%%%%%%%%%%%%%%%%%%%%%%%%%%%%%%%%%%%%%%%%%%%%%%%%%%%%%%%%%%%%%%%%%%%%%%%%%%%%%%
%  TEMPLATE INSA CENTRE VAL DE LOIRE : main.tex
%
%  - Chaque partie du document est un fichier séparé, importé plus bas avec
%    \include{...} : pagetitre, introduction, chapitre1 ... chapitre6, qcm, conclusion.
%    Ces fichiers ne contiennent ni préambule ni \begin{document}. Un chapitre
%    commence par \section{...} (la classe est article : un chapitre = une \section) ;
%    l'introduction et la conclusion sont des \section* ajoutées au sommaire.
%  - Nouveau document : modifier le bloc « Informations du document »,
%    ajouter ou retirer des \include, décommenter les sections utiles
%    (résumés, table des illustrations, annexes).
%  - Compilation : pdflatex main, biber main, pdflatex main, pdflatex main.
%    (deux passes de pdflatex au minimum : sommaire, renvois, page de titre TikZ)
%  - Couleurs de la charte : insarouge (principale) et insagris (secondaire),
%    relevées sur le logo INSA. À utiliser dans les chapitres (figures, tableaux).
%  - Glossaire : entrées dans glossaire.tex, appel \gls{qubit} dans le texte.
%    Bibliographie : références dans sources.bib, appel \parencite{cle}.
%%%%%%%%%%%%%%%%%%%%%%%%%%%%%%%%%%%%%%%%%%%%%%%%%%%%%%%%%%%%%%%%%%%%%%%%%%%%%%
\documentclass[10pt]{article}

%%%%% Informations du document (à modifier pour chaque document) %%%%%
\newcommand{\doctitre}{Informatique quantique}
\newcommand{\docsoustitre}{Notes de cours -- 4A FISE SF}
\newcommand{\doctagline}{De l'onde de Schr\"odinger au circuit quantique}
\newcommand{\docauteur}{Niels LOPEZ}
\newcommand{\docprof}{M.~Adell}
\newcommand{\docecole}{INSA Centre Val de Loire}
\newcommand{\docdate}{\today}
% Nombre de pages de garde au début du document : la page de titre (pagetitre.tex) ou des pages
% importées en PDF (\includepdf). Elles sont comptées dans la numérotation (le numéro affiché est
% celui de la page du PDF, sans décalage) mais n'ont ni en-tête ni numéro : le logo n'apparaît
% qu'après elles. 1 = la page de titre seule ; 0 = aucune page de garde.
\newcounter{pagesdegarde}
\setcounter{pagesdegarde}{1}

%%%%% Package de base %%%%%
\usepackage{graphicx} % pour insérer des images
\graphicspath{{images/}}
\usepackage{float} % pour placer les images ([H])

\usepackage{amsmath}% pour les équations
\usepackage{amssymb}% pour les symboles mathématiques (\mathbb, ...)
\numberwithin{figure}{section}%pour les numérotations des figures
\numberwithin{equation}{section}%idem pour les équations
\numberwithin{table}{section}%idem pour les tableaux

\usepackage{array}%pour les tableaux
\usepackage{makecell}  % idem

\usepackage[utf8]{inputenc}%pour glossaire
\usepackage[T1]{fontenc}
\usepackage{lmodern} % police Latin Modern (équivalent de Computer Modern en T1)
\usepackage[french]{babel}

%%%%% Charte graphique INSA CVL %%%%%
\usepackage{xcolor}
\definecolor{insarouge}{HTML}{E42618} % rouge du logo INSA (couleur principale)
\definecolor{insagris}{HTML}{5F5F5F}  % gris du logo INSA (couleur secondaire)

%%%%% Pour la biblio %%%%%
\usepackage{csquotes}
\usepackage[backend=biber,style=apa]{biblatex} % pour la biblio
\addbibresource{sources.bib} % fichier contenant les références
% autorise les coupures de ligne dans les URL et les DOI de la bibliographie
\setcounter{biburlnumpenalty}{100}
\setcounter{biburlucpenalty}{100}
\setcounter{biburllcpenalty}{100}

%%%%% Pour les lignes de code %%%%%
\usepackage{listings}
\lstset{
  basicstyle=\ttfamily\small,
  backgroundcolor=\color{insagris!8},
  frame=single,
  rulecolor=\color{insagris!60},
  breaklines=true
}
\lstdefinestyle{terminal}{
  backgroundcolor=\color{black!5},
  basicstyle=\ttfamily\small\color{black},
  frame=single,
  framerule=0pt,
  xleftmargin=1em,
  showstringspaces=false,
  breaklines=true,
  escapeinside={(*@}{@*)},
}
\newcommand{\prompt}{\textcolor{gray}{\texttt{user@linux:~\$}} }
\lstdefinestyle{darkterminal}{
  backgroundcolor=\color{black},
  basicstyle=\ttfamily\small\color{green},
  frame=single,
  framerule=0pt,
  xleftmargin=1em,
  breaklines=true,
  showstringspaces=false
}
\lstdefinestyle{verbatimlike}{
    basicstyle=\ttfamily\small,     % police type machine à écrire, comme verbatim
    breaklines=true,                % retour automatique à la ligne
    breakatwhitespace=false,        % casse les mots si nécessaire
    showstringspaces=false,         % pas de soulignement des espaces dans les chaînes
    numbers=none,                   % pas de numérotation
    frame=none,                     % pas de cadre
    framerule=0pt,                  % épaisseur de la bordure à 0
    rulecolor=\color{black},        % couleur du cadre (inutile ici)
    backgroundcolor=,               % pas de fond
    tabsize=4,                      % largeur d'une tabulation
    keepspaces=true,                % préserve les espaces consécutifs
    columns=flexible                % adapte les colonnes pour éviter les débordements
}

\usepackage{fancyvrb} %pour les retours à la ligne dans verbatim

\usepackage{enumitem} %pour les listes de puces
\setlist[itemize,1]{label=\textcolor{insarouge}{\textemdash}} % puces rouges

\usepackage{changepage} %pour dépasser dans les marges
\usepackage{subcaption}    % pour gérer des figures côte à côte

\usepackage[version=4]{mhchem}   % pour les équations chimiques
\usepackage{booktabs}            % tableau plus soigné (optionnel, sinon \hline suffit)

\usepackage{tikz}
\usetikzlibrary{positioning, arrows.meta, fit, backgrounds, calc, shapes.geometric, babel,
                decorations.pathmorphing, decorations.pathreplacing, patterns} % pour les graphs
\usepackage{adjustbox}

%%%%% Symboles de circuits quantiques (utilisés dans tous les chapitres) %%%%%
\tikzset{
  ctl/.pic   = {\fill[insarouge] (0,0) circle (2.6pt);},                       % contrôle
  cible/.pic = {\draw[insarouge, thick, fill=white] (0,0) circle (0.28);       % cible (NOT)
                \draw[insarouge, thick] (-0.28,0) -- (0.28,0);
                \draw[insarouge, thick] (0,-0.28) -- (0,0.28);},
  croix/.pic = {\draw[insarouge, thick] (-0.17,-0.17) -- (0.17,0.17);          % croix du SWAP
                \draw[insarouge, thick] (-0.17,0.17) -- (0.17,-0.17);},
  mesure/.pic = {\draw[thick, fill=white] (-0.3,-0.27) rectangle (0.3,0.27);   % mesure
                 \draw[thick] (-0.19,-0.13) arc[start angle=155, end angle=25, radius=0.21];
                 \draw[thick, -{Stealth[length=1.3mm]}] (0,-0.17) -- (0.14,0.12);},
  porte/.style = {draw=insarouge, fill=insarouge!8, minimum size=7mm, thick},  % porte à un qubit
}

%%%%% Encadrés (définitions, résultats, remarques, récapitulatifs, exercices) %%%%%
\usepackage[skins,breakable]{tcolorbox}
\tcbset{boxrule=0.6pt,arc=2pt,left=6pt,right=6pt,top=4pt,bottom=4pt}
\newtcolorbox{definition}[1]{colback=insarouge!4,colframe=insarouge,
  fonttitle=\bfseries,title={#1}}
\newtcolorbox{resultat}[1]{colback=insarouge!10,colframe=insarouge!65!black,boxrule=1pt,
  fonttitle=\bfseries,title={#1}}
\newtcolorbox{remarque}[1][Remarque]{colback=insagris!7,colframe=insagris,
  fonttitle=\bfseries,title={#1}}
% Récapitulatifs de fin de chapitre : points à retenir et savoir-faire
\newtcolorbox{aretenir}{colback=insarouge!6,colframe=insarouge,boxrule=1pt,
  fonttitle=\bfseries,title={À retenir}}
\newtcolorbox{savoirfaire}{colback=insagris!6,colframe=insagris,
  fonttitle=\bfseries,title={Savoir-faire}}
% Exercices (numérotés dans la section, cités par \ref{...}) et leur correction
\newcounter{exercice}[section]
\renewcommand{\theexercice}{\thesection.\arabic{exercice}}
\newtcolorbox{exercicebox}[1]{colback=white,
  colframe=insarouge,boxrule=0.8pt,fonttitle=\bfseries,coltitle=white,
  colbacktitle=insarouge,title={Exercice~\theexercice\ --- #1}}
% \refstepcounter est appelé avant la boîte : \label{...} placé dans l'énoncé donne le bon numéro
\newenvironment{exercice}[1]{\refstepcounter{exercice}\begin{exercicebox}{#1}}{\end{exercicebox}}
% Début d'une correction : le mot « Correction. » reste sur la même page que la première question
\makeatletter
\newcommand{\corrstart}{%
  \captionsetup{hypcap=false}%
  \textbf{\color{insagris}Correction.}\par\nobreak
  \let\corr@oldparagraph\paragraph
  \renewcommand{\paragraph}[1]{\corr@oldparagraph{##1}\@secpenalty=-300\relax}%
  \@secpenalty=10000\relax}
\makeatother
\newtcolorbox{correction}{breakable,colback=insagris!5,colframe=insagris,boxrule=0pt,
  toprule=0pt,bottomrule=0pt,rightrule=0pt,leftrule=2.5pt,sharp corners,
  before upper={\corrstart}}

%%%%% Macros de mathématiques / physique %%%%%
\newcommand{\ii}{\mathrm{i}}   % unité imaginaire
\newcommand{\ee}{\mathrm{e}}   % nombre e
\newcommand{\dd}{\mathrm{d}}   % différentielle droite
\newcommand{\pd}[2]{\dfrac{\partial #1}{\partial #2}}
\newcommand{\pdd}[2]{\dfrac{\partial^{2} #1}{\partial #2^{2}}}
\newcommand{\Hop}{\hat{H}}     % opérateur hamiltonien
% Notation de Dirac : \ket{a} = |a>, \bra{a} = <a|, \braket{a}{b} = <a|b>
\newcommand{\ket}[1]{|#1\rangle}
\newcommand{\bra}[1]{\langle #1|}
\newcommand{\braket}[2]{\langle #1|#2\rangle}

%%%%% Titres de sections (rouge INSA) %%%%%
\usepackage{titlesec}
\titleformat{\section}
  {\fontsize{15}{18}\selectfont\bfseries\color{insarouge}}{\thesection}{0.7em}{}
  [{\vspace{0.5ex}\color{insarouge}\titlerule[0.8pt]}]
% Titres non numérotés (\section*, sommaire, résumés) : même style. Le filet n'est tracé que si le titre
% n'est pas vide (les titres vides servent à masquer le titre automatique de la liste des figures ou du glossaire).
\newsavebox{\insatitbox}
\newcommand{\insatitrenn}[1]{%
  \sbox{\insatitbox}{\ignorespaces#1\unskip}%
  \ifdim\wd\insatitbox=0pt\else #1\par\nobreak\vspace{0.5ex}{\color{insarouge}\hrule height 0.8pt}\fi}
\titleformat{name=\section,numberless}
  {\fontsize{15}{18}\selectfont\bfseries\color{insarouge}}{}{0pt}{\insatitrenn}
\titleformat{\subsection}
  {\fontsize{12}{14.5}\selectfont\bfseries\color{insarouge}}{\thesubsection}{0.7em}{}
\titleformat{\subsubsection}
  {\fontsize{10.5}{13}\selectfont\bfseries\color{insagris}}{\thesubsubsection}{0.7em}{}
\titleformat{\paragraph}[runin]
  {\normalfont\normalsize\bfseries\color{insagris}}{\theparagraph}{0.7em}{}
\titlespacing*{\paragraph}{0pt}{1.5ex plus .5ex minus .2ex}{0.6em}
\titlespacing*{\section}{0pt}{3ex plus 1ex minus .2ex}{1.8ex plus .2ex}
\titlespacing*{\subsection}{0pt}{2.5ex plus 1ex minus .2ex}{1.2ex plus .2ex}
\titlespacing*{\subsubsection}{0pt}{2ex plus .5ex minus .2ex}{1ex plus .2ex}

%%%%% Pour le glossaire %%%%%
\usepackage[hidelinks]{hyperref} %pour rendre les liens cliquables
\hypersetup{pdftitle={\doctitre}, pdfauthor={\docauteur}, pdfsubject={\docsoustitre}}
\usepackage[nonumberlist]{glossaries}
\renewcommand{\glstextformat}[1]{\textcolor{insarouge}{#1}} %pour colorer les mots glossariés
%\makeglossaries
\makenoidxglossaries % pas de makeindex : deux passes de pdflatex suffisent
% =====================================================================
%  glossaire.tex : entrées du glossaire (chargé par main.tex via % =====================================================================
%  glossaire.tex : entrées du glossaire (chargé par main.tex via % =====================================================================
%  glossaire.tex : entrées du glossaire (chargé par main.tex via \input{glossaire})
%  Utilisation dans le texte : \gls{qubit} (le mot apparaît en rouge).
%  Toutes les entrées sont imprimées en fin de document (\glsaddall dans main.tex).
%  La clé sort={...} donne le texte brut utilisé pour le tri (accents, commandes).
%  Compilation : deux passes de pdflatex suffisent (\makenoidxglossaries).
% =====================================================================

\newglossaryentry{amplitude}{name={amplitude de probabilité}, sort={amplitude},
  description={Coefficient complexe $c_i$ devant un état de base dans une superposition ; la probabilité de mesurer cet état est $|c_i|^2$}}

\newglossaryentry{basecalcul}{name={base de calcul}, sort={base de calcul},
  description={Base orthonormée $\{\ket{0},\ket{1}\}$ d'un qubit (base $Z$) ; pour $n$ qubits, les $2^n$ états $\ket{i_1\dots i_n}$}}

\newglossaryentry{bell}{name={Bell (état de)}, sort={Bell},
  description={L'un des quatre états intriqués maximaux de deux qubits, $\ket{\Phi^\pm}$ et $\ket{\Psi^\pm}$ ; ils forment une base orthonormée}}

\newglossaryentry{bloch}{name={Bloch (sphère de)}, sort={Bloch},
  description={Sphère unité dont les points représentent les états purs d'un qubit ; l'intérieur de la boule représente les états mixtes}}

\newglossaryentry{born}{name={Born (règle de)}, sort={Born},
  description={La probabilité de trouver la particule en $x$ est $|\Psi(x,t)|^2\,\dd x$ ; pour un qubit, $P(i)=|\braket{i}{\Psi}|^2$}}

\newglossaryentry{bra}{name={bra}, sort={bra},
  description={Vecteur ligne $\bra{a}=(\ket{a})^\dagger$, transposé conjugué du ket $\ket{a}$}}

\newglossaryentry{cnot}{name={CNOT}, sort={CNOT},
  description={Porte à deux qubits : la cible est inversée ($X$) si le contrôle vaut $1$ ; avec la convention du cours (contrôle $B$, cible $A$), $\ket{a,b}\to\ket{a\oplus b,b}$}}

\newglossaryentry{coherence}{name={cohérence}, sort={coherence},
  description={Terme hors diagonale de la matrice de densité ; il contient la phase relative entre les composantes de l'état}}

\newglossaryentry{densite}{name={densité (matrice de)}, sort={densite},
  description={Opérateur $\rho$ hermitien, positif, de trace $1$, qui décrit un état pur ($\rho=\ket{\Psi}\bra{\Psi}$) ou mixte ($\rho=\sum_ip_i\ket{\Psi_i}\bra{\Psi_i}$)}}

\newglossaryentry{dirac}{name={Dirac (notation de)}, sort={Dirac},
  description={Notation $\ket{a}$ (ket), $\bra{a}$ (bra) et $\braket{a}{b}$ (produit scalaire) pour les vecteurs d'un espace de Hilbert}}

\newglossaryentry{ebit}{name={e-bit}, sort={e-bit},
  description={Paire de qubits maximalement intriqués, par exemple $\ket{\Phi^+}$, partagée entre deux parties ; ressource de la téléportation}}

\newglossaryentry{entropie}{name={entropie de Shannon}, sort={entropie de Shannon},
  description={$H(A)=-\sum_ap(a)\log_2p(a)$, incertitude moyenne (en bits) sur une variable aléatoire $A$}}

\newglossaryentry{vonneumann}{name={entropie de von Neumann}, sort={entropie de von Neumann},
  description={$S(\rho)=-\operatorname{Tr}(\rho\log_2\rho)=-\sum_i\lambda_i\log_2\lambda_i$ ; elle vaut $0$ pour un état pur}}

\newglossaryentry{hilbert}{name={Hilbert (espace de)}, sort={Hilbert},
  description={Espace vectoriel complexe muni d'un produit scalaire et complet ; c'est l'espace des états d'un système quantique}}

\newglossaryentry{etatmixte}{name={état mixte}, sort={etat mixte},
  description={État décrit par une matrice de densité avec $\operatorname{Tr}\rho^2<1$ : mélange statistique d'états purs}}

\newglossaryentry{etatpur}{name={état pur}, sort={etat pur},
  description={État décrit par un vecteur $\ket{\Psi}$, soit $\rho=\ket{\Psi}\bra{\Psi}$ et $\operatorname{Tr}\rho^2=1$}}

\newglossaryentry{stationnaire}{name={état stationnaire}, sort={etat stationnaire},
  description={Solution de l'équation de Schrödinger de la forme $\psi(x)\,\ee^{-\ii Et/\hbar}$ : énergie certaine et densité de probabilité indépendante du temps}}

\newglossaryentry{fonctiononde}{name={fonction d'onde}, sort={fonction d'onde},
  description={Fonction complexe $\Psi(x,t)$ décrivant l'état d'une particule ; $|\Psi|^2$ est une densité de probabilité de présence}}

\newglossaryentry{fredkin}{name={Fredkin (porte de)}, sort={Fredkin},
  description={Porte à trois qubits (C-SWAP) : échange les deux cibles si le contrôle vaut $1$}}

\newglossaryentry{hadamard}{name={Hadamard (porte de)}, sort={Hadamard},
  description={Porte $H=\frac{1}{\sqrt2}\begin{pmatrix}1&1\\1&-1\end{pmatrix}$ ; elle crée la superposition $\ket{+}$ à partir de $\ket{0}$ et change de base $Z\leftrightarrow X$}}

\newglossaryentry{hamiltonien}{name={hamiltonien}, sort={hamiltonien},
  description={Opérateur $\hat H=-\frac{\hbar^2}{2m}\partial_x^2+V(x)$ associé à l'énergie totale}}

\newglossaryentry{hermitien}{name={hermitien}, sort={hermitien},
  description={Se dit d'un opérateur égal à son adjoint, $A^\dagger=A$ ; ses valeurs propres sont réelles}}

\newglossaryentry{informationmutuelle}{name={information mutuelle}, sort={information mutuelle},
  description={$I(A;B)=H(A)+H(B)-H(A,B)$ : ce que la connaissance de $B$ apprend sur $A$}}

\newglossaryentry{intrication}{name={intrication}, sort={intrication},
  description={Propriété d'un état de plusieurs qubits qui n'est pas un produit d'états de chaque qubit}}

\newglossaryentry{ket}{name={ket}, sort={ket},
  description={Vecteur colonne $\ket{a}$ d'un espace de Hilbert ; il représente un état quantique}}

\newglossaryentry{mesure}{name={mesure}, sort={mesure},
  description={Opération qui donne l'un des résultats possibles avec la probabilité de Born et projette l'état sur l'état associé}}

\newglossaryentry{nonclonage}{name={non-clonage (théorème du)}, sort={non-clonage},
  description={Aucune opération ne copie un état quantique inconnu ; la téléportation déplace l'état sans le copier}}

\newglossaryentry{operateur}{name={opérateur}, sort={operateur},
  description={Application linéaire d'un espace de Hilbert dans lui-même ; en dimension finie, une matrice}}

\newglossaryentry{orthonormee}{name={orthonormée (base)}, sort={orthonormee},
  description={Base $\{\ket{e_i}\}$ de vecteurs de norme $1$ et deux à deux orthogonaux : $\braket{e_i}{e_j}=\delta_{ij}$}}

\newglossaryentry{pauli}{name={Pauli (matrices de)}, sort={Pauli},
  description={Les trois matrices $X=\sigma_x$, $Y=\sigma_y$, $Z=\sigma_z$ ; ce sont des portes à un qubit}}

\newglossaryentry{population}{name={population}, sort={population},
  description={Terme diagonal $\rho_{ii}$ de la matrice de densité : probabilité de mesurer l'état de base $\ket{i}$}}

\newglossaryentry{tenseur}{name={produit tensoriel}, sort={produit tensoriel},
  description={Construction $\mathcal{H}_A\otimes\mathcal{H}_B$ de l'espace de deux systèmes ; $\dim=\dim\mathcal{H}_A\times\dim\mathcal{H}_B$}}

\newglossaryentry{projecteur}{name={projecteur}, sort={projecteur},
  description={Opérateur $\hat P$ tel que $\hat P^2=\hat P$ et $\hat P^\dagger=\hat P$ ; par exemple $\hat P_i=\ket{i}\bra{i}$}}

\newglossaryentry{qubit}{name={qubit}, sort={qubit},
  description={Bit quantique : système à deux niveaux d'état $\alpha\ket{0}+\beta\ket{1}$, avec $|\alpha|^2+|\beta|^2=1$}}

\newglossaryentry{schrodinger}{name={Schrödinger (équation de)}, sort={Schrodinger},
  description={Équation $\ii\hbar\,\partial_t\Psi=\hat H\Psi$ qui gouverne l'évolution de la fonction d'onde}}

\newglossaryentry{separable}{name={séparable (état)}, sort={separable},
  description={État de deux qubits qui s'écrit $\ket{\Psi}_A\otimes\ket{\Psi}_B$ ; critère : $c_{00}c_{11}-c_{01}c_{10}=0$}}

\newglossaryentry{superposition}{name={superposition}, sort={superposition},
  description={Combinaison linéaire $\sum_ic_i\ket{i}$ d'états ; conséquence de la linéarité de l'équation de Schrödinger}}

\newglossaryentry{swap}{name={SWAP}, sort={SWAP},
  description={Porte à deux qubits qui échange leurs états : $\ket{a,b}\to\ket{b,a}$}}

\newglossaryentry{teleportation}{name={téléportation quantique}, sort={teleportation},
  description={Protocole qui transmet un état de qubit inconnu avec un e-bit et deux bits classiques, sans déplacer le qubit}}

\newglossaryentry{toffoli}{name={Toffoli (porte de)}, sort={Toffoli},
  description={Porte à trois qubits (CCNOT) : la cible est inversée si les deux contrôles valent $1$ ; $\ket{x,y,z}\to\ket{x,y,z\oplus xy}$}}

\newglossaryentry{trace}{name={trace (partielle)}, sort={trace},
  description={$\operatorname{Tr}_B$ : opération qui «~oublie~» le sous-système $B$ et donne l'état $\rho_A$ du sous-système $A$}}

\newglossaryentry{unitaire}{name={unitaire}, sort={unitaire},
  description={Matrice $U$ telle que $U^\dagger U=\mathrm{Id}$ ; elle conserve la norme et est réversible ; c'est une porte quantique}}

\newglossaryentry{observable}{name={observable}, sort={observable},
  description={Opérateur hermitien qui représente une grandeur mesurable ; ses valeurs propres sont les résultats possibles d'une mesure}}

\newglossaryentry{effondrement}{name={effondrement (de l'état)}, sort={effondrement},
  description={Changement brusque de l'état lors d'une mesure : il devient l'état propre associé au résultat obtenu}}

\newglossaryentry{valeurpropre}{name={valeur propre}, sort={valeur propre},
  description={Nombre $\lambda$ tel que $A\ket{v}=\lambda\ket{v}$ avec $\ket{v}\neq0$ (vecteur propre) ; pour une observable, résultat possible d'une mesure}}

\newglossaryentry{phase}{name={phase (relative)}, sort={phase},
  description={Angle $\varphi$ entre les amplitudes de $\ket{0}$ et $\ket{1}$ ; observable, contrairement à la phase globale}}

\newglossaryentry{purete}{name={pureté}, sort={purete},
  description={Nombre $\operatorname{Tr}\rho^2$, compris entre $\frac1d$ et $1$ (dimension $d$) ; il vaut $1$ si et seulement si l'état est pur}}

\newglossaryentry{porte}{name={porte quantique}, sort={porte quantique},
  description={Opération unitaire appliquée à un ou plusieurs qubits ; dessinée comme une boîte sur un circuit}}
)
%  Utilisation dans le texte : \gls{qubit} (le mot apparaît en rouge).
%  Toutes les entrées sont imprimées en fin de document (\glsaddall dans main.tex).
%  La clé sort={...} donne le texte brut utilisé pour le tri (accents, commandes).
%  Compilation : deux passes de pdflatex suffisent (\makenoidxglossaries).
% =====================================================================

\newglossaryentry{amplitude}{name={amplitude de probabilité}, sort={amplitude},
  description={Coefficient complexe $c_i$ devant un état de base dans une superposition ; la probabilité de mesurer cet état est $|c_i|^2$}}

\newglossaryentry{basecalcul}{name={base de calcul}, sort={base de calcul},
  description={Base orthonormée $\{\ket{0},\ket{1}\}$ d'un qubit (base $Z$) ; pour $n$ qubits, les $2^n$ états $\ket{i_1\dots i_n}$}}

\newglossaryentry{bell}{name={Bell (état de)}, sort={Bell},
  description={L'un des quatre états intriqués maximaux de deux qubits, $\ket{\Phi^\pm}$ et $\ket{\Psi^\pm}$ ; ils forment une base orthonormée}}

\newglossaryentry{bloch}{name={Bloch (sphère de)}, sort={Bloch},
  description={Sphère unité dont les points représentent les états purs d'un qubit ; l'intérieur de la boule représente les états mixtes}}

\newglossaryentry{born}{name={Born (règle de)}, sort={Born},
  description={La probabilité de trouver la particule en $x$ est $|\Psi(x,t)|^2\,\dd x$ ; pour un qubit, $P(i)=|\braket{i}{\Psi}|^2$}}

\newglossaryentry{bra}{name={bra}, sort={bra},
  description={Vecteur ligne $\bra{a}=(\ket{a})^\dagger$, transposé conjugué du ket $\ket{a}$}}

\newglossaryentry{cnot}{name={CNOT}, sort={CNOT},
  description={Porte à deux qubits : la cible est inversée ($X$) si le contrôle vaut $1$ ; avec la convention du cours (contrôle $B$, cible $A$), $\ket{a,b}\to\ket{a\oplus b,b}$}}

\newglossaryentry{coherence}{name={cohérence}, sort={coherence},
  description={Terme hors diagonale de la matrice de densité ; il contient la phase relative entre les composantes de l'état}}

\newglossaryentry{densite}{name={densité (matrice de)}, sort={densite},
  description={Opérateur $\rho$ hermitien, positif, de trace $1$, qui décrit un état pur ($\rho=\ket{\Psi}\bra{\Psi}$) ou mixte ($\rho=\sum_ip_i\ket{\Psi_i}\bra{\Psi_i}$)}}

\newglossaryentry{dirac}{name={Dirac (notation de)}, sort={Dirac},
  description={Notation $\ket{a}$ (ket), $\bra{a}$ (bra) et $\braket{a}{b}$ (produit scalaire) pour les vecteurs d'un espace de Hilbert}}

\newglossaryentry{ebit}{name={e-bit}, sort={e-bit},
  description={Paire de qubits maximalement intriqués, par exemple $\ket{\Phi^+}$, partagée entre deux parties ; ressource de la téléportation}}

\newglossaryentry{entropie}{name={entropie de Shannon}, sort={entropie de Shannon},
  description={$H(A)=-\sum_ap(a)\log_2p(a)$, incertitude moyenne (en bits) sur une variable aléatoire $A$}}

\newglossaryentry{vonneumann}{name={entropie de von Neumann}, sort={entropie de von Neumann},
  description={$S(\rho)=-\operatorname{Tr}(\rho\log_2\rho)=-\sum_i\lambda_i\log_2\lambda_i$ ; elle vaut $0$ pour un état pur}}

\newglossaryentry{hilbert}{name={Hilbert (espace de)}, sort={Hilbert},
  description={Espace vectoriel complexe muni d'un produit scalaire et complet ; c'est l'espace des états d'un système quantique}}

\newglossaryentry{etatmixte}{name={état mixte}, sort={etat mixte},
  description={État décrit par une matrice de densité avec $\operatorname{Tr}\rho^2<1$ : mélange statistique d'états purs}}

\newglossaryentry{etatpur}{name={état pur}, sort={etat pur},
  description={État décrit par un vecteur $\ket{\Psi}$, soit $\rho=\ket{\Psi}\bra{\Psi}$ et $\operatorname{Tr}\rho^2=1$}}

\newglossaryentry{stationnaire}{name={état stationnaire}, sort={etat stationnaire},
  description={Solution de l'équation de Schrödinger de la forme $\psi(x)\,\ee^{-\ii Et/\hbar}$ : énergie certaine et densité de probabilité indépendante du temps}}

\newglossaryentry{fonctiononde}{name={fonction d'onde}, sort={fonction d'onde},
  description={Fonction complexe $\Psi(x,t)$ décrivant l'état d'une particule ; $|\Psi|^2$ est une densité de probabilité de présence}}

\newglossaryentry{fredkin}{name={Fredkin (porte de)}, sort={Fredkin},
  description={Porte à trois qubits (C-SWAP) : échange les deux cibles si le contrôle vaut $1$}}

\newglossaryentry{hadamard}{name={Hadamard (porte de)}, sort={Hadamard},
  description={Porte $H=\frac{1}{\sqrt2}\begin{pmatrix}1&1\\1&-1\end{pmatrix}$ ; elle crée la superposition $\ket{+}$ à partir de $\ket{0}$ et change de base $Z\leftrightarrow X$}}

\newglossaryentry{hamiltonien}{name={hamiltonien}, sort={hamiltonien},
  description={Opérateur $\hat H=-\frac{\hbar^2}{2m}\partial_x^2+V(x)$ associé à l'énergie totale}}

\newglossaryentry{hermitien}{name={hermitien}, sort={hermitien},
  description={Se dit d'un opérateur égal à son adjoint, $A^\dagger=A$ ; ses valeurs propres sont réelles}}

\newglossaryentry{informationmutuelle}{name={information mutuelle}, sort={information mutuelle},
  description={$I(A;B)=H(A)+H(B)-H(A,B)$ : ce que la connaissance de $B$ apprend sur $A$}}

\newglossaryentry{intrication}{name={intrication}, sort={intrication},
  description={Propriété d'un état de plusieurs qubits qui n'est pas un produit d'états de chaque qubit}}

\newglossaryentry{ket}{name={ket}, sort={ket},
  description={Vecteur colonne $\ket{a}$ d'un espace de Hilbert ; il représente un état quantique}}

\newglossaryentry{mesure}{name={mesure}, sort={mesure},
  description={Opération qui donne l'un des résultats possibles avec la probabilité de Born et projette l'état sur l'état associé}}

\newglossaryentry{nonclonage}{name={non-clonage (théorème du)}, sort={non-clonage},
  description={Aucune opération ne copie un état quantique inconnu ; la téléportation déplace l'état sans le copier}}

\newglossaryentry{operateur}{name={opérateur}, sort={operateur},
  description={Application linéaire d'un espace de Hilbert dans lui-même ; en dimension finie, une matrice}}

\newglossaryentry{orthonormee}{name={orthonormée (base)}, sort={orthonormee},
  description={Base $\{\ket{e_i}\}$ de vecteurs de norme $1$ et deux à deux orthogonaux : $\braket{e_i}{e_j}=\delta_{ij}$}}

\newglossaryentry{pauli}{name={Pauli (matrices de)}, sort={Pauli},
  description={Les trois matrices $X=\sigma_x$, $Y=\sigma_y$, $Z=\sigma_z$ ; ce sont des portes à un qubit}}

\newglossaryentry{population}{name={population}, sort={population},
  description={Terme diagonal $\rho_{ii}$ de la matrice de densité : probabilité de mesurer l'état de base $\ket{i}$}}

\newglossaryentry{tenseur}{name={produit tensoriel}, sort={produit tensoriel},
  description={Construction $\mathcal{H}_A\otimes\mathcal{H}_B$ de l'espace de deux systèmes ; $\dim=\dim\mathcal{H}_A\times\dim\mathcal{H}_B$}}

\newglossaryentry{projecteur}{name={projecteur}, sort={projecteur},
  description={Opérateur $\hat P$ tel que $\hat P^2=\hat P$ et $\hat P^\dagger=\hat P$ ; par exemple $\hat P_i=\ket{i}\bra{i}$}}

\newglossaryentry{qubit}{name={qubit}, sort={qubit},
  description={Bit quantique : système à deux niveaux d'état $\alpha\ket{0}+\beta\ket{1}$, avec $|\alpha|^2+|\beta|^2=1$}}

\newglossaryentry{schrodinger}{name={Schrödinger (équation de)}, sort={Schrodinger},
  description={Équation $\ii\hbar\,\partial_t\Psi=\hat H\Psi$ qui gouverne l'évolution de la fonction d'onde}}

\newglossaryentry{separable}{name={séparable (état)}, sort={separable},
  description={État de deux qubits qui s'écrit $\ket{\Psi}_A\otimes\ket{\Psi}_B$ ; critère : $c_{00}c_{11}-c_{01}c_{10}=0$}}

\newglossaryentry{superposition}{name={superposition}, sort={superposition},
  description={Combinaison linéaire $\sum_ic_i\ket{i}$ d'états ; conséquence de la linéarité de l'équation de Schrödinger}}

\newglossaryentry{swap}{name={SWAP}, sort={SWAP},
  description={Porte à deux qubits qui échange leurs états : $\ket{a,b}\to\ket{b,a}$}}

\newglossaryentry{teleportation}{name={téléportation quantique}, sort={teleportation},
  description={Protocole qui transmet un état de qubit inconnu avec un e-bit et deux bits classiques, sans déplacer le qubit}}

\newglossaryentry{toffoli}{name={Toffoli (porte de)}, sort={Toffoli},
  description={Porte à trois qubits (CCNOT) : la cible est inversée si les deux contrôles valent $1$ ; $\ket{x,y,z}\to\ket{x,y,z\oplus xy}$}}

\newglossaryentry{trace}{name={trace (partielle)}, sort={trace},
  description={$\operatorname{Tr}_B$ : opération qui «~oublie~» le sous-système $B$ et donne l'état $\rho_A$ du sous-système $A$}}

\newglossaryentry{unitaire}{name={unitaire}, sort={unitaire},
  description={Matrice $U$ telle que $U^\dagger U=\mathrm{Id}$ ; elle conserve la norme et est réversible ; c'est une porte quantique}}

\newglossaryentry{observable}{name={observable}, sort={observable},
  description={Opérateur hermitien qui représente une grandeur mesurable ; ses valeurs propres sont les résultats possibles d'une mesure}}

\newglossaryentry{effondrement}{name={effondrement (de l'état)}, sort={effondrement},
  description={Changement brusque de l'état lors d'une mesure : il devient l'état propre associé au résultat obtenu}}

\newglossaryentry{valeurpropre}{name={valeur propre}, sort={valeur propre},
  description={Nombre $\lambda$ tel que $A\ket{v}=\lambda\ket{v}$ avec $\ket{v}\neq0$ (vecteur propre) ; pour une observable, résultat possible d'une mesure}}

\newglossaryentry{phase}{name={phase (relative)}, sort={phase},
  description={Angle $\varphi$ entre les amplitudes de $\ket{0}$ et $\ket{1}$ ; observable, contrairement à la phase globale}}

\newglossaryentry{purete}{name={pureté}, sort={purete},
  description={Nombre $\operatorname{Tr}\rho^2$, compris entre $\frac1d$ et $1$ (dimension $d$) ; il vaut $1$ si et seulement si l'état est pur}}

\newglossaryentry{porte}{name={porte quantique}, sort={porte quantique},
  description={Opération unitaire appliquée à un ou plusieurs qubits ; dessinée comme une boîte sur un circuit}}
)
%  Utilisation dans le texte : \gls{qubit} (le mot apparaît en rouge).
%  Toutes les entrées sont imprimées en fin de document (\glsaddall dans main.tex).
%  La clé sort={...} donne le texte brut utilisé pour le tri (accents, commandes).
%  Compilation : deux passes de pdflatex suffisent (\makenoidxglossaries).
% =====================================================================

\newglossaryentry{amplitude}{name={amplitude de probabilité}, sort={amplitude},
  description={Coefficient complexe $c_i$ devant un état de base dans une superposition ; la probabilité de mesurer cet état est $|c_i|^2$}}

\newglossaryentry{basecalcul}{name={base de calcul}, sort={base de calcul},
  description={Base orthonormée $\{\ket{0},\ket{1}\}$ d'un qubit (base $Z$) ; pour $n$ qubits, les $2^n$ états $\ket{i_1\dots i_n}$}}

\newglossaryentry{bell}{name={Bell (état de)}, sort={Bell},
  description={L'un des quatre états intriqués maximaux de deux qubits, $\ket{\Phi^\pm}$ et $\ket{\Psi^\pm}$ ; ils forment une base orthonormée}}

\newglossaryentry{bloch}{name={Bloch (sphère de)}, sort={Bloch},
  description={Sphère unité dont les points représentent les états purs d'un qubit ; l'intérieur de la boule représente les états mixtes}}

\newglossaryentry{born}{name={Born (règle de)}, sort={Born},
  description={La probabilité de trouver la particule en $x$ est $|\Psi(x,t)|^2\,\dd x$ ; pour un qubit, $P(i)=|\braket{i}{\Psi}|^2$}}

\newglossaryentry{bra}{name={bra}, sort={bra},
  description={Vecteur ligne $\bra{a}=(\ket{a})^\dagger$, transposé conjugué du ket $\ket{a}$}}

\newglossaryentry{cnot}{name={CNOT}, sort={CNOT},
  description={Porte à deux qubits : la cible est inversée ($X$) si le contrôle vaut $1$ ; avec la convention du cours (contrôle $B$, cible $A$), $\ket{a,b}\to\ket{a\oplus b,b}$}}

\newglossaryentry{coherence}{name={cohérence}, sort={coherence},
  description={Terme hors diagonale de la matrice de densité ; il contient la phase relative entre les composantes de l'état}}

\newglossaryentry{densite}{name={densité (matrice de)}, sort={densite},
  description={Opérateur $\rho$ hermitien, positif, de trace $1$, qui décrit un état pur ($\rho=\ket{\Psi}\bra{\Psi}$) ou mixte ($\rho=\sum_ip_i\ket{\Psi_i}\bra{\Psi_i}$)}}

\newglossaryentry{dirac}{name={Dirac (notation de)}, sort={Dirac},
  description={Notation $\ket{a}$ (ket), $\bra{a}$ (bra) et $\braket{a}{b}$ (produit scalaire) pour les vecteurs d'un espace de Hilbert}}

\newglossaryentry{ebit}{name={e-bit}, sort={e-bit},
  description={Paire de qubits maximalement intriqués, par exemple $\ket{\Phi^+}$, partagée entre deux parties ; ressource de la téléportation}}

\newglossaryentry{entropie}{name={entropie de Shannon}, sort={entropie de Shannon},
  description={$H(A)=-\sum_ap(a)\log_2p(a)$, incertitude moyenne (en bits) sur une variable aléatoire $A$}}

\newglossaryentry{vonneumann}{name={entropie de von Neumann}, sort={entropie de von Neumann},
  description={$S(\rho)=-\operatorname{Tr}(\rho\log_2\rho)=-\sum_i\lambda_i\log_2\lambda_i$ ; elle vaut $0$ pour un état pur}}

\newglossaryentry{hilbert}{name={Hilbert (espace de)}, sort={Hilbert},
  description={Espace vectoriel complexe muni d'un produit scalaire et complet ; c'est l'espace des états d'un système quantique}}

\newglossaryentry{etatmixte}{name={état mixte}, sort={etat mixte},
  description={État décrit par une matrice de densité avec $\operatorname{Tr}\rho^2<1$ : mélange statistique d'états purs}}

\newglossaryentry{etatpur}{name={état pur}, sort={etat pur},
  description={État décrit par un vecteur $\ket{\Psi}$, soit $\rho=\ket{\Psi}\bra{\Psi}$ et $\operatorname{Tr}\rho^2=1$}}

\newglossaryentry{stationnaire}{name={état stationnaire}, sort={etat stationnaire},
  description={Solution de l'équation de Schrödinger de la forme $\psi(x)\,\ee^{-\ii Et/\hbar}$ : énergie certaine et densité de probabilité indépendante du temps}}

\newglossaryentry{fonctiononde}{name={fonction d'onde}, sort={fonction d'onde},
  description={Fonction complexe $\Psi(x,t)$ décrivant l'état d'une particule ; $|\Psi|^2$ est une densité de probabilité de présence}}

\newglossaryentry{fredkin}{name={Fredkin (porte de)}, sort={Fredkin},
  description={Porte à trois qubits (C-SWAP) : échange les deux cibles si le contrôle vaut $1$}}

\newglossaryentry{hadamard}{name={Hadamard (porte de)}, sort={Hadamard},
  description={Porte $H=\frac{1}{\sqrt2}\begin{pmatrix}1&1\\1&-1\end{pmatrix}$ ; elle crée la superposition $\ket{+}$ à partir de $\ket{0}$ et change de base $Z\leftrightarrow X$}}

\newglossaryentry{hamiltonien}{name={hamiltonien}, sort={hamiltonien},
  description={Opérateur $\hat H=-\frac{\hbar^2}{2m}\partial_x^2+V(x)$ associé à l'énergie totale}}

\newglossaryentry{hermitien}{name={hermitien}, sort={hermitien},
  description={Se dit d'un opérateur égal à son adjoint, $A^\dagger=A$ ; ses valeurs propres sont réelles}}

\newglossaryentry{informationmutuelle}{name={information mutuelle}, sort={information mutuelle},
  description={$I(A;B)=H(A)+H(B)-H(A,B)$ : ce que la connaissance de $B$ apprend sur $A$}}

\newglossaryentry{intrication}{name={intrication}, sort={intrication},
  description={Propriété d'un état de plusieurs qubits qui n'est pas un produit d'états de chaque qubit}}

\newglossaryentry{ket}{name={ket}, sort={ket},
  description={Vecteur colonne $\ket{a}$ d'un espace de Hilbert ; il représente un état quantique}}

\newglossaryentry{mesure}{name={mesure}, sort={mesure},
  description={Opération qui donne l'un des résultats possibles avec la probabilité de Born et projette l'état sur l'état associé}}

\newglossaryentry{nonclonage}{name={non-clonage (théorème du)}, sort={non-clonage},
  description={Aucune opération ne copie un état quantique inconnu ; la téléportation déplace l'état sans le copier}}

\newglossaryentry{operateur}{name={opérateur}, sort={operateur},
  description={Application linéaire d'un espace de Hilbert dans lui-même ; en dimension finie, une matrice}}

\newglossaryentry{orthonormee}{name={orthonormée (base)}, sort={orthonormee},
  description={Base $\{\ket{e_i}\}$ de vecteurs de norme $1$ et deux à deux orthogonaux : $\braket{e_i}{e_j}=\delta_{ij}$}}

\newglossaryentry{pauli}{name={Pauli (matrices de)}, sort={Pauli},
  description={Les trois matrices $X=\sigma_x$, $Y=\sigma_y$, $Z=\sigma_z$ ; ce sont des portes à un qubit}}

\newglossaryentry{population}{name={population}, sort={population},
  description={Terme diagonal $\rho_{ii}$ de la matrice de densité : probabilité de mesurer l'état de base $\ket{i}$}}

\newglossaryentry{tenseur}{name={produit tensoriel}, sort={produit tensoriel},
  description={Construction $\mathcal{H}_A\otimes\mathcal{H}_B$ de l'espace de deux systèmes ; $\dim=\dim\mathcal{H}_A\times\dim\mathcal{H}_B$}}

\newglossaryentry{projecteur}{name={projecteur}, sort={projecteur},
  description={Opérateur $\hat P$ tel que $\hat P^2=\hat P$ et $\hat P^\dagger=\hat P$ ; par exemple $\hat P_i=\ket{i}\bra{i}$}}

\newglossaryentry{qubit}{name={qubit}, sort={qubit},
  description={Bit quantique : système à deux niveaux d'état $\alpha\ket{0}+\beta\ket{1}$, avec $|\alpha|^2+|\beta|^2=1$}}

\newglossaryentry{schrodinger}{name={Schrödinger (équation de)}, sort={Schrodinger},
  description={Équation $\ii\hbar\,\partial_t\Psi=\hat H\Psi$ qui gouverne l'évolution de la fonction d'onde}}

\newglossaryentry{separable}{name={séparable (état)}, sort={separable},
  description={État de deux qubits qui s'écrit $\ket{\Psi}_A\otimes\ket{\Psi}_B$ ; critère : $c_{00}c_{11}-c_{01}c_{10}=0$}}

\newglossaryentry{superposition}{name={superposition}, sort={superposition},
  description={Combinaison linéaire $\sum_ic_i\ket{i}$ d'états ; conséquence de la linéarité de l'équation de Schrödinger}}

\newglossaryentry{swap}{name={SWAP}, sort={SWAP},
  description={Porte à deux qubits qui échange leurs états : $\ket{a,b}\to\ket{b,a}$}}

\newglossaryentry{teleportation}{name={téléportation quantique}, sort={teleportation},
  description={Protocole qui transmet un état de qubit inconnu avec un e-bit et deux bits classiques, sans déplacer le qubit}}

\newglossaryentry{toffoli}{name={Toffoli (porte de)}, sort={Toffoli},
  description={Porte à trois qubits (CCNOT) : la cible est inversée si les deux contrôles valent $1$ ; $\ket{x,y,z}\to\ket{x,y,z\oplus xy}$}}

\newglossaryentry{trace}{name={trace (partielle)}, sort={trace},
  description={$\operatorname{Tr}_B$ : opération qui «~oublie~» le sous-système $B$ et donne l'état $\rho_A$ du sous-système $A$}}

\newglossaryentry{unitaire}{name={unitaire}, sort={unitaire},
  description={Matrice $U$ telle que $U^\dagger U=\mathrm{Id}$ ; elle conserve la norme et est réversible ; c'est une porte quantique}}

\newglossaryentry{observable}{name={observable}, sort={observable},
  description={Opérateur hermitien qui représente une grandeur mesurable ; ses valeurs propres sont les résultats possibles d'une mesure}}

\newglossaryentry{effondrement}{name={effondrement (de l'état)}, sort={effondrement},
  description={Changement brusque de l'état lors d'une mesure : il devient l'état propre associé au résultat obtenu}}

\newglossaryentry{valeurpropre}{name={valeur propre}, sort={valeur propre},
  description={Nombre $\lambda$ tel que $A\ket{v}=\lambda\ket{v}$ avec $\ket{v}\neq0$ (vecteur propre) ; pour une observable, résultat possible d'une mesure}}

\newglossaryentry{phase}{name={phase (relative)}, sort={phase},
  description={Angle $\varphi$ entre les amplitudes de $\ket{0}$ et $\ket{1}$ ; observable, contrairement à la phase globale}}

\newglossaryentry{purete}{name={pureté}, sort={purete},
  description={Nombre $\operatorname{Tr}\rho^2$, compris entre $\frac1d$ et $1$ (dimension $d$) ; il vaut $1$ si et seulement si l'état est pur}}

\newglossaryentry{porte}{name={porte quantique}, sort={porte quantique},
  description={Opération unitaire appliquée à un ou plusieurs qubits ; dessinée comme une boîte sur un circuit}}
 % charge les entrées depuis glossaire.tex

%%%%% Pour les corner images : logo INSA, titre du document, filet rouge %%%%%
\usepackage{eso-pic} % for corner picutres
\usepackage{ifthen}
\AddToShipoutPictureBG{%
 \ifthenelse{\value{page}>\value{pagesdegarde}}{%
  % logo INSA en haut à gauche
  \put(0.5cm,\paperheight-1.5cm){\includegraphics[width=6cm]{images/logo_insa.png}}%
  % titre du document en haut à droite
  \put(\paperwidth-0.5cm,\paperheight-0.9cm){\makebox(0,0)[r]{%
    \normalfont\normalsize% (évite d'hériter du gras ou de la taille du texte en cours de page)
    \renewcommand{\arraystretch}{1}% (évite d'hériter d'un \arraystretch réglé dans un tableau en cours de page)
    \begin{tabular}{@{}r@{}}
      {\bfseries\color{insarouge}\doctitre}\\
      {\footnotesize\color{insagris}\docsoustitre}
    \end{tabular}}}%
  % filet rouge sous l'en-tête
  \put(0.5cm,\paperheight-1.8cm){\color{insarouge}\rule{\dimexpr\paperwidth-1cm\relax}{1pt}}%
 }{}%
}

%%%%% Numéro de page (rouge, centré en bas) %%%%%
\usepackage{fancyhdr}
\pagestyle{fancy}
\fancyhf{}
\renewcommand{\headrulewidth}{0pt}
\fancyfoot[C]{\small\bfseries\textcolor{insarouge}{\thepage}}
\fancypagestyle{plain}{%
  \fancyhf{}%
  \renewcommand{\headrulewidth}{0pt}%
  \fancyfoot[C]{\small\bfseries\textcolor{insarouge}{\thepage}}%
}

%%%%% Coupures de ligne : évite les débordements dans les marges %%%%%
\setlength{\emergencystretch}{1.5em}

%%%%% Marges %%%%%
\usepackage[a4paper,
            top=3.5cm,
            bottom=2.5cm,
            left=2.2cm,
            right=2.2cm,
            headheight=14pt]{geometry}

%%%%% Import de pages PDF (par exemple une page de garde) %%%%%
% Import de pages PDF, par exemple une page de garde :
%   \includepdf[pages={1,2}]{page1.pdf}   (puis régler le compteur pagesdegarde en haut du fichier)
\usepackage{pdfpages}

\title{\doctitre}
\author{\docauteur}
\date{\docdate}

% QCM (chapitre 7) : bascules. Valeurs par défaut : tout est affiché.
% Décommenter une ligne pour changer le comportement de qcm.tex.
\newif\ifqcmrefs    \qcmrefstrue      % références « Chap. · § » sous chaque énoncé
\newif\ifqcmcorrige \qcmcorrigetrue   % corrigé (clé, explications, pièges, fiche de score)
%\qcmrefsfalse      % masque les références sous les énoncés (le corrigé les garde)
%\qcmcorrigefalse   % supprime le corrigé : document avec énoncés seuls

% Pour ne compiler que certains chapitres (la numérotation est conservée) :
%\includeonly{chapitre1}

\begin{document}

%%%%% Page de titre : page 1 du PDF, sans en-tête ni numéro (le sommaire est la page 2) %%%%%
% =====================================================================
%  Page de titre : importée par main.tex avec % =====================================================================
%  Page de titre : importée par main.tex avec % =====================================================================
%  Page de titre : importée par main.tex avec \input{pagetitre}.
%  Elle utilise \doctitre, \docsoustitre, \doctagline, \docauteur, \docprof,
%  \docecole et \docdate, définis dans le bloc « Informations du document »
%  de main.tex. Cette page est la page 1 du document : elle est comptée dans la
%  numérotation, sans décalage entre le numéro affiché et la page du PDF. Elle
%  n'a ni numéro (\thispagestyle{empty}) ni en-tête (compteur pagesdegarde de
%  main.tex). Le \setcounter qui suit \end{titlepage} donne le numéro 2 au sommaire,
%  car l'environnement titlepage remet sinon le compteur à 1.
%  (TikZ « remember picture » : deux passes de compilation.)
% =====================================================================
\begin{titlepage}
\thispagestyle{empty}
\setcounter{page}{1}
\begin{tikzpicture}[remember picture, overlay]
  % logo en haut à gauche et filet rouge dessous (mêmes positions que l'en-tête des autres pages)
  \node[anchor=south west, inner sep=0]
    at ($(current page.north west)+(0.5cm,-1.5cm)$)
    {\includegraphics[width=6cm]{images/logo_insa.png}};
  \draw[insarouge, line width=1pt]
    ($(current page.north west)+(0.5cm,-1.8cm)$) -- ($(current page.north east)+(-0.5cm,-1.8cm)$);

  % titre du document sur un bandeau rouge
  \fill[insarouge]
    ($(current page.north west)+(0,-6.6cm)$) rectangle ($(current page.north east)+(0,-11.4cm)$);
  \node[text=white, align=center, text width=17cm, font=\fontsize{38}{44}\selectfont\bfseries]
    at ($(current page.north)+(0,-8.3cm)$) {\doctitre};
  \node[text=white, align=center, text width=17cm, font=\Large]
    at ($(current page.north)+(0,-10.1cm)$) {\docsoustitre};

  % accroche et cours
  \node[text=insagris, align=center, text width=16cm, font=\Large\itshape]
    at ($(current page.north)+(0,-13.0cm)$) {\doctagline};
  \node[text=insagris, align=center, text width=16cm, font=\large]
    at ($(current page.north)+(0,-14.4cm)$) {Cours de \docprof{} -- \docecole};

  % motif : de l'onde au circuit quantique
  \begin{scope}[shift={($(current page.north)+(0,-18.4cm)$)}, >=Stealth, font=\small]
    % paquet d'ondes (mécanique ondulatoire)
    \draw[insagris, thick, domain=-7.6:-1.6, samples=240]
      plot (\x, {0.65*sin(\x*115)*exp(-((\x+4.6)^2)/2.2)});
    \draw[insagris!60, dashed, domain=-7.6:-1.6, samples=80]
      plot (\x, {0.65*exp(-((\x+4.6)^2)/2.2)});
    \draw[insagris!60, dashed, domain=-7.6:-1.6, samples=80]
      plot (\x, {-0.65*exp(-((\x+4.6)^2)/2.2)});
    \node[text=insagris, font=\footnotesize] at (-4.6,-1.15) {$\Psi(x,t)$};
    % flèche
    \draw[->, insarouge, very thick] (-1.1,0) -- (0.3,0);
    % circuit à deux qubits : H, CNOT, mesure
    \draw (0.8,0.45) -- (7.6,0.45);
    \draw (0.8,-0.45) -- (7.6,-0.45);
    \node[anchor=east, font=\footnotesize] at (0.8,0.45) {$\ket{0}$};
    \node[anchor=east, font=\footnotesize] at (0.8,-0.45) {$\ket{0}$};
    \node[porte] at (2.4,0.45) {$H$};
    \draw[insarouge, thick] (4.2,0.45) -- (4.2,-0.45);
    \pic at (4.2,0.45) {ctl};
    \pic at (4.2,-0.45) {cible};
    \pic at (6.2,0.45) {mesure};
    \pic at (6.2,-0.45) {mesure};
    \node[text=insagris, font=\footnotesize] at (4.2,-1.15) {$\ket{\Phi^+}=\tfrac{1}{\sqrt2}(\ket{00}+\ket{11})$};
  \end{scope}

  % auteur et date
  \node[align=center, font=\Large]
    at ($(current page.south)+(0,6.4cm)$) {\textbf{\docauteur}};
  \draw[insarouge, line width=1pt]
    ($(current page.south)+(-3cm,5.3cm)$) -- ($(current page.south)+(3cm,5.3cm)$);
  \node[text=insagris, align=center, font=\large]
    at ($(current page.south)+(0,4.4cm)$) {\docdate};
\end{tikzpicture}
\end{titlepage}
\setcounter{page}{2}
.
%  Elle utilise \doctitre, \docsoustitre, \doctagline, \docauteur, \docprof,
%  \docecole et \docdate, définis dans le bloc « Informations du document »
%  de main.tex. Cette page est la page 1 du document : elle est comptée dans la
%  numérotation, sans décalage entre le numéro affiché et la page du PDF. Elle
%  n'a ni numéro (\thispagestyle{empty}) ni en-tête (compteur pagesdegarde de
%  main.tex). Le \setcounter qui suit \end{titlepage} donne le numéro 2 au sommaire,
%  car l'environnement titlepage remet sinon le compteur à 1.
%  (TikZ « remember picture » : deux passes de compilation.)
% =====================================================================
\begin{titlepage}
\thispagestyle{empty}
\setcounter{page}{1}
\begin{tikzpicture}[remember picture, overlay]
  % logo en haut à gauche et filet rouge dessous (mêmes positions que l'en-tête des autres pages)
  \node[anchor=south west, inner sep=0]
    at ($(current page.north west)+(0.5cm,-1.5cm)$)
    {\includegraphics[width=6cm]{images/logo_insa.png}};
  \draw[insarouge, line width=1pt]
    ($(current page.north west)+(0.5cm,-1.8cm)$) -- ($(current page.north east)+(-0.5cm,-1.8cm)$);

  % titre du document sur un bandeau rouge
  \fill[insarouge]
    ($(current page.north west)+(0,-6.6cm)$) rectangle ($(current page.north east)+(0,-11.4cm)$);
  \node[text=white, align=center, text width=17cm, font=\fontsize{38}{44}\selectfont\bfseries]
    at ($(current page.north)+(0,-8.3cm)$) {\doctitre};
  \node[text=white, align=center, text width=17cm, font=\Large]
    at ($(current page.north)+(0,-10.1cm)$) {\docsoustitre};

  % accroche et cours
  \node[text=insagris, align=center, text width=16cm, font=\Large\itshape]
    at ($(current page.north)+(0,-13.0cm)$) {\doctagline};
  \node[text=insagris, align=center, text width=16cm, font=\large]
    at ($(current page.north)+(0,-14.4cm)$) {Cours de \docprof{} -- \docecole};

  % motif : de l'onde au circuit quantique
  \begin{scope}[shift={($(current page.north)+(0,-18.4cm)$)}, >=Stealth, font=\small]
    % paquet d'ondes (mécanique ondulatoire)
    \draw[insagris, thick, domain=-7.6:-1.6, samples=240]
      plot (\x, {0.65*sin(\x*115)*exp(-((\x+4.6)^2)/2.2)});
    \draw[insagris!60, dashed, domain=-7.6:-1.6, samples=80]
      plot (\x, {0.65*exp(-((\x+4.6)^2)/2.2)});
    \draw[insagris!60, dashed, domain=-7.6:-1.6, samples=80]
      plot (\x, {-0.65*exp(-((\x+4.6)^2)/2.2)});
    \node[text=insagris, font=\footnotesize] at (-4.6,-1.15) {$\Psi(x,t)$};
    % flèche
    \draw[->, insarouge, very thick] (-1.1,0) -- (0.3,0);
    % circuit à deux qubits : H, CNOT, mesure
    \draw (0.8,0.45) -- (7.6,0.45);
    \draw (0.8,-0.45) -- (7.6,-0.45);
    \node[anchor=east, font=\footnotesize] at (0.8,0.45) {$\ket{0}$};
    \node[anchor=east, font=\footnotesize] at (0.8,-0.45) {$\ket{0}$};
    \node[porte] at (2.4,0.45) {$H$};
    \draw[insarouge, thick] (4.2,0.45) -- (4.2,-0.45);
    \pic at (4.2,0.45) {ctl};
    \pic at (4.2,-0.45) {cible};
    \pic at (6.2,0.45) {mesure};
    \pic at (6.2,-0.45) {mesure};
    \node[text=insagris, font=\footnotesize] at (4.2,-1.15) {$\ket{\Phi^+}=\tfrac{1}{\sqrt2}(\ket{00}+\ket{11})$};
  \end{scope}

  % auteur et date
  \node[align=center, font=\Large]
    at ($(current page.south)+(0,6.4cm)$) {\textbf{\docauteur}};
  \draw[insarouge, line width=1pt]
    ($(current page.south)+(-3cm,5.3cm)$) -- ($(current page.south)+(3cm,5.3cm)$);
  \node[text=insagris, align=center, font=\large]
    at ($(current page.south)+(0,4.4cm)$) {\docdate};
\end{tikzpicture}
\end{titlepage}
\setcounter{page}{2}
.
%  Elle utilise \doctitre, \docsoustitre, \doctagline, \docauteur, \docprof,
%  \docecole et \docdate, définis dans le bloc « Informations du document »
%  de main.tex. Cette page est la page 1 du document : elle est comptée dans la
%  numérotation, sans décalage entre le numéro affiché et la page du PDF. Elle
%  n'a ni numéro (\thispagestyle{empty}) ni en-tête (compteur pagesdegarde de
%  main.tex). Le \setcounter qui suit \end{titlepage} donne le numéro 2 au sommaire,
%  car l'environnement titlepage remet sinon le compteur à 1.
%  (TikZ « remember picture » : deux passes de compilation.)
% =====================================================================
\begin{titlepage}
\thispagestyle{empty}
\setcounter{page}{1}
\begin{tikzpicture}[remember picture, overlay]
  % logo en haut à gauche et filet rouge dessous (mêmes positions que l'en-tête des autres pages)
  \node[anchor=south west, inner sep=0]
    at ($(current page.north west)+(0.5cm,-1.5cm)$)
    {\includegraphics[width=6cm]{images/logo_insa.png}};
  \draw[insarouge, line width=1pt]
    ($(current page.north west)+(0.5cm,-1.8cm)$) -- ($(current page.north east)+(-0.5cm,-1.8cm)$);

  % titre du document sur un bandeau rouge
  \fill[insarouge]
    ($(current page.north west)+(0,-6.6cm)$) rectangle ($(current page.north east)+(0,-11.4cm)$);
  \node[text=white, align=center, text width=17cm, font=\fontsize{38}{44}\selectfont\bfseries]
    at ($(current page.north)+(0,-8.3cm)$) {\doctitre};
  \node[text=white, align=center, text width=17cm, font=\Large]
    at ($(current page.north)+(0,-10.1cm)$) {\docsoustitre};

  % accroche et cours
  \node[text=insagris, align=center, text width=16cm, font=\Large\itshape]
    at ($(current page.north)+(0,-13.0cm)$) {\doctagline};
  \node[text=insagris, align=center, text width=16cm, font=\large]
    at ($(current page.north)+(0,-14.4cm)$) {Cours de \docprof{} -- \docecole};

  % motif : de l'onde au circuit quantique
  \begin{scope}[shift={($(current page.north)+(0,-18.4cm)$)}, >=Stealth, font=\small]
    % paquet d'ondes (mécanique ondulatoire)
    \draw[insagris, thick, domain=-7.6:-1.6, samples=240]
      plot (\x, {0.65*sin(\x*115)*exp(-((\x+4.6)^2)/2.2)});
    \draw[insagris!60, dashed, domain=-7.6:-1.6, samples=80]
      plot (\x, {0.65*exp(-((\x+4.6)^2)/2.2)});
    \draw[insagris!60, dashed, domain=-7.6:-1.6, samples=80]
      plot (\x, {-0.65*exp(-((\x+4.6)^2)/2.2)});
    \node[text=insagris, font=\footnotesize] at (-4.6,-1.15) {$\Psi(x,t)$};
    % flèche
    \draw[->, insarouge, very thick] (-1.1,0) -- (0.3,0);
    % circuit à deux qubits : H, CNOT, mesure
    \draw (0.8,0.45) -- (7.6,0.45);
    \draw (0.8,-0.45) -- (7.6,-0.45);
    \node[anchor=east, font=\footnotesize] at (0.8,0.45) {$\ket{0}$};
    \node[anchor=east, font=\footnotesize] at (0.8,-0.45) {$\ket{0}$};
    \node[porte] at (2.4,0.45) {$H$};
    \draw[insarouge, thick] (4.2,0.45) -- (4.2,-0.45);
    \pic at (4.2,0.45) {ctl};
    \pic at (4.2,-0.45) {cible};
    \pic at (6.2,0.45) {mesure};
    \pic at (6.2,-0.45) {mesure};
    \node[text=insagris, font=\footnotesize] at (4.2,-1.15) {$\ket{\Phi^+}=\tfrac{1}{\sqrt2}(\ket{00}+\ket{11})$};
  \end{scope}

  % auteur et date
  \node[align=center, font=\Large]
    at ($(current page.south)+(0,6.4cm)$) {\textbf{\docauteur}};
  \draw[insarouge, line width=1pt]
    ($(current page.south)+(-3cm,5.3cm)$) -- ($(current page.south)+(3cm,5.3cm)$);
  \node[text=insagris, align=center, font=\large]
    at ($(current page.south)+(0,4.4cm)$) {\docdate};
\end{tikzpicture}
\end{titlepage}
\setcounter{page}{2}


%%%%% Sommaire %%%%%
\renewcommand{\contentsname}{Sommaire}
\setcounter{tocdepth}{2}
\begingroup
\fontsize{8pt}{9.4pt}\selectfont
\tableofcontents
\endgroup

%%%%% Résumés %%%%%
%\newpage
%\section*{Résumés}
%\addcontentsline{toc}{section}{Résumés}
%\input{sections/0_summary}

%%%%% Introduction (section non numérotée, ajoutée au sommaire) %%%%%
% =====================================================================
%  Introduction : importée par main.tex avec % =====================================================================
%  Introduction : importée par main.tex avec % =====================================================================
%  Introduction : importée par main.tex avec \include{introduction}.
%  Section non numérotée, ajoutée au sommaire. Pas de figure flottante ici
%  (la numérotation des figures suit celle des chapitres) : le plan est un
%  dessin TikZ en place.
% =====================================================================

% Macros de Dirac (déjà définies dans main.tex ; ce bloc évite l'erreur
% « Undefined control sequence \ket » si main.tex est une ancienne version).
\providecommand{\ket}[1]{|#1\rangle}
\providecommand{\bra}[1]{\langle #1|}
\providecommand{\braket}[2]{\langle #1|#2\rangle}

\phantomsection
\section*{Introduction}
\addcontentsline{toc}{section}{Introduction}

Ce document rassemble les notes du cours d'informatique quantique de \docprof{}
(4A FISE SF). Il part de la mécanique ondulatoire (fonction d'onde de Schrödinger,
interprétation de Born) et en tire, pas à pas, les outils de l'informatique
quantique : états d'un \gls{qubit}, \glslink{hilbert}{espace de Hilbert}, opérateurs,
matrice de densité, portes quantiques, systèmes de plusieurs qubits,
\gls{intrication} et téléportation. Tous les calculs sont détaillés.

\begin{center}
\begin{tikzpicture}[>=Stealth, font=\scriptsize,
    boite/.style={draw=insarouge, fill=insarouge!6, thick, rounded corners=2pt,
                  text width=2.02cm, minimum height=1.55cm, align=center},
    num/.style={circle, fill=insarouge, text=white, font=\small\bfseries, minimum size=5.5mm, inner sep=0pt}]
  \foreach \i/\txt/\idee in {
      1/{Histoire et mécanique quantique}/{onde de matière},
      2/{Notation de Dirac et espace de Hilbert}/{états = vecteurs},
      3/{Opérateurs, matrice de densité, sphère de Bloch, produit tensoriel}/{états mixtes, Bloch},
      4/{Portes quantiques et systèmes multi-qubits}/{portes unitaires},
      5/{États quantiques, intrication, états de Bell}/{corrélations, téléportation},
      6/{Exemples et exercices corrigés}/{s'entraîner},
      7/{QCM sur le cours}/{s'évaluer}} {
    \node[boite] (b\i) at ({(\i-1)*2.36},0) {\txt};
    \node[num] at ({(\i-1)*2.36},1.2) {\i};
    \node[text=insagris, font=\scriptsize\itshape, text width=2.1cm, align=center] at ({(\i-1)*2.36},-1.25) {\idee};
  }
  \foreach \i/\j in {1/2, 2/3, 3/4, 4/5, 5/6, 6/7}
    \draw[->, insarouge, thick] (b\i.east) -- (b\j.west);
\end{tikzpicture}
\end{center}

Chaque chapitre se termine par un \emph{récapitulatif} : schéma de synthèse, points à retenir,
tableau de formules renvoyant aux sections, et liste de savoir-faire. Le chapitre~\ref{chap:exercices}
rassemble des exemples et exercices corrigés, et le chapitre~\ref{chap:qcm} un QCM de 100 questions
avec corrigé détaillé. Les termes techniques sont définis dans
le glossaire, en fin de document ; les ouvrages de référence figurent dans la bibliographie
\parencite[par exemple][]{nielsen2010,cohen1973}.

\paragraph{Prérequis et conventions.} Algèbre linéaire (matrices, valeurs propres),
nombres complexes, notions d'ondes. Le symbole $\ii$ désigne l'unité imaginaire, $\ket{0}$
et $\ket{1}$ la base de calcul d'un qubit, et $\hbar=h/2\pi$.
.
%  Section non numérotée, ajoutée au sommaire. Pas de figure flottante ici
%  (la numérotation des figures suit celle des chapitres) : le plan est un
%  dessin TikZ en place.
% =====================================================================

% Macros de Dirac (déjà définies dans main.tex ; ce bloc évite l'erreur
% « Undefined control sequence \ket » si main.tex est une ancienne version).
\providecommand{\ket}[1]{|#1\rangle}
\providecommand{\bra}[1]{\langle #1|}
\providecommand{\braket}[2]{\langle #1|#2\rangle}

\phantomsection
\section*{Introduction}
\addcontentsline{toc}{section}{Introduction}

Ce document rassemble les notes du cours d'informatique quantique de \docprof{}
(4A FISE SF). Il part de la mécanique ondulatoire (fonction d'onde de Schrödinger,
interprétation de Born) et en tire, pas à pas, les outils de l'informatique
quantique : états d'un \gls{qubit}, \glslink{hilbert}{espace de Hilbert}, opérateurs,
matrice de densité, portes quantiques, systèmes de plusieurs qubits,
\gls{intrication} et téléportation. Tous les calculs sont détaillés.

\begin{center}
\begin{tikzpicture}[>=Stealth, font=\scriptsize,
    boite/.style={draw=insarouge, fill=insarouge!6, thick, rounded corners=2pt,
                  text width=2.02cm, minimum height=1.55cm, align=center},
    num/.style={circle, fill=insarouge, text=white, font=\small\bfseries, minimum size=5.5mm, inner sep=0pt}]
  \foreach \i/\txt/\idee in {
      1/{Histoire et mécanique quantique}/{onde de matière},
      2/{Notation de Dirac et espace de Hilbert}/{états = vecteurs},
      3/{Opérateurs, matrice de densité, sphère de Bloch, produit tensoriel}/{états mixtes, Bloch},
      4/{Portes quantiques et systèmes multi-qubits}/{portes unitaires},
      5/{États quantiques, intrication, états de Bell}/{corrélations, téléportation},
      6/{Exemples et exercices corrigés}/{s'entraîner},
      7/{QCM sur le cours}/{s'évaluer}} {
    \node[boite] (b\i) at ({(\i-1)*2.36},0) {\txt};
    \node[num] at ({(\i-1)*2.36},1.2) {\i};
    \node[text=insagris, font=\scriptsize\itshape, text width=2.1cm, align=center] at ({(\i-1)*2.36},-1.25) {\idee};
  }
  \foreach \i/\j in {1/2, 2/3, 3/4, 4/5, 5/6, 6/7}
    \draw[->, insarouge, thick] (b\i.east) -- (b\j.west);
\end{tikzpicture}
\end{center}

Chaque chapitre se termine par un \emph{récapitulatif} : schéma de synthèse, points à retenir,
tableau de formules renvoyant aux sections, et liste de savoir-faire. Le chapitre~\ref{chap:exercices}
rassemble des exemples et exercices corrigés, et le chapitre~\ref{chap:qcm} un QCM de 100 questions
avec corrigé détaillé. Les termes techniques sont définis dans
le glossaire, en fin de document ; les ouvrages de référence figurent dans la bibliographie
\parencite[par exemple][]{nielsen2010,cohen1973}.

\paragraph{Prérequis et conventions.} Algèbre linéaire (matrices, valeurs propres),
nombres complexes, notions d'ondes. Le symbole $\ii$ désigne l'unité imaginaire, $\ket{0}$
et $\ket{1}$ la base de calcul d'un qubit, et $\hbar=h/2\pi$.
.
%  Section non numérotée, ajoutée au sommaire. Pas de figure flottante ici
%  (la numérotation des figures suit celle des chapitres) : le plan est un
%  dessin TikZ en place.
% =====================================================================

% Macros de Dirac (déjà définies dans main.tex ; ce bloc évite l'erreur
% « Undefined control sequence \ket » si main.tex est une ancienne version).
\providecommand{\ket}[1]{|#1\rangle}
\providecommand{\bra}[1]{\langle #1|}
\providecommand{\braket}[2]{\langle #1|#2\rangle}

\phantomsection
\section*{Introduction}
\addcontentsline{toc}{section}{Introduction}

Ce document rassemble les notes du cours d'informatique quantique de \docprof{}
(4A FISE SF). Il part de la mécanique ondulatoire (fonction d'onde de Schrödinger,
interprétation de Born) et en tire, pas à pas, les outils de l'informatique
quantique : états d'un \gls{qubit}, \glslink{hilbert}{espace de Hilbert}, opérateurs,
matrice de densité, portes quantiques, systèmes de plusieurs qubits,
\gls{intrication} et téléportation. Tous les calculs sont détaillés.

\begin{center}
\begin{tikzpicture}[>=Stealth, font=\scriptsize,
    boite/.style={draw=insarouge, fill=insarouge!6, thick, rounded corners=2pt,
                  text width=2.02cm, minimum height=1.55cm, align=center},
    num/.style={circle, fill=insarouge, text=white, font=\small\bfseries, minimum size=5.5mm, inner sep=0pt}]
  \foreach \i/\txt/\idee in {
      1/{Histoire et mécanique quantique}/{onde de matière},
      2/{Notation de Dirac et espace de Hilbert}/{états = vecteurs},
      3/{Opérateurs, matrice de densité, sphère de Bloch, produit tensoriel}/{états mixtes, Bloch},
      4/{Portes quantiques et systèmes multi-qubits}/{portes unitaires},
      5/{États quantiques, intrication, états de Bell}/{corrélations, téléportation},
      6/{Exemples et exercices corrigés}/{s'entraîner},
      7/{QCM sur le cours}/{s'évaluer}} {
    \node[boite] (b\i) at ({(\i-1)*2.36},0) {\txt};
    \node[num] at ({(\i-1)*2.36},1.2) {\i};
    \node[text=insagris, font=\scriptsize\itshape, text width=2.1cm, align=center] at ({(\i-1)*2.36},-1.25) {\idee};
  }
  \foreach \i/\j in {1/2, 2/3, 3/4, 4/5, 5/6, 6/7}
    \draw[->, insarouge, thick] (b\i.east) -- (b\j.west);
\end{tikzpicture}
\end{center}

Chaque chapitre se termine par un \emph{récapitulatif} : schéma de synthèse, points à retenir,
tableau de formules renvoyant aux sections, et liste de savoir-faire. Le chapitre~\ref{chap:exercices}
rassemble des exemples et exercices corrigés, et le chapitre~\ref{chap:qcm} un QCM de 100 questions
avec corrigé détaillé. Les termes techniques sont définis dans
le glossaire, en fin de document ; les ouvrages de référence figurent dans la bibliographie
\parencite[par exemple][]{nielsen2010,cohen1973}.

\paragraph{Prérequis et conventions.} Algèbre linéaire (matrices, valeurs propres),
nombres complexes, notions d'ondes. Le symbole $\ii$ désigne l'unité imaginaire, $\ket{0}$
et $\ket{1}$ la base de calcul d'un qubit, et $\hbar=h/2\pi$.


%%%%% Chapitres (un fichier par chapitre, chacun commence par \section) %%%%%
% =====================================================================
%  Chapitre 1 : Fondements physiques de la mécanique quantique
%  Fichier importé par main.tex avec % =====================================================================
%  Chapitre 1 : Fondements physiques de la mécanique quantique
%  Fichier importé par main.tex avec % =====================================================================
%  Chapitre 1 : Fondements physiques de la mécanique quantique
%  Fichier importé par main.tex avec \include{chapitre1}.
%  Il ne contient ni préambule ni \begin{document} : les environnements
%  (definition, resultat, remarque), les macros (\ii, \ee, \dd, \pd, \pdd,
%  \Hop) et les couleurs (insarouge, insagris) sont définis dans main.tex.
% =====================================================================

% =====================================================================
\section{Fondements physiques de la mécanique quantique}\label{sec:fondements}
% =====================================================================

\begin{remarque}[Plan général du cours]
\begin{enumerate}[nosep]
  \item Mécanique quantique (ce chapitre) ;
  \item Espace de Hilbert $\mathcal{H}$ ;
  \item Notation de Dirac ;
  \item Espace produit tensoriel ;
  \item États quantiques, intrication quantique, états de Bell.
\end{enumerate}
\end{remarque}

% ---------------------------------------------------------------------
\subsection{Dualité onde--corpuscule}
% ---------------------------------------------------------------------

\subsubsection{Einstein : la lumière est faite de photons}

À la fin du \textsc{xix}\textsuperscript{e} siècle, l'effet photo-électrique
(émission d'électrons par un métal éclairé) contredit la théorie purement
ondulatoire de la lumière : l'émission dépend de la \emph{fréquence} de la
lumière, et non de son intensité. En 1905, Einstein l'explique en supposant que
la lumière est formée de quanta d'énergie, les \emph{photons}. Un électron n'est
éjecté que si l'énergie du photon dépasse le travail d'extraction $W$ du métal :
$E_{c,\max} = hf - W$. Une onde lumineuse se comporte donc aussi comme une
particule.

\begin{definition}{Photon (Planck--Einstein)}
\[
  E = hf = h\nu = \hbar\omega ,
  \qquad
  p = \frac{h}{\lambda} = \hbar k ,
\]
avec $h \approx 6{,}63\times10^{-34}\ \mathrm{J\cdot s}$ (constante de Planck),
$\hbar = \dfrac{h}{2\pi}$, $\omega = 2\pi f$ (pulsation) et
$k = \dfrac{2\pi}{\lambda}$ (nombre d'onde).
\end{definition}

\subsubsection{de Broglie : la matière est une onde}

En 1924, de Broglie inverse le raisonnement : si une onde peut se comporter
comme une particule, une \emph{particule matérielle} peut se comporter comme une
onde. À une particule de masse $m$ et de vitesse $v$, donc de quantité de
mouvement $p = mv$, on associe une \emph{onde de matière}.

\begin{definition}{Relations de de Broglie}
\[
  \lambda = \frac{h}{p} = \frac{h}{mv} ,
  \qquad
  f = \frac{E}{h}
  \qquad\Longleftrightarrow\qquad
  p = \hbar k ,\quad E = \hbar\omega .
\]
\end{definition}

\begin{figure}[!ht]
\centering
\begin{tikzpicture}[>=Stealth, thick]
% --- Ligne 1 : Einstein (onde -> particule)
\begin{scope}[yshift=2.3cm]
  \draw[insarouge, domain=0:2.4, samples=100] plot (\x, {0.35*sin(360*\x/0.8)});
  \node at (1.2,-0.8) {\small onde lumineuse};
  \draw[<->] (3.0,0) -- (5.4,0)
    node[midway, above] {\small Einstein}
    node[midway, below] {\small $E = hf = \hbar\omega$};
  \shade[ball color=insarouge!85] (6.8,0) circle (0.3);
  \node at (6.8,0) {\small $\gamma$};
  \node at (6.8,-0.8) {\small photon};
\end{scope}
% --- Ligne 2 : de Broglie (particule -> onde)
\begin{scope}
  \shade[ball color=insagris!80] (0.3,0) circle (0.3);
  \draw[->] (0.7,0) -- (2.0,0) node[midway, above] {\small $p = mv$};
  \node at (1.0,-0.8) {\small particule};
  \draw[<->] (3.0,0) -- (5.4,0)
    node[midway, above] {\small de Broglie}
    node[midway, below] {\small $\lambda = h/p$};
  \draw[insagris, domain=0:2.4, samples=100, shift={(6.0,0)}]
    plot (\x, {0.35*sin(360*\x/0.8)});
  \node at (7.2,-0.8) {\small onde de matière};
\end{scope}
\end{tikzpicture}
\caption{Dualité onde--corpuscule : Einstein associe une particule à l'onde
lumineuse, de Broglie associe une onde à la particule matérielle.}
\end{figure}

\begin{remarque}[Ordres de grandeur]
\begin{itemize}[nosep]
  \item Balle de $0{,}1$ kg à $10$ m/s : $\lambda = \dfrac{6{,}63\times10^{-34}}{0{,}1\times 10}
        \approx 6{,}6\times10^{-34}$ m, inobservable.
  \item Électron accéléré sous $100$ V : $p=\sqrt{2m_e\,eU}$ donne
        $\lambda \approx 0{,}12$ nm, de l'ordre de la distance entre atomes d'un cristal.
        La diffraction d'électrons par un cristal (Davisson--Germer, 1927) confirme
        l'existence de l'onde de matière.
\end{itemize}
\end{remarque}

\subsubsection{Le problème posé à Schrödinger}

Si toute particule est associée à une onde, existe-t-il une \emph{équation
d'onde} qui gouverne cette onde de matière ? C'est la question à laquelle
répond l'équation de Schrödinger (1926).

% ---------------------------------------------------------------------
\subsection{Rappel : l'onde classique}
% ---------------------------------------------------------------------

Une onde plane progressive sinusoïdale se propageant selon $x$ s'écrit
\[
  y(x,t) = A\cos(kx - \omega t),
\]
où $A$ est l'amplitude, $k = 2\pi/\lambda$ le nombre d'onde, $\omega = 2\pi f$ la
pulsation et $v = \omega/k = \lambda f$ la vitesse de phase (un maximum vérifie
$kx-\omega t = \text{cste}$, donc $\dd x/\dd t = \omega/k$).

\begin{figure}[!ht]
\centering
\begin{tikzpicture}[>=Stealth]
  \draw[->] (-0.3,0) -- (8.7,0) node[right] {$x$};
  \draw[->] (0,-1.6) -- (0,2.0) node[above] {$y$};
  \draw[thick, insagris, domain=0:8.4, samples=200]
    plot (\x, {cos(360*\x/4)});
  \draw[thick, dashed, insarouge, domain=0:8.4, samples=200]
    plot (\x, {cos(360*(\x-0.9)/4)});
  \draw[dotted] (0,1) -- (4,1);
  \node[left] at (0,1) {$A$};
  \node[left] at (0,-1) {$-A$};
  \draw[dotted] (4,1) -- (4,1.55);
  \draw[<->] (0,1.45) -- (4,1.45) node[midway, above] {$\lambda$};
  \draw[->, insarouge] (6,-1.2) -- (6.9,-1.2)
    node[midway, below] {\small $v\,\Delta t$};
  \draw[thick, insagris] (6.0,1.75) -- (6.6,1.75)
    node[right, black] {\small $t$};
  \draw[thick, dashed, insarouge] (7.4,1.75) -- (8.0,1.75)
    node[right, black] {\small $t+\Delta t$};
\end{tikzpicture}
\caption{Onde plane progressive $y(x,t)=A\cos(kx-\omega t)$ à deux instants :
le profil se translate de $v\,\Delta t$.}
\end{figure}

Elle vérifie l'\emph{équation des ondes}
\[
  \pdd{y}{x} = \frac{1}{v^2}\,\pdd{y}{t}
  \qquad\text{(courbure $\leftrightarrow$ accélération).}
\]
En effet,
\[
  \pdd{y}{x} = -Ak^2\cos(kx-\omega t) = -k^2\,y ,
  \qquad
  \pdd{y}{t} = -A\omega^2\cos(kx-\omega t) = -\omega^2\,y ,
\]
donc $-k^2 y = -\dfrac{\omega^2}{v^2}\,y$, soit $v = \dfrac{\omega}{k}$.

\begin{remarque}
Cette équation est du second ordre en temps et $y$ est réelle. L'équation de
Schrödinger sera du \emph{premier} ordre en temps, avec une fonction d'onde
\emph{complexe}.
\end{remarque}

% ---------------------------------------------------------------------
\subsection{Construction de l'équation de Schrödinger}
% ---------------------------------------------------------------------

\subsubsection{Particule libre et particule dans un potentiel}

La force qui s'exerce sur la particule dérive de son énergie potentielle
$V(x,t)$ : $F = -\partial V/\partial x$, et $\dd p/\dd t = F$.

\begin{itemize}
  \item \textbf{Particule libre} ($V = V_0$ constante) : $F = 0$, donc $p$ est
        constante, ainsi que $\lambda = h/p$ et $f = E/h$. On s'attend à une onde
        plane sinusoïdale $\Psi(x,t)\propto\cos(kx-\omega t)$.
  \item \textbf{Particule dans un champ de potentiel} : $F \neq 0$, donc $p$
        varie, et $\lambda$ varie dans l'espace. Si $E$ est conservée,
        $p(x)=\sqrt{2m\,(E-V(x))}$, donc $\lambda(x)=h\big/\sqrt{2m\,(E-V(x))}$.
\end{itemize}

\begin{figure}[!ht]
\centering
\begin{minipage}{0.48\textwidth}
\centering
\begin{tikzpicture}[>=Stealth, scale=0.7]
  \draw[->] (0,0) -- (6.4,0) node[right] {\small $x$};
  \draw[thick, insarouge] (0,0.5) -- (6,0.5);
  \node[insarouge, above] at (5.2,0.5) {\small $V_0$};
  \draw[thick, insagris, domain=0:6, samples=200]
    plot (\x, {1.9+0.6*cos(360*\x/1.5)});
  \node at (3,-0.7) {\small $F=0$ : $p$ et $\lambda$ constants};
\end{tikzpicture}
\end{minipage}\hfill
\begin{minipage}{0.48\textwidth}
\centering
\begin{tikzpicture}[>=Stealth, scale=0.7]
  \draw[->] (0,0) -- (6.4,0) node[right] {\small $x$};
  \draw[thick, insarouge, domain=0:6, samples=60]
    plot (\x, {1.2*exp(-0.5*\x)+0.1});
  \node[insarouge, above] at (4.8,0.3) {\small $V(x)$};
  \draw[thick, insagris, domain=0:6, samples=300]
    plot (\x, {1.9+0.6*cos(360*(0.5*\x+0.07*\x*\x))});
  \node at (3,-0.7) {\small $F\neq0$ : $p$ et $\lambda$ varient};
\end{tikzpicture}
\end{minipage}
\caption{À gauche : particule libre, onde sinusoïdale de $\lambda$ constante.
À droite : particule dans un potentiel décroissant, l'énergie cinétique augmente
et $\lambda$ diminue.}
\end{figure}

\subsubsection{Hypothèses raisonnables}

On cherche une équation vérifiant les quatre hypothèses suivantes.

\begin{enumerate}[label=(\alph*)]
  \item \textbf{Compatibilité avec Einstein--de Broglie :}
        $\lambda = h/p$ et $f = E/h$, soit $p = \hbar k$ et $E = \hbar\omega$.
  \item \textbf{Compatibilité avec l'énergie totale :}
        $E = E_c + E_p = \tfrac12 mv^2 + V = (mv)^2/2m + V = p^2/2m + V$.
        Avec $p = \hbar k$ et $E = \hbar\omega$ :
        \begin{equation}\label{eq:R}
          \hbar\omega = \frac{\hbar^2k^2}{2m} + V(x,t).
        \end{equation}
  \item \textbf{Linéarité :} si $\Psi_1$ et $\Psi_2$ sont solutions, alors
        $\Psi = c_1\Psi_1 + c_2\Psi_2$ l'est aussi. C'est le principe de
        superposition, source des interférences constructives et destructives.
  \item \textbf{Particule libre :} si $V = V_0$, la solution est une onde plane
        progressive sinusoïdale.
\end{enumerate}

\subsubsection{Forme de l'équation}

D'après (b), l'équation doit relier $E \leftrightarrow \omega$ (une dérivée
première en temps), $p^2 \leftrightarrow k^2$ (une dérivée seconde en espace) et
$V$ (un terme $V\Psi$). D'après (c), elle est linéaire en $\Psi$. On cherche donc
\begin{equation}\label{eq:forme}
  \alpha\,\pd{\Psi}{t} = \beta\,\pdd{\Psi}{x} + V(x,t)\,\Psi ,
\end{equation}
où $\alpha$ et $\beta$ sont des constantes à déterminer.

\subsubsection{Un cosinus réel ne convient pas}

Prenons une particule libre ($V=V_0$) et $\Psi = \cos\theta$ avec
$\theta = kx-\omega t$. Alors
\[
  \pd{\Psi}{t} = \omega\sin\theta,
  \qquad
  \pdd{\Psi}{x} = -k^2\cos\theta ,
\]
et l'équation~\eqref{eq:forme} devient
\[
  \alpha\omega\,\sin\theta + (\beta k^2 - V_0)\cos\theta = 0
  \qquad \forall\,(x,t).
\]
Comme $\sin\theta$ et $\cos\theta$ sont indépendants, il faut
$\alpha\omega = 0$ et $\beta k^2 = V_0$. Or $\alpha = 0$ supprime la dérivée
temporelle : l'équation est triviale. On ne peut donc pas travailler avec des
fonctions réelles.

\subsubsection{Passage aux nombres complexes}

Cherchons $\Psi = \cos\theta + \gamma\sin\theta$ avec $\gamma\in\mathbb{C}$
constant. Alors
\[
  \pd{\Psi}{t} = \omega\,(\sin\theta - \gamma\cos\theta),
  \qquad
  \pdd{\Psi}{x} = -k^2\,\Psi .
\]
L'équation~\eqref{eq:forme} s'écrit
\[
  \alpha\omega\,(\sin\theta - \gamma\cos\theta) = (V_0 - \beta k^2)(\cos\theta + \gamma\sin\theta).
\]
On identifie les coefficients de $\cos\theta$ et de $\sin\theta$ :
\begin{align}
  \cos\theta &:\quad -\alpha\omega\gamma = V_0 - \beta k^2, \label{eq:cos}\\
  \sin\theta &:\quad \alpha\omega = \gamma\,(V_0 - \beta k^2). \label{eq:sin}
\end{align}
En reportant~\eqref{eq:cos} dans~\eqref{eq:sin}, on obtient
$\alpha\omega = -\gamma^2\alpha\omega$, donc
\[
  \gamma^2 = -1 \quad\Longrightarrow\quad \gamma = \pm\ii .
\]
Les nombres complexes s'imposent. Avec le choix usuel $\gamma = +\ii$,
\[
  \Psi(x,t) = \cos(kx-\omega t) + \ii\sin(kx-\omega t) = \ee^{\,\ii(kx-\omega t)} .
\]
La relation~\eqref{eq:cos} devient $(-\ii\alpha)\,\omega = V_0 - \beta k^2$. Elle
doit coïncider avec~\eqref{eq:R} écrite sous la forme
$\hbar\,\omega = V_0 + \tfrac{\hbar^2}{2m}\,k^2$. Par identification,
\[
  -\ii\alpha = \hbar \;\Longrightarrow\; \alpha = \ii\hbar,
  \qquad
  -\beta = \frac{\hbar^2}{2m} \;\Longrightarrow\; \beta = -\frac{\hbar^2}{2m}.
\]

\begin{resultat}{Équation de Schrödinger dépendante du temps (1D)}
\[
  \ii\hbar\,\pd{\Psi(x,t)}{t}
  = -\frac{\hbar^2}{2m}\,\pdd{\Psi(x,t)}{x} + V(x,t)\,\Psi(x,t)
\]
En 3D, avec $\nabla^2 = \dfrac{\partial^2}{\partial x^2}+\dfrac{\partial^2}{\partial y^2}+\dfrac{\partial^2}{\partial z^2}$ :
\[
  \ii\hbar\,\pd{\Psi(\mathbf{r},t)}{t} = \Hop\,\Psi(\mathbf{r},t),
  \qquad
  \Hop = -\frac{\hbar^2}{2m}\,\nabla^2 + V(\mathbf{r},t),
\]
où $\Hop$ est l'\emph{opérateur hamiltonien} (énergie totale).
\end{resultat}

\begin{remarque}[Portée de la construction]
Ce n'est pas une démonstration : on a construit l'équation la plus simple
compatible avec (a)--(d) pour un potentiel constant, puis on postule qu'elle
reste valable pour $V(x,t)$ quelconque. Sa validité repose sur l'accord avec
l'expérience. Le choix $\gamma=-\ii$ conduirait à la convention complexe
conjuguée ($\alpha = -\ii\hbar$), physiquement équivalente.
\end{remarque}

\begin{remarque}[Lien avec les opérateurs]
Sur l'onde plane $\ee^{\ii(kx-\omega t)}$ :
$\ii\hbar\,\partial_t\Psi = \hbar\omega\,\Psi = E\,\Psi$ et
$-\ii\hbar\,\partial_x\Psi = \hbar k\,\Psi = p\,\Psi$. L'équation de Schrödinger
traduit donc $E = \frac{p^2}{2m} + V$ avec $E\to\ii\hbar\,\partial_t$ et
$p\to-\ii\hbar\,\partial_x$ : c'est le point de départ du chapitre sur les
opérateurs quantiques.
\end{remarque}

% ---------------------------------------------------------------------
\subsection{Interprétation de Max Born}
% ---------------------------------------------------------------------

La fonction d'onde $\Psi:\mathbb{R}\times\mathbb{R}\to\mathbb{C}$ est à valeurs
complexes : elle n'est ni observable, ni mesurable directement. En revanche, son
module au carré est un réel positif, donc mesurable :
\[
  |\Psi(x,t)|^2 = \Psi^*(x,t)\,\Psi(x,t) \in \mathbb{R}_+ .
\]

\begin{definition}{Interprétation de Born (1926)}
$|\Psi(x,t)|^2$ est la \emph{densité de probabilité de présence} de la particule
en $x$ à l'instant $t$ :
\[
  \dd P\big([x,x+\dd x],t\big) = |\Psi(x,t)|^2\,\dd x ,
  \qquad
  P(a<x<b) = \int_a^b \Psi^*\Psi\,\dd x .
\]
En 3D : $\dd P = |\Psi(\mathbf{r},t)|^2\,\dd^3 r$.
\end{definition}

La particule se trouve forcément quelque part, d'où la \emph{condition de
normalisation}
\[
  \int_{-\infty}^{+\infty} |\Psi(x,t)|^2\,\dd x = 1 .
\]

\begin{figure}[!ht]
\centering
\begin{tikzpicture}[>=Stealth]
  \draw[->] (-3.3,0) -- (3.6,0) node[right] {$x$};
  \draw[->] (0,0) -- (0,2.5) node[above] {$|\Psi(x,t)|^2$};
  \fill[insarouge!20] (0.3,0) -- plot[domain=0.3:1.6, samples=40]
    (\x, {2*exp(-(\x)*(\x)/1.2)}) -- (1.6,0) -- cycle;
  \draw[thick, insagris, domain=-3:3, samples=120]
    plot (\x, {2*exp(-(\x)*(\x)/1.2)});
  \draw[dashed] (0.3,0) -- (0.3,1.86);
  \draw[dashed] (1.6,0) -- (1.6,0.24);
  \node[below] at (0.3,0) {$a$};
  \node[below] at (1.6,0) {$b$};
  \node[anchor=west, align=left] at (1.9,1.7)
    {\small $P(a<x<b)$\\[1mm] \small $=\displaystyle\int_a^b|\Psi|^2\,\dd x$};
  \draw[->] (2.3,1.2) -- (1.15,0.45);
\end{tikzpicture}
\caption{Densité de probabilité de présence ; la probabilité de trouver la
particule entre $a$ et $b$ est l'aire colorée.}
\end{figure}

\begin{remarque}
\begin{itemize}[nosep]
  \item Multiplier $\Psi$ par une phase globale $\ee^{\ii\varphi}$ ne change pas
        $|\Psi|^2$ : la phase globale n'est pas observable.
  \item L'onde plane $\ee^{\ii(kx-\omega t)}$ a $|\Psi|^2 = 1$ partout : elle
        n'est pas normalisable, la particule serait équiprobable à l'infini. Un
        état physique est une \emph{superposition} d'ondes planes (paquet
        d'ondes), possible grâce à la linéarité (c).
  \item Avec la notation de Dirac (chapitres suivants), $|\Psi\rangle$ désigne
        l'état et la normalisation s'écrit $\int\Psi^*\Psi\,\dd x = \langle\Psi|\Psi\rangle = 1$.
\end{itemize}
\end{remarque}

% ---------------------------------------------------------------------
\subsection{Récapitulatif}
% ---------------------------------------------------------------------

\begin{center}
\renewcommand{\arraystretch}{1.5}
\begin{tabular}{@{}ll@{}}
\toprule
\textbf{Relation} & \textbf{Formule} \\
\midrule
Planck--Einstein & $E = hf = \hbar\omega$ \\
de Broglie & $\lambda = \dfrac{h}{p}$, \quad $p = \hbar k$ \\
Énergie totale & $E = \dfrac{p^2}{2m} + V(x,t)$ \\
Onde plane complexe & $\Psi = \ee^{\,\ii(kx-\omega t)}$ \\
Schrödinger (dépendante du temps) & $\ii\hbar\,\partial_t\Psi = \Hop\Psi$, \quad
   $\Hop = -\dfrac{\hbar^2}{2m}\partial_x^2 + V$ \\
Born & $\dd P = |\Psi|^2\,\dd x$, \quad $\displaystyle\int_{-\infty}^{+\infty}|\Psi|^2\,\dd x = 1$ \\
\bottomrule
\end{tabular}
\end{center}
.
%  Il ne contient ni préambule ni \begin{document} : les environnements
%  (definition, resultat, remarque), les macros (\ii, \ee, \dd, \pd, \pdd,
%  \Hop) et les couleurs (insarouge, insagris) sont définis dans main.tex.
% =====================================================================

% =====================================================================
\section{Fondements physiques de la mécanique quantique}\label{sec:fondements}
% =====================================================================

\begin{remarque}[Plan général du cours]
\begin{enumerate}[nosep]
  \item Mécanique quantique (ce chapitre) ;
  \item Espace de Hilbert $\mathcal{H}$ ;
  \item Notation de Dirac ;
  \item Espace produit tensoriel ;
  \item États quantiques, intrication quantique, états de Bell.
\end{enumerate}
\end{remarque}

% ---------------------------------------------------------------------
\subsection{Dualité onde--corpuscule}
% ---------------------------------------------------------------------

\subsubsection{Einstein : la lumière est faite de photons}

À la fin du \textsc{xix}\textsuperscript{e} siècle, l'effet photo-électrique
(émission d'électrons par un métal éclairé) contredit la théorie purement
ondulatoire de la lumière : l'émission dépend de la \emph{fréquence} de la
lumière, et non de son intensité. En 1905, Einstein l'explique en supposant que
la lumière est formée de quanta d'énergie, les \emph{photons}. Un électron n'est
éjecté que si l'énergie du photon dépasse le travail d'extraction $W$ du métal :
$E_{c,\max} = hf - W$. Une onde lumineuse se comporte donc aussi comme une
particule.

\begin{definition}{Photon (Planck--Einstein)}
\[
  E = hf = h\nu = \hbar\omega ,
  \qquad
  p = \frac{h}{\lambda} = \hbar k ,
\]
avec $h \approx 6{,}63\times10^{-34}\ \mathrm{J\cdot s}$ (constante de Planck),
$\hbar = \dfrac{h}{2\pi}$, $\omega = 2\pi f$ (pulsation) et
$k = \dfrac{2\pi}{\lambda}$ (nombre d'onde).
\end{definition}

\subsubsection{de Broglie : la matière est une onde}

En 1924, de Broglie inverse le raisonnement : si une onde peut se comporter
comme une particule, une \emph{particule matérielle} peut se comporter comme une
onde. À une particule de masse $m$ et de vitesse $v$, donc de quantité de
mouvement $p = mv$, on associe une \emph{onde de matière}.

\begin{definition}{Relations de de Broglie}
\[
  \lambda = \frac{h}{p} = \frac{h}{mv} ,
  \qquad
  f = \frac{E}{h}
  \qquad\Longleftrightarrow\qquad
  p = \hbar k ,\quad E = \hbar\omega .
\]
\end{definition}

\begin{figure}[!ht]
\centering
\begin{tikzpicture}[>=Stealth, thick]
% --- Ligne 1 : Einstein (onde -> particule)
\begin{scope}[yshift=2.3cm]
  \draw[insarouge, domain=0:2.4, samples=100] plot (\x, {0.35*sin(360*\x/0.8)});
  \node at (1.2,-0.8) {\small onde lumineuse};
  \draw[<->] (3.0,0) -- (5.4,0)
    node[midway, above] {\small Einstein}
    node[midway, below] {\small $E = hf = \hbar\omega$};
  \shade[ball color=insarouge!85] (6.8,0) circle (0.3);
  \node at (6.8,0) {\small $\gamma$};
  \node at (6.8,-0.8) {\small photon};
\end{scope}
% --- Ligne 2 : de Broglie (particule -> onde)
\begin{scope}
  \shade[ball color=insagris!80] (0.3,0) circle (0.3);
  \draw[->] (0.7,0) -- (2.0,0) node[midway, above] {\small $p = mv$};
  \node at (1.0,-0.8) {\small particule};
  \draw[<->] (3.0,0) -- (5.4,0)
    node[midway, above] {\small de Broglie}
    node[midway, below] {\small $\lambda = h/p$};
  \draw[insagris, domain=0:2.4, samples=100, shift={(6.0,0)}]
    plot (\x, {0.35*sin(360*\x/0.8)});
  \node at (7.2,-0.8) {\small onde de matière};
\end{scope}
\end{tikzpicture}
\caption{Dualité onde--corpuscule : Einstein associe une particule à l'onde
lumineuse, de Broglie associe une onde à la particule matérielle.}
\end{figure}

\begin{remarque}[Ordres de grandeur]
\begin{itemize}[nosep]
  \item Balle de $0{,}1$ kg à $10$ m/s : $\lambda = \dfrac{6{,}63\times10^{-34}}{0{,}1\times 10}
        \approx 6{,}6\times10^{-34}$ m, inobservable.
  \item Électron accéléré sous $100$ V : $p=\sqrt{2m_e\,eU}$ donne
        $\lambda \approx 0{,}12$ nm, de l'ordre de la distance entre atomes d'un cristal.
        La diffraction d'électrons par un cristal (Davisson--Germer, 1927) confirme
        l'existence de l'onde de matière.
\end{itemize}
\end{remarque}

\subsubsection{Le problème posé à Schrödinger}

Si toute particule est associée à une onde, existe-t-il une \emph{équation
d'onde} qui gouverne cette onde de matière ? C'est la question à laquelle
répond l'équation de Schrödinger (1926).

% ---------------------------------------------------------------------
\subsection{Rappel : l'onde classique}
% ---------------------------------------------------------------------

Une onde plane progressive sinusoïdale se propageant selon $x$ s'écrit
\[
  y(x,t) = A\cos(kx - \omega t),
\]
où $A$ est l'amplitude, $k = 2\pi/\lambda$ le nombre d'onde, $\omega = 2\pi f$ la
pulsation et $v = \omega/k = \lambda f$ la vitesse de phase (un maximum vérifie
$kx-\omega t = \text{cste}$, donc $\dd x/\dd t = \omega/k$).

\begin{figure}[!ht]
\centering
\begin{tikzpicture}[>=Stealth]
  \draw[->] (-0.3,0) -- (8.7,0) node[right] {$x$};
  \draw[->] (0,-1.6) -- (0,2.0) node[above] {$y$};
  \draw[thick, insagris, domain=0:8.4, samples=200]
    plot (\x, {cos(360*\x/4)});
  \draw[thick, dashed, insarouge, domain=0:8.4, samples=200]
    plot (\x, {cos(360*(\x-0.9)/4)});
  \draw[dotted] (0,1) -- (4,1);
  \node[left] at (0,1) {$A$};
  \node[left] at (0,-1) {$-A$};
  \draw[dotted] (4,1) -- (4,1.55);
  \draw[<->] (0,1.45) -- (4,1.45) node[midway, above] {$\lambda$};
  \draw[->, insarouge] (6,-1.2) -- (6.9,-1.2)
    node[midway, below] {\small $v\,\Delta t$};
  \draw[thick, insagris] (6.0,1.75) -- (6.6,1.75)
    node[right, black] {\small $t$};
  \draw[thick, dashed, insarouge] (7.4,1.75) -- (8.0,1.75)
    node[right, black] {\small $t+\Delta t$};
\end{tikzpicture}
\caption{Onde plane progressive $y(x,t)=A\cos(kx-\omega t)$ à deux instants :
le profil se translate de $v\,\Delta t$.}
\end{figure}

Elle vérifie l'\emph{équation des ondes}
\[
  \pdd{y}{x} = \frac{1}{v^2}\,\pdd{y}{t}
  \qquad\text{(courbure $\leftrightarrow$ accélération).}
\]
En effet,
\[
  \pdd{y}{x} = -Ak^2\cos(kx-\omega t) = -k^2\,y ,
  \qquad
  \pdd{y}{t} = -A\omega^2\cos(kx-\omega t) = -\omega^2\,y ,
\]
donc $-k^2 y = -\dfrac{\omega^2}{v^2}\,y$, soit $v = \dfrac{\omega}{k}$.

\begin{remarque}
Cette équation est du second ordre en temps et $y$ est réelle. L'équation de
Schrödinger sera du \emph{premier} ordre en temps, avec une fonction d'onde
\emph{complexe}.
\end{remarque}

% ---------------------------------------------------------------------
\subsection{Construction de l'équation de Schrödinger}
% ---------------------------------------------------------------------

\subsubsection{Particule libre et particule dans un potentiel}

La force qui s'exerce sur la particule dérive de son énergie potentielle
$V(x,t)$ : $F = -\partial V/\partial x$, et $\dd p/\dd t = F$.

\begin{itemize}
  \item \textbf{Particule libre} ($V = V_0$ constante) : $F = 0$, donc $p$ est
        constante, ainsi que $\lambda = h/p$ et $f = E/h$. On s'attend à une onde
        plane sinusoïdale $\Psi(x,t)\propto\cos(kx-\omega t)$.
  \item \textbf{Particule dans un champ de potentiel} : $F \neq 0$, donc $p$
        varie, et $\lambda$ varie dans l'espace. Si $E$ est conservée,
        $p(x)=\sqrt{2m\,(E-V(x))}$, donc $\lambda(x)=h\big/\sqrt{2m\,(E-V(x))}$.
\end{itemize}

\begin{figure}[!ht]
\centering
\begin{minipage}{0.48\textwidth}
\centering
\begin{tikzpicture}[>=Stealth, scale=0.7]
  \draw[->] (0,0) -- (6.4,0) node[right] {\small $x$};
  \draw[thick, insarouge] (0,0.5) -- (6,0.5);
  \node[insarouge, above] at (5.2,0.5) {\small $V_0$};
  \draw[thick, insagris, domain=0:6, samples=200]
    plot (\x, {1.9+0.6*cos(360*\x/1.5)});
  \node at (3,-0.7) {\small $F=0$ : $p$ et $\lambda$ constants};
\end{tikzpicture}
\end{minipage}\hfill
\begin{minipage}{0.48\textwidth}
\centering
\begin{tikzpicture}[>=Stealth, scale=0.7]
  \draw[->] (0,0) -- (6.4,0) node[right] {\small $x$};
  \draw[thick, insarouge, domain=0:6, samples=60]
    plot (\x, {1.2*exp(-0.5*\x)+0.1});
  \node[insarouge, above] at (4.8,0.3) {\small $V(x)$};
  \draw[thick, insagris, domain=0:6, samples=300]
    plot (\x, {1.9+0.6*cos(360*(0.5*\x+0.07*\x*\x))});
  \node at (3,-0.7) {\small $F\neq0$ : $p$ et $\lambda$ varient};
\end{tikzpicture}
\end{minipage}
\caption{À gauche : particule libre, onde sinusoïdale de $\lambda$ constante.
À droite : particule dans un potentiel décroissant, l'énergie cinétique augmente
et $\lambda$ diminue.}
\end{figure}

\subsubsection{Hypothèses raisonnables}

On cherche une équation vérifiant les quatre hypothèses suivantes.

\begin{enumerate}[label=(\alph*)]
  \item \textbf{Compatibilité avec Einstein--de Broglie :}
        $\lambda = h/p$ et $f = E/h$, soit $p = \hbar k$ et $E = \hbar\omega$.
  \item \textbf{Compatibilité avec l'énergie totale :}
        $E = E_c + E_p = \tfrac12 mv^2 + V = (mv)^2/2m + V = p^2/2m + V$.
        Avec $p = \hbar k$ et $E = \hbar\omega$ :
        \begin{equation}\label{eq:R}
          \hbar\omega = \frac{\hbar^2k^2}{2m} + V(x,t).
        \end{equation}
  \item \textbf{Linéarité :} si $\Psi_1$ et $\Psi_2$ sont solutions, alors
        $\Psi = c_1\Psi_1 + c_2\Psi_2$ l'est aussi. C'est le principe de
        superposition, source des interférences constructives et destructives.
  \item \textbf{Particule libre :} si $V = V_0$, la solution est une onde plane
        progressive sinusoïdale.
\end{enumerate}

\subsubsection{Forme de l'équation}

D'après (b), l'équation doit relier $E \leftrightarrow \omega$ (une dérivée
première en temps), $p^2 \leftrightarrow k^2$ (une dérivée seconde en espace) et
$V$ (un terme $V\Psi$). D'après (c), elle est linéaire en $\Psi$. On cherche donc
\begin{equation}\label{eq:forme}
  \alpha\,\pd{\Psi}{t} = \beta\,\pdd{\Psi}{x} + V(x,t)\,\Psi ,
\end{equation}
où $\alpha$ et $\beta$ sont des constantes à déterminer.

\subsubsection{Un cosinus réel ne convient pas}

Prenons une particule libre ($V=V_0$) et $\Psi = \cos\theta$ avec
$\theta = kx-\omega t$. Alors
\[
  \pd{\Psi}{t} = \omega\sin\theta,
  \qquad
  \pdd{\Psi}{x} = -k^2\cos\theta ,
\]
et l'équation~\eqref{eq:forme} devient
\[
  \alpha\omega\,\sin\theta + (\beta k^2 - V_0)\cos\theta = 0
  \qquad \forall\,(x,t).
\]
Comme $\sin\theta$ et $\cos\theta$ sont indépendants, il faut
$\alpha\omega = 0$ et $\beta k^2 = V_0$. Or $\alpha = 0$ supprime la dérivée
temporelle : l'équation est triviale. On ne peut donc pas travailler avec des
fonctions réelles.

\subsubsection{Passage aux nombres complexes}

Cherchons $\Psi = \cos\theta + \gamma\sin\theta$ avec $\gamma\in\mathbb{C}$
constant. Alors
\[
  \pd{\Psi}{t} = \omega\,(\sin\theta - \gamma\cos\theta),
  \qquad
  \pdd{\Psi}{x} = -k^2\,\Psi .
\]
L'équation~\eqref{eq:forme} s'écrit
\[
  \alpha\omega\,(\sin\theta - \gamma\cos\theta) = (V_0 - \beta k^2)(\cos\theta + \gamma\sin\theta).
\]
On identifie les coefficients de $\cos\theta$ et de $\sin\theta$ :
\begin{align}
  \cos\theta &:\quad -\alpha\omega\gamma = V_0 - \beta k^2, \label{eq:cos}\\
  \sin\theta &:\quad \alpha\omega = \gamma\,(V_0 - \beta k^2). \label{eq:sin}
\end{align}
En reportant~\eqref{eq:cos} dans~\eqref{eq:sin}, on obtient
$\alpha\omega = -\gamma^2\alpha\omega$, donc
\[
  \gamma^2 = -1 \quad\Longrightarrow\quad \gamma = \pm\ii .
\]
Les nombres complexes s'imposent. Avec le choix usuel $\gamma = +\ii$,
\[
  \Psi(x,t) = \cos(kx-\omega t) + \ii\sin(kx-\omega t) = \ee^{\,\ii(kx-\omega t)} .
\]
La relation~\eqref{eq:cos} devient $(-\ii\alpha)\,\omega = V_0 - \beta k^2$. Elle
doit coïncider avec~\eqref{eq:R} écrite sous la forme
$\hbar\,\omega = V_0 + \tfrac{\hbar^2}{2m}\,k^2$. Par identification,
\[
  -\ii\alpha = \hbar \;\Longrightarrow\; \alpha = \ii\hbar,
  \qquad
  -\beta = \frac{\hbar^2}{2m} \;\Longrightarrow\; \beta = -\frac{\hbar^2}{2m}.
\]

\begin{resultat}{Équation de Schrödinger dépendante du temps (1D)}
\[
  \ii\hbar\,\pd{\Psi(x,t)}{t}
  = -\frac{\hbar^2}{2m}\,\pdd{\Psi(x,t)}{x} + V(x,t)\,\Psi(x,t)
\]
En 3D, avec $\nabla^2 = \dfrac{\partial^2}{\partial x^2}+\dfrac{\partial^2}{\partial y^2}+\dfrac{\partial^2}{\partial z^2}$ :
\[
  \ii\hbar\,\pd{\Psi(\mathbf{r},t)}{t} = \Hop\,\Psi(\mathbf{r},t),
  \qquad
  \Hop = -\frac{\hbar^2}{2m}\,\nabla^2 + V(\mathbf{r},t),
\]
où $\Hop$ est l'\emph{opérateur hamiltonien} (énergie totale).
\end{resultat}

\begin{remarque}[Portée de la construction]
Ce n'est pas une démonstration : on a construit l'équation la plus simple
compatible avec (a)--(d) pour un potentiel constant, puis on postule qu'elle
reste valable pour $V(x,t)$ quelconque. Sa validité repose sur l'accord avec
l'expérience. Le choix $\gamma=-\ii$ conduirait à la convention complexe
conjuguée ($\alpha = -\ii\hbar$), physiquement équivalente.
\end{remarque}

\begin{remarque}[Lien avec les opérateurs]
Sur l'onde plane $\ee^{\ii(kx-\omega t)}$ :
$\ii\hbar\,\partial_t\Psi = \hbar\omega\,\Psi = E\,\Psi$ et
$-\ii\hbar\,\partial_x\Psi = \hbar k\,\Psi = p\,\Psi$. L'équation de Schrödinger
traduit donc $E = \frac{p^2}{2m} + V$ avec $E\to\ii\hbar\,\partial_t$ et
$p\to-\ii\hbar\,\partial_x$ : c'est le point de départ du chapitre sur les
opérateurs quantiques.
\end{remarque}

% ---------------------------------------------------------------------
\subsection{Interprétation de Max Born}
% ---------------------------------------------------------------------

La fonction d'onde $\Psi:\mathbb{R}\times\mathbb{R}\to\mathbb{C}$ est à valeurs
complexes : elle n'est ni observable, ni mesurable directement. En revanche, son
module au carré est un réel positif, donc mesurable :
\[
  |\Psi(x,t)|^2 = \Psi^*(x,t)\,\Psi(x,t) \in \mathbb{R}_+ .
\]

\begin{definition}{Interprétation de Born (1926)}
$|\Psi(x,t)|^2$ est la \emph{densité de probabilité de présence} de la particule
en $x$ à l'instant $t$ :
\[
  \dd P\big([x,x+\dd x],t\big) = |\Psi(x,t)|^2\,\dd x ,
  \qquad
  P(a<x<b) = \int_a^b \Psi^*\Psi\,\dd x .
\]
En 3D : $\dd P = |\Psi(\mathbf{r},t)|^2\,\dd^3 r$.
\end{definition}

La particule se trouve forcément quelque part, d'où la \emph{condition de
normalisation}
\[
  \int_{-\infty}^{+\infty} |\Psi(x,t)|^2\,\dd x = 1 .
\]

\begin{figure}[!ht]
\centering
\begin{tikzpicture}[>=Stealth]
  \draw[->] (-3.3,0) -- (3.6,0) node[right] {$x$};
  \draw[->] (0,0) -- (0,2.5) node[above] {$|\Psi(x,t)|^2$};
  \fill[insarouge!20] (0.3,0) -- plot[domain=0.3:1.6, samples=40]
    (\x, {2*exp(-(\x)*(\x)/1.2)}) -- (1.6,0) -- cycle;
  \draw[thick, insagris, domain=-3:3, samples=120]
    plot (\x, {2*exp(-(\x)*(\x)/1.2)});
  \draw[dashed] (0.3,0) -- (0.3,1.86);
  \draw[dashed] (1.6,0) -- (1.6,0.24);
  \node[below] at (0.3,0) {$a$};
  \node[below] at (1.6,0) {$b$};
  \node[anchor=west, align=left] at (1.9,1.7)
    {\small $P(a<x<b)$\\[1mm] \small $=\displaystyle\int_a^b|\Psi|^2\,\dd x$};
  \draw[->] (2.3,1.2) -- (1.15,0.45);
\end{tikzpicture}
\caption{Densité de probabilité de présence ; la probabilité de trouver la
particule entre $a$ et $b$ est l'aire colorée.}
\end{figure}

\begin{remarque}
\begin{itemize}[nosep]
  \item Multiplier $\Psi$ par une phase globale $\ee^{\ii\varphi}$ ne change pas
        $|\Psi|^2$ : la phase globale n'est pas observable.
  \item L'onde plane $\ee^{\ii(kx-\omega t)}$ a $|\Psi|^2 = 1$ partout : elle
        n'est pas normalisable, la particule serait équiprobable à l'infini. Un
        état physique est une \emph{superposition} d'ondes planes (paquet
        d'ondes), possible grâce à la linéarité (c).
  \item Avec la notation de Dirac (chapitres suivants), $|\Psi\rangle$ désigne
        l'état et la normalisation s'écrit $\int\Psi^*\Psi\,\dd x = \langle\Psi|\Psi\rangle = 1$.
\end{itemize}
\end{remarque}

% ---------------------------------------------------------------------
\subsection{Récapitulatif}
% ---------------------------------------------------------------------

\begin{center}
\renewcommand{\arraystretch}{1.5}
\begin{tabular}{@{}ll@{}}
\toprule
\textbf{Relation} & \textbf{Formule} \\
\midrule
Planck--Einstein & $E = hf = \hbar\omega$ \\
de Broglie & $\lambda = \dfrac{h}{p}$, \quad $p = \hbar k$ \\
Énergie totale & $E = \dfrac{p^2}{2m} + V(x,t)$ \\
Onde plane complexe & $\Psi = \ee^{\,\ii(kx-\omega t)}$ \\
Schrödinger (dépendante du temps) & $\ii\hbar\,\partial_t\Psi = \Hop\Psi$, \quad
   $\Hop = -\dfrac{\hbar^2}{2m}\partial_x^2 + V$ \\
Born & $\dd P = |\Psi|^2\,\dd x$, \quad $\displaystyle\int_{-\infty}^{+\infty}|\Psi|^2\,\dd x = 1$ \\
\bottomrule
\end{tabular}
\end{center}
.
%  Il ne contient ni préambule ni \begin{document} : les environnements
%  (definition, resultat, remarque), les macros (\ii, \ee, \dd, \pd, \pdd,
%  \Hop) et les couleurs (insarouge, insagris) sont définis dans main.tex.
% =====================================================================

% =====================================================================
\section{Fondements physiques de la mécanique quantique}\label{sec:fondements}
% =====================================================================

\begin{remarque}[Plan général du cours]
\begin{enumerate}[nosep]
  \item Mécanique quantique (ce chapitre) ;
  \item Espace de Hilbert $\mathcal{H}$ ;
  \item Notation de Dirac ;
  \item Espace produit tensoriel ;
  \item États quantiques, intrication quantique, états de Bell.
\end{enumerate}
\end{remarque}

% ---------------------------------------------------------------------
\subsection{Dualité onde--corpuscule}
% ---------------------------------------------------------------------

\subsubsection{Einstein : la lumière est faite de photons}

À la fin du \textsc{xix}\textsuperscript{e} siècle, l'effet photo-électrique
(émission d'électrons par un métal éclairé) contredit la théorie purement
ondulatoire de la lumière : l'émission dépend de la \emph{fréquence} de la
lumière, et non de son intensité. En 1905, Einstein l'explique en supposant que
la lumière est formée de quanta d'énergie, les \emph{photons}. Un électron n'est
éjecté que si l'énergie du photon dépasse le travail d'extraction $W$ du métal :
$E_{c,\max} = hf - W$. Une onde lumineuse se comporte donc aussi comme une
particule.

\begin{definition}{Photon (Planck--Einstein)}
\[
  E = hf = h\nu = \hbar\omega ,
  \qquad
  p = \frac{h}{\lambda} = \hbar k ,
\]
avec $h \approx 6{,}63\times10^{-34}\ \mathrm{J\cdot s}$ (constante de Planck),
$\hbar = \dfrac{h}{2\pi}$, $\omega = 2\pi f$ (pulsation) et
$k = \dfrac{2\pi}{\lambda}$ (nombre d'onde).
\end{definition}

\subsubsection{de Broglie : la matière est une onde}

En 1924, de Broglie inverse le raisonnement : si une onde peut se comporter
comme une particule, une \emph{particule matérielle} peut se comporter comme une
onde. À une particule de masse $m$ et de vitesse $v$, donc de quantité de
mouvement $p = mv$, on associe une \emph{onde de matière}.

\begin{definition}{Relations de de Broglie}
\[
  \lambda = \frac{h}{p} = \frac{h}{mv} ,
  \qquad
  f = \frac{E}{h}
  \qquad\Longleftrightarrow\qquad
  p = \hbar k ,\quad E = \hbar\omega .
\]
\end{definition}

\begin{figure}[!ht]
\centering
\begin{tikzpicture}[>=Stealth, thick]
% --- Ligne 1 : Einstein (onde -> particule)
\begin{scope}[yshift=2.3cm]
  \draw[insarouge, domain=0:2.4, samples=100] plot (\x, {0.35*sin(360*\x/0.8)});
  \node at (1.2,-0.8) {\small onde lumineuse};
  \draw[<->] (3.0,0) -- (5.4,0)
    node[midway, above] {\small Einstein}
    node[midway, below] {\small $E = hf = \hbar\omega$};
  \shade[ball color=insarouge!85] (6.8,0) circle (0.3);
  \node at (6.8,0) {\small $\gamma$};
  \node at (6.8,-0.8) {\small photon};
\end{scope}
% --- Ligne 2 : de Broglie (particule -> onde)
\begin{scope}
  \shade[ball color=insagris!80] (0.3,0) circle (0.3);
  \draw[->] (0.7,0) -- (2.0,0) node[midway, above] {\small $p = mv$};
  \node at (1.0,-0.8) {\small particule};
  \draw[<->] (3.0,0) -- (5.4,0)
    node[midway, above] {\small de Broglie}
    node[midway, below] {\small $\lambda = h/p$};
  \draw[insagris, domain=0:2.4, samples=100, shift={(6.0,0)}]
    plot (\x, {0.35*sin(360*\x/0.8)});
  \node at (7.2,-0.8) {\small onde de matière};
\end{scope}
\end{tikzpicture}
\caption{Dualité onde--corpuscule : Einstein associe une particule à l'onde
lumineuse, de Broglie associe une onde à la particule matérielle.}
\end{figure}

\begin{remarque}[Ordres de grandeur]
\begin{itemize}[nosep]
  \item Balle de $0{,}1$ kg à $10$ m/s : $\lambda = \dfrac{6{,}63\times10^{-34}}{0{,}1\times 10}
        \approx 6{,}6\times10^{-34}$ m, inobservable.
  \item Électron accéléré sous $100$ V : $p=\sqrt{2m_e\,eU}$ donne
        $\lambda \approx 0{,}12$ nm, de l'ordre de la distance entre atomes d'un cristal.
        La diffraction d'électrons par un cristal (Davisson--Germer, 1927) confirme
        l'existence de l'onde de matière.
\end{itemize}
\end{remarque}

\subsubsection{Le problème posé à Schrödinger}

Si toute particule est associée à une onde, existe-t-il une \emph{équation
d'onde} qui gouverne cette onde de matière ? C'est la question à laquelle
répond l'équation de Schrödinger (1926).

% ---------------------------------------------------------------------
\subsection{Rappel : l'onde classique}
% ---------------------------------------------------------------------

Une onde plane progressive sinusoïdale se propageant selon $x$ s'écrit
\[
  y(x,t) = A\cos(kx - \omega t),
\]
où $A$ est l'amplitude, $k = 2\pi/\lambda$ le nombre d'onde, $\omega = 2\pi f$ la
pulsation et $v = \omega/k = \lambda f$ la vitesse de phase (un maximum vérifie
$kx-\omega t = \text{cste}$, donc $\dd x/\dd t = \omega/k$).

\begin{figure}[!ht]
\centering
\begin{tikzpicture}[>=Stealth]
  \draw[->] (-0.3,0) -- (8.7,0) node[right] {$x$};
  \draw[->] (0,-1.6) -- (0,2.0) node[above] {$y$};
  \draw[thick, insagris, domain=0:8.4, samples=200]
    plot (\x, {cos(360*\x/4)});
  \draw[thick, dashed, insarouge, domain=0:8.4, samples=200]
    plot (\x, {cos(360*(\x-0.9)/4)});
  \draw[dotted] (0,1) -- (4,1);
  \node[left] at (0,1) {$A$};
  \node[left] at (0,-1) {$-A$};
  \draw[dotted] (4,1) -- (4,1.55);
  \draw[<->] (0,1.45) -- (4,1.45) node[midway, above] {$\lambda$};
  \draw[->, insarouge] (6,-1.2) -- (6.9,-1.2)
    node[midway, below] {\small $v\,\Delta t$};
  \draw[thick, insagris] (6.0,1.75) -- (6.6,1.75)
    node[right, black] {\small $t$};
  \draw[thick, dashed, insarouge] (7.4,1.75) -- (8.0,1.75)
    node[right, black] {\small $t+\Delta t$};
\end{tikzpicture}
\caption{Onde plane progressive $y(x,t)=A\cos(kx-\omega t)$ à deux instants :
le profil se translate de $v\,\Delta t$.}
\end{figure}

Elle vérifie l'\emph{équation des ondes}
\[
  \pdd{y}{x} = \frac{1}{v^2}\,\pdd{y}{t}
  \qquad\text{(courbure $\leftrightarrow$ accélération).}
\]
En effet,
\[
  \pdd{y}{x} = -Ak^2\cos(kx-\omega t) = -k^2\,y ,
  \qquad
  \pdd{y}{t} = -A\omega^2\cos(kx-\omega t) = -\omega^2\,y ,
\]
donc $-k^2 y = -\dfrac{\omega^2}{v^2}\,y$, soit $v = \dfrac{\omega}{k}$.

\begin{remarque}
Cette équation est du second ordre en temps et $y$ est réelle. L'équation de
Schrödinger sera du \emph{premier} ordre en temps, avec une fonction d'onde
\emph{complexe}.
\end{remarque}

% ---------------------------------------------------------------------
\subsection{Construction de l'équation de Schrödinger}
% ---------------------------------------------------------------------

\subsubsection{Particule libre et particule dans un potentiel}

La force qui s'exerce sur la particule dérive de son énergie potentielle
$V(x,t)$ : $F = -\partial V/\partial x$, et $\dd p/\dd t = F$.

\begin{itemize}
  \item \textbf{Particule libre} ($V = V_0$ constante) : $F = 0$, donc $p$ est
        constante, ainsi que $\lambda = h/p$ et $f = E/h$. On s'attend à une onde
        plane sinusoïdale $\Psi(x,t)\propto\cos(kx-\omega t)$.
  \item \textbf{Particule dans un champ de potentiel} : $F \neq 0$, donc $p$
        varie, et $\lambda$ varie dans l'espace. Si $E$ est conservée,
        $p(x)=\sqrt{2m\,(E-V(x))}$, donc $\lambda(x)=h\big/\sqrt{2m\,(E-V(x))}$.
\end{itemize}

\begin{figure}[!ht]
\centering
\begin{minipage}{0.48\textwidth}
\centering
\begin{tikzpicture}[>=Stealth, scale=0.7]
  \draw[->] (0,0) -- (6.4,0) node[right] {\small $x$};
  \draw[thick, insarouge] (0,0.5) -- (6,0.5);
  \node[insarouge, above] at (5.2,0.5) {\small $V_0$};
  \draw[thick, insagris, domain=0:6, samples=200]
    plot (\x, {1.9+0.6*cos(360*\x/1.5)});
  \node at (3,-0.7) {\small $F=0$ : $p$ et $\lambda$ constants};
\end{tikzpicture}
\end{minipage}\hfill
\begin{minipage}{0.48\textwidth}
\centering
\begin{tikzpicture}[>=Stealth, scale=0.7]
  \draw[->] (0,0) -- (6.4,0) node[right] {\small $x$};
  \draw[thick, insarouge, domain=0:6, samples=60]
    plot (\x, {1.2*exp(-0.5*\x)+0.1});
  \node[insarouge, above] at (4.8,0.3) {\small $V(x)$};
  \draw[thick, insagris, domain=0:6, samples=300]
    plot (\x, {1.9+0.6*cos(360*(0.5*\x+0.07*\x*\x))});
  \node at (3,-0.7) {\small $F\neq0$ : $p$ et $\lambda$ varient};
\end{tikzpicture}
\end{minipage}
\caption{À gauche : particule libre, onde sinusoïdale de $\lambda$ constante.
À droite : particule dans un potentiel décroissant, l'énergie cinétique augmente
et $\lambda$ diminue.}
\end{figure}

\subsubsection{Hypothèses raisonnables}

On cherche une équation vérifiant les quatre hypothèses suivantes.

\begin{enumerate}[label=(\alph*)]
  \item \textbf{Compatibilité avec Einstein--de Broglie :}
        $\lambda = h/p$ et $f = E/h$, soit $p = \hbar k$ et $E = \hbar\omega$.
  \item \textbf{Compatibilité avec l'énergie totale :}
        $E = E_c + E_p = \tfrac12 mv^2 + V = (mv)^2/2m + V = p^2/2m + V$.
        Avec $p = \hbar k$ et $E = \hbar\omega$ :
        \begin{equation}\label{eq:R}
          \hbar\omega = \frac{\hbar^2k^2}{2m} + V(x,t).
        \end{equation}
  \item \textbf{Linéarité :} si $\Psi_1$ et $\Psi_2$ sont solutions, alors
        $\Psi = c_1\Psi_1 + c_2\Psi_2$ l'est aussi. C'est le principe de
        superposition, source des interférences constructives et destructives.
  \item \textbf{Particule libre :} si $V = V_0$, la solution est une onde plane
        progressive sinusoïdale.
\end{enumerate}

\subsubsection{Forme de l'équation}

D'après (b), l'équation doit relier $E \leftrightarrow \omega$ (une dérivée
première en temps), $p^2 \leftrightarrow k^2$ (une dérivée seconde en espace) et
$V$ (un terme $V\Psi$). D'après (c), elle est linéaire en $\Psi$. On cherche donc
\begin{equation}\label{eq:forme}
  \alpha\,\pd{\Psi}{t} = \beta\,\pdd{\Psi}{x} + V(x,t)\,\Psi ,
\end{equation}
où $\alpha$ et $\beta$ sont des constantes à déterminer.

\subsubsection{Un cosinus réel ne convient pas}

Prenons une particule libre ($V=V_0$) et $\Psi = \cos\theta$ avec
$\theta = kx-\omega t$. Alors
\[
  \pd{\Psi}{t} = \omega\sin\theta,
  \qquad
  \pdd{\Psi}{x} = -k^2\cos\theta ,
\]
et l'équation~\eqref{eq:forme} devient
\[
  \alpha\omega\,\sin\theta + (\beta k^2 - V_0)\cos\theta = 0
  \qquad \forall\,(x,t).
\]
Comme $\sin\theta$ et $\cos\theta$ sont indépendants, il faut
$\alpha\omega = 0$ et $\beta k^2 = V_0$. Or $\alpha = 0$ supprime la dérivée
temporelle : l'équation est triviale. On ne peut donc pas travailler avec des
fonctions réelles.

\subsubsection{Passage aux nombres complexes}

Cherchons $\Psi = \cos\theta + \gamma\sin\theta$ avec $\gamma\in\mathbb{C}$
constant. Alors
\[
  \pd{\Psi}{t} = \omega\,(\sin\theta - \gamma\cos\theta),
  \qquad
  \pdd{\Psi}{x} = -k^2\,\Psi .
\]
L'équation~\eqref{eq:forme} s'écrit
\[
  \alpha\omega\,(\sin\theta - \gamma\cos\theta) = (V_0 - \beta k^2)(\cos\theta + \gamma\sin\theta).
\]
On identifie les coefficients de $\cos\theta$ et de $\sin\theta$ :
\begin{align}
  \cos\theta &:\quad -\alpha\omega\gamma = V_0 - \beta k^2, \label{eq:cos}\\
  \sin\theta &:\quad \alpha\omega = \gamma\,(V_0 - \beta k^2). \label{eq:sin}
\end{align}
En reportant~\eqref{eq:cos} dans~\eqref{eq:sin}, on obtient
$\alpha\omega = -\gamma^2\alpha\omega$, donc
\[
  \gamma^2 = -1 \quad\Longrightarrow\quad \gamma = \pm\ii .
\]
Les nombres complexes s'imposent. Avec le choix usuel $\gamma = +\ii$,
\[
  \Psi(x,t) = \cos(kx-\omega t) + \ii\sin(kx-\omega t) = \ee^{\,\ii(kx-\omega t)} .
\]
La relation~\eqref{eq:cos} devient $(-\ii\alpha)\,\omega = V_0 - \beta k^2$. Elle
doit coïncider avec~\eqref{eq:R} écrite sous la forme
$\hbar\,\omega = V_0 + \tfrac{\hbar^2}{2m}\,k^2$. Par identification,
\[
  -\ii\alpha = \hbar \;\Longrightarrow\; \alpha = \ii\hbar,
  \qquad
  -\beta = \frac{\hbar^2}{2m} \;\Longrightarrow\; \beta = -\frac{\hbar^2}{2m}.
\]

\begin{resultat}{Équation de Schrödinger dépendante du temps (1D)}
\[
  \ii\hbar\,\pd{\Psi(x,t)}{t}
  = -\frac{\hbar^2}{2m}\,\pdd{\Psi(x,t)}{x} + V(x,t)\,\Psi(x,t)
\]
En 3D, avec $\nabla^2 = \dfrac{\partial^2}{\partial x^2}+\dfrac{\partial^2}{\partial y^2}+\dfrac{\partial^2}{\partial z^2}$ :
\[
  \ii\hbar\,\pd{\Psi(\mathbf{r},t)}{t} = \Hop\,\Psi(\mathbf{r},t),
  \qquad
  \Hop = -\frac{\hbar^2}{2m}\,\nabla^2 + V(\mathbf{r},t),
\]
où $\Hop$ est l'\emph{opérateur hamiltonien} (énergie totale).
\end{resultat}

\begin{remarque}[Portée de la construction]
Ce n'est pas une démonstration : on a construit l'équation la plus simple
compatible avec (a)--(d) pour un potentiel constant, puis on postule qu'elle
reste valable pour $V(x,t)$ quelconque. Sa validité repose sur l'accord avec
l'expérience. Le choix $\gamma=-\ii$ conduirait à la convention complexe
conjuguée ($\alpha = -\ii\hbar$), physiquement équivalente.
\end{remarque}

\begin{remarque}[Lien avec les opérateurs]
Sur l'onde plane $\ee^{\ii(kx-\omega t)}$ :
$\ii\hbar\,\partial_t\Psi = \hbar\omega\,\Psi = E\,\Psi$ et
$-\ii\hbar\,\partial_x\Psi = \hbar k\,\Psi = p\,\Psi$. L'équation de Schrödinger
traduit donc $E = \frac{p^2}{2m} + V$ avec $E\to\ii\hbar\,\partial_t$ et
$p\to-\ii\hbar\,\partial_x$ : c'est le point de départ du chapitre sur les
opérateurs quantiques.
\end{remarque}

% ---------------------------------------------------------------------
\subsection{Interprétation de Max Born}
% ---------------------------------------------------------------------

La fonction d'onde $\Psi:\mathbb{R}\times\mathbb{R}\to\mathbb{C}$ est à valeurs
complexes : elle n'est ni observable, ni mesurable directement. En revanche, son
module au carré est un réel positif, donc mesurable :
\[
  |\Psi(x,t)|^2 = \Psi^*(x,t)\,\Psi(x,t) \in \mathbb{R}_+ .
\]

\begin{definition}{Interprétation de Born (1926)}
$|\Psi(x,t)|^2$ est la \emph{densité de probabilité de présence} de la particule
en $x$ à l'instant $t$ :
\[
  \dd P\big([x,x+\dd x],t\big) = |\Psi(x,t)|^2\,\dd x ,
  \qquad
  P(a<x<b) = \int_a^b \Psi^*\Psi\,\dd x .
\]
En 3D : $\dd P = |\Psi(\mathbf{r},t)|^2\,\dd^3 r$.
\end{definition}

La particule se trouve forcément quelque part, d'où la \emph{condition de
normalisation}
\[
  \int_{-\infty}^{+\infty} |\Psi(x,t)|^2\,\dd x = 1 .
\]

\begin{figure}[!ht]
\centering
\begin{tikzpicture}[>=Stealth]
  \draw[->] (-3.3,0) -- (3.6,0) node[right] {$x$};
  \draw[->] (0,0) -- (0,2.5) node[above] {$|\Psi(x,t)|^2$};
  \fill[insarouge!20] (0.3,0) -- plot[domain=0.3:1.6, samples=40]
    (\x, {2*exp(-(\x)*(\x)/1.2)}) -- (1.6,0) -- cycle;
  \draw[thick, insagris, domain=-3:3, samples=120]
    plot (\x, {2*exp(-(\x)*(\x)/1.2)});
  \draw[dashed] (0.3,0) -- (0.3,1.86);
  \draw[dashed] (1.6,0) -- (1.6,0.24);
  \node[below] at (0.3,0) {$a$};
  \node[below] at (1.6,0) {$b$};
  \node[anchor=west, align=left] at (1.9,1.7)
    {\small $P(a<x<b)$\\[1mm] \small $=\displaystyle\int_a^b|\Psi|^2\,\dd x$};
  \draw[->] (2.3,1.2) -- (1.15,0.45);
\end{tikzpicture}
\caption{Densité de probabilité de présence ; la probabilité de trouver la
particule entre $a$ et $b$ est l'aire colorée.}
\end{figure}

\begin{remarque}
\begin{itemize}[nosep]
  \item Multiplier $\Psi$ par une phase globale $\ee^{\ii\varphi}$ ne change pas
        $|\Psi|^2$ : la phase globale n'est pas observable.
  \item L'onde plane $\ee^{\ii(kx-\omega t)}$ a $|\Psi|^2 = 1$ partout : elle
        n'est pas normalisable, la particule serait équiprobable à l'infini. Un
        état physique est une \emph{superposition} d'ondes planes (paquet
        d'ondes), possible grâce à la linéarité (c).
  \item Avec la notation de Dirac (chapitres suivants), $|\Psi\rangle$ désigne
        l'état et la normalisation s'écrit $\int\Psi^*\Psi\,\dd x = \langle\Psi|\Psi\rangle = 1$.
\end{itemize}
\end{remarque}

% ---------------------------------------------------------------------
\subsection{Récapitulatif}
% ---------------------------------------------------------------------

\begin{center}
\renewcommand{\arraystretch}{1.5}
\begin{tabular}{@{}ll@{}}
\toprule
\textbf{Relation} & \textbf{Formule} \\
\midrule
Planck--Einstein & $E = hf = \hbar\omega$ \\
de Broglie & $\lambda = \dfrac{h}{p}$, \quad $p = \hbar k$ \\
Énergie totale & $E = \dfrac{p^2}{2m} + V(x,t)$ \\
Onde plane complexe & $\Psi = \ee^{\,\ii(kx-\omega t)}$ \\
Schrödinger (dépendante du temps) & $\ii\hbar\,\partial_t\Psi = \Hop\Psi$, \quad
   $\Hop = -\dfrac{\hbar^2}{2m}\partial_x^2 + V$ \\
Born & $\dd P = |\Psi|^2\,\dd x$, \quad $\displaystyle\int_{-\infty}^{+\infty}|\Psi|^2\,\dd x = 1$ \\
\bottomrule
\end{tabular}
\end{center}
 % 1. Histoire et mécanique quantique
% =====================================================================
%  Chapitre 2 : suite de la section 1 (mécanique quantique)
%  Fichier importé par main.tex avec % =====================================================================
%  Chapitre 2 : suite de la section 1 (mécanique quantique)
%  Fichier importé par main.tex avec % =====================================================================
%  Chapitre 2 : suite de la section 1 (mécanique quantique)
%  Fichier importé par main.tex avec \include{chapitre2}.
%  Il PROLONGE chapitre1.tex : pas de \section ici, les sous-sections
%  1.6 à 1.9 font suite à 1.1 à 1.5 (la suite est dans chapitre3.tex).
%  Pas de préambule non plus : les
%  environnements (definition, resultat, remarque), les macros (\ii, \ee,
%  \dd, \pd, \pdd, \Hop, \ket, \bra, \braket) et les couleurs (insarouge,
%  insagris) sont définis dans main.tex.
%  Convention de signe : p = -i hbar d/dx, cohérente avec l'onde plane
%  e^{i(kx-wt)} de la section 1.3 (les notes de cours donnaient +i hbar).
% =====================================================================

% Macros de Dirac (déjà définies dans main.tex récent ; ce bloc évite l'erreur
% « Undefined control sequence \ket » si main.tex est une ancienne version).
\providecommand{\ket}[1]{|#1\rangle}
\providecommand{\bra}[1]{\langle #1|}
\providecommand{\braket}[2]{\langle #1|#2\rangle}

% ---------------------------------------------------------------------
\subsection{Notation de Dirac}\label{sec:dirac}
% ---------------------------------------------------------------------

Un état s'écrit comme un vecteur colonne, par exemple
$\begin{pmatrix} a_1\\ a_2\end{pmatrix}\in\mathbb{C}^2$ : c'est un élément d'un
espace de Hilbert de dimension~2, noté $\mathcal{H}^2$. Dirac le note simplement
$\ket{a}$ et évite d'écrire les composantes à chaque ligne.

\begin{definition}{Notation de Dirac : ket, bra et bra-ket}
\begin{center}
\begin{tabular}{@{}r@{\hspace{0.5em}}l@{\hspace{2em}}l@{}}
  {\Large\color{insarouge}$\ket{a}$}
    & $=\begin{pmatrix}a_1\\ a_2\end{pmatrix}\in\mathcal{H}^2$
    & \textbf{ket} : vecteur colonne \\[1.6ex]
  {\Large\color{insarouge}$\bra{a}$}
    & $=\big(\ket{a}\big)^\dagger=\begin{pmatrix}a_1^* & a_2^*\end{pmatrix}$
    & \textbf{bra} : vecteur ligne \\[1.6ex]
  {\Large\color{insarouge}$\braket{a}{b}$}
    & $=\begin{pmatrix}a_1^* & a_2^*\end{pmatrix}\begin{pmatrix}b_1\\ b_2\end{pmatrix}
       =a_1^*b_1+a_2^*b_2$
    & \textbf{bra-ket} : produit scalaire \\
\end{tabular}
\end{center}
Le bra est le transposé \emph{conjugué} du ket. Le nom vient de l'anglais
\emph{bracket} (crochet) : en accolant $\bra{a}$~(\emph{bra}) et $\ket{b}$~(\emph{ket})
on obtient $\braket{a}{b}$~(\emph{bracket}), un nombre complexe.

Propriétés : $\braket{b}{a}=\braket{a}{b}^*$ et
$\braket{a}{a}=|a_1|^2+|a_2|^2\in\mathbb{R}_+$ (norme au carré).
\end{definition}

Dans l'autre ordre, le produit d'un ket par un bra est une matrice (un opérateur) :
\[
  \ket{a}\bra{b}
  =\begin{pmatrix}a_1\\ a_2\end{pmatrix}\begin{pmatrix}b_1^* & b_2^*\end{pmatrix}
  =\begin{pmatrix}a_1b_1^* & a_1b_2^*\\ a_2b_1^* & a_2b_2^*\end{pmatrix}.
\]

\begin{remarque}[Pourquoi le conjugué ?]
Pour $\ket{a}=\dfrac{1}{\sqrt2}\begin{pmatrix}1\\ \ii\end{pmatrix}$, on a
$\bra{a}=\dfrac{1}{\sqrt2}\begin{pmatrix}1 & -\ii\end{pmatrix}$ et
$\braket{a}{a}=\tfrac12\,(1+1)=1$. Sans conjugaison on trouverait
$\tfrac12\,(1+\ii^2)=0$, ce qui n'est pas une norme.
\end{remarque}

Un état physique est \emph{normalisé} :
\[
  \braket{\Psi}{\Psi}=1 .
\]
Pour une particule dans l'espace, $\ket{\Psi}$ est une fonction de $x$ (dimension
infinie) et le produit scalaire devient une intégrale,
$\braket{\varphi}{\psi}=\int\varphi^*\psi\,\dd x$. On retrouve la condition de
normalisation de Born (section~\ref{sec:born}) :
\[
  \braket{\Psi}{\Psi} = \int_{-\infty}^{+\infty}\Psi^*\Psi\,\dd x = 1 .
\]
L'espace de Hilbert et le formalisme complet sont repris plus loin dans le cours.

% ---------------------------------------------------------------------
\subsection{Opérateurs quantiques}\label{sec:operateurs}
% ---------------------------------------------------------------------

À chaque grandeur physique (position, quantité de mouvement, énergie) on associe
un \emph{opérateur} $\hat O$ : une règle qui transforme une fonction d'onde en une
autre, $\Psi\mapsto\hat O\Psi$. Le point de départ est l'onde plane de la
section~\ref{sec:construction} : $\Psi=\ee^{\ii(kx-\omega t)}$ vérifie
$-\ii\hbar\,\partial_x\Psi=\hbar k\,\Psi=p\,\Psi$.

\begin{definition}{Opérateur quantité de mouvement}
\[
  \hat p = -\ii\hbar\,\frac{\partial}{\partial x},
  \qquad
  \hat p\,\Psi = -\ii\hbar\,\frac{\partial\Psi}{\partial x}.
\]
L'onde plane est fonction propre de $\hat p$, de valeur propre $p=\hbar k$.
\end{definition}

Les autres opérateurs s'obtiennent en remplaçant $p$ par $\hat p$ dans
l'expression classique. Avec $E_c=\tfrac12 mv^2=p^2/2m$ :
\[
  \hat E_c = \frac{\hat p^{\,2}}{2m}
  = \frac{1}{2m}\Big(-\ii\hbar\frac{\partial}{\partial x}\Big)
                \Big(-\ii\hbar\frac{\partial}{\partial x}\Big)
  = \frac{(-\ii\hbar)^2}{2m}\,\frac{\partial^2}{\partial x^2}
  = -\frac{\hbar^2}{2m}\,\frac{\partial^2}{\partial x^2},
  \qquad
  \hat v = \frac{\hat p}{m} = -\frac{\ii\hbar}{m}\,\frac{\partial}{\partial x}.
\]

\begin{center}
\renewcommand{\arraystretch}{1.6}
\begin{tabular}{@{}lll@{}}
\toprule
\textbf{Grandeur} & \textbf{Classique} & \textbf{Opérateur} \\
\midrule
Position & $x$ & $\hat x = x\times$ \\
Quantité de mouvement & $p$ & $\hat p = -\ii\hbar\,\partial_x$ \\
Vitesse & $v=\dfrac{p}{m}$ & $\hat v = -\dfrac{\ii\hbar}{m}\,\partial_x$ \\
Énergie cinétique & $E_c=\dfrac{p^2}{2m}$ & $\hat E_c = -\dfrac{\hbar^2}{2m}\,\partial_x^2$ \\
Énergie potentielle & $V(x)$ & $\hat V = V(x)\times$ \\
Énergie totale & $E=E_c+V$ & $\Hop = \hat E_c+\hat V = -\dfrac{\hbar^2}{2m}\,\partial_x^2+V(x)$ \\
\bottomrule
\end{tabular}
\end{center}

L'opérateur hamiltonien $\Hop$ représente l'énergie totale. Avec $E\to\ii\hbar\,\partial_t$,
l'équation de Schrödinger de la section~\ref{sec:construction} n'est que la relation
$E=E_c+V$ écrite avec des opérateurs : $\ii\hbar\,\partial_t\Psi=\Hop\Psi$.

\begin{remarque}[Convention de signe]
Le signe de $\hat p$ dépend du choix de l'onde plane $\ee^{\ii(kx-\omega t)}$. La
convention conjuguée $\ee^{-\ii(kx-\omega t)}$ inverse les deux signes
($\hat p=+\ii\hbar\,\partial_x$ et $E\to-\ii\hbar\,\partial_t$) sans changer les
résultats physiques. Comme $\hat p$ intervient au carré, $\hat E_c$ est le même dans
les deux conventions.
\end{remarque}

\begin{definition}{Valeur moyenne d'un opérateur}
Pour une fonction d'onde normalisée $\Psi$,
\[
  \langle\hat O\rangle
  = \int_{-\infty}^{+\infty}\Psi^*(x,t)\;\hat O\,\Psi(x,t)\,\dd x
  = \bra{\Psi}\hat O\ket{\Psi}.
\]
C'est la moyenne des résultats obtenus en mesurant la grandeur sur un grand nombre
de systèmes identiques, tous préparés dans l'état $\Psi$. L'opérateur agit sur
$\Psi$ seul, à droite de $\Psi^*$. Pour la position,
$\langle\hat x\rangle=\int x\,|\Psi|^2\,\dd x$ : c'est l'espérance de la densité
de probabilité de Born.

L'\emph{incertitude} (écart type) sur la grandeur associée à $\hat O$ mesure la
dispersion des résultats autour de cette moyenne :
\[
  \sigma_O=\sqrt{\langle\hat O^{\,2}\rangle-\langle\hat O\rangle^2}\,,
  \qquad
  \langle\hat O^{\,2}\rangle=\int\Psi^*\,\hat O^{\,2}\,\Psi\,\dd x .
\]
Si $\Psi$ est une fonction propre de $\hat O$ ($\hat O\Psi=\lambda\Psi$), alors
$\langle\hat O\rangle=\lambda$, $\langle\hat O^{\,2}\rangle=\lambda^2$ et
$\sigma_O=0$ : la grandeur a une valeur certaine.
\end{definition}

% ---------------------------------------------------------------------
\subsection{Solutions de l'équation de Schrödinger}\label{sec:solutions}
% ---------------------------------------------------------------------

\subsubsection{Du potentiel à la fonction d'onde}

En mécanique classique, l'énergie potentielle $V(x,t)$ fixe la force
$F=-\partial V/\partial x$, puis $\dd p/\dd t=F$ donne la trajectoire. En mécanique
quantique, c'est l'équation de Schrödinger qui joue ce rôle : connaissant $V$, on la
résout pour obtenir $\Psi(x,t)$.

\begin{resultat}{ESDT : équation de Schrödinger dépendante du temps}
\[
  \ii\hbar\,\pd{\Psi(x,t)}{t} = \Hop\,\Psi(x,t),
  \qquad
  \Hop = -\frac{\hbar^2}{2m}\,\pdd{}{x} + V(x,t).
\]
\end{resultat}

Si l'énergie potentielle ne dépend pas du temps, $V(x,t)=V(x)$, la force
$F=-\dd V/\dd x$ (dérivée ordinaire) est constante dans le temps et l'énergie se
conserve. L'équation se résout alors par séparation des variables.

\subsubsection{Séparation des variables}

On cherche des solutions produit d'une fonction de $x$ seul et d'une fonction de $t$
seul :
\[
  \Psi(x,t)=\psi(x)\,\varphi(t).
\]
Dans l'ESDT, les dérivées partielles deviennent des dérivées ordinaires :
\[
  \ii\hbar\,\psi(x)\,\frac{\dd\varphi}{\dd t}
  = \varphi(t)\left[-\frac{\hbar^2}{2m}\,\frac{\dd^2\psi}{\dd x^2}+V(x)\,\psi(x)\right].
\]
On divise par $\psi(x)\,\varphi(t)$ :
\[
  \underbrace{\ii\hbar\,\frac{1}{\varphi}\,\frac{\dd\varphi}{\dd t}}_{\text{ne dépend que de }t}
  \;=\;
  \underbrace{\frac{1}{\psi}\left[-\frac{\hbar^2}{2m}\,\frac{\dd^2\psi}{\dd x^2}+V(x)\,\psi\right]}_{\text{ne dépend que de }x}.
\]
Les deux membres dépendent de variables indépendantes et sont pourtant égaux pour
tout $(x,t)$ : ils sont égaux à une même constante, notée $E$. On obtient deux
équations.

\paragraph{Partie temporelle.} $\ii\hbar\,\dd\varphi/\dd t=E\,\varphi$ donne
\[
  \varphi(t)=\ee^{-\ii Et/\hbar}=\ee^{-\ii\omega t},
  \qquad \omega=\frac{E}{\hbar},
\]
la constante multiplicative étant absorbée dans $\psi$. On retrouve $E=\hbar\omega$ :
la constante $E$ est bien l'énergie de la particule.

\paragraph{Partie spatiale.} C'est l'ESIT ci-dessous.

\begin{resultat}{ESIT : équation de Schrödinger indépendante du temps}
\[
  -\frac{\hbar^2}{2m}\,\frac{\dd^2\psi(x)}{\dd x^2}+V(x)\,\psi(x)=E\,\psi(x)
  \qquad\Longleftrightarrow\qquad
  \Hop\,\psi(x)=E\,\psi(x).
\]
C'est un \emph{problème aux valeurs propres} : $E$ est une valeur propre de $\Hop$
(énergie permise) et $\psi$ la fonction propre associée. Chaque solution
$(E_n,\psi_n)$ donne un \emph{état stationnaire}
\[
  \Psi_n(x,t)=\psi_n(x)\,\ee^{-\ii E_nt/\hbar}.
\]
\end{resultat}

Dans un état stationnaire, $|\Psi_n(x,t)|^2=|\psi_n(x)|^2$ ne dépend pas du temps :
le facteur temporel est une phase globale, non observable. Les conditions physiques
($\psi$ continue et normalisable) ne sont satisfaites que pour certaines valeurs de
$E$ : c'est l'origine de la quantification de l'énergie, illustrée ci-dessous.

Comme l'ESDT est linéaire (hypothèse~(c) de la section~\ref{sec:construction}), toute
superposition $\Psi=\sum_n c_n\,\Psi_n$ est aussi solution, normalisée si
$\sum_n|c_n|^2=1$ et si les $\psi_n$ sont orthonormées. Une telle superposition
n'est en général pas stationnaire : $|\Psi|^2$ contient des termes croisés en
$\ee^{\ii(E_m-E_n)t/\hbar}$.

\subsubsection{Exemple : puits de potentiel infini}

Une particule de masse $m$ (un électron par exemple) est enfermée entre deux parois
impénétrables, en $x=0$ et $x=a$ :
\[
  V(x)=
  \begin{cases}
    0 & \text{si } 0<x<a,\\
    +\infty & \text{ailleurs.}
  \end{cases}
\]
La particule ne peut pas se trouver hors du puits : $\psi=0$ pour $x\le0$ et
$x\ge a$. La fonction $\psi$ est continue (mais pas sa dérivée, à cause du saut
infini de $V$), d'où les conditions aux limites
\[
  \psi(0)=\psi(a)=0 .
\]

\paragraph{Opérateur hamiltonien.} Dans le puits $V=0$, donc $\Hop$ se réduit à
l'opérateur énergie cinétique :
\[
  \Hop=\hat E_c=-\frac{\hbar^2}{2m}\,\frac{\dd^2}{\dd x^2}\qquad(0<x<a).
\]

\paragraph{Résolution de l'ESIT.} Avec $V=0$, l'ESIT s'écrit
$\psi''=-k^2\psi$ avec $k=\sqrt{2mE}/\hbar$. On prend $E>0$ : pour $E\le0$, la seule
solution vérifiant $\psi(0)=\psi(a)=0$ est $\psi\equiv0$. La solution générale est
\[
  \psi(x)=A\sin(kx)+B\cos(kx).
\]
\begin{itemize}[nosep]
  \item $\psi(0)=0$ donne $B=0$ ;
  \item $\psi(a)=A\sin(ka)=0$ avec $A\neq0$ (sinon $\psi\equiv0$ : pas de
        particule) donne $ka=n\pi$, soit
        $k_n=\dfrac{n\pi}{a}$ avec $n=1,2,3,\dots$
        ($n=0$ donne $\psi\equiv0$, et $n<0$ redonne les mêmes états au signe près).
\end{itemize}
Cela revient à $\lambda_n=2\pi/k_n=2a/n$, soit $a=n\,\lambda_n/2$ : $n$ demi-longueurs
d'onde tiennent dans le puits, comme pour une onde stationnaire.

\begin{resultat}{Énergie quantifiée}
\[
  E_n=\frac{\hbar^2k_n^2}{2m}=\frac{n^2\pi^2\hbar^2}{2ma^2}=\frac{n^2h^2}{8ma^2}=n^2E_1,
  \qquad n=1,2,3,\dots
\]
\textbf{Énergie fondamentale} (niveau $n=1$, le plus bas) :
$\displaystyle E_1=\frac{\pi^2\hbar^2}{2ma^2}=\frac{h^2}{8ma^2}>0$.
\end{resultat}

\begin{remarque}[Lecture du résultat]
\begin{itemize}[nosep]
  \item Les énergies sont \emph{discrètes} : $E_{n+1}-E_n=(2n+1)E_1$. Comme
        $E_n\propto1/a^2$, plus le puits est étroit, plus les énergies sont grandes.
  \item $E_1\neq0$ : une particule confinée ne peut pas être au repos (énergie de
        point zéro), conséquence du principe d'incertitude de Heisenberg.
  \item Ordre de grandeur : pour un électron avec $a=1$~nm,
        $E_1\approx0{,}38$~eV et $E_2=4E_1\approx1{,}5$~eV. Pour $a=1$~cm,
        $E_1\approx4\times10^{-15}$~eV : la quantification est indétectable.
\end{itemize}
\end{remarque}

\paragraph{Normalisation.} Le module de $\ee^{-\ii E_nt/\hbar}$ vaut 1, donc
$|\Psi_n|^2=A^2\sin^2(n\pi x/a)$ et la condition de Born impose
\[
  \int_0^a|\Psi_n|^2\,\dd x
  = A^2\int_0^a\frac{1-\cos\!\big(2n\pi x/a\big)}{2}\,\dd x
  = A^2\,\frac{a}{2}=1
  \quad\Longrightarrow\quad A=\sqrt{\frac{2}{a}},
\]
car l'intégrale du cosinus sur $n$ périodes est nulle.

\begin{resultat}{États stationnaires du puits infini}
\[
  \Psi_n(x,t)=\sqrt{\frac{2}{a}}\;\sin\!\Big(\frac{n\pi x}{a}\Big)\,\ee^{-\ii E_nt/\hbar}
  \quad(0<x<a),
  \qquad
  |\Psi_n|^2=\frac{2}{a}\,\sin^2\!\Big(\frac{n\pi x}{a}\Big).
\]
$\psi_n$ possède $n-1$ nœuds à l'intérieur du puits. Pour $n=1$ (état fondamental,
d'énergie $E_1$) :
\[
  \Psi_1(x,t)=\sqrt{\frac{2}{a}}\;\sin\!\Big(\frac{\pi x}{a}\Big)\,\ee^{-\ii E_1t/\hbar}.
\]
\end{resultat}

\begin{figure}[!ht]
\centering
\begin{tikzpicture}[>=Stealth, x=1cm, y=0.95cm]
  % ---------- (a) niveaux d'énergie et fonctions propres ----------
  \begin{scope}
    \node[above, font=\small] at (2,5.0) {(a) niveaux $E_n$ et fonctions propres $\psi_n$};
    \draw[->] (-2.5,0) -- (-2.5,4.8) node[above] {$E$};
    \draw[very thick, insarouge] (0,4.4) -- (0,0) -- (4,0) -- (4,4.4);
    \node[insarouge, right, font=\small] at (4.05,3.9) {$V=+\infty$};
    \node[insarouge, below, font=\small] at (2,0) {$V=0$};
    \node[below, font=\small] at (0,-0.35) {$0$};
    \node[below, font=\small] at (4,-0.35) {$a$};
    \node[left, font=\small] at (-0.1,0.4) {$E_1$};
    \node[left, font=\small] at (-0.1,1.6) {$E_2=4E_1$};
    \node[left, font=\small] at (-0.1,3.6) {$E_3=9E_1$};
    \foreach \n/\lev in {1/0.4, 2/1.6, 3/3.6} {
      \draw[dashed, insagris!70] (0,\lev) -- (4,\lev);
      \draw[thick, insagris, domain=0:1, samples=80]
        plot ({4*\x}, {\lev+0.5*sin(\n*180*\x)});
    }
  \end{scope}
  % ---------- (b) densités de probabilité ----------
  \begin{scope}[xshift=7cm]
    \node[above, font=\small] at (2,5.0) {(b) densités de probabilité $|\psi_n|^2$};
    \draw[->] (-0.3,0) -- (4.6,0) node[right] {$x$};
    \draw[thick, insarouge] (0,0) -- (0,4.4) (4,0) -- (4,4.4);
    \node[below, font=\small] at (0,0) {$0$};
    \node[below, font=\small] at (4,0) {$a$};
    \node[below, font=\small] at (2,0) {$a/2$};
    \foreach \n/\lev in {1/0.4, 2/1.6, 3/3.6} {
      \draw[dashed, insagris!70] (0,\lev) -- (4,\lev);
      \fill[insarouge!20] (0,\lev) -- plot[domain=0:1, samples=80]
        ({4*\x}, {\lev+0.6*(sin(\n*180*\x))^2}) -- (4,\lev) -- cycle;
      \draw[thick, insarouge, domain=0:1, samples=80]
        plot ({4*\x}, {\lev+0.6*(sin(\n*180*\x))^2});
    }
    \draw[dotted, thick, insagris] (2,0) -- (2,4.4);
    \node[above, font=\small] at (2,4.4) {$\langle\hat x\rangle=a/2$};
  \end{scope}
\end{tikzpicture}
\caption{Puits de potentiel infini de largeur $a$. (a) Niveaux d'énergie
$E_n=n^2E_1$ et fonctions propres $\psi_n$. (b) Densités de probabilité
$|\psi_n|^2$, symétriques par rapport à $a/2$.}
\end{figure}

\subsubsection{Position moyenne dans le puits}

D'après la définition de la valeur moyenne, avec $\hat x=x\times$ :
\[
  \langle\hat x\rangle
  = \int_0^a\Psi_n^*(x,t)\;x\;\Psi_n(x,t)\,\dd x
  = \frac{2}{a}\int_0^a x\,\sin^2\!\Big(\frac{n\pi x}{a}\Big)\dd x
  = \frac{1}{a}\int_0^a x\Big(1-\cos\frac{2n\pi x}{a}\Big)\dd x .
\]
Le premier terme vaut $\dfrac1a\cdot\dfrac{a^2}{2}=\dfrac a2$. Pour le second, une
intégration par parties donne
\[
  \int_0^a x\cos\frac{2n\pi x}{a}\,\dd x
  =\Big[\frac{a\,x}{2n\pi}\sin\frac{2n\pi x}{a}\Big]_0^a
   +\frac{a^2}{4n^2\pi^2}\Big[\cos\frac{2n\pi x}{a}\Big]_0^a = 0,
\]
puisque $\sin(2n\pi)=0$ et $\cos(2n\pi)=1$.

\begin{resultat}{Position moyenne}
\[
  \langle\hat x\rangle=\frac{a}{2}\qquad\text{pour tout }n\text{ et à tout instant.}
\]
\end{resultat}

Les phases $\ee^{\pm\ii E_nt/\hbar}$ se compensent dans $\Psi^*\Psi$. Le résultat se
voit d'ailleurs sans calcul, car $|\Psi_n|^2$ est symétrique par rapport à $a/2$.

\begin{remarque}[Moyenne et position la plus probable]
Pour $n=2$, la densité de probabilité s'annule en $x=a/2$ alors que
$\langle\hat x\rangle=a/2$ : la valeur moyenne n'est pas la position la plus
probable, et la particule n'est jamais trouvée là.
\end{remarque}

\paragraph{Incertitude sur l'énergie.} Comme $\Hop\Psi_n=E_n\Psi_n$, on a aussi
$\Hop^{2}\Psi_n=E_n^2\Psi_n$, donc $\langle\Hop\rangle=E_n$, $\langle\Hop^{2}\rangle=E_n^2$
et
\[
  \sigma_E=\sqrt{\langle\Hop^{2}\rangle-\langle\Hop\rangle^2}=\sqrt{E_n^2-E_n^2}=0 .
\]
Dans un état propre (stationnaire) l'énergie est parfaitement déterminée, alors que la
position ne l'est pas. Que se passe-t-il si la particule n'est pas dans un état propre ?
C'est l'objet de la section~\ref{sec:superposition}.

% ---------------------------------------------------------------------
\subsection{Bilan : opérateurs et puits infini}
% ---------------------------------------------------------------------

\begin{center}
\small\renewcommand{\arraystretch}{1.25}
\begin{tabular}{@{}ll@{}}
\toprule
\textbf{Notion} & \textbf{Formule} \\
\midrule
Bra-ket & $\braket{a}{b}=a_1^*b_1+a_2^*b_2$, \quad
          $\braket{\Psi}{\Psi}=\displaystyle\int\Psi^*\Psi\,\dd x=1$ \\
Opérateurs & $\hat p=-\ii\hbar\,\partial_x$, \quad
             $\hat E_c=\dfrac{\hat p^{\,2}}{2m}=-\dfrac{\hbar^2}{2m}\,\partial_x^2$, \quad
             $\Hop=\hat E_c+V(x)$ \\
Valeur moyenne, incertitude & $\langle\hat O\rangle=\displaystyle\int\Psi^*\,\hat O\,\Psi\,\dd x$,
          \quad $\sigma_O=\sqrt{\langle\hat O^{\,2}\rangle-\langle\hat O\rangle^2}$ \\
ESDT & $\ii\hbar\,\partial_t\Psi=\Hop\,\Psi$ \\
Séparation ($V=V(x)$) & $\Psi=\psi(x)\,\ee^{-\ii Et/\hbar}$, \quad ESIT : $\Hop\,\psi=E\,\psi$ \\
Puits infini & $\psi_n=\sqrt{2/a}\,\sin(n\pi x/a)$, \quad
               $E_n=\dfrac{n^2\pi^2\hbar^2}{2ma^2}$, \quad $E_1=\dfrac{\pi^2\hbar^2}{2ma^2}$ \\
Moyennes dans $\Psi_n$ & $\langle\hat x\rangle=a/2$, \quad $\langle\Hop\rangle=E_n$, \quad $\sigma_E=0$ \\
\bottomrule
\end{tabular}
\end{center}
.
%  Il PROLONGE chapitre1.tex : pas de \section ici, les sous-sections
%  1.6 à 1.9 font suite à 1.1 à 1.5 (la suite est dans chapitre3.tex).
%  Pas de préambule non plus : les
%  environnements (definition, resultat, remarque), les macros (\ii, \ee,
%  \dd, \pd, \pdd, \Hop, \ket, \bra, \braket) et les couleurs (insarouge,
%  insagris) sont définis dans main.tex.
%  Convention de signe : p = -i hbar d/dx, cohérente avec l'onde plane
%  e^{i(kx-wt)} de la section 1.3 (les notes de cours donnaient +i hbar).
% =====================================================================

% Macros de Dirac (déjà définies dans main.tex récent ; ce bloc évite l'erreur
% « Undefined control sequence \ket » si main.tex est une ancienne version).
\providecommand{\ket}[1]{|#1\rangle}
\providecommand{\bra}[1]{\langle #1|}
\providecommand{\braket}[2]{\langle #1|#2\rangle}

% ---------------------------------------------------------------------
\subsection{Notation de Dirac}\label{sec:dirac}
% ---------------------------------------------------------------------

Un état s'écrit comme un vecteur colonne, par exemple
$\begin{pmatrix} a_1\\ a_2\end{pmatrix}\in\mathbb{C}^2$ : c'est un élément d'un
espace de Hilbert de dimension~2, noté $\mathcal{H}^2$. Dirac le note simplement
$\ket{a}$ et évite d'écrire les composantes à chaque ligne.

\begin{definition}{Notation de Dirac : ket, bra et bra-ket}
\begin{center}
\begin{tabular}{@{}r@{\hspace{0.5em}}l@{\hspace{2em}}l@{}}
  {\Large\color{insarouge}$\ket{a}$}
    & $=\begin{pmatrix}a_1\\ a_2\end{pmatrix}\in\mathcal{H}^2$
    & \textbf{ket} : vecteur colonne \\[1.6ex]
  {\Large\color{insarouge}$\bra{a}$}
    & $=\big(\ket{a}\big)^\dagger=\begin{pmatrix}a_1^* & a_2^*\end{pmatrix}$
    & \textbf{bra} : vecteur ligne \\[1.6ex]
  {\Large\color{insarouge}$\braket{a}{b}$}
    & $=\begin{pmatrix}a_1^* & a_2^*\end{pmatrix}\begin{pmatrix}b_1\\ b_2\end{pmatrix}
       =a_1^*b_1+a_2^*b_2$
    & \textbf{bra-ket} : produit scalaire \\
\end{tabular}
\end{center}
Le bra est le transposé \emph{conjugué} du ket. Le nom vient de l'anglais
\emph{bracket} (crochet) : en accolant $\bra{a}$~(\emph{bra}) et $\ket{b}$~(\emph{ket})
on obtient $\braket{a}{b}$~(\emph{bracket}), un nombre complexe.

Propriétés : $\braket{b}{a}=\braket{a}{b}^*$ et
$\braket{a}{a}=|a_1|^2+|a_2|^2\in\mathbb{R}_+$ (norme au carré).
\end{definition}

Dans l'autre ordre, le produit d'un ket par un bra est une matrice (un opérateur) :
\[
  \ket{a}\bra{b}
  =\begin{pmatrix}a_1\\ a_2\end{pmatrix}\begin{pmatrix}b_1^* & b_2^*\end{pmatrix}
  =\begin{pmatrix}a_1b_1^* & a_1b_2^*\\ a_2b_1^* & a_2b_2^*\end{pmatrix}.
\]

\begin{remarque}[Pourquoi le conjugué ?]
Pour $\ket{a}=\dfrac{1}{\sqrt2}\begin{pmatrix}1\\ \ii\end{pmatrix}$, on a
$\bra{a}=\dfrac{1}{\sqrt2}\begin{pmatrix}1 & -\ii\end{pmatrix}$ et
$\braket{a}{a}=\tfrac12\,(1+1)=1$. Sans conjugaison on trouverait
$\tfrac12\,(1+\ii^2)=0$, ce qui n'est pas une norme.
\end{remarque}

Un état physique est \emph{normalisé} :
\[
  \braket{\Psi}{\Psi}=1 .
\]
Pour une particule dans l'espace, $\ket{\Psi}$ est une fonction de $x$ (dimension
infinie) et le produit scalaire devient une intégrale,
$\braket{\varphi}{\psi}=\int\varphi^*\psi\,\dd x$. On retrouve la condition de
normalisation de Born (section~\ref{sec:born}) :
\[
  \braket{\Psi}{\Psi} = \int_{-\infty}^{+\infty}\Psi^*\Psi\,\dd x = 1 .
\]
L'espace de Hilbert et le formalisme complet sont repris plus loin dans le cours.

% ---------------------------------------------------------------------
\subsection{Opérateurs quantiques}\label{sec:operateurs}
% ---------------------------------------------------------------------

À chaque grandeur physique (position, quantité de mouvement, énergie) on associe
un \emph{opérateur} $\hat O$ : une règle qui transforme une fonction d'onde en une
autre, $\Psi\mapsto\hat O\Psi$. Le point de départ est l'onde plane de la
section~\ref{sec:construction} : $\Psi=\ee^{\ii(kx-\omega t)}$ vérifie
$-\ii\hbar\,\partial_x\Psi=\hbar k\,\Psi=p\,\Psi$.

\begin{definition}{Opérateur quantité de mouvement}
\[
  \hat p = -\ii\hbar\,\frac{\partial}{\partial x},
  \qquad
  \hat p\,\Psi = -\ii\hbar\,\frac{\partial\Psi}{\partial x}.
\]
L'onde plane est fonction propre de $\hat p$, de valeur propre $p=\hbar k$.
\end{definition}

Les autres opérateurs s'obtiennent en remplaçant $p$ par $\hat p$ dans
l'expression classique. Avec $E_c=\tfrac12 mv^2=p^2/2m$ :
\[
  \hat E_c = \frac{\hat p^{\,2}}{2m}
  = \frac{1}{2m}\Big(-\ii\hbar\frac{\partial}{\partial x}\Big)
                \Big(-\ii\hbar\frac{\partial}{\partial x}\Big)
  = \frac{(-\ii\hbar)^2}{2m}\,\frac{\partial^2}{\partial x^2}
  = -\frac{\hbar^2}{2m}\,\frac{\partial^2}{\partial x^2},
  \qquad
  \hat v = \frac{\hat p}{m} = -\frac{\ii\hbar}{m}\,\frac{\partial}{\partial x}.
\]

\begin{center}
\renewcommand{\arraystretch}{1.6}
\begin{tabular}{@{}lll@{}}
\toprule
\textbf{Grandeur} & \textbf{Classique} & \textbf{Opérateur} \\
\midrule
Position & $x$ & $\hat x = x\times$ \\
Quantité de mouvement & $p$ & $\hat p = -\ii\hbar\,\partial_x$ \\
Vitesse & $v=\dfrac{p}{m}$ & $\hat v = -\dfrac{\ii\hbar}{m}\,\partial_x$ \\
Énergie cinétique & $E_c=\dfrac{p^2}{2m}$ & $\hat E_c = -\dfrac{\hbar^2}{2m}\,\partial_x^2$ \\
Énergie potentielle & $V(x)$ & $\hat V = V(x)\times$ \\
Énergie totale & $E=E_c+V$ & $\Hop = \hat E_c+\hat V = -\dfrac{\hbar^2}{2m}\,\partial_x^2+V(x)$ \\
\bottomrule
\end{tabular}
\end{center}

L'opérateur hamiltonien $\Hop$ représente l'énergie totale. Avec $E\to\ii\hbar\,\partial_t$,
l'équation de Schrödinger de la section~\ref{sec:construction} n'est que la relation
$E=E_c+V$ écrite avec des opérateurs : $\ii\hbar\,\partial_t\Psi=\Hop\Psi$.

\begin{remarque}[Convention de signe]
Le signe de $\hat p$ dépend du choix de l'onde plane $\ee^{\ii(kx-\omega t)}$. La
convention conjuguée $\ee^{-\ii(kx-\omega t)}$ inverse les deux signes
($\hat p=+\ii\hbar\,\partial_x$ et $E\to-\ii\hbar\,\partial_t$) sans changer les
résultats physiques. Comme $\hat p$ intervient au carré, $\hat E_c$ est le même dans
les deux conventions.
\end{remarque}

\begin{definition}{Valeur moyenne d'un opérateur}
Pour une fonction d'onde normalisée $\Psi$,
\[
  \langle\hat O\rangle
  = \int_{-\infty}^{+\infty}\Psi^*(x,t)\;\hat O\,\Psi(x,t)\,\dd x
  = \bra{\Psi}\hat O\ket{\Psi}.
\]
C'est la moyenne des résultats obtenus en mesurant la grandeur sur un grand nombre
de systèmes identiques, tous préparés dans l'état $\Psi$. L'opérateur agit sur
$\Psi$ seul, à droite de $\Psi^*$. Pour la position,
$\langle\hat x\rangle=\int x\,|\Psi|^2\,\dd x$ : c'est l'espérance de la densité
de probabilité de Born.

L'\emph{incertitude} (écart type) sur la grandeur associée à $\hat O$ mesure la
dispersion des résultats autour de cette moyenne :
\[
  \sigma_O=\sqrt{\langle\hat O^{\,2}\rangle-\langle\hat O\rangle^2}\,,
  \qquad
  \langle\hat O^{\,2}\rangle=\int\Psi^*\,\hat O^{\,2}\,\Psi\,\dd x .
\]
Si $\Psi$ est une fonction propre de $\hat O$ ($\hat O\Psi=\lambda\Psi$), alors
$\langle\hat O\rangle=\lambda$, $\langle\hat O^{\,2}\rangle=\lambda^2$ et
$\sigma_O=0$ : la grandeur a une valeur certaine.
\end{definition}

% ---------------------------------------------------------------------
\subsection{Solutions de l'équation de Schrödinger}\label{sec:solutions}
% ---------------------------------------------------------------------

\subsubsection{Du potentiel à la fonction d'onde}

En mécanique classique, l'énergie potentielle $V(x,t)$ fixe la force
$F=-\partial V/\partial x$, puis $\dd p/\dd t=F$ donne la trajectoire. En mécanique
quantique, c'est l'équation de Schrödinger qui joue ce rôle : connaissant $V$, on la
résout pour obtenir $\Psi(x,t)$.

\begin{resultat}{ESDT : équation de Schrödinger dépendante du temps}
\[
  \ii\hbar\,\pd{\Psi(x,t)}{t} = \Hop\,\Psi(x,t),
  \qquad
  \Hop = -\frac{\hbar^2}{2m}\,\pdd{}{x} + V(x,t).
\]
\end{resultat}

Si l'énergie potentielle ne dépend pas du temps, $V(x,t)=V(x)$, la force
$F=-\dd V/\dd x$ (dérivée ordinaire) est constante dans le temps et l'énergie se
conserve. L'équation se résout alors par séparation des variables.

\subsubsection{Séparation des variables}

On cherche des solutions produit d'une fonction de $x$ seul et d'une fonction de $t$
seul :
\[
  \Psi(x,t)=\psi(x)\,\varphi(t).
\]
Dans l'ESDT, les dérivées partielles deviennent des dérivées ordinaires :
\[
  \ii\hbar\,\psi(x)\,\frac{\dd\varphi}{\dd t}
  = \varphi(t)\left[-\frac{\hbar^2}{2m}\,\frac{\dd^2\psi}{\dd x^2}+V(x)\,\psi(x)\right].
\]
On divise par $\psi(x)\,\varphi(t)$ :
\[
  \underbrace{\ii\hbar\,\frac{1}{\varphi}\,\frac{\dd\varphi}{\dd t}}_{\text{ne dépend que de }t}
  \;=\;
  \underbrace{\frac{1}{\psi}\left[-\frac{\hbar^2}{2m}\,\frac{\dd^2\psi}{\dd x^2}+V(x)\,\psi\right]}_{\text{ne dépend que de }x}.
\]
Les deux membres dépendent de variables indépendantes et sont pourtant égaux pour
tout $(x,t)$ : ils sont égaux à une même constante, notée $E$. On obtient deux
équations.

\paragraph{Partie temporelle.} $\ii\hbar\,\dd\varphi/\dd t=E\,\varphi$ donne
\[
  \varphi(t)=\ee^{-\ii Et/\hbar}=\ee^{-\ii\omega t},
  \qquad \omega=\frac{E}{\hbar},
\]
la constante multiplicative étant absorbée dans $\psi$. On retrouve $E=\hbar\omega$ :
la constante $E$ est bien l'énergie de la particule.

\paragraph{Partie spatiale.} C'est l'ESIT ci-dessous.

\begin{resultat}{ESIT : équation de Schrödinger indépendante du temps}
\[
  -\frac{\hbar^2}{2m}\,\frac{\dd^2\psi(x)}{\dd x^2}+V(x)\,\psi(x)=E\,\psi(x)
  \qquad\Longleftrightarrow\qquad
  \Hop\,\psi(x)=E\,\psi(x).
\]
C'est un \emph{problème aux valeurs propres} : $E$ est une valeur propre de $\Hop$
(énergie permise) et $\psi$ la fonction propre associée. Chaque solution
$(E_n,\psi_n)$ donne un \emph{état stationnaire}
\[
  \Psi_n(x,t)=\psi_n(x)\,\ee^{-\ii E_nt/\hbar}.
\]
\end{resultat}

Dans un état stationnaire, $|\Psi_n(x,t)|^2=|\psi_n(x)|^2$ ne dépend pas du temps :
le facteur temporel est une phase globale, non observable. Les conditions physiques
($\psi$ continue et normalisable) ne sont satisfaites que pour certaines valeurs de
$E$ : c'est l'origine de la quantification de l'énergie, illustrée ci-dessous.

Comme l'ESDT est linéaire (hypothèse~(c) de la section~\ref{sec:construction}), toute
superposition $\Psi=\sum_n c_n\,\Psi_n$ est aussi solution, normalisée si
$\sum_n|c_n|^2=1$ et si les $\psi_n$ sont orthonormées. Une telle superposition
n'est en général pas stationnaire : $|\Psi|^2$ contient des termes croisés en
$\ee^{\ii(E_m-E_n)t/\hbar}$.

\subsubsection{Exemple : puits de potentiel infini}

Une particule de masse $m$ (un électron par exemple) est enfermée entre deux parois
impénétrables, en $x=0$ et $x=a$ :
\[
  V(x)=
  \begin{cases}
    0 & \text{si } 0<x<a,\\
    +\infty & \text{ailleurs.}
  \end{cases}
\]
La particule ne peut pas se trouver hors du puits : $\psi=0$ pour $x\le0$ et
$x\ge a$. La fonction $\psi$ est continue (mais pas sa dérivée, à cause du saut
infini de $V$), d'où les conditions aux limites
\[
  \psi(0)=\psi(a)=0 .
\]

\paragraph{Opérateur hamiltonien.} Dans le puits $V=0$, donc $\Hop$ se réduit à
l'opérateur énergie cinétique :
\[
  \Hop=\hat E_c=-\frac{\hbar^2}{2m}\,\frac{\dd^2}{\dd x^2}\qquad(0<x<a).
\]

\paragraph{Résolution de l'ESIT.} Avec $V=0$, l'ESIT s'écrit
$\psi''=-k^2\psi$ avec $k=\sqrt{2mE}/\hbar$. On prend $E>0$ : pour $E\le0$, la seule
solution vérifiant $\psi(0)=\psi(a)=0$ est $\psi\equiv0$. La solution générale est
\[
  \psi(x)=A\sin(kx)+B\cos(kx).
\]
\begin{itemize}[nosep]
  \item $\psi(0)=0$ donne $B=0$ ;
  \item $\psi(a)=A\sin(ka)=0$ avec $A\neq0$ (sinon $\psi\equiv0$ : pas de
        particule) donne $ka=n\pi$, soit
        $k_n=\dfrac{n\pi}{a}$ avec $n=1,2,3,\dots$
        ($n=0$ donne $\psi\equiv0$, et $n<0$ redonne les mêmes états au signe près).
\end{itemize}
Cela revient à $\lambda_n=2\pi/k_n=2a/n$, soit $a=n\,\lambda_n/2$ : $n$ demi-longueurs
d'onde tiennent dans le puits, comme pour une onde stationnaire.

\begin{resultat}{Énergie quantifiée}
\[
  E_n=\frac{\hbar^2k_n^2}{2m}=\frac{n^2\pi^2\hbar^2}{2ma^2}=\frac{n^2h^2}{8ma^2}=n^2E_1,
  \qquad n=1,2,3,\dots
\]
\textbf{Énergie fondamentale} (niveau $n=1$, le plus bas) :
$\displaystyle E_1=\frac{\pi^2\hbar^2}{2ma^2}=\frac{h^2}{8ma^2}>0$.
\end{resultat}

\begin{remarque}[Lecture du résultat]
\begin{itemize}[nosep]
  \item Les énergies sont \emph{discrètes} : $E_{n+1}-E_n=(2n+1)E_1$. Comme
        $E_n\propto1/a^2$, plus le puits est étroit, plus les énergies sont grandes.
  \item $E_1\neq0$ : une particule confinée ne peut pas être au repos (énergie de
        point zéro), conséquence du principe d'incertitude de Heisenberg.
  \item Ordre de grandeur : pour un électron avec $a=1$~nm,
        $E_1\approx0{,}38$~eV et $E_2=4E_1\approx1{,}5$~eV. Pour $a=1$~cm,
        $E_1\approx4\times10^{-15}$~eV : la quantification est indétectable.
\end{itemize}
\end{remarque}

\paragraph{Normalisation.} Le module de $\ee^{-\ii E_nt/\hbar}$ vaut 1, donc
$|\Psi_n|^2=A^2\sin^2(n\pi x/a)$ et la condition de Born impose
\[
  \int_0^a|\Psi_n|^2\,\dd x
  = A^2\int_0^a\frac{1-\cos\!\big(2n\pi x/a\big)}{2}\,\dd x
  = A^2\,\frac{a}{2}=1
  \quad\Longrightarrow\quad A=\sqrt{\frac{2}{a}},
\]
car l'intégrale du cosinus sur $n$ périodes est nulle.

\begin{resultat}{États stationnaires du puits infini}
\[
  \Psi_n(x,t)=\sqrt{\frac{2}{a}}\;\sin\!\Big(\frac{n\pi x}{a}\Big)\,\ee^{-\ii E_nt/\hbar}
  \quad(0<x<a),
  \qquad
  |\Psi_n|^2=\frac{2}{a}\,\sin^2\!\Big(\frac{n\pi x}{a}\Big).
\]
$\psi_n$ possède $n-1$ nœuds à l'intérieur du puits. Pour $n=1$ (état fondamental,
d'énergie $E_1$) :
\[
  \Psi_1(x,t)=\sqrt{\frac{2}{a}}\;\sin\!\Big(\frac{\pi x}{a}\Big)\,\ee^{-\ii E_1t/\hbar}.
\]
\end{resultat}

\begin{figure}[!ht]
\centering
\begin{tikzpicture}[>=Stealth, x=1cm, y=0.95cm]
  % ---------- (a) niveaux d'énergie et fonctions propres ----------
  \begin{scope}
    \node[above, font=\small] at (2,5.0) {(a) niveaux $E_n$ et fonctions propres $\psi_n$};
    \draw[->] (-2.5,0) -- (-2.5,4.8) node[above] {$E$};
    \draw[very thick, insarouge] (0,4.4) -- (0,0) -- (4,0) -- (4,4.4);
    \node[insarouge, right, font=\small] at (4.05,3.9) {$V=+\infty$};
    \node[insarouge, below, font=\small] at (2,0) {$V=0$};
    \node[below, font=\small] at (0,-0.35) {$0$};
    \node[below, font=\small] at (4,-0.35) {$a$};
    \node[left, font=\small] at (-0.1,0.4) {$E_1$};
    \node[left, font=\small] at (-0.1,1.6) {$E_2=4E_1$};
    \node[left, font=\small] at (-0.1,3.6) {$E_3=9E_1$};
    \foreach \n/\lev in {1/0.4, 2/1.6, 3/3.6} {
      \draw[dashed, insagris!70] (0,\lev) -- (4,\lev);
      \draw[thick, insagris, domain=0:1, samples=80]
        plot ({4*\x}, {\lev+0.5*sin(\n*180*\x)});
    }
  \end{scope}
  % ---------- (b) densités de probabilité ----------
  \begin{scope}[xshift=7cm]
    \node[above, font=\small] at (2,5.0) {(b) densités de probabilité $|\psi_n|^2$};
    \draw[->] (-0.3,0) -- (4.6,0) node[right] {$x$};
    \draw[thick, insarouge] (0,0) -- (0,4.4) (4,0) -- (4,4.4);
    \node[below, font=\small] at (0,0) {$0$};
    \node[below, font=\small] at (4,0) {$a$};
    \node[below, font=\small] at (2,0) {$a/2$};
    \foreach \n/\lev in {1/0.4, 2/1.6, 3/3.6} {
      \draw[dashed, insagris!70] (0,\lev) -- (4,\lev);
      \fill[insarouge!20] (0,\lev) -- plot[domain=0:1, samples=80]
        ({4*\x}, {\lev+0.6*(sin(\n*180*\x))^2}) -- (4,\lev) -- cycle;
      \draw[thick, insarouge, domain=0:1, samples=80]
        plot ({4*\x}, {\lev+0.6*(sin(\n*180*\x))^2});
    }
    \draw[dotted, thick, insagris] (2,0) -- (2,4.4);
    \node[above, font=\small] at (2,4.4) {$\langle\hat x\rangle=a/2$};
  \end{scope}
\end{tikzpicture}
\caption{Puits de potentiel infini de largeur $a$. (a) Niveaux d'énergie
$E_n=n^2E_1$ et fonctions propres $\psi_n$. (b) Densités de probabilité
$|\psi_n|^2$, symétriques par rapport à $a/2$.}
\end{figure}

\subsubsection{Position moyenne dans le puits}

D'après la définition de la valeur moyenne, avec $\hat x=x\times$ :
\[
  \langle\hat x\rangle
  = \int_0^a\Psi_n^*(x,t)\;x\;\Psi_n(x,t)\,\dd x
  = \frac{2}{a}\int_0^a x\,\sin^2\!\Big(\frac{n\pi x}{a}\Big)\dd x
  = \frac{1}{a}\int_0^a x\Big(1-\cos\frac{2n\pi x}{a}\Big)\dd x .
\]
Le premier terme vaut $\dfrac1a\cdot\dfrac{a^2}{2}=\dfrac a2$. Pour le second, une
intégration par parties donne
\[
  \int_0^a x\cos\frac{2n\pi x}{a}\,\dd x
  =\Big[\frac{a\,x}{2n\pi}\sin\frac{2n\pi x}{a}\Big]_0^a
   +\frac{a^2}{4n^2\pi^2}\Big[\cos\frac{2n\pi x}{a}\Big]_0^a = 0,
\]
puisque $\sin(2n\pi)=0$ et $\cos(2n\pi)=1$.

\begin{resultat}{Position moyenne}
\[
  \langle\hat x\rangle=\frac{a}{2}\qquad\text{pour tout }n\text{ et à tout instant.}
\]
\end{resultat}

Les phases $\ee^{\pm\ii E_nt/\hbar}$ se compensent dans $\Psi^*\Psi$. Le résultat se
voit d'ailleurs sans calcul, car $|\Psi_n|^2$ est symétrique par rapport à $a/2$.

\begin{remarque}[Moyenne et position la plus probable]
Pour $n=2$, la densité de probabilité s'annule en $x=a/2$ alors que
$\langle\hat x\rangle=a/2$ : la valeur moyenne n'est pas la position la plus
probable, et la particule n'est jamais trouvée là.
\end{remarque}

\paragraph{Incertitude sur l'énergie.} Comme $\Hop\Psi_n=E_n\Psi_n$, on a aussi
$\Hop^{2}\Psi_n=E_n^2\Psi_n$, donc $\langle\Hop\rangle=E_n$, $\langle\Hop^{2}\rangle=E_n^2$
et
\[
  \sigma_E=\sqrt{\langle\Hop^{2}\rangle-\langle\Hop\rangle^2}=\sqrt{E_n^2-E_n^2}=0 .
\]
Dans un état propre (stationnaire) l'énergie est parfaitement déterminée, alors que la
position ne l'est pas. Que se passe-t-il si la particule n'est pas dans un état propre ?
C'est l'objet de la section~\ref{sec:superposition}.

% ---------------------------------------------------------------------
\subsection{Bilan : opérateurs et puits infini}
% ---------------------------------------------------------------------

\begin{center}
\small\renewcommand{\arraystretch}{1.25}
\begin{tabular}{@{}ll@{}}
\toprule
\textbf{Notion} & \textbf{Formule} \\
\midrule
Bra-ket & $\braket{a}{b}=a_1^*b_1+a_2^*b_2$, \quad
          $\braket{\Psi}{\Psi}=\displaystyle\int\Psi^*\Psi\,\dd x=1$ \\
Opérateurs & $\hat p=-\ii\hbar\,\partial_x$, \quad
             $\hat E_c=\dfrac{\hat p^{\,2}}{2m}=-\dfrac{\hbar^2}{2m}\,\partial_x^2$, \quad
             $\Hop=\hat E_c+V(x)$ \\
Valeur moyenne, incertitude & $\langle\hat O\rangle=\displaystyle\int\Psi^*\,\hat O\,\Psi\,\dd x$,
          \quad $\sigma_O=\sqrt{\langle\hat O^{\,2}\rangle-\langle\hat O\rangle^2}$ \\
ESDT & $\ii\hbar\,\partial_t\Psi=\Hop\,\Psi$ \\
Séparation ($V=V(x)$) & $\Psi=\psi(x)\,\ee^{-\ii Et/\hbar}$, \quad ESIT : $\Hop\,\psi=E\,\psi$ \\
Puits infini & $\psi_n=\sqrt{2/a}\,\sin(n\pi x/a)$, \quad
               $E_n=\dfrac{n^2\pi^2\hbar^2}{2ma^2}$, \quad $E_1=\dfrac{\pi^2\hbar^2}{2ma^2}$ \\
Moyennes dans $\Psi_n$ & $\langle\hat x\rangle=a/2$, \quad $\langle\Hop\rangle=E_n$, \quad $\sigma_E=0$ \\
\bottomrule
\end{tabular}
\end{center}
.
%  Il PROLONGE chapitre1.tex : pas de \section ici, les sous-sections
%  1.6 à 1.9 font suite à 1.1 à 1.5 (la suite est dans chapitre3.tex).
%  Pas de préambule non plus : les
%  environnements (definition, resultat, remarque), les macros (\ii, \ee,
%  \dd, \pd, \pdd, \Hop, \ket, \bra, \braket) et les couleurs (insarouge,
%  insagris) sont définis dans main.tex.
%  Convention de signe : p = -i hbar d/dx, cohérente avec l'onde plane
%  e^{i(kx-wt)} de la section 1.3 (les notes de cours donnaient +i hbar).
% =====================================================================

% Macros de Dirac (déjà définies dans main.tex récent ; ce bloc évite l'erreur
% « Undefined control sequence \ket » si main.tex est une ancienne version).
\providecommand{\ket}[1]{|#1\rangle}
\providecommand{\bra}[1]{\langle #1|}
\providecommand{\braket}[2]{\langle #1|#2\rangle}

% ---------------------------------------------------------------------
\subsection{Notation de Dirac}\label{sec:dirac}
% ---------------------------------------------------------------------

Un état s'écrit comme un vecteur colonne, par exemple
$\begin{pmatrix} a_1\\ a_2\end{pmatrix}\in\mathbb{C}^2$ : c'est un élément d'un
espace de Hilbert de dimension~2, noté $\mathcal{H}^2$. Dirac le note simplement
$\ket{a}$ et évite d'écrire les composantes à chaque ligne.

\begin{definition}{Notation de Dirac : ket, bra et bra-ket}
\begin{center}
\begin{tabular}{@{}r@{\hspace{0.5em}}l@{\hspace{2em}}l@{}}
  {\Large\color{insarouge}$\ket{a}$}
    & $=\begin{pmatrix}a_1\\ a_2\end{pmatrix}\in\mathcal{H}^2$
    & \textbf{ket} : vecteur colonne \\[1.6ex]
  {\Large\color{insarouge}$\bra{a}$}
    & $=\big(\ket{a}\big)^\dagger=\begin{pmatrix}a_1^* & a_2^*\end{pmatrix}$
    & \textbf{bra} : vecteur ligne \\[1.6ex]
  {\Large\color{insarouge}$\braket{a}{b}$}
    & $=\begin{pmatrix}a_1^* & a_2^*\end{pmatrix}\begin{pmatrix}b_1\\ b_2\end{pmatrix}
       =a_1^*b_1+a_2^*b_2$
    & \textbf{bra-ket} : produit scalaire \\
\end{tabular}
\end{center}
Le bra est le transposé \emph{conjugué} du ket. Le nom vient de l'anglais
\emph{bracket} (crochet) : en accolant $\bra{a}$~(\emph{bra}) et $\ket{b}$~(\emph{ket})
on obtient $\braket{a}{b}$~(\emph{bracket}), un nombre complexe.

Propriétés : $\braket{b}{a}=\braket{a}{b}^*$ et
$\braket{a}{a}=|a_1|^2+|a_2|^2\in\mathbb{R}_+$ (norme au carré).
\end{definition}

Dans l'autre ordre, le produit d'un ket par un bra est une matrice (un opérateur) :
\[
  \ket{a}\bra{b}
  =\begin{pmatrix}a_1\\ a_2\end{pmatrix}\begin{pmatrix}b_1^* & b_2^*\end{pmatrix}
  =\begin{pmatrix}a_1b_1^* & a_1b_2^*\\ a_2b_1^* & a_2b_2^*\end{pmatrix}.
\]

\begin{remarque}[Pourquoi le conjugué ?]
Pour $\ket{a}=\dfrac{1}{\sqrt2}\begin{pmatrix}1\\ \ii\end{pmatrix}$, on a
$\bra{a}=\dfrac{1}{\sqrt2}\begin{pmatrix}1 & -\ii\end{pmatrix}$ et
$\braket{a}{a}=\tfrac12\,(1+1)=1$. Sans conjugaison on trouverait
$\tfrac12\,(1+\ii^2)=0$, ce qui n'est pas une norme.
\end{remarque}

Un état physique est \emph{normalisé} :
\[
  \braket{\Psi}{\Psi}=1 .
\]
Pour une particule dans l'espace, $\ket{\Psi}$ est une fonction de $x$ (dimension
infinie) et le produit scalaire devient une intégrale,
$\braket{\varphi}{\psi}=\int\varphi^*\psi\,\dd x$. On retrouve la condition de
normalisation de Born (section~\ref{sec:born}) :
\[
  \braket{\Psi}{\Psi} = \int_{-\infty}^{+\infty}\Psi^*\Psi\,\dd x = 1 .
\]
L'espace de Hilbert et le formalisme complet sont repris plus loin dans le cours.

% ---------------------------------------------------------------------
\subsection{Opérateurs quantiques}\label{sec:operateurs}
% ---------------------------------------------------------------------

À chaque grandeur physique (position, quantité de mouvement, énergie) on associe
un \emph{opérateur} $\hat O$ : une règle qui transforme une fonction d'onde en une
autre, $\Psi\mapsto\hat O\Psi$. Le point de départ est l'onde plane de la
section~\ref{sec:construction} : $\Psi=\ee^{\ii(kx-\omega t)}$ vérifie
$-\ii\hbar\,\partial_x\Psi=\hbar k\,\Psi=p\,\Psi$.

\begin{definition}{Opérateur quantité de mouvement}
\[
  \hat p = -\ii\hbar\,\frac{\partial}{\partial x},
  \qquad
  \hat p\,\Psi = -\ii\hbar\,\frac{\partial\Psi}{\partial x}.
\]
L'onde plane est fonction propre de $\hat p$, de valeur propre $p=\hbar k$.
\end{definition}

Les autres opérateurs s'obtiennent en remplaçant $p$ par $\hat p$ dans
l'expression classique. Avec $E_c=\tfrac12 mv^2=p^2/2m$ :
\[
  \hat E_c = \frac{\hat p^{\,2}}{2m}
  = \frac{1}{2m}\Big(-\ii\hbar\frac{\partial}{\partial x}\Big)
                \Big(-\ii\hbar\frac{\partial}{\partial x}\Big)
  = \frac{(-\ii\hbar)^2}{2m}\,\frac{\partial^2}{\partial x^2}
  = -\frac{\hbar^2}{2m}\,\frac{\partial^2}{\partial x^2},
  \qquad
  \hat v = \frac{\hat p}{m} = -\frac{\ii\hbar}{m}\,\frac{\partial}{\partial x}.
\]

\begin{center}
\renewcommand{\arraystretch}{1.6}
\begin{tabular}{@{}lll@{}}
\toprule
\textbf{Grandeur} & \textbf{Classique} & \textbf{Opérateur} \\
\midrule
Position & $x$ & $\hat x = x\times$ \\
Quantité de mouvement & $p$ & $\hat p = -\ii\hbar\,\partial_x$ \\
Vitesse & $v=\dfrac{p}{m}$ & $\hat v = -\dfrac{\ii\hbar}{m}\,\partial_x$ \\
Énergie cinétique & $E_c=\dfrac{p^2}{2m}$ & $\hat E_c = -\dfrac{\hbar^2}{2m}\,\partial_x^2$ \\
Énergie potentielle & $V(x)$ & $\hat V = V(x)\times$ \\
Énergie totale & $E=E_c+V$ & $\Hop = \hat E_c+\hat V = -\dfrac{\hbar^2}{2m}\,\partial_x^2+V(x)$ \\
\bottomrule
\end{tabular}
\end{center}

L'opérateur hamiltonien $\Hop$ représente l'énergie totale. Avec $E\to\ii\hbar\,\partial_t$,
l'équation de Schrödinger de la section~\ref{sec:construction} n'est que la relation
$E=E_c+V$ écrite avec des opérateurs : $\ii\hbar\,\partial_t\Psi=\Hop\Psi$.

\begin{remarque}[Convention de signe]
Le signe de $\hat p$ dépend du choix de l'onde plane $\ee^{\ii(kx-\omega t)}$. La
convention conjuguée $\ee^{-\ii(kx-\omega t)}$ inverse les deux signes
($\hat p=+\ii\hbar\,\partial_x$ et $E\to-\ii\hbar\,\partial_t$) sans changer les
résultats physiques. Comme $\hat p$ intervient au carré, $\hat E_c$ est le même dans
les deux conventions.
\end{remarque}

\begin{definition}{Valeur moyenne d'un opérateur}
Pour une fonction d'onde normalisée $\Psi$,
\[
  \langle\hat O\rangle
  = \int_{-\infty}^{+\infty}\Psi^*(x,t)\;\hat O\,\Psi(x,t)\,\dd x
  = \bra{\Psi}\hat O\ket{\Psi}.
\]
C'est la moyenne des résultats obtenus en mesurant la grandeur sur un grand nombre
de systèmes identiques, tous préparés dans l'état $\Psi$. L'opérateur agit sur
$\Psi$ seul, à droite de $\Psi^*$. Pour la position,
$\langle\hat x\rangle=\int x\,|\Psi|^2\,\dd x$ : c'est l'espérance de la densité
de probabilité de Born.

L'\emph{incertitude} (écart type) sur la grandeur associée à $\hat O$ mesure la
dispersion des résultats autour de cette moyenne :
\[
  \sigma_O=\sqrt{\langle\hat O^{\,2}\rangle-\langle\hat O\rangle^2}\,,
  \qquad
  \langle\hat O^{\,2}\rangle=\int\Psi^*\,\hat O^{\,2}\,\Psi\,\dd x .
\]
Si $\Psi$ est une fonction propre de $\hat O$ ($\hat O\Psi=\lambda\Psi$), alors
$\langle\hat O\rangle=\lambda$, $\langle\hat O^{\,2}\rangle=\lambda^2$ et
$\sigma_O=0$ : la grandeur a une valeur certaine.
\end{definition}

% ---------------------------------------------------------------------
\subsection{Solutions de l'équation de Schrödinger}\label{sec:solutions}
% ---------------------------------------------------------------------

\subsubsection{Du potentiel à la fonction d'onde}

En mécanique classique, l'énergie potentielle $V(x,t)$ fixe la force
$F=-\partial V/\partial x$, puis $\dd p/\dd t=F$ donne la trajectoire. En mécanique
quantique, c'est l'équation de Schrödinger qui joue ce rôle : connaissant $V$, on la
résout pour obtenir $\Psi(x,t)$.

\begin{resultat}{ESDT : équation de Schrödinger dépendante du temps}
\[
  \ii\hbar\,\pd{\Psi(x,t)}{t} = \Hop\,\Psi(x,t),
  \qquad
  \Hop = -\frac{\hbar^2}{2m}\,\pdd{}{x} + V(x,t).
\]
\end{resultat}

Si l'énergie potentielle ne dépend pas du temps, $V(x,t)=V(x)$, la force
$F=-\dd V/\dd x$ (dérivée ordinaire) est constante dans le temps et l'énergie se
conserve. L'équation se résout alors par séparation des variables.

\subsubsection{Séparation des variables}

On cherche des solutions produit d'une fonction de $x$ seul et d'une fonction de $t$
seul :
\[
  \Psi(x,t)=\psi(x)\,\varphi(t).
\]
Dans l'ESDT, les dérivées partielles deviennent des dérivées ordinaires :
\[
  \ii\hbar\,\psi(x)\,\frac{\dd\varphi}{\dd t}
  = \varphi(t)\left[-\frac{\hbar^2}{2m}\,\frac{\dd^2\psi}{\dd x^2}+V(x)\,\psi(x)\right].
\]
On divise par $\psi(x)\,\varphi(t)$ :
\[
  \underbrace{\ii\hbar\,\frac{1}{\varphi}\,\frac{\dd\varphi}{\dd t}}_{\text{ne dépend que de }t}
  \;=\;
  \underbrace{\frac{1}{\psi}\left[-\frac{\hbar^2}{2m}\,\frac{\dd^2\psi}{\dd x^2}+V(x)\,\psi\right]}_{\text{ne dépend que de }x}.
\]
Les deux membres dépendent de variables indépendantes et sont pourtant égaux pour
tout $(x,t)$ : ils sont égaux à une même constante, notée $E$. On obtient deux
équations.

\paragraph{Partie temporelle.} $\ii\hbar\,\dd\varphi/\dd t=E\,\varphi$ donne
\[
  \varphi(t)=\ee^{-\ii Et/\hbar}=\ee^{-\ii\omega t},
  \qquad \omega=\frac{E}{\hbar},
\]
la constante multiplicative étant absorbée dans $\psi$. On retrouve $E=\hbar\omega$ :
la constante $E$ est bien l'énergie de la particule.

\paragraph{Partie spatiale.} C'est l'ESIT ci-dessous.

\begin{resultat}{ESIT : équation de Schrödinger indépendante du temps}
\[
  -\frac{\hbar^2}{2m}\,\frac{\dd^2\psi(x)}{\dd x^2}+V(x)\,\psi(x)=E\,\psi(x)
  \qquad\Longleftrightarrow\qquad
  \Hop\,\psi(x)=E\,\psi(x).
\]
C'est un \emph{problème aux valeurs propres} : $E$ est une valeur propre de $\Hop$
(énergie permise) et $\psi$ la fonction propre associée. Chaque solution
$(E_n,\psi_n)$ donne un \emph{état stationnaire}
\[
  \Psi_n(x,t)=\psi_n(x)\,\ee^{-\ii E_nt/\hbar}.
\]
\end{resultat}

Dans un état stationnaire, $|\Psi_n(x,t)|^2=|\psi_n(x)|^2$ ne dépend pas du temps :
le facteur temporel est une phase globale, non observable. Les conditions physiques
($\psi$ continue et normalisable) ne sont satisfaites que pour certaines valeurs de
$E$ : c'est l'origine de la quantification de l'énergie, illustrée ci-dessous.

Comme l'ESDT est linéaire (hypothèse~(c) de la section~\ref{sec:construction}), toute
superposition $\Psi=\sum_n c_n\,\Psi_n$ est aussi solution, normalisée si
$\sum_n|c_n|^2=1$ et si les $\psi_n$ sont orthonormées. Une telle superposition
n'est en général pas stationnaire : $|\Psi|^2$ contient des termes croisés en
$\ee^{\ii(E_m-E_n)t/\hbar}$.

\subsubsection{Exemple : puits de potentiel infini}

Une particule de masse $m$ (un électron par exemple) est enfermée entre deux parois
impénétrables, en $x=0$ et $x=a$ :
\[
  V(x)=
  \begin{cases}
    0 & \text{si } 0<x<a,\\
    +\infty & \text{ailleurs.}
  \end{cases}
\]
La particule ne peut pas se trouver hors du puits : $\psi=0$ pour $x\le0$ et
$x\ge a$. La fonction $\psi$ est continue (mais pas sa dérivée, à cause du saut
infini de $V$), d'où les conditions aux limites
\[
  \psi(0)=\psi(a)=0 .
\]

\paragraph{Opérateur hamiltonien.} Dans le puits $V=0$, donc $\Hop$ se réduit à
l'opérateur énergie cinétique :
\[
  \Hop=\hat E_c=-\frac{\hbar^2}{2m}\,\frac{\dd^2}{\dd x^2}\qquad(0<x<a).
\]

\paragraph{Résolution de l'ESIT.} Avec $V=0$, l'ESIT s'écrit
$\psi''=-k^2\psi$ avec $k=\sqrt{2mE}/\hbar$. On prend $E>0$ : pour $E\le0$, la seule
solution vérifiant $\psi(0)=\psi(a)=0$ est $\psi\equiv0$. La solution générale est
\[
  \psi(x)=A\sin(kx)+B\cos(kx).
\]
\begin{itemize}[nosep]
  \item $\psi(0)=0$ donne $B=0$ ;
  \item $\psi(a)=A\sin(ka)=0$ avec $A\neq0$ (sinon $\psi\equiv0$ : pas de
        particule) donne $ka=n\pi$, soit
        $k_n=\dfrac{n\pi}{a}$ avec $n=1,2,3,\dots$
        ($n=0$ donne $\psi\equiv0$, et $n<0$ redonne les mêmes états au signe près).
\end{itemize}
Cela revient à $\lambda_n=2\pi/k_n=2a/n$, soit $a=n\,\lambda_n/2$ : $n$ demi-longueurs
d'onde tiennent dans le puits, comme pour une onde stationnaire.

\begin{resultat}{Énergie quantifiée}
\[
  E_n=\frac{\hbar^2k_n^2}{2m}=\frac{n^2\pi^2\hbar^2}{2ma^2}=\frac{n^2h^2}{8ma^2}=n^2E_1,
  \qquad n=1,2,3,\dots
\]
\textbf{Énergie fondamentale} (niveau $n=1$, le plus bas) :
$\displaystyle E_1=\frac{\pi^2\hbar^2}{2ma^2}=\frac{h^2}{8ma^2}>0$.
\end{resultat}

\begin{remarque}[Lecture du résultat]
\begin{itemize}[nosep]
  \item Les énergies sont \emph{discrètes} : $E_{n+1}-E_n=(2n+1)E_1$. Comme
        $E_n\propto1/a^2$, plus le puits est étroit, plus les énergies sont grandes.
  \item $E_1\neq0$ : une particule confinée ne peut pas être au repos (énergie de
        point zéro), conséquence du principe d'incertitude de Heisenberg.
  \item Ordre de grandeur : pour un électron avec $a=1$~nm,
        $E_1\approx0{,}38$~eV et $E_2=4E_1\approx1{,}5$~eV. Pour $a=1$~cm,
        $E_1\approx4\times10^{-15}$~eV : la quantification est indétectable.
\end{itemize}
\end{remarque}

\paragraph{Normalisation.} Le module de $\ee^{-\ii E_nt/\hbar}$ vaut 1, donc
$|\Psi_n|^2=A^2\sin^2(n\pi x/a)$ et la condition de Born impose
\[
  \int_0^a|\Psi_n|^2\,\dd x
  = A^2\int_0^a\frac{1-\cos\!\big(2n\pi x/a\big)}{2}\,\dd x
  = A^2\,\frac{a}{2}=1
  \quad\Longrightarrow\quad A=\sqrt{\frac{2}{a}},
\]
car l'intégrale du cosinus sur $n$ périodes est nulle.

\begin{resultat}{États stationnaires du puits infini}
\[
  \Psi_n(x,t)=\sqrt{\frac{2}{a}}\;\sin\!\Big(\frac{n\pi x}{a}\Big)\,\ee^{-\ii E_nt/\hbar}
  \quad(0<x<a),
  \qquad
  |\Psi_n|^2=\frac{2}{a}\,\sin^2\!\Big(\frac{n\pi x}{a}\Big).
\]
$\psi_n$ possède $n-1$ nœuds à l'intérieur du puits. Pour $n=1$ (état fondamental,
d'énergie $E_1$) :
\[
  \Psi_1(x,t)=\sqrt{\frac{2}{a}}\;\sin\!\Big(\frac{\pi x}{a}\Big)\,\ee^{-\ii E_1t/\hbar}.
\]
\end{resultat}

\begin{figure}[!ht]
\centering
\begin{tikzpicture}[>=Stealth, x=1cm, y=0.95cm]
  % ---------- (a) niveaux d'énergie et fonctions propres ----------
  \begin{scope}
    \node[above, font=\small] at (2,5.0) {(a) niveaux $E_n$ et fonctions propres $\psi_n$};
    \draw[->] (-2.5,0) -- (-2.5,4.8) node[above] {$E$};
    \draw[very thick, insarouge] (0,4.4) -- (0,0) -- (4,0) -- (4,4.4);
    \node[insarouge, right, font=\small] at (4.05,3.9) {$V=+\infty$};
    \node[insarouge, below, font=\small] at (2,0) {$V=0$};
    \node[below, font=\small] at (0,-0.35) {$0$};
    \node[below, font=\small] at (4,-0.35) {$a$};
    \node[left, font=\small] at (-0.1,0.4) {$E_1$};
    \node[left, font=\small] at (-0.1,1.6) {$E_2=4E_1$};
    \node[left, font=\small] at (-0.1,3.6) {$E_3=9E_1$};
    \foreach \n/\lev in {1/0.4, 2/1.6, 3/3.6} {
      \draw[dashed, insagris!70] (0,\lev) -- (4,\lev);
      \draw[thick, insagris, domain=0:1, samples=80]
        plot ({4*\x}, {\lev+0.5*sin(\n*180*\x)});
    }
  \end{scope}
  % ---------- (b) densités de probabilité ----------
  \begin{scope}[xshift=7cm]
    \node[above, font=\small] at (2,5.0) {(b) densités de probabilité $|\psi_n|^2$};
    \draw[->] (-0.3,0) -- (4.6,0) node[right] {$x$};
    \draw[thick, insarouge] (0,0) -- (0,4.4) (4,0) -- (4,4.4);
    \node[below, font=\small] at (0,0) {$0$};
    \node[below, font=\small] at (4,0) {$a$};
    \node[below, font=\small] at (2,0) {$a/2$};
    \foreach \n/\lev in {1/0.4, 2/1.6, 3/3.6} {
      \draw[dashed, insagris!70] (0,\lev) -- (4,\lev);
      \fill[insarouge!20] (0,\lev) -- plot[domain=0:1, samples=80]
        ({4*\x}, {\lev+0.6*(sin(\n*180*\x))^2}) -- (4,\lev) -- cycle;
      \draw[thick, insarouge, domain=0:1, samples=80]
        plot ({4*\x}, {\lev+0.6*(sin(\n*180*\x))^2});
    }
    \draw[dotted, thick, insagris] (2,0) -- (2,4.4);
    \node[above, font=\small] at (2,4.4) {$\langle\hat x\rangle=a/2$};
  \end{scope}
\end{tikzpicture}
\caption{Puits de potentiel infini de largeur $a$. (a) Niveaux d'énergie
$E_n=n^2E_1$ et fonctions propres $\psi_n$. (b) Densités de probabilité
$|\psi_n|^2$, symétriques par rapport à $a/2$.}
\end{figure}

\subsubsection{Position moyenne dans le puits}

D'après la définition de la valeur moyenne, avec $\hat x=x\times$ :
\[
  \langle\hat x\rangle
  = \int_0^a\Psi_n^*(x,t)\;x\;\Psi_n(x,t)\,\dd x
  = \frac{2}{a}\int_0^a x\,\sin^2\!\Big(\frac{n\pi x}{a}\Big)\dd x
  = \frac{1}{a}\int_0^a x\Big(1-\cos\frac{2n\pi x}{a}\Big)\dd x .
\]
Le premier terme vaut $\dfrac1a\cdot\dfrac{a^2}{2}=\dfrac a2$. Pour le second, une
intégration par parties donne
\[
  \int_0^a x\cos\frac{2n\pi x}{a}\,\dd x
  =\Big[\frac{a\,x}{2n\pi}\sin\frac{2n\pi x}{a}\Big]_0^a
   +\frac{a^2}{4n^2\pi^2}\Big[\cos\frac{2n\pi x}{a}\Big]_0^a = 0,
\]
puisque $\sin(2n\pi)=0$ et $\cos(2n\pi)=1$.

\begin{resultat}{Position moyenne}
\[
  \langle\hat x\rangle=\frac{a}{2}\qquad\text{pour tout }n\text{ et à tout instant.}
\]
\end{resultat}

Les phases $\ee^{\pm\ii E_nt/\hbar}$ se compensent dans $\Psi^*\Psi$. Le résultat se
voit d'ailleurs sans calcul, car $|\Psi_n|^2$ est symétrique par rapport à $a/2$.

\begin{remarque}[Moyenne et position la plus probable]
Pour $n=2$, la densité de probabilité s'annule en $x=a/2$ alors que
$\langle\hat x\rangle=a/2$ : la valeur moyenne n'est pas la position la plus
probable, et la particule n'est jamais trouvée là.
\end{remarque}

\paragraph{Incertitude sur l'énergie.} Comme $\Hop\Psi_n=E_n\Psi_n$, on a aussi
$\Hop^{2}\Psi_n=E_n^2\Psi_n$, donc $\langle\Hop\rangle=E_n$, $\langle\Hop^{2}\rangle=E_n^2$
et
\[
  \sigma_E=\sqrt{\langle\Hop^{2}\rangle-\langle\Hop\rangle^2}=\sqrt{E_n^2-E_n^2}=0 .
\]
Dans un état propre (stationnaire) l'énergie est parfaitement déterminée, alors que la
position ne l'est pas. Que se passe-t-il si la particule n'est pas dans un état propre ?
C'est l'objet de la section~\ref{sec:superposition}.

% ---------------------------------------------------------------------
\subsection{Bilan : opérateurs et puits infini}
% ---------------------------------------------------------------------

\begin{center}
\small\renewcommand{\arraystretch}{1.25}
\begin{tabular}{@{}ll@{}}
\toprule
\textbf{Notion} & \textbf{Formule} \\
\midrule
Bra-ket & $\braket{a}{b}=a_1^*b_1+a_2^*b_2$, \quad
          $\braket{\Psi}{\Psi}=\displaystyle\int\Psi^*\Psi\,\dd x=1$ \\
Opérateurs & $\hat p=-\ii\hbar\,\partial_x$, \quad
             $\hat E_c=\dfrac{\hat p^{\,2}}{2m}=-\dfrac{\hbar^2}{2m}\,\partial_x^2$, \quad
             $\Hop=\hat E_c+V(x)$ \\
Valeur moyenne, incertitude & $\langle\hat O\rangle=\displaystyle\int\Psi^*\,\hat O\,\Psi\,\dd x$,
          \quad $\sigma_O=\sqrt{\langle\hat O^{\,2}\rangle-\langle\hat O\rangle^2}$ \\
ESDT & $\ii\hbar\,\partial_t\Psi=\Hop\,\Psi$ \\
Séparation ($V=V(x)$) & $\Psi=\psi(x)\,\ee^{-\ii Et/\hbar}$, \quad ESIT : $\Hop\,\psi=E\,\psi$ \\
Puits infini & $\psi_n=\sqrt{2/a}\,\sin(n\pi x/a)$, \quad
               $E_n=\dfrac{n^2\pi^2\hbar^2}{2ma^2}$, \quad $E_1=\dfrac{\pi^2\hbar^2}{2ma^2}$ \\
Moyennes dans $\Psi_n$ & $\langle\hat x\rangle=a/2$, \quad $\langle\Hop\rangle=E_n$, \quad $\sigma_E=0$ \\
\bottomrule
\end{tabular}
\end{center}
 % 2. Notation de Dirac et espace de Hilbert
% =====================================================================
%  Chapitre 3 : suite de la section 1 (mécanique quantique)
%  Fichier importé par main.tex avec % =====================================================================
%  Chapitre 3 : suite de la section 1 (mécanique quantique)
%  Fichier importé par main.tex avec % =====================================================================
%  Chapitre 3 : suite de la section 1 (mécanique quantique)
%  Fichier importé par main.tex avec \include{chapitre3}.
%  Il PROLONGE chapitre2.tex : pas de \section ici, les sous-sections
%  1.10 à 1.13 font suite à 1.1 à 1.9 (la suite est le chapitre4.tex, section 2).
%  Pas de préambule non plus : les
%  environnements (definition, resultat, remarque), les macros (\ii, \ee,
%  \dd, \Hop, \ket, \bra, \braket) et les couleurs (insarouge, insagris)
%  sont définis dans main.tex.
% =====================================================================

% Macros de Dirac (déjà définies dans main.tex récent ; ce bloc évite l'erreur
% « Undefined control sequence \ket » si main.tex est une ancienne version).
\providecommand{\ket}[1]{|#1\rangle}
\providecommand{\bra}[1]{\langle #1|}
\providecommand{\braket}[2]{\langle #1|#2\rangle}

% ---------------------------------------------------------------------
\subsection{Superposition d'états stationnaires}\label{sec:superposition}
% ---------------------------------------------------------------------

\subsubsection{Un état qui n'est pas un état propre}

Dans un état propre $\Psi_n$ du puits infini, l'énergie est certaine
($\sigma_E=0$, section~\ref{sec:solutions}). Rien n'oblige pourtant la particule à
être dans un état propre. Par linéarité de l'équation de Schrödinger
(3\textsuperscript{e} hypothèse de la section~\ref{sec:construction}), toute
combinaison de solutions est encore solution.

\begin{definition}{Superposition quantique}
\[
  \Psi(x,t)=\sum_{i=1}^{N}c_i\,\Psi_i(x,t),
  \qquad c_i\in\mathbb{C},
\]
où les $\Psi_i$ sont des états stationnaires d'énergies $E_i$ ($i$ est le nombre
quantique de l'état $\Psi_i$). L'état n'a plus une énergie unique : il est
\emph{probabiliste}, l'énergie mesurée pouvant valoir chacune des $E_i$.
\end{definition}

\subsubsection{Exemple : deux états du puits}

À $t=0$, on prépare la particule dans l'état
\[
  \Psi(x,0)=A\left[\sqrt{\frac{2}{a}}\,\sin\!\Big(\frac{\pi x}{a}\Big)
                  +\sqrt{\frac{2}{a}}\,\sin\!\Big(\frac{2\pi x}{a}\Big)\right]
  =A\,\big[\psi_1(x)+\psi_2(x)\big].
\]
La normalisation $\int_0^a\Psi^*\Psi\,\dd x=1$ donne, en développant,
\[
  A^2\left[\int_0^a\psi_1^2\,\dd x+2\int_0^a\psi_1\psi_2\,\dd x+\int_0^a\psi_2^2\,\dd x\right]
  =A^2\,(1+0+1)=2A^2=1
  \quad\Longrightarrow\quad A=\frac{1}{\sqrt2}.
\]
Les $\psi_n$ étant normalisées, $\int\psi_n^2\,\dd x=1$. Le terme croisé s'annule :
\[
  \int_0^a\psi_1\psi_2\,\dd x
  =\frac{1}{a}\int_0^a\Big[\cos\frac{\pi x}{a}-\cos\frac{3\pi x}{a}\Big]\dd x=0 .
\]
Chaque terme évolue avec sa propre phase (chacun est une solution stationnaire), donc
\[
  \Psi(x,t)=\frac{1}{\sqrt2}\Big[\psi_1(x)\,\ee^{-\ii E_1t/\hbar}
                                 +\psi_2(x)\,\ee^{-\ii E_2t/\hbar}\Big],
\]
\[
  |\Psi(x,t)|^2=\frac12\big[\psi_1^2+\psi_2^2\big]
                +\psi_1\psi_2\,\cos\!\Big(\frac{E_2-E_1}{\hbar}\,t\Big).
\]
Le terme croisé oscille à la pulsation $(E_2-E_1)/\hbar$ : contrairement à un état
stationnaire, la densité de probabilité dépend du temps.

\begin{figure}[!ht]
\centering
\begin{tikzpicture}[>=Stealth, x=1cm, y=1cm]
  \draw[->] (-0.3,0) -- (6.0,0) node[right] {$x$};
  \draw[->] (0,0) -- (0,4.2) node[above] {$|\Psi|^2$};
  \draw[thick, insagris, dashed, domain=0:1, samples=100]
    plot ({5*\x}, {1.0*(sin(180*\x)-sin(360*\x))^2});
  \draw[thick, insarouge, domain=0:1, samples=100]
    plot ({5*\x}, {1.0*(sin(180*\x)+sin(360*\x))^2});
  \node[below] at (0,0) {$0$};
  \node[below] at (5,0) {$a$};
  \draw[insarouge, thick] (0.6,3.8) -- (1.2,3.8) node[right, black, font=\small] {$t=0$};
  \draw[insagris, thick, dashed] (3.2,3.8) -- (3.8,3.8)
    node[right, black, font=\small] {$t=T/2$};
\end{tikzpicture}
\caption{Densité de probabilité de $(\psi_1+\psi_2)/\sqrt2$ à $t=0$ et à une
demi-période $T/2$ (avec $T=2\pi\hbar/(E_2-E_1)$) : le maximum passe de la gauche à
la droite du puits. Unité verticale : $1/a$.}
\end{figure}

\subsubsection{Cas général : normalisation}

Pour $\Psi=\sum_i c_i\Psi_i$, la condition $\int\Psi^*\Psi\,\dd x=1$ s'écrit
\[
  \int_0^a\Big[\sum_i c_i^*\Psi_i^*\Big]\Big[\sum_j c_j\Psi_j\Big]\dd x
  =|c_1|^2\!\int\!\Psi_1^*\Psi_1\,\dd x+c_1^*c_2\!\int\!\Psi_1^*\Psi_2\,\dd x+\cdots=1 .
\]
Les $\Psi_i$ sont orthonormées, comme des vecteurs unitaires : le produit scalaire d'un
vecteur avec lui-même vaut 1, celui de deux vecteurs orthogonaux vaut 0.

\begin{resultat}{Orthonormalité et normalisation}
\[
  \int_0^a\Psi_i^*\Psi_j\,\dd x=\delta_{ij}=
  \begin{cases}1 & \text{si } i=j,\\ 0 & \text{si } i\neq j,\end{cases}
  \qquad\Longrightarrow\qquad
  |c_1|^2+|c_2|^2+\cdots+|c_N|^2=\sum_{i=1}^{N}|c_i|^2=1
\]
\end{resultat}

L'orthogonalité tient à tout instant, car les phases $\ee^{\ii(E_i-E_j)t/\hbar}$
multiplient une intégrale nulle. En multipliant $\Psi=\sum_j c_j\Psi_j$ par $\Psi_i^*$
et en intégrant, on obtient aussi $c_i=\int\Psi_i^*\,\Psi\,\dd x$.

\subsubsection{Probabilités et énergie moyenne}

L'énergie moyenne est la valeur moyenne de $\Hop$, avec $\Hop\Psi_j=E_j\Psi_j$ :
\[
  \langle E\rangle=\int_0^a\Psi^*\,\Hop\,\Psi\,\dd x
  =\sum_{i,j}c_i^*c_jE_j\!\int\!\Psi_i^*\Psi_j\,\dd x
  =\sum_{i=1}^{N}E_i\,|c_i|^2,
\]
et de même, avec $\Hop^{2}\Psi_j=E_j^2\Psi_j$,
\[
  \langle E^2\rangle=\int_0^a\Psi^*\,\Hop^{2}\,\Psi\,\dd x=\sum_{i=1}^{N}E_i^2\,|c_i|^2 ,
  \qquad
  \sigma_E=\sqrt{\langle E^2\rangle-\langle E\rangle^2}.
\]

\begin{resultat}{Probabilité de mesurer l'énergie $E_i$}
\[
  P(E_i)=|c_i|^2 ,\qquad
  \langle E\rangle=\sum_i E_i\,P(E_i).
\]
$c_i$ est l'\emph{amplitude de probabilité} de l'état $\Psi_i$ et $|c_i|^2$ sa
probabilité. Les $|c_i|^2$ ne dépendent pas du temps.
\end{resultat}

La moyenne pondérée garde un sens même si les $E_i$ sont discrets. Elle n'est
d'ailleurs pas forcément une valeur mesurable : une mesure ne donne jamais que l'un des
$E_i$.

\paragraph{Application à l'exemple.} Avec $c_1=c_2=1/\sqrt2$ et $E_2=4E_1$ :
$P(E_1)=P(E_2)=\tfrac12$,
\[
  \langle E\rangle=\frac{E_1+E_2}{2}=\frac{5}{2}\,E_1,
  \qquad
  \langle E^2\rangle=\frac{E_1^2+E_2^2}{2}=\frac{17}{2}\,E_1^2,
  \qquad
  \sigma_E=\sqrt{\frac{17}{2}-\frac{25}{4}}\;E_1=\frac{3}{2}\,E_1\neq0 .
\]
Ici l'énergie est incertaine, et la valeur moyenne $\tfrac52E_1$ n'est jamais obtenue :
on mesure $E_1$ ou $E_2$, chacune avec la probabilité $1/2$.

\begin{remarque}[Effondrement]
Si la mesure donne $E_i$, l'état devient $\Psi_i$ : une seconde mesure faite
aussitôt redonne $E_i$ avec certitude. La probabilité $|c_i|^2$ ne se lit donc que sur
un grand nombre de systèmes identiques préparés dans le même état.
\end{remarque}

% ---------------------------------------------------------------------
\subsection{Retour à la notation de Dirac : le qubit}\label{sec:qubit}
% ---------------------------------------------------------------------

On ne garde que deux états du puits, notés comme un bit quantique (\emph{qubit}) :
$\ket{0}$ pour $\Psi_1$ et $\ket{1}$ pour $\Psi_2$. L'état de l'exemple s'écrit
\[
  \ket{\Psi}=\frac{1}{\sqrt2}\,\ket{0}+\frac{1}{\sqrt2}\,\ket{1}.
\]

\begin{center}
\renewcommand{\arraystretch}{1.5}
\begin{tabular}{@{}ll@{}}
\toprule
\textbf{Fonctions d'onde} & \textbf{Notation de Dirac} \\
\midrule
$\Psi=c_1\Psi_1+c_2\Psi_2$ & $\ket{\Psi}=c_1\ket{0}+c_2\ket{1}$ \\
$\displaystyle\int\Psi_i^*\Psi_j\,\dd x=\delta_{ij}$ & $\braket{i}{j}=\delta_{ij}$ \\
$\displaystyle\int\Psi^*\Psi\,\dd x=\sum_i|c_i|^2=1$ & $\braket{\Psi}{\Psi}=1$ \\
$P(E_i)=|c_i|^2$ & $P(i)=|c_i|^2=|\braket{i}{\Psi}|^2$ \\
\bottomrule
\end{tabular}
\end{center}

\paragraph{Probabilité.} La probabilité que $\ket{\Psi}$ soit mesuré dans l'état
$\ket{0}$ est
\[
  P(\ket{\Psi}\to\ket{0})=P(0)=\left|\frac{1}{\sqrt2}\right|^2=\frac12 .
\]

\paragraph{Normalisation.}
\[
  \braket{\Psi}{\Psi}
  =\Big(\tfrac{1}{\sqrt2}\bra{0}+\tfrac{1}{\sqrt2}\bra{1}\Big)
   \Big(\tfrac{1}{\sqrt2}\ket{0}+\tfrac{1}{\sqrt2}\ket{1}\Big)
  =\tfrac12\braket{0}{0}+\tfrac12\braket{0}{1}+\tfrac12\braket{1}{0}+\tfrac12\braket{1}{1}
  =\tfrac12\,(1+0+0+1)=1 .
\]

\begin{definition}{Base de calcul (base $Z$)}
$\{\ket{0},\ket{1}\}$ est une base \emph{orthonormée} de $\mathcal{H}^2=\mathbb{C}^2$ :
\[
  \ket{0}=\begin{pmatrix}1\\ 0\end{pmatrix},\qquad
  \ket{1}=\begin{pmatrix}0\\ 1\end{pmatrix},\qquad
  \braket{0}{0}=\begin{pmatrix}1 & 0\end{pmatrix}\begin{pmatrix}1\\ 0\end{pmatrix}=1,\qquad
  \braket{0}{1}=\begin{pmatrix}1 & 0\end{pmatrix}\begin{pmatrix}0\\ 1\end{pmatrix}=0 .
\]
Tout état s'écrit $\ket{\Psi}=\alpha\ket{0}+\beta\ket{1}=\begin{pmatrix}\alpha\\ \beta\end{pmatrix}$
avec $|\alpha|^2+|\beta|^2=1$ ; alors $P(0)=|\alpha|^2$ et $P(1)=|\beta|^2$.
\end{definition}

\paragraph{Exemple.} Soit $\ket{\Psi}=\dfrac{1}{\sqrt5}\,\ket{0}+\dfrac{2}{\sqrt5}\,\ket{1}$
dans $\mathcal{H}^2$.
\begin{enumerate}[nosep]
  \item \emph{$\ket{\Psi}$ est-il normalisé ?}
        $\braket{\Psi}{\Psi}=\dfrac15+\dfrac45=1$ : oui.
  \item $P(0)=\left|\dfrac{1}{\sqrt5}\right|^2=\dfrac15$.
  \item $P(1)=\left|\dfrac{2}{\sqrt5}\right|^2=\dfrac45$.
\end{enumerate}

\begin{remarque}[État non normalisé]
Si $\braket{\Psi}{\Psi}\neq1$, on divise l'état par sa norme $\sqrt{\braket{\Psi}{\Psi}}$.
Par exemple $\ket{0}+2\ket{1}$ a pour norme au carré $1+4=5$ : l'état normalisé est
$\frac{1}{\sqrt5}\ket{0}+\frac{2}{\sqrt5}\ket{1}$, celui de l'exemple.
\end{remarque}

% ---------------------------------------------------------------------
\subsection{Projecteurs et mesure}\label{sec:projecteurs}
% ---------------------------------------------------------------------

Le produit d'un ket par un bra est un opérateur (section~\ref{sec:dirac}). Avec les
états de la base de calcul, on obtient deux \emph{projecteurs}.

\begin{definition}{Projecteurs sur $\ket{0}$ et sur $\ket{1}$}
\[
  \hat P_0=\ket{0}\bra{0}=\begin{pmatrix}1\\ 0\end{pmatrix}\begin{pmatrix}1 & 0\end{pmatrix}
  =\begin{pmatrix}1 & 0\\ 0 & 0\end{pmatrix},
  \qquad
  \hat P_1=\ket{1}\bra{1}=\begin{pmatrix}0\\ 1\end{pmatrix}\begin{pmatrix}0 & 1\end{pmatrix}
  =\begin{pmatrix}0 & 0\\ 0 & 1\end{pmatrix}.
\]
Ils vérifient
\[
  \hat P_i^{\,2}=\hat P_i,
  \qquad
  \hat P_0+\hat P_1=\begin{pmatrix}1 & 0\\ 0 & 1\end{pmatrix}=\mathrm{Id}
  \quad(\text{complétude}).
\]
\end{definition}

En effet $\hat P_0^{\,2}=\ket{0}\braket{0}{0}\bra{0}=\ket{0}\bra{0}$ car
$\braket{0}{0}=1$. Appliqué à $\ket{\Psi}=\alpha\ket{0}+\beta\ket{1}$, $\hat P_0$
ne garde que la composante sur $\ket{0}$ : $\hat P_0\ket{\Psi}=\alpha\ket{0}$.

\begin{resultat}{Probabilité de mesure}
\[
  P(\ket{\Psi}\to\ket{i})=\bra{\Psi}\hat P_i\ket{\Psi}
  =\braket{\Psi}{i}\braket{i}{\Psi}=|\braket{i}{\Psi}|^2 .
\]
La complétude donne bien $\sum_i P(i)=\bra{\Psi}\big(\sum_i\hat P_i\big)\ket{\Psi}=\braket{\Psi}{\Psi}=1$.
\end{resultat}

\paragraph{Exemple.} Pour $\ket{\Psi}=\dfrac{1}{\sqrt5}\ket{0}+\dfrac{2}{\sqrt5}\ket{1}
=\begin{pmatrix}1/\sqrt5\\ 2/\sqrt5\end{pmatrix}$ :
\[
  \hat P_0\ket{\Psi}=\begin{pmatrix}1 & 0\\ 0 & 0\end{pmatrix}
  \begin{pmatrix}1/\sqrt5\\ 2/\sqrt5\end{pmatrix}
  =\begin{pmatrix}1/\sqrt5\\ 0\end{pmatrix},
  \qquad
  P(\ket{\Psi}\to\ket{0})=\bra{\Psi}\hat P_0\ket{\Psi}
  =\begin{pmatrix}\frac{1}{\sqrt5} & \frac{2}{\sqrt5}\end{pmatrix}
   \begin{pmatrix}\frac{1}{\sqrt5}\\ 0\end{pmatrix}=\frac15 ,
\]
\[
  P(\ket{\Psi}\to\ket{1})=\bra{\Psi}\hat P_1\ket{\Psi}
  =\begin{pmatrix}\frac{1}{\sqrt5} & \frac{2}{\sqrt5}\end{pmatrix}
   \begin{pmatrix}0 & 0\\ 0 & 1\end{pmatrix}
   \begin{pmatrix}\frac{1}{\sqrt5}\\ \frac{2}{\sqrt5}\end{pmatrix}=\frac45 .
\]

\paragraph{Effet d'une phase relative.} Soit $\ket{\Psi'}=\dfrac{1}{\sqrt5}\ket{0}
+\dfrac{2\ii}{\sqrt5}\ket{1}$. Son bra est le \emph{conjugué} :
$\bra{\Psi'}=\begin{pmatrix}\frac{1}{\sqrt5} & -\frac{2\ii}{\sqrt5}\end{pmatrix}$, donc
\[
  \braket{\Psi'}{\Psi'}=\frac15+\Big(-\frac{2\ii}{\sqrt5}\Big)\Big(\frac{2\ii}{\sqrt5}\Big)
  =\frac15+\frac45=1 ,
\]
et $P(0)=1/5$, $P(1)=4/5$ comme pour $\ket{\Psi}$. Les deux états sont pourtant
différents : le facteur $\ii$ est une phase \emph{relative} entre les deux composantes,
que la mesure dans la base $Z$ ne voit pas (contrairement à une phase globale, jamais
observable).

\begin{remarque}[Rappels]
$\hat P^{\,2}=\hat P$ définit un projecteur ; $\braket{a}{b}=\braket{b}{a}^*$ ;
$\hat P_i$ est hermitien : $\hat P_i^{\,\dagger}=\hat P_i$.
\end{remarque}

% ---------------------------------------------------------------------
\subsection{Récapitulatif du chapitre}
% ---------------------------------------------------------------------

\begin{center}
\renewcommand{\arraystretch}{1.7}
\begin{tabular}{@{}ll@{}}
\toprule
\textbf{Notion} & \textbf{Formule} \\
\midrule
Superposition & $\Psi=\displaystyle\sum_i c_i\Psi_i$, \quad $c_i\in\mathbb{C}$ \\
Orthonormalité & $\displaystyle\int\Psi_i^*\Psi_j\,\dd x=\delta_{ij}$ \\
Normalisation & $\displaystyle\sum_i|c_i|^2=1$ \\
Probabilité d'une énergie & $P(E_i)=|c_i|^2$ \\
Énergie moyenne & $\langle E\rangle=\displaystyle\sum_iE_i|c_i|^2$ \\
Incertitude & $\sigma_E=\sqrt{\langle E^2\rangle-\langle E\rangle^2}$, \quad
              $\langle E^2\rangle=\displaystyle\sum_iE_i^2|c_i|^2$ \\
Qubit & $\ket{\Psi}=\alpha\ket{0}+\beta\ket{1}$, \quad $|\alpha|^2+|\beta|^2=1$ \\
Mesure dans la base $Z$ & $P(0)=|\alpha|^2$, \quad $P(1)=|\beta|^2$ \\
Projecteurs & $\hat P_i=\ket{i}\bra{i}$, \quad $\hat P_i^{\,2}=\hat P_i$, \quad $\hat P_0+\hat P_1=\mathrm{Id}$ \\
Probabilité de mesure & $P(i)=\bra{\Psi}\hat P_i\ket{\Psi}=|\braket{i}{\Psi}|^2$ \\
\bottomrule
\end{tabular}
\end{center}
.
%  Il PROLONGE chapitre2.tex : pas de \section ici, les sous-sections
%  1.10 à 1.13 font suite à 1.1 à 1.9 (la suite est le chapitre4.tex, section 2).
%  Pas de préambule non plus : les
%  environnements (definition, resultat, remarque), les macros (\ii, \ee,
%  \dd, \Hop, \ket, \bra, \braket) et les couleurs (insarouge, insagris)
%  sont définis dans main.tex.
% =====================================================================

% Macros de Dirac (déjà définies dans main.tex récent ; ce bloc évite l'erreur
% « Undefined control sequence \ket » si main.tex est une ancienne version).
\providecommand{\ket}[1]{|#1\rangle}
\providecommand{\bra}[1]{\langle #1|}
\providecommand{\braket}[2]{\langle #1|#2\rangle}

% ---------------------------------------------------------------------
\subsection{Superposition d'états stationnaires}\label{sec:superposition}
% ---------------------------------------------------------------------

\subsubsection{Un état qui n'est pas un état propre}

Dans un état propre $\Psi_n$ du puits infini, l'énergie est certaine
($\sigma_E=0$, section~\ref{sec:solutions}). Rien n'oblige pourtant la particule à
être dans un état propre. Par linéarité de l'équation de Schrödinger
(3\textsuperscript{e} hypothèse de la section~\ref{sec:construction}), toute
combinaison de solutions est encore solution.

\begin{definition}{Superposition quantique}
\[
  \Psi(x,t)=\sum_{i=1}^{N}c_i\,\Psi_i(x,t),
  \qquad c_i\in\mathbb{C},
\]
où les $\Psi_i$ sont des états stationnaires d'énergies $E_i$ ($i$ est le nombre
quantique de l'état $\Psi_i$). L'état n'a plus une énergie unique : il est
\emph{probabiliste}, l'énergie mesurée pouvant valoir chacune des $E_i$.
\end{definition}

\subsubsection{Exemple : deux états du puits}

À $t=0$, on prépare la particule dans l'état
\[
  \Psi(x,0)=A\left[\sqrt{\frac{2}{a}}\,\sin\!\Big(\frac{\pi x}{a}\Big)
                  +\sqrt{\frac{2}{a}}\,\sin\!\Big(\frac{2\pi x}{a}\Big)\right]
  =A\,\big[\psi_1(x)+\psi_2(x)\big].
\]
La normalisation $\int_0^a\Psi^*\Psi\,\dd x=1$ donne, en développant,
\[
  A^2\left[\int_0^a\psi_1^2\,\dd x+2\int_0^a\psi_1\psi_2\,\dd x+\int_0^a\psi_2^2\,\dd x\right]
  =A^2\,(1+0+1)=2A^2=1
  \quad\Longrightarrow\quad A=\frac{1}{\sqrt2}.
\]
Les $\psi_n$ étant normalisées, $\int\psi_n^2\,\dd x=1$. Le terme croisé s'annule :
\[
  \int_0^a\psi_1\psi_2\,\dd x
  =\frac{1}{a}\int_0^a\Big[\cos\frac{\pi x}{a}-\cos\frac{3\pi x}{a}\Big]\dd x=0 .
\]
Chaque terme évolue avec sa propre phase (chacun est une solution stationnaire), donc
\[
  \Psi(x,t)=\frac{1}{\sqrt2}\Big[\psi_1(x)\,\ee^{-\ii E_1t/\hbar}
                                 +\psi_2(x)\,\ee^{-\ii E_2t/\hbar}\Big],
\]
\[
  |\Psi(x,t)|^2=\frac12\big[\psi_1^2+\psi_2^2\big]
                +\psi_1\psi_2\,\cos\!\Big(\frac{E_2-E_1}{\hbar}\,t\Big).
\]
Le terme croisé oscille à la pulsation $(E_2-E_1)/\hbar$ : contrairement à un état
stationnaire, la densité de probabilité dépend du temps.

\begin{figure}[!ht]
\centering
\begin{tikzpicture}[>=Stealth, x=1cm, y=1cm]
  \draw[->] (-0.3,0) -- (6.0,0) node[right] {$x$};
  \draw[->] (0,0) -- (0,4.2) node[above] {$|\Psi|^2$};
  \draw[thick, insagris, dashed, domain=0:1, samples=100]
    plot ({5*\x}, {1.0*(sin(180*\x)-sin(360*\x))^2});
  \draw[thick, insarouge, domain=0:1, samples=100]
    plot ({5*\x}, {1.0*(sin(180*\x)+sin(360*\x))^2});
  \node[below] at (0,0) {$0$};
  \node[below] at (5,0) {$a$};
  \draw[insarouge, thick] (0.6,3.8) -- (1.2,3.8) node[right, black, font=\small] {$t=0$};
  \draw[insagris, thick, dashed] (3.2,3.8) -- (3.8,3.8)
    node[right, black, font=\small] {$t=T/2$};
\end{tikzpicture}
\caption{Densité de probabilité de $(\psi_1+\psi_2)/\sqrt2$ à $t=0$ et à une
demi-période $T/2$ (avec $T=2\pi\hbar/(E_2-E_1)$) : le maximum passe de la gauche à
la droite du puits. Unité verticale : $1/a$.}
\end{figure}

\subsubsection{Cas général : normalisation}

Pour $\Psi=\sum_i c_i\Psi_i$, la condition $\int\Psi^*\Psi\,\dd x=1$ s'écrit
\[
  \int_0^a\Big[\sum_i c_i^*\Psi_i^*\Big]\Big[\sum_j c_j\Psi_j\Big]\dd x
  =|c_1|^2\!\int\!\Psi_1^*\Psi_1\,\dd x+c_1^*c_2\!\int\!\Psi_1^*\Psi_2\,\dd x+\cdots=1 .
\]
Les $\Psi_i$ sont orthonormées, comme des vecteurs unitaires : le produit scalaire d'un
vecteur avec lui-même vaut 1, celui de deux vecteurs orthogonaux vaut 0.

\begin{resultat}{Orthonormalité et normalisation}
\[
  \int_0^a\Psi_i^*\Psi_j\,\dd x=\delta_{ij}=
  \begin{cases}1 & \text{si } i=j,\\ 0 & \text{si } i\neq j,\end{cases}
  \qquad\Longrightarrow\qquad
  |c_1|^2+|c_2|^2+\cdots+|c_N|^2=\sum_{i=1}^{N}|c_i|^2=1
\]
\end{resultat}

L'orthogonalité tient à tout instant, car les phases $\ee^{\ii(E_i-E_j)t/\hbar}$
multiplient une intégrale nulle. En multipliant $\Psi=\sum_j c_j\Psi_j$ par $\Psi_i^*$
et en intégrant, on obtient aussi $c_i=\int\Psi_i^*\,\Psi\,\dd x$.

\subsubsection{Probabilités et énergie moyenne}

L'énergie moyenne est la valeur moyenne de $\Hop$, avec $\Hop\Psi_j=E_j\Psi_j$ :
\[
  \langle E\rangle=\int_0^a\Psi^*\,\Hop\,\Psi\,\dd x
  =\sum_{i,j}c_i^*c_jE_j\!\int\!\Psi_i^*\Psi_j\,\dd x
  =\sum_{i=1}^{N}E_i\,|c_i|^2,
\]
et de même, avec $\Hop^{2}\Psi_j=E_j^2\Psi_j$,
\[
  \langle E^2\rangle=\int_0^a\Psi^*\,\Hop^{2}\,\Psi\,\dd x=\sum_{i=1}^{N}E_i^2\,|c_i|^2 ,
  \qquad
  \sigma_E=\sqrt{\langle E^2\rangle-\langle E\rangle^2}.
\]

\begin{resultat}{Probabilité de mesurer l'énergie $E_i$}
\[
  P(E_i)=|c_i|^2 ,\qquad
  \langle E\rangle=\sum_i E_i\,P(E_i).
\]
$c_i$ est l'\emph{amplitude de probabilité} de l'état $\Psi_i$ et $|c_i|^2$ sa
probabilité. Les $|c_i|^2$ ne dépendent pas du temps.
\end{resultat}

La moyenne pondérée garde un sens même si les $E_i$ sont discrets. Elle n'est
d'ailleurs pas forcément une valeur mesurable : une mesure ne donne jamais que l'un des
$E_i$.

\paragraph{Application à l'exemple.} Avec $c_1=c_2=1/\sqrt2$ et $E_2=4E_1$ :
$P(E_1)=P(E_2)=\tfrac12$,
\[
  \langle E\rangle=\frac{E_1+E_2}{2}=\frac{5}{2}\,E_1,
  \qquad
  \langle E^2\rangle=\frac{E_1^2+E_2^2}{2}=\frac{17}{2}\,E_1^2,
  \qquad
  \sigma_E=\sqrt{\frac{17}{2}-\frac{25}{4}}\;E_1=\frac{3}{2}\,E_1\neq0 .
\]
Ici l'énergie est incertaine, et la valeur moyenne $\tfrac52E_1$ n'est jamais obtenue :
on mesure $E_1$ ou $E_2$, chacune avec la probabilité $1/2$.

\begin{remarque}[Effondrement]
Si la mesure donne $E_i$, l'état devient $\Psi_i$ : une seconde mesure faite
aussitôt redonne $E_i$ avec certitude. La probabilité $|c_i|^2$ ne se lit donc que sur
un grand nombre de systèmes identiques préparés dans le même état.
\end{remarque}

% ---------------------------------------------------------------------
\subsection{Retour à la notation de Dirac : le qubit}\label{sec:qubit}
% ---------------------------------------------------------------------

On ne garde que deux états du puits, notés comme un bit quantique (\emph{qubit}) :
$\ket{0}$ pour $\Psi_1$ et $\ket{1}$ pour $\Psi_2$. L'état de l'exemple s'écrit
\[
  \ket{\Psi}=\frac{1}{\sqrt2}\,\ket{0}+\frac{1}{\sqrt2}\,\ket{1}.
\]

\begin{center}
\renewcommand{\arraystretch}{1.5}
\begin{tabular}{@{}ll@{}}
\toprule
\textbf{Fonctions d'onde} & \textbf{Notation de Dirac} \\
\midrule
$\Psi=c_1\Psi_1+c_2\Psi_2$ & $\ket{\Psi}=c_1\ket{0}+c_2\ket{1}$ \\
$\displaystyle\int\Psi_i^*\Psi_j\,\dd x=\delta_{ij}$ & $\braket{i}{j}=\delta_{ij}$ \\
$\displaystyle\int\Psi^*\Psi\,\dd x=\sum_i|c_i|^2=1$ & $\braket{\Psi}{\Psi}=1$ \\
$P(E_i)=|c_i|^2$ & $P(i)=|c_i|^2=|\braket{i}{\Psi}|^2$ \\
\bottomrule
\end{tabular}
\end{center}

\paragraph{Probabilité.} La probabilité que $\ket{\Psi}$ soit mesuré dans l'état
$\ket{0}$ est
\[
  P(\ket{\Psi}\to\ket{0})=P(0)=\left|\frac{1}{\sqrt2}\right|^2=\frac12 .
\]

\paragraph{Normalisation.}
\[
  \braket{\Psi}{\Psi}
  =\Big(\tfrac{1}{\sqrt2}\bra{0}+\tfrac{1}{\sqrt2}\bra{1}\Big)
   \Big(\tfrac{1}{\sqrt2}\ket{0}+\tfrac{1}{\sqrt2}\ket{1}\Big)
  =\tfrac12\braket{0}{0}+\tfrac12\braket{0}{1}+\tfrac12\braket{1}{0}+\tfrac12\braket{1}{1}
  =\tfrac12\,(1+0+0+1)=1 .
\]

\begin{definition}{Base de calcul (base $Z$)}
$\{\ket{0},\ket{1}\}$ est une base \emph{orthonormée} de $\mathcal{H}^2=\mathbb{C}^2$ :
\[
  \ket{0}=\begin{pmatrix}1\\ 0\end{pmatrix},\qquad
  \ket{1}=\begin{pmatrix}0\\ 1\end{pmatrix},\qquad
  \braket{0}{0}=\begin{pmatrix}1 & 0\end{pmatrix}\begin{pmatrix}1\\ 0\end{pmatrix}=1,\qquad
  \braket{0}{1}=\begin{pmatrix}1 & 0\end{pmatrix}\begin{pmatrix}0\\ 1\end{pmatrix}=0 .
\]
Tout état s'écrit $\ket{\Psi}=\alpha\ket{0}+\beta\ket{1}=\begin{pmatrix}\alpha\\ \beta\end{pmatrix}$
avec $|\alpha|^2+|\beta|^2=1$ ; alors $P(0)=|\alpha|^2$ et $P(1)=|\beta|^2$.
\end{definition}

\paragraph{Exemple.} Soit $\ket{\Psi}=\dfrac{1}{\sqrt5}\,\ket{0}+\dfrac{2}{\sqrt5}\,\ket{1}$
dans $\mathcal{H}^2$.
\begin{enumerate}[nosep]
  \item \emph{$\ket{\Psi}$ est-il normalisé ?}
        $\braket{\Psi}{\Psi}=\dfrac15+\dfrac45=1$ : oui.
  \item $P(0)=\left|\dfrac{1}{\sqrt5}\right|^2=\dfrac15$.
  \item $P(1)=\left|\dfrac{2}{\sqrt5}\right|^2=\dfrac45$.
\end{enumerate}

\begin{remarque}[État non normalisé]
Si $\braket{\Psi}{\Psi}\neq1$, on divise l'état par sa norme $\sqrt{\braket{\Psi}{\Psi}}$.
Par exemple $\ket{0}+2\ket{1}$ a pour norme au carré $1+4=5$ : l'état normalisé est
$\frac{1}{\sqrt5}\ket{0}+\frac{2}{\sqrt5}\ket{1}$, celui de l'exemple.
\end{remarque}

% ---------------------------------------------------------------------
\subsection{Projecteurs et mesure}\label{sec:projecteurs}
% ---------------------------------------------------------------------

Le produit d'un ket par un bra est un opérateur (section~\ref{sec:dirac}). Avec les
états de la base de calcul, on obtient deux \emph{projecteurs}.

\begin{definition}{Projecteurs sur $\ket{0}$ et sur $\ket{1}$}
\[
  \hat P_0=\ket{0}\bra{0}=\begin{pmatrix}1\\ 0\end{pmatrix}\begin{pmatrix}1 & 0\end{pmatrix}
  =\begin{pmatrix}1 & 0\\ 0 & 0\end{pmatrix},
  \qquad
  \hat P_1=\ket{1}\bra{1}=\begin{pmatrix}0\\ 1\end{pmatrix}\begin{pmatrix}0 & 1\end{pmatrix}
  =\begin{pmatrix}0 & 0\\ 0 & 1\end{pmatrix}.
\]
Ils vérifient
\[
  \hat P_i^{\,2}=\hat P_i,
  \qquad
  \hat P_0+\hat P_1=\begin{pmatrix}1 & 0\\ 0 & 1\end{pmatrix}=\mathrm{Id}
  \quad(\text{complétude}).
\]
\end{definition}

En effet $\hat P_0^{\,2}=\ket{0}\braket{0}{0}\bra{0}=\ket{0}\bra{0}$ car
$\braket{0}{0}=1$. Appliqué à $\ket{\Psi}=\alpha\ket{0}+\beta\ket{1}$, $\hat P_0$
ne garde que la composante sur $\ket{0}$ : $\hat P_0\ket{\Psi}=\alpha\ket{0}$.

\begin{resultat}{Probabilité de mesure}
\[
  P(\ket{\Psi}\to\ket{i})=\bra{\Psi}\hat P_i\ket{\Psi}
  =\braket{\Psi}{i}\braket{i}{\Psi}=|\braket{i}{\Psi}|^2 .
\]
La complétude donne bien $\sum_i P(i)=\bra{\Psi}\big(\sum_i\hat P_i\big)\ket{\Psi}=\braket{\Psi}{\Psi}=1$.
\end{resultat}

\paragraph{Exemple.} Pour $\ket{\Psi}=\dfrac{1}{\sqrt5}\ket{0}+\dfrac{2}{\sqrt5}\ket{1}
=\begin{pmatrix}1/\sqrt5\\ 2/\sqrt5\end{pmatrix}$ :
\[
  \hat P_0\ket{\Psi}=\begin{pmatrix}1 & 0\\ 0 & 0\end{pmatrix}
  \begin{pmatrix}1/\sqrt5\\ 2/\sqrt5\end{pmatrix}
  =\begin{pmatrix}1/\sqrt5\\ 0\end{pmatrix},
  \qquad
  P(\ket{\Psi}\to\ket{0})=\bra{\Psi}\hat P_0\ket{\Psi}
  =\begin{pmatrix}\frac{1}{\sqrt5} & \frac{2}{\sqrt5}\end{pmatrix}
   \begin{pmatrix}\frac{1}{\sqrt5}\\ 0\end{pmatrix}=\frac15 ,
\]
\[
  P(\ket{\Psi}\to\ket{1})=\bra{\Psi}\hat P_1\ket{\Psi}
  =\begin{pmatrix}\frac{1}{\sqrt5} & \frac{2}{\sqrt5}\end{pmatrix}
   \begin{pmatrix}0 & 0\\ 0 & 1\end{pmatrix}
   \begin{pmatrix}\frac{1}{\sqrt5}\\ \frac{2}{\sqrt5}\end{pmatrix}=\frac45 .
\]

\paragraph{Effet d'une phase relative.} Soit $\ket{\Psi'}=\dfrac{1}{\sqrt5}\ket{0}
+\dfrac{2\ii}{\sqrt5}\ket{1}$. Son bra est le \emph{conjugué} :
$\bra{\Psi'}=\begin{pmatrix}\frac{1}{\sqrt5} & -\frac{2\ii}{\sqrt5}\end{pmatrix}$, donc
\[
  \braket{\Psi'}{\Psi'}=\frac15+\Big(-\frac{2\ii}{\sqrt5}\Big)\Big(\frac{2\ii}{\sqrt5}\Big)
  =\frac15+\frac45=1 ,
\]
et $P(0)=1/5$, $P(1)=4/5$ comme pour $\ket{\Psi}$. Les deux états sont pourtant
différents : le facteur $\ii$ est une phase \emph{relative} entre les deux composantes,
que la mesure dans la base $Z$ ne voit pas (contrairement à une phase globale, jamais
observable).

\begin{remarque}[Rappels]
$\hat P^{\,2}=\hat P$ définit un projecteur ; $\braket{a}{b}=\braket{b}{a}^*$ ;
$\hat P_i$ est hermitien : $\hat P_i^{\,\dagger}=\hat P_i$.
\end{remarque}

% ---------------------------------------------------------------------
\subsection{Récapitulatif du chapitre}
% ---------------------------------------------------------------------

\begin{center}
\renewcommand{\arraystretch}{1.7}
\begin{tabular}{@{}ll@{}}
\toprule
\textbf{Notion} & \textbf{Formule} \\
\midrule
Superposition & $\Psi=\displaystyle\sum_i c_i\Psi_i$, \quad $c_i\in\mathbb{C}$ \\
Orthonormalité & $\displaystyle\int\Psi_i^*\Psi_j\,\dd x=\delta_{ij}$ \\
Normalisation & $\displaystyle\sum_i|c_i|^2=1$ \\
Probabilité d'une énergie & $P(E_i)=|c_i|^2$ \\
Énergie moyenne & $\langle E\rangle=\displaystyle\sum_iE_i|c_i|^2$ \\
Incertitude & $\sigma_E=\sqrt{\langle E^2\rangle-\langle E\rangle^2}$, \quad
              $\langle E^2\rangle=\displaystyle\sum_iE_i^2|c_i|^2$ \\
Qubit & $\ket{\Psi}=\alpha\ket{0}+\beta\ket{1}$, \quad $|\alpha|^2+|\beta|^2=1$ \\
Mesure dans la base $Z$ & $P(0)=|\alpha|^2$, \quad $P(1)=|\beta|^2$ \\
Projecteurs & $\hat P_i=\ket{i}\bra{i}$, \quad $\hat P_i^{\,2}=\hat P_i$, \quad $\hat P_0+\hat P_1=\mathrm{Id}$ \\
Probabilité de mesure & $P(i)=\bra{\Psi}\hat P_i\ket{\Psi}=|\braket{i}{\Psi}|^2$ \\
\bottomrule
\end{tabular}
\end{center}
.
%  Il PROLONGE chapitre2.tex : pas de \section ici, les sous-sections
%  1.10 à 1.13 font suite à 1.1 à 1.9 (la suite est le chapitre4.tex, section 2).
%  Pas de préambule non plus : les
%  environnements (definition, resultat, remarque), les macros (\ii, \ee,
%  \dd, \Hop, \ket, \bra, \braket) et les couleurs (insarouge, insagris)
%  sont définis dans main.tex.
% =====================================================================

% Macros de Dirac (déjà définies dans main.tex récent ; ce bloc évite l'erreur
% « Undefined control sequence \ket » si main.tex est une ancienne version).
\providecommand{\ket}[1]{|#1\rangle}
\providecommand{\bra}[1]{\langle #1|}
\providecommand{\braket}[2]{\langle #1|#2\rangle}

% ---------------------------------------------------------------------
\subsection{Superposition d'états stationnaires}\label{sec:superposition}
% ---------------------------------------------------------------------

\subsubsection{Un état qui n'est pas un état propre}

Dans un état propre $\Psi_n$ du puits infini, l'énergie est certaine
($\sigma_E=0$, section~\ref{sec:solutions}). Rien n'oblige pourtant la particule à
être dans un état propre. Par linéarité de l'équation de Schrödinger
(3\textsuperscript{e} hypothèse de la section~\ref{sec:construction}), toute
combinaison de solutions est encore solution.

\begin{definition}{Superposition quantique}
\[
  \Psi(x,t)=\sum_{i=1}^{N}c_i\,\Psi_i(x,t),
  \qquad c_i\in\mathbb{C},
\]
où les $\Psi_i$ sont des états stationnaires d'énergies $E_i$ ($i$ est le nombre
quantique de l'état $\Psi_i$). L'état n'a plus une énergie unique : il est
\emph{probabiliste}, l'énergie mesurée pouvant valoir chacune des $E_i$.
\end{definition}

\subsubsection{Exemple : deux états du puits}

À $t=0$, on prépare la particule dans l'état
\[
  \Psi(x,0)=A\left[\sqrt{\frac{2}{a}}\,\sin\!\Big(\frac{\pi x}{a}\Big)
                  +\sqrt{\frac{2}{a}}\,\sin\!\Big(\frac{2\pi x}{a}\Big)\right]
  =A\,\big[\psi_1(x)+\psi_2(x)\big].
\]
La normalisation $\int_0^a\Psi^*\Psi\,\dd x=1$ donne, en développant,
\[
  A^2\left[\int_0^a\psi_1^2\,\dd x+2\int_0^a\psi_1\psi_2\,\dd x+\int_0^a\psi_2^2\,\dd x\right]
  =A^2\,(1+0+1)=2A^2=1
  \quad\Longrightarrow\quad A=\frac{1}{\sqrt2}.
\]
Les $\psi_n$ étant normalisées, $\int\psi_n^2\,\dd x=1$. Le terme croisé s'annule :
\[
  \int_0^a\psi_1\psi_2\,\dd x
  =\frac{1}{a}\int_0^a\Big[\cos\frac{\pi x}{a}-\cos\frac{3\pi x}{a}\Big]\dd x=0 .
\]
Chaque terme évolue avec sa propre phase (chacun est une solution stationnaire), donc
\[
  \Psi(x,t)=\frac{1}{\sqrt2}\Big[\psi_1(x)\,\ee^{-\ii E_1t/\hbar}
                                 +\psi_2(x)\,\ee^{-\ii E_2t/\hbar}\Big],
\]
\[
  |\Psi(x,t)|^2=\frac12\big[\psi_1^2+\psi_2^2\big]
                +\psi_1\psi_2\,\cos\!\Big(\frac{E_2-E_1}{\hbar}\,t\Big).
\]
Le terme croisé oscille à la pulsation $(E_2-E_1)/\hbar$ : contrairement à un état
stationnaire, la densité de probabilité dépend du temps.

\begin{figure}[!ht]
\centering
\begin{tikzpicture}[>=Stealth, x=1cm, y=1cm]
  \draw[->] (-0.3,0) -- (6.0,0) node[right] {$x$};
  \draw[->] (0,0) -- (0,4.2) node[above] {$|\Psi|^2$};
  \draw[thick, insagris, dashed, domain=0:1, samples=100]
    plot ({5*\x}, {1.0*(sin(180*\x)-sin(360*\x))^2});
  \draw[thick, insarouge, domain=0:1, samples=100]
    plot ({5*\x}, {1.0*(sin(180*\x)+sin(360*\x))^2});
  \node[below] at (0,0) {$0$};
  \node[below] at (5,0) {$a$};
  \draw[insarouge, thick] (0.6,3.8) -- (1.2,3.8) node[right, black, font=\small] {$t=0$};
  \draw[insagris, thick, dashed] (3.2,3.8) -- (3.8,3.8)
    node[right, black, font=\small] {$t=T/2$};
\end{tikzpicture}
\caption{Densité de probabilité de $(\psi_1+\psi_2)/\sqrt2$ à $t=0$ et à une
demi-période $T/2$ (avec $T=2\pi\hbar/(E_2-E_1)$) : le maximum passe de la gauche à
la droite du puits. Unité verticale : $1/a$.}
\end{figure}

\subsubsection{Cas général : normalisation}

Pour $\Psi=\sum_i c_i\Psi_i$, la condition $\int\Psi^*\Psi\,\dd x=1$ s'écrit
\[
  \int_0^a\Big[\sum_i c_i^*\Psi_i^*\Big]\Big[\sum_j c_j\Psi_j\Big]\dd x
  =|c_1|^2\!\int\!\Psi_1^*\Psi_1\,\dd x+c_1^*c_2\!\int\!\Psi_1^*\Psi_2\,\dd x+\cdots=1 .
\]
Les $\Psi_i$ sont orthonormées, comme des vecteurs unitaires : le produit scalaire d'un
vecteur avec lui-même vaut 1, celui de deux vecteurs orthogonaux vaut 0.

\begin{resultat}{Orthonormalité et normalisation}
\[
  \int_0^a\Psi_i^*\Psi_j\,\dd x=\delta_{ij}=
  \begin{cases}1 & \text{si } i=j,\\ 0 & \text{si } i\neq j,\end{cases}
  \qquad\Longrightarrow\qquad
  |c_1|^2+|c_2|^2+\cdots+|c_N|^2=\sum_{i=1}^{N}|c_i|^2=1
\]
\end{resultat}

L'orthogonalité tient à tout instant, car les phases $\ee^{\ii(E_i-E_j)t/\hbar}$
multiplient une intégrale nulle. En multipliant $\Psi=\sum_j c_j\Psi_j$ par $\Psi_i^*$
et en intégrant, on obtient aussi $c_i=\int\Psi_i^*\,\Psi\,\dd x$.

\subsubsection{Probabilités et énergie moyenne}

L'énergie moyenne est la valeur moyenne de $\Hop$, avec $\Hop\Psi_j=E_j\Psi_j$ :
\[
  \langle E\rangle=\int_0^a\Psi^*\,\Hop\,\Psi\,\dd x
  =\sum_{i,j}c_i^*c_jE_j\!\int\!\Psi_i^*\Psi_j\,\dd x
  =\sum_{i=1}^{N}E_i\,|c_i|^2,
\]
et de même, avec $\Hop^{2}\Psi_j=E_j^2\Psi_j$,
\[
  \langle E^2\rangle=\int_0^a\Psi^*\,\Hop^{2}\,\Psi\,\dd x=\sum_{i=1}^{N}E_i^2\,|c_i|^2 ,
  \qquad
  \sigma_E=\sqrt{\langle E^2\rangle-\langle E\rangle^2}.
\]

\begin{resultat}{Probabilité de mesurer l'énergie $E_i$}
\[
  P(E_i)=|c_i|^2 ,\qquad
  \langle E\rangle=\sum_i E_i\,P(E_i).
\]
$c_i$ est l'\emph{amplitude de probabilité} de l'état $\Psi_i$ et $|c_i|^2$ sa
probabilité. Les $|c_i|^2$ ne dépendent pas du temps.
\end{resultat}

La moyenne pondérée garde un sens même si les $E_i$ sont discrets. Elle n'est
d'ailleurs pas forcément une valeur mesurable : une mesure ne donne jamais que l'un des
$E_i$.

\paragraph{Application à l'exemple.} Avec $c_1=c_2=1/\sqrt2$ et $E_2=4E_1$ :
$P(E_1)=P(E_2)=\tfrac12$,
\[
  \langle E\rangle=\frac{E_1+E_2}{2}=\frac{5}{2}\,E_1,
  \qquad
  \langle E^2\rangle=\frac{E_1^2+E_2^2}{2}=\frac{17}{2}\,E_1^2,
  \qquad
  \sigma_E=\sqrt{\frac{17}{2}-\frac{25}{4}}\;E_1=\frac{3}{2}\,E_1\neq0 .
\]
Ici l'énergie est incertaine, et la valeur moyenne $\tfrac52E_1$ n'est jamais obtenue :
on mesure $E_1$ ou $E_2$, chacune avec la probabilité $1/2$.

\begin{remarque}[Effondrement]
Si la mesure donne $E_i$, l'état devient $\Psi_i$ : une seconde mesure faite
aussitôt redonne $E_i$ avec certitude. La probabilité $|c_i|^2$ ne se lit donc que sur
un grand nombre de systèmes identiques préparés dans le même état.
\end{remarque}

% ---------------------------------------------------------------------
\subsection{Retour à la notation de Dirac : le qubit}\label{sec:qubit}
% ---------------------------------------------------------------------

On ne garde que deux états du puits, notés comme un bit quantique (\emph{qubit}) :
$\ket{0}$ pour $\Psi_1$ et $\ket{1}$ pour $\Psi_2$. L'état de l'exemple s'écrit
\[
  \ket{\Psi}=\frac{1}{\sqrt2}\,\ket{0}+\frac{1}{\sqrt2}\,\ket{1}.
\]

\begin{center}
\renewcommand{\arraystretch}{1.5}
\begin{tabular}{@{}ll@{}}
\toprule
\textbf{Fonctions d'onde} & \textbf{Notation de Dirac} \\
\midrule
$\Psi=c_1\Psi_1+c_2\Psi_2$ & $\ket{\Psi}=c_1\ket{0}+c_2\ket{1}$ \\
$\displaystyle\int\Psi_i^*\Psi_j\,\dd x=\delta_{ij}$ & $\braket{i}{j}=\delta_{ij}$ \\
$\displaystyle\int\Psi^*\Psi\,\dd x=\sum_i|c_i|^2=1$ & $\braket{\Psi}{\Psi}=1$ \\
$P(E_i)=|c_i|^2$ & $P(i)=|c_i|^2=|\braket{i}{\Psi}|^2$ \\
\bottomrule
\end{tabular}
\end{center}

\paragraph{Probabilité.} La probabilité que $\ket{\Psi}$ soit mesuré dans l'état
$\ket{0}$ est
\[
  P(\ket{\Psi}\to\ket{0})=P(0)=\left|\frac{1}{\sqrt2}\right|^2=\frac12 .
\]

\paragraph{Normalisation.}
\[
  \braket{\Psi}{\Psi}
  =\Big(\tfrac{1}{\sqrt2}\bra{0}+\tfrac{1}{\sqrt2}\bra{1}\Big)
   \Big(\tfrac{1}{\sqrt2}\ket{0}+\tfrac{1}{\sqrt2}\ket{1}\Big)
  =\tfrac12\braket{0}{0}+\tfrac12\braket{0}{1}+\tfrac12\braket{1}{0}+\tfrac12\braket{1}{1}
  =\tfrac12\,(1+0+0+1)=1 .
\]

\begin{definition}{Base de calcul (base $Z$)}
$\{\ket{0},\ket{1}\}$ est une base \emph{orthonormée} de $\mathcal{H}^2=\mathbb{C}^2$ :
\[
  \ket{0}=\begin{pmatrix}1\\ 0\end{pmatrix},\qquad
  \ket{1}=\begin{pmatrix}0\\ 1\end{pmatrix},\qquad
  \braket{0}{0}=\begin{pmatrix}1 & 0\end{pmatrix}\begin{pmatrix}1\\ 0\end{pmatrix}=1,\qquad
  \braket{0}{1}=\begin{pmatrix}1 & 0\end{pmatrix}\begin{pmatrix}0\\ 1\end{pmatrix}=0 .
\]
Tout état s'écrit $\ket{\Psi}=\alpha\ket{0}+\beta\ket{1}=\begin{pmatrix}\alpha\\ \beta\end{pmatrix}$
avec $|\alpha|^2+|\beta|^2=1$ ; alors $P(0)=|\alpha|^2$ et $P(1)=|\beta|^2$.
\end{definition}

\paragraph{Exemple.} Soit $\ket{\Psi}=\dfrac{1}{\sqrt5}\,\ket{0}+\dfrac{2}{\sqrt5}\,\ket{1}$
dans $\mathcal{H}^2$.
\begin{enumerate}[nosep]
  \item \emph{$\ket{\Psi}$ est-il normalisé ?}
        $\braket{\Psi}{\Psi}=\dfrac15+\dfrac45=1$ : oui.
  \item $P(0)=\left|\dfrac{1}{\sqrt5}\right|^2=\dfrac15$.
  \item $P(1)=\left|\dfrac{2}{\sqrt5}\right|^2=\dfrac45$.
\end{enumerate}

\begin{remarque}[État non normalisé]
Si $\braket{\Psi}{\Psi}\neq1$, on divise l'état par sa norme $\sqrt{\braket{\Psi}{\Psi}}$.
Par exemple $\ket{0}+2\ket{1}$ a pour norme au carré $1+4=5$ : l'état normalisé est
$\frac{1}{\sqrt5}\ket{0}+\frac{2}{\sqrt5}\ket{1}$, celui de l'exemple.
\end{remarque}

% ---------------------------------------------------------------------
\subsection{Projecteurs et mesure}\label{sec:projecteurs}
% ---------------------------------------------------------------------

Le produit d'un ket par un bra est un opérateur (section~\ref{sec:dirac}). Avec les
états de la base de calcul, on obtient deux \emph{projecteurs}.

\begin{definition}{Projecteurs sur $\ket{0}$ et sur $\ket{1}$}
\[
  \hat P_0=\ket{0}\bra{0}=\begin{pmatrix}1\\ 0\end{pmatrix}\begin{pmatrix}1 & 0\end{pmatrix}
  =\begin{pmatrix}1 & 0\\ 0 & 0\end{pmatrix},
  \qquad
  \hat P_1=\ket{1}\bra{1}=\begin{pmatrix}0\\ 1\end{pmatrix}\begin{pmatrix}0 & 1\end{pmatrix}
  =\begin{pmatrix}0 & 0\\ 0 & 1\end{pmatrix}.
\]
Ils vérifient
\[
  \hat P_i^{\,2}=\hat P_i,
  \qquad
  \hat P_0+\hat P_1=\begin{pmatrix}1 & 0\\ 0 & 1\end{pmatrix}=\mathrm{Id}
  \quad(\text{complétude}).
\]
\end{definition}

En effet $\hat P_0^{\,2}=\ket{0}\braket{0}{0}\bra{0}=\ket{0}\bra{0}$ car
$\braket{0}{0}=1$. Appliqué à $\ket{\Psi}=\alpha\ket{0}+\beta\ket{1}$, $\hat P_0$
ne garde que la composante sur $\ket{0}$ : $\hat P_0\ket{\Psi}=\alpha\ket{0}$.

\begin{resultat}{Probabilité de mesure}
\[
  P(\ket{\Psi}\to\ket{i})=\bra{\Psi}\hat P_i\ket{\Psi}
  =\braket{\Psi}{i}\braket{i}{\Psi}=|\braket{i}{\Psi}|^2 .
\]
La complétude donne bien $\sum_i P(i)=\bra{\Psi}\big(\sum_i\hat P_i\big)\ket{\Psi}=\braket{\Psi}{\Psi}=1$.
\end{resultat}

\paragraph{Exemple.} Pour $\ket{\Psi}=\dfrac{1}{\sqrt5}\ket{0}+\dfrac{2}{\sqrt5}\ket{1}
=\begin{pmatrix}1/\sqrt5\\ 2/\sqrt5\end{pmatrix}$ :
\[
  \hat P_0\ket{\Psi}=\begin{pmatrix}1 & 0\\ 0 & 0\end{pmatrix}
  \begin{pmatrix}1/\sqrt5\\ 2/\sqrt5\end{pmatrix}
  =\begin{pmatrix}1/\sqrt5\\ 0\end{pmatrix},
  \qquad
  P(\ket{\Psi}\to\ket{0})=\bra{\Psi}\hat P_0\ket{\Psi}
  =\begin{pmatrix}\frac{1}{\sqrt5} & \frac{2}{\sqrt5}\end{pmatrix}
   \begin{pmatrix}\frac{1}{\sqrt5}\\ 0\end{pmatrix}=\frac15 ,
\]
\[
  P(\ket{\Psi}\to\ket{1})=\bra{\Psi}\hat P_1\ket{\Psi}
  =\begin{pmatrix}\frac{1}{\sqrt5} & \frac{2}{\sqrt5}\end{pmatrix}
   \begin{pmatrix}0 & 0\\ 0 & 1\end{pmatrix}
   \begin{pmatrix}\frac{1}{\sqrt5}\\ \frac{2}{\sqrt5}\end{pmatrix}=\frac45 .
\]

\paragraph{Effet d'une phase relative.} Soit $\ket{\Psi'}=\dfrac{1}{\sqrt5}\ket{0}
+\dfrac{2\ii}{\sqrt5}\ket{1}$. Son bra est le \emph{conjugué} :
$\bra{\Psi'}=\begin{pmatrix}\frac{1}{\sqrt5} & -\frac{2\ii}{\sqrt5}\end{pmatrix}$, donc
\[
  \braket{\Psi'}{\Psi'}=\frac15+\Big(-\frac{2\ii}{\sqrt5}\Big)\Big(\frac{2\ii}{\sqrt5}\Big)
  =\frac15+\frac45=1 ,
\]
et $P(0)=1/5$, $P(1)=4/5$ comme pour $\ket{\Psi}$. Les deux états sont pourtant
différents : le facteur $\ii$ est une phase \emph{relative} entre les deux composantes,
que la mesure dans la base $Z$ ne voit pas (contrairement à une phase globale, jamais
observable).

\begin{remarque}[Rappels]
$\hat P^{\,2}=\hat P$ définit un projecteur ; $\braket{a}{b}=\braket{b}{a}^*$ ;
$\hat P_i$ est hermitien : $\hat P_i^{\,\dagger}=\hat P_i$.
\end{remarque}

% ---------------------------------------------------------------------
\subsection{Récapitulatif du chapitre}
% ---------------------------------------------------------------------

\begin{center}
\renewcommand{\arraystretch}{1.7}
\begin{tabular}{@{}ll@{}}
\toprule
\textbf{Notion} & \textbf{Formule} \\
\midrule
Superposition & $\Psi=\displaystyle\sum_i c_i\Psi_i$, \quad $c_i\in\mathbb{C}$ \\
Orthonormalité & $\displaystyle\int\Psi_i^*\Psi_j\,\dd x=\delta_{ij}$ \\
Normalisation & $\displaystyle\sum_i|c_i|^2=1$ \\
Probabilité d'une énergie & $P(E_i)=|c_i|^2$ \\
Énergie moyenne & $\langle E\rangle=\displaystyle\sum_iE_i|c_i|^2$ \\
Incertitude & $\sigma_E=\sqrt{\langle E^2\rangle-\langle E\rangle^2}$, \quad
              $\langle E^2\rangle=\displaystyle\sum_iE_i^2|c_i|^2$ \\
Qubit & $\ket{\Psi}=\alpha\ket{0}+\beta\ket{1}$, \quad $|\alpha|^2+|\beta|^2=1$ \\
Mesure dans la base $Z$ & $P(0)=|\alpha|^2$, \quad $P(1)=|\beta|^2$ \\
Projecteurs & $\hat P_i=\ket{i}\bra{i}$, \quad $\hat P_i^{\,2}=\hat P_i$, \quad $\hat P_0+\hat P_1=\mathrm{Id}$ \\
Probabilité de mesure & $P(i)=\bra{\Psi}\hat P_i\ket{\Psi}=|\braket{i}{\Psi}|^2$ \\
\bottomrule
\end{tabular}
\end{center}
 % 3. Opérateurs, matrice de densité, sphère de Bloch, produit tensoriel
% =====================================================================
%  Chapitre 4 : Portes quantiques et systèmes multi-qubits
%  Fichier importé par main.tex avec % =====================================================================
%  Chapitre 4 : Portes quantiques et systèmes multi-qubits
%  Fichier importé par main.tex avec % =====================================================================
%  Chapitre 4 : Portes quantiques et systèmes multi-qubits
%  Fichier importé par main.tex avec \include{chapitre4}.
%  Il ne contient ni préambule ni \begin{document} : les environnements
%  (definition, resultat, remarque, aretenir, savoirfaire, exercice, correction),
%  les macros (\ii, \ee, \dd, \pd, \pdd, \Hop, \ket, \bra, \braket), les symboles
%  de circuits (ctl, cible, croix, mesure, porte) et les couleurs (insarouge,
%  insagris) sont définis dans main.tex. Un chapitre = une \section.
% =====================================================================

% Macros de Dirac (déjà définies dans main.tex ; ce bloc évite l'erreur
% « Undefined control sequence \ket » si main.tex est une ancienne version).
\providecommand{\ket}[1]{|#1\rangle}
\providecommand{\bra}[1]{\langle #1|}
\providecommand{\braket}[2]{\langle #1|#2\rangle}

% =====================================================================
\section{Portes quantiques et systèmes multi-qubits}\label{chap:portes}
% =====================================================================

Un calcul quantique est une suite de transformations unitaires appliquées à des qubits. Ce
chapitre décrit les portes à un qubit (Pauli, Hadamard, phase) et leur effet sur la sphère de
Bloch, la lecture d'un circuit, puis les portes qui agissent sur plusieurs qubits : CNOT, SWAP et
portes contrôlées.

% ---------------------------------------------------------------------
\subsection{Portes quantiques}\label{sec:portes}
% ---------------------------------------------------------------------

Un calcul quantique transforme des états. Les opérations sur un seul qubit sont appelées
\emph{portes} ; les sections suivantes passent à plusieurs qubits.

\begin{resultat}{Postulat 4 : évolution d'un système isolé}
L'évolution d'un système isolé est décrite par un opérateur unitaire : $\ket{\Psi}\to U\ket{\Psi}$
avec $U^\dagger U=\mathrm{Id}$.
\end{resultat}

C'est la forme de l'équation de Schrödinger du chapitre~\ref{chap:histoire}. Pour un hamiltonien
$\Hop$ indépendant du temps, $\ii\hbar\,\partial_t\ket{\Psi}=\Hop\ket{\Psi}$ a pour solution
$\ket{\Psi(t)}=U(t)\ket{\Psi(0)}$ avec $U(t)=\ee^{-\ii\Hop t/\hbar}$. Comme $\Hop$ est hermitien,
$U(t)^\dagger=\ee^{+\ii\Hop t/\hbar}$ et $U^\dagger U=\mathrm{Id}$ : l'évolution est unitaire. Une
porte est une telle évolution, obtenue en contrôlant $\Hop$ pendant une durée donnée.

\begin{definition}{Porte quantique}
Une porte à un qubit est un opérateur $U$ (matrice $2\times2$) \emph{unitaire} :
$U^\dagger U=\mathrm{Id}$. Elle transforme $\ket{\Psi}\to U\ket{\Psi}$, conserve la norme
($\bra{\Psi}U^\dagger U\ket{\Psi}=\braket{\Psi}{\Psi}=1$) et est réversible
($U^{-1}=U^\dagger$). Sur la sphère de Bloch (section~\ref{sec:bloch}), c'est une
\emph{rotation}.
\end{definition}

% ---------------------------------------------------------------------
\subsection{Portes de Pauli}\label{sec:pauli}
% ---------------------------------------------------------------------

Les matrices de Pauli de la section~\ref{sec:op-dirac} sont trois portes.

\begin{definition}{Portes $X$, $Y$, $Z$}
\begin{center}
\renewcommand{\arraystretch}{1}
\begin{tabular}{@{}lccl@{}}
\textbf{Porte} & \textbf{Matrice} & \textbf{Ket-bra} & \textbf{Nom} \\[0.8ex]
$X=\sigma_x$ & $\begin{pmatrix}0 & 1\\ 1 & 0\end{pmatrix}$
  & $\ket{0}\bra{1}+\ket{1}\bra{0}$ & bit flip \\[2.2ex]
$Y=\sigma_y$ & $\begin{pmatrix}0 & -\ii\\ \ii & 0\end{pmatrix}$
  & $-\ii\ket{0}\bra{1}+\ii\ket{1}\bra{0}$ &  \\[2.2ex]
$Z=\sigma_z$ & $\begin{pmatrix}1 & 0\\ 0 & -1\end{pmatrix}$
  & $\ket{0}\bra{0}-\ket{1}\bra{1}$ & phase flip \\
\end{tabular}
\end{center}
\end{definition}

\paragraph{Porte $X$ (inverseur).} Avec $\braket{0}{0}=\braket{1}{1}=1$ et
$\braket{0}{1}=\braket{1}{0}=0$ :
\begin{align*}
  X\ket{0}&=\ket{0}\braket{1}{0}+\ket{1}\braket{0}{0}=0\cdot\ket{0}+1\cdot\ket{1}=\ket{1},\\
  X\ket{1}&=\ket{0}\braket{1}{1}+\ket{1}\braket{0}{1}=1\cdot\ket{0}+0\cdot\ket{1}=\ket{0}.
\end{align*}
$X$ échange $\ket{0}$ et $\ket{1}$ : c'est le NOT quantique (\emph{bit flip}). Sur la
sphère de Bloch, c'est une rotation de $\pi$ autour de l'axe $x$ : le pôle nord devient
le pôle sud. Les états de l'axe $x$ ne bougent pas : $X\ket{+}=\ket{+}$ et
$X\ket{-}=-\ket{-}$.

\paragraph{Porte $Z$ (changement de phase).}
\begin{align*}
  Z\ket{0}&=\ket{0}\braket{0}{0}-\ket{1}\braket{1}{0}=\ket{0},\\
  Z\ket{1}&=\ket{0}\braket{0}{1}-\ket{1}\braket{1}{1}=-\ket{1},
\end{align*}
donc $Z$ change le signe de la composante sur $\ket{1}$ (\emph{phase flip}). Sur les
états de la base $X$, par linéarité :
\begin{align*}
  Z\ket{+}&=\tfrac{1}{\sqrt2}\big(Z\ket{0}+Z\ket{1}\big)=\tfrac{1}{\sqrt2}\big(\ket{0}-\ket{1}\big)=\ket{-},\\
  Z\ket{-}&=\tfrac{1}{\sqrt2}\big(Z\ket{0}-Z\ket{1}\big)=\tfrac{1}{\sqrt2}\big(\ket{0}+\ket{1}\big)=\ket{+}.
\end{align*}
Sous forme matricielle, $Z\ket{+}=\begin{pmatrix}1 & 0\\ 0 & -1\end{pmatrix}
\frac{1}{\sqrt2}\begin{pmatrix}1\\ 1\end{pmatrix}=\frac{1}{\sqrt2}\begin{pmatrix}1\\ -1\end{pmatrix}=\ket{-}$.
$Z$ est une rotation de $\pi$ autour de l'axe $z$ : elle laisse $\ket{0}$ et $\ket{1}$ en
place et échange $\ket{+}$ et $\ket{-}$.

\paragraph{Porte $Y$.} Avec $Y=-\ii\ket{0}\bra{1}+\ii\ket{1}\bra{0}$ :
\begin{align*}
  Y\ket{0}&=-\ii\ket{0}\braket{1}{0}+\ii\ket{1}\braket{0}{0}=\ii\ket{1},\\
  Y\ket{1}&=-\ii\ket{0}\braket{1}{1}+\ii\ket{1}\braket{0}{1}=-\ii\ket{0}.
\end{align*}
Sur les états de la base $X$ :
\begin{align*}
  Y\ket{+}&=\tfrac{1}{\sqrt2}\big(\ii\ket{1}-\ii\ket{0}\big)=-\ii\,\tfrac{1}{\sqrt2}\big(\ket{0}-\ket{1}\big)=-\ii\ket{-},\\
  Y\ket{-}&=\tfrac{1}{\sqrt2}\big(\ii\ket{1}+\ii\ket{0}\big)=\ii\,\tfrac{1}{\sqrt2}\big(\ket{0}+\ket{1}\big)=\ii\ket{+}.
\end{align*}
Un facteur global ($-\ii$ ou $\ii$) ne change pas l'état physique : $Y$ échange bien
$\ket{+}$ et $\ket{-}$. Sur les états de l'axe $y$ :
\begin{align*}
  Y\ket{{+}\ii}&=\tfrac{1}{\sqrt2}\big(Y\ket{0}+\ii Y\ket{1}\big)
   =\tfrac{1}{\sqrt2}\big(\ii\ket{1}+\ket{0}\big)=\ket{{+}\ii},\\
  Y\ket{{-}\ii}&=\tfrac{1}{\sqrt2}\big(Y\ket{0}-\ii Y\ket{1}\big)
   =\tfrac{1}{\sqrt2}\big(\ii\ket{1}-\ket{0}\big)=-\ket{{-}\ii}.
\end{align*}
$Y$ est une rotation de $\pi$ autour de l'axe $y$.

\begin{remarque}[Chaque Pauli fixe son axe]
$Z\ket{0}=\ket{0}$, $Z\ket{1}=-\ket{1}$ ; $X\ket{\pm}=\pm\ket{\pm}$ ; $Y\ket{{\pm}\ii}=\pm\ket{{\pm}\ii}$.
Les bases $Z$, $X$, $Y$ de la section~\ref{sec:bloch} sont les bases propres de
$Z$, $X$, $Y$. Une porte de Pauli laisse donc son propre axe invariant et retourne les
deux autres ; c'est pourquoi les exemples utilisent les états des deux autres axes.
\end{remarque}

% ---------------------------------------------------------------------
\subsection{Porte de Hadamard}\label{sec:hadamard}
% ---------------------------------------------------------------------

\begin{definition}{Porte de Hadamard}
\[
  H=\frac{1}{\sqrt2}\begin{pmatrix}1 & 1\\ 1 & -1\end{pmatrix}
  =\ket{+}\bra{0}+\ket{-}\bra{1}.
\]
\end{definition}

\begin{align*}
  H\ket{0}&=\frac{1}{\sqrt2}\begin{pmatrix}1 & 1\\ 1 & -1\end{pmatrix}\begin{pmatrix}1\\ 0\end{pmatrix}
    =\frac{1}{\sqrt2}\begin{pmatrix}1\\ 1\end{pmatrix}=\ket{+},\\[0.6ex]
  H\ket{1}&=\frac{1}{\sqrt2}\begin{pmatrix}1 & 1\\ 1 & -1\end{pmatrix}\begin{pmatrix}0\\ 1\end{pmatrix}
    =\frac{1}{\sqrt2}\begin{pmatrix}1\\ -1\end{pmatrix}=\ket{-},\\[0.6ex]
  H\ket{+}&=\frac{1}{2}\begin{pmatrix}1 & 1\\ 1 & -1\end{pmatrix}\begin{pmatrix}1\\ 1\end{pmatrix}
    =\frac{1}{2}\begin{pmatrix}2\\ 0\end{pmatrix}=\ket{0},\qquad
  H\ket{-}=\frac{1}{2}\begin{pmatrix}0\\ 2\end{pmatrix}=\ket{1}.
\end{align*}
Comme $H\ket{+}=\ket{0}$ et $H\ket{-}=\ket{1}$, on a $H^2=\mathrm{Id}$ ($H$ est son propre
inverse). Sur la sphère de Bloch, $H$ est une rotation de $\pi$ autour de l'axe
$(x+z)/\sqrt2$ : elle échange l'axe $x$ et l'axe $z$.

\begin{remarque}[À quoi sert $H$ ?]
\begin{itemize}[nosep]
  \item Elle \emph{crée une superposition} à partir d'un état de base :
        $H\ket{0}=\ket{+}$ donne $P(0)=P(1)=\frac12$. C'est la première étape de
        la plupart des algorithmes quantiques.
  \item Elle \emph{change de base} : elle envoie la base $Z$ sur la base $X$ et
        inversement. Mesurer dans la base $X$ revient à appliquer $H$ puis à mesurer
        dans la base $Z$ (figure~\ref{fig:mesurex}).
  \item Elle est réversible ($H^2=\mathrm{Id}$).
\end{itemize}
\end{remarque}

% ---------------------------------------------------------------------
\subsection{Portes de phase}\label{sec:phase}
% ---------------------------------------------------------------------

\begin{definition}{Portes de phase $P_\theta$, $S$ et $T$}
\[
  P_\theta=\begin{pmatrix}1 & 0\\ 0 & \ee^{\ii\theta}\end{pmatrix}
  =\ket{0}\bra{0}+\ee^{\ii\theta}\ket{1}\bra{1},
  \qquad
  S=P_{\pi/2}=\begin{pmatrix}1 & 0\\ 0 & \ii\end{pmatrix},
  \qquad
  T=P_{\pi/4}=\begin{pmatrix}1 & 0\\ 0 & \dfrac{1+\ii}{\sqrt2}\end{pmatrix}.
\]
\end{definition}

Ici $\ee^{\ii\pi/2}=\ii$ et $\ee^{\ii\pi/4}=\cos\frac\pi4+\ii\sin\frac\pi4=\frac{1+\ii}{\sqrt2}$.
$P_\theta$ laisse $\ket{0}$ inchangé et multiplie $\ket{1}$ par $\ee^{\ii\theta}$ :
\[
  P_\theta\big(\alpha\ket{0}+\beta\ket{1}\big)=\alpha\ket{0}+\ee^{\ii\theta}\beta\ket{1}.
\]
Dans la paramétrisation de Bloch, $\varphi\to\varphi+\theta$ : $P_\theta$ est une rotation
d'angle $\theta$ autour de l'axe $z$. On retrouve $P_\pi=Z$, et $S^2=Z$, $T^2=S$.

\paragraph{Exemples.}
\begin{align*}
  S\ket{+}&=\tfrac{1}{\sqrt2}\big(\ket{0}+\ii\ket{1}\big)=\ket{{+}\ii},\\
  S\ket{{+}\ii}&=\tfrac{1}{\sqrt2}\big(\ket{0}+\ii\cdot\ii\ket{1}\big)=\tfrac{1}{\sqrt2}\big(\ket{0}-\ket{1}\big)=\ket{-},\\
  T\ket{+}&=\tfrac{1}{\sqrt2}\begin{pmatrix}1 & 0\\ 0 & \frac{1+\ii}{\sqrt2}\end{pmatrix}\begin{pmatrix}1\\ 1\end{pmatrix}
    =\tfrac{1}{\sqrt2}\begin{pmatrix}1\\ \frac{1+\ii}{\sqrt2}\end{pmatrix}
    =\tfrac{1}{\sqrt2}\big(\ket{0}+\ee^{\ii\pi/4}\ket{1}\big)\quad\big(\text{angle polaire }\tfrac\pi2,\ \text{azimut }\tfrac\pi4\big).
\end{align*}
$S$ tourne donc $\ket{+}$ d'un quart de tour autour de $z$ (de l'axe $x$ à l'axe $y$),
puis d'un second quart de tour jusqu'à $\ket{-}$ ; $T$ fait un huitième de tour.

\begin{remarque}[Une porte de phase est une évolution libre]
Pour $\Hop=\frac{\hbar\omega}{2}Z$, $U(t)=\ee^{-\ii\Hop t/\hbar}=\mathrm{diag}\big(\ee^{-\ii\omega t/2},\ee^{\ii\omega t/2}\big)
=\ee^{-\ii\omega t/2}\,P_{\omega t}$. Au facteur global près, laisser évoluer le qubit pendant la
durée $t$ applique la porte de phase $P_{\omega t}$ : une rotation d'angle $\omega t$ autour de
l'axe $z$.
\end{remarque}

% ---------------------------------------------------------------------
\subsection{Portes et sphère de Bloch}\label{sec:portes-bloch}
% ---------------------------------------------------------------------

Toute porte à un qubit est une rotation de la sphère de Bloch. La figure~\ref{fig:portes} montre
les axes des portes usuelles et le tableau qui suit les résume.

\begin{figure}[!ht]
\centering
\begin{tikzpicture}[x={(-1.63cm,-0.56cm)}, y={(1.63cm,-0.56cm)}, z={(0cm,2.16cm)}, >=Stealth, scale=0.9]
  % sphère et grands cercles
  \shade[ball color=insagris!15, opacity=0.55] (0,0,0) circle[radius=2.3cm];
  \begin{scope}[insagris!75, thin]
    \draw[domain=-45:135, samples=60, variable=\t] plot ({cos(\t)},{sin(\t)},0);
    \draw[domain=135:315, samples=60, variable=\t, dashed] plot ({cos(\t)},{sin(\t)},0);
    \draw[domain=-62.7:117.3, samples=60, variable=\t] plot ({cos(\t)},0,{sin(\t)});
    \draw[domain=117.3:297.3, samples=60, variable=\t, dashed] plot ({cos(\t)},0,{sin(\t)});
    \draw[domain=-62.7:117.3, samples=60, variable=\t] plot (0,{cos(\t)},{sin(\t)});
    \draw[domain=117.3:297.3, samples=60, variable=\t, dashed] plot (0,{cos(\t)},{sin(\t)});
  \end{scope}
  % axes x, y, z (portes X, Y, Z)
  \draw[insagris, dashed] (0,0,0) -- (-1,0,0);
  \draw[insagris, dashed] (0,0,0) -- (0,-1,0);
  \draw[insagris, dashed] (0,0,0) -- (0,0,-1);
  \draw[insagris] (-1,0,0) -- (-1.3,0,0);
  \draw[insagris] (0,-1,0) -- (0,-1.3,0);
  \draw[insagris] (0,0,-1) -- (0,0,-1.3);
  \draw[insagris, -{Stealth[length=2.2mm]}] (0,0,0) -- (1.3,0,0);
  \draw[insagris, -{Stealth[length=2.2mm]}] (0,0,0) -- (0,1.3,0);
  \draw[insagris, -{Stealth[length=2.2mm]}] (0,0,0) -- (0,0,1.3);
  \node[anchor=east]  at (1.33,0,0) {$x$};
  \node[anchor=west]  at (0,1.33,0) {$y$};
  \node[anchor=south] at (0,0,1.33) {$z$};
  % axe de la porte H : (x+z)/sqrt(2)
  \draw[insarouge, dashed] (0,0,0) -- (-0.707,0,-0.707);
  \draw[insarouge, very thick, -{Stealth[length=2.6mm]}] (0,0,0) -- (0.98,0,0.98);
  \node[insarouge, anchor=south east] at (0.98,0,0.98) {axe de $H$};
  % rotation autour de z (portes de phase)
  \draw[insarouge, thick, -{Stealth[length=2.2mm]}, domain=15:105, samples=30, variable=\t]
    plot ({1.15*cos(\t)},{1.15*sin(\t)},0);
  \node[insarouge, anchor=north west, font=\small] at ({1.15*cos(60)},{1.15*sin(60)},0) {$P_\theta$};
\end{tikzpicture}
\caption{Portes et rotations sur la sphère de Bloch. Les portes $X$, $Y$, $Z$ sont des
rotations de $\pi$ autour des axes $x$, $y$, $z$ ; $H$ est une rotation de $\pi$ autour de
l'axe $(x+z)/\sqrt2$ (rouge) ; $P_\theta$ (donc $Z$, $S$, $T$) tourne d'un angle $\theta$
autour de $z$.}
\label{fig:portes}
\end{figure}

\begin{center}
\renewcommand{\arraystretch}{1}
\begin{tabular}{@{}lccc@{}}
\toprule
\textbf{Porte} & \textbf{Matrice} & \textbf{Effet sur la base} & \textbf{Rotation de Bloch} \\
\midrule
$X$ & $\begin{pmatrix}0 & 1\\ 1 & 0\end{pmatrix}$ & $\ket{0}\leftrightarrow\ket{1}$ & $\pi$ autour de $x$ \\[2.2ex]
$Y$ & $\begin{pmatrix}0 & -\ii\\ \ii & 0\end{pmatrix}$ & $\ket{0}\to\ii\ket{1},\ \ket{1}\to-\ii\ket{0}$ & $\pi$ autour de $y$ \\[2.2ex]
$Z$ & $\begin{pmatrix}1 & 0\\ 0 & -1\end{pmatrix}$ & $\ket{1}\to-\ket{1},\ \ket{\pm}\leftrightarrow\ket{\mp}$ & $\pi$ autour de $z$ \\[2.2ex]
$H$ & $\dfrac{1}{\sqrt2}\begin{pmatrix}1 & 1\\ 1 & -1\end{pmatrix}$ & $\ket{0}\leftrightarrow\ket{+},\ \ket{1}\leftrightarrow\ket{-}$ & $\pi$ autour de $\frac{x+z}{\sqrt2}$ \\[2.2ex]
$S$ & $\begin{pmatrix}1 & 0\\ 0 & \ii\end{pmatrix}$ & $\ket{1}\to\ii\ket{1}$ & $\frac\pi2$ autour de $z$ \\[2.2ex]
$T$ & $\begin{pmatrix}1 & 0\\ 0 & \frac{1+\ii}{\sqrt2}\end{pmatrix}$ & $\ket{1}\to\ee^{\ii\pi/4}\ket{1}$ & $\frac\pi4$ autour de $z$ \\
\bottomrule
\end{tabular}
\end{center}

\begin{figure}[!ht]
\centering
\begin{tikzpicture}[x={(-1.63cm,-0.56cm)}, y={(1.63cm,-0.56cm)}, z={(0cm,2.16cm)}, >=Stealth, scale=0.9]
  % sphère et grands cercles
  \shade[ball color=insagris!15, opacity=0.55] (0,0,0) circle[radius=2.3cm];
  \begin{scope}[insagris!75, thin]
    \draw[domain=-45:135, samples=60, variable=\t] plot ({cos(\t)},{sin(\t)},0);
    \draw[domain=135:315, samples=60, variable=\t, dashed] plot ({cos(\t)},{sin(\t)},0);
    \draw[domain=-62.7:117.3, samples=60, variable=\t] plot ({cos(\t)},0,{sin(\t)});
    \draw[domain=117.3:297.3, samples=60, variable=\t, dashed] plot ({cos(\t)},0,{sin(\t)});
    \draw[domain=-62.7:117.3, samples=60, variable=\t] plot (0,{cos(\t)},{sin(\t)});
    \draw[domain=117.3:297.3, samples=60, variable=\t, dashed] plot (0,{cos(\t)},{sin(\t)});
  \end{scope}
  % axes x, y, z (portes X, Y, Z)
  \draw[insagris, dashed] (0,0,0) -- (-1,0,0);
  \draw[insagris, dashed] (0,0,0) -- (0,-1,0);
  \draw[insagris, dashed] (0,0,0) -- (0,0,-1);
  \draw[insagris] (-1,0,0) -- (-1.3,0,0);
  \draw[insagris] (0,-1,0) -- (0,-1.3,0);
  \draw[insagris] (0,0,-1) -- (0,0,-1.3);
  \draw[insagris, -{Stealth[length=2.2mm]}] (0,0,0) -- (1.3,0,0);
  \draw[insagris, -{Stealth[length=2.2mm]}] (0,0,0) -- (0,1.3,0);
  \draw[insagris, -{Stealth[length=2.2mm]}] (0,0,0) -- (0,0,1.3);
  % trajet : |0> --H--> |+> --S--> |+i>
  \draw[insarouge, very thick, -{Stealth[length=2.6mm]}, domain=3:87, samples=30, variable=\t]
    plot ({sin(\t)},0,{cos(\t)});
  \draw[insarouge, very thick, -{Stealth[length=2.6mm]}, domain=3:87, samples=30, variable=\t]
    plot ({cos(\t)},{sin(\t)},0);
  \foreach \ax/\ay/\az in {0/0/1, 1/0/0, 0/1/0} \fill[insarouge] (\ax,\ay,\az) circle[radius=2pt];
  \node[insarouge, anchor=south] at (0,0,1.33) {$\ket{0}$};
  \node[insarouge, anchor=east] at (1.33,0,0) {$\ket{+}$};
  \node[insarouge, anchor=west] at (0,1.33,0) {$\ket{{+}\ii}$};
  \node[insarouge, font=\small, anchor=south west] at ({sin(45)},0,{cos(45)}) {$H$};
  \node[insarouge, font=\small, anchor=north west] at ({cos(45)},{sin(45)},0) {$S$};
\end{tikzpicture}
\caption{Trajet d'un état sous le circuit $H$ puis $S$ (flèches schématiques). $H$ envoie
$\ket{0}$ sur $\ket{+}$ ; $S$, rotation d'un quart de tour autour de $z$, envoie $\ket{+}$ sur
$\ket{{+}\ii}$. Les trois états sont alignés sur les axes $z$, $x$ et $y$.}
\label{fig:trajet}
\end{figure}

% ---------------------------------------------------------------------
\subsection{Circuits quantiques}\label{sec:circuits}
% ---------------------------------------------------------------------

Un \emph{circuit} représente un calcul : chaque qubit est un fil horizontal, le temps va de gauche
à droite, et chaque porte est un symbole posé sur les fils qu'elle concerne. La
figure~\ref{fig:symboles} donne les symboles utilisés dans ce document.

\begin{figure}[!ht]
\centering
\begin{adjustbox}{max width=\linewidth}
\begin{tikzpicture}[>=Stealth, font=\small,
    nom/.style={font=\footnotesize, align=center, anchor=north}]
  % (a) porte à un qubit
  \draw (0,0) -- (1.8,0);
  \node[porte] at (0.9,0) {$U$};
  \node[nom] at (0.9,-1.15) {porte $U$\\ à un qubit};
  % (b) CNOT
  \begin{scope}[xshift=3.2cm]
    \draw (0,0.6) -- (1.8,0.6);
    \draw (0,-0.6) -- (1.8,-0.6);
    \draw[insarouge, thick] (0.9,0.6) -- (0.9,-0.6);
    \pic at (0.9,0.6) {ctl};
    \pic at (0.9,-0.6) {cible};
    \node[nom] at (0.9,-1.15) {CNOT\\ (contrôle $\bullet$, cible $\oplus$)};
  \end{scope}
  % (c) SWAP
  \begin{scope}[xshift=6.4cm]
    \draw (0,0.6) -- (1.8,0.6);
    \draw (0,-0.6) -- (1.8,-0.6);
    \draw[insarouge, thick] (0.9,0.6) -- (0.9,-0.6);
    \pic at (0.9,0.6) {croix};
    \pic at (0.9,-0.6) {croix};
    \node[nom] at (0.9,-1.15) {SWAP\\ (échange)};
  \end{scope}
  % (d) contrôlé-U
  \begin{scope}[xshift=9.6cm]
    \draw (0,0.6) -- (1.8,0.6);
    \draw (0,-0.6) -- (1.8,-0.6);
    \draw[insarouge, thick] (0.9,-0.6) -- (0.9,0.3);
    \pic at (0.9,-0.6) {ctl};
    \node[porte] at (0.9,0.6) {$U$};
    \node[nom] at (0.9,-1.15) {contrôlé-$U$\\ (contrôle en bas)};
  \end{scope}
  % (e) mesure
  \begin{scope}[xshift=12.8cm]
    \draw (0,0) -- (0.6,0);
    \pic at (0.9,0) {mesure};
    \draw (1.2,0.03) -- (2.0,0.03);
    \draw (1.2,-0.03) -- (2.0,-0.03);
    \node[nom] at (1.0,-1.15) {mesure\\ (fil double : bit)};
  \end{scope}
\end{tikzpicture}
\end{adjustbox}
\caption{Symboles des circuits quantiques. Le point noir $\bullet$ marque un qubit de contrôle,
le symbole $\oplus$ une cible qui subit $X$, les deux croix un échange. Après une mesure, le fil
double transporte un bit classique.}
\label{fig:symboles}
\end{figure}

\begin{remarque}[Conventions de lecture]
\begin{itemize}[nosep]
  \item Le temps va de gauche à droite. Deux portes successives $U_1$ puis $U_2$ se multiplient
        dans l'ordre inverse : la matrice du circuit est $U_2U_1$, car $U_1$ agit d'abord sur le ket.
  \item Avec plusieurs fils, le fil du \emph{haut} porte le \emph{dernier} chiffre du ket et le
        fil du \emph{bas} le premier : si $A$ est le premier qubit et $B$ le second,
        $\ket{a,b}=\ket{a}_A\otimes\ket{b}_B$ et $B$ est dessiné au-dessus de $A$.
  \item Deux portes placées dans la même colonne, sur deux fils, agissent en parallèle : c'est un
        produit tensoriel $U_A\otimes U_B$ (section~\ref{sec:op2}).
\end{itemize}
\end{remarque}

\paragraph{Exemple.} Le circuit de la figure~\ref{fig:circuit-hs} applique $H$ puis $S$ à
$\ket{0}$. Sa matrice est $SH$ :
\[
  SH=\begin{pmatrix}1&0\\0&\ii\end{pmatrix}\frac{1}{\sqrt2}\begin{pmatrix}1&1\\1&-1\end{pmatrix}
  =\frac{1}{\sqrt2}\begin{pmatrix}1&1\\ \ii&-\ii\end{pmatrix},
  \qquad
  SH\ket{0}=\frac{1}{\sqrt2}\begin{pmatrix}1\\ \ii\end{pmatrix}=\ket{{+}\ii}.
\]
Le circuit prépare donc l'état $\ket{{+}\ii}$ de la base $Y$ à partir de $\ket{0}$
(figure~\ref{fig:trajet}). L'ordre compte : $HS\ket{0}=H\ket{0}=\ket{+}$, car $S\ket{0}=\ket{0}$.

\begin{figure}[!ht]
\centering
\begin{tikzpicture}[>=Stealth, font=\small]
  \draw (0,0) -- (6.6,0);
  \node[anchor=east] at (0,0) {$\ket{0}$};
  \node[porte] at (2.0,0) {$H$};
  \node[porte] at (4.3,0) {$S$};
  \node[anchor=west] at (6.6,0) {$\ket{{+}\ii}$};
  \foreach \x in {0.7, 3.15, 5.5} \draw[dashed, insagris!70] (\x,0.6) -- (\x,-0.95);
  \node[font=\footnotesize, anchor=north] at (0.7,-0.95) {$\ket{0}$};
  \node[font=\footnotesize, anchor=north] at (3.15,-0.95) {$\ket{+}$};
  \node[font=\footnotesize, anchor=north] at (5.5,-0.95) {$\ket{{+}\ii}$};
\end{tikzpicture}
\caption{Circuit $H$ puis $S$ appliqué à $\ket{0}$. Les états intermédiaires sont indiqués sous
les traits pointillés ; la matrice du circuit est $SH$.}
\label{fig:circuit-hs}
\end{figure}

\paragraph{Mesure.} Le symbole de mesure (compteur) donne le résultat $0$ ou $1$ dans la base $Z$.
Mesurer dans la base $X$ revient à appliquer $H$ puis à mesurer dans la base $Z$
(figure~\ref{fig:mesurex}) : comme $H\ket{\pm}=\ket{0}$ ou $\ket{1}$, le résultat $0$ correspond à
$\ket{+}$ et le résultat $1$ à $\ket{-}$, avec $P(\pm)=|\braket{\pm}{\Psi}|^2$.

\begin{figure}[!ht]
\centering
\begin{tikzpicture}[>=Stealth, font=\small]
  % (a) base Z
  \draw (0,0) -- (1.2,0);
  \node[anchor=east] at (0,0) {$\ket{\Psi}$};
  \pic at (1.5,0) {mesure};
  \draw (1.8,0.03) -- (2.6,0.03);
  \draw (1.8,-0.03) -- (2.6,-0.03);
  \node[anchor=west] at (2.7,0) {$0$ ou $1$};
  \node[font=\footnotesize] at (1.5,-0.95) {(a) mesure en base $Z$};
  % (b) base X
  \begin{scope}[xshift=7.5cm]
    \draw (0,0) -- (2.6,0);
    \node[anchor=east] at (0,0) {$\ket{\Psi}$};
    \node[porte] at (1.0,0) {$H$};
    \pic at (2.3,0) {mesure};
    \draw (2.6,0.03) -- (3.4,0.03);
    \draw (2.6,-0.03) -- (3.4,-0.03);
    \node[anchor=west] at (3.5,0) {$+$ ou $-$};
    \node[font=\footnotesize] at (1.9,-0.95) {(b) mesure en base $X$};
  \end{scope}
\end{tikzpicture}
\caption{(a) Mesure dans la base $Z$. (b) Mesure dans la base $X$ : porte de Hadamard, puis
mesure dans la base $Z$ ; les résultats $0$ et $1$ se lisent $+$ et $-$.}
\label{fig:mesurex}
\end{figure}

% ---------------------------------------------------------------------
\subsection{Portes locales sur deux qubits}\label{sec:op2}
% ---------------------------------------------------------------------

Une porte appliquée à un seul qubit d'un système de deux s'écrit avec un produit tensoriel
(section~\ref{sec:tenseur}) : $U\otimes I$ agit sur $A$ (premier facteur, fil du bas), $I\otimes U$
agit sur $B$, et $U_A\otimes U_B$ agit sur les deux en parallèle (figure~\ref{fig:local}).

\begin{figure}[!ht]
\centering
\begin{tikzpicture}[>=Stealth, font=\small]
  \foreach \xs/\haut/\bas/\tit in {0/{}/{$X$}/{(a) $X\otimes I$},
                                   5.4/{$X$}/{}/{(b) $I\otimes X$},
                                   10.8/{$H$}/{$H$}/{(c) $H\otimes H$}} {
    \begin{scope}[xshift=\xs cm]
      \draw (0,0.75) -- (2.6,0.75);
      \draw (0,0) -- (2.6,0);
      \node[anchor=east] at (0,0.75) {$B$};
      \node[anchor=east] at (0,0) {$A$};
      \ifx\haut\empty\else \node[porte] at (1.3,0.75) {\haut}; \fi
      \ifx\bas\empty\else \node[porte] at (1.3,0) {\bas}; \fi
      \node[font=\footnotesize] at (1.3,-0.85) {\tit};
    \end{scope}
  }
\end{tikzpicture}
\caption{Portes en parallèle sur deux qubits. Le qubit $A$, premier facteur du produit tensoriel,
est sur le fil du bas. (a) $X$ sur $A$ seulement. (b) $X$ sur $B$ seulement. (c) $H$ sur les deux.}
\label{fig:local}
\end{figure}

\subsubsection{Portes agissant sur un seul des deux qubits}

On note $I=\mathrm{Id}$ ($2\times2$). $X\otimes I$ agit sur $A$, $I\otimes X$ agit sur $B$ :
\begin{align*}
  I\otimes X&=\begin{pmatrix}1\cdot X & 0\cdot X\\ 0\cdot X & 1\cdot X\end{pmatrix}
    =\begin{pmatrix}0 & 1 & 0 & 0\\ 1 & 0 & 0 & 0\\ 0 & 0 & 0 & 1\\ 0 & 0 & 1 & 0\end{pmatrix},\\[1ex]
  X\otimes I&=\begin{pmatrix}0\cdot I & 1\cdot I\\ 1\cdot I & 0\cdot I\end{pmatrix}
    =\begin{pmatrix}0 & 0 & 1 & 0\\ 0 & 0 & 0 & 1\\ 1 & 0 & 0 & 0\\ 0 & 1 & 0 & 0\end{pmatrix}.
\end{align*}
Sur $\ket{00}=(1,0,0,0)^{\mathsf T}$ (première colonne de la matrice) :
$(I\otimes X)\ket{00}=(0,1,0,0)^{\mathsf T}=\ket{01}$ et
$(X\otimes I)\ket{00}=(0,0,1,0)^{\mathsf T}=\ket{10}$. Avec la règle :
$(I\otimes X)\ket{00}=I\ket{0}\otimes X\ket{0}=\ket{0}\otimes\ket{1}=\ket{01}$.

Pour la porte de Hadamard, avec $H=\frac{1}{\sqrt2}\begin{pmatrix}1 & 1\\ 1 & -1\end{pmatrix}$ :
\[
  H\otimes I=\frac{1}{\sqrt2}\begin{pmatrix}1\cdot I & 1\cdot I\\ 1\cdot I & -1\cdot I\end{pmatrix}
  =\frac{1}{\sqrt2}\begin{pmatrix}1 & 0 & 1 & 0\\ 0 & 1 & 0 & 1\\ 1 & 0 & -1 & 0\\ 0 & 1 & 0 & -1\end{pmatrix}.
\]
Ses colonnes donnent l'image des quatre états de base :
\begin{align*}
  (H\otimes I)\ket{00}&=\tfrac{1}{\sqrt2}(1,0,1,0)^{\mathsf T}=\ket{+}\otimes\ket{0},&
  (H\otimes I)\ket{01}&=\tfrac{1}{\sqrt2}(0,1,0,1)^{\mathsf T}=\ket{+}\otimes\ket{1},\\
  (H\otimes I)\ket{10}&=\tfrac{1}{\sqrt2}(1,0,-1,0)^{\mathsf T}=\ket{-}\otimes\ket{0},&
  (H\otimes I)\ket{11}&=\tfrac{1}{\sqrt2}(0,1,0,-1)^{\mathsf T}=\ket{-}\otimes\ket{1}.
\end{align*}

\subsubsection{Hadamard sur les deux qubits}

\[
  H\otimes H=\frac12\begin{pmatrix}1\cdot H' & 1\cdot H'\\ 1\cdot H' & -1\cdot H'\end{pmatrix},
  \quad H'=\begin{pmatrix}1 & 1\\ 1 & -1\end{pmatrix},
  \qquad
  H\otimes H=\frac12\begin{pmatrix}1 & 1 & 1 & 1\\ 1 & -1 & 1 & -1\\ 1 & 1 & -1 & -1\\ 1 & -1 & -1 & 1\end{pmatrix}.
\]
La première colonne donne $(H\otimes H)\ket{00}=\frac12(1,1,1,1)^{\mathsf T}=\ket{+}\otimes\ket{+}$
(cohérent avec $H\ket{0}\otimes H\ket{0}$) : les quatre résultats $00,01,10,11$ sont
équiprobables, chacun avec la probabilité $\frac14$. Sur $n$ qubits, $H^{\otimes n}\ket{0\ldots0}$
donne une superposition uniforme de $2^n$ états.

% ---------------------------------------------------------------------
\subsection{Porte CNOT}\label{sec:cnot}
% ---------------------------------------------------------------------

\begin{definition}{Porte CNOT (NOT contrôlé)}
La porte CNOT agit sur deux qubits : le qubit $B$ est le \emph{contrôle}, le qubit $A$ est la
\emph{cible}. La cible est inversée (porte $X$) si le contrôle vaut $1$ et reste inchangée
s'il vaut $0$. Avec $\ket{a,b}=\ket{a}_A\otimes\ket{b}_B$ :
\[
  \mathrm{CNOT}\,\ket{a,b}=\ket{a\oplus b,\,b},
  \qquad
  \mathrm{CNOT}=I\otimes\ket{0}\bra{0}+X\otimes\ket{1}\bra{1},
\]
où $\oplus$ est l'addition modulo $2$ ($0\oplus0=1\oplus1=0$ et $0\oplus1=1\oplus0=1$).
\end{definition}

Sur le circuit de la figure~\ref{fig:cnot}, le point noir posé sur le fil de $B$ (en haut) marque
le contrôle et le symbole $\oplus$ posé sur le fil de $A$ (en bas) marque la cible.

\begin{figure}[!ht]
\centering
\begin{tikzpicture}[>=Stealth, font=\small]
  \draw (0,1.2) -- (6,1.2);
  \draw (0,0) -- (6,0);
  \node[anchor=east] at (0,1.2) {$\ket{b}_B$};
  \node[anchor=east] at (0,0) {$\ket{a}_A$};
  \node[anchor=west] at (6,1.2) {$\ket{b}_B$};
  \node[anchor=west] at (6,0) {$\ket{a\oplus b}_A$};
  \draw[insarouge, thick] (3,1.2) -- (3,0);
  \pic at (3,1.2) {ctl};
  \pic at (3,0) {cible};
  \node[anchor=south, text=insagris] at (3,1.5) {qubit de contrôle};
  \node[anchor=north, text=insagris] at (3,-0.5) {qubit cible};
\end{tikzpicture}
\caption{Porte CNOT : le fil du haut ($B$) est le contrôle, le fil du bas ($A$) est la cible. Le
ket s'écrit $\ket{a,b}$ : le fil du haut donne le dernier chiffre.}
\label{fig:cnot}
\end{figure}

\subsubsection{Table de vérité}

\begin{center}
\renewcommand{\arraystretch}{1.3}
\begin{tabular}{@{}cccc@{}}
\toprule
Entrée $\ket{a,b}$ & Contrôle $b$ & Sortie $\ket{a\oplus b,b}$ & Effet sur la cible $A$ \\
\midrule
$\ket{00}$ & $0$ & $\ket{00}$ & aucun \\
$\ket{01}$ & $1$ & $\ket{11}$ & $0\to1$ \\
$\ket{10}$ & $0$ & $\ket{10}$ & aucun \\
$\ket{11}$ & $1$ & $\ket{01}$ & $1\to0$ \\
\bottomrule
\end{tabular}
\end{center}

\subsubsection{Matrice du CNOT}

La $j$-ième colonne d'une matrice $4\times4$ est l'image du $j$-ième vecteur de base
($\ket{00},\ket{01},\ket{10},\ket{11}$, dans cet ordre). La table donne :
\begin{align*}
  \mathrm{CNOT}\ket{00}&=\ket{00}=\begin{pmatrix}1\\0\\0\\0\end{pmatrix},&
  \mathrm{CNOT}\ket{01}&=\ket{11}=\begin{pmatrix}0\\0\\0\\1\end{pmatrix},\\[1.5ex]
  \mathrm{CNOT}\ket{10}&=\ket{10}=\begin{pmatrix}0\\0\\1\\0\end{pmatrix},&
  \mathrm{CNOT}\ket{11}&=\ket{01}=\begin{pmatrix}0\\1\\0\\0\end{pmatrix},
\end{align*}
d'où, en juxtaposant ces quatre colonnes,
\[
  \mathrm{CNOT}=\begin{pmatrix}
    1 & 0 & 0 & 0\\ 0 & 0 & 0 & 1\\ 0 & 0 & 1 & 0\\ 0 & 1 & 0 & 0
  \end{pmatrix}.
\]

On retrouve cette matrice avec la formule $I\otimes\ket{0}\bra{0}+X\otimes\ket{1}\bra{1}$ et
la règle $A\otimes B=(A_{ij}B)$ (section~\ref{sec:tenseur}), où
$\ket{0}\bra{0}=\begin{pmatrix}1&0\\0&0\end{pmatrix}$ et
$\ket{1}\bra{1}=\begin{pmatrix}0&0\\0&1\end{pmatrix}$ :
\[
  I\otimes\ket{0}\bra{0}=\begin{pmatrix}1&0&0&0\\0&0&0&0\\0&0&1&0\\0&0&0&0\end{pmatrix},
  \qquad
  X\otimes\ket{1}\bra{1}=\begin{pmatrix}0&0&0&0\\0&0&0&1\\0&0&0&0\\0&1&0&0\end{pmatrix},
\]
et la somme de ces deux matrices est la matrice du CNOT.

\subsubsection{Action sur un état quelconque}

Pour $\ket{\Psi}=c_{00}\ket{00}+c_{01}\ket{01}+c_{10}\ket{10}+c_{11}\ket{11}$ :
\[
  \mathrm{CNOT}\begin{pmatrix}c_{00}\\c_{01}\\c_{10}\\c_{11}\end{pmatrix}
  =\begin{pmatrix}
    1\cdot c_{00}+0\cdot c_{01}+0\cdot c_{10}+0\cdot c_{11}\\
    0\cdot c_{00}+0\cdot c_{01}+0\cdot c_{10}+1\cdot c_{11}\\
    0\cdot c_{00}+0\cdot c_{01}+1\cdot c_{10}+0\cdot c_{11}\\
    0\cdot c_{00}+1\cdot c_{01}+0\cdot c_{10}+0\cdot c_{11}
  \end{pmatrix}
  =\begin{pmatrix}c_{00}\\c_{11}\\c_{10}\\c_{01}\end{pmatrix}.
\]
Le CNOT échange les amplitudes de $\ket{01}$ et de $\ket{11}$, c'est-à-dire des deux états
dont le contrôle vaut $1$, et laisse les autres.

\paragraph{Exemple : contrôle en superposition.} Soit $\ket{0}_A\otimes\ket{+}_B
=\frac{1}{\sqrt2}\big(\ket{00}+\ket{01}\big)=\frac{1}{\sqrt2}(1,1,0,0)^{\mathsf T}$,
un état séparable. Alors
\[
  \mathrm{CNOT}\,\frac{1}{\sqrt2}\begin{pmatrix}1\\1\\0\\0\end{pmatrix}
  =\frac{1}{\sqrt2}\begin{pmatrix}1\\0\\0\\1\end{pmatrix}
  =\frac{1}{\sqrt2}\big(\ket{00}+\ket{11}\big).
\]
Le résultat est intriqué ($c_{00}c_{11}-c_{01}c_{10}=\frac12\neq0$ ; section~\ref{sec:separable}) :
un CNOT dont le contrôle est en superposition crée de l'intrication. C'est l'état de Bell
$\ket{\Phi^+}$ du chapitre~\ref{chap:bell}.

\paragraph{Réversibilité.} Deux CNOT successifs redonnent l'état initial :
$\ket{a,b}\to\ket{a\oplus b,b}\to\ket{a\oplus b\oplus b,b}=\ket{a,b}$, donc
$\mathrm{CNOT}^2=\mathrm{Id}$. La matrice est réelle et symétrique, donc
$\mathrm{CNOT}^\dagger\mathrm{CNOT}=\mathrm{CNOT}^2=\mathrm{Id}$ : elle est bien unitaire.

\begin{remarque}[CNOT et copie d'un bit]
Avec la cible dans l'état $\ket{0}$, $\mathrm{CNOT}\ket{0,b}=\ket{b,b}$ pour $b=0$ et $b=1$ : le CNOT
copie le bit classique $b$. En revanche $\mathrm{CNOT}\big(\ket{0}\otimes\ket{+}\big)$ est l'état
ci-dessus, et non $\ket{+}\otimes\ket{+}$ : un état quantique inconnu ne peut pas être copié
(théorème de non-clonage, section~\ref{sec:teleportation}).
\end{remarque}

\begin{remarque}[Contrôle en dernier ou en premier]
La matrice dépend de l'ordre d'écriture du ket. Notons $\mathrm{CNOT}_{21}$ le CNOT dont le
contrôle est le \emph{second} chiffre du ket : c'est la matrice ci-dessus, adoptée dans ce
document (fil du haut). Beaucoup d'ouvrages mettent le contrôle en \emph{premier} ; le CNOT
correspondant est
\[
  \mathrm{CNOT}_{12}=\ket{0}\bra{0}\otimes I+\ket{1}\bra{1}\otimes X
  =\begin{pmatrix}1&0&0&0\\0&1&0&0\\0&0&0&1\\0&0&1&0\end{pmatrix}.
\]
Le même circuit a donc deux matrices, selon que le contrôle est écrit en dernier ou en premier.
Les deux sont reliées par le SWAP (section~\ref{sec:swap}). La section~\ref{sec:controlees} sur les
portes contrôlées met le contrôle en premier, avec la matrice $\mathrm{diag}(I,U)$.
\end{remarque}

% ---------------------------------------------------------------------
\subsection{Porte SWAP}\label{sec:swap}
% ---------------------------------------------------------------------

\begin{definition}{Porte SWAP (échange)}
La porte SWAP échange les états des deux qubits :
\[
  \mathrm{SWAP}\big(\ket{\psi}\otimes\ket{\phi}\big)=\ket{\phi}\otimes\ket{\psi}.
\]
Sur les états de base, $\mathrm{SWAP}\ket{a,b}=\ket{b,a}$ ; on prolonge par linéarité.
\end{definition}

Son symbole est formé de deux croix reliées par un trait vertical
(figure~\ref{fig:swap}).

\begin{figure}[!ht]
\centering
\begin{tikzpicture}[>=Stealth, font=\small]
  \draw (0,1.2) -- (6,1.2);
  \draw (0,0) -- (6,0);
  \node[anchor=east] at (0,1.2) {$\ket{a}$};
  \node[anchor=east] at (0,0) {$\ket{b}$};
  \node[anchor=west] at (6,1.2) {$\ket{b}$};
  \node[anchor=west] at (6,0) {$\ket{a}$};
  \draw[insarouge, thick] (3,1.2) -- (3,0);
  \pic at (3,1.2) {croix};
  \pic at (3,0) {croix};
\end{tikzpicture}
\caption{Porte SWAP : les deux qubits échangent leurs états.}
\label{fig:swap}
\end{figure}

\subsubsection{Table et matrice}

\begin{center}
\renewcommand{\arraystretch}{1.3}
\begin{tabular}{@{}cc@{}}
\toprule
Entrée & Sortie \\
\midrule
$\ket{00}$ & $\ket{00}$ \\
$\ket{01}$ & $\ket{10}$ \\
$\ket{10}$ & $\ket{01}$ \\
$\ket{11}$ & $\ket{11}$ \\
\bottomrule
\end{tabular}
\end{center}

Les images des quatre vecteurs de base sont les colonnes de la matrice :
\[
  \mathrm{SWAP}\ket{00}=\begin{pmatrix}1\\0\\0\\0\end{pmatrix},\quad
  \mathrm{SWAP}\ket{01}=\begin{pmatrix}0\\0\\1\\0\end{pmatrix},\quad
  \mathrm{SWAP}\ket{10}=\begin{pmatrix}0\\1\\0\\0\end{pmatrix},\quad
  \mathrm{SWAP}\ket{11}=\begin{pmatrix}0\\0\\0\\1\end{pmatrix},
\]
\[
  \mathrm{SWAP}=\begin{pmatrix}
    1 & 0 & 0 & 0\\ 0 & 0 & 1 & 0\\ 0 & 1 & 0 & 0\\ 0 & 0 & 0 & 1
  \end{pmatrix}.
\]
Appliquée à $\ket{01}=(0,1,0,0)^{\mathsf T}$, cette matrice sélectionne sa deuxième colonne :
\[
  \begin{pmatrix}
    1 & 0 & 0 & 0\\ 0 & 0 & 1 & 0\\ 0 & 1 & 0 & 0\\ 0 & 0 & 0 & 1
  \end{pmatrix}
  \begin{pmatrix}0\\1\\0\\0\end{pmatrix}
  =\begin{pmatrix}0\\0\\1\\0\end{pmatrix}=\ket{10}.
\]
Sur un état quelconque, $\mathrm{SWAP}\,(c_{00},c_{01},c_{10},c_{11})^{\mathsf T}
=(c_{00},c_{10},c_{01},c_{11})^{\mathsf T}$ : les amplitudes de $\ket{01}$ et $\ket{10}$
sont échangées.

\subsubsection{Accord avec la définition}

Pour deux états $\ket{\psi}=(a_0,a_1)^{\mathsf T}$ et $\ket{\phi}=(b_0,b_1)^{\mathsf T}$ :
\[
  \mathrm{SWAP}\,(\ket{\psi}\otimes\ket{\phi})
  =\mathrm{SWAP}\begin{pmatrix}a_0b_0\\a_0b_1\\a_1b_0\\a_1b_1\end{pmatrix}
  =\begin{pmatrix}a_0b_0\\a_1b_0\\a_0b_1\\a_1b_1\end{pmatrix}
  =\begin{pmatrix}b_0\\b_1\end{pmatrix}\otimes\begin{pmatrix}a_0\\a_1\end{pmatrix}
  =\ket{\phi}\otimes\ket{\psi},
\]
ce qui est bien la définition. La porte s'écrit aussi avec des projecteurs :
\[
  \mathrm{SWAP}=\sum_{a,b\in\{0,1\}}\ket{a}\bra{b}\otimes\ket{b}\bra{a}.
\]
En effet, pour $\ket{c}\otimes\ket{d}$ on trouve
$\sum_{a,b}\ket{a}\braket{b}{c}\otimes\ket{b}\braket{a}{d}$, et seul le terme $b=c$, $a=d$
est non nul : le résultat est $\ket{d}\otimes\ket{c}$.

\paragraph{Exemples.}
\begin{itemize}[nosep]
  \item Deux états différents : $\mathrm{SWAP}\big(\ket{0}\otimes\ket{+}\big)
        =\mathrm{SWAP}\,\frac{1}{\sqrt2}(1,1,0,0)^{\mathsf T}=\frac{1}{\sqrt2}(1,0,1,0)^{\mathsf T}
        =\ket{+}\otimes\ket{0}$.
  \item L'état $\ket{\psi}=\frac{1}{\sqrt2}\ket{00}+\frac{1}{\sqrt2}\ket{11}$ ne change pas :
        $\mathrm{SWAP}\ket{\psi}=\frac{1}{\sqrt2}\ket{00}+\frac{1}{\sqrt2}\ket{11}=\ket{\psi}$,
        car $\ket{00}$ et $\ket{11}$ sont invariants. De même
        $\frac{1}{\sqrt2}(\ket{00}-\ket{11})$ et $\frac{1}{\sqrt2}(\ket{01}+\ket{10})$ sont invariants,
        alors que $\mathrm{SWAP}\,\frac{1}{\sqrt2}(\ket{01}-\ket{10})=\frac{1}{\sqrt2}(\ket{10}-\ket{01})$
        est l'opposé de l'état de départ. Ce sont les quatre états de Bell (chapitre~\ref{chap:bell}).
\end{itemize}

\paragraph{Propriétés.} Échanger deux fois redonne l'état de départ :
$\mathrm{SWAP}^2=\mathrm{Id}$ ; la matrice est réelle et symétrique, donc
$\mathrm{SWAP}^\dagger\mathrm{SWAP}=\mathrm{SWAP}^2=\mathrm{Id}$ : elle est unitaire.

\subsubsection{Lien avec le CNOT}

\begin{resultat}{SWAP échange les rôles de contrôle et de cible}
\[
  \mathrm{CNOT}_{21}=\mathrm{SWAP}\;\mathrm{CNOT}_{12}\;\mathrm{SWAP}.
\]
\end{resultat}

\emph{Sur les états de base.}
$\mathrm{SWAP}\,\mathrm{CNOT}_{12}\,\mathrm{SWAP}\ket{a,b}
=\mathrm{SWAP}\,\mathrm{CNOT}_{12}\ket{b,a}
=\mathrm{SWAP}\ket{b,\,b\oplus a}=\ket{b\oplus a,\,b}=\mathrm{CNOT}_{21}\ket{a,b}$.

\emph{Par les matrices.} Multiplier à gauche par $\mathrm{SWAP}$ échange les lignes $2$ et
$3$, multiplier à droite l'échange les colonnes $2$ et $3$ :
\[
  \mathrm{SWAP}\,\mathrm{CNOT}_{12}
  =\begin{pmatrix}1&0&0&0\\0&0&0&1\\0&1&0&0\\0&0&1&0\end{pmatrix},
  \qquad
  \mathrm{SWAP}\,\mathrm{CNOT}_{12}\,\mathrm{SWAP}
  =\begin{pmatrix}1&0&0&0\\0&0&0&1\\0&0&1&0\\0&1&0&0\end{pmatrix}=\mathrm{CNOT}_{21}.
\]

\begin{remarque}[SWAP avec trois CNOT]
$\mathrm{SWAP}=\mathrm{CNOT}_{12}\,\mathrm{CNOT}_{21}\,\mathrm{CNOT}_{12}$. En effet, sur
$\ket{a,b}$ (on applique les portes de droite à gauche) :
\[
  \ket{a,b}\ \xrightarrow{\ \mathrm{CNOT}_{12}\ }\ \ket{a,\,a\oplus b}
  \ \xrightarrow{\ \mathrm{CNOT}_{21}\ }\ \ket{a\oplus(a\oplus b),\,a\oplus b}=\ket{b,\,a\oplus b}
  \ \xrightarrow{\ \mathrm{CNOT}_{12}\ }\ \ket{b,\,b\oplus(a\oplus b)}=\ket{b,a}.
\]
Les deux membres coïncident sur les états de base, donc sur tous les états par linéarité.
\end{remarque}

% ---------------------------------------------------------------------
\subsection{Portes contrôlées : contrôlé-\texorpdfstring{$U$}{U}, Fredkin et Toffoli}\label{sec:controlees}
% ---------------------------------------------------------------------

Le CNOT est le cas le plus simple d'une idée générale : appliquer une opération à un ou
plusieurs qubits \emph{cibles} seulement si un ou plusieurs qubits de \emph{contrôle} valent $1$.
Dans cette section, le contrôle est le \emph{premier} chiffre
du ket (remarque de la section~\ref{sec:cnot}). Les qubits sont notés $x$, $y$, $z$ (le qubit $x$ est le contrôle) et
$\ket{x,y,z}=\ket{x}\otimes\ket{y}\otimes\ket{z}$ ; d'après la convention de la section~\ref{sec:circuits}, le
contrôle est dessiné en bas (figure~\ref{fig:controlees}). Pour éviter la confusion avec les
matrices de Pauli, on écrit ici $\sigma_x=X$ pour la porte NOT.

\begin{figure}[!ht]
\centering
\begin{adjustbox}{max width=\linewidth}
\begin{tikzpicture}[>=Stealth, font=\small]
  % (a) contrôlé-U : contrôle x en bas, cible y en haut
  \draw (0,1.2) -- (3.4,1.2);
  \draw (0,0) -- (3.4,0);
  \node[anchor=east] at (0,1.2) {$y$};
  \node[anchor=east] at (0,0) {$x$};
  \draw[insarouge, thick] (1.7,0) -- (1.7,0.85);
  \pic at (1.7,0) {ctl};
  \node[porte] at (1.7,1.2) {$U$};
  \node[anchor=north, font=\footnotesize] at (1.7,-0.5) {(a) contrôlé-$U$};
  % (b) contrôlé-SWAP (Fredkin) : x contrôle, y et z échangés
  \draw (5,2.4) -- (8.4,2.4);
  \draw (5,1.2) -- (8.4,1.2);
  \draw (5,0) -- (8.4,0);
  \node[anchor=east] at (5,2.4) {$z$};
  \node[anchor=east] at (5,1.2) {$y$};
  \node[anchor=east] at (5,0) {$x$};
  \draw[insarouge, thick] (6.7,0) -- (6.7,2.4);
  \pic at (6.7,0) {ctl};
  \pic at (6.7,1.2) {croix};
  \pic at (6.7,2.4) {croix};
  \node[anchor=north, font=\footnotesize] at (6.7,-0.5) {(b) Fredkin (C-SWAP)};
  % (c) Toffoli (CCNOT) : x et y contrôlent, z cible
  \draw (10,2.4) -- (13.4,2.4);
  \draw (10,1.2) -- (13.4,1.2);
  \draw (10,0) -- (13.4,0);
  \node[anchor=east] at (10,2.4) {$z$};
  \node[anchor=east] at (10,1.2) {$y$};
  \node[anchor=east] at (10,0) {$x$};
  \draw[insarouge, thick] (11.7,0) -- (11.7,2.4);
  \pic at (11.7,0) {ctl};
  \pic at (11.7,1.2) {ctl};
  \pic at (11.7,2.4) {cible};
  \node[anchor=north, font=\footnotesize] at (11.7,-0.5) {(c) Toffoli (CCNOT)};
\end{tikzpicture}
\end{adjustbox}
\caption{Portes contrôlées. Le contrôle $x$ est le premier chiffre du ket, donc le fil du bas.
(a) Le qubit $y$ subit $U$ si $x=1$ ; pour $U=X$ on retrouve le CNOT. (b) Si $x=1$, les qubits
$y$ et $z$ sont échangés. (c) Si $x=y=1$, le qubit $z$ est inversé.}
\label{fig:controlees}
\end{figure}

\subsubsection{Contrôlé-\texorpdfstring{$U$}{U}}

\begin{definition}{Porte contrôlée-$U$}
Soit $U$ une matrice unitaire $2\times2$. Sur le système $(x,y)$, de contrôle $x$ et de cible $y$,
\[
  \mathrm{C}U=\ket{0}\bra{0}_x\otimes I_y+\ket{1}\bra{1}_x\otimes U_y .
\]
Le premier terme agit quand le contrôle est $\ket{0}$ : la cible ne change pas ($I$). Le second
agit quand le contrôle est $\ket{1}$ : la cible subit $U$.
\end{definition}

Comme $\ket{0}\bra{0}\ket{x}=\delta_{x0}\ket{0}$ et $\ket{1}\bra{1}\ket{x}=\delta_{x1}\ket{1}$,
\[
  \mathrm{C}U\big(\ket{x}\otimes\ket{y}\big)=\ket{x}\otimes U^x\ket{y},\qquad U^0=I,\ U^1=U .
\]
Par la règle $A\otimes B=(A_{ij}B)$, $\ket{0}\bra{0}\otimes I=\begin{pmatrix}I&0\\0&0\end{pmatrix}$
et $\ket{1}\bra{1}\otimes U=\begin{pmatrix}0&0\\0&U\end{pmatrix}$ (blocs $2\times2$), donc
\[
  \mathrm{C}U=\begin{pmatrix}I&0\\0&U\end{pmatrix}
  =\begin{pmatrix}1&0&0&0\\0&1&0&0\\0&0&u_{00}&u_{01}\\0&0&u_{10}&u_{11}\end{pmatrix}
  \qquad\text{pour }U=\begin{pmatrix}u_{00}&u_{01}\\u_{10}&u_{11}\end{pmatrix}.
\]
Cette matrice est unitaire car $I^\dagger I=I$ et $U^\dagger U=I$.

\paragraph{Le CNOT est un cas particulier.} Pour $U=\sigma_x$,
$\mathrm{C}\sigma_x=\ket{0}\bra{0}_x\otimes I_y+\ket{1}\bra{1}_x\otimes\sigma_x$ est le CNOT de
contrôle $x$ (premier chiffre), c'est-à-dire $\mathrm{CNOT}_{12}$ de la section~\ref{sec:cnot}.
On calcule l'action sur les quatre états de base, avec $\braket{0}{0}=\braket{1}{1}=1$,
$\braket{0}{1}=\braket{1}{0}=0$, $\sigma_x\ket{0}=\ket{1}$ et $\sigma_x\ket{1}=\ket{0}$ :
\begin{align*}
  \mathrm{C}\sigma_x\ket{00}&=\big[\ket{0}\bra{0}_x\otimes I_y\big]\ket{0}_x\ket{0}_y
      +\big[\ket{1}\bra{1}_x\otimes\sigma_x\big]\ket{0}_x\ket{0}_y\\
    &=\ket{0}\braket{0}{0}\otimes I\ket{0}+\ket{1}\braket{1}{0}\otimes\sigma_x\ket{0}
     =\ket{0}\otimes\ket{0}+0=\ket{00},\\[1.2ex]
  \mathrm{C}\sigma_x\ket{01}&=\ket{0}\braket{0}{0}\otimes I\ket{1}+\ket{1}\braket{1}{0}\otimes\sigma_x\ket{1}
     =\ket{0}\otimes\ket{1}+0=\ket{01},\\[1.2ex]
  \mathrm{C}\sigma_x\ket{10}&=\ket{0}\braket{0}{1}\otimes I\ket{0}+\ket{1}\braket{1}{1}\otimes\sigma_x\ket{0}
     =0+\ket{1}\otimes\ket{1}=\ket{11},\\[1.2ex]
  \mathrm{C}\sigma_x\ket{11}&=\ket{0}\braket{0}{1}\otimes I\ket{1}+\ket{1}\braket{1}{1}\otimes\sigma_x\ket{1}
     =0+\ket{1}\otimes\ket{0}=\ket{10}.
\end{align*}
On retrouve la table du CNOT : la cible n'est inversée que si le contrôle vaut $1$.

\paragraph{Un autre exemple : $U=Z$.} Ici
\[
  \mathrm{C}Z=\begin{pmatrix}I&0\\0&Z\end{pmatrix}=\mathrm{diag}(1,1,1,-1) :
\]
la porte multiplie $\ket{11}$ par $-1$ et ne touche pas les autres états de base. Elle est
symétrique entre les deux qubits. Elle se déduit du CNOT par
$\mathrm{C}Z=(I\otimes H)\,\mathrm{C}X\,(I\otimes H)$, où $\mathrm{C}X$ est le contrôlé-$X$.
En effet, $I\otimes H=\mathrm{diag}(H,H)$ (deux blocs $2\times2$ égaux à $H$), et le produit de
matrices diagonales par blocs se fait bloc par bloc :
\[
  (I\otimes H)\begin{pmatrix}I&0\\0&X\end{pmatrix}(I\otimes H)
  =\begin{pmatrix}H\,I\,H&0\\0&H\,X\,H\end{pmatrix}
  =\begin{pmatrix}I&0\\0&HXH\end{pmatrix},
\]
car $H^2=I$, et
\[
  HXH=\frac12\begin{pmatrix}1&1\\1&-1\end{pmatrix}\begin{pmatrix}0&1\\1&0\end{pmatrix}\begin{pmatrix}1&1\\1&-1\end{pmatrix}
  =\frac12\begin{pmatrix}1&1\\1&-1\end{pmatrix}\begin{pmatrix}1&-1\\1&1\end{pmatrix}
  =\frac12\begin{pmatrix}2&0\\0&-2\end{pmatrix}=Z
\]
(section~\ref{sec:hadamard} : $H$ échange les axes $x$ et $z$).

\subsubsection{Contrôlé-SWAP : la porte de Fredkin}

\begin{definition}{Porte de Fredkin (C-SWAP)}
Sur trois qubits $(x,y,z)$ (espace de dimension $2^3=8$), avec $x$ pour contrôle,
\[
  \mathrm{C\text{-}SWAP}=\ket{0}\bra{0}_x\otimes I_y\otimes I_z+\ket{1}\bra{1}_x\otimes\mathrm{SWAP}_{yz}.
\]
Si $x=0$, rien ne change ; si $x=1$, les qubits $y$ et $z$ sont échangés.
\end{definition}

Calcul sur quatre états de base (avec $I\ket{ab}=\ket{ab}$ et $\mathrm{SWAP}\ket{ab}=\ket{ba}$) :
\begin{align*}
  \mathrm{C\text{-}SWAP}\ket{000}&=\ket{0}\braket{0}{0}\otimes I\ket{00}+\ket{1}\braket{1}{0}\otimes\mathrm{SWAP}\ket{00}
     =\ket{0}\otimes\ket{00}+0=\ket{000},\\[1.2ex]
  \mathrm{C\text{-}SWAP}\ket{101}&=\ket{0}\braket{0}{1}\otimes I\ket{01}+\ket{1}\braket{1}{1}\otimes\mathrm{SWAP}\ket{01}
     =0+\ket{1}\otimes\ket{10}=\ket{110},\\[1.2ex]
  \mathrm{C\text{-}SWAP}\ket{110}&=\ket{0}\braket{0}{1}\otimes I\ket{10}+\ket{1}\braket{1}{1}\otimes\mathrm{SWAP}\ket{10}
     =0+\ket{1}\otimes\ket{01}=\ket{101},\\[1.2ex]
  \mathrm{C\text{-}SWAP}\ket{010}&=\ket{0}\braket{0}{0}\otimes I\ket{10}+\ket{1}\braket{1}{0}\otimes\mathrm{SWAP}\ket{10}
     =\ket{0}\otimes\ket{10}+0=\ket{010}.
\end{align*}
Les huit états de base donnent :
\begin{center}
\renewcommand{\arraystretch}{1.3}
\begin{tabular}{@{}ccccccccc@{}}
\toprule
Entrée & $\ket{000}$ & $\ket{001}$ & $\ket{010}$ & $\ket{011}$ & $\ket{100}$ & $\ket{101}$ & $\ket{110}$ & $\ket{111}$ \\
\midrule
Sortie & $\ket{000}$ & $\ket{001}$ & $\ket{010}$ & $\ket{011}$ & $\ket{100}$ & $\ket{110}$ & $\ket{101}$ & $\ket{111}$ \\
\bottomrule
\end{tabular}
\end{center}
Seuls $\ket{101}$ et $\ket{110}$ changent (contrôle à $1$ et $y\neq z$) ; pour $\ket{100}$ et
$\ket{111}$ l'échange de deux qubits égaux ne change rien. Dans la base
$\ket{000},\ket{001},\dots,\ket{111}$ (ordre binaire), la matrice est
$\mathrm{diag}(I_4,\mathrm{SWAP})$ (deux blocs $4\times4$) :
\[
  \mathrm{C\text{-}SWAP}=\begingroup\setlength{\arraycolsep}{3pt}\begin{pmatrix}
    1&0&0&0&0&0&0&0\\
    0&1&0&0&0&0&0&0\\
    0&0&1&0&0&0&0&0\\
    0&0&0&1&0&0&0&0\\
    0&0&0&0&1&0&0&0\\
    0&0&0&0&0&0&1&0\\
    0&0&0&0&0&1&0&0\\
    0&0&0&0&0&0&0&1
  \end{pmatrix}\endgroup .
\]
C'est l'identité dont les lignes $6$ et $7$ sont échangées.

\paragraph{Contrôle en superposition.} Avec $\ket{+}_x\otimes\ket{0}_y\otimes\ket{1}_z
=\frac{1}{\sqrt2}\big(\ket{001}+\ket{101}\big)$ :
\[
  \mathrm{C\text{-}SWAP}\,\frac{1}{\sqrt2}\big(\ket{001}+\ket{101}\big)
  =\frac{1}{\sqrt2}\big(\ket{001}+\ket{110}\big).
\]
Le résultat n'est pas un produit : les coefficients de $\ket{x}\otimes\ket{yz}$ forment la
matrice (lignes $x=0,1$, colonnes $yz=00,01,10,11$)
\[
  \frac{1}{\sqrt2}\begin{pmatrix}0&1&0&0\\0&0&1&0\end{pmatrix},
\]
dont les deux lignes ne sont pas proportionnelles. Le qubit $x$ est intriqué avec la paire
$(y,z)$.

\subsubsection{Contrôlé-contrôlé-NOT : la porte de Toffoli}

\begin{definition}{Porte de Toffoli (CCNOT)}
Sur trois qubits $(x,y,z)$, avec $x$ et $y$ pour contrôles et $z$ pour cible,
\[
  \mathrm{CCNOT}=\ket{0}\bra{0}_x\otimes I_y\otimes I_z
   +\ket{1}\bra{1}_x\otimes\Big[\ket{0}\bra{0}_y\otimes I_z+\ket{1}\bra{1}_y\otimes\sigma_x\Big].
\]
Le crochet est un CNOT (contrôle $y$, cible $z$), appliqué seulement si $x=1$ : si $x=0$, rien ;
si $x=1$ et $y=0$, rien ; si $x=1$ et $y=1$, le qubit $z$ est inversé.
\end{definition}

Calcul de $\mathrm{CCNOT}\ket{110}$ :
\begin{align*}
  \mathrm{CCNOT}\ket{110}&=\ket{0}\braket{0}{1}\otimes I\ket{1}\otimes I\ket{0}
     +\ket{1}\braket{1}{1}\otimes\Big[\ket{0}\braket{0}{1}\otimes I\ket{0}
       +\ket{1}\braket{1}{1}\otimes\sigma_x\ket{0}\Big]\\
  &=0+\ket{1}\otimes\Big[0+\ket{1}\otimes\ket{1}\Big]=\ket{111}.
\end{align*}
Pour les huit états de base, $\mathrm{CCNOT}\ket{x,y,z}=\ket{x,\,y,\,z\oplus xy}$ :
\begin{center}
\renewcommand{\arraystretch}{1.3}
\begin{tabular}{@{}ccccccccc@{}}
\toprule
Entrée & $\ket{000}$ & $\ket{001}$ & $\ket{010}$ & $\ket{011}$ & $\ket{100}$ & $\ket{101}$ & $\ket{110}$ & $\ket{111}$ \\
\midrule
Sortie & $\ket{000}$ & $\ket{001}$ & $\ket{010}$ & $\ket{011}$ & $\ket{100}$ & $\ket{101}$ & $\ket{111}$ & $\ket{110}$ \\
\bottomrule
\end{tabular}
\end{center}
Seuls $\ket{110}$ et $\ket{111}$ sont échangés : la matrice est $\mathrm{diag}(I_6,X)$
(un bloc $6\times6$ égal à l'identité, puis $X$),
\[
  \mathrm{CCNOT}=\begingroup\setlength{\arraycolsep}{3pt}\begin{pmatrix}
    1&0&0&0&0&0&0&0\\
    0&1&0&0&0&0&0&0\\
    0&0&1&0&0&0&0&0\\
    0&0&0&1&0&0&0&0\\
    0&0&0&0&1&0&0&0\\
    0&0&0&0&0&1&0&0\\
    0&0&0&0&0&0&0&1\\
    0&0&0&0&0&0&1&0
  \end{pmatrix}\endgroup .
\]

\begin{remarque}[Portes réversibles]
Les portes de Fredkin et de Toffoli sont des permutations des états de base qui échangent
seulement deux d'entre eux : elles sont leurs propres inverses ($\mathrm{C\text{-}SWAP}^2=\mathrm{CCNOT}^2=\mathrm{Id}$)
et unitaires. Avec $z=0$, la Toffoli calcule le ET logique $x\wedge y$ sans rien effacer :
$\ket{x,y,0}\to\ket{x,y,xy}$ ; avec $z=1$, elle calcule le NON-ET. Elle sert à écrire tout calcul
classique de façon réversible.
\end{remarque}

\clearpage
% ---------------------------------------------------------------------
\subsection{Récapitulatif du chapitre}\label{sec:recap4}
% ---------------------------------------------------------------------

\begin{center}
\adjustbox{max width=\linewidth}{%
\begin{tikzpicture}[>=Stealth, font=\footnotesize,
    boite/.style={draw=insarouge, fill=insarouge!6, rounded corners=2pt, thick,
                  minimum width=4.4cm, minimum height=1.5cm, align=center, inner sep=3pt},
    lien/.style={->, very thick, insagris},
    etiq/.style={font=\scriptsize\itshape, text=insagris, inner sep=2pt}]
  \node[boite] (uni) at (0,0)
    {\textbf{Porte unitaire}\\ $U^\dagger U=\mathrm{Id}$, \ $\ket{\Psi}\to U\ket{\Psi}$\\ {\scriptsize\S\,\ref{sec:portes}}};
  \node[boite] (un) at (5.9,0)
    {\textbf{Portes à un qubit}\\ $X$, $Y$, $Z$, $H$, $S$, $T$\\ {\scriptsize\S\,\ref{sec:pauli}--\ref{sec:phase}}};
  \node[boite] (blo) at (11.8,0)
    {\textbf{Rotations de Bloch}\\ $X,Y,Z,H$ : angle $\pi$ ; $P_\theta$ autour de $z$\\
     {\scriptsize\S\,\ref{sec:portes-bloch}}};
  \node[boite] (cir) at (11.8,-2.5)
    {\textbf{Circuits}\\ fils, boîtes, $U_2U_1$\\ {\scriptsize\S\,\ref{sec:circuits}, \ref{sec:op2}}};
  \node[boite] (cn) at (5.9,-2.5)
    {\textbf{CNOT et SWAP}\\ $\ket{a,b}\to\ket{a\oplus b,b}$\\ {\scriptsize\S\,\ref{sec:cnot}, \ref{sec:swap}}};
  \node[boite] (cu) at (0,-2.5)
    {\textbf{Portes contrôlées}\\ $\mathrm{C}U=\mathrm{diag}(I,U)$, Fredkin, Toffoli\\
     {\scriptsize\S\,\ref{sec:controlees}}};
  \draw[lien] (uni) -- (un);
  \draw[lien] (un) -- (blo);
  \draw[lien] (blo) -- (cir);
  \draw[lien] (cir) -- (cn) node[etiq, midway, above=1pt] {2 qubits};
  \draw[lien] (cn) -- (cu) node[etiq, midway, above=1pt] {contrôle};
\end{tikzpicture}}
\end{center}

\begin{aretenir}
\begin{itemize}
  \item Une porte est un opérateur unitaire : elle conserve la norme, elle est réversible
        ($U^{-1}=U^\dagger$) et agit sur la sphère de Bloch comme une rotation. C'est l'évolution de
        Schrödinger, $U=\ee^{-\ii\Hop t/\hbar}$.
  \item $X$ échange $\ket{0}$ et $\ket{1}$, $Z$ change le signe de $\ket{1}$, $Y=\ii XZ$ fait les deux ;
        $H$ échange les bases $Z$ et $X$ et crée des superpositions ; $P_\theta$, $S=P_{\pi/2}$ et
        $T=P_{\pi/4}$ font tourner autour de $z$ ($T^2=S$, $S^2=Z$).
  \item Dans un circuit, le temps va de gauche à droite et la matrice est le produit dans l'ordre
        inverse, $U_2U_1$. Le fil du haut porte le dernier chiffre du ket. Des portes en parallèle
        forment un produit tensoriel $U_A\otimes U_B$.
  \item Le CNOT inverse la cible si le contrôle vaut $1$ : $\ket{a,b}\to\ket{a\oplus b,b}$. Le SWAP
        échange deux qubits ; $\mathrm{CNOT}_{21}=\mathrm{SWAP}\,\mathrm{CNOT}_{12}\,\mathrm{SWAP}$.
  \item Une porte contrôlée a pour matrice $\mathrm{diag}(I,U)$. Fredkin (contrôlé-SWAP) et
        Toffoli (CCNOT) sont des permutations des états de base : réversibles et universelles
        pour le calcul classique.
\end{itemize}
\end{aretenir}

\begin{center}
\renewcommand{\arraystretch}{1.45}
\begin{tabular}{@{}l l r@{}}
\toprule
\textbf{Notion} & \textbf{Formule} & \textbf{§} \\
\midrule
Porte unitaire & $U^\dagger U=\mathrm{Id}$, \quad $\bra{\Psi}U^\dagger U\ket{\Psi}=1$ & \ref{sec:portes} \\
Pauli & $X\ket{0}=\ket{1}$, \quad $Z\ket{1}=-\ket{1}$, \quad $X^2=Y^2=Z^2=\mathrm{Id}$ & \ref{sec:pauli} \\
Hadamard & $H\ket{0}=\ket{+}$, \quad $H\ket{1}=\ket{-}$, \quad $H^2=\mathrm{Id}$ & \ref{sec:hadamard} \\
Phase & $P_\theta=\mathrm{diag}(1,\ee^{\ii\theta})$, \quad $S=P_{\pi/2}$, \quad $T=P_{\pi/4}$ & \ref{sec:phase} \\
Circuit & $U_1$ puis $U_2$ : \quad $U_2U_1$ & \ref{sec:circuits} \\
Portes locales & $(U\otimes I)(\ket{u}\otimes\ket{v})=U\ket{u}\otimes\ket{v}$ & \ref{sec:op2} \\
CNOT & $\ket{a,b}\to\ket{a\oplus b,b}$, \quad $\mathrm{CNOT}=I\otimes\ket{0}\bra{0}+X\otimes\ket{1}\bra{1}$
  & \ref{sec:cnot} \\
SWAP & $\mathrm{SWAP}\ket{a,b}=\ket{b,a}$, \quad $\mathrm{SWAP}=\mathrm{CNOT}_{12}\mathrm{CNOT}_{21}\mathrm{CNOT}_{12}$
  & \ref{sec:swap} \\
Contrôlé-$U$ & $\mathrm{C}U=\ket{0}\bra{0}\otimes I+\ket{1}\bra{1}\otimes U=\mathrm{diag}(I,U)$ & \ref{sec:controlees} \\
Toffoli & $\mathrm{CCNOT}\ket{x,y,z}=\ket{x,y,z\oplus xy}$ & \ref{sec:controlees} \\
\bottomrule
\end{tabular}
\end{center}

\begin{savoirfaire}
\begin{itemize}
  \item Appliquer une porte à un état (par la matrice ou par les ket-bra) et identifier sa rotation
        de Bloch (exercices~\ref{ex:conj} et~\ref{ex:phases}).
  \item Lire un circuit, écrire sa matrice dans le bon ordre et suivre l'état à chaque étape
        (exercice~\ref{ex:kickback}).
  \item Écrire la matrice de $U\otimes I$, du CNOT, du SWAP ou d'une porte contrôlée, et calculer
        leur action sur les états de base.
  \item Reconnaître une porte réversible qui calcule une fonction booléenne
        (exercice~\ref{ex:toffoli}).
\end{itemize}
\end{savoirfaire}
.
%  Il ne contient ni préambule ni \begin{document} : les environnements
%  (definition, resultat, remarque, aretenir, savoirfaire, exercice, correction),
%  les macros (\ii, \ee, \dd, \pd, \pdd, \Hop, \ket, \bra, \braket), les symboles
%  de circuits (ctl, cible, croix, mesure, porte) et les couleurs (insarouge,
%  insagris) sont définis dans main.tex. Un chapitre = une \section.
% =====================================================================

% Macros de Dirac (déjà définies dans main.tex ; ce bloc évite l'erreur
% « Undefined control sequence \ket » si main.tex est une ancienne version).
\providecommand{\ket}[1]{|#1\rangle}
\providecommand{\bra}[1]{\langle #1|}
\providecommand{\braket}[2]{\langle #1|#2\rangle}

% =====================================================================
\section{Portes quantiques et systèmes multi-qubits}\label{chap:portes}
% =====================================================================

Un calcul quantique est une suite de transformations unitaires appliquées à des qubits. Ce
chapitre décrit les portes à un qubit (Pauli, Hadamard, phase) et leur effet sur la sphère de
Bloch, la lecture d'un circuit, puis les portes qui agissent sur plusieurs qubits : CNOT, SWAP et
portes contrôlées.

% ---------------------------------------------------------------------
\subsection{Portes quantiques}\label{sec:portes}
% ---------------------------------------------------------------------

Un calcul quantique transforme des états. Les opérations sur un seul qubit sont appelées
\emph{portes} ; les sections suivantes passent à plusieurs qubits.

\begin{resultat}{Postulat 4 : évolution d'un système isolé}
L'évolution d'un système isolé est décrite par un opérateur unitaire : $\ket{\Psi}\to U\ket{\Psi}$
avec $U^\dagger U=\mathrm{Id}$.
\end{resultat}

C'est la forme de l'équation de Schrödinger du chapitre~\ref{chap:histoire}. Pour un hamiltonien
$\Hop$ indépendant du temps, $\ii\hbar\,\partial_t\ket{\Psi}=\Hop\ket{\Psi}$ a pour solution
$\ket{\Psi(t)}=U(t)\ket{\Psi(0)}$ avec $U(t)=\ee^{-\ii\Hop t/\hbar}$. Comme $\Hop$ est hermitien,
$U(t)^\dagger=\ee^{+\ii\Hop t/\hbar}$ et $U^\dagger U=\mathrm{Id}$ : l'évolution est unitaire. Une
porte est une telle évolution, obtenue en contrôlant $\Hop$ pendant une durée donnée.

\begin{definition}{Porte quantique}
Une porte à un qubit est un opérateur $U$ (matrice $2\times2$) \emph{unitaire} :
$U^\dagger U=\mathrm{Id}$. Elle transforme $\ket{\Psi}\to U\ket{\Psi}$, conserve la norme
($\bra{\Psi}U^\dagger U\ket{\Psi}=\braket{\Psi}{\Psi}=1$) et est réversible
($U^{-1}=U^\dagger$). Sur la sphère de Bloch (section~\ref{sec:bloch}), c'est une
\emph{rotation}.
\end{definition}

% ---------------------------------------------------------------------
\subsection{Portes de Pauli}\label{sec:pauli}
% ---------------------------------------------------------------------

Les matrices de Pauli de la section~\ref{sec:op-dirac} sont trois portes.

\begin{definition}{Portes $X$, $Y$, $Z$}
\begin{center}
\renewcommand{\arraystretch}{1}
\begin{tabular}{@{}lccl@{}}
\textbf{Porte} & \textbf{Matrice} & \textbf{Ket-bra} & \textbf{Nom} \\[0.8ex]
$X=\sigma_x$ & $\begin{pmatrix}0 & 1\\ 1 & 0\end{pmatrix}$
  & $\ket{0}\bra{1}+\ket{1}\bra{0}$ & bit flip \\[2.2ex]
$Y=\sigma_y$ & $\begin{pmatrix}0 & -\ii\\ \ii & 0\end{pmatrix}$
  & $-\ii\ket{0}\bra{1}+\ii\ket{1}\bra{0}$ &  \\[2.2ex]
$Z=\sigma_z$ & $\begin{pmatrix}1 & 0\\ 0 & -1\end{pmatrix}$
  & $\ket{0}\bra{0}-\ket{1}\bra{1}$ & phase flip \\
\end{tabular}
\end{center}
\end{definition}

\paragraph{Porte $X$ (inverseur).} Avec $\braket{0}{0}=\braket{1}{1}=1$ et
$\braket{0}{1}=\braket{1}{0}=0$ :
\begin{align*}
  X\ket{0}&=\ket{0}\braket{1}{0}+\ket{1}\braket{0}{0}=0\cdot\ket{0}+1\cdot\ket{1}=\ket{1},\\
  X\ket{1}&=\ket{0}\braket{1}{1}+\ket{1}\braket{0}{1}=1\cdot\ket{0}+0\cdot\ket{1}=\ket{0}.
\end{align*}
$X$ échange $\ket{0}$ et $\ket{1}$ : c'est le NOT quantique (\emph{bit flip}). Sur la
sphère de Bloch, c'est une rotation de $\pi$ autour de l'axe $x$ : le pôle nord devient
le pôle sud. Les états de l'axe $x$ ne bougent pas : $X\ket{+}=\ket{+}$ et
$X\ket{-}=-\ket{-}$.

\paragraph{Porte $Z$ (changement de phase).}
\begin{align*}
  Z\ket{0}&=\ket{0}\braket{0}{0}-\ket{1}\braket{1}{0}=\ket{0},\\
  Z\ket{1}&=\ket{0}\braket{0}{1}-\ket{1}\braket{1}{1}=-\ket{1},
\end{align*}
donc $Z$ change le signe de la composante sur $\ket{1}$ (\emph{phase flip}). Sur les
états de la base $X$, par linéarité :
\begin{align*}
  Z\ket{+}&=\tfrac{1}{\sqrt2}\big(Z\ket{0}+Z\ket{1}\big)=\tfrac{1}{\sqrt2}\big(\ket{0}-\ket{1}\big)=\ket{-},\\
  Z\ket{-}&=\tfrac{1}{\sqrt2}\big(Z\ket{0}-Z\ket{1}\big)=\tfrac{1}{\sqrt2}\big(\ket{0}+\ket{1}\big)=\ket{+}.
\end{align*}
Sous forme matricielle, $Z\ket{+}=\begin{pmatrix}1 & 0\\ 0 & -1\end{pmatrix}
\frac{1}{\sqrt2}\begin{pmatrix}1\\ 1\end{pmatrix}=\frac{1}{\sqrt2}\begin{pmatrix}1\\ -1\end{pmatrix}=\ket{-}$.
$Z$ est une rotation de $\pi$ autour de l'axe $z$ : elle laisse $\ket{0}$ et $\ket{1}$ en
place et échange $\ket{+}$ et $\ket{-}$.

\paragraph{Porte $Y$.} Avec $Y=-\ii\ket{0}\bra{1}+\ii\ket{1}\bra{0}$ :
\begin{align*}
  Y\ket{0}&=-\ii\ket{0}\braket{1}{0}+\ii\ket{1}\braket{0}{0}=\ii\ket{1},\\
  Y\ket{1}&=-\ii\ket{0}\braket{1}{1}+\ii\ket{1}\braket{0}{1}=-\ii\ket{0}.
\end{align*}
Sur les états de la base $X$ :
\begin{align*}
  Y\ket{+}&=\tfrac{1}{\sqrt2}\big(\ii\ket{1}-\ii\ket{0}\big)=-\ii\,\tfrac{1}{\sqrt2}\big(\ket{0}-\ket{1}\big)=-\ii\ket{-},\\
  Y\ket{-}&=\tfrac{1}{\sqrt2}\big(\ii\ket{1}+\ii\ket{0}\big)=\ii\,\tfrac{1}{\sqrt2}\big(\ket{0}+\ket{1}\big)=\ii\ket{+}.
\end{align*}
Un facteur global ($-\ii$ ou $\ii$) ne change pas l'état physique : $Y$ échange bien
$\ket{+}$ et $\ket{-}$. Sur les états de l'axe $y$ :
\begin{align*}
  Y\ket{{+}\ii}&=\tfrac{1}{\sqrt2}\big(Y\ket{0}+\ii Y\ket{1}\big)
   =\tfrac{1}{\sqrt2}\big(\ii\ket{1}+\ket{0}\big)=\ket{{+}\ii},\\
  Y\ket{{-}\ii}&=\tfrac{1}{\sqrt2}\big(Y\ket{0}-\ii Y\ket{1}\big)
   =\tfrac{1}{\sqrt2}\big(\ii\ket{1}-\ket{0}\big)=-\ket{{-}\ii}.
\end{align*}
$Y$ est une rotation de $\pi$ autour de l'axe $y$.

\begin{remarque}[Chaque Pauli fixe son axe]
$Z\ket{0}=\ket{0}$, $Z\ket{1}=-\ket{1}$ ; $X\ket{\pm}=\pm\ket{\pm}$ ; $Y\ket{{\pm}\ii}=\pm\ket{{\pm}\ii}$.
Les bases $Z$, $X$, $Y$ de la section~\ref{sec:bloch} sont les bases propres de
$Z$, $X$, $Y$. Une porte de Pauli laisse donc son propre axe invariant et retourne les
deux autres ; c'est pourquoi les exemples utilisent les états des deux autres axes.
\end{remarque}

% ---------------------------------------------------------------------
\subsection{Porte de Hadamard}\label{sec:hadamard}
% ---------------------------------------------------------------------

\begin{definition}{Porte de Hadamard}
\[
  H=\frac{1}{\sqrt2}\begin{pmatrix}1 & 1\\ 1 & -1\end{pmatrix}
  =\ket{+}\bra{0}+\ket{-}\bra{1}.
\]
\end{definition}

\begin{align*}
  H\ket{0}&=\frac{1}{\sqrt2}\begin{pmatrix}1 & 1\\ 1 & -1\end{pmatrix}\begin{pmatrix}1\\ 0\end{pmatrix}
    =\frac{1}{\sqrt2}\begin{pmatrix}1\\ 1\end{pmatrix}=\ket{+},\\[0.6ex]
  H\ket{1}&=\frac{1}{\sqrt2}\begin{pmatrix}1 & 1\\ 1 & -1\end{pmatrix}\begin{pmatrix}0\\ 1\end{pmatrix}
    =\frac{1}{\sqrt2}\begin{pmatrix}1\\ -1\end{pmatrix}=\ket{-},\\[0.6ex]
  H\ket{+}&=\frac{1}{2}\begin{pmatrix}1 & 1\\ 1 & -1\end{pmatrix}\begin{pmatrix}1\\ 1\end{pmatrix}
    =\frac{1}{2}\begin{pmatrix}2\\ 0\end{pmatrix}=\ket{0},\qquad
  H\ket{-}=\frac{1}{2}\begin{pmatrix}0\\ 2\end{pmatrix}=\ket{1}.
\end{align*}
Comme $H\ket{+}=\ket{0}$ et $H\ket{-}=\ket{1}$, on a $H^2=\mathrm{Id}$ ($H$ est son propre
inverse). Sur la sphère de Bloch, $H$ est une rotation de $\pi$ autour de l'axe
$(x+z)/\sqrt2$ : elle échange l'axe $x$ et l'axe $z$.

\begin{remarque}[À quoi sert $H$ ?]
\begin{itemize}[nosep]
  \item Elle \emph{crée une superposition} à partir d'un état de base :
        $H\ket{0}=\ket{+}$ donne $P(0)=P(1)=\frac12$. C'est la première étape de
        la plupart des algorithmes quantiques.
  \item Elle \emph{change de base} : elle envoie la base $Z$ sur la base $X$ et
        inversement. Mesurer dans la base $X$ revient à appliquer $H$ puis à mesurer
        dans la base $Z$ (figure~\ref{fig:mesurex}).
  \item Elle est réversible ($H^2=\mathrm{Id}$).
\end{itemize}
\end{remarque}

% ---------------------------------------------------------------------
\subsection{Portes de phase}\label{sec:phase}
% ---------------------------------------------------------------------

\begin{definition}{Portes de phase $P_\theta$, $S$ et $T$}
\[
  P_\theta=\begin{pmatrix}1 & 0\\ 0 & \ee^{\ii\theta}\end{pmatrix}
  =\ket{0}\bra{0}+\ee^{\ii\theta}\ket{1}\bra{1},
  \qquad
  S=P_{\pi/2}=\begin{pmatrix}1 & 0\\ 0 & \ii\end{pmatrix},
  \qquad
  T=P_{\pi/4}=\begin{pmatrix}1 & 0\\ 0 & \dfrac{1+\ii}{\sqrt2}\end{pmatrix}.
\]
\end{definition}

Ici $\ee^{\ii\pi/2}=\ii$ et $\ee^{\ii\pi/4}=\cos\frac\pi4+\ii\sin\frac\pi4=\frac{1+\ii}{\sqrt2}$.
$P_\theta$ laisse $\ket{0}$ inchangé et multiplie $\ket{1}$ par $\ee^{\ii\theta}$ :
\[
  P_\theta\big(\alpha\ket{0}+\beta\ket{1}\big)=\alpha\ket{0}+\ee^{\ii\theta}\beta\ket{1}.
\]
Dans la paramétrisation de Bloch, $\varphi\to\varphi+\theta$ : $P_\theta$ est une rotation
d'angle $\theta$ autour de l'axe $z$. On retrouve $P_\pi=Z$, et $S^2=Z$, $T^2=S$.

\paragraph{Exemples.}
\begin{align*}
  S\ket{+}&=\tfrac{1}{\sqrt2}\big(\ket{0}+\ii\ket{1}\big)=\ket{{+}\ii},\\
  S\ket{{+}\ii}&=\tfrac{1}{\sqrt2}\big(\ket{0}+\ii\cdot\ii\ket{1}\big)=\tfrac{1}{\sqrt2}\big(\ket{0}-\ket{1}\big)=\ket{-},\\
  T\ket{+}&=\tfrac{1}{\sqrt2}\begin{pmatrix}1 & 0\\ 0 & \frac{1+\ii}{\sqrt2}\end{pmatrix}\begin{pmatrix}1\\ 1\end{pmatrix}
    =\tfrac{1}{\sqrt2}\begin{pmatrix}1\\ \frac{1+\ii}{\sqrt2}\end{pmatrix}
    =\tfrac{1}{\sqrt2}\big(\ket{0}+\ee^{\ii\pi/4}\ket{1}\big)\quad\big(\text{angle polaire }\tfrac\pi2,\ \text{azimut }\tfrac\pi4\big).
\end{align*}
$S$ tourne donc $\ket{+}$ d'un quart de tour autour de $z$ (de l'axe $x$ à l'axe $y$),
puis d'un second quart de tour jusqu'à $\ket{-}$ ; $T$ fait un huitième de tour.

\begin{remarque}[Une porte de phase est une évolution libre]
Pour $\Hop=\frac{\hbar\omega}{2}Z$, $U(t)=\ee^{-\ii\Hop t/\hbar}=\mathrm{diag}\big(\ee^{-\ii\omega t/2},\ee^{\ii\omega t/2}\big)
=\ee^{-\ii\omega t/2}\,P_{\omega t}$. Au facteur global près, laisser évoluer le qubit pendant la
durée $t$ applique la porte de phase $P_{\omega t}$ : une rotation d'angle $\omega t$ autour de
l'axe $z$.
\end{remarque}

% ---------------------------------------------------------------------
\subsection{Portes et sphère de Bloch}\label{sec:portes-bloch}
% ---------------------------------------------------------------------

Toute porte à un qubit est une rotation de la sphère de Bloch. La figure~\ref{fig:portes} montre
les axes des portes usuelles et le tableau qui suit les résume.

\begin{figure}[!ht]
\centering
\begin{tikzpicture}[x={(-1.63cm,-0.56cm)}, y={(1.63cm,-0.56cm)}, z={(0cm,2.16cm)}, >=Stealth, scale=0.9]
  % sphère et grands cercles
  \shade[ball color=insagris!15, opacity=0.55] (0,0,0) circle[radius=2.3cm];
  \begin{scope}[insagris!75, thin]
    \draw[domain=-45:135, samples=60, variable=\t] plot ({cos(\t)},{sin(\t)},0);
    \draw[domain=135:315, samples=60, variable=\t, dashed] plot ({cos(\t)},{sin(\t)},0);
    \draw[domain=-62.7:117.3, samples=60, variable=\t] plot ({cos(\t)},0,{sin(\t)});
    \draw[domain=117.3:297.3, samples=60, variable=\t, dashed] plot ({cos(\t)},0,{sin(\t)});
    \draw[domain=-62.7:117.3, samples=60, variable=\t] plot (0,{cos(\t)},{sin(\t)});
    \draw[domain=117.3:297.3, samples=60, variable=\t, dashed] plot (0,{cos(\t)},{sin(\t)});
  \end{scope}
  % axes x, y, z (portes X, Y, Z)
  \draw[insagris, dashed] (0,0,0) -- (-1,0,0);
  \draw[insagris, dashed] (0,0,0) -- (0,-1,0);
  \draw[insagris, dashed] (0,0,0) -- (0,0,-1);
  \draw[insagris] (-1,0,0) -- (-1.3,0,0);
  \draw[insagris] (0,-1,0) -- (0,-1.3,0);
  \draw[insagris] (0,0,-1) -- (0,0,-1.3);
  \draw[insagris, -{Stealth[length=2.2mm]}] (0,0,0) -- (1.3,0,0);
  \draw[insagris, -{Stealth[length=2.2mm]}] (0,0,0) -- (0,1.3,0);
  \draw[insagris, -{Stealth[length=2.2mm]}] (0,0,0) -- (0,0,1.3);
  \node[anchor=east]  at (1.33,0,0) {$x$};
  \node[anchor=west]  at (0,1.33,0) {$y$};
  \node[anchor=south] at (0,0,1.33) {$z$};
  % axe de la porte H : (x+z)/sqrt(2)
  \draw[insarouge, dashed] (0,0,0) -- (-0.707,0,-0.707);
  \draw[insarouge, very thick, -{Stealth[length=2.6mm]}] (0,0,0) -- (0.98,0,0.98);
  \node[insarouge, anchor=south east] at (0.98,0,0.98) {axe de $H$};
  % rotation autour de z (portes de phase)
  \draw[insarouge, thick, -{Stealth[length=2.2mm]}, domain=15:105, samples=30, variable=\t]
    plot ({1.15*cos(\t)},{1.15*sin(\t)},0);
  \node[insarouge, anchor=north west, font=\small] at ({1.15*cos(60)},{1.15*sin(60)},0) {$P_\theta$};
\end{tikzpicture}
\caption{Portes et rotations sur la sphère de Bloch. Les portes $X$, $Y$, $Z$ sont des
rotations de $\pi$ autour des axes $x$, $y$, $z$ ; $H$ est une rotation de $\pi$ autour de
l'axe $(x+z)/\sqrt2$ (rouge) ; $P_\theta$ (donc $Z$, $S$, $T$) tourne d'un angle $\theta$
autour de $z$.}
\label{fig:portes}
\end{figure}

\begin{center}
\renewcommand{\arraystretch}{1}
\begin{tabular}{@{}lccc@{}}
\toprule
\textbf{Porte} & \textbf{Matrice} & \textbf{Effet sur la base} & \textbf{Rotation de Bloch} \\
\midrule
$X$ & $\begin{pmatrix}0 & 1\\ 1 & 0\end{pmatrix}$ & $\ket{0}\leftrightarrow\ket{1}$ & $\pi$ autour de $x$ \\[2.2ex]
$Y$ & $\begin{pmatrix}0 & -\ii\\ \ii & 0\end{pmatrix}$ & $\ket{0}\to\ii\ket{1},\ \ket{1}\to-\ii\ket{0}$ & $\pi$ autour de $y$ \\[2.2ex]
$Z$ & $\begin{pmatrix}1 & 0\\ 0 & -1\end{pmatrix}$ & $\ket{1}\to-\ket{1},\ \ket{\pm}\leftrightarrow\ket{\mp}$ & $\pi$ autour de $z$ \\[2.2ex]
$H$ & $\dfrac{1}{\sqrt2}\begin{pmatrix}1 & 1\\ 1 & -1\end{pmatrix}$ & $\ket{0}\leftrightarrow\ket{+},\ \ket{1}\leftrightarrow\ket{-}$ & $\pi$ autour de $\frac{x+z}{\sqrt2}$ \\[2.2ex]
$S$ & $\begin{pmatrix}1 & 0\\ 0 & \ii\end{pmatrix}$ & $\ket{1}\to\ii\ket{1}$ & $\frac\pi2$ autour de $z$ \\[2.2ex]
$T$ & $\begin{pmatrix}1 & 0\\ 0 & \frac{1+\ii}{\sqrt2}\end{pmatrix}$ & $\ket{1}\to\ee^{\ii\pi/4}\ket{1}$ & $\frac\pi4$ autour de $z$ \\
\bottomrule
\end{tabular}
\end{center}

\begin{figure}[!ht]
\centering
\begin{tikzpicture}[x={(-1.63cm,-0.56cm)}, y={(1.63cm,-0.56cm)}, z={(0cm,2.16cm)}, >=Stealth, scale=0.9]
  % sphère et grands cercles
  \shade[ball color=insagris!15, opacity=0.55] (0,0,0) circle[radius=2.3cm];
  \begin{scope}[insagris!75, thin]
    \draw[domain=-45:135, samples=60, variable=\t] plot ({cos(\t)},{sin(\t)},0);
    \draw[domain=135:315, samples=60, variable=\t, dashed] plot ({cos(\t)},{sin(\t)},0);
    \draw[domain=-62.7:117.3, samples=60, variable=\t] plot ({cos(\t)},0,{sin(\t)});
    \draw[domain=117.3:297.3, samples=60, variable=\t, dashed] plot ({cos(\t)},0,{sin(\t)});
    \draw[domain=-62.7:117.3, samples=60, variable=\t] plot (0,{cos(\t)},{sin(\t)});
    \draw[domain=117.3:297.3, samples=60, variable=\t, dashed] plot (0,{cos(\t)},{sin(\t)});
  \end{scope}
  % axes x, y, z (portes X, Y, Z)
  \draw[insagris, dashed] (0,0,0) -- (-1,0,0);
  \draw[insagris, dashed] (0,0,0) -- (0,-1,0);
  \draw[insagris, dashed] (0,0,0) -- (0,0,-1);
  \draw[insagris] (-1,0,0) -- (-1.3,0,0);
  \draw[insagris] (0,-1,0) -- (0,-1.3,0);
  \draw[insagris] (0,0,-1) -- (0,0,-1.3);
  \draw[insagris, -{Stealth[length=2.2mm]}] (0,0,0) -- (1.3,0,0);
  \draw[insagris, -{Stealth[length=2.2mm]}] (0,0,0) -- (0,1.3,0);
  \draw[insagris, -{Stealth[length=2.2mm]}] (0,0,0) -- (0,0,1.3);
  % trajet : |0> --H--> |+> --S--> |+i>
  \draw[insarouge, very thick, -{Stealth[length=2.6mm]}, domain=3:87, samples=30, variable=\t]
    plot ({sin(\t)},0,{cos(\t)});
  \draw[insarouge, very thick, -{Stealth[length=2.6mm]}, domain=3:87, samples=30, variable=\t]
    plot ({cos(\t)},{sin(\t)},0);
  \foreach \ax/\ay/\az in {0/0/1, 1/0/0, 0/1/0} \fill[insarouge] (\ax,\ay,\az) circle[radius=2pt];
  \node[insarouge, anchor=south] at (0,0,1.33) {$\ket{0}$};
  \node[insarouge, anchor=east] at (1.33,0,0) {$\ket{+}$};
  \node[insarouge, anchor=west] at (0,1.33,0) {$\ket{{+}\ii}$};
  \node[insarouge, font=\small, anchor=south west] at ({sin(45)},0,{cos(45)}) {$H$};
  \node[insarouge, font=\small, anchor=north west] at ({cos(45)},{sin(45)},0) {$S$};
\end{tikzpicture}
\caption{Trajet d'un état sous le circuit $H$ puis $S$ (flèches schématiques). $H$ envoie
$\ket{0}$ sur $\ket{+}$ ; $S$, rotation d'un quart de tour autour de $z$, envoie $\ket{+}$ sur
$\ket{{+}\ii}$. Les trois états sont alignés sur les axes $z$, $x$ et $y$.}
\label{fig:trajet}
\end{figure}

% ---------------------------------------------------------------------
\subsection{Circuits quantiques}\label{sec:circuits}
% ---------------------------------------------------------------------

Un \emph{circuit} représente un calcul : chaque qubit est un fil horizontal, le temps va de gauche
à droite, et chaque porte est un symbole posé sur les fils qu'elle concerne. La
figure~\ref{fig:symboles} donne les symboles utilisés dans ce document.

\begin{figure}[!ht]
\centering
\begin{adjustbox}{max width=\linewidth}
\begin{tikzpicture}[>=Stealth, font=\small,
    nom/.style={font=\footnotesize, align=center, anchor=north}]
  % (a) porte à un qubit
  \draw (0,0) -- (1.8,0);
  \node[porte] at (0.9,0) {$U$};
  \node[nom] at (0.9,-1.15) {porte $U$\\ à un qubit};
  % (b) CNOT
  \begin{scope}[xshift=3.2cm]
    \draw (0,0.6) -- (1.8,0.6);
    \draw (0,-0.6) -- (1.8,-0.6);
    \draw[insarouge, thick] (0.9,0.6) -- (0.9,-0.6);
    \pic at (0.9,0.6) {ctl};
    \pic at (0.9,-0.6) {cible};
    \node[nom] at (0.9,-1.15) {CNOT\\ (contrôle $\bullet$, cible $\oplus$)};
  \end{scope}
  % (c) SWAP
  \begin{scope}[xshift=6.4cm]
    \draw (0,0.6) -- (1.8,0.6);
    \draw (0,-0.6) -- (1.8,-0.6);
    \draw[insarouge, thick] (0.9,0.6) -- (0.9,-0.6);
    \pic at (0.9,0.6) {croix};
    \pic at (0.9,-0.6) {croix};
    \node[nom] at (0.9,-1.15) {SWAP\\ (échange)};
  \end{scope}
  % (d) contrôlé-U
  \begin{scope}[xshift=9.6cm]
    \draw (0,0.6) -- (1.8,0.6);
    \draw (0,-0.6) -- (1.8,-0.6);
    \draw[insarouge, thick] (0.9,-0.6) -- (0.9,0.3);
    \pic at (0.9,-0.6) {ctl};
    \node[porte] at (0.9,0.6) {$U$};
    \node[nom] at (0.9,-1.15) {contrôlé-$U$\\ (contrôle en bas)};
  \end{scope}
  % (e) mesure
  \begin{scope}[xshift=12.8cm]
    \draw (0,0) -- (0.6,0);
    \pic at (0.9,0) {mesure};
    \draw (1.2,0.03) -- (2.0,0.03);
    \draw (1.2,-0.03) -- (2.0,-0.03);
    \node[nom] at (1.0,-1.15) {mesure\\ (fil double : bit)};
  \end{scope}
\end{tikzpicture}
\end{adjustbox}
\caption{Symboles des circuits quantiques. Le point noir $\bullet$ marque un qubit de contrôle,
le symbole $\oplus$ une cible qui subit $X$, les deux croix un échange. Après une mesure, le fil
double transporte un bit classique.}
\label{fig:symboles}
\end{figure}

\begin{remarque}[Conventions de lecture]
\begin{itemize}[nosep]
  \item Le temps va de gauche à droite. Deux portes successives $U_1$ puis $U_2$ se multiplient
        dans l'ordre inverse : la matrice du circuit est $U_2U_1$, car $U_1$ agit d'abord sur le ket.
  \item Avec plusieurs fils, le fil du \emph{haut} porte le \emph{dernier} chiffre du ket et le
        fil du \emph{bas} le premier : si $A$ est le premier qubit et $B$ le second,
        $\ket{a,b}=\ket{a}_A\otimes\ket{b}_B$ et $B$ est dessiné au-dessus de $A$.
  \item Deux portes placées dans la même colonne, sur deux fils, agissent en parallèle : c'est un
        produit tensoriel $U_A\otimes U_B$ (section~\ref{sec:op2}).
\end{itemize}
\end{remarque}

\paragraph{Exemple.} Le circuit de la figure~\ref{fig:circuit-hs} applique $H$ puis $S$ à
$\ket{0}$. Sa matrice est $SH$ :
\[
  SH=\begin{pmatrix}1&0\\0&\ii\end{pmatrix}\frac{1}{\sqrt2}\begin{pmatrix}1&1\\1&-1\end{pmatrix}
  =\frac{1}{\sqrt2}\begin{pmatrix}1&1\\ \ii&-\ii\end{pmatrix},
  \qquad
  SH\ket{0}=\frac{1}{\sqrt2}\begin{pmatrix}1\\ \ii\end{pmatrix}=\ket{{+}\ii}.
\]
Le circuit prépare donc l'état $\ket{{+}\ii}$ de la base $Y$ à partir de $\ket{0}$
(figure~\ref{fig:trajet}). L'ordre compte : $HS\ket{0}=H\ket{0}=\ket{+}$, car $S\ket{0}=\ket{0}$.

\begin{figure}[!ht]
\centering
\begin{tikzpicture}[>=Stealth, font=\small]
  \draw (0,0) -- (6.6,0);
  \node[anchor=east] at (0,0) {$\ket{0}$};
  \node[porte] at (2.0,0) {$H$};
  \node[porte] at (4.3,0) {$S$};
  \node[anchor=west] at (6.6,0) {$\ket{{+}\ii}$};
  \foreach \x in {0.7, 3.15, 5.5} \draw[dashed, insagris!70] (\x,0.6) -- (\x,-0.95);
  \node[font=\footnotesize, anchor=north] at (0.7,-0.95) {$\ket{0}$};
  \node[font=\footnotesize, anchor=north] at (3.15,-0.95) {$\ket{+}$};
  \node[font=\footnotesize, anchor=north] at (5.5,-0.95) {$\ket{{+}\ii}$};
\end{tikzpicture}
\caption{Circuit $H$ puis $S$ appliqué à $\ket{0}$. Les états intermédiaires sont indiqués sous
les traits pointillés ; la matrice du circuit est $SH$.}
\label{fig:circuit-hs}
\end{figure}

\paragraph{Mesure.} Le symbole de mesure (compteur) donne le résultat $0$ ou $1$ dans la base $Z$.
Mesurer dans la base $X$ revient à appliquer $H$ puis à mesurer dans la base $Z$
(figure~\ref{fig:mesurex}) : comme $H\ket{\pm}=\ket{0}$ ou $\ket{1}$, le résultat $0$ correspond à
$\ket{+}$ et le résultat $1$ à $\ket{-}$, avec $P(\pm)=|\braket{\pm}{\Psi}|^2$.

\begin{figure}[!ht]
\centering
\begin{tikzpicture}[>=Stealth, font=\small]
  % (a) base Z
  \draw (0,0) -- (1.2,0);
  \node[anchor=east] at (0,0) {$\ket{\Psi}$};
  \pic at (1.5,0) {mesure};
  \draw (1.8,0.03) -- (2.6,0.03);
  \draw (1.8,-0.03) -- (2.6,-0.03);
  \node[anchor=west] at (2.7,0) {$0$ ou $1$};
  \node[font=\footnotesize] at (1.5,-0.95) {(a) mesure en base $Z$};
  % (b) base X
  \begin{scope}[xshift=7.5cm]
    \draw (0,0) -- (2.6,0);
    \node[anchor=east] at (0,0) {$\ket{\Psi}$};
    \node[porte] at (1.0,0) {$H$};
    \pic at (2.3,0) {mesure};
    \draw (2.6,0.03) -- (3.4,0.03);
    \draw (2.6,-0.03) -- (3.4,-0.03);
    \node[anchor=west] at (3.5,0) {$+$ ou $-$};
    \node[font=\footnotesize] at (1.9,-0.95) {(b) mesure en base $X$};
  \end{scope}
\end{tikzpicture}
\caption{(a) Mesure dans la base $Z$. (b) Mesure dans la base $X$ : porte de Hadamard, puis
mesure dans la base $Z$ ; les résultats $0$ et $1$ se lisent $+$ et $-$.}
\label{fig:mesurex}
\end{figure}

% ---------------------------------------------------------------------
\subsection{Portes locales sur deux qubits}\label{sec:op2}
% ---------------------------------------------------------------------

Une porte appliquée à un seul qubit d'un système de deux s'écrit avec un produit tensoriel
(section~\ref{sec:tenseur}) : $U\otimes I$ agit sur $A$ (premier facteur, fil du bas), $I\otimes U$
agit sur $B$, et $U_A\otimes U_B$ agit sur les deux en parallèle (figure~\ref{fig:local}).

\begin{figure}[!ht]
\centering
\begin{tikzpicture}[>=Stealth, font=\small]
  \foreach \xs/\haut/\bas/\tit in {0/{}/{$X$}/{(a) $X\otimes I$},
                                   5.4/{$X$}/{}/{(b) $I\otimes X$},
                                   10.8/{$H$}/{$H$}/{(c) $H\otimes H$}} {
    \begin{scope}[xshift=\xs cm]
      \draw (0,0.75) -- (2.6,0.75);
      \draw (0,0) -- (2.6,0);
      \node[anchor=east] at (0,0.75) {$B$};
      \node[anchor=east] at (0,0) {$A$};
      \ifx\haut\empty\else \node[porte] at (1.3,0.75) {\haut}; \fi
      \ifx\bas\empty\else \node[porte] at (1.3,0) {\bas}; \fi
      \node[font=\footnotesize] at (1.3,-0.85) {\tit};
    \end{scope}
  }
\end{tikzpicture}
\caption{Portes en parallèle sur deux qubits. Le qubit $A$, premier facteur du produit tensoriel,
est sur le fil du bas. (a) $X$ sur $A$ seulement. (b) $X$ sur $B$ seulement. (c) $H$ sur les deux.}
\label{fig:local}
\end{figure}

\subsubsection{Portes agissant sur un seul des deux qubits}

On note $I=\mathrm{Id}$ ($2\times2$). $X\otimes I$ agit sur $A$, $I\otimes X$ agit sur $B$ :
\begin{align*}
  I\otimes X&=\begin{pmatrix}1\cdot X & 0\cdot X\\ 0\cdot X & 1\cdot X\end{pmatrix}
    =\begin{pmatrix}0 & 1 & 0 & 0\\ 1 & 0 & 0 & 0\\ 0 & 0 & 0 & 1\\ 0 & 0 & 1 & 0\end{pmatrix},\\[1ex]
  X\otimes I&=\begin{pmatrix}0\cdot I & 1\cdot I\\ 1\cdot I & 0\cdot I\end{pmatrix}
    =\begin{pmatrix}0 & 0 & 1 & 0\\ 0 & 0 & 0 & 1\\ 1 & 0 & 0 & 0\\ 0 & 1 & 0 & 0\end{pmatrix}.
\end{align*}
Sur $\ket{00}=(1,0,0,0)^{\mathsf T}$ (première colonne de la matrice) :
$(I\otimes X)\ket{00}=(0,1,0,0)^{\mathsf T}=\ket{01}$ et
$(X\otimes I)\ket{00}=(0,0,1,0)^{\mathsf T}=\ket{10}$. Avec la règle :
$(I\otimes X)\ket{00}=I\ket{0}\otimes X\ket{0}=\ket{0}\otimes\ket{1}=\ket{01}$.

Pour la porte de Hadamard, avec $H=\frac{1}{\sqrt2}\begin{pmatrix}1 & 1\\ 1 & -1\end{pmatrix}$ :
\[
  H\otimes I=\frac{1}{\sqrt2}\begin{pmatrix}1\cdot I & 1\cdot I\\ 1\cdot I & -1\cdot I\end{pmatrix}
  =\frac{1}{\sqrt2}\begin{pmatrix}1 & 0 & 1 & 0\\ 0 & 1 & 0 & 1\\ 1 & 0 & -1 & 0\\ 0 & 1 & 0 & -1\end{pmatrix}.
\]
Ses colonnes donnent l'image des quatre états de base :
\begin{align*}
  (H\otimes I)\ket{00}&=\tfrac{1}{\sqrt2}(1,0,1,0)^{\mathsf T}=\ket{+}\otimes\ket{0},&
  (H\otimes I)\ket{01}&=\tfrac{1}{\sqrt2}(0,1,0,1)^{\mathsf T}=\ket{+}\otimes\ket{1},\\
  (H\otimes I)\ket{10}&=\tfrac{1}{\sqrt2}(1,0,-1,0)^{\mathsf T}=\ket{-}\otimes\ket{0},&
  (H\otimes I)\ket{11}&=\tfrac{1}{\sqrt2}(0,1,0,-1)^{\mathsf T}=\ket{-}\otimes\ket{1}.
\end{align*}

\subsubsection{Hadamard sur les deux qubits}

\[
  H\otimes H=\frac12\begin{pmatrix}1\cdot H' & 1\cdot H'\\ 1\cdot H' & -1\cdot H'\end{pmatrix},
  \quad H'=\begin{pmatrix}1 & 1\\ 1 & -1\end{pmatrix},
  \qquad
  H\otimes H=\frac12\begin{pmatrix}1 & 1 & 1 & 1\\ 1 & -1 & 1 & -1\\ 1 & 1 & -1 & -1\\ 1 & -1 & -1 & 1\end{pmatrix}.
\]
La première colonne donne $(H\otimes H)\ket{00}=\frac12(1,1,1,1)^{\mathsf T}=\ket{+}\otimes\ket{+}$
(cohérent avec $H\ket{0}\otimes H\ket{0}$) : les quatre résultats $00,01,10,11$ sont
équiprobables, chacun avec la probabilité $\frac14$. Sur $n$ qubits, $H^{\otimes n}\ket{0\ldots0}$
donne une superposition uniforme de $2^n$ états.

% ---------------------------------------------------------------------
\subsection{Porte CNOT}\label{sec:cnot}
% ---------------------------------------------------------------------

\begin{definition}{Porte CNOT (NOT contrôlé)}
La porte CNOT agit sur deux qubits : le qubit $B$ est le \emph{contrôle}, le qubit $A$ est la
\emph{cible}. La cible est inversée (porte $X$) si le contrôle vaut $1$ et reste inchangée
s'il vaut $0$. Avec $\ket{a,b}=\ket{a}_A\otimes\ket{b}_B$ :
\[
  \mathrm{CNOT}\,\ket{a,b}=\ket{a\oplus b,\,b},
  \qquad
  \mathrm{CNOT}=I\otimes\ket{0}\bra{0}+X\otimes\ket{1}\bra{1},
\]
où $\oplus$ est l'addition modulo $2$ ($0\oplus0=1\oplus1=0$ et $0\oplus1=1\oplus0=1$).
\end{definition}

Sur le circuit de la figure~\ref{fig:cnot}, le point noir posé sur le fil de $B$ (en haut) marque
le contrôle et le symbole $\oplus$ posé sur le fil de $A$ (en bas) marque la cible.

\begin{figure}[!ht]
\centering
\begin{tikzpicture}[>=Stealth, font=\small]
  \draw (0,1.2) -- (6,1.2);
  \draw (0,0) -- (6,0);
  \node[anchor=east] at (0,1.2) {$\ket{b}_B$};
  \node[anchor=east] at (0,0) {$\ket{a}_A$};
  \node[anchor=west] at (6,1.2) {$\ket{b}_B$};
  \node[anchor=west] at (6,0) {$\ket{a\oplus b}_A$};
  \draw[insarouge, thick] (3,1.2) -- (3,0);
  \pic at (3,1.2) {ctl};
  \pic at (3,0) {cible};
  \node[anchor=south, text=insagris] at (3,1.5) {qubit de contrôle};
  \node[anchor=north, text=insagris] at (3,-0.5) {qubit cible};
\end{tikzpicture}
\caption{Porte CNOT : le fil du haut ($B$) est le contrôle, le fil du bas ($A$) est la cible. Le
ket s'écrit $\ket{a,b}$ : le fil du haut donne le dernier chiffre.}
\label{fig:cnot}
\end{figure}

\subsubsection{Table de vérité}

\begin{center}
\renewcommand{\arraystretch}{1.3}
\begin{tabular}{@{}cccc@{}}
\toprule
Entrée $\ket{a,b}$ & Contrôle $b$ & Sortie $\ket{a\oplus b,b}$ & Effet sur la cible $A$ \\
\midrule
$\ket{00}$ & $0$ & $\ket{00}$ & aucun \\
$\ket{01}$ & $1$ & $\ket{11}$ & $0\to1$ \\
$\ket{10}$ & $0$ & $\ket{10}$ & aucun \\
$\ket{11}$ & $1$ & $\ket{01}$ & $1\to0$ \\
\bottomrule
\end{tabular}
\end{center}

\subsubsection{Matrice du CNOT}

La $j$-ième colonne d'une matrice $4\times4$ est l'image du $j$-ième vecteur de base
($\ket{00},\ket{01},\ket{10},\ket{11}$, dans cet ordre). La table donne :
\begin{align*}
  \mathrm{CNOT}\ket{00}&=\ket{00}=\begin{pmatrix}1\\0\\0\\0\end{pmatrix},&
  \mathrm{CNOT}\ket{01}&=\ket{11}=\begin{pmatrix}0\\0\\0\\1\end{pmatrix},\\[1.5ex]
  \mathrm{CNOT}\ket{10}&=\ket{10}=\begin{pmatrix}0\\0\\1\\0\end{pmatrix},&
  \mathrm{CNOT}\ket{11}&=\ket{01}=\begin{pmatrix}0\\1\\0\\0\end{pmatrix},
\end{align*}
d'où, en juxtaposant ces quatre colonnes,
\[
  \mathrm{CNOT}=\begin{pmatrix}
    1 & 0 & 0 & 0\\ 0 & 0 & 0 & 1\\ 0 & 0 & 1 & 0\\ 0 & 1 & 0 & 0
  \end{pmatrix}.
\]

On retrouve cette matrice avec la formule $I\otimes\ket{0}\bra{0}+X\otimes\ket{1}\bra{1}$ et
la règle $A\otimes B=(A_{ij}B)$ (section~\ref{sec:tenseur}), où
$\ket{0}\bra{0}=\begin{pmatrix}1&0\\0&0\end{pmatrix}$ et
$\ket{1}\bra{1}=\begin{pmatrix}0&0\\0&1\end{pmatrix}$ :
\[
  I\otimes\ket{0}\bra{0}=\begin{pmatrix}1&0&0&0\\0&0&0&0\\0&0&1&0\\0&0&0&0\end{pmatrix},
  \qquad
  X\otimes\ket{1}\bra{1}=\begin{pmatrix}0&0&0&0\\0&0&0&1\\0&0&0&0\\0&1&0&0\end{pmatrix},
\]
et la somme de ces deux matrices est la matrice du CNOT.

\subsubsection{Action sur un état quelconque}

Pour $\ket{\Psi}=c_{00}\ket{00}+c_{01}\ket{01}+c_{10}\ket{10}+c_{11}\ket{11}$ :
\[
  \mathrm{CNOT}\begin{pmatrix}c_{00}\\c_{01}\\c_{10}\\c_{11}\end{pmatrix}
  =\begin{pmatrix}
    1\cdot c_{00}+0\cdot c_{01}+0\cdot c_{10}+0\cdot c_{11}\\
    0\cdot c_{00}+0\cdot c_{01}+0\cdot c_{10}+1\cdot c_{11}\\
    0\cdot c_{00}+0\cdot c_{01}+1\cdot c_{10}+0\cdot c_{11}\\
    0\cdot c_{00}+1\cdot c_{01}+0\cdot c_{10}+0\cdot c_{11}
  \end{pmatrix}
  =\begin{pmatrix}c_{00}\\c_{11}\\c_{10}\\c_{01}\end{pmatrix}.
\]
Le CNOT échange les amplitudes de $\ket{01}$ et de $\ket{11}$, c'est-à-dire des deux états
dont le contrôle vaut $1$, et laisse les autres.

\paragraph{Exemple : contrôle en superposition.} Soit $\ket{0}_A\otimes\ket{+}_B
=\frac{1}{\sqrt2}\big(\ket{00}+\ket{01}\big)=\frac{1}{\sqrt2}(1,1,0,0)^{\mathsf T}$,
un état séparable. Alors
\[
  \mathrm{CNOT}\,\frac{1}{\sqrt2}\begin{pmatrix}1\\1\\0\\0\end{pmatrix}
  =\frac{1}{\sqrt2}\begin{pmatrix}1\\0\\0\\1\end{pmatrix}
  =\frac{1}{\sqrt2}\big(\ket{00}+\ket{11}\big).
\]
Le résultat est intriqué ($c_{00}c_{11}-c_{01}c_{10}=\frac12\neq0$ ; section~\ref{sec:separable}) :
un CNOT dont le contrôle est en superposition crée de l'intrication. C'est l'état de Bell
$\ket{\Phi^+}$ du chapitre~\ref{chap:bell}.

\paragraph{Réversibilité.} Deux CNOT successifs redonnent l'état initial :
$\ket{a,b}\to\ket{a\oplus b,b}\to\ket{a\oplus b\oplus b,b}=\ket{a,b}$, donc
$\mathrm{CNOT}^2=\mathrm{Id}$. La matrice est réelle et symétrique, donc
$\mathrm{CNOT}^\dagger\mathrm{CNOT}=\mathrm{CNOT}^2=\mathrm{Id}$ : elle est bien unitaire.

\begin{remarque}[CNOT et copie d'un bit]
Avec la cible dans l'état $\ket{0}$, $\mathrm{CNOT}\ket{0,b}=\ket{b,b}$ pour $b=0$ et $b=1$ : le CNOT
copie le bit classique $b$. En revanche $\mathrm{CNOT}\big(\ket{0}\otimes\ket{+}\big)$ est l'état
ci-dessus, et non $\ket{+}\otimes\ket{+}$ : un état quantique inconnu ne peut pas être copié
(théorème de non-clonage, section~\ref{sec:teleportation}).
\end{remarque}

\begin{remarque}[Contrôle en dernier ou en premier]
La matrice dépend de l'ordre d'écriture du ket. Notons $\mathrm{CNOT}_{21}$ le CNOT dont le
contrôle est le \emph{second} chiffre du ket : c'est la matrice ci-dessus, adoptée dans ce
document (fil du haut). Beaucoup d'ouvrages mettent le contrôle en \emph{premier} ; le CNOT
correspondant est
\[
  \mathrm{CNOT}_{12}=\ket{0}\bra{0}\otimes I+\ket{1}\bra{1}\otimes X
  =\begin{pmatrix}1&0&0&0\\0&1&0&0\\0&0&0&1\\0&0&1&0\end{pmatrix}.
\]
Le même circuit a donc deux matrices, selon que le contrôle est écrit en dernier ou en premier.
Les deux sont reliées par le SWAP (section~\ref{sec:swap}). La section~\ref{sec:controlees} sur les
portes contrôlées met le contrôle en premier, avec la matrice $\mathrm{diag}(I,U)$.
\end{remarque}

% ---------------------------------------------------------------------
\subsection{Porte SWAP}\label{sec:swap}
% ---------------------------------------------------------------------

\begin{definition}{Porte SWAP (échange)}
La porte SWAP échange les états des deux qubits :
\[
  \mathrm{SWAP}\big(\ket{\psi}\otimes\ket{\phi}\big)=\ket{\phi}\otimes\ket{\psi}.
\]
Sur les états de base, $\mathrm{SWAP}\ket{a,b}=\ket{b,a}$ ; on prolonge par linéarité.
\end{definition}

Son symbole est formé de deux croix reliées par un trait vertical
(figure~\ref{fig:swap}).

\begin{figure}[!ht]
\centering
\begin{tikzpicture}[>=Stealth, font=\small]
  \draw (0,1.2) -- (6,1.2);
  \draw (0,0) -- (6,0);
  \node[anchor=east] at (0,1.2) {$\ket{a}$};
  \node[anchor=east] at (0,0) {$\ket{b}$};
  \node[anchor=west] at (6,1.2) {$\ket{b}$};
  \node[anchor=west] at (6,0) {$\ket{a}$};
  \draw[insarouge, thick] (3,1.2) -- (3,0);
  \pic at (3,1.2) {croix};
  \pic at (3,0) {croix};
\end{tikzpicture}
\caption{Porte SWAP : les deux qubits échangent leurs états.}
\label{fig:swap}
\end{figure}

\subsubsection{Table et matrice}

\begin{center}
\renewcommand{\arraystretch}{1.3}
\begin{tabular}{@{}cc@{}}
\toprule
Entrée & Sortie \\
\midrule
$\ket{00}$ & $\ket{00}$ \\
$\ket{01}$ & $\ket{10}$ \\
$\ket{10}$ & $\ket{01}$ \\
$\ket{11}$ & $\ket{11}$ \\
\bottomrule
\end{tabular}
\end{center}

Les images des quatre vecteurs de base sont les colonnes de la matrice :
\[
  \mathrm{SWAP}\ket{00}=\begin{pmatrix}1\\0\\0\\0\end{pmatrix},\quad
  \mathrm{SWAP}\ket{01}=\begin{pmatrix}0\\0\\1\\0\end{pmatrix},\quad
  \mathrm{SWAP}\ket{10}=\begin{pmatrix}0\\1\\0\\0\end{pmatrix},\quad
  \mathrm{SWAP}\ket{11}=\begin{pmatrix}0\\0\\0\\1\end{pmatrix},
\]
\[
  \mathrm{SWAP}=\begin{pmatrix}
    1 & 0 & 0 & 0\\ 0 & 0 & 1 & 0\\ 0 & 1 & 0 & 0\\ 0 & 0 & 0 & 1
  \end{pmatrix}.
\]
Appliquée à $\ket{01}=(0,1,0,0)^{\mathsf T}$, cette matrice sélectionne sa deuxième colonne :
\[
  \begin{pmatrix}
    1 & 0 & 0 & 0\\ 0 & 0 & 1 & 0\\ 0 & 1 & 0 & 0\\ 0 & 0 & 0 & 1
  \end{pmatrix}
  \begin{pmatrix}0\\1\\0\\0\end{pmatrix}
  =\begin{pmatrix}0\\0\\1\\0\end{pmatrix}=\ket{10}.
\]
Sur un état quelconque, $\mathrm{SWAP}\,(c_{00},c_{01},c_{10},c_{11})^{\mathsf T}
=(c_{00},c_{10},c_{01},c_{11})^{\mathsf T}$ : les amplitudes de $\ket{01}$ et $\ket{10}$
sont échangées.

\subsubsection{Accord avec la définition}

Pour deux états $\ket{\psi}=(a_0,a_1)^{\mathsf T}$ et $\ket{\phi}=(b_0,b_1)^{\mathsf T}$ :
\[
  \mathrm{SWAP}\,(\ket{\psi}\otimes\ket{\phi})
  =\mathrm{SWAP}\begin{pmatrix}a_0b_0\\a_0b_1\\a_1b_0\\a_1b_1\end{pmatrix}
  =\begin{pmatrix}a_0b_0\\a_1b_0\\a_0b_1\\a_1b_1\end{pmatrix}
  =\begin{pmatrix}b_0\\b_1\end{pmatrix}\otimes\begin{pmatrix}a_0\\a_1\end{pmatrix}
  =\ket{\phi}\otimes\ket{\psi},
\]
ce qui est bien la définition. La porte s'écrit aussi avec des projecteurs :
\[
  \mathrm{SWAP}=\sum_{a,b\in\{0,1\}}\ket{a}\bra{b}\otimes\ket{b}\bra{a}.
\]
En effet, pour $\ket{c}\otimes\ket{d}$ on trouve
$\sum_{a,b}\ket{a}\braket{b}{c}\otimes\ket{b}\braket{a}{d}$, et seul le terme $b=c$, $a=d$
est non nul : le résultat est $\ket{d}\otimes\ket{c}$.

\paragraph{Exemples.}
\begin{itemize}[nosep]
  \item Deux états différents : $\mathrm{SWAP}\big(\ket{0}\otimes\ket{+}\big)
        =\mathrm{SWAP}\,\frac{1}{\sqrt2}(1,1,0,0)^{\mathsf T}=\frac{1}{\sqrt2}(1,0,1,0)^{\mathsf T}
        =\ket{+}\otimes\ket{0}$.
  \item L'état $\ket{\psi}=\frac{1}{\sqrt2}\ket{00}+\frac{1}{\sqrt2}\ket{11}$ ne change pas :
        $\mathrm{SWAP}\ket{\psi}=\frac{1}{\sqrt2}\ket{00}+\frac{1}{\sqrt2}\ket{11}=\ket{\psi}$,
        car $\ket{00}$ et $\ket{11}$ sont invariants. De même
        $\frac{1}{\sqrt2}(\ket{00}-\ket{11})$ et $\frac{1}{\sqrt2}(\ket{01}+\ket{10})$ sont invariants,
        alors que $\mathrm{SWAP}\,\frac{1}{\sqrt2}(\ket{01}-\ket{10})=\frac{1}{\sqrt2}(\ket{10}-\ket{01})$
        est l'opposé de l'état de départ. Ce sont les quatre états de Bell (chapitre~\ref{chap:bell}).
\end{itemize}

\paragraph{Propriétés.} Échanger deux fois redonne l'état de départ :
$\mathrm{SWAP}^2=\mathrm{Id}$ ; la matrice est réelle et symétrique, donc
$\mathrm{SWAP}^\dagger\mathrm{SWAP}=\mathrm{SWAP}^2=\mathrm{Id}$ : elle est unitaire.

\subsubsection{Lien avec le CNOT}

\begin{resultat}{SWAP échange les rôles de contrôle et de cible}
\[
  \mathrm{CNOT}_{21}=\mathrm{SWAP}\;\mathrm{CNOT}_{12}\;\mathrm{SWAP}.
\]
\end{resultat}

\emph{Sur les états de base.}
$\mathrm{SWAP}\,\mathrm{CNOT}_{12}\,\mathrm{SWAP}\ket{a,b}
=\mathrm{SWAP}\,\mathrm{CNOT}_{12}\ket{b,a}
=\mathrm{SWAP}\ket{b,\,b\oplus a}=\ket{b\oplus a,\,b}=\mathrm{CNOT}_{21}\ket{a,b}$.

\emph{Par les matrices.} Multiplier à gauche par $\mathrm{SWAP}$ échange les lignes $2$ et
$3$, multiplier à droite l'échange les colonnes $2$ et $3$ :
\[
  \mathrm{SWAP}\,\mathrm{CNOT}_{12}
  =\begin{pmatrix}1&0&0&0\\0&0&0&1\\0&1&0&0\\0&0&1&0\end{pmatrix},
  \qquad
  \mathrm{SWAP}\,\mathrm{CNOT}_{12}\,\mathrm{SWAP}
  =\begin{pmatrix}1&0&0&0\\0&0&0&1\\0&0&1&0\\0&1&0&0\end{pmatrix}=\mathrm{CNOT}_{21}.
\]

\begin{remarque}[SWAP avec trois CNOT]
$\mathrm{SWAP}=\mathrm{CNOT}_{12}\,\mathrm{CNOT}_{21}\,\mathrm{CNOT}_{12}$. En effet, sur
$\ket{a,b}$ (on applique les portes de droite à gauche) :
\[
  \ket{a,b}\ \xrightarrow{\ \mathrm{CNOT}_{12}\ }\ \ket{a,\,a\oplus b}
  \ \xrightarrow{\ \mathrm{CNOT}_{21}\ }\ \ket{a\oplus(a\oplus b),\,a\oplus b}=\ket{b,\,a\oplus b}
  \ \xrightarrow{\ \mathrm{CNOT}_{12}\ }\ \ket{b,\,b\oplus(a\oplus b)}=\ket{b,a}.
\]
Les deux membres coïncident sur les états de base, donc sur tous les états par linéarité.
\end{remarque}

% ---------------------------------------------------------------------
\subsection{Portes contrôlées : contrôlé-\texorpdfstring{$U$}{U}, Fredkin et Toffoli}\label{sec:controlees}
% ---------------------------------------------------------------------

Le CNOT est le cas le plus simple d'une idée générale : appliquer une opération à un ou
plusieurs qubits \emph{cibles} seulement si un ou plusieurs qubits de \emph{contrôle} valent $1$.
Dans cette section, le contrôle est le \emph{premier} chiffre
du ket (remarque de la section~\ref{sec:cnot}). Les qubits sont notés $x$, $y$, $z$ (le qubit $x$ est le contrôle) et
$\ket{x,y,z}=\ket{x}\otimes\ket{y}\otimes\ket{z}$ ; d'après la convention de la section~\ref{sec:circuits}, le
contrôle est dessiné en bas (figure~\ref{fig:controlees}). Pour éviter la confusion avec les
matrices de Pauli, on écrit ici $\sigma_x=X$ pour la porte NOT.

\begin{figure}[!ht]
\centering
\begin{adjustbox}{max width=\linewidth}
\begin{tikzpicture}[>=Stealth, font=\small]
  % (a) contrôlé-U : contrôle x en bas, cible y en haut
  \draw (0,1.2) -- (3.4,1.2);
  \draw (0,0) -- (3.4,0);
  \node[anchor=east] at (0,1.2) {$y$};
  \node[anchor=east] at (0,0) {$x$};
  \draw[insarouge, thick] (1.7,0) -- (1.7,0.85);
  \pic at (1.7,0) {ctl};
  \node[porte] at (1.7,1.2) {$U$};
  \node[anchor=north, font=\footnotesize] at (1.7,-0.5) {(a) contrôlé-$U$};
  % (b) contrôlé-SWAP (Fredkin) : x contrôle, y et z échangés
  \draw (5,2.4) -- (8.4,2.4);
  \draw (5,1.2) -- (8.4,1.2);
  \draw (5,0) -- (8.4,0);
  \node[anchor=east] at (5,2.4) {$z$};
  \node[anchor=east] at (5,1.2) {$y$};
  \node[anchor=east] at (5,0) {$x$};
  \draw[insarouge, thick] (6.7,0) -- (6.7,2.4);
  \pic at (6.7,0) {ctl};
  \pic at (6.7,1.2) {croix};
  \pic at (6.7,2.4) {croix};
  \node[anchor=north, font=\footnotesize] at (6.7,-0.5) {(b) Fredkin (C-SWAP)};
  % (c) Toffoli (CCNOT) : x et y contrôlent, z cible
  \draw (10,2.4) -- (13.4,2.4);
  \draw (10,1.2) -- (13.4,1.2);
  \draw (10,0) -- (13.4,0);
  \node[anchor=east] at (10,2.4) {$z$};
  \node[anchor=east] at (10,1.2) {$y$};
  \node[anchor=east] at (10,0) {$x$};
  \draw[insarouge, thick] (11.7,0) -- (11.7,2.4);
  \pic at (11.7,0) {ctl};
  \pic at (11.7,1.2) {ctl};
  \pic at (11.7,2.4) {cible};
  \node[anchor=north, font=\footnotesize] at (11.7,-0.5) {(c) Toffoli (CCNOT)};
\end{tikzpicture}
\end{adjustbox}
\caption{Portes contrôlées. Le contrôle $x$ est le premier chiffre du ket, donc le fil du bas.
(a) Le qubit $y$ subit $U$ si $x=1$ ; pour $U=X$ on retrouve le CNOT. (b) Si $x=1$, les qubits
$y$ et $z$ sont échangés. (c) Si $x=y=1$, le qubit $z$ est inversé.}
\label{fig:controlees}
\end{figure}

\subsubsection{Contrôlé-\texorpdfstring{$U$}{U}}

\begin{definition}{Porte contrôlée-$U$}
Soit $U$ une matrice unitaire $2\times2$. Sur le système $(x,y)$, de contrôle $x$ et de cible $y$,
\[
  \mathrm{C}U=\ket{0}\bra{0}_x\otimes I_y+\ket{1}\bra{1}_x\otimes U_y .
\]
Le premier terme agit quand le contrôle est $\ket{0}$ : la cible ne change pas ($I$). Le second
agit quand le contrôle est $\ket{1}$ : la cible subit $U$.
\end{definition}

Comme $\ket{0}\bra{0}\ket{x}=\delta_{x0}\ket{0}$ et $\ket{1}\bra{1}\ket{x}=\delta_{x1}\ket{1}$,
\[
  \mathrm{C}U\big(\ket{x}\otimes\ket{y}\big)=\ket{x}\otimes U^x\ket{y},\qquad U^0=I,\ U^1=U .
\]
Par la règle $A\otimes B=(A_{ij}B)$, $\ket{0}\bra{0}\otimes I=\begin{pmatrix}I&0\\0&0\end{pmatrix}$
et $\ket{1}\bra{1}\otimes U=\begin{pmatrix}0&0\\0&U\end{pmatrix}$ (blocs $2\times2$), donc
\[
  \mathrm{C}U=\begin{pmatrix}I&0\\0&U\end{pmatrix}
  =\begin{pmatrix}1&0&0&0\\0&1&0&0\\0&0&u_{00}&u_{01}\\0&0&u_{10}&u_{11}\end{pmatrix}
  \qquad\text{pour }U=\begin{pmatrix}u_{00}&u_{01}\\u_{10}&u_{11}\end{pmatrix}.
\]
Cette matrice est unitaire car $I^\dagger I=I$ et $U^\dagger U=I$.

\paragraph{Le CNOT est un cas particulier.} Pour $U=\sigma_x$,
$\mathrm{C}\sigma_x=\ket{0}\bra{0}_x\otimes I_y+\ket{1}\bra{1}_x\otimes\sigma_x$ est le CNOT de
contrôle $x$ (premier chiffre), c'est-à-dire $\mathrm{CNOT}_{12}$ de la section~\ref{sec:cnot}.
On calcule l'action sur les quatre états de base, avec $\braket{0}{0}=\braket{1}{1}=1$,
$\braket{0}{1}=\braket{1}{0}=0$, $\sigma_x\ket{0}=\ket{1}$ et $\sigma_x\ket{1}=\ket{0}$ :
\begin{align*}
  \mathrm{C}\sigma_x\ket{00}&=\big[\ket{0}\bra{0}_x\otimes I_y\big]\ket{0}_x\ket{0}_y
      +\big[\ket{1}\bra{1}_x\otimes\sigma_x\big]\ket{0}_x\ket{0}_y\\
    &=\ket{0}\braket{0}{0}\otimes I\ket{0}+\ket{1}\braket{1}{0}\otimes\sigma_x\ket{0}
     =\ket{0}\otimes\ket{0}+0=\ket{00},\\[1.2ex]
  \mathrm{C}\sigma_x\ket{01}&=\ket{0}\braket{0}{0}\otimes I\ket{1}+\ket{1}\braket{1}{0}\otimes\sigma_x\ket{1}
     =\ket{0}\otimes\ket{1}+0=\ket{01},\\[1.2ex]
  \mathrm{C}\sigma_x\ket{10}&=\ket{0}\braket{0}{1}\otimes I\ket{0}+\ket{1}\braket{1}{1}\otimes\sigma_x\ket{0}
     =0+\ket{1}\otimes\ket{1}=\ket{11},\\[1.2ex]
  \mathrm{C}\sigma_x\ket{11}&=\ket{0}\braket{0}{1}\otimes I\ket{1}+\ket{1}\braket{1}{1}\otimes\sigma_x\ket{1}
     =0+\ket{1}\otimes\ket{0}=\ket{10}.
\end{align*}
On retrouve la table du CNOT : la cible n'est inversée que si le contrôle vaut $1$.

\paragraph{Un autre exemple : $U=Z$.} Ici
\[
  \mathrm{C}Z=\begin{pmatrix}I&0\\0&Z\end{pmatrix}=\mathrm{diag}(1,1,1,-1) :
\]
la porte multiplie $\ket{11}$ par $-1$ et ne touche pas les autres états de base. Elle est
symétrique entre les deux qubits. Elle se déduit du CNOT par
$\mathrm{C}Z=(I\otimes H)\,\mathrm{C}X\,(I\otimes H)$, où $\mathrm{C}X$ est le contrôlé-$X$.
En effet, $I\otimes H=\mathrm{diag}(H,H)$ (deux blocs $2\times2$ égaux à $H$), et le produit de
matrices diagonales par blocs se fait bloc par bloc :
\[
  (I\otimes H)\begin{pmatrix}I&0\\0&X\end{pmatrix}(I\otimes H)
  =\begin{pmatrix}H\,I\,H&0\\0&H\,X\,H\end{pmatrix}
  =\begin{pmatrix}I&0\\0&HXH\end{pmatrix},
\]
car $H^2=I$, et
\[
  HXH=\frac12\begin{pmatrix}1&1\\1&-1\end{pmatrix}\begin{pmatrix}0&1\\1&0\end{pmatrix}\begin{pmatrix}1&1\\1&-1\end{pmatrix}
  =\frac12\begin{pmatrix}1&1\\1&-1\end{pmatrix}\begin{pmatrix}1&-1\\1&1\end{pmatrix}
  =\frac12\begin{pmatrix}2&0\\0&-2\end{pmatrix}=Z
\]
(section~\ref{sec:hadamard} : $H$ échange les axes $x$ et $z$).

\subsubsection{Contrôlé-SWAP : la porte de Fredkin}

\begin{definition}{Porte de Fredkin (C-SWAP)}
Sur trois qubits $(x,y,z)$ (espace de dimension $2^3=8$), avec $x$ pour contrôle,
\[
  \mathrm{C\text{-}SWAP}=\ket{0}\bra{0}_x\otimes I_y\otimes I_z+\ket{1}\bra{1}_x\otimes\mathrm{SWAP}_{yz}.
\]
Si $x=0$, rien ne change ; si $x=1$, les qubits $y$ et $z$ sont échangés.
\end{definition}

Calcul sur quatre états de base (avec $I\ket{ab}=\ket{ab}$ et $\mathrm{SWAP}\ket{ab}=\ket{ba}$) :
\begin{align*}
  \mathrm{C\text{-}SWAP}\ket{000}&=\ket{0}\braket{0}{0}\otimes I\ket{00}+\ket{1}\braket{1}{0}\otimes\mathrm{SWAP}\ket{00}
     =\ket{0}\otimes\ket{00}+0=\ket{000},\\[1.2ex]
  \mathrm{C\text{-}SWAP}\ket{101}&=\ket{0}\braket{0}{1}\otimes I\ket{01}+\ket{1}\braket{1}{1}\otimes\mathrm{SWAP}\ket{01}
     =0+\ket{1}\otimes\ket{10}=\ket{110},\\[1.2ex]
  \mathrm{C\text{-}SWAP}\ket{110}&=\ket{0}\braket{0}{1}\otimes I\ket{10}+\ket{1}\braket{1}{1}\otimes\mathrm{SWAP}\ket{10}
     =0+\ket{1}\otimes\ket{01}=\ket{101},\\[1.2ex]
  \mathrm{C\text{-}SWAP}\ket{010}&=\ket{0}\braket{0}{0}\otimes I\ket{10}+\ket{1}\braket{1}{0}\otimes\mathrm{SWAP}\ket{10}
     =\ket{0}\otimes\ket{10}+0=\ket{010}.
\end{align*}
Les huit états de base donnent :
\begin{center}
\renewcommand{\arraystretch}{1.3}
\begin{tabular}{@{}ccccccccc@{}}
\toprule
Entrée & $\ket{000}$ & $\ket{001}$ & $\ket{010}$ & $\ket{011}$ & $\ket{100}$ & $\ket{101}$ & $\ket{110}$ & $\ket{111}$ \\
\midrule
Sortie & $\ket{000}$ & $\ket{001}$ & $\ket{010}$ & $\ket{011}$ & $\ket{100}$ & $\ket{110}$ & $\ket{101}$ & $\ket{111}$ \\
\bottomrule
\end{tabular}
\end{center}
Seuls $\ket{101}$ et $\ket{110}$ changent (contrôle à $1$ et $y\neq z$) ; pour $\ket{100}$ et
$\ket{111}$ l'échange de deux qubits égaux ne change rien. Dans la base
$\ket{000},\ket{001},\dots,\ket{111}$ (ordre binaire), la matrice est
$\mathrm{diag}(I_4,\mathrm{SWAP})$ (deux blocs $4\times4$) :
\[
  \mathrm{C\text{-}SWAP}=\begingroup\setlength{\arraycolsep}{3pt}\begin{pmatrix}
    1&0&0&0&0&0&0&0\\
    0&1&0&0&0&0&0&0\\
    0&0&1&0&0&0&0&0\\
    0&0&0&1&0&0&0&0\\
    0&0&0&0&1&0&0&0\\
    0&0&0&0&0&0&1&0\\
    0&0&0&0&0&1&0&0\\
    0&0&0&0&0&0&0&1
  \end{pmatrix}\endgroup .
\]
C'est l'identité dont les lignes $6$ et $7$ sont échangées.

\paragraph{Contrôle en superposition.} Avec $\ket{+}_x\otimes\ket{0}_y\otimes\ket{1}_z
=\frac{1}{\sqrt2}\big(\ket{001}+\ket{101}\big)$ :
\[
  \mathrm{C\text{-}SWAP}\,\frac{1}{\sqrt2}\big(\ket{001}+\ket{101}\big)
  =\frac{1}{\sqrt2}\big(\ket{001}+\ket{110}\big).
\]
Le résultat n'est pas un produit : les coefficients de $\ket{x}\otimes\ket{yz}$ forment la
matrice (lignes $x=0,1$, colonnes $yz=00,01,10,11$)
\[
  \frac{1}{\sqrt2}\begin{pmatrix}0&1&0&0\\0&0&1&0\end{pmatrix},
\]
dont les deux lignes ne sont pas proportionnelles. Le qubit $x$ est intriqué avec la paire
$(y,z)$.

\subsubsection{Contrôlé-contrôlé-NOT : la porte de Toffoli}

\begin{definition}{Porte de Toffoli (CCNOT)}
Sur trois qubits $(x,y,z)$, avec $x$ et $y$ pour contrôles et $z$ pour cible,
\[
  \mathrm{CCNOT}=\ket{0}\bra{0}_x\otimes I_y\otimes I_z
   +\ket{1}\bra{1}_x\otimes\Big[\ket{0}\bra{0}_y\otimes I_z+\ket{1}\bra{1}_y\otimes\sigma_x\Big].
\]
Le crochet est un CNOT (contrôle $y$, cible $z$), appliqué seulement si $x=1$ : si $x=0$, rien ;
si $x=1$ et $y=0$, rien ; si $x=1$ et $y=1$, le qubit $z$ est inversé.
\end{definition}

Calcul de $\mathrm{CCNOT}\ket{110}$ :
\begin{align*}
  \mathrm{CCNOT}\ket{110}&=\ket{0}\braket{0}{1}\otimes I\ket{1}\otimes I\ket{0}
     +\ket{1}\braket{1}{1}\otimes\Big[\ket{0}\braket{0}{1}\otimes I\ket{0}
       +\ket{1}\braket{1}{1}\otimes\sigma_x\ket{0}\Big]\\
  &=0+\ket{1}\otimes\Big[0+\ket{1}\otimes\ket{1}\Big]=\ket{111}.
\end{align*}
Pour les huit états de base, $\mathrm{CCNOT}\ket{x,y,z}=\ket{x,\,y,\,z\oplus xy}$ :
\begin{center}
\renewcommand{\arraystretch}{1.3}
\begin{tabular}{@{}ccccccccc@{}}
\toprule
Entrée & $\ket{000}$ & $\ket{001}$ & $\ket{010}$ & $\ket{011}$ & $\ket{100}$ & $\ket{101}$ & $\ket{110}$ & $\ket{111}$ \\
\midrule
Sortie & $\ket{000}$ & $\ket{001}$ & $\ket{010}$ & $\ket{011}$ & $\ket{100}$ & $\ket{101}$ & $\ket{111}$ & $\ket{110}$ \\
\bottomrule
\end{tabular}
\end{center}
Seuls $\ket{110}$ et $\ket{111}$ sont échangés : la matrice est $\mathrm{diag}(I_6,X)$
(un bloc $6\times6$ égal à l'identité, puis $X$),
\[
  \mathrm{CCNOT}=\begingroup\setlength{\arraycolsep}{3pt}\begin{pmatrix}
    1&0&0&0&0&0&0&0\\
    0&1&0&0&0&0&0&0\\
    0&0&1&0&0&0&0&0\\
    0&0&0&1&0&0&0&0\\
    0&0&0&0&1&0&0&0\\
    0&0&0&0&0&1&0&0\\
    0&0&0&0&0&0&0&1\\
    0&0&0&0&0&0&1&0
  \end{pmatrix}\endgroup .
\]

\begin{remarque}[Portes réversibles]
Les portes de Fredkin et de Toffoli sont des permutations des états de base qui échangent
seulement deux d'entre eux : elles sont leurs propres inverses ($\mathrm{C\text{-}SWAP}^2=\mathrm{CCNOT}^2=\mathrm{Id}$)
et unitaires. Avec $z=0$, la Toffoli calcule le ET logique $x\wedge y$ sans rien effacer :
$\ket{x,y,0}\to\ket{x,y,xy}$ ; avec $z=1$, elle calcule le NON-ET. Elle sert à écrire tout calcul
classique de façon réversible.
\end{remarque}

\clearpage
% ---------------------------------------------------------------------
\subsection{Récapitulatif du chapitre}\label{sec:recap4}
% ---------------------------------------------------------------------

\begin{center}
\adjustbox{max width=\linewidth}{%
\begin{tikzpicture}[>=Stealth, font=\footnotesize,
    boite/.style={draw=insarouge, fill=insarouge!6, rounded corners=2pt, thick,
                  minimum width=4.4cm, minimum height=1.5cm, align=center, inner sep=3pt},
    lien/.style={->, very thick, insagris},
    etiq/.style={font=\scriptsize\itshape, text=insagris, inner sep=2pt}]
  \node[boite] (uni) at (0,0)
    {\textbf{Porte unitaire}\\ $U^\dagger U=\mathrm{Id}$, \ $\ket{\Psi}\to U\ket{\Psi}$\\ {\scriptsize\S\,\ref{sec:portes}}};
  \node[boite] (un) at (5.9,0)
    {\textbf{Portes à un qubit}\\ $X$, $Y$, $Z$, $H$, $S$, $T$\\ {\scriptsize\S\,\ref{sec:pauli}--\ref{sec:phase}}};
  \node[boite] (blo) at (11.8,0)
    {\textbf{Rotations de Bloch}\\ $X,Y,Z,H$ : angle $\pi$ ; $P_\theta$ autour de $z$\\
     {\scriptsize\S\,\ref{sec:portes-bloch}}};
  \node[boite] (cir) at (11.8,-2.5)
    {\textbf{Circuits}\\ fils, boîtes, $U_2U_1$\\ {\scriptsize\S\,\ref{sec:circuits}, \ref{sec:op2}}};
  \node[boite] (cn) at (5.9,-2.5)
    {\textbf{CNOT et SWAP}\\ $\ket{a,b}\to\ket{a\oplus b,b}$\\ {\scriptsize\S\,\ref{sec:cnot}, \ref{sec:swap}}};
  \node[boite] (cu) at (0,-2.5)
    {\textbf{Portes contrôlées}\\ $\mathrm{C}U=\mathrm{diag}(I,U)$, Fredkin, Toffoli\\
     {\scriptsize\S\,\ref{sec:controlees}}};
  \draw[lien] (uni) -- (un);
  \draw[lien] (un) -- (blo);
  \draw[lien] (blo) -- (cir);
  \draw[lien] (cir) -- (cn) node[etiq, midway, above=1pt] {2 qubits};
  \draw[lien] (cn) -- (cu) node[etiq, midway, above=1pt] {contrôle};
\end{tikzpicture}}
\end{center}

\begin{aretenir}
\begin{itemize}
  \item Une porte est un opérateur unitaire : elle conserve la norme, elle est réversible
        ($U^{-1}=U^\dagger$) et agit sur la sphère de Bloch comme une rotation. C'est l'évolution de
        Schrödinger, $U=\ee^{-\ii\Hop t/\hbar}$.
  \item $X$ échange $\ket{0}$ et $\ket{1}$, $Z$ change le signe de $\ket{1}$, $Y=\ii XZ$ fait les deux ;
        $H$ échange les bases $Z$ et $X$ et crée des superpositions ; $P_\theta$, $S=P_{\pi/2}$ et
        $T=P_{\pi/4}$ font tourner autour de $z$ ($T^2=S$, $S^2=Z$).
  \item Dans un circuit, le temps va de gauche à droite et la matrice est le produit dans l'ordre
        inverse, $U_2U_1$. Le fil du haut porte le dernier chiffre du ket. Des portes en parallèle
        forment un produit tensoriel $U_A\otimes U_B$.
  \item Le CNOT inverse la cible si le contrôle vaut $1$ : $\ket{a,b}\to\ket{a\oplus b,b}$. Le SWAP
        échange deux qubits ; $\mathrm{CNOT}_{21}=\mathrm{SWAP}\,\mathrm{CNOT}_{12}\,\mathrm{SWAP}$.
  \item Une porte contrôlée a pour matrice $\mathrm{diag}(I,U)$. Fredkin (contrôlé-SWAP) et
        Toffoli (CCNOT) sont des permutations des états de base : réversibles et universelles
        pour le calcul classique.
\end{itemize}
\end{aretenir}

\begin{center}
\renewcommand{\arraystretch}{1.45}
\begin{tabular}{@{}l l r@{}}
\toprule
\textbf{Notion} & \textbf{Formule} & \textbf{§} \\
\midrule
Porte unitaire & $U^\dagger U=\mathrm{Id}$, \quad $\bra{\Psi}U^\dagger U\ket{\Psi}=1$ & \ref{sec:portes} \\
Pauli & $X\ket{0}=\ket{1}$, \quad $Z\ket{1}=-\ket{1}$, \quad $X^2=Y^2=Z^2=\mathrm{Id}$ & \ref{sec:pauli} \\
Hadamard & $H\ket{0}=\ket{+}$, \quad $H\ket{1}=\ket{-}$, \quad $H^2=\mathrm{Id}$ & \ref{sec:hadamard} \\
Phase & $P_\theta=\mathrm{diag}(1,\ee^{\ii\theta})$, \quad $S=P_{\pi/2}$, \quad $T=P_{\pi/4}$ & \ref{sec:phase} \\
Circuit & $U_1$ puis $U_2$ : \quad $U_2U_1$ & \ref{sec:circuits} \\
Portes locales & $(U\otimes I)(\ket{u}\otimes\ket{v})=U\ket{u}\otimes\ket{v}$ & \ref{sec:op2} \\
CNOT & $\ket{a,b}\to\ket{a\oplus b,b}$, \quad $\mathrm{CNOT}=I\otimes\ket{0}\bra{0}+X\otimes\ket{1}\bra{1}$
  & \ref{sec:cnot} \\
SWAP & $\mathrm{SWAP}\ket{a,b}=\ket{b,a}$, \quad $\mathrm{SWAP}=\mathrm{CNOT}_{12}\mathrm{CNOT}_{21}\mathrm{CNOT}_{12}$
  & \ref{sec:swap} \\
Contrôlé-$U$ & $\mathrm{C}U=\ket{0}\bra{0}\otimes I+\ket{1}\bra{1}\otimes U=\mathrm{diag}(I,U)$ & \ref{sec:controlees} \\
Toffoli & $\mathrm{CCNOT}\ket{x,y,z}=\ket{x,y,z\oplus xy}$ & \ref{sec:controlees} \\
\bottomrule
\end{tabular}
\end{center}

\begin{savoirfaire}
\begin{itemize}
  \item Appliquer une porte à un état (par la matrice ou par les ket-bra) et identifier sa rotation
        de Bloch (exercices~\ref{ex:conj} et~\ref{ex:phases}).
  \item Lire un circuit, écrire sa matrice dans le bon ordre et suivre l'état à chaque étape
        (exercice~\ref{ex:kickback}).
  \item Écrire la matrice de $U\otimes I$, du CNOT, du SWAP ou d'une porte contrôlée, et calculer
        leur action sur les états de base.
  \item Reconnaître une porte réversible qui calcule une fonction booléenne
        (exercice~\ref{ex:toffoli}).
\end{itemize}
\end{savoirfaire}
.
%  Il ne contient ni préambule ni \begin{document} : les environnements
%  (definition, resultat, remarque, aretenir, savoirfaire, exercice, correction),
%  les macros (\ii, \ee, \dd, \pd, \pdd, \Hop, \ket, \bra, \braket), les symboles
%  de circuits (ctl, cible, croix, mesure, porte) et les couleurs (insarouge,
%  insagris) sont définis dans main.tex. Un chapitre = une \section.
% =====================================================================

% Macros de Dirac (déjà définies dans main.tex ; ce bloc évite l'erreur
% « Undefined control sequence \ket » si main.tex est une ancienne version).
\providecommand{\ket}[1]{|#1\rangle}
\providecommand{\bra}[1]{\langle #1|}
\providecommand{\braket}[2]{\langle #1|#2\rangle}

% =====================================================================
\section{Portes quantiques et systèmes multi-qubits}\label{chap:portes}
% =====================================================================

Un calcul quantique est une suite de transformations unitaires appliquées à des qubits. Ce
chapitre décrit les portes à un qubit (Pauli, Hadamard, phase) et leur effet sur la sphère de
Bloch, la lecture d'un circuit, puis les portes qui agissent sur plusieurs qubits : CNOT, SWAP et
portes contrôlées.

% ---------------------------------------------------------------------
\subsection{Portes quantiques}\label{sec:portes}
% ---------------------------------------------------------------------

Un calcul quantique transforme des états. Les opérations sur un seul qubit sont appelées
\emph{portes} ; les sections suivantes passent à plusieurs qubits.

\begin{resultat}{Postulat 4 : évolution d'un système isolé}
L'évolution d'un système isolé est décrite par un opérateur unitaire : $\ket{\Psi}\to U\ket{\Psi}$
avec $U^\dagger U=\mathrm{Id}$.
\end{resultat}

C'est la forme de l'équation de Schrödinger du chapitre~\ref{chap:histoire}. Pour un hamiltonien
$\Hop$ indépendant du temps, $\ii\hbar\,\partial_t\ket{\Psi}=\Hop\ket{\Psi}$ a pour solution
$\ket{\Psi(t)}=U(t)\ket{\Psi(0)}$ avec $U(t)=\ee^{-\ii\Hop t/\hbar}$. Comme $\Hop$ est hermitien,
$U(t)^\dagger=\ee^{+\ii\Hop t/\hbar}$ et $U^\dagger U=\mathrm{Id}$ : l'évolution est unitaire. Une
porte est une telle évolution, obtenue en contrôlant $\Hop$ pendant une durée donnée.

\begin{definition}{Porte quantique}
Une porte à un qubit est un opérateur $U$ (matrice $2\times2$) \emph{unitaire} :
$U^\dagger U=\mathrm{Id}$. Elle transforme $\ket{\Psi}\to U\ket{\Psi}$, conserve la norme
($\bra{\Psi}U^\dagger U\ket{\Psi}=\braket{\Psi}{\Psi}=1$) et est réversible
($U^{-1}=U^\dagger$). Sur la sphère de Bloch (section~\ref{sec:bloch}), c'est une
\emph{rotation}.
\end{definition}

% ---------------------------------------------------------------------
\subsection{Portes de Pauli}\label{sec:pauli}
% ---------------------------------------------------------------------

Les matrices de Pauli de la section~\ref{sec:op-dirac} sont trois portes.

\begin{definition}{Portes $X$, $Y$, $Z$}
\begin{center}
\renewcommand{\arraystretch}{1}
\begin{tabular}{@{}lccl@{}}
\textbf{Porte} & \textbf{Matrice} & \textbf{Ket-bra} & \textbf{Nom} \\[0.8ex]
$X=\sigma_x$ & $\begin{pmatrix}0 & 1\\ 1 & 0\end{pmatrix}$
  & $\ket{0}\bra{1}+\ket{1}\bra{0}$ & bit flip \\[2.2ex]
$Y=\sigma_y$ & $\begin{pmatrix}0 & -\ii\\ \ii & 0\end{pmatrix}$
  & $-\ii\ket{0}\bra{1}+\ii\ket{1}\bra{0}$ &  \\[2.2ex]
$Z=\sigma_z$ & $\begin{pmatrix}1 & 0\\ 0 & -1\end{pmatrix}$
  & $\ket{0}\bra{0}-\ket{1}\bra{1}$ & phase flip \\
\end{tabular}
\end{center}
\end{definition}

\paragraph{Porte $X$ (inverseur).} Avec $\braket{0}{0}=\braket{1}{1}=1$ et
$\braket{0}{1}=\braket{1}{0}=0$ :
\begin{align*}
  X\ket{0}&=\ket{0}\braket{1}{0}+\ket{1}\braket{0}{0}=0\cdot\ket{0}+1\cdot\ket{1}=\ket{1},\\
  X\ket{1}&=\ket{0}\braket{1}{1}+\ket{1}\braket{0}{1}=1\cdot\ket{0}+0\cdot\ket{1}=\ket{0}.
\end{align*}
$X$ échange $\ket{0}$ et $\ket{1}$ : c'est le NOT quantique (\emph{bit flip}). Sur la
sphère de Bloch, c'est une rotation de $\pi$ autour de l'axe $x$ : le pôle nord devient
le pôle sud. Les états de l'axe $x$ ne bougent pas : $X\ket{+}=\ket{+}$ et
$X\ket{-}=-\ket{-}$.

\paragraph{Porte $Z$ (changement de phase).}
\begin{align*}
  Z\ket{0}&=\ket{0}\braket{0}{0}-\ket{1}\braket{1}{0}=\ket{0},\\
  Z\ket{1}&=\ket{0}\braket{0}{1}-\ket{1}\braket{1}{1}=-\ket{1},
\end{align*}
donc $Z$ change le signe de la composante sur $\ket{1}$ (\emph{phase flip}). Sur les
états de la base $X$, par linéarité :
\begin{align*}
  Z\ket{+}&=\tfrac{1}{\sqrt2}\big(Z\ket{0}+Z\ket{1}\big)=\tfrac{1}{\sqrt2}\big(\ket{0}-\ket{1}\big)=\ket{-},\\
  Z\ket{-}&=\tfrac{1}{\sqrt2}\big(Z\ket{0}-Z\ket{1}\big)=\tfrac{1}{\sqrt2}\big(\ket{0}+\ket{1}\big)=\ket{+}.
\end{align*}
Sous forme matricielle, $Z\ket{+}=\begin{pmatrix}1 & 0\\ 0 & -1\end{pmatrix}
\frac{1}{\sqrt2}\begin{pmatrix}1\\ 1\end{pmatrix}=\frac{1}{\sqrt2}\begin{pmatrix}1\\ -1\end{pmatrix}=\ket{-}$.
$Z$ est une rotation de $\pi$ autour de l'axe $z$ : elle laisse $\ket{0}$ et $\ket{1}$ en
place et échange $\ket{+}$ et $\ket{-}$.

\paragraph{Porte $Y$.} Avec $Y=-\ii\ket{0}\bra{1}+\ii\ket{1}\bra{0}$ :
\begin{align*}
  Y\ket{0}&=-\ii\ket{0}\braket{1}{0}+\ii\ket{1}\braket{0}{0}=\ii\ket{1},\\
  Y\ket{1}&=-\ii\ket{0}\braket{1}{1}+\ii\ket{1}\braket{0}{1}=-\ii\ket{0}.
\end{align*}
Sur les états de la base $X$ :
\begin{align*}
  Y\ket{+}&=\tfrac{1}{\sqrt2}\big(\ii\ket{1}-\ii\ket{0}\big)=-\ii\,\tfrac{1}{\sqrt2}\big(\ket{0}-\ket{1}\big)=-\ii\ket{-},\\
  Y\ket{-}&=\tfrac{1}{\sqrt2}\big(\ii\ket{1}+\ii\ket{0}\big)=\ii\,\tfrac{1}{\sqrt2}\big(\ket{0}+\ket{1}\big)=\ii\ket{+}.
\end{align*}
Un facteur global ($-\ii$ ou $\ii$) ne change pas l'état physique : $Y$ échange bien
$\ket{+}$ et $\ket{-}$. Sur les états de l'axe $y$ :
\begin{align*}
  Y\ket{{+}\ii}&=\tfrac{1}{\sqrt2}\big(Y\ket{0}+\ii Y\ket{1}\big)
   =\tfrac{1}{\sqrt2}\big(\ii\ket{1}+\ket{0}\big)=\ket{{+}\ii},\\
  Y\ket{{-}\ii}&=\tfrac{1}{\sqrt2}\big(Y\ket{0}-\ii Y\ket{1}\big)
   =\tfrac{1}{\sqrt2}\big(\ii\ket{1}-\ket{0}\big)=-\ket{{-}\ii}.
\end{align*}
$Y$ est une rotation de $\pi$ autour de l'axe $y$.

\begin{remarque}[Chaque Pauli fixe son axe]
$Z\ket{0}=\ket{0}$, $Z\ket{1}=-\ket{1}$ ; $X\ket{\pm}=\pm\ket{\pm}$ ; $Y\ket{{\pm}\ii}=\pm\ket{{\pm}\ii}$.
Les bases $Z$, $X$, $Y$ de la section~\ref{sec:bloch} sont les bases propres de
$Z$, $X$, $Y$. Une porte de Pauli laisse donc son propre axe invariant et retourne les
deux autres ; c'est pourquoi les exemples utilisent les états des deux autres axes.
\end{remarque}

% ---------------------------------------------------------------------
\subsection{Porte de Hadamard}\label{sec:hadamard}
% ---------------------------------------------------------------------

\begin{definition}{Porte de Hadamard}
\[
  H=\frac{1}{\sqrt2}\begin{pmatrix}1 & 1\\ 1 & -1\end{pmatrix}
  =\ket{+}\bra{0}+\ket{-}\bra{1}.
\]
\end{definition}

\begin{align*}
  H\ket{0}&=\frac{1}{\sqrt2}\begin{pmatrix}1 & 1\\ 1 & -1\end{pmatrix}\begin{pmatrix}1\\ 0\end{pmatrix}
    =\frac{1}{\sqrt2}\begin{pmatrix}1\\ 1\end{pmatrix}=\ket{+},\\[0.6ex]
  H\ket{1}&=\frac{1}{\sqrt2}\begin{pmatrix}1 & 1\\ 1 & -1\end{pmatrix}\begin{pmatrix}0\\ 1\end{pmatrix}
    =\frac{1}{\sqrt2}\begin{pmatrix}1\\ -1\end{pmatrix}=\ket{-},\\[0.6ex]
  H\ket{+}&=\frac{1}{2}\begin{pmatrix}1 & 1\\ 1 & -1\end{pmatrix}\begin{pmatrix}1\\ 1\end{pmatrix}
    =\frac{1}{2}\begin{pmatrix}2\\ 0\end{pmatrix}=\ket{0},\qquad
  H\ket{-}=\frac{1}{2}\begin{pmatrix}0\\ 2\end{pmatrix}=\ket{1}.
\end{align*}
Comme $H\ket{+}=\ket{0}$ et $H\ket{-}=\ket{1}$, on a $H^2=\mathrm{Id}$ ($H$ est son propre
inverse). Sur la sphère de Bloch, $H$ est une rotation de $\pi$ autour de l'axe
$(x+z)/\sqrt2$ : elle échange l'axe $x$ et l'axe $z$.

\begin{remarque}[À quoi sert $H$ ?]
\begin{itemize}[nosep]
  \item Elle \emph{crée une superposition} à partir d'un état de base :
        $H\ket{0}=\ket{+}$ donne $P(0)=P(1)=\frac12$. C'est la première étape de
        la plupart des algorithmes quantiques.
  \item Elle \emph{change de base} : elle envoie la base $Z$ sur la base $X$ et
        inversement. Mesurer dans la base $X$ revient à appliquer $H$ puis à mesurer
        dans la base $Z$ (figure~\ref{fig:mesurex}).
  \item Elle est réversible ($H^2=\mathrm{Id}$).
\end{itemize}
\end{remarque}

% ---------------------------------------------------------------------
\subsection{Portes de phase}\label{sec:phase}
% ---------------------------------------------------------------------

\begin{definition}{Portes de phase $P_\theta$, $S$ et $T$}
\[
  P_\theta=\begin{pmatrix}1 & 0\\ 0 & \ee^{\ii\theta}\end{pmatrix}
  =\ket{0}\bra{0}+\ee^{\ii\theta}\ket{1}\bra{1},
  \qquad
  S=P_{\pi/2}=\begin{pmatrix}1 & 0\\ 0 & \ii\end{pmatrix},
  \qquad
  T=P_{\pi/4}=\begin{pmatrix}1 & 0\\ 0 & \dfrac{1+\ii}{\sqrt2}\end{pmatrix}.
\]
\end{definition}

Ici $\ee^{\ii\pi/2}=\ii$ et $\ee^{\ii\pi/4}=\cos\frac\pi4+\ii\sin\frac\pi4=\frac{1+\ii}{\sqrt2}$.
$P_\theta$ laisse $\ket{0}$ inchangé et multiplie $\ket{1}$ par $\ee^{\ii\theta}$ :
\[
  P_\theta\big(\alpha\ket{0}+\beta\ket{1}\big)=\alpha\ket{0}+\ee^{\ii\theta}\beta\ket{1}.
\]
Dans la paramétrisation de Bloch, $\varphi\to\varphi+\theta$ : $P_\theta$ est une rotation
d'angle $\theta$ autour de l'axe $z$. On retrouve $P_\pi=Z$, et $S^2=Z$, $T^2=S$.

\paragraph{Exemples.}
\begin{align*}
  S\ket{+}&=\tfrac{1}{\sqrt2}\big(\ket{0}+\ii\ket{1}\big)=\ket{{+}\ii},\\
  S\ket{{+}\ii}&=\tfrac{1}{\sqrt2}\big(\ket{0}+\ii\cdot\ii\ket{1}\big)=\tfrac{1}{\sqrt2}\big(\ket{0}-\ket{1}\big)=\ket{-},\\
  T\ket{+}&=\tfrac{1}{\sqrt2}\begin{pmatrix}1 & 0\\ 0 & \frac{1+\ii}{\sqrt2}\end{pmatrix}\begin{pmatrix}1\\ 1\end{pmatrix}
    =\tfrac{1}{\sqrt2}\begin{pmatrix}1\\ \frac{1+\ii}{\sqrt2}\end{pmatrix}
    =\tfrac{1}{\sqrt2}\big(\ket{0}+\ee^{\ii\pi/4}\ket{1}\big)\quad\big(\text{angle polaire }\tfrac\pi2,\ \text{azimut }\tfrac\pi4\big).
\end{align*}
$S$ tourne donc $\ket{+}$ d'un quart de tour autour de $z$ (de l'axe $x$ à l'axe $y$),
puis d'un second quart de tour jusqu'à $\ket{-}$ ; $T$ fait un huitième de tour.

\begin{remarque}[Une porte de phase est une évolution libre]
Pour $\Hop=\frac{\hbar\omega}{2}Z$, $U(t)=\ee^{-\ii\Hop t/\hbar}=\mathrm{diag}\big(\ee^{-\ii\omega t/2},\ee^{\ii\omega t/2}\big)
=\ee^{-\ii\omega t/2}\,P_{\omega t}$. Au facteur global près, laisser évoluer le qubit pendant la
durée $t$ applique la porte de phase $P_{\omega t}$ : une rotation d'angle $\omega t$ autour de
l'axe $z$.
\end{remarque}

% ---------------------------------------------------------------------
\subsection{Portes et sphère de Bloch}\label{sec:portes-bloch}
% ---------------------------------------------------------------------

Toute porte à un qubit est une rotation de la sphère de Bloch. La figure~\ref{fig:portes} montre
les axes des portes usuelles et le tableau qui suit les résume.

\begin{figure}[!ht]
\centering
\begin{tikzpicture}[x={(-1.63cm,-0.56cm)}, y={(1.63cm,-0.56cm)}, z={(0cm,2.16cm)}, >=Stealth, scale=0.9]
  % sphère et grands cercles
  \shade[ball color=insagris!15, opacity=0.55] (0,0,0) circle[radius=2.3cm];
  \begin{scope}[insagris!75, thin]
    \draw[domain=-45:135, samples=60, variable=\t] plot ({cos(\t)},{sin(\t)},0);
    \draw[domain=135:315, samples=60, variable=\t, dashed] plot ({cos(\t)},{sin(\t)},0);
    \draw[domain=-62.7:117.3, samples=60, variable=\t] plot ({cos(\t)},0,{sin(\t)});
    \draw[domain=117.3:297.3, samples=60, variable=\t, dashed] plot ({cos(\t)},0,{sin(\t)});
    \draw[domain=-62.7:117.3, samples=60, variable=\t] plot (0,{cos(\t)},{sin(\t)});
    \draw[domain=117.3:297.3, samples=60, variable=\t, dashed] plot (0,{cos(\t)},{sin(\t)});
  \end{scope}
  % axes x, y, z (portes X, Y, Z)
  \draw[insagris, dashed] (0,0,0) -- (-1,0,0);
  \draw[insagris, dashed] (0,0,0) -- (0,-1,0);
  \draw[insagris, dashed] (0,0,0) -- (0,0,-1);
  \draw[insagris] (-1,0,0) -- (-1.3,0,0);
  \draw[insagris] (0,-1,0) -- (0,-1.3,0);
  \draw[insagris] (0,0,-1) -- (0,0,-1.3);
  \draw[insagris, -{Stealth[length=2.2mm]}] (0,0,0) -- (1.3,0,0);
  \draw[insagris, -{Stealth[length=2.2mm]}] (0,0,0) -- (0,1.3,0);
  \draw[insagris, -{Stealth[length=2.2mm]}] (0,0,0) -- (0,0,1.3);
  \node[anchor=east]  at (1.33,0,0) {$x$};
  \node[anchor=west]  at (0,1.33,0) {$y$};
  \node[anchor=south] at (0,0,1.33) {$z$};
  % axe de la porte H : (x+z)/sqrt(2)
  \draw[insarouge, dashed] (0,0,0) -- (-0.707,0,-0.707);
  \draw[insarouge, very thick, -{Stealth[length=2.6mm]}] (0,0,0) -- (0.98,0,0.98);
  \node[insarouge, anchor=south east] at (0.98,0,0.98) {axe de $H$};
  % rotation autour de z (portes de phase)
  \draw[insarouge, thick, -{Stealth[length=2.2mm]}, domain=15:105, samples=30, variable=\t]
    plot ({1.15*cos(\t)},{1.15*sin(\t)},0);
  \node[insarouge, anchor=north west, font=\small] at ({1.15*cos(60)},{1.15*sin(60)},0) {$P_\theta$};
\end{tikzpicture}
\caption{Portes et rotations sur la sphère de Bloch. Les portes $X$, $Y$, $Z$ sont des
rotations de $\pi$ autour des axes $x$, $y$, $z$ ; $H$ est une rotation de $\pi$ autour de
l'axe $(x+z)/\sqrt2$ (rouge) ; $P_\theta$ (donc $Z$, $S$, $T$) tourne d'un angle $\theta$
autour de $z$.}
\label{fig:portes}
\end{figure}

\begin{center}
\renewcommand{\arraystretch}{1}
\begin{tabular}{@{}lccc@{}}
\toprule
\textbf{Porte} & \textbf{Matrice} & \textbf{Effet sur la base} & \textbf{Rotation de Bloch} \\
\midrule
$X$ & $\begin{pmatrix}0 & 1\\ 1 & 0\end{pmatrix}$ & $\ket{0}\leftrightarrow\ket{1}$ & $\pi$ autour de $x$ \\[2.2ex]
$Y$ & $\begin{pmatrix}0 & -\ii\\ \ii & 0\end{pmatrix}$ & $\ket{0}\to\ii\ket{1},\ \ket{1}\to-\ii\ket{0}$ & $\pi$ autour de $y$ \\[2.2ex]
$Z$ & $\begin{pmatrix}1 & 0\\ 0 & -1\end{pmatrix}$ & $\ket{1}\to-\ket{1},\ \ket{\pm}\leftrightarrow\ket{\mp}$ & $\pi$ autour de $z$ \\[2.2ex]
$H$ & $\dfrac{1}{\sqrt2}\begin{pmatrix}1 & 1\\ 1 & -1\end{pmatrix}$ & $\ket{0}\leftrightarrow\ket{+},\ \ket{1}\leftrightarrow\ket{-}$ & $\pi$ autour de $\frac{x+z}{\sqrt2}$ \\[2.2ex]
$S$ & $\begin{pmatrix}1 & 0\\ 0 & \ii\end{pmatrix}$ & $\ket{1}\to\ii\ket{1}$ & $\frac\pi2$ autour de $z$ \\[2.2ex]
$T$ & $\begin{pmatrix}1 & 0\\ 0 & \frac{1+\ii}{\sqrt2}\end{pmatrix}$ & $\ket{1}\to\ee^{\ii\pi/4}\ket{1}$ & $\frac\pi4$ autour de $z$ \\
\bottomrule
\end{tabular}
\end{center}

\begin{figure}[!ht]
\centering
\begin{tikzpicture}[x={(-1.63cm,-0.56cm)}, y={(1.63cm,-0.56cm)}, z={(0cm,2.16cm)}, >=Stealth, scale=0.9]
  % sphère et grands cercles
  \shade[ball color=insagris!15, opacity=0.55] (0,0,0) circle[radius=2.3cm];
  \begin{scope}[insagris!75, thin]
    \draw[domain=-45:135, samples=60, variable=\t] plot ({cos(\t)},{sin(\t)},0);
    \draw[domain=135:315, samples=60, variable=\t, dashed] plot ({cos(\t)},{sin(\t)},0);
    \draw[domain=-62.7:117.3, samples=60, variable=\t] plot ({cos(\t)},0,{sin(\t)});
    \draw[domain=117.3:297.3, samples=60, variable=\t, dashed] plot ({cos(\t)},0,{sin(\t)});
    \draw[domain=-62.7:117.3, samples=60, variable=\t] plot (0,{cos(\t)},{sin(\t)});
    \draw[domain=117.3:297.3, samples=60, variable=\t, dashed] plot (0,{cos(\t)},{sin(\t)});
  \end{scope}
  % axes x, y, z (portes X, Y, Z)
  \draw[insagris, dashed] (0,0,0) -- (-1,0,0);
  \draw[insagris, dashed] (0,0,0) -- (0,-1,0);
  \draw[insagris, dashed] (0,0,0) -- (0,0,-1);
  \draw[insagris] (-1,0,0) -- (-1.3,0,0);
  \draw[insagris] (0,-1,0) -- (0,-1.3,0);
  \draw[insagris] (0,0,-1) -- (0,0,-1.3);
  \draw[insagris, -{Stealth[length=2.2mm]}] (0,0,0) -- (1.3,0,0);
  \draw[insagris, -{Stealth[length=2.2mm]}] (0,0,0) -- (0,1.3,0);
  \draw[insagris, -{Stealth[length=2.2mm]}] (0,0,0) -- (0,0,1.3);
  % trajet : |0> --H--> |+> --S--> |+i>
  \draw[insarouge, very thick, -{Stealth[length=2.6mm]}, domain=3:87, samples=30, variable=\t]
    plot ({sin(\t)},0,{cos(\t)});
  \draw[insarouge, very thick, -{Stealth[length=2.6mm]}, domain=3:87, samples=30, variable=\t]
    plot ({cos(\t)},{sin(\t)},0);
  \foreach \ax/\ay/\az in {0/0/1, 1/0/0, 0/1/0} \fill[insarouge] (\ax,\ay,\az) circle[radius=2pt];
  \node[insarouge, anchor=south] at (0,0,1.33) {$\ket{0}$};
  \node[insarouge, anchor=east] at (1.33,0,0) {$\ket{+}$};
  \node[insarouge, anchor=west] at (0,1.33,0) {$\ket{{+}\ii}$};
  \node[insarouge, font=\small, anchor=south west] at ({sin(45)},0,{cos(45)}) {$H$};
  \node[insarouge, font=\small, anchor=north west] at ({cos(45)},{sin(45)},0) {$S$};
\end{tikzpicture}
\caption{Trajet d'un état sous le circuit $H$ puis $S$ (flèches schématiques). $H$ envoie
$\ket{0}$ sur $\ket{+}$ ; $S$, rotation d'un quart de tour autour de $z$, envoie $\ket{+}$ sur
$\ket{{+}\ii}$. Les trois états sont alignés sur les axes $z$, $x$ et $y$.}
\label{fig:trajet}
\end{figure}

% ---------------------------------------------------------------------
\subsection{Circuits quantiques}\label{sec:circuits}
% ---------------------------------------------------------------------

Un \emph{circuit} représente un calcul : chaque qubit est un fil horizontal, le temps va de gauche
à droite, et chaque porte est un symbole posé sur les fils qu'elle concerne. La
figure~\ref{fig:symboles} donne les symboles utilisés dans ce document.

\begin{figure}[!ht]
\centering
\begin{adjustbox}{max width=\linewidth}
\begin{tikzpicture}[>=Stealth, font=\small,
    nom/.style={font=\footnotesize, align=center, anchor=north}]
  % (a) porte à un qubit
  \draw (0,0) -- (1.8,0);
  \node[porte] at (0.9,0) {$U$};
  \node[nom] at (0.9,-1.15) {porte $U$\\ à un qubit};
  % (b) CNOT
  \begin{scope}[xshift=3.2cm]
    \draw (0,0.6) -- (1.8,0.6);
    \draw (0,-0.6) -- (1.8,-0.6);
    \draw[insarouge, thick] (0.9,0.6) -- (0.9,-0.6);
    \pic at (0.9,0.6) {ctl};
    \pic at (0.9,-0.6) {cible};
    \node[nom] at (0.9,-1.15) {CNOT\\ (contrôle $\bullet$, cible $\oplus$)};
  \end{scope}
  % (c) SWAP
  \begin{scope}[xshift=6.4cm]
    \draw (0,0.6) -- (1.8,0.6);
    \draw (0,-0.6) -- (1.8,-0.6);
    \draw[insarouge, thick] (0.9,0.6) -- (0.9,-0.6);
    \pic at (0.9,0.6) {croix};
    \pic at (0.9,-0.6) {croix};
    \node[nom] at (0.9,-1.15) {SWAP\\ (échange)};
  \end{scope}
  % (d) contrôlé-U
  \begin{scope}[xshift=9.6cm]
    \draw (0,0.6) -- (1.8,0.6);
    \draw (0,-0.6) -- (1.8,-0.6);
    \draw[insarouge, thick] (0.9,-0.6) -- (0.9,0.3);
    \pic at (0.9,-0.6) {ctl};
    \node[porte] at (0.9,0.6) {$U$};
    \node[nom] at (0.9,-1.15) {contrôlé-$U$\\ (contrôle en bas)};
  \end{scope}
  % (e) mesure
  \begin{scope}[xshift=12.8cm]
    \draw (0,0) -- (0.6,0);
    \pic at (0.9,0) {mesure};
    \draw (1.2,0.03) -- (2.0,0.03);
    \draw (1.2,-0.03) -- (2.0,-0.03);
    \node[nom] at (1.0,-1.15) {mesure\\ (fil double : bit)};
  \end{scope}
\end{tikzpicture}
\end{adjustbox}
\caption{Symboles des circuits quantiques. Le point noir $\bullet$ marque un qubit de contrôle,
le symbole $\oplus$ une cible qui subit $X$, les deux croix un échange. Après une mesure, le fil
double transporte un bit classique.}
\label{fig:symboles}
\end{figure}

\begin{remarque}[Conventions de lecture]
\begin{itemize}[nosep]
  \item Le temps va de gauche à droite. Deux portes successives $U_1$ puis $U_2$ se multiplient
        dans l'ordre inverse : la matrice du circuit est $U_2U_1$, car $U_1$ agit d'abord sur le ket.
  \item Avec plusieurs fils, le fil du \emph{haut} porte le \emph{dernier} chiffre du ket et le
        fil du \emph{bas} le premier : si $A$ est le premier qubit et $B$ le second,
        $\ket{a,b}=\ket{a}_A\otimes\ket{b}_B$ et $B$ est dessiné au-dessus de $A$.
  \item Deux portes placées dans la même colonne, sur deux fils, agissent en parallèle : c'est un
        produit tensoriel $U_A\otimes U_B$ (section~\ref{sec:op2}).
\end{itemize}
\end{remarque}

\paragraph{Exemple.} Le circuit de la figure~\ref{fig:circuit-hs} applique $H$ puis $S$ à
$\ket{0}$. Sa matrice est $SH$ :
\[
  SH=\begin{pmatrix}1&0\\0&\ii\end{pmatrix}\frac{1}{\sqrt2}\begin{pmatrix}1&1\\1&-1\end{pmatrix}
  =\frac{1}{\sqrt2}\begin{pmatrix}1&1\\ \ii&-\ii\end{pmatrix},
  \qquad
  SH\ket{0}=\frac{1}{\sqrt2}\begin{pmatrix}1\\ \ii\end{pmatrix}=\ket{{+}\ii}.
\]
Le circuit prépare donc l'état $\ket{{+}\ii}$ de la base $Y$ à partir de $\ket{0}$
(figure~\ref{fig:trajet}). L'ordre compte : $HS\ket{0}=H\ket{0}=\ket{+}$, car $S\ket{0}=\ket{0}$.

\begin{figure}[!ht]
\centering
\begin{tikzpicture}[>=Stealth, font=\small]
  \draw (0,0) -- (6.6,0);
  \node[anchor=east] at (0,0) {$\ket{0}$};
  \node[porte] at (2.0,0) {$H$};
  \node[porte] at (4.3,0) {$S$};
  \node[anchor=west] at (6.6,0) {$\ket{{+}\ii}$};
  \foreach \x in {0.7, 3.15, 5.5} \draw[dashed, insagris!70] (\x,0.6) -- (\x,-0.95);
  \node[font=\footnotesize, anchor=north] at (0.7,-0.95) {$\ket{0}$};
  \node[font=\footnotesize, anchor=north] at (3.15,-0.95) {$\ket{+}$};
  \node[font=\footnotesize, anchor=north] at (5.5,-0.95) {$\ket{{+}\ii}$};
\end{tikzpicture}
\caption{Circuit $H$ puis $S$ appliqué à $\ket{0}$. Les états intermédiaires sont indiqués sous
les traits pointillés ; la matrice du circuit est $SH$.}
\label{fig:circuit-hs}
\end{figure}

\paragraph{Mesure.} Le symbole de mesure (compteur) donne le résultat $0$ ou $1$ dans la base $Z$.
Mesurer dans la base $X$ revient à appliquer $H$ puis à mesurer dans la base $Z$
(figure~\ref{fig:mesurex}) : comme $H\ket{\pm}=\ket{0}$ ou $\ket{1}$, le résultat $0$ correspond à
$\ket{+}$ et le résultat $1$ à $\ket{-}$, avec $P(\pm)=|\braket{\pm}{\Psi}|^2$.

\begin{figure}[!ht]
\centering
\begin{tikzpicture}[>=Stealth, font=\small]
  % (a) base Z
  \draw (0,0) -- (1.2,0);
  \node[anchor=east] at (0,0) {$\ket{\Psi}$};
  \pic at (1.5,0) {mesure};
  \draw (1.8,0.03) -- (2.6,0.03);
  \draw (1.8,-0.03) -- (2.6,-0.03);
  \node[anchor=west] at (2.7,0) {$0$ ou $1$};
  \node[font=\footnotesize] at (1.5,-0.95) {(a) mesure en base $Z$};
  % (b) base X
  \begin{scope}[xshift=7.5cm]
    \draw (0,0) -- (2.6,0);
    \node[anchor=east] at (0,0) {$\ket{\Psi}$};
    \node[porte] at (1.0,0) {$H$};
    \pic at (2.3,0) {mesure};
    \draw (2.6,0.03) -- (3.4,0.03);
    \draw (2.6,-0.03) -- (3.4,-0.03);
    \node[anchor=west] at (3.5,0) {$+$ ou $-$};
    \node[font=\footnotesize] at (1.9,-0.95) {(b) mesure en base $X$};
  \end{scope}
\end{tikzpicture}
\caption{(a) Mesure dans la base $Z$. (b) Mesure dans la base $X$ : porte de Hadamard, puis
mesure dans la base $Z$ ; les résultats $0$ et $1$ se lisent $+$ et $-$.}
\label{fig:mesurex}
\end{figure}

% ---------------------------------------------------------------------
\subsection{Portes locales sur deux qubits}\label{sec:op2}
% ---------------------------------------------------------------------

Une porte appliquée à un seul qubit d'un système de deux s'écrit avec un produit tensoriel
(section~\ref{sec:tenseur}) : $U\otimes I$ agit sur $A$ (premier facteur, fil du bas), $I\otimes U$
agit sur $B$, et $U_A\otimes U_B$ agit sur les deux en parallèle (figure~\ref{fig:local}).

\begin{figure}[!ht]
\centering
\begin{tikzpicture}[>=Stealth, font=\small]
  \foreach \xs/\haut/\bas/\tit in {0/{}/{$X$}/{(a) $X\otimes I$},
                                   5.4/{$X$}/{}/{(b) $I\otimes X$},
                                   10.8/{$H$}/{$H$}/{(c) $H\otimes H$}} {
    \begin{scope}[xshift=\xs cm]
      \draw (0,0.75) -- (2.6,0.75);
      \draw (0,0) -- (2.6,0);
      \node[anchor=east] at (0,0.75) {$B$};
      \node[anchor=east] at (0,0) {$A$};
      \ifx\haut\empty\else \node[porte] at (1.3,0.75) {\haut}; \fi
      \ifx\bas\empty\else \node[porte] at (1.3,0) {\bas}; \fi
      \node[font=\footnotesize] at (1.3,-0.85) {\tit};
    \end{scope}
  }
\end{tikzpicture}
\caption{Portes en parallèle sur deux qubits. Le qubit $A$, premier facteur du produit tensoriel,
est sur le fil du bas. (a) $X$ sur $A$ seulement. (b) $X$ sur $B$ seulement. (c) $H$ sur les deux.}
\label{fig:local}
\end{figure}

\subsubsection{Portes agissant sur un seul des deux qubits}

On note $I=\mathrm{Id}$ ($2\times2$). $X\otimes I$ agit sur $A$, $I\otimes X$ agit sur $B$ :
\begin{align*}
  I\otimes X&=\begin{pmatrix}1\cdot X & 0\cdot X\\ 0\cdot X & 1\cdot X\end{pmatrix}
    =\begin{pmatrix}0 & 1 & 0 & 0\\ 1 & 0 & 0 & 0\\ 0 & 0 & 0 & 1\\ 0 & 0 & 1 & 0\end{pmatrix},\\[1ex]
  X\otimes I&=\begin{pmatrix}0\cdot I & 1\cdot I\\ 1\cdot I & 0\cdot I\end{pmatrix}
    =\begin{pmatrix}0 & 0 & 1 & 0\\ 0 & 0 & 0 & 1\\ 1 & 0 & 0 & 0\\ 0 & 1 & 0 & 0\end{pmatrix}.
\end{align*}
Sur $\ket{00}=(1,0,0,0)^{\mathsf T}$ (première colonne de la matrice) :
$(I\otimes X)\ket{00}=(0,1,0,0)^{\mathsf T}=\ket{01}$ et
$(X\otimes I)\ket{00}=(0,0,1,0)^{\mathsf T}=\ket{10}$. Avec la règle :
$(I\otimes X)\ket{00}=I\ket{0}\otimes X\ket{0}=\ket{0}\otimes\ket{1}=\ket{01}$.

Pour la porte de Hadamard, avec $H=\frac{1}{\sqrt2}\begin{pmatrix}1 & 1\\ 1 & -1\end{pmatrix}$ :
\[
  H\otimes I=\frac{1}{\sqrt2}\begin{pmatrix}1\cdot I & 1\cdot I\\ 1\cdot I & -1\cdot I\end{pmatrix}
  =\frac{1}{\sqrt2}\begin{pmatrix}1 & 0 & 1 & 0\\ 0 & 1 & 0 & 1\\ 1 & 0 & -1 & 0\\ 0 & 1 & 0 & -1\end{pmatrix}.
\]
Ses colonnes donnent l'image des quatre états de base :
\begin{align*}
  (H\otimes I)\ket{00}&=\tfrac{1}{\sqrt2}(1,0,1,0)^{\mathsf T}=\ket{+}\otimes\ket{0},&
  (H\otimes I)\ket{01}&=\tfrac{1}{\sqrt2}(0,1,0,1)^{\mathsf T}=\ket{+}\otimes\ket{1},\\
  (H\otimes I)\ket{10}&=\tfrac{1}{\sqrt2}(1,0,-1,0)^{\mathsf T}=\ket{-}\otimes\ket{0},&
  (H\otimes I)\ket{11}&=\tfrac{1}{\sqrt2}(0,1,0,-1)^{\mathsf T}=\ket{-}\otimes\ket{1}.
\end{align*}

\subsubsection{Hadamard sur les deux qubits}

\[
  H\otimes H=\frac12\begin{pmatrix}1\cdot H' & 1\cdot H'\\ 1\cdot H' & -1\cdot H'\end{pmatrix},
  \quad H'=\begin{pmatrix}1 & 1\\ 1 & -1\end{pmatrix},
  \qquad
  H\otimes H=\frac12\begin{pmatrix}1 & 1 & 1 & 1\\ 1 & -1 & 1 & -1\\ 1 & 1 & -1 & -1\\ 1 & -1 & -1 & 1\end{pmatrix}.
\]
La première colonne donne $(H\otimes H)\ket{00}=\frac12(1,1,1,1)^{\mathsf T}=\ket{+}\otimes\ket{+}$
(cohérent avec $H\ket{0}\otimes H\ket{0}$) : les quatre résultats $00,01,10,11$ sont
équiprobables, chacun avec la probabilité $\frac14$. Sur $n$ qubits, $H^{\otimes n}\ket{0\ldots0}$
donne une superposition uniforme de $2^n$ états.

% ---------------------------------------------------------------------
\subsection{Porte CNOT}\label{sec:cnot}
% ---------------------------------------------------------------------

\begin{definition}{Porte CNOT (NOT contrôlé)}
La porte CNOT agit sur deux qubits : le qubit $B$ est le \emph{contrôle}, le qubit $A$ est la
\emph{cible}. La cible est inversée (porte $X$) si le contrôle vaut $1$ et reste inchangée
s'il vaut $0$. Avec $\ket{a,b}=\ket{a}_A\otimes\ket{b}_B$ :
\[
  \mathrm{CNOT}\,\ket{a,b}=\ket{a\oplus b,\,b},
  \qquad
  \mathrm{CNOT}=I\otimes\ket{0}\bra{0}+X\otimes\ket{1}\bra{1},
\]
où $\oplus$ est l'addition modulo $2$ ($0\oplus0=1\oplus1=0$ et $0\oplus1=1\oplus0=1$).
\end{definition}

Sur le circuit de la figure~\ref{fig:cnot}, le point noir posé sur le fil de $B$ (en haut) marque
le contrôle et le symbole $\oplus$ posé sur le fil de $A$ (en bas) marque la cible.

\begin{figure}[!ht]
\centering
\begin{tikzpicture}[>=Stealth, font=\small]
  \draw (0,1.2) -- (6,1.2);
  \draw (0,0) -- (6,0);
  \node[anchor=east] at (0,1.2) {$\ket{b}_B$};
  \node[anchor=east] at (0,0) {$\ket{a}_A$};
  \node[anchor=west] at (6,1.2) {$\ket{b}_B$};
  \node[anchor=west] at (6,0) {$\ket{a\oplus b}_A$};
  \draw[insarouge, thick] (3,1.2) -- (3,0);
  \pic at (3,1.2) {ctl};
  \pic at (3,0) {cible};
  \node[anchor=south, text=insagris] at (3,1.5) {qubit de contrôle};
  \node[anchor=north, text=insagris] at (3,-0.5) {qubit cible};
\end{tikzpicture}
\caption{Porte CNOT : le fil du haut ($B$) est le contrôle, le fil du bas ($A$) est la cible. Le
ket s'écrit $\ket{a,b}$ : le fil du haut donne le dernier chiffre.}
\label{fig:cnot}
\end{figure}

\subsubsection{Table de vérité}

\begin{center}
\renewcommand{\arraystretch}{1.3}
\begin{tabular}{@{}cccc@{}}
\toprule
Entrée $\ket{a,b}$ & Contrôle $b$ & Sortie $\ket{a\oplus b,b}$ & Effet sur la cible $A$ \\
\midrule
$\ket{00}$ & $0$ & $\ket{00}$ & aucun \\
$\ket{01}$ & $1$ & $\ket{11}$ & $0\to1$ \\
$\ket{10}$ & $0$ & $\ket{10}$ & aucun \\
$\ket{11}$ & $1$ & $\ket{01}$ & $1\to0$ \\
\bottomrule
\end{tabular}
\end{center}

\subsubsection{Matrice du CNOT}

La $j$-ième colonne d'une matrice $4\times4$ est l'image du $j$-ième vecteur de base
($\ket{00},\ket{01},\ket{10},\ket{11}$, dans cet ordre). La table donne :
\begin{align*}
  \mathrm{CNOT}\ket{00}&=\ket{00}=\begin{pmatrix}1\\0\\0\\0\end{pmatrix},&
  \mathrm{CNOT}\ket{01}&=\ket{11}=\begin{pmatrix}0\\0\\0\\1\end{pmatrix},\\[1.5ex]
  \mathrm{CNOT}\ket{10}&=\ket{10}=\begin{pmatrix}0\\0\\1\\0\end{pmatrix},&
  \mathrm{CNOT}\ket{11}&=\ket{01}=\begin{pmatrix}0\\1\\0\\0\end{pmatrix},
\end{align*}
d'où, en juxtaposant ces quatre colonnes,
\[
  \mathrm{CNOT}=\begin{pmatrix}
    1 & 0 & 0 & 0\\ 0 & 0 & 0 & 1\\ 0 & 0 & 1 & 0\\ 0 & 1 & 0 & 0
  \end{pmatrix}.
\]

On retrouve cette matrice avec la formule $I\otimes\ket{0}\bra{0}+X\otimes\ket{1}\bra{1}$ et
la règle $A\otimes B=(A_{ij}B)$ (section~\ref{sec:tenseur}), où
$\ket{0}\bra{0}=\begin{pmatrix}1&0\\0&0\end{pmatrix}$ et
$\ket{1}\bra{1}=\begin{pmatrix}0&0\\0&1\end{pmatrix}$ :
\[
  I\otimes\ket{0}\bra{0}=\begin{pmatrix}1&0&0&0\\0&0&0&0\\0&0&1&0\\0&0&0&0\end{pmatrix},
  \qquad
  X\otimes\ket{1}\bra{1}=\begin{pmatrix}0&0&0&0\\0&0&0&1\\0&0&0&0\\0&1&0&0\end{pmatrix},
\]
et la somme de ces deux matrices est la matrice du CNOT.

\subsubsection{Action sur un état quelconque}

Pour $\ket{\Psi}=c_{00}\ket{00}+c_{01}\ket{01}+c_{10}\ket{10}+c_{11}\ket{11}$ :
\[
  \mathrm{CNOT}\begin{pmatrix}c_{00}\\c_{01}\\c_{10}\\c_{11}\end{pmatrix}
  =\begin{pmatrix}
    1\cdot c_{00}+0\cdot c_{01}+0\cdot c_{10}+0\cdot c_{11}\\
    0\cdot c_{00}+0\cdot c_{01}+0\cdot c_{10}+1\cdot c_{11}\\
    0\cdot c_{00}+0\cdot c_{01}+1\cdot c_{10}+0\cdot c_{11}\\
    0\cdot c_{00}+1\cdot c_{01}+0\cdot c_{10}+0\cdot c_{11}
  \end{pmatrix}
  =\begin{pmatrix}c_{00}\\c_{11}\\c_{10}\\c_{01}\end{pmatrix}.
\]
Le CNOT échange les amplitudes de $\ket{01}$ et de $\ket{11}$, c'est-à-dire des deux états
dont le contrôle vaut $1$, et laisse les autres.

\paragraph{Exemple : contrôle en superposition.} Soit $\ket{0}_A\otimes\ket{+}_B
=\frac{1}{\sqrt2}\big(\ket{00}+\ket{01}\big)=\frac{1}{\sqrt2}(1,1,0,0)^{\mathsf T}$,
un état séparable. Alors
\[
  \mathrm{CNOT}\,\frac{1}{\sqrt2}\begin{pmatrix}1\\1\\0\\0\end{pmatrix}
  =\frac{1}{\sqrt2}\begin{pmatrix}1\\0\\0\\1\end{pmatrix}
  =\frac{1}{\sqrt2}\big(\ket{00}+\ket{11}\big).
\]
Le résultat est intriqué ($c_{00}c_{11}-c_{01}c_{10}=\frac12\neq0$ ; section~\ref{sec:separable}) :
un CNOT dont le contrôle est en superposition crée de l'intrication. C'est l'état de Bell
$\ket{\Phi^+}$ du chapitre~\ref{chap:bell}.

\paragraph{Réversibilité.} Deux CNOT successifs redonnent l'état initial :
$\ket{a,b}\to\ket{a\oplus b,b}\to\ket{a\oplus b\oplus b,b}=\ket{a,b}$, donc
$\mathrm{CNOT}^2=\mathrm{Id}$. La matrice est réelle et symétrique, donc
$\mathrm{CNOT}^\dagger\mathrm{CNOT}=\mathrm{CNOT}^2=\mathrm{Id}$ : elle est bien unitaire.

\begin{remarque}[CNOT et copie d'un bit]
Avec la cible dans l'état $\ket{0}$, $\mathrm{CNOT}\ket{0,b}=\ket{b,b}$ pour $b=0$ et $b=1$ : le CNOT
copie le bit classique $b$. En revanche $\mathrm{CNOT}\big(\ket{0}\otimes\ket{+}\big)$ est l'état
ci-dessus, et non $\ket{+}\otimes\ket{+}$ : un état quantique inconnu ne peut pas être copié
(théorème de non-clonage, section~\ref{sec:teleportation}).
\end{remarque}

\begin{remarque}[Contrôle en dernier ou en premier]
La matrice dépend de l'ordre d'écriture du ket. Notons $\mathrm{CNOT}_{21}$ le CNOT dont le
contrôle est le \emph{second} chiffre du ket : c'est la matrice ci-dessus, adoptée dans ce
document (fil du haut). Beaucoup d'ouvrages mettent le contrôle en \emph{premier} ; le CNOT
correspondant est
\[
  \mathrm{CNOT}_{12}=\ket{0}\bra{0}\otimes I+\ket{1}\bra{1}\otimes X
  =\begin{pmatrix}1&0&0&0\\0&1&0&0\\0&0&0&1\\0&0&1&0\end{pmatrix}.
\]
Le même circuit a donc deux matrices, selon que le contrôle est écrit en dernier ou en premier.
Les deux sont reliées par le SWAP (section~\ref{sec:swap}). La section~\ref{sec:controlees} sur les
portes contrôlées met le contrôle en premier, avec la matrice $\mathrm{diag}(I,U)$.
\end{remarque}

% ---------------------------------------------------------------------
\subsection{Porte SWAP}\label{sec:swap}
% ---------------------------------------------------------------------

\begin{definition}{Porte SWAP (échange)}
La porte SWAP échange les états des deux qubits :
\[
  \mathrm{SWAP}\big(\ket{\psi}\otimes\ket{\phi}\big)=\ket{\phi}\otimes\ket{\psi}.
\]
Sur les états de base, $\mathrm{SWAP}\ket{a,b}=\ket{b,a}$ ; on prolonge par linéarité.
\end{definition}

Son symbole est formé de deux croix reliées par un trait vertical
(figure~\ref{fig:swap}).

\begin{figure}[!ht]
\centering
\begin{tikzpicture}[>=Stealth, font=\small]
  \draw (0,1.2) -- (6,1.2);
  \draw (0,0) -- (6,0);
  \node[anchor=east] at (0,1.2) {$\ket{a}$};
  \node[anchor=east] at (0,0) {$\ket{b}$};
  \node[anchor=west] at (6,1.2) {$\ket{b}$};
  \node[anchor=west] at (6,0) {$\ket{a}$};
  \draw[insarouge, thick] (3,1.2) -- (3,0);
  \pic at (3,1.2) {croix};
  \pic at (3,0) {croix};
\end{tikzpicture}
\caption{Porte SWAP : les deux qubits échangent leurs états.}
\label{fig:swap}
\end{figure}

\subsubsection{Table et matrice}

\begin{center}
\renewcommand{\arraystretch}{1.3}
\begin{tabular}{@{}cc@{}}
\toprule
Entrée & Sortie \\
\midrule
$\ket{00}$ & $\ket{00}$ \\
$\ket{01}$ & $\ket{10}$ \\
$\ket{10}$ & $\ket{01}$ \\
$\ket{11}$ & $\ket{11}$ \\
\bottomrule
\end{tabular}
\end{center}

Les images des quatre vecteurs de base sont les colonnes de la matrice :
\[
  \mathrm{SWAP}\ket{00}=\begin{pmatrix}1\\0\\0\\0\end{pmatrix},\quad
  \mathrm{SWAP}\ket{01}=\begin{pmatrix}0\\0\\1\\0\end{pmatrix},\quad
  \mathrm{SWAP}\ket{10}=\begin{pmatrix}0\\1\\0\\0\end{pmatrix},\quad
  \mathrm{SWAP}\ket{11}=\begin{pmatrix}0\\0\\0\\1\end{pmatrix},
\]
\[
  \mathrm{SWAP}=\begin{pmatrix}
    1 & 0 & 0 & 0\\ 0 & 0 & 1 & 0\\ 0 & 1 & 0 & 0\\ 0 & 0 & 0 & 1
  \end{pmatrix}.
\]
Appliquée à $\ket{01}=(0,1,0,0)^{\mathsf T}$, cette matrice sélectionne sa deuxième colonne :
\[
  \begin{pmatrix}
    1 & 0 & 0 & 0\\ 0 & 0 & 1 & 0\\ 0 & 1 & 0 & 0\\ 0 & 0 & 0 & 1
  \end{pmatrix}
  \begin{pmatrix}0\\1\\0\\0\end{pmatrix}
  =\begin{pmatrix}0\\0\\1\\0\end{pmatrix}=\ket{10}.
\]
Sur un état quelconque, $\mathrm{SWAP}\,(c_{00},c_{01},c_{10},c_{11})^{\mathsf T}
=(c_{00},c_{10},c_{01},c_{11})^{\mathsf T}$ : les amplitudes de $\ket{01}$ et $\ket{10}$
sont échangées.

\subsubsection{Accord avec la définition}

Pour deux états $\ket{\psi}=(a_0,a_1)^{\mathsf T}$ et $\ket{\phi}=(b_0,b_1)^{\mathsf T}$ :
\[
  \mathrm{SWAP}\,(\ket{\psi}\otimes\ket{\phi})
  =\mathrm{SWAP}\begin{pmatrix}a_0b_0\\a_0b_1\\a_1b_0\\a_1b_1\end{pmatrix}
  =\begin{pmatrix}a_0b_0\\a_1b_0\\a_0b_1\\a_1b_1\end{pmatrix}
  =\begin{pmatrix}b_0\\b_1\end{pmatrix}\otimes\begin{pmatrix}a_0\\a_1\end{pmatrix}
  =\ket{\phi}\otimes\ket{\psi},
\]
ce qui est bien la définition. La porte s'écrit aussi avec des projecteurs :
\[
  \mathrm{SWAP}=\sum_{a,b\in\{0,1\}}\ket{a}\bra{b}\otimes\ket{b}\bra{a}.
\]
En effet, pour $\ket{c}\otimes\ket{d}$ on trouve
$\sum_{a,b}\ket{a}\braket{b}{c}\otimes\ket{b}\braket{a}{d}$, et seul le terme $b=c$, $a=d$
est non nul : le résultat est $\ket{d}\otimes\ket{c}$.

\paragraph{Exemples.}
\begin{itemize}[nosep]
  \item Deux états différents : $\mathrm{SWAP}\big(\ket{0}\otimes\ket{+}\big)
        =\mathrm{SWAP}\,\frac{1}{\sqrt2}(1,1,0,0)^{\mathsf T}=\frac{1}{\sqrt2}(1,0,1,0)^{\mathsf T}
        =\ket{+}\otimes\ket{0}$.
  \item L'état $\ket{\psi}=\frac{1}{\sqrt2}\ket{00}+\frac{1}{\sqrt2}\ket{11}$ ne change pas :
        $\mathrm{SWAP}\ket{\psi}=\frac{1}{\sqrt2}\ket{00}+\frac{1}{\sqrt2}\ket{11}=\ket{\psi}$,
        car $\ket{00}$ et $\ket{11}$ sont invariants. De même
        $\frac{1}{\sqrt2}(\ket{00}-\ket{11})$ et $\frac{1}{\sqrt2}(\ket{01}+\ket{10})$ sont invariants,
        alors que $\mathrm{SWAP}\,\frac{1}{\sqrt2}(\ket{01}-\ket{10})=\frac{1}{\sqrt2}(\ket{10}-\ket{01})$
        est l'opposé de l'état de départ. Ce sont les quatre états de Bell (chapitre~\ref{chap:bell}).
\end{itemize}

\paragraph{Propriétés.} Échanger deux fois redonne l'état de départ :
$\mathrm{SWAP}^2=\mathrm{Id}$ ; la matrice est réelle et symétrique, donc
$\mathrm{SWAP}^\dagger\mathrm{SWAP}=\mathrm{SWAP}^2=\mathrm{Id}$ : elle est unitaire.

\subsubsection{Lien avec le CNOT}

\begin{resultat}{SWAP échange les rôles de contrôle et de cible}
\[
  \mathrm{CNOT}_{21}=\mathrm{SWAP}\;\mathrm{CNOT}_{12}\;\mathrm{SWAP}.
\]
\end{resultat}

\emph{Sur les états de base.}
$\mathrm{SWAP}\,\mathrm{CNOT}_{12}\,\mathrm{SWAP}\ket{a,b}
=\mathrm{SWAP}\,\mathrm{CNOT}_{12}\ket{b,a}
=\mathrm{SWAP}\ket{b,\,b\oplus a}=\ket{b\oplus a,\,b}=\mathrm{CNOT}_{21}\ket{a,b}$.

\emph{Par les matrices.} Multiplier à gauche par $\mathrm{SWAP}$ échange les lignes $2$ et
$3$, multiplier à droite l'échange les colonnes $2$ et $3$ :
\[
  \mathrm{SWAP}\,\mathrm{CNOT}_{12}
  =\begin{pmatrix}1&0&0&0\\0&0&0&1\\0&1&0&0\\0&0&1&0\end{pmatrix},
  \qquad
  \mathrm{SWAP}\,\mathrm{CNOT}_{12}\,\mathrm{SWAP}
  =\begin{pmatrix}1&0&0&0\\0&0&0&1\\0&0&1&0\\0&1&0&0\end{pmatrix}=\mathrm{CNOT}_{21}.
\]

\begin{remarque}[SWAP avec trois CNOT]
$\mathrm{SWAP}=\mathrm{CNOT}_{12}\,\mathrm{CNOT}_{21}\,\mathrm{CNOT}_{12}$. En effet, sur
$\ket{a,b}$ (on applique les portes de droite à gauche) :
\[
  \ket{a,b}\ \xrightarrow{\ \mathrm{CNOT}_{12}\ }\ \ket{a,\,a\oplus b}
  \ \xrightarrow{\ \mathrm{CNOT}_{21}\ }\ \ket{a\oplus(a\oplus b),\,a\oplus b}=\ket{b,\,a\oplus b}
  \ \xrightarrow{\ \mathrm{CNOT}_{12}\ }\ \ket{b,\,b\oplus(a\oplus b)}=\ket{b,a}.
\]
Les deux membres coïncident sur les états de base, donc sur tous les états par linéarité.
\end{remarque}

% ---------------------------------------------------------------------
\subsection{Portes contrôlées : contrôlé-\texorpdfstring{$U$}{U}, Fredkin et Toffoli}\label{sec:controlees}
% ---------------------------------------------------------------------

Le CNOT est le cas le plus simple d'une idée générale : appliquer une opération à un ou
plusieurs qubits \emph{cibles} seulement si un ou plusieurs qubits de \emph{contrôle} valent $1$.
Dans cette section, le contrôle est le \emph{premier} chiffre
du ket (remarque de la section~\ref{sec:cnot}). Les qubits sont notés $x$, $y$, $z$ (le qubit $x$ est le contrôle) et
$\ket{x,y,z}=\ket{x}\otimes\ket{y}\otimes\ket{z}$ ; d'après la convention de la section~\ref{sec:circuits}, le
contrôle est dessiné en bas (figure~\ref{fig:controlees}). Pour éviter la confusion avec les
matrices de Pauli, on écrit ici $\sigma_x=X$ pour la porte NOT.

\begin{figure}[!ht]
\centering
\begin{adjustbox}{max width=\linewidth}
\begin{tikzpicture}[>=Stealth, font=\small]
  % (a) contrôlé-U : contrôle x en bas, cible y en haut
  \draw (0,1.2) -- (3.4,1.2);
  \draw (0,0) -- (3.4,0);
  \node[anchor=east] at (0,1.2) {$y$};
  \node[anchor=east] at (0,0) {$x$};
  \draw[insarouge, thick] (1.7,0) -- (1.7,0.85);
  \pic at (1.7,0) {ctl};
  \node[porte] at (1.7,1.2) {$U$};
  \node[anchor=north, font=\footnotesize] at (1.7,-0.5) {(a) contrôlé-$U$};
  % (b) contrôlé-SWAP (Fredkin) : x contrôle, y et z échangés
  \draw (5,2.4) -- (8.4,2.4);
  \draw (5,1.2) -- (8.4,1.2);
  \draw (5,0) -- (8.4,0);
  \node[anchor=east] at (5,2.4) {$z$};
  \node[anchor=east] at (5,1.2) {$y$};
  \node[anchor=east] at (5,0) {$x$};
  \draw[insarouge, thick] (6.7,0) -- (6.7,2.4);
  \pic at (6.7,0) {ctl};
  \pic at (6.7,1.2) {croix};
  \pic at (6.7,2.4) {croix};
  \node[anchor=north, font=\footnotesize] at (6.7,-0.5) {(b) Fredkin (C-SWAP)};
  % (c) Toffoli (CCNOT) : x et y contrôlent, z cible
  \draw (10,2.4) -- (13.4,2.4);
  \draw (10,1.2) -- (13.4,1.2);
  \draw (10,0) -- (13.4,0);
  \node[anchor=east] at (10,2.4) {$z$};
  \node[anchor=east] at (10,1.2) {$y$};
  \node[anchor=east] at (10,0) {$x$};
  \draw[insarouge, thick] (11.7,0) -- (11.7,2.4);
  \pic at (11.7,0) {ctl};
  \pic at (11.7,1.2) {ctl};
  \pic at (11.7,2.4) {cible};
  \node[anchor=north, font=\footnotesize] at (11.7,-0.5) {(c) Toffoli (CCNOT)};
\end{tikzpicture}
\end{adjustbox}
\caption{Portes contrôlées. Le contrôle $x$ est le premier chiffre du ket, donc le fil du bas.
(a) Le qubit $y$ subit $U$ si $x=1$ ; pour $U=X$ on retrouve le CNOT. (b) Si $x=1$, les qubits
$y$ et $z$ sont échangés. (c) Si $x=y=1$, le qubit $z$ est inversé.}
\label{fig:controlees}
\end{figure}

\subsubsection{Contrôlé-\texorpdfstring{$U$}{U}}

\begin{definition}{Porte contrôlée-$U$}
Soit $U$ une matrice unitaire $2\times2$. Sur le système $(x,y)$, de contrôle $x$ et de cible $y$,
\[
  \mathrm{C}U=\ket{0}\bra{0}_x\otimes I_y+\ket{1}\bra{1}_x\otimes U_y .
\]
Le premier terme agit quand le contrôle est $\ket{0}$ : la cible ne change pas ($I$). Le second
agit quand le contrôle est $\ket{1}$ : la cible subit $U$.
\end{definition}

Comme $\ket{0}\bra{0}\ket{x}=\delta_{x0}\ket{0}$ et $\ket{1}\bra{1}\ket{x}=\delta_{x1}\ket{1}$,
\[
  \mathrm{C}U\big(\ket{x}\otimes\ket{y}\big)=\ket{x}\otimes U^x\ket{y},\qquad U^0=I,\ U^1=U .
\]
Par la règle $A\otimes B=(A_{ij}B)$, $\ket{0}\bra{0}\otimes I=\begin{pmatrix}I&0\\0&0\end{pmatrix}$
et $\ket{1}\bra{1}\otimes U=\begin{pmatrix}0&0\\0&U\end{pmatrix}$ (blocs $2\times2$), donc
\[
  \mathrm{C}U=\begin{pmatrix}I&0\\0&U\end{pmatrix}
  =\begin{pmatrix}1&0&0&0\\0&1&0&0\\0&0&u_{00}&u_{01}\\0&0&u_{10}&u_{11}\end{pmatrix}
  \qquad\text{pour }U=\begin{pmatrix}u_{00}&u_{01}\\u_{10}&u_{11}\end{pmatrix}.
\]
Cette matrice est unitaire car $I^\dagger I=I$ et $U^\dagger U=I$.

\paragraph{Le CNOT est un cas particulier.} Pour $U=\sigma_x$,
$\mathrm{C}\sigma_x=\ket{0}\bra{0}_x\otimes I_y+\ket{1}\bra{1}_x\otimes\sigma_x$ est le CNOT de
contrôle $x$ (premier chiffre), c'est-à-dire $\mathrm{CNOT}_{12}$ de la section~\ref{sec:cnot}.
On calcule l'action sur les quatre états de base, avec $\braket{0}{0}=\braket{1}{1}=1$,
$\braket{0}{1}=\braket{1}{0}=0$, $\sigma_x\ket{0}=\ket{1}$ et $\sigma_x\ket{1}=\ket{0}$ :
\begin{align*}
  \mathrm{C}\sigma_x\ket{00}&=\big[\ket{0}\bra{0}_x\otimes I_y\big]\ket{0}_x\ket{0}_y
      +\big[\ket{1}\bra{1}_x\otimes\sigma_x\big]\ket{0}_x\ket{0}_y\\
    &=\ket{0}\braket{0}{0}\otimes I\ket{0}+\ket{1}\braket{1}{0}\otimes\sigma_x\ket{0}
     =\ket{0}\otimes\ket{0}+0=\ket{00},\\[1.2ex]
  \mathrm{C}\sigma_x\ket{01}&=\ket{0}\braket{0}{0}\otimes I\ket{1}+\ket{1}\braket{1}{0}\otimes\sigma_x\ket{1}
     =\ket{0}\otimes\ket{1}+0=\ket{01},\\[1.2ex]
  \mathrm{C}\sigma_x\ket{10}&=\ket{0}\braket{0}{1}\otimes I\ket{0}+\ket{1}\braket{1}{1}\otimes\sigma_x\ket{0}
     =0+\ket{1}\otimes\ket{1}=\ket{11},\\[1.2ex]
  \mathrm{C}\sigma_x\ket{11}&=\ket{0}\braket{0}{1}\otimes I\ket{1}+\ket{1}\braket{1}{1}\otimes\sigma_x\ket{1}
     =0+\ket{1}\otimes\ket{0}=\ket{10}.
\end{align*}
On retrouve la table du CNOT : la cible n'est inversée que si le contrôle vaut $1$.

\paragraph{Un autre exemple : $U=Z$.} Ici
\[
  \mathrm{C}Z=\begin{pmatrix}I&0\\0&Z\end{pmatrix}=\mathrm{diag}(1,1,1,-1) :
\]
la porte multiplie $\ket{11}$ par $-1$ et ne touche pas les autres états de base. Elle est
symétrique entre les deux qubits. Elle se déduit du CNOT par
$\mathrm{C}Z=(I\otimes H)\,\mathrm{C}X\,(I\otimes H)$, où $\mathrm{C}X$ est le contrôlé-$X$.
En effet, $I\otimes H=\mathrm{diag}(H,H)$ (deux blocs $2\times2$ égaux à $H$), et le produit de
matrices diagonales par blocs se fait bloc par bloc :
\[
  (I\otimes H)\begin{pmatrix}I&0\\0&X\end{pmatrix}(I\otimes H)
  =\begin{pmatrix}H\,I\,H&0\\0&H\,X\,H\end{pmatrix}
  =\begin{pmatrix}I&0\\0&HXH\end{pmatrix},
\]
car $H^2=I$, et
\[
  HXH=\frac12\begin{pmatrix}1&1\\1&-1\end{pmatrix}\begin{pmatrix}0&1\\1&0\end{pmatrix}\begin{pmatrix}1&1\\1&-1\end{pmatrix}
  =\frac12\begin{pmatrix}1&1\\1&-1\end{pmatrix}\begin{pmatrix}1&-1\\1&1\end{pmatrix}
  =\frac12\begin{pmatrix}2&0\\0&-2\end{pmatrix}=Z
\]
(section~\ref{sec:hadamard} : $H$ échange les axes $x$ et $z$).

\subsubsection{Contrôlé-SWAP : la porte de Fredkin}

\begin{definition}{Porte de Fredkin (C-SWAP)}
Sur trois qubits $(x,y,z)$ (espace de dimension $2^3=8$), avec $x$ pour contrôle,
\[
  \mathrm{C\text{-}SWAP}=\ket{0}\bra{0}_x\otimes I_y\otimes I_z+\ket{1}\bra{1}_x\otimes\mathrm{SWAP}_{yz}.
\]
Si $x=0$, rien ne change ; si $x=1$, les qubits $y$ et $z$ sont échangés.
\end{definition}

Calcul sur quatre états de base (avec $I\ket{ab}=\ket{ab}$ et $\mathrm{SWAP}\ket{ab}=\ket{ba}$) :
\begin{align*}
  \mathrm{C\text{-}SWAP}\ket{000}&=\ket{0}\braket{0}{0}\otimes I\ket{00}+\ket{1}\braket{1}{0}\otimes\mathrm{SWAP}\ket{00}
     =\ket{0}\otimes\ket{00}+0=\ket{000},\\[1.2ex]
  \mathrm{C\text{-}SWAP}\ket{101}&=\ket{0}\braket{0}{1}\otimes I\ket{01}+\ket{1}\braket{1}{1}\otimes\mathrm{SWAP}\ket{01}
     =0+\ket{1}\otimes\ket{10}=\ket{110},\\[1.2ex]
  \mathrm{C\text{-}SWAP}\ket{110}&=\ket{0}\braket{0}{1}\otimes I\ket{10}+\ket{1}\braket{1}{1}\otimes\mathrm{SWAP}\ket{10}
     =0+\ket{1}\otimes\ket{01}=\ket{101},\\[1.2ex]
  \mathrm{C\text{-}SWAP}\ket{010}&=\ket{0}\braket{0}{0}\otimes I\ket{10}+\ket{1}\braket{1}{0}\otimes\mathrm{SWAP}\ket{10}
     =\ket{0}\otimes\ket{10}+0=\ket{010}.
\end{align*}
Les huit états de base donnent :
\begin{center}
\renewcommand{\arraystretch}{1.3}
\begin{tabular}{@{}ccccccccc@{}}
\toprule
Entrée & $\ket{000}$ & $\ket{001}$ & $\ket{010}$ & $\ket{011}$ & $\ket{100}$ & $\ket{101}$ & $\ket{110}$ & $\ket{111}$ \\
\midrule
Sortie & $\ket{000}$ & $\ket{001}$ & $\ket{010}$ & $\ket{011}$ & $\ket{100}$ & $\ket{110}$ & $\ket{101}$ & $\ket{111}$ \\
\bottomrule
\end{tabular}
\end{center}
Seuls $\ket{101}$ et $\ket{110}$ changent (contrôle à $1$ et $y\neq z$) ; pour $\ket{100}$ et
$\ket{111}$ l'échange de deux qubits égaux ne change rien. Dans la base
$\ket{000},\ket{001},\dots,\ket{111}$ (ordre binaire), la matrice est
$\mathrm{diag}(I_4,\mathrm{SWAP})$ (deux blocs $4\times4$) :
\[
  \mathrm{C\text{-}SWAP}=\begingroup\setlength{\arraycolsep}{3pt}\begin{pmatrix}
    1&0&0&0&0&0&0&0\\
    0&1&0&0&0&0&0&0\\
    0&0&1&0&0&0&0&0\\
    0&0&0&1&0&0&0&0\\
    0&0&0&0&1&0&0&0\\
    0&0&0&0&0&0&1&0\\
    0&0&0&0&0&1&0&0\\
    0&0&0&0&0&0&0&1
  \end{pmatrix}\endgroup .
\]
C'est l'identité dont les lignes $6$ et $7$ sont échangées.

\paragraph{Contrôle en superposition.} Avec $\ket{+}_x\otimes\ket{0}_y\otimes\ket{1}_z
=\frac{1}{\sqrt2}\big(\ket{001}+\ket{101}\big)$ :
\[
  \mathrm{C\text{-}SWAP}\,\frac{1}{\sqrt2}\big(\ket{001}+\ket{101}\big)
  =\frac{1}{\sqrt2}\big(\ket{001}+\ket{110}\big).
\]
Le résultat n'est pas un produit : les coefficients de $\ket{x}\otimes\ket{yz}$ forment la
matrice (lignes $x=0,1$, colonnes $yz=00,01,10,11$)
\[
  \frac{1}{\sqrt2}\begin{pmatrix}0&1&0&0\\0&0&1&0\end{pmatrix},
\]
dont les deux lignes ne sont pas proportionnelles. Le qubit $x$ est intriqué avec la paire
$(y,z)$.

\subsubsection{Contrôlé-contrôlé-NOT : la porte de Toffoli}

\begin{definition}{Porte de Toffoli (CCNOT)}
Sur trois qubits $(x,y,z)$, avec $x$ et $y$ pour contrôles et $z$ pour cible,
\[
  \mathrm{CCNOT}=\ket{0}\bra{0}_x\otimes I_y\otimes I_z
   +\ket{1}\bra{1}_x\otimes\Big[\ket{0}\bra{0}_y\otimes I_z+\ket{1}\bra{1}_y\otimes\sigma_x\Big].
\]
Le crochet est un CNOT (contrôle $y$, cible $z$), appliqué seulement si $x=1$ : si $x=0$, rien ;
si $x=1$ et $y=0$, rien ; si $x=1$ et $y=1$, le qubit $z$ est inversé.
\end{definition}

Calcul de $\mathrm{CCNOT}\ket{110}$ :
\begin{align*}
  \mathrm{CCNOT}\ket{110}&=\ket{0}\braket{0}{1}\otimes I\ket{1}\otimes I\ket{0}
     +\ket{1}\braket{1}{1}\otimes\Big[\ket{0}\braket{0}{1}\otimes I\ket{0}
       +\ket{1}\braket{1}{1}\otimes\sigma_x\ket{0}\Big]\\
  &=0+\ket{1}\otimes\Big[0+\ket{1}\otimes\ket{1}\Big]=\ket{111}.
\end{align*}
Pour les huit états de base, $\mathrm{CCNOT}\ket{x,y,z}=\ket{x,\,y,\,z\oplus xy}$ :
\begin{center}
\renewcommand{\arraystretch}{1.3}
\begin{tabular}{@{}ccccccccc@{}}
\toprule
Entrée & $\ket{000}$ & $\ket{001}$ & $\ket{010}$ & $\ket{011}$ & $\ket{100}$ & $\ket{101}$ & $\ket{110}$ & $\ket{111}$ \\
\midrule
Sortie & $\ket{000}$ & $\ket{001}$ & $\ket{010}$ & $\ket{011}$ & $\ket{100}$ & $\ket{101}$ & $\ket{111}$ & $\ket{110}$ \\
\bottomrule
\end{tabular}
\end{center}
Seuls $\ket{110}$ et $\ket{111}$ sont échangés : la matrice est $\mathrm{diag}(I_6,X)$
(un bloc $6\times6$ égal à l'identité, puis $X$),
\[
  \mathrm{CCNOT}=\begingroup\setlength{\arraycolsep}{3pt}\begin{pmatrix}
    1&0&0&0&0&0&0&0\\
    0&1&0&0&0&0&0&0\\
    0&0&1&0&0&0&0&0\\
    0&0&0&1&0&0&0&0\\
    0&0&0&0&1&0&0&0\\
    0&0&0&0&0&1&0&0\\
    0&0&0&0&0&0&0&1\\
    0&0&0&0&0&0&1&0
  \end{pmatrix}\endgroup .
\]

\begin{remarque}[Portes réversibles]
Les portes de Fredkin et de Toffoli sont des permutations des états de base qui échangent
seulement deux d'entre eux : elles sont leurs propres inverses ($\mathrm{C\text{-}SWAP}^2=\mathrm{CCNOT}^2=\mathrm{Id}$)
et unitaires. Avec $z=0$, la Toffoli calcule le ET logique $x\wedge y$ sans rien effacer :
$\ket{x,y,0}\to\ket{x,y,xy}$ ; avec $z=1$, elle calcule le NON-ET. Elle sert à écrire tout calcul
classique de façon réversible.
\end{remarque}

\clearpage
% ---------------------------------------------------------------------
\subsection{Récapitulatif du chapitre}\label{sec:recap4}
% ---------------------------------------------------------------------

\begin{center}
\adjustbox{max width=\linewidth}{%
\begin{tikzpicture}[>=Stealth, font=\footnotesize,
    boite/.style={draw=insarouge, fill=insarouge!6, rounded corners=2pt, thick,
                  minimum width=4.4cm, minimum height=1.5cm, align=center, inner sep=3pt},
    lien/.style={->, very thick, insagris},
    etiq/.style={font=\scriptsize\itshape, text=insagris, inner sep=2pt}]
  \node[boite] (uni) at (0,0)
    {\textbf{Porte unitaire}\\ $U^\dagger U=\mathrm{Id}$, \ $\ket{\Psi}\to U\ket{\Psi}$\\ {\scriptsize\S\,\ref{sec:portes}}};
  \node[boite] (un) at (5.9,0)
    {\textbf{Portes à un qubit}\\ $X$, $Y$, $Z$, $H$, $S$, $T$\\ {\scriptsize\S\,\ref{sec:pauli}--\ref{sec:phase}}};
  \node[boite] (blo) at (11.8,0)
    {\textbf{Rotations de Bloch}\\ $X,Y,Z,H$ : angle $\pi$ ; $P_\theta$ autour de $z$\\
     {\scriptsize\S\,\ref{sec:portes-bloch}}};
  \node[boite] (cir) at (11.8,-2.5)
    {\textbf{Circuits}\\ fils, boîtes, $U_2U_1$\\ {\scriptsize\S\,\ref{sec:circuits}, \ref{sec:op2}}};
  \node[boite] (cn) at (5.9,-2.5)
    {\textbf{CNOT et SWAP}\\ $\ket{a,b}\to\ket{a\oplus b,b}$\\ {\scriptsize\S\,\ref{sec:cnot}, \ref{sec:swap}}};
  \node[boite] (cu) at (0,-2.5)
    {\textbf{Portes contrôlées}\\ $\mathrm{C}U=\mathrm{diag}(I,U)$, Fredkin, Toffoli\\
     {\scriptsize\S\,\ref{sec:controlees}}};
  \draw[lien] (uni) -- (un);
  \draw[lien] (un) -- (blo);
  \draw[lien] (blo) -- (cir);
  \draw[lien] (cir) -- (cn) node[etiq, midway, above=1pt] {2 qubits};
  \draw[lien] (cn) -- (cu) node[etiq, midway, above=1pt] {contrôle};
\end{tikzpicture}}
\end{center}

\begin{aretenir}
\begin{itemize}
  \item Une porte est un opérateur unitaire : elle conserve la norme, elle est réversible
        ($U^{-1}=U^\dagger$) et agit sur la sphère de Bloch comme une rotation. C'est l'évolution de
        Schrödinger, $U=\ee^{-\ii\Hop t/\hbar}$.
  \item $X$ échange $\ket{0}$ et $\ket{1}$, $Z$ change le signe de $\ket{1}$, $Y=\ii XZ$ fait les deux ;
        $H$ échange les bases $Z$ et $X$ et crée des superpositions ; $P_\theta$, $S=P_{\pi/2}$ et
        $T=P_{\pi/4}$ font tourner autour de $z$ ($T^2=S$, $S^2=Z$).
  \item Dans un circuit, le temps va de gauche à droite et la matrice est le produit dans l'ordre
        inverse, $U_2U_1$. Le fil du haut porte le dernier chiffre du ket. Des portes en parallèle
        forment un produit tensoriel $U_A\otimes U_B$.
  \item Le CNOT inverse la cible si le contrôle vaut $1$ : $\ket{a,b}\to\ket{a\oplus b,b}$. Le SWAP
        échange deux qubits ; $\mathrm{CNOT}_{21}=\mathrm{SWAP}\,\mathrm{CNOT}_{12}\,\mathrm{SWAP}$.
  \item Une porte contrôlée a pour matrice $\mathrm{diag}(I,U)$. Fredkin (contrôlé-SWAP) et
        Toffoli (CCNOT) sont des permutations des états de base : réversibles et universelles
        pour le calcul classique.
\end{itemize}
\end{aretenir}

\begin{center}
\renewcommand{\arraystretch}{1.45}
\begin{tabular}{@{}l l r@{}}
\toprule
\textbf{Notion} & \textbf{Formule} & \textbf{§} \\
\midrule
Porte unitaire & $U^\dagger U=\mathrm{Id}$, \quad $\bra{\Psi}U^\dagger U\ket{\Psi}=1$ & \ref{sec:portes} \\
Pauli & $X\ket{0}=\ket{1}$, \quad $Z\ket{1}=-\ket{1}$, \quad $X^2=Y^2=Z^2=\mathrm{Id}$ & \ref{sec:pauli} \\
Hadamard & $H\ket{0}=\ket{+}$, \quad $H\ket{1}=\ket{-}$, \quad $H^2=\mathrm{Id}$ & \ref{sec:hadamard} \\
Phase & $P_\theta=\mathrm{diag}(1,\ee^{\ii\theta})$, \quad $S=P_{\pi/2}$, \quad $T=P_{\pi/4}$ & \ref{sec:phase} \\
Circuit & $U_1$ puis $U_2$ : \quad $U_2U_1$ & \ref{sec:circuits} \\
Portes locales & $(U\otimes I)(\ket{u}\otimes\ket{v})=U\ket{u}\otimes\ket{v}$ & \ref{sec:op2} \\
CNOT & $\ket{a,b}\to\ket{a\oplus b,b}$, \quad $\mathrm{CNOT}=I\otimes\ket{0}\bra{0}+X\otimes\ket{1}\bra{1}$
  & \ref{sec:cnot} \\
SWAP & $\mathrm{SWAP}\ket{a,b}=\ket{b,a}$, \quad $\mathrm{SWAP}=\mathrm{CNOT}_{12}\mathrm{CNOT}_{21}\mathrm{CNOT}_{12}$
  & \ref{sec:swap} \\
Contrôlé-$U$ & $\mathrm{C}U=\ket{0}\bra{0}\otimes I+\ket{1}\bra{1}\otimes U=\mathrm{diag}(I,U)$ & \ref{sec:controlees} \\
Toffoli & $\mathrm{CCNOT}\ket{x,y,z}=\ket{x,y,z\oplus xy}$ & \ref{sec:controlees} \\
\bottomrule
\end{tabular}
\end{center}

\begin{savoirfaire}
\begin{itemize}
  \item Appliquer une porte à un état (par la matrice ou par les ket-bra) et identifier sa rotation
        de Bloch (exercices~\ref{ex:conj} et~\ref{ex:phases}).
  \item Lire un circuit, écrire sa matrice dans le bon ordre et suivre l'état à chaque étape
        (exercice~\ref{ex:kickback}).
  \item Écrire la matrice de $U\otimes I$, du CNOT, du SWAP ou d'une porte contrôlée, et calculer
        leur action sur les états de base.
  \item Reconnaître une porte réversible qui calcule une fonction booléenne
        (exercice~\ref{ex:toffoli}).
\end{itemize}
\end{savoirfaire}
 % 4. Portes quantiques et systèmes multi-qubits
% =====================================================================
%  Chapitre 5 : États quantiques, intrication et états de Bell
%  Fichier importé par main.tex avec % =====================================================================
%  Chapitre 5 : États quantiques, intrication et états de Bell
%  Fichier importé par main.tex avec % =====================================================================
%  Chapitre 5 : États quantiques, intrication et états de Bell
%  Fichier importé par main.tex avec \include{chapitre5}.
%  Il ne contient ni préambule ni \begin{document} : les environnements
%  (definition, resultat, remarque, aretenir, savoirfaire, exercice, correction),
%  les macros (\ii, \ee, \dd, \pd, \pdd, \Hop, \ket, \bra, \braket), les symboles
%  de circuits (ctl, cible, croix, mesure, porte) et les couleurs (insarouge,
%  insagris) sont définis dans main.tex. Un chapitre = une \section.
% =====================================================================

% Macros de Dirac (déjà définies dans main.tex ; ce bloc évite l'erreur
% « Undefined control sequence \ket » si main.tex est une ancienne version).
\providecommand{\ket}[1]{|#1\rangle}
\providecommand{\bra}[1]{\langle #1|}
\providecommand{\braket}[2]{\langle #1|#2\rangle}

% =====================================================================
\section{États quantiques, intrication et états de Bell}\label{chap:bell}
% =====================================================================

Les états de deux qubits forment un espace de dimension~$4$, et la plupart d'entre eux ne se
décomposent pas en «~un état pour $A$, un état pour $B$~» : ils sont \emph{intriqués}. Ce chapitre
apprend à mesurer un couple de qubits, à reconnaître l'intrication (critère de séparabilité), à la
produire avec un circuit (états de Bell), à décrire un qubit pris seul (trace partielle), à
quantifier l'information (entropies) et à s'en servir pour téléporter un état.

% ---------------------------------------------------------------------
\subsection{États de deux qubits et mesure}\label{sec:mesure2}
% ---------------------------------------------------------------------

Un état quelconque de $\mathcal{H}^4$ se décompose sur la base de calcul :
\[
  \ket{\Psi}=c_{00}\ket{00}+c_{01}\ket{01}+c_{10}\ket{10}+c_{11}\ket{11}
  =\begin{pmatrix}c_{00}\\ c_{01}\\ c_{10}\\ c_{11}\end{pmatrix},
  \qquad \sum_{i,j}|c_{ij}|^2=1 .
\]

\begin{resultat}{Mesure des deux qubits dans la base de calcul}
\[
  P(ij)=|\braket{ij}{\Psi}|^2=|c_{ij}|^2 .
\]
Pour un seul qubit : $P(A{=}0)=P(00)+P(01)$, $P(A{=}1)=P(10)+P(11)$,
$P(B{=}0)=P(00)+P(10)$, $P(B{=}1)=P(01)+P(11)$.
\end{resultat}

En effet $\braket{ij}{\Psi}=\sum_{k,l}c_{kl}\braket{ij}{kl}=\sum_{k,l}c_{kl}\delta_{ik}\delta_{jl}=c_{ij}$.

\paragraph{Exemple 1.} $\ket{\Psi}=\frac{1}{\sqrt2}\ket{00}+\frac{1}{\sqrt2}\ket{01}
=\big(\tfrac{1}{\sqrt2},\tfrac{1}{\sqrt2},0,0\big)^{\mathsf T}$. Il est normalisé :
$\frac12+\frac12+0+0=1$. Les probabilités jointes sont
$P(00)=P(01)=\frac12$ et $P(10)=P(11)=0$. Pour chaque qubit :
\[
  P(A{=}0)=\tfrac12+\tfrac12=1,\quad P(A{=}1)=0+0=0,\quad
  P(B{=}0)=\tfrac12+0=\tfrac12,\quad P(B{=}1)=\tfrac12+0=\tfrac12 .
\]
$A$ donne toujours $0$ ; $B$ donne $0$ ou $1$ au hasard.

\paragraph{Exemple 2 (mesure de $A$ seul).} Soit
$\ket{\Psi}=\frac12\ket{00}+\frac12\ket{01}+\frac{1}{\sqrt2}\ket{11}$. Normalisation :
$\frac14+\frac14+0+\frac12=1$. Probabilités jointes : $P(00)=\frac14$, $P(01)=\frac14$,
$P(10)=0$, $P(11)=\frac12$. Donc
\[
  P(A{=}0)=\tfrac14+\tfrac14=\tfrac12,\quad P(A{=}1)=0+\tfrac12=\tfrac12,\quad
  P(B{=}0)=\tfrac14+0=\tfrac14,\quad P(B{=}1)=\tfrac14+\tfrac12=\tfrac34 .
\]

\begin{remarque}[Effondrement après la mesure de $A$]
Si $A$ donne le résultat $a$, on ne garde que les termes de $\ket{\Psi}$ dont le premier
chiffre vaut $a$, puis on renormalise :
\[
  \ket{\Psi}\ \to\ \frac{c_{a0}\ket{a0}+c_{a1}\ket{a1}}{\sqrt{|c_{a0}|^2+|c_{a1}|^2}},
  \qquad\text{probabilité }|c_{a0}|^2+|c_{a1}|^2 .
\]
\end{remarque}

Pour l'exemple 2 :
\begin{itemize}[nosep]
  \item $A=0$ (probabilité $\frac12$) : on garde $\frac12\ket{00}+\frac12\ket{01}$, de norme
        au carré $\frac12$ ; renormalisé, $\frac{1}{\sqrt2}\big(\ket{00}+\ket{01}\big)
        =\ket{0}\otimes\ket{+}$. Ensuite $B$ vaut $0$ ou $1$ avec la probabilité $\frac12$.
  \item $A=1$ (probabilité $\frac12$) : on garde $\frac{1}{\sqrt2}\ket{11}$, de norme au carré
        $\frac12$ ; renormalisé, $\ket{11}=\ket{1}\otimes\ket{1}$. Ensuite $B$ vaut $1$
        avec certitude.
\end{itemize}
Contrôle : $P(B{=}0)=\frac12\cdot\frac12+\frac12\cdot0=\frac14$, comme calculé plus haut.
Ici la mesure de $A$ \emph{modifie} l'état de $B$ ($\ket{+}$ ou $\ket{1}$) : c'est la
signature d'un état intriqué (section~\ref{sec:separable}). La figure~\ref{fig:mesure2} résume
les deux mesures successives.

\begin{figure}[!ht]
\centering
\adjustbox{max width=\linewidth}{%
\begin{tikzpicture}[>=Stealth, font=\small,
    noeud/.style={draw=insarouge, fill=insarouge!6, rounded corners=2pt, thick, inner sep=4pt, align=center},
    fin/.style={draw=insagris, fill=insagris!10, rounded corners=2pt, inner sep=4pt, align=center},
    lien/.style={->, thick, insagris},
    proba/.style={font=\footnotesize, fill=white, inner sep=1.5pt, text=black}]
  \node[noeud] (psi) at (0,0) {$\ket{\Psi}$};
  \node[noeud] (a0) at (4.4,1.6) {$A=0$\\ {\footnotesize état $\ket{0}\otimes\ket{+}$}};
  \node[noeud] (a1) at (4.4,-1.6) {$A=1$\\ {\footnotesize état $\ket{1}\otimes\ket{1}$}};
  \node[fin] (b00) at (9.2,2.5) {$B=0$ : $\ket{00}$};
  \node[fin] (b01) at (9.2,0.7) {$B=1$ : $\ket{01}$};
  \node[fin] (b11) at (9.2,-1.6) {$B=1$ : $\ket{11}$};
  \draw[lien] (psi) -- (a0) node[proba, midway] {$\frac12$};
  \draw[lien] (psi) -- (a1) node[proba, midway] {$\frac12$};
  \draw[lien] (a0) -- (b00) node[proba, midway] {$\frac12$};
  \draw[lien] (a0) -- (b01) node[proba, midway] {$\frac12$};
  \draw[lien] (a1) -- (b11) node[proba, midway] {$1$};
  \node[anchor=west] at (11.0,2.5) {$P(00)=\frac12\cdot\frac12=\frac14$};
  \node[anchor=west] at (11.0,0.7) {$P(01)=\frac12\cdot\frac12=\frac14$};
  \node[anchor=west] at (11.0,-1.6) {$P(11)=\frac12\cdot1=\frac12$};
  \node[text=insagris, font=\footnotesize\itshape] at (2.2,3.0) {mesure de $A$};
  \node[text=insagris, font=\footnotesize\itshape] at (6.8,3.6) {mesure de $B$};
  \node[text=insagris, font=\footnotesize\itshape, anchor=west] at (11.0,3.6) {probabilité jointe};
\end{tikzpicture}}
\caption{Mesure de $A$ puis de $B$ pour l'état de l'exemple~2,
$\ket{\Psi}=\frac12\ket{00}+\frac12\ket{01}+\frac{1}{\sqrt2}\ket{11}$. Chaque flèche porte une
probabilité, et le produit des probabilités le long d'un chemin est la probabilité jointe. Après la
mesure de $A$, l'état de $B$ dépend du résultat ($\ket{+}$ ou $\ket{1}$).}
\label{fig:mesure2}
\end{figure}

% ---------------------------------------------------------------------
\subsection{États séparables et états intriqués}\label{sec:separable}
% ---------------------------------------------------------------------

\begin{definition}{Séparable ou intriqué}
Un état $\ket{\Psi}_{AB}$ est \emph{séparable} s'il s'écrit comme un produit
$\ket{\Psi}_{AB}=\ket{\Psi}_A\otimes\ket{\Psi}_B$. Sinon il est \emph{intriqué} : on ne
peut pas attribuer d'état propre à chaque qubit.
\end{definition}

\paragraph{Exemple séparable (exemple 1 ci-dessus).} Par linéarité du produit tensoriel :
\begin{align*}
  \ket{\Psi}&=\tfrac{1}{\sqrt2}\ket{00}_{AB}+\tfrac{1}{\sqrt2}\ket{01}_{AB}
    =\tfrac{1}{\sqrt2}\Big(\ket{0}_A\otimes\ket{0}_B+\ket{0}_A\otimes\ket{1}_B\Big)\\
  &=\ket{0}_A\otimes\Big(\tfrac{1}{\sqrt2}\big(\ket{0}+\ket{1}\big)\Big)_B
    =\ket{0}_A\otimes\ket{+}_B ,
\end{align*}
donc $\ket{\Psi}_A=\ket{0}$ et $\ket{\Psi}_B=\ket{+}$. Contrôle par les composantes :
$\begin{pmatrix}1\\ 0\end{pmatrix}\otimes\tfrac{1}{\sqrt2}\begin{pmatrix}1\\ 1\end{pmatrix}
=\tfrac{1}{\sqrt2}\begin{pmatrix}1\cdot1\\ 1\cdot1\\ 0\cdot1\\ 0\cdot1\end{pmatrix}
=\tfrac{1}{\sqrt2}\begin{pmatrix}1\\ 1\\ 0\\ 0\end{pmatrix}$.

\subsubsection{Un critère de séparabilité}

Un produit quelconque a pour coefficients $c_{00}=ac$, $c_{01}=ad$, $c_{10}=bc$,
$c_{11}=bd$ (formule de la section~\ref{sec:tenseur}, avec $a,b$ pour $A$ et $c,d$ pour $B$), donc
\[
  c_{00}c_{11}=(ac)(bd)=abcd=(ad)(bc)=c_{01}c_{10}.
\]

\begin{resultat}{Test de séparabilité (deux qubits)}
$\ket{\Psi}=\sum c_{ij}\ket{ij}$ est séparable si et seulement si
\[
  D=c_{00}c_{11}-c_{01}c_{10}=0 .
\]
\end{resultat}

\emph{Démonstration.} ($\Rightarrow$) Calcul ci-dessus. ($\Leftarrow$) On regroupe les
termes selon l'état de $A$ :
\[
  \ket{\Psi}=\ket{0}\otimes\big(c_{00}\ket{0}+c_{01}\ket{1}\big)+\ket{1}\otimes\big(c_{10}\ket{0}+c_{11}\ket{1}\big).
\]
Si $D=0$, les lignes $(c_{00},c_{01})$ et $(c_{10},c_{11})$ sont proportionnelles. Si
$(c_{00},c_{01})\neq(0,0)$, il existe $\lambda$ avec $(c_{10},c_{11})=\lambda(c_{00},c_{01})$
et $\ket{\Psi}=\big(\ket{0}+\lambda\ket{1}\big)\otimes\big(c_{00}\ket{0}+c_{01}\ket{1}\big)$.
Si $(c_{00},c_{01})=(0,0)$, alors $\ket{\Psi}=\ket{1}\otimes\big(c_{10}\ket{0}+c_{11}\ket{1}\big)$.
Dans les deux cas $\ket{\Psi}$ est un produit ; on renormalise chaque facteur.
\hfill$\square$

\subsubsection{Application à plusieurs états}

\begin{center}
\renewcommand{\arraystretch}{1.5}
\begin{tabular}{@{}lccl@{}}
\toprule
\textbf{État} & $(c_{00},c_{01},c_{10},c_{11})$ & $D$ & \textbf{Nature} \\
\midrule
$\ket{00}$ & $(1,0,0,0)$ & $1\cdot0-0\cdot0=0$ & séparable \\
$\ket{+}\otimes\ket{+}$ & $\frac12(1,1,1,1)$ & $\frac14-\frac14=0$ & séparable \\
$\ket{+}\otimes\ket{-}$ & $\frac12(1,-1,1,-1)$ & $-\frac14-(-\frac14)=0$ & séparable \\
$\frac{1}{\sqrt2}(\ket{00}+\ket{01})$ & $\frac{1}{\sqrt2}(1,1,0,0)$ & $0-\frac12\cdot0=0$ & séparable \\
$\ket{\Phi^+}=\frac{1}{\sqrt2}(\ket{00}+\ket{11})$ & $\frac{1}{\sqrt2}(1,0,0,1)$ & $\frac12-0=\frac12$ & intriqué \\
$\ket{\Phi^-}=\frac{1}{\sqrt2}(\ket{00}-\ket{11})$ & $\frac{1}{\sqrt2}(1,0,0,-1)$ & $-\frac12-0=-\frac12$ & intriqué \\
$\ket{\Psi^+}=\frac{1}{\sqrt2}(\ket{01}+\ket{10})$ & $\frac{1}{\sqrt2}(0,1,1,0)$ & $0-\frac12=-\frac12$ & intriqué \\
$\ket{\Psi^-}=\frac{1}{\sqrt2}(\ket{01}-\ket{10})$ & $\frac{1}{\sqrt2}(0,1,-1,0)$ & $0-(-\frac12)=\frac12$ & intriqué \\
$\frac12\ket{00}+\frac12\ket{01}+\frac{1}{\sqrt2}\ket{11}$ & $(\frac12,\frac12,0,\frac{1}{\sqrt2})$ & $\frac{1}{2\sqrt2}-0$ & intriqué \\
\bottomrule
\end{tabular}
\end{center}

Les quatre états de Bell $\ket{\Phi^\pm}$ et $\ket{\Psi^\pm}$ de la table sont tous intriqués
($D=\pm\frac12$) ; leurs propriétés font l'objet de la section~\ref{sec:bell}.

\paragraph{Preuve directe que $\ket{\Phi^+}$ est intriqué.} Supposons
$\ket{\Phi^+}=(a\ket{0}+b\ket{1})\otimes(c\ket{0}+d\ket{1})$. En identifiant les
composantes : $ac=\frac{1}{\sqrt2}$, $ad=0$, $bc=0$, $bd=\frac{1}{\sqrt2}$. La condition
$ad=0$ impose $a=0$ (impossible car $ac\neq0$) ou $d=0$ (impossible car $bd\neq0$).
Contradiction : aucun produit ne convient.

\paragraph{Corrélations.} Pour un produit $\ket{\Psi}_A\otimes\ket{\Psi}_B$ avec
$\ket{\Psi}_A=a_0\ket{0}+a_1\ket{1}$ et $\ket{\Psi}_B=b_0\ket{0}+b_1\ket{1}$, on a
$P(ij)=|a_ib_j|^2=|a_i|^2|b_j|^2=P(A{=}i)\,P(B{=}j)$ : les deux qubits sont
\emph{indépendants}. Vérification avec $\ket{+}\otimes\ket{+}$ : $P(00)=\frac14=\frac12\cdot\frac12$.
Pour $\ket{\Phi^+}$ : $P(00)=P(11)=\frac12$ et $P(01)=P(10)=0$, donc
$P(A{=}0)=P(B{=}0)=\frac12$, mais
\[
  P(00)=\tfrac12\neq\tfrac12\cdot\tfrac12=P(A{=}0)\,P(B{=}0) .
\]
Les deux qubits donnent toujours le même résultat, aléatoire : ils sont corrélés.
La figure~\ref{fig:epr} illustre la différence entre les deux cas.

\begin{figure}[!ht]
\centering
\adjustbox{max width=\linewidth}{%
\begin{tikzpicture}[>=Stealth, font=\small,
    labo/.style={draw=insagris, fill=insagris!8, rounded corners=2pt, thick, minimum width=2.6cm,
                 minimum height=1.0cm, align=center},
    source/.style={draw=insarouge, fill=insarouge!10, rounded corners=2pt, thick, minimum width=3.6cm,
                   minimum height=1.0cm, align=center},
    droit/.style={->, thick, insagris},
    onde/.style={->, thick, insarouge, decorate,
                 decoration={snake, amplitude=1.3pt, segment length=7pt, post length=2mm}},
    res/.style={font=\small, align=center}]
  % (a) état séparable
  \node[labo] (Aa) at (0,0) {Alice\\ mesure $A$};
  \node[source] (Sa) at (6,0) {source\\ $\ket{0}_A\otimes\ket{+}_B$};
  \node[labo] (Ba) at (12,0) {Bob\\ mesure $B$};
  \draw[droit] (Sa.west) -- (Aa.east);
  \draw[droit] (Sa.east) -- (Ba.west);
  \node[res, below=3pt of Aa] {$0\ \ 0\ \ 0\ \ 0\ \ 0\ \ 0$};
  \node[res, below=3pt of Ba] {$1\ \ 0\ \ 0\ \ 1\ \ 0\ \ 1$};
  \node[res, text=insagris, font=\footnotesize] at (6,-1.0) {$A$ donne toujours $0$ ; $B$ est aléatoire.\\
        Les deux résultats sont indépendants.};
  \node[anchor=south west, font=\small\bfseries] at (-1.3,0.75) {(a) état séparable};
  % (b) état intriqué
  \begin{scope}[yshift=-3.6cm]
    \node[labo] (Ab) at (0,0) {Alice\\ mesure $A$};
    \node[source] (Sb) at (6,0) {source\\ $\ket{\Phi^+}_{AB}$};
    \node[labo] (Bb) at (12,0) {Bob\\ mesure $B$};
    \draw[onde] (Sb.west) -- (Ab.east);
    \draw[onde] (Sb.east) -- (Bb.west);
    \node[res, below=3pt of Ab] {$1\ \ 0\ \ 0\ \ 1\ \ 1\ \ 0$};
    \node[res, below=3pt of Bb] {$1\ \ 0\ \ 0\ \ 1\ \ 1\ \ 0$};
    \node[res, text=insagris, font=\footnotesize] at (6,-1.0) {$A$ et $B$ sont chacun aléatoires,\\
          mais toujours égaux : ils sont corrélés.};
    \node[anchor=south west, font=\small\bfseries] at (-1.3,0.75) {(b) état intriqué};
  \end{scope}
\end{tikzpicture}}
\caption{Six mesures successives sur des paires préparées dans le même état (les résultats sont
des exemples). (a) Pour un produit, chaque qubit a son propre comportement. (b) Pour
$\ket{\Phi^+}=\frac{1}{\sqrt2}(\ket{00}+\ket{11})$, chaque résultat pris seul est un tirage à pile ou
face, mais les deux résultats coïncident toujours ; la ligne ondulée symbolise l'intrication.}
\label{fig:epr}
\end{figure}

\paragraph{Et un qubit pris seul ?} Dans un état séparable, chaque qubit garde son propre état pur ;
dans un état intriqué, aucun ket ne décrit un qubit isolé. On le décrit alors par une matrice de
densité, obtenue par la trace partielle (section~\ref{sec:trace}).

% ---------------------------------------------------------------------
\subsection{Les états de Bell}\label{sec:bell}
% ---------------------------------------------------------------------

\begin{definition}{Les quatre états de Bell}
\[
  \ket{\Phi^{\pm}}=\frac{1}{\sqrt2}\big(\ket{00}\pm\ket{11}\big),
  \qquad
  \ket{\Psi^{\pm}}=\frac{1}{\sqrt2}\big(\ket{01}\pm\ket{10}\big).
\]
Ils sont intriqués ($|D|=\frac12$, section~\ref{sec:separable}) et le sont au maximum :
$|D|\leq\frac12$ pour tout état normalisé de deux qubits, et chaque qubit d'un état de Bell,
pris seul, est complètement aléatoire (section~\ref{sec:trace}).
\end{definition}

\paragraph{Une base de $\mathcal{H}^4$.} Chaque état est de norme $\frac12+\frac12=1$. Les
$\ket{\Phi^\pm}$ n'ont de composantes que sur $\ket{00}$ et $\ket{11}$, les $\ket{\Psi^\pm}$ que sur
$\ket{01}$ et $\ket{10}$, donc $\braket{\Phi^\pm}{\Psi^\pm}=0$ pour toutes les combinaisons de signes ;
enfin $\braket{\Phi^+}{\Phi^-}=\frac12(1-1)=0$ et $\braket{\Psi^+}{\Psi^-}=\frac12(1-1)=0$. Les quatre
états forment donc une base orthonormée de $\mathcal{H}^4$, la \emph{base de Bell}, constituée
uniquement d'états intriqués. Le circuit de la section~\ref{sec:circuit-bell} fait passer de la
base de calcul à cette base.

\paragraph{Ce qui les distingue : deux parités.} Sur la base de calcul,
$Z\otimes Z\ket{ij}=(-1)^{i+j}\ket{ij}$ (valeur propre $+1$ si les deux chiffres sont égaux, $-1$
sinon) et $X\otimes X\ket{ij}=\ket{\bar i\bar j}$ (les deux chiffres sont inversés, avec
$\bar i=1-i$). On en déduit, par exemple,
\[
  X\otimes X\ket{\Phi^-}=\tfrac{1}{\sqrt2}\big(\ket{11}-\ket{00}\big)=-\ket{\Phi^-},
  \qquad
  Z\otimes Z\ket{\Psi^+}=\tfrac{1}{\sqrt2}\big(-\ket{01}-\ket{10}\big)=-\ket{\Psi^+},
\]
et le tableau complet des valeurs propres :
\begin{center}
\renewcommand{\arraystretch}{1.3}
\begin{tabular}{@{}lcccc@{}}
\toprule
 & $\ket{\Phi^+}$ & $\ket{\Phi^-}$ & $\ket{\Psi^+}$ & $\ket{\Psi^-}$ \\
\midrule
$Z\otimes Z$ & $+1$ & $+1$ & $-1$ & $-1$ \\
$X\otimes X$ & $+1$ & $-1$ & $+1$ & $-1$ \\
\bottomrule
\end{tabular}
\end{center}
Les deux opérateurs commutent, car $ZX=-XZ$ donne
$(Z\otimes Z)(X\otimes X)=ZX\otimes ZX=XZ\otimes XZ=(X\otimes X)(Z\otimes Z)$ ; les quatre états de
Bell sont leurs états propres communs, et chaque couple de valeurs propres apparaît une fois
(figure~\ref{fig:grille-bell}). Le même calcul montre que $X\otimes Z$ et $Z\otimes X$
commutent : $(X\otimes Z)(Z\otimes X)=XZ\otimes ZX=(-ZX)\otimes(-XZ)=ZX\otimes XZ=(Z\otimes X)(X\otimes Z)$.
Plus généralement, $P\otimes Q$ et $P'\otimes Q'$ (matrices de Pauli) commutent si $P,P'$
anticommutent et si $Q,Q'$ anticommutent : les deux signes $-1$ se compensent.

\begin{figure}[!ht]
\centering
\adjustbox{max width=\linewidth}{%
\begin{tikzpicture}[>=Stealth, font=\small,
    cell/.style={draw=insarouge, fill=insarouge!6, rounded corners=2pt, thick, minimum width=5.4cm,
                 minimum height=1.35cm, align=center},
    tete/.style={text=insagris, align=center}]
  \node[tete] at (3.0,2.0) {$X\otimes X=+1$\\ {\footnotesize mêmes résultats en base $X$}};
  \node[tete] at (8.8,2.0) {$X\otimes X=-1$\\ {\footnotesize résultats opposés en base $X$}};
  \node[tete, anchor=east] at (-0.1,0.75) {$Z\otimes Z=+1$\\ {\footnotesize mêmes résultats}\\ {\footnotesize en base $Z$}};
  \node[tete, anchor=east] at (-0.1,-0.85) {$Z\otimes Z=-1$\\ {\footnotesize résultats opposés}\\ {\footnotesize en base $Z$}};
  \node[cell] at (3.0,0.75) {$\ket{\Phi^+}=\frac{1}{\sqrt2}\big(\ket{00}+\ket{11}\big)$};
  \node[cell] at (8.8,0.75) {$\ket{\Phi^-}=\frac{1}{\sqrt2}\big(\ket{00}-\ket{11}\big)$};
  \node[cell] at (3.0,-0.85) {$\ket{\Psi^+}=\frac{1}{\sqrt2}\big(\ket{01}+\ket{10}\big)$};
  \node[cell] at (8.8,-0.85) {$\ket{\Psi^-}=\frac{1}{\sqrt2}\big(\ket{01}-\ket{10}\big)$};
\end{tikzpicture}}
\caption{Les quatre états de Bell, classés par les valeurs propres de $Z\otimes Z$ (lignes) et de
$X\otimes X$ (colonnes). La valeur propre est le produit des deux résultats $\pm1$ d'une mesure des
deux qubits dans la base correspondante : $+1$ si les résultats sont égaux, $-1$ s'ils sont opposés.}
\label{fig:grille-bell}
\end{figure}

Les corrélations ne se limitent donc pas à la base $Z$ : $\ket{\Phi^+}$ s'écrit aussi
$\frac{1}{\sqrt2}\big(\ket{++}+\ket{--}\big)$ et donne des résultats égaux en base $X$ comme en base $Z$
(exercice~\ref{ex:bellX}).

\paragraph{D'un état de Bell à l'autre.} Une porte agissant sur un seul qubit suffit. Avec
$X\ket{0}=\ket{1}$, $X\ket{1}=\ket{0}$ et $Z\ket{1}=-\ket{1}$, en agissant sur le second qubit :
\begin{align*}
  (I\otimes X)\ket{\Phi^+}&=\tfrac{1}{\sqrt2}\big(\ket{01}+\ket{10}\big)=\ket{\Psi^+},\\
  (I\otimes Z)\ket{\Phi^+}&=\tfrac{1}{\sqrt2}\big(\ket{00}-\ket{11}\big)=\ket{\Phi^-},\\
  (I\otimes XZ)\ket{\Phi^+}&=(I\otimes X)\tfrac{1}{\sqrt2}\big(\ket{00}-\ket{11}\big)
    =\tfrac{1}{\sqrt2}\big(\ket{01}-\ket{10}\big)=\ket{\Psi^-}.
\end{align*}
Quatre opérations locales ($I$, $X$, $Z$, $XZ$) sur \emph{un} qubit d'une paire de Bell mènent donc
aux quatre états de la base ; c'est le principe du codage superdense
(exercice~\ref{ex:superdense}).

\begin{remarque}[Le singulet est le même dans toutes les bases]\label{rem:singulet}
Pour toute matrice unitaire $U=\begin{pmatrix}a&b\\c&d\end{pmatrix}$, on a $U\ket{0}=a\ket{0}+c\ket{1}$
et $U\ket{1}=b\ket{0}+d\ket{1}$, donc
\begin{align*}
  (U\otimes U)\big(\ket{01}-\ket{10}\big)
  &=\big(a\ket{0}+c\ket{1}\big)\big(b\ket{0}+d\ket{1}\big)-\big(b\ket{0}+d\ket{1}\big)\big(a\ket{0}+c\ket{1}\big)\\
  &=(ad-bc)\big(\ket{01}-\ket{10}\big),
\end{align*}
c'est-à-dire $(U\otimes U)\ket{\Psi^-}=\det(U)\,\ket{\Psi^-}$. Comme $|\det U|=1$, ce facteur est une
phase globale, sans effet physique : $\ket{\Psi^-}$ s'écrit \emph{de la même façon} dans toutes les
bases orthonormées (avec $U=H$, de déterminant $-1$ : $\ket{\Psi^-}=\frac{1}{\sqrt2}\big(\ket{-+}-\ket{+-}\big)$).
Si Alice et Bob mesurent dans la \emph{même} base, quelle qu'elle soit, leurs résultats sont
toujours opposés : c'est le «~singulet~». Parmi les états de Bell, seul $\ket{\Psi^-}$ a cette propriété.
\end{remarque}

% ---------------------------------------------------------------------
\subsection{Un circuit quantique : le générateur d'états de Bell}\label{sec:circuit-bell}
% ---------------------------------------------------------------------

On assemble une porte de Hadamard et un CNOT. Le système est $(A,B)$ avec le ket
$\ket{a,b}=\ket{a}_A\otimes\ket{b}_B$ ; d'après la convention de la section~\ref{sec:circuits}, $B$ est dessiné
en haut et $A$ en bas (figure~\ref{fig:circuit-bell}). Le circuit a deux étapes :
\begin{enumerate}[nosep]
  \item la porte $I\otimes H$ : identité sur $A$, Hadamard sur $B$ ;
  \item un CNOT de contrôle $B$ et de cible $A$ (celui de la section~\ref{sec:cnot}, dont le
        contrôle est le second chiffre).
\end{enumerate}
Les traits pointillés délimitent les états $\ket{\Pi_0}$, $\ket{\Pi_1}$ et $\ket{\Pi_2}$ :
avant les portes, après $I\otimes H$, après le CNOT. Le circuit transforme la base de calcul
en base de Bell (section~\ref{sec:bell}).

\begin{figure}[!ht]
\centering
\begin{tikzpicture}[>=Stealth, font=\small]
  % fils : B en haut, A en bas
  \draw (0,1.4) -- (7.6,1.4);
  \draw (0,0) -- (7.6,0);
  \node[anchor=east] at (0,1.4) {$\ket{0}_B$};
  \node[anchor=east] at (0,0) {$\ket{0}_A$};
  % première étape : H sur B, identité (pointillée) sur A
  \node[porte] at (2.2,1.4) {$H$};
  \node[draw=insagris, dashed, fill=white, minimum size=7mm] at (2.2,0) {$I$};
  % CNOT : contrôle B (haut), cible A (bas)
  \draw[insarouge, thick] (5,1.4) -- (5,0);
  \pic at (5,1.4) {ctl};
  \pic at (5,0) {cible};
  % accolade de sortie
  \draw[insagris, thick] (7.75,1.6) -- (7.9,1.6) -- (7.9,-0.2) -- (7.75,-0.2);
  \node[anchor=west] at (7.9,0.7) {$\ket{\beta^{00}}_{AB}$};
  % états intermédiaires
  \foreach \x in {0.7,3.6,6.4}
    \draw[dashed, insagris!70] (\x,1.9) -- (\x,-0.7);
  \node[anchor=north, font=\footnotesize] at (0.7,-0.7) {$\ket{\Pi_0}=\ket{00}$};
  \node[anchor=north, font=\footnotesize] at (3.6,-0.7) {$\ket{\Pi_1}=\ket{0}\otimes\ket{+}$};
  \node[anchor=north, font=\footnotesize] at (6.4,-0.7) {$\ket{\Pi_2}=\ket{\beta^{00}}$};
\end{tikzpicture}
\caption{Générateur d'états de Bell : $I\otimes H$, puis un CNOT de contrôle $B$ (fil du haut)
et de cible $A$ (fil du bas). L'entrée représentée est $\ket{00}$.}
\label{fig:circuit-bell}
\end{figure}

\subsubsection{Entrée \texorpdfstring{$\ket{00}$}{|00>} : calcul pas à pas}

\begin{align*}
  \ket{\Pi_0}&=\ket{0}_A\otimes\ket{0}_B=\ket{00},\\[1ex]
  \ket{\Pi_1}&=(I\otimes H)\ket{\Pi_0}=I\ket{0}_A\otimes H\ket{0}_B=\ket{0}\otimes\ket{+}
    =\frac{\ket{00}+\ket{01}}{\sqrt2},\\[1ex]
  \ket{\Pi_2}&=\mathrm{CNOT}\ket{\Pi_1}=\frac{\mathrm{CNOT}\ket{00}+\mathrm{CNOT}\ket{01}}{\sqrt2}
    =\frac{\ket{00}+\ket{11}}{\sqrt2}=\ket{\beta^{00}}=\ket{\Phi^+}.
\end{align*}
Dans $\ket{\Pi_1}$, le premier chiffre est celui de $A$ et le second celui de $B$. Le CNOT a
pour contrôle $B$ : $\ket{00}$ ($B=0$) ne change pas, $\ket{01}$ ($B=1$) devient $\ket{11}$
puisque $A$ est inversé. Le résultat est un \emph{état de Bell}.

Avec les matrices, la règle $A\otimes B=(A_{ij}B)$ donne
\[
  I\otimes H=\begin{pmatrix}1\cdot H&0\\0&1\cdot H\end{pmatrix}
  =\frac{1}{\sqrt2}\begin{pmatrix}1&1&0&0\\1&-1&0&0\\0&0&1&1\\0&0&1&-1\end{pmatrix}.
\]
Sa première colonne est $\ket{\Pi_1}=\frac{1}{\sqrt2}(1,1,0,0)^{\mathsf T}$ (image de $\ket{00}$).
Le CNOT échange les amplitudes de $\ket{01}$ et de $\ket{11}$ (section~\ref{sec:cnot}), donc
\[
  \ket{\Pi_2}=\mathrm{CNOT}\,\frac{1}{\sqrt2}\begin{pmatrix}1\\1\\0\\0\end{pmatrix}
  =\frac{1}{\sqrt2}\begin{pmatrix}1\\0\\0\\1\end{pmatrix}.
\]

\subsubsection{Les quatre entrées}

Prenons $B$ dans l'état $\ket{x}$ et $A$ dans l'état $\ket{y}$, avec $x,y\in\{0,1\}$. Alors
$\ket{\Pi_0}=\ket{y}_A\otimes\ket{x}_B=\ket{y,x}$. Comme
$H\ket{x}=\frac{1}{\sqrt2}\big(\ket{0}+(-1)^x\ket{1}\big)$ (c'est $\ket{+}$ pour $x=0$ et
$\ket{-}$ pour $x=1$),
\[
  \ket{\Pi_1}=\ket{y}\otimes H\ket{x}=\frac{1}{\sqrt2}\Big(\ket{y,0}+(-1)^x\ket{y,1}\Big).
\]
Le CNOT (contrôle $B$) ne change pas $\ket{y,0}$ et transforme $\ket{y,1}$ en
$\ket{\bar y,1}$, avec $\bar y=1-y$ :
\[
  \ket{\Pi_2}=\ket{\beta^{xy}}=\frac{1}{\sqrt2}\Big(\ket{y,0}+(-1)^x\ket{\bar y,1}\Big).
\]
Le premier indice de $\ket{\beta^{xy}}$ est l'entrée de $B$ (le qubit qui reçoit $H$), le
second celle de $A$ (la cible). Les quatre cas :

\begin{center}
\renewcommand{\arraystretch}{1.7}
\begin{tabular}{@{}ccccc@{}}
\toprule
$(B,A)=(x,y)$ & $\ket{\Pi_0}=\ket{y,x}$ & $\ket{\Pi_1}$ & $\ket{\Pi_2}=\ket{\beta^{xy}}$ & Nom \\
\midrule
$(0,0)$ & $\ket{00}$ & $\frac{1}{\sqrt2}(\ket{00}+\ket{01})$ & $\frac{1}{\sqrt2}(\ket{00}+\ket{11})$ & $\ket{\Phi^+}$ \\
$(1,0)$ & $\ket{01}$ & $\frac{1}{\sqrt2}(\ket{00}-\ket{01})$ & $\frac{1}{\sqrt2}(\ket{00}-\ket{11})$ & $\ket{\Phi^-}$ \\
$(0,1)$ & $\ket{10}$ & $\frac{1}{\sqrt2}(\ket{10}+\ket{11})$ & $\frac{1}{\sqrt2}(\ket{10}+\ket{01})$ & $\ket{\Psi^+}$ \\
$(1,1)$ & $\ket{11}$ & $\frac{1}{\sqrt2}(\ket{10}-\ket{11})$ & $\frac{1}{\sqrt2}(\ket{10}-\ket{01})$ & $-\ket{\Psi^-}$ \\
\bottomrule
\end{tabular}
\end{center}
Détail de la dernière ligne : $\ket{\Pi_1}=\ket{1}\otimes\ket{-}=\frac{1}{\sqrt2}(\ket{10}-\ket{11})$ ;
le CNOT laisse $\ket{10}$ et envoie $\ket{11}$ sur $\ket{01}$, d'où
$\frac{1}{\sqrt2}(\ket{10}-\ket{01})=-\frac{1}{\sqrt2}(\ket{01}-\ket{10})=-\ket{\Psi^-}$.

\paragraph{Forme matricielle.} Le CNOT échange les lignes $2$ et $4$ de
$I\otimes H$ (il échange les amplitudes de $\ket{01}$ et $\ket{11}$), donc
\[
  \mathrm{CNOT}\,(I\otimes H)=\frac{1}{\sqrt2}\begin{pmatrix}
    1 & 1 & 0 & 0\\ 0 & 0 & 1 & -1\\ 0 & 0 & 1 & 1\\ 1 & -1 & 0 & 0
  \end{pmatrix}.
\]
Ses colonnes sont les images de $\ket{00},\ket{01},\ket{10},\ket{11}$, c'est-à-dire de
$(B,A)=(0,0),(1,0),(0,1),(1,1)$ : ce sont $\ket{\beta^{00}}$, $\ket{\beta^{10}}$,
$\ket{\beta^{01}}$, $\ket{\beta^{11}}$, soit $\ket{\Phi^+}$, $\ket{\Phi^-}$, $\ket{\Psi^+}$
et $-\ket{\Psi^-}$.

\subsubsection{Matrice de densité de l'état de sortie}

Pour $\ket{\Psi}=\ket{\beta^{00}}=\frac{1}{\sqrt2}\ket{00}+\frac{1}{\sqrt2}\ket{11}$ :
\[
  \rho_{AB}=\ket{\Psi}\bra{\Psi}
  =\begin{pmatrix}\frac{1}{\sqrt2}\\0\\0\\\frac{1}{\sqrt2}\end{pmatrix}
   \begin{pmatrix}\frac{1}{\sqrt2}&0&0&\frac{1}{\sqrt2}\end{pmatrix}
  =\frac12\begin{pmatrix}1&0&0&1\\0&0&0&0\\0&0&0&0\\1&0&0&1\end{pmatrix}.
\]
Sous forme de ket-bra :
$\rho_{AB}=\frac12\big(\ket{00}\bra{00}+\ket{00}\bra{11}+\ket{11}\bra{00}+\ket{11}\bra{11}\big)$.
On vérifie $\operatorname{Tr}\rho_{AB}=\frac12(1+1)=1$ et
\[
  \rho_{AB}^2=\frac14\begin{pmatrix}2&0&0&2\\0&0&0&0\\0&0&0&0\\2&0&0&2\end{pmatrix}=\rho_{AB},
\]
donc $\operatorname{Tr}\rho_{AB}^2=1$ : le couple $(A,B)$ est dans un état \emph{pur}.

\begin{remarque}[Porte intriquante]
Les états d'entrée $\ket{ij}$ sont séparables ($D=0$) ; les sorties sont intriquées
($D=\pm\frac12$). Aucune porte agissant séparément sur $A$ et sur $B$ ($U_A\otimes U_B$,
section~\ref{sec:op2}) ne peut créer de l'intrication à partir d'un produit, puisque
$(U_A\otimes U_B)(\ket{u}\otimes\ket{v})=U_A\ket{u}\otimes U_B\ket{v}$ reste un produit.
Une porte à deux qubits comme le CNOT est indispensable.
\end{remarque}

% ---------------------------------------------------------------------
\subsection{Trace partielle : l'état d'un qubit pris seul}\label{sec:trace}
% ---------------------------------------------------------------------

Un état intriqué n'a pas de ket pour $A$ seul (section~\ref{sec:separable}). On peut pourtant
décrire $A$ par une matrice de densité, obtenue en «~oubliant~» $B$.

\begin{definition}{Trace partielle}
Soit $\rho_{AB}$ la matrice de densité du couple $(A,B)$. L'état du qubit $A$ pris seul est
$\rho_A=\operatorname{Tr}_B(\rho_{AB})$, celui de $B$ est $\rho_B=\operatorname{Tr}_A(\rho_{AB})$. La
trace partielle est linéaire et remplace la partie tracée de chaque terme par sa trace :
\[
  \operatorname{Tr}_B\big(\ket{a}\bra{a'}\otimes\ket{b}\bra{b'}\big)=\ket{a}\bra{a'}\;\braket{b'}{b},
  \qquad
  \operatorname{Tr}_A\big(\ket{a}\bra{a'}\otimes\ket{b}\bra{b'}\big)=\braket{a'}{a}\;\ket{b}\bra{b'}.
\]
\end{definition}

\paragraph{En matrices.} On découpe $\rho_{AB}$ (matrice $4\times4$, lignes et colonnes ordonnées
$00,01,10,11$) en quatre blocs $2\times2$ :
\[
  \rho_{AB}=\begin{pmatrix}R_{00}&R_{01}\\R_{10}&R_{11}\end{pmatrix}
  \quad\Longrightarrow\quad
  \rho_A=\begin{pmatrix}\operatorname{Tr}R_{00}&\operatorname{Tr}R_{01}\\
                        \operatorname{Tr}R_{10}&\operatorname{Tr}R_{11}\end{pmatrix},
  \qquad
  \rho_B=R_{00}+R_{11}.
\]
Autrement dit $(\rho_A)_{ij}=\rho_{(i0),(j0)}+\rho_{(i1),(j1)}$ : $\rho_A$ est la matrice des traces
des blocs, et $\rho_B$ la somme des deux blocs diagonaux.

\begin{remarque}[Pourquoi cette définition]
Toutes les prédictions qui portent sur $A$ seul viennent d'opérateurs de la forme $O_A\otimes I$, et
\[
  \operatorname{Tr}\big(\rho_{AB}\,(O_A\otimes I)\big)=\operatorname{Tr}\big(\rho_A\,O_A\big)
\]
(on le vérifie sur $\ket{a}\bra{a'}\otimes\ket{b}\bra{b'}$, puis par linéarité). La matrice $\rho_A$
contient donc exactement ce que l'on peut prédire sur $A$ sans regarder $B$.
\end{remarque}

\paragraph{Exemple 1 : état séparable $\ket{0}\otimes\ket{+}$.} Ici
$\rho_{AB}=\ket{0}\bra{0}\otimes\ket{+}\bra{+}$, donc
\[
  \rho_A=\ket{0}\bra{0}\,\braket{+}{+}=\ket{0}\bra{0},
  \qquad
  \rho_B=\braket{0}{0}\,\ket{+}\bra{+}=\ket{+}\bra{+}.
\]
Avec les matrices,
\[
  \rho_{AB}=\frac12\begin{pmatrix}1&1&0&0\\1&1&0&0\\0&0&0&0\\0&0&0&0\end{pmatrix},
  \qquad
  R_{00}=\frac12\begin{pmatrix}1&1\\1&1\end{pmatrix},\quad R_{01}=R_{10}=R_{11}=0,
\]
d'où
\[
  \rho_A=\begin{pmatrix}\operatorname{Tr}R_{00}&0\\0&0\end{pmatrix}=\begin{pmatrix}1&0\\0&0\end{pmatrix},
  \qquad
  \rho_B=R_{00}+R_{11}=\frac12\begin{pmatrix}1&1\\1&1\end{pmatrix}.
\]
Un qubit d'un état séparable garde son propre état pur.

\paragraph{Exemple 2 : état de Bell $\ket{\Phi^+}$.} On écrit
$\rho_{AB}=\ket{\Phi^+}\bra{\Phi^+}=\frac12\big(\ket{00}\bra{00}+\ket{00}\bra{11}+\ket{11}\bra{00}
+\ket{11}\bra{11}\big)$. Chaque terme est un produit d'opérateurs à un qubit, par exemple
$\ket{00}\bra{11}=\ket{0}\bra{1}\otimes\ket{0}\bra{1}$, donc
\begin{align*}
  \rho_A=\operatorname{Tr}_B(\rho_{AB})
  &=\tfrac12\Big(\ket{0}\bra{0}\,\braket{0}{0}+\ket{0}\bra{1}\,\braket{1}{0}
      +\ket{1}\bra{0}\,\braket{0}{1}+\ket{1}\bra{1}\,\braket{1}{1}\Big)\\
  &=\tfrac12\Big(\ket{0}\bra{0}\cdot1+\ket{0}\bra{1}\cdot0+\ket{1}\bra{0}\cdot0+\ket{1}\bra{1}\cdot1\Big)
   =\tfrac12\big(\ket{0}\bra{0}+\ket{1}\bra{1}\big)=\tfrac12\,\mathrm{Id}.
\end{align*}
Par les blocs,
\[
  \rho_{AB}=\frac12\begin{pmatrix}1&0&0&1\\0&0&0&0\\0&0&0&0\\1&0&0&1\end{pmatrix},
  \qquad
  R_{00}=\frac12\begin{pmatrix}1&0\\0&0\end{pmatrix},\quad
  R_{11}=\frac12\begin{pmatrix}0&0\\0&1\end{pmatrix},\quad
  \operatorname{Tr}R_{01}=\operatorname{Tr}R_{10}=0,
\]
d'où $\rho_A=\frac12\mathrm{Id}$ et $\rho_B=R_{00}+R_{11}=\frac12\mathrm{Id}$ : le même résultat.

Le qubit $A$ est donc dans l'état \emph{complètement mixte} ($\operatorname{Tr}\rho_A^2=\frac14+\frac14=\frac12$),
au centre de la boule de Bloch (section~\ref{sec:bloch}), alors que le couple est dans un état pur
($\operatorname{Tr}\rho_{AB}^2=1$, section~\ref{sec:circuit-bell}). On connaît parfaitement l'état du
couple mais pas celui d'un qubit pris seul : l'information est entièrement contenue dans les
corrélations entre $A$ et $B$. Mesuré seul, $A$ donne $0$ ou $1$ avec la probabilité $\frac12$, dans
toutes les bases.

\paragraph{Formule pratique pour un état pur.} Pour $\ket{\Psi}=\sum_{ij}c_{ij}\ket{ij}$, notons
$C=\begin{pmatrix}c_{00}&c_{01}\\c_{10}&c_{11}\end{pmatrix}$ la matrice $2\times2$ des coefficients
(ligne $i$ : état de $A$, colonne $j$ : état de $B$). Comme $\rho_{(ij),(i'j')}=c_{ij}\,c_{i'j'}^*$,
\[
  (\rho_A)_{ii'}=\sum_jc_{ij}\,c_{i'j}^*,\qquad\text{soit}\qquad \rho_A=C\,C^\dagger .
\]
\emph{Exemple 3.} Pour $\ket{\Psi}=\frac12\ket{00}+\frac12\ket{01}+\frac{1}{\sqrt2}\ket{11}$ (exemple~2
de la section~\ref{sec:mesure2}), $C=\begin{pmatrix}\frac12&\frac12\\0&\frac{1}{\sqrt2}\end{pmatrix}$ et
\[
  \rho_A=CC^\dagger=\begin{pmatrix}\frac14+\frac14&0+\frac{1}{2\sqrt2}\\[0.4ex]\frac{1}{2\sqrt2}&\frac12\end{pmatrix}
  =\begin{pmatrix}\frac12&\frac{\sqrt2}{4}\\[0.4ex]\frac{\sqrt2}{4}&\frac12\end{pmatrix},
  \qquad
  \operatorname{Tr}\rho_A^2=\tfrac14+\tfrac14+2\cdot\tfrac18=\tfrac34 .
\]
Le qubit $A$ est mixte, mais pas complètement : cet état est intriqué sans être un état de Bell.

\begin{resultat}{Pureté d'un qubit et séparabilité}
Pour un état pur $\ket{\Psi}=\sum c_{ij}\ket{ij}$ de deux qubits, avec $D=c_{00}c_{11}-c_{01}c_{10}$,
\[
  \operatorname{Tr}\rho_A^2=\operatorname{Tr}\rho_B^2=1-2|D|^2 .
\]
Donc $\ket{\Psi}$ est séparable ($D=0$) si et seulement si $\rho_A$ est pur. Plus $|D|$ est grand,
plus $\rho_A$ est mixte, jusqu'à $|D|=\frac12$ et $\rho_A=\mathrm{Id}/2$ pour les états de Bell.
\end{resultat}

\emph{Démonstration.} La matrice $\rho_A=CC^\dagger$ a pour trace $\sum|c_{ij}|^2=1$ et pour
déterminant $\det C\,\overline{\det C}=|D|^2$. Si $\lambda_1,\lambda_2$ sont ses valeurs propres,
$\lambda_1+\lambda_2=1$ et $\operatorname{Tr}\rho_A^2=\lambda_1^2+\lambda_2^2=1-2\lambda_1\lambda_2
=1-2|D|^2$ ; le calcul est le même pour $\rho_B$. \hfill$\square$

Comme $\operatorname{Tr}\rho_A^2\geq\frac12$ pour un qubit, on en déduit $|D|\leq\frac12$. Avec
$\operatorname{Tr}\rho_A^2=\frac{1+|\vec r|^2}{2}$ (section~\ref{sec:bloch}), le vecteur de Bloch de $A$
a pour norme $|\vec r|=\sqrt{1-4|D|^2}$ : il part de la sphère pour un état séparable et se rapproche
du centre quand l'intrication croît (figure~\ref{fig:trace-bloch}).

\begin{center}
\renewcommand{\arraystretch}{1.5}
\begin{tabular}{@{}lcccl@{}}
\toprule
\textbf{État du couple} & $D$ & $\operatorname{Tr}\rho_A^2$ & $|\vec r|$ & \textbf{Qubit $A$} \\
\midrule
$\ket{0}\otimes\ket{+}$ & $0$ & $1$ & $1$ & pur, sur la sphère \\
$\frac12\ket{00}+\frac12\ket{01}+\frac{1}{\sqrt2}\ket{11}$ & $\frac{\sqrt2}{4}$ & $\frac34$
  & $\frac{1}{\sqrt2}$ & mixte, dans la boule \\
$\ket{\Phi^+}$ (comme les trois autres états de Bell) & $\pm\frac12$ & $\frac12$ & $0$
  & complètement mixte, au centre \\
\bottomrule
\end{tabular}
\end{center}

\begin{figure}[!ht]
\centering
\adjustbox{max width=\linewidth}{%
\begin{tikzpicture}[>=Stealth, font=\small]
  \def\R{2.3}
  \fill[insagris!12] (0,0) circle (\R);
  \draw[insagris, dotted] (-\R,0) -- (\R,0);
  \draw[insagris, dotted] (0,-\R) -- (0,\R);
  \draw[insagris, dashed] (0,0) circle ({\R*0.70711});
  \draw[insarouge, very thick] (0,0) circle (\R);
  \node[font=\footnotesize, text=insagris, anchor=south west] at (\R,0) {$x$};
  \node[font=\footnotesize, text=insagris, anchor=south west] at (0.03,\R) {$z$};
  % vecteurs de Bloch de rho_A pour trois états du couple
  \fill[insarouge] (0,\R) circle (2.4pt);
  \fill[insagris!80!black] ({\R*0.70711},0) circle (2.4pt);
  \fill[insagris!80!black] (0,0) circle (2.4pt);
  \draw[insagris] (0,\R) -- (3.1,\R);
  \node[anchor=west, font=\footnotesize, align=left] at (3.1,\R)
    {$\ket{0}\otimes\ket{+}$ : $\rho_A=\ket{0}\bra{0}$\\ $D=0$, \ $|\vec r|=1$ (état pur)};
  \draw[insagris] ({\R*0.70711},0) -- (3.1,0);
  \node[anchor=west, font=\footnotesize, align=left] at (3.1,0)
    {$\frac12\ket{00}+\frac12\ket{01}+\frac{1}{\sqrt2}\ket{11}$\\
     $D=\frac{\sqrt2}{4}$, \ $|\vec r|=\frac{1}{\sqrt2}$ (mixte)};
  \draw[insagris] (0,0) -- (3.1,-\R);
  \node[anchor=west, font=\footnotesize, align=left] at (3.1,-\R)
    {$\ket{\Phi^+}$ : $\rho_A=\mathrm{Id}/2$\\ $D=\frac12$, \ $\vec r=\vec 0$ (centre)};
\end{tikzpicture}}
\caption{Le vecteur de Bloch de $\rho_A$ (boule de Bloch vue en coupe) pour trois états du couple
$(A,B)$. Le cercle pointillé est la sphère de rayon $1/\sqrt2$. Plus l'état est intriqué, plus le
point s'éloigne de la sphère des états purs : $|\vec r|=\sqrt{1-4|D|^2}$.}
\label{fig:trace-bloch}
\end{figure}

% ---------------------------------------------------------------------
\subsection{Information : variables aléatoires et états quantiques}\label{sec:info}
% ---------------------------------------------------------------------

\begin{definition}{Entropie, information mutuelle, entropie conditionnelle}
Soit $A$ une variable aléatoire de loi $p(a)$. Son \emph{entropie de Shannon}, en bits, est
\[
  H(A)=-\sum_a p(a)\log_2p(a)\qquad(\text{avec }0\log_20=0).
\]
Elle mesure l'incertitude moyenne sur $A$ : une pièce équilibrée a $H=1$ bit. Pour un couple
$(A,B)$ de loi $p(a,b)$, $H(A,B)=-\sum_{a,b}p(a,b)\log_2p(a,b)$, et
\begin{align*}
  I(A;B)&=H(A)+H(B)-H(A,B) &&\text{(information mutuelle)},\\
  H(A|B)&=H(A,B)-H(B)=H(A)-I(A;B) &&\text{(entropie conditionnelle)}.
\end{align*}
\end{definition}

Les deux écritures de $H(A|B)$ sont égales : la définition de $I(A;B)$ donne
$H(A,B)=H(A)+H(B)-I(A;B)$, donc $H(A,B)-H(B)=H(A)-I(A;B)$. L'information mutuelle $I(A;B)$
mesure ce que la connaissance de $B$ apprend sur $A$ ; $H(A|B)$ est l'incertitude qui reste sur
$A$ quand on connaît $B$.

\subsubsection{Deux variables aléatoires indépendantes}

Alice et Bob lancent chacun une pièce équilibrée, sans lien entre les deux lancers :
$p(a,b)=\frac14$ pour les quatre couples. Chaque pièce donne $0$ ou $1$ avec la probabilité $\frac12$, donc
\begin{align*}
  H(A)&=-\Big(\tfrac12\log_2\tfrac12+\tfrac12\log_2\tfrac12\Big)=1\text{ bit},\\
  H(B)&=1\text{ bit (même calcul)},\\
  H(A,B)&=-4\cdot\tfrac14\log_2\tfrac14=2\text{ bits},\\
  I(A;B)&=H(A)+H(B)-H(A,B)=1+1-2=0,\\
  H(A|B)&=H(A)-I(A;B)=1-0=H(A)=1\text{ bit}.
\end{align*}
Connaître $B$ n'apprend rien sur $A$.

\subsubsection{Deux variables solidaires}

Les deux pièces tombent toujours du même côté : $p(0,0)=p(1,1)=\frac12$ et
$p(0,1)=p(1,0)=0$. Les lois de chaque pièce sont inchangées ($p(a{=}0)=p(0,0)+p(0,1)=\frac12$,
etc.), donc $H(A)=H(B)=1$ bit, mais
\begin{align*}
  H(A,B)&=-2\cdot\tfrac12\log_2\tfrac12-2\cdot0\log_20=1\text{ bit},\\
  I(A;B)&=1+1-1=1\text{ bit},\qquad H(A|B)=1-1=0.
\end{align*}
Connaître $B$ détermine $A$ : plus aucune incertitude ne reste.

\subsubsection{Deux qubits dans un état de Bell}

Pour un système quantique, on remplace la loi de probabilité par la matrice de densité.

\begin{definition}{Entropie de von Neumann}
L'\emph{entropie de von Neumann} d'un état $\rho$, de valeurs propres $\lambda_i$, est
\[
  S(\rho)=-\operatorname{Tr}(\rho\log_2\rho)=-\sum_i\lambda_i\log_2\lambda_i .
\]
Elle vaut $0$ pour un état pur et $1$ bit pour le qubit complètement mélangé $\mathrm{Id}/2$.
En quantique, $H(A)=S(\rho_A)$, $H(B)=S(\rho_B)$ et $H(A,B)=S(\rho_{AB})$ ; les définitions de
$I(A;B)$ et de $H(A|B)$ sont les mêmes.
\end{definition}

On prend l'état de Bell $\ket{\Phi^+}$ produit par le circuit de la
section~\ref{sec:circuit-bell}. Le couple est dans un état pur,
$\rho_{AB}=\ket{\Phi^+}\bra{\Phi^+}$ ; chaque qubit pris seul est dans l'état
$\rho_A=\rho_B=\frac12\mathrm{Id}$ (trace partielle, section~\ref{sec:trace}).

Calcul des entropies :
\begin{itemize}[nosep]
  \item $\rho_{AB}$ est pur : ses valeurs propres sont $1,0,0,0$, donc
        $H(A,B)=S(\rho_{AB})=-1\cdot\log_21=0$ ;
  \item $\rho_A$ et $\rho_B$ ont pour valeurs propres $\frac12,\frac12$, donc
        $H(A)=H(B)=-\big(\frac12\log_2\frac12+\frac12\log_2\frac12\big)=1$ bit.
\end{itemize}
D'où
\begin{align*}
  I(A;B)&=H(A)+H(B)-H(A,B)=1+1-0=2\text{ bits},\\
  H(A|B)&=H(A,B)-H(B)=0-1=-1\text{ bit}\qquad(=H(A)-I(A;B)=1-2).
\end{align*}

\begin{remarque}[Une entropie conditionnelle négative]
Pour des variables classiques, $p(a,b)\leq p(b)$, donc $-\log_2p(a,b)\geq-\log_2p(b)$ et
$H(A,B)\geq H(B)$ : on a toujours $H(A|B)\geq0$, et $I(A;B)\leq H(A)$. Ici $H(A|B)=-1$ et
$I(A;B)=2>H(A)$ : ces valeurs n'ont pas d'équivalent classique, c'est une signature de
l'intrication. Le résultat est surprenant mais mathématiquement correct.
\end{remarque}

\begin{center}
\renewcommand{\arraystretch}{1.4}
\begin{tabular}{@{}lccc@{}}
\toprule
 & indépendantes & solidaires & état de Bell \\
\midrule
$H(A)$ & $1$ & $1$ & $1$ \\
$H(B)$ & $1$ & $1$ & $1$ \\
$H(A,B)$ & $2$ & $1$ & $0$ \\
$I(A;B)$ & $0$ & $1$ & $2$ \\
$H(A|B)$ & $1$ & $0$ & $-1$ \\
\bottomrule
\end{tabular}
\end{center}
(en bits ; les deux premières colonnes sont classiques, la troisième est quantique).

\begin{figure}[!ht]
\centering
\adjustbox{max width=\linewidth}{%
\begin{tikzpicture}[>=Stealth, font=\small]
  \def\bw{0.8}
  % axe des ordonnées (bits)
  \draw[->] (0,-1.5) -- (0,2.8) node[above] {bits};
  \foreach \v in {-1,0,1,2} {
    \draw (-0.07,\v) -- (0.07,\v);
    \node[left, font=\footnotesize] at (-0.1,\v) {$\v$};
  }
  \draw[insagris!60] (0,0) -- (11.4,0);
  % limite classique de I(A;B) : H(A) = 1
  \draw[dashed, insarouge] (0,1) -- (11.4,1);
  \node[anchor=west, font=\scriptsize, text=insarouge, align=left] at (11.4,1)
    {limite classique\\ $I(A;B)\leq H(A)=1$};
  \node[anchor=west, font=\scriptsize, text=insagris, align=left] at (11.4,0)
    {limite classique\\ $H(A|B)\geq0$};
  % barres : (groupe, H(A,B), I(A;B), H(A|B))
  \foreach \g/\hab/\mi/\hc in {0/2/0/1, 1/1/1/0, 2/0/2/-1} {
    \pgfmathsetmacro{\x}{0.9+3.7*\g}
    \fill[insagris!65] (\x,0) rectangle (\x+\bw,\hab);
    \fill[insarouge] (\x+\bw+0.1,0) rectangle (\x+2*\bw+0.1,\mi);
    \fill[black!80] (\x+2*\bw+0.2,0) rectangle (\x+3*\bw+0.2,\hc);
    \node[above, font=\footnotesize] at (\x+0.5*\bw,{max(\hab,0)}) {$\hab$};
    \node[above, font=\footnotesize] at (\x+1.5*\bw+0.1,{max(\mi,0)}) {$\mi$};
    \ifnum\hc<0
      \node[below, font=\footnotesize] at (\x+2.5*\bw+0.2,\hc) {$\hc$};
    \else
      \node[above, font=\footnotesize] at (\x+2.5*\bw+0.2,\hc) {$\hc$};
    \fi
  }
  \node[align=center, anchor=north] at (2.15,-1.65) {variables\\ indépendantes};
  \node[align=center, anchor=north] at (5.85,-1.65) {variables\\ solidaires};
  \node[align=center, anchor=north] at (9.55,-1.65) {état de Bell\\ $\ket{\Phi^+}$};
  % légende
  \fill[insagris!65] (0.9,3.35) rectangle (1.2,3.65);
  \node[anchor=west, font=\footnotesize] at (1.25,3.5) {$H(A,B)$};
  \fill[insarouge] (3.6,3.35) rectangle (3.9,3.65);
  \node[anchor=west, font=\footnotesize] at (3.95,3.5) {$I(A;B)$};
  \fill[black!80] (6.3,3.35) rectangle (6.6,3.65);
  \node[anchor=west, font=\footnotesize] at (6.65,3.5) {$H(A|B)$};
\end{tikzpicture}}
\caption{Entropies (en bits) pour deux pièces indépendantes, deux pièces solidaires et un état de Bell.
Dans les trois cas $H(A)=H(B)=1$. Les deux premiers groupes respectent les bornes classiques
$I(A;B)\leq H(A)$ et $H(A|B)\geq0$ ; l'état de Bell les dépasse ($I=2$, $H(A|B)=-1$).}
\label{fig:entropies}
\end{figure}

% ---------------------------------------------------------------------
\subsection{Téléportation quantique}\label{sec:teleportation}
% ---------------------------------------------------------------------

\emph{Problème.} Alice détient un qubit $Q$ dans un état qu'elle ne connaît pas,
$\ket{\varphi}=\alpha\ket{0}+\beta\ket{1}$ avec $|\alpha|^2+|\beta|^2=1$, et elle veut que Bob
obtienne cet état sur l'un de ses qubits. Elle ne peut pas mesurer $Q$ pour connaître
$\alpha$ et $\beta$ (une mesure ne donne qu'un bit et détruit l'état), elle ne peut pas le
copier (théorème de non-clonage, remarque en fin de section) et elle n'envoie pas le qubit
lui-même. Le protocole est dû à \textcite{bennett1993}.

\begin{definition}{Protocole de téléportation quantique}
\textbf{Ressources.}
\begin{itemize}[nosep]
  \item Un \emph{e-bit} : Alice et Bob partagent une paire de qubits $(A,B)$ dans l'état de Bell
        $\ket{\Phi^+}_{BA}=\frac{1}{\sqrt2}\big(\ket{00}+\ket{11}\big)$ ; Alice détient $A$,
        Bob détient $B$.
  \item Le qubit $Q$ à téléporter, chez Alice.
  \item Un canal classique (deux bits).
\end{itemize}
Le système est écrit $(B,A,Q)$ : $\ket{b,a,q}=\ket{b}_B\otimes\ket{a}_A\otimes\ket{q}_Q$.

\textbf{Étapes.}
\begin{enumerate}[nosep]
  \item Alice applique un CNOT de contrôle $Q$ et de cible $A$.
  \item Alice applique une porte $H$ sur $Q$.
  \item Alice mesure $A$ et $Q$ dans la base de calcul, obtient deux bits $(a,q)$ et les envoie
        à Bob.
  \item Bob applique $X$ si $a=1$, puis $Z$ si $q=1$, à son qubit $B$.
\end{enumerate}
\end{definition}

La figure~\ref{fig:tele-schema} montre l'organisation du protocole et la
figure~\ref{fig:teleportation} son circuit.

\begin{figure}[!ht]
\centering
\adjustbox{max width=\linewidth}{%
\begin{tikzpicture}[>=Stealth, font=\small,
    qb/.style={circle, draw=insarouge, fill=insarouge!8, thick, minimum size=8mm, inner sep=0pt},
    zone/.style={draw=insagris, rounded corners=4pt, thick, dashed, inner sep=8pt},
    num/.style={circle, fill=insarouge, text=white, font=\footnotesize\bfseries, inner sep=0pt,
                minimum size=4.6mm},
    etape/.style={anchor=west, font=\footnotesize, text=insarouge},
    onde/.style={thick, insarouge, decorate,
                 decoration={snake, amplitude=1.3pt, segment length=7pt}}]
  % Alice : Q (état inconnu) et A
  \node[qb] (Q) at (0,1.3) {$Q$};
  \node[qb] (A) at (0,-0.5) {$A$};
  \node[zone, fit=(Q)(A), label={[text=insagris]above:Alice}] (alice) {};
  \node[anchor=east] at (-0.95,1.3) {$\ket{\varphi}=\alpha\ket{0}+\beta\ket{1}$};
  \draw[insarouge, thick] (Q) -- (A);
  % Bob : B
  \node[qb] (B) at (9,-0.5) {$B$};
  \node[zone, fit=(B), label={[text=insagris]above:Bob}] (bob) {};
  % source de la paire de Bell
  \node[draw=insarouge, fill=insarouge!10, rounded corners=2pt, thick, minimum width=3.2cm,
        minimum height=0.9cm] (src) at (4.5,-2.3) {source $\ket{\Phi^+}_{BA}$};
  \draw[onde] (src.west) -| (A.south);
  \draw[onde] (src.east) -| (B.south);
  % étapes d'Alice : pastilles sur le bord droit de sa zone
  \node[num] (n1) at (alice.east |- Q) {1};
  \node[num] (n2) at (alice.east) {2};
  \node[num] (n3) at (alice.east |- A) {3};
  \node[etape] at (n1.east) {\ CNOT de $Q$ vers $A$};
  \node[etape] at (n2.east) {\ porte $H$ sur $Q$};
  % canal classique
  \draw[double, double distance=1.2pt, ->, insagris] (n3.east) -- (bob.west);
  \node[above, font=\footnotesize] at (4.6,-0.45) {deux bits classiques $(a,q)$};
  \node[below, font=\footnotesize, text=insarouge] at (4.6,-0.55) {mesure de $A$ et $Q$ (étape 3)};
  % correction de Bob
  \node[num] (n4) at (bob.east) {4};
  \node[etape, align=left] at (n4.east) {\ $X$ si $a=1$, puis $Z$ si $q=1$};
  \node[anchor=west, font=\footnotesize, align=left] at ($(n4.east)+(0.05,-0.8)$)
    {\ $B$ est alors dans l'état $\ket{\varphi}$};
\end{tikzpicture}}
\caption{Organisation de la téléportation. Alice possède le qubit $Q$ à téléporter et la moitié $A$
d'une paire de Bell, Bob possède l'autre moitié $B$. (1)~CNOT de $Q$ vers $A$, (2)~porte $H$ sur $Q$,
(3)~mesure de $A$ et $Q$, dont les résultats $(a,q)$ partent par un canal classique, (4)~corrections de
Bob. Aucun qubit ne voyage entre Alice et Bob.}
\label{fig:tele-schema}
\end{figure}

\begin{figure}[!ht]
\centering
\begin{adjustbox}{max width=\linewidth}
\begin{tikzpicture}[>=Stealth, font=\small]
  % fils quantiques : Q en haut, A au milieu, B en bas
  \draw (0,2.4) -- (5.3,2.4);
  \draw (0,1.2) -- (5.3,1.2);
  \draw (0,0) -- (10.6,0);
  \node[anchor=east] at (0,2.4) {$\ket{\varphi}_Q$};
  \node[anchor=south west, text=insagris, font=\footnotesize] at (0.05,1.2) {$A$};
  \node[anchor=south west, text=insagris, font=\footnotesize] at (0.05,0) {$B$};
  % paire de Bell partagée
  \draw[insagris, thick] (-0.1,1.5) -- (-0.25,1.5) -- (-0.25,-0.3) -- (-0.1,-0.3);
  \node[anchor=east] at (-0.35,0.6) {$\ket{\Phi^+}_{BA}$};
  % CNOT : contrôle Q, cible A
  \draw[insarouge, thick] (2.6,2.4) -- (2.6,1.2);
  \pic at (2.6,2.4) {ctl};
  \pic at (2.6,1.2) {cible};
  % Hadamard sur Q
  \node[porte] at (4.2,2.4) {$H$};
  % mesures d'Alice
  \pic at (5.6,2.4) {mesure};
  \pic at (5.6,1.2) {mesure};
  % fils classiques (doubles)
  \draw[double, double distance=1.2pt, insagris] (5.9,1.2) -- (7.6,1.2) -- (7.6,0.35);
  \draw[double, double distance=1.2pt, insagris] (5.9,2.4) -- (9.6,2.4) -- (9.6,0.35);
  \node[above, font=\footnotesize] at (6.5,1.2) {$a$};
  \node[above, font=\footnotesize] at (6.5,2.4) {$q$};
  % corrections de Bob
  \node[porte] at (7.6,0) {$X$};
  \node[porte] at (9.6,0) {$Z$};
  \node[anchor=west] at (10.6,0) {$\ket{\varphi}_B$};
  % états intermédiaires
  \foreach \x in {1.4,3.4,4.9}
    \draw[dashed, insagris!70] (\x,2.8) -- (\x,-0.5);
  \node[anchor=north, font=\footnotesize] at (1.4,-0.5) {$\ket{\Pi_0}$};
  \node[anchor=north, font=\footnotesize] at (3.4,-0.5) {$\ket{\Pi_1}$};
  \node[anchor=north, font=\footnotesize] at (4.9,-0.5) {$\ket{\Pi_2}$};
\end{tikzpicture}
\end{adjustbox}
\caption{Circuit de téléportation. Alice détient $Q$ (dans $\ket{\varphi}=\alpha\ket{0}+\beta\ket{1}$)
et $A$, Bob détient $B$ ; $A$ et $B$ forment la paire $\ket{\Phi^+}_{BA}$. Les traits doubles
sont des fils \emph{classiques} qui portent les résultats $a$ (mesure de $A$) et $q$ (mesure de
$Q$) ; $X$ est appliqué si $a=1$, $Z$ si $q=1$. Le ket s'écrit $\ket{b,a,q}$ : le fil du haut
($Q$) est le dernier chiffre.}
\label{fig:teleportation}
\end{figure}

\subsubsection{Calcul pas à pas}

\paragraph{État initial $\ket{\Pi_0}$.} Alice et Bob partagent la paire, et $Q$ est indépendant
d'elle :
\begin{align*}
  \ket{\Pi_0}&=\ket{\Phi^+}_{BA}\otimes\ket{\varphi}_Q
    =\frac{1}{\sqrt2}\big(\ket{00}+\ket{11}\big)\otimes\big(\alpha\ket{0}+\beta\ket{1}\big)\\
  &=\frac{1}{\sqrt2}\Big(\alpha\ket{00}\ket{0}+\beta\ket{00}\ket{1}+\alpha\ket{11}\ket{0}+\beta\ket{11}\ket{1}\Big)\\
  &=\frac{\alpha\ket{000}+\beta\ket{001}+\alpha\ket{110}+\beta\ket{111}}{\sqrt2}.
\end{align*}

\paragraph{Après le CNOT $Q\to A$ : $\ket{\Pi_1}$.} Le CNOT (contrôle $Q$, dernier chiffre ;
cible $A$, chiffre du milieu) agit ainsi : $\ket{b,a,q}\to\ket{b,\,a\oplus q,\,q}$. Terme par terme :
\begin{center}
\renewcommand{\arraystretch}{1.3}
\begin{tabular}{@{}cccc@{}}
\toprule
Terme & $q$ & Effet sur $A$ & Résultat \\
\midrule
$\alpha\ket{000}$ & $0$ & aucun & $\alpha\ket{000}$ \\
$\beta\ket{001}$ & $1$ & $0\to1$ & $\beta\ket{011}$ \\
$\alpha\ket{110}$ & $0$ & aucun & $\alpha\ket{110}$ \\
$\beta\ket{111}$ & $1$ & $1\to0$ & $\beta\ket{101}$ \\
\bottomrule
\end{tabular}
\end{center}
donc
\[
  \ket{\Pi_1}=\frac{\alpha\ket{000}+\beta\ket{011}+\alpha\ket{110}+\beta\ket{101}}{\sqrt2}.
\]
On isole ensuite le dernier chiffre, celui de $Q$, qui va changer à l'étape suivante
(le premier ket est celui du couple $\ket{b,a}$) :
\[
  \ket{\Pi_1}=\frac{\alpha\ket{00}\otimes\ket{0}+\beta\ket{01}\otimes\ket{1}
    +\alpha\ket{11}\otimes\ket{0}+\beta\ket{10}\otimes\ket{1}}{\sqrt2}.
\]

\paragraph{Après la porte $H$ sur $Q$ : $\ket{\Pi_2}$.} On remplace $\ket{0}$ par
$H\ket{0}=\ket{+}$ et $\ket{1}$ par $H\ket{1}=\ket{-}$ dans le facteur $Q$ :
\[
  \ket{\Pi_2}=\frac{1}{\sqrt2}\Big(\alpha\ket{00}\otimes\ket{+}+\beta\ket{01}\otimes\ket{-}
    +\alpha\ket{11}\otimes\ket{+}+\beta\ket{10}\otimes\ket{-}\Big).
\]
Avec $\ket{\pm}=\frac{1}{\sqrt2}(\ket{0}\pm\ket{1})$, le préfacteur devient
$\frac{1}{\sqrt2}\cdot\frac{1}{\sqrt2}=\frac12$ :
\begin{align*}
  \ket{\Pi_2}&=\frac12\Big(\alpha\ket{00}\big(\ket{0}+\ket{1}\big)+\beta\ket{01}\big(\ket{0}-\ket{1}\big)
     +\alpha\ket{11}\big(\ket{0}+\ket{1}\big)+\beta\ket{10}\big(\ket{0}-\ket{1}\big)\Big)\\
  &=\frac12\Big(\alpha\ket{000}+\alpha\ket{001}+\beta\ket{010}-\beta\ket{011}
     +\alpha\ket{110}+\alpha\ket{111}+\beta\ket{100}-\beta\ket{101}\Big).
\end{align*}

\paragraph{Regroupement selon le résultat d'Alice.} Alice mesurera $A$ et $Q$, donc les deux
derniers chiffres. On regroupe les huit termes qui ont les mêmes chiffres $(a,q)$, en mettant
le ket de Bob en premier :
\begin{center}
\renewcommand{\arraystretch}{1.3}
\begin{tabular}{@{}cll@{}}
\toprule
$(a,q)$ & Termes de $\ket{\Pi_2}$ (facteur $\frac12$ omis) & Ket de Bob \\
\midrule
$(0,0)$ & $\alpha\ket{000}+\beta\ket{100}$ & $(\alpha\ket{0}+\beta\ket{1})\ket{00}$ \\
$(0,1)$ & $\alpha\ket{001}-\beta\ket{101}$ & $(\alpha\ket{0}-\beta\ket{1})\ket{01}$ \\
$(1,0)$ & $\beta\ket{010}+\alpha\ket{110}$ & $(\beta\ket{0}+\alpha\ket{1})\ket{10}$ \\
$(1,1)$ & $-\beta\ket{011}+\alpha\ket{111}$ & $(-\beta\ket{0}+\alpha\ket{1})\ket{11}$ \\
\bottomrule
\end{tabular}
\end{center}
d'où
\[
  \ket{\Pi_2}=\frac12\Big[\big(\alpha\ket{0}+\beta\ket{1}\big)\ket{00}
    +\big(\alpha\ket{0}-\beta\ket{1}\big)\ket{01}
    +\big(\beta\ket{0}+\alpha\ket{1}\big)\ket{10}
    +\big(\alpha\ket{1}-\beta\ket{0}\big)\ket{11}\Big].
\]
Dans chaque terme, le ket de gauche est l'état de $B$ et le ket de droite $\ket{a,q}$ celui du
couple $(A,Q)$ d'Alice.

\subsubsection{Mesure d'Alice et corrections de Bob}

Alice mesure $A$ et $Q$. Comme à la section~\ref{sec:mesure2}, la probabilité d'un résultat
$(a,q)$ est le carré de la norme du terme correspondant, et l'état de Bob devient le ket
associé (déjà normé). Ici
\[
  P(a,q)=\Big\|\tfrac12\ket{\phi_{aq}}\Big\|^2=\tfrac14\big(|\alpha|^2+|\beta|^2\big)=\tfrac14
  \quad\text{pour les quatre résultats,}
\]
qui sont donc équiprobables et \emph{indépendants} de $\alpha$ et $\beta$ : les deux bits
envoyés ne révèlent rien sur $\ket{\varphi}$. Les kets de Bob sont
$\ket{\phi_{00}}=\alpha\ket{0}+\beta\ket{1}$, $\ket{\phi_{01}}=\alpha\ket{0}-\beta\ket{1}$,
$\ket{\phi_{10}}=\beta\ket{0}+\alpha\ket{1}$ et $\ket{\phi_{11}}=-\beta\ket{0}+\alpha\ket{1}$.

\begin{center}
\renewcommand{\arraystretch}{1.5}
\begin{tabular}{@{}cccll@{}}
\toprule
Résultat $(a,q)$ & Probabilité & État de Bob & Lien avec $\ket{\varphi}$ & Correction de Bob \\
\midrule
$(0,0)$ & $\frac14$ & $\alpha\ket{0}+\beta\ket{1}$ & $\ket{\varphi}$ & $I$ (rien) \\
$(0,1)$ & $\frac14$ & $\alpha\ket{0}-\beta\ket{1}$ & $Z\ket{\varphi}$ & $Z$ \\
$(1,0)$ & $\frac14$ & $\beta\ket{0}+\alpha\ket{1}$ & $X\ket{\varphi}$ & $X$ \\
$(1,1)$ & $\frac14$ & $-\beta\ket{0}+\alpha\ket{1}$ & $XZ\ket{\varphi}$ & $X$ puis $Z$ (opérateur $ZX$) \\
\bottomrule
\end{tabular}
\end{center}

\paragraph{Vérification des quatre corrections.} Avec $X\ket{0}=\ket{1}$, $X\ket{1}=\ket{0}$,
$Z\ket{0}=\ket{0}$ et $Z\ket{1}=-\ket{1}$ :
\begin{align*}
  (0,0):\quad & I\big(\alpha\ket{0}+\beta\ket{1}\big)=\alpha\ket{0}+\beta\ket{1},\\
  (0,1):\quad & Z\big(\alpha\ket{0}-\beta\ket{1}\big)=\alpha\ket{0}-\beta Z\ket{1}=\alpha\ket{0}+\beta\ket{1},\\
  (1,0):\quad & X\big(\beta\ket{0}+\alpha\ket{1}\big)=\beta\ket{1}+\alpha\ket{0}=\alpha\ket{0}+\beta\ket{1},\\
  (1,1):\quad & X\big(-\beta\ket{0}+\alpha\ket{1}\big)=-\beta\ket{1}+\alpha\ket{0}=\alpha\ket{0}-\beta\ket{1},
     \ \text{puis}\ Z\big(\alpha\ket{0}-\beta\ket{1}\big)=\alpha\ket{0}+\beta\ket{1}.
\end{align*}
Dans le dernier cas, avec les matrices,
$ZX=\begin{pmatrix}1&0\\0&-1\end{pmatrix}\begin{pmatrix}0&1\\1&0\end{pmatrix}
=\begin{pmatrix}0&1\\-1&0\end{pmatrix}$ et
$ZX\begin{pmatrix}-\beta\\ \alpha\end{pmatrix}=\begin{pmatrix}\alpha\\ \beta\end{pmatrix}$.
L'ordre compte : l'état de Bob est $XZ\ket{\varphi}$ ($Z\ket{\varphi}=\alpha\ket{0}-\beta\ket{1}$,
puis $X$ donne $\alpha\ket{1}-\beta\ket{0}$), et on l'annule avec
$(XZ)^{-1}=Z^{-1}X^{-1}=ZX$ : on applique d'abord $X$, puis $Z$, comme dans le circuit.
Dans les quatre cas, Bob retrouve exactement $\ket{\varphi}=\alpha\ket{0}+\beta\ket{1}$.

\subsubsection{Ce que Bob sait avant de recevoir les bits}

Tant que Bob ignore $(a,q)$, son qubit est dans le mélange des quatre kets $\ket{\phi_{aq}}$, chacun
avec la probabilité $\frac14$ :
\begin{align*}
  \rho_B&=\frac14\sum_{a,q}\ket{\phi_{aq}}\bra{\phi_{aq}}\\
  &=\frac14\Bigg[\begin{pmatrix}|\alpha|^2&\alpha\beta^*\\ \alpha^*\beta&|\beta|^2\end{pmatrix}
    +\begin{pmatrix}|\alpha|^2&-\alpha\beta^*\\ -\alpha^*\beta&|\beta|^2\end{pmatrix}
    +\begin{pmatrix}|\beta|^2&\alpha^*\beta\\ \alpha\beta^*&|\alpha|^2\end{pmatrix}
    +\begin{pmatrix}|\beta|^2&-\alpha^*\beta\\ -\alpha\beta^*&|\alpha|^2\end{pmatrix}\Bigg]\\
  &=\frac14\begin{pmatrix}2(|\alpha|^2+|\beta|^2)&0\\0&2(|\alpha|^2+|\beta|^2)\end{pmatrix}
   =\frac12\begin{pmatrix}1&0\\0&1\end{pmatrix}=\frac12\,\mathrm{Id}.
\end{align*}
Sans les deux bits, Bob a donc un qubit complètement mélangé, qui ne dépend pas de $\alpha$ ni de
$\beta$ : il n'a reçu aucune information sur $\ket{\varphi}$. C'est l'envoi classique de $(a,q)$
qui permet de la retrouver, donc rien ne voyage plus vite que la lumière.

\begin{remarque}[Bilan de la téléportation]
\begin{itemize}[nosep]
  \item Ressources : $1$ e-bit (la paire de Bell) et $2$ bits classiques pour transmettre
        \emph{un qubit} d'état quelconque, sans qu'Alice connaisse $\alpha$ ni $\beta$.
  \item Pas de clonage : après la mesure, $A$ et $Q$ sont dans l'état de base $\ket{a,q}$.
        Alice n'a plus $\ket{\varphi}$ : l'état a été \emph{déplacé}, pas copié. L'intrication
        entre $A$ et $B$ a été consommée.
  \item Les quatre corrections $I,Z,X,ZX$ sont des portes contrôlées par des bits classiques
        (section~\ref{sec:controlees}).
\end{itemize}
\end{remarque}

\begin{remarque}[Les gestes d'Alice forment une mesure de Bell]
Le CNOT de contrôle $Q$ sur $A$ suivi de $H$ sur $Q$ est l'\emph{inverse} du générateur de la
section~\ref{sec:circuit-bell} (les portes $H$ et CNOT sont leurs propres inverses, et l'ordre
s'inverse) : il envoie la base de Bell du couple $(A,Q)$ sur la base de calcul. Mesurer ensuite dans
la base de calcul revient donc à mesurer $(A,Q)$ dans la base de Bell. Les deux bits $(a,q)$
disent \emph{quel} état de Bell Alice a trouvé, ce qui explique les quatre corrections de Bob.
\end{remarque}

\begin{remarque}[Non-clonage]
Alice ne peut pas simplement copier $\ket{\varphi}$ \parencite{wootters1982}. Supposons une porte
unitaire $U$ telle que $U\big(\ket{\psi}\otimes\ket{0}\big)=\ket{\psi}\otimes\ket{\psi}$ pour tout
$\ket{\psi}$. Elle copie $\ket{0}$ et $\ket{1}$, mais par linéarité
\[
  U\big(\ket{+}\otimes\ket{0}\big)=\tfrac{1}{\sqrt2}\big(\ket{00}+\ket{11}\big)
  \neq\ket{+}\otimes\ket{+}=\tfrac12\big(\ket{00}+\ket{01}+\ket{10}+\ket{11}\big),
\]
ce qui contredit l'hypothèse. Le CNOT copie un bit classique, pas une superposition (remarque de la
section~\ref{sec:cnot}). La téléportation respecte cette interdiction : l'état est déplacé, il n'est
jamais présent en deux endroits.
\end{remarque}

% ---------------------------------------------------------------------
\subsection{Corrélations non locales : le jeu CHSH}\label{sec:chsh}
% ---------------------------------------------------------------------

Les états de Bell donnent des résultats corrélés en base $Z$ comme en base $X$
(section~\ref{sec:bell}). Mais deux enveloppes contenant la même carte, préparées à l'avance, donnent
aussi des résultats égaux : des corrélations dans une seule base ne prouvent pas l'intrication. Bell
\parencite{bell1964} a montré qu'on peut trancher en \emph{changeant de base de mesure} ; la version de
Clauser, Horne, Shimony et Holt \parencite{clauser1969} se présente comme un jeu.

\begin{definition}{Le jeu CHSH}
Alice et Bob, éloignés, ne communiquent plus une fois le jeu commencé. Un arbitre tire deux bits $x$ et
$y$ de façon uniforme et indépendante, donne $x$ à Alice et $y$ à Bob. Chacun répond par un bit : $a$ pour
Alice, $b$ pour Bob. Ils \emph{gagnent} si
\[
  a\oplus b=x\cdot y,
\]
c'est-à-dire si $a=b$ pour $(x,y)\in\{(0,0),(0,1),(1,0)\}$ et si $a\neq b$ pour $(x,y)=(1,1)$.
\end{definition}

\paragraph{Stratégie classique : au plus $\frac34$.} Une stratégie déterministe est $a=f(x)$,
$b=g(y)$. Si elle gagnait les quatre parties, on aurait
\[
  f(0)\oplus g(0)=0,\qquad f(0)\oplus g(1)=0,\qquad f(1)\oplus g(0)=0,\qquad f(1)\oplus g(1)=1 .
\]
En additionnant ces quatre égalités modulo $2$, chacune des valeurs $f(0),f(1),g(0),g(1)$ apparaît
deux fois : le membre de gauche vaut $0$, celui de droite vaut $1$. C'est absurde, donc au moins une
partie sur quatre est perdue et $P(\text{gain})\leq\frac34$. La stratégie «~toujours répondre $0$~»
atteint $\frac34$ (elle ne perd que pour $(1,1)$). Un hasard partagé ne change rien : il mélange des
stratégies déterministes, dont aucune ne dépasse $\frac34$.

\paragraph{Stratégie quantique.} Alice et Bob partagent avant le jeu une paire $\ket{\Phi^+}$ (un
e-bit), puis chacun mesure son qubit sans rien dire à l'autre. Pour un angle $\theta$, notons
\[
  \ket{\theta_+}=\cos\tfrac\theta2\,\ket{0}+\sin\tfrac\theta2\,\ket{1},
  \qquad
  \ket{\theta_-}=-\sin\tfrac\theta2\,\ket{0}+\cos\tfrac\theta2\,\ket{1}
\]
la base de mesure d'angle $\theta$ : c'est la direction $\theta$ dans le plan $xz$ de la sphère de
Bloch ($\theta=0$ : base $Z$ ; $\theta=\frac\pi2$ : base $X$). Le résultat vaut $0$ pour
$\ket{\theta_+}$ et $1$ pour $\ket{\theta_-}$. La stratégie est
\begin{itemize}[nosep]
  \item Alice : $\theta=0$ si $x=0$ (base $Z$), $\theta=\frac\pi2$ si $x=1$ (base $X$) ;
  \item Bob : $\theta=\frac\pi4$ si $y=0$, $\theta=-\frac\pi4$ si $y=1$ (bases propres de
        $\frac{Z\pm X}{\sqrt2}$).
\end{itemize}

\begin{figure}[!ht]
\centering
\adjustbox{max width=\linewidth}{%
\begin{tikzpicture}[>=Stealth, font=\footnotesize,
    joueur/.style={draw=insarouge, fill=insarouge!6, rounded corners=2pt, thick,
                   minimum width=2.1cm, minimum height=1.1cm, align=center},
    arb/.style={draw=insagris, fill=insagris!8, rounded corners=2pt, thick,
                minimum width=3.6cm, minimum height=0.8cm, align=center}]
  % --- le jeu ---
  \node[arb] (arb) at (3.2,2.7) {arbitre : $x,y$ au hasard};
  \node[joueur] (al) at (0,0) {Alice};
  \node[joueur] (bo) at (6.4,0) {Bob};
  \draw[->, thick, insagris] (arb.west) -| node[pos=0.2, above] {$x$} (al.north);
  \draw[->, thick, insagris] (arb.east) -| node[pos=0.2, above] {$y$} (bo.north);
  \draw[insarouge, thick, decorate, decoration={snake, amplitude=1.1pt, segment length=6pt}]
    (al.east) -- (bo.west);
  \node[insarouge, above=1pt] at (3.2,0.05) {$\ket{\Phi^+}$};
  \node[text=insagris, below=2pt, font=\scriptsize] at (3.2,-0.05) {paire intriquée};
  \draw[->, thick] (al.south) -- ++(0,-1.0) node[below] {$a$};
  \draw[->, thick] (bo.south) -- ++(0,-1.0) node[below] {$b$};
  \node[draw=insagris, rounded corners=2pt, inner sep=3pt] at (3.2,-1.75) {gain si $a\oplus b=x\cdot y$};
  % --- directions de mesure dans le plan xz de la sphère de Bloch ---
  \begin{scope}[shift={(11.9,0.3)}]
    \draw[insagris!70] (0,0) circle (1.7);
    \draw[insagris!50, dashed] (-1.7,0) -- (1.7,0);
    \draw[insagris!50, dashed] (0,-1.7) -- (0,1.7);
    \node[font=\scriptsize, text=insagris, anchor=north] at (0,-1.72) {$\ket{1}$};
    \node[font=\scriptsize, text=insagris, anchor=east] at (-1.72,-0.02) {$\ket{-}$};
    \draw[->, very thick, insarouge] (0,0) -- (90:1.7);
    \draw[->, very thick, insarouge] (0,0) -- (0:1.7);
    \draw[->, very thick, insagris] (0,0) -- (45:1.7);
    \draw[->, very thick, insagris] (0,0) -- (135:1.7);
    \node[insarouge, anchor=south, inner sep=2pt] at (0,1.72) {$x=0$ : $Z$ $(\ket{0})$};
    \node[insarouge, anchor=west, inner sep=2pt] at (1.72,0) {$x=1$ : $X$ $(\ket{+})$};
    \node[text=insagris, anchor=south west, align=left, inner sep=2pt] at (1.25,1.25) {$y=0$\\ $\frac{Z+X}{\sqrt2}$};
    \node[text=insagris, anchor=south east, align=right, inner sep=2pt] at (-1.25,1.25) {$y=1$\\ $\frac{Z-X}{\sqrt2}$};
    \draw[insagris] (0.85,0) arc (0:45:0.85);
    \node[font=\scriptsize, text=insagris] at (22.5:1.12) {$\frac\pi4$};
    \draw[insagris] (0,0.85) arc (90:135:0.85);
    \node[font=\scriptsize, text=insagris] at (112.5:1.12) {$\frac\pi4$};
    \draw[insagris] (0.6,0.6) arc (45:90:0.85);
    \node[font=\scriptsize, text=insagris] at (67.5:1.12) {$\frac\pi4$};
  \end{scope}
\end{tikzpicture}}
\caption{Le jeu CHSH (à gauche) et les quatre directions de mesure de la stratégie quantique (à
droite), dans le plan $xz$ de la sphère de Bloch : Alice en rouge, Bob en gris. Pour
$(x,y)=(0,0),(0,1),(1,0)$, les directions font un angle $\frac\pi4$ ; pour $(1,1)$, un angle
$\frac{3\pi}4$.}
\label{fig:chsh}
\end{figure}

\begin{resultat}{Corrélations de $\ket{\Phi^+}$ en bases tournées}
Si Alice mesure dans la base d'angle $\alpha$ et Bob dans la base d'angle $\beta$, alors
\[
  P(a=b)=\cos^2\frac{\alpha-\beta}{2}.
\]
\end{resultat}

\noindent\emph{Preuve.} Avec $c=\cos\frac\alpha2$ et $s=\sin\frac\alpha2$,
$\ket{\alpha_+\alpha_+}+\ket{\alpha_-\alpha_-}=(c^2+s^2)\ket{00}+(cs-sc)\big(\ket{01}+\ket{10}\big)+(s^2+c^2)\ket{11}
=\ket{00}+\ket{11}$ : on peut donc écrire $\ket{\Phi^+}=\frac{1}{\sqrt2}\big(\ket{\alpha_+\alpha_+}+\ket{\alpha_-\alpha_-}\big)$.
Si Alice obtient $a=0$ (probabilité $\frac12$), le qubit de Bob est dans $\ket{\alpha_+}$ et Bob obtient $b=0$
avec la probabilité $|\braket{\beta_+}{\alpha_+}|^2$, où $\braket{\beta_+}{\alpha_+}=\cos\frac\beta2\cos\frac\alpha2+\sin\frac\beta2\sin\frac\alpha2=\cos\frac{\alpha-\beta}{2}$.
Si Alice obtient $a=1$, Bob est dans $\ket{\alpha_-}$ et obtient $b=1$ avec la même probabilité, car
$\braket{\beta_-}{\alpha_-}=\braket{\beta_+}{\alpha_+}$. Donc
$P(a=b)=\frac12\cos^2\frac{\alpha-\beta}2+\frac12\cos^2\frac{\alpha-\beta}2$. \hfill$\square$

\medskip
\noindent Pour les quatre couples $(x,y)$ :
\begin{center}
\renewcommand{\arraystretch}{1.3}
\begin{tabular}{@{}c c c c c c@{}}
\toprule
$(x,y)$ & $\alpha$ & $\beta$ & $P(a=b)$ & condition de gain & $P(\text{gain})$\\
\midrule
$(0,0)$ & $0$ & $\frac\pi4$ & $\cos^2\frac\pi8$ & $a=b$ & $\cos^2\frac\pi8$\\
$(0,1)$ & $0$ & $-\frac\pi4$ & $\cos^2\frac\pi8$ & $a=b$ & $\cos^2\frac\pi8$\\
$(1,0)$ & $\frac\pi2$ & $\frac\pi4$ & $\cos^2\frac\pi8$ & $a=b$ & $\cos^2\frac\pi8$\\
$(1,1)$ & $\frac\pi2$ & $-\frac\pi4$ & $\cos^2\frac{3\pi}8=\sin^2\frac\pi8$ & $a\neq b$ &
  $1-\sin^2\frac\pi8=\cos^2\frac\pi8$\\
\bottomrule
\end{tabular}
\end{center}
Alice et Bob gagnent donc chaque partie avec la même probabilité :
\[
  P(\text{gain})=\cos^2\frac\pi8=\frac{2+\sqrt2}{4}\approx0{,}854>\frac34 .
\]

\begin{remarque}[Ce que montre le jeu CHSH]
\begin{itemize}[nosep]
  \item Aucune stratégie classique, même avec du hasard partagé, ne dépasse $\frac34$ : l'intrication
        produit des corrélations qu'aucun modèle «~local~» (résultats fixés à l'avance) ne reproduit.
        Les expériences sur des photons intriqués confirment la valeur quantique.
  \item $\cos^2\frac\pi8$ est le maximum quantique (\emph{borne de Tsirelson}, résultat admis,
        \cite{tsirelson1980}) : on ne peut pas gagner à tous les coups. Avec
        $E_{xy}=P(a=b)-P(a\neq b)=\cos(\alpha-\beta)$ et $S=E_{00}+E_{01}+E_{10}-E_{11}$, on a
        $P(\text{gain})=\frac12+\frac S8$ ; l'inégalité de Bell donne $|S|\leq2$ en classique, et la
        stratégie ci-dessus atteint $S=2\sqrt2$.
  \item Il n'y a pas de communication : la réponse de Bob est uniforme ($P(b=0)=\frac12$) quelle que
        soit la base choisie par Alice. L'intrication ne transmet pas d'information ; elle ne fait
        que corréler les résultats.
\end{itemize}
\end{remarque}

\clearpage
% ---------------------------------------------------------------------
\subsection{Récapitulatif du chapitre}\label{sec:recap5}
% ---------------------------------------------------------------------

\begin{center}
\adjustbox{max width=\linewidth}{%
\begin{tikzpicture}[>=Stealth, font=\footnotesize,
    boite/.style={draw=insarouge, fill=insarouge!6, rounded corners=2pt, thick,
                  minimum width=4.4cm, minimum height=1.45cm, align=center, inner sep=3pt},
    lien/.style={->, very thick, insagris},
    etiq/.style={font=\scriptsize\itshape, text=insagris, inner sep=2pt}]
  \node[boite] (mes) at (0,0)
    {\textbf{Mesure à deux qubits}\\ $P(ij)=|c_{ij}|^2$\\ {\scriptsize\S\,\ref{sec:mesure2}}};
  \node[boite] (sep) at (6.3,0)
    {\textbf{Séparable ou intriqué}\\ $D=c_{00}c_{11}-c_{01}c_{10}$\\ {\scriptsize\S\,\ref{sec:separable}}};
  \node[boite] (bel) at (12.6,0)
    {\textbf{États de Bell}\\ $\ket{\Phi^\pm}$, $\ket{\Psi^\pm}$ ; $\mathrm{CNOT}\,(I\otimes H)$\\
     {\scriptsize\S\,\ref{sec:bell}, \ref{sec:circuit-bell}}};
  \node[boite] (tra) at (12.6,-2.3)
    {\textbf{Trace partielle, entropies}\\ $\rho_A=\operatorname{Tr}_B\rho_{AB}$ ; $S(\rho)$, $I(A;B)$\\
     {\scriptsize\S\,\ref{sec:trace}, \ref{sec:info}}};
  \node[boite] (tel) at (6.3,-2.3)
    {\textbf{Téléportation}\\ e-bit + 2 bits classiques\\ {\scriptsize\S\,\ref{sec:teleportation}}};
  \node[boite] (chs) at (0,-2.3)
    {\textbf{Jeu CHSH}\\ $\frac34$ classique, $\cos^2\frac\pi8$ quantique\\ {\scriptsize\S\,\ref{sec:chsh}}};
  \draw[lien] (mes) -- (sep) node[etiq, midway, above=1pt] {corrélations};
  \draw[lien] (sep) -- (bel) node[etiq, midway, above=1pt] {cas extrême};
  \draw[lien] (bel) -- (tra) node[etiq, midway, right=1pt] {un qubit seul};
  \draw[lien] (tra) -- (tel) node[etiq, midway, above=1pt] {ressource};
  \draw[lien] (tel) -- (chs) node[etiq, midway, above=1pt] {non-localité};
\end{tikzpicture}}
\end{center}

\begin{aretenir}
\begin{itemize}
  \item Deux qubits : quatre amplitudes $c_{ij}$, $P(ij)=|c_{ij}|^2$. Mesurer $A$ supprime les termes
        incompatibles et renormalise : l'état de $B$ peut changer.
  \item Séparable $\Longleftrightarrow D=c_{00}c_{11}-c_{01}c_{10}=0$. Sinon l'état est intriqué et
        aucun qubit n'a d'état propre (pour $\ket{\Phi^+}$, $P(00)\neq P(A{=}0)P(B{=}0)$).
  \item États de Bell : base orthonormée d'états maximalement intriqués, états propres de
        $Z\otimes Z$ et $X\otimes X$ ; $\mathrm{CNOT}\,(I\otimes H)$ les produit à partir de la base de
        calcul. Des portes locales ne créent pas d'intrication.
  \item Qubit seul : $\rho_A=\operatorname{Tr}_B\rho_{AB}$, avec $\operatorname{Tr}\rho_A^2=1-2|D|^2$ pour un état pur.
        Pour $\ket{\Phi^+}$ : $\rho_A=\mathrm{Id}/2$ alors que le couple est pur ; $S(\rho_{AB})=0$,
        $S(\rho_A)=1$ bit, $I(A;B)=2$ bits, $H(A|B)=-1$ (impossible en classique).
  \item Téléportation : un e-bit et deux bits classiques ; $P(a,q)=\frac14$ ; Bob applique $X$ si $a=1$
        puis $Z$ si $q=1$. L'état est déplacé, pas copié ; sans les bits, $\rho_B=\mathrm{Id}/2$.
  \item Jeu CHSH : toute stratégie classique gagne avec une probabilité $\leq\frac34$ ; avec $\ket{\Phi^+}$
        et des bases tournées, $\cos^2\frac\pi8\approx0{,}854$ (borne de Tsirelson), sans aucune
        communication.
\end{itemize}
\end{aretenir}

\begin{center}
\renewcommand{\arraystretch}{1.2}
\begin{tabular}{@{}l l r@{}}
\toprule
\textbf{Notion} & \textbf{Formule} & \textbf{§} \\
\midrule
Mesure de deux qubits & $P(ij)=|c_{ij}|^2$, \quad $P(A{=}a)=P(a0)+P(a1)$ & \ref{sec:mesure2} \\
Séparabilité & $\ket{\Psi}=\ket{\Psi}_A\otimes\ket{\Psi}_B\ \Longleftrightarrow\ D=0$ & \ref{sec:separable} \\
États de Bell & $\ket{\Phi^\pm}=\dfrac{\ket{00}\pm\ket{11}}{\sqrt2}$, \quad
  $\ket{\Psi^\pm}=\dfrac{\ket{01}\pm\ket{10}}{\sqrt2}$ & \ref{sec:bell} \\
Trace partielle & $\operatorname{Tr}_B\big(\ket{a}\bra{a'}\otimes\ket{b}\bra{b'}\big)
  =\ket{a}\bra{a'}\,\braket{b'}{b}$ & \ref{sec:trace} \\
Qubit d'un état pur & $\rho_A=CC^\dagger$, \quad $\operatorname{Tr}\rho_A^2=1-2|D|^2$ & \ref{sec:trace} \\
Entropies & $S(\rho)=-\sum_i\lambda_i\log_2\lambda_i$, \quad $I(A;B)=H(A)+H(B)-H(A,B)$ & \ref{sec:info} \\
Téléportation & $P(a,q)=\frac14$, \quad correction $Z^{q}X^{a}$ (d'abord $X$) & \ref{sec:teleportation} \\
Singulet & $(U\otimes U)\ket{\Psi^-}=\det(U)\,\ket{\Psi^-}$ & \ref{sec:bell} \\
Jeu CHSH & $P_{\mathrm{cl}}\leq\frac34$, \quad $P_{\mathrm{quant}}=\cos^2\frac\pi8=\frac{2+\sqrt2}{4}$ & \ref{sec:chsh} \\
\bottomrule
\end{tabular}
\end{center}

\begin{savoirfaire}
\begin{itemize}
  \item Calculer $P(ij)$, les probabilités d'un seul qubit et l'état après la mesure de l'un d'eux ;
        tester la séparabilité avec $D$ et calculer la pureté de $\rho_A$ (exercice~\ref{ex:separable3}).
  \item Suivre un état dans le générateur de Bell, écrire $\rho_{AB}$ et reconnaître les corrélations
        d'un état de Bell dans deux bases (exercice~\ref{ex:bellX}).
  \item Calculer $H(A)$, $H(A,B)$, $I(A;B)$ et $H(A|B)$ pour une loi ou un état (exercice~\ref{ex:entropies}).
  \item Dérouler la téléportation pour un état donné et utiliser les états de Bell pour le codage
        superdense (exercices~\ref{ex:telenum} et~\ref{ex:superdense}).
  \item Prouver la borne $\frac34$ du jeu CHSH classique et calculer la probabilité de gain d'une
        stratégie quantique (section~\ref{sec:chsh}).
\end{itemize}
\end{savoirfaire}
.
%  Il ne contient ni préambule ni \begin{document} : les environnements
%  (definition, resultat, remarque, aretenir, savoirfaire, exercice, correction),
%  les macros (\ii, \ee, \dd, \pd, \pdd, \Hop, \ket, \bra, \braket), les symboles
%  de circuits (ctl, cible, croix, mesure, porte) et les couleurs (insarouge,
%  insagris) sont définis dans main.tex. Un chapitre = une \section.
% =====================================================================

% Macros de Dirac (déjà définies dans main.tex ; ce bloc évite l'erreur
% « Undefined control sequence \ket » si main.tex est une ancienne version).
\providecommand{\ket}[1]{|#1\rangle}
\providecommand{\bra}[1]{\langle #1|}
\providecommand{\braket}[2]{\langle #1|#2\rangle}

% =====================================================================
\section{États quantiques, intrication et états de Bell}\label{chap:bell}
% =====================================================================

Les états de deux qubits forment un espace de dimension~$4$, et la plupart d'entre eux ne se
décomposent pas en «~un état pour $A$, un état pour $B$~» : ils sont \emph{intriqués}. Ce chapitre
apprend à mesurer un couple de qubits, à reconnaître l'intrication (critère de séparabilité), à la
produire avec un circuit (états de Bell), à décrire un qubit pris seul (trace partielle), à
quantifier l'information (entropies) et à s'en servir pour téléporter un état.

% ---------------------------------------------------------------------
\subsection{États de deux qubits et mesure}\label{sec:mesure2}
% ---------------------------------------------------------------------

Un état quelconque de $\mathcal{H}^4$ se décompose sur la base de calcul :
\[
  \ket{\Psi}=c_{00}\ket{00}+c_{01}\ket{01}+c_{10}\ket{10}+c_{11}\ket{11}
  =\begin{pmatrix}c_{00}\\ c_{01}\\ c_{10}\\ c_{11}\end{pmatrix},
  \qquad \sum_{i,j}|c_{ij}|^2=1 .
\]

\begin{resultat}{Mesure des deux qubits dans la base de calcul}
\[
  P(ij)=|\braket{ij}{\Psi}|^2=|c_{ij}|^2 .
\]
Pour un seul qubit : $P(A{=}0)=P(00)+P(01)$, $P(A{=}1)=P(10)+P(11)$,
$P(B{=}0)=P(00)+P(10)$, $P(B{=}1)=P(01)+P(11)$.
\end{resultat}

En effet $\braket{ij}{\Psi}=\sum_{k,l}c_{kl}\braket{ij}{kl}=\sum_{k,l}c_{kl}\delta_{ik}\delta_{jl}=c_{ij}$.

\paragraph{Exemple 1.} $\ket{\Psi}=\frac{1}{\sqrt2}\ket{00}+\frac{1}{\sqrt2}\ket{01}
=\big(\tfrac{1}{\sqrt2},\tfrac{1}{\sqrt2},0,0\big)^{\mathsf T}$. Il est normalisé :
$\frac12+\frac12+0+0=1$. Les probabilités jointes sont
$P(00)=P(01)=\frac12$ et $P(10)=P(11)=0$. Pour chaque qubit :
\[
  P(A{=}0)=\tfrac12+\tfrac12=1,\quad P(A{=}1)=0+0=0,\quad
  P(B{=}0)=\tfrac12+0=\tfrac12,\quad P(B{=}1)=\tfrac12+0=\tfrac12 .
\]
$A$ donne toujours $0$ ; $B$ donne $0$ ou $1$ au hasard.

\paragraph{Exemple 2 (mesure de $A$ seul).} Soit
$\ket{\Psi}=\frac12\ket{00}+\frac12\ket{01}+\frac{1}{\sqrt2}\ket{11}$. Normalisation :
$\frac14+\frac14+0+\frac12=1$. Probabilités jointes : $P(00)=\frac14$, $P(01)=\frac14$,
$P(10)=0$, $P(11)=\frac12$. Donc
\[
  P(A{=}0)=\tfrac14+\tfrac14=\tfrac12,\quad P(A{=}1)=0+\tfrac12=\tfrac12,\quad
  P(B{=}0)=\tfrac14+0=\tfrac14,\quad P(B{=}1)=\tfrac14+\tfrac12=\tfrac34 .
\]

\begin{remarque}[Effondrement après la mesure de $A$]
Si $A$ donne le résultat $a$, on ne garde que les termes de $\ket{\Psi}$ dont le premier
chiffre vaut $a$, puis on renormalise :
\[
  \ket{\Psi}\ \to\ \frac{c_{a0}\ket{a0}+c_{a1}\ket{a1}}{\sqrt{|c_{a0}|^2+|c_{a1}|^2}},
  \qquad\text{probabilité }|c_{a0}|^2+|c_{a1}|^2 .
\]
\end{remarque}

Pour l'exemple 2 :
\begin{itemize}[nosep]
  \item $A=0$ (probabilité $\frac12$) : on garde $\frac12\ket{00}+\frac12\ket{01}$, de norme
        au carré $\frac12$ ; renormalisé, $\frac{1}{\sqrt2}\big(\ket{00}+\ket{01}\big)
        =\ket{0}\otimes\ket{+}$. Ensuite $B$ vaut $0$ ou $1$ avec la probabilité $\frac12$.
  \item $A=1$ (probabilité $\frac12$) : on garde $\frac{1}{\sqrt2}\ket{11}$, de norme au carré
        $\frac12$ ; renormalisé, $\ket{11}=\ket{1}\otimes\ket{1}$. Ensuite $B$ vaut $1$
        avec certitude.
\end{itemize}
Contrôle : $P(B{=}0)=\frac12\cdot\frac12+\frac12\cdot0=\frac14$, comme calculé plus haut.
Ici la mesure de $A$ \emph{modifie} l'état de $B$ ($\ket{+}$ ou $\ket{1}$) : c'est la
signature d'un état intriqué (section~\ref{sec:separable}). La figure~\ref{fig:mesure2} résume
les deux mesures successives.

\begin{figure}[!ht]
\centering
\adjustbox{max width=\linewidth}{%
\begin{tikzpicture}[>=Stealth, font=\small,
    noeud/.style={draw=insarouge, fill=insarouge!6, rounded corners=2pt, thick, inner sep=4pt, align=center},
    fin/.style={draw=insagris, fill=insagris!10, rounded corners=2pt, inner sep=4pt, align=center},
    lien/.style={->, thick, insagris},
    proba/.style={font=\footnotesize, fill=white, inner sep=1.5pt, text=black}]
  \node[noeud] (psi) at (0,0) {$\ket{\Psi}$};
  \node[noeud] (a0) at (4.4,1.6) {$A=0$\\ {\footnotesize état $\ket{0}\otimes\ket{+}$}};
  \node[noeud] (a1) at (4.4,-1.6) {$A=1$\\ {\footnotesize état $\ket{1}\otimes\ket{1}$}};
  \node[fin] (b00) at (9.2,2.5) {$B=0$ : $\ket{00}$};
  \node[fin] (b01) at (9.2,0.7) {$B=1$ : $\ket{01}$};
  \node[fin] (b11) at (9.2,-1.6) {$B=1$ : $\ket{11}$};
  \draw[lien] (psi) -- (a0) node[proba, midway] {$\frac12$};
  \draw[lien] (psi) -- (a1) node[proba, midway] {$\frac12$};
  \draw[lien] (a0) -- (b00) node[proba, midway] {$\frac12$};
  \draw[lien] (a0) -- (b01) node[proba, midway] {$\frac12$};
  \draw[lien] (a1) -- (b11) node[proba, midway] {$1$};
  \node[anchor=west] at (11.0,2.5) {$P(00)=\frac12\cdot\frac12=\frac14$};
  \node[anchor=west] at (11.0,0.7) {$P(01)=\frac12\cdot\frac12=\frac14$};
  \node[anchor=west] at (11.0,-1.6) {$P(11)=\frac12\cdot1=\frac12$};
  \node[text=insagris, font=\footnotesize\itshape] at (2.2,3.0) {mesure de $A$};
  \node[text=insagris, font=\footnotesize\itshape] at (6.8,3.6) {mesure de $B$};
  \node[text=insagris, font=\footnotesize\itshape, anchor=west] at (11.0,3.6) {probabilité jointe};
\end{tikzpicture}}
\caption{Mesure de $A$ puis de $B$ pour l'état de l'exemple~2,
$\ket{\Psi}=\frac12\ket{00}+\frac12\ket{01}+\frac{1}{\sqrt2}\ket{11}$. Chaque flèche porte une
probabilité, et le produit des probabilités le long d'un chemin est la probabilité jointe. Après la
mesure de $A$, l'état de $B$ dépend du résultat ($\ket{+}$ ou $\ket{1}$).}
\label{fig:mesure2}
\end{figure}

% ---------------------------------------------------------------------
\subsection{États séparables et états intriqués}\label{sec:separable}
% ---------------------------------------------------------------------

\begin{definition}{Séparable ou intriqué}
Un état $\ket{\Psi}_{AB}$ est \emph{séparable} s'il s'écrit comme un produit
$\ket{\Psi}_{AB}=\ket{\Psi}_A\otimes\ket{\Psi}_B$. Sinon il est \emph{intriqué} : on ne
peut pas attribuer d'état propre à chaque qubit.
\end{definition}

\paragraph{Exemple séparable (exemple 1 ci-dessus).} Par linéarité du produit tensoriel :
\begin{align*}
  \ket{\Psi}&=\tfrac{1}{\sqrt2}\ket{00}_{AB}+\tfrac{1}{\sqrt2}\ket{01}_{AB}
    =\tfrac{1}{\sqrt2}\Big(\ket{0}_A\otimes\ket{0}_B+\ket{0}_A\otimes\ket{1}_B\Big)\\
  &=\ket{0}_A\otimes\Big(\tfrac{1}{\sqrt2}\big(\ket{0}+\ket{1}\big)\Big)_B
    =\ket{0}_A\otimes\ket{+}_B ,
\end{align*}
donc $\ket{\Psi}_A=\ket{0}$ et $\ket{\Psi}_B=\ket{+}$. Contrôle par les composantes :
$\begin{pmatrix}1\\ 0\end{pmatrix}\otimes\tfrac{1}{\sqrt2}\begin{pmatrix}1\\ 1\end{pmatrix}
=\tfrac{1}{\sqrt2}\begin{pmatrix}1\cdot1\\ 1\cdot1\\ 0\cdot1\\ 0\cdot1\end{pmatrix}
=\tfrac{1}{\sqrt2}\begin{pmatrix}1\\ 1\\ 0\\ 0\end{pmatrix}$.

\subsubsection{Un critère de séparabilité}

Un produit quelconque a pour coefficients $c_{00}=ac$, $c_{01}=ad$, $c_{10}=bc$,
$c_{11}=bd$ (formule de la section~\ref{sec:tenseur}, avec $a,b$ pour $A$ et $c,d$ pour $B$), donc
\[
  c_{00}c_{11}=(ac)(bd)=abcd=(ad)(bc)=c_{01}c_{10}.
\]

\begin{resultat}{Test de séparabilité (deux qubits)}
$\ket{\Psi}=\sum c_{ij}\ket{ij}$ est séparable si et seulement si
\[
  D=c_{00}c_{11}-c_{01}c_{10}=0 .
\]
\end{resultat}

\emph{Démonstration.} ($\Rightarrow$) Calcul ci-dessus. ($\Leftarrow$) On regroupe les
termes selon l'état de $A$ :
\[
  \ket{\Psi}=\ket{0}\otimes\big(c_{00}\ket{0}+c_{01}\ket{1}\big)+\ket{1}\otimes\big(c_{10}\ket{0}+c_{11}\ket{1}\big).
\]
Si $D=0$, les lignes $(c_{00},c_{01})$ et $(c_{10},c_{11})$ sont proportionnelles. Si
$(c_{00},c_{01})\neq(0,0)$, il existe $\lambda$ avec $(c_{10},c_{11})=\lambda(c_{00},c_{01})$
et $\ket{\Psi}=\big(\ket{0}+\lambda\ket{1}\big)\otimes\big(c_{00}\ket{0}+c_{01}\ket{1}\big)$.
Si $(c_{00},c_{01})=(0,0)$, alors $\ket{\Psi}=\ket{1}\otimes\big(c_{10}\ket{0}+c_{11}\ket{1}\big)$.
Dans les deux cas $\ket{\Psi}$ est un produit ; on renormalise chaque facteur.
\hfill$\square$

\subsubsection{Application à plusieurs états}

\begin{center}
\renewcommand{\arraystretch}{1.5}
\begin{tabular}{@{}lccl@{}}
\toprule
\textbf{État} & $(c_{00},c_{01},c_{10},c_{11})$ & $D$ & \textbf{Nature} \\
\midrule
$\ket{00}$ & $(1,0,0,0)$ & $1\cdot0-0\cdot0=0$ & séparable \\
$\ket{+}\otimes\ket{+}$ & $\frac12(1,1,1,1)$ & $\frac14-\frac14=0$ & séparable \\
$\ket{+}\otimes\ket{-}$ & $\frac12(1,-1,1,-1)$ & $-\frac14-(-\frac14)=0$ & séparable \\
$\frac{1}{\sqrt2}(\ket{00}+\ket{01})$ & $\frac{1}{\sqrt2}(1,1,0,0)$ & $0-\frac12\cdot0=0$ & séparable \\
$\ket{\Phi^+}=\frac{1}{\sqrt2}(\ket{00}+\ket{11})$ & $\frac{1}{\sqrt2}(1,0,0,1)$ & $\frac12-0=\frac12$ & intriqué \\
$\ket{\Phi^-}=\frac{1}{\sqrt2}(\ket{00}-\ket{11})$ & $\frac{1}{\sqrt2}(1,0,0,-1)$ & $-\frac12-0=-\frac12$ & intriqué \\
$\ket{\Psi^+}=\frac{1}{\sqrt2}(\ket{01}+\ket{10})$ & $\frac{1}{\sqrt2}(0,1,1,0)$ & $0-\frac12=-\frac12$ & intriqué \\
$\ket{\Psi^-}=\frac{1}{\sqrt2}(\ket{01}-\ket{10})$ & $\frac{1}{\sqrt2}(0,1,-1,0)$ & $0-(-\frac12)=\frac12$ & intriqué \\
$\frac12\ket{00}+\frac12\ket{01}+\frac{1}{\sqrt2}\ket{11}$ & $(\frac12,\frac12,0,\frac{1}{\sqrt2})$ & $\frac{1}{2\sqrt2}-0$ & intriqué \\
\bottomrule
\end{tabular}
\end{center}

Les quatre états de Bell $\ket{\Phi^\pm}$ et $\ket{\Psi^\pm}$ de la table sont tous intriqués
($D=\pm\frac12$) ; leurs propriétés font l'objet de la section~\ref{sec:bell}.

\paragraph{Preuve directe que $\ket{\Phi^+}$ est intriqué.} Supposons
$\ket{\Phi^+}=(a\ket{0}+b\ket{1})\otimes(c\ket{0}+d\ket{1})$. En identifiant les
composantes : $ac=\frac{1}{\sqrt2}$, $ad=0$, $bc=0$, $bd=\frac{1}{\sqrt2}$. La condition
$ad=0$ impose $a=0$ (impossible car $ac\neq0$) ou $d=0$ (impossible car $bd\neq0$).
Contradiction : aucun produit ne convient.

\paragraph{Corrélations.} Pour un produit $\ket{\Psi}_A\otimes\ket{\Psi}_B$ avec
$\ket{\Psi}_A=a_0\ket{0}+a_1\ket{1}$ et $\ket{\Psi}_B=b_0\ket{0}+b_1\ket{1}$, on a
$P(ij)=|a_ib_j|^2=|a_i|^2|b_j|^2=P(A{=}i)\,P(B{=}j)$ : les deux qubits sont
\emph{indépendants}. Vérification avec $\ket{+}\otimes\ket{+}$ : $P(00)=\frac14=\frac12\cdot\frac12$.
Pour $\ket{\Phi^+}$ : $P(00)=P(11)=\frac12$ et $P(01)=P(10)=0$, donc
$P(A{=}0)=P(B{=}0)=\frac12$, mais
\[
  P(00)=\tfrac12\neq\tfrac12\cdot\tfrac12=P(A{=}0)\,P(B{=}0) .
\]
Les deux qubits donnent toujours le même résultat, aléatoire : ils sont corrélés.
La figure~\ref{fig:epr} illustre la différence entre les deux cas.

\begin{figure}[!ht]
\centering
\adjustbox{max width=\linewidth}{%
\begin{tikzpicture}[>=Stealth, font=\small,
    labo/.style={draw=insagris, fill=insagris!8, rounded corners=2pt, thick, minimum width=2.6cm,
                 minimum height=1.0cm, align=center},
    source/.style={draw=insarouge, fill=insarouge!10, rounded corners=2pt, thick, minimum width=3.6cm,
                   minimum height=1.0cm, align=center},
    droit/.style={->, thick, insagris},
    onde/.style={->, thick, insarouge, decorate,
                 decoration={snake, amplitude=1.3pt, segment length=7pt, post length=2mm}},
    res/.style={font=\small, align=center}]
  % (a) état séparable
  \node[labo] (Aa) at (0,0) {Alice\\ mesure $A$};
  \node[source] (Sa) at (6,0) {source\\ $\ket{0}_A\otimes\ket{+}_B$};
  \node[labo] (Ba) at (12,0) {Bob\\ mesure $B$};
  \draw[droit] (Sa.west) -- (Aa.east);
  \draw[droit] (Sa.east) -- (Ba.west);
  \node[res, below=3pt of Aa] {$0\ \ 0\ \ 0\ \ 0\ \ 0\ \ 0$};
  \node[res, below=3pt of Ba] {$1\ \ 0\ \ 0\ \ 1\ \ 0\ \ 1$};
  \node[res, text=insagris, font=\footnotesize] at (6,-1.0) {$A$ donne toujours $0$ ; $B$ est aléatoire.\\
        Les deux résultats sont indépendants.};
  \node[anchor=south west, font=\small\bfseries] at (-1.3,0.75) {(a) état séparable};
  % (b) état intriqué
  \begin{scope}[yshift=-3.6cm]
    \node[labo] (Ab) at (0,0) {Alice\\ mesure $A$};
    \node[source] (Sb) at (6,0) {source\\ $\ket{\Phi^+}_{AB}$};
    \node[labo] (Bb) at (12,0) {Bob\\ mesure $B$};
    \draw[onde] (Sb.west) -- (Ab.east);
    \draw[onde] (Sb.east) -- (Bb.west);
    \node[res, below=3pt of Ab] {$1\ \ 0\ \ 0\ \ 1\ \ 1\ \ 0$};
    \node[res, below=3pt of Bb] {$1\ \ 0\ \ 0\ \ 1\ \ 1\ \ 0$};
    \node[res, text=insagris, font=\footnotesize] at (6,-1.0) {$A$ et $B$ sont chacun aléatoires,\\
          mais toujours égaux : ils sont corrélés.};
    \node[anchor=south west, font=\small\bfseries] at (-1.3,0.75) {(b) état intriqué};
  \end{scope}
\end{tikzpicture}}
\caption{Six mesures successives sur des paires préparées dans le même état (les résultats sont
des exemples). (a) Pour un produit, chaque qubit a son propre comportement. (b) Pour
$\ket{\Phi^+}=\frac{1}{\sqrt2}(\ket{00}+\ket{11})$, chaque résultat pris seul est un tirage à pile ou
face, mais les deux résultats coïncident toujours ; la ligne ondulée symbolise l'intrication.}
\label{fig:epr}
\end{figure}

\paragraph{Et un qubit pris seul ?} Dans un état séparable, chaque qubit garde son propre état pur ;
dans un état intriqué, aucun ket ne décrit un qubit isolé. On le décrit alors par une matrice de
densité, obtenue par la trace partielle (section~\ref{sec:trace}).

% ---------------------------------------------------------------------
\subsection{Les états de Bell}\label{sec:bell}
% ---------------------------------------------------------------------

\begin{definition}{Les quatre états de Bell}
\[
  \ket{\Phi^{\pm}}=\frac{1}{\sqrt2}\big(\ket{00}\pm\ket{11}\big),
  \qquad
  \ket{\Psi^{\pm}}=\frac{1}{\sqrt2}\big(\ket{01}\pm\ket{10}\big).
\]
Ils sont intriqués ($|D|=\frac12$, section~\ref{sec:separable}) et le sont au maximum :
$|D|\leq\frac12$ pour tout état normalisé de deux qubits, et chaque qubit d'un état de Bell,
pris seul, est complètement aléatoire (section~\ref{sec:trace}).
\end{definition}

\paragraph{Une base de $\mathcal{H}^4$.} Chaque état est de norme $\frac12+\frac12=1$. Les
$\ket{\Phi^\pm}$ n'ont de composantes que sur $\ket{00}$ et $\ket{11}$, les $\ket{\Psi^\pm}$ que sur
$\ket{01}$ et $\ket{10}$, donc $\braket{\Phi^\pm}{\Psi^\pm}=0$ pour toutes les combinaisons de signes ;
enfin $\braket{\Phi^+}{\Phi^-}=\frac12(1-1)=0$ et $\braket{\Psi^+}{\Psi^-}=\frac12(1-1)=0$. Les quatre
états forment donc une base orthonormée de $\mathcal{H}^4$, la \emph{base de Bell}, constituée
uniquement d'états intriqués. Le circuit de la section~\ref{sec:circuit-bell} fait passer de la
base de calcul à cette base.

\paragraph{Ce qui les distingue : deux parités.} Sur la base de calcul,
$Z\otimes Z\ket{ij}=(-1)^{i+j}\ket{ij}$ (valeur propre $+1$ si les deux chiffres sont égaux, $-1$
sinon) et $X\otimes X\ket{ij}=\ket{\bar i\bar j}$ (les deux chiffres sont inversés, avec
$\bar i=1-i$). On en déduit, par exemple,
\[
  X\otimes X\ket{\Phi^-}=\tfrac{1}{\sqrt2}\big(\ket{11}-\ket{00}\big)=-\ket{\Phi^-},
  \qquad
  Z\otimes Z\ket{\Psi^+}=\tfrac{1}{\sqrt2}\big(-\ket{01}-\ket{10}\big)=-\ket{\Psi^+},
\]
et le tableau complet des valeurs propres :
\begin{center}
\renewcommand{\arraystretch}{1.3}
\begin{tabular}{@{}lcccc@{}}
\toprule
 & $\ket{\Phi^+}$ & $\ket{\Phi^-}$ & $\ket{\Psi^+}$ & $\ket{\Psi^-}$ \\
\midrule
$Z\otimes Z$ & $+1$ & $+1$ & $-1$ & $-1$ \\
$X\otimes X$ & $+1$ & $-1$ & $+1$ & $-1$ \\
\bottomrule
\end{tabular}
\end{center}
Les deux opérateurs commutent, car $ZX=-XZ$ donne
$(Z\otimes Z)(X\otimes X)=ZX\otimes ZX=XZ\otimes XZ=(X\otimes X)(Z\otimes Z)$ ; les quatre états de
Bell sont leurs états propres communs, et chaque couple de valeurs propres apparaît une fois
(figure~\ref{fig:grille-bell}). Le même calcul montre que $X\otimes Z$ et $Z\otimes X$
commutent : $(X\otimes Z)(Z\otimes X)=XZ\otimes ZX=(-ZX)\otimes(-XZ)=ZX\otimes XZ=(Z\otimes X)(X\otimes Z)$.
Plus généralement, $P\otimes Q$ et $P'\otimes Q'$ (matrices de Pauli) commutent si $P,P'$
anticommutent et si $Q,Q'$ anticommutent : les deux signes $-1$ se compensent.

\begin{figure}[!ht]
\centering
\adjustbox{max width=\linewidth}{%
\begin{tikzpicture}[>=Stealth, font=\small,
    cell/.style={draw=insarouge, fill=insarouge!6, rounded corners=2pt, thick, minimum width=5.4cm,
                 minimum height=1.35cm, align=center},
    tete/.style={text=insagris, align=center}]
  \node[tete] at (3.0,2.0) {$X\otimes X=+1$\\ {\footnotesize mêmes résultats en base $X$}};
  \node[tete] at (8.8,2.0) {$X\otimes X=-1$\\ {\footnotesize résultats opposés en base $X$}};
  \node[tete, anchor=east] at (-0.1,0.75) {$Z\otimes Z=+1$\\ {\footnotesize mêmes résultats}\\ {\footnotesize en base $Z$}};
  \node[tete, anchor=east] at (-0.1,-0.85) {$Z\otimes Z=-1$\\ {\footnotesize résultats opposés}\\ {\footnotesize en base $Z$}};
  \node[cell] at (3.0,0.75) {$\ket{\Phi^+}=\frac{1}{\sqrt2}\big(\ket{00}+\ket{11}\big)$};
  \node[cell] at (8.8,0.75) {$\ket{\Phi^-}=\frac{1}{\sqrt2}\big(\ket{00}-\ket{11}\big)$};
  \node[cell] at (3.0,-0.85) {$\ket{\Psi^+}=\frac{1}{\sqrt2}\big(\ket{01}+\ket{10}\big)$};
  \node[cell] at (8.8,-0.85) {$\ket{\Psi^-}=\frac{1}{\sqrt2}\big(\ket{01}-\ket{10}\big)$};
\end{tikzpicture}}
\caption{Les quatre états de Bell, classés par les valeurs propres de $Z\otimes Z$ (lignes) et de
$X\otimes X$ (colonnes). La valeur propre est le produit des deux résultats $\pm1$ d'une mesure des
deux qubits dans la base correspondante : $+1$ si les résultats sont égaux, $-1$ s'ils sont opposés.}
\label{fig:grille-bell}
\end{figure}

Les corrélations ne se limitent donc pas à la base $Z$ : $\ket{\Phi^+}$ s'écrit aussi
$\frac{1}{\sqrt2}\big(\ket{++}+\ket{--}\big)$ et donne des résultats égaux en base $X$ comme en base $Z$
(exercice~\ref{ex:bellX}).

\paragraph{D'un état de Bell à l'autre.} Une porte agissant sur un seul qubit suffit. Avec
$X\ket{0}=\ket{1}$, $X\ket{1}=\ket{0}$ et $Z\ket{1}=-\ket{1}$, en agissant sur le second qubit :
\begin{align*}
  (I\otimes X)\ket{\Phi^+}&=\tfrac{1}{\sqrt2}\big(\ket{01}+\ket{10}\big)=\ket{\Psi^+},\\
  (I\otimes Z)\ket{\Phi^+}&=\tfrac{1}{\sqrt2}\big(\ket{00}-\ket{11}\big)=\ket{\Phi^-},\\
  (I\otimes XZ)\ket{\Phi^+}&=(I\otimes X)\tfrac{1}{\sqrt2}\big(\ket{00}-\ket{11}\big)
    =\tfrac{1}{\sqrt2}\big(\ket{01}-\ket{10}\big)=\ket{\Psi^-}.
\end{align*}
Quatre opérations locales ($I$, $X$, $Z$, $XZ$) sur \emph{un} qubit d'une paire de Bell mènent donc
aux quatre états de la base ; c'est le principe du codage superdense
(exercice~\ref{ex:superdense}).

\begin{remarque}[Le singulet est le même dans toutes les bases]\label{rem:singulet}
Pour toute matrice unitaire $U=\begin{pmatrix}a&b\\c&d\end{pmatrix}$, on a $U\ket{0}=a\ket{0}+c\ket{1}$
et $U\ket{1}=b\ket{0}+d\ket{1}$, donc
\begin{align*}
  (U\otimes U)\big(\ket{01}-\ket{10}\big)
  &=\big(a\ket{0}+c\ket{1}\big)\big(b\ket{0}+d\ket{1}\big)-\big(b\ket{0}+d\ket{1}\big)\big(a\ket{0}+c\ket{1}\big)\\
  &=(ad-bc)\big(\ket{01}-\ket{10}\big),
\end{align*}
c'est-à-dire $(U\otimes U)\ket{\Psi^-}=\det(U)\,\ket{\Psi^-}$. Comme $|\det U|=1$, ce facteur est une
phase globale, sans effet physique : $\ket{\Psi^-}$ s'écrit \emph{de la même façon} dans toutes les
bases orthonormées (avec $U=H$, de déterminant $-1$ : $\ket{\Psi^-}=\frac{1}{\sqrt2}\big(\ket{-+}-\ket{+-}\big)$).
Si Alice et Bob mesurent dans la \emph{même} base, quelle qu'elle soit, leurs résultats sont
toujours opposés : c'est le «~singulet~». Parmi les états de Bell, seul $\ket{\Psi^-}$ a cette propriété.
\end{remarque}

% ---------------------------------------------------------------------
\subsection{Un circuit quantique : le générateur d'états de Bell}\label{sec:circuit-bell}
% ---------------------------------------------------------------------

On assemble une porte de Hadamard et un CNOT. Le système est $(A,B)$ avec le ket
$\ket{a,b}=\ket{a}_A\otimes\ket{b}_B$ ; d'après la convention de la section~\ref{sec:circuits}, $B$ est dessiné
en haut et $A$ en bas (figure~\ref{fig:circuit-bell}). Le circuit a deux étapes :
\begin{enumerate}[nosep]
  \item la porte $I\otimes H$ : identité sur $A$, Hadamard sur $B$ ;
  \item un CNOT de contrôle $B$ et de cible $A$ (celui de la section~\ref{sec:cnot}, dont le
        contrôle est le second chiffre).
\end{enumerate}
Les traits pointillés délimitent les états $\ket{\Pi_0}$, $\ket{\Pi_1}$ et $\ket{\Pi_2}$ :
avant les portes, après $I\otimes H$, après le CNOT. Le circuit transforme la base de calcul
en base de Bell (section~\ref{sec:bell}).

\begin{figure}[!ht]
\centering
\begin{tikzpicture}[>=Stealth, font=\small]
  % fils : B en haut, A en bas
  \draw (0,1.4) -- (7.6,1.4);
  \draw (0,0) -- (7.6,0);
  \node[anchor=east] at (0,1.4) {$\ket{0}_B$};
  \node[anchor=east] at (0,0) {$\ket{0}_A$};
  % première étape : H sur B, identité (pointillée) sur A
  \node[porte] at (2.2,1.4) {$H$};
  \node[draw=insagris, dashed, fill=white, minimum size=7mm] at (2.2,0) {$I$};
  % CNOT : contrôle B (haut), cible A (bas)
  \draw[insarouge, thick] (5,1.4) -- (5,0);
  \pic at (5,1.4) {ctl};
  \pic at (5,0) {cible};
  % accolade de sortie
  \draw[insagris, thick] (7.75,1.6) -- (7.9,1.6) -- (7.9,-0.2) -- (7.75,-0.2);
  \node[anchor=west] at (7.9,0.7) {$\ket{\beta^{00}}_{AB}$};
  % états intermédiaires
  \foreach \x in {0.7,3.6,6.4}
    \draw[dashed, insagris!70] (\x,1.9) -- (\x,-0.7);
  \node[anchor=north, font=\footnotesize] at (0.7,-0.7) {$\ket{\Pi_0}=\ket{00}$};
  \node[anchor=north, font=\footnotesize] at (3.6,-0.7) {$\ket{\Pi_1}=\ket{0}\otimes\ket{+}$};
  \node[anchor=north, font=\footnotesize] at (6.4,-0.7) {$\ket{\Pi_2}=\ket{\beta^{00}}$};
\end{tikzpicture}
\caption{Générateur d'états de Bell : $I\otimes H$, puis un CNOT de contrôle $B$ (fil du haut)
et de cible $A$ (fil du bas). L'entrée représentée est $\ket{00}$.}
\label{fig:circuit-bell}
\end{figure}

\subsubsection{Entrée \texorpdfstring{$\ket{00}$}{|00>} : calcul pas à pas}

\begin{align*}
  \ket{\Pi_0}&=\ket{0}_A\otimes\ket{0}_B=\ket{00},\\[1ex]
  \ket{\Pi_1}&=(I\otimes H)\ket{\Pi_0}=I\ket{0}_A\otimes H\ket{0}_B=\ket{0}\otimes\ket{+}
    =\frac{\ket{00}+\ket{01}}{\sqrt2},\\[1ex]
  \ket{\Pi_2}&=\mathrm{CNOT}\ket{\Pi_1}=\frac{\mathrm{CNOT}\ket{00}+\mathrm{CNOT}\ket{01}}{\sqrt2}
    =\frac{\ket{00}+\ket{11}}{\sqrt2}=\ket{\beta^{00}}=\ket{\Phi^+}.
\end{align*}
Dans $\ket{\Pi_1}$, le premier chiffre est celui de $A$ et le second celui de $B$. Le CNOT a
pour contrôle $B$ : $\ket{00}$ ($B=0$) ne change pas, $\ket{01}$ ($B=1$) devient $\ket{11}$
puisque $A$ est inversé. Le résultat est un \emph{état de Bell}.

Avec les matrices, la règle $A\otimes B=(A_{ij}B)$ donne
\[
  I\otimes H=\begin{pmatrix}1\cdot H&0\\0&1\cdot H\end{pmatrix}
  =\frac{1}{\sqrt2}\begin{pmatrix}1&1&0&0\\1&-1&0&0\\0&0&1&1\\0&0&1&-1\end{pmatrix}.
\]
Sa première colonne est $\ket{\Pi_1}=\frac{1}{\sqrt2}(1,1,0,0)^{\mathsf T}$ (image de $\ket{00}$).
Le CNOT échange les amplitudes de $\ket{01}$ et de $\ket{11}$ (section~\ref{sec:cnot}), donc
\[
  \ket{\Pi_2}=\mathrm{CNOT}\,\frac{1}{\sqrt2}\begin{pmatrix}1\\1\\0\\0\end{pmatrix}
  =\frac{1}{\sqrt2}\begin{pmatrix}1\\0\\0\\1\end{pmatrix}.
\]

\subsubsection{Les quatre entrées}

Prenons $B$ dans l'état $\ket{x}$ et $A$ dans l'état $\ket{y}$, avec $x,y\in\{0,1\}$. Alors
$\ket{\Pi_0}=\ket{y}_A\otimes\ket{x}_B=\ket{y,x}$. Comme
$H\ket{x}=\frac{1}{\sqrt2}\big(\ket{0}+(-1)^x\ket{1}\big)$ (c'est $\ket{+}$ pour $x=0$ et
$\ket{-}$ pour $x=1$),
\[
  \ket{\Pi_1}=\ket{y}\otimes H\ket{x}=\frac{1}{\sqrt2}\Big(\ket{y,0}+(-1)^x\ket{y,1}\Big).
\]
Le CNOT (contrôle $B$) ne change pas $\ket{y,0}$ et transforme $\ket{y,1}$ en
$\ket{\bar y,1}$, avec $\bar y=1-y$ :
\[
  \ket{\Pi_2}=\ket{\beta^{xy}}=\frac{1}{\sqrt2}\Big(\ket{y,0}+(-1)^x\ket{\bar y,1}\Big).
\]
Le premier indice de $\ket{\beta^{xy}}$ est l'entrée de $B$ (le qubit qui reçoit $H$), le
second celle de $A$ (la cible). Les quatre cas :

\begin{center}
\renewcommand{\arraystretch}{1.7}
\begin{tabular}{@{}ccccc@{}}
\toprule
$(B,A)=(x,y)$ & $\ket{\Pi_0}=\ket{y,x}$ & $\ket{\Pi_1}$ & $\ket{\Pi_2}=\ket{\beta^{xy}}$ & Nom \\
\midrule
$(0,0)$ & $\ket{00}$ & $\frac{1}{\sqrt2}(\ket{00}+\ket{01})$ & $\frac{1}{\sqrt2}(\ket{00}+\ket{11})$ & $\ket{\Phi^+}$ \\
$(1,0)$ & $\ket{01}$ & $\frac{1}{\sqrt2}(\ket{00}-\ket{01})$ & $\frac{1}{\sqrt2}(\ket{00}-\ket{11})$ & $\ket{\Phi^-}$ \\
$(0,1)$ & $\ket{10}$ & $\frac{1}{\sqrt2}(\ket{10}+\ket{11})$ & $\frac{1}{\sqrt2}(\ket{10}+\ket{01})$ & $\ket{\Psi^+}$ \\
$(1,1)$ & $\ket{11}$ & $\frac{1}{\sqrt2}(\ket{10}-\ket{11})$ & $\frac{1}{\sqrt2}(\ket{10}-\ket{01})$ & $-\ket{\Psi^-}$ \\
\bottomrule
\end{tabular}
\end{center}
Détail de la dernière ligne : $\ket{\Pi_1}=\ket{1}\otimes\ket{-}=\frac{1}{\sqrt2}(\ket{10}-\ket{11})$ ;
le CNOT laisse $\ket{10}$ et envoie $\ket{11}$ sur $\ket{01}$, d'où
$\frac{1}{\sqrt2}(\ket{10}-\ket{01})=-\frac{1}{\sqrt2}(\ket{01}-\ket{10})=-\ket{\Psi^-}$.

\paragraph{Forme matricielle.} Le CNOT échange les lignes $2$ et $4$ de
$I\otimes H$ (il échange les amplitudes de $\ket{01}$ et $\ket{11}$), donc
\[
  \mathrm{CNOT}\,(I\otimes H)=\frac{1}{\sqrt2}\begin{pmatrix}
    1 & 1 & 0 & 0\\ 0 & 0 & 1 & -1\\ 0 & 0 & 1 & 1\\ 1 & -1 & 0 & 0
  \end{pmatrix}.
\]
Ses colonnes sont les images de $\ket{00},\ket{01},\ket{10},\ket{11}$, c'est-à-dire de
$(B,A)=(0,0),(1,0),(0,1),(1,1)$ : ce sont $\ket{\beta^{00}}$, $\ket{\beta^{10}}$,
$\ket{\beta^{01}}$, $\ket{\beta^{11}}$, soit $\ket{\Phi^+}$, $\ket{\Phi^-}$, $\ket{\Psi^+}$
et $-\ket{\Psi^-}$.

\subsubsection{Matrice de densité de l'état de sortie}

Pour $\ket{\Psi}=\ket{\beta^{00}}=\frac{1}{\sqrt2}\ket{00}+\frac{1}{\sqrt2}\ket{11}$ :
\[
  \rho_{AB}=\ket{\Psi}\bra{\Psi}
  =\begin{pmatrix}\frac{1}{\sqrt2}\\0\\0\\\frac{1}{\sqrt2}\end{pmatrix}
   \begin{pmatrix}\frac{1}{\sqrt2}&0&0&\frac{1}{\sqrt2}\end{pmatrix}
  =\frac12\begin{pmatrix}1&0&0&1\\0&0&0&0\\0&0&0&0\\1&0&0&1\end{pmatrix}.
\]
Sous forme de ket-bra :
$\rho_{AB}=\frac12\big(\ket{00}\bra{00}+\ket{00}\bra{11}+\ket{11}\bra{00}+\ket{11}\bra{11}\big)$.
On vérifie $\operatorname{Tr}\rho_{AB}=\frac12(1+1)=1$ et
\[
  \rho_{AB}^2=\frac14\begin{pmatrix}2&0&0&2\\0&0&0&0\\0&0&0&0\\2&0&0&2\end{pmatrix}=\rho_{AB},
\]
donc $\operatorname{Tr}\rho_{AB}^2=1$ : le couple $(A,B)$ est dans un état \emph{pur}.

\begin{remarque}[Porte intriquante]
Les états d'entrée $\ket{ij}$ sont séparables ($D=0$) ; les sorties sont intriquées
($D=\pm\frac12$). Aucune porte agissant séparément sur $A$ et sur $B$ ($U_A\otimes U_B$,
section~\ref{sec:op2}) ne peut créer de l'intrication à partir d'un produit, puisque
$(U_A\otimes U_B)(\ket{u}\otimes\ket{v})=U_A\ket{u}\otimes U_B\ket{v}$ reste un produit.
Une porte à deux qubits comme le CNOT est indispensable.
\end{remarque}

% ---------------------------------------------------------------------
\subsection{Trace partielle : l'état d'un qubit pris seul}\label{sec:trace}
% ---------------------------------------------------------------------

Un état intriqué n'a pas de ket pour $A$ seul (section~\ref{sec:separable}). On peut pourtant
décrire $A$ par une matrice de densité, obtenue en «~oubliant~» $B$.

\begin{definition}{Trace partielle}
Soit $\rho_{AB}$ la matrice de densité du couple $(A,B)$. L'état du qubit $A$ pris seul est
$\rho_A=\operatorname{Tr}_B(\rho_{AB})$, celui de $B$ est $\rho_B=\operatorname{Tr}_A(\rho_{AB})$. La
trace partielle est linéaire et remplace la partie tracée de chaque terme par sa trace :
\[
  \operatorname{Tr}_B\big(\ket{a}\bra{a'}\otimes\ket{b}\bra{b'}\big)=\ket{a}\bra{a'}\;\braket{b'}{b},
  \qquad
  \operatorname{Tr}_A\big(\ket{a}\bra{a'}\otimes\ket{b}\bra{b'}\big)=\braket{a'}{a}\;\ket{b}\bra{b'}.
\]
\end{definition}

\paragraph{En matrices.} On découpe $\rho_{AB}$ (matrice $4\times4$, lignes et colonnes ordonnées
$00,01,10,11$) en quatre blocs $2\times2$ :
\[
  \rho_{AB}=\begin{pmatrix}R_{00}&R_{01}\\R_{10}&R_{11}\end{pmatrix}
  \quad\Longrightarrow\quad
  \rho_A=\begin{pmatrix}\operatorname{Tr}R_{00}&\operatorname{Tr}R_{01}\\
                        \operatorname{Tr}R_{10}&\operatorname{Tr}R_{11}\end{pmatrix},
  \qquad
  \rho_B=R_{00}+R_{11}.
\]
Autrement dit $(\rho_A)_{ij}=\rho_{(i0),(j0)}+\rho_{(i1),(j1)}$ : $\rho_A$ est la matrice des traces
des blocs, et $\rho_B$ la somme des deux blocs diagonaux.

\begin{remarque}[Pourquoi cette définition]
Toutes les prédictions qui portent sur $A$ seul viennent d'opérateurs de la forme $O_A\otimes I$, et
\[
  \operatorname{Tr}\big(\rho_{AB}\,(O_A\otimes I)\big)=\operatorname{Tr}\big(\rho_A\,O_A\big)
\]
(on le vérifie sur $\ket{a}\bra{a'}\otimes\ket{b}\bra{b'}$, puis par linéarité). La matrice $\rho_A$
contient donc exactement ce que l'on peut prédire sur $A$ sans regarder $B$.
\end{remarque}

\paragraph{Exemple 1 : état séparable $\ket{0}\otimes\ket{+}$.} Ici
$\rho_{AB}=\ket{0}\bra{0}\otimes\ket{+}\bra{+}$, donc
\[
  \rho_A=\ket{0}\bra{0}\,\braket{+}{+}=\ket{0}\bra{0},
  \qquad
  \rho_B=\braket{0}{0}\,\ket{+}\bra{+}=\ket{+}\bra{+}.
\]
Avec les matrices,
\[
  \rho_{AB}=\frac12\begin{pmatrix}1&1&0&0\\1&1&0&0\\0&0&0&0\\0&0&0&0\end{pmatrix},
  \qquad
  R_{00}=\frac12\begin{pmatrix}1&1\\1&1\end{pmatrix},\quad R_{01}=R_{10}=R_{11}=0,
\]
d'où
\[
  \rho_A=\begin{pmatrix}\operatorname{Tr}R_{00}&0\\0&0\end{pmatrix}=\begin{pmatrix}1&0\\0&0\end{pmatrix},
  \qquad
  \rho_B=R_{00}+R_{11}=\frac12\begin{pmatrix}1&1\\1&1\end{pmatrix}.
\]
Un qubit d'un état séparable garde son propre état pur.

\paragraph{Exemple 2 : état de Bell $\ket{\Phi^+}$.} On écrit
$\rho_{AB}=\ket{\Phi^+}\bra{\Phi^+}=\frac12\big(\ket{00}\bra{00}+\ket{00}\bra{11}+\ket{11}\bra{00}
+\ket{11}\bra{11}\big)$. Chaque terme est un produit d'opérateurs à un qubit, par exemple
$\ket{00}\bra{11}=\ket{0}\bra{1}\otimes\ket{0}\bra{1}$, donc
\begin{align*}
  \rho_A=\operatorname{Tr}_B(\rho_{AB})
  &=\tfrac12\Big(\ket{0}\bra{0}\,\braket{0}{0}+\ket{0}\bra{1}\,\braket{1}{0}
      +\ket{1}\bra{0}\,\braket{0}{1}+\ket{1}\bra{1}\,\braket{1}{1}\Big)\\
  &=\tfrac12\Big(\ket{0}\bra{0}\cdot1+\ket{0}\bra{1}\cdot0+\ket{1}\bra{0}\cdot0+\ket{1}\bra{1}\cdot1\Big)
   =\tfrac12\big(\ket{0}\bra{0}+\ket{1}\bra{1}\big)=\tfrac12\,\mathrm{Id}.
\end{align*}
Par les blocs,
\[
  \rho_{AB}=\frac12\begin{pmatrix}1&0&0&1\\0&0&0&0\\0&0&0&0\\1&0&0&1\end{pmatrix},
  \qquad
  R_{00}=\frac12\begin{pmatrix}1&0\\0&0\end{pmatrix},\quad
  R_{11}=\frac12\begin{pmatrix}0&0\\0&1\end{pmatrix},\quad
  \operatorname{Tr}R_{01}=\operatorname{Tr}R_{10}=0,
\]
d'où $\rho_A=\frac12\mathrm{Id}$ et $\rho_B=R_{00}+R_{11}=\frac12\mathrm{Id}$ : le même résultat.

Le qubit $A$ est donc dans l'état \emph{complètement mixte} ($\operatorname{Tr}\rho_A^2=\frac14+\frac14=\frac12$),
au centre de la boule de Bloch (section~\ref{sec:bloch}), alors que le couple est dans un état pur
($\operatorname{Tr}\rho_{AB}^2=1$, section~\ref{sec:circuit-bell}). On connaît parfaitement l'état du
couple mais pas celui d'un qubit pris seul : l'information est entièrement contenue dans les
corrélations entre $A$ et $B$. Mesuré seul, $A$ donne $0$ ou $1$ avec la probabilité $\frac12$, dans
toutes les bases.

\paragraph{Formule pratique pour un état pur.} Pour $\ket{\Psi}=\sum_{ij}c_{ij}\ket{ij}$, notons
$C=\begin{pmatrix}c_{00}&c_{01}\\c_{10}&c_{11}\end{pmatrix}$ la matrice $2\times2$ des coefficients
(ligne $i$ : état de $A$, colonne $j$ : état de $B$). Comme $\rho_{(ij),(i'j')}=c_{ij}\,c_{i'j'}^*$,
\[
  (\rho_A)_{ii'}=\sum_jc_{ij}\,c_{i'j}^*,\qquad\text{soit}\qquad \rho_A=C\,C^\dagger .
\]
\emph{Exemple 3.} Pour $\ket{\Psi}=\frac12\ket{00}+\frac12\ket{01}+\frac{1}{\sqrt2}\ket{11}$ (exemple~2
de la section~\ref{sec:mesure2}), $C=\begin{pmatrix}\frac12&\frac12\\0&\frac{1}{\sqrt2}\end{pmatrix}$ et
\[
  \rho_A=CC^\dagger=\begin{pmatrix}\frac14+\frac14&0+\frac{1}{2\sqrt2}\\[0.4ex]\frac{1}{2\sqrt2}&\frac12\end{pmatrix}
  =\begin{pmatrix}\frac12&\frac{\sqrt2}{4}\\[0.4ex]\frac{\sqrt2}{4}&\frac12\end{pmatrix},
  \qquad
  \operatorname{Tr}\rho_A^2=\tfrac14+\tfrac14+2\cdot\tfrac18=\tfrac34 .
\]
Le qubit $A$ est mixte, mais pas complètement : cet état est intriqué sans être un état de Bell.

\begin{resultat}{Pureté d'un qubit et séparabilité}
Pour un état pur $\ket{\Psi}=\sum c_{ij}\ket{ij}$ de deux qubits, avec $D=c_{00}c_{11}-c_{01}c_{10}$,
\[
  \operatorname{Tr}\rho_A^2=\operatorname{Tr}\rho_B^2=1-2|D|^2 .
\]
Donc $\ket{\Psi}$ est séparable ($D=0$) si et seulement si $\rho_A$ est pur. Plus $|D|$ est grand,
plus $\rho_A$ est mixte, jusqu'à $|D|=\frac12$ et $\rho_A=\mathrm{Id}/2$ pour les états de Bell.
\end{resultat}

\emph{Démonstration.} La matrice $\rho_A=CC^\dagger$ a pour trace $\sum|c_{ij}|^2=1$ et pour
déterminant $\det C\,\overline{\det C}=|D|^2$. Si $\lambda_1,\lambda_2$ sont ses valeurs propres,
$\lambda_1+\lambda_2=1$ et $\operatorname{Tr}\rho_A^2=\lambda_1^2+\lambda_2^2=1-2\lambda_1\lambda_2
=1-2|D|^2$ ; le calcul est le même pour $\rho_B$. \hfill$\square$

Comme $\operatorname{Tr}\rho_A^2\geq\frac12$ pour un qubit, on en déduit $|D|\leq\frac12$. Avec
$\operatorname{Tr}\rho_A^2=\frac{1+|\vec r|^2}{2}$ (section~\ref{sec:bloch}), le vecteur de Bloch de $A$
a pour norme $|\vec r|=\sqrt{1-4|D|^2}$ : il part de la sphère pour un état séparable et se rapproche
du centre quand l'intrication croît (figure~\ref{fig:trace-bloch}).

\begin{center}
\renewcommand{\arraystretch}{1.5}
\begin{tabular}{@{}lcccl@{}}
\toprule
\textbf{État du couple} & $D$ & $\operatorname{Tr}\rho_A^2$ & $|\vec r|$ & \textbf{Qubit $A$} \\
\midrule
$\ket{0}\otimes\ket{+}$ & $0$ & $1$ & $1$ & pur, sur la sphère \\
$\frac12\ket{00}+\frac12\ket{01}+\frac{1}{\sqrt2}\ket{11}$ & $\frac{\sqrt2}{4}$ & $\frac34$
  & $\frac{1}{\sqrt2}$ & mixte, dans la boule \\
$\ket{\Phi^+}$ (comme les trois autres états de Bell) & $\pm\frac12$ & $\frac12$ & $0$
  & complètement mixte, au centre \\
\bottomrule
\end{tabular}
\end{center}

\begin{figure}[!ht]
\centering
\adjustbox{max width=\linewidth}{%
\begin{tikzpicture}[>=Stealth, font=\small]
  \def\R{2.3}
  \fill[insagris!12] (0,0) circle (\R);
  \draw[insagris, dotted] (-\R,0) -- (\R,0);
  \draw[insagris, dotted] (0,-\R) -- (0,\R);
  \draw[insagris, dashed] (0,0) circle ({\R*0.70711});
  \draw[insarouge, very thick] (0,0) circle (\R);
  \node[font=\footnotesize, text=insagris, anchor=south west] at (\R,0) {$x$};
  \node[font=\footnotesize, text=insagris, anchor=south west] at (0.03,\R) {$z$};
  % vecteurs de Bloch de rho_A pour trois états du couple
  \fill[insarouge] (0,\R) circle (2.4pt);
  \fill[insagris!80!black] ({\R*0.70711},0) circle (2.4pt);
  \fill[insagris!80!black] (0,0) circle (2.4pt);
  \draw[insagris] (0,\R) -- (3.1,\R);
  \node[anchor=west, font=\footnotesize, align=left] at (3.1,\R)
    {$\ket{0}\otimes\ket{+}$ : $\rho_A=\ket{0}\bra{0}$\\ $D=0$, \ $|\vec r|=1$ (état pur)};
  \draw[insagris] ({\R*0.70711},0) -- (3.1,0);
  \node[anchor=west, font=\footnotesize, align=left] at (3.1,0)
    {$\frac12\ket{00}+\frac12\ket{01}+\frac{1}{\sqrt2}\ket{11}$\\
     $D=\frac{\sqrt2}{4}$, \ $|\vec r|=\frac{1}{\sqrt2}$ (mixte)};
  \draw[insagris] (0,0) -- (3.1,-\R);
  \node[anchor=west, font=\footnotesize, align=left] at (3.1,-\R)
    {$\ket{\Phi^+}$ : $\rho_A=\mathrm{Id}/2$\\ $D=\frac12$, \ $\vec r=\vec 0$ (centre)};
\end{tikzpicture}}
\caption{Le vecteur de Bloch de $\rho_A$ (boule de Bloch vue en coupe) pour trois états du couple
$(A,B)$. Le cercle pointillé est la sphère de rayon $1/\sqrt2$. Plus l'état est intriqué, plus le
point s'éloigne de la sphère des états purs : $|\vec r|=\sqrt{1-4|D|^2}$.}
\label{fig:trace-bloch}
\end{figure}

% ---------------------------------------------------------------------
\subsection{Information : variables aléatoires et états quantiques}\label{sec:info}
% ---------------------------------------------------------------------

\begin{definition}{Entropie, information mutuelle, entropie conditionnelle}
Soit $A$ une variable aléatoire de loi $p(a)$. Son \emph{entropie de Shannon}, en bits, est
\[
  H(A)=-\sum_a p(a)\log_2p(a)\qquad(\text{avec }0\log_20=0).
\]
Elle mesure l'incertitude moyenne sur $A$ : une pièce équilibrée a $H=1$ bit. Pour un couple
$(A,B)$ de loi $p(a,b)$, $H(A,B)=-\sum_{a,b}p(a,b)\log_2p(a,b)$, et
\begin{align*}
  I(A;B)&=H(A)+H(B)-H(A,B) &&\text{(information mutuelle)},\\
  H(A|B)&=H(A,B)-H(B)=H(A)-I(A;B) &&\text{(entropie conditionnelle)}.
\end{align*}
\end{definition}

Les deux écritures de $H(A|B)$ sont égales : la définition de $I(A;B)$ donne
$H(A,B)=H(A)+H(B)-I(A;B)$, donc $H(A,B)-H(B)=H(A)-I(A;B)$. L'information mutuelle $I(A;B)$
mesure ce que la connaissance de $B$ apprend sur $A$ ; $H(A|B)$ est l'incertitude qui reste sur
$A$ quand on connaît $B$.

\subsubsection{Deux variables aléatoires indépendantes}

Alice et Bob lancent chacun une pièce équilibrée, sans lien entre les deux lancers :
$p(a,b)=\frac14$ pour les quatre couples. Chaque pièce donne $0$ ou $1$ avec la probabilité $\frac12$, donc
\begin{align*}
  H(A)&=-\Big(\tfrac12\log_2\tfrac12+\tfrac12\log_2\tfrac12\Big)=1\text{ bit},\\
  H(B)&=1\text{ bit (même calcul)},\\
  H(A,B)&=-4\cdot\tfrac14\log_2\tfrac14=2\text{ bits},\\
  I(A;B)&=H(A)+H(B)-H(A,B)=1+1-2=0,\\
  H(A|B)&=H(A)-I(A;B)=1-0=H(A)=1\text{ bit}.
\end{align*}
Connaître $B$ n'apprend rien sur $A$.

\subsubsection{Deux variables solidaires}

Les deux pièces tombent toujours du même côté : $p(0,0)=p(1,1)=\frac12$ et
$p(0,1)=p(1,0)=0$. Les lois de chaque pièce sont inchangées ($p(a{=}0)=p(0,0)+p(0,1)=\frac12$,
etc.), donc $H(A)=H(B)=1$ bit, mais
\begin{align*}
  H(A,B)&=-2\cdot\tfrac12\log_2\tfrac12-2\cdot0\log_20=1\text{ bit},\\
  I(A;B)&=1+1-1=1\text{ bit},\qquad H(A|B)=1-1=0.
\end{align*}
Connaître $B$ détermine $A$ : plus aucune incertitude ne reste.

\subsubsection{Deux qubits dans un état de Bell}

Pour un système quantique, on remplace la loi de probabilité par la matrice de densité.

\begin{definition}{Entropie de von Neumann}
L'\emph{entropie de von Neumann} d'un état $\rho$, de valeurs propres $\lambda_i$, est
\[
  S(\rho)=-\operatorname{Tr}(\rho\log_2\rho)=-\sum_i\lambda_i\log_2\lambda_i .
\]
Elle vaut $0$ pour un état pur et $1$ bit pour le qubit complètement mélangé $\mathrm{Id}/2$.
En quantique, $H(A)=S(\rho_A)$, $H(B)=S(\rho_B)$ et $H(A,B)=S(\rho_{AB})$ ; les définitions de
$I(A;B)$ et de $H(A|B)$ sont les mêmes.
\end{definition}

On prend l'état de Bell $\ket{\Phi^+}$ produit par le circuit de la
section~\ref{sec:circuit-bell}. Le couple est dans un état pur,
$\rho_{AB}=\ket{\Phi^+}\bra{\Phi^+}$ ; chaque qubit pris seul est dans l'état
$\rho_A=\rho_B=\frac12\mathrm{Id}$ (trace partielle, section~\ref{sec:trace}).

Calcul des entropies :
\begin{itemize}[nosep]
  \item $\rho_{AB}$ est pur : ses valeurs propres sont $1,0,0,0$, donc
        $H(A,B)=S(\rho_{AB})=-1\cdot\log_21=0$ ;
  \item $\rho_A$ et $\rho_B$ ont pour valeurs propres $\frac12,\frac12$, donc
        $H(A)=H(B)=-\big(\frac12\log_2\frac12+\frac12\log_2\frac12\big)=1$ bit.
\end{itemize}
D'où
\begin{align*}
  I(A;B)&=H(A)+H(B)-H(A,B)=1+1-0=2\text{ bits},\\
  H(A|B)&=H(A,B)-H(B)=0-1=-1\text{ bit}\qquad(=H(A)-I(A;B)=1-2).
\end{align*}

\begin{remarque}[Une entropie conditionnelle négative]
Pour des variables classiques, $p(a,b)\leq p(b)$, donc $-\log_2p(a,b)\geq-\log_2p(b)$ et
$H(A,B)\geq H(B)$ : on a toujours $H(A|B)\geq0$, et $I(A;B)\leq H(A)$. Ici $H(A|B)=-1$ et
$I(A;B)=2>H(A)$ : ces valeurs n'ont pas d'équivalent classique, c'est une signature de
l'intrication. Le résultat est surprenant mais mathématiquement correct.
\end{remarque}

\begin{center}
\renewcommand{\arraystretch}{1.4}
\begin{tabular}{@{}lccc@{}}
\toprule
 & indépendantes & solidaires & état de Bell \\
\midrule
$H(A)$ & $1$ & $1$ & $1$ \\
$H(B)$ & $1$ & $1$ & $1$ \\
$H(A,B)$ & $2$ & $1$ & $0$ \\
$I(A;B)$ & $0$ & $1$ & $2$ \\
$H(A|B)$ & $1$ & $0$ & $-1$ \\
\bottomrule
\end{tabular}
\end{center}
(en bits ; les deux premières colonnes sont classiques, la troisième est quantique).

\begin{figure}[!ht]
\centering
\adjustbox{max width=\linewidth}{%
\begin{tikzpicture}[>=Stealth, font=\small]
  \def\bw{0.8}
  % axe des ordonnées (bits)
  \draw[->] (0,-1.5) -- (0,2.8) node[above] {bits};
  \foreach \v in {-1,0,1,2} {
    \draw (-0.07,\v) -- (0.07,\v);
    \node[left, font=\footnotesize] at (-0.1,\v) {$\v$};
  }
  \draw[insagris!60] (0,0) -- (11.4,0);
  % limite classique de I(A;B) : H(A) = 1
  \draw[dashed, insarouge] (0,1) -- (11.4,1);
  \node[anchor=west, font=\scriptsize, text=insarouge, align=left] at (11.4,1)
    {limite classique\\ $I(A;B)\leq H(A)=1$};
  \node[anchor=west, font=\scriptsize, text=insagris, align=left] at (11.4,0)
    {limite classique\\ $H(A|B)\geq0$};
  % barres : (groupe, H(A,B), I(A;B), H(A|B))
  \foreach \g/\hab/\mi/\hc in {0/2/0/1, 1/1/1/0, 2/0/2/-1} {
    \pgfmathsetmacro{\x}{0.9+3.7*\g}
    \fill[insagris!65] (\x,0) rectangle (\x+\bw,\hab);
    \fill[insarouge] (\x+\bw+0.1,0) rectangle (\x+2*\bw+0.1,\mi);
    \fill[black!80] (\x+2*\bw+0.2,0) rectangle (\x+3*\bw+0.2,\hc);
    \node[above, font=\footnotesize] at (\x+0.5*\bw,{max(\hab,0)}) {$\hab$};
    \node[above, font=\footnotesize] at (\x+1.5*\bw+0.1,{max(\mi,0)}) {$\mi$};
    \ifnum\hc<0
      \node[below, font=\footnotesize] at (\x+2.5*\bw+0.2,\hc) {$\hc$};
    \else
      \node[above, font=\footnotesize] at (\x+2.5*\bw+0.2,\hc) {$\hc$};
    \fi
  }
  \node[align=center, anchor=north] at (2.15,-1.65) {variables\\ indépendantes};
  \node[align=center, anchor=north] at (5.85,-1.65) {variables\\ solidaires};
  \node[align=center, anchor=north] at (9.55,-1.65) {état de Bell\\ $\ket{\Phi^+}$};
  % légende
  \fill[insagris!65] (0.9,3.35) rectangle (1.2,3.65);
  \node[anchor=west, font=\footnotesize] at (1.25,3.5) {$H(A,B)$};
  \fill[insarouge] (3.6,3.35) rectangle (3.9,3.65);
  \node[anchor=west, font=\footnotesize] at (3.95,3.5) {$I(A;B)$};
  \fill[black!80] (6.3,3.35) rectangle (6.6,3.65);
  \node[anchor=west, font=\footnotesize] at (6.65,3.5) {$H(A|B)$};
\end{tikzpicture}}
\caption{Entropies (en bits) pour deux pièces indépendantes, deux pièces solidaires et un état de Bell.
Dans les trois cas $H(A)=H(B)=1$. Les deux premiers groupes respectent les bornes classiques
$I(A;B)\leq H(A)$ et $H(A|B)\geq0$ ; l'état de Bell les dépasse ($I=2$, $H(A|B)=-1$).}
\label{fig:entropies}
\end{figure}

% ---------------------------------------------------------------------
\subsection{Téléportation quantique}\label{sec:teleportation}
% ---------------------------------------------------------------------

\emph{Problème.} Alice détient un qubit $Q$ dans un état qu'elle ne connaît pas,
$\ket{\varphi}=\alpha\ket{0}+\beta\ket{1}$ avec $|\alpha|^2+|\beta|^2=1$, et elle veut que Bob
obtienne cet état sur l'un de ses qubits. Elle ne peut pas mesurer $Q$ pour connaître
$\alpha$ et $\beta$ (une mesure ne donne qu'un bit et détruit l'état), elle ne peut pas le
copier (théorème de non-clonage, remarque en fin de section) et elle n'envoie pas le qubit
lui-même. Le protocole est dû à \textcite{bennett1993}.

\begin{definition}{Protocole de téléportation quantique}
\textbf{Ressources.}
\begin{itemize}[nosep]
  \item Un \emph{e-bit} : Alice et Bob partagent une paire de qubits $(A,B)$ dans l'état de Bell
        $\ket{\Phi^+}_{BA}=\frac{1}{\sqrt2}\big(\ket{00}+\ket{11}\big)$ ; Alice détient $A$,
        Bob détient $B$.
  \item Le qubit $Q$ à téléporter, chez Alice.
  \item Un canal classique (deux bits).
\end{itemize}
Le système est écrit $(B,A,Q)$ : $\ket{b,a,q}=\ket{b}_B\otimes\ket{a}_A\otimes\ket{q}_Q$.

\textbf{Étapes.}
\begin{enumerate}[nosep]
  \item Alice applique un CNOT de contrôle $Q$ et de cible $A$.
  \item Alice applique une porte $H$ sur $Q$.
  \item Alice mesure $A$ et $Q$ dans la base de calcul, obtient deux bits $(a,q)$ et les envoie
        à Bob.
  \item Bob applique $X$ si $a=1$, puis $Z$ si $q=1$, à son qubit $B$.
\end{enumerate}
\end{definition}

La figure~\ref{fig:tele-schema} montre l'organisation du protocole et la
figure~\ref{fig:teleportation} son circuit.

\begin{figure}[!ht]
\centering
\adjustbox{max width=\linewidth}{%
\begin{tikzpicture}[>=Stealth, font=\small,
    qb/.style={circle, draw=insarouge, fill=insarouge!8, thick, minimum size=8mm, inner sep=0pt},
    zone/.style={draw=insagris, rounded corners=4pt, thick, dashed, inner sep=8pt},
    num/.style={circle, fill=insarouge, text=white, font=\footnotesize\bfseries, inner sep=0pt,
                minimum size=4.6mm},
    etape/.style={anchor=west, font=\footnotesize, text=insarouge},
    onde/.style={thick, insarouge, decorate,
                 decoration={snake, amplitude=1.3pt, segment length=7pt}}]
  % Alice : Q (état inconnu) et A
  \node[qb] (Q) at (0,1.3) {$Q$};
  \node[qb] (A) at (0,-0.5) {$A$};
  \node[zone, fit=(Q)(A), label={[text=insagris]above:Alice}] (alice) {};
  \node[anchor=east] at (-0.95,1.3) {$\ket{\varphi}=\alpha\ket{0}+\beta\ket{1}$};
  \draw[insarouge, thick] (Q) -- (A);
  % Bob : B
  \node[qb] (B) at (9,-0.5) {$B$};
  \node[zone, fit=(B), label={[text=insagris]above:Bob}] (bob) {};
  % source de la paire de Bell
  \node[draw=insarouge, fill=insarouge!10, rounded corners=2pt, thick, minimum width=3.2cm,
        minimum height=0.9cm] (src) at (4.5,-2.3) {source $\ket{\Phi^+}_{BA}$};
  \draw[onde] (src.west) -| (A.south);
  \draw[onde] (src.east) -| (B.south);
  % étapes d'Alice : pastilles sur le bord droit de sa zone
  \node[num] (n1) at (alice.east |- Q) {1};
  \node[num] (n2) at (alice.east) {2};
  \node[num] (n3) at (alice.east |- A) {3};
  \node[etape] at (n1.east) {\ CNOT de $Q$ vers $A$};
  \node[etape] at (n2.east) {\ porte $H$ sur $Q$};
  % canal classique
  \draw[double, double distance=1.2pt, ->, insagris] (n3.east) -- (bob.west);
  \node[above, font=\footnotesize] at (4.6,-0.45) {deux bits classiques $(a,q)$};
  \node[below, font=\footnotesize, text=insarouge] at (4.6,-0.55) {mesure de $A$ et $Q$ (étape 3)};
  % correction de Bob
  \node[num] (n4) at (bob.east) {4};
  \node[etape, align=left] at (n4.east) {\ $X$ si $a=1$, puis $Z$ si $q=1$};
  \node[anchor=west, font=\footnotesize, align=left] at ($(n4.east)+(0.05,-0.8)$)
    {\ $B$ est alors dans l'état $\ket{\varphi}$};
\end{tikzpicture}}
\caption{Organisation de la téléportation. Alice possède le qubit $Q$ à téléporter et la moitié $A$
d'une paire de Bell, Bob possède l'autre moitié $B$. (1)~CNOT de $Q$ vers $A$, (2)~porte $H$ sur $Q$,
(3)~mesure de $A$ et $Q$, dont les résultats $(a,q)$ partent par un canal classique, (4)~corrections de
Bob. Aucun qubit ne voyage entre Alice et Bob.}
\label{fig:tele-schema}
\end{figure}

\begin{figure}[!ht]
\centering
\begin{adjustbox}{max width=\linewidth}
\begin{tikzpicture}[>=Stealth, font=\small]
  % fils quantiques : Q en haut, A au milieu, B en bas
  \draw (0,2.4) -- (5.3,2.4);
  \draw (0,1.2) -- (5.3,1.2);
  \draw (0,0) -- (10.6,0);
  \node[anchor=east] at (0,2.4) {$\ket{\varphi}_Q$};
  \node[anchor=south west, text=insagris, font=\footnotesize] at (0.05,1.2) {$A$};
  \node[anchor=south west, text=insagris, font=\footnotesize] at (0.05,0) {$B$};
  % paire de Bell partagée
  \draw[insagris, thick] (-0.1,1.5) -- (-0.25,1.5) -- (-0.25,-0.3) -- (-0.1,-0.3);
  \node[anchor=east] at (-0.35,0.6) {$\ket{\Phi^+}_{BA}$};
  % CNOT : contrôle Q, cible A
  \draw[insarouge, thick] (2.6,2.4) -- (2.6,1.2);
  \pic at (2.6,2.4) {ctl};
  \pic at (2.6,1.2) {cible};
  % Hadamard sur Q
  \node[porte] at (4.2,2.4) {$H$};
  % mesures d'Alice
  \pic at (5.6,2.4) {mesure};
  \pic at (5.6,1.2) {mesure};
  % fils classiques (doubles)
  \draw[double, double distance=1.2pt, insagris] (5.9,1.2) -- (7.6,1.2) -- (7.6,0.35);
  \draw[double, double distance=1.2pt, insagris] (5.9,2.4) -- (9.6,2.4) -- (9.6,0.35);
  \node[above, font=\footnotesize] at (6.5,1.2) {$a$};
  \node[above, font=\footnotesize] at (6.5,2.4) {$q$};
  % corrections de Bob
  \node[porte] at (7.6,0) {$X$};
  \node[porte] at (9.6,0) {$Z$};
  \node[anchor=west] at (10.6,0) {$\ket{\varphi}_B$};
  % états intermédiaires
  \foreach \x in {1.4,3.4,4.9}
    \draw[dashed, insagris!70] (\x,2.8) -- (\x,-0.5);
  \node[anchor=north, font=\footnotesize] at (1.4,-0.5) {$\ket{\Pi_0}$};
  \node[anchor=north, font=\footnotesize] at (3.4,-0.5) {$\ket{\Pi_1}$};
  \node[anchor=north, font=\footnotesize] at (4.9,-0.5) {$\ket{\Pi_2}$};
\end{tikzpicture}
\end{adjustbox}
\caption{Circuit de téléportation. Alice détient $Q$ (dans $\ket{\varphi}=\alpha\ket{0}+\beta\ket{1}$)
et $A$, Bob détient $B$ ; $A$ et $B$ forment la paire $\ket{\Phi^+}_{BA}$. Les traits doubles
sont des fils \emph{classiques} qui portent les résultats $a$ (mesure de $A$) et $q$ (mesure de
$Q$) ; $X$ est appliqué si $a=1$, $Z$ si $q=1$. Le ket s'écrit $\ket{b,a,q}$ : le fil du haut
($Q$) est le dernier chiffre.}
\label{fig:teleportation}
\end{figure}

\subsubsection{Calcul pas à pas}

\paragraph{État initial $\ket{\Pi_0}$.} Alice et Bob partagent la paire, et $Q$ est indépendant
d'elle :
\begin{align*}
  \ket{\Pi_0}&=\ket{\Phi^+}_{BA}\otimes\ket{\varphi}_Q
    =\frac{1}{\sqrt2}\big(\ket{00}+\ket{11}\big)\otimes\big(\alpha\ket{0}+\beta\ket{1}\big)\\
  &=\frac{1}{\sqrt2}\Big(\alpha\ket{00}\ket{0}+\beta\ket{00}\ket{1}+\alpha\ket{11}\ket{0}+\beta\ket{11}\ket{1}\Big)\\
  &=\frac{\alpha\ket{000}+\beta\ket{001}+\alpha\ket{110}+\beta\ket{111}}{\sqrt2}.
\end{align*}

\paragraph{Après le CNOT $Q\to A$ : $\ket{\Pi_1}$.} Le CNOT (contrôle $Q$, dernier chiffre ;
cible $A$, chiffre du milieu) agit ainsi : $\ket{b,a,q}\to\ket{b,\,a\oplus q,\,q}$. Terme par terme :
\begin{center}
\renewcommand{\arraystretch}{1.3}
\begin{tabular}{@{}cccc@{}}
\toprule
Terme & $q$ & Effet sur $A$ & Résultat \\
\midrule
$\alpha\ket{000}$ & $0$ & aucun & $\alpha\ket{000}$ \\
$\beta\ket{001}$ & $1$ & $0\to1$ & $\beta\ket{011}$ \\
$\alpha\ket{110}$ & $0$ & aucun & $\alpha\ket{110}$ \\
$\beta\ket{111}$ & $1$ & $1\to0$ & $\beta\ket{101}$ \\
\bottomrule
\end{tabular}
\end{center}
donc
\[
  \ket{\Pi_1}=\frac{\alpha\ket{000}+\beta\ket{011}+\alpha\ket{110}+\beta\ket{101}}{\sqrt2}.
\]
On isole ensuite le dernier chiffre, celui de $Q$, qui va changer à l'étape suivante
(le premier ket est celui du couple $\ket{b,a}$) :
\[
  \ket{\Pi_1}=\frac{\alpha\ket{00}\otimes\ket{0}+\beta\ket{01}\otimes\ket{1}
    +\alpha\ket{11}\otimes\ket{0}+\beta\ket{10}\otimes\ket{1}}{\sqrt2}.
\]

\paragraph{Après la porte $H$ sur $Q$ : $\ket{\Pi_2}$.} On remplace $\ket{0}$ par
$H\ket{0}=\ket{+}$ et $\ket{1}$ par $H\ket{1}=\ket{-}$ dans le facteur $Q$ :
\[
  \ket{\Pi_2}=\frac{1}{\sqrt2}\Big(\alpha\ket{00}\otimes\ket{+}+\beta\ket{01}\otimes\ket{-}
    +\alpha\ket{11}\otimes\ket{+}+\beta\ket{10}\otimes\ket{-}\Big).
\]
Avec $\ket{\pm}=\frac{1}{\sqrt2}(\ket{0}\pm\ket{1})$, le préfacteur devient
$\frac{1}{\sqrt2}\cdot\frac{1}{\sqrt2}=\frac12$ :
\begin{align*}
  \ket{\Pi_2}&=\frac12\Big(\alpha\ket{00}\big(\ket{0}+\ket{1}\big)+\beta\ket{01}\big(\ket{0}-\ket{1}\big)
     +\alpha\ket{11}\big(\ket{0}+\ket{1}\big)+\beta\ket{10}\big(\ket{0}-\ket{1}\big)\Big)\\
  &=\frac12\Big(\alpha\ket{000}+\alpha\ket{001}+\beta\ket{010}-\beta\ket{011}
     +\alpha\ket{110}+\alpha\ket{111}+\beta\ket{100}-\beta\ket{101}\Big).
\end{align*}

\paragraph{Regroupement selon le résultat d'Alice.} Alice mesurera $A$ et $Q$, donc les deux
derniers chiffres. On regroupe les huit termes qui ont les mêmes chiffres $(a,q)$, en mettant
le ket de Bob en premier :
\begin{center}
\renewcommand{\arraystretch}{1.3}
\begin{tabular}{@{}cll@{}}
\toprule
$(a,q)$ & Termes de $\ket{\Pi_2}$ (facteur $\frac12$ omis) & Ket de Bob \\
\midrule
$(0,0)$ & $\alpha\ket{000}+\beta\ket{100}$ & $(\alpha\ket{0}+\beta\ket{1})\ket{00}$ \\
$(0,1)$ & $\alpha\ket{001}-\beta\ket{101}$ & $(\alpha\ket{0}-\beta\ket{1})\ket{01}$ \\
$(1,0)$ & $\beta\ket{010}+\alpha\ket{110}$ & $(\beta\ket{0}+\alpha\ket{1})\ket{10}$ \\
$(1,1)$ & $-\beta\ket{011}+\alpha\ket{111}$ & $(-\beta\ket{0}+\alpha\ket{1})\ket{11}$ \\
\bottomrule
\end{tabular}
\end{center}
d'où
\[
  \ket{\Pi_2}=\frac12\Big[\big(\alpha\ket{0}+\beta\ket{1}\big)\ket{00}
    +\big(\alpha\ket{0}-\beta\ket{1}\big)\ket{01}
    +\big(\beta\ket{0}+\alpha\ket{1}\big)\ket{10}
    +\big(\alpha\ket{1}-\beta\ket{0}\big)\ket{11}\Big].
\]
Dans chaque terme, le ket de gauche est l'état de $B$ et le ket de droite $\ket{a,q}$ celui du
couple $(A,Q)$ d'Alice.

\subsubsection{Mesure d'Alice et corrections de Bob}

Alice mesure $A$ et $Q$. Comme à la section~\ref{sec:mesure2}, la probabilité d'un résultat
$(a,q)$ est le carré de la norme du terme correspondant, et l'état de Bob devient le ket
associé (déjà normé). Ici
\[
  P(a,q)=\Big\|\tfrac12\ket{\phi_{aq}}\Big\|^2=\tfrac14\big(|\alpha|^2+|\beta|^2\big)=\tfrac14
  \quad\text{pour les quatre résultats,}
\]
qui sont donc équiprobables et \emph{indépendants} de $\alpha$ et $\beta$ : les deux bits
envoyés ne révèlent rien sur $\ket{\varphi}$. Les kets de Bob sont
$\ket{\phi_{00}}=\alpha\ket{0}+\beta\ket{1}$, $\ket{\phi_{01}}=\alpha\ket{0}-\beta\ket{1}$,
$\ket{\phi_{10}}=\beta\ket{0}+\alpha\ket{1}$ et $\ket{\phi_{11}}=-\beta\ket{0}+\alpha\ket{1}$.

\begin{center}
\renewcommand{\arraystretch}{1.5}
\begin{tabular}{@{}cccll@{}}
\toprule
Résultat $(a,q)$ & Probabilité & État de Bob & Lien avec $\ket{\varphi}$ & Correction de Bob \\
\midrule
$(0,0)$ & $\frac14$ & $\alpha\ket{0}+\beta\ket{1}$ & $\ket{\varphi}$ & $I$ (rien) \\
$(0,1)$ & $\frac14$ & $\alpha\ket{0}-\beta\ket{1}$ & $Z\ket{\varphi}$ & $Z$ \\
$(1,0)$ & $\frac14$ & $\beta\ket{0}+\alpha\ket{1}$ & $X\ket{\varphi}$ & $X$ \\
$(1,1)$ & $\frac14$ & $-\beta\ket{0}+\alpha\ket{1}$ & $XZ\ket{\varphi}$ & $X$ puis $Z$ (opérateur $ZX$) \\
\bottomrule
\end{tabular}
\end{center}

\paragraph{Vérification des quatre corrections.} Avec $X\ket{0}=\ket{1}$, $X\ket{1}=\ket{0}$,
$Z\ket{0}=\ket{0}$ et $Z\ket{1}=-\ket{1}$ :
\begin{align*}
  (0,0):\quad & I\big(\alpha\ket{0}+\beta\ket{1}\big)=\alpha\ket{0}+\beta\ket{1},\\
  (0,1):\quad & Z\big(\alpha\ket{0}-\beta\ket{1}\big)=\alpha\ket{0}-\beta Z\ket{1}=\alpha\ket{0}+\beta\ket{1},\\
  (1,0):\quad & X\big(\beta\ket{0}+\alpha\ket{1}\big)=\beta\ket{1}+\alpha\ket{0}=\alpha\ket{0}+\beta\ket{1},\\
  (1,1):\quad & X\big(-\beta\ket{0}+\alpha\ket{1}\big)=-\beta\ket{1}+\alpha\ket{0}=\alpha\ket{0}-\beta\ket{1},
     \ \text{puis}\ Z\big(\alpha\ket{0}-\beta\ket{1}\big)=\alpha\ket{0}+\beta\ket{1}.
\end{align*}
Dans le dernier cas, avec les matrices,
$ZX=\begin{pmatrix}1&0\\0&-1\end{pmatrix}\begin{pmatrix}0&1\\1&0\end{pmatrix}
=\begin{pmatrix}0&1\\-1&0\end{pmatrix}$ et
$ZX\begin{pmatrix}-\beta\\ \alpha\end{pmatrix}=\begin{pmatrix}\alpha\\ \beta\end{pmatrix}$.
L'ordre compte : l'état de Bob est $XZ\ket{\varphi}$ ($Z\ket{\varphi}=\alpha\ket{0}-\beta\ket{1}$,
puis $X$ donne $\alpha\ket{1}-\beta\ket{0}$), et on l'annule avec
$(XZ)^{-1}=Z^{-1}X^{-1}=ZX$ : on applique d'abord $X$, puis $Z$, comme dans le circuit.
Dans les quatre cas, Bob retrouve exactement $\ket{\varphi}=\alpha\ket{0}+\beta\ket{1}$.

\subsubsection{Ce que Bob sait avant de recevoir les bits}

Tant que Bob ignore $(a,q)$, son qubit est dans le mélange des quatre kets $\ket{\phi_{aq}}$, chacun
avec la probabilité $\frac14$ :
\begin{align*}
  \rho_B&=\frac14\sum_{a,q}\ket{\phi_{aq}}\bra{\phi_{aq}}\\
  &=\frac14\Bigg[\begin{pmatrix}|\alpha|^2&\alpha\beta^*\\ \alpha^*\beta&|\beta|^2\end{pmatrix}
    +\begin{pmatrix}|\alpha|^2&-\alpha\beta^*\\ -\alpha^*\beta&|\beta|^2\end{pmatrix}
    +\begin{pmatrix}|\beta|^2&\alpha^*\beta\\ \alpha\beta^*&|\alpha|^2\end{pmatrix}
    +\begin{pmatrix}|\beta|^2&-\alpha^*\beta\\ -\alpha\beta^*&|\alpha|^2\end{pmatrix}\Bigg]\\
  &=\frac14\begin{pmatrix}2(|\alpha|^2+|\beta|^2)&0\\0&2(|\alpha|^2+|\beta|^2)\end{pmatrix}
   =\frac12\begin{pmatrix}1&0\\0&1\end{pmatrix}=\frac12\,\mathrm{Id}.
\end{align*}
Sans les deux bits, Bob a donc un qubit complètement mélangé, qui ne dépend pas de $\alpha$ ni de
$\beta$ : il n'a reçu aucune information sur $\ket{\varphi}$. C'est l'envoi classique de $(a,q)$
qui permet de la retrouver, donc rien ne voyage plus vite que la lumière.

\begin{remarque}[Bilan de la téléportation]
\begin{itemize}[nosep]
  \item Ressources : $1$ e-bit (la paire de Bell) et $2$ bits classiques pour transmettre
        \emph{un qubit} d'état quelconque, sans qu'Alice connaisse $\alpha$ ni $\beta$.
  \item Pas de clonage : après la mesure, $A$ et $Q$ sont dans l'état de base $\ket{a,q}$.
        Alice n'a plus $\ket{\varphi}$ : l'état a été \emph{déplacé}, pas copié. L'intrication
        entre $A$ et $B$ a été consommée.
  \item Les quatre corrections $I,Z,X,ZX$ sont des portes contrôlées par des bits classiques
        (section~\ref{sec:controlees}).
\end{itemize}
\end{remarque}

\begin{remarque}[Les gestes d'Alice forment une mesure de Bell]
Le CNOT de contrôle $Q$ sur $A$ suivi de $H$ sur $Q$ est l'\emph{inverse} du générateur de la
section~\ref{sec:circuit-bell} (les portes $H$ et CNOT sont leurs propres inverses, et l'ordre
s'inverse) : il envoie la base de Bell du couple $(A,Q)$ sur la base de calcul. Mesurer ensuite dans
la base de calcul revient donc à mesurer $(A,Q)$ dans la base de Bell. Les deux bits $(a,q)$
disent \emph{quel} état de Bell Alice a trouvé, ce qui explique les quatre corrections de Bob.
\end{remarque}

\begin{remarque}[Non-clonage]
Alice ne peut pas simplement copier $\ket{\varphi}$ \parencite{wootters1982}. Supposons une porte
unitaire $U$ telle que $U\big(\ket{\psi}\otimes\ket{0}\big)=\ket{\psi}\otimes\ket{\psi}$ pour tout
$\ket{\psi}$. Elle copie $\ket{0}$ et $\ket{1}$, mais par linéarité
\[
  U\big(\ket{+}\otimes\ket{0}\big)=\tfrac{1}{\sqrt2}\big(\ket{00}+\ket{11}\big)
  \neq\ket{+}\otimes\ket{+}=\tfrac12\big(\ket{00}+\ket{01}+\ket{10}+\ket{11}\big),
\]
ce qui contredit l'hypothèse. Le CNOT copie un bit classique, pas une superposition (remarque de la
section~\ref{sec:cnot}). La téléportation respecte cette interdiction : l'état est déplacé, il n'est
jamais présent en deux endroits.
\end{remarque}

% ---------------------------------------------------------------------
\subsection{Corrélations non locales : le jeu CHSH}\label{sec:chsh}
% ---------------------------------------------------------------------

Les états de Bell donnent des résultats corrélés en base $Z$ comme en base $X$
(section~\ref{sec:bell}). Mais deux enveloppes contenant la même carte, préparées à l'avance, donnent
aussi des résultats égaux : des corrélations dans une seule base ne prouvent pas l'intrication. Bell
\parencite{bell1964} a montré qu'on peut trancher en \emph{changeant de base de mesure} ; la version de
Clauser, Horne, Shimony et Holt \parencite{clauser1969} se présente comme un jeu.

\begin{definition}{Le jeu CHSH}
Alice et Bob, éloignés, ne communiquent plus une fois le jeu commencé. Un arbitre tire deux bits $x$ et
$y$ de façon uniforme et indépendante, donne $x$ à Alice et $y$ à Bob. Chacun répond par un bit : $a$ pour
Alice, $b$ pour Bob. Ils \emph{gagnent} si
\[
  a\oplus b=x\cdot y,
\]
c'est-à-dire si $a=b$ pour $(x,y)\in\{(0,0),(0,1),(1,0)\}$ et si $a\neq b$ pour $(x,y)=(1,1)$.
\end{definition}

\paragraph{Stratégie classique : au plus $\frac34$.} Une stratégie déterministe est $a=f(x)$,
$b=g(y)$. Si elle gagnait les quatre parties, on aurait
\[
  f(0)\oplus g(0)=0,\qquad f(0)\oplus g(1)=0,\qquad f(1)\oplus g(0)=0,\qquad f(1)\oplus g(1)=1 .
\]
En additionnant ces quatre égalités modulo $2$, chacune des valeurs $f(0),f(1),g(0),g(1)$ apparaît
deux fois : le membre de gauche vaut $0$, celui de droite vaut $1$. C'est absurde, donc au moins une
partie sur quatre est perdue et $P(\text{gain})\leq\frac34$. La stratégie «~toujours répondre $0$~»
atteint $\frac34$ (elle ne perd que pour $(1,1)$). Un hasard partagé ne change rien : il mélange des
stratégies déterministes, dont aucune ne dépasse $\frac34$.

\paragraph{Stratégie quantique.} Alice et Bob partagent avant le jeu une paire $\ket{\Phi^+}$ (un
e-bit), puis chacun mesure son qubit sans rien dire à l'autre. Pour un angle $\theta$, notons
\[
  \ket{\theta_+}=\cos\tfrac\theta2\,\ket{0}+\sin\tfrac\theta2\,\ket{1},
  \qquad
  \ket{\theta_-}=-\sin\tfrac\theta2\,\ket{0}+\cos\tfrac\theta2\,\ket{1}
\]
la base de mesure d'angle $\theta$ : c'est la direction $\theta$ dans le plan $xz$ de la sphère de
Bloch ($\theta=0$ : base $Z$ ; $\theta=\frac\pi2$ : base $X$). Le résultat vaut $0$ pour
$\ket{\theta_+}$ et $1$ pour $\ket{\theta_-}$. La stratégie est
\begin{itemize}[nosep]
  \item Alice : $\theta=0$ si $x=0$ (base $Z$), $\theta=\frac\pi2$ si $x=1$ (base $X$) ;
  \item Bob : $\theta=\frac\pi4$ si $y=0$, $\theta=-\frac\pi4$ si $y=1$ (bases propres de
        $\frac{Z\pm X}{\sqrt2}$).
\end{itemize}

\begin{figure}[!ht]
\centering
\adjustbox{max width=\linewidth}{%
\begin{tikzpicture}[>=Stealth, font=\footnotesize,
    joueur/.style={draw=insarouge, fill=insarouge!6, rounded corners=2pt, thick,
                   minimum width=2.1cm, minimum height=1.1cm, align=center},
    arb/.style={draw=insagris, fill=insagris!8, rounded corners=2pt, thick,
                minimum width=3.6cm, minimum height=0.8cm, align=center}]
  % --- le jeu ---
  \node[arb] (arb) at (3.2,2.7) {arbitre : $x,y$ au hasard};
  \node[joueur] (al) at (0,0) {Alice};
  \node[joueur] (bo) at (6.4,0) {Bob};
  \draw[->, thick, insagris] (arb.west) -| node[pos=0.2, above] {$x$} (al.north);
  \draw[->, thick, insagris] (arb.east) -| node[pos=0.2, above] {$y$} (bo.north);
  \draw[insarouge, thick, decorate, decoration={snake, amplitude=1.1pt, segment length=6pt}]
    (al.east) -- (bo.west);
  \node[insarouge, above=1pt] at (3.2,0.05) {$\ket{\Phi^+}$};
  \node[text=insagris, below=2pt, font=\scriptsize] at (3.2,-0.05) {paire intriquée};
  \draw[->, thick] (al.south) -- ++(0,-1.0) node[below] {$a$};
  \draw[->, thick] (bo.south) -- ++(0,-1.0) node[below] {$b$};
  \node[draw=insagris, rounded corners=2pt, inner sep=3pt] at (3.2,-1.75) {gain si $a\oplus b=x\cdot y$};
  % --- directions de mesure dans le plan xz de la sphère de Bloch ---
  \begin{scope}[shift={(11.9,0.3)}]
    \draw[insagris!70] (0,0) circle (1.7);
    \draw[insagris!50, dashed] (-1.7,0) -- (1.7,0);
    \draw[insagris!50, dashed] (0,-1.7) -- (0,1.7);
    \node[font=\scriptsize, text=insagris, anchor=north] at (0,-1.72) {$\ket{1}$};
    \node[font=\scriptsize, text=insagris, anchor=east] at (-1.72,-0.02) {$\ket{-}$};
    \draw[->, very thick, insarouge] (0,0) -- (90:1.7);
    \draw[->, very thick, insarouge] (0,0) -- (0:1.7);
    \draw[->, very thick, insagris] (0,0) -- (45:1.7);
    \draw[->, very thick, insagris] (0,0) -- (135:1.7);
    \node[insarouge, anchor=south, inner sep=2pt] at (0,1.72) {$x=0$ : $Z$ $(\ket{0})$};
    \node[insarouge, anchor=west, inner sep=2pt] at (1.72,0) {$x=1$ : $X$ $(\ket{+})$};
    \node[text=insagris, anchor=south west, align=left, inner sep=2pt] at (1.25,1.25) {$y=0$\\ $\frac{Z+X}{\sqrt2}$};
    \node[text=insagris, anchor=south east, align=right, inner sep=2pt] at (-1.25,1.25) {$y=1$\\ $\frac{Z-X}{\sqrt2}$};
    \draw[insagris] (0.85,0) arc (0:45:0.85);
    \node[font=\scriptsize, text=insagris] at (22.5:1.12) {$\frac\pi4$};
    \draw[insagris] (0,0.85) arc (90:135:0.85);
    \node[font=\scriptsize, text=insagris] at (112.5:1.12) {$\frac\pi4$};
    \draw[insagris] (0.6,0.6) arc (45:90:0.85);
    \node[font=\scriptsize, text=insagris] at (67.5:1.12) {$\frac\pi4$};
  \end{scope}
\end{tikzpicture}}
\caption{Le jeu CHSH (à gauche) et les quatre directions de mesure de la stratégie quantique (à
droite), dans le plan $xz$ de la sphère de Bloch : Alice en rouge, Bob en gris. Pour
$(x,y)=(0,0),(0,1),(1,0)$, les directions font un angle $\frac\pi4$ ; pour $(1,1)$, un angle
$\frac{3\pi}4$.}
\label{fig:chsh}
\end{figure}

\begin{resultat}{Corrélations de $\ket{\Phi^+}$ en bases tournées}
Si Alice mesure dans la base d'angle $\alpha$ et Bob dans la base d'angle $\beta$, alors
\[
  P(a=b)=\cos^2\frac{\alpha-\beta}{2}.
\]
\end{resultat}

\noindent\emph{Preuve.} Avec $c=\cos\frac\alpha2$ et $s=\sin\frac\alpha2$,
$\ket{\alpha_+\alpha_+}+\ket{\alpha_-\alpha_-}=(c^2+s^2)\ket{00}+(cs-sc)\big(\ket{01}+\ket{10}\big)+(s^2+c^2)\ket{11}
=\ket{00}+\ket{11}$ : on peut donc écrire $\ket{\Phi^+}=\frac{1}{\sqrt2}\big(\ket{\alpha_+\alpha_+}+\ket{\alpha_-\alpha_-}\big)$.
Si Alice obtient $a=0$ (probabilité $\frac12$), le qubit de Bob est dans $\ket{\alpha_+}$ et Bob obtient $b=0$
avec la probabilité $|\braket{\beta_+}{\alpha_+}|^2$, où $\braket{\beta_+}{\alpha_+}=\cos\frac\beta2\cos\frac\alpha2+\sin\frac\beta2\sin\frac\alpha2=\cos\frac{\alpha-\beta}{2}$.
Si Alice obtient $a=1$, Bob est dans $\ket{\alpha_-}$ et obtient $b=1$ avec la même probabilité, car
$\braket{\beta_-}{\alpha_-}=\braket{\beta_+}{\alpha_+}$. Donc
$P(a=b)=\frac12\cos^2\frac{\alpha-\beta}2+\frac12\cos^2\frac{\alpha-\beta}2$. \hfill$\square$

\medskip
\noindent Pour les quatre couples $(x,y)$ :
\begin{center}
\renewcommand{\arraystretch}{1.3}
\begin{tabular}{@{}c c c c c c@{}}
\toprule
$(x,y)$ & $\alpha$ & $\beta$ & $P(a=b)$ & condition de gain & $P(\text{gain})$\\
\midrule
$(0,0)$ & $0$ & $\frac\pi4$ & $\cos^2\frac\pi8$ & $a=b$ & $\cos^2\frac\pi8$\\
$(0,1)$ & $0$ & $-\frac\pi4$ & $\cos^2\frac\pi8$ & $a=b$ & $\cos^2\frac\pi8$\\
$(1,0)$ & $\frac\pi2$ & $\frac\pi4$ & $\cos^2\frac\pi8$ & $a=b$ & $\cos^2\frac\pi8$\\
$(1,1)$ & $\frac\pi2$ & $-\frac\pi4$ & $\cos^2\frac{3\pi}8=\sin^2\frac\pi8$ & $a\neq b$ &
  $1-\sin^2\frac\pi8=\cos^2\frac\pi8$\\
\bottomrule
\end{tabular}
\end{center}
Alice et Bob gagnent donc chaque partie avec la même probabilité :
\[
  P(\text{gain})=\cos^2\frac\pi8=\frac{2+\sqrt2}{4}\approx0{,}854>\frac34 .
\]

\begin{remarque}[Ce que montre le jeu CHSH]
\begin{itemize}[nosep]
  \item Aucune stratégie classique, même avec du hasard partagé, ne dépasse $\frac34$ : l'intrication
        produit des corrélations qu'aucun modèle «~local~» (résultats fixés à l'avance) ne reproduit.
        Les expériences sur des photons intriqués confirment la valeur quantique.
  \item $\cos^2\frac\pi8$ est le maximum quantique (\emph{borne de Tsirelson}, résultat admis,
        \cite{tsirelson1980}) : on ne peut pas gagner à tous les coups. Avec
        $E_{xy}=P(a=b)-P(a\neq b)=\cos(\alpha-\beta)$ et $S=E_{00}+E_{01}+E_{10}-E_{11}$, on a
        $P(\text{gain})=\frac12+\frac S8$ ; l'inégalité de Bell donne $|S|\leq2$ en classique, et la
        stratégie ci-dessus atteint $S=2\sqrt2$.
  \item Il n'y a pas de communication : la réponse de Bob est uniforme ($P(b=0)=\frac12$) quelle que
        soit la base choisie par Alice. L'intrication ne transmet pas d'information ; elle ne fait
        que corréler les résultats.
\end{itemize}
\end{remarque}

\clearpage
% ---------------------------------------------------------------------
\subsection{Récapitulatif du chapitre}\label{sec:recap5}
% ---------------------------------------------------------------------

\begin{center}
\adjustbox{max width=\linewidth}{%
\begin{tikzpicture}[>=Stealth, font=\footnotesize,
    boite/.style={draw=insarouge, fill=insarouge!6, rounded corners=2pt, thick,
                  minimum width=4.4cm, minimum height=1.45cm, align=center, inner sep=3pt},
    lien/.style={->, very thick, insagris},
    etiq/.style={font=\scriptsize\itshape, text=insagris, inner sep=2pt}]
  \node[boite] (mes) at (0,0)
    {\textbf{Mesure à deux qubits}\\ $P(ij)=|c_{ij}|^2$\\ {\scriptsize\S\,\ref{sec:mesure2}}};
  \node[boite] (sep) at (6.3,0)
    {\textbf{Séparable ou intriqué}\\ $D=c_{00}c_{11}-c_{01}c_{10}$\\ {\scriptsize\S\,\ref{sec:separable}}};
  \node[boite] (bel) at (12.6,0)
    {\textbf{États de Bell}\\ $\ket{\Phi^\pm}$, $\ket{\Psi^\pm}$ ; $\mathrm{CNOT}\,(I\otimes H)$\\
     {\scriptsize\S\,\ref{sec:bell}, \ref{sec:circuit-bell}}};
  \node[boite] (tra) at (12.6,-2.3)
    {\textbf{Trace partielle, entropies}\\ $\rho_A=\operatorname{Tr}_B\rho_{AB}$ ; $S(\rho)$, $I(A;B)$\\
     {\scriptsize\S\,\ref{sec:trace}, \ref{sec:info}}};
  \node[boite] (tel) at (6.3,-2.3)
    {\textbf{Téléportation}\\ e-bit + 2 bits classiques\\ {\scriptsize\S\,\ref{sec:teleportation}}};
  \node[boite] (chs) at (0,-2.3)
    {\textbf{Jeu CHSH}\\ $\frac34$ classique, $\cos^2\frac\pi8$ quantique\\ {\scriptsize\S\,\ref{sec:chsh}}};
  \draw[lien] (mes) -- (sep) node[etiq, midway, above=1pt] {corrélations};
  \draw[lien] (sep) -- (bel) node[etiq, midway, above=1pt] {cas extrême};
  \draw[lien] (bel) -- (tra) node[etiq, midway, right=1pt] {un qubit seul};
  \draw[lien] (tra) -- (tel) node[etiq, midway, above=1pt] {ressource};
  \draw[lien] (tel) -- (chs) node[etiq, midway, above=1pt] {non-localité};
\end{tikzpicture}}
\end{center}

\begin{aretenir}
\begin{itemize}
  \item Deux qubits : quatre amplitudes $c_{ij}$, $P(ij)=|c_{ij}|^2$. Mesurer $A$ supprime les termes
        incompatibles et renormalise : l'état de $B$ peut changer.
  \item Séparable $\Longleftrightarrow D=c_{00}c_{11}-c_{01}c_{10}=0$. Sinon l'état est intriqué et
        aucun qubit n'a d'état propre (pour $\ket{\Phi^+}$, $P(00)\neq P(A{=}0)P(B{=}0)$).
  \item États de Bell : base orthonormée d'états maximalement intriqués, états propres de
        $Z\otimes Z$ et $X\otimes X$ ; $\mathrm{CNOT}\,(I\otimes H)$ les produit à partir de la base de
        calcul. Des portes locales ne créent pas d'intrication.
  \item Qubit seul : $\rho_A=\operatorname{Tr}_B\rho_{AB}$, avec $\operatorname{Tr}\rho_A^2=1-2|D|^2$ pour un état pur.
        Pour $\ket{\Phi^+}$ : $\rho_A=\mathrm{Id}/2$ alors que le couple est pur ; $S(\rho_{AB})=0$,
        $S(\rho_A)=1$ bit, $I(A;B)=2$ bits, $H(A|B)=-1$ (impossible en classique).
  \item Téléportation : un e-bit et deux bits classiques ; $P(a,q)=\frac14$ ; Bob applique $X$ si $a=1$
        puis $Z$ si $q=1$. L'état est déplacé, pas copié ; sans les bits, $\rho_B=\mathrm{Id}/2$.
  \item Jeu CHSH : toute stratégie classique gagne avec une probabilité $\leq\frac34$ ; avec $\ket{\Phi^+}$
        et des bases tournées, $\cos^2\frac\pi8\approx0{,}854$ (borne de Tsirelson), sans aucune
        communication.
\end{itemize}
\end{aretenir}

\begin{center}
\renewcommand{\arraystretch}{1.2}
\begin{tabular}{@{}l l r@{}}
\toprule
\textbf{Notion} & \textbf{Formule} & \textbf{§} \\
\midrule
Mesure de deux qubits & $P(ij)=|c_{ij}|^2$, \quad $P(A{=}a)=P(a0)+P(a1)$ & \ref{sec:mesure2} \\
Séparabilité & $\ket{\Psi}=\ket{\Psi}_A\otimes\ket{\Psi}_B\ \Longleftrightarrow\ D=0$ & \ref{sec:separable} \\
États de Bell & $\ket{\Phi^\pm}=\dfrac{\ket{00}\pm\ket{11}}{\sqrt2}$, \quad
  $\ket{\Psi^\pm}=\dfrac{\ket{01}\pm\ket{10}}{\sqrt2}$ & \ref{sec:bell} \\
Trace partielle & $\operatorname{Tr}_B\big(\ket{a}\bra{a'}\otimes\ket{b}\bra{b'}\big)
  =\ket{a}\bra{a'}\,\braket{b'}{b}$ & \ref{sec:trace} \\
Qubit d'un état pur & $\rho_A=CC^\dagger$, \quad $\operatorname{Tr}\rho_A^2=1-2|D|^2$ & \ref{sec:trace} \\
Entropies & $S(\rho)=-\sum_i\lambda_i\log_2\lambda_i$, \quad $I(A;B)=H(A)+H(B)-H(A,B)$ & \ref{sec:info} \\
Téléportation & $P(a,q)=\frac14$, \quad correction $Z^{q}X^{a}$ (d'abord $X$) & \ref{sec:teleportation} \\
Singulet & $(U\otimes U)\ket{\Psi^-}=\det(U)\,\ket{\Psi^-}$ & \ref{sec:bell} \\
Jeu CHSH & $P_{\mathrm{cl}}\leq\frac34$, \quad $P_{\mathrm{quant}}=\cos^2\frac\pi8=\frac{2+\sqrt2}{4}$ & \ref{sec:chsh} \\
\bottomrule
\end{tabular}
\end{center}

\begin{savoirfaire}
\begin{itemize}
  \item Calculer $P(ij)$, les probabilités d'un seul qubit et l'état après la mesure de l'un d'eux ;
        tester la séparabilité avec $D$ et calculer la pureté de $\rho_A$ (exercice~\ref{ex:separable3}).
  \item Suivre un état dans le générateur de Bell, écrire $\rho_{AB}$ et reconnaître les corrélations
        d'un état de Bell dans deux bases (exercice~\ref{ex:bellX}).
  \item Calculer $H(A)$, $H(A,B)$, $I(A;B)$ et $H(A|B)$ pour une loi ou un état (exercice~\ref{ex:entropies}).
  \item Dérouler la téléportation pour un état donné et utiliser les états de Bell pour le codage
        superdense (exercices~\ref{ex:telenum} et~\ref{ex:superdense}).
  \item Prouver la borne $\frac34$ du jeu CHSH classique et calculer la probabilité de gain d'une
        stratégie quantique (section~\ref{sec:chsh}).
\end{itemize}
\end{savoirfaire}
.
%  Il ne contient ni préambule ni \begin{document} : les environnements
%  (definition, resultat, remarque, aretenir, savoirfaire, exercice, correction),
%  les macros (\ii, \ee, \dd, \pd, \pdd, \Hop, \ket, \bra, \braket), les symboles
%  de circuits (ctl, cible, croix, mesure, porte) et les couleurs (insarouge,
%  insagris) sont définis dans main.tex. Un chapitre = une \section.
% =====================================================================

% Macros de Dirac (déjà définies dans main.tex ; ce bloc évite l'erreur
% « Undefined control sequence \ket » si main.tex est une ancienne version).
\providecommand{\ket}[1]{|#1\rangle}
\providecommand{\bra}[1]{\langle #1|}
\providecommand{\braket}[2]{\langle #1|#2\rangle}

% =====================================================================
\section{États quantiques, intrication et états de Bell}\label{chap:bell}
% =====================================================================

Les états de deux qubits forment un espace de dimension~$4$, et la plupart d'entre eux ne se
décomposent pas en «~un état pour $A$, un état pour $B$~» : ils sont \emph{intriqués}. Ce chapitre
apprend à mesurer un couple de qubits, à reconnaître l'intrication (critère de séparabilité), à la
produire avec un circuit (états de Bell), à décrire un qubit pris seul (trace partielle), à
quantifier l'information (entropies) et à s'en servir pour téléporter un état.

% ---------------------------------------------------------------------
\subsection{États de deux qubits et mesure}\label{sec:mesure2}
% ---------------------------------------------------------------------

Un état quelconque de $\mathcal{H}^4$ se décompose sur la base de calcul :
\[
  \ket{\Psi}=c_{00}\ket{00}+c_{01}\ket{01}+c_{10}\ket{10}+c_{11}\ket{11}
  =\begin{pmatrix}c_{00}\\ c_{01}\\ c_{10}\\ c_{11}\end{pmatrix},
  \qquad \sum_{i,j}|c_{ij}|^2=1 .
\]

\begin{resultat}{Mesure des deux qubits dans la base de calcul}
\[
  P(ij)=|\braket{ij}{\Psi}|^2=|c_{ij}|^2 .
\]
Pour un seul qubit : $P(A{=}0)=P(00)+P(01)$, $P(A{=}1)=P(10)+P(11)$,
$P(B{=}0)=P(00)+P(10)$, $P(B{=}1)=P(01)+P(11)$.
\end{resultat}

En effet $\braket{ij}{\Psi}=\sum_{k,l}c_{kl}\braket{ij}{kl}=\sum_{k,l}c_{kl}\delta_{ik}\delta_{jl}=c_{ij}$.

\paragraph{Exemple 1.} $\ket{\Psi}=\frac{1}{\sqrt2}\ket{00}+\frac{1}{\sqrt2}\ket{01}
=\big(\tfrac{1}{\sqrt2},\tfrac{1}{\sqrt2},0,0\big)^{\mathsf T}$. Il est normalisé :
$\frac12+\frac12+0+0=1$. Les probabilités jointes sont
$P(00)=P(01)=\frac12$ et $P(10)=P(11)=0$. Pour chaque qubit :
\[
  P(A{=}0)=\tfrac12+\tfrac12=1,\quad P(A{=}1)=0+0=0,\quad
  P(B{=}0)=\tfrac12+0=\tfrac12,\quad P(B{=}1)=\tfrac12+0=\tfrac12 .
\]
$A$ donne toujours $0$ ; $B$ donne $0$ ou $1$ au hasard.

\paragraph{Exemple 2 (mesure de $A$ seul).} Soit
$\ket{\Psi}=\frac12\ket{00}+\frac12\ket{01}+\frac{1}{\sqrt2}\ket{11}$. Normalisation :
$\frac14+\frac14+0+\frac12=1$. Probabilités jointes : $P(00)=\frac14$, $P(01)=\frac14$,
$P(10)=0$, $P(11)=\frac12$. Donc
\[
  P(A{=}0)=\tfrac14+\tfrac14=\tfrac12,\quad P(A{=}1)=0+\tfrac12=\tfrac12,\quad
  P(B{=}0)=\tfrac14+0=\tfrac14,\quad P(B{=}1)=\tfrac14+\tfrac12=\tfrac34 .
\]

\begin{remarque}[Effondrement après la mesure de $A$]
Si $A$ donne le résultat $a$, on ne garde que les termes de $\ket{\Psi}$ dont le premier
chiffre vaut $a$, puis on renormalise :
\[
  \ket{\Psi}\ \to\ \frac{c_{a0}\ket{a0}+c_{a1}\ket{a1}}{\sqrt{|c_{a0}|^2+|c_{a1}|^2}},
  \qquad\text{probabilité }|c_{a0}|^2+|c_{a1}|^2 .
\]
\end{remarque}

Pour l'exemple 2 :
\begin{itemize}[nosep]
  \item $A=0$ (probabilité $\frac12$) : on garde $\frac12\ket{00}+\frac12\ket{01}$, de norme
        au carré $\frac12$ ; renormalisé, $\frac{1}{\sqrt2}\big(\ket{00}+\ket{01}\big)
        =\ket{0}\otimes\ket{+}$. Ensuite $B$ vaut $0$ ou $1$ avec la probabilité $\frac12$.
  \item $A=1$ (probabilité $\frac12$) : on garde $\frac{1}{\sqrt2}\ket{11}$, de norme au carré
        $\frac12$ ; renormalisé, $\ket{11}=\ket{1}\otimes\ket{1}$. Ensuite $B$ vaut $1$
        avec certitude.
\end{itemize}
Contrôle : $P(B{=}0)=\frac12\cdot\frac12+\frac12\cdot0=\frac14$, comme calculé plus haut.
Ici la mesure de $A$ \emph{modifie} l'état de $B$ ($\ket{+}$ ou $\ket{1}$) : c'est la
signature d'un état intriqué (section~\ref{sec:separable}). La figure~\ref{fig:mesure2} résume
les deux mesures successives.

\begin{figure}[!ht]
\centering
\adjustbox{max width=\linewidth}{%
\begin{tikzpicture}[>=Stealth, font=\small,
    noeud/.style={draw=insarouge, fill=insarouge!6, rounded corners=2pt, thick, inner sep=4pt, align=center},
    fin/.style={draw=insagris, fill=insagris!10, rounded corners=2pt, inner sep=4pt, align=center},
    lien/.style={->, thick, insagris},
    proba/.style={font=\footnotesize, fill=white, inner sep=1.5pt, text=black}]
  \node[noeud] (psi) at (0,0) {$\ket{\Psi}$};
  \node[noeud] (a0) at (4.4,1.6) {$A=0$\\ {\footnotesize état $\ket{0}\otimes\ket{+}$}};
  \node[noeud] (a1) at (4.4,-1.6) {$A=1$\\ {\footnotesize état $\ket{1}\otimes\ket{1}$}};
  \node[fin] (b00) at (9.2,2.5) {$B=0$ : $\ket{00}$};
  \node[fin] (b01) at (9.2,0.7) {$B=1$ : $\ket{01}$};
  \node[fin] (b11) at (9.2,-1.6) {$B=1$ : $\ket{11}$};
  \draw[lien] (psi) -- (a0) node[proba, midway] {$\frac12$};
  \draw[lien] (psi) -- (a1) node[proba, midway] {$\frac12$};
  \draw[lien] (a0) -- (b00) node[proba, midway] {$\frac12$};
  \draw[lien] (a0) -- (b01) node[proba, midway] {$\frac12$};
  \draw[lien] (a1) -- (b11) node[proba, midway] {$1$};
  \node[anchor=west] at (11.0,2.5) {$P(00)=\frac12\cdot\frac12=\frac14$};
  \node[anchor=west] at (11.0,0.7) {$P(01)=\frac12\cdot\frac12=\frac14$};
  \node[anchor=west] at (11.0,-1.6) {$P(11)=\frac12\cdot1=\frac12$};
  \node[text=insagris, font=\footnotesize\itshape] at (2.2,3.0) {mesure de $A$};
  \node[text=insagris, font=\footnotesize\itshape] at (6.8,3.6) {mesure de $B$};
  \node[text=insagris, font=\footnotesize\itshape, anchor=west] at (11.0,3.6) {probabilité jointe};
\end{tikzpicture}}
\caption{Mesure de $A$ puis de $B$ pour l'état de l'exemple~2,
$\ket{\Psi}=\frac12\ket{00}+\frac12\ket{01}+\frac{1}{\sqrt2}\ket{11}$. Chaque flèche porte une
probabilité, et le produit des probabilités le long d'un chemin est la probabilité jointe. Après la
mesure de $A$, l'état de $B$ dépend du résultat ($\ket{+}$ ou $\ket{1}$).}
\label{fig:mesure2}
\end{figure}

% ---------------------------------------------------------------------
\subsection{États séparables et états intriqués}\label{sec:separable}
% ---------------------------------------------------------------------

\begin{definition}{Séparable ou intriqué}
Un état $\ket{\Psi}_{AB}$ est \emph{séparable} s'il s'écrit comme un produit
$\ket{\Psi}_{AB}=\ket{\Psi}_A\otimes\ket{\Psi}_B$. Sinon il est \emph{intriqué} : on ne
peut pas attribuer d'état propre à chaque qubit.
\end{definition}

\paragraph{Exemple séparable (exemple 1 ci-dessus).} Par linéarité du produit tensoriel :
\begin{align*}
  \ket{\Psi}&=\tfrac{1}{\sqrt2}\ket{00}_{AB}+\tfrac{1}{\sqrt2}\ket{01}_{AB}
    =\tfrac{1}{\sqrt2}\Big(\ket{0}_A\otimes\ket{0}_B+\ket{0}_A\otimes\ket{1}_B\Big)\\
  &=\ket{0}_A\otimes\Big(\tfrac{1}{\sqrt2}\big(\ket{0}+\ket{1}\big)\Big)_B
    =\ket{0}_A\otimes\ket{+}_B ,
\end{align*}
donc $\ket{\Psi}_A=\ket{0}$ et $\ket{\Psi}_B=\ket{+}$. Contrôle par les composantes :
$\begin{pmatrix}1\\ 0\end{pmatrix}\otimes\tfrac{1}{\sqrt2}\begin{pmatrix}1\\ 1\end{pmatrix}
=\tfrac{1}{\sqrt2}\begin{pmatrix}1\cdot1\\ 1\cdot1\\ 0\cdot1\\ 0\cdot1\end{pmatrix}
=\tfrac{1}{\sqrt2}\begin{pmatrix}1\\ 1\\ 0\\ 0\end{pmatrix}$.

\subsubsection{Un critère de séparabilité}

Un produit quelconque a pour coefficients $c_{00}=ac$, $c_{01}=ad$, $c_{10}=bc$,
$c_{11}=bd$ (formule de la section~\ref{sec:tenseur}, avec $a,b$ pour $A$ et $c,d$ pour $B$), donc
\[
  c_{00}c_{11}=(ac)(bd)=abcd=(ad)(bc)=c_{01}c_{10}.
\]

\begin{resultat}{Test de séparabilité (deux qubits)}
$\ket{\Psi}=\sum c_{ij}\ket{ij}$ est séparable si et seulement si
\[
  D=c_{00}c_{11}-c_{01}c_{10}=0 .
\]
\end{resultat}

\emph{Démonstration.} ($\Rightarrow$) Calcul ci-dessus. ($\Leftarrow$) On regroupe les
termes selon l'état de $A$ :
\[
  \ket{\Psi}=\ket{0}\otimes\big(c_{00}\ket{0}+c_{01}\ket{1}\big)+\ket{1}\otimes\big(c_{10}\ket{0}+c_{11}\ket{1}\big).
\]
Si $D=0$, les lignes $(c_{00},c_{01})$ et $(c_{10},c_{11})$ sont proportionnelles. Si
$(c_{00},c_{01})\neq(0,0)$, il existe $\lambda$ avec $(c_{10},c_{11})=\lambda(c_{00},c_{01})$
et $\ket{\Psi}=\big(\ket{0}+\lambda\ket{1}\big)\otimes\big(c_{00}\ket{0}+c_{01}\ket{1}\big)$.
Si $(c_{00},c_{01})=(0,0)$, alors $\ket{\Psi}=\ket{1}\otimes\big(c_{10}\ket{0}+c_{11}\ket{1}\big)$.
Dans les deux cas $\ket{\Psi}$ est un produit ; on renormalise chaque facteur.
\hfill$\square$

\subsubsection{Application à plusieurs états}

\begin{center}
\renewcommand{\arraystretch}{1.5}
\begin{tabular}{@{}lccl@{}}
\toprule
\textbf{État} & $(c_{00},c_{01},c_{10},c_{11})$ & $D$ & \textbf{Nature} \\
\midrule
$\ket{00}$ & $(1,0,0,0)$ & $1\cdot0-0\cdot0=0$ & séparable \\
$\ket{+}\otimes\ket{+}$ & $\frac12(1,1,1,1)$ & $\frac14-\frac14=0$ & séparable \\
$\ket{+}\otimes\ket{-}$ & $\frac12(1,-1,1,-1)$ & $-\frac14-(-\frac14)=0$ & séparable \\
$\frac{1}{\sqrt2}(\ket{00}+\ket{01})$ & $\frac{1}{\sqrt2}(1,1,0,0)$ & $0-\frac12\cdot0=0$ & séparable \\
$\ket{\Phi^+}=\frac{1}{\sqrt2}(\ket{00}+\ket{11})$ & $\frac{1}{\sqrt2}(1,0,0,1)$ & $\frac12-0=\frac12$ & intriqué \\
$\ket{\Phi^-}=\frac{1}{\sqrt2}(\ket{00}-\ket{11})$ & $\frac{1}{\sqrt2}(1,0,0,-1)$ & $-\frac12-0=-\frac12$ & intriqué \\
$\ket{\Psi^+}=\frac{1}{\sqrt2}(\ket{01}+\ket{10})$ & $\frac{1}{\sqrt2}(0,1,1,0)$ & $0-\frac12=-\frac12$ & intriqué \\
$\ket{\Psi^-}=\frac{1}{\sqrt2}(\ket{01}-\ket{10})$ & $\frac{1}{\sqrt2}(0,1,-1,0)$ & $0-(-\frac12)=\frac12$ & intriqué \\
$\frac12\ket{00}+\frac12\ket{01}+\frac{1}{\sqrt2}\ket{11}$ & $(\frac12,\frac12,0,\frac{1}{\sqrt2})$ & $\frac{1}{2\sqrt2}-0$ & intriqué \\
\bottomrule
\end{tabular}
\end{center}

Les quatre états de Bell $\ket{\Phi^\pm}$ et $\ket{\Psi^\pm}$ de la table sont tous intriqués
($D=\pm\frac12$) ; leurs propriétés font l'objet de la section~\ref{sec:bell}.

\paragraph{Preuve directe que $\ket{\Phi^+}$ est intriqué.} Supposons
$\ket{\Phi^+}=(a\ket{0}+b\ket{1})\otimes(c\ket{0}+d\ket{1})$. En identifiant les
composantes : $ac=\frac{1}{\sqrt2}$, $ad=0$, $bc=0$, $bd=\frac{1}{\sqrt2}$. La condition
$ad=0$ impose $a=0$ (impossible car $ac\neq0$) ou $d=0$ (impossible car $bd\neq0$).
Contradiction : aucun produit ne convient.

\paragraph{Corrélations.} Pour un produit $\ket{\Psi}_A\otimes\ket{\Psi}_B$ avec
$\ket{\Psi}_A=a_0\ket{0}+a_1\ket{1}$ et $\ket{\Psi}_B=b_0\ket{0}+b_1\ket{1}$, on a
$P(ij)=|a_ib_j|^2=|a_i|^2|b_j|^2=P(A{=}i)\,P(B{=}j)$ : les deux qubits sont
\emph{indépendants}. Vérification avec $\ket{+}\otimes\ket{+}$ : $P(00)=\frac14=\frac12\cdot\frac12$.
Pour $\ket{\Phi^+}$ : $P(00)=P(11)=\frac12$ et $P(01)=P(10)=0$, donc
$P(A{=}0)=P(B{=}0)=\frac12$, mais
\[
  P(00)=\tfrac12\neq\tfrac12\cdot\tfrac12=P(A{=}0)\,P(B{=}0) .
\]
Les deux qubits donnent toujours le même résultat, aléatoire : ils sont corrélés.
La figure~\ref{fig:epr} illustre la différence entre les deux cas.

\begin{figure}[!ht]
\centering
\adjustbox{max width=\linewidth}{%
\begin{tikzpicture}[>=Stealth, font=\small,
    labo/.style={draw=insagris, fill=insagris!8, rounded corners=2pt, thick, minimum width=2.6cm,
                 minimum height=1.0cm, align=center},
    source/.style={draw=insarouge, fill=insarouge!10, rounded corners=2pt, thick, minimum width=3.6cm,
                   minimum height=1.0cm, align=center},
    droit/.style={->, thick, insagris},
    onde/.style={->, thick, insarouge, decorate,
                 decoration={snake, amplitude=1.3pt, segment length=7pt, post length=2mm}},
    res/.style={font=\small, align=center}]
  % (a) état séparable
  \node[labo] (Aa) at (0,0) {Alice\\ mesure $A$};
  \node[source] (Sa) at (6,0) {source\\ $\ket{0}_A\otimes\ket{+}_B$};
  \node[labo] (Ba) at (12,0) {Bob\\ mesure $B$};
  \draw[droit] (Sa.west) -- (Aa.east);
  \draw[droit] (Sa.east) -- (Ba.west);
  \node[res, below=3pt of Aa] {$0\ \ 0\ \ 0\ \ 0\ \ 0\ \ 0$};
  \node[res, below=3pt of Ba] {$1\ \ 0\ \ 0\ \ 1\ \ 0\ \ 1$};
  \node[res, text=insagris, font=\footnotesize] at (6,-1.0) {$A$ donne toujours $0$ ; $B$ est aléatoire.\\
        Les deux résultats sont indépendants.};
  \node[anchor=south west, font=\small\bfseries] at (-1.3,0.75) {(a) état séparable};
  % (b) état intriqué
  \begin{scope}[yshift=-3.6cm]
    \node[labo] (Ab) at (0,0) {Alice\\ mesure $A$};
    \node[source] (Sb) at (6,0) {source\\ $\ket{\Phi^+}_{AB}$};
    \node[labo] (Bb) at (12,0) {Bob\\ mesure $B$};
    \draw[onde] (Sb.west) -- (Ab.east);
    \draw[onde] (Sb.east) -- (Bb.west);
    \node[res, below=3pt of Ab] {$1\ \ 0\ \ 0\ \ 1\ \ 1\ \ 0$};
    \node[res, below=3pt of Bb] {$1\ \ 0\ \ 0\ \ 1\ \ 1\ \ 0$};
    \node[res, text=insagris, font=\footnotesize] at (6,-1.0) {$A$ et $B$ sont chacun aléatoires,\\
          mais toujours égaux : ils sont corrélés.};
    \node[anchor=south west, font=\small\bfseries] at (-1.3,0.75) {(b) état intriqué};
  \end{scope}
\end{tikzpicture}}
\caption{Six mesures successives sur des paires préparées dans le même état (les résultats sont
des exemples). (a) Pour un produit, chaque qubit a son propre comportement. (b) Pour
$\ket{\Phi^+}=\frac{1}{\sqrt2}(\ket{00}+\ket{11})$, chaque résultat pris seul est un tirage à pile ou
face, mais les deux résultats coïncident toujours ; la ligne ondulée symbolise l'intrication.}
\label{fig:epr}
\end{figure}

\paragraph{Et un qubit pris seul ?} Dans un état séparable, chaque qubit garde son propre état pur ;
dans un état intriqué, aucun ket ne décrit un qubit isolé. On le décrit alors par une matrice de
densité, obtenue par la trace partielle (section~\ref{sec:trace}).

% ---------------------------------------------------------------------
\subsection{Les états de Bell}\label{sec:bell}
% ---------------------------------------------------------------------

\begin{definition}{Les quatre états de Bell}
\[
  \ket{\Phi^{\pm}}=\frac{1}{\sqrt2}\big(\ket{00}\pm\ket{11}\big),
  \qquad
  \ket{\Psi^{\pm}}=\frac{1}{\sqrt2}\big(\ket{01}\pm\ket{10}\big).
\]
Ils sont intriqués ($|D|=\frac12$, section~\ref{sec:separable}) et le sont au maximum :
$|D|\leq\frac12$ pour tout état normalisé de deux qubits, et chaque qubit d'un état de Bell,
pris seul, est complètement aléatoire (section~\ref{sec:trace}).
\end{definition}

\paragraph{Une base de $\mathcal{H}^4$.} Chaque état est de norme $\frac12+\frac12=1$. Les
$\ket{\Phi^\pm}$ n'ont de composantes que sur $\ket{00}$ et $\ket{11}$, les $\ket{\Psi^\pm}$ que sur
$\ket{01}$ et $\ket{10}$, donc $\braket{\Phi^\pm}{\Psi^\pm}=0$ pour toutes les combinaisons de signes ;
enfin $\braket{\Phi^+}{\Phi^-}=\frac12(1-1)=0$ et $\braket{\Psi^+}{\Psi^-}=\frac12(1-1)=0$. Les quatre
états forment donc une base orthonormée de $\mathcal{H}^4$, la \emph{base de Bell}, constituée
uniquement d'états intriqués. Le circuit de la section~\ref{sec:circuit-bell} fait passer de la
base de calcul à cette base.

\paragraph{Ce qui les distingue : deux parités.} Sur la base de calcul,
$Z\otimes Z\ket{ij}=(-1)^{i+j}\ket{ij}$ (valeur propre $+1$ si les deux chiffres sont égaux, $-1$
sinon) et $X\otimes X\ket{ij}=\ket{\bar i\bar j}$ (les deux chiffres sont inversés, avec
$\bar i=1-i$). On en déduit, par exemple,
\[
  X\otimes X\ket{\Phi^-}=\tfrac{1}{\sqrt2}\big(\ket{11}-\ket{00}\big)=-\ket{\Phi^-},
  \qquad
  Z\otimes Z\ket{\Psi^+}=\tfrac{1}{\sqrt2}\big(-\ket{01}-\ket{10}\big)=-\ket{\Psi^+},
\]
et le tableau complet des valeurs propres :
\begin{center}
\renewcommand{\arraystretch}{1.3}
\begin{tabular}{@{}lcccc@{}}
\toprule
 & $\ket{\Phi^+}$ & $\ket{\Phi^-}$ & $\ket{\Psi^+}$ & $\ket{\Psi^-}$ \\
\midrule
$Z\otimes Z$ & $+1$ & $+1$ & $-1$ & $-1$ \\
$X\otimes X$ & $+1$ & $-1$ & $+1$ & $-1$ \\
\bottomrule
\end{tabular}
\end{center}
Les deux opérateurs commutent, car $ZX=-XZ$ donne
$(Z\otimes Z)(X\otimes X)=ZX\otimes ZX=XZ\otimes XZ=(X\otimes X)(Z\otimes Z)$ ; les quatre états de
Bell sont leurs états propres communs, et chaque couple de valeurs propres apparaît une fois
(figure~\ref{fig:grille-bell}). Le même calcul montre que $X\otimes Z$ et $Z\otimes X$
commutent : $(X\otimes Z)(Z\otimes X)=XZ\otimes ZX=(-ZX)\otimes(-XZ)=ZX\otimes XZ=(Z\otimes X)(X\otimes Z)$.
Plus généralement, $P\otimes Q$ et $P'\otimes Q'$ (matrices de Pauli) commutent si $P,P'$
anticommutent et si $Q,Q'$ anticommutent : les deux signes $-1$ se compensent.

\begin{figure}[!ht]
\centering
\adjustbox{max width=\linewidth}{%
\begin{tikzpicture}[>=Stealth, font=\small,
    cell/.style={draw=insarouge, fill=insarouge!6, rounded corners=2pt, thick, minimum width=5.4cm,
                 minimum height=1.35cm, align=center},
    tete/.style={text=insagris, align=center}]
  \node[tete] at (3.0,2.0) {$X\otimes X=+1$\\ {\footnotesize mêmes résultats en base $X$}};
  \node[tete] at (8.8,2.0) {$X\otimes X=-1$\\ {\footnotesize résultats opposés en base $X$}};
  \node[tete, anchor=east] at (-0.1,0.75) {$Z\otimes Z=+1$\\ {\footnotesize mêmes résultats}\\ {\footnotesize en base $Z$}};
  \node[tete, anchor=east] at (-0.1,-0.85) {$Z\otimes Z=-1$\\ {\footnotesize résultats opposés}\\ {\footnotesize en base $Z$}};
  \node[cell] at (3.0,0.75) {$\ket{\Phi^+}=\frac{1}{\sqrt2}\big(\ket{00}+\ket{11}\big)$};
  \node[cell] at (8.8,0.75) {$\ket{\Phi^-}=\frac{1}{\sqrt2}\big(\ket{00}-\ket{11}\big)$};
  \node[cell] at (3.0,-0.85) {$\ket{\Psi^+}=\frac{1}{\sqrt2}\big(\ket{01}+\ket{10}\big)$};
  \node[cell] at (8.8,-0.85) {$\ket{\Psi^-}=\frac{1}{\sqrt2}\big(\ket{01}-\ket{10}\big)$};
\end{tikzpicture}}
\caption{Les quatre états de Bell, classés par les valeurs propres de $Z\otimes Z$ (lignes) et de
$X\otimes X$ (colonnes). La valeur propre est le produit des deux résultats $\pm1$ d'une mesure des
deux qubits dans la base correspondante : $+1$ si les résultats sont égaux, $-1$ s'ils sont opposés.}
\label{fig:grille-bell}
\end{figure}

Les corrélations ne se limitent donc pas à la base $Z$ : $\ket{\Phi^+}$ s'écrit aussi
$\frac{1}{\sqrt2}\big(\ket{++}+\ket{--}\big)$ et donne des résultats égaux en base $X$ comme en base $Z$
(exercice~\ref{ex:bellX}).

\paragraph{D'un état de Bell à l'autre.} Une porte agissant sur un seul qubit suffit. Avec
$X\ket{0}=\ket{1}$, $X\ket{1}=\ket{0}$ et $Z\ket{1}=-\ket{1}$, en agissant sur le second qubit :
\begin{align*}
  (I\otimes X)\ket{\Phi^+}&=\tfrac{1}{\sqrt2}\big(\ket{01}+\ket{10}\big)=\ket{\Psi^+},\\
  (I\otimes Z)\ket{\Phi^+}&=\tfrac{1}{\sqrt2}\big(\ket{00}-\ket{11}\big)=\ket{\Phi^-},\\
  (I\otimes XZ)\ket{\Phi^+}&=(I\otimes X)\tfrac{1}{\sqrt2}\big(\ket{00}-\ket{11}\big)
    =\tfrac{1}{\sqrt2}\big(\ket{01}-\ket{10}\big)=\ket{\Psi^-}.
\end{align*}
Quatre opérations locales ($I$, $X$, $Z$, $XZ$) sur \emph{un} qubit d'une paire de Bell mènent donc
aux quatre états de la base ; c'est le principe du codage superdense
(exercice~\ref{ex:superdense}).

\begin{remarque}[Le singulet est le même dans toutes les bases]\label{rem:singulet}
Pour toute matrice unitaire $U=\begin{pmatrix}a&b\\c&d\end{pmatrix}$, on a $U\ket{0}=a\ket{0}+c\ket{1}$
et $U\ket{1}=b\ket{0}+d\ket{1}$, donc
\begin{align*}
  (U\otimes U)\big(\ket{01}-\ket{10}\big)
  &=\big(a\ket{0}+c\ket{1}\big)\big(b\ket{0}+d\ket{1}\big)-\big(b\ket{0}+d\ket{1}\big)\big(a\ket{0}+c\ket{1}\big)\\
  &=(ad-bc)\big(\ket{01}-\ket{10}\big),
\end{align*}
c'est-à-dire $(U\otimes U)\ket{\Psi^-}=\det(U)\,\ket{\Psi^-}$. Comme $|\det U|=1$, ce facteur est une
phase globale, sans effet physique : $\ket{\Psi^-}$ s'écrit \emph{de la même façon} dans toutes les
bases orthonormées (avec $U=H$, de déterminant $-1$ : $\ket{\Psi^-}=\frac{1}{\sqrt2}\big(\ket{-+}-\ket{+-}\big)$).
Si Alice et Bob mesurent dans la \emph{même} base, quelle qu'elle soit, leurs résultats sont
toujours opposés : c'est le «~singulet~». Parmi les états de Bell, seul $\ket{\Psi^-}$ a cette propriété.
\end{remarque}

% ---------------------------------------------------------------------
\subsection{Un circuit quantique : le générateur d'états de Bell}\label{sec:circuit-bell}
% ---------------------------------------------------------------------

On assemble une porte de Hadamard et un CNOT. Le système est $(A,B)$ avec le ket
$\ket{a,b}=\ket{a}_A\otimes\ket{b}_B$ ; d'après la convention de la section~\ref{sec:circuits}, $B$ est dessiné
en haut et $A$ en bas (figure~\ref{fig:circuit-bell}). Le circuit a deux étapes :
\begin{enumerate}[nosep]
  \item la porte $I\otimes H$ : identité sur $A$, Hadamard sur $B$ ;
  \item un CNOT de contrôle $B$ et de cible $A$ (celui de la section~\ref{sec:cnot}, dont le
        contrôle est le second chiffre).
\end{enumerate}
Les traits pointillés délimitent les états $\ket{\Pi_0}$, $\ket{\Pi_1}$ et $\ket{\Pi_2}$ :
avant les portes, après $I\otimes H$, après le CNOT. Le circuit transforme la base de calcul
en base de Bell (section~\ref{sec:bell}).

\begin{figure}[!ht]
\centering
\begin{tikzpicture}[>=Stealth, font=\small]
  % fils : B en haut, A en bas
  \draw (0,1.4) -- (7.6,1.4);
  \draw (0,0) -- (7.6,0);
  \node[anchor=east] at (0,1.4) {$\ket{0}_B$};
  \node[anchor=east] at (0,0) {$\ket{0}_A$};
  % première étape : H sur B, identité (pointillée) sur A
  \node[porte] at (2.2,1.4) {$H$};
  \node[draw=insagris, dashed, fill=white, minimum size=7mm] at (2.2,0) {$I$};
  % CNOT : contrôle B (haut), cible A (bas)
  \draw[insarouge, thick] (5,1.4) -- (5,0);
  \pic at (5,1.4) {ctl};
  \pic at (5,0) {cible};
  % accolade de sortie
  \draw[insagris, thick] (7.75,1.6) -- (7.9,1.6) -- (7.9,-0.2) -- (7.75,-0.2);
  \node[anchor=west] at (7.9,0.7) {$\ket{\beta^{00}}_{AB}$};
  % états intermédiaires
  \foreach \x in {0.7,3.6,6.4}
    \draw[dashed, insagris!70] (\x,1.9) -- (\x,-0.7);
  \node[anchor=north, font=\footnotesize] at (0.7,-0.7) {$\ket{\Pi_0}=\ket{00}$};
  \node[anchor=north, font=\footnotesize] at (3.6,-0.7) {$\ket{\Pi_1}=\ket{0}\otimes\ket{+}$};
  \node[anchor=north, font=\footnotesize] at (6.4,-0.7) {$\ket{\Pi_2}=\ket{\beta^{00}}$};
\end{tikzpicture}
\caption{Générateur d'états de Bell : $I\otimes H$, puis un CNOT de contrôle $B$ (fil du haut)
et de cible $A$ (fil du bas). L'entrée représentée est $\ket{00}$.}
\label{fig:circuit-bell}
\end{figure}

\subsubsection{Entrée \texorpdfstring{$\ket{00}$}{|00>} : calcul pas à pas}

\begin{align*}
  \ket{\Pi_0}&=\ket{0}_A\otimes\ket{0}_B=\ket{00},\\[1ex]
  \ket{\Pi_1}&=(I\otimes H)\ket{\Pi_0}=I\ket{0}_A\otimes H\ket{0}_B=\ket{0}\otimes\ket{+}
    =\frac{\ket{00}+\ket{01}}{\sqrt2},\\[1ex]
  \ket{\Pi_2}&=\mathrm{CNOT}\ket{\Pi_1}=\frac{\mathrm{CNOT}\ket{00}+\mathrm{CNOT}\ket{01}}{\sqrt2}
    =\frac{\ket{00}+\ket{11}}{\sqrt2}=\ket{\beta^{00}}=\ket{\Phi^+}.
\end{align*}
Dans $\ket{\Pi_1}$, le premier chiffre est celui de $A$ et le second celui de $B$. Le CNOT a
pour contrôle $B$ : $\ket{00}$ ($B=0$) ne change pas, $\ket{01}$ ($B=1$) devient $\ket{11}$
puisque $A$ est inversé. Le résultat est un \emph{état de Bell}.

Avec les matrices, la règle $A\otimes B=(A_{ij}B)$ donne
\[
  I\otimes H=\begin{pmatrix}1\cdot H&0\\0&1\cdot H\end{pmatrix}
  =\frac{1}{\sqrt2}\begin{pmatrix}1&1&0&0\\1&-1&0&0\\0&0&1&1\\0&0&1&-1\end{pmatrix}.
\]
Sa première colonne est $\ket{\Pi_1}=\frac{1}{\sqrt2}(1,1,0,0)^{\mathsf T}$ (image de $\ket{00}$).
Le CNOT échange les amplitudes de $\ket{01}$ et de $\ket{11}$ (section~\ref{sec:cnot}), donc
\[
  \ket{\Pi_2}=\mathrm{CNOT}\,\frac{1}{\sqrt2}\begin{pmatrix}1\\1\\0\\0\end{pmatrix}
  =\frac{1}{\sqrt2}\begin{pmatrix}1\\0\\0\\1\end{pmatrix}.
\]

\subsubsection{Les quatre entrées}

Prenons $B$ dans l'état $\ket{x}$ et $A$ dans l'état $\ket{y}$, avec $x,y\in\{0,1\}$. Alors
$\ket{\Pi_0}=\ket{y}_A\otimes\ket{x}_B=\ket{y,x}$. Comme
$H\ket{x}=\frac{1}{\sqrt2}\big(\ket{0}+(-1)^x\ket{1}\big)$ (c'est $\ket{+}$ pour $x=0$ et
$\ket{-}$ pour $x=1$),
\[
  \ket{\Pi_1}=\ket{y}\otimes H\ket{x}=\frac{1}{\sqrt2}\Big(\ket{y,0}+(-1)^x\ket{y,1}\Big).
\]
Le CNOT (contrôle $B$) ne change pas $\ket{y,0}$ et transforme $\ket{y,1}$ en
$\ket{\bar y,1}$, avec $\bar y=1-y$ :
\[
  \ket{\Pi_2}=\ket{\beta^{xy}}=\frac{1}{\sqrt2}\Big(\ket{y,0}+(-1)^x\ket{\bar y,1}\Big).
\]
Le premier indice de $\ket{\beta^{xy}}$ est l'entrée de $B$ (le qubit qui reçoit $H$), le
second celle de $A$ (la cible). Les quatre cas :

\begin{center}
\renewcommand{\arraystretch}{1.7}
\begin{tabular}{@{}ccccc@{}}
\toprule
$(B,A)=(x,y)$ & $\ket{\Pi_0}=\ket{y,x}$ & $\ket{\Pi_1}$ & $\ket{\Pi_2}=\ket{\beta^{xy}}$ & Nom \\
\midrule
$(0,0)$ & $\ket{00}$ & $\frac{1}{\sqrt2}(\ket{00}+\ket{01})$ & $\frac{1}{\sqrt2}(\ket{00}+\ket{11})$ & $\ket{\Phi^+}$ \\
$(1,0)$ & $\ket{01}$ & $\frac{1}{\sqrt2}(\ket{00}-\ket{01})$ & $\frac{1}{\sqrt2}(\ket{00}-\ket{11})$ & $\ket{\Phi^-}$ \\
$(0,1)$ & $\ket{10}$ & $\frac{1}{\sqrt2}(\ket{10}+\ket{11})$ & $\frac{1}{\sqrt2}(\ket{10}+\ket{01})$ & $\ket{\Psi^+}$ \\
$(1,1)$ & $\ket{11}$ & $\frac{1}{\sqrt2}(\ket{10}-\ket{11})$ & $\frac{1}{\sqrt2}(\ket{10}-\ket{01})$ & $-\ket{\Psi^-}$ \\
\bottomrule
\end{tabular}
\end{center}
Détail de la dernière ligne : $\ket{\Pi_1}=\ket{1}\otimes\ket{-}=\frac{1}{\sqrt2}(\ket{10}-\ket{11})$ ;
le CNOT laisse $\ket{10}$ et envoie $\ket{11}$ sur $\ket{01}$, d'où
$\frac{1}{\sqrt2}(\ket{10}-\ket{01})=-\frac{1}{\sqrt2}(\ket{01}-\ket{10})=-\ket{\Psi^-}$.

\paragraph{Forme matricielle.} Le CNOT échange les lignes $2$ et $4$ de
$I\otimes H$ (il échange les amplitudes de $\ket{01}$ et $\ket{11}$), donc
\[
  \mathrm{CNOT}\,(I\otimes H)=\frac{1}{\sqrt2}\begin{pmatrix}
    1 & 1 & 0 & 0\\ 0 & 0 & 1 & -1\\ 0 & 0 & 1 & 1\\ 1 & -1 & 0 & 0
  \end{pmatrix}.
\]
Ses colonnes sont les images de $\ket{00},\ket{01},\ket{10},\ket{11}$, c'est-à-dire de
$(B,A)=(0,0),(1,0),(0,1),(1,1)$ : ce sont $\ket{\beta^{00}}$, $\ket{\beta^{10}}$,
$\ket{\beta^{01}}$, $\ket{\beta^{11}}$, soit $\ket{\Phi^+}$, $\ket{\Phi^-}$, $\ket{\Psi^+}$
et $-\ket{\Psi^-}$.

\subsubsection{Matrice de densité de l'état de sortie}

Pour $\ket{\Psi}=\ket{\beta^{00}}=\frac{1}{\sqrt2}\ket{00}+\frac{1}{\sqrt2}\ket{11}$ :
\[
  \rho_{AB}=\ket{\Psi}\bra{\Psi}
  =\begin{pmatrix}\frac{1}{\sqrt2}\\0\\0\\\frac{1}{\sqrt2}\end{pmatrix}
   \begin{pmatrix}\frac{1}{\sqrt2}&0&0&\frac{1}{\sqrt2}\end{pmatrix}
  =\frac12\begin{pmatrix}1&0&0&1\\0&0&0&0\\0&0&0&0\\1&0&0&1\end{pmatrix}.
\]
Sous forme de ket-bra :
$\rho_{AB}=\frac12\big(\ket{00}\bra{00}+\ket{00}\bra{11}+\ket{11}\bra{00}+\ket{11}\bra{11}\big)$.
On vérifie $\operatorname{Tr}\rho_{AB}=\frac12(1+1)=1$ et
\[
  \rho_{AB}^2=\frac14\begin{pmatrix}2&0&0&2\\0&0&0&0\\0&0&0&0\\2&0&0&2\end{pmatrix}=\rho_{AB},
\]
donc $\operatorname{Tr}\rho_{AB}^2=1$ : le couple $(A,B)$ est dans un état \emph{pur}.

\begin{remarque}[Porte intriquante]
Les états d'entrée $\ket{ij}$ sont séparables ($D=0$) ; les sorties sont intriquées
($D=\pm\frac12$). Aucune porte agissant séparément sur $A$ et sur $B$ ($U_A\otimes U_B$,
section~\ref{sec:op2}) ne peut créer de l'intrication à partir d'un produit, puisque
$(U_A\otimes U_B)(\ket{u}\otimes\ket{v})=U_A\ket{u}\otimes U_B\ket{v}$ reste un produit.
Une porte à deux qubits comme le CNOT est indispensable.
\end{remarque}

% ---------------------------------------------------------------------
\subsection{Trace partielle : l'état d'un qubit pris seul}\label{sec:trace}
% ---------------------------------------------------------------------

Un état intriqué n'a pas de ket pour $A$ seul (section~\ref{sec:separable}). On peut pourtant
décrire $A$ par une matrice de densité, obtenue en «~oubliant~» $B$.

\begin{definition}{Trace partielle}
Soit $\rho_{AB}$ la matrice de densité du couple $(A,B)$. L'état du qubit $A$ pris seul est
$\rho_A=\operatorname{Tr}_B(\rho_{AB})$, celui de $B$ est $\rho_B=\operatorname{Tr}_A(\rho_{AB})$. La
trace partielle est linéaire et remplace la partie tracée de chaque terme par sa trace :
\[
  \operatorname{Tr}_B\big(\ket{a}\bra{a'}\otimes\ket{b}\bra{b'}\big)=\ket{a}\bra{a'}\;\braket{b'}{b},
  \qquad
  \operatorname{Tr}_A\big(\ket{a}\bra{a'}\otimes\ket{b}\bra{b'}\big)=\braket{a'}{a}\;\ket{b}\bra{b'}.
\]
\end{definition}

\paragraph{En matrices.} On découpe $\rho_{AB}$ (matrice $4\times4$, lignes et colonnes ordonnées
$00,01,10,11$) en quatre blocs $2\times2$ :
\[
  \rho_{AB}=\begin{pmatrix}R_{00}&R_{01}\\R_{10}&R_{11}\end{pmatrix}
  \quad\Longrightarrow\quad
  \rho_A=\begin{pmatrix}\operatorname{Tr}R_{00}&\operatorname{Tr}R_{01}\\
                        \operatorname{Tr}R_{10}&\operatorname{Tr}R_{11}\end{pmatrix},
  \qquad
  \rho_B=R_{00}+R_{11}.
\]
Autrement dit $(\rho_A)_{ij}=\rho_{(i0),(j0)}+\rho_{(i1),(j1)}$ : $\rho_A$ est la matrice des traces
des blocs, et $\rho_B$ la somme des deux blocs diagonaux.

\begin{remarque}[Pourquoi cette définition]
Toutes les prédictions qui portent sur $A$ seul viennent d'opérateurs de la forme $O_A\otimes I$, et
\[
  \operatorname{Tr}\big(\rho_{AB}\,(O_A\otimes I)\big)=\operatorname{Tr}\big(\rho_A\,O_A\big)
\]
(on le vérifie sur $\ket{a}\bra{a'}\otimes\ket{b}\bra{b'}$, puis par linéarité). La matrice $\rho_A$
contient donc exactement ce que l'on peut prédire sur $A$ sans regarder $B$.
\end{remarque}

\paragraph{Exemple 1 : état séparable $\ket{0}\otimes\ket{+}$.} Ici
$\rho_{AB}=\ket{0}\bra{0}\otimes\ket{+}\bra{+}$, donc
\[
  \rho_A=\ket{0}\bra{0}\,\braket{+}{+}=\ket{0}\bra{0},
  \qquad
  \rho_B=\braket{0}{0}\,\ket{+}\bra{+}=\ket{+}\bra{+}.
\]
Avec les matrices,
\[
  \rho_{AB}=\frac12\begin{pmatrix}1&1&0&0\\1&1&0&0\\0&0&0&0\\0&0&0&0\end{pmatrix},
  \qquad
  R_{00}=\frac12\begin{pmatrix}1&1\\1&1\end{pmatrix},\quad R_{01}=R_{10}=R_{11}=0,
\]
d'où
\[
  \rho_A=\begin{pmatrix}\operatorname{Tr}R_{00}&0\\0&0\end{pmatrix}=\begin{pmatrix}1&0\\0&0\end{pmatrix},
  \qquad
  \rho_B=R_{00}+R_{11}=\frac12\begin{pmatrix}1&1\\1&1\end{pmatrix}.
\]
Un qubit d'un état séparable garde son propre état pur.

\paragraph{Exemple 2 : état de Bell $\ket{\Phi^+}$.} On écrit
$\rho_{AB}=\ket{\Phi^+}\bra{\Phi^+}=\frac12\big(\ket{00}\bra{00}+\ket{00}\bra{11}+\ket{11}\bra{00}
+\ket{11}\bra{11}\big)$. Chaque terme est un produit d'opérateurs à un qubit, par exemple
$\ket{00}\bra{11}=\ket{0}\bra{1}\otimes\ket{0}\bra{1}$, donc
\begin{align*}
  \rho_A=\operatorname{Tr}_B(\rho_{AB})
  &=\tfrac12\Big(\ket{0}\bra{0}\,\braket{0}{0}+\ket{0}\bra{1}\,\braket{1}{0}
      +\ket{1}\bra{0}\,\braket{0}{1}+\ket{1}\bra{1}\,\braket{1}{1}\Big)\\
  &=\tfrac12\Big(\ket{0}\bra{0}\cdot1+\ket{0}\bra{1}\cdot0+\ket{1}\bra{0}\cdot0+\ket{1}\bra{1}\cdot1\Big)
   =\tfrac12\big(\ket{0}\bra{0}+\ket{1}\bra{1}\big)=\tfrac12\,\mathrm{Id}.
\end{align*}
Par les blocs,
\[
  \rho_{AB}=\frac12\begin{pmatrix}1&0&0&1\\0&0&0&0\\0&0&0&0\\1&0&0&1\end{pmatrix},
  \qquad
  R_{00}=\frac12\begin{pmatrix}1&0\\0&0\end{pmatrix},\quad
  R_{11}=\frac12\begin{pmatrix}0&0\\0&1\end{pmatrix},\quad
  \operatorname{Tr}R_{01}=\operatorname{Tr}R_{10}=0,
\]
d'où $\rho_A=\frac12\mathrm{Id}$ et $\rho_B=R_{00}+R_{11}=\frac12\mathrm{Id}$ : le même résultat.

Le qubit $A$ est donc dans l'état \emph{complètement mixte} ($\operatorname{Tr}\rho_A^2=\frac14+\frac14=\frac12$),
au centre de la boule de Bloch (section~\ref{sec:bloch}), alors que le couple est dans un état pur
($\operatorname{Tr}\rho_{AB}^2=1$, section~\ref{sec:circuit-bell}). On connaît parfaitement l'état du
couple mais pas celui d'un qubit pris seul : l'information est entièrement contenue dans les
corrélations entre $A$ et $B$. Mesuré seul, $A$ donne $0$ ou $1$ avec la probabilité $\frac12$, dans
toutes les bases.

\paragraph{Formule pratique pour un état pur.} Pour $\ket{\Psi}=\sum_{ij}c_{ij}\ket{ij}$, notons
$C=\begin{pmatrix}c_{00}&c_{01}\\c_{10}&c_{11}\end{pmatrix}$ la matrice $2\times2$ des coefficients
(ligne $i$ : état de $A$, colonne $j$ : état de $B$). Comme $\rho_{(ij),(i'j')}=c_{ij}\,c_{i'j'}^*$,
\[
  (\rho_A)_{ii'}=\sum_jc_{ij}\,c_{i'j}^*,\qquad\text{soit}\qquad \rho_A=C\,C^\dagger .
\]
\emph{Exemple 3.} Pour $\ket{\Psi}=\frac12\ket{00}+\frac12\ket{01}+\frac{1}{\sqrt2}\ket{11}$ (exemple~2
de la section~\ref{sec:mesure2}), $C=\begin{pmatrix}\frac12&\frac12\\0&\frac{1}{\sqrt2}\end{pmatrix}$ et
\[
  \rho_A=CC^\dagger=\begin{pmatrix}\frac14+\frac14&0+\frac{1}{2\sqrt2}\\[0.4ex]\frac{1}{2\sqrt2}&\frac12\end{pmatrix}
  =\begin{pmatrix}\frac12&\frac{\sqrt2}{4}\\[0.4ex]\frac{\sqrt2}{4}&\frac12\end{pmatrix},
  \qquad
  \operatorname{Tr}\rho_A^2=\tfrac14+\tfrac14+2\cdot\tfrac18=\tfrac34 .
\]
Le qubit $A$ est mixte, mais pas complètement : cet état est intriqué sans être un état de Bell.

\begin{resultat}{Pureté d'un qubit et séparabilité}
Pour un état pur $\ket{\Psi}=\sum c_{ij}\ket{ij}$ de deux qubits, avec $D=c_{00}c_{11}-c_{01}c_{10}$,
\[
  \operatorname{Tr}\rho_A^2=\operatorname{Tr}\rho_B^2=1-2|D|^2 .
\]
Donc $\ket{\Psi}$ est séparable ($D=0$) si et seulement si $\rho_A$ est pur. Plus $|D|$ est grand,
plus $\rho_A$ est mixte, jusqu'à $|D|=\frac12$ et $\rho_A=\mathrm{Id}/2$ pour les états de Bell.
\end{resultat}

\emph{Démonstration.} La matrice $\rho_A=CC^\dagger$ a pour trace $\sum|c_{ij}|^2=1$ et pour
déterminant $\det C\,\overline{\det C}=|D|^2$. Si $\lambda_1,\lambda_2$ sont ses valeurs propres,
$\lambda_1+\lambda_2=1$ et $\operatorname{Tr}\rho_A^2=\lambda_1^2+\lambda_2^2=1-2\lambda_1\lambda_2
=1-2|D|^2$ ; le calcul est le même pour $\rho_B$. \hfill$\square$

Comme $\operatorname{Tr}\rho_A^2\geq\frac12$ pour un qubit, on en déduit $|D|\leq\frac12$. Avec
$\operatorname{Tr}\rho_A^2=\frac{1+|\vec r|^2}{2}$ (section~\ref{sec:bloch}), le vecteur de Bloch de $A$
a pour norme $|\vec r|=\sqrt{1-4|D|^2}$ : il part de la sphère pour un état séparable et se rapproche
du centre quand l'intrication croît (figure~\ref{fig:trace-bloch}).

\begin{center}
\renewcommand{\arraystretch}{1.5}
\begin{tabular}{@{}lcccl@{}}
\toprule
\textbf{État du couple} & $D$ & $\operatorname{Tr}\rho_A^2$ & $|\vec r|$ & \textbf{Qubit $A$} \\
\midrule
$\ket{0}\otimes\ket{+}$ & $0$ & $1$ & $1$ & pur, sur la sphère \\
$\frac12\ket{00}+\frac12\ket{01}+\frac{1}{\sqrt2}\ket{11}$ & $\frac{\sqrt2}{4}$ & $\frac34$
  & $\frac{1}{\sqrt2}$ & mixte, dans la boule \\
$\ket{\Phi^+}$ (comme les trois autres états de Bell) & $\pm\frac12$ & $\frac12$ & $0$
  & complètement mixte, au centre \\
\bottomrule
\end{tabular}
\end{center}

\begin{figure}[!ht]
\centering
\adjustbox{max width=\linewidth}{%
\begin{tikzpicture}[>=Stealth, font=\small]
  \def\R{2.3}
  \fill[insagris!12] (0,0) circle (\R);
  \draw[insagris, dotted] (-\R,0) -- (\R,0);
  \draw[insagris, dotted] (0,-\R) -- (0,\R);
  \draw[insagris, dashed] (0,0) circle ({\R*0.70711});
  \draw[insarouge, very thick] (0,0) circle (\R);
  \node[font=\footnotesize, text=insagris, anchor=south west] at (\R,0) {$x$};
  \node[font=\footnotesize, text=insagris, anchor=south west] at (0.03,\R) {$z$};
  % vecteurs de Bloch de rho_A pour trois états du couple
  \fill[insarouge] (0,\R) circle (2.4pt);
  \fill[insagris!80!black] ({\R*0.70711},0) circle (2.4pt);
  \fill[insagris!80!black] (0,0) circle (2.4pt);
  \draw[insagris] (0,\R) -- (3.1,\R);
  \node[anchor=west, font=\footnotesize, align=left] at (3.1,\R)
    {$\ket{0}\otimes\ket{+}$ : $\rho_A=\ket{0}\bra{0}$\\ $D=0$, \ $|\vec r|=1$ (état pur)};
  \draw[insagris] ({\R*0.70711},0) -- (3.1,0);
  \node[anchor=west, font=\footnotesize, align=left] at (3.1,0)
    {$\frac12\ket{00}+\frac12\ket{01}+\frac{1}{\sqrt2}\ket{11}$\\
     $D=\frac{\sqrt2}{4}$, \ $|\vec r|=\frac{1}{\sqrt2}$ (mixte)};
  \draw[insagris] (0,0) -- (3.1,-\R);
  \node[anchor=west, font=\footnotesize, align=left] at (3.1,-\R)
    {$\ket{\Phi^+}$ : $\rho_A=\mathrm{Id}/2$\\ $D=\frac12$, \ $\vec r=\vec 0$ (centre)};
\end{tikzpicture}}
\caption{Le vecteur de Bloch de $\rho_A$ (boule de Bloch vue en coupe) pour trois états du couple
$(A,B)$. Le cercle pointillé est la sphère de rayon $1/\sqrt2$. Plus l'état est intriqué, plus le
point s'éloigne de la sphère des états purs : $|\vec r|=\sqrt{1-4|D|^2}$.}
\label{fig:trace-bloch}
\end{figure}

% ---------------------------------------------------------------------
\subsection{Information : variables aléatoires et états quantiques}\label{sec:info}
% ---------------------------------------------------------------------

\begin{definition}{Entropie, information mutuelle, entropie conditionnelle}
Soit $A$ une variable aléatoire de loi $p(a)$. Son \emph{entropie de Shannon}, en bits, est
\[
  H(A)=-\sum_a p(a)\log_2p(a)\qquad(\text{avec }0\log_20=0).
\]
Elle mesure l'incertitude moyenne sur $A$ : une pièce équilibrée a $H=1$ bit. Pour un couple
$(A,B)$ de loi $p(a,b)$, $H(A,B)=-\sum_{a,b}p(a,b)\log_2p(a,b)$, et
\begin{align*}
  I(A;B)&=H(A)+H(B)-H(A,B) &&\text{(information mutuelle)},\\
  H(A|B)&=H(A,B)-H(B)=H(A)-I(A;B) &&\text{(entropie conditionnelle)}.
\end{align*}
\end{definition}

Les deux écritures de $H(A|B)$ sont égales : la définition de $I(A;B)$ donne
$H(A,B)=H(A)+H(B)-I(A;B)$, donc $H(A,B)-H(B)=H(A)-I(A;B)$. L'information mutuelle $I(A;B)$
mesure ce que la connaissance de $B$ apprend sur $A$ ; $H(A|B)$ est l'incertitude qui reste sur
$A$ quand on connaît $B$.

\subsubsection{Deux variables aléatoires indépendantes}

Alice et Bob lancent chacun une pièce équilibrée, sans lien entre les deux lancers :
$p(a,b)=\frac14$ pour les quatre couples. Chaque pièce donne $0$ ou $1$ avec la probabilité $\frac12$, donc
\begin{align*}
  H(A)&=-\Big(\tfrac12\log_2\tfrac12+\tfrac12\log_2\tfrac12\Big)=1\text{ bit},\\
  H(B)&=1\text{ bit (même calcul)},\\
  H(A,B)&=-4\cdot\tfrac14\log_2\tfrac14=2\text{ bits},\\
  I(A;B)&=H(A)+H(B)-H(A,B)=1+1-2=0,\\
  H(A|B)&=H(A)-I(A;B)=1-0=H(A)=1\text{ bit}.
\end{align*}
Connaître $B$ n'apprend rien sur $A$.

\subsubsection{Deux variables solidaires}

Les deux pièces tombent toujours du même côté : $p(0,0)=p(1,1)=\frac12$ et
$p(0,1)=p(1,0)=0$. Les lois de chaque pièce sont inchangées ($p(a{=}0)=p(0,0)+p(0,1)=\frac12$,
etc.), donc $H(A)=H(B)=1$ bit, mais
\begin{align*}
  H(A,B)&=-2\cdot\tfrac12\log_2\tfrac12-2\cdot0\log_20=1\text{ bit},\\
  I(A;B)&=1+1-1=1\text{ bit},\qquad H(A|B)=1-1=0.
\end{align*}
Connaître $B$ détermine $A$ : plus aucune incertitude ne reste.

\subsubsection{Deux qubits dans un état de Bell}

Pour un système quantique, on remplace la loi de probabilité par la matrice de densité.

\begin{definition}{Entropie de von Neumann}
L'\emph{entropie de von Neumann} d'un état $\rho$, de valeurs propres $\lambda_i$, est
\[
  S(\rho)=-\operatorname{Tr}(\rho\log_2\rho)=-\sum_i\lambda_i\log_2\lambda_i .
\]
Elle vaut $0$ pour un état pur et $1$ bit pour le qubit complètement mélangé $\mathrm{Id}/2$.
En quantique, $H(A)=S(\rho_A)$, $H(B)=S(\rho_B)$ et $H(A,B)=S(\rho_{AB})$ ; les définitions de
$I(A;B)$ et de $H(A|B)$ sont les mêmes.
\end{definition}

On prend l'état de Bell $\ket{\Phi^+}$ produit par le circuit de la
section~\ref{sec:circuit-bell}. Le couple est dans un état pur,
$\rho_{AB}=\ket{\Phi^+}\bra{\Phi^+}$ ; chaque qubit pris seul est dans l'état
$\rho_A=\rho_B=\frac12\mathrm{Id}$ (trace partielle, section~\ref{sec:trace}).

Calcul des entropies :
\begin{itemize}[nosep]
  \item $\rho_{AB}$ est pur : ses valeurs propres sont $1,0,0,0$, donc
        $H(A,B)=S(\rho_{AB})=-1\cdot\log_21=0$ ;
  \item $\rho_A$ et $\rho_B$ ont pour valeurs propres $\frac12,\frac12$, donc
        $H(A)=H(B)=-\big(\frac12\log_2\frac12+\frac12\log_2\frac12\big)=1$ bit.
\end{itemize}
D'où
\begin{align*}
  I(A;B)&=H(A)+H(B)-H(A,B)=1+1-0=2\text{ bits},\\
  H(A|B)&=H(A,B)-H(B)=0-1=-1\text{ bit}\qquad(=H(A)-I(A;B)=1-2).
\end{align*}

\begin{remarque}[Une entropie conditionnelle négative]
Pour des variables classiques, $p(a,b)\leq p(b)$, donc $-\log_2p(a,b)\geq-\log_2p(b)$ et
$H(A,B)\geq H(B)$ : on a toujours $H(A|B)\geq0$, et $I(A;B)\leq H(A)$. Ici $H(A|B)=-1$ et
$I(A;B)=2>H(A)$ : ces valeurs n'ont pas d'équivalent classique, c'est une signature de
l'intrication. Le résultat est surprenant mais mathématiquement correct.
\end{remarque}

\begin{center}
\renewcommand{\arraystretch}{1.4}
\begin{tabular}{@{}lccc@{}}
\toprule
 & indépendantes & solidaires & état de Bell \\
\midrule
$H(A)$ & $1$ & $1$ & $1$ \\
$H(B)$ & $1$ & $1$ & $1$ \\
$H(A,B)$ & $2$ & $1$ & $0$ \\
$I(A;B)$ & $0$ & $1$ & $2$ \\
$H(A|B)$ & $1$ & $0$ & $-1$ \\
\bottomrule
\end{tabular}
\end{center}
(en bits ; les deux premières colonnes sont classiques, la troisième est quantique).

\begin{figure}[!ht]
\centering
\adjustbox{max width=\linewidth}{%
\begin{tikzpicture}[>=Stealth, font=\small]
  \def\bw{0.8}
  % axe des ordonnées (bits)
  \draw[->] (0,-1.5) -- (0,2.8) node[above] {bits};
  \foreach \v in {-1,0,1,2} {
    \draw (-0.07,\v) -- (0.07,\v);
    \node[left, font=\footnotesize] at (-0.1,\v) {$\v$};
  }
  \draw[insagris!60] (0,0) -- (11.4,0);
  % limite classique de I(A;B) : H(A) = 1
  \draw[dashed, insarouge] (0,1) -- (11.4,1);
  \node[anchor=west, font=\scriptsize, text=insarouge, align=left] at (11.4,1)
    {limite classique\\ $I(A;B)\leq H(A)=1$};
  \node[anchor=west, font=\scriptsize, text=insagris, align=left] at (11.4,0)
    {limite classique\\ $H(A|B)\geq0$};
  % barres : (groupe, H(A,B), I(A;B), H(A|B))
  \foreach \g/\hab/\mi/\hc in {0/2/0/1, 1/1/1/0, 2/0/2/-1} {
    \pgfmathsetmacro{\x}{0.9+3.7*\g}
    \fill[insagris!65] (\x,0) rectangle (\x+\bw,\hab);
    \fill[insarouge] (\x+\bw+0.1,0) rectangle (\x+2*\bw+0.1,\mi);
    \fill[black!80] (\x+2*\bw+0.2,0) rectangle (\x+3*\bw+0.2,\hc);
    \node[above, font=\footnotesize] at (\x+0.5*\bw,{max(\hab,0)}) {$\hab$};
    \node[above, font=\footnotesize] at (\x+1.5*\bw+0.1,{max(\mi,0)}) {$\mi$};
    \ifnum\hc<0
      \node[below, font=\footnotesize] at (\x+2.5*\bw+0.2,\hc) {$\hc$};
    \else
      \node[above, font=\footnotesize] at (\x+2.5*\bw+0.2,\hc) {$\hc$};
    \fi
  }
  \node[align=center, anchor=north] at (2.15,-1.65) {variables\\ indépendantes};
  \node[align=center, anchor=north] at (5.85,-1.65) {variables\\ solidaires};
  \node[align=center, anchor=north] at (9.55,-1.65) {état de Bell\\ $\ket{\Phi^+}$};
  % légende
  \fill[insagris!65] (0.9,3.35) rectangle (1.2,3.65);
  \node[anchor=west, font=\footnotesize] at (1.25,3.5) {$H(A,B)$};
  \fill[insarouge] (3.6,3.35) rectangle (3.9,3.65);
  \node[anchor=west, font=\footnotesize] at (3.95,3.5) {$I(A;B)$};
  \fill[black!80] (6.3,3.35) rectangle (6.6,3.65);
  \node[anchor=west, font=\footnotesize] at (6.65,3.5) {$H(A|B)$};
\end{tikzpicture}}
\caption{Entropies (en bits) pour deux pièces indépendantes, deux pièces solidaires et un état de Bell.
Dans les trois cas $H(A)=H(B)=1$. Les deux premiers groupes respectent les bornes classiques
$I(A;B)\leq H(A)$ et $H(A|B)\geq0$ ; l'état de Bell les dépasse ($I=2$, $H(A|B)=-1$).}
\label{fig:entropies}
\end{figure}

% ---------------------------------------------------------------------
\subsection{Téléportation quantique}\label{sec:teleportation}
% ---------------------------------------------------------------------

\emph{Problème.} Alice détient un qubit $Q$ dans un état qu'elle ne connaît pas,
$\ket{\varphi}=\alpha\ket{0}+\beta\ket{1}$ avec $|\alpha|^2+|\beta|^2=1$, et elle veut que Bob
obtienne cet état sur l'un de ses qubits. Elle ne peut pas mesurer $Q$ pour connaître
$\alpha$ et $\beta$ (une mesure ne donne qu'un bit et détruit l'état), elle ne peut pas le
copier (théorème de non-clonage, remarque en fin de section) et elle n'envoie pas le qubit
lui-même. Le protocole est dû à \textcite{bennett1993}.

\begin{definition}{Protocole de téléportation quantique}
\textbf{Ressources.}
\begin{itemize}[nosep]
  \item Un \emph{e-bit} : Alice et Bob partagent une paire de qubits $(A,B)$ dans l'état de Bell
        $\ket{\Phi^+}_{BA}=\frac{1}{\sqrt2}\big(\ket{00}+\ket{11}\big)$ ; Alice détient $A$,
        Bob détient $B$.
  \item Le qubit $Q$ à téléporter, chez Alice.
  \item Un canal classique (deux bits).
\end{itemize}
Le système est écrit $(B,A,Q)$ : $\ket{b,a,q}=\ket{b}_B\otimes\ket{a}_A\otimes\ket{q}_Q$.

\textbf{Étapes.}
\begin{enumerate}[nosep]
  \item Alice applique un CNOT de contrôle $Q$ et de cible $A$.
  \item Alice applique une porte $H$ sur $Q$.
  \item Alice mesure $A$ et $Q$ dans la base de calcul, obtient deux bits $(a,q)$ et les envoie
        à Bob.
  \item Bob applique $X$ si $a=1$, puis $Z$ si $q=1$, à son qubit $B$.
\end{enumerate}
\end{definition}

La figure~\ref{fig:tele-schema} montre l'organisation du protocole et la
figure~\ref{fig:teleportation} son circuit.

\begin{figure}[!ht]
\centering
\adjustbox{max width=\linewidth}{%
\begin{tikzpicture}[>=Stealth, font=\small,
    qb/.style={circle, draw=insarouge, fill=insarouge!8, thick, minimum size=8mm, inner sep=0pt},
    zone/.style={draw=insagris, rounded corners=4pt, thick, dashed, inner sep=8pt},
    num/.style={circle, fill=insarouge, text=white, font=\footnotesize\bfseries, inner sep=0pt,
                minimum size=4.6mm},
    etape/.style={anchor=west, font=\footnotesize, text=insarouge},
    onde/.style={thick, insarouge, decorate,
                 decoration={snake, amplitude=1.3pt, segment length=7pt}}]
  % Alice : Q (état inconnu) et A
  \node[qb] (Q) at (0,1.3) {$Q$};
  \node[qb] (A) at (0,-0.5) {$A$};
  \node[zone, fit=(Q)(A), label={[text=insagris]above:Alice}] (alice) {};
  \node[anchor=east] at (-0.95,1.3) {$\ket{\varphi}=\alpha\ket{0}+\beta\ket{1}$};
  \draw[insarouge, thick] (Q) -- (A);
  % Bob : B
  \node[qb] (B) at (9,-0.5) {$B$};
  \node[zone, fit=(B), label={[text=insagris]above:Bob}] (bob) {};
  % source de la paire de Bell
  \node[draw=insarouge, fill=insarouge!10, rounded corners=2pt, thick, minimum width=3.2cm,
        minimum height=0.9cm] (src) at (4.5,-2.3) {source $\ket{\Phi^+}_{BA}$};
  \draw[onde] (src.west) -| (A.south);
  \draw[onde] (src.east) -| (B.south);
  % étapes d'Alice : pastilles sur le bord droit de sa zone
  \node[num] (n1) at (alice.east |- Q) {1};
  \node[num] (n2) at (alice.east) {2};
  \node[num] (n3) at (alice.east |- A) {3};
  \node[etape] at (n1.east) {\ CNOT de $Q$ vers $A$};
  \node[etape] at (n2.east) {\ porte $H$ sur $Q$};
  % canal classique
  \draw[double, double distance=1.2pt, ->, insagris] (n3.east) -- (bob.west);
  \node[above, font=\footnotesize] at (4.6,-0.45) {deux bits classiques $(a,q)$};
  \node[below, font=\footnotesize, text=insarouge] at (4.6,-0.55) {mesure de $A$ et $Q$ (étape 3)};
  % correction de Bob
  \node[num] (n4) at (bob.east) {4};
  \node[etape, align=left] at (n4.east) {\ $X$ si $a=1$, puis $Z$ si $q=1$};
  \node[anchor=west, font=\footnotesize, align=left] at ($(n4.east)+(0.05,-0.8)$)
    {\ $B$ est alors dans l'état $\ket{\varphi}$};
\end{tikzpicture}}
\caption{Organisation de la téléportation. Alice possède le qubit $Q$ à téléporter et la moitié $A$
d'une paire de Bell, Bob possède l'autre moitié $B$. (1)~CNOT de $Q$ vers $A$, (2)~porte $H$ sur $Q$,
(3)~mesure de $A$ et $Q$, dont les résultats $(a,q)$ partent par un canal classique, (4)~corrections de
Bob. Aucun qubit ne voyage entre Alice et Bob.}
\label{fig:tele-schema}
\end{figure}

\begin{figure}[!ht]
\centering
\begin{adjustbox}{max width=\linewidth}
\begin{tikzpicture}[>=Stealth, font=\small]
  % fils quantiques : Q en haut, A au milieu, B en bas
  \draw (0,2.4) -- (5.3,2.4);
  \draw (0,1.2) -- (5.3,1.2);
  \draw (0,0) -- (10.6,0);
  \node[anchor=east] at (0,2.4) {$\ket{\varphi}_Q$};
  \node[anchor=south west, text=insagris, font=\footnotesize] at (0.05,1.2) {$A$};
  \node[anchor=south west, text=insagris, font=\footnotesize] at (0.05,0) {$B$};
  % paire de Bell partagée
  \draw[insagris, thick] (-0.1,1.5) -- (-0.25,1.5) -- (-0.25,-0.3) -- (-0.1,-0.3);
  \node[anchor=east] at (-0.35,0.6) {$\ket{\Phi^+}_{BA}$};
  % CNOT : contrôle Q, cible A
  \draw[insarouge, thick] (2.6,2.4) -- (2.6,1.2);
  \pic at (2.6,2.4) {ctl};
  \pic at (2.6,1.2) {cible};
  % Hadamard sur Q
  \node[porte] at (4.2,2.4) {$H$};
  % mesures d'Alice
  \pic at (5.6,2.4) {mesure};
  \pic at (5.6,1.2) {mesure};
  % fils classiques (doubles)
  \draw[double, double distance=1.2pt, insagris] (5.9,1.2) -- (7.6,1.2) -- (7.6,0.35);
  \draw[double, double distance=1.2pt, insagris] (5.9,2.4) -- (9.6,2.4) -- (9.6,0.35);
  \node[above, font=\footnotesize] at (6.5,1.2) {$a$};
  \node[above, font=\footnotesize] at (6.5,2.4) {$q$};
  % corrections de Bob
  \node[porte] at (7.6,0) {$X$};
  \node[porte] at (9.6,0) {$Z$};
  \node[anchor=west] at (10.6,0) {$\ket{\varphi}_B$};
  % états intermédiaires
  \foreach \x in {1.4,3.4,4.9}
    \draw[dashed, insagris!70] (\x,2.8) -- (\x,-0.5);
  \node[anchor=north, font=\footnotesize] at (1.4,-0.5) {$\ket{\Pi_0}$};
  \node[anchor=north, font=\footnotesize] at (3.4,-0.5) {$\ket{\Pi_1}$};
  \node[anchor=north, font=\footnotesize] at (4.9,-0.5) {$\ket{\Pi_2}$};
\end{tikzpicture}
\end{adjustbox}
\caption{Circuit de téléportation. Alice détient $Q$ (dans $\ket{\varphi}=\alpha\ket{0}+\beta\ket{1}$)
et $A$, Bob détient $B$ ; $A$ et $B$ forment la paire $\ket{\Phi^+}_{BA}$. Les traits doubles
sont des fils \emph{classiques} qui portent les résultats $a$ (mesure de $A$) et $q$ (mesure de
$Q$) ; $X$ est appliqué si $a=1$, $Z$ si $q=1$. Le ket s'écrit $\ket{b,a,q}$ : le fil du haut
($Q$) est le dernier chiffre.}
\label{fig:teleportation}
\end{figure}

\subsubsection{Calcul pas à pas}

\paragraph{État initial $\ket{\Pi_0}$.} Alice et Bob partagent la paire, et $Q$ est indépendant
d'elle :
\begin{align*}
  \ket{\Pi_0}&=\ket{\Phi^+}_{BA}\otimes\ket{\varphi}_Q
    =\frac{1}{\sqrt2}\big(\ket{00}+\ket{11}\big)\otimes\big(\alpha\ket{0}+\beta\ket{1}\big)\\
  &=\frac{1}{\sqrt2}\Big(\alpha\ket{00}\ket{0}+\beta\ket{00}\ket{1}+\alpha\ket{11}\ket{0}+\beta\ket{11}\ket{1}\Big)\\
  &=\frac{\alpha\ket{000}+\beta\ket{001}+\alpha\ket{110}+\beta\ket{111}}{\sqrt2}.
\end{align*}

\paragraph{Après le CNOT $Q\to A$ : $\ket{\Pi_1}$.} Le CNOT (contrôle $Q$, dernier chiffre ;
cible $A$, chiffre du milieu) agit ainsi : $\ket{b,a,q}\to\ket{b,\,a\oplus q,\,q}$. Terme par terme :
\begin{center}
\renewcommand{\arraystretch}{1.3}
\begin{tabular}{@{}cccc@{}}
\toprule
Terme & $q$ & Effet sur $A$ & Résultat \\
\midrule
$\alpha\ket{000}$ & $0$ & aucun & $\alpha\ket{000}$ \\
$\beta\ket{001}$ & $1$ & $0\to1$ & $\beta\ket{011}$ \\
$\alpha\ket{110}$ & $0$ & aucun & $\alpha\ket{110}$ \\
$\beta\ket{111}$ & $1$ & $1\to0$ & $\beta\ket{101}$ \\
\bottomrule
\end{tabular}
\end{center}
donc
\[
  \ket{\Pi_1}=\frac{\alpha\ket{000}+\beta\ket{011}+\alpha\ket{110}+\beta\ket{101}}{\sqrt2}.
\]
On isole ensuite le dernier chiffre, celui de $Q$, qui va changer à l'étape suivante
(le premier ket est celui du couple $\ket{b,a}$) :
\[
  \ket{\Pi_1}=\frac{\alpha\ket{00}\otimes\ket{0}+\beta\ket{01}\otimes\ket{1}
    +\alpha\ket{11}\otimes\ket{0}+\beta\ket{10}\otimes\ket{1}}{\sqrt2}.
\]

\paragraph{Après la porte $H$ sur $Q$ : $\ket{\Pi_2}$.} On remplace $\ket{0}$ par
$H\ket{0}=\ket{+}$ et $\ket{1}$ par $H\ket{1}=\ket{-}$ dans le facteur $Q$ :
\[
  \ket{\Pi_2}=\frac{1}{\sqrt2}\Big(\alpha\ket{00}\otimes\ket{+}+\beta\ket{01}\otimes\ket{-}
    +\alpha\ket{11}\otimes\ket{+}+\beta\ket{10}\otimes\ket{-}\Big).
\]
Avec $\ket{\pm}=\frac{1}{\sqrt2}(\ket{0}\pm\ket{1})$, le préfacteur devient
$\frac{1}{\sqrt2}\cdot\frac{1}{\sqrt2}=\frac12$ :
\begin{align*}
  \ket{\Pi_2}&=\frac12\Big(\alpha\ket{00}\big(\ket{0}+\ket{1}\big)+\beta\ket{01}\big(\ket{0}-\ket{1}\big)
     +\alpha\ket{11}\big(\ket{0}+\ket{1}\big)+\beta\ket{10}\big(\ket{0}-\ket{1}\big)\Big)\\
  &=\frac12\Big(\alpha\ket{000}+\alpha\ket{001}+\beta\ket{010}-\beta\ket{011}
     +\alpha\ket{110}+\alpha\ket{111}+\beta\ket{100}-\beta\ket{101}\Big).
\end{align*}

\paragraph{Regroupement selon le résultat d'Alice.} Alice mesurera $A$ et $Q$, donc les deux
derniers chiffres. On regroupe les huit termes qui ont les mêmes chiffres $(a,q)$, en mettant
le ket de Bob en premier :
\begin{center}
\renewcommand{\arraystretch}{1.3}
\begin{tabular}{@{}cll@{}}
\toprule
$(a,q)$ & Termes de $\ket{\Pi_2}$ (facteur $\frac12$ omis) & Ket de Bob \\
\midrule
$(0,0)$ & $\alpha\ket{000}+\beta\ket{100}$ & $(\alpha\ket{0}+\beta\ket{1})\ket{00}$ \\
$(0,1)$ & $\alpha\ket{001}-\beta\ket{101}$ & $(\alpha\ket{0}-\beta\ket{1})\ket{01}$ \\
$(1,0)$ & $\beta\ket{010}+\alpha\ket{110}$ & $(\beta\ket{0}+\alpha\ket{1})\ket{10}$ \\
$(1,1)$ & $-\beta\ket{011}+\alpha\ket{111}$ & $(-\beta\ket{0}+\alpha\ket{1})\ket{11}$ \\
\bottomrule
\end{tabular}
\end{center}
d'où
\[
  \ket{\Pi_2}=\frac12\Big[\big(\alpha\ket{0}+\beta\ket{1}\big)\ket{00}
    +\big(\alpha\ket{0}-\beta\ket{1}\big)\ket{01}
    +\big(\beta\ket{0}+\alpha\ket{1}\big)\ket{10}
    +\big(\alpha\ket{1}-\beta\ket{0}\big)\ket{11}\Big].
\]
Dans chaque terme, le ket de gauche est l'état de $B$ et le ket de droite $\ket{a,q}$ celui du
couple $(A,Q)$ d'Alice.

\subsubsection{Mesure d'Alice et corrections de Bob}

Alice mesure $A$ et $Q$. Comme à la section~\ref{sec:mesure2}, la probabilité d'un résultat
$(a,q)$ est le carré de la norme du terme correspondant, et l'état de Bob devient le ket
associé (déjà normé). Ici
\[
  P(a,q)=\Big\|\tfrac12\ket{\phi_{aq}}\Big\|^2=\tfrac14\big(|\alpha|^2+|\beta|^2\big)=\tfrac14
  \quad\text{pour les quatre résultats,}
\]
qui sont donc équiprobables et \emph{indépendants} de $\alpha$ et $\beta$ : les deux bits
envoyés ne révèlent rien sur $\ket{\varphi}$. Les kets de Bob sont
$\ket{\phi_{00}}=\alpha\ket{0}+\beta\ket{1}$, $\ket{\phi_{01}}=\alpha\ket{0}-\beta\ket{1}$,
$\ket{\phi_{10}}=\beta\ket{0}+\alpha\ket{1}$ et $\ket{\phi_{11}}=-\beta\ket{0}+\alpha\ket{1}$.

\begin{center}
\renewcommand{\arraystretch}{1.5}
\begin{tabular}{@{}cccll@{}}
\toprule
Résultat $(a,q)$ & Probabilité & État de Bob & Lien avec $\ket{\varphi}$ & Correction de Bob \\
\midrule
$(0,0)$ & $\frac14$ & $\alpha\ket{0}+\beta\ket{1}$ & $\ket{\varphi}$ & $I$ (rien) \\
$(0,1)$ & $\frac14$ & $\alpha\ket{0}-\beta\ket{1}$ & $Z\ket{\varphi}$ & $Z$ \\
$(1,0)$ & $\frac14$ & $\beta\ket{0}+\alpha\ket{1}$ & $X\ket{\varphi}$ & $X$ \\
$(1,1)$ & $\frac14$ & $-\beta\ket{0}+\alpha\ket{1}$ & $XZ\ket{\varphi}$ & $X$ puis $Z$ (opérateur $ZX$) \\
\bottomrule
\end{tabular}
\end{center}

\paragraph{Vérification des quatre corrections.} Avec $X\ket{0}=\ket{1}$, $X\ket{1}=\ket{0}$,
$Z\ket{0}=\ket{0}$ et $Z\ket{1}=-\ket{1}$ :
\begin{align*}
  (0,0):\quad & I\big(\alpha\ket{0}+\beta\ket{1}\big)=\alpha\ket{0}+\beta\ket{1},\\
  (0,1):\quad & Z\big(\alpha\ket{0}-\beta\ket{1}\big)=\alpha\ket{0}-\beta Z\ket{1}=\alpha\ket{0}+\beta\ket{1},\\
  (1,0):\quad & X\big(\beta\ket{0}+\alpha\ket{1}\big)=\beta\ket{1}+\alpha\ket{0}=\alpha\ket{0}+\beta\ket{1},\\
  (1,1):\quad & X\big(-\beta\ket{0}+\alpha\ket{1}\big)=-\beta\ket{1}+\alpha\ket{0}=\alpha\ket{0}-\beta\ket{1},
     \ \text{puis}\ Z\big(\alpha\ket{0}-\beta\ket{1}\big)=\alpha\ket{0}+\beta\ket{1}.
\end{align*}
Dans le dernier cas, avec les matrices,
$ZX=\begin{pmatrix}1&0\\0&-1\end{pmatrix}\begin{pmatrix}0&1\\1&0\end{pmatrix}
=\begin{pmatrix}0&1\\-1&0\end{pmatrix}$ et
$ZX\begin{pmatrix}-\beta\\ \alpha\end{pmatrix}=\begin{pmatrix}\alpha\\ \beta\end{pmatrix}$.
L'ordre compte : l'état de Bob est $XZ\ket{\varphi}$ ($Z\ket{\varphi}=\alpha\ket{0}-\beta\ket{1}$,
puis $X$ donne $\alpha\ket{1}-\beta\ket{0}$), et on l'annule avec
$(XZ)^{-1}=Z^{-1}X^{-1}=ZX$ : on applique d'abord $X$, puis $Z$, comme dans le circuit.
Dans les quatre cas, Bob retrouve exactement $\ket{\varphi}=\alpha\ket{0}+\beta\ket{1}$.

\subsubsection{Ce que Bob sait avant de recevoir les bits}

Tant que Bob ignore $(a,q)$, son qubit est dans le mélange des quatre kets $\ket{\phi_{aq}}$, chacun
avec la probabilité $\frac14$ :
\begin{align*}
  \rho_B&=\frac14\sum_{a,q}\ket{\phi_{aq}}\bra{\phi_{aq}}\\
  &=\frac14\Bigg[\begin{pmatrix}|\alpha|^2&\alpha\beta^*\\ \alpha^*\beta&|\beta|^2\end{pmatrix}
    +\begin{pmatrix}|\alpha|^2&-\alpha\beta^*\\ -\alpha^*\beta&|\beta|^2\end{pmatrix}
    +\begin{pmatrix}|\beta|^2&\alpha^*\beta\\ \alpha\beta^*&|\alpha|^2\end{pmatrix}
    +\begin{pmatrix}|\beta|^2&-\alpha^*\beta\\ -\alpha\beta^*&|\alpha|^2\end{pmatrix}\Bigg]\\
  &=\frac14\begin{pmatrix}2(|\alpha|^2+|\beta|^2)&0\\0&2(|\alpha|^2+|\beta|^2)\end{pmatrix}
   =\frac12\begin{pmatrix}1&0\\0&1\end{pmatrix}=\frac12\,\mathrm{Id}.
\end{align*}
Sans les deux bits, Bob a donc un qubit complètement mélangé, qui ne dépend pas de $\alpha$ ni de
$\beta$ : il n'a reçu aucune information sur $\ket{\varphi}$. C'est l'envoi classique de $(a,q)$
qui permet de la retrouver, donc rien ne voyage plus vite que la lumière.

\begin{remarque}[Bilan de la téléportation]
\begin{itemize}[nosep]
  \item Ressources : $1$ e-bit (la paire de Bell) et $2$ bits classiques pour transmettre
        \emph{un qubit} d'état quelconque, sans qu'Alice connaisse $\alpha$ ni $\beta$.
  \item Pas de clonage : après la mesure, $A$ et $Q$ sont dans l'état de base $\ket{a,q}$.
        Alice n'a plus $\ket{\varphi}$ : l'état a été \emph{déplacé}, pas copié. L'intrication
        entre $A$ et $B$ a été consommée.
  \item Les quatre corrections $I,Z,X,ZX$ sont des portes contrôlées par des bits classiques
        (section~\ref{sec:controlees}).
\end{itemize}
\end{remarque}

\begin{remarque}[Les gestes d'Alice forment une mesure de Bell]
Le CNOT de contrôle $Q$ sur $A$ suivi de $H$ sur $Q$ est l'\emph{inverse} du générateur de la
section~\ref{sec:circuit-bell} (les portes $H$ et CNOT sont leurs propres inverses, et l'ordre
s'inverse) : il envoie la base de Bell du couple $(A,Q)$ sur la base de calcul. Mesurer ensuite dans
la base de calcul revient donc à mesurer $(A,Q)$ dans la base de Bell. Les deux bits $(a,q)$
disent \emph{quel} état de Bell Alice a trouvé, ce qui explique les quatre corrections de Bob.
\end{remarque}

\begin{remarque}[Non-clonage]
Alice ne peut pas simplement copier $\ket{\varphi}$ \parencite{wootters1982}. Supposons une porte
unitaire $U$ telle que $U\big(\ket{\psi}\otimes\ket{0}\big)=\ket{\psi}\otimes\ket{\psi}$ pour tout
$\ket{\psi}$. Elle copie $\ket{0}$ et $\ket{1}$, mais par linéarité
\[
  U\big(\ket{+}\otimes\ket{0}\big)=\tfrac{1}{\sqrt2}\big(\ket{00}+\ket{11}\big)
  \neq\ket{+}\otimes\ket{+}=\tfrac12\big(\ket{00}+\ket{01}+\ket{10}+\ket{11}\big),
\]
ce qui contredit l'hypothèse. Le CNOT copie un bit classique, pas une superposition (remarque de la
section~\ref{sec:cnot}). La téléportation respecte cette interdiction : l'état est déplacé, il n'est
jamais présent en deux endroits.
\end{remarque}

% ---------------------------------------------------------------------
\subsection{Corrélations non locales : le jeu CHSH}\label{sec:chsh}
% ---------------------------------------------------------------------

Les états de Bell donnent des résultats corrélés en base $Z$ comme en base $X$
(section~\ref{sec:bell}). Mais deux enveloppes contenant la même carte, préparées à l'avance, donnent
aussi des résultats égaux : des corrélations dans une seule base ne prouvent pas l'intrication. Bell
\parencite{bell1964} a montré qu'on peut trancher en \emph{changeant de base de mesure} ; la version de
Clauser, Horne, Shimony et Holt \parencite{clauser1969} se présente comme un jeu.

\begin{definition}{Le jeu CHSH}
Alice et Bob, éloignés, ne communiquent plus une fois le jeu commencé. Un arbitre tire deux bits $x$ et
$y$ de façon uniforme et indépendante, donne $x$ à Alice et $y$ à Bob. Chacun répond par un bit : $a$ pour
Alice, $b$ pour Bob. Ils \emph{gagnent} si
\[
  a\oplus b=x\cdot y,
\]
c'est-à-dire si $a=b$ pour $(x,y)\in\{(0,0),(0,1),(1,0)\}$ et si $a\neq b$ pour $(x,y)=(1,1)$.
\end{definition}

\paragraph{Stratégie classique : au plus $\frac34$.} Une stratégie déterministe est $a=f(x)$,
$b=g(y)$. Si elle gagnait les quatre parties, on aurait
\[
  f(0)\oplus g(0)=0,\qquad f(0)\oplus g(1)=0,\qquad f(1)\oplus g(0)=0,\qquad f(1)\oplus g(1)=1 .
\]
En additionnant ces quatre égalités modulo $2$, chacune des valeurs $f(0),f(1),g(0),g(1)$ apparaît
deux fois : le membre de gauche vaut $0$, celui de droite vaut $1$. C'est absurde, donc au moins une
partie sur quatre est perdue et $P(\text{gain})\leq\frac34$. La stratégie «~toujours répondre $0$~»
atteint $\frac34$ (elle ne perd que pour $(1,1)$). Un hasard partagé ne change rien : il mélange des
stratégies déterministes, dont aucune ne dépasse $\frac34$.

\paragraph{Stratégie quantique.} Alice et Bob partagent avant le jeu une paire $\ket{\Phi^+}$ (un
e-bit), puis chacun mesure son qubit sans rien dire à l'autre. Pour un angle $\theta$, notons
\[
  \ket{\theta_+}=\cos\tfrac\theta2\,\ket{0}+\sin\tfrac\theta2\,\ket{1},
  \qquad
  \ket{\theta_-}=-\sin\tfrac\theta2\,\ket{0}+\cos\tfrac\theta2\,\ket{1}
\]
la base de mesure d'angle $\theta$ : c'est la direction $\theta$ dans le plan $xz$ de la sphère de
Bloch ($\theta=0$ : base $Z$ ; $\theta=\frac\pi2$ : base $X$). Le résultat vaut $0$ pour
$\ket{\theta_+}$ et $1$ pour $\ket{\theta_-}$. La stratégie est
\begin{itemize}[nosep]
  \item Alice : $\theta=0$ si $x=0$ (base $Z$), $\theta=\frac\pi2$ si $x=1$ (base $X$) ;
  \item Bob : $\theta=\frac\pi4$ si $y=0$, $\theta=-\frac\pi4$ si $y=1$ (bases propres de
        $\frac{Z\pm X}{\sqrt2}$).
\end{itemize}

\begin{figure}[!ht]
\centering
\adjustbox{max width=\linewidth}{%
\begin{tikzpicture}[>=Stealth, font=\footnotesize,
    joueur/.style={draw=insarouge, fill=insarouge!6, rounded corners=2pt, thick,
                   minimum width=2.1cm, minimum height=1.1cm, align=center},
    arb/.style={draw=insagris, fill=insagris!8, rounded corners=2pt, thick,
                minimum width=3.6cm, minimum height=0.8cm, align=center}]
  % --- le jeu ---
  \node[arb] (arb) at (3.2,2.7) {arbitre : $x,y$ au hasard};
  \node[joueur] (al) at (0,0) {Alice};
  \node[joueur] (bo) at (6.4,0) {Bob};
  \draw[->, thick, insagris] (arb.west) -| node[pos=0.2, above] {$x$} (al.north);
  \draw[->, thick, insagris] (arb.east) -| node[pos=0.2, above] {$y$} (bo.north);
  \draw[insarouge, thick, decorate, decoration={snake, amplitude=1.1pt, segment length=6pt}]
    (al.east) -- (bo.west);
  \node[insarouge, above=1pt] at (3.2,0.05) {$\ket{\Phi^+}$};
  \node[text=insagris, below=2pt, font=\scriptsize] at (3.2,-0.05) {paire intriquée};
  \draw[->, thick] (al.south) -- ++(0,-1.0) node[below] {$a$};
  \draw[->, thick] (bo.south) -- ++(0,-1.0) node[below] {$b$};
  \node[draw=insagris, rounded corners=2pt, inner sep=3pt] at (3.2,-1.75) {gain si $a\oplus b=x\cdot y$};
  % --- directions de mesure dans le plan xz de la sphère de Bloch ---
  \begin{scope}[shift={(11.9,0.3)}]
    \draw[insagris!70] (0,0) circle (1.7);
    \draw[insagris!50, dashed] (-1.7,0) -- (1.7,0);
    \draw[insagris!50, dashed] (0,-1.7) -- (0,1.7);
    \node[font=\scriptsize, text=insagris, anchor=north] at (0,-1.72) {$\ket{1}$};
    \node[font=\scriptsize, text=insagris, anchor=east] at (-1.72,-0.02) {$\ket{-}$};
    \draw[->, very thick, insarouge] (0,0) -- (90:1.7);
    \draw[->, very thick, insarouge] (0,0) -- (0:1.7);
    \draw[->, very thick, insagris] (0,0) -- (45:1.7);
    \draw[->, very thick, insagris] (0,0) -- (135:1.7);
    \node[insarouge, anchor=south, inner sep=2pt] at (0,1.72) {$x=0$ : $Z$ $(\ket{0})$};
    \node[insarouge, anchor=west, inner sep=2pt] at (1.72,0) {$x=1$ : $X$ $(\ket{+})$};
    \node[text=insagris, anchor=south west, align=left, inner sep=2pt] at (1.25,1.25) {$y=0$\\ $\frac{Z+X}{\sqrt2}$};
    \node[text=insagris, anchor=south east, align=right, inner sep=2pt] at (-1.25,1.25) {$y=1$\\ $\frac{Z-X}{\sqrt2}$};
    \draw[insagris] (0.85,0) arc (0:45:0.85);
    \node[font=\scriptsize, text=insagris] at (22.5:1.12) {$\frac\pi4$};
    \draw[insagris] (0,0.85) arc (90:135:0.85);
    \node[font=\scriptsize, text=insagris] at (112.5:1.12) {$\frac\pi4$};
    \draw[insagris] (0.6,0.6) arc (45:90:0.85);
    \node[font=\scriptsize, text=insagris] at (67.5:1.12) {$\frac\pi4$};
  \end{scope}
\end{tikzpicture}}
\caption{Le jeu CHSH (à gauche) et les quatre directions de mesure de la stratégie quantique (à
droite), dans le plan $xz$ de la sphère de Bloch : Alice en rouge, Bob en gris. Pour
$(x,y)=(0,0),(0,1),(1,0)$, les directions font un angle $\frac\pi4$ ; pour $(1,1)$, un angle
$\frac{3\pi}4$.}
\label{fig:chsh}
\end{figure}

\begin{resultat}{Corrélations de $\ket{\Phi^+}$ en bases tournées}
Si Alice mesure dans la base d'angle $\alpha$ et Bob dans la base d'angle $\beta$, alors
\[
  P(a=b)=\cos^2\frac{\alpha-\beta}{2}.
\]
\end{resultat}

\noindent\emph{Preuve.} Avec $c=\cos\frac\alpha2$ et $s=\sin\frac\alpha2$,
$\ket{\alpha_+\alpha_+}+\ket{\alpha_-\alpha_-}=(c^2+s^2)\ket{00}+(cs-sc)\big(\ket{01}+\ket{10}\big)+(s^2+c^2)\ket{11}
=\ket{00}+\ket{11}$ : on peut donc écrire $\ket{\Phi^+}=\frac{1}{\sqrt2}\big(\ket{\alpha_+\alpha_+}+\ket{\alpha_-\alpha_-}\big)$.
Si Alice obtient $a=0$ (probabilité $\frac12$), le qubit de Bob est dans $\ket{\alpha_+}$ et Bob obtient $b=0$
avec la probabilité $|\braket{\beta_+}{\alpha_+}|^2$, où $\braket{\beta_+}{\alpha_+}=\cos\frac\beta2\cos\frac\alpha2+\sin\frac\beta2\sin\frac\alpha2=\cos\frac{\alpha-\beta}{2}$.
Si Alice obtient $a=1$, Bob est dans $\ket{\alpha_-}$ et obtient $b=1$ avec la même probabilité, car
$\braket{\beta_-}{\alpha_-}=\braket{\beta_+}{\alpha_+}$. Donc
$P(a=b)=\frac12\cos^2\frac{\alpha-\beta}2+\frac12\cos^2\frac{\alpha-\beta}2$. \hfill$\square$

\medskip
\noindent Pour les quatre couples $(x,y)$ :
\begin{center}
\renewcommand{\arraystretch}{1.3}
\begin{tabular}{@{}c c c c c c@{}}
\toprule
$(x,y)$ & $\alpha$ & $\beta$ & $P(a=b)$ & condition de gain & $P(\text{gain})$\\
\midrule
$(0,0)$ & $0$ & $\frac\pi4$ & $\cos^2\frac\pi8$ & $a=b$ & $\cos^2\frac\pi8$\\
$(0,1)$ & $0$ & $-\frac\pi4$ & $\cos^2\frac\pi8$ & $a=b$ & $\cos^2\frac\pi8$\\
$(1,0)$ & $\frac\pi2$ & $\frac\pi4$ & $\cos^2\frac\pi8$ & $a=b$ & $\cos^2\frac\pi8$\\
$(1,1)$ & $\frac\pi2$ & $-\frac\pi4$ & $\cos^2\frac{3\pi}8=\sin^2\frac\pi8$ & $a\neq b$ &
  $1-\sin^2\frac\pi8=\cos^2\frac\pi8$\\
\bottomrule
\end{tabular}
\end{center}
Alice et Bob gagnent donc chaque partie avec la même probabilité :
\[
  P(\text{gain})=\cos^2\frac\pi8=\frac{2+\sqrt2}{4}\approx0{,}854>\frac34 .
\]

\begin{remarque}[Ce que montre le jeu CHSH]
\begin{itemize}[nosep]
  \item Aucune stratégie classique, même avec du hasard partagé, ne dépasse $\frac34$ : l'intrication
        produit des corrélations qu'aucun modèle «~local~» (résultats fixés à l'avance) ne reproduit.
        Les expériences sur des photons intriqués confirment la valeur quantique.
  \item $\cos^2\frac\pi8$ est le maximum quantique (\emph{borne de Tsirelson}, résultat admis,
        \cite{tsirelson1980}) : on ne peut pas gagner à tous les coups. Avec
        $E_{xy}=P(a=b)-P(a\neq b)=\cos(\alpha-\beta)$ et $S=E_{00}+E_{01}+E_{10}-E_{11}$, on a
        $P(\text{gain})=\frac12+\frac S8$ ; l'inégalité de Bell donne $|S|\leq2$ en classique, et la
        stratégie ci-dessus atteint $S=2\sqrt2$.
  \item Il n'y a pas de communication : la réponse de Bob est uniforme ($P(b=0)=\frac12$) quelle que
        soit la base choisie par Alice. L'intrication ne transmet pas d'information ; elle ne fait
        que corréler les résultats.
\end{itemize}
\end{remarque}

\clearpage
% ---------------------------------------------------------------------
\subsection{Récapitulatif du chapitre}\label{sec:recap5}
% ---------------------------------------------------------------------

\begin{center}
\adjustbox{max width=\linewidth}{%
\begin{tikzpicture}[>=Stealth, font=\footnotesize,
    boite/.style={draw=insarouge, fill=insarouge!6, rounded corners=2pt, thick,
                  minimum width=4.4cm, minimum height=1.45cm, align=center, inner sep=3pt},
    lien/.style={->, very thick, insagris},
    etiq/.style={font=\scriptsize\itshape, text=insagris, inner sep=2pt}]
  \node[boite] (mes) at (0,0)
    {\textbf{Mesure à deux qubits}\\ $P(ij)=|c_{ij}|^2$\\ {\scriptsize\S\,\ref{sec:mesure2}}};
  \node[boite] (sep) at (6.3,0)
    {\textbf{Séparable ou intriqué}\\ $D=c_{00}c_{11}-c_{01}c_{10}$\\ {\scriptsize\S\,\ref{sec:separable}}};
  \node[boite] (bel) at (12.6,0)
    {\textbf{États de Bell}\\ $\ket{\Phi^\pm}$, $\ket{\Psi^\pm}$ ; $\mathrm{CNOT}\,(I\otimes H)$\\
     {\scriptsize\S\,\ref{sec:bell}, \ref{sec:circuit-bell}}};
  \node[boite] (tra) at (12.6,-2.3)
    {\textbf{Trace partielle, entropies}\\ $\rho_A=\operatorname{Tr}_B\rho_{AB}$ ; $S(\rho)$, $I(A;B)$\\
     {\scriptsize\S\,\ref{sec:trace}, \ref{sec:info}}};
  \node[boite] (tel) at (6.3,-2.3)
    {\textbf{Téléportation}\\ e-bit + 2 bits classiques\\ {\scriptsize\S\,\ref{sec:teleportation}}};
  \node[boite] (chs) at (0,-2.3)
    {\textbf{Jeu CHSH}\\ $\frac34$ classique, $\cos^2\frac\pi8$ quantique\\ {\scriptsize\S\,\ref{sec:chsh}}};
  \draw[lien] (mes) -- (sep) node[etiq, midway, above=1pt] {corrélations};
  \draw[lien] (sep) -- (bel) node[etiq, midway, above=1pt] {cas extrême};
  \draw[lien] (bel) -- (tra) node[etiq, midway, right=1pt] {un qubit seul};
  \draw[lien] (tra) -- (tel) node[etiq, midway, above=1pt] {ressource};
  \draw[lien] (tel) -- (chs) node[etiq, midway, above=1pt] {non-localité};
\end{tikzpicture}}
\end{center}

\begin{aretenir}
\begin{itemize}
  \item Deux qubits : quatre amplitudes $c_{ij}$, $P(ij)=|c_{ij}|^2$. Mesurer $A$ supprime les termes
        incompatibles et renormalise : l'état de $B$ peut changer.
  \item Séparable $\Longleftrightarrow D=c_{00}c_{11}-c_{01}c_{10}=0$. Sinon l'état est intriqué et
        aucun qubit n'a d'état propre (pour $\ket{\Phi^+}$, $P(00)\neq P(A{=}0)P(B{=}0)$).
  \item États de Bell : base orthonormée d'états maximalement intriqués, états propres de
        $Z\otimes Z$ et $X\otimes X$ ; $\mathrm{CNOT}\,(I\otimes H)$ les produit à partir de la base de
        calcul. Des portes locales ne créent pas d'intrication.
  \item Qubit seul : $\rho_A=\operatorname{Tr}_B\rho_{AB}$, avec $\operatorname{Tr}\rho_A^2=1-2|D|^2$ pour un état pur.
        Pour $\ket{\Phi^+}$ : $\rho_A=\mathrm{Id}/2$ alors que le couple est pur ; $S(\rho_{AB})=0$,
        $S(\rho_A)=1$ bit, $I(A;B)=2$ bits, $H(A|B)=-1$ (impossible en classique).
  \item Téléportation : un e-bit et deux bits classiques ; $P(a,q)=\frac14$ ; Bob applique $X$ si $a=1$
        puis $Z$ si $q=1$. L'état est déplacé, pas copié ; sans les bits, $\rho_B=\mathrm{Id}/2$.
  \item Jeu CHSH : toute stratégie classique gagne avec une probabilité $\leq\frac34$ ; avec $\ket{\Phi^+}$
        et des bases tournées, $\cos^2\frac\pi8\approx0{,}854$ (borne de Tsirelson), sans aucune
        communication.
\end{itemize}
\end{aretenir}

\begin{center}
\renewcommand{\arraystretch}{1.2}
\begin{tabular}{@{}l l r@{}}
\toprule
\textbf{Notion} & \textbf{Formule} & \textbf{§} \\
\midrule
Mesure de deux qubits & $P(ij)=|c_{ij}|^2$, \quad $P(A{=}a)=P(a0)+P(a1)$ & \ref{sec:mesure2} \\
Séparabilité & $\ket{\Psi}=\ket{\Psi}_A\otimes\ket{\Psi}_B\ \Longleftrightarrow\ D=0$ & \ref{sec:separable} \\
États de Bell & $\ket{\Phi^\pm}=\dfrac{\ket{00}\pm\ket{11}}{\sqrt2}$, \quad
  $\ket{\Psi^\pm}=\dfrac{\ket{01}\pm\ket{10}}{\sqrt2}$ & \ref{sec:bell} \\
Trace partielle & $\operatorname{Tr}_B\big(\ket{a}\bra{a'}\otimes\ket{b}\bra{b'}\big)
  =\ket{a}\bra{a'}\,\braket{b'}{b}$ & \ref{sec:trace} \\
Qubit d'un état pur & $\rho_A=CC^\dagger$, \quad $\operatorname{Tr}\rho_A^2=1-2|D|^2$ & \ref{sec:trace} \\
Entropies & $S(\rho)=-\sum_i\lambda_i\log_2\lambda_i$, \quad $I(A;B)=H(A)+H(B)-H(A,B)$ & \ref{sec:info} \\
Téléportation & $P(a,q)=\frac14$, \quad correction $Z^{q}X^{a}$ (d'abord $X$) & \ref{sec:teleportation} \\
Singulet & $(U\otimes U)\ket{\Psi^-}=\det(U)\,\ket{\Psi^-}$ & \ref{sec:bell} \\
Jeu CHSH & $P_{\mathrm{cl}}\leq\frac34$, \quad $P_{\mathrm{quant}}=\cos^2\frac\pi8=\frac{2+\sqrt2}{4}$ & \ref{sec:chsh} \\
\bottomrule
\end{tabular}
\end{center}

\begin{savoirfaire}
\begin{itemize}
  \item Calculer $P(ij)$, les probabilités d'un seul qubit et l'état après la mesure de l'un d'eux ;
        tester la séparabilité avec $D$ et calculer la pureté de $\rho_A$ (exercice~\ref{ex:separable3}).
  \item Suivre un état dans le générateur de Bell, écrire $\rho_{AB}$ et reconnaître les corrélations
        d'un état de Bell dans deux bases (exercice~\ref{ex:bellX}).
  \item Calculer $H(A)$, $H(A,B)$, $I(A;B)$ et $H(A|B)$ pour une loi ou un état (exercice~\ref{ex:entropies}).
  \item Dérouler la téléportation pour un état donné et utiliser les états de Bell pour le codage
        superdense (exercices~\ref{ex:telenum} et~\ref{ex:superdense}).
  \item Prouver la borne $\frac34$ du jeu CHSH classique et calculer la probabilité de gain d'une
        stratégie quantique (section~\ref{sec:chsh}).
\end{itemize}
\end{savoirfaire}
 % 5. États quantiques, intrication, états de Bell
% =====================================================================
%  Chapitre 6 : systèmes multi-qubits et intrication
%  Fichier importé par main.tex avec % =====================================================================
%  Chapitre 6 : systèmes multi-qubits et intrication
%  Fichier importé par main.tex avec % =====================================================================
%  Chapitre 6 : systèmes multi-qubits et intrication
%  Fichier importé par main.tex avec \include{chapitre6}.
%  Comme chapitre4.tex et chapitre5.tex, il ouvre une NOUVELLE SECTION
%  (numéro 4). Pour le rattacher à la section précédente, remplacer
%  \section par \subsection et chaque \subsection par \subsubsection.
%  Pas de préambule : environnements (definition, resultat, remarque),
%  macros (\ii, \ee, \dd, \ket, \bra, \braket) et couleurs (insarouge,
%  insagris) sont définis dans main.tex.
% =====================================================================

% Macros de Dirac (déjà définies dans main.tex récent ; ce bloc évite l'erreur
% « Undefined control sequence \ket » si main.tex est une ancienne version).
\providecommand{\ket}[1]{|#1\rangle}
\providecommand{\bra}[1]{\langle #1|}
\providecommand{\braket}[2]{\langle #1|#2\rangle}

% =====================================================================
\section{Systèmes multi-qubits et intrication}\label{sec:multi}
% =====================================================================

Un ordinateur quantique utilise plusieurs qubits. Cette section décrit l'espace des états
de deux qubits (le produit tensoriel), la mesure, la distinction entre états séparables
et états intriqués, puis les opérateurs qui agissent sur deux qubits. Tous les calculs
sont détaillés ; les qubits sont notés $A$ et $B$.

% ---------------------------------------------------------------------
\subsection{L'espace de deux qubits}\label{sec:tenseur}
% ---------------------------------------------------------------------

Chaque qubit vit dans $\mathcal{H}^2=\mathbb{C}^2$. L'espace du système $(A,B)$ est le
\emph{produit tensoriel}
\[
  \mathcal{H}^2\otimes\mathcal{H}^2=\mathcal{H}^4=\mathbb{C}^4 ,
\]
de dimension $2\times2=4$. Pour $n$ qubits la dimension est $2^n$ :

\begin{center}
\renewcommand{\arraystretch}{1.3}
\begin{tabular}{@{}lccccc@{}}
\toprule
Nombre de qubits $n$ & $1$ & $2$ & $3$ & $10$ & $50$ \\
\midrule
Dimension $2^n$ & $2$ & $4$ & $8$ & $1024$ & $\approx1{,}1\times10^{15}$ \\
\bottomrule
\end{tabular}
\end{center}

\begin{definition}{Produit tensoriel de deux vecteurs}
\[
  \begin{pmatrix}a_1\\ a_2\end{pmatrix}\otimes\begin{pmatrix}b_1\\ b_2\end{pmatrix}
  =\begin{pmatrix}a_1\begin{pmatrix}b_1\\ b_2\end{pmatrix}\\[2ex]
                  a_2\begin{pmatrix}b_1\\ b_2\end{pmatrix}\end{pmatrix}
  =\begin{pmatrix}a_1b_1\\ a_1b_2\\ a_2b_1\\ a_2b_2\end{pmatrix}.
\]
Chaque composante du premier vecteur multiplie le second vecteur en entier. On note
$\ket{ij}_{AB}=\ket{i}_A\otimes\ket{j}_B$ : le premier chiffre est l'état de $A$, le second
celui de $B$.
\end{definition}

\subsubsection{Les quatre états de base}

On développe chacun des quatre produits avec $\ket{0}=\begin{pmatrix}1\\ 0\end{pmatrix}$
et $\ket{1}=\begin{pmatrix}0\\ 1\end{pmatrix}$ :
\begin{align*}
  \ket{0}\otimes\ket{0}&=\begin{pmatrix}1\\ 0\end{pmatrix}\otimes\begin{pmatrix}1\\ 0\end{pmatrix}
    =\begin{pmatrix}1\cdot\begin{pmatrix}1\\ 0\end{pmatrix}\\[1.5ex]0\cdot\begin{pmatrix}1\\ 0\end{pmatrix}\end{pmatrix}
    =\begin{pmatrix}1\\ 0\\ 0\\ 0\end{pmatrix}=\ket{00},\\[1.5ex]
  \ket{0}\otimes\ket{1}&=\begin{pmatrix}1\\ 0\end{pmatrix}\otimes\begin{pmatrix}0\\ 1\end{pmatrix}
    =\begin{pmatrix}1\cdot\begin{pmatrix}0\\ 1\end{pmatrix}\\[1.5ex]0\cdot\begin{pmatrix}0\\ 1\end{pmatrix}\end{pmatrix}
    =\begin{pmatrix}0\\ 1\\ 0\\ 0\end{pmatrix}=\ket{01},\\[1.5ex]
  \ket{1}\otimes\ket{0}&=\begin{pmatrix}0\\ 1\end{pmatrix}\otimes\begin{pmatrix}1\\ 0\end{pmatrix}
    =\begin{pmatrix}0\cdot\begin{pmatrix}1\\ 0\end{pmatrix}\\[1.5ex]1\cdot\begin{pmatrix}1\\ 0\end{pmatrix}\end{pmatrix}
    =\begin{pmatrix}0\\ 0\\ 1\\ 0\end{pmatrix}=\ket{10},\\[1.5ex]
  \ket{1}\otimes\ket{1}&=\begin{pmatrix}0\\ 1\end{pmatrix}\otimes\begin{pmatrix}0\\ 1\end{pmatrix}
    =\begin{pmatrix}0\cdot\begin{pmatrix}0\\ 1\end{pmatrix}\\[1.5ex]1\cdot\begin{pmatrix}0\\ 1\end{pmatrix}\end{pmatrix}
    =\begin{pmatrix}0\\ 0\\ 0\\ 1\end{pmatrix}=\ket{11}.
\end{align*}
Le vecteur $\ket{ij}$ a son $1$ en position $2i+j$ (en comptant à partir de $0$) : c'est
la lecture de $ij$ en binaire. Les quatre vecteurs forment la base de calcul de
$\mathcal{H}^4$.

\subsubsection{Produit scalaire et orthonormalité}

\begin{resultat}{Produit scalaire d'un produit tensoriel}
\[
  \big(\bra{a}\otimes\bra{b}\big)\big(\ket{c}\otimes\ket{d}\big)=\braket{a}{c}\,\braket{b}{d}.
\]
\end{resultat}

\emph{Démonstration.} Les composantes de $\ket{a}\otimes\ket{b}$ sont $a_ib_j$, celles de
$\ket{c}\otimes\ket{d}$ sont $c_id_j$. Le produit scalaire est
$\sum_{i,j}(a_ib_j)^*c_id_j=\big(\sum_ia_i^*c_i\big)\big(\sum_jb_j^*d_j\big)=\braket{a}{c}\braket{b}{d}$.
\hfill$\square$

On en déduit $\braket{ij}{kl}=\braket{i}{k}\braket{j}{l}$. Les seize produits scalaires
des états de base sont :
\begin{center}
\renewcommand{\arraystretch}{1.3}
\begin{tabular}{@{}c|cccc@{}}
 & $\ket{00}$ & $\ket{01}$ & $\ket{10}$ & $\ket{11}$ \\
\hline
$\bra{00}$ & $1\cdot1=1$ & $1\cdot0=0$ & $0\cdot1=0$ & $0\cdot0=0$ \\
$\bra{01}$ & $1\cdot0=0$ & $1\cdot1=1$ & $0\cdot0=0$ & $0\cdot1=0$ \\
$\bra{10}$ & $0\cdot1=0$ & $0\cdot0=0$ & $1\cdot1=1$ & $1\cdot0=0$ \\
$\bra{11}$ & $0\cdot0=0$ & $0\cdot1=0$ & $1\cdot0=0$ & $1\cdot1=1$ \\
\end{tabular}
\end{center}
(chaque case est $\braket{i}{k}\cdot\braket{j}{l}$). On lit $\braket{ij}{kl}=\delta_{ik}\delta_{jl}$ :
la base est \emph{orthonormée}.

\subsubsection{Produit tensoriel d'états quelconques}

Le produit tensoriel est linéaire en chaque facteur :
$(\lambda\ket{a}+\mu\ket{a'})\otimes\ket{b}=\lambda\ket{a}\otimes\ket{b}+\mu\ket{a'}\otimes\ket{b}$,
et de même pour le second facteur. Pour $\alpha\ket{0}+\beta\ket{1}$ et
$\gamma\ket{0}+\delta\ket{1}$ :
\begin{align*}
  \big(\alpha\ket{0}+\beta\ket{1}\big)\otimes\big(\gamma\ket{0}+\delta\ket{1}\big)
  &=\alpha\ket{0}\otimes\big(\gamma\ket{0}+\delta\ket{1}\big)
    +\beta\ket{1}\otimes\big(\gamma\ket{0}+\delta\ket{1}\big)\\
  &=\alpha\gamma\ket{00}+\alpha\delta\ket{01}+\beta\gamma\ket{10}+\beta\delta\ket{11}.
\end{align*}
Par les composantes, $\begin{pmatrix}\alpha\\ \beta\end{pmatrix}\otimes\begin{pmatrix}\gamma\\ \delta\end{pmatrix}
=\begin{pmatrix}\alpha\gamma\\ \alpha\delta\\ \beta\gamma\\ \beta\delta\end{pmatrix}$ : c'est le même
résultat. La norme est conservée :
\[
  |\alpha\gamma|^2+|\alpha\delta|^2+|\beta\gamma|^2+|\beta\delta|^2
  =\big(|\alpha|^2+|\beta|^2\big)\big(|\gamma|^2+|\delta|^2\big)=1\cdot1=1 .
\]

\paragraph{Exemples.}
\begin{align*}
  \ket{+}\otimes\ket{+}&=\tfrac12\big(\ket{0}+\ket{1}\big)\otimes\big(\ket{0}+\ket{1}\big)
    =\tfrac12\big(\ket{00}+\ket{01}+\ket{10}+\ket{11}\big)
    =\tfrac12\begin{pmatrix}1\\ 1\\ 1\\ 1\end{pmatrix},\\[1ex]
  \ket{+}\otimes\ket{-}&=\tfrac12\big(\ket{0}+\ket{1}\big)\otimes\big(\ket{0}-\ket{1}\big)
    =\tfrac12\big(\ket{00}-\ket{01}+\ket{10}-\ket{11}\big)
    =\tfrac12\begin{pmatrix}1\\ -1\\ 1\\ -1\end{pmatrix},\\[1ex]
  \ket{0}\otimes\ket{+}&=\tfrac{1}{\sqrt2}\big(\ket{00}+\ket{01}\big)
    =\tfrac{1}{\sqrt2}\begin{pmatrix}1\\ 1\\ 0\\ 0\end{pmatrix},\qquad
  \ket{+}\otimes\ket{0}=\tfrac{1}{\sqrt2}\big(\ket{00}+\ket{10}\big)
    =\tfrac{1}{\sqrt2}\begin{pmatrix}1\\ 0\\ 1\\ 0\end{pmatrix}.
\end{align*}
Les deux derniers vecteurs sont différents : \emph{l'ordre des facteurs compte}, le premier
est toujours celui de $A$.

% ---------------------------------------------------------------------
\subsection{États à deux qubits et mesure}\label{sec:mesure2}
% ---------------------------------------------------------------------

Un état quelconque de $\mathcal{H}^4$ se décompose sur la base de calcul :
\[
  \ket{\Psi}=c_{00}\ket{00}+c_{01}\ket{01}+c_{10}\ket{10}+c_{11}\ket{11}
  =\begin{pmatrix}c_{00}\\ c_{01}\\ c_{10}\\ c_{11}\end{pmatrix},
  \qquad \sum_{i,j}|c_{ij}|^2=1 .
\]

\begin{resultat}{Mesure des deux qubits dans la base de calcul}
\[
  P(ij)=|\braket{ij}{\Psi}|^2=|c_{ij}|^2 .
\]
Pour un seul qubit : $P(A{=}0)=P(00)+P(01)$, $P(A{=}1)=P(10)+P(11)$,
$P(B{=}0)=P(00)+P(10)$, $P(B{=}1)=P(01)+P(11)$.
\end{resultat}

En effet $\braket{ij}{\Psi}=\sum_{k,l}c_{kl}\braket{ij}{kl}=\sum_{k,l}c_{kl}\delta_{ik}\delta_{jl}=c_{ij}$.

\paragraph{Exemple 1.} $\ket{\Psi}=\frac{1}{\sqrt2}\ket{00}+\frac{1}{\sqrt2}\ket{01}
=\big(\tfrac{1}{\sqrt2},\tfrac{1}{\sqrt2},0,0\big)^{\mathsf T}$. Il est normalisé :
$\frac12+\frac12+0+0=1$. Les probabilités jointes sont
$P(00)=P(01)=\frac12$ et $P(10)=P(11)=0$. Pour chaque qubit :
\[
  P(A{=}0)=\tfrac12+\tfrac12=1,\quad P(A{=}1)=0+0=0,\quad
  P(B{=}0)=\tfrac12+0=\tfrac12,\quad P(B{=}1)=\tfrac12+0=\tfrac12 .
\]
$A$ donne toujours $0$ ; $B$ donne $0$ ou $1$ au hasard.

\paragraph{Exemple 2 (mesure de $A$ seul).} Soit
$\ket{\Psi}=\frac12\ket{00}+\frac12\ket{01}+\frac{1}{\sqrt2}\ket{11}$. Normalisation :
$\frac14+\frac14+0+\frac12=1$. Probabilités jointes : $P(00)=\frac14$, $P(01)=\frac14$,
$P(10)=0$, $P(11)=\frac12$. Donc
\[
  P(A{=}0)=\tfrac14+\tfrac14=\tfrac12,\quad P(A{=}1)=0+\tfrac12=\tfrac12,\quad
  P(B{=}0)=\tfrac14+0=\tfrac14,\quad P(B{=}1)=\tfrac14+\tfrac12=\tfrac34 .
\]

\begin{remarque}[Effondrement après la mesure de $A$]
Si $A$ donne le résultat $a$, on ne garde que les termes de $\ket{\Psi}$ dont le premier
chiffre vaut $a$, puis on renormalise :
\[
  \ket{\Psi}\ \to\ \frac{c_{a0}\ket{a0}+c_{a1}\ket{a1}}{\sqrt{|c_{a0}|^2+|c_{a1}|^2}},
  \qquad\text{probabilité }|c_{a0}|^2+|c_{a1}|^2 .
\]
\end{remarque}

Pour l'exemple 2 :
\begin{itemize}[nosep]
  \item $A=0$ (probabilité $\frac12$) : on garde $\frac12\ket{00}+\frac12\ket{01}$, de norme
        au carré $\frac12$ ; renormalisé, $\frac{1}{\sqrt2}\big(\ket{00}+\ket{01}\big)
        =\ket{0}\otimes\ket{+}$. Ensuite $B$ vaut $0$ ou $1$ avec la probabilité $\frac12$.
  \item $A=1$ (probabilité $\frac12$) : on garde $\frac{1}{\sqrt2}\ket{11}$, de norme au carré
        $\frac12$ ; renormalisé, $\ket{11}=\ket{1}\otimes\ket{1}$. Ensuite $B$ vaut $1$
        avec certitude.
\end{itemize}
Contrôle : $P(B{=}0)=\frac12\cdot\frac12+\frac12\cdot0=\frac14$, comme calculé plus haut.
Ici la mesure de $A$ \emph{modifie} l'état de $B$ ($\ket{+}$ ou $\ket{1}$) : c'est la
signature d'un état intriqué (section suivante).

% ---------------------------------------------------------------------
\subsection{États séparables et états intriqués}\label{sec:separable}
% ---------------------------------------------------------------------

\begin{definition}{Séparable ou intriqué}
Un état $\ket{\Psi}_{AB}$ est \emph{séparable} s'il s'écrit comme un produit
$\ket{\Psi}_{AB}=\ket{\Psi}_A\otimes\ket{\Psi}_B$. Sinon il est \emph{intriqué} : on ne
peut pas attribuer d'état propre à chaque qubit.
\end{definition}

\paragraph{Exemple séparable (exemple 1 ci-dessus).} Par linéarité du produit tensoriel :
\begin{align*}
  \ket{\Psi}&=\tfrac{1}{\sqrt2}\ket{00}_{AB}+\tfrac{1}{\sqrt2}\ket{01}_{AB}
    =\tfrac{1}{\sqrt2}\Big(\ket{0}_A\otimes\ket{0}_B+\ket{0}_A\otimes\ket{1}_B\Big)\\
  &=\ket{0}_A\otimes\Big(\tfrac{1}{\sqrt2}\big(\ket{0}+\ket{1}\big)\Big)_B
    =\ket{0}_A\otimes\ket{+}_B ,
\end{align*}
donc $\ket{\Psi}_A=\ket{0}$ et $\ket{\Psi}_B=\ket{+}$. Contrôle par les composantes :
$\begin{pmatrix}1\\ 0\end{pmatrix}\otimes\tfrac{1}{\sqrt2}\begin{pmatrix}1\\ 1\end{pmatrix}
=\tfrac{1}{\sqrt2}\begin{pmatrix}1\cdot1\\ 1\cdot1\\ 0\cdot1\\ 0\cdot1\end{pmatrix}
=\tfrac{1}{\sqrt2}\begin{pmatrix}1\\ 1\\ 0\\ 0\end{pmatrix}$.

\subsubsection{Un critère de séparabilité}

Un produit quelconque a pour coefficients $c_{00}=ac$, $c_{01}=ad$, $c_{10}=bc$,
$c_{11}=bd$ (formule de la section~\ref{sec:tenseur}, avec $a,b$ pour $A$ et $c,d$ pour $B$), donc
\[
  c_{00}c_{11}=(ac)(bd)=abcd=(ad)(bc)=c_{01}c_{10}.
\]

\begin{resultat}{Test de séparabilité (deux qubits)}
$\ket{\Psi}=\sum c_{ij}\ket{ij}$ est séparable si et seulement si
\[
  D=c_{00}c_{11}-c_{01}c_{10}=0 .
\]
\end{resultat}

\emph{Démonstration.} ($\Rightarrow$) Calcul ci-dessus. ($\Leftarrow$) On regroupe les
termes selon l'état de $A$ :
\[
  \ket{\Psi}=\ket{0}\otimes\big(c_{00}\ket{0}+c_{01}\ket{1}\big)+\ket{1}\otimes\big(c_{10}\ket{0}+c_{11}\ket{1}\big).
\]
Si $D=0$, les lignes $(c_{00},c_{01})$ et $(c_{10},c_{11})$ sont proportionnelles. Si
$(c_{00},c_{01})\neq(0,0)$, il existe $\lambda$ avec $(c_{10},c_{11})=\lambda(c_{00},c_{01})$
et $\ket{\Psi}=\big(\ket{0}+\lambda\ket{1}\big)\otimes\big(c_{00}\ket{0}+c_{01}\ket{1}\big)$.
Si $(c_{00},c_{01})=(0,0)$, alors $\ket{\Psi}=\ket{1}\otimes\big(c_{10}\ket{0}+c_{11}\ket{1}\big)$.
Dans les deux cas $\ket{\Psi}$ est un produit ; on renormalise chaque facteur.
\hfill$\square$

\subsubsection{Application à plusieurs états}

\begin{center}
\renewcommand{\arraystretch}{1.5}
\begin{tabular}{@{}lccl@{}}
\toprule
\textbf{État} & $(c_{00},c_{01},c_{10},c_{11})$ & $D$ & \textbf{Nature} \\
\midrule
$\ket{00}$ & $(1,0,0,0)$ & $1\cdot0-0\cdot0=0$ & séparable \\
$\ket{+}\otimes\ket{+}$ & $\frac12(1,1,1,1)$ & $\frac14-\frac14=0$ & séparable \\
$\ket{+}\otimes\ket{-}$ & $\frac12(1,-1,1,-1)$ & $-\frac14-(-\frac14)=0$ & séparable \\
$\frac{1}{\sqrt2}(\ket{00}+\ket{01})$ & $\frac{1}{\sqrt2}(1,1,0,0)$ & $0-\frac12\cdot0=0$ & séparable \\
$\ket{\Phi^+}=\frac{1}{\sqrt2}(\ket{00}+\ket{11})$ & $\frac{1}{\sqrt2}(1,0,0,1)$ & $\frac12-0=\frac12$ & intriqué \\
$\ket{\Phi^-}=\frac{1}{\sqrt2}(\ket{00}-\ket{11})$ & $\frac{1}{\sqrt2}(1,0,0,-1)$ & $-\frac12-0=-\frac12$ & intriqué \\
$\ket{\Psi^+}=\frac{1}{\sqrt2}(\ket{01}+\ket{10})$ & $\frac{1}{\sqrt2}(0,1,1,0)$ & $0-\frac12=-\frac12$ & intriqué \\
$\ket{\Psi^-}=\frac{1}{\sqrt2}(\ket{01}-\ket{10})$ & $\frac{1}{\sqrt2}(0,1,-1,0)$ & $0-(-\frac12)=\frac12$ & intriqué \\
$\frac12\ket{00}+\frac12\ket{01}+\frac{1}{\sqrt2}\ket{11}$ & $(\frac12,\frac12,0,\frac{1}{\sqrt2})$ & $\frac{1}{2\sqrt2}-0$ & intriqué \\
\bottomrule
\end{tabular}
\end{center}

Les quatre états de Bell $\ket{\Phi^\pm}$, $\ket{\Psi^\pm}$ sont \emph{orthonormés} :
chacun est de norme $\frac12+\frac12=1$ ; $\ket{\Phi^\pm}$ n'ont des composantes que sur
$\ket{00}$ et $\ket{11}$, $\ket{\Psi^\pm}$ que sur $\ket{01}$ et $\ket{10}$, donc
$\braket{\Phi^\pm}{\Psi^\pm}=0$, et
$\braket{\Phi^+}{\Phi^-}=\frac12(1-1)=0$, $\braket{\Psi^+}{\Psi^-}=\frac12(1-1)=0$.
Ils forment une base de $\mathcal{H}^4$ constituée uniquement d'états intriqués.

\paragraph{Preuve directe que $\ket{\Phi^+}$ est intriqué.} Supposons
$\ket{\Phi^+}=(a\ket{0}+b\ket{1})\otimes(c\ket{0}+d\ket{1})$. En identifiant les
composantes : $ac=\frac{1}{\sqrt2}$, $ad=0$, $bc=0$, $bd=\frac{1}{\sqrt2}$. La condition
$ad=0$ impose $a=0$ (impossible car $ac\neq0$) ou $d=0$ (impossible car $bd\neq0$).
Contradiction : aucun produit ne convient.

\paragraph{Corrélations.} Pour un produit $\ket{\Psi}_A\otimes\ket{\Psi}_B$ avec
$\ket{\Psi}_A=a_0\ket{0}+a_1\ket{1}$ et $\ket{\Psi}_B=b_0\ket{0}+b_1\ket{1}$, on a
$P(ij)=|a_ib_j|^2=|a_i|^2|b_j|^2=P(A{=}i)\,P(B{=}j)$ : les deux qubits sont
\emph{indépendants}. Vérification avec $\ket{+}\otimes\ket{+}$ : $P(00)=\frac14=\frac12\cdot\frac12$.
Pour $\ket{\Phi^+}$ : $P(00)=P(11)=\frac12$ et $P(01)=P(10)=0$, donc
$P(A{=}0)=P(B{=}0)=\frac12$, mais
\[
  P(00)=\tfrac12\neq\tfrac12\cdot\tfrac12=P(A{=}0)\,P(B{=}0) .
\]
Les deux qubits donnent toujours le même résultat, aléatoire : ils sont corrélés.

\paragraph{État d'un qubit pris seul.} On prend la matrice de densité
$\rho=\ket{\Psi}\bra{\Psi}$ ($4\times4$, lignes et colonnes ordonnées $00,01,10,11$) et on
trace sur $B$ : $(\rho_A)_{ij}=\rho_{(i0),(j0)}+\rho_{(i1),(j1)}$.

\emph{État de Bell $\ket{\Phi^+}$.}
\[
  \rho=\frac12\begin{pmatrix}1 & 0 & 0 & 1\\ 0 & 0 & 0 & 0\\ 0 & 0 & 0 & 0\\ 1 & 0 & 0 & 1\end{pmatrix},
  \qquad
  \begin{aligned}
    (\rho_A)_{00}&=\rho_{00,00}+\rho_{01,01}=\tfrac12+0=\tfrac12,\\
    (\rho_A)_{01}&=\rho_{00,10}+\rho_{01,11}=0+0=0,\\
    (\rho_A)_{11}&=\rho_{10,10}+\rho_{11,11}=0+\tfrac12=\tfrac12 .
  \end{aligned}
\]
Donc $\rho_A=\frac12\mathrm{Id}$ : état \emph{mixte} ($\operatorname{Tr}\rho_A^2=\frac12$), au centre
de la boule de Bloch (section~\ref{sec:bloch}).

\emph{État séparable $\ket{0}\otimes\ket{+}$.}
\[
  \rho=\frac12\begin{pmatrix}1 & 1 & 0 & 0\\ 1 & 1 & 0 & 0\\ 0 & 0 & 0 & 0\\ 0 & 0 & 0 & 0\end{pmatrix},
  \qquad
  \begin{aligned}
    (\rho_A)_{00}&=\rho_{00,00}+\rho_{01,01}=\tfrac12+\tfrac12=1,\\
    (\rho_A)_{01}&=\rho_{00,10}+\rho_{01,11}=0+0=0,\\
    (\rho_A)_{11}&=\rho_{10,10}+\rho_{11,11}=0+0=0 .
  \end{aligned}
\]
Donc $\rho_A=\ket{0}\bra{0}$ : état \emph{pur}, celui du facteur $\ket{\Psi}_A$.

Un qubit d'un état séparable garde son propre état pur ; un qubit d'un état de Bell est
complètement aléatoire alors que le couple est parfaitement défini.

% ---------------------------------------------------------------------
\subsection{Opérateurs sur deux qubits}\label{sec:op2}
% ---------------------------------------------------------------------

\begin{definition}{Produit tensoriel de deux matrices}
\[
  A\otimes B=\begin{pmatrix}A_{11} & A_{12}\\ A_{21} & A_{22}\end{pmatrix}\otimes B
  =\begin{pmatrix}A_{11}B & A_{12}B\\ A_{21}B & A_{22}B\end{pmatrix}
  \qquad(\text{matrice }4\times4).
\]
\end{definition}

\begin{resultat}{Action sur un produit}
\[
  (A\otimes B)\big(\ket{u}\otimes\ket{v}\big)=A\ket{u}\otimes B\ket{v}.
\]
\end{resultat}

\emph{Démonstration.} Avec $\ket{u}=(u_1,u_2)^{\mathsf T}$, le vecteur $\ket{u}\otimes\ket{v}$ a
pour blocs $u_k\ket{v}$ ($k=1,2$). Le bloc de lignes $i$ de $(A\otimes B)(\ket{u}\otimes\ket{v})$ est
donc
\[
  \sum_kA_{ik}B\,\big(u_k\ket{v}\big)=\Big(\sum_kA_{ik}u_k\Big)B\ket{v}=(A\ket{u})_i\,B\ket{v},
\]
qui est le bloc $i$ de $A\ket{u}\otimes B\ket{v}$. \hfill$\square$

\subsubsection{Portes agissant sur un seul des deux qubits}

On note $I=\mathrm{Id}$ ($2\times2$). $X\otimes I$ agit sur $A$, $I\otimes X$ agit sur $B$ :
\begin{align*}
  I\otimes X&=\begin{pmatrix}1\cdot X & 0\cdot X\\ 0\cdot X & 1\cdot X\end{pmatrix}
    =\begin{pmatrix}0 & 1 & 0 & 0\\ 1 & 0 & 0 & 0\\ 0 & 0 & 0 & 1\\ 0 & 0 & 1 & 0\end{pmatrix},\\[1ex]
  X\otimes I&=\begin{pmatrix}0\cdot I & 1\cdot I\\ 1\cdot I & 0\cdot I\end{pmatrix}
    =\begin{pmatrix}0 & 0 & 1 & 0\\ 0 & 0 & 0 & 1\\ 1 & 0 & 0 & 0\\ 0 & 1 & 0 & 0\end{pmatrix}.
\end{align*}
Sur $\ket{00}=(1,0,0,0)^{\mathsf T}$ (première colonne de la matrice) :
$(I\otimes X)\ket{00}=(0,1,0,0)^{\mathsf T}=\ket{01}$ et
$(X\otimes I)\ket{00}=(0,0,1,0)^{\mathsf T}=\ket{10}$. Avec la règle :
$(I\otimes X)\ket{00}=I\ket{0}\otimes X\ket{0}=\ket{0}\otimes\ket{1}=\ket{01}$.

Pour la porte de Hadamard, avec $H=\frac{1}{\sqrt2}\begin{pmatrix}1 & 1\\ 1 & -1\end{pmatrix}$ :
\[
  H\otimes I=\frac{1}{\sqrt2}\begin{pmatrix}1\cdot I & 1\cdot I\\ 1\cdot I & -1\cdot I\end{pmatrix}
  =\frac{1}{\sqrt2}\begin{pmatrix}1 & 0 & 1 & 0\\ 0 & 1 & 0 & 1\\ 1 & 0 & -1 & 0\\ 0 & 1 & 0 & -1\end{pmatrix}.
\]
Ses colonnes donnent l'image des quatre états de base :
\begin{align*}
  (H\otimes I)\ket{00}&=\tfrac{1}{\sqrt2}(1,0,1,0)^{\mathsf T}=\ket{+}\otimes\ket{0},&
  (H\otimes I)\ket{01}&=\tfrac{1}{\sqrt2}(0,1,0,1)^{\mathsf T}=\ket{+}\otimes\ket{1},\\
  (H\otimes I)\ket{10}&=\tfrac{1}{\sqrt2}(1,0,-1,0)^{\mathsf T}=\ket{-}\otimes\ket{0},&
  (H\otimes I)\ket{11}&=\tfrac{1}{\sqrt2}(0,1,0,-1)^{\mathsf T}=\ket{-}\otimes\ket{1}.
\end{align*}

\subsubsection{Hadamard sur les deux qubits}

\[
  H\otimes H=\frac12\begin{pmatrix}1\cdot H' & 1\cdot H'\\ 1\cdot H' & -1\cdot H'\end{pmatrix},
  \quad H'=\begin{pmatrix}1 & 1\\ 1 & -1\end{pmatrix},
  \qquad
  H\otimes H=\frac12\begin{pmatrix}1 & 1 & 1 & 1\\ 1 & -1 & 1 & -1\\ 1 & 1 & -1 & -1\\ 1 & -1 & -1 & 1\end{pmatrix}.
\]
La première colonne donne $(H\otimes H)\ket{00}=\frac12(1,1,1,1)^{\mathsf T}=\ket{+}\otimes\ket{+}$
(cohérent avec $H\ket{0}\otimes H\ket{0}$) : les quatre résultats $00,01,10,11$ sont
équiprobables, chacun avec la probabilité $\frac14$. Sur $n$ qubits, $H^{\otimes n}\ket{0\ldots0}$
donne une superposition uniforme de $2^n$ états.

\subsubsection{La porte CNOT et les états de Bell}

\begin{definition}{Porte CNOT (NOT contrôlé)}
Le qubit $A$ est le \emph{contrôle}, $B$ la \emph{cible} : $B$ est inversé si $A=1$.
\[
  \mathrm{CNOT}=\ket{0}\bra{0}\otimes I+\ket{1}\bra{1}\otimes X
  =\begin{pmatrix}I & 0\\ 0 & X\end{pmatrix}
  =\begin{pmatrix}1 & 0 & 0 & 0\\ 0 & 1 & 0 & 0\\ 0 & 0 & 0 & 1\\ 0 & 0 & 1 & 0\end{pmatrix}.
\]
\end{definition}

Sa matrice échange les composantes $10$ et $11$ :
$\mathrm{CNOT}\ket{00}=\ket{00}$, $\mathrm{CNOT}\ket{01}=\ket{01}$,
$\mathrm{CNOT}\ket{10}=\ket{11}$, $\mathrm{CNOT}\ket{11}=\ket{10}$, soit
$\ket{i,j}\to\ket{i,i\oplus j}$ (addition modulo $2$).

On applique $H\otimes I$ puis $\mathrm{CNOT}$ à $\ket{00}$ (figure~\ref{fig:bell}) :
\[
  \ket{00}\ \xrightarrow{\ H\otimes I\ }\ \tfrac{1}{\sqrt2}\begin{pmatrix}1\\ 0\\ 1\\ 0\end{pmatrix}
  =\tfrac{1}{\sqrt2}\big(\ket{00}+\ket{10}\big)
  \ \xrightarrow{\ \mathrm{CNOT}\ }\ \tfrac{1}{\sqrt2}\begin{pmatrix}1\\ 0\\ 0\\ 1\end{pmatrix}
  =\tfrac{1}{\sqrt2}\big(\ket{00}+\ket{11}\big)=\ket{\Phi^+}.
\]

\begin{figure}[!ht]
\centering
\begin{tikzpicture}[>=Stealth, font=\small]
  % fils : B en haut (dernier chiffre du ket), A en bas (premier chiffre)
  \draw (0,1) -- (8.6,1);
  \draw (0,0) -- (8.6,0);
  \node[anchor=east] at (0,1) {$\ket{0}_B$};
  \node[anchor=east] at (0,0) {$\ket{0}_A$};
  % Hadamard sur A
  \node[draw=insarouge, fill=insarouge!8, minimum size=7mm, thick] at (2.2,0) {$H$};
  % CNOT : contrôle sur A (bas), cible sur B (haut)
  \draw[insarouge, thick] (5.4,0) -- (5.4,1.3);
  \fill[insarouge] (5.4,0) circle (2.6pt);
  \draw[insarouge, thick, fill=white] (5.4,1) circle (0.3);
  \draw[insarouge, thick] (5.1,1) -- (5.7,1);
  \draw[insarouge, thick] (5.4,0.7) -- (5.4,1.3);
  % accolade de sortie
  \draw[insagris, thick] (8.75,1.2) -- (8.9,1.2) -- (8.9,-0.2) -- (8.75,-0.2);
  \node[anchor=west] at (8.9,0.5) {$\ket{\Phi^+}_{AB}$};
  % états intermédiaires
  \foreach \x in {0.6,3.8,7.2}
    \draw[dashed, insagris!70] (\x,1.35) -- (\x,-0.95);
  \node[anchor=north, font=\footnotesize] at (0.6,-0.95) {$\ket{00}$};
  \node[anchor=north, font=\footnotesize] at (3.8,-0.95) {$\tfrac{1}{\sqrt2}(\ket{00}+\ket{10})$};
  \node[anchor=north, font=\footnotesize] at (7.2,-0.95) {$\tfrac{1}{\sqrt2}(\ket{00}+\ket{11})$};
\end{tikzpicture}
\caption{Circuit qui produit l'état de Bell $\ket{\Phi^+}$ : une porte de Hadamard sur $A$,
puis un CNOT ($A$ contrôle, $B$ cible). Les états intermédiaires sont indiqués sous les
traits pointillés. Convention : les fils sont dessinés dans l'ordre inverse de l'écriture
du ket $\ket{ij}_{AB}$, si bien que le premier chiffre (celui de $A$) est sur le fil du bas.}
\label{fig:bell}
\end{figure}

Les quatre états de base d'entrée donnent les quatre états de Bell :
\begin{center}
\renewcommand{\arraystretch}{1.6}
\begin{tabular}{@{}cccc@{}}
\toprule
Entrée & après $H\otimes I$ & après CNOT & Sortie \\
\midrule
$\ket{00}$ & $\frac{1}{\sqrt2}(\ket{00}+\ket{10})$ & $\frac{1}{\sqrt2}(\ket{00}+\ket{11})$ & $\ket{\Phi^+}$ \\
$\ket{01}$ & $\frac{1}{\sqrt2}(\ket{01}+\ket{11})$ & $\frac{1}{\sqrt2}(\ket{01}+\ket{10})$ & $\ket{\Psi^+}$ \\
$\ket{10}$ & $\frac{1}{\sqrt2}(\ket{00}-\ket{10})$ & $\frac{1}{\sqrt2}(\ket{00}-\ket{11})$ & $\ket{\Phi^-}$ \\
$\ket{11}$ & $\frac{1}{\sqrt2}(\ket{01}-\ket{11})$ & $\frac{1}{\sqrt2}(\ket{01}-\ket{10})$ & $\ket{\Psi^-}$ \\
\bottomrule
\end{tabular}
\end{center}
Sous forme matricielle, le produit $\mathrm{CNOT}\cdot(H\otimes I)$ échange les lignes $3$ et $4$ de
$H\otimes I$ :
\[
  \mathrm{CNOT}\,(H\otimes I)=\frac{1}{\sqrt2}\begin{pmatrix}1 & 0 & 1 & 0\\ 0 & 1 & 0 & 1\\ 0 & 1 & 0 & -1\\ 1 & 0 & -1 & 0\end{pmatrix},
\]
et ses quatre colonnes sont $\ket{\Phi^+}$, $\ket{\Psi^+}$, $\ket{\Phi^-}$, $\ket{\Psi^-}$.

\begin{remarque}[Porte intriquante]
Les états d'entrée $\ket{ij}$ sont séparables ($D=0$) ; les sorties sont intriquées
($D=\pm\frac12$). Aucune porte agissant séparément sur $A$ et sur $B$ ($U_A\otimes U_B$)
ne peut créer de l'intrication à partir d'un produit, puisque
$(U_A\otimes U_B)(\ket{u}\otimes\ket{v})=U_A\ket{u}\otimes U_B\ket{v}$ reste un produit.
Une porte à deux qubits comme le CNOT est indispensable.
\end{remarque}

% ---------------------------------------------------------------------
\subsection{Récapitulatif du chapitre}
% ---------------------------------------------------------------------

\begin{center}
\renewcommand{\arraystretch}{1.4}
\begin{tabular}{@{}ll@{}}
\toprule
\textbf{Notion} & \textbf{Formule} \\
\midrule
Espace de deux qubits & $\mathcal{H}^2\otimes\mathcal{H}^2=\mathcal{H}^4$, \quad dimension $2^n$ pour $n$ qubits \\
Produit de vecteurs & $(a_1,a_2)^{\mathsf T}\otimes(b_1,b_2)^{\mathsf T}=(a_1b_1,a_1b_2,a_2b_1,a_2b_2)^{\mathsf T}$ \\
Base de calcul & $\ket{ij}_{AB}=\ket{i}_A\otimes\ket{j}_B$, \quad $\braket{ij}{kl}=\delta_{ik}\delta_{jl}$ \\
État général & $\ket{\Psi}=\sum c_{ij}\ket{ij}$, \quad $\sum|c_{ij}|^2=1$, \quad $P(ij)=|c_{ij}|^2$ \\
Séparable & $\ket{\Psi}_A\otimes\ket{\Psi}_B$ \quad $\Leftrightarrow$ \quad $c_{00}c_{11}-c_{01}c_{10}=0$ \\
Intriqué & non séparable, \quad par exemple $\ket{\Phi^+}=\frac{1}{\sqrt2}(\ket{00}+\ket{11})$ \\
Matrices & $A\otimes B=(A_{ij}B)$, \quad $(A\otimes B)(\ket{u}\otimes\ket{v})=A\ket{u}\otimes B\ket{v}$ \\
CNOT & $\ket{i,j}\to\ket{i,i\oplus j}$, \quad $\mathrm{CNOT}\,(H\otimes I)\ket{00}=\ket{\Phi^+}$ \\
\bottomrule
\end{tabular}
\end{center}
.
%  Comme chapitre4.tex et chapitre5.tex, il ouvre une NOUVELLE SECTION
%  (numéro 4). Pour le rattacher à la section précédente, remplacer
%  \section par \subsection et chaque \subsection par \subsubsection.
%  Pas de préambule : environnements (definition, resultat, remarque),
%  macros (\ii, \ee, \dd, \ket, \bra, \braket) et couleurs (insarouge,
%  insagris) sont définis dans main.tex.
% =====================================================================

% Macros de Dirac (déjà définies dans main.tex récent ; ce bloc évite l'erreur
% « Undefined control sequence \ket » si main.tex est une ancienne version).
\providecommand{\ket}[1]{|#1\rangle}
\providecommand{\bra}[1]{\langle #1|}
\providecommand{\braket}[2]{\langle #1|#2\rangle}

% =====================================================================
\section{Systèmes multi-qubits et intrication}\label{sec:multi}
% =====================================================================

Un ordinateur quantique utilise plusieurs qubits. Cette section décrit l'espace des états
de deux qubits (le produit tensoriel), la mesure, la distinction entre états séparables
et états intriqués, puis les opérateurs qui agissent sur deux qubits. Tous les calculs
sont détaillés ; les qubits sont notés $A$ et $B$.

% ---------------------------------------------------------------------
\subsection{L'espace de deux qubits}\label{sec:tenseur}
% ---------------------------------------------------------------------

Chaque qubit vit dans $\mathcal{H}^2=\mathbb{C}^2$. L'espace du système $(A,B)$ est le
\emph{produit tensoriel}
\[
  \mathcal{H}^2\otimes\mathcal{H}^2=\mathcal{H}^4=\mathbb{C}^4 ,
\]
de dimension $2\times2=4$. Pour $n$ qubits la dimension est $2^n$ :

\begin{center}
\renewcommand{\arraystretch}{1.3}
\begin{tabular}{@{}lccccc@{}}
\toprule
Nombre de qubits $n$ & $1$ & $2$ & $3$ & $10$ & $50$ \\
\midrule
Dimension $2^n$ & $2$ & $4$ & $8$ & $1024$ & $\approx1{,}1\times10^{15}$ \\
\bottomrule
\end{tabular}
\end{center}

\begin{definition}{Produit tensoriel de deux vecteurs}
\[
  \begin{pmatrix}a_1\\ a_2\end{pmatrix}\otimes\begin{pmatrix}b_1\\ b_2\end{pmatrix}
  =\begin{pmatrix}a_1\begin{pmatrix}b_1\\ b_2\end{pmatrix}\\[2ex]
                  a_2\begin{pmatrix}b_1\\ b_2\end{pmatrix}\end{pmatrix}
  =\begin{pmatrix}a_1b_1\\ a_1b_2\\ a_2b_1\\ a_2b_2\end{pmatrix}.
\]
Chaque composante du premier vecteur multiplie le second vecteur en entier. On note
$\ket{ij}_{AB}=\ket{i}_A\otimes\ket{j}_B$ : le premier chiffre est l'état de $A$, le second
celui de $B$.
\end{definition}

\subsubsection{Les quatre états de base}

On développe chacun des quatre produits avec $\ket{0}=\begin{pmatrix}1\\ 0\end{pmatrix}$
et $\ket{1}=\begin{pmatrix}0\\ 1\end{pmatrix}$ :
\begin{align*}
  \ket{0}\otimes\ket{0}&=\begin{pmatrix}1\\ 0\end{pmatrix}\otimes\begin{pmatrix}1\\ 0\end{pmatrix}
    =\begin{pmatrix}1\cdot\begin{pmatrix}1\\ 0\end{pmatrix}\\[1.5ex]0\cdot\begin{pmatrix}1\\ 0\end{pmatrix}\end{pmatrix}
    =\begin{pmatrix}1\\ 0\\ 0\\ 0\end{pmatrix}=\ket{00},\\[1.5ex]
  \ket{0}\otimes\ket{1}&=\begin{pmatrix}1\\ 0\end{pmatrix}\otimes\begin{pmatrix}0\\ 1\end{pmatrix}
    =\begin{pmatrix}1\cdot\begin{pmatrix}0\\ 1\end{pmatrix}\\[1.5ex]0\cdot\begin{pmatrix}0\\ 1\end{pmatrix}\end{pmatrix}
    =\begin{pmatrix}0\\ 1\\ 0\\ 0\end{pmatrix}=\ket{01},\\[1.5ex]
  \ket{1}\otimes\ket{0}&=\begin{pmatrix}0\\ 1\end{pmatrix}\otimes\begin{pmatrix}1\\ 0\end{pmatrix}
    =\begin{pmatrix}0\cdot\begin{pmatrix}1\\ 0\end{pmatrix}\\[1.5ex]1\cdot\begin{pmatrix}1\\ 0\end{pmatrix}\end{pmatrix}
    =\begin{pmatrix}0\\ 0\\ 1\\ 0\end{pmatrix}=\ket{10},\\[1.5ex]
  \ket{1}\otimes\ket{1}&=\begin{pmatrix}0\\ 1\end{pmatrix}\otimes\begin{pmatrix}0\\ 1\end{pmatrix}
    =\begin{pmatrix}0\cdot\begin{pmatrix}0\\ 1\end{pmatrix}\\[1.5ex]1\cdot\begin{pmatrix}0\\ 1\end{pmatrix}\end{pmatrix}
    =\begin{pmatrix}0\\ 0\\ 0\\ 1\end{pmatrix}=\ket{11}.
\end{align*}
Le vecteur $\ket{ij}$ a son $1$ en position $2i+j$ (en comptant à partir de $0$) : c'est
la lecture de $ij$ en binaire. Les quatre vecteurs forment la base de calcul de
$\mathcal{H}^4$.

\subsubsection{Produit scalaire et orthonormalité}

\begin{resultat}{Produit scalaire d'un produit tensoriel}
\[
  \big(\bra{a}\otimes\bra{b}\big)\big(\ket{c}\otimes\ket{d}\big)=\braket{a}{c}\,\braket{b}{d}.
\]
\end{resultat}

\emph{Démonstration.} Les composantes de $\ket{a}\otimes\ket{b}$ sont $a_ib_j$, celles de
$\ket{c}\otimes\ket{d}$ sont $c_id_j$. Le produit scalaire est
$\sum_{i,j}(a_ib_j)^*c_id_j=\big(\sum_ia_i^*c_i\big)\big(\sum_jb_j^*d_j\big)=\braket{a}{c}\braket{b}{d}$.
\hfill$\square$

On en déduit $\braket{ij}{kl}=\braket{i}{k}\braket{j}{l}$. Les seize produits scalaires
des états de base sont :
\begin{center}
\renewcommand{\arraystretch}{1.3}
\begin{tabular}{@{}c|cccc@{}}
 & $\ket{00}$ & $\ket{01}$ & $\ket{10}$ & $\ket{11}$ \\
\hline
$\bra{00}$ & $1\cdot1=1$ & $1\cdot0=0$ & $0\cdot1=0$ & $0\cdot0=0$ \\
$\bra{01}$ & $1\cdot0=0$ & $1\cdot1=1$ & $0\cdot0=0$ & $0\cdot1=0$ \\
$\bra{10}$ & $0\cdot1=0$ & $0\cdot0=0$ & $1\cdot1=1$ & $1\cdot0=0$ \\
$\bra{11}$ & $0\cdot0=0$ & $0\cdot1=0$ & $1\cdot0=0$ & $1\cdot1=1$ \\
\end{tabular}
\end{center}
(chaque case est $\braket{i}{k}\cdot\braket{j}{l}$). On lit $\braket{ij}{kl}=\delta_{ik}\delta_{jl}$ :
la base est \emph{orthonormée}.

\subsubsection{Produit tensoriel d'états quelconques}

Le produit tensoriel est linéaire en chaque facteur :
$(\lambda\ket{a}+\mu\ket{a'})\otimes\ket{b}=\lambda\ket{a}\otimes\ket{b}+\mu\ket{a'}\otimes\ket{b}$,
et de même pour le second facteur. Pour $\alpha\ket{0}+\beta\ket{1}$ et
$\gamma\ket{0}+\delta\ket{1}$ :
\begin{align*}
  \big(\alpha\ket{0}+\beta\ket{1}\big)\otimes\big(\gamma\ket{0}+\delta\ket{1}\big)
  &=\alpha\ket{0}\otimes\big(\gamma\ket{0}+\delta\ket{1}\big)
    +\beta\ket{1}\otimes\big(\gamma\ket{0}+\delta\ket{1}\big)\\
  &=\alpha\gamma\ket{00}+\alpha\delta\ket{01}+\beta\gamma\ket{10}+\beta\delta\ket{11}.
\end{align*}
Par les composantes, $\begin{pmatrix}\alpha\\ \beta\end{pmatrix}\otimes\begin{pmatrix}\gamma\\ \delta\end{pmatrix}
=\begin{pmatrix}\alpha\gamma\\ \alpha\delta\\ \beta\gamma\\ \beta\delta\end{pmatrix}$ : c'est le même
résultat. La norme est conservée :
\[
  |\alpha\gamma|^2+|\alpha\delta|^2+|\beta\gamma|^2+|\beta\delta|^2
  =\big(|\alpha|^2+|\beta|^2\big)\big(|\gamma|^2+|\delta|^2\big)=1\cdot1=1 .
\]

\paragraph{Exemples.}
\begin{align*}
  \ket{+}\otimes\ket{+}&=\tfrac12\big(\ket{0}+\ket{1}\big)\otimes\big(\ket{0}+\ket{1}\big)
    =\tfrac12\big(\ket{00}+\ket{01}+\ket{10}+\ket{11}\big)
    =\tfrac12\begin{pmatrix}1\\ 1\\ 1\\ 1\end{pmatrix},\\[1ex]
  \ket{+}\otimes\ket{-}&=\tfrac12\big(\ket{0}+\ket{1}\big)\otimes\big(\ket{0}-\ket{1}\big)
    =\tfrac12\big(\ket{00}-\ket{01}+\ket{10}-\ket{11}\big)
    =\tfrac12\begin{pmatrix}1\\ -1\\ 1\\ -1\end{pmatrix},\\[1ex]
  \ket{0}\otimes\ket{+}&=\tfrac{1}{\sqrt2}\big(\ket{00}+\ket{01}\big)
    =\tfrac{1}{\sqrt2}\begin{pmatrix}1\\ 1\\ 0\\ 0\end{pmatrix},\qquad
  \ket{+}\otimes\ket{0}=\tfrac{1}{\sqrt2}\big(\ket{00}+\ket{10}\big)
    =\tfrac{1}{\sqrt2}\begin{pmatrix}1\\ 0\\ 1\\ 0\end{pmatrix}.
\end{align*}
Les deux derniers vecteurs sont différents : \emph{l'ordre des facteurs compte}, le premier
est toujours celui de $A$.

% ---------------------------------------------------------------------
\subsection{États à deux qubits et mesure}\label{sec:mesure2}
% ---------------------------------------------------------------------

Un état quelconque de $\mathcal{H}^4$ se décompose sur la base de calcul :
\[
  \ket{\Psi}=c_{00}\ket{00}+c_{01}\ket{01}+c_{10}\ket{10}+c_{11}\ket{11}
  =\begin{pmatrix}c_{00}\\ c_{01}\\ c_{10}\\ c_{11}\end{pmatrix},
  \qquad \sum_{i,j}|c_{ij}|^2=1 .
\]

\begin{resultat}{Mesure des deux qubits dans la base de calcul}
\[
  P(ij)=|\braket{ij}{\Psi}|^2=|c_{ij}|^2 .
\]
Pour un seul qubit : $P(A{=}0)=P(00)+P(01)$, $P(A{=}1)=P(10)+P(11)$,
$P(B{=}0)=P(00)+P(10)$, $P(B{=}1)=P(01)+P(11)$.
\end{resultat}

En effet $\braket{ij}{\Psi}=\sum_{k,l}c_{kl}\braket{ij}{kl}=\sum_{k,l}c_{kl}\delta_{ik}\delta_{jl}=c_{ij}$.

\paragraph{Exemple 1.} $\ket{\Psi}=\frac{1}{\sqrt2}\ket{00}+\frac{1}{\sqrt2}\ket{01}
=\big(\tfrac{1}{\sqrt2},\tfrac{1}{\sqrt2},0,0\big)^{\mathsf T}$. Il est normalisé :
$\frac12+\frac12+0+0=1$. Les probabilités jointes sont
$P(00)=P(01)=\frac12$ et $P(10)=P(11)=0$. Pour chaque qubit :
\[
  P(A{=}0)=\tfrac12+\tfrac12=1,\quad P(A{=}1)=0+0=0,\quad
  P(B{=}0)=\tfrac12+0=\tfrac12,\quad P(B{=}1)=\tfrac12+0=\tfrac12 .
\]
$A$ donne toujours $0$ ; $B$ donne $0$ ou $1$ au hasard.

\paragraph{Exemple 2 (mesure de $A$ seul).} Soit
$\ket{\Psi}=\frac12\ket{00}+\frac12\ket{01}+\frac{1}{\sqrt2}\ket{11}$. Normalisation :
$\frac14+\frac14+0+\frac12=1$. Probabilités jointes : $P(00)=\frac14$, $P(01)=\frac14$,
$P(10)=0$, $P(11)=\frac12$. Donc
\[
  P(A{=}0)=\tfrac14+\tfrac14=\tfrac12,\quad P(A{=}1)=0+\tfrac12=\tfrac12,\quad
  P(B{=}0)=\tfrac14+0=\tfrac14,\quad P(B{=}1)=\tfrac14+\tfrac12=\tfrac34 .
\]

\begin{remarque}[Effondrement après la mesure de $A$]
Si $A$ donne le résultat $a$, on ne garde que les termes de $\ket{\Psi}$ dont le premier
chiffre vaut $a$, puis on renormalise :
\[
  \ket{\Psi}\ \to\ \frac{c_{a0}\ket{a0}+c_{a1}\ket{a1}}{\sqrt{|c_{a0}|^2+|c_{a1}|^2}},
  \qquad\text{probabilité }|c_{a0}|^2+|c_{a1}|^2 .
\]
\end{remarque}

Pour l'exemple 2 :
\begin{itemize}[nosep]
  \item $A=0$ (probabilité $\frac12$) : on garde $\frac12\ket{00}+\frac12\ket{01}$, de norme
        au carré $\frac12$ ; renormalisé, $\frac{1}{\sqrt2}\big(\ket{00}+\ket{01}\big)
        =\ket{0}\otimes\ket{+}$. Ensuite $B$ vaut $0$ ou $1$ avec la probabilité $\frac12$.
  \item $A=1$ (probabilité $\frac12$) : on garde $\frac{1}{\sqrt2}\ket{11}$, de norme au carré
        $\frac12$ ; renormalisé, $\ket{11}=\ket{1}\otimes\ket{1}$. Ensuite $B$ vaut $1$
        avec certitude.
\end{itemize}
Contrôle : $P(B{=}0)=\frac12\cdot\frac12+\frac12\cdot0=\frac14$, comme calculé plus haut.
Ici la mesure de $A$ \emph{modifie} l'état de $B$ ($\ket{+}$ ou $\ket{1}$) : c'est la
signature d'un état intriqué (section suivante).

% ---------------------------------------------------------------------
\subsection{États séparables et états intriqués}\label{sec:separable}
% ---------------------------------------------------------------------

\begin{definition}{Séparable ou intriqué}
Un état $\ket{\Psi}_{AB}$ est \emph{séparable} s'il s'écrit comme un produit
$\ket{\Psi}_{AB}=\ket{\Psi}_A\otimes\ket{\Psi}_B$. Sinon il est \emph{intriqué} : on ne
peut pas attribuer d'état propre à chaque qubit.
\end{definition}

\paragraph{Exemple séparable (exemple 1 ci-dessus).} Par linéarité du produit tensoriel :
\begin{align*}
  \ket{\Psi}&=\tfrac{1}{\sqrt2}\ket{00}_{AB}+\tfrac{1}{\sqrt2}\ket{01}_{AB}
    =\tfrac{1}{\sqrt2}\Big(\ket{0}_A\otimes\ket{0}_B+\ket{0}_A\otimes\ket{1}_B\Big)\\
  &=\ket{0}_A\otimes\Big(\tfrac{1}{\sqrt2}\big(\ket{0}+\ket{1}\big)\Big)_B
    =\ket{0}_A\otimes\ket{+}_B ,
\end{align*}
donc $\ket{\Psi}_A=\ket{0}$ et $\ket{\Psi}_B=\ket{+}$. Contrôle par les composantes :
$\begin{pmatrix}1\\ 0\end{pmatrix}\otimes\tfrac{1}{\sqrt2}\begin{pmatrix}1\\ 1\end{pmatrix}
=\tfrac{1}{\sqrt2}\begin{pmatrix}1\cdot1\\ 1\cdot1\\ 0\cdot1\\ 0\cdot1\end{pmatrix}
=\tfrac{1}{\sqrt2}\begin{pmatrix}1\\ 1\\ 0\\ 0\end{pmatrix}$.

\subsubsection{Un critère de séparabilité}

Un produit quelconque a pour coefficients $c_{00}=ac$, $c_{01}=ad$, $c_{10}=bc$,
$c_{11}=bd$ (formule de la section~\ref{sec:tenseur}, avec $a,b$ pour $A$ et $c,d$ pour $B$), donc
\[
  c_{00}c_{11}=(ac)(bd)=abcd=(ad)(bc)=c_{01}c_{10}.
\]

\begin{resultat}{Test de séparabilité (deux qubits)}
$\ket{\Psi}=\sum c_{ij}\ket{ij}$ est séparable si et seulement si
\[
  D=c_{00}c_{11}-c_{01}c_{10}=0 .
\]
\end{resultat}

\emph{Démonstration.} ($\Rightarrow$) Calcul ci-dessus. ($\Leftarrow$) On regroupe les
termes selon l'état de $A$ :
\[
  \ket{\Psi}=\ket{0}\otimes\big(c_{00}\ket{0}+c_{01}\ket{1}\big)+\ket{1}\otimes\big(c_{10}\ket{0}+c_{11}\ket{1}\big).
\]
Si $D=0$, les lignes $(c_{00},c_{01})$ et $(c_{10},c_{11})$ sont proportionnelles. Si
$(c_{00},c_{01})\neq(0,0)$, il existe $\lambda$ avec $(c_{10},c_{11})=\lambda(c_{00},c_{01})$
et $\ket{\Psi}=\big(\ket{0}+\lambda\ket{1}\big)\otimes\big(c_{00}\ket{0}+c_{01}\ket{1}\big)$.
Si $(c_{00},c_{01})=(0,0)$, alors $\ket{\Psi}=\ket{1}\otimes\big(c_{10}\ket{0}+c_{11}\ket{1}\big)$.
Dans les deux cas $\ket{\Psi}$ est un produit ; on renormalise chaque facteur.
\hfill$\square$

\subsubsection{Application à plusieurs états}

\begin{center}
\renewcommand{\arraystretch}{1.5}
\begin{tabular}{@{}lccl@{}}
\toprule
\textbf{État} & $(c_{00},c_{01},c_{10},c_{11})$ & $D$ & \textbf{Nature} \\
\midrule
$\ket{00}$ & $(1,0,0,0)$ & $1\cdot0-0\cdot0=0$ & séparable \\
$\ket{+}\otimes\ket{+}$ & $\frac12(1,1,1,1)$ & $\frac14-\frac14=0$ & séparable \\
$\ket{+}\otimes\ket{-}$ & $\frac12(1,-1,1,-1)$ & $-\frac14-(-\frac14)=0$ & séparable \\
$\frac{1}{\sqrt2}(\ket{00}+\ket{01})$ & $\frac{1}{\sqrt2}(1,1,0,0)$ & $0-\frac12\cdot0=0$ & séparable \\
$\ket{\Phi^+}=\frac{1}{\sqrt2}(\ket{00}+\ket{11})$ & $\frac{1}{\sqrt2}(1,0,0,1)$ & $\frac12-0=\frac12$ & intriqué \\
$\ket{\Phi^-}=\frac{1}{\sqrt2}(\ket{00}-\ket{11})$ & $\frac{1}{\sqrt2}(1,0,0,-1)$ & $-\frac12-0=-\frac12$ & intriqué \\
$\ket{\Psi^+}=\frac{1}{\sqrt2}(\ket{01}+\ket{10})$ & $\frac{1}{\sqrt2}(0,1,1,0)$ & $0-\frac12=-\frac12$ & intriqué \\
$\ket{\Psi^-}=\frac{1}{\sqrt2}(\ket{01}-\ket{10})$ & $\frac{1}{\sqrt2}(0,1,-1,0)$ & $0-(-\frac12)=\frac12$ & intriqué \\
$\frac12\ket{00}+\frac12\ket{01}+\frac{1}{\sqrt2}\ket{11}$ & $(\frac12,\frac12,0,\frac{1}{\sqrt2})$ & $\frac{1}{2\sqrt2}-0$ & intriqué \\
\bottomrule
\end{tabular}
\end{center}

Les quatre états de Bell $\ket{\Phi^\pm}$, $\ket{\Psi^\pm}$ sont \emph{orthonormés} :
chacun est de norme $\frac12+\frac12=1$ ; $\ket{\Phi^\pm}$ n'ont des composantes que sur
$\ket{00}$ et $\ket{11}$, $\ket{\Psi^\pm}$ que sur $\ket{01}$ et $\ket{10}$, donc
$\braket{\Phi^\pm}{\Psi^\pm}=0$, et
$\braket{\Phi^+}{\Phi^-}=\frac12(1-1)=0$, $\braket{\Psi^+}{\Psi^-}=\frac12(1-1)=0$.
Ils forment une base de $\mathcal{H}^4$ constituée uniquement d'états intriqués.

\paragraph{Preuve directe que $\ket{\Phi^+}$ est intriqué.} Supposons
$\ket{\Phi^+}=(a\ket{0}+b\ket{1})\otimes(c\ket{0}+d\ket{1})$. En identifiant les
composantes : $ac=\frac{1}{\sqrt2}$, $ad=0$, $bc=0$, $bd=\frac{1}{\sqrt2}$. La condition
$ad=0$ impose $a=0$ (impossible car $ac\neq0$) ou $d=0$ (impossible car $bd\neq0$).
Contradiction : aucun produit ne convient.

\paragraph{Corrélations.} Pour un produit $\ket{\Psi}_A\otimes\ket{\Psi}_B$ avec
$\ket{\Psi}_A=a_0\ket{0}+a_1\ket{1}$ et $\ket{\Psi}_B=b_0\ket{0}+b_1\ket{1}$, on a
$P(ij)=|a_ib_j|^2=|a_i|^2|b_j|^2=P(A{=}i)\,P(B{=}j)$ : les deux qubits sont
\emph{indépendants}. Vérification avec $\ket{+}\otimes\ket{+}$ : $P(00)=\frac14=\frac12\cdot\frac12$.
Pour $\ket{\Phi^+}$ : $P(00)=P(11)=\frac12$ et $P(01)=P(10)=0$, donc
$P(A{=}0)=P(B{=}0)=\frac12$, mais
\[
  P(00)=\tfrac12\neq\tfrac12\cdot\tfrac12=P(A{=}0)\,P(B{=}0) .
\]
Les deux qubits donnent toujours le même résultat, aléatoire : ils sont corrélés.

\paragraph{État d'un qubit pris seul.} On prend la matrice de densité
$\rho=\ket{\Psi}\bra{\Psi}$ ($4\times4$, lignes et colonnes ordonnées $00,01,10,11$) et on
trace sur $B$ : $(\rho_A)_{ij}=\rho_{(i0),(j0)}+\rho_{(i1),(j1)}$.

\emph{État de Bell $\ket{\Phi^+}$.}
\[
  \rho=\frac12\begin{pmatrix}1 & 0 & 0 & 1\\ 0 & 0 & 0 & 0\\ 0 & 0 & 0 & 0\\ 1 & 0 & 0 & 1\end{pmatrix},
  \qquad
  \begin{aligned}
    (\rho_A)_{00}&=\rho_{00,00}+\rho_{01,01}=\tfrac12+0=\tfrac12,\\
    (\rho_A)_{01}&=\rho_{00,10}+\rho_{01,11}=0+0=0,\\
    (\rho_A)_{11}&=\rho_{10,10}+\rho_{11,11}=0+\tfrac12=\tfrac12 .
  \end{aligned}
\]
Donc $\rho_A=\frac12\mathrm{Id}$ : état \emph{mixte} ($\operatorname{Tr}\rho_A^2=\frac12$), au centre
de la boule de Bloch (section~\ref{sec:bloch}).

\emph{État séparable $\ket{0}\otimes\ket{+}$.}
\[
  \rho=\frac12\begin{pmatrix}1 & 1 & 0 & 0\\ 1 & 1 & 0 & 0\\ 0 & 0 & 0 & 0\\ 0 & 0 & 0 & 0\end{pmatrix},
  \qquad
  \begin{aligned}
    (\rho_A)_{00}&=\rho_{00,00}+\rho_{01,01}=\tfrac12+\tfrac12=1,\\
    (\rho_A)_{01}&=\rho_{00,10}+\rho_{01,11}=0+0=0,\\
    (\rho_A)_{11}&=\rho_{10,10}+\rho_{11,11}=0+0=0 .
  \end{aligned}
\]
Donc $\rho_A=\ket{0}\bra{0}$ : état \emph{pur}, celui du facteur $\ket{\Psi}_A$.

Un qubit d'un état séparable garde son propre état pur ; un qubit d'un état de Bell est
complètement aléatoire alors que le couple est parfaitement défini.

% ---------------------------------------------------------------------
\subsection{Opérateurs sur deux qubits}\label{sec:op2}
% ---------------------------------------------------------------------

\begin{definition}{Produit tensoriel de deux matrices}
\[
  A\otimes B=\begin{pmatrix}A_{11} & A_{12}\\ A_{21} & A_{22}\end{pmatrix}\otimes B
  =\begin{pmatrix}A_{11}B & A_{12}B\\ A_{21}B & A_{22}B\end{pmatrix}
  \qquad(\text{matrice }4\times4).
\]
\end{definition}

\begin{resultat}{Action sur un produit}
\[
  (A\otimes B)\big(\ket{u}\otimes\ket{v}\big)=A\ket{u}\otimes B\ket{v}.
\]
\end{resultat}

\emph{Démonstration.} Avec $\ket{u}=(u_1,u_2)^{\mathsf T}$, le vecteur $\ket{u}\otimes\ket{v}$ a
pour blocs $u_k\ket{v}$ ($k=1,2$). Le bloc de lignes $i$ de $(A\otimes B)(\ket{u}\otimes\ket{v})$ est
donc
\[
  \sum_kA_{ik}B\,\big(u_k\ket{v}\big)=\Big(\sum_kA_{ik}u_k\Big)B\ket{v}=(A\ket{u})_i\,B\ket{v},
\]
qui est le bloc $i$ de $A\ket{u}\otimes B\ket{v}$. \hfill$\square$

\subsubsection{Portes agissant sur un seul des deux qubits}

On note $I=\mathrm{Id}$ ($2\times2$). $X\otimes I$ agit sur $A$, $I\otimes X$ agit sur $B$ :
\begin{align*}
  I\otimes X&=\begin{pmatrix}1\cdot X & 0\cdot X\\ 0\cdot X & 1\cdot X\end{pmatrix}
    =\begin{pmatrix}0 & 1 & 0 & 0\\ 1 & 0 & 0 & 0\\ 0 & 0 & 0 & 1\\ 0 & 0 & 1 & 0\end{pmatrix},\\[1ex]
  X\otimes I&=\begin{pmatrix}0\cdot I & 1\cdot I\\ 1\cdot I & 0\cdot I\end{pmatrix}
    =\begin{pmatrix}0 & 0 & 1 & 0\\ 0 & 0 & 0 & 1\\ 1 & 0 & 0 & 0\\ 0 & 1 & 0 & 0\end{pmatrix}.
\end{align*}
Sur $\ket{00}=(1,0,0,0)^{\mathsf T}$ (première colonne de la matrice) :
$(I\otimes X)\ket{00}=(0,1,0,0)^{\mathsf T}=\ket{01}$ et
$(X\otimes I)\ket{00}=(0,0,1,0)^{\mathsf T}=\ket{10}$. Avec la règle :
$(I\otimes X)\ket{00}=I\ket{0}\otimes X\ket{0}=\ket{0}\otimes\ket{1}=\ket{01}$.

Pour la porte de Hadamard, avec $H=\frac{1}{\sqrt2}\begin{pmatrix}1 & 1\\ 1 & -1\end{pmatrix}$ :
\[
  H\otimes I=\frac{1}{\sqrt2}\begin{pmatrix}1\cdot I & 1\cdot I\\ 1\cdot I & -1\cdot I\end{pmatrix}
  =\frac{1}{\sqrt2}\begin{pmatrix}1 & 0 & 1 & 0\\ 0 & 1 & 0 & 1\\ 1 & 0 & -1 & 0\\ 0 & 1 & 0 & -1\end{pmatrix}.
\]
Ses colonnes donnent l'image des quatre états de base :
\begin{align*}
  (H\otimes I)\ket{00}&=\tfrac{1}{\sqrt2}(1,0,1,0)^{\mathsf T}=\ket{+}\otimes\ket{0},&
  (H\otimes I)\ket{01}&=\tfrac{1}{\sqrt2}(0,1,0,1)^{\mathsf T}=\ket{+}\otimes\ket{1},\\
  (H\otimes I)\ket{10}&=\tfrac{1}{\sqrt2}(1,0,-1,0)^{\mathsf T}=\ket{-}\otimes\ket{0},&
  (H\otimes I)\ket{11}&=\tfrac{1}{\sqrt2}(0,1,0,-1)^{\mathsf T}=\ket{-}\otimes\ket{1}.
\end{align*}

\subsubsection{Hadamard sur les deux qubits}

\[
  H\otimes H=\frac12\begin{pmatrix}1\cdot H' & 1\cdot H'\\ 1\cdot H' & -1\cdot H'\end{pmatrix},
  \quad H'=\begin{pmatrix}1 & 1\\ 1 & -1\end{pmatrix},
  \qquad
  H\otimes H=\frac12\begin{pmatrix}1 & 1 & 1 & 1\\ 1 & -1 & 1 & -1\\ 1 & 1 & -1 & -1\\ 1 & -1 & -1 & 1\end{pmatrix}.
\]
La première colonne donne $(H\otimes H)\ket{00}=\frac12(1,1,1,1)^{\mathsf T}=\ket{+}\otimes\ket{+}$
(cohérent avec $H\ket{0}\otimes H\ket{0}$) : les quatre résultats $00,01,10,11$ sont
équiprobables, chacun avec la probabilité $\frac14$. Sur $n$ qubits, $H^{\otimes n}\ket{0\ldots0}$
donne une superposition uniforme de $2^n$ états.

\subsubsection{La porte CNOT et les états de Bell}

\begin{definition}{Porte CNOT (NOT contrôlé)}
Le qubit $A$ est le \emph{contrôle}, $B$ la \emph{cible} : $B$ est inversé si $A=1$.
\[
  \mathrm{CNOT}=\ket{0}\bra{0}\otimes I+\ket{1}\bra{1}\otimes X
  =\begin{pmatrix}I & 0\\ 0 & X\end{pmatrix}
  =\begin{pmatrix}1 & 0 & 0 & 0\\ 0 & 1 & 0 & 0\\ 0 & 0 & 0 & 1\\ 0 & 0 & 1 & 0\end{pmatrix}.
\]
\end{definition}

Sa matrice échange les composantes $10$ et $11$ :
$\mathrm{CNOT}\ket{00}=\ket{00}$, $\mathrm{CNOT}\ket{01}=\ket{01}$,
$\mathrm{CNOT}\ket{10}=\ket{11}$, $\mathrm{CNOT}\ket{11}=\ket{10}$, soit
$\ket{i,j}\to\ket{i,i\oplus j}$ (addition modulo $2$).

On applique $H\otimes I$ puis $\mathrm{CNOT}$ à $\ket{00}$ (figure~\ref{fig:bell}) :
\[
  \ket{00}\ \xrightarrow{\ H\otimes I\ }\ \tfrac{1}{\sqrt2}\begin{pmatrix}1\\ 0\\ 1\\ 0\end{pmatrix}
  =\tfrac{1}{\sqrt2}\big(\ket{00}+\ket{10}\big)
  \ \xrightarrow{\ \mathrm{CNOT}\ }\ \tfrac{1}{\sqrt2}\begin{pmatrix}1\\ 0\\ 0\\ 1\end{pmatrix}
  =\tfrac{1}{\sqrt2}\big(\ket{00}+\ket{11}\big)=\ket{\Phi^+}.
\]

\begin{figure}[!ht]
\centering
\begin{tikzpicture}[>=Stealth, font=\small]
  % fils : B en haut (dernier chiffre du ket), A en bas (premier chiffre)
  \draw (0,1) -- (8.6,1);
  \draw (0,0) -- (8.6,0);
  \node[anchor=east] at (0,1) {$\ket{0}_B$};
  \node[anchor=east] at (0,0) {$\ket{0}_A$};
  % Hadamard sur A
  \node[draw=insarouge, fill=insarouge!8, minimum size=7mm, thick] at (2.2,0) {$H$};
  % CNOT : contrôle sur A (bas), cible sur B (haut)
  \draw[insarouge, thick] (5.4,0) -- (5.4,1.3);
  \fill[insarouge] (5.4,0) circle (2.6pt);
  \draw[insarouge, thick, fill=white] (5.4,1) circle (0.3);
  \draw[insarouge, thick] (5.1,1) -- (5.7,1);
  \draw[insarouge, thick] (5.4,0.7) -- (5.4,1.3);
  % accolade de sortie
  \draw[insagris, thick] (8.75,1.2) -- (8.9,1.2) -- (8.9,-0.2) -- (8.75,-0.2);
  \node[anchor=west] at (8.9,0.5) {$\ket{\Phi^+}_{AB}$};
  % états intermédiaires
  \foreach \x in {0.6,3.8,7.2}
    \draw[dashed, insagris!70] (\x,1.35) -- (\x,-0.95);
  \node[anchor=north, font=\footnotesize] at (0.6,-0.95) {$\ket{00}$};
  \node[anchor=north, font=\footnotesize] at (3.8,-0.95) {$\tfrac{1}{\sqrt2}(\ket{00}+\ket{10})$};
  \node[anchor=north, font=\footnotesize] at (7.2,-0.95) {$\tfrac{1}{\sqrt2}(\ket{00}+\ket{11})$};
\end{tikzpicture}
\caption{Circuit qui produit l'état de Bell $\ket{\Phi^+}$ : une porte de Hadamard sur $A$,
puis un CNOT ($A$ contrôle, $B$ cible). Les états intermédiaires sont indiqués sous les
traits pointillés. Convention : les fils sont dessinés dans l'ordre inverse de l'écriture
du ket $\ket{ij}_{AB}$, si bien que le premier chiffre (celui de $A$) est sur le fil du bas.}
\label{fig:bell}
\end{figure}

Les quatre états de base d'entrée donnent les quatre états de Bell :
\begin{center}
\renewcommand{\arraystretch}{1.6}
\begin{tabular}{@{}cccc@{}}
\toprule
Entrée & après $H\otimes I$ & après CNOT & Sortie \\
\midrule
$\ket{00}$ & $\frac{1}{\sqrt2}(\ket{00}+\ket{10})$ & $\frac{1}{\sqrt2}(\ket{00}+\ket{11})$ & $\ket{\Phi^+}$ \\
$\ket{01}$ & $\frac{1}{\sqrt2}(\ket{01}+\ket{11})$ & $\frac{1}{\sqrt2}(\ket{01}+\ket{10})$ & $\ket{\Psi^+}$ \\
$\ket{10}$ & $\frac{1}{\sqrt2}(\ket{00}-\ket{10})$ & $\frac{1}{\sqrt2}(\ket{00}-\ket{11})$ & $\ket{\Phi^-}$ \\
$\ket{11}$ & $\frac{1}{\sqrt2}(\ket{01}-\ket{11})$ & $\frac{1}{\sqrt2}(\ket{01}-\ket{10})$ & $\ket{\Psi^-}$ \\
\bottomrule
\end{tabular}
\end{center}
Sous forme matricielle, le produit $\mathrm{CNOT}\cdot(H\otimes I)$ échange les lignes $3$ et $4$ de
$H\otimes I$ :
\[
  \mathrm{CNOT}\,(H\otimes I)=\frac{1}{\sqrt2}\begin{pmatrix}1 & 0 & 1 & 0\\ 0 & 1 & 0 & 1\\ 0 & 1 & 0 & -1\\ 1 & 0 & -1 & 0\end{pmatrix},
\]
et ses quatre colonnes sont $\ket{\Phi^+}$, $\ket{\Psi^+}$, $\ket{\Phi^-}$, $\ket{\Psi^-}$.

\begin{remarque}[Porte intriquante]
Les états d'entrée $\ket{ij}$ sont séparables ($D=0$) ; les sorties sont intriquées
($D=\pm\frac12$). Aucune porte agissant séparément sur $A$ et sur $B$ ($U_A\otimes U_B$)
ne peut créer de l'intrication à partir d'un produit, puisque
$(U_A\otimes U_B)(\ket{u}\otimes\ket{v})=U_A\ket{u}\otimes U_B\ket{v}$ reste un produit.
Une porte à deux qubits comme le CNOT est indispensable.
\end{remarque}

% ---------------------------------------------------------------------
\subsection{Récapitulatif du chapitre}
% ---------------------------------------------------------------------

\begin{center}
\renewcommand{\arraystretch}{1.4}
\begin{tabular}{@{}ll@{}}
\toprule
\textbf{Notion} & \textbf{Formule} \\
\midrule
Espace de deux qubits & $\mathcal{H}^2\otimes\mathcal{H}^2=\mathcal{H}^4$, \quad dimension $2^n$ pour $n$ qubits \\
Produit de vecteurs & $(a_1,a_2)^{\mathsf T}\otimes(b_1,b_2)^{\mathsf T}=(a_1b_1,a_1b_2,a_2b_1,a_2b_2)^{\mathsf T}$ \\
Base de calcul & $\ket{ij}_{AB}=\ket{i}_A\otimes\ket{j}_B$, \quad $\braket{ij}{kl}=\delta_{ik}\delta_{jl}$ \\
État général & $\ket{\Psi}=\sum c_{ij}\ket{ij}$, \quad $\sum|c_{ij}|^2=1$, \quad $P(ij)=|c_{ij}|^2$ \\
Séparable & $\ket{\Psi}_A\otimes\ket{\Psi}_B$ \quad $\Leftrightarrow$ \quad $c_{00}c_{11}-c_{01}c_{10}=0$ \\
Intriqué & non séparable, \quad par exemple $\ket{\Phi^+}=\frac{1}{\sqrt2}(\ket{00}+\ket{11})$ \\
Matrices & $A\otimes B=(A_{ij}B)$, \quad $(A\otimes B)(\ket{u}\otimes\ket{v})=A\ket{u}\otimes B\ket{v}$ \\
CNOT & $\ket{i,j}\to\ket{i,i\oplus j}$, \quad $\mathrm{CNOT}\,(H\otimes I)\ket{00}=\ket{\Phi^+}$ \\
\bottomrule
\end{tabular}
\end{center}
.
%  Comme chapitre4.tex et chapitre5.tex, il ouvre une NOUVELLE SECTION
%  (numéro 4). Pour le rattacher à la section précédente, remplacer
%  \section par \subsection et chaque \subsection par \subsubsection.
%  Pas de préambule : environnements (definition, resultat, remarque),
%  macros (\ii, \ee, \dd, \ket, \bra, \braket) et couleurs (insarouge,
%  insagris) sont définis dans main.tex.
% =====================================================================

% Macros de Dirac (déjà définies dans main.tex récent ; ce bloc évite l'erreur
% « Undefined control sequence \ket » si main.tex est une ancienne version).
\providecommand{\ket}[1]{|#1\rangle}
\providecommand{\bra}[1]{\langle #1|}
\providecommand{\braket}[2]{\langle #1|#2\rangle}

% =====================================================================
\section{Systèmes multi-qubits et intrication}\label{sec:multi}
% =====================================================================

Un ordinateur quantique utilise plusieurs qubits. Cette section décrit l'espace des états
de deux qubits (le produit tensoriel), la mesure, la distinction entre états séparables
et états intriqués, puis les opérateurs qui agissent sur deux qubits. Tous les calculs
sont détaillés ; les qubits sont notés $A$ et $B$.

% ---------------------------------------------------------------------
\subsection{L'espace de deux qubits}\label{sec:tenseur}
% ---------------------------------------------------------------------

Chaque qubit vit dans $\mathcal{H}^2=\mathbb{C}^2$. L'espace du système $(A,B)$ est le
\emph{produit tensoriel}
\[
  \mathcal{H}^2\otimes\mathcal{H}^2=\mathcal{H}^4=\mathbb{C}^4 ,
\]
de dimension $2\times2=4$. Pour $n$ qubits la dimension est $2^n$ :

\begin{center}
\renewcommand{\arraystretch}{1.3}
\begin{tabular}{@{}lccccc@{}}
\toprule
Nombre de qubits $n$ & $1$ & $2$ & $3$ & $10$ & $50$ \\
\midrule
Dimension $2^n$ & $2$ & $4$ & $8$ & $1024$ & $\approx1{,}1\times10^{15}$ \\
\bottomrule
\end{tabular}
\end{center}

\begin{definition}{Produit tensoriel de deux vecteurs}
\[
  \begin{pmatrix}a_1\\ a_2\end{pmatrix}\otimes\begin{pmatrix}b_1\\ b_2\end{pmatrix}
  =\begin{pmatrix}a_1\begin{pmatrix}b_1\\ b_2\end{pmatrix}\\[2ex]
                  a_2\begin{pmatrix}b_1\\ b_2\end{pmatrix}\end{pmatrix}
  =\begin{pmatrix}a_1b_1\\ a_1b_2\\ a_2b_1\\ a_2b_2\end{pmatrix}.
\]
Chaque composante du premier vecteur multiplie le second vecteur en entier. On note
$\ket{ij}_{AB}=\ket{i}_A\otimes\ket{j}_B$ : le premier chiffre est l'état de $A$, le second
celui de $B$.
\end{definition}

\subsubsection{Les quatre états de base}

On développe chacun des quatre produits avec $\ket{0}=\begin{pmatrix}1\\ 0\end{pmatrix}$
et $\ket{1}=\begin{pmatrix}0\\ 1\end{pmatrix}$ :
\begin{align*}
  \ket{0}\otimes\ket{0}&=\begin{pmatrix}1\\ 0\end{pmatrix}\otimes\begin{pmatrix}1\\ 0\end{pmatrix}
    =\begin{pmatrix}1\cdot\begin{pmatrix}1\\ 0\end{pmatrix}\\[1.5ex]0\cdot\begin{pmatrix}1\\ 0\end{pmatrix}\end{pmatrix}
    =\begin{pmatrix}1\\ 0\\ 0\\ 0\end{pmatrix}=\ket{00},\\[1.5ex]
  \ket{0}\otimes\ket{1}&=\begin{pmatrix}1\\ 0\end{pmatrix}\otimes\begin{pmatrix}0\\ 1\end{pmatrix}
    =\begin{pmatrix}1\cdot\begin{pmatrix}0\\ 1\end{pmatrix}\\[1.5ex]0\cdot\begin{pmatrix}0\\ 1\end{pmatrix}\end{pmatrix}
    =\begin{pmatrix}0\\ 1\\ 0\\ 0\end{pmatrix}=\ket{01},\\[1.5ex]
  \ket{1}\otimes\ket{0}&=\begin{pmatrix}0\\ 1\end{pmatrix}\otimes\begin{pmatrix}1\\ 0\end{pmatrix}
    =\begin{pmatrix}0\cdot\begin{pmatrix}1\\ 0\end{pmatrix}\\[1.5ex]1\cdot\begin{pmatrix}1\\ 0\end{pmatrix}\end{pmatrix}
    =\begin{pmatrix}0\\ 0\\ 1\\ 0\end{pmatrix}=\ket{10},\\[1.5ex]
  \ket{1}\otimes\ket{1}&=\begin{pmatrix}0\\ 1\end{pmatrix}\otimes\begin{pmatrix}0\\ 1\end{pmatrix}
    =\begin{pmatrix}0\cdot\begin{pmatrix}0\\ 1\end{pmatrix}\\[1.5ex]1\cdot\begin{pmatrix}0\\ 1\end{pmatrix}\end{pmatrix}
    =\begin{pmatrix}0\\ 0\\ 0\\ 1\end{pmatrix}=\ket{11}.
\end{align*}
Le vecteur $\ket{ij}$ a son $1$ en position $2i+j$ (en comptant à partir de $0$) : c'est
la lecture de $ij$ en binaire. Les quatre vecteurs forment la base de calcul de
$\mathcal{H}^4$.

\subsubsection{Produit scalaire et orthonormalité}

\begin{resultat}{Produit scalaire d'un produit tensoriel}
\[
  \big(\bra{a}\otimes\bra{b}\big)\big(\ket{c}\otimes\ket{d}\big)=\braket{a}{c}\,\braket{b}{d}.
\]
\end{resultat}

\emph{Démonstration.} Les composantes de $\ket{a}\otimes\ket{b}$ sont $a_ib_j$, celles de
$\ket{c}\otimes\ket{d}$ sont $c_id_j$. Le produit scalaire est
$\sum_{i,j}(a_ib_j)^*c_id_j=\big(\sum_ia_i^*c_i\big)\big(\sum_jb_j^*d_j\big)=\braket{a}{c}\braket{b}{d}$.
\hfill$\square$

On en déduit $\braket{ij}{kl}=\braket{i}{k}\braket{j}{l}$. Les seize produits scalaires
des états de base sont :
\begin{center}
\renewcommand{\arraystretch}{1.3}
\begin{tabular}{@{}c|cccc@{}}
 & $\ket{00}$ & $\ket{01}$ & $\ket{10}$ & $\ket{11}$ \\
\hline
$\bra{00}$ & $1\cdot1=1$ & $1\cdot0=0$ & $0\cdot1=0$ & $0\cdot0=0$ \\
$\bra{01}$ & $1\cdot0=0$ & $1\cdot1=1$ & $0\cdot0=0$ & $0\cdot1=0$ \\
$\bra{10}$ & $0\cdot1=0$ & $0\cdot0=0$ & $1\cdot1=1$ & $1\cdot0=0$ \\
$\bra{11}$ & $0\cdot0=0$ & $0\cdot1=0$ & $1\cdot0=0$ & $1\cdot1=1$ \\
\end{tabular}
\end{center}
(chaque case est $\braket{i}{k}\cdot\braket{j}{l}$). On lit $\braket{ij}{kl}=\delta_{ik}\delta_{jl}$ :
la base est \emph{orthonormée}.

\subsubsection{Produit tensoriel d'états quelconques}

Le produit tensoriel est linéaire en chaque facteur :
$(\lambda\ket{a}+\mu\ket{a'})\otimes\ket{b}=\lambda\ket{a}\otimes\ket{b}+\mu\ket{a'}\otimes\ket{b}$,
et de même pour le second facteur. Pour $\alpha\ket{0}+\beta\ket{1}$ et
$\gamma\ket{0}+\delta\ket{1}$ :
\begin{align*}
  \big(\alpha\ket{0}+\beta\ket{1}\big)\otimes\big(\gamma\ket{0}+\delta\ket{1}\big)
  &=\alpha\ket{0}\otimes\big(\gamma\ket{0}+\delta\ket{1}\big)
    +\beta\ket{1}\otimes\big(\gamma\ket{0}+\delta\ket{1}\big)\\
  &=\alpha\gamma\ket{00}+\alpha\delta\ket{01}+\beta\gamma\ket{10}+\beta\delta\ket{11}.
\end{align*}
Par les composantes, $\begin{pmatrix}\alpha\\ \beta\end{pmatrix}\otimes\begin{pmatrix}\gamma\\ \delta\end{pmatrix}
=\begin{pmatrix}\alpha\gamma\\ \alpha\delta\\ \beta\gamma\\ \beta\delta\end{pmatrix}$ : c'est le même
résultat. La norme est conservée :
\[
  |\alpha\gamma|^2+|\alpha\delta|^2+|\beta\gamma|^2+|\beta\delta|^2
  =\big(|\alpha|^2+|\beta|^2\big)\big(|\gamma|^2+|\delta|^2\big)=1\cdot1=1 .
\]

\paragraph{Exemples.}
\begin{align*}
  \ket{+}\otimes\ket{+}&=\tfrac12\big(\ket{0}+\ket{1}\big)\otimes\big(\ket{0}+\ket{1}\big)
    =\tfrac12\big(\ket{00}+\ket{01}+\ket{10}+\ket{11}\big)
    =\tfrac12\begin{pmatrix}1\\ 1\\ 1\\ 1\end{pmatrix},\\[1ex]
  \ket{+}\otimes\ket{-}&=\tfrac12\big(\ket{0}+\ket{1}\big)\otimes\big(\ket{0}-\ket{1}\big)
    =\tfrac12\big(\ket{00}-\ket{01}+\ket{10}-\ket{11}\big)
    =\tfrac12\begin{pmatrix}1\\ -1\\ 1\\ -1\end{pmatrix},\\[1ex]
  \ket{0}\otimes\ket{+}&=\tfrac{1}{\sqrt2}\big(\ket{00}+\ket{01}\big)
    =\tfrac{1}{\sqrt2}\begin{pmatrix}1\\ 1\\ 0\\ 0\end{pmatrix},\qquad
  \ket{+}\otimes\ket{0}=\tfrac{1}{\sqrt2}\big(\ket{00}+\ket{10}\big)
    =\tfrac{1}{\sqrt2}\begin{pmatrix}1\\ 0\\ 1\\ 0\end{pmatrix}.
\end{align*}
Les deux derniers vecteurs sont différents : \emph{l'ordre des facteurs compte}, le premier
est toujours celui de $A$.

% ---------------------------------------------------------------------
\subsection{États à deux qubits et mesure}\label{sec:mesure2}
% ---------------------------------------------------------------------

Un état quelconque de $\mathcal{H}^4$ se décompose sur la base de calcul :
\[
  \ket{\Psi}=c_{00}\ket{00}+c_{01}\ket{01}+c_{10}\ket{10}+c_{11}\ket{11}
  =\begin{pmatrix}c_{00}\\ c_{01}\\ c_{10}\\ c_{11}\end{pmatrix},
  \qquad \sum_{i,j}|c_{ij}|^2=1 .
\]

\begin{resultat}{Mesure des deux qubits dans la base de calcul}
\[
  P(ij)=|\braket{ij}{\Psi}|^2=|c_{ij}|^2 .
\]
Pour un seul qubit : $P(A{=}0)=P(00)+P(01)$, $P(A{=}1)=P(10)+P(11)$,
$P(B{=}0)=P(00)+P(10)$, $P(B{=}1)=P(01)+P(11)$.
\end{resultat}

En effet $\braket{ij}{\Psi}=\sum_{k,l}c_{kl}\braket{ij}{kl}=\sum_{k,l}c_{kl}\delta_{ik}\delta_{jl}=c_{ij}$.

\paragraph{Exemple 1.} $\ket{\Psi}=\frac{1}{\sqrt2}\ket{00}+\frac{1}{\sqrt2}\ket{01}
=\big(\tfrac{1}{\sqrt2},\tfrac{1}{\sqrt2},0,0\big)^{\mathsf T}$. Il est normalisé :
$\frac12+\frac12+0+0=1$. Les probabilités jointes sont
$P(00)=P(01)=\frac12$ et $P(10)=P(11)=0$. Pour chaque qubit :
\[
  P(A{=}0)=\tfrac12+\tfrac12=1,\quad P(A{=}1)=0+0=0,\quad
  P(B{=}0)=\tfrac12+0=\tfrac12,\quad P(B{=}1)=\tfrac12+0=\tfrac12 .
\]
$A$ donne toujours $0$ ; $B$ donne $0$ ou $1$ au hasard.

\paragraph{Exemple 2 (mesure de $A$ seul).} Soit
$\ket{\Psi}=\frac12\ket{00}+\frac12\ket{01}+\frac{1}{\sqrt2}\ket{11}$. Normalisation :
$\frac14+\frac14+0+\frac12=1$. Probabilités jointes : $P(00)=\frac14$, $P(01)=\frac14$,
$P(10)=0$, $P(11)=\frac12$. Donc
\[
  P(A{=}0)=\tfrac14+\tfrac14=\tfrac12,\quad P(A{=}1)=0+\tfrac12=\tfrac12,\quad
  P(B{=}0)=\tfrac14+0=\tfrac14,\quad P(B{=}1)=\tfrac14+\tfrac12=\tfrac34 .
\]

\begin{remarque}[Effondrement après la mesure de $A$]
Si $A$ donne le résultat $a$, on ne garde que les termes de $\ket{\Psi}$ dont le premier
chiffre vaut $a$, puis on renormalise :
\[
  \ket{\Psi}\ \to\ \frac{c_{a0}\ket{a0}+c_{a1}\ket{a1}}{\sqrt{|c_{a0}|^2+|c_{a1}|^2}},
  \qquad\text{probabilité }|c_{a0}|^2+|c_{a1}|^2 .
\]
\end{remarque}

Pour l'exemple 2 :
\begin{itemize}[nosep]
  \item $A=0$ (probabilité $\frac12$) : on garde $\frac12\ket{00}+\frac12\ket{01}$, de norme
        au carré $\frac12$ ; renormalisé, $\frac{1}{\sqrt2}\big(\ket{00}+\ket{01}\big)
        =\ket{0}\otimes\ket{+}$. Ensuite $B$ vaut $0$ ou $1$ avec la probabilité $\frac12$.
  \item $A=1$ (probabilité $\frac12$) : on garde $\frac{1}{\sqrt2}\ket{11}$, de norme au carré
        $\frac12$ ; renormalisé, $\ket{11}=\ket{1}\otimes\ket{1}$. Ensuite $B$ vaut $1$
        avec certitude.
\end{itemize}
Contrôle : $P(B{=}0)=\frac12\cdot\frac12+\frac12\cdot0=\frac14$, comme calculé plus haut.
Ici la mesure de $A$ \emph{modifie} l'état de $B$ ($\ket{+}$ ou $\ket{1}$) : c'est la
signature d'un état intriqué (section suivante).

% ---------------------------------------------------------------------
\subsection{États séparables et états intriqués}\label{sec:separable}
% ---------------------------------------------------------------------

\begin{definition}{Séparable ou intriqué}
Un état $\ket{\Psi}_{AB}$ est \emph{séparable} s'il s'écrit comme un produit
$\ket{\Psi}_{AB}=\ket{\Psi}_A\otimes\ket{\Psi}_B$. Sinon il est \emph{intriqué} : on ne
peut pas attribuer d'état propre à chaque qubit.
\end{definition}

\paragraph{Exemple séparable (exemple 1 ci-dessus).} Par linéarité du produit tensoriel :
\begin{align*}
  \ket{\Psi}&=\tfrac{1}{\sqrt2}\ket{00}_{AB}+\tfrac{1}{\sqrt2}\ket{01}_{AB}
    =\tfrac{1}{\sqrt2}\Big(\ket{0}_A\otimes\ket{0}_B+\ket{0}_A\otimes\ket{1}_B\Big)\\
  &=\ket{0}_A\otimes\Big(\tfrac{1}{\sqrt2}\big(\ket{0}+\ket{1}\big)\Big)_B
    =\ket{0}_A\otimes\ket{+}_B ,
\end{align*}
donc $\ket{\Psi}_A=\ket{0}$ et $\ket{\Psi}_B=\ket{+}$. Contrôle par les composantes :
$\begin{pmatrix}1\\ 0\end{pmatrix}\otimes\tfrac{1}{\sqrt2}\begin{pmatrix}1\\ 1\end{pmatrix}
=\tfrac{1}{\sqrt2}\begin{pmatrix}1\cdot1\\ 1\cdot1\\ 0\cdot1\\ 0\cdot1\end{pmatrix}
=\tfrac{1}{\sqrt2}\begin{pmatrix}1\\ 1\\ 0\\ 0\end{pmatrix}$.

\subsubsection{Un critère de séparabilité}

Un produit quelconque a pour coefficients $c_{00}=ac$, $c_{01}=ad$, $c_{10}=bc$,
$c_{11}=bd$ (formule de la section~\ref{sec:tenseur}, avec $a,b$ pour $A$ et $c,d$ pour $B$), donc
\[
  c_{00}c_{11}=(ac)(bd)=abcd=(ad)(bc)=c_{01}c_{10}.
\]

\begin{resultat}{Test de séparabilité (deux qubits)}
$\ket{\Psi}=\sum c_{ij}\ket{ij}$ est séparable si et seulement si
\[
  D=c_{00}c_{11}-c_{01}c_{10}=0 .
\]
\end{resultat}

\emph{Démonstration.} ($\Rightarrow$) Calcul ci-dessus. ($\Leftarrow$) On regroupe les
termes selon l'état de $A$ :
\[
  \ket{\Psi}=\ket{0}\otimes\big(c_{00}\ket{0}+c_{01}\ket{1}\big)+\ket{1}\otimes\big(c_{10}\ket{0}+c_{11}\ket{1}\big).
\]
Si $D=0$, les lignes $(c_{00},c_{01})$ et $(c_{10},c_{11})$ sont proportionnelles. Si
$(c_{00},c_{01})\neq(0,0)$, il existe $\lambda$ avec $(c_{10},c_{11})=\lambda(c_{00},c_{01})$
et $\ket{\Psi}=\big(\ket{0}+\lambda\ket{1}\big)\otimes\big(c_{00}\ket{0}+c_{01}\ket{1}\big)$.
Si $(c_{00},c_{01})=(0,0)$, alors $\ket{\Psi}=\ket{1}\otimes\big(c_{10}\ket{0}+c_{11}\ket{1}\big)$.
Dans les deux cas $\ket{\Psi}$ est un produit ; on renormalise chaque facteur.
\hfill$\square$

\subsubsection{Application à plusieurs états}

\begin{center}
\renewcommand{\arraystretch}{1.5}
\begin{tabular}{@{}lccl@{}}
\toprule
\textbf{État} & $(c_{00},c_{01},c_{10},c_{11})$ & $D$ & \textbf{Nature} \\
\midrule
$\ket{00}$ & $(1,0,0,0)$ & $1\cdot0-0\cdot0=0$ & séparable \\
$\ket{+}\otimes\ket{+}$ & $\frac12(1,1,1,1)$ & $\frac14-\frac14=0$ & séparable \\
$\ket{+}\otimes\ket{-}$ & $\frac12(1,-1,1,-1)$ & $-\frac14-(-\frac14)=0$ & séparable \\
$\frac{1}{\sqrt2}(\ket{00}+\ket{01})$ & $\frac{1}{\sqrt2}(1,1,0,0)$ & $0-\frac12\cdot0=0$ & séparable \\
$\ket{\Phi^+}=\frac{1}{\sqrt2}(\ket{00}+\ket{11})$ & $\frac{1}{\sqrt2}(1,0,0,1)$ & $\frac12-0=\frac12$ & intriqué \\
$\ket{\Phi^-}=\frac{1}{\sqrt2}(\ket{00}-\ket{11})$ & $\frac{1}{\sqrt2}(1,0,0,-1)$ & $-\frac12-0=-\frac12$ & intriqué \\
$\ket{\Psi^+}=\frac{1}{\sqrt2}(\ket{01}+\ket{10})$ & $\frac{1}{\sqrt2}(0,1,1,0)$ & $0-\frac12=-\frac12$ & intriqué \\
$\ket{\Psi^-}=\frac{1}{\sqrt2}(\ket{01}-\ket{10})$ & $\frac{1}{\sqrt2}(0,1,-1,0)$ & $0-(-\frac12)=\frac12$ & intriqué \\
$\frac12\ket{00}+\frac12\ket{01}+\frac{1}{\sqrt2}\ket{11}$ & $(\frac12,\frac12,0,\frac{1}{\sqrt2})$ & $\frac{1}{2\sqrt2}-0$ & intriqué \\
\bottomrule
\end{tabular}
\end{center}

Les quatre états de Bell $\ket{\Phi^\pm}$, $\ket{\Psi^\pm}$ sont \emph{orthonormés} :
chacun est de norme $\frac12+\frac12=1$ ; $\ket{\Phi^\pm}$ n'ont des composantes que sur
$\ket{00}$ et $\ket{11}$, $\ket{\Psi^\pm}$ que sur $\ket{01}$ et $\ket{10}$, donc
$\braket{\Phi^\pm}{\Psi^\pm}=0$, et
$\braket{\Phi^+}{\Phi^-}=\frac12(1-1)=0$, $\braket{\Psi^+}{\Psi^-}=\frac12(1-1)=0$.
Ils forment une base de $\mathcal{H}^4$ constituée uniquement d'états intriqués.

\paragraph{Preuve directe que $\ket{\Phi^+}$ est intriqué.} Supposons
$\ket{\Phi^+}=(a\ket{0}+b\ket{1})\otimes(c\ket{0}+d\ket{1})$. En identifiant les
composantes : $ac=\frac{1}{\sqrt2}$, $ad=0$, $bc=0$, $bd=\frac{1}{\sqrt2}$. La condition
$ad=0$ impose $a=0$ (impossible car $ac\neq0$) ou $d=0$ (impossible car $bd\neq0$).
Contradiction : aucun produit ne convient.

\paragraph{Corrélations.} Pour un produit $\ket{\Psi}_A\otimes\ket{\Psi}_B$ avec
$\ket{\Psi}_A=a_0\ket{0}+a_1\ket{1}$ et $\ket{\Psi}_B=b_0\ket{0}+b_1\ket{1}$, on a
$P(ij)=|a_ib_j|^2=|a_i|^2|b_j|^2=P(A{=}i)\,P(B{=}j)$ : les deux qubits sont
\emph{indépendants}. Vérification avec $\ket{+}\otimes\ket{+}$ : $P(00)=\frac14=\frac12\cdot\frac12$.
Pour $\ket{\Phi^+}$ : $P(00)=P(11)=\frac12$ et $P(01)=P(10)=0$, donc
$P(A{=}0)=P(B{=}0)=\frac12$, mais
\[
  P(00)=\tfrac12\neq\tfrac12\cdot\tfrac12=P(A{=}0)\,P(B{=}0) .
\]
Les deux qubits donnent toujours le même résultat, aléatoire : ils sont corrélés.

\paragraph{État d'un qubit pris seul.} On prend la matrice de densité
$\rho=\ket{\Psi}\bra{\Psi}$ ($4\times4$, lignes et colonnes ordonnées $00,01,10,11$) et on
trace sur $B$ : $(\rho_A)_{ij}=\rho_{(i0),(j0)}+\rho_{(i1),(j1)}$.

\emph{État de Bell $\ket{\Phi^+}$.}
\[
  \rho=\frac12\begin{pmatrix}1 & 0 & 0 & 1\\ 0 & 0 & 0 & 0\\ 0 & 0 & 0 & 0\\ 1 & 0 & 0 & 1\end{pmatrix},
  \qquad
  \begin{aligned}
    (\rho_A)_{00}&=\rho_{00,00}+\rho_{01,01}=\tfrac12+0=\tfrac12,\\
    (\rho_A)_{01}&=\rho_{00,10}+\rho_{01,11}=0+0=0,\\
    (\rho_A)_{11}&=\rho_{10,10}+\rho_{11,11}=0+\tfrac12=\tfrac12 .
  \end{aligned}
\]
Donc $\rho_A=\frac12\mathrm{Id}$ : état \emph{mixte} ($\operatorname{Tr}\rho_A^2=\frac12$), au centre
de la boule de Bloch (section~\ref{sec:bloch}).

\emph{État séparable $\ket{0}\otimes\ket{+}$.}
\[
  \rho=\frac12\begin{pmatrix}1 & 1 & 0 & 0\\ 1 & 1 & 0 & 0\\ 0 & 0 & 0 & 0\\ 0 & 0 & 0 & 0\end{pmatrix},
  \qquad
  \begin{aligned}
    (\rho_A)_{00}&=\rho_{00,00}+\rho_{01,01}=\tfrac12+\tfrac12=1,\\
    (\rho_A)_{01}&=\rho_{00,10}+\rho_{01,11}=0+0=0,\\
    (\rho_A)_{11}&=\rho_{10,10}+\rho_{11,11}=0+0=0 .
  \end{aligned}
\]
Donc $\rho_A=\ket{0}\bra{0}$ : état \emph{pur}, celui du facteur $\ket{\Psi}_A$.

Un qubit d'un état séparable garde son propre état pur ; un qubit d'un état de Bell est
complètement aléatoire alors que le couple est parfaitement défini.

% ---------------------------------------------------------------------
\subsection{Opérateurs sur deux qubits}\label{sec:op2}
% ---------------------------------------------------------------------

\begin{definition}{Produit tensoriel de deux matrices}
\[
  A\otimes B=\begin{pmatrix}A_{11} & A_{12}\\ A_{21} & A_{22}\end{pmatrix}\otimes B
  =\begin{pmatrix}A_{11}B & A_{12}B\\ A_{21}B & A_{22}B\end{pmatrix}
  \qquad(\text{matrice }4\times4).
\]
\end{definition}

\begin{resultat}{Action sur un produit}
\[
  (A\otimes B)\big(\ket{u}\otimes\ket{v}\big)=A\ket{u}\otimes B\ket{v}.
\]
\end{resultat}

\emph{Démonstration.} Avec $\ket{u}=(u_1,u_2)^{\mathsf T}$, le vecteur $\ket{u}\otimes\ket{v}$ a
pour blocs $u_k\ket{v}$ ($k=1,2$). Le bloc de lignes $i$ de $(A\otimes B)(\ket{u}\otimes\ket{v})$ est
donc
\[
  \sum_kA_{ik}B\,\big(u_k\ket{v}\big)=\Big(\sum_kA_{ik}u_k\Big)B\ket{v}=(A\ket{u})_i\,B\ket{v},
\]
qui est le bloc $i$ de $A\ket{u}\otimes B\ket{v}$. \hfill$\square$

\subsubsection{Portes agissant sur un seul des deux qubits}

On note $I=\mathrm{Id}$ ($2\times2$). $X\otimes I$ agit sur $A$, $I\otimes X$ agit sur $B$ :
\begin{align*}
  I\otimes X&=\begin{pmatrix}1\cdot X & 0\cdot X\\ 0\cdot X & 1\cdot X\end{pmatrix}
    =\begin{pmatrix}0 & 1 & 0 & 0\\ 1 & 0 & 0 & 0\\ 0 & 0 & 0 & 1\\ 0 & 0 & 1 & 0\end{pmatrix},\\[1ex]
  X\otimes I&=\begin{pmatrix}0\cdot I & 1\cdot I\\ 1\cdot I & 0\cdot I\end{pmatrix}
    =\begin{pmatrix}0 & 0 & 1 & 0\\ 0 & 0 & 0 & 1\\ 1 & 0 & 0 & 0\\ 0 & 1 & 0 & 0\end{pmatrix}.
\end{align*}
Sur $\ket{00}=(1,0,0,0)^{\mathsf T}$ (première colonne de la matrice) :
$(I\otimes X)\ket{00}=(0,1,0,0)^{\mathsf T}=\ket{01}$ et
$(X\otimes I)\ket{00}=(0,0,1,0)^{\mathsf T}=\ket{10}$. Avec la règle :
$(I\otimes X)\ket{00}=I\ket{0}\otimes X\ket{0}=\ket{0}\otimes\ket{1}=\ket{01}$.

Pour la porte de Hadamard, avec $H=\frac{1}{\sqrt2}\begin{pmatrix}1 & 1\\ 1 & -1\end{pmatrix}$ :
\[
  H\otimes I=\frac{1}{\sqrt2}\begin{pmatrix}1\cdot I & 1\cdot I\\ 1\cdot I & -1\cdot I\end{pmatrix}
  =\frac{1}{\sqrt2}\begin{pmatrix}1 & 0 & 1 & 0\\ 0 & 1 & 0 & 1\\ 1 & 0 & -1 & 0\\ 0 & 1 & 0 & -1\end{pmatrix}.
\]
Ses colonnes donnent l'image des quatre états de base :
\begin{align*}
  (H\otimes I)\ket{00}&=\tfrac{1}{\sqrt2}(1,0,1,0)^{\mathsf T}=\ket{+}\otimes\ket{0},&
  (H\otimes I)\ket{01}&=\tfrac{1}{\sqrt2}(0,1,0,1)^{\mathsf T}=\ket{+}\otimes\ket{1},\\
  (H\otimes I)\ket{10}&=\tfrac{1}{\sqrt2}(1,0,-1,0)^{\mathsf T}=\ket{-}\otimes\ket{0},&
  (H\otimes I)\ket{11}&=\tfrac{1}{\sqrt2}(0,1,0,-1)^{\mathsf T}=\ket{-}\otimes\ket{1}.
\end{align*}

\subsubsection{Hadamard sur les deux qubits}

\[
  H\otimes H=\frac12\begin{pmatrix}1\cdot H' & 1\cdot H'\\ 1\cdot H' & -1\cdot H'\end{pmatrix},
  \quad H'=\begin{pmatrix}1 & 1\\ 1 & -1\end{pmatrix},
  \qquad
  H\otimes H=\frac12\begin{pmatrix}1 & 1 & 1 & 1\\ 1 & -1 & 1 & -1\\ 1 & 1 & -1 & -1\\ 1 & -1 & -1 & 1\end{pmatrix}.
\]
La première colonne donne $(H\otimes H)\ket{00}=\frac12(1,1,1,1)^{\mathsf T}=\ket{+}\otimes\ket{+}$
(cohérent avec $H\ket{0}\otimes H\ket{0}$) : les quatre résultats $00,01,10,11$ sont
équiprobables, chacun avec la probabilité $\frac14$. Sur $n$ qubits, $H^{\otimes n}\ket{0\ldots0}$
donne une superposition uniforme de $2^n$ états.

\subsubsection{La porte CNOT et les états de Bell}

\begin{definition}{Porte CNOT (NOT contrôlé)}
Le qubit $A$ est le \emph{contrôle}, $B$ la \emph{cible} : $B$ est inversé si $A=1$.
\[
  \mathrm{CNOT}=\ket{0}\bra{0}\otimes I+\ket{1}\bra{1}\otimes X
  =\begin{pmatrix}I & 0\\ 0 & X\end{pmatrix}
  =\begin{pmatrix}1 & 0 & 0 & 0\\ 0 & 1 & 0 & 0\\ 0 & 0 & 0 & 1\\ 0 & 0 & 1 & 0\end{pmatrix}.
\]
\end{definition}

Sa matrice échange les composantes $10$ et $11$ :
$\mathrm{CNOT}\ket{00}=\ket{00}$, $\mathrm{CNOT}\ket{01}=\ket{01}$,
$\mathrm{CNOT}\ket{10}=\ket{11}$, $\mathrm{CNOT}\ket{11}=\ket{10}$, soit
$\ket{i,j}\to\ket{i,i\oplus j}$ (addition modulo $2$).

On applique $H\otimes I$ puis $\mathrm{CNOT}$ à $\ket{00}$ (figure~\ref{fig:bell}) :
\[
  \ket{00}\ \xrightarrow{\ H\otimes I\ }\ \tfrac{1}{\sqrt2}\begin{pmatrix}1\\ 0\\ 1\\ 0\end{pmatrix}
  =\tfrac{1}{\sqrt2}\big(\ket{00}+\ket{10}\big)
  \ \xrightarrow{\ \mathrm{CNOT}\ }\ \tfrac{1}{\sqrt2}\begin{pmatrix}1\\ 0\\ 0\\ 1\end{pmatrix}
  =\tfrac{1}{\sqrt2}\big(\ket{00}+\ket{11}\big)=\ket{\Phi^+}.
\]

\begin{figure}[!ht]
\centering
\begin{tikzpicture}[>=Stealth, font=\small]
  % fils : B en haut (dernier chiffre du ket), A en bas (premier chiffre)
  \draw (0,1) -- (8.6,1);
  \draw (0,0) -- (8.6,0);
  \node[anchor=east] at (0,1) {$\ket{0}_B$};
  \node[anchor=east] at (0,0) {$\ket{0}_A$};
  % Hadamard sur A
  \node[draw=insarouge, fill=insarouge!8, minimum size=7mm, thick] at (2.2,0) {$H$};
  % CNOT : contrôle sur A (bas), cible sur B (haut)
  \draw[insarouge, thick] (5.4,0) -- (5.4,1.3);
  \fill[insarouge] (5.4,0) circle (2.6pt);
  \draw[insarouge, thick, fill=white] (5.4,1) circle (0.3);
  \draw[insarouge, thick] (5.1,1) -- (5.7,1);
  \draw[insarouge, thick] (5.4,0.7) -- (5.4,1.3);
  % accolade de sortie
  \draw[insagris, thick] (8.75,1.2) -- (8.9,1.2) -- (8.9,-0.2) -- (8.75,-0.2);
  \node[anchor=west] at (8.9,0.5) {$\ket{\Phi^+}_{AB}$};
  % états intermédiaires
  \foreach \x in {0.6,3.8,7.2}
    \draw[dashed, insagris!70] (\x,1.35) -- (\x,-0.95);
  \node[anchor=north, font=\footnotesize] at (0.6,-0.95) {$\ket{00}$};
  \node[anchor=north, font=\footnotesize] at (3.8,-0.95) {$\tfrac{1}{\sqrt2}(\ket{00}+\ket{10})$};
  \node[anchor=north, font=\footnotesize] at (7.2,-0.95) {$\tfrac{1}{\sqrt2}(\ket{00}+\ket{11})$};
\end{tikzpicture}
\caption{Circuit qui produit l'état de Bell $\ket{\Phi^+}$ : une porte de Hadamard sur $A$,
puis un CNOT ($A$ contrôle, $B$ cible). Les états intermédiaires sont indiqués sous les
traits pointillés. Convention : les fils sont dessinés dans l'ordre inverse de l'écriture
du ket $\ket{ij}_{AB}$, si bien que le premier chiffre (celui de $A$) est sur le fil du bas.}
\label{fig:bell}
\end{figure}

Les quatre états de base d'entrée donnent les quatre états de Bell :
\begin{center}
\renewcommand{\arraystretch}{1.6}
\begin{tabular}{@{}cccc@{}}
\toprule
Entrée & après $H\otimes I$ & après CNOT & Sortie \\
\midrule
$\ket{00}$ & $\frac{1}{\sqrt2}(\ket{00}+\ket{10})$ & $\frac{1}{\sqrt2}(\ket{00}+\ket{11})$ & $\ket{\Phi^+}$ \\
$\ket{01}$ & $\frac{1}{\sqrt2}(\ket{01}+\ket{11})$ & $\frac{1}{\sqrt2}(\ket{01}+\ket{10})$ & $\ket{\Psi^+}$ \\
$\ket{10}$ & $\frac{1}{\sqrt2}(\ket{00}-\ket{10})$ & $\frac{1}{\sqrt2}(\ket{00}-\ket{11})$ & $\ket{\Phi^-}$ \\
$\ket{11}$ & $\frac{1}{\sqrt2}(\ket{01}-\ket{11})$ & $\frac{1}{\sqrt2}(\ket{01}-\ket{10})$ & $\ket{\Psi^-}$ \\
\bottomrule
\end{tabular}
\end{center}
Sous forme matricielle, le produit $\mathrm{CNOT}\cdot(H\otimes I)$ échange les lignes $3$ et $4$ de
$H\otimes I$ :
\[
  \mathrm{CNOT}\,(H\otimes I)=\frac{1}{\sqrt2}\begin{pmatrix}1 & 0 & 1 & 0\\ 0 & 1 & 0 & 1\\ 0 & 1 & 0 & -1\\ 1 & 0 & -1 & 0\end{pmatrix},
\]
et ses quatre colonnes sont $\ket{\Phi^+}$, $\ket{\Psi^+}$, $\ket{\Phi^-}$, $\ket{\Psi^-}$.

\begin{remarque}[Porte intriquante]
Les états d'entrée $\ket{ij}$ sont séparables ($D=0$) ; les sorties sont intriquées
($D=\pm\frac12$). Aucune porte agissant séparément sur $A$ et sur $B$ ($U_A\otimes U_B$)
ne peut créer de l'intrication à partir d'un produit, puisque
$(U_A\otimes U_B)(\ket{u}\otimes\ket{v})=U_A\ket{u}\otimes U_B\ket{v}$ reste un produit.
Une porte à deux qubits comme le CNOT est indispensable.
\end{remarque}

% ---------------------------------------------------------------------
\subsection{Récapitulatif du chapitre}
% ---------------------------------------------------------------------

\begin{center}
\renewcommand{\arraystretch}{1.4}
\begin{tabular}{@{}ll@{}}
\toprule
\textbf{Notion} & \textbf{Formule} \\
\midrule
Espace de deux qubits & $\mathcal{H}^2\otimes\mathcal{H}^2=\mathcal{H}^4$, \quad dimension $2^n$ pour $n$ qubits \\
Produit de vecteurs & $(a_1,a_2)^{\mathsf T}\otimes(b_1,b_2)^{\mathsf T}=(a_1b_1,a_1b_2,a_2b_1,a_2b_2)^{\mathsf T}$ \\
Base de calcul & $\ket{ij}_{AB}=\ket{i}_A\otimes\ket{j}_B$, \quad $\braket{ij}{kl}=\delta_{ik}\delta_{jl}$ \\
État général & $\ket{\Psi}=\sum c_{ij}\ket{ij}$, \quad $\sum|c_{ij}|^2=1$, \quad $P(ij)=|c_{ij}|^2$ \\
Séparable & $\ket{\Psi}_A\otimes\ket{\Psi}_B$ \quad $\Leftrightarrow$ \quad $c_{00}c_{11}-c_{01}c_{10}=0$ \\
Intriqué & non séparable, \quad par exemple $\ket{\Phi^+}=\frac{1}{\sqrt2}(\ket{00}+\ket{11})$ \\
Matrices & $A\otimes B=(A_{ij}B)$, \quad $(A\otimes B)(\ket{u}\otimes\ket{v})=A\ket{u}\otimes B\ket{v}$ \\
CNOT & $\ket{i,j}\to\ket{i,i\oplus j}$, \quad $\mathrm{CNOT}\,(H\otimes I)\ket{00}=\ket{\Phi^+}$ \\
\bottomrule
\end{tabular}
\end{center}
 % 6. Exemples et exercices corrigés
% =====================================================================
%  Chapitre 7 : QCM sur le cours (100 questions, 4 niveaux).
%  Fichier importé par main.tex avec % =====================================================================
%  Chapitre 7 : QCM sur le cours (100 questions, 4 niveaux).
%  Fichier importé par main.tex avec % =====================================================================
%  Chapitre 7 : QCM sur le cours (100 questions, 4 niveaux).
%  Fichier importé par main.tex avec \include{qcm}. Un chapitre = une \section.
%
%  Organisation :
%    - les énoncés (niveaux 1 à 4) occupent des pages à part ;
%    - le corrigé (clé des réponses, explications, pièges, fiche de score)
%      commence sur une nouvelle page : on peut imprimer les énoncés seuls ;
%    - chaque question est écrite avec \qcmq puis \qcmx (voir plus bas) : le même
%      code produit l'énoncé (cases à cocher) et, plus loin, le corrigé.
%
%  Bascules (définies dans main.tex ; à placer avant \include{qcm}) :
%    \qcmrefsfalse     masque la référence « Chap. · § » sous chaque énoncé ;
%    \qcmcorrigefalse  supprime le corrigé du document (QCM à distribuer seul).
%
%  Ce fichier est produit par outils_qcm/gen_qcm.py à partir de la banque de
%  questions (fichiers texte), mais il reste un fichier LaTeX ordinaire que l'on
%  peut modifier à la main. Pour ajouter une question à la main : insérer un
%  couple \qcmq ... \qcmx dans le bon niveau (la numérotation est automatique),
%  puis mettre à jour (1) les bornes « questions a à b » des titres de niveaux,
%  (2) les lignes \qcmcorr{n} du corrigé, (3) le tableau « Où revoir ? ».
% =====================================================================

% Macros de Dirac (déjà définies dans main.tex ; ce bloc évite l'erreur
% « Undefined control sequence \ket » si main.tex est une ancienne version).
\providecommand{\ket}[1]{|#1\rangle}
\providecommand{\bra}[1]{\langle #1|}
\providecommand{\braket}[2]{\langle #1|#2\rangle}

\makeatletter
% bascules : normalement déjà définies dans main.tex (valeurs par défaut : vrai)
\@ifundefined{ifqcmrefs}{\newif\ifqcmrefs\qcmrefstrue}{}
\@ifundefined{ifqcmcorrige}{\newif\ifqcmcorrige\qcmcorrigetrue}{}
\newcounter{qcmq}

% case à cocher (3 mm) pour le crayon
\newcommand{\qcmbox}{{\setlength{\fboxsep}{0pt}\setlength{\fboxrule}{0.5pt}%
  \raisebox{-0.4mm}{\fbox{\rule{0pt}{3mm}\rule{3mm}{0pt}}}}}

% une proposition : case, lettre, texte (retrait suspendu)
\newcommand{\qcm@opt}[2]{\noindent\hangindent=3.2em \hangafter=1
  \hspace*{0.3em}\qcmbox\hspace{0.45em}\textbf{#1)}\hspace{0.4em}#2\par\vspace{2.5pt}}

% \qcmq[colonnes]{ref}{énoncé}{A}{B}{C}{D}{bonne lettre}   (TeX limite les macros à 9 paramètres)
% \qcmx{com. A}{com. B}{com. C}{com. D}  : commentaires, placés juste après \qcmq
%   - com. de la bonne lettre = explication ; com. des autres = piège ;
%   - colonnes = 1 (une proposition par ligne) ou 2 (grille 2x2).
\newcommand{\qcmq}[8][1]{%
  \stepcounter{qcmq}%
  \expandafter\gdef\csname qcm@ref@\theqcmq\endcsname{#2}%
  \expandafter\gdef\csname qcm@rep@\theqcmq\endcsname{#8}%
  \expandafter\gdef\csname qcm@o@\theqcmq @A\endcsname{#4}%
  \expandafter\gdef\csname qcm@o@\theqcmq @B\endcsname{#5}%
  \expandafter\gdef\csname qcm@o@\theqcmq @C\endcsname{#6}%
  \expandafter\gdef\csname qcm@o@\theqcmq @D\endcsname{#7}%
  \par\vspace{7pt plus 3pt minus 2pt}\noindent
  \begin{minipage}{\linewidth}%
    \noindent{\color{insarouge}\bfseries Question~\theqcmq}%
    \ifqcmrefs\hfill{\footnotesize\color{insagris}#2}\fi\par\nopagebreak
    \vspace{2pt}\noindent #3\par\nopagebreak\vspace{3pt}%
    \ifnum#1=2
      \begin{minipage}[t]{0.495\linewidth}\qcm@opt{A}{#4}\end{minipage}\hfill
      \begin{minipage}[t]{0.495\linewidth}\qcm@opt{B}{#5}\end{minipage}\par\vspace{2pt}
      \begin{minipage}[t]{0.495\linewidth}\qcm@opt{C}{#6}\end{minipage}\hfill
      \begin{minipage}[t]{0.495\linewidth}\qcm@opt{D}{#7}\end{minipage}\par
    \else
      \qcm@opt{A}{#4}\qcm@opt{B}{#5}\qcm@opt{C}{#6}\qcm@opt{D}{#7}%
    \fi
  \end{minipage}\par}

\newcommand{\qcmx}[4]{%
  \expandafter\gdef\csname qcm@c@\theqcmq @A\endcsname{#1}%
  \expandafter\gdef\csname qcm@c@\theqcmq @B\endcsname{#2}%
  \expandafter\gdef\csname qcm@c@\theqcmq @C\endcsname{#3}%
  \expandafter\gdef\csname qcm@c@\theqcmq @D\endcsname{#4}}

% lettre de la bonne réponse de la question n (« ? » si la question n'existe pas)
\newcommand{\qcmkey}[1]{\@ifundefined{qcm@rep@#1}{?}{\csname qcm@rep@#1\endcsname}}

% une mauvaise réponse du corrigé (n° de question, lettre) : ignorée si c'est la bonne.
% (pas de \@for : « := » ne fonctionne pas avec le « : » rendu actif par babel-french)
\newcommand{\qcm@wrongone}[2]{%
  \edef\qcm@a{#2}\edef\qcm@b{\csname qcm@rep@#1\endcsname}%
  \ifx\qcm@a\qcm@b\else
    \noindent\hangindent=1.4em\hangafter=1 {\color{insagris}\textbf{#2)}}\enspace
    \csname qcm@o@#1@#2\endcsname\enspace--\enspace\emph{\csname qcm@c@#1@#2\endcsname}\par
  \fi}

% bloc de corrigé de la question n
\newcommand{\qcmcorr}[1]{%
  \par\vspace{6pt plus 2pt minus 2pt}\noindent
  \begin{minipage}{\linewidth}\small
    \noindent{\color{insarouge}\bfseries Question~#1}\enspace
    \textbf{Réponse : \csname qcm@rep@#1\endcsname}\hfill
    {\footnotesize\color{insagris}\csname qcm@ref@#1\endcsname}\par\nopagebreak
    \vspace{1pt}\noindent\hangindent=1.4em\hangafter=1
    \textbf{\csname qcm@rep@#1\endcsname)}\enspace
    \csname qcm@o@#1@\csname qcm@rep@#1\endcsname\endcsname\enspace---\enspace
    \csname qcm@c@#1@\csname qcm@rep@#1\endcsname\endcsname\par
    \qcm@wrongone{#1}{A}\qcm@wrongone{#1}{B}\qcm@wrongone{#1}{C}\qcm@wrongone{#1}{D}%
  \end{minipage}\par}
\makeatother

% =====================================================================
\section{QCM sur le cours}\label{chap:qcm}
% =====================================================================

\subsection{Mode d'emploi}\label{sec:qcm-mode}

Ce QCM compte \textbf{100 questions} sur les chapitres~\ref{chap:histoire} à~\ref{chap:bell},
classées en quatre niveaux. Chaque question propose quatre réponses dont \textbf{une seule} est
exacte : cochez au crayon la case de la proposition choisie.
\ifqcmrefs À droite du numéro de chaque question, une référence (chapitre, section) indique la
partie du cours concernée.\else\ifqcmcorrige{} La référence de chaque question (chapitre, section)
est donnée dans le corrigé.\fi\fi

\begin{center}
\renewcommand{\arraystretch}{1.3}
\begin{tabular}{@{}llc>{\raggedright\arraybackslash}p{8.2cm}@{}}
\toprule
\textbf{Niveau} & \textbf{Questions} & \textbf{Nombre} & \textbf{Ce qui est demandé}\\
\midrule
1. Débutant & 1 à 30 & 30 & définitions, notations, résultats du cours en une étape\\
2. Intermédiaire & 31 à 60 & 30 & petits calculs, propriétés, application directe d'un résultat\\
3. Avancé & 61 à 90 & 30 & calculs en plusieurs étapes, portes, circuits et protocoles\\
4. Maîtrise & 91 à 100 & 10 & raisonnements fins et pièges conceptuels, qui combinent plusieurs notions\\
\bottomrule
\end{tabular}
\end{center}

\begin{itemize}[nosep]
  \item \textbf{Sur papier} : imprimez seulement les pages~\pageref{chap:qcm} à~\pageref{qcm:fin}
        (mode d'emploi et énoncés, sections~\ref{sec:qcm-n1} à~\ref{sec:qcm-n4}).%
        \ifqcmcorrige{} Le corrigé commence sur une page distincte, à la page~\pageref{sec:qcm-cle}.\fi
  \item \textbf{Barème} : 1 point par bonne réponse, 0 sinon (pas de point négatif). Calculatrice
        autorisée, mais rarement utile (quelques logarithmes et racines).
  \item \textbf{Correction} :%
        \ifqcmcorrige{} comparez avec la clé des réponses, puis lisez l'explication de chaque
        question, même réussie. Les « pièges » décrivent l'erreur qui mène à chaque mauvaise réponse.%
        \else{} le corrigé n'est pas inclus dans cette version du document.\fi
\end{itemize}

\bigskip
\noindent Nom : \makebox[6cm]{\dotfill}\hfill Date : \makebox[3.5cm]{\dotfill}\hfill Score : \makebox[1.6cm]{\dotfill}\,/\,100
\bigskip

\begin{remarque}[Conventions des énoncés]
\begin{itemize}[nosep]
  \item $\ket{ab}=\ket{a}_A\otimes\ket{b}_B$ : le premier chiffre est le qubit $A$, le second le
        qubit $B$ ; dans un circuit, le fil du haut est le dernier chiffre.
  \item Le CNOT de la section~\ref{sec:cnot} a pour contrôle le \emph{second} chiffre ; les portes
        contrôlées de la section~\ref{sec:controlees} ont pour contrôle le \emph{premier}. Chaque
        énoncé précise le contrôle.
  \item Sans précision, une mesure est faite dans la base de calcul $Z$. Les logarithmes des
        entropies sont en base $2$ (résultats en bits).
  \item $\ket{\pm}=\frac{1}{\sqrt2}(\ket0\pm\ket1)$ et $\ket{{\pm}\ii}=\frac{1}{\sqrt2}(\ket0\pm\ii\ket1)$.
\end{itemize}
\end{remarque}

\clearpage
%% ======================================================================
\subsection{Niveau 1 : débutant (questions 1 à 30)}\label{sec:qcm-n1}
%% ======================================================================
\noindent{\color{insagris}\small Niveau de difficulté : Débutant\quad $\bullet\circ\circ\circ$}\par\smallskip
% ---- Question 1 (niveau 1, correct : D)
\qcmq[2]{Chap.~\ref{chap:histoire}\,$\cdot$\,\S\,\ref{sec:dualite}}
  {En 1905, qui explique l'effet photo-électrique en postulant que la lumière est faite de quanta d'énergie $E=hf$, les photons ?}
  {Louis de Broglie}
  {Erwin Schrödinger}
  {Max Planck}
  {Albert Einstein}
  {D}
\qcmx
  {Il propose en 1924 une onde de matière ($\lambda=h/p$) pour les particules comme l'électron ; l'idée du photon date de 1905.}
  {Il écrit en 1926 l'équation d'onde de la matière, plus de vingt ans après l'article sur les photons.}
  {Il introduit le quantum d'énergie en 1900 pour le rayonnement du corps noir ; c'est Einstein qui en fait une propriété de la lumière elle-même.}
  {Il postule que la lumière échange son énergie par paquets $E=hf=\hbar\omega$ : c'est la première face de la dualité onde--corpuscule.}

% ---- Question 2 (niveau 1, correct : A)
\qcmq[1]{Chap.~\ref{chap:histoire}\,$\cdot$\,\S\,\ref{sec:born}}
  {Selon l'interprétation de Max Born, que représente $|\Psi(x,t)|^2$ ?}
  {La densité de probabilité de présence de la particule en $x$ à l'instant $t$}
  {L'énergie cinétique de la particule}
  {L'amplitude de l'onde de matière}
  {La vitesse de la particule}
  {A}
\qcmx
  {$|\Psi(x,t)|^2\,\dd x$ est la probabilité de trouver la particule entre $x$ et $x+\dd x$ ; $\Psi$ lui-même n'est pas mesurable.}
  {L'énergie est associée à l'opérateur hamiltonien $\Hop$ ; $|\Psi|^2$ n'a pas la dimension d'une énergie.}
  {$\Psi$ est l'amplitude (complexe) ; $|\Psi|^2$ en est le carré du module.}
  {La vitesse s'obtient avec l'opérateur $\hat p=-\ii\hbar\,\partial_x$, pas avec $|\Psi|^2$.}

% ---- Question 3 (niveau 1, correct : D)
\qcmq[2]{Chap.~\ref{chap:histoire}\,$\cdot$\,\S\,\ref{sec:dualite}}
  {Une particule de quantité de mouvement $p$ est associée, d'après de Broglie, à une onde de longueur d'onde $\lambda$. Quelle relation est correcte ?}
  {$\lambda=hp$}
  {$\lambda=p/h$}
  {$\lambda=\hbar/p$}
  {$\lambda=h/p$}
  {D}
\qcmx
  {Produit au lieu du quotient : $\lambda$ grandirait avec $p$, le contraire de ce qu'on observe.}
  {Quotient inversé : $p/h=1/\lambda$ s'exprime en m$^{-1}$, pas en mètres.}
  {Confusion entre $h$ et $\hbar=h/2\pi$ : $\hbar/p=\lambda/2\pi$ est l'inverse du nombre d'onde $k$.}
  {Relation de de Broglie : plus la particule est rapide ou lourde, plus $\lambda$ est petite (balle de $0{,}1$~kg à $10$~m/s : $\lambda\approx10^{-34}$~m, inobservable).}

% ---- Question 4 (niveau 1, correct : C)
\qcmq[2]{Chap.~\ref{chap:histoire}\,$\cdot$\,\S\,\ref{sec:construction}}
  {Quelle est l'équation de Schrödinger dépendante du temps pour un état $\Psi$ et un hamiltonien $\Hop$ ?}
  {$-\ii\hbar\,\partial_t\Psi=\Hop\Psi$}
  {$\hbar\,\partial_t\Psi=\Hop\Psi$}
  {$\ii\hbar\,\partial_t\Psi=\Hop\Psi$}
  {$\ii\hbar\,\partial_x\Psi=\Hop\Psi$}
  {C}
\qcmx
  {Erreur de signe : la solution serait $\ee^{+\ii\Hop t/\hbar}\Psi(0)$, qui décrit l'évolution à rebours du temps.}
  {Sans le facteur $\ii$, les solutions sont des exponentielles réelles (croissantes ou décroissantes) et la norme n'est pas conservée.}
  {Elle relie l'évolution de l'état à l'énergie ; pour un $\Hop$ indépendant du temps, sa solution est $\Psi(t)=\ee^{-\ii\Hop t/\hbar}\Psi(0)$ (chapitre~\ref{chap:portes}).}
  {C'est la dérivée en $t$ qui décrit l'évolution ; la dérivée en $x$ sert à construire $\hat p=-\ii\hbar\,\partial_x$.}

% ---- Question 5 (niveau 1, correct : A)
\qcmq[1]{Chap.~\ref{chap:histoire}\,$\cdot$\,\S\,\ref{sec:construction}}
  {Quel opérateur représente l'énergie totale d'un système quantique ?}
  {L'hamiltonien $\Hop$}
  {L'opérateur quantité de mouvement $\hat p$}
  {L'opérateur position $\hat x$}
  {L'opérateur identité}
  {A}
\qcmx
  {Somme de l'énergie cinétique $-\frac{\hbar^2}{2m}\partial_x^2$ et de l'énergie potentielle $V$ ; ses valeurs propres sont les énergies mesurables.}
  {Il représente l'impulsion ; l'énergie cinétique est $\hat p^2/2m$, et il faut y ajouter $V$.}
  {C'est la multiplication par $x$ : il donne la position, pas l'énergie.}
  {Il vaut $1$ pour tout état : il ne contient aucune information sur l'énergie.}

% ---- Question 6 (niveau 1, correct : C)
\qcmq[1]{Chap.~\ref{chap:histoire}\,$\cdot$\,\S\,\ref{sec:histoire}}
  {Dans quel ordre chronologique ces avancées ont-elles eu lieu ?}
  {de Broglie $\to$ Planck $\to$ Schrödinger $\to$ Bell}
  {Planck $\to$ Schrödinger $\to$ de Broglie $\to$ Bell}
  {Planck (quantum d'énergie) $\to$ de Broglie (onde de matière) $\to$ Schrödinger (équation d'onde) $\to$ Bell (inégalités)}
  {Planck $\to$ de Broglie $\to$ Bell $\to$ Schrödinger}
  {C}
\qcmx
  {Le quantum d'énergie de Planck (1900) précède l'onde de matière de de Broglie (1924).}
  {L'équation de Schrödinger (1926) généralise l'idée d'onde de matière de de Broglie (1924) : elle vient après.}
  {Planck 1900, de Broglie 1924, Schrödinger 1926, Bell 1964.}
  {Les inégalités de Bell (1964) viennent près de quarante ans après l'équation de Schrödinger.}

% ---- Question 7 (niveau 1, correct : C)
\qcmq[1]{Chap.~\ref{chap:dirac}\,$\cdot$\,\S\,\ref{sec:dirac}}
  {En notation de Dirac, comment appelle-t-on $\ket{\Psi}$ et comment le représente-t-on pour un qubit ?}
  {Un bra, représenté par un vecteur ligne à deux composantes complexes}
  {Un ket, représenté par un nombre complexe}
  {Un ket, représenté par un vecteur colonne à deux composantes complexes}
  {Un bra, représenté par une matrice $2\times2$}
  {C}
\qcmx
  {Le bra $\bra{\Psi}$ est le conjugué transposé du ket : un vecteur ligne, mais ce n'est pas la notation $\ket{\Psi}$.}
  {Un scalaire ne peut pas décrire une superposition : il faut deux amplitudes $\alpha$ et $\beta$.}
  {Le ket est le vecteur colonne $\binom{\alpha}{\beta}$ ; son bra $\bra{\Psi}$ est le vecteur ligne conjugué.}
  {Une matrice est un opérateur ; le bra est une matrice ligne $1\times2$.}

% ---- Question 8 (niveau 1, correct : B)
\qcmq[2]{Chap.~\ref{chap:dirac}\,$\cdot$\,\S\,\ref{sec:qubit}}
  {Quelle condition les amplitudes $\alpha$ et $\beta$ de $\alpha\ket{0}+\beta\ket{1}$ doivent-elles vérifier pour que l'état soit normalisé ?}
  {$\alpha+\beta=1$}
  {$|\alpha|^2+|\beta|^2=1$}
  {$\alpha^2+\beta^2=1$}
  {$|\alpha|+|\beta|=1$}
  {B}
\qcmx
  {On somme les amplitudes (complexes, parfois négatives) ; ce sont leurs modules au carré qui sont des probabilités.}
  {Les probabilités de mesure $P(0)=|\alpha|^2$ et $P(1)=|\beta|^2$ doivent sommer à $1$.}
  {Oubli du module : pour $\alpha=\frac{1}{\sqrt2}$ et $\beta=\frac{\ii}{\sqrt2}$, $\alpha^2+\beta^2=0$ alors que l'état est normalisé.}
  {Somme des modules au lieu des carrés : $\alpha=\beta=\frac{1}{\sqrt2}$ est normalisé, mais $|\alpha|+|\beta|=\sqrt2$.}

% ---- Question 9 (niveau 1, correct : C)
\qcmq[1]{Chap.~\ref{chap:dirac}\,$\cdot$\,\S\,\ref{sec:projecteurs}}
  {On mesure dans la base $Z$ un qubit préparé dans $\ket{+}=\frac{1}{\sqrt2}(\ket{0}+\ket{1})$. Que se passe-t-il ?}
  {On lit $+$ et le qubit reste dans $\ket{+}$}
  {On lit toujours $0$}
  {On lit $0$ ou $1$, chacun avec la probabilité $\frac12$, et l'état devient $\ket{0}$ ou $\ket{1}$}
  {On lit $0$ et $1$ à la fois}
  {C}
\qcmx
  {Une mesure en base $Z$ ne donne que $0$ ou $1$ ; elle modifie l'état.}
  {Les deux amplitudes ont le même module : $P(0)=P(1)=\frac12$.}
  {Postulat de la mesure : la superposition s'effondre sur l'état propre correspondant au résultat.}
  {Chaque mesure donne un seul résultat ; la superposition ne s'observe qu'à travers les statistiques.}

% ---- Question 10 (niveau 1, correct : B)
\qcmq[2]{Chap.~\ref{chap:dirac}\,$\cdot$\,\S\,\ref{sec:qubit}}
  {Que vaut le produit scalaire $\braket{0}{1}$ ?}
  {$1$}
  {$0$}
  {$\frac{1}{\sqrt2}$}
  {$\ket{01}$}
  {B}
\qcmx
  {C'est la valeur de $\braket{0}{0}$ et de $\braket{1}{1}$ (vecteurs normés), pas de $\braket{0}{1}$.}
  {$\ket{0}$ et $\ket{1}$ sont orthogonaux : $\begin{pmatrix}1&0\end{pmatrix}\begin{pmatrix}0\\1\end{pmatrix}=0$.}
  {C'est la valeur de $\braket{0}{+}$, pas de $\braket{0}{1}$.}
  {Confusion avec le produit tensoriel : un produit scalaire est un nombre, pas un ket.}

% ---- Question 11 (niveau 1, correct : D)
\qcmq[2]{Chap.~\ref{chap:dirac}\,$\cdot$\,\S\,\ref{sec:qubit}}
  {Un qubit est dans l'état $\ket{-}=\frac{1}{\sqrt2}(\ket{0}-\ket{1})$. Quelle est la probabilité de mesurer $1$ en base $Z$ ?}
  {$0$}
  {$-\frac12$}
  {$\frac{1}{\sqrt2}$}
  {$\frac12$}
  {D}
\qcmx
  {Le signe $-$ ne « retire » pas $\ket{1}$ de la superposition.}
  {Une probabilité est toujours positive : on prend le module au carré de l'amplitude.}
  {C'est l'amplitude en module, pas la probabilité (il faut élever au carré).}
  {$P(1)=\left|-\frac{1}{\sqrt2}\right|^2=\frac12$ : le signe est une phase relative, qui disparaît dans le module.}

% ---- Question 12 (niveau 1, correct : B)
\qcmq[2]{Chap.~\ref{chap:dirac}\,$\cdot$\,\S\,\ref{sec:projecteurs}}
  {Quel est le projecteur sur l'état $\ket{0}$ ?}
  {$\ket{0}\bra{1}=\begin{pmatrix}0&1\\0&0\end{pmatrix}$}
  {$\hat P_0=\ket{0}\bra{0}=\begin{pmatrix}1&0\\0&0\end{pmatrix}$}
  {$\braket{0}{0}=1$}
  {$\begin{pmatrix}0&0\\0&1\end{pmatrix}$}
  {B}
\qcmx
  {Cet opérateur envoie $\ket{1}$ sur $\ket{0}$ ; son carré est nul : ce n'est pas un projecteur.}
  {$\hat P_0\ket{\Psi}=\alpha\ket{0}$ : il garde la composante sur $\ket{0}$ ; $\hat P_0^2=\hat P_0$.}
  {C'est un nombre, pas un opérateur.}
  {C'est $\ket{1}\bra{1}$, le projecteur sur $\ket{1}$.}

% ---- Question 13 (niveau 1, correct : A)
\qcmq[1]{Chap.~\ref{chap:operateurs}\,$\cdot$\,\S\,\ref{sec:bloch}}
  {Où se trouvent les états $\ket{0}$ et $\ket{1}$ sur la sphère de Bloch ?}
  {Aux deux pôles, Nord et Sud}
  {Aux deux extrémités de l'axe $x$}
  {Aux deux extrémités de l'axe $y$}
  {En deux points voisins de l'équateur}
  {A}
\qcmx
  {$\ket{0}$ au pôle Nord ($\theta=0$), $\ket{1}$ au pôle Sud ($\theta=\pi$) : deux états orthogonaux sont diamétralement opposés.}
  {Ce sont les places de $\ket{+}$ et $\ket{-}$.}
  {Ce sont les places de $\ket{{+}\ii}$ et $\ket{{-}\ii}$.}
  {Des états orthogonaux sont antipodaux, jamais voisins.}

% ---- Question 14 (niveau 1, correct : A)
\qcmq[2]{Chap.~\ref{chap:operateurs}\,$\cdot$\,\S\,\ref{sec:tenseur}}
  {Quelle est la dimension de l'espace de Hilbert de $n$ qubits ?}
  {$2^n$}
  {$2n$}
  {$n^2$}
  {$2^n-1$}
  {A}
\qcmx
  {Le produit tensoriel multiplie les dimensions : $2\times2\times\cdots\times2$ ; deux qubits : $4$, trois qubits : $8$.}
  {On additionne les dimensions au lieu de les multiplier ; pour $n=3$ on trouverait $6$ au lieu de $8$.}
  {Juste pour $n=2$ par coïncidence seulement : pour $n=3$ on trouverait $9$ au lieu de $8$.}
  {Il y a bien $2^n$ vecteurs dans la base de calcul, $\ket{0\cdots0}$ compris.}

% ---- Question 15 (niveau 1, correct : C)
\qcmq[1]{Chap.~\ref{chap:operateurs}\,$\cdot$\,\S\,\ref{sec:rho}}
  {Que vaut toujours la trace d'une matrice de densité $\rho$ ?}
  {$0$}
  {$1$ pour un état pur, et moins de $1$ pour un état mixte}
  {$1$}
  {$2$, la dimension de l'espace}
  {C}
\qcmx
  {C'est la trace des matrices de Pauli, pas d'une matrice de densité.}
  {Confusion avec $\operatorname{Tr}(\rho^2)$ : c'est la pureté qui distingue les deux cas, la trace vaut toujours $1$.}
  {C'est la somme des probabilités de mesure (les populations), que l'état soit pur ou mixte.}
  {C'est la trace de l'identité $\operatorname{Tr}\mathrm{Id}=2$, pas celle de $\rho$.}

% ---- Question 16 (niveau 1, correct : B)
\qcmq[1]{Chap.~\ref{chap:operateurs}\,$\cdot$\,\S\,\ref{sec:pur-mixte}}
  {Quelle est la valeur de $\operatorname{Tr}(\rho^2)$ pour un état pur ?}
  {$0$}
  {$1$}
  {$\frac12$}
  {Une valeur strictement inférieure à $1$}
  {B}
\qcmx
  {$\rho^2=0$ imposerait $\rho=0$, ce qui contredit $\operatorname{Tr}\rho=1$.}
  {Un état pur s'écrit $\rho=\ket{\Psi}\bra{\Psi}$, donc $\rho^2=\rho$ et $\operatorname{Tr}\rho^2=\operatorname{Tr}\rho=1$.}
  {C'est la valeur minimale pour un qubit, atteinte par l'état maximalement mixte $\mathrm{Id}/2$.}
  {C'est la caractérisation d'un état mixte.}

% ---- Question 17 (niveau 1, correct : B)
\qcmq[1]{Chap.~\ref{chap:operateurs}\,$\cdot$\,\S\,\ref{sec:rho}}
  {Que représentent les éléments diagonaux $\rho_{00}$ et $\rho_{11}$ de la matrice de densité d'un qubit ?}
  {Les cohérences entre $\ket{0}$ et $\ket{1}$}
  {Les probabilités de mesurer $\ket{0}$ et $\ket{1}$ (les populations)}
  {Les valeurs propres de $\rho$}
  {Les amplitudes $\alpha$ et $\beta$}
  {B}
\qcmx
  {Les cohérences sont les éléments hors diagonale $\rho_{01}$ et $\rho_{10}$.}
  {$\rho_{00}=\bra{0}\rho\ket{0}=P(0)$ et $\rho_{11}=P(1)$, d'où $\operatorname{Tr}\rho=1$.}
  {Seulement si $\rho$ est diagonale dans la base de calcul : $\ket{+}\bra{+}$ a pour diagonale $(\frac12,\frac12)$ mais pour valeurs propres $(1,0)$.}
  {Pour un état pur, $\rho_{00}=|\alpha|^2$ : le module au carré, pas l'amplitude.}

% ---- Question 18 (niveau 1, correct : D)
\qcmq[1]{Chap.~\ref{chap:operateurs}\,$\cdot$\,\S\,\ref{sec:op-dirac}}
  {Pourquoi associe-t-on aux grandeurs mesurables (observables) des opérateurs hermitiens ?}
  {Parce qu'ils conservent la norme des vecteurs}
  {Parce qu'ils ont tous une trace nulle}
  {Parce que leurs valeurs propres valent toutes $\pm1$}
  {Parce que leurs valeurs propres sont réelles, comme les résultats d'une mesure}
  {D}
\qcmx
  {C'est la propriété des opérateurs unitaires (les portes), pas des opérateurs hermitiens.}
  {Vrai pour les matrices de Pauli, faux en général (l'identité est hermitienne, de trace $2$), et sans rapport avec le caractère mesurable.}
  {Vrai pour $X$, $Y$, $Z$ seulement ; l'hamiltonien du puits a les valeurs propres $E_n=n^2E_1$.}
  {Un opérateur hermitien ($A^\dagger=A$) a des valeurs propres réelles et des vecteurs propres orthogonaux ; les résultats possibles sont ses valeurs propres.}

% ---- Question 19 (niveau 1, correct : C)
\qcmq[2]{Chap.~\ref{chap:portes}\,$\cdot$\,\S\,\ref{sec:pauli}}
  {Quelle porte de Pauli réalise l'inversion de bit $\ket{0}\leftrightarrow\ket{1}$ ?}
  {$Z$}
  {$Y$}
  {$X$}
  {$H$}
  {C}
\qcmx
  {$Z$ laisse $\ket{0}$ et change $\ket{1}$ en $-\ket{1}$ : c'est une inversion de phase.}
  {$Y\ket{0}=\ii\ket{1}$ : elle inverse le bit mais ajoute une phase ; la porte d'inversion « pure » est $X$.}
  {$X\ket{0}=\ket{1}$ et $X\ket{1}=\ket{0}$ : c'est le NOT quantique.}
  {$H$ crée des superpositions : $H\ket{0}=\ket{+}$.}

% ---- Question 20 (niveau 1, correct : B)
\qcmq[2]{Chap.~\ref{chap:portes}\,$\cdot$\,\S\,\ref{sec:hadamard}}
  {Que donne la porte de Hadamard appliquée à $\ket{0}$ ?}
  {$\ket{1}$}
  {$\ket{+}=\frac{1}{\sqrt2}(\ket{0}+\ket{1})$}
  {$\ket{-}=\frac{1}{\sqrt2}(\ket{0}-\ket{1})$}
  {$\ket{0}$}
  {B}
\qcmx
  {C'est l'action de $X$.}
  {Superposition à poids égaux : $P(0)=P(1)=\frac12$ ; c'est la première étape de la plupart des algorithmes.}
  {C'est $H\ket{1}$ ; le signe $-$ n'apparaît qu'à partir de $\ket{1}$.}
  {$H$ n'est pas l'identité ; elle ne laisse inchangés ni $\ket{0}$ ni $\ket{1}$ (elle envoie $\ket{+}$ sur $\ket{0}$).}

% ---- Question 21 (niveau 1, correct : D)
\qcmq[2]{Chap.~\ref{chap:portes}\,$\cdot$\,\S\,\ref{sec:portes}}
  {Quelle propriété doit vérifier la matrice $U$ d'une porte quantique ?}
  {Elle est hermitienne : $U^\dagger=U$}
  {Sa trace vaut $1$}
  {Ses coefficients sont réels}
  {Elle est unitaire : $U^\dagger U=\mathrm{Id}$}
  {D}
\qcmx
  {Vrai pour $X$, $Y$, $Z$, $H$ mais faux pour $S$ et $T$ : ce n'est pas la condition générale.}
  {C'est une propriété des matrices de densité ; $\operatorname{Tr}X=0$ et $\operatorname{Tr}\mathrm{Id}=2$.}
  {$Y$ et $S$ sont des portes à coefficients complexes.}
  {Elle conserve la norme ($\braket{\Psi}{\Psi}=1$ reste vrai) et elle est réversible.}

% ---- Question 22 (niveau 1, correct : A)
\qcmq[2]{Chap.~\ref{chap:portes}\,$\cdot$\,\S\,\ref{sec:portes}}
  {Quel est l'inverse d'une porte quantique $U$ ?}
  {$U^{-1}=U^\dagger$ (transposée conjuguée)}
  {$U^{-1}=U^{\mathsf T}$ (transposée seule)}
  {$U^{-1}=U^{*}$ (conjuguée seule)}
  {$U^{-1}=-U$}
  {A}
\qcmx
  {Conséquence de $U^\dagger U=\mathrm{Id}$ ; par exemple $S^{-1}=S^\dagger=\mathrm{diag}(1,-\ii)$.}
  {Oubli de la conjugaison : $S^{\mathsf T}=S=\mathrm{diag}(1,\ii)$ alors que $S^{-1}=\mathrm{diag}(1,-\ii)$.}
  {Faux en général : $Y^*=-Y$ alors que $Y^{-1}=Y$.}
  {Faux : $X^{-1}=X\neq-X$.}

% ---- Question 23 (niveau 1, correct : C)
\qcmq[1]{Chap.~\ref{chap:portes}\,$\cdot$\,\S\,\ref{sec:cnot}}
  {Dans une porte CNOT, que devient le qubit cible lorsque le qubit de contrôle est dans l'état $\ket{0}$ ?}
  {Il est inversé par la porte $X$}
  {Il est remis à $\ket{0}$}
  {Il reste inchangé}
  {Il subit un changement de signe, comme avec $Z$}
  {C}
\qcmx
  {C'est ce qui se passe quand le contrôle vaut $\ket{1}$.}
  {Le CNOT est réversible : remettre la cible à zéro détruirait de l'information.}
  {Avec le contrôle à $\ket{0}$, la porte agit comme l'identité sur la cible.}
  {Ce serait la porte $Z$ contrôlée (CZ), pas le CNOT.}

% ---- Question 24 (niveau 1, correct : A)
\qcmq[2]{Chap.~\ref{chap:portes}\,$\cdot$\,\S\,\ref{sec:cnot}}
  {Le CNOT du cours a pour contrôle le second chiffre du ket et pour cible le premier : $\mathrm{CNOT}\ket{a,b}=\ket{a\oplus b,\,b}$. Que donne-t-il sur $\ket{0,1}$ ?}
  {$\ket{1,1}$}
  {$\ket{0,1}$}
  {$\ket{0,0}$}
  {$\ket{1,0}$}
  {A}
\qcmx
  {Le contrôle $b=1$ inverse la cible : $0\oplus1=1$.}
  {Le contrôle vaut $1$ : la cible est bien inversée.}
  {C'est le contrôle qui serait inversé : le contrôle n'est jamais modifié.}
  {Les deux qubits seraient inversés : seule la cible change.}

% ---- Question 25 (niveau 1, correct : D)
\qcmq[1]{Chap.~\ref{chap:bell}\,$\cdot$\,\S\,\ref{sec:bell}}
  {Quelle propriété caractérise les quatre états de Bell ?}
  {Ils sont séparables : chaque qubit a son propre état}
  {Ils décrivent un seul qubit}
  {Ils ne diffèrent que par une phase globale}
  {Ils sont maximalement intriqués et forment une base orthonormée de l'espace de deux qubits}
  {D}
\qcmx
  {$D=c_{00}c_{11}-c_{01}c_{10}=\pm\frac12\neq0$ : ils ne se factorisent pas.}
  {Ce sont des états de deux qubits (quatre amplitudes sur $\ket{00},\ket{01},\ket{10},\ket{11}$).}
  {Ils sont orthogonaux deux à deux : ce sont quatre états distincts, alors qu'une phase globale ne change pas l'état.}
  {Aucun n'est un produit d'états à un qubit ($|D|=\frac12$, valeur maximale) et ils sont orthogonaux deux à deux.}

% ---- Question 26 (niveau 1, correct : D)
\qcmq[1]{Chap.~\ref{chap:bell}\,$\cdot$\,\S\,\ref{sec:teleportation}}
  {Quel est l'objectif du protocole de téléportation quantique ?}
  {Copier l'état d'un qubit inconnu sur le qubit de Bob}
  {Transmettre une information plus vite que la lumière}
  {Déplacer la particule physique d'Alice vers Bob}
  {Transférer l'état d'un qubit inconnu d'Alice à Bob, avec une paire intriquée et deux bits classiques}
  {D}
\qcmx
  {La copie est interdite (non-clonage) : après le protocole, l'état n'existe plus que chez Bob.}
  {Sans les deux bits classiques, le qubit de Bob est dans l'état $\mathrm{Id}/2$ : il n'a reçu aucune information.}
  {C'est l'état quantique qui est transféré, pas la particule qui le porte.}
  {Le qubit lui-même n'est pas envoyé : son état est reconstruit chez Bob, et détruit chez Alice par la mesure.}

% ---- Question 27 (niveau 1, correct : B)
\qcmq[1]{Chap.~\ref{chap:bell}\,$\cdot$\,\S\,\ref{sec:teleportation}}
  {Dans la téléportation, que doit envoyer Alice à Bob sur un canal classique pour qu'il retrouve l'état ?}
  {Un qubit intriqué}
  {Deux bits classiques, le résultat de sa mesure}
  {Les amplitudes exactes $\alpha$ et $\beta$ de l'état}
  {Un seul bit}
  {B}
\qcmx
  {La paire intriquée est partagée avant le protocole ; on ne l'envoie pas après.}
  {Les quatre résultats $(a,q)$ correspondent aux quatre corrections possibles ($I$, $Z$, $X$, $X$ puis $Z$) : il faut deux bits.}
  {Alice ne connaît pas l'état à téléporter, et ces nombres complexes n'ont pas de description finie en bits.}
  {Un bit ne distingue que deux cas, alors que Bob doit choisir parmi quatre corrections.}

% ---- Question 28 (niveau 1, correct : A)
\qcmq[1]{Chap.~\ref{chap:bell}\,$\cdot$\,\S\,\ref{sec:bell} ; ex.\,\ref{ex:superdense}}
  {Quel est le but du codage superdense ?}
  {Transmettre deux bits classiques en envoyant un seul qubit, grâce à une paire intriquée partagée}
  {Transmettre l'état d'un qubit avec deux bits classiques}
  {Compresser plusieurs qubits en un seul}
  {Chiffrer un message de façon inviolable}
  {A}
\qcmx
  {Alice applique l'une des quatre portes $I$, $Z$, $X$, $XZ$ à son qubit de la paire, l'envoie, et Bob identifie l'état de Bell obtenu.}
  {C'est la téléportation, le protocole inverse.}
  {Le protocole ne compresse pas de qubits : il fait passer deux bits classiques par un canal à un qubit.}
  {Ce n'est pas un protocole de cryptographie : la paire intriquée sert à doubler la capacité du canal.}

% ---- Question 29 (niveau 1, correct : A)
\qcmq[1]{Chap.~\ref{chap:bell}\,$\cdot$\,\S\,\ref{sec:teleportation}}
  {Que dit le théorème de non-clonage ?}
  {Aucune opération ne peut copier parfaitement un état quantique inconnu quelconque}
  {Il est impossible de mesurer un qubit sans le perturber}
  {Il est impossible de créer de l'intrication}
  {Il est impossible de copier un bit classique}
  {A}
\qcmx
  {Une porte $U$ qui copierait $\ket{0}$ et $\ket{1}$ enverrait $\ket{+}\ket{0}$ sur un état intriqué, par linéarité, et non sur $\ket{+}\ket{+}$.}
  {C'est l'effondrement de l'état, un autre principe.}
  {L'intrication se crée très bien (Hadamard puis CNOT).}
  {Le CNOT copie un bit classique : $\ket{0,b}\to\ket{b,b}$.}

% ---- Question 30 (niveau 1, correct : B)
\qcmq[1]{Chap.~\ref{chap:bell}\,$\cdot$\,\S\,\ref{sec:circuit-bell}}
  {Quel enchaînement de portes transforme $\ket{00}$ en l'état de Bell $\ket{\Phi^+}=\frac{1}{\sqrt2}(\ket{00}+\ket{11})$ ?}
  {CNOT, puis Hadamard sur le contrôle}
  {Hadamard sur l'un des qubits, puis CNOT dont le contrôle est ce même qubit}
  {Hadamard sur chacun des deux qubits}
  {SWAP, puis $X$ sur un qubit}
  {B}
\qcmx
  {Le CNOT ne fait rien sur $\ket{00}$ (contrôle à $0$) ; $H$ donne ensuite un état séparable.}
  {$H$ crée $\ket{0}\otimes\ket{+}$, puis le CNOT recopie la superposition : $\frac{1}{\sqrt2}(\ket{00}+\ket{11})$.}
  {On obtient $\ket{+}\otimes\ket{+}$, séparable : des portes locales ne créent pas d'intrication.}
  {$\ket{00}$ devient $\ket{01}$ ou $\ket{10}$, un état de base séparable.}

\clearpage
%% ======================================================================
\subsection{Niveau 2 : intermédiaire (questions 31 à 60)}\label{sec:qcm-n2}
%% ======================================================================
\noindent{\color{insagris}\small Niveau de difficulté : Intermédiaire\quad $\bullet\bullet\circ\circ$}\par\smallskip
% ---- Question 31 (niveau 2, correct : A)
\qcmq[1]{Chap.~\ref{chap:histoire}\,$\cdot$\,\S\,\ref{sec:born}}
  {Quelle condition la fonction d'onde d'une particule doit-elle vérifier sur tout l'espace ?}
  {$\displaystyle\int_{-\infty}^{+\infty}|\Psi(x,t)|^2\,\dd x=1$}
  {$\displaystyle\int_{-\infty}^{+\infty}\Psi(x,t)\,\dd x=1$}
  {$|\Psi(x,t)|^2=1$ pour tout $x$}
  {$\partial_x\Psi(x,t)=0$ en tout point}
  {A}
\qcmx
  {La particule se trouve forcément quelque part : la probabilité totale vaut $1$ (condition de normalisation).}
  {Sans le module au carré : $\Psi$ est complexe et son intégrale n'est pas une probabilité (elle peut être nulle ou complexe).}
  {$|\Psi|^2$ est une densité : si elle valait $1$ partout, son intégrale divergerait.}
  {Une fonction constante n'est pas normalisable ($\int|\Psi|^2\dd x$ diverge) ; la condition porte sur l'intégrale de $|\Psi|^2$.}

% ---- Question 32 (niveau 2, correct : D)
\qcmq[2]{Chap.~\ref{chap:histoire}\,$\cdot$\,\S\,\ref{sec:operateurs}}
  {Quelle est l'expression de l'opérateur quantité de mouvement $\hat p$ en une dimension ?}
  {$\hat p=m\,v$}
  {$\hat p=-\dfrac{\hbar^2}{2m}\dfrac{\partial^2}{\partial x^2}$}
  {$\hat p=+\ii\hbar\,\dfrac{\partial}{\partial x}$}
  {$\hat p=-\ii\hbar\,\dfrac{\partial}{\partial x}$}
  {D}
\qcmx
  {C'est la grandeur classique ; en mécanique quantique $p$ est un opérateur différentiel.}
  {C'est l'énergie cinétique $\hat p^2/2m$.}
  {Erreur de signe : on trouverait $p=-\hbar k$ pour l'onde $\ee^{\ii kx}$.}
  {Convention du cours : $\hat p\,\ee^{\ii kx}=\hbar k\,\ee^{\ii kx}$, donc une onde vers les $x$ croissants a $p=\hbar k>0$.}

% ---- Question 33 (niveau 2, correct : B)
\qcmq[2]{Chap.~\ref{chap:histoire}\,$\cdot$\,\S\,\ref{sec:solutions}}
  {Pour une particule dans un puits de potentiel infini de largeur $a$, comment varie l'énergie $E_n$ avec le numéro $n$ du niveau ?}
  {$E_n=E_1/n$}
  {$E_n=n^2E_1$, avec $E_1=\dfrac{\pi^2\hbar^2}{2ma^2}$}
  {$E_n=nE_1$}
  {$E_n=E_1$ pour tout $n$}
  {B}
\qcmx
  {Les niveaux se resserreraient quand $n$ augmente ; c'est l'inverse pour le puits infini.}
  {$k_n=\frac{n\pi}{a}$ et $E_n=\frac{\hbar^2k_n^2}{2m}$ : les niveaux s'écartent quand $n$ croît ($E_{n+1}-E_n=(2n+1)E_1$).}
  {Niveaux équidistants : ce n'est pas le puits infini, où $E_2=4E_1$ et $E_3=9E_1$.}
  {Il n'y aurait pas de quantification : tous les états auraient la même énergie.}

% ---- Question 34 (niveau 2, correct : C)
\qcmq[2]{Chap.~\ref{chap:histoire}\,$\cdot$\,\S\,\ref{sec:solutions}}
  {Une particule est dans un état stationnaire $\Psi_n$ de l'hamiltonien ($\Hop\Psi_n=E_n\Psi_n$). Que vaut l'écart-type $\sigma_E$ de son énergie ?}
  {$E_n$}
  {$E_n^2$}
  {$0$}
  {$\sqrt{E_n}$}
  {C}
\qcmx
  {C'est la valeur moyenne de l'énergie, pas son écart-type.}
  {C'est la valeur de $\langle\Hop^2\rangle$, avant soustraction de $\langle\Hop\rangle^2$.}
  {$\langle\Hop\rangle=E_n$ et $\langle\Hop^2\rangle=E_n^2$, donc $\sigma_E=\sqrt{E_n^2-E_n^2}=0$ : l'énergie est parfaitement déterminée.}
  {Mélange avec la formule $\sigma=\sqrt{\langle\Hop^2\rangle-\langle\Hop\rangle^2}$ : les deux termes sont égaux à $E_n^2$.}

% ---- Question 35 (niveau 2, correct : C)
\qcmq[2]{Chap.~\ref{chap:histoire}\,$\cdot$\,\S\,\ref{sec:solutions}}
  {Pour un état stationnaire $\Psi_n$ du puits infini $]0,a[$, que vaut la position moyenne $\langle\hat x\rangle$ ?}
  {$a/n$}
  {$a/4$}
  {$a/2$ pour tout $n$}
  {$0$}
  {C}
\qcmx
  {Elle dépendrait de $n$, ce que la symétrie de $|\Psi_n|^2$ interdit.}
  {C'est la position d'un maximum de $|\Psi_2|^2$ ; la moyenne des deux maxima $\frac a4$ et $\frac{3a}{4}$ donne $\frac a2$.}
  {$|\Psi_n|^2=\frac2a\sin^2\frac{n\pi x}{a}$ est symétrique par rapport au centre du puits.}
  {$\Psi_n$ s'annule sur la paroi, mais $\langle\hat x\rangle$ est une moyenne pondérée sur tout le puits.}

% ---- Question 36 (niveau 2, correct : D)
\qcmq[2]{Chap.~\ref{chap:dirac}\,$\cdot$\,\S\,\ref{sec:qubit}}
  {Un qubit est dans l'état $\ket{\Psi}=\frac{1}{\sqrt5}\ket{0}+\frac{2\ii}{\sqrt5}\ket{1}$. Quelle est la probabilité de mesurer $1$ en base $Z$ ?}
  {$\frac25$}
  {$-\frac45$}
  {$\frac15$}
  {$\frac45$}
  {D}
\qcmx
  {On divise le coefficient $2$ par $5$ : on oublie à la fois la racine et le carré.}
  {On élève $\frac{2\ii}{\sqrt5}$ au carré sans conjuguer : $\left(\frac{2\ii}{\sqrt5}\right)^2=-\frac45$ ; une probabilité est $|\cdot|^2$.}
  {C'est $P(0)=\left|\frac{1}{\sqrt5}\right|^2$.}
  {$P(1)=\left|\frac{2\ii}{\sqrt5}\right|^2=\frac45$ : le facteur $\ii$ est une phase relative, invisible en base $Z$.}

% ---- Question 37 (niveau 2, correct : C)
\qcmq[2]{Chap.~\ref{chap:dirac}\,$\cdot$\,\S\,\ref{sec:projecteurs}}
  {Un qubit préparé dans $\ket{0}$ est mesuré dans la base $X$ $\{\ket{+},\ket{-}\}$. Quelle est la probabilité d'obtenir $-$ ?}
  {$0$}
  {$1$}
  {$\frac12$}
  {$\frac{1}{\sqrt2}$}
  {C}
\qcmx
  {$\ket{0}$ est certain en base $Z$, mais pas en base $X$ : une autre base donne une autre loi.}
  {Même confusion dans l'autre sens : $\ket{0}$ n'est pas l'état $\ket{-}$.}
  {$P(-)=|\braket{-}{0}|^2=\left|\frac{1}{\sqrt2}\right|^2=\frac12$.}
  {C'est l'amplitude $\braket{-}{0}$ ; il faut élever son module au carré.}

% ---- Question 38 (niveau 2, correct : C)
\qcmq[2]{Chap.~\ref{chap:dirac}\,$\cdot$\,\S\,\ref{sec:qubit}}
  {Quel facteur $N>0$ normalise l'état $N\big(\ket{0}+(1-\ii)\ket{1}\big)$ ?}
  {$N=\frac{1}{\sqrt2}$}
  {$N=\frac13$}
  {$N=\frac{1}{\sqrt3}$}
  {$N=\frac{1}{\sqrt{1-2\ii}}$}
  {C}
\qcmx
  {On oublie la composante $\ket{0}$ : $|1-\ii|^2=2$, mais la norme au carré est $1+2=3$.}
  {On oublie la racine carrée : $3N^2=1$ donne $N=\frac{1}{\sqrt3}$, pas $\frac13$.}
  {$\braket{\Psi}{\Psi}=N^2\big(1+|1-\ii|^2\big)=N^2(1+2)=3N^2$.}
  {On élève $1-\ii$ au carré sans prendre le module : $(1-\ii)^2=-2\ii$ n'est pas réel ; il faut $|1-\ii|^2=2$.}

% ---- Question 39 (niveau 2, correct : D)
\qcmq[2]{Chap.~\ref{chap:dirac}\,$\cdot$\,\S\,\ref{sec:hilbert}}
  {Soient $\ket{\psi}=\frac{1}{\sqrt2}(\ket{0}+\ii\ket{1})$ et $\ket{\varphi}=\frac{1}{\sqrt2}(\ket{0}-\ket{1})$. Que vaut $\braket{\psi}{\varphi}$ ?}
  {$\frac{1-\ii}{2}$}
  {$\frac12$}
  {$0$}
  {$\frac{1+\ii}{2}$}
  {D}
\qcmx
  {On oublie de conjuguer l'amplitude $\ii$ du bra : on calcule $\frac12\big(1+\ii\cdot(-1)\big)$.}
  {Seul le terme en $\ket{0}$ est compté ; le terme en $\ket{1}$ vaut $\frac12(-\ii)(-1)=\frac{\ii}{2}$, non nul.}
  {Confusion avec $\braket{+}{-}=0$ ; ici $|\braket{\psi}{\varphi}|^2=\frac12$.}
  {Le bra conjugue les amplitudes : $\bra{\psi}=\frac{1}{\sqrt2}(\bra{0}-\ii\bra{1})$, d'où $\frac12\big(1+(-\ii)(-1)\big)=\frac{1+\ii}{2}$.}

% ---- Question 40 (niveau 2, correct : A)
\qcmq[2]{Chap.~\ref{chap:dirac}\,$\cdot$\,\S\,\ref{sec:qubit}}
  {Lequel de ces états est orthogonal à $\ket{{+}\ii}=\frac{1}{\sqrt2}(\ket{0}+\ii\ket{1})$ ?}
  {$\ket{{-}\ii}=\frac{1}{\sqrt2}(\ket{0}-\ii\ket{1})$}
  {$\ket{-}=\frac{1}{\sqrt2}(\ket{0}-\ket{1})$}
  {$\ket{+}=\frac{1}{\sqrt2}(\ket{0}+\ket{1})$}
  {$\ket{1}$}
  {A}
\qcmx
  {$\braket{{+}\ii}{{-}\ii}=\frac12\big(1+(-\ii)(-\ii)\big)=\frac12(1-1)=0$, en conjuguant le bra.}
  {$\braket{{+}\ii}{-}=\frac{1+\ii}{2}\neq0$ : $\ket{-}$ est orthogonal à $\ket{+}$, pas à $\ket{{+}\ii}$.}
  {$\braket{{+}\ii}{+}=\frac{1-\ii}{2}\neq0$.}
  {$\braket{{+}\ii}{1}=\frac{-\ii}{\sqrt2}\neq0$ : $\ket{1}$ est orthogonal à $\ket{0}$ seulement.}

% ---- Question 41 (niveau 2, correct : D)
\qcmq[1]{Chap.~\ref{chap:operateurs}\,$\cdot$\,\S\,\ref{sec:pur-mixte}}
  {Une matrice de densité de qubit a pour pureté $\operatorname{Tr}(\rho^2)=\frac58$. Que peut-on en conclure ?}
  {L'état est pur}
  {La matrice n'est pas une matrice de densité valide}
  {L'état est maximalement mixte, $\rho=\mathrm{Id}/2$}
  {L'état est mixte (mélange statistique)}
  {D}
\qcmx
  {Il faudrait $\operatorname{Tr}\rho^2=1$.}
  {$\frac58$ est compris entre $\frac12$ et $1$ : valeur admissible (exemple du cours : $\mathrm{diag}(\frac34,\frac14)$).}
  {Ce cas donnerait $\operatorname{Tr}\rho^2=\frac12$.}
  {$\operatorname{Tr}\rho^2<1$ ; pour un qubit, $\frac12\leq\operatorname{Tr}\rho^2\leq1$ et $\frac58$ est une valeur admissible.}

% ---- Question 42 (niveau 2, correct : B)
\qcmq[2]{Chap.~\ref{chap:operateurs}\,$\cdot$\,\S\,\ref{sec:rho}}
  {Que vaut $\langle Z\rangle=\operatorname{Tr}(\rho Z)$ pour un qubit dans l'état $\ket{+}$ ?}
  {$1$}
  {$0$}
  {$-1$}
  {$\frac{1}{\sqrt2}$}
  {B}
\qcmx
  {C'est la valeur pour $\ket{0}$.}
  {$\rho=\frac12\begin{pmatrix}1&1\\1&1\end{pmatrix}$ et $\operatorname{Tr}(\rho Z)=\frac12-\frac12=0$ : les résultats $+1$ et $-1$ sont équiprobables.}
  {C'est la valeur pour $\ket{1}$.}
  {C'est l'amplitude de $\ket{0}$ dans $\ket{+}$, pas une valeur moyenne de $Z$.}

% ---- Question 43 (niveau 2, correct : A)
\qcmq[1]{Chap.~\ref{chap:operateurs}\,$\cdot$\,\S\,\ref{sec:bloch}}
  {Sur la boule de Bloch, que représente un point situé à l'intérieur de la sphère ($|\vec r|<1$) ?}
  {Un état mixte, de pureté $\operatorname{Tr}\rho^2=\frac{1+|\vec r|^2}{2}<1$}
  {Un état pur dont la norme n'est pas $1$}
  {Un état pur en superposition}
  {Un point sans signification physique}
  {A}
\qcmx
  {Les états purs sont sur la sphère ; le centre $\vec r=\vec0$ est l'état maximalement mixte $\mathrm{Id}/2$.}
  {Tout état pur est normalisé et situé sur la sphère.}
  {Les états purs, superpositions comprises, sont tous \emph{sur} la sphère.}
  {Tout point de la boule est une matrice de densité valide.}

% ---- Question 44 (niveau 2, correct : B)
\qcmq[1]{Chap.~\ref{chap:operateurs}\,$\cdot$\,\S\,\ref{sec:pur-mixte}}
  {Par quoi la matrice de densité distingue-t-elle l'état pur $\ket{+}$ du mélange équiprobable de $\ket{0}$ et $\ket{1}$ ?}
  {Par leur trace, égale à $1$ pour l'état pur seulement}
  {Par les éléments hors diagonale (cohérences), non nuls pour $\ket{+}$ et nuls pour le mélange}
  {Par leurs éléments diagonaux, qui sont différents}
  {Ils ne se distinguent pas : c'est la même matrice}
  {B}
\qcmx
  {La trace vaut $1$ pour toute matrice de densité.}
  {$\ket{+}\bra{+}=\frac12\begin{pmatrix}1&1\\1&1\end{pmatrix}$ et le mélange vaut $\frac12\mathrm{Id}$ ; une mesure en base $X$ les distingue.}
  {Les deux ont la diagonale $(\frac12,\frac12)$ : $P(0)=P(1)=\frac12$ dans les deux cas.}
  {Les matrices sont différentes : en base $X$, $\ket{+}$ donne $+$ avec certitude, le mélange donne $\pm$ à pile ou face.}

% ---- Question 45 (niveau 2, correct : C)
\qcmq[2]{Chap.~\ref{chap:operateurs}\,$\cdot$\,\S\,\ref{sec:rho}}
  {Un qubit est préparé dans $\ket{0}$ avec la probabilité $\frac34$ et dans $\ket{1}$ avec la probabilité $\frac14$. Que vaut $\langle Z\rangle=\operatorname{Tr}(\rho Z)$ ?}
  {$\frac34$}
  {$1$}
  {$\frac12$}
  {$0$}
  {C}
\qcmx
  {C'est la probabilité $\rho_{00}$ de lire $0$ ; or $Z$ vaut $+1$ pour $0$ et $-1$ pour $1$.}
  {Valeur pour $\rho=\ket{0}\bra{0}$ (préparation certaine de $\ket{0}$).}
  {$\rho=\mathrm{diag}(\frac34,\frac14)$, donc $\operatorname{Tr}(\rho Z)=\frac34\cdot1+\frac14\cdot(-1)=\frac12$ : c'est la coordonnée $r_z$ de la sphère de Bloch.}
  {Valeur pour le mélange équiprobable $\mathrm{Id}/2$.}

% ---- Question 46 (niveau 2, correct : B)
\qcmq[2]{Chap.~\ref{chap:operateurs}\,$\cdot$\,\S\,\ref{sec:tenseur}}
  {Dans la base ordonnée $\{\ket{00},\ket{01},\ket{10},\ket{11}\}$, quel vecteur colonne représente $\ket{1}\otimes\ket{0}=\ket{10}$ ?}
  {$(0,1,0,0)^{\mathsf T}$}
  {$(0,0,1,0)^{\mathsf T}$}
  {$(1,0,0,0)^{\mathsf T}$}
  {$(0,0,0,1)^{\mathsf T}$}
  {B}
\qcmx
  {C'est $\ket{01}$ : l'ordre des facteurs est inversé.}
  {$\binom{0}{1}\otimes\binom{1}{0}=\big(0\binom{1}{0},\,1\binom{1}{0}\big)=(0,0,1,0)^{\mathsf T}$ : c'est le troisième vecteur de la base.}
  {C'est $\ket{00}$.}
  {C'est $\ket{11}$.}

% ---- Question 47 (niveau 2, correct : D)
\qcmq[2]{Chap.~\ref{chap:portes}\,$\cdot$\,\S\,\ref{sec:phase}}
  {Quelle relation lie la porte de phase $S=\mathrm{diag}(1,\ii)$ à la porte $Z$ ?}
  {$S=Z$}
  {$S^2=X$}
  {$S^4=Z$}
  {$S^2=Z$}
  {D}
\qcmx
  {$Z=\mathrm{diag}(1,-1)$ : $S$ est un quart de tour, $Z$ un demi-tour.}
  {$S^2$ est diagonale alors que $X$ ne l'est pas ; c'est $HSH$ qui est une racine carrée de $X$.}
  {$S^4=\mathrm{diag}(1,\ii^4)=\mathrm{Id}$.}
  {$S^2=\mathrm{diag}(1,\ii^2)=\mathrm{diag}(1,-1)=Z$ : deux quarts de tour autour de $z$ font un demi-tour.}

% ---- Question 48 (niveau 2, correct : A)
\qcmq[2]{Chap.~\ref{chap:portes}\,$\cdot$\,\S\,\ref{sec:controlees}}
  {Dans la porte de Toffoli (CCNOT) à trois qubits $x,y,z$, à quelle condition la cible $z$ est-elle inversée ?}
  {Seulement si $x=1$ et $y=1$}
  {Si $x=1$ ou $y=1$}
  {Si $x\neq y$}
  {Si $x=0$ et $y=0$}
  {A}
\qcmx
  {$\ket{x,y,z}\to\ket{x,y,z\oplus xy}$ : la cible est inversée si et seulement si le produit $xy$ vaut $1$.}
  {Un seul contrôle à $1$ ne suffit pas : pour $(x,y)=(1,0)$, $xy=0$ et la cible ne change pas.}
  {Ce serait deux CNOT ($z\to z\oplus x\oplus y$) ; pour $(1,0)$ le Toffoli n'inverse rien.}
  {Pour $(0,0)$ on a $xy=0$ : la cible ne change pas.}

% ---- Question 49 (niveau 2, correct : A)
\qcmq[1]{Chap.~\ref{chap:portes}\,$\cdot$\,\S\,\ref{sec:portes}}
  {Pour un hamiltonien $\Hop$ indépendant du temps, quel opérateur d'évolution $U(t,t_0)$ vérifie $\ket{\Psi(t)}=U(t,t_0)\ket{\Psi(t_0)}$ ?}
  {$U(t,t_0)=\ee^{-\ii\Hop(t-t_0)/\hbar}$}
  {$U(t,t_0)=\ee^{+\ii\Hop(t-t_0)/\hbar}$}
  {$U(t,t_0)=-\ii\hbar\,\partial_t$}
  {$U(t,t_0)=\Hop\,(t-t_0)/\hbar$}
  {A}
\qcmx
  {Solution de $\ii\hbar\,\partial_t\ket{\Psi}=\Hop\ket{\Psi}$ ; comme $\Hop$ est hermitien, $U^\dagger U=\mathrm{Id}$ : l'évolution est unitaire.}
  {Erreur de signe : c'est $U^\dagger$, qui remonte le temps.}
  {C'est un opérateur différentiel, pas un opérateur d'évolution ; l'équation de Schrödinger s'écrit $\ii\hbar\,\partial_t\ket{\Psi}=\Hop\ket{\Psi}$.}
  {On oublie l'exponentielle et le facteur $\ii$ : cet opérateur est hermitien, pas unitaire.}

% ---- Question 50 (niveau 2, correct : D)
\qcmq[2]{Chap.~\ref{chap:portes}\,$\cdot$\,\S\,\ref{sec:controlees}}
  {La porte de Fredkin (SWAP contrôlé) a pour contrôle le premier qubit. Que donne-t-elle sur $\ket{1,0,1}$ ?}
  {$\ket{1,0,1}$}
  {$\ket{0,1,1}$}
  {$\ket{1,1,1}$}
  {$\ket{1,1,0}$}
  {D}
\qcmx
  {C'est le résultat si le contrôle valait $0$ ; ici il vaut $1$ et les deux autres qubits sont différents.}
  {Le SWAP porte sur les qubits $2$ et $3$, pas sur le contrôle.}
  {Un SWAP ne change pas le nombre de $1$, il permute seulement ; $\ket{1,1,1}$ viendrait d'un NOT.}
  {Le contrôle vaut $1$ : les deux autres qubits sont échangés, $\ket{0,1}\to\ket{1,0}$.}

% ---- Question 51 (niveau 2, correct : A)
\qcmq[1]{Chap.~\ref{chap:portes}\,$\cdot$\,\S\,\ref{sec:op2}}
  {Que donne $(H\otimes I)\ket{00}$, où $H$ agit sur le premier facteur du produit tensoriel $\ket{a}\otimes\ket{b}$ ?}
  {$\ket{+}\otimes\ket{0}=\frac{1}{\sqrt2}(\ket{00}+\ket{10})$}
  {$\frac12(\ket{00}+\ket{01}+\ket{10}+\ket{11})$}
  {$\ket{0}\otimes\ket{+}=\frac{1}{\sqrt2}(\ket{00}+\ket{01})$}
  {$\frac{1}{\sqrt2}(\ket{00}+\ket{11})$}
  {A}
\qcmx
  {$(A\otimes B)(\ket{u}\otimes\ket{v})=A\ket{u}\otimes B\ket{v}$ : $H\ket{0}=\ket{+}$ sur le premier facteur, $I\ket{0}=\ket{0}$ sur le second.}
  {C'est $(H\otimes H)\ket{00}$ : ici le second qubit n'est pas touché.}
  {$H$ agit sur le premier qubit, pas sur le second : les facteurs sont inversés.}
  {C'est $\ket{\Phi^+}$ : il faut un CNOT pour intriquer ; des portes locales laissent l'état séparable.}

% ---- Question 52 (niveau 2, correct : C)
\qcmq[1]{Chap.~\ref{chap:bell}\,$\cdot$\,\S\,\ref{sec:separable}}
  {Pourquoi l'état de Bell $\ket{\Phi^+}=\frac{1}{\sqrt2}(\ket{00}+\ket{11})$ est-il dit intriqué ?}
  {Parce que $P(01)=0$}
  {Parce que chaque qubit, pris seul, est dans un état pur}
  {Il ne peut pas s'écrire comme un produit $\ket{\varphi}_A\otimes\ket{\chi}_B$ : $D=c_{00}c_{11}-c_{01}c_{10}=\frac12\neq0$}
  {Parce qu'il est orthogonal à $\ket{00}$}
  {C}
\qcmx
  {C'est vrai, mais $\ket{00}$ a aussi $P(01)=0$ et il est séparable : cette condition ne prouve rien.}
  {C'est l'inverse : pris seul, chaque qubit est complètement mixte ($\rho_A=\mathrm{Id}/2$).}
  {Un produit imposerait $D=0$ ; aucun des deux qubits n'a d'état propre.}
  {$\braket{00}{\Phi^+}=\frac{1}{\sqrt2}\neq0$ ; l'orthogonalité n'a pas de lien avec l'intrication.}

% ---- Question 53 (niveau 2, correct : B)
\qcmq[2]{Chap.~\ref{chap:bell}\,$\cdot$\,\S\,\ref{sec:teleportation}}
  {Dans la téléportation du cours, Alice obtient les résultats $a=0$ et $q=1$. Quelle correction Bob doit-il appliquer à son qubit ?}
  {La porte $X$}
  {La porte $Z$}
  {$X$ puis $Z$}
  {Aucune, l'état est déjà $\ket{\varphi}$}
  {B}
\qcmx
  {$X$ corrige le cas $a=1$ (inversion de bit) ; ici $a=0$.}
  {Bob détient $\alpha\ket{0}-\beta\ket{1}=Z\ket{\varphi}$ ; $Z$ rétablit le signe : $Z(\alpha\ket{0}-\beta\ket{1})=\alpha\ket{0}+\beta\ket{1}$.}
  {Les deux corrections ne sont nécessaires que pour $(a,q)=(1,1)$.}
  {Ce n'est vrai que pour $(a,q)=(0,0)$.}

% ---- Question 54 (niveau 2, correct : B)
\qcmq[2]{Chap.~\ref{chap:bell}\,$\cdot$\,\S\,\ref{sec:bell} ; ex.\,\ref{ex:superdense}}
  {Dans le codage superdense, Alice et Bob partagent $\ket{\Phi^+}$ et Alice agit sur son qubit. Quelle porte transforme $\ket{\Phi^+}$ en $\ket{\Psi^+}=\frac{1}{\sqrt2}(\ket{01}+\ket{10})$ ?}
  {La porte $Z$}
  {La porte $X$}
  {La porte identité}
  {La porte $H$}
  {B}
\qcmx
  {$Z$ donne $\ket{\Phi^-}=\frac{1}{\sqrt2}(\ket{00}-\ket{11})$ : elle change un signe, pas les chiffres.}
  {$X$ échange $\ket{0}$ et $\ket{1}$ sur un qubit : $\ket{00}+\ket{11}\to\ket{01}+\ket{10}$ ; c'est le message $10$.}
  {Elle laisse $\ket{\Phi^+}$ : c'est le message $00$.}
  {$(H\otimes I)\ket{\Phi^+}$ est une superposition de deux états de Bell, pas un état de Bell.}

% ---- Question 55 (niveau 2, correct : D)
\qcmq[2]{Chap.~\ref{chap:bell}\,$\cdot$\,\S\,\ref{sec:info}}
  {Pour deux qubits dans l'état de Bell $\ket{\Phi^+}$, que vaut l'entropie conditionnelle $H(A|B)=H(A,B)-H(B)$ ?}
  {$0$ bit}
  {$+1$ bit}
  {$+2$ bits}
  {$-1$ bit}
  {D}
\qcmx
  {C'est la valeur classique quand $A$ est entièrement déterminé par $B$ (deux bits identiques) ; ici $H(A,B)=0$, pas $H(B)$.}
  {C'est la valeur quand $A$ et $B$ sont indépendants : $H(A|B)=H(A)=1$.}
  {C'est l'information mutuelle $I(A;B)=1+1-0=2$ bits, pas l'entropie conditionnelle.}
  {$H(A,B)=S(\rho_{AB})=0$ (état pur) et $H(B)=1$ bit ($\rho_B=\mathrm{Id}/2$) : $0-1=-1$, valeur impossible en classique.}

% ---- Question 56 (niveau 2, correct : B)
\qcmq[1]{Chap.~\ref{chap:bell}\,$\cdot$\,\S\,\ref{sec:trace}}
  {Quelle est la matrice de densité réduite $\rho_A=\operatorname{Tr}_B\rho_{AB}$ d'un qubit d'une paire dans l'état $\ket{\Phi^+}$ ?}
  {$\rho_A=\ket{0}\bra{0}$}
  {$\rho_A=\frac12\mathrm{Id}$, de pureté $\operatorname{Tr}\rho_A^2=\frac12$}
  {$\rho_A=\ket{+}\bra{+}=\frac12\begin{pmatrix}1&1\\1&1\end{pmatrix}$}
  {$\rho_A=\ket{\Phi^+}\bra{\Phi^+}$}
  {B}
\qcmx
  {Le qubit aurait un état défini, ce qui est faux pour un état intriqué : $P(0)=P(1)=\frac12$.}
  {$\rho_A=\frac12\big(\ket{0}\bra{0}+\ket{1}\bra{1}\big)$ : les termes croisés $\ket{00}\bra{11}$ s'annulent ($\braket{1}{0}=0$). Le couple est pur, mais chaque qubit est complètement mixte.}
  {Les termes croisés sont multipliés par $\braket{1}{0}=0$ dans la trace partielle : ils disparaissent.}
  {C'est la matrice $4\times4$ du couple ; la trace partielle donne une matrice $2\times2$.}

% ---- Question 57 (niveau 2, correct : C)
\qcmq[1]{Chap.~\ref{chap:bell}\,$\cdot$\,\S\,\ref{sec:separable}}
  {Pour $\ket{\Psi}=c_{00}\ket{00}+c_{01}\ket{01}+c_{10}\ket{10}+c_{11}\ket{11}$, quel est le critère de séparabilité ?}
  {$c_{00}c_{11}+c_{01}c_{10}=0$}
  {$c_{00}c_{01}-c_{10}c_{11}=0$}
  {$D=c_{00}c_{11}-c_{01}c_{10}=0$}
  {$c_{00}=c_{11}$ et $c_{01}=c_{10}$}
  {C}
\qcmx
  {Mauvais signe : $\ket{+}\otimes\ket{-}$ est séparable mais donnerait $-\frac12\neq0$.}
  {Mauvais appariement : $\ket{\Psi^+}$ (intriqué) la vérifie, puisque $c_{00}=c_{11}=0$.}
  {Un produit $(a_0\ket{0}+a_1\ket{1})\otimes(b_0\ket{0}+b_1\ket{1})$ a $c_{ij}=a_ib_j$, donc $c_{00}c_{11}=a_0a_1b_0b_1=c_{01}c_{10}$ ; la réciproque est vraie aussi.}
  {Ni nécessaire ($\ket{00}$ est séparable avec $c_{00}\neq c_{11}$) ni suffisant ($\ket{\Phi^+}$ la vérifie et est intriqué).}

% ---- Question 58 (niveau 2, correct : B)
\qcmq[2]{Chap.~\ref{chap:bell}\,$\cdot$\,\S\,\ref{sec:bell}}
  {Quelle est la valeur propre de $X\otimes X$ pour l'état $\ket{\Phi^+}$ ?}
  {$-1$}
  {$+1$}
  {$0$}
  {$\ket{\Phi^+}$ n'est pas un état propre de $X\otimes X$}
  {B}
\qcmx
  {C'est la valeur pour $\ket{\Phi^-}$ et $\ket{\Psi^-}$ : le signe $-$ de l'état se retrouve dans la valeur propre.}
  {$(X\otimes X)\ket{00}=\ket{11}$ et $(X\otimes X)\ket{11}=\ket{00}$, donc $(X\otimes X)\ket{\Phi^+}=\ket{\Phi^+}$.}
  {$X\otimes X$ est unitaire : ses valeurs propres sont de module $1$, jamais nulles.}
  {Les quatre états de Bell sont des états propres communs de $Z\otimes Z$ et de $X\otimes X$.}

% ---- Question 59 (niveau 2, correct : A)
\qcmq[2]{Chap.~\ref{chap:bell}\,$\cdot$\,\S\,\ref{sec:info}}
  {Que vaut l'entropie de von Neumann $S(\rho)=-\sum_i\lambda_i\log_2\lambda_i$ d'un état pur ?}
  {$0$ bit}
  {$1$ bit}
  {$-1$ bit}
  {$+\infty$}
  {A}
\qcmx
  {Les valeurs propres sont $1$ et $0$ : $1\cdot\log_21=0$ et $0\cdot\log_20=0$ (par continuité, $x\log_2x\to0$).}
  {C'est la valeur pour l'état maximalement mixte $\mathrm{Id}/2$ d'un qubit.}
  {L'entropie de $\rho$ est positive ou nulle ; seule l'entropie conditionnelle quantique peut être négative.}
  {$\log_20$ diverge, mais le terme $0\cdot\log_20$ vaut $0$ par continuité.}

% ---- Question 60 (niveau 2, correct : A)
\qcmq[2]{Chap.~\ref{chap:bell}\,$\cdot$\,\S\,\ref{sec:bell}}
  {Lequel des états de Bell a pour valeur propre $-1$ à la fois pour $Z\otimes Z$ et pour $X\otimes X$ ?}
  {$\ket{\Psi^-}=\frac{1}{\sqrt2}(\ket{01}-\ket{10})$}
  {$\ket{\Phi^+}$}
  {$\ket{\Phi^-}$}
  {$\ket{\Psi^+}$}
  {A}
\qcmx
  {Résultats opposés dans les deux bases : $Z\otimes Z=-1$ (chiffres différents) et $X\otimes X=-1$ (signe $-$ dans l'état).}
  {Valeurs propres $(Z\otimes Z,\,X\otimes X)=(+1,+1)$.}
  {Valeurs propres $(Z\otimes Z,\,X\otimes X)=(+1,-1)$.}
  {Valeurs propres $(Z\otimes Z,\,X\otimes X)=(-1,+1)$.}

\clearpage
%% ======================================================================
\subsection{Niveau 3 : avancé (questions 61 à 90)}\label{sec:qcm-n3}
%% ======================================================================
\noindent{\color{insagris}\small Niveau de difficulté : Avancé\quad $\bullet\bullet\bullet\circ$}\par\smallskip
% ---- Question 61 (niveau 3, correct : A)
\qcmq[2]{Chap.~\ref{chap:histoire}\,$\cdot$\,\S\,\ref{sec:solutions}}
  {Un électron d'un puits infini passe du niveau $n=3$ au niveau $n=2$ en émettant un photon. Quelle est l'énergie du photon, en fonction de $E_1$ ?}
  {$5E_1$}
  {$E_1$}
  {$\frac94E_1$}
  {$13E_1$}
  {A}
\qcmx
  {$E_3-E_2=9E_1-4E_1=5E_1$ ; le photon a la fréquence $f=(E_3-E_2)/h$.}
  {On a soustrait les numéros de niveau ($3-2=1$) au lieu de leurs carrés.}
  {C'est le rapport $E_3/E_2$, pas la différence $E_3-E_2$.}
  {On a additionné $9E_1+4E_1$ : l'énergie du photon est une différence de niveaux.}

% ---- Question 62 (niveau 3, correct : A)
\qcmq[2]{Chap.~\ref{chap:histoire}\,$\cdot$\,\S\,\ref{sec:superposition}}
  {À $t=0$, une particule du puits infini est dans l'état $\Psi=\frac{\sqrt3}{2}\,\psi_1+\frac12\,\psi_3$ ($\psi_n$ : états stationnaires normalisés, d'énergies $E_n=n^2E_1$). Quelle est son énergie moyenne $\langle E\rangle$ ?}
  {$3E_1$}
  {$5E_1$}
  {$\left(\frac{\sqrt3}{2}+\frac92\right)E_1$}
  {$7E_1$}
  {A}
\qcmx
  {$P(E_1)=\frac34$ et $P(E_3)=\frac14$, donc $\langle E\rangle=\frac34E_1+\frac14\cdot9E_1=3E_1$.}
  {On a pondéré également les deux niveaux, $\frac{E_1+9E_1}{2}$ : les poids sont les $|c_i|^2$, pas $\frac12$.}
  {On a pondéré par les amplitudes $c_i$ au lieu de leurs carrés $|c_i|^2$.}
  {Poids permutés : $\frac14E_1+\frac34\cdot9E_1=7E_1$ ; la probabilité de $E_1$ est $\left|\frac{\sqrt3}{2}\right|^2=\frac34$.}

% ---- Question 63 (niveau 3, correct : C)
\qcmq[1]{Chap.~\ref{chap:histoire}\,$\cdot$\,\S\,\ref{sec:solutions}}
  {Pour l'état stationnaire $n=2$ du puits infini $]0,a[$, que valent la densité de probabilité de présence en $x=a/2$ et la position moyenne $\langle\hat x\rangle$ ?}
  {Densité maximale en $a/2$, et $\langle\hat x\rangle=a/2$}
  {Densité nulle en $a/2$, et $\langle\hat x\rangle=0$}
  {Densité nulle en $a/2$, et $\langle\hat x\rangle=a/2$}
  {Densité maximale en $a/2$, et $\langle\hat x\rangle=a/4$}
  {C}
\qcmx
  {C'est le cas de $n=1$ ; pour $n=2$ il y a un nœud au centre.}
  {Une densité nulle au centre ne donne pas une moyenne nulle : $\langle\hat x\rangle$ pondère tout le puits et vaut $\frac a2$ par symétrie.}
  {$|\Psi_2|^2=\frac2a\sin^2\frac{2\pi x}{a}$ s'annule en $x=\frac a2$ (nœud), alors que la symétrie impose $\langle\hat x\rangle=\frac a2$ : la moyenne n'est pas la position la plus probable.}
  {Ni l'un ni l'autre : maxima en $\frac a4$ et $\frac{3a}4$, nœud en $\frac a2$, moyenne $\frac a2$.}

% ---- Question 64 (niveau 3, correct : C)
\qcmq[2]{Chap.~\ref{chap:dirac}\,$\cdot$\,\S\,\ref{sec:projecteurs}}
  {Quelle est la matrice du projecteur $\hat P_+=\ket{+}\bra{+}$ dans la base $\{\ket{0},\ket{1}\}$ ?}
  {$\frac12\begin{pmatrix}1&-1\\-1&1\end{pmatrix}$}
  {$\begin{pmatrix}1&0\\0&0\end{pmatrix}$}
  {$\frac12\begin{pmatrix}1&1\\1&1\end{pmatrix}$}
  {$\frac{1}{\sqrt2}\begin{pmatrix}1&1\\1&-1\end{pmatrix}$}
  {C}
\qcmx
  {C'est $\ket{-}\bra{-}$, le projecteur sur $\ket{-}$.}
  {C'est $\ket{0}\bra{0}$ : projecteur de la base $Z$, pas de la base $X$.}
  {$\ket{+}\bra{+}=\frac12(\ket0+\ket1)(\bra0+\bra1)$ ; on vérifie $\hat P_+^2=\hat P_+$ et $\operatorname{Tr}\hat P_+=1$.}
  {C'est la matrice de Hadamard, unitaire ($H^2=\mathrm{Id}$), et non un projecteur (un projecteur vérifie $P^2=P$).}

% ---- Question 65 (niveau 3, correct : B)
\qcmq[2]{Chap.~\ref{chap:dirac}\,$\cdot$\,\S\,\ref{sec:projecteurs}}
  {Un qubit préparé dans $\ket{0}$ est mesuré en base $X$, puis aussitôt en base $Z$. Quelle est la probabilité d'obtenir $1$ à la seconde mesure ?}
  {$0$}
  {$\frac12$}
  {$1$}
  {$\frac14$}
  {B}
\qcmx
  {On garde $\ket{0}$ en mémoire : la mesure en base $X$ a effondré l'état sur $\ket{\pm}$ et effacé $\ket{0}$.}
  {La première mesure donne $\ket{+}$ ou $\ket{-}$ (probabilité $\frac12$ chacun) ; dans les deux cas $P(1)=\frac12$ en base $Z$ : $\frac12\cdot\frac12+\frac12\cdot\frac12=\frac12$.}
  {Sans rapport : $\ket{0}$ n'est plus l'état après la mesure en base $X$.}
  {On n'a compté qu'une branche de la première mesure ; il faut sommer les deux ($\frac14+\frac14$).}

% ---- Question 66 (niveau 3, correct : D)
\qcmq[2]{Chap.~\ref{chap:dirac}\,$\cdot$\,\S\,\ref{sec:projecteurs}}
  {Un qubit est dans l'état $\ket{\psi}=\frac{3}{5}\ket{0}+\frac{4}{5}\ket{1}$. Quelle est la probabilité d'obtenir $-$ en mesurant dans la base $X$ ?}
  {$\frac{16}{25}$}
  {$\frac{1}{25}$}
  {$\frac{49}{50}$}
  {$\frac{1}{50}$}
  {D}
\qcmx
  {C'est $P(1)$ en base $Z$ ($|\frac45|^2$) : on a mesuré dans la mauvaise base.}
  {On a oublié le facteur $\frac{1}{\sqrt2}$ de $\ket{-}$ : $\left|\frac35-\frac45\right|^2=\frac1{25}$, il manque le $\frac12$.}
  {C'est $P(+)$ ; la probabilité demandée est celle du résultat $-$.}
  {$\braket{-}{\psi}=\frac{1}{\sqrt2}\left(\frac35-\frac45\right)=-\frac{1}{5\sqrt2}$, donc $P(-)=\frac1{50}$ ; et $P(+)=\frac{49}{50}$.}

% ---- Question 67 (niveau 3, correct : B)
\qcmq[1]{Chap.~\ref{chap:dirac}\,$\cdot$\,\S\,\ref{sec:projecteurs}}
  {Lequel de ces kets décrit le même état physique que $\ket{+}=\frac{1}{\sqrt2}(\ket{0}+\ket{1})$ ?}
  {$\frac{1}{\sqrt2}\big(\ket{0}+\ii\ket{1}\big)$}
  {$\ee^{\ii\pi/3}\ket{+}$}
  {$\frac{1}{\sqrt2}\big(\ket{0}+\ee^{\ii\pi}\ket{1}\big)$}
  {$\frac{1}{\sqrt2}\big(\ket{0}+\ee^{\ii\pi/3}\ket{1}\big)$}
  {B}
\qcmx
  {La phase $\ii$ ne porte que sur $\ket{1}$ : c'est une phase \emph{relative}, et l'état est $\ket{{+}\ii}$ (base $Y$).}
  {Un facteur de phase \emph{globale} multiplie tout le ket : toutes les probabilités de mesure (et $\rho=\ket{+}\bra{+}$) sont inchangées.}
  {$\ee^{\ii\pi}=-1$ : c'est $\ket{-}$, orthogonal à $\ket{+}$.}
  {Phase relative $\frac\pi3$ : c'est un autre état, sur l'équateur de la sphère de Bloch, à l'azimut $\frac\pi3$.}

% ---- Question 68 (niveau 3, correct : A)
\qcmq[2]{Chap.~\ref{chap:dirac}\,$\cdot$\,\S\,\ref{sec:qubit}}
  {Laquelle de ces paires est une base orthonormée de $\mathbb{C}^2$ ?}
  {$\{\ket{{+}\ii},\ket{{-}\ii}\}$}
  {$\{\ket{0},\ket{+}\}$}
  {$\{\ket{+},\ket{{-}\ii}\}$}
  {$\{\ket{0},\ket{0}+\ket{1}\}$}
  {A}
\qcmx
  {Ces deux états sont normés et orthogonaux : $\braket{{+}\ii}{{-}\ii}=\frac12(1-1)=0$.}
  {Normés, mais $\braket{0}{+}=\frac{1}{\sqrt2}\neq0$ : pas orthogonaux.}
  {$\braket{+}{{-}\ii}=\frac12(1-\ii)\neq0$ : pas orthogonaux.}
  {Le second ket n'est pas normé (norme au carré $2$) ni orthogonal à $\ket{0}$.}

% ---- Question 69 (niveau 3, correct : C)
\qcmq[2]{Chap.~\ref{chap:operateurs}\,$\cdot$\,\S\,\ref{sec:rho}}
  {Laquelle de ces matrices peut être la matrice de densité d'un qubit ?}
  {$\begin{pmatrix}\frac12&\frac34\\\frac34&\frac12\end{pmatrix}$}
  {$\begin{pmatrix}\frac12&\frac{\ii}{2}\\\frac{\ii}{2}&\frac12\end{pmatrix}$}
  {$\begin{pmatrix}\frac23&\frac{\ii}{3}\\-\frac{\ii}{3}&\frac13\end{pmatrix}$}
  {$\begin{pmatrix}\frac34&0\\0&\frac12\end{pmatrix}$}
  {C}
\qcmx
  {Hermitienne et de trace $1$, mais pas positive : $\det=\frac14-\frac9{16}<0$, valeurs propres $\frac54$ et $-\frac14$ (une « probabilité » négative).}
  {De trace $1$ mais pas hermitienne : il faudrait $\rho_{10}=\overline{\rho_{01}}=-\frac{\ii}{2}$.}
  {Hermitienne ($\rho_{10}=\overline{\rho_{01}}$), de trace $1$ et positive : $\det\rho=\frac29-\frac19=\frac19>0$ avec $\operatorname{Tr}\rho>0$ donne deux valeurs propres positives, $\frac{3\pm\sqrt5}{6}$.}
  {Hermitienne et positive, mais de trace $\frac54\neq1$ : les probabilités diagonales ne somment pas à $1$.}

% ---- Question 70 (niveau 3, correct : B)
\qcmq[2]{Chap.~\ref{chap:operateurs}\,$\cdot$\,\S\,\ref{sec:op-dirac}}
  {Un qubit est dans l'état $\ket{\psi}=\frac{1}{\sqrt5}\big(\ket{0}+2\ii\ket{1}\big)$. Que vaut la valeur moyenne $\langle Y\rangle=\bra{\psi}Y\ket{\psi}$, avec $Y=\begin{pmatrix}0&-\ii\\\ii&0\end{pmatrix}$ ?}
  {$0$}
  {$\frac45$}
  {$-\frac45$}
  {$-\frac35$}
  {B}
\qcmx
  {On a oublié de conjuguer les amplitudes du bra : $\psi^{\mathsf T}Y\psi=\frac15(2+2\ii\cdot\ii)=0$ (c'est d'ailleurs la valeur de $\langle X\rangle$ pour cet état).}
  {$Y\ket\psi=\frac1{\sqrt5}(2\ket0+\ii\ket1)$, puis $\bra\psi Y\ket\psi=\frac15\big(1\cdot2+(-2\ii)(\ii)\big)=\frac45$ ; en général $\langle Y\rangle=2\,\mathrm{Im}(\bar\alpha\beta)$.}
  {Erreur de signe : on a utilisé la transposée $\begin{pmatrix}0&\ii\\-\ii&0\end{pmatrix}=-Y$.}
  {C'est $\langle Z\rangle=\frac15-\frac45$ : on a calculé la valeur moyenne d'une autre observable.}

% ---- Question 71 (niveau 3, correct : D)
\qcmq[1]{Chap.~\ref{chap:operateurs}\,$\cdot$\,\S\,\ref{sec:spectral}}
  {On mesure l'observable $A=\begin{pmatrix}2&1\\1&2\end{pmatrix}$ sur un qubit préparé dans $\ket{0}$. Quelle affirmation est correcte ?}
  {On obtient toujours $2$}
  {On obtient $+1$ ou $-1$, avec la probabilité $\frac12$ chacun}
  {On obtient $2$ ou $1$, avec les probabilités $\frac45$ et $\frac15$}
  {On obtient $3$ ou $1$, avec la probabilité $\frac12$ chacun (moyenne $2$)}
  {D}
\qcmx
  {On a lu le coefficient diagonal : $A\ket0=2\ket0+\ket1$ n'est pas proportionnel à $\ket0$, donc $\ket0$ n'est pas vecteur propre de $A$.}
  {Ce sont les valeurs propres de $X$, pas de $A$ : le décalage $2\,\mathrm{Id}$ translate le spectre ($2\pm1$).}
  {On a lu $A\ket0=2\ket0+\ket1$ comme un état : $A$ n'est pas unitaire, ce n'est pas une évolution ; les résultats possibles d'une mesure sont les valeurs propres de $A$.}
  {$A=2\,\mathrm{Id}+X$ a les vecteurs propres de $X$ : valeur propre $3$ pour $\ket+$, $1$ pour $\ket-$. Comme $\ket0=\frac{1}{\sqrt2}(\ket++\ket-)$, les deux résultats sont équiprobables, et $\langle A\rangle=\bra0A\ket0=2$.}

% ---- Question 72 (niveau 3, correct : D)
\qcmq[2]{Chap.~\ref{chap:operateurs}\,$\cdot$\,\S\,\ref{sec:pur-mixte}}
  {Un qubit est préparé dans $\ket{0}$ avec la probabilité $\frac12$ et dans $\ket{+}$ avec la probabilité $\frac12$. Quelle est la pureté $\operatorname{Tr}(\rho^2)$ de son état ?}
  {$1$}
  {$\frac12$}
  {$\frac58$}
  {$\frac34$}
  {D}
\qcmx
  {Un mélange de deux états purs n'est pur que s'ils sont identiques ; ici $\ket0\neq\ket+$, donc $\operatorname{Tr}\rho^2<1$.}
  {Valeur de l'état maximalement mixte $\frac{\mathrm{Id}}2$, atteinte pour deux états \emph{orthogonaux} ; ici ils ne le sont pas (on a aussi pu oublier les termes croisés).}
  {On n'a additionné que les carrés des éléments diagonaux ($\frac9{16}+\frac1{16}$) : il manque $2\rho_{01}\rho_{10}=\frac2{16}$.}
  {$\rho=\frac12\ket0\bra0+\frac12\ket+\bra+=\begin{pmatrix}\frac34&\frac14\\\frac14&\frac14\end{pmatrix}$, donc $\operatorname{Tr}\rho^2=\frac9{16}+\frac1{16}+2\cdot\frac1{16}=\frac34$ : état mixte, mais moins que $\frac{\mathrm{Id}}2$ car $\ket0$ et $\ket+$ ne sont pas orthogonaux.}

% ---- Question 73 (niveau 3, correct : B)
\qcmq[2]{Chap.~\ref{chap:operateurs}\,$\cdot$\,\S\,\ref{sec:bloch}}
  {Quel est le vecteur de Bloch $(r_x,r_y,r_z)$ de l'état $\ket{\psi}=\cos\frac\pi6\,\ket0+\sin\frac\pi6\,\ket1$ ?}
  {$\left(\frac12,\,0,\,\frac{\sqrt3}{2}\right)$}
  {$\left(\frac{\sqrt3}{2},\,0,\,\frac12\right)$}
  {$\left(\frac{\sqrt3}{4},\,0,\,\frac12\right)$}
  {$\left(\frac{\sqrt3}{2},\,0,\,\frac34\right)$}
  {B}
\qcmx
  {On a pris $\theta=\frac\pi6$ : l'angle de Bloch est le \emph{double} de l'angle qui figure dans les amplitudes.}
  {Ici $\frac\theta2=\frac\pi6$, donc $\theta=\frac\pi3$ et $\varphi=0$ : $\vec r=(\sin\theta\cos\varphi,\sin\theta\sin\varphi,\cos\theta)$. Contrôle : $r_z=|\alpha|^2-|\beta|^2=\frac34-\frac14$ et $r_x=2\alpha\beta=\frac{\sqrt3}2$.}
  {On a écrit $r_x=\alpha\beta$ en oubliant le facteur $2$ ; on le voit à $|\vec r|^2=\frac7{16}<1$, impossible pour un état pur.}
  {On a pris $r_z=|\alpha|^2$ au lieu de $|\alpha|^2-|\beta|^2$ ; or $|\vec r|^2=\frac{21}{16}>1$ est impossible.}

% ---- Question 74 (niveau 3, correct : A)
\qcmq[2]{Chap.~\ref{chap:operateurs}\,$\cdot$\,\S\,\ref{sec:bloch}}
  {Un qubit a pour matrice de densité $\rho=\frac12\begin{pmatrix}1&-\ii\\\ii&1\end{pmatrix}$. Quel est son vecteur de Bloch ?}
  {$(0,\,1,\,0)$ : c'est l'état pur $\ket{{+}\ii}$}
  {$(0,\,-1,\,0)$}
  {$\left(0,\,\frac12,\,0\right)$}
  {$(0,\,0,\,0)$}
  {A}
\qcmx
  {Avec $\rho=\frac12\begin{pmatrix}1+r_z&r_x-\ii r_y\\r_x+\ii r_y&1-r_z\end{pmatrix}$ : $r_z=0$ et $r_x-\ii r_y=-\ii$, donc $r_y=1$, $r_x=0$. Comme $|\vec r|=1$, l'état est pur : $\ket{{+}\ii}\bra{{+}\ii}$.}
  {Erreur de signe : $\rho_{01}=\frac{r_x-\ii r_y}{2}$, donc $\rho_{01}=-\frac\ii2$ correspond à $r_y=+1$ ($r_y=-1$ pour $\ket{{-}\ii}$).}
  {On a lu $\rho_{01}=-\frac\ii2$ comme $-\ii r_y$ : il faut multiplier par $2$ ($\rho=\frac12(\mathrm{Id}+\vec r\cdot\vec\sigma)$).}
  {Les diagonales $\frac12,\frac12$ donnent $r_z=0$, mais les éléments non diagonaux (cohérences) ne sont pas nuls : l'état est pur, et non $\frac{\mathrm{Id}}2$.}

% ---- Question 75 (niveau 3, correct : C)
\qcmq[2]{Chap.~\ref{chap:operateurs}\,$\cdot$\,\S\,\ref{sec:tenseur}}
  {Que donne $(X\otimes Z)\,\big(\ket{+}\otimes\ket{-}\big)$ ?}
  {$-\ket{-}\otimes\ket{-}$}
  {$\ket{-}\otimes\ket{+}$}
  {$\ket{+}\otimes\ket{+}$}
  {$\ket{+}\otimes\ket{-}$}
  {C}
\qcmx
  {On a échangé les facteurs ($Z\otimes X$) : $Z\ket+=\ket-$ et $X\ket-=-\ket-$. Le premier opérateur agit sur le premier qubit.}
  {On a cru que $X$ transforme $\ket+$ en $\ket-$ ; or $X$ échange $\ket0$ et $\ket1$ et laisse $\ket+$ inchangé.}
  {$(A\otimes B)(\ket u\otimes\ket v)=A\ket u\otimes B\ket v$ : $X\ket+=\ket+$ ($\ket+$ est propre pour $X$, valeur propre $+1$) et $Z\ket-=\ket+$ ($Z$ échange $\ket+$ et $\ket-$).}
  {On a cru que $Z$ laisse $\ket-$ inchangé ; $Z$ ne laisse invariants que $\ket0$ et $\ket1$, et échange $\ket+$ et $\ket-$.}

% ---- Question 76 (niveau 3, correct : C)
\qcmq[2]{Chap.~\ref{chap:portes}\,$\cdot$\,\S\,\ref{sec:phase}}
  {Le qubit $\ket{+}$ traverse la porte $T=\mathrm{diag}\big(1,\ee^{\ii\pi/4}\big)$, puis est mesuré dans la base $X$ $\{\ket+,\ket-\}$. Quelle est la probabilité d'obtenir $+$ ?}
  {$\frac12$}
  {$\frac{2-\sqrt2}{4}\approx0{,}15$}
  {$\frac{2+\sqrt2}{4}\approx0{,}85$}
  {$\frac{\sqrt2}{2}\approx0{,}71$}
  {C}
\qcmx
  {On a cru que la phase ne change pas les probabilités : c'est vrai en base $Z$, mais la base $X$ est sensible à la phase \emph{relative} (le point de Bloch est à l'azimut $\frac\pi4$ de l'axe $+x$).}
  {C'est $P(-)=\sin^2\frac\pi8$ : on a permuté les résultats $+$ et $-$.}
  {$T\ket+=\frac1{\sqrt2}(\ket0+\ee^{\ii\pi/4}\ket1)$, donc $\braket{+}{\psi}=\frac12\big(1+\ee^{\ii\pi/4}\big)$ et $P(+)=\frac14\big|1+\ee^{\ii\pi/4}\big|^2=\frac{2+2\cos\frac\pi4}{4}=\frac{2+\sqrt2}{4}$ (soit $\cos^2\frac\pi8$).}
  {On a pris $\cos\frac\pi4$ comme probabilité ; il faut le \emph{demi}-angle et le carré : $P=\cos^2\frac{\pi/4}{2}=\cos^2\frac\pi8$.}

% ---- Question 77 (niveau 3, correct : B)
\qcmq[2]{Chap.~\ref{chap:portes}\,$\cdot$\,\S\,\ref{sec:hadamard}}
  {Que vaut le produit de matrices $HYH$, où $Y=\begin{pmatrix}0&-\ii\\\ii&0\end{pmatrix}$ ?}
  {$Y$}
  {$-Y$}
  {$Z$}
  {$X$}
  {B}
\qcmx
  {Par analogie avec $HXH=Z$, on a cru que $H$ ne modifie pas $Y$ ; en fait l'axe $y$ est retourné (signe $-$).}
  {$H$ échange les axes $x$ et $z$ et retourne l'axe $y$ : $HXH=Z$, $HZH=X$, $HYH=-Y$. Direct : $HY=\frac{1}{\sqrt2}\begin{pmatrix}\ii&-\ii\\-\ii&-\ii\end{pmatrix}$, puis $HYH=\begin{pmatrix}0&\ii\\-\ii&0\end{pmatrix}=-Y$.}
  {C'est $HXH$ : $H$ n'échange que les axes $x$ et $z$, l'axe $y$ reste l'axe $y$.}
  {C'est $HZH$ : $Y$ ne peut pas devenir $X$, car $H$ laisse l'axe $y$ à sa place (au signe près).}

% ---- Question 78 (niveau 3, correct : D)
\qcmq[1]{Chap.~\ref{chap:portes}\,$\cdot$\,\S\,\ref{sec:op2}}
  {Que donne $(H\otimes H)\ket{01}$, où $\ket{01}=\ket0\otimes\ket1$ ?}
  {$\frac12\big(\ket{00}+\ket{01}-\ket{10}-\ket{11}\big)$}
  {$\frac12\big(\ket{00}-\ket{01}-\ket{10}+\ket{11}\big)$}
  {$\frac12\big(\ket{00}+\ket{01}+\ket{10}+\ket{11}\big)$}
  {$\frac12\big(\ket{00}-\ket{01}+\ket{10}-\ket{11}\big)$}
  {D}
\qcmx
  {C'est $\ket-\otimes\ket+$ : on a attribué le $\ket-$ au premier qubit, alors que c'est le second qui vaut $\ket1$.}
  {C'est $\ket-\otimes\ket-$ : $H\ket0=\ket+$, et non $\ket-$.}
  {C'est $\ket+\otimes\ket+$, le résultat pour $\ket{00}$ : on a oublié que $H\ket1=\ket-$ introduit des signes.}
  {$H\ket0\otimes H\ket1=\ket+\otimes\ket-=\frac12(\ket0+\ket1)\otimes(\ket0-\ket1)$. Le signe $-$ apparaît quand le \emph{second} chiffre vaut $1$ (formule générale $\frac12\sum_y(-1)^{x\cdot y}\ket y$ avec $x=01$).}

% ---- Question 79 (niveau 3, correct : D)
\qcmq[2]{Chap.~\ref{chap:portes}\,$\cdot$\,\S\,\ref{sec:cnot}}
  {Avec le CNOT du cours (contrôle : second qubit $B$ ; cible : premier qubit $A$ ; $\mathrm{CNOT}\ket{a,b}=\ket{a\oplus b,\,b}$), quel état obtient-on à partir de $\ket{1}_A\otimes\ket{+}_B$ ?}
  {$\ket{1}\otimes\ket{+}$ (inchangé)}
  {$\frac{1}{\sqrt2}\big(\ket{00}+\ket{11}\big)=\ket{\Phi^+}$}
  {$\frac{1}{\sqrt2}\big(\ket{01}-\ket{10}\big)=\ket{\Psi^-}$}
  {$\frac{1}{\sqrt2}\big(\ket{01}+\ket{10}\big)=\ket{\Psi^+}$}
  {D}
\qcmx
  {On a pris le premier qubit comme contrôle : le CNOT agirait alors sur $B$, et $X\ket+=\ket+$ laisserait l'état inchangé. Mais ici le contrôle est $B$.}
  {C'est le résultat pour $\ket{0}_A\ket{+}_B$ ; ici $A$ vaut $\ket1$, ce qui échange les deux termes.}
  {Aucun signe $-$ n'apparaît : le CNOT permute des états de base sans leur ajouter de phase.}
  {$\ket1\ket+=\frac1{\sqrt2}(\ket{10}+\ket{11})$ ; $\ket{10}\mapsto\ket{1\oplus0,0}=\ket{10}$ et $\ket{11}\mapsto\ket{1\oplus1,1}=\ket{01}$. C'est un état de Bell, donc intriqué.}

% ---- Question 80 (niveau 3, correct : C)
\qcmq[2]{Chap.~\ref{chap:portes}\,$\cdot$\,\S\,\ref{sec:swap}}
  {Que vaut $\mathrm{SWAP}\,(X\otimes\mathrm{Id})\,\mathrm{SWAP}$ ?}
  {$X\otimes\mathrm{Id}$}
  {$\mathrm{Id}\otimes\mathrm{Id}$}
  {$\mathrm{Id}\otimes X$}
  {$X\otimes X$}
  {C}
\qcmx
  {On a supposé que $\mathrm{SWAP}$ commute avec $X\otimes\mathrm{Id}$ ; ce n'est pas le cas, puisque cet opérateur n'agit pas de la même façon sur les deux qubits.}
  {On a simplifié $\mathrm{SWAP}^2=\mathrm{Id}$ en oubliant que $X\otimes\mathrm{Id}$ se trouve \emph{entre} les deux SWAP : on ne peut pas les rapprocher.}
  {$\mathrm{SWAP}\,(X\otimes\mathrm{Id})\,\mathrm{SWAP}\ket{a,b}=\mathrm{SWAP}\,(X\otimes\mathrm{Id})\ket{b,a}=\mathrm{SWAP}\ket{b\oplus1,a}=\ket{a,b\oplus1}$ : le SWAP transporte l'opérateur d'un qubit à l'autre (même mécanisme que $\mathrm{CNOT}_{21}=\mathrm{SWAP}\,\mathrm{CNOT}_{12}\,\mathrm{SWAP}$).}
  {Cet opérateur inverserait les deux bits ; ici un seul bit est inversé, celui du \emph{second} qubit.}

% ---- Question 81 (niveau 3, correct : C)
\qcmq[1]{Chap.~\ref{chap:portes}\,$\cdot$\,\S\,\ref{sec:controlees}}
  {La porte $Z$ contrôlée ($\mathrm{CZ}=\mathrm{diag}(1,1,1,-1)$, contrôle : premier qubit) est appliquée à $\ket{+}\otimes\ket{+}$. Quel état obtient-on ?}
  {$\ket{+}\otimes\ket{+}$}
  {$\ket{-}\otimes\ket{-}$}
  {$\frac12\big(\ket{00}+\ket{01}+\ket{10}-\ket{11}\big)$, un état intriqué}
  {$\ket{+}\otimes\ket{-}$}
  {C}
\qcmx
  {On a cru qu'un simple signe est une phase globale ; mais le signe $-$ ne touche que $\ket{11}$ : c'est une phase \emph{relative}.}
  {C'est $Z\otimes Z$ ($Z$ sur les deux qubits), et non une porte contrôlée.}
  {CZ ne change que le signe de $\ket{11}$. Critère de séparabilité : $c_{00}c_{11}-c_{01}c_{10}=-\frac14-\frac14=-\frac12\neq0$ ; on peut l'écrire $\frac{1}{\sqrt2}\big(\ket0\ket++\ket1\ket-\big)$.}
  {On a appliqué $Z$ au second qubit sans condition ; le contrôle exige que le premier qubit soit dans $\ket1$.}

% ---- Question 82 (niveau 3, correct : D)
\qcmq[2]{Chap.~\ref{chap:portes}\,$\cdot$\,\S\,\ref{sec:phase}, \ref{sec:portes-bloch}}
  {Un qubit préparé dans $\ket{0}$ traverse successivement $H$, puis $S$, puis $H$ (dans cet ordre). Quel est l'état final, à une phase globale près ?}
  {$\ket{{+}\ii}$}
  {$\ket{0}$}
  {$\ket{1}$}
  {$\ket{{-}\ii}=\frac{1}{\sqrt2}\big(\ket0-\ii\ket1\big)$}
  {D}
\qcmx
  {Mauvais sens de rotation : on a perdu un signe en calculant $H\ket{{+}\ii}$ (la phase relative vaut $-\ii$, pas $+\ii$).}
  {On a « simplifié » $HSH$ en $S$ (puis $S\ket0=\ket0$) : mais $H$ ne commute pas avec $S$, on ne peut pas supprimer les deux $H$.}
  {C'est le résultat pour $HZH=X$ ; or $S$ n'est qu'un \emph{quart} de tour autour de $z$ (et non un demi-tour comme $Z$), donc $HSH$ est un quart de tour autour de $x$.}
  {$HSH=\sqrt X$ est un quart de tour autour de $x$ (§\,\ref{sec:phase}). Directement : $H\ket0=\ket+$, $S\ket+=\ket{{+}\ii}$, $H\ket{{+}\ii}=\frac{1+\ii}{2}\ket0+\frac{1-\ii}{2}\ket1=\frac{1+\ii}{2}\big(\ket0-\ii\ket1\big)$.}

% ---- Question 83 (niveau 3, correct : A)
\qcmq[1]{Chap.~\ref{chap:bell}\,$\cdot$\,\S\,\ref{sec:mesure2}}
  {Deux qubits sont dans l'état $\ket{\Psi}=\frac{1}{\sqrt3}\big(\ket{00}+\ket{01}+\ket{11}\big)$. On mesure le \emph{second} qubit $B$ en base $Z$ et on obtient $1$. Que valait la probabilité de ce résultat, et dans quel état se trouve alors $A$ ?}
  {$P(B{=}1)=\frac23$ ; $A$ est dans l'état $\ket{+}$}
  {$P(B{=}1)=\frac13$ ; $A$ est dans l'état $\ket{0}$}
  {$P(B{=}1)=\frac23$ ; $A$ est dans l'état $\ket{1}$}
  {$P(B{=}1)=\frac13$ ; $A$ est dans l'état $\ket{+}$}
  {A}
\qcmx
  {$P(B{=}1)=|c_{01}|^2+|c_{11}|^2=\frac13+\frac13$. On garde les termes dont le second chiffre vaut $1$ : $\frac{1}{\sqrt3}\big(\ket{01}+\ket{11}\big)=\frac{1}{\sqrt3}(\ket0+\ket1)\otimes\ket1$, qui renormalisé donne $\ket{+}_A\otimes\ket{1}_B$.}
  {C'est le résultat $B=0$ : seul le terme $\ket{00}$ a un second chiffre nul, d'où la probabilité $\frac13$ et $A$ dans $\ket0$.}
  {On a projeté sur le seul terme $\ket{11}$ ; or $\ket{01}$ a aussi son second chiffre égal à $1$ : $A$ vaut $0$ ou $1$ (état $\ket+$).}
  {L'état de $A$ est juste, mais on n'a compté qu'un seul terme pour la probabilité ; il faut sommer $|c_{01}|^2+|c_{11}|^2$.}

% ---- Question 84 (niveau 3, correct : A)
\qcmq[2]{Chap.~\ref{chap:bell}\,$\cdot$\,\S\,\ref{sec:separable}}
  {Pour quelle valeur de $x$ l'état (non normalisé) $\ket{00}+2\ket{01}+3\ket{10}+x\ket{11}$ est-il séparable ?}
  {$x=6$}
  {$x=-6$}
  {$x=\frac16$}
  {$x=5$}
  {A}
\qcmx
  {Critère $D=c_{00}c_{11}-c_{01}c_{10}=x-2\cdot3=0$. On obtient bien un produit : $(\ket0+3\ket1)\otimes(\ket0+2\ket1)=\ket{00}+2\ket{01}+3\ket{10}+6\ket{11}$.}
  {Erreur de signe : $D=x-6$ s'annule pour $x=+6$ ; pour $x=-6$, $D=-12\neq0$ et l'état est intriqué.}
  {On a inversé le rapport : $c_{00}c_{11}=c_{01}c_{10}$ donne $x=\frac{c_{01}c_{10}}{c_{00}}=6$, et non $\frac{c_{00}}{c_{01}c_{10}}$.}
  {On a additionné ($2+3$) au lieu de multiplier : le critère de séparabilité porte sur des \emph{produits} de coefficients.}

% ---- Question 85 (niveau 3, correct : B)
\qcmq[2]{Chap.~\ref{chap:bell}\,$\cdot$\,\S\,\ref{sec:bell}}
  {Comment s'écrit $\ket{00}$ dans la base de Bell ?}
  {$\frac{1}{\sqrt2}\big(\ket{\Phi^+}-\ket{\Phi^-}\big)$}
  {$\frac{1}{\sqrt2}\big(\ket{\Phi^+}+\ket{\Phi^-}\big)$}
  {$\frac{1}{\sqrt2}\big(\ket{\Phi^+}+\ket{\Psi^+}\big)$}
  {$\ket{\Phi^+}$}
  {B}
\qcmx
  {Erreur de signe : cette combinaison vaut $\ket{11}$, puisque $\ket{\Phi^+}-\ket{\Phi^-}=\sqrt2\,\ket{11}$.}
  {$\ket{\Phi^+}+\ket{\Phi^-}=\frac{1}{\sqrt2}\cdot2\ket{00}=\sqrt2\,\ket{00}$. Une mesure de Bell sur $\ket{00}$ donne donc $\Phi^+$ ou $\Phi^-$ avec la probabilité $\frac12$ chacun (jamais $\Psi^\pm$).}
  {Cette combinaison vaut $\ket{+}\otimes\ket{+}=\frac12(\ket{00}+\ket{01}+\ket{10}+\ket{11})$ : les deux états de Bell « pairs » et « impairs » s'y ajoutent.}
  {On a confondu $\ket{00}$, état séparable, avec $\ket{\Phi^+}=\frac{1}{\sqrt2}(\ket{00}+\ket{11})$, qui contient aussi $\ket{11}$ et est intriqué.}

% ---- Question 86 (niveau 3, correct : A)
\qcmq[1]{Chap.~\ref{chap:bell}\,$\cdot$\,\S\,\ref{sec:circuit-bell}}
  {Le générateur de Bell du cours est $G=\mathrm{CNOT}\cdot(\mathrm{Id}\otimes H)$ : $H$ sur le second qubit $B$, puis CNOT de contrôle $B$ et de cible $A$. Quel circuit permet de \emph{mesurer} une paire dans la base de Bell (on le place avant une mesure en base $Z$) ?}
  {D'abord le CNOT (contrôle $B$, cible $A$), puis $H$ sur $B$}
  {D'abord $H$ sur $B$, puis le CNOT}
  {D'abord le CNOT, puis $H$ sur $A$ (la cible)}
  {Aucun circuit : on mesure directement les deux qubits en base $Z$}
  {A}
\qcmx
  {C'est l'inverse $G^\dagger=(\mathrm{Id}\otimes H)\,\mathrm{CNOT}$ : les portes sont leurs propres inverses et l'ordre s'inverse. Il envoie $\Phi^+,\Phi^-,\Psi^+,\Psi^-$ sur $\ket{00},\ket{01},\ket{10},\ket{11}$ (à une phase près).}
  {C'est $G$ lui-même, qui \emph{fabrique} des états de Bell : on garde le même ordre au lieu de l'inverser, et la mesure ne distingue plus les quatre états.}
  {Mauvais qubit : pour $\Phi^+$ on trouve $\ket{+}\otimes\ket{+}$, dont les deux bits sont aléatoires. Le $H$ doit agir sur $B$, qui a reçu le $H$ dans $G$.}
  {$\Phi^+$ et $\Phi^-$ (comme $\Psi^+$ et $\Psi^-$) donnent les mêmes résultats en base $Z$ : cette base ne voit pas le signe relatif, il faut d'abord défaire l'intrication.}

% ---- Question 87 (niveau 3, correct : B)
\qcmq[2]{Chap.~\ref{chap:bell}\,$\cdot$\,\S\,\ref{sec:trace}}
  {Quelle est la matrice de densité réduite $\rho_A=\operatorname{Tr}_B\ket\Psi\bra\Psi$ du premier qubit pour $\ket{\Psi}=\frac12\ket{00}+\frac{\sqrt3}{2}\ket{11}$ ?}
  {$\begin{pmatrix}\frac14&\frac{\sqrt3}{4}\\\frac{\sqrt3}{4}&\frac34\end{pmatrix}$}
  {$\begin{pmatrix}\frac14&0\\0&\frac34\end{pmatrix}$}
  {$\begin{pmatrix}\frac12&0\\0&\frac12\end{pmatrix}$}
  {$\begin{pmatrix}\frac12&0\\0&\frac{\sqrt3}{2}\end{pmatrix}$}
  {B}
\qcmx
  {C'est la matrice $\ket\chi\bra\chi$ de l'état pur $\ket\chi=\frac12\ket0+\frac{\sqrt3}2\ket1$ : on a gardé les cohérences, comme si $A$ était seul ; l'intrication avec $B$ les détruit.}
  {Avec $C=\begin{pmatrix}c_{00}&c_{01}\\c_{10}&c_{11}\end{pmatrix}=\begin{pmatrix}\frac12&0\\0&\frac{\sqrt3}2\end{pmatrix}$, $\rho_A=CC^\dagger=\mathrm{diag}\big(\frac14,\frac34\big)$ : état mixte, de pureté $\frac{10}{16}=1-2|D|^2$ avec $D=\frac{\sqrt3}{4}$.}
  {C'est la réponse pour un état de Bell (intrication maximale) ; ici les poids $\frac14$ et $\frac34$ diffèrent : l'intrication est partielle.}
  {On a mis les amplitudes au lieu de leurs carrés : la trace vaudrait $\frac{1+\sqrt3}{2}\neq1$.}

% ---- Question 88 (niveau 3, correct : A)
\qcmq[2]{Chap.~\ref{chap:bell}\,$\cdot$\,\S\,\ref{sec:info}}
  {Deux qubits sont dans le mélange statistique $\rho_{AB}=\frac12\ket{00}\bra{00}+\frac12\ket{11}\bra{11}$ (corrélations purement classiques). Que vaut $H(A|B)=S(\rho_{AB})-S(\rho_B)$, en bits ?}
  {$0$}
  {$-1$}
  {$1$}
  {$2$}
  {A}
\qcmx
  {$\rho_{AB}$ a pour valeurs propres $\frac12,\frac12,0,0$, donc $S(\rho_{AB})=1$ ; et $\rho_B=\frac12\mathrm{Id}$ donne $S(\rho_B)=1$. Ainsi $H(A|B)=0$ : connaître $B$ fixe $A$. Une valeur négative est propre à l'intrication.}
  {C'est la valeur pour l'état de Bell \emph{pur} $\ket{\Phi^+}$ : mêmes résultats en base $Z$, mais $S(\rho_{AB})=0$ car l'état global est pur. Ici l'état est mixte.}
  {On a pris $H(A|B)=H(A)$, comme si $B$ ne disait rien sur $A$ ; or les deux bits sont toujours égaux.}
  {On a additionné les entropies des deux qubits ($1+1$), valeur de deux bits indépendants : ce n'est pas une entropie conditionnelle.}

% ---- Question 89 (niveau 3, correct : B)
\qcmq[1]{Chap.~\ref{chap:bell}\,$\cdot$\,\S\,\ref{sec:teleportation}}
  {Dans la téléportation, avant de recevoir les deux bits $(a,q)$ d'Alice, dans quel état se trouve le qubit de Bob ?}
  {Dans l'état pur $\alpha\ket0+\beta\ket1$, déjà téléporté}
  {Dans l'état complètement mélangé $\rho_B=\frac12\mathrm{Id}$, indépendant de $\alpha$ et $\beta$}
  {Dans le mélange $|\alpha|^2\ket0\bra0+|\beta|^2\ket1\bra1$}
  {Dans l'état pur $\ket{0}$, son état initial}
  {B}
\qcmx
  {Ce serait une communication instantanée. L'état $\alpha\ket0+\beta\ket1$ n'apparaît qu'\emph{après} la correction, qui dépend de $(a,q)$.}
  {Les quatre kets $\ket{\phi_{aq}}$ sont équiprobables ($\frac14$) et $\frac14\sum_{a,q}\ket{\phi_{aq}}\bra{\phi_{aq}}=\frac12\mathrm{Id}$. Bob n'apprend rien avant le message classique : rien ne voyage plus vite que la lumière.}
  {Bob connaîtrait alors les statistiques de $\ket\varphi$ sans aucun message : or les probabilités d'Alice ($\frac14$ pour chaque résultat) ne dépendent pas de $\alpha,\beta$.}
  {Le qubit de Bob est la moitié de la paire $\ket{\Phi^+}$ : pris seul, il n'a pas d'état pur ($\rho_B=\frac12\mathrm{Id}$ avant comme après la mesure d'Alice, tant que Bob ignore le résultat).}

% ---- Question 90 (niveau 3, correct : D)
\qcmq[1]{Chap.~\ref{chap:bell}\,$\cdot$\,\S\,\ref{sec:chsh}}
  {Au jeu CHSH (gain si $a\oplus b=x\cdot y$), quelles sont les meilleures probabilités de gain avec une stratégie classique, même avec du hasard partagé, et avec une paire $\ket{\Phi^+}$ mesurée dans des bases bien choisies ?}
  {$\frac12$ en classique ; $1$ avec l'intrication}
  {$\frac34$ dans les deux cas}
  {$1$ en classique si Alice et Bob s'accordent à l'avance ; aucun avantage quantique}
  {$\frac34$ en classique ; $\cos^2\frac\pi8\approx0{,}85$ avec l'intrication}
  {D}
\qcmx
  {Classique : répondre toujours $0$ donne déjà $\frac34$. Quantique : l'intrication ne permet pas de gagner à tous les coups (borne de Tsirelson $\cos^2\frac\pi8$).}
  {Si l'intrication n'apportait rien, on n'aurait pas de violation de l'inégalité de Bell ; l'avantage quantique ($0{,}85>0{,}75$) est précisément ce que l'expérience de Bell met en évidence.}
  {S'accorder à l'avance (ou partager du hasard) ne suffit pas : aucune stratégie déterministe ne gagne les quatre parties, d'après l'argument de parité.}
  {Une stratégie déterministe perd au moins une des quatre parties (la somme modulo $2$ des conditions donne $0=1$) ; avec $\Phi^+$, chaque partie est gagnée avec la probabilité $\cos^2\frac\pi8=\frac{2+\sqrt2}{4}$.}

\clearpage
%% ======================================================================
\subsection{Niveau 4 : maîtrise (questions 91 à 100)}\label{sec:qcm-n4}
%% ======================================================================
\noindent{\color{insagris}\small Niveau de difficulté : Maîtrise\quad $\bullet\bullet\bullet\bullet$}\par\smallskip
% ---- Question 91 (niveau 4, correct : C)
\qcmq[2]{Chap.~\ref{chap:histoire}\,$\cdot$\,\S\,\ref{sec:superposition}}
  {Une particule du puits infini est préparée à $t=0$ dans l'état $\Psi=\frac{1}{\sqrt2}\big(\psi_1+\psi_2\big)$, avec $E_n=n^2E_1$. Au bout de quelle durée minimale la densité de probabilité $|\Psi(x,t)|^2$ retrouve-t-elle sa forme initiale ?}
  {$T=\dfrac{2\pi\hbar}{E_1}$}
  {$T=\dfrac{2\pi\hbar}{5E_1}$}
  {$T=\dfrac{2\pi\hbar}{3E_1}$}
  {$T=\dfrac{\pi\hbar}{3E_1}$}
  {C}
\qcmx
  {C'est la période de la phase $\ee^{-\ii E_1t/\hbar}$ de $\psi_1$ seul ; or les phases globales ne se voient pas dans $|\Psi|^2$ : seul compte l'écart $E_2-E_1=3E_1$.}
  {On a pris la somme $E_1+E_2=5E_1$ ; c'est la \emph{différence} des énergies (phase relative) qui intervient.}
  {$|\Psi|^2=\frac12(\psi_1^2+\psi_2^2)+\psi_1\psi_2\cos\frac{(E_2-E_1)t}{\hbar}$ : seul le terme croisé dépend du temps, à la pulsation $\omega=\frac{E_2-E_1}{\hbar}=\frac{3E_1}{\hbar}$, donc $T=\frac{2\pi}{\omega}$.}
  {C'est la demi-période : à $t=\frac T2$ le cosinus vaut $-1$ et la densité est le symétrique de l'état initial ($x\to a-x$), pas encore sa forme de départ.}

% ---- Question 92 (niveau 4, correct : D)
\qcmq[1]{Chap.~\ref{chap:dirac}\,$\cdot$\,\S\,\ref{sec:projecteurs}, \ref{sec:qubit}}
  {Un expérimentateur reçoit un qubit préparé soit dans $\ket{0}$, soit dans $\ket{+}$. Laquelle de ces affirmations est correcte ?}
  {Une mesure en base $Z$ les distingue : $\ket0$ donne $0$ et $\ket+$ donne $1$}
  {Une mesure en base $X$ les distingue : $\ket+$ donne $+$ et $\ket0$ donne $-$}
  {On les distingue en répétant la mesure sur le même qubit}
  {Aucune mesure ne permet de les distinguer à coup sûr, car $\braket{0}{+}=\frac{1}{\sqrt2}\neq0$}
  {D}
\qcmx
  {$\ket+$ donne $0$ ou $1$ avec la probabilité $\frac12$ : le résultat $0$ est possible pour les deux états.}
  {$\ket0=\frac{1}{\sqrt2}(\ket++\ket-)$ donne $+$ ou $-$ avec la probabilité $\frac12$ : le résultat $+$ est possible pour les deux états.}
  {Après une première mesure l'état est effondré : répéter ne donne aucune information nouvelle, et on ne peut pas copier le qubit pour recommencer (non-clonage).}
  {Seuls des états \emph{orthogonaux} sont parfaitement discernables. En base $Z$, $\ket0$ donne toujours $0$ et $\ket+$ donne $0$ une fois sur deux : obtenir $0$ ne tranche pas (obtenir $1$ prouve seulement qu'on avait $\ket+$).}

% ---- Question 93 (niveau 4, correct : C)
\qcmq[1]{Chap.~\ref{chap:operateurs}\,$\cdot$\,\S\,\ref{sec:pur-mixte}}
  {Une source envoie au hasard $\ket{0}$ ou $\ket{1}$ (probabilité $\frac12$ chacun) ; une autre envoie au hasard $\ket{+}$ ou $\ket{-}$ (probabilité $\frac12$ chacun). Peut-on distinguer ces deux sources en mesurant les qubits émis ?}
  {Oui : en base $Z$, la première source donne des résultats prévisibles}
  {Oui : en base $X$, la seconde source donne des résultats prévisibles}
  {Non : les deux ont la même matrice de densité $\frac{\mathrm{Id}}{2}$, donc les mêmes statistiques pour toute mesure}
  {Oui : les états purs envoyés ne sont pas les mêmes}
  {C}
\qcmx
  {Chaque qubit est bien $\ket0$ ou $\ket1$, mais on ignore lequel : le résultat vaut $0$ ou $1$ avec la probabilité $\frac12$, exactement comme avec la seconde source ($|\braket{0}{\pm}|^2=\frac12$).}
  {Même erreur : chaque qubit est $\ket+$ ou $\ket-$, mais on ignore lequel ; le résultat en base $X$ est aléatoire, et la première source donne aussi $+$ ou $-$ avec la probabilité $\frac12$.}
  {$\frac12\ket0\bra0+\frac12\ket1\bra1=\frac12\ket+\bra++\frac12\ket-\bra-=\frac{\mathrm{Id}}2$ (relation de fermeture dans chaque base). Les probabilités de toute mesure ne dépendent que de $\rho$.}
  {Ce qui compte n'est pas la liste des états purs envoyés mais la matrice de densité moyenne : des préparations différentes peuvent avoir la même $\rho$.}

% ---- Question 94 (niveau 4, correct : B)
\qcmq[1]{Chap.~\ref{chap:operateurs}\,$\cdot$\,\S\,\ref{sec:pur-mixte} ; Chap.~\ref{chap:portes}\,$\cdot$\,\S\,\ref{sec:portes-bloch}}
  {Un qubit est dans l'état complètement mélangé $\rho=\frac{\mathrm{Id}}{2}$. Existe-t-il une porte quantique $U$ qui le transforme en un état pur, par exemple $\ket{0}\bra{0}$ ?}
  {Oui : la porte de Hadamard, qui crée des superpositions}
  {Non : $\rho\mapsto U\rho U^\dagger$ conserve les valeurs propres, donc la pureté ; ici $U\frac{\mathrm{Id}}{2}U^\dagger=\frac{\mathrm{Id}}{2}$}
  {Oui : une rotation de Bloch qui amène le centre de la boule au pôle nord}
  {Oui : une mesure en base $Z$, qui donne $\ket0$}
  {B}
\qcmx
  {$H$ envoie un état pur sur un état pur ($\ket0\mapsto\ket+$) et laisse $\frac{\mathrm{Id}}{2}$ invariant : $H\frac{\mathrm{Id}}2H=\frac{\mathrm{Id}}2$ car $H^2=\mathrm{Id}$.}
  {$\operatorname{Tr}\big[(U\rho U^\dagger)^2\big]=\operatorname{Tr}\rho^2$ ($=\frac12$ ici). Sur la boule de Bloch, une porte est une rotation : elle laisse le centre fixe.}
  {Une rotation tourne la boule autour de son centre, qui reste fixe : elle ne peut pas l'amener sur la sphère.}
  {Une mesure n'est pas une porte (elle est non unitaire) ; de plus elle donne $\ket0$ ou $\ket1$ au hasard, pas $\ket0$ à coup sûr.}

% ---- Question 95 (niveau 4, correct : D)
\qcmq[1]{Chap.~\ref{chap:portes}\,$\cdot$\,\S\,\ref{sec:cnot}}
  {Que vaut $(H\otimes H)\,\mathrm{CNOT}_{12}\,(H\otimes H)$, où $\mathrm{CNOT}_{12}$ a pour contrôle le \emph{premier} qubit et pour cible le second ?}
  {$\mathrm{CNOT}_{12}$, inchangé}
  {$\mathrm{CZ}$ (porte $Z$ contrôlée)}
  {$\mathrm{SWAP}$}
  {$\mathrm{CNOT}_{21}$ (contrôle : second qubit ; cible : premier qubit)}
  {D}
\qcmx
  {On a cru que les deux $H$ s'annulent ($H^2=\mathrm{Id}$) ; mais ils encadrent le CNOT au lieu d'être côte à côte : on ne peut pas les simplifier.}
  {C'est le résultat quand $H$ n'encadre que la \emph{cible} : $(\mathrm{Id}\otimes H)\,\mathrm{CNOT}_{12}\,(\mathrm{Id}\otimes H)=\mathrm{CZ}$, car $HXH=Z$. Ici $H$ agit aussi sur le contrôle.}
  {Le SWAP s'obtient avec \emph{trois} CNOT alternés, pas avec un seul CNOT conjugué par $H\otimes H$ ; il n'y a pas de contrôle dans un SWAP.}
  {$H\ket0\bra0H=\ket+\bra+$, $H\ket1\bra1H=\ket-\bra-$ et $HXH=Z$ donnent $\ket+\bra+\otimes I+\ket-\bra-\otimes Z=I\otimes\ket0\bra0+X\otimes\ket1\bra1=\mathrm{CNOT}_{21}$ : dans la base $\{\ket+,\ket-\}$, contrôle et cible échangent leurs rôles (rebond de phase).}

% ---- Question 96 (niveau 4, correct : A)
\qcmq[1]{Chap.~\ref{chap:portes}\,$\cdot$\,\S\,\ref{sec:controlees}}
  {On applique la porte de Toffoli (contrôles : qubits 1 et 2 ; cible : qubit 3) à $\ket{+}\otimes\ket{+}\otimes\ket{0}$, puis on mesure le troisième qubit en base $Z$. Que peut-on dire si l'on obtient $1$ ?}
  {Ce résultat a la probabilité $\frac14$, et les deux premiers qubits sont alors dans l'état $\ket{11}$}
  {Ce résultat a la probabilité $\frac12$, et les deux premiers qubits restent dans $\ket{+}\ket{+}$}
  {Ce résultat a la probabilité $\frac14$, mais les deux premiers qubits restent dans $\ket{+}\ket{+}$}
  {Ce résultat a la probabilité $\frac34$, et les deux premiers qubits sont dans l'état $\frac{1}{\sqrt3}\big(\ket{00}+\ket{01}+\ket{10}\big)$}
  {A}
\qcmx
  {$\mathrm{CCNOT}\ket{x,y,0}=\ket{x,y,xy}$ donne $\frac12\big(\ket{000}+\ket{010}+\ket{100}+\ket{111}\big)$ : seul $\ket{111}$ a un troisième chiffre égal à $1$ (probabilité $\frac14$), et il impose $x=y=1$.}
  {La cible n'est inversée que si $x=y=1$ (probabilité $\frac14$) et non une fois sur deux ; de plus l'état est intriqué, donc mesurer le qubit 3 modifie les qubits 1 et 2.}
  {La probabilité est juste, mais la mesure de la cible n'est pas sans effet sur les contrôles : l'état après la Toffoli est intriqué, donc le résultat $1$ les projette sur $\ket{11}$.}
  {C'est le cas du résultat $0$ (probabilité $\frac34$ : $x,y\in\{00,01,10\}$). Le résultat $1$ est le cas rare, celui du ET vrai ($x=y=1$).}

% ---- Question 97 (niveau 4, correct : A)
\qcmq[1]{Chap.~\ref{chap:bell}\,$\cdot$\,\S\,\ref{sec:bell}}
  {Que donne $(H\otimes H)\ket{\Psi^-}$, avec $\ket{\Psi^-}=\frac{1}{\sqrt2}\big(\ket{01}-\ket{10}\big)$ ?}
  {$-\ket{\Psi^-}$, c'est-à-dire le même état physique}
  {$\ket{\Phi^-}$}
  {$\ket{\Psi^+}$}
  {$\ket{\Phi^+}$}
  {A}
\qcmx
  {Le singulet est invariant par $U\otimes U$ à un facteur près : $(U\otimes U)\ket{\Psi^-}=\det(U)\ket{\Psi^-}$ avec $\det H=-1$. Directement : $\frac{1}{\sqrt2}\big(\ket+\ket--\ket-\ket+\big)=-\ket{\Psi^-}$.}
  {C'est l'image de $\ket{\Psi^+}$, pas de $\ket{\Psi^-}$ : $(H\otimes H)\ket{\Psi^+}=\ket{\Phi^-}$. Le signe $-$ du singulet change le calcul (les termes $\ket{00}$ et $\ket{11}$ se compensent).}
  {$\ket{\Psi^+}$ est symétrique par échange des deux qubits, $\ket{\Psi^-}$ antisymétrique ; $H\otimes H$ commute avec le SWAP et conserve donc cette symétrie.}
  {C'est l'état que $H\otimes H$ laisse invariant ($\ket{\Phi^+}=\frac{1}{\sqrt2}(\ket{++}+\ket{--})$) ; le singulet n'a pas ce comportement.}

% ---- Question 98 (niveau 4, correct : B)
\qcmq[1]{Chap.~\ref{chap:bell}\,$\cdot$\,\S\,\ref{sec:chsh}}
  {Une expérience de type CHSH (avec $S=E_{00}+E_{01}+E_{10}-E_{11}$) mesure $S=2{,}6$. Que peut-on en conclure ?}
  {C'est impossible : on a toujours $|S|\leq2$}
  {Aucun modèle local à résultats prédéterminés ne l'explique, mais c'est compatible avec la mécanique quantique}
  {C'est impossible en mécanique quantique, car $2\sqrt2\approx2{,}4<2{,}6$}
  {Les particules s'échangent de l'information plus vite que la lumière}
  {B}
\qcmx
  {Cette borne ne vaut que pour les modèles classiques (locaux) ; l'expérience et la mécanique quantique la dépassent, jusqu'à $2\sqrt2$.}
  {Classique : $|S|\leq2$ ; quantique : $|S|\leq2\sqrt2\approx2{,}83$ (borne de Tsirelson). Comme $2<2{,}6<2\sqrt2$, l'inégalité de Bell est violée sans dépasser le maximum quantique ; le gain correspondant est $\frac12+\frac{2{,}6}{8}\approx0{,}825$.}
  {Erreur numérique : $2\sqrt2\approx2{,}83$, donc $2{,}6$ est permis.}
  {Les corrélations violent l'inégalité de Bell mais ne permettent aucun signal : la réponse de Bob est uniforme quelle que soit la base choisie par Alice.}

% ---- Question 99 (niveau 4, correct : A)
\qcmq[2]{Chap.~\ref{chap:bell}\,$\cdot$\,\S\,\ref{sec:info}}
  {Pour l'état $\ket{\Psi}=\frac{1}{\sqrt3}\big(\ket{00}+\sqrt2\,\ket{11}\big)$, quelle est l'entropie de von Neumann $S(\rho_A)$ du premier qubit, en bits ?}
  {$\log_2 3-\frac23\approx0{,}92$}
  {$1$}
  {$0$}
  {$0{,}64$}
  {A}
\qcmx
  {$\rho_A=\mathrm{diag}\big(\frac13,\frac23\big)$ (trace partielle), donc $S=-\frac13\log_2\frac13-\frac23\log_2\frac23=\log_23-\frac23\approx1{,}585-0{,}667$. C'est moins que le maximum d'un bit (état de Bell) : l'intrication est partielle.}
  {C'est la valeur pour un état de Bell ($\rho_A=\frac{\mathrm{Id}}{2}$) ; ici les poids $\frac13$ et $\frac23$ diffèrent, donc $S<1$.}
  {$S(\rho_A)=0$ vaut pour un état séparable ; ici $D=\frac{\sqrt2}{3}\neq0$ : l'état est intriqué, $\rho_A$ est mixte et $S>0$ (le couple $AB$ est pur, pas $A$ seul).}
  {On a utilisé le logarithme \emph{népérien} (résultat en nats) ; en bits il faut $\log_2$, ce qui donne $0{,}92$.}

% ---- Question 100 (niveau 4, correct : B)
\qcmq[1]{Chap.~\ref{chap:bell}\,$\cdot$\,\S\,\ref{sec:bell}}
  {Pourquoi peut-on mesurer simultanément les deux « parités » $Z\otimes Z$ et $X\otimes X$ de deux qubits, alors qu'on ne peut pas mesurer simultanément $Z$ et $X$ sur un seul qubit ?}
  {Parce que l'incertitude de Heisenberg ne s'applique pas à deux qubits}
  {Parce que $Z\otimes Z$ et $X\otimes X$ commutent : les deux signes $-1$ de $ZX=-XZ$ se compensent}
  {Parce que les deux qubits sont indépendants}
  {Parce que $Z\otimes Z$ et $X\otimes X$ agissent sur des qubits différents}
  {B}
\qcmx
  {Elle s'applique toujours : $Z\otimes\mathrm{Id}$ et $X\otimes\mathrm{Id}$ ne commutent pas. C'est la commutation de ces deux opérateurs précis qui permet la mesure simultanée.}
  {$(Z\otimes Z)(X\otimes X)=ZX\otimes ZX=(-XZ)\otimes(-XZ)=XZ\otimes XZ=(X\otimes X)(Z\otimes Z)$. Ils ont donc une base propre commune : la base de Bell.}
  {Dans un état de Bell les qubits sont intriqués, donc loin d'être indépendants ; la commutation est une propriété des opérateurs, pas de l'état.}
  {Chacun des deux opérateurs agit sur \emph{les deux} qubits à la fois. (Ce sont $Z\otimes\mathrm{Id}$ et $\mathrm{Id}\otimes X$, sur des qubits différents, qui commutent pour cette raison.)}

\label{qcm:fin}

\clearpage
\ifqcmcorrige

% ======================================================================
\subsection{Clé des réponses et fiche de score}\label{sec:qcm-cle}
% ======================================================================
\noindent\textbf{Cette partie donne les réponses : ne la lisez qu'après avoir répondu à toutes les
questions.}

\paragraph{Clé des réponses.} La lettre de la bonne réponse de la question $n$ se lit à la ligne
qui contient $n$, dans la colonne du rang de $n$ dans la ligne (question~24 : ligne 21--30,
colonne 4).
\begin{center}
\renewcommand{\arraystretch}{1.25}
\setlength{\tabcolsep}{5.5pt}
\begin{tabular}{@{}l*{10}{c}@{}}
\toprule
 & 1 & 2 & 3 & 4 & 5 & 6 & 7 & 8 & 9 & 10 \\
\midrule
1--10 & \qcmkey{1} & \qcmkey{2} & \qcmkey{3} & \qcmkey{4} & \qcmkey{5} & \qcmkey{6} & \qcmkey{7} & \qcmkey{8} & \qcmkey{9} & \qcmkey{10} \\
11--20 & \qcmkey{11} & \qcmkey{12} & \qcmkey{13} & \qcmkey{14} & \qcmkey{15} & \qcmkey{16} & \qcmkey{17} & \qcmkey{18} & \qcmkey{19} & \qcmkey{20} \\
21--30 & \qcmkey{21} & \qcmkey{22} & \qcmkey{23} & \qcmkey{24} & \qcmkey{25} & \qcmkey{26} & \qcmkey{27} & \qcmkey{28} & \qcmkey{29} & \qcmkey{30} \\
31--40 & \qcmkey{31} & \qcmkey{32} & \qcmkey{33} & \qcmkey{34} & \qcmkey{35} & \qcmkey{36} & \qcmkey{37} & \qcmkey{38} & \qcmkey{39} & \qcmkey{40} \\
41--50 & \qcmkey{41} & \qcmkey{42} & \qcmkey{43} & \qcmkey{44} & \qcmkey{45} & \qcmkey{46} & \qcmkey{47} & \qcmkey{48} & \qcmkey{49} & \qcmkey{50} \\
51--60 & \qcmkey{51} & \qcmkey{52} & \qcmkey{53} & \qcmkey{54} & \qcmkey{55} & \qcmkey{56} & \qcmkey{57} & \qcmkey{58} & \qcmkey{59} & \qcmkey{60} \\
61--70 & \qcmkey{61} & \qcmkey{62} & \qcmkey{63} & \qcmkey{64} & \qcmkey{65} & \qcmkey{66} & \qcmkey{67} & \qcmkey{68} & \qcmkey{69} & \qcmkey{70} \\
71--80 & \qcmkey{71} & \qcmkey{72} & \qcmkey{73} & \qcmkey{74} & \qcmkey{75} & \qcmkey{76} & \qcmkey{77} & \qcmkey{78} & \qcmkey{79} & \qcmkey{80} \\
81--90 & \qcmkey{81} & \qcmkey{82} & \qcmkey{83} & \qcmkey{84} & \qcmkey{85} & \qcmkey{86} & \qcmkey{87} & \qcmkey{88} & \qcmkey{89} & \qcmkey{90} \\
91--100 & \qcmkey{91} & \qcmkey{92} & \qcmkey{93} & \qcmkey{94} & \qcmkey{95} & \qcmkey{96} & \qcmkey{97} & \qcmkey{98} & \qcmkey{99} & \qcmkey{100} \\
\bottomrule
\end{tabular}
\end{center}

\paragraph{Fiche de score.} Reportez le nombre de bonnes réponses de chaque niveau.
\begin{center}
\renewcommand{\arraystretch}{1.6}
\begin{tabular}{@{}l l c c c@{}}
\toprule
\textbf{Niveau} & \textbf{Questions} & \textbf{Bonnes réponses} & \textbf{Sur} & \textbf{Pourcentage}\\
\midrule
1. Débutant & 1--30 & \makebox[2.6cm]{\dotfill} & 30 & \makebox[2cm]{\dotfill}\,\%\\
2. Intermédiaire & 31--60 & \makebox[2.6cm]{\dotfill} & 30 & \makebox[2cm]{\dotfill}\,\%\\
3. Avancé & 61--90 & \makebox[2.6cm]{\dotfill} & 30 & \makebox[2cm]{\dotfill}\,\%\\
4. Maîtrise & 91--100 & \makebox[2.6cm]{\dotfill} & 10 & \makebox[2cm]{\dotfill}\,\%\\
\midrule
\textbf{Total} & 1--100 & \makebox[2.6cm]{\dotfill} & 100 & \makebox[2cm]{\dotfill}\,\%\\
\bottomrule
\end{tabular}
\end{center}
Repères : au moins $80\,\%$ à un niveau, il est acquis et on peut passer au suivant ; entre $50$ et
$80\,\%$, revoyez les sections des questions manquées ; en dessous, relisez le chapitre concerné et
son récapitulatif avant de recommencer.

\paragraph{Où revoir ?} Questions classées par chapitre du cours.
\begin{center}
\renewcommand{\arraystretch}{1.3}
\begin{tabular}{@{}l >{\raggedright\arraybackslash}p{5.9cm} >{\raggedright\arraybackslash}p{6.9cm}@{}}
\toprule
\textbf{Chapitre} & \textbf{Titre} & \textbf{Questions}\\
\midrule
Chap.~\ref{chap:histoire} & Histoire et mécanique quantique & 1, 2, 3, 4, 5, 6, 31, 32, 33, 34, 35, 61, 62, 63, 91 \\
Chap.~\ref{chap:dirac} & Notation de Dirac et espace de Hilbert & 7, 8, 9, 10, 11, 12, 36, 37, 38, 39, 40, 64, 65, 66, 67, 68, 92 \\
Chap.~\ref{chap:operateurs} & Opérateurs, matrice de densité, Bloch, produit tensoriel & 13, 14, 15, 16, 17, 18, 41, 42, 43, 44, 45, 46, 69, 70, 71, 72, 73, 74, 75, 93, 94 \\
Chap.~\ref{chap:portes} & Portes quantiques et systèmes multi-qubits & 19, 20, 21, 22, 23, 24, 47, 48, 49, 50, 51, 76, 77, 78, 79, 80, 81, 82, 95, 96 \\
Chap.~\ref{chap:bell} & États quantiques, intrication, Bell & 25, 26, 27, 28, 29, 30, 52, 53, 54, 55, 56, 57, 58, 59, 60, 83, 84, 85, 86, 87, 88, 89, 90, 97, 98, 99, 100 \\
\bottomrule
\end{tabular}
\end{center}

\clearpage
\subsection{Corrigé du niveau 1 : débutant (questions 1 à 30)}\label{sec:qcm-c1}
\qcmcorr{1}
\qcmcorr{2}
\qcmcorr{3}
\qcmcorr{4}
\qcmcorr{5}
\qcmcorr{6}
\qcmcorr{7}
\qcmcorr{8}
\qcmcorr{9}
\qcmcorr{10}
\qcmcorr{11}
\qcmcorr{12}
\qcmcorr{13}
\qcmcorr{14}
\qcmcorr{15}
\qcmcorr{16}
\qcmcorr{17}
\qcmcorr{18}
\qcmcorr{19}
\qcmcorr{20}
\qcmcorr{21}
\qcmcorr{22}
\qcmcorr{23}
\qcmcorr{24}
\qcmcorr{25}
\qcmcorr{26}
\qcmcorr{27}
\qcmcorr{28}
\qcmcorr{29}
\qcmcorr{30}

\clearpage
\subsection{Corrigé du niveau 2 : intermédiaire (questions 31 à 60)}\label{sec:qcm-c2}
\qcmcorr{31}
\qcmcorr{32}
\qcmcorr{33}
\qcmcorr{34}
\qcmcorr{35}
\qcmcorr{36}
\qcmcorr{37}
\qcmcorr{38}
\qcmcorr{39}
\qcmcorr{40}
\qcmcorr{41}
\qcmcorr{42}
\qcmcorr{43}
\qcmcorr{44}
\qcmcorr{45}
\qcmcorr{46}
\qcmcorr{47}
\qcmcorr{48}
\qcmcorr{49}
\qcmcorr{50}
\qcmcorr{51}
\qcmcorr{52}
\qcmcorr{53}
\qcmcorr{54}
\qcmcorr{55}
\qcmcorr{56}
\qcmcorr{57}
\qcmcorr{58}
\qcmcorr{59}
\qcmcorr{60}

\clearpage
\subsection{Corrigé du niveau 3 : avancé (questions 61 à 90)}\label{sec:qcm-c3}
\qcmcorr{61}
\qcmcorr{62}
\qcmcorr{63}
\qcmcorr{64}
\qcmcorr{65}
\qcmcorr{66}
\qcmcorr{67}
\qcmcorr{68}
\qcmcorr{69}
\qcmcorr{70}
\qcmcorr{71}
\qcmcorr{72}
\qcmcorr{73}
\qcmcorr{74}
\qcmcorr{75}
\qcmcorr{76}
\qcmcorr{77}
\qcmcorr{78}
\qcmcorr{79}
\qcmcorr{80}
\qcmcorr{81}
\qcmcorr{82}
\qcmcorr{83}
\qcmcorr{84}
\qcmcorr{85}
\qcmcorr{86}
\qcmcorr{87}
\qcmcorr{88}
\qcmcorr{89}
\qcmcorr{90}

\clearpage
\subsection{Corrigé du niveau 4 : maîtrise (questions 91 à 100)}\label{sec:qcm-c4}
\qcmcorr{91}
\qcmcorr{92}
\qcmcorr{93}
\qcmcorr{94}
\qcmcorr{95}
\qcmcorr{96}
\qcmcorr{97}
\qcmcorr{98}
\qcmcorr{99}
\qcmcorr{100}
\fi
. Un chapitre = une \section.
%
%  Organisation :
%    - les énoncés (niveaux 1 à 4) occupent des pages à part ;
%    - le corrigé (clé des réponses, explications, pièges, fiche de score)
%      commence sur une nouvelle page : on peut imprimer les énoncés seuls ;
%    - chaque question est écrite avec \qcmq puis \qcmx (voir plus bas) : le même
%      code produit l'énoncé (cases à cocher) et, plus loin, le corrigé.
%
%  Bascules (définies dans main.tex ; à placer avant % =====================================================================
%  Chapitre 7 : QCM sur le cours (100 questions, 4 niveaux).
%  Fichier importé par main.tex avec \include{qcm}. Un chapitre = une \section.
%
%  Organisation :
%    - les énoncés (niveaux 1 à 4) occupent des pages à part ;
%    - le corrigé (clé des réponses, explications, pièges, fiche de score)
%      commence sur une nouvelle page : on peut imprimer les énoncés seuls ;
%    - chaque question est écrite avec \qcmq puis \qcmx (voir plus bas) : le même
%      code produit l'énoncé (cases à cocher) et, plus loin, le corrigé.
%
%  Bascules (définies dans main.tex ; à placer avant \include{qcm}) :
%    \qcmrefsfalse     masque la référence « Chap. · § » sous chaque énoncé ;
%    \qcmcorrigefalse  supprime le corrigé du document (QCM à distribuer seul).
%
%  Ce fichier est produit par outils_qcm/gen_qcm.py à partir de la banque de
%  questions (fichiers texte), mais il reste un fichier LaTeX ordinaire que l'on
%  peut modifier à la main. Pour ajouter une question à la main : insérer un
%  couple \qcmq ... \qcmx dans le bon niveau (la numérotation est automatique),
%  puis mettre à jour (1) les bornes « questions a à b » des titres de niveaux,
%  (2) les lignes \qcmcorr{n} du corrigé, (3) le tableau « Où revoir ? ».
% =====================================================================

% Macros de Dirac (déjà définies dans main.tex ; ce bloc évite l'erreur
% « Undefined control sequence \ket » si main.tex est une ancienne version).
\providecommand{\ket}[1]{|#1\rangle}
\providecommand{\bra}[1]{\langle #1|}
\providecommand{\braket}[2]{\langle #1|#2\rangle}

\makeatletter
% bascules : normalement déjà définies dans main.tex (valeurs par défaut : vrai)
\@ifundefined{ifqcmrefs}{\newif\ifqcmrefs\qcmrefstrue}{}
\@ifundefined{ifqcmcorrige}{\newif\ifqcmcorrige\qcmcorrigetrue}{}
\newcounter{qcmq}

% case à cocher (3 mm) pour le crayon
\newcommand{\qcmbox}{{\setlength{\fboxsep}{0pt}\setlength{\fboxrule}{0.5pt}%
  \raisebox{-0.4mm}{\fbox{\rule{0pt}{3mm}\rule{3mm}{0pt}}}}}

% une proposition : case, lettre, texte (retrait suspendu)
\newcommand{\qcm@opt}[2]{\noindent\hangindent=3.2em \hangafter=1
  \hspace*{0.3em}\qcmbox\hspace{0.45em}\textbf{#1)}\hspace{0.4em}#2\par\vspace{2.5pt}}

% \qcmq[colonnes]{ref}{énoncé}{A}{B}{C}{D}{bonne lettre}   (TeX limite les macros à 9 paramètres)
% \qcmx{com. A}{com. B}{com. C}{com. D}  : commentaires, placés juste après \qcmq
%   - com. de la bonne lettre = explication ; com. des autres = piège ;
%   - colonnes = 1 (une proposition par ligne) ou 2 (grille 2x2).
\newcommand{\qcmq}[8][1]{%
  \stepcounter{qcmq}%
  \expandafter\gdef\csname qcm@ref@\theqcmq\endcsname{#2}%
  \expandafter\gdef\csname qcm@rep@\theqcmq\endcsname{#8}%
  \expandafter\gdef\csname qcm@o@\theqcmq @A\endcsname{#4}%
  \expandafter\gdef\csname qcm@o@\theqcmq @B\endcsname{#5}%
  \expandafter\gdef\csname qcm@o@\theqcmq @C\endcsname{#6}%
  \expandafter\gdef\csname qcm@o@\theqcmq @D\endcsname{#7}%
  \par\vspace{7pt plus 3pt minus 2pt}\noindent
  \begin{minipage}{\linewidth}%
    \noindent{\color{insarouge}\bfseries Question~\theqcmq}%
    \ifqcmrefs\hfill{\footnotesize\color{insagris}#2}\fi\par\nopagebreak
    \vspace{2pt}\noindent #3\par\nopagebreak\vspace{3pt}%
    \ifnum#1=2
      \begin{minipage}[t]{0.495\linewidth}\qcm@opt{A}{#4}\end{minipage}\hfill
      \begin{minipage}[t]{0.495\linewidth}\qcm@opt{B}{#5}\end{minipage}\par\vspace{2pt}
      \begin{minipage}[t]{0.495\linewidth}\qcm@opt{C}{#6}\end{minipage}\hfill
      \begin{minipage}[t]{0.495\linewidth}\qcm@opt{D}{#7}\end{minipage}\par
    \else
      \qcm@opt{A}{#4}\qcm@opt{B}{#5}\qcm@opt{C}{#6}\qcm@opt{D}{#7}%
    \fi
  \end{minipage}\par}

\newcommand{\qcmx}[4]{%
  \expandafter\gdef\csname qcm@c@\theqcmq @A\endcsname{#1}%
  \expandafter\gdef\csname qcm@c@\theqcmq @B\endcsname{#2}%
  \expandafter\gdef\csname qcm@c@\theqcmq @C\endcsname{#3}%
  \expandafter\gdef\csname qcm@c@\theqcmq @D\endcsname{#4}}

% lettre de la bonne réponse de la question n (« ? » si la question n'existe pas)
\newcommand{\qcmkey}[1]{\@ifundefined{qcm@rep@#1}{?}{\csname qcm@rep@#1\endcsname}}

% une mauvaise réponse du corrigé (n° de question, lettre) : ignorée si c'est la bonne.
% (pas de \@for : « := » ne fonctionne pas avec le « : » rendu actif par babel-french)
\newcommand{\qcm@wrongone}[2]{%
  \edef\qcm@a{#2}\edef\qcm@b{\csname qcm@rep@#1\endcsname}%
  \ifx\qcm@a\qcm@b\else
    \noindent\hangindent=1.4em\hangafter=1 {\color{insagris}\textbf{#2)}}\enspace
    \csname qcm@o@#1@#2\endcsname\enspace--\enspace\emph{\csname qcm@c@#1@#2\endcsname}\par
  \fi}

% bloc de corrigé de la question n
\newcommand{\qcmcorr}[1]{%
  \par\vspace{6pt plus 2pt minus 2pt}\noindent
  \begin{minipage}{\linewidth}\small
    \noindent{\color{insarouge}\bfseries Question~#1}\enspace
    \textbf{Réponse : \csname qcm@rep@#1\endcsname}\hfill
    {\footnotesize\color{insagris}\csname qcm@ref@#1\endcsname}\par\nopagebreak
    \vspace{1pt}\noindent\hangindent=1.4em\hangafter=1
    \textbf{\csname qcm@rep@#1\endcsname)}\enspace
    \csname qcm@o@#1@\csname qcm@rep@#1\endcsname\endcsname\enspace---\enspace
    \csname qcm@c@#1@\csname qcm@rep@#1\endcsname\endcsname\par
    \qcm@wrongone{#1}{A}\qcm@wrongone{#1}{B}\qcm@wrongone{#1}{C}\qcm@wrongone{#1}{D}%
  \end{minipage}\par}
\makeatother

% =====================================================================
\section{QCM sur le cours}\label{chap:qcm}
% =====================================================================

\subsection{Mode d'emploi}\label{sec:qcm-mode}

Ce QCM compte \textbf{100 questions} sur les chapitres~\ref{chap:histoire} à~\ref{chap:bell},
classées en quatre niveaux. Chaque question propose quatre réponses dont \textbf{une seule} est
exacte : cochez au crayon la case de la proposition choisie.
\ifqcmrefs À droite du numéro de chaque question, une référence (chapitre, section) indique la
partie du cours concernée.\else\ifqcmcorrige{} La référence de chaque question (chapitre, section)
est donnée dans le corrigé.\fi\fi

\begin{center}
\renewcommand{\arraystretch}{1.3}
\begin{tabular}{@{}llc>{\raggedright\arraybackslash}p{8.2cm}@{}}
\toprule
\textbf{Niveau} & \textbf{Questions} & \textbf{Nombre} & \textbf{Ce qui est demandé}\\
\midrule
1. Débutant & 1 à 30 & 30 & définitions, notations, résultats du cours en une étape\\
2. Intermédiaire & 31 à 60 & 30 & petits calculs, propriétés, application directe d'un résultat\\
3. Avancé & 61 à 90 & 30 & calculs en plusieurs étapes, portes, circuits et protocoles\\
4. Maîtrise & 91 à 100 & 10 & raisonnements fins et pièges conceptuels, qui combinent plusieurs notions\\
\bottomrule
\end{tabular}
\end{center}

\begin{itemize}[nosep]
  \item \textbf{Sur papier} : imprimez seulement les pages~\pageref{chap:qcm} à~\pageref{qcm:fin}
        (mode d'emploi et énoncés, sections~\ref{sec:qcm-n1} à~\ref{sec:qcm-n4}).%
        \ifqcmcorrige{} Le corrigé commence sur une page distincte, à la page~\pageref{sec:qcm-cle}.\fi
  \item \textbf{Barème} : 1 point par bonne réponse, 0 sinon (pas de point négatif). Calculatrice
        autorisée, mais rarement utile (quelques logarithmes et racines).
  \item \textbf{Correction} :%
        \ifqcmcorrige{} comparez avec la clé des réponses, puis lisez l'explication de chaque
        question, même réussie. Les « pièges » décrivent l'erreur qui mène à chaque mauvaise réponse.%
        \else{} le corrigé n'est pas inclus dans cette version du document.\fi
\end{itemize}

\bigskip
\noindent Nom : \makebox[6cm]{\dotfill}\hfill Date : \makebox[3.5cm]{\dotfill}\hfill Score : \makebox[1.6cm]{\dotfill}\,/\,100
\bigskip

\begin{remarque}[Conventions des énoncés]
\begin{itemize}[nosep]
  \item $\ket{ab}=\ket{a}_A\otimes\ket{b}_B$ : le premier chiffre est le qubit $A$, le second le
        qubit $B$ ; dans un circuit, le fil du haut est le dernier chiffre.
  \item Le CNOT de la section~\ref{sec:cnot} a pour contrôle le \emph{second} chiffre ; les portes
        contrôlées de la section~\ref{sec:controlees} ont pour contrôle le \emph{premier}. Chaque
        énoncé précise le contrôle.
  \item Sans précision, une mesure est faite dans la base de calcul $Z$. Les logarithmes des
        entropies sont en base $2$ (résultats en bits).
  \item $\ket{\pm}=\frac{1}{\sqrt2}(\ket0\pm\ket1)$ et $\ket{{\pm}\ii}=\frac{1}{\sqrt2}(\ket0\pm\ii\ket1)$.
\end{itemize}
\end{remarque}

\clearpage
%% ======================================================================
\subsection{Niveau 1 : débutant (questions 1 à 30)}\label{sec:qcm-n1}
%% ======================================================================
\noindent{\color{insagris}\small Niveau de difficulté : Débutant\quad $\bullet\circ\circ\circ$}\par\smallskip
% ---- Question 1 (niveau 1, correct : D)
\qcmq[2]{Chap.~\ref{chap:histoire}\,$\cdot$\,\S\,\ref{sec:dualite}}
  {En 1905, qui explique l'effet photo-électrique en postulant que la lumière est faite de quanta d'énergie $E=hf$, les photons ?}
  {Louis de Broglie}
  {Erwin Schrödinger}
  {Max Planck}
  {Albert Einstein}
  {D}
\qcmx
  {Il propose en 1924 une onde de matière ($\lambda=h/p$) pour les particules comme l'électron ; l'idée du photon date de 1905.}
  {Il écrit en 1926 l'équation d'onde de la matière, plus de vingt ans après l'article sur les photons.}
  {Il introduit le quantum d'énergie en 1900 pour le rayonnement du corps noir ; c'est Einstein qui en fait une propriété de la lumière elle-même.}
  {Il postule que la lumière échange son énergie par paquets $E=hf=\hbar\omega$ : c'est la première face de la dualité onde--corpuscule.}

% ---- Question 2 (niveau 1, correct : A)
\qcmq[1]{Chap.~\ref{chap:histoire}\,$\cdot$\,\S\,\ref{sec:born}}
  {Selon l'interprétation de Max Born, que représente $|\Psi(x,t)|^2$ ?}
  {La densité de probabilité de présence de la particule en $x$ à l'instant $t$}
  {L'énergie cinétique de la particule}
  {L'amplitude de l'onde de matière}
  {La vitesse de la particule}
  {A}
\qcmx
  {$|\Psi(x,t)|^2\,\dd x$ est la probabilité de trouver la particule entre $x$ et $x+\dd x$ ; $\Psi$ lui-même n'est pas mesurable.}
  {L'énergie est associée à l'opérateur hamiltonien $\Hop$ ; $|\Psi|^2$ n'a pas la dimension d'une énergie.}
  {$\Psi$ est l'amplitude (complexe) ; $|\Psi|^2$ en est le carré du module.}
  {La vitesse s'obtient avec l'opérateur $\hat p=-\ii\hbar\,\partial_x$, pas avec $|\Psi|^2$.}

% ---- Question 3 (niveau 1, correct : D)
\qcmq[2]{Chap.~\ref{chap:histoire}\,$\cdot$\,\S\,\ref{sec:dualite}}
  {Une particule de quantité de mouvement $p$ est associée, d'après de Broglie, à une onde de longueur d'onde $\lambda$. Quelle relation est correcte ?}
  {$\lambda=hp$}
  {$\lambda=p/h$}
  {$\lambda=\hbar/p$}
  {$\lambda=h/p$}
  {D}
\qcmx
  {Produit au lieu du quotient : $\lambda$ grandirait avec $p$, le contraire de ce qu'on observe.}
  {Quotient inversé : $p/h=1/\lambda$ s'exprime en m$^{-1}$, pas en mètres.}
  {Confusion entre $h$ et $\hbar=h/2\pi$ : $\hbar/p=\lambda/2\pi$ est l'inverse du nombre d'onde $k$.}
  {Relation de de Broglie : plus la particule est rapide ou lourde, plus $\lambda$ est petite (balle de $0{,}1$~kg à $10$~m/s : $\lambda\approx10^{-34}$~m, inobservable).}

% ---- Question 4 (niveau 1, correct : C)
\qcmq[2]{Chap.~\ref{chap:histoire}\,$\cdot$\,\S\,\ref{sec:construction}}
  {Quelle est l'équation de Schrödinger dépendante du temps pour un état $\Psi$ et un hamiltonien $\Hop$ ?}
  {$-\ii\hbar\,\partial_t\Psi=\Hop\Psi$}
  {$\hbar\,\partial_t\Psi=\Hop\Psi$}
  {$\ii\hbar\,\partial_t\Psi=\Hop\Psi$}
  {$\ii\hbar\,\partial_x\Psi=\Hop\Psi$}
  {C}
\qcmx
  {Erreur de signe : la solution serait $\ee^{+\ii\Hop t/\hbar}\Psi(0)$, qui décrit l'évolution à rebours du temps.}
  {Sans le facteur $\ii$, les solutions sont des exponentielles réelles (croissantes ou décroissantes) et la norme n'est pas conservée.}
  {Elle relie l'évolution de l'état à l'énergie ; pour un $\Hop$ indépendant du temps, sa solution est $\Psi(t)=\ee^{-\ii\Hop t/\hbar}\Psi(0)$ (chapitre~\ref{chap:portes}).}
  {C'est la dérivée en $t$ qui décrit l'évolution ; la dérivée en $x$ sert à construire $\hat p=-\ii\hbar\,\partial_x$.}

% ---- Question 5 (niveau 1, correct : A)
\qcmq[1]{Chap.~\ref{chap:histoire}\,$\cdot$\,\S\,\ref{sec:construction}}
  {Quel opérateur représente l'énergie totale d'un système quantique ?}
  {L'hamiltonien $\Hop$}
  {L'opérateur quantité de mouvement $\hat p$}
  {L'opérateur position $\hat x$}
  {L'opérateur identité}
  {A}
\qcmx
  {Somme de l'énergie cinétique $-\frac{\hbar^2}{2m}\partial_x^2$ et de l'énergie potentielle $V$ ; ses valeurs propres sont les énergies mesurables.}
  {Il représente l'impulsion ; l'énergie cinétique est $\hat p^2/2m$, et il faut y ajouter $V$.}
  {C'est la multiplication par $x$ : il donne la position, pas l'énergie.}
  {Il vaut $1$ pour tout état : il ne contient aucune information sur l'énergie.}

% ---- Question 6 (niveau 1, correct : C)
\qcmq[1]{Chap.~\ref{chap:histoire}\,$\cdot$\,\S\,\ref{sec:histoire}}
  {Dans quel ordre chronologique ces avancées ont-elles eu lieu ?}
  {de Broglie $\to$ Planck $\to$ Schrödinger $\to$ Bell}
  {Planck $\to$ Schrödinger $\to$ de Broglie $\to$ Bell}
  {Planck (quantum d'énergie) $\to$ de Broglie (onde de matière) $\to$ Schrödinger (équation d'onde) $\to$ Bell (inégalités)}
  {Planck $\to$ de Broglie $\to$ Bell $\to$ Schrödinger}
  {C}
\qcmx
  {Le quantum d'énergie de Planck (1900) précède l'onde de matière de de Broglie (1924).}
  {L'équation de Schrödinger (1926) généralise l'idée d'onde de matière de de Broglie (1924) : elle vient après.}
  {Planck 1900, de Broglie 1924, Schrödinger 1926, Bell 1964.}
  {Les inégalités de Bell (1964) viennent près de quarante ans après l'équation de Schrödinger.}

% ---- Question 7 (niveau 1, correct : C)
\qcmq[1]{Chap.~\ref{chap:dirac}\,$\cdot$\,\S\,\ref{sec:dirac}}
  {En notation de Dirac, comment appelle-t-on $\ket{\Psi}$ et comment le représente-t-on pour un qubit ?}
  {Un bra, représenté par un vecteur ligne à deux composantes complexes}
  {Un ket, représenté par un nombre complexe}
  {Un ket, représenté par un vecteur colonne à deux composantes complexes}
  {Un bra, représenté par une matrice $2\times2$}
  {C}
\qcmx
  {Le bra $\bra{\Psi}$ est le conjugué transposé du ket : un vecteur ligne, mais ce n'est pas la notation $\ket{\Psi}$.}
  {Un scalaire ne peut pas décrire une superposition : il faut deux amplitudes $\alpha$ et $\beta$.}
  {Le ket est le vecteur colonne $\binom{\alpha}{\beta}$ ; son bra $\bra{\Psi}$ est le vecteur ligne conjugué.}
  {Une matrice est un opérateur ; le bra est une matrice ligne $1\times2$.}

% ---- Question 8 (niveau 1, correct : B)
\qcmq[2]{Chap.~\ref{chap:dirac}\,$\cdot$\,\S\,\ref{sec:qubit}}
  {Quelle condition les amplitudes $\alpha$ et $\beta$ de $\alpha\ket{0}+\beta\ket{1}$ doivent-elles vérifier pour que l'état soit normalisé ?}
  {$\alpha+\beta=1$}
  {$|\alpha|^2+|\beta|^2=1$}
  {$\alpha^2+\beta^2=1$}
  {$|\alpha|+|\beta|=1$}
  {B}
\qcmx
  {On somme les amplitudes (complexes, parfois négatives) ; ce sont leurs modules au carré qui sont des probabilités.}
  {Les probabilités de mesure $P(0)=|\alpha|^2$ et $P(1)=|\beta|^2$ doivent sommer à $1$.}
  {Oubli du module : pour $\alpha=\frac{1}{\sqrt2}$ et $\beta=\frac{\ii}{\sqrt2}$, $\alpha^2+\beta^2=0$ alors que l'état est normalisé.}
  {Somme des modules au lieu des carrés : $\alpha=\beta=\frac{1}{\sqrt2}$ est normalisé, mais $|\alpha|+|\beta|=\sqrt2$.}

% ---- Question 9 (niveau 1, correct : C)
\qcmq[1]{Chap.~\ref{chap:dirac}\,$\cdot$\,\S\,\ref{sec:projecteurs}}
  {On mesure dans la base $Z$ un qubit préparé dans $\ket{+}=\frac{1}{\sqrt2}(\ket{0}+\ket{1})$. Que se passe-t-il ?}
  {On lit $+$ et le qubit reste dans $\ket{+}$}
  {On lit toujours $0$}
  {On lit $0$ ou $1$, chacun avec la probabilité $\frac12$, et l'état devient $\ket{0}$ ou $\ket{1}$}
  {On lit $0$ et $1$ à la fois}
  {C}
\qcmx
  {Une mesure en base $Z$ ne donne que $0$ ou $1$ ; elle modifie l'état.}
  {Les deux amplitudes ont le même module : $P(0)=P(1)=\frac12$.}
  {Postulat de la mesure : la superposition s'effondre sur l'état propre correspondant au résultat.}
  {Chaque mesure donne un seul résultat ; la superposition ne s'observe qu'à travers les statistiques.}

% ---- Question 10 (niveau 1, correct : B)
\qcmq[2]{Chap.~\ref{chap:dirac}\,$\cdot$\,\S\,\ref{sec:qubit}}
  {Que vaut le produit scalaire $\braket{0}{1}$ ?}
  {$1$}
  {$0$}
  {$\frac{1}{\sqrt2}$}
  {$\ket{01}$}
  {B}
\qcmx
  {C'est la valeur de $\braket{0}{0}$ et de $\braket{1}{1}$ (vecteurs normés), pas de $\braket{0}{1}$.}
  {$\ket{0}$ et $\ket{1}$ sont orthogonaux : $\begin{pmatrix}1&0\end{pmatrix}\begin{pmatrix}0\\1\end{pmatrix}=0$.}
  {C'est la valeur de $\braket{0}{+}$, pas de $\braket{0}{1}$.}
  {Confusion avec le produit tensoriel : un produit scalaire est un nombre, pas un ket.}

% ---- Question 11 (niveau 1, correct : D)
\qcmq[2]{Chap.~\ref{chap:dirac}\,$\cdot$\,\S\,\ref{sec:qubit}}
  {Un qubit est dans l'état $\ket{-}=\frac{1}{\sqrt2}(\ket{0}-\ket{1})$. Quelle est la probabilité de mesurer $1$ en base $Z$ ?}
  {$0$}
  {$-\frac12$}
  {$\frac{1}{\sqrt2}$}
  {$\frac12$}
  {D}
\qcmx
  {Le signe $-$ ne « retire » pas $\ket{1}$ de la superposition.}
  {Une probabilité est toujours positive : on prend le module au carré de l'amplitude.}
  {C'est l'amplitude en module, pas la probabilité (il faut élever au carré).}
  {$P(1)=\left|-\frac{1}{\sqrt2}\right|^2=\frac12$ : le signe est une phase relative, qui disparaît dans le module.}

% ---- Question 12 (niveau 1, correct : B)
\qcmq[2]{Chap.~\ref{chap:dirac}\,$\cdot$\,\S\,\ref{sec:projecteurs}}
  {Quel est le projecteur sur l'état $\ket{0}$ ?}
  {$\ket{0}\bra{1}=\begin{pmatrix}0&1\\0&0\end{pmatrix}$}
  {$\hat P_0=\ket{0}\bra{0}=\begin{pmatrix}1&0\\0&0\end{pmatrix}$}
  {$\braket{0}{0}=1$}
  {$\begin{pmatrix}0&0\\0&1\end{pmatrix}$}
  {B}
\qcmx
  {Cet opérateur envoie $\ket{1}$ sur $\ket{0}$ ; son carré est nul : ce n'est pas un projecteur.}
  {$\hat P_0\ket{\Psi}=\alpha\ket{0}$ : il garde la composante sur $\ket{0}$ ; $\hat P_0^2=\hat P_0$.}
  {C'est un nombre, pas un opérateur.}
  {C'est $\ket{1}\bra{1}$, le projecteur sur $\ket{1}$.}

% ---- Question 13 (niveau 1, correct : A)
\qcmq[1]{Chap.~\ref{chap:operateurs}\,$\cdot$\,\S\,\ref{sec:bloch}}
  {Où se trouvent les états $\ket{0}$ et $\ket{1}$ sur la sphère de Bloch ?}
  {Aux deux pôles, Nord et Sud}
  {Aux deux extrémités de l'axe $x$}
  {Aux deux extrémités de l'axe $y$}
  {En deux points voisins de l'équateur}
  {A}
\qcmx
  {$\ket{0}$ au pôle Nord ($\theta=0$), $\ket{1}$ au pôle Sud ($\theta=\pi$) : deux états orthogonaux sont diamétralement opposés.}
  {Ce sont les places de $\ket{+}$ et $\ket{-}$.}
  {Ce sont les places de $\ket{{+}\ii}$ et $\ket{{-}\ii}$.}
  {Des états orthogonaux sont antipodaux, jamais voisins.}

% ---- Question 14 (niveau 1, correct : A)
\qcmq[2]{Chap.~\ref{chap:operateurs}\,$\cdot$\,\S\,\ref{sec:tenseur}}
  {Quelle est la dimension de l'espace de Hilbert de $n$ qubits ?}
  {$2^n$}
  {$2n$}
  {$n^2$}
  {$2^n-1$}
  {A}
\qcmx
  {Le produit tensoriel multiplie les dimensions : $2\times2\times\cdots\times2$ ; deux qubits : $4$, trois qubits : $8$.}
  {On additionne les dimensions au lieu de les multiplier ; pour $n=3$ on trouverait $6$ au lieu de $8$.}
  {Juste pour $n=2$ par coïncidence seulement : pour $n=3$ on trouverait $9$ au lieu de $8$.}
  {Il y a bien $2^n$ vecteurs dans la base de calcul, $\ket{0\cdots0}$ compris.}

% ---- Question 15 (niveau 1, correct : C)
\qcmq[1]{Chap.~\ref{chap:operateurs}\,$\cdot$\,\S\,\ref{sec:rho}}
  {Que vaut toujours la trace d'une matrice de densité $\rho$ ?}
  {$0$}
  {$1$ pour un état pur, et moins de $1$ pour un état mixte}
  {$1$}
  {$2$, la dimension de l'espace}
  {C}
\qcmx
  {C'est la trace des matrices de Pauli, pas d'une matrice de densité.}
  {Confusion avec $\operatorname{Tr}(\rho^2)$ : c'est la pureté qui distingue les deux cas, la trace vaut toujours $1$.}
  {C'est la somme des probabilités de mesure (les populations), que l'état soit pur ou mixte.}
  {C'est la trace de l'identité $\operatorname{Tr}\mathrm{Id}=2$, pas celle de $\rho$.}

% ---- Question 16 (niveau 1, correct : B)
\qcmq[1]{Chap.~\ref{chap:operateurs}\,$\cdot$\,\S\,\ref{sec:pur-mixte}}
  {Quelle est la valeur de $\operatorname{Tr}(\rho^2)$ pour un état pur ?}
  {$0$}
  {$1$}
  {$\frac12$}
  {Une valeur strictement inférieure à $1$}
  {B}
\qcmx
  {$\rho^2=0$ imposerait $\rho=0$, ce qui contredit $\operatorname{Tr}\rho=1$.}
  {Un état pur s'écrit $\rho=\ket{\Psi}\bra{\Psi}$, donc $\rho^2=\rho$ et $\operatorname{Tr}\rho^2=\operatorname{Tr}\rho=1$.}
  {C'est la valeur minimale pour un qubit, atteinte par l'état maximalement mixte $\mathrm{Id}/2$.}
  {C'est la caractérisation d'un état mixte.}

% ---- Question 17 (niveau 1, correct : B)
\qcmq[1]{Chap.~\ref{chap:operateurs}\,$\cdot$\,\S\,\ref{sec:rho}}
  {Que représentent les éléments diagonaux $\rho_{00}$ et $\rho_{11}$ de la matrice de densité d'un qubit ?}
  {Les cohérences entre $\ket{0}$ et $\ket{1}$}
  {Les probabilités de mesurer $\ket{0}$ et $\ket{1}$ (les populations)}
  {Les valeurs propres de $\rho$}
  {Les amplitudes $\alpha$ et $\beta$}
  {B}
\qcmx
  {Les cohérences sont les éléments hors diagonale $\rho_{01}$ et $\rho_{10}$.}
  {$\rho_{00}=\bra{0}\rho\ket{0}=P(0)$ et $\rho_{11}=P(1)$, d'où $\operatorname{Tr}\rho=1$.}
  {Seulement si $\rho$ est diagonale dans la base de calcul : $\ket{+}\bra{+}$ a pour diagonale $(\frac12,\frac12)$ mais pour valeurs propres $(1,0)$.}
  {Pour un état pur, $\rho_{00}=|\alpha|^2$ : le module au carré, pas l'amplitude.}

% ---- Question 18 (niveau 1, correct : D)
\qcmq[1]{Chap.~\ref{chap:operateurs}\,$\cdot$\,\S\,\ref{sec:op-dirac}}
  {Pourquoi associe-t-on aux grandeurs mesurables (observables) des opérateurs hermitiens ?}
  {Parce qu'ils conservent la norme des vecteurs}
  {Parce qu'ils ont tous une trace nulle}
  {Parce que leurs valeurs propres valent toutes $\pm1$}
  {Parce que leurs valeurs propres sont réelles, comme les résultats d'une mesure}
  {D}
\qcmx
  {C'est la propriété des opérateurs unitaires (les portes), pas des opérateurs hermitiens.}
  {Vrai pour les matrices de Pauli, faux en général (l'identité est hermitienne, de trace $2$), et sans rapport avec le caractère mesurable.}
  {Vrai pour $X$, $Y$, $Z$ seulement ; l'hamiltonien du puits a les valeurs propres $E_n=n^2E_1$.}
  {Un opérateur hermitien ($A^\dagger=A$) a des valeurs propres réelles et des vecteurs propres orthogonaux ; les résultats possibles sont ses valeurs propres.}

% ---- Question 19 (niveau 1, correct : C)
\qcmq[2]{Chap.~\ref{chap:portes}\,$\cdot$\,\S\,\ref{sec:pauli}}
  {Quelle porte de Pauli réalise l'inversion de bit $\ket{0}\leftrightarrow\ket{1}$ ?}
  {$Z$}
  {$Y$}
  {$X$}
  {$H$}
  {C}
\qcmx
  {$Z$ laisse $\ket{0}$ et change $\ket{1}$ en $-\ket{1}$ : c'est une inversion de phase.}
  {$Y\ket{0}=\ii\ket{1}$ : elle inverse le bit mais ajoute une phase ; la porte d'inversion « pure » est $X$.}
  {$X\ket{0}=\ket{1}$ et $X\ket{1}=\ket{0}$ : c'est le NOT quantique.}
  {$H$ crée des superpositions : $H\ket{0}=\ket{+}$.}

% ---- Question 20 (niveau 1, correct : B)
\qcmq[2]{Chap.~\ref{chap:portes}\,$\cdot$\,\S\,\ref{sec:hadamard}}
  {Que donne la porte de Hadamard appliquée à $\ket{0}$ ?}
  {$\ket{1}$}
  {$\ket{+}=\frac{1}{\sqrt2}(\ket{0}+\ket{1})$}
  {$\ket{-}=\frac{1}{\sqrt2}(\ket{0}-\ket{1})$}
  {$\ket{0}$}
  {B}
\qcmx
  {C'est l'action de $X$.}
  {Superposition à poids égaux : $P(0)=P(1)=\frac12$ ; c'est la première étape de la plupart des algorithmes.}
  {C'est $H\ket{1}$ ; le signe $-$ n'apparaît qu'à partir de $\ket{1}$.}
  {$H$ n'est pas l'identité ; elle ne laisse inchangés ni $\ket{0}$ ni $\ket{1}$ (elle envoie $\ket{+}$ sur $\ket{0}$).}

% ---- Question 21 (niveau 1, correct : D)
\qcmq[2]{Chap.~\ref{chap:portes}\,$\cdot$\,\S\,\ref{sec:portes}}
  {Quelle propriété doit vérifier la matrice $U$ d'une porte quantique ?}
  {Elle est hermitienne : $U^\dagger=U$}
  {Sa trace vaut $1$}
  {Ses coefficients sont réels}
  {Elle est unitaire : $U^\dagger U=\mathrm{Id}$}
  {D}
\qcmx
  {Vrai pour $X$, $Y$, $Z$, $H$ mais faux pour $S$ et $T$ : ce n'est pas la condition générale.}
  {C'est une propriété des matrices de densité ; $\operatorname{Tr}X=0$ et $\operatorname{Tr}\mathrm{Id}=2$.}
  {$Y$ et $S$ sont des portes à coefficients complexes.}
  {Elle conserve la norme ($\braket{\Psi}{\Psi}=1$ reste vrai) et elle est réversible.}

% ---- Question 22 (niveau 1, correct : A)
\qcmq[2]{Chap.~\ref{chap:portes}\,$\cdot$\,\S\,\ref{sec:portes}}
  {Quel est l'inverse d'une porte quantique $U$ ?}
  {$U^{-1}=U^\dagger$ (transposée conjuguée)}
  {$U^{-1}=U^{\mathsf T}$ (transposée seule)}
  {$U^{-1}=U^{*}$ (conjuguée seule)}
  {$U^{-1}=-U$}
  {A}
\qcmx
  {Conséquence de $U^\dagger U=\mathrm{Id}$ ; par exemple $S^{-1}=S^\dagger=\mathrm{diag}(1,-\ii)$.}
  {Oubli de la conjugaison : $S^{\mathsf T}=S=\mathrm{diag}(1,\ii)$ alors que $S^{-1}=\mathrm{diag}(1,-\ii)$.}
  {Faux en général : $Y^*=-Y$ alors que $Y^{-1}=Y$.}
  {Faux : $X^{-1}=X\neq-X$.}

% ---- Question 23 (niveau 1, correct : C)
\qcmq[1]{Chap.~\ref{chap:portes}\,$\cdot$\,\S\,\ref{sec:cnot}}
  {Dans une porte CNOT, que devient le qubit cible lorsque le qubit de contrôle est dans l'état $\ket{0}$ ?}
  {Il est inversé par la porte $X$}
  {Il est remis à $\ket{0}$}
  {Il reste inchangé}
  {Il subit un changement de signe, comme avec $Z$}
  {C}
\qcmx
  {C'est ce qui se passe quand le contrôle vaut $\ket{1}$.}
  {Le CNOT est réversible : remettre la cible à zéro détruirait de l'information.}
  {Avec le contrôle à $\ket{0}$, la porte agit comme l'identité sur la cible.}
  {Ce serait la porte $Z$ contrôlée (CZ), pas le CNOT.}

% ---- Question 24 (niveau 1, correct : A)
\qcmq[2]{Chap.~\ref{chap:portes}\,$\cdot$\,\S\,\ref{sec:cnot}}
  {Le CNOT du cours a pour contrôle le second chiffre du ket et pour cible le premier : $\mathrm{CNOT}\ket{a,b}=\ket{a\oplus b,\,b}$. Que donne-t-il sur $\ket{0,1}$ ?}
  {$\ket{1,1}$}
  {$\ket{0,1}$}
  {$\ket{0,0}$}
  {$\ket{1,0}$}
  {A}
\qcmx
  {Le contrôle $b=1$ inverse la cible : $0\oplus1=1$.}
  {Le contrôle vaut $1$ : la cible est bien inversée.}
  {C'est le contrôle qui serait inversé : le contrôle n'est jamais modifié.}
  {Les deux qubits seraient inversés : seule la cible change.}

% ---- Question 25 (niveau 1, correct : D)
\qcmq[1]{Chap.~\ref{chap:bell}\,$\cdot$\,\S\,\ref{sec:bell}}
  {Quelle propriété caractérise les quatre états de Bell ?}
  {Ils sont séparables : chaque qubit a son propre état}
  {Ils décrivent un seul qubit}
  {Ils ne diffèrent que par une phase globale}
  {Ils sont maximalement intriqués et forment une base orthonormée de l'espace de deux qubits}
  {D}
\qcmx
  {$D=c_{00}c_{11}-c_{01}c_{10}=\pm\frac12\neq0$ : ils ne se factorisent pas.}
  {Ce sont des états de deux qubits (quatre amplitudes sur $\ket{00},\ket{01},\ket{10},\ket{11}$).}
  {Ils sont orthogonaux deux à deux : ce sont quatre états distincts, alors qu'une phase globale ne change pas l'état.}
  {Aucun n'est un produit d'états à un qubit ($|D|=\frac12$, valeur maximale) et ils sont orthogonaux deux à deux.}

% ---- Question 26 (niveau 1, correct : D)
\qcmq[1]{Chap.~\ref{chap:bell}\,$\cdot$\,\S\,\ref{sec:teleportation}}
  {Quel est l'objectif du protocole de téléportation quantique ?}
  {Copier l'état d'un qubit inconnu sur le qubit de Bob}
  {Transmettre une information plus vite que la lumière}
  {Déplacer la particule physique d'Alice vers Bob}
  {Transférer l'état d'un qubit inconnu d'Alice à Bob, avec une paire intriquée et deux bits classiques}
  {D}
\qcmx
  {La copie est interdite (non-clonage) : après le protocole, l'état n'existe plus que chez Bob.}
  {Sans les deux bits classiques, le qubit de Bob est dans l'état $\mathrm{Id}/2$ : il n'a reçu aucune information.}
  {C'est l'état quantique qui est transféré, pas la particule qui le porte.}
  {Le qubit lui-même n'est pas envoyé : son état est reconstruit chez Bob, et détruit chez Alice par la mesure.}

% ---- Question 27 (niveau 1, correct : B)
\qcmq[1]{Chap.~\ref{chap:bell}\,$\cdot$\,\S\,\ref{sec:teleportation}}
  {Dans la téléportation, que doit envoyer Alice à Bob sur un canal classique pour qu'il retrouve l'état ?}
  {Un qubit intriqué}
  {Deux bits classiques, le résultat de sa mesure}
  {Les amplitudes exactes $\alpha$ et $\beta$ de l'état}
  {Un seul bit}
  {B}
\qcmx
  {La paire intriquée est partagée avant le protocole ; on ne l'envoie pas après.}
  {Les quatre résultats $(a,q)$ correspondent aux quatre corrections possibles ($I$, $Z$, $X$, $X$ puis $Z$) : il faut deux bits.}
  {Alice ne connaît pas l'état à téléporter, et ces nombres complexes n'ont pas de description finie en bits.}
  {Un bit ne distingue que deux cas, alors que Bob doit choisir parmi quatre corrections.}

% ---- Question 28 (niveau 1, correct : A)
\qcmq[1]{Chap.~\ref{chap:bell}\,$\cdot$\,\S\,\ref{sec:bell} ; ex.\,\ref{ex:superdense}}
  {Quel est le but du codage superdense ?}
  {Transmettre deux bits classiques en envoyant un seul qubit, grâce à une paire intriquée partagée}
  {Transmettre l'état d'un qubit avec deux bits classiques}
  {Compresser plusieurs qubits en un seul}
  {Chiffrer un message de façon inviolable}
  {A}
\qcmx
  {Alice applique l'une des quatre portes $I$, $Z$, $X$, $XZ$ à son qubit de la paire, l'envoie, et Bob identifie l'état de Bell obtenu.}
  {C'est la téléportation, le protocole inverse.}
  {Le protocole ne compresse pas de qubits : il fait passer deux bits classiques par un canal à un qubit.}
  {Ce n'est pas un protocole de cryptographie : la paire intriquée sert à doubler la capacité du canal.}

% ---- Question 29 (niveau 1, correct : A)
\qcmq[1]{Chap.~\ref{chap:bell}\,$\cdot$\,\S\,\ref{sec:teleportation}}
  {Que dit le théorème de non-clonage ?}
  {Aucune opération ne peut copier parfaitement un état quantique inconnu quelconque}
  {Il est impossible de mesurer un qubit sans le perturber}
  {Il est impossible de créer de l'intrication}
  {Il est impossible de copier un bit classique}
  {A}
\qcmx
  {Une porte $U$ qui copierait $\ket{0}$ et $\ket{1}$ enverrait $\ket{+}\ket{0}$ sur un état intriqué, par linéarité, et non sur $\ket{+}\ket{+}$.}
  {C'est l'effondrement de l'état, un autre principe.}
  {L'intrication se crée très bien (Hadamard puis CNOT).}
  {Le CNOT copie un bit classique : $\ket{0,b}\to\ket{b,b}$.}

% ---- Question 30 (niveau 1, correct : B)
\qcmq[1]{Chap.~\ref{chap:bell}\,$\cdot$\,\S\,\ref{sec:circuit-bell}}
  {Quel enchaînement de portes transforme $\ket{00}$ en l'état de Bell $\ket{\Phi^+}=\frac{1}{\sqrt2}(\ket{00}+\ket{11})$ ?}
  {CNOT, puis Hadamard sur le contrôle}
  {Hadamard sur l'un des qubits, puis CNOT dont le contrôle est ce même qubit}
  {Hadamard sur chacun des deux qubits}
  {SWAP, puis $X$ sur un qubit}
  {B}
\qcmx
  {Le CNOT ne fait rien sur $\ket{00}$ (contrôle à $0$) ; $H$ donne ensuite un état séparable.}
  {$H$ crée $\ket{0}\otimes\ket{+}$, puis le CNOT recopie la superposition : $\frac{1}{\sqrt2}(\ket{00}+\ket{11})$.}
  {On obtient $\ket{+}\otimes\ket{+}$, séparable : des portes locales ne créent pas d'intrication.}
  {$\ket{00}$ devient $\ket{01}$ ou $\ket{10}$, un état de base séparable.}

\clearpage
%% ======================================================================
\subsection{Niveau 2 : intermédiaire (questions 31 à 60)}\label{sec:qcm-n2}
%% ======================================================================
\noindent{\color{insagris}\small Niveau de difficulté : Intermédiaire\quad $\bullet\bullet\circ\circ$}\par\smallskip
% ---- Question 31 (niveau 2, correct : A)
\qcmq[1]{Chap.~\ref{chap:histoire}\,$\cdot$\,\S\,\ref{sec:born}}
  {Quelle condition la fonction d'onde d'une particule doit-elle vérifier sur tout l'espace ?}
  {$\displaystyle\int_{-\infty}^{+\infty}|\Psi(x,t)|^2\,\dd x=1$}
  {$\displaystyle\int_{-\infty}^{+\infty}\Psi(x,t)\,\dd x=1$}
  {$|\Psi(x,t)|^2=1$ pour tout $x$}
  {$\partial_x\Psi(x,t)=0$ en tout point}
  {A}
\qcmx
  {La particule se trouve forcément quelque part : la probabilité totale vaut $1$ (condition de normalisation).}
  {Sans le module au carré : $\Psi$ est complexe et son intégrale n'est pas une probabilité (elle peut être nulle ou complexe).}
  {$|\Psi|^2$ est une densité : si elle valait $1$ partout, son intégrale divergerait.}
  {Une fonction constante n'est pas normalisable ($\int|\Psi|^2\dd x$ diverge) ; la condition porte sur l'intégrale de $|\Psi|^2$.}

% ---- Question 32 (niveau 2, correct : D)
\qcmq[2]{Chap.~\ref{chap:histoire}\,$\cdot$\,\S\,\ref{sec:operateurs}}
  {Quelle est l'expression de l'opérateur quantité de mouvement $\hat p$ en une dimension ?}
  {$\hat p=m\,v$}
  {$\hat p=-\dfrac{\hbar^2}{2m}\dfrac{\partial^2}{\partial x^2}$}
  {$\hat p=+\ii\hbar\,\dfrac{\partial}{\partial x}$}
  {$\hat p=-\ii\hbar\,\dfrac{\partial}{\partial x}$}
  {D}
\qcmx
  {C'est la grandeur classique ; en mécanique quantique $p$ est un opérateur différentiel.}
  {C'est l'énergie cinétique $\hat p^2/2m$.}
  {Erreur de signe : on trouverait $p=-\hbar k$ pour l'onde $\ee^{\ii kx}$.}
  {Convention du cours : $\hat p\,\ee^{\ii kx}=\hbar k\,\ee^{\ii kx}$, donc une onde vers les $x$ croissants a $p=\hbar k>0$.}

% ---- Question 33 (niveau 2, correct : B)
\qcmq[2]{Chap.~\ref{chap:histoire}\,$\cdot$\,\S\,\ref{sec:solutions}}
  {Pour une particule dans un puits de potentiel infini de largeur $a$, comment varie l'énergie $E_n$ avec le numéro $n$ du niveau ?}
  {$E_n=E_1/n$}
  {$E_n=n^2E_1$, avec $E_1=\dfrac{\pi^2\hbar^2}{2ma^2}$}
  {$E_n=nE_1$}
  {$E_n=E_1$ pour tout $n$}
  {B}
\qcmx
  {Les niveaux se resserreraient quand $n$ augmente ; c'est l'inverse pour le puits infini.}
  {$k_n=\frac{n\pi}{a}$ et $E_n=\frac{\hbar^2k_n^2}{2m}$ : les niveaux s'écartent quand $n$ croît ($E_{n+1}-E_n=(2n+1)E_1$).}
  {Niveaux équidistants : ce n'est pas le puits infini, où $E_2=4E_1$ et $E_3=9E_1$.}
  {Il n'y aurait pas de quantification : tous les états auraient la même énergie.}

% ---- Question 34 (niveau 2, correct : C)
\qcmq[2]{Chap.~\ref{chap:histoire}\,$\cdot$\,\S\,\ref{sec:solutions}}
  {Une particule est dans un état stationnaire $\Psi_n$ de l'hamiltonien ($\Hop\Psi_n=E_n\Psi_n$). Que vaut l'écart-type $\sigma_E$ de son énergie ?}
  {$E_n$}
  {$E_n^2$}
  {$0$}
  {$\sqrt{E_n}$}
  {C}
\qcmx
  {C'est la valeur moyenne de l'énergie, pas son écart-type.}
  {C'est la valeur de $\langle\Hop^2\rangle$, avant soustraction de $\langle\Hop\rangle^2$.}
  {$\langle\Hop\rangle=E_n$ et $\langle\Hop^2\rangle=E_n^2$, donc $\sigma_E=\sqrt{E_n^2-E_n^2}=0$ : l'énergie est parfaitement déterminée.}
  {Mélange avec la formule $\sigma=\sqrt{\langle\Hop^2\rangle-\langle\Hop\rangle^2}$ : les deux termes sont égaux à $E_n^2$.}

% ---- Question 35 (niveau 2, correct : C)
\qcmq[2]{Chap.~\ref{chap:histoire}\,$\cdot$\,\S\,\ref{sec:solutions}}
  {Pour un état stationnaire $\Psi_n$ du puits infini $]0,a[$, que vaut la position moyenne $\langle\hat x\rangle$ ?}
  {$a/n$}
  {$a/4$}
  {$a/2$ pour tout $n$}
  {$0$}
  {C}
\qcmx
  {Elle dépendrait de $n$, ce que la symétrie de $|\Psi_n|^2$ interdit.}
  {C'est la position d'un maximum de $|\Psi_2|^2$ ; la moyenne des deux maxima $\frac a4$ et $\frac{3a}{4}$ donne $\frac a2$.}
  {$|\Psi_n|^2=\frac2a\sin^2\frac{n\pi x}{a}$ est symétrique par rapport au centre du puits.}
  {$\Psi_n$ s'annule sur la paroi, mais $\langle\hat x\rangle$ est une moyenne pondérée sur tout le puits.}

% ---- Question 36 (niveau 2, correct : D)
\qcmq[2]{Chap.~\ref{chap:dirac}\,$\cdot$\,\S\,\ref{sec:qubit}}
  {Un qubit est dans l'état $\ket{\Psi}=\frac{1}{\sqrt5}\ket{0}+\frac{2\ii}{\sqrt5}\ket{1}$. Quelle est la probabilité de mesurer $1$ en base $Z$ ?}
  {$\frac25$}
  {$-\frac45$}
  {$\frac15$}
  {$\frac45$}
  {D}
\qcmx
  {On divise le coefficient $2$ par $5$ : on oublie à la fois la racine et le carré.}
  {On élève $\frac{2\ii}{\sqrt5}$ au carré sans conjuguer : $\left(\frac{2\ii}{\sqrt5}\right)^2=-\frac45$ ; une probabilité est $|\cdot|^2$.}
  {C'est $P(0)=\left|\frac{1}{\sqrt5}\right|^2$.}
  {$P(1)=\left|\frac{2\ii}{\sqrt5}\right|^2=\frac45$ : le facteur $\ii$ est une phase relative, invisible en base $Z$.}

% ---- Question 37 (niveau 2, correct : C)
\qcmq[2]{Chap.~\ref{chap:dirac}\,$\cdot$\,\S\,\ref{sec:projecteurs}}
  {Un qubit préparé dans $\ket{0}$ est mesuré dans la base $X$ $\{\ket{+},\ket{-}\}$. Quelle est la probabilité d'obtenir $-$ ?}
  {$0$}
  {$1$}
  {$\frac12$}
  {$\frac{1}{\sqrt2}$}
  {C}
\qcmx
  {$\ket{0}$ est certain en base $Z$, mais pas en base $X$ : une autre base donne une autre loi.}
  {Même confusion dans l'autre sens : $\ket{0}$ n'est pas l'état $\ket{-}$.}
  {$P(-)=|\braket{-}{0}|^2=\left|\frac{1}{\sqrt2}\right|^2=\frac12$.}
  {C'est l'amplitude $\braket{-}{0}$ ; il faut élever son module au carré.}

% ---- Question 38 (niveau 2, correct : C)
\qcmq[2]{Chap.~\ref{chap:dirac}\,$\cdot$\,\S\,\ref{sec:qubit}}
  {Quel facteur $N>0$ normalise l'état $N\big(\ket{0}+(1-\ii)\ket{1}\big)$ ?}
  {$N=\frac{1}{\sqrt2}$}
  {$N=\frac13$}
  {$N=\frac{1}{\sqrt3}$}
  {$N=\frac{1}{\sqrt{1-2\ii}}$}
  {C}
\qcmx
  {On oublie la composante $\ket{0}$ : $|1-\ii|^2=2$, mais la norme au carré est $1+2=3$.}
  {On oublie la racine carrée : $3N^2=1$ donne $N=\frac{1}{\sqrt3}$, pas $\frac13$.}
  {$\braket{\Psi}{\Psi}=N^2\big(1+|1-\ii|^2\big)=N^2(1+2)=3N^2$.}
  {On élève $1-\ii$ au carré sans prendre le module : $(1-\ii)^2=-2\ii$ n'est pas réel ; il faut $|1-\ii|^2=2$.}

% ---- Question 39 (niveau 2, correct : D)
\qcmq[2]{Chap.~\ref{chap:dirac}\,$\cdot$\,\S\,\ref{sec:hilbert}}
  {Soient $\ket{\psi}=\frac{1}{\sqrt2}(\ket{0}+\ii\ket{1})$ et $\ket{\varphi}=\frac{1}{\sqrt2}(\ket{0}-\ket{1})$. Que vaut $\braket{\psi}{\varphi}$ ?}
  {$\frac{1-\ii}{2}$}
  {$\frac12$}
  {$0$}
  {$\frac{1+\ii}{2}$}
  {D}
\qcmx
  {On oublie de conjuguer l'amplitude $\ii$ du bra : on calcule $\frac12\big(1+\ii\cdot(-1)\big)$.}
  {Seul le terme en $\ket{0}$ est compté ; le terme en $\ket{1}$ vaut $\frac12(-\ii)(-1)=\frac{\ii}{2}$, non nul.}
  {Confusion avec $\braket{+}{-}=0$ ; ici $|\braket{\psi}{\varphi}|^2=\frac12$.}
  {Le bra conjugue les amplitudes : $\bra{\psi}=\frac{1}{\sqrt2}(\bra{0}-\ii\bra{1})$, d'où $\frac12\big(1+(-\ii)(-1)\big)=\frac{1+\ii}{2}$.}

% ---- Question 40 (niveau 2, correct : A)
\qcmq[2]{Chap.~\ref{chap:dirac}\,$\cdot$\,\S\,\ref{sec:qubit}}
  {Lequel de ces états est orthogonal à $\ket{{+}\ii}=\frac{1}{\sqrt2}(\ket{0}+\ii\ket{1})$ ?}
  {$\ket{{-}\ii}=\frac{1}{\sqrt2}(\ket{0}-\ii\ket{1})$}
  {$\ket{-}=\frac{1}{\sqrt2}(\ket{0}-\ket{1})$}
  {$\ket{+}=\frac{1}{\sqrt2}(\ket{0}+\ket{1})$}
  {$\ket{1}$}
  {A}
\qcmx
  {$\braket{{+}\ii}{{-}\ii}=\frac12\big(1+(-\ii)(-\ii)\big)=\frac12(1-1)=0$, en conjuguant le bra.}
  {$\braket{{+}\ii}{-}=\frac{1+\ii}{2}\neq0$ : $\ket{-}$ est orthogonal à $\ket{+}$, pas à $\ket{{+}\ii}$.}
  {$\braket{{+}\ii}{+}=\frac{1-\ii}{2}\neq0$.}
  {$\braket{{+}\ii}{1}=\frac{-\ii}{\sqrt2}\neq0$ : $\ket{1}$ est orthogonal à $\ket{0}$ seulement.}

% ---- Question 41 (niveau 2, correct : D)
\qcmq[1]{Chap.~\ref{chap:operateurs}\,$\cdot$\,\S\,\ref{sec:pur-mixte}}
  {Une matrice de densité de qubit a pour pureté $\operatorname{Tr}(\rho^2)=\frac58$. Que peut-on en conclure ?}
  {L'état est pur}
  {La matrice n'est pas une matrice de densité valide}
  {L'état est maximalement mixte, $\rho=\mathrm{Id}/2$}
  {L'état est mixte (mélange statistique)}
  {D}
\qcmx
  {Il faudrait $\operatorname{Tr}\rho^2=1$.}
  {$\frac58$ est compris entre $\frac12$ et $1$ : valeur admissible (exemple du cours : $\mathrm{diag}(\frac34,\frac14)$).}
  {Ce cas donnerait $\operatorname{Tr}\rho^2=\frac12$.}
  {$\operatorname{Tr}\rho^2<1$ ; pour un qubit, $\frac12\leq\operatorname{Tr}\rho^2\leq1$ et $\frac58$ est une valeur admissible.}

% ---- Question 42 (niveau 2, correct : B)
\qcmq[2]{Chap.~\ref{chap:operateurs}\,$\cdot$\,\S\,\ref{sec:rho}}
  {Que vaut $\langle Z\rangle=\operatorname{Tr}(\rho Z)$ pour un qubit dans l'état $\ket{+}$ ?}
  {$1$}
  {$0$}
  {$-1$}
  {$\frac{1}{\sqrt2}$}
  {B}
\qcmx
  {C'est la valeur pour $\ket{0}$.}
  {$\rho=\frac12\begin{pmatrix}1&1\\1&1\end{pmatrix}$ et $\operatorname{Tr}(\rho Z)=\frac12-\frac12=0$ : les résultats $+1$ et $-1$ sont équiprobables.}
  {C'est la valeur pour $\ket{1}$.}
  {C'est l'amplitude de $\ket{0}$ dans $\ket{+}$, pas une valeur moyenne de $Z$.}

% ---- Question 43 (niveau 2, correct : A)
\qcmq[1]{Chap.~\ref{chap:operateurs}\,$\cdot$\,\S\,\ref{sec:bloch}}
  {Sur la boule de Bloch, que représente un point situé à l'intérieur de la sphère ($|\vec r|<1$) ?}
  {Un état mixte, de pureté $\operatorname{Tr}\rho^2=\frac{1+|\vec r|^2}{2}<1$}
  {Un état pur dont la norme n'est pas $1$}
  {Un état pur en superposition}
  {Un point sans signification physique}
  {A}
\qcmx
  {Les états purs sont sur la sphère ; le centre $\vec r=\vec0$ est l'état maximalement mixte $\mathrm{Id}/2$.}
  {Tout état pur est normalisé et situé sur la sphère.}
  {Les états purs, superpositions comprises, sont tous \emph{sur} la sphère.}
  {Tout point de la boule est une matrice de densité valide.}

% ---- Question 44 (niveau 2, correct : B)
\qcmq[1]{Chap.~\ref{chap:operateurs}\,$\cdot$\,\S\,\ref{sec:pur-mixte}}
  {Par quoi la matrice de densité distingue-t-elle l'état pur $\ket{+}$ du mélange équiprobable de $\ket{0}$ et $\ket{1}$ ?}
  {Par leur trace, égale à $1$ pour l'état pur seulement}
  {Par les éléments hors diagonale (cohérences), non nuls pour $\ket{+}$ et nuls pour le mélange}
  {Par leurs éléments diagonaux, qui sont différents}
  {Ils ne se distinguent pas : c'est la même matrice}
  {B}
\qcmx
  {La trace vaut $1$ pour toute matrice de densité.}
  {$\ket{+}\bra{+}=\frac12\begin{pmatrix}1&1\\1&1\end{pmatrix}$ et le mélange vaut $\frac12\mathrm{Id}$ ; une mesure en base $X$ les distingue.}
  {Les deux ont la diagonale $(\frac12,\frac12)$ : $P(0)=P(1)=\frac12$ dans les deux cas.}
  {Les matrices sont différentes : en base $X$, $\ket{+}$ donne $+$ avec certitude, le mélange donne $\pm$ à pile ou face.}

% ---- Question 45 (niveau 2, correct : C)
\qcmq[2]{Chap.~\ref{chap:operateurs}\,$\cdot$\,\S\,\ref{sec:rho}}
  {Un qubit est préparé dans $\ket{0}$ avec la probabilité $\frac34$ et dans $\ket{1}$ avec la probabilité $\frac14$. Que vaut $\langle Z\rangle=\operatorname{Tr}(\rho Z)$ ?}
  {$\frac34$}
  {$1$}
  {$\frac12$}
  {$0$}
  {C}
\qcmx
  {C'est la probabilité $\rho_{00}$ de lire $0$ ; or $Z$ vaut $+1$ pour $0$ et $-1$ pour $1$.}
  {Valeur pour $\rho=\ket{0}\bra{0}$ (préparation certaine de $\ket{0}$).}
  {$\rho=\mathrm{diag}(\frac34,\frac14)$, donc $\operatorname{Tr}(\rho Z)=\frac34\cdot1+\frac14\cdot(-1)=\frac12$ : c'est la coordonnée $r_z$ de la sphère de Bloch.}
  {Valeur pour le mélange équiprobable $\mathrm{Id}/2$.}

% ---- Question 46 (niveau 2, correct : B)
\qcmq[2]{Chap.~\ref{chap:operateurs}\,$\cdot$\,\S\,\ref{sec:tenseur}}
  {Dans la base ordonnée $\{\ket{00},\ket{01},\ket{10},\ket{11}\}$, quel vecteur colonne représente $\ket{1}\otimes\ket{0}=\ket{10}$ ?}
  {$(0,1,0,0)^{\mathsf T}$}
  {$(0,0,1,0)^{\mathsf T}$}
  {$(1,0,0,0)^{\mathsf T}$}
  {$(0,0,0,1)^{\mathsf T}$}
  {B}
\qcmx
  {C'est $\ket{01}$ : l'ordre des facteurs est inversé.}
  {$\binom{0}{1}\otimes\binom{1}{0}=\big(0\binom{1}{0},\,1\binom{1}{0}\big)=(0,0,1,0)^{\mathsf T}$ : c'est le troisième vecteur de la base.}
  {C'est $\ket{00}$.}
  {C'est $\ket{11}$.}

% ---- Question 47 (niveau 2, correct : D)
\qcmq[2]{Chap.~\ref{chap:portes}\,$\cdot$\,\S\,\ref{sec:phase}}
  {Quelle relation lie la porte de phase $S=\mathrm{diag}(1,\ii)$ à la porte $Z$ ?}
  {$S=Z$}
  {$S^2=X$}
  {$S^4=Z$}
  {$S^2=Z$}
  {D}
\qcmx
  {$Z=\mathrm{diag}(1,-1)$ : $S$ est un quart de tour, $Z$ un demi-tour.}
  {$S^2$ est diagonale alors que $X$ ne l'est pas ; c'est $HSH$ qui est une racine carrée de $X$.}
  {$S^4=\mathrm{diag}(1,\ii^4)=\mathrm{Id}$.}
  {$S^2=\mathrm{diag}(1,\ii^2)=\mathrm{diag}(1,-1)=Z$ : deux quarts de tour autour de $z$ font un demi-tour.}

% ---- Question 48 (niveau 2, correct : A)
\qcmq[2]{Chap.~\ref{chap:portes}\,$\cdot$\,\S\,\ref{sec:controlees}}
  {Dans la porte de Toffoli (CCNOT) à trois qubits $x,y,z$, à quelle condition la cible $z$ est-elle inversée ?}
  {Seulement si $x=1$ et $y=1$}
  {Si $x=1$ ou $y=1$}
  {Si $x\neq y$}
  {Si $x=0$ et $y=0$}
  {A}
\qcmx
  {$\ket{x,y,z}\to\ket{x,y,z\oplus xy}$ : la cible est inversée si et seulement si le produit $xy$ vaut $1$.}
  {Un seul contrôle à $1$ ne suffit pas : pour $(x,y)=(1,0)$, $xy=0$ et la cible ne change pas.}
  {Ce serait deux CNOT ($z\to z\oplus x\oplus y$) ; pour $(1,0)$ le Toffoli n'inverse rien.}
  {Pour $(0,0)$ on a $xy=0$ : la cible ne change pas.}

% ---- Question 49 (niveau 2, correct : A)
\qcmq[1]{Chap.~\ref{chap:portes}\,$\cdot$\,\S\,\ref{sec:portes}}
  {Pour un hamiltonien $\Hop$ indépendant du temps, quel opérateur d'évolution $U(t,t_0)$ vérifie $\ket{\Psi(t)}=U(t,t_0)\ket{\Psi(t_0)}$ ?}
  {$U(t,t_0)=\ee^{-\ii\Hop(t-t_0)/\hbar}$}
  {$U(t,t_0)=\ee^{+\ii\Hop(t-t_0)/\hbar}$}
  {$U(t,t_0)=-\ii\hbar\,\partial_t$}
  {$U(t,t_0)=\Hop\,(t-t_0)/\hbar$}
  {A}
\qcmx
  {Solution de $\ii\hbar\,\partial_t\ket{\Psi}=\Hop\ket{\Psi}$ ; comme $\Hop$ est hermitien, $U^\dagger U=\mathrm{Id}$ : l'évolution est unitaire.}
  {Erreur de signe : c'est $U^\dagger$, qui remonte le temps.}
  {C'est un opérateur différentiel, pas un opérateur d'évolution ; l'équation de Schrödinger s'écrit $\ii\hbar\,\partial_t\ket{\Psi}=\Hop\ket{\Psi}$.}
  {On oublie l'exponentielle et le facteur $\ii$ : cet opérateur est hermitien, pas unitaire.}

% ---- Question 50 (niveau 2, correct : D)
\qcmq[2]{Chap.~\ref{chap:portes}\,$\cdot$\,\S\,\ref{sec:controlees}}
  {La porte de Fredkin (SWAP contrôlé) a pour contrôle le premier qubit. Que donne-t-elle sur $\ket{1,0,1}$ ?}
  {$\ket{1,0,1}$}
  {$\ket{0,1,1}$}
  {$\ket{1,1,1}$}
  {$\ket{1,1,0}$}
  {D}
\qcmx
  {C'est le résultat si le contrôle valait $0$ ; ici il vaut $1$ et les deux autres qubits sont différents.}
  {Le SWAP porte sur les qubits $2$ et $3$, pas sur le contrôle.}
  {Un SWAP ne change pas le nombre de $1$, il permute seulement ; $\ket{1,1,1}$ viendrait d'un NOT.}
  {Le contrôle vaut $1$ : les deux autres qubits sont échangés, $\ket{0,1}\to\ket{1,0}$.}

% ---- Question 51 (niveau 2, correct : A)
\qcmq[1]{Chap.~\ref{chap:portes}\,$\cdot$\,\S\,\ref{sec:op2}}
  {Que donne $(H\otimes I)\ket{00}$, où $H$ agit sur le premier facteur du produit tensoriel $\ket{a}\otimes\ket{b}$ ?}
  {$\ket{+}\otimes\ket{0}=\frac{1}{\sqrt2}(\ket{00}+\ket{10})$}
  {$\frac12(\ket{00}+\ket{01}+\ket{10}+\ket{11})$}
  {$\ket{0}\otimes\ket{+}=\frac{1}{\sqrt2}(\ket{00}+\ket{01})$}
  {$\frac{1}{\sqrt2}(\ket{00}+\ket{11})$}
  {A}
\qcmx
  {$(A\otimes B)(\ket{u}\otimes\ket{v})=A\ket{u}\otimes B\ket{v}$ : $H\ket{0}=\ket{+}$ sur le premier facteur, $I\ket{0}=\ket{0}$ sur le second.}
  {C'est $(H\otimes H)\ket{00}$ : ici le second qubit n'est pas touché.}
  {$H$ agit sur le premier qubit, pas sur le second : les facteurs sont inversés.}
  {C'est $\ket{\Phi^+}$ : il faut un CNOT pour intriquer ; des portes locales laissent l'état séparable.}

% ---- Question 52 (niveau 2, correct : C)
\qcmq[1]{Chap.~\ref{chap:bell}\,$\cdot$\,\S\,\ref{sec:separable}}
  {Pourquoi l'état de Bell $\ket{\Phi^+}=\frac{1}{\sqrt2}(\ket{00}+\ket{11})$ est-il dit intriqué ?}
  {Parce que $P(01)=0$}
  {Parce que chaque qubit, pris seul, est dans un état pur}
  {Il ne peut pas s'écrire comme un produit $\ket{\varphi}_A\otimes\ket{\chi}_B$ : $D=c_{00}c_{11}-c_{01}c_{10}=\frac12\neq0$}
  {Parce qu'il est orthogonal à $\ket{00}$}
  {C}
\qcmx
  {C'est vrai, mais $\ket{00}$ a aussi $P(01)=0$ et il est séparable : cette condition ne prouve rien.}
  {C'est l'inverse : pris seul, chaque qubit est complètement mixte ($\rho_A=\mathrm{Id}/2$).}
  {Un produit imposerait $D=0$ ; aucun des deux qubits n'a d'état propre.}
  {$\braket{00}{\Phi^+}=\frac{1}{\sqrt2}\neq0$ ; l'orthogonalité n'a pas de lien avec l'intrication.}

% ---- Question 53 (niveau 2, correct : B)
\qcmq[2]{Chap.~\ref{chap:bell}\,$\cdot$\,\S\,\ref{sec:teleportation}}
  {Dans la téléportation du cours, Alice obtient les résultats $a=0$ et $q=1$. Quelle correction Bob doit-il appliquer à son qubit ?}
  {La porte $X$}
  {La porte $Z$}
  {$X$ puis $Z$}
  {Aucune, l'état est déjà $\ket{\varphi}$}
  {B}
\qcmx
  {$X$ corrige le cas $a=1$ (inversion de bit) ; ici $a=0$.}
  {Bob détient $\alpha\ket{0}-\beta\ket{1}=Z\ket{\varphi}$ ; $Z$ rétablit le signe : $Z(\alpha\ket{0}-\beta\ket{1})=\alpha\ket{0}+\beta\ket{1}$.}
  {Les deux corrections ne sont nécessaires que pour $(a,q)=(1,1)$.}
  {Ce n'est vrai que pour $(a,q)=(0,0)$.}

% ---- Question 54 (niveau 2, correct : B)
\qcmq[2]{Chap.~\ref{chap:bell}\,$\cdot$\,\S\,\ref{sec:bell} ; ex.\,\ref{ex:superdense}}
  {Dans le codage superdense, Alice et Bob partagent $\ket{\Phi^+}$ et Alice agit sur son qubit. Quelle porte transforme $\ket{\Phi^+}$ en $\ket{\Psi^+}=\frac{1}{\sqrt2}(\ket{01}+\ket{10})$ ?}
  {La porte $Z$}
  {La porte $X$}
  {La porte identité}
  {La porte $H$}
  {B}
\qcmx
  {$Z$ donne $\ket{\Phi^-}=\frac{1}{\sqrt2}(\ket{00}-\ket{11})$ : elle change un signe, pas les chiffres.}
  {$X$ échange $\ket{0}$ et $\ket{1}$ sur un qubit : $\ket{00}+\ket{11}\to\ket{01}+\ket{10}$ ; c'est le message $10$.}
  {Elle laisse $\ket{\Phi^+}$ : c'est le message $00$.}
  {$(H\otimes I)\ket{\Phi^+}$ est une superposition de deux états de Bell, pas un état de Bell.}

% ---- Question 55 (niveau 2, correct : D)
\qcmq[2]{Chap.~\ref{chap:bell}\,$\cdot$\,\S\,\ref{sec:info}}
  {Pour deux qubits dans l'état de Bell $\ket{\Phi^+}$, que vaut l'entropie conditionnelle $H(A|B)=H(A,B)-H(B)$ ?}
  {$0$ bit}
  {$+1$ bit}
  {$+2$ bits}
  {$-1$ bit}
  {D}
\qcmx
  {C'est la valeur classique quand $A$ est entièrement déterminé par $B$ (deux bits identiques) ; ici $H(A,B)=0$, pas $H(B)$.}
  {C'est la valeur quand $A$ et $B$ sont indépendants : $H(A|B)=H(A)=1$.}
  {C'est l'information mutuelle $I(A;B)=1+1-0=2$ bits, pas l'entropie conditionnelle.}
  {$H(A,B)=S(\rho_{AB})=0$ (état pur) et $H(B)=1$ bit ($\rho_B=\mathrm{Id}/2$) : $0-1=-1$, valeur impossible en classique.}

% ---- Question 56 (niveau 2, correct : B)
\qcmq[1]{Chap.~\ref{chap:bell}\,$\cdot$\,\S\,\ref{sec:trace}}
  {Quelle est la matrice de densité réduite $\rho_A=\operatorname{Tr}_B\rho_{AB}$ d'un qubit d'une paire dans l'état $\ket{\Phi^+}$ ?}
  {$\rho_A=\ket{0}\bra{0}$}
  {$\rho_A=\frac12\mathrm{Id}$, de pureté $\operatorname{Tr}\rho_A^2=\frac12$}
  {$\rho_A=\ket{+}\bra{+}=\frac12\begin{pmatrix}1&1\\1&1\end{pmatrix}$}
  {$\rho_A=\ket{\Phi^+}\bra{\Phi^+}$}
  {B}
\qcmx
  {Le qubit aurait un état défini, ce qui est faux pour un état intriqué : $P(0)=P(1)=\frac12$.}
  {$\rho_A=\frac12\big(\ket{0}\bra{0}+\ket{1}\bra{1}\big)$ : les termes croisés $\ket{00}\bra{11}$ s'annulent ($\braket{1}{0}=0$). Le couple est pur, mais chaque qubit est complètement mixte.}
  {Les termes croisés sont multipliés par $\braket{1}{0}=0$ dans la trace partielle : ils disparaissent.}
  {C'est la matrice $4\times4$ du couple ; la trace partielle donne une matrice $2\times2$.}

% ---- Question 57 (niveau 2, correct : C)
\qcmq[1]{Chap.~\ref{chap:bell}\,$\cdot$\,\S\,\ref{sec:separable}}
  {Pour $\ket{\Psi}=c_{00}\ket{00}+c_{01}\ket{01}+c_{10}\ket{10}+c_{11}\ket{11}$, quel est le critère de séparabilité ?}
  {$c_{00}c_{11}+c_{01}c_{10}=0$}
  {$c_{00}c_{01}-c_{10}c_{11}=0$}
  {$D=c_{00}c_{11}-c_{01}c_{10}=0$}
  {$c_{00}=c_{11}$ et $c_{01}=c_{10}$}
  {C}
\qcmx
  {Mauvais signe : $\ket{+}\otimes\ket{-}$ est séparable mais donnerait $-\frac12\neq0$.}
  {Mauvais appariement : $\ket{\Psi^+}$ (intriqué) la vérifie, puisque $c_{00}=c_{11}=0$.}
  {Un produit $(a_0\ket{0}+a_1\ket{1})\otimes(b_0\ket{0}+b_1\ket{1})$ a $c_{ij}=a_ib_j$, donc $c_{00}c_{11}=a_0a_1b_0b_1=c_{01}c_{10}$ ; la réciproque est vraie aussi.}
  {Ni nécessaire ($\ket{00}$ est séparable avec $c_{00}\neq c_{11}$) ni suffisant ($\ket{\Phi^+}$ la vérifie et est intriqué).}

% ---- Question 58 (niveau 2, correct : B)
\qcmq[2]{Chap.~\ref{chap:bell}\,$\cdot$\,\S\,\ref{sec:bell}}
  {Quelle est la valeur propre de $X\otimes X$ pour l'état $\ket{\Phi^+}$ ?}
  {$-1$}
  {$+1$}
  {$0$}
  {$\ket{\Phi^+}$ n'est pas un état propre de $X\otimes X$}
  {B}
\qcmx
  {C'est la valeur pour $\ket{\Phi^-}$ et $\ket{\Psi^-}$ : le signe $-$ de l'état se retrouve dans la valeur propre.}
  {$(X\otimes X)\ket{00}=\ket{11}$ et $(X\otimes X)\ket{11}=\ket{00}$, donc $(X\otimes X)\ket{\Phi^+}=\ket{\Phi^+}$.}
  {$X\otimes X$ est unitaire : ses valeurs propres sont de module $1$, jamais nulles.}
  {Les quatre états de Bell sont des états propres communs de $Z\otimes Z$ et de $X\otimes X$.}

% ---- Question 59 (niveau 2, correct : A)
\qcmq[2]{Chap.~\ref{chap:bell}\,$\cdot$\,\S\,\ref{sec:info}}
  {Que vaut l'entropie de von Neumann $S(\rho)=-\sum_i\lambda_i\log_2\lambda_i$ d'un état pur ?}
  {$0$ bit}
  {$1$ bit}
  {$-1$ bit}
  {$+\infty$}
  {A}
\qcmx
  {Les valeurs propres sont $1$ et $0$ : $1\cdot\log_21=0$ et $0\cdot\log_20=0$ (par continuité, $x\log_2x\to0$).}
  {C'est la valeur pour l'état maximalement mixte $\mathrm{Id}/2$ d'un qubit.}
  {L'entropie de $\rho$ est positive ou nulle ; seule l'entropie conditionnelle quantique peut être négative.}
  {$\log_20$ diverge, mais le terme $0\cdot\log_20$ vaut $0$ par continuité.}

% ---- Question 60 (niveau 2, correct : A)
\qcmq[2]{Chap.~\ref{chap:bell}\,$\cdot$\,\S\,\ref{sec:bell}}
  {Lequel des états de Bell a pour valeur propre $-1$ à la fois pour $Z\otimes Z$ et pour $X\otimes X$ ?}
  {$\ket{\Psi^-}=\frac{1}{\sqrt2}(\ket{01}-\ket{10})$}
  {$\ket{\Phi^+}$}
  {$\ket{\Phi^-}$}
  {$\ket{\Psi^+}$}
  {A}
\qcmx
  {Résultats opposés dans les deux bases : $Z\otimes Z=-1$ (chiffres différents) et $X\otimes X=-1$ (signe $-$ dans l'état).}
  {Valeurs propres $(Z\otimes Z,\,X\otimes X)=(+1,+1)$.}
  {Valeurs propres $(Z\otimes Z,\,X\otimes X)=(+1,-1)$.}
  {Valeurs propres $(Z\otimes Z,\,X\otimes X)=(-1,+1)$.}

\clearpage
%% ======================================================================
\subsection{Niveau 3 : avancé (questions 61 à 90)}\label{sec:qcm-n3}
%% ======================================================================
\noindent{\color{insagris}\small Niveau de difficulté : Avancé\quad $\bullet\bullet\bullet\circ$}\par\smallskip
% ---- Question 61 (niveau 3, correct : A)
\qcmq[2]{Chap.~\ref{chap:histoire}\,$\cdot$\,\S\,\ref{sec:solutions}}
  {Un électron d'un puits infini passe du niveau $n=3$ au niveau $n=2$ en émettant un photon. Quelle est l'énergie du photon, en fonction de $E_1$ ?}
  {$5E_1$}
  {$E_1$}
  {$\frac94E_1$}
  {$13E_1$}
  {A}
\qcmx
  {$E_3-E_2=9E_1-4E_1=5E_1$ ; le photon a la fréquence $f=(E_3-E_2)/h$.}
  {On a soustrait les numéros de niveau ($3-2=1$) au lieu de leurs carrés.}
  {C'est le rapport $E_3/E_2$, pas la différence $E_3-E_2$.}
  {On a additionné $9E_1+4E_1$ : l'énergie du photon est une différence de niveaux.}

% ---- Question 62 (niveau 3, correct : A)
\qcmq[2]{Chap.~\ref{chap:histoire}\,$\cdot$\,\S\,\ref{sec:superposition}}
  {À $t=0$, une particule du puits infini est dans l'état $\Psi=\frac{\sqrt3}{2}\,\psi_1+\frac12\,\psi_3$ ($\psi_n$ : états stationnaires normalisés, d'énergies $E_n=n^2E_1$). Quelle est son énergie moyenne $\langle E\rangle$ ?}
  {$3E_1$}
  {$5E_1$}
  {$\left(\frac{\sqrt3}{2}+\frac92\right)E_1$}
  {$7E_1$}
  {A}
\qcmx
  {$P(E_1)=\frac34$ et $P(E_3)=\frac14$, donc $\langle E\rangle=\frac34E_1+\frac14\cdot9E_1=3E_1$.}
  {On a pondéré également les deux niveaux, $\frac{E_1+9E_1}{2}$ : les poids sont les $|c_i|^2$, pas $\frac12$.}
  {On a pondéré par les amplitudes $c_i$ au lieu de leurs carrés $|c_i|^2$.}
  {Poids permutés : $\frac14E_1+\frac34\cdot9E_1=7E_1$ ; la probabilité de $E_1$ est $\left|\frac{\sqrt3}{2}\right|^2=\frac34$.}

% ---- Question 63 (niveau 3, correct : C)
\qcmq[1]{Chap.~\ref{chap:histoire}\,$\cdot$\,\S\,\ref{sec:solutions}}
  {Pour l'état stationnaire $n=2$ du puits infini $]0,a[$, que valent la densité de probabilité de présence en $x=a/2$ et la position moyenne $\langle\hat x\rangle$ ?}
  {Densité maximale en $a/2$, et $\langle\hat x\rangle=a/2$}
  {Densité nulle en $a/2$, et $\langle\hat x\rangle=0$}
  {Densité nulle en $a/2$, et $\langle\hat x\rangle=a/2$}
  {Densité maximale en $a/2$, et $\langle\hat x\rangle=a/4$}
  {C}
\qcmx
  {C'est le cas de $n=1$ ; pour $n=2$ il y a un nœud au centre.}
  {Une densité nulle au centre ne donne pas une moyenne nulle : $\langle\hat x\rangle$ pondère tout le puits et vaut $\frac a2$ par symétrie.}
  {$|\Psi_2|^2=\frac2a\sin^2\frac{2\pi x}{a}$ s'annule en $x=\frac a2$ (nœud), alors que la symétrie impose $\langle\hat x\rangle=\frac a2$ : la moyenne n'est pas la position la plus probable.}
  {Ni l'un ni l'autre : maxima en $\frac a4$ et $\frac{3a}4$, nœud en $\frac a2$, moyenne $\frac a2$.}

% ---- Question 64 (niveau 3, correct : C)
\qcmq[2]{Chap.~\ref{chap:dirac}\,$\cdot$\,\S\,\ref{sec:projecteurs}}
  {Quelle est la matrice du projecteur $\hat P_+=\ket{+}\bra{+}$ dans la base $\{\ket{0},\ket{1}\}$ ?}
  {$\frac12\begin{pmatrix}1&-1\\-1&1\end{pmatrix}$}
  {$\begin{pmatrix}1&0\\0&0\end{pmatrix}$}
  {$\frac12\begin{pmatrix}1&1\\1&1\end{pmatrix}$}
  {$\frac{1}{\sqrt2}\begin{pmatrix}1&1\\1&-1\end{pmatrix}$}
  {C}
\qcmx
  {C'est $\ket{-}\bra{-}$, le projecteur sur $\ket{-}$.}
  {C'est $\ket{0}\bra{0}$ : projecteur de la base $Z$, pas de la base $X$.}
  {$\ket{+}\bra{+}=\frac12(\ket0+\ket1)(\bra0+\bra1)$ ; on vérifie $\hat P_+^2=\hat P_+$ et $\operatorname{Tr}\hat P_+=1$.}
  {C'est la matrice de Hadamard, unitaire ($H^2=\mathrm{Id}$), et non un projecteur (un projecteur vérifie $P^2=P$).}

% ---- Question 65 (niveau 3, correct : B)
\qcmq[2]{Chap.~\ref{chap:dirac}\,$\cdot$\,\S\,\ref{sec:projecteurs}}
  {Un qubit préparé dans $\ket{0}$ est mesuré en base $X$, puis aussitôt en base $Z$. Quelle est la probabilité d'obtenir $1$ à la seconde mesure ?}
  {$0$}
  {$\frac12$}
  {$1$}
  {$\frac14$}
  {B}
\qcmx
  {On garde $\ket{0}$ en mémoire : la mesure en base $X$ a effondré l'état sur $\ket{\pm}$ et effacé $\ket{0}$.}
  {La première mesure donne $\ket{+}$ ou $\ket{-}$ (probabilité $\frac12$ chacun) ; dans les deux cas $P(1)=\frac12$ en base $Z$ : $\frac12\cdot\frac12+\frac12\cdot\frac12=\frac12$.}
  {Sans rapport : $\ket{0}$ n'est plus l'état après la mesure en base $X$.}
  {On n'a compté qu'une branche de la première mesure ; il faut sommer les deux ($\frac14+\frac14$).}

% ---- Question 66 (niveau 3, correct : D)
\qcmq[2]{Chap.~\ref{chap:dirac}\,$\cdot$\,\S\,\ref{sec:projecteurs}}
  {Un qubit est dans l'état $\ket{\psi}=\frac{3}{5}\ket{0}+\frac{4}{5}\ket{1}$. Quelle est la probabilité d'obtenir $-$ en mesurant dans la base $X$ ?}
  {$\frac{16}{25}$}
  {$\frac{1}{25}$}
  {$\frac{49}{50}$}
  {$\frac{1}{50}$}
  {D}
\qcmx
  {C'est $P(1)$ en base $Z$ ($|\frac45|^2$) : on a mesuré dans la mauvaise base.}
  {On a oublié le facteur $\frac{1}{\sqrt2}$ de $\ket{-}$ : $\left|\frac35-\frac45\right|^2=\frac1{25}$, il manque le $\frac12$.}
  {C'est $P(+)$ ; la probabilité demandée est celle du résultat $-$.}
  {$\braket{-}{\psi}=\frac{1}{\sqrt2}\left(\frac35-\frac45\right)=-\frac{1}{5\sqrt2}$, donc $P(-)=\frac1{50}$ ; et $P(+)=\frac{49}{50}$.}

% ---- Question 67 (niveau 3, correct : B)
\qcmq[1]{Chap.~\ref{chap:dirac}\,$\cdot$\,\S\,\ref{sec:projecteurs}}
  {Lequel de ces kets décrit le même état physique que $\ket{+}=\frac{1}{\sqrt2}(\ket{0}+\ket{1})$ ?}
  {$\frac{1}{\sqrt2}\big(\ket{0}+\ii\ket{1}\big)$}
  {$\ee^{\ii\pi/3}\ket{+}$}
  {$\frac{1}{\sqrt2}\big(\ket{0}+\ee^{\ii\pi}\ket{1}\big)$}
  {$\frac{1}{\sqrt2}\big(\ket{0}+\ee^{\ii\pi/3}\ket{1}\big)$}
  {B}
\qcmx
  {La phase $\ii$ ne porte que sur $\ket{1}$ : c'est une phase \emph{relative}, et l'état est $\ket{{+}\ii}$ (base $Y$).}
  {Un facteur de phase \emph{globale} multiplie tout le ket : toutes les probabilités de mesure (et $\rho=\ket{+}\bra{+}$) sont inchangées.}
  {$\ee^{\ii\pi}=-1$ : c'est $\ket{-}$, orthogonal à $\ket{+}$.}
  {Phase relative $\frac\pi3$ : c'est un autre état, sur l'équateur de la sphère de Bloch, à l'azimut $\frac\pi3$.}

% ---- Question 68 (niveau 3, correct : A)
\qcmq[2]{Chap.~\ref{chap:dirac}\,$\cdot$\,\S\,\ref{sec:qubit}}
  {Laquelle de ces paires est une base orthonormée de $\mathbb{C}^2$ ?}
  {$\{\ket{{+}\ii},\ket{{-}\ii}\}$}
  {$\{\ket{0},\ket{+}\}$}
  {$\{\ket{+},\ket{{-}\ii}\}$}
  {$\{\ket{0},\ket{0}+\ket{1}\}$}
  {A}
\qcmx
  {Ces deux états sont normés et orthogonaux : $\braket{{+}\ii}{{-}\ii}=\frac12(1-1)=0$.}
  {Normés, mais $\braket{0}{+}=\frac{1}{\sqrt2}\neq0$ : pas orthogonaux.}
  {$\braket{+}{{-}\ii}=\frac12(1-\ii)\neq0$ : pas orthogonaux.}
  {Le second ket n'est pas normé (norme au carré $2$) ni orthogonal à $\ket{0}$.}

% ---- Question 69 (niveau 3, correct : C)
\qcmq[2]{Chap.~\ref{chap:operateurs}\,$\cdot$\,\S\,\ref{sec:rho}}
  {Laquelle de ces matrices peut être la matrice de densité d'un qubit ?}
  {$\begin{pmatrix}\frac12&\frac34\\\frac34&\frac12\end{pmatrix}$}
  {$\begin{pmatrix}\frac12&\frac{\ii}{2}\\\frac{\ii}{2}&\frac12\end{pmatrix}$}
  {$\begin{pmatrix}\frac23&\frac{\ii}{3}\\-\frac{\ii}{3}&\frac13\end{pmatrix}$}
  {$\begin{pmatrix}\frac34&0\\0&\frac12\end{pmatrix}$}
  {C}
\qcmx
  {Hermitienne et de trace $1$, mais pas positive : $\det=\frac14-\frac9{16}<0$, valeurs propres $\frac54$ et $-\frac14$ (une « probabilité » négative).}
  {De trace $1$ mais pas hermitienne : il faudrait $\rho_{10}=\overline{\rho_{01}}=-\frac{\ii}{2}$.}
  {Hermitienne ($\rho_{10}=\overline{\rho_{01}}$), de trace $1$ et positive : $\det\rho=\frac29-\frac19=\frac19>0$ avec $\operatorname{Tr}\rho>0$ donne deux valeurs propres positives, $\frac{3\pm\sqrt5}{6}$.}
  {Hermitienne et positive, mais de trace $\frac54\neq1$ : les probabilités diagonales ne somment pas à $1$.}

% ---- Question 70 (niveau 3, correct : B)
\qcmq[2]{Chap.~\ref{chap:operateurs}\,$\cdot$\,\S\,\ref{sec:op-dirac}}
  {Un qubit est dans l'état $\ket{\psi}=\frac{1}{\sqrt5}\big(\ket{0}+2\ii\ket{1}\big)$. Que vaut la valeur moyenne $\langle Y\rangle=\bra{\psi}Y\ket{\psi}$, avec $Y=\begin{pmatrix}0&-\ii\\\ii&0\end{pmatrix}$ ?}
  {$0$}
  {$\frac45$}
  {$-\frac45$}
  {$-\frac35$}
  {B}
\qcmx
  {On a oublié de conjuguer les amplitudes du bra : $\psi^{\mathsf T}Y\psi=\frac15(2+2\ii\cdot\ii)=0$ (c'est d'ailleurs la valeur de $\langle X\rangle$ pour cet état).}
  {$Y\ket\psi=\frac1{\sqrt5}(2\ket0+\ii\ket1)$, puis $\bra\psi Y\ket\psi=\frac15\big(1\cdot2+(-2\ii)(\ii)\big)=\frac45$ ; en général $\langle Y\rangle=2\,\mathrm{Im}(\bar\alpha\beta)$.}
  {Erreur de signe : on a utilisé la transposée $\begin{pmatrix}0&\ii\\-\ii&0\end{pmatrix}=-Y$.}
  {C'est $\langle Z\rangle=\frac15-\frac45$ : on a calculé la valeur moyenne d'une autre observable.}

% ---- Question 71 (niveau 3, correct : D)
\qcmq[1]{Chap.~\ref{chap:operateurs}\,$\cdot$\,\S\,\ref{sec:spectral}}
  {On mesure l'observable $A=\begin{pmatrix}2&1\\1&2\end{pmatrix}$ sur un qubit préparé dans $\ket{0}$. Quelle affirmation est correcte ?}
  {On obtient toujours $2$}
  {On obtient $+1$ ou $-1$, avec la probabilité $\frac12$ chacun}
  {On obtient $2$ ou $1$, avec les probabilités $\frac45$ et $\frac15$}
  {On obtient $3$ ou $1$, avec la probabilité $\frac12$ chacun (moyenne $2$)}
  {D}
\qcmx
  {On a lu le coefficient diagonal : $A\ket0=2\ket0+\ket1$ n'est pas proportionnel à $\ket0$, donc $\ket0$ n'est pas vecteur propre de $A$.}
  {Ce sont les valeurs propres de $X$, pas de $A$ : le décalage $2\,\mathrm{Id}$ translate le spectre ($2\pm1$).}
  {On a lu $A\ket0=2\ket0+\ket1$ comme un état : $A$ n'est pas unitaire, ce n'est pas une évolution ; les résultats possibles d'une mesure sont les valeurs propres de $A$.}
  {$A=2\,\mathrm{Id}+X$ a les vecteurs propres de $X$ : valeur propre $3$ pour $\ket+$, $1$ pour $\ket-$. Comme $\ket0=\frac{1}{\sqrt2}(\ket++\ket-)$, les deux résultats sont équiprobables, et $\langle A\rangle=\bra0A\ket0=2$.}

% ---- Question 72 (niveau 3, correct : D)
\qcmq[2]{Chap.~\ref{chap:operateurs}\,$\cdot$\,\S\,\ref{sec:pur-mixte}}
  {Un qubit est préparé dans $\ket{0}$ avec la probabilité $\frac12$ et dans $\ket{+}$ avec la probabilité $\frac12$. Quelle est la pureté $\operatorname{Tr}(\rho^2)$ de son état ?}
  {$1$}
  {$\frac12$}
  {$\frac58$}
  {$\frac34$}
  {D}
\qcmx
  {Un mélange de deux états purs n'est pur que s'ils sont identiques ; ici $\ket0\neq\ket+$, donc $\operatorname{Tr}\rho^2<1$.}
  {Valeur de l'état maximalement mixte $\frac{\mathrm{Id}}2$, atteinte pour deux états \emph{orthogonaux} ; ici ils ne le sont pas (on a aussi pu oublier les termes croisés).}
  {On n'a additionné que les carrés des éléments diagonaux ($\frac9{16}+\frac1{16}$) : il manque $2\rho_{01}\rho_{10}=\frac2{16}$.}
  {$\rho=\frac12\ket0\bra0+\frac12\ket+\bra+=\begin{pmatrix}\frac34&\frac14\\\frac14&\frac14\end{pmatrix}$, donc $\operatorname{Tr}\rho^2=\frac9{16}+\frac1{16}+2\cdot\frac1{16}=\frac34$ : état mixte, mais moins que $\frac{\mathrm{Id}}2$ car $\ket0$ et $\ket+$ ne sont pas orthogonaux.}

% ---- Question 73 (niveau 3, correct : B)
\qcmq[2]{Chap.~\ref{chap:operateurs}\,$\cdot$\,\S\,\ref{sec:bloch}}
  {Quel est le vecteur de Bloch $(r_x,r_y,r_z)$ de l'état $\ket{\psi}=\cos\frac\pi6\,\ket0+\sin\frac\pi6\,\ket1$ ?}
  {$\left(\frac12,\,0,\,\frac{\sqrt3}{2}\right)$}
  {$\left(\frac{\sqrt3}{2},\,0,\,\frac12\right)$}
  {$\left(\frac{\sqrt3}{4},\,0,\,\frac12\right)$}
  {$\left(\frac{\sqrt3}{2},\,0,\,\frac34\right)$}
  {B}
\qcmx
  {On a pris $\theta=\frac\pi6$ : l'angle de Bloch est le \emph{double} de l'angle qui figure dans les amplitudes.}
  {Ici $\frac\theta2=\frac\pi6$, donc $\theta=\frac\pi3$ et $\varphi=0$ : $\vec r=(\sin\theta\cos\varphi,\sin\theta\sin\varphi,\cos\theta)$. Contrôle : $r_z=|\alpha|^2-|\beta|^2=\frac34-\frac14$ et $r_x=2\alpha\beta=\frac{\sqrt3}2$.}
  {On a écrit $r_x=\alpha\beta$ en oubliant le facteur $2$ ; on le voit à $|\vec r|^2=\frac7{16}<1$, impossible pour un état pur.}
  {On a pris $r_z=|\alpha|^2$ au lieu de $|\alpha|^2-|\beta|^2$ ; or $|\vec r|^2=\frac{21}{16}>1$ est impossible.}

% ---- Question 74 (niveau 3, correct : A)
\qcmq[2]{Chap.~\ref{chap:operateurs}\,$\cdot$\,\S\,\ref{sec:bloch}}
  {Un qubit a pour matrice de densité $\rho=\frac12\begin{pmatrix}1&-\ii\\\ii&1\end{pmatrix}$. Quel est son vecteur de Bloch ?}
  {$(0,\,1,\,0)$ : c'est l'état pur $\ket{{+}\ii}$}
  {$(0,\,-1,\,0)$}
  {$\left(0,\,\frac12,\,0\right)$}
  {$(0,\,0,\,0)$}
  {A}
\qcmx
  {Avec $\rho=\frac12\begin{pmatrix}1+r_z&r_x-\ii r_y\\r_x+\ii r_y&1-r_z\end{pmatrix}$ : $r_z=0$ et $r_x-\ii r_y=-\ii$, donc $r_y=1$, $r_x=0$. Comme $|\vec r|=1$, l'état est pur : $\ket{{+}\ii}\bra{{+}\ii}$.}
  {Erreur de signe : $\rho_{01}=\frac{r_x-\ii r_y}{2}$, donc $\rho_{01}=-\frac\ii2$ correspond à $r_y=+1$ ($r_y=-1$ pour $\ket{{-}\ii}$).}
  {On a lu $\rho_{01}=-\frac\ii2$ comme $-\ii r_y$ : il faut multiplier par $2$ ($\rho=\frac12(\mathrm{Id}+\vec r\cdot\vec\sigma)$).}
  {Les diagonales $\frac12,\frac12$ donnent $r_z=0$, mais les éléments non diagonaux (cohérences) ne sont pas nuls : l'état est pur, et non $\frac{\mathrm{Id}}2$.}

% ---- Question 75 (niveau 3, correct : C)
\qcmq[2]{Chap.~\ref{chap:operateurs}\,$\cdot$\,\S\,\ref{sec:tenseur}}
  {Que donne $(X\otimes Z)\,\big(\ket{+}\otimes\ket{-}\big)$ ?}
  {$-\ket{-}\otimes\ket{-}$}
  {$\ket{-}\otimes\ket{+}$}
  {$\ket{+}\otimes\ket{+}$}
  {$\ket{+}\otimes\ket{-}$}
  {C}
\qcmx
  {On a échangé les facteurs ($Z\otimes X$) : $Z\ket+=\ket-$ et $X\ket-=-\ket-$. Le premier opérateur agit sur le premier qubit.}
  {On a cru que $X$ transforme $\ket+$ en $\ket-$ ; or $X$ échange $\ket0$ et $\ket1$ et laisse $\ket+$ inchangé.}
  {$(A\otimes B)(\ket u\otimes\ket v)=A\ket u\otimes B\ket v$ : $X\ket+=\ket+$ ($\ket+$ est propre pour $X$, valeur propre $+1$) et $Z\ket-=\ket+$ ($Z$ échange $\ket+$ et $\ket-$).}
  {On a cru que $Z$ laisse $\ket-$ inchangé ; $Z$ ne laisse invariants que $\ket0$ et $\ket1$, et échange $\ket+$ et $\ket-$.}

% ---- Question 76 (niveau 3, correct : C)
\qcmq[2]{Chap.~\ref{chap:portes}\,$\cdot$\,\S\,\ref{sec:phase}}
  {Le qubit $\ket{+}$ traverse la porte $T=\mathrm{diag}\big(1,\ee^{\ii\pi/4}\big)$, puis est mesuré dans la base $X$ $\{\ket+,\ket-\}$. Quelle est la probabilité d'obtenir $+$ ?}
  {$\frac12$}
  {$\frac{2-\sqrt2}{4}\approx0{,}15$}
  {$\frac{2+\sqrt2}{4}\approx0{,}85$}
  {$\frac{\sqrt2}{2}\approx0{,}71$}
  {C}
\qcmx
  {On a cru que la phase ne change pas les probabilités : c'est vrai en base $Z$, mais la base $X$ est sensible à la phase \emph{relative} (le point de Bloch est à l'azimut $\frac\pi4$ de l'axe $+x$).}
  {C'est $P(-)=\sin^2\frac\pi8$ : on a permuté les résultats $+$ et $-$.}
  {$T\ket+=\frac1{\sqrt2}(\ket0+\ee^{\ii\pi/4}\ket1)$, donc $\braket{+}{\psi}=\frac12\big(1+\ee^{\ii\pi/4}\big)$ et $P(+)=\frac14\big|1+\ee^{\ii\pi/4}\big|^2=\frac{2+2\cos\frac\pi4}{4}=\frac{2+\sqrt2}{4}$ (soit $\cos^2\frac\pi8$).}
  {On a pris $\cos\frac\pi4$ comme probabilité ; il faut le \emph{demi}-angle et le carré : $P=\cos^2\frac{\pi/4}{2}=\cos^2\frac\pi8$.}

% ---- Question 77 (niveau 3, correct : B)
\qcmq[2]{Chap.~\ref{chap:portes}\,$\cdot$\,\S\,\ref{sec:hadamard}}
  {Que vaut le produit de matrices $HYH$, où $Y=\begin{pmatrix}0&-\ii\\\ii&0\end{pmatrix}$ ?}
  {$Y$}
  {$-Y$}
  {$Z$}
  {$X$}
  {B}
\qcmx
  {Par analogie avec $HXH=Z$, on a cru que $H$ ne modifie pas $Y$ ; en fait l'axe $y$ est retourné (signe $-$).}
  {$H$ échange les axes $x$ et $z$ et retourne l'axe $y$ : $HXH=Z$, $HZH=X$, $HYH=-Y$. Direct : $HY=\frac{1}{\sqrt2}\begin{pmatrix}\ii&-\ii\\-\ii&-\ii\end{pmatrix}$, puis $HYH=\begin{pmatrix}0&\ii\\-\ii&0\end{pmatrix}=-Y$.}
  {C'est $HXH$ : $H$ n'échange que les axes $x$ et $z$, l'axe $y$ reste l'axe $y$.}
  {C'est $HZH$ : $Y$ ne peut pas devenir $X$, car $H$ laisse l'axe $y$ à sa place (au signe près).}

% ---- Question 78 (niveau 3, correct : D)
\qcmq[1]{Chap.~\ref{chap:portes}\,$\cdot$\,\S\,\ref{sec:op2}}
  {Que donne $(H\otimes H)\ket{01}$, où $\ket{01}=\ket0\otimes\ket1$ ?}
  {$\frac12\big(\ket{00}+\ket{01}-\ket{10}-\ket{11}\big)$}
  {$\frac12\big(\ket{00}-\ket{01}-\ket{10}+\ket{11}\big)$}
  {$\frac12\big(\ket{00}+\ket{01}+\ket{10}+\ket{11}\big)$}
  {$\frac12\big(\ket{00}-\ket{01}+\ket{10}-\ket{11}\big)$}
  {D}
\qcmx
  {C'est $\ket-\otimes\ket+$ : on a attribué le $\ket-$ au premier qubit, alors que c'est le second qui vaut $\ket1$.}
  {C'est $\ket-\otimes\ket-$ : $H\ket0=\ket+$, et non $\ket-$.}
  {C'est $\ket+\otimes\ket+$, le résultat pour $\ket{00}$ : on a oublié que $H\ket1=\ket-$ introduit des signes.}
  {$H\ket0\otimes H\ket1=\ket+\otimes\ket-=\frac12(\ket0+\ket1)\otimes(\ket0-\ket1)$. Le signe $-$ apparaît quand le \emph{second} chiffre vaut $1$ (formule générale $\frac12\sum_y(-1)^{x\cdot y}\ket y$ avec $x=01$).}

% ---- Question 79 (niveau 3, correct : D)
\qcmq[2]{Chap.~\ref{chap:portes}\,$\cdot$\,\S\,\ref{sec:cnot}}
  {Avec le CNOT du cours (contrôle : second qubit $B$ ; cible : premier qubit $A$ ; $\mathrm{CNOT}\ket{a,b}=\ket{a\oplus b,\,b}$), quel état obtient-on à partir de $\ket{1}_A\otimes\ket{+}_B$ ?}
  {$\ket{1}\otimes\ket{+}$ (inchangé)}
  {$\frac{1}{\sqrt2}\big(\ket{00}+\ket{11}\big)=\ket{\Phi^+}$}
  {$\frac{1}{\sqrt2}\big(\ket{01}-\ket{10}\big)=\ket{\Psi^-}$}
  {$\frac{1}{\sqrt2}\big(\ket{01}+\ket{10}\big)=\ket{\Psi^+}$}
  {D}
\qcmx
  {On a pris le premier qubit comme contrôle : le CNOT agirait alors sur $B$, et $X\ket+=\ket+$ laisserait l'état inchangé. Mais ici le contrôle est $B$.}
  {C'est le résultat pour $\ket{0}_A\ket{+}_B$ ; ici $A$ vaut $\ket1$, ce qui échange les deux termes.}
  {Aucun signe $-$ n'apparaît : le CNOT permute des états de base sans leur ajouter de phase.}
  {$\ket1\ket+=\frac1{\sqrt2}(\ket{10}+\ket{11})$ ; $\ket{10}\mapsto\ket{1\oplus0,0}=\ket{10}$ et $\ket{11}\mapsto\ket{1\oplus1,1}=\ket{01}$. C'est un état de Bell, donc intriqué.}

% ---- Question 80 (niveau 3, correct : C)
\qcmq[2]{Chap.~\ref{chap:portes}\,$\cdot$\,\S\,\ref{sec:swap}}
  {Que vaut $\mathrm{SWAP}\,(X\otimes\mathrm{Id})\,\mathrm{SWAP}$ ?}
  {$X\otimes\mathrm{Id}$}
  {$\mathrm{Id}\otimes\mathrm{Id}$}
  {$\mathrm{Id}\otimes X$}
  {$X\otimes X$}
  {C}
\qcmx
  {On a supposé que $\mathrm{SWAP}$ commute avec $X\otimes\mathrm{Id}$ ; ce n'est pas le cas, puisque cet opérateur n'agit pas de la même façon sur les deux qubits.}
  {On a simplifié $\mathrm{SWAP}^2=\mathrm{Id}$ en oubliant que $X\otimes\mathrm{Id}$ se trouve \emph{entre} les deux SWAP : on ne peut pas les rapprocher.}
  {$\mathrm{SWAP}\,(X\otimes\mathrm{Id})\,\mathrm{SWAP}\ket{a,b}=\mathrm{SWAP}\,(X\otimes\mathrm{Id})\ket{b,a}=\mathrm{SWAP}\ket{b\oplus1,a}=\ket{a,b\oplus1}$ : le SWAP transporte l'opérateur d'un qubit à l'autre (même mécanisme que $\mathrm{CNOT}_{21}=\mathrm{SWAP}\,\mathrm{CNOT}_{12}\,\mathrm{SWAP}$).}
  {Cet opérateur inverserait les deux bits ; ici un seul bit est inversé, celui du \emph{second} qubit.}

% ---- Question 81 (niveau 3, correct : C)
\qcmq[1]{Chap.~\ref{chap:portes}\,$\cdot$\,\S\,\ref{sec:controlees}}
  {La porte $Z$ contrôlée ($\mathrm{CZ}=\mathrm{diag}(1,1,1,-1)$, contrôle : premier qubit) est appliquée à $\ket{+}\otimes\ket{+}$. Quel état obtient-on ?}
  {$\ket{+}\otimes\ket{+}$}
  {$\ket{-}\otimes\ket{-}$}
  {$\frac12\big(\ket{00}+\ket{01}+\ket{10}-\ket{11}\big)$, un état intriqué}
  {$\ket{+}\otimes\ket{-}$}
  {C}
\qcmx
  {On a cru qu'un simple signe est une phase globale ; mais le signe $-$ ne touche que $\ket{11}$ : c'est une phase \emph{relative}.}
  {C'est $Z\otimes Z$ ($Z$ sur les deux qubits), et non une porte contrôlée.}
  {CZ ne change que le signe de $\ket{11}$. Critère de séparabilité : $c_{00}c_{11}-c_{01}c_{10}=-\frac14-\frac14=-\frac12\neq0$ ; on peut l'écrire $\frac{1}{\sqrt2}\big(\ket0\ket++\ket1\ket-\big)$.}
  {On a appliqué $Z$ au second qubit sans condition ; le contrôle exige que le premier qubit soit dans $\ket1$.}

% ---- Question 82 (niveau 3, correct : D)
\qcmq[2]{Chap.~\ref{chap:portes}\,$\cdot$\,\S\,\ref{sec:phase}, \ref{sec:portes-bloch}}
  {Un qubit préparé dans $\ket{0}$ traverse successivement $H$, puis $S$, puis $H$ (dans cet ordre). Quel est l'état final, à une phase globale près ?}
  {$\ket{{+}\ii}$}
  {$\ket{0}$}
  {$\ket{1}$}
  {$\ket{{-}\ii}=\frac{1}{\sqrt2}\big(\ket0-\ii\ket1\big)$}
  {D}
\qcmx
  {Mauvais sens de rotation : on a perdu un signe en calculant $H\ket{{+}\ii}$ (la phase relative vaut $-\ii$, pas $+\ii$).}
  {On a « simplifié » $HSH$ en $S$ (puis $S\ket0=\ket0$) : mais $H$ ne commute pas avec $S$, on ne peut pas supprimer les deux $H$.}
  {C'est le résultat pour $HZH=X$ ; or $S$ n'est qu'un \emph{quart} de tour autour de $z$ (et non un demi-tour comme $Z$), donc $HSH$ est un quart de tour autour de $x$.}
  {$HSH=\sqrt X$ est un quart de tour autour de $x$ (§\,\ref{sec:phase}). Directement : $H\ket0=\ket+$, $S\ket+=\ket{{+}\ii}$, $H\ket{{+}\ii}=\frac{1+\ii}{2}\ket0+\frac{1-\ii}{2}\ket1=\frac{1+\ii}{2}\big(\ket0-\ii\ket1\big)$.}

% ---- Question 83 (niveau 3, correct : A)
\qcmq[1]{Chap.~\ref{chap:bell}\,$\cdot$\,\S\,\ref{sec:mesure2}}
  {Deux qubits sont dans l'état $\ket{\Psi}=\frac{1}{\sqrt3}\big(\ket{00}+\ket{01}+\ket{11}\big)$. On mesure le \emph{second} qubit $B$ en base $Z$ et on obtient $1$. Que valait la probabilité de ce résultat, et dans quel état se trouve alors $A$ ?}
  {$P(B{=}1)=\frac23$ ; $A$ est dans l'état $\ket{+}$}
  {$P(B{=}1)=\frac13$ ; $A$ est dans l'état $\ket{0}$}
  {$P(B{=}1)=\frac23$ ; $A$ est dans l'état $\ket{1}$}
  {$P(B{=}1)=\frac13$ ; $A$ est dans l'état $\ket{+}$}
  {A}
\qcmx
  {$P(B{=}1)=|c_{01}|^2+|c_{11}|^2=\frac13+\frac13$. On garde les termes dont le second chiffre vaut $1$ : $\frac{1}{\sqrt3}\big(\ket{01}+\ket{11}\big)=\frac{1}{\sqrt3}(\ket0+\ket1)\otimes\ket1$, qui renormalisé donne $\ket{+}_A\otimes\ket{1}_B$.}
  {C'est le résultat $B=0$ : seul le terme $\ket{00}$ a un second chiffre nul, d'où la probabilité $\frac13$ et $A$ dans $\ket0$.}
  {On a projeté sur le seul terme $\ket{11}$ ; or $\ket{01}$ a aussi son second chiffre égal à $1$ : $A$ vaut $0$ ou $1$ (état $\ket+$).}
  {L'état de $A$ est juste, mais on n'a compté qu'un seul terme pour la probabilité ; il faut sommer $|c_{01}|^2+|c_{11}|^2$.}

% ---- Question 84 (niveau 3, correct : A)
\qcmq[2]{Chap.~\ref{chap:bell}\,$\cdot$\,\S\,\ref{sec:separable}}
  {Pour quelle valeur de $x$ l'état (non normalisé) $\ket{00}+2\ket{01}+3\ket{10}+x\ket{11}$ est-il séparable ?}
  {$x=6$}
  {$x=-6$}
  {$x=\frac16$}
  {$x=5$}
  {A}
\qcmx
  {Critère $D=c_{00}c_{11}-c_{01}c_{10}=x-2\cdot3=0$. On obtient bien un produit : $(\ket0+3\ket1)\otimes(\ket0+2\ket1)=\ket{00}+2\ket{01}+3\ket{10}+6\ket{11}$.}
  {Erreur de signe : $D=x-6$ s'annule pour $x=+6$ ; pour $x=-6$, $D=-12\neq0$ et l'état est intriqué.}
  {On a inversé le rapport : $c_{00}c_{11}=c_{01}c_{10}$ donne $x=\frac{c_{01}c_{10}}{c_{00}}=6$, et non $\frac{c_{00}}{c_{01}c_{10}}$.}
  {On a additionné ($2+3$) au lieu de multiplier : le critère de séparabilité porte sur des \emph{produits} de coefficients.}

% ---- Question 85 (niveau 3, correct : B)
\qcmq[2]{Chap.~\ref{chap:bell}\,$\cdot$\,\S\,\ref{sec:bell}}
  {Comment s'écrit $\ket{00}$ dans la base de Bell ?}
  {$\frac{1}{\sqrt2}\big(\ket{\Phi^+}-\ket{\Phi^-}\big)$}
  {$\frac{1}{\sqrt2}\big(\ket{\Phi^+}+\ket{\Phi^-}\big)$}
  {$\frac{1}{\sqrt2}\big(\ket{\Phi^+}+\ket{\Psi^+}\big)$}
  {$\ket{\Phi^+}$}
  {B}
\qcmx
  {Erreur de signe : cette combinaison vaut $\ket{11}$, puisque $\ket{\Phi^+}-\ket{\Phi^-}=\sqrt2\,\ket{11}$.}
  {$\ket{\Phi^+}+\ket{\Phi^-}=\frac{1}{\sqrt2}\cdot2\ket{00}=\sqrt2\,\ket{00}$. Une mesure de Bell sur $\ket{00}$ donne donc $\Phi^+$ ou $\Phi^-$ avec la probabilité $\frac12$ chacun (jamais $\Psi^\pm$).}
  {Cette combinaison vaut $\ket{+}\otimes\ket{+}=\frac12(\ket{00}+\ket{01}+\ket{10}+\ket{11})$ : les deux états de Bell « pairs » et « impairs » s'y ajoutent.}
  {On a confondu $\ket{00}$, état séparable, avec $\ket{\Phi^+}=\frac{1}{\sqrt2}(\ket{00}+\ket{11})$, qui contient aussi $\ket{11}$ et est intriqué.}

% ---- Question 86 (niveau 3, correct : A)
\qcmq[1]{Chap.~\ref{chap:bell}\,$\cdot$\,\S\,\ref{sec:circuit-bell}}
  {Le générateur de Bell du cours est $G=\mathrm{CNOT}\cdot(\mathrm{Id}\otimes H)$ : $H$ sur le second qubit $B$, puis CNOT de contrôle $B$ et de cible $A$. Quel circuit permet de \emph{mesurer} une paire dans la base de Bell (on le place avant une mesure en base $Z$) ?}
  {D'abord le CNOT (contrôle $B$, cible $A$), puis $H$ sur $B$}
  {D'abord $H$ sur $B$, puis le CNOT}
  {D'abord le CNOT, puis $H$ sur $A$ (la cible)}
  {Aucun circuit : on mesure directement les deux qubits en base $Z$}
  {A}
\qcmx
  {C'est l'inverse $G^\dagger=(\mathrm{Id}\otimes H)\,\mathrm{CNOT}$ : les portes sont leurs propres inverses et l'ordre s'inverse. Il envoie $\Phi^+,\Phi^-,\Psi^+,\Psi^-$ sur $\ket{00},\ket{01},\ket{10},\ket{11}$ (à une phase près).}
  {C'est $G$ lui-même, qui \emph{fabrique} des états de Bell : on garde le même ordre au lieu de l'inverser, et la mesure ne distingue plus les quatre états.}
  {Mauvais qubit : pour $\Phi^+$ on trouve $\ket{+}\otimes\ket{+}$, dont les deux bits sont aléatoires. Le $H$ doit agir sur $B$, qui a reçu le $H$ dans $G$.}
  {$\Phi^+$ et $\Phi^-$ (comme $\Psi^+$ et $\Psi^-$) donnent les mêmes résultats en base $Z$ : cette base ne voit pas le signe relatif, il faut d'abord défaire l'intrication.}

% ---- Question 87 (niveau 3, correct : B)
\qcmq[2]{Chap.~\ref{chap:bell}\,$\cdot$\,\S\,\ref{sec:trace}}
  {Quelle est la matrice de densité réduite $\rho_A=\operatorname{Tr}_B\ket\Psi\bra\Psi$ du premier qubit pour $\ket{\Psi}=\frac12\ket{00}+\frac{\sqrt3}{2}\ket{11}$ ?}
  {$\begin{pmatrix}\frac14&\frac{\sqrt3}{4}\\\frac{\sqrt3}{4}&\frac34\end{pmatrix}$}
  {$\begin{pmatrix}\frac14&0\\0&\frac34\end{pmatrix}$}
  {$\begin{pmatrix}\frac12&0\\0&\frac12\end{pmatrix}$}
  {$\begin{pmatrix}\frac12&0\\0&\frac{\sqrt3}{2}\end{pmatrix}$}
  {B}
\qcmx
  {C'est la matrice $\ket\chi\bra\chi$ de l'état pur $\ket\chi=\frac12\ket0+\frac{\sqrt3}2\ket1$ : on a gardé les cohérences, comme si $A$ était seul ; l'intrication avec $B$ les détruit.}
  {Avec $C=\begin{pmatrix}c_{00}&c_{01}\\c_{10}&c_{11}\end{pmatrix}=\begin{pmatrix}\frac12&0\\0&\frac{\sqrt3}2\end{pmatrix}$, $\rho_A=CC^\dagger=\mathrm{diag}\big(\frac14,\frac34\big)$ : état mixte, de pureté $\frac{10}{16}=1-2|D|^2$ avec $D=\frac{\sqrt3}{4}$.}
  {C'est la réponse pour un état de Bell (intrication maximale) ; ici les poids $\frac14$ et $\frac34$ diffèrent : l'intrication est partielle.}
  {On a mis les amplitudes au lieu de leurs carrés : la trace vaudrait $\frac{1+\sqrt3}{2}\neq1$.}

% ---- Question 88 (niveau 3, correct : A)
\qcmq[2]{Chap.~\ref{chap:bell}\,$\cdot$\,\S\,\ref{sec:info}}
  {Deux qubits sont dans le mélange statistique $\rho_{AB}=\frac12\ket{00}\bra{00}+\frac12\ket{11}\bra{11}$ (corrélations purement classiques). Que vaut $H(A|B)=S(\rho_{AB})-S(\rho_B)$, en bits ?}
  {$0$}
  {$-1$}
  {$1$}
  {$2$}
  {A}
\qcmx
  {$\rho_{AB}$ a pour valeurs propres $\frac12,\frac12,0,0$, donc $S(\rho_{AB})=1$ ; et $\rho_B=\frac12\mathrm{Id}$ donne $S(\rho_B)=1$. Ainsi $H(A|B)=0$ : connaître $B$ fixe $A$. Une valeur négative est propre à l'intrication.}
  {C'est la valeur pour l'état de Bell \emph{pur} $\ket{\Phi^+}$ : mêmes résultats en base $Z$, mais $S(\rho_{AB})=0$ car l'état global est pur. Ici l'état est mixte.}
  {On a pris $H(A|B)=H(A)$, comme si $B$ ne disait rien sur $A$ ; or les deux bits sont toujours égaux.}
  {On a additionné les entropies des deux qubits ($1+1$), valeur de deux bits indépendants : ce n'est pas une entropie conditionnelle.}

% ---- Question 89 (niveau 3, correct : B)
\qcmq[1]{Chap.~\ref{chap:bell}\,$\cdot$\,\S\,\ref{sec:teleportation}}
  {Dans la téléportation, avant de recevoir les deux bits $(a,q)$ d'Alice, dans quel état se trouve le qubit de Bob ?}
  {Dans l'état pur $\alpha\ket0+\beta\ket1$, déjà téléporté}
  {Dans l'état complètement mélangé $\rho_B=\frac12\mathrm{Id}$, indépendant de $\alpha$ et $\beta$}
  {Dans le mélange $|\alpha|^2\ket0\bra0+|\beta|^2\ket1\bra1$}
  {Dans l'état pur $\ket{0}$, son état initial}
  {B}
\qcmx
  {Ce serait une communication instantanée. L'état $\alpha\ket0+\beta\ket1$ n'apparaît qu'\emph{après} la correction, qui dépend de $(a,q)$.}
  {Les quatre kets $\ket{\phi_{aq}}$ sont équiprobables ($\frac14$) et $\frac14\sum_{a,q}\ket{\phi_{aq}}\bra{\phi_{aq}}=\frac12\mathrm{Id}$. Bob n'apprend rien avant le message classique : rien ne voyage plus vite que la lumière.}
  {Bob connaîtrait alors les statistiques de $\ket\varphi$ sans aucun message : or les probabilités d'Alice ($\frac14$ pour chaque résultat) ne dépendent pas de $\alpha,\beta$.}
  {Le qubit de Bob est la moitié de la paire $\ket{\Phi^+}$ : pris seul, il n'a pas d'état pur ($\rho_B=\frac12\mathrm{Id}$ avant comme après la mesure d'Alice, tant que Bob ignore le résultat).}

% ---- Question 90 (niveau 3, correct : D)
\qcmq[1]{Chap.~\ref{chap:bell}\,$\cdot$\,\S\,\ref{sec:chsh}}
  {Au jeu CHSH (gain si $a\oplus b=x\cdot y$), quelles sont les meilleures probabilités de gain avec une stratégie classique, même avec du hasard partagé, et avec une paire $\ket{\Phi^+}$ mesurée dans des bases bien choisies ?}
  {$\frac12$ en classique ; $1$ avec l'intrication}
  {$\frac34$ dans les deux cas}
  {$1$ en classique si Alice et Bob s'accordent à l'avance ; aucun avantage quantique}
  {$\frac34$ en classique ; $\cos^2\frac\pi8\approx0{,}85$ avec l'intrication}
  {D}
\qcmx
  {Classique : répondre toujours $0$ donne déjà $\frac34$. Quantique : l'intrication ne permet pas de gagner à tous les coups (borne de Tsirelson $\cos^2\frac\pi8$).}
  {Si l'intrication n'apportait rien, on n'aurait pas de violation de l'inégalité de Bell ; l'avantage quantique ($0{,}85>0{,}75$) est précisément ce que l'expérience de Bell met en évidence.}
  {S'accorder à l'avance (ou partager du hasard) ne suffit pas : aucune stratégie déterministe ne gagne les quatre parties, d'après l'argument de parité.}
  {Une stratégie déterministe perd au moins une des quatre parties (la somme modulo $2$ des conditions donne $0=1$) ; avec $\Phi^+$, chaque partie est gagnée avec la probabilité $\cos^2\frac\pi8=\frac{2+\sqrt2}{4}$.}

\clearpage
%% ======================================================================
\subsection{Niveau 4 : maîtrise (questions 91 à 100)}\label{sec:qcm-n4}
%% ======================================================================
\noindent{\color{insagris}\small Niveau de difficulté : Maîtrise\quad $\bullet\bullet\bullet\bullet$}\par\smallskip
% ---- Question 91 (niveau 4, correct : C)
\qcmq[2]{Chap.~\ref{chap:histoire}\,$\cdot$\,\S\,\ref{sec:superposition}}
  {Une particule du puits infini est préparée à $t=0$ dans l'état $\Psi=\frac{1}{\sqrt2}\big(\psi_1+\psi_2\big)$, avec $E_n=n^2E_1$. Au bout de quelle durée minimale la densité de probabilité $|\Psi(x,t)|^2$ retrouve-t-elle sa forme initiale ?}
  {$T=\dfrac{2\pi\hbar}{E_1}$}
  {$T=\dfrac{2\pi\hbar}{5E_1}$}
  {$T=\dfrac{2\pi\hbar}{3E_1}$}
  {$T=\dfrac{\pi\hbar}{3E_1}$}
  {C}
\qcmx
  {C'est la période de la phase $\ee^{-\ii E_1t/\hbar}$ de $\psi_1$ seul ; or les phases globales ne se voient pas dans $|\Psi|^2$ : seul compte l'écart $E_2-E_1=3E_1$.}
  {On a pris la somme $E_1+E_2=5E_1$ ; c'est la \emph{différence} des énergies (phase relative) qui intervient.}
  {$|\Psi|^2=\frac12(\psi_1^2+\psi_2^2)+\psi_1\psi_2\cos\frac{(E_2-E_1)t}{\hbar}$ : seul le terme croisé dépend du temps, à la pulsation $\omega=\frac{E_2-E_1}{\hbar}=\frac{3E_1}{\hbar}$, donc $T=\frac{2\pi}{\omega}$.}
  {C'est la demi-période : à $t=\frac T2$ le cosinus vaut $-1$ et la densité est le symétrique de l'état initial ($x\to a-x$), pas encore sa forme de départ.}

% ---- Question 92 (niveau 4, correct : D)
\qcmq[1]{Chap.~\ref{chap:dirac}\,$\cdot$\,\S\,\ref{sec:projecteurs}, \ref{sec:qubit}}
  {Un expérimentateur reçoit un qubit préparé soit dans $\ket{0}$, soit dans $\ket{+}$. Laquelle de ces affirmations est correcte ?}
  {Une mesure en base $Z$ les distingue : $\ket0$ donne $0$ et $\ket+$ donne $1$}
  {Une mesure en base $X$ les distingue : $\ket+$ donne $+$ et $\ket0$ donne $-$}
  {On les distingue en répétant la mesure sur le même qubit}
  {Aucune mesure ne permet de les distinguer à coup sûr, car $\braket{0}{+}=\frac{1}{\sqrt2}\neq0$}
  {D}
\qcmx
  {$\ket+$ donne $0$ ou $1$ avec la probabilité $\frac12$ : le résultat $0$ est possible pour les deux états.}
  {$\ket0=\frac{1}{\sqrt2}(\ket++\ket-)$ donne $+$ ou $-$ avec la probabilité $\frac12$ : le résultat $+$ est possible pour les deux états.}
  {Après une première mesure l'état est effondré : répéter ne donne aucune information nouvelle, et on ne peut pas copier le qubit pour recommencer (non-clonage).}
  {Seuls des états \emph{orthogonaux} sont parfaitement discernables. En base $Z$, $\ket0$ donne toujours $0$ et $\ket+$ donne $0$ une fois sur deux : obtenir $0$ ne tranche pas (obtenir $1$ prouve seulement qu'on avait $\ket+$).}

% ---- Question 93 (niveau 4, correct : C)
\qcmq[1]{Chap.~\ref{chap:operateurs}\,$\cdot$\,\S\,\ref{sec:pur-mixte}}
  {Une source envoie au hasard $\ket{0}$ ou $\ket{1}$ (probabilité $\frac12$ chacun) ; une autre envoie au hasard $\ket{+}$ ou $\ket{-}$ (probabilité $\frac12$ chacun). Peut-on distinguer ces deux sources en mesurant les qubits émis ?}
  {Oui : en base $Z$, la première source donne des résultats prévisibles}
  {Oui : en base $X$, la seconde source donne des résultats prévisibles}
  {Non : les deux ont la même matrice de densité $\frac{\mathrm{Id}}{2}$, donc les mêmes statistiques pour toute mesure}
  {Oui : les états purs envoyés ne sont pas les mêmes}
  {C}
\qcmx
  {Chaque qubit est bien $\ket0$ ou $\ket1$, mais on ignore lequel : le résultat vaut $0$ ou $1$ avec la probabilité $\frac12$, exactement comme avec la seconde source ($|\braket{0}{\pm}|^2=\frac12$).}
  {Même erreur : chaque qubit est $\ket+$ ou $\ket-$, mais on ignore lequel ; le résultat en base $X$ est aléatoire, et la première source donne aussi $+$ ou $-$ avec la probabilité $\frac12$.}
  {$\frac12\ket0\bra0+\frac12\ket1\bra1=\frac12\ket+\bra++\frac12\ket-\bra-=\frac{\mathrm{Id}}2$ (relation de fermeture dans chaque base). Les probabilités de toute mesure ne dépendent que de $\rho$.}
  {Ce qui compte n'est pas la liste des états purs envoyés mais la matrice de densité moyenne : des préparations différentes peuvent avoir la même $\rho$.}

% ---- Question 94 (niveau 4, correct : B)
\qcmq[1]{Chap.~\ref{chap:operateurs}\,$\cdot$\,\S\,\ref{sec:pur-mixte} ; Chap.~\ref{chap:portes}\,$\cdot$\,\S\,\ref{sec:portes-bloch}}
  {Un qubit est dans l'état complètement mélangé $\rho=\frac{\mathrm{Id}}{2}$. Existe-t-il une porte quantique $U$ qui le transforme en un état pur, par exemple $\ket{0}\bra{0}$ ?}
  {Oui : la porte de Hadamard, qui crée des superpositions}
  {Non : $\rho\mapsto U\rho U^\dagger$ conserve les valeurs propres, donc la pureté ; ici $U\frac{\mathrm{Id}}{2}U^\dagger=\frac{\mathrm{Id}}{2}$}
  {Oui : une rotation de Bloch qui amène le centre de la boule au pôle nord}
  {Oui : une mesure en base $Z$, qui donne $\ket0$}
  {B}
\qcmx
  {$H$ envoie un état pur sur un état pur ($\ket0\mapsto\ket+$) et laisse $\frac{\mathrm{Id}}{2}$ invariant : $H\frac{\mathrm{Id}}2H=\frac{\mathrm{Id}}2$ car $H^2=\mathrm{Id}$.}
  {$\operatorname{Tr}\big[(U\rho U^\dagger)^2\big]=\operatorname{Tr}\rho^2$ ($=\frac12$ ici). Sur la boule de Bloch, une porte est une rotation : elle laisse le centre fixe.}
  {Une rotation tourne la boule autour de son centre, qui reste fixe : elle ne peut pas l'amener sur la sphère.}
  {Une mesure n'est pas une porte (elle est non unitaire) ; de plus elle donne $\ket0$ ou $\ket1$ au hasard, pas $\ket0$ à coup sûr.}

% ---- Question 95 (niveau 4, correct : D)
\qcmq[1]{Chap.~\ref{chap:portes}\,$\cdot$\,\S\,\ref{sec:cnot}}
  {Que vaut $(H\otimes H)\,\mathrm{CNOT}_{12}\,(H\otimes H)$, où $\mathrm{CNOT}_{12}$ a pour contrôle le \emph{premier} qubit et pour cible le second ?}
  {$\mathrm{CNOT}_{12}$, inchangé}
  {$\mathrm{CZ}$ (porte $Z$ contrôlée)}
  {$\mathrm{SWAP}$}
  {$\mathrm{CNOT}_{21}$ (contrôle : second qubit ; cible : premier qubit)}
  {D}
\qcmx
  {On a cru que les deux $H$ s'annulent ($H^2=\mathrm{Id}$) ; mais ils encadrent le CNOT au lieu d'être côte à côte : on ne peut pas les simplifier.}
  {C'est le résultat quand $H$ n'encadre que la \emph{cible} : $(\mathrm{Id}\otimes H)\,\mathrm{CNOT}_{12}\,(\mathrm{Id}\otimes H)=\mathrm{CZ}$, car $HXH=Z$. Ici $H$ agit aussi sur le contrôle.}
  {Le SWAP s'obtient avec \emph{trois} CNOT alternés, pas avec un seul CNOT conjugué par $H\otimes H$ ; il n'y a pas de contrôle dans un SWAP.}
  {$H\ket0\bra0H=\ket+\bra+$, $H\ket1\bra1H=\ket-\bra-$ et $HXH=Z$ donnent $\ket+\bra+\otimes I+\ket-\bra-\otimes Z=I\otimes\ket0\bra0+X\otimes\ket1\bra1=\mathrm{CNOT}_{21}$ : dans la base $\{\ket+,\ket-\}$, contrôle et cible échangent leurs rôles (rebond de phase).}

% ---- Question 96 (niveau 4, correct : A)
\qcmq[1]{Chap.~\ref{chap:portes}\,$\cdot$\,\S\,\ref{sec:controlees}}
  {On applique la porte de Toffoli (contrôles : qubits 1 et 2 ; cible : qubit 3) à $\ket{+}\otimes\ket{+}\otimes\ket{0}$, puis on mesure le troisième qubit en base $Z$. Que peut-on dire si l'on obtient $1$ ?}
  {Ce résultat a la probabilité $\frac14$, et les deux premiers qubits sont alors dans l'état $\ket{11}$}
  {Ce résultat a la probabilité $\frac12$, et les deux premiers qubits restent dans $\ket{+}\ket{+}$}
  {Ce résultat a la probabilité $\frac14$, mais les deux premiers qubits restent dans $\ket{+}\ket{+}$}
  {Ce résultat a la probabilité $\frac34$, et les deux premiers qubits sont dans l'état $\frac{1}{\sqrt3}\big(\ket{00}+\ket{01}+\ket{10}\big)$}
  {A}
\qcmx
  {$\mathrm{CCNOT}\ket{x,y,0}=\ket{x,y,xy}$ donne $\frac12\big(\ket{000}+\ket{010}+\ket{100}+\ket{111}\big)$ : seul $\ket{111}$ a un troisième chiffre égal à $1$ (probabilité $\frac14$), et il impose $x=y=1$.}
  {La cible n'est inversée que si $x=y=1$ (probabilité $\frac14$) et non une fois sur deux ; de plus l'état est intriqué, donc mesurer le qubit 3 modifie les qubits 1 et 2.}
  {La probabilité est juste, mais la mesure de la cible n'est pas sans effet sur les contrôles : l'état après la Toffoli est intriqué, donc le résultat $1$ les projette sur $\ket{11}$.}
  {C'est le cas du résultat $0$ (probabilité $\frac34$ : $x,y\in\{00,01,10\}$). Le résultat $1$ est le cas rare, celui du ET vrai ($x=y=1$).}

% ---- Question 97 (niveau 4, correct : A)
\qcmq[1]{Chap.~\ref{chap:bell}\,$\cdot$\,\S\,\ref{sec:bell}}
  {Que donne $(H\otimes H)\ket{\Psi^-}$, avec $\ket{\Psi^-}=\frac{1}{\sqrt2}\big(\ket{01}-\ket{10}\big)$ ?}
  {$-\ket{\Psi^-}$, c'est-à-dire le même état physique}
  {$\ket{\Phi^-}$}
  {$\ket{\Psi^+}$}
  {$\ket{\Phi^+}$}
  {A}
\qcmx
  {Le singulet est invariant par $U\otimes U$ à un facteur près : $(U\otimes U)\ket{\Psi^-}=\det(U)\ket{\Psi^-}$ avec $\det H=-1$. Directement : $\frac{1}{\sqrt2}\big(\ket+\ket--\ket-\ket+\big)=-\ket{\Psi^-}$.}
  {C'est l'image de $\ket{\Psi^+}$, pas de $\ket{\Psi^-}$ : $(H\otimes H)\ket{\Psi^+}=\ket{\Phi^-}$. Le signe $-$ du singulet change le calcul (les termes $\ket{00}$ et $\ket{11}$ se compensent).}
  {$\ket{\Psi^+}$ est symétrique par échange des deux qubits, $\ket{\Psi^-}$ antisymétrique ; $H\otimes H$ commute avec le SWAP et conserve donc cette symétrie.}
  {C'est l'état que $H\otimes H$ laisse invariant ($\ket{\Phi^+}=\frac{1}{\sqrt2}(\ket{++}+\ket{--})$) ; le singulet n'a pas ce comportement.}

% ---- Question 98 (niveau 4, correct : B)
\qcmq[1]{Chap.~\ref{chap:bell}\,$\cdot$\,\S\,\ref{sec:chsh}}
  {Une expérience de type CHSH (avec $S=E_{00}+E_{01}+E_{10}-E_{11}$) mesure $S=2{,}6$. Que peut-on en conclure ?}
  {C'est impossible : on a toujours $|S|\leq2$}
  {Aucun modèle local à résultats prédéterminés ne l'explique, mais c'est compatible avec la mécanique quantique}
  {C'est impossible en mécanique quantique, car $2\sqrt2\approx2{,}4<2{,}6$}
  {Les particules s'échangent de l'information plus vite que la lumière}
  {B}
\qcmx
  {Cette borne ne vaut que pour les modèles classiques (locaux) ; l'expérience et la mécanique quantique la dépassent, jusqu'à $2\sqrt2$.}
  {Classique : $|S|\leq2$ ; quantique : $|S|\leq2\sqrt2\approx2{,}83$ (borne de Tsirelson). Comme $2<2{,}6<2\sqrt2$, l'inégalité de Bell est violée sans dépasser le maximum quantique ; le gain correspondant est $\frac12+\frac{2{,}6}{8}\approx0{,}825$.}
  {Erreur numérique : $2\sqrt2\approx2{,}83$, donc $2{,}6$ est permis.}
  {Les corrélations violent l'inégalité de Bell mais ne permettent aucun signal : la réponse de Bob est uniforme quelle que soit la base choisie par Alice.}

% ---- Question 99 (niveau 4, correct : A)
\qcmq[2]{Chap.~\ref{chap:bell}\,$\cdot$\,\S\,\ref{sec:info}}
  {Pour l'état $\ket{\Psi}=\frac{1}{\sqrt3}\big(\ket{00}+\sqrt2\,\ket{11}\big)$, quelle est l'entropie de von Neumann $S(\rho_A)$ du premier qubit, en bits ?}
  {$\log_2 3-\frac23\approx0{,}92$}
  {$1$}
  {$0$}
  {$0{,}64$}
  {A}
\qcmx
  {$\rho_A=\mathrm{diag}\big(\frac13,\frac23\big)$ (trace partielle), donc $S=-\frac13\log_2\frac13-\frac23\log_2\frac23=\log_23-\frac23\approx1{,}585-0{,}667$. C'est moins que le maximum d'un bit (état de Bell) : l'intrication est partielle.}
  {C'est la valeur pour un état de Bell ($\rho_A=\frac{\mathrm{Id}}{2}$) ; ici les poids $\frac13$ et $\frac23$ diffèrent, donc $S<1$.}
  {$S(\rho_A)=0$ vaut pour un état séparable ; ici $D=\frac{\sqrt2}{3}\neq0$ : l'état est intriqué, $\rho_A$ est mixte et $S>0$ (le couple $AB$ est pur, pas $A$ seul).}
  {On a utilisé le logarithme \emph{népérien} (résultat en nats) ; en bits il faut $\log_2$, ce qui donne $0{,}92$.}

% ---- Question 100 (niveau 4, correct : B)
\qcmq[1]{Chap.~\ref{chap:bell}\,$\cdot$\,\S\,\ref{sec:bell}}
  {Pourquoi peut-on mesurer simultanément les deux « parités » $Z\otimes Z$ et $X\otimes X$ de deux qubits, alors qu'on ne peut pas mesurer simultanément $Z$ et $X$ sur un seul qubit ?}
  {Parce que l'incertitude de Heisenberg ne s'applique pas à deux qubits}
  {Parce que $Z\otimes Z$ et $X\otimes X$ commutent : les deux signes $-1$ de $ZX=-XZ$ se compensent}
  {Parce que les deux qubits sont indépendants}
  {Parce que $Z\otimes Z$ et $X\otimes X$ agissent sur des qubits différents}
  {B}
\qcmx
  {Elle s'applique toujours : $Z\otimes\mathrm{Id}$ et $X\otimes\mathrm{Id}$ ne commutent pas. C'est la commutation de ces deux opérateurs précis qui permet la mesure simultanée.}
  {$(Z\otimes Z)(X\otimes X)=ZX\otimes ZX=(-XZ)\otimes(-XZ)=XZ\otimes XZ=(X\otimes X)(Z\otimes Z)$. Ils ont donc une base propre commune : la base de Bell.}
  {Dans un état de Bell les qubits sont intriqués, donc loin d'être indépendants ; la commutation est une propriété des opérateurs, pas de l'état.}
  {Chacun des deux opérateurs agit sur \emph{les deux} qubits à la fois. (Ce sont $Z\otimes\mathrm{Id}$ et $\mathrm{Id}\otimes X$, sur des qubits différents, qui commutent pour cette raison.)}

\label{qcm:fin}

\clearpage
\ifqcmcorrige

% ======================================================================
\subsection{Clé des réponses et fiche de score}\label{sec:qcm-cle}
% ======================================================================
\noindent\textbf{Cette partie donne les réponses : ne la lisez qu'après avoir répondu à toutes les
questions.}

\paragraph{Clé des réponses.} La lettre de la bonne réponse de la question $n$ se lit à la ligne
qui contient $n$, dans la colonne du rang de $n$ dans la ligne (question~24 : ligne 21--30,
colonne 4).
\begin{center}
\renewcommand{\arraystretch}{1.25}
\setlength{\tabcolsep}{5.5pt}
\begin{tabular}{@{}l*{10}{c}@{}}
\toprule
 & 1 & 2 & 3 & 4 & 5 & 6 & 7 & 8 & 9 & 10 \\
\midrule
1--10 & \qcmkey{1} & \qcmkey{2} & \qcmkey{3} & \qcmkey{4} & \qcmkey{5} & \qcmkey{6} & \qcmkey{7} & \qcmkey{8} & \qcmkey{9} & \qcmkey{10} \\
11--20 & \qcmkey{11} & \qcmkey{12} & \qcmkey{13} & \qcmkey{14} & \qcmkey{15} & \qcmkey{16} & \qcmkey{17} & \qcmkey{18} & \qcmkey{19} & \qcmkey{20} \\
21--30 & \qcmkey{21} & \qcmkey{22} & \qcmkey{23} & \qcmkey{24} & \qcmkey{25} & \qcmkey{26} & \qcmkey{27} & \qcmkey{28} & \qcmkey{29} & \qcmkey{30} \\
31--40 & \qcmkey{31} & \qcmkey{32} & \qcmkey{33} & \qcmkey{34} & \qcmkey{35} & \qcmkey{36} & \qcmkey{37} & \qcmkey{38} & \qcmkey{39} & \qcmkey{40} \\
41--50 & \qcmkey{41} & \qcmkey{42} & \qcmkey{43} & \qcmkey{44} & \qcmkey{45} & \qcmkey{46} & \qcmkey{47} & \qcmkey{48} & \qcmkey{49} & \qcmkey{50} \\
51--60 & \qcmkey{51} & \qcmkey{52} & \qcmkey{53} & \qcmkey{54} & \qcmkey{55} & \qcmkey{56} & \qcmkey{57} & \qcmkey{58} & \qcmkey{59} & \qcmkey{60} \\
61--70 & \qcmkey{61} & \qcmkey{62} & \qcmkey{63} & \qcmkey{64} & \qcmkey{65} & \qcmkey{66} & \qcmkey{67} & \qcmkey{68} & \qcmkey{69} & \qcmkey{70} \\
71--80 & \qcmkey{71} & \qcmkey{72} & \qcmkey{73} & \qcmkey{74} & \qcmkey{75} & \qcmkey{76} & \qcmkey{77} & \qcmkey{78} & \qcmkey{79} & \qcmkey{80} \\
81--90 & \qcmkey{81} & \qcmkey{82} & \qcmkey{83} & \qcmkey{84} & \qcmkey{85} & \qcmkey{86} & \qcmkey{87} & \qcmkey{88} & \qcmkey{89} & \qcmkey{90} \\
91--100 & \qcmkey{91} & \qcmkey{92} & \qcmkey{93} & \qcmkey{94} & \qcmkey{95} & \qcmkey{96} & \qcmkey{97} & \qcmkey{98} & \qcmkey{99} & \qcmkey{100} \\
\bottomrule
\end{tabular}
\end{center}

\paragraph{Fiche de score.} Reportez le nombre de bonnes réponses de chaque niveau.
\begin{center}
\renewcommand{\arraystretch}{1.6}
\begin{tabular}{@{}l l c c c@{}}
\toprule
\textbf{Niveau} & \textbf{Questions} & \textbf{Bonnes réponses} & \textbf{Sur} & \textbf{Pourcentage}\\
\midrule
1. Débutant & 1--30 & \makebox[2.6cm]{\dotfill} & 30 & \makebox[2cm]{\dotfill}\,\%\\
2. Intermédiaire & 31--60 & \makebox[2.6cm]{\dotfill} & 30 & \makebox[2cm]{\dotfill}\,\%\\
3. Avancé & 61--90 & \makebox[2.6cm]{\dotfill} & 30 & \makebox[2cm]{\dotfill}\,\%\\
4. Maîtrise & 91--100 & \makebox[2.6cm]{\dotfill} & 10 & \makebox[2cm]{\dotfill}\,\%\\
\midrule
\textbf{Total} & 1--100 & \makebox[2.6cm]{\dotfill} & 100 & \makebox[2cm]{\dotfill}\,\%\\
\bottomrule
\end{tabular}
\end{center}
Repères : au moins $80\,\%$ à un niveau, il est acquis et on peut passer au suivant ; entre $50$ et
$80\,\%$, revoyez les sections des questions manquées ; en dessous, relisez le chapitre concerné et
son récapitulatif avant de recommencer.

\paragraph{Où revoir ?} Questions classées par chapitre du cours.
\begin{center}
\renewcommand{\arraystretch}{1.3}
\begin{tabular}{@{}l >{\raggedright\arraybackslash}p{5.9cm} >{\raggedright\arraybackslash}p{6.9cm}@{}}
\toprule
\textbf{Chapitre} & \textbf{Titre} & \textbf{Questions}\\
\midrule
Chap.~\ref{chap:histoire} & Histoire et mécanique quantique & 1, 2, 3, 4, 5, 6, 31, 32, 33, 34, 35, 61, 62, 63, 91 \\
Chap.~\ref{chap:dirac} & Notation de Dirac et espace de Hilbert & 7, 8, 9, 10, 11, 12, 36, 37, 38, 39, 40, 64, 65, 66, 67, 68, 92 \\
Chap.~\ref{chap:operateurs} & Opérateurs, matrice de densité, Bloch, produit tensoriel & 13, 14, 15, 16, 17, 18, 41, 42, 43, 44, 45, 46, 69, 70, 71, 72, 73, 74, 75, 93, 94 \\
Chap.~\ref{chap:portes} & Portes quantiques et systèmes multi-qubits & 19, 20, 21, 22, 23, 24, 47, 48, 49, 50, 51, 76, 77, 78, 79, 80, 81, 82, 95, 96 \\
Chap.~\ref{chap:bell} & États quantiques, intrication, Bell & 25, 26, 27, 28, 29, 30, 52, 53, 54, 55, 56, 57, 58, 59, 60, 83, 84, 85, 86, 87, 88, 89, 90, 97, 98, 99, 100 \\
\bottomrule
\end{tabular}
\end{center}

\clearpage
\subsection{Corrigé du niveau 1 : débutant (questions 1 à 30)}\label{sec:qcm-c1}
\qcmcorr{1}
\qcmcorr{2}
\qcmcorr{3}
\qcmcorr{4}
\qcmcorr{5}
\qcmcorr{6}
\qcmcorr{7}
\qcmcorr{8}
\qcmcorr{9}
\qcmcorr{10}
\qcmcorr{11}
\qcmcorr{12}
\qcmcorr{13}
\qcmcorr{14}
\qcmcorr{15}
\qcmcorr{16}
\qcmcorr{17}
\qcmcorr{18}
\qcmcorr{19}
\qcmcorr{20}
\qcmcorr{21}
\qcmcorr{22}
\qcmcorr{23}
\qcmcorr{24}
\qcmcorr{25}
\qcmcorr{26}
\qcmcorr{27}
\qcmcorr{28}
\qcmcorr{29}
\qcmcorr{30}

\clearpage
\subsection{Corrigé du niveau 2 : intermédiaire (questions 31 à 60)}\label{sec:qcm-c2}
\qcmcorr{31}
\qcmcorr{32}
\qcmcorr{33}
\qcmcorr{34}
\qcmcorr{35}
\qcmcorr{36}
\qcmcorr{37}
\qcmcorr{38}
\qcmcorr{39}
\qcmcorr{40}
\qcmcorr{41}
\qcmcorr{42}
\qcmcorr{43}
\qcmcorr{44}
\qcmcorr{45}
\qcmcorr{46}
\qcmcorr{47}
\qcmcorr{48}
\qcmcorr{49}
\qcmcorr{50}
\qcmcorr{51}
\qcmcorr{52}
\qcmcorr{53}
\qcmcorr{54}
\qcmcorr{55}
\qcmcorr{56}
\qcmcorr{57}
\qcmcorr{58}
\qcmcorr{59}
\qcmcorr{60}

\clearpage
\subsection{Corrigé du niveau 3 : avancé (questions 61 à 90)}\label{sec:qcm-c3}
\qcmcorr{61}
\qcmcorr{62}
\qcmcorr{63}
\qcmcorr{64}
\qcmcorr{65}
\qcmcorr{66}
\qcmcorr{67}
\qcmcorr{68}
\qcmcorr{69}
\qcmcorr{70}
\qcmcorr{71}
\qcmcorr{72}
\qcmcorr{73}
\qcmcorr{74}
\qcmcorr{75}
\qcmcorr{76}
\qcmcorr{77}
\qcmcorr{78}
\qcmcorr{79}
\qcmcorr{80}
\qcmcorr{81}
\qcmcorr{82}
\qcmcorr{83}
\qcmcorr{84}
\qcmcorr{85}
\qcmcorr{86}
\qcmcorr{87}
\qcmcorr{88}
\qcmcorr{89}
\qcmcorr{90}

\clearpage
\subsection{Corrigé du niveau 4 : maîtrise (questions 91 à 100)}\label{sec:qcm-c4}
\qcmcorr{91}
\qcmcorr{92}
\qcmcorr{93}
\qcmcorr{94}
\qcmcorr{95}
\qcmcorr{96}
\qcmcorr{97}
\qcmcorr{98}
\qcmcorr{99}
\qcmcorr{100}
\fi
) :
%    \qcmrefsfalse     masque la référence « Chap. · § » sous chaque énoncé ;
%    \qcmcorrigefalse  supprime le corrigé du document (QCM à distribuer seul).
%
%  Ce fichier est produit par outils_qcm/gen_qcm.py à partir de la banque de
%  questions (fichiers texte), mais il reste un fichier LaTeX ordinaire que l'on
%  peut modifier à la main. Pour ajouter une question à la main : insérer un
%  couple \qcmq ... \qcmx dans le bon niveau (la numérotation est automatique),
%  puis mettre à jour (1) les bornes « questions a à b » des titres de niveaux,
%  (2) les lignes \qcmcorr{n} du corrigé, (3) le tableau « Où revoir ? ».
% =====================================================================

% Macros de Dirac (déjà définies dans main.tex ; ce bloc évite l'erreur
% « Undefined control sequence \ket » si main.tex est une ancienne version).
\providecommand{\ket}[1]{|#1\rangle}
\providecommand{\bra}[1]{\langle #1|}
\providecommand{\braket}[2]{\langle #1|#2\rangle}

\makeatletter
% bascules : normalement déjà définies dans main.tex (valeurs par défaut : vrai)
\@ifundefined{ifqcmrefs}{\newif\ifqcmrefs\qcmrefstrue}{}
\@ifundefined{ifqcmcorrige}{\newif\ifqcmcorrige\qcmcorrigetrue}{}
\newcounter{qcmq}

% case à cocher (3 mm) pour le crayon
\newcommand{\qcmbox}{{\setlength{\fboxsep}{0pt}\setlength{\fboxrule}{0.5pt}%
  \raisebox{-0.4mm}{\fbox{\rule{0pt}{3mm}\rule{3mm}{0pt}}}}}

% une proposition : case, lettre, texte (retrait suspendu)
\newcommand{\qcm@opt}[2]{\noindent\hangindent=3.2em \hangafter=1
  \hspace*{0.3em}\qcmbox\hspace{0.45em}\textbf{#1)}\hspace{0.4em}#2\par\vspace{2.5pt}}

% \qcmq[colonnes]{ref}{énoncé}{A}{B}{C}{D}{bonne lettre}   (TeX limite les macros à 9 paramètres)
% \qcmx{com. A}{com. B}{com. C}{com. D}  : commentaires, placés juste après \qcmq
%   - com. de la bonne lettre = explication ; com. des autres = piège ;
%   - colonnes = 1 (une proposition par ligne) ou 2 (grille 2x2).
\newcommand{\qcmq}[8][1]{%
  \stepcounter{qcmq}%
  \expandafter\gdef\csname qcm@ref@\theqcmq\endcsname{#2}%
  \expandafter\gdef\csname qcm@rep@\theqcmq\endcsname{#8}%
  \expandafter\gdef\csname qcm@o@\theqcmq @A\endcsname{#4}%
  \expandafter\gdef\csname qcm@o@\theqcmq @B\endcsname{#5}%
  \expandafter\gdef\csname qcm@o@\theqcmq @C\endcsname{#6}%
  \expandafter\gdef\csname qcm@o@\theqcmq @D\endcsname{#7}%
  \par\vspace{7pt plus 3pt minus 2pt}\noindent
  \begin{minipage}{\linewidth}%
    \noindent{\color{insarouge}\bfseries Question~\theqcmq}%
    \ifqcmrefs\hfill{\footnotesize\color{insagris}#2}\fi\par\nopagebreak
    \vspace{2pt}\noindent #3\par\nopagebreak\vspace{3pt}%
    \ifnum#1=2
      \begin{minipage}[t]{0.495\linewidth}\qcm@opt{A}{#4}\end{minipage}\hfill
      \begin{minipage}[t]{0.495\linewidth}\qcm@opt{B}{#5}\end{minipage}\par\vspace{2pt}
      \begin{minipage}[t]{0.495\linewidth}\qcm@opt{C}{#6}\end{minipage}\hfill
      \begin{minipage}[t]{0.495\linewidth}\qcm@opt{D}{#7}\end{minipage}\par
    \else
      \qcm@opt{A}{#4}\qcm@opt{B}{#5}\qcm@opt{C}{#6}\qcm@opt{D}{#7}%
    \fi
  \end{minipage}\par}

\newcommand{\qcmx}[4]{%
  \expandafter\gdef\csname qcm@c@\theqcmq @A\endcsname{#1}%
  \expandafter\gdef\csname qcm@c@\theqcmq @B\endcsname{#2}%
  \expandafter\gdef\csname qcm@c@\theqcmq @C\endcsname{#3}%
  \expandafter\gdef\csname qcm@c@\theqcmq @D\endcsname{#4}}

% lettre de la bonne réponse de la question n (« ? » si la question n'existe pas)
\newcommand{\qcmkey}[1]{\@ifundefined{qcm@rep@#1}{?}{\csname qcm@rep@#1\endcsname}}

% une mauvaise réponse du corrigé (n° de question, lettre) : ignorée si c'est la bonne.
% (pas de \@for : « := » ne fonctionne pas avec le « : » rendu actif par babel-french)
\newcommand{\qcm@wrongone}[2]{%
  \edef\qcm@a{#2}\edef\qcm@b{\csname qcm@rep@#1\endcsname}%
  \ifx\qcm@a\qcm@b\else
    \noindent\hangindent=1.4em\hangafter=1 {\color{insagris}\textbf{#2)}}\enspace
    \csname qcm@o@#1@#2\endcsname\enspace--\enspace\emph{\csname qcm@c@#1@#2\endcsname}\par
  \fi}

% bloc de corrigé de la question n
\newcommand{\qcmcorr}[1]{%
  \par\vspace{6pt plus 2pt minus 2pt}\noindent
  \begin{minipage}{\linewidth}\small
    \noindent{\color{insarouge}\bfseries Question~#1}\enspace
    \textbf{Réponse : \csname qcm@rep@#1\endcsname}\hfill
    {\footnotesize\color{insagris}\csname qcm@ref@#1\endcsname}\par\nopagebreak
    \vspace{1pt}\noindent\hangindent=1.4em\hangafter=1
    \textbf{\csname qcm@rep@#1\endcsname)}\enspace
    \csname qcm@o@#1@\csname qcm@rep@#1\endcsname\endcsname\enspace---\enspace
    \csname qcm@c@#1@\csname qcm@rep@#1\endcsname\endcsname\par
    \qcm@wrongone{#1}{A}\qcm@wrongone{#1}{B}\qcm@wrongone{#1}{C}\qcm@wrongone{#1}{D}%
  \end{minipage}\par}
\makeatother

% =====================================================================
\section{QCM sur le cours}\label{chap:qcm}
% =====================================================================

\subsection{Mode d'emploi}\label{sec:qcm-mode}

Ce QCM compte \textbf{100 questions} sur les chapitres~\ref{chap:histoire} à~\ref{chap:bell},
classées en quatre niveaux. Chaque question propose quatre réponses dont \textbf{une seule} est
exacte : cochez au crayon la case de la proposition choisie.
\ifqcmrefs À droite du numéro de chaque question, une référence (chapitre, section) indique la
partie du cours concernée.\else\ifqcmcorrige{} La référence de chaque question (chapitre, section)
est donnée dans le corrigé.\fi\fi

\begin{center}
\renewcommand{\arraystretch}{1.3}
\begin{tabular}{@{}llc>{\raggedright\arraybackslash}p{8.2cm}@{}}
\toprule
\textbf{Niveau} & \textbf{Questions} & \textbf{Nombre} & \textbf{Ce qui est demandé}\\
\midrule
1. Débutant & 1 à 30 & 30 & définitions, notations, résultats du cours en une étape\\
2. Intermédiaire & 31 à 60 & 30 & petits calculs, propriétés, application directe d'un résultat\\
3. Avancé & 61 à 90 & 30 & calculs en plusieurs étapes, portes, circuits et protocoles\\
4. Maîtrise & 91 à 100 & 10 & raisonnements fins et pièges conceptuels, qui combinent plusieurs notions\\
\bottomrule
\end{tabular}
\end{center}

\begin{itemize}[nosep]
  \item \textbf{Sur papier} : imprimez seulement les pages~\pageref{chap:qcm} à~\pageref{qcm:fin}
        (mode d'emploi et énoncés, sections~\ref{sec:qcm-n1} à~\ref{sec:qcm-n4}).%
        \ifqcmcorrige{} Le corrigé commence sur une page distincte, à la page~\pageref{sec:qcm-cle}.\fi
  \item \textbf{Barème} : 1 point par bonne réponse, 0 sinon (pas de point négatif). Calculatrice
        autorisée, mais rarement utile (quelques logarithmes et racines).
  \item \textbf{Correction} :%
        \ifqcmcorrige{} comparez avec la clé des réponses, puis lisez l'explication de chaque
        question, même réussie. Les « pièges » décrivent l'erreur qui mène à chaque mauvaise réponse.%
        \else{} le corrigé n'est pas inclus dans cette version du document.\fi
\end{itemize}

\bigskip
\noindent Nom : \makebox[6cm]{\dotfill}\hfill Date : \makebox[3.5cm]{\dotfill}\hfill Score : \makebox[1.6cm]{\dotfill}\,/\,100
\bigskip

\begin{remarque}[Conventions des énoncés]
\begin{itemize}[nosep]
  \item $\ket{ab}=\ket{a}_A\otimes\ket{b}_B$ : le premier chiffre est le qubit $A$, le second le
        qubit $B$ ; dans un circuit, le fil du haut est le dernier chiffre.
  \item Le CNOT de la section~\ref{sec:cnot} a pour contrôle le \emph{second} chiffre ; les portes
        contrôlées de la section~\ref{sec:controlees} ont pour contrôle le \emph{premier}. Chaque
        énoncé précise le contrôle.
  \item Sans précision, une mesure est faite dans la base de calcul $Z$. Les logarithmes des
        entropies sont en base $2$ (résultats en bits).
  \item $\ket{\pm}=\frac{1}{\sqrt2}(\ket0\pm\ket1)$ et $\ket{{\pm}\ii}=\frac{1}{\sqrt2}(\ket0\pm\ii\ket1)$.
\end{itemize}
\end{remarque}

\clearpage
%% ======================================================================
\subsection{Niveau 1 : débutant (questions 1 à 30)}\label{sec:qcm-n1}
%% ======================================================================
\noindent{\color{insagris}\small Niveau de difficulté : Débutant\quad $\bullet\circ\circ\circ$}\par\smallskip
% ---- Question 1 (niveau 1, correct : D)
\qcmq[2]{Chap.~\ref{chap:histoire}\,$\cdot$\,\S\,\ref{sec:dualite}}
  {En 1905, qui explique l'effet photo-électrique en postulant que la lumière est faite de quanta d'énergie $E=hf$, les photons ?}
  {Louis de Broglie}
  {Erwin Schrödinger}
  {Max Planck}
  {Albert Einstein}
  {D}
\qcmx
  {Il propose en 1924 une onde de matière ($\lambda=h/p$) pour les particules comme l'électron ; l'idée du photon date de 1905.}
  {Il écrit en 1926 l'équation d'onde de la matière, plus de vingt ans après l'article sur les photons.}
  {Il introduit le quantum d'énergie en 1900 pour le rayonnement du corps noir ; c'est Einstein qui en fait une propriété de la lumière elle-même.}
  {Il postule que la lumière échange son énergie par paquets $E=hf=\hbar\omega$ : c'est la première face de la dualité onde--corpuscule.}

% ---- Question 2 (niveau 1, correct : A)
\qcmq[1]{Chap.~\ref{chap:histoire}\,$\cdot$\,\S\,\ref{sec:born}}
  {Selon l'interprétation de Max Born, que représente $|\Psi(x,t)|^2$ ?}
  {La densité de probabilité de présence de la particule en $x$ à l'instant $t$}
  {L'énergie cinétique de la particule}
  {L'amplitude de l'onde de matière}
  {La vitesse de la particule}
  {A}
\qcmx
  {$|\Psi(x,t)|^2\,\dd x$ est la probabilité de trouver la particule entre $x$ et $x+\dd x$ ; $\Psi$ lui-même n'est pas mesurable.}
  {L'énergie est associée à l'opérateur hamiltonien $\Hop$ ; $|\Psi|^2$ n'a pas la dimension d'une énergie.}
  {$\Psi$ est l'amplitude (complexe) ; $|\Psi|^2$ en est le carré du module.}
  {La vitesse s'obtient avec l'opérateur $\hat p=-\ii\hbar\,\partial_x$, pas avec $|\Psi|^2$.}

% ---- Question 3 (niveau 1, correct : D)
\qcmq[2]{Chap.~\ref{chap:histoire}\,$\cdot$\,\S\,\ref{sec:dualite}}
  {Une particule de quantité de mouvement $p$ est associée, d'après de Broglie, à une onde de longueur d'onde $\lambda$. Quelle relation est correcte ?}
  {$\lambda=hp$}
  {$\lambda=p/h$}
  {$\lambda=\hbar/p$}
  {$\lambda=h/p$}
  {D}
\qcmx
  {Produit au lieu du quotient : $\lambda$ grandirait avec $p$, le contraire de ce qu'on observe.}
  {Quotient inversé : $p/h=1/\lambda$ s'exprime en m$^{-1}$, pas en mètres.}
  {Confusion entre $h$ et $\hbar=h/2\pi$ : $\hbar/p=\lambda/2\pi$ est l'inverse du nombre d'onde $k$.}
  {Relation de de Broglie : plus la particule est rapide ou lourde, plus $\lambda$ est petite (balle de $0{,}1$~kg à $10$~m/s : $\lambda\approx10^{-34}$~m, inobservable).}

% ---- Question 4 (niveau 1, correct : C)
\qcmq[2]{Chap.~\ref{chap:histoire}\,$\cdot$\,\S\,\ref{sec:construction}}
  {Quelle est l'équation de Schrödinger dépendante du temps pour un état $\Psi$ et un hamiltonien $\Hop$ ?}
  {$-\ii\hbar\,\partial_t\Psi=\Hop\Psi$}
  {$\hbar\,\partial_t\Psi=\Hop\Psi$}
  {$\ii\hbar\,\partial_t\Psi=\Hop\Psi$}
  {$\ii\hbar\,\partial_x\Psi=\Hop\Psi$}
  {C}
\qcmx
  {Erreur de signe : la solution serait $\ee^{+\ii\Hop t/\hbar}\Psi(0)$, qui décrit l'évolution à rebours du temps.}
  {Sans le facteur $\ii$, les solutions sont des exponentielles réelles (croissantes ou décroissantes) et la norme n'est pas conservée.}
  {Elle relie l'évolution de l'état à l'énergie ; pour un $\Hop$ indépendant du temps, sa solution est $\Psi(t)=\ee^{-\ii\Hop t/\hbar}\Psi(0)$ (chapitre~\ref{chap:portes}).}
  {C'est la dérivée en $t$ qui décrit l'évolution ; la dérivée en $x$ sert à construire $\hat p=-\ii\hbar\,\partial_x$.}

% ---- Question 5 (niveau 1, correct : A)
\qcmq[1]{Chap.~\ref{chap:histoire}\,$\cdot$\,\S\,\ref{sec:construction}}
  {Quel opérateur représente l'énergie totale d'un système quantique ?}
  {L'hamiltonien $\Hop$}
  {L'opérateur quantité de mouvement $\hat p$}
  {L'opérateur position $\hat x$}
  {L'opérateur identité}
  {A}
\qcmx
  {Somme de l'énergie cinétique $-\frac{\hbar^2}{2m}\partial_x^2$ et de l'énergie potentielle $V$ ; ses valeurs propres sont les énergies mesurables.}
  {Il représente l'impulsion ; l'énergie cinétique est $\hat p^2/2m$, et il faut y ajouter $V$.}
  {C'est la multiplication par $x$ : il donne la position, pas l'énergie.}
  {Il vaut $1$ pour tout état : il ne contient aucune information sur l'énergie.}

% ---- Question 6 (niveau 1, correct : C)
\qcmq[1]{Chap.~\ref{chap:histoire}\,$\cdot$\,\S\,\ref{sec:histoire}}
  {Dans quel ordre chronologique ces avancées ont-elles eu lieu ?}
  {de Broglie $\to$ Planck $\to$ Schrödinger $\to$ Bell}
  {Planck $\to$ Schrödinger $\to$ de Broglie $\to$ Bell}
  {Planck (quantum d'énergie) $\to$ de Broglie (onde de matière) $\to$ Schrödinger (équation d'onde) $\to$ Bell (inégalités)}
  {Planck $\to$ de Broglie $\to$ Bell $\to$ Schrödinger}
  {C}
\qcmx
  {Le quantum d'énergie de Planck (1900) précède l'onde de matière de de Broglie (1924).}
  {L'équation de Schrödinger (1926) généralise l'idée d'onde de matière de de Broglie (1924) : elle vient après.}
  {Planck 1900, de Broglie 1924, Schrödinger 1926, Bell 1964.}
  {Les inégalités de Bell (1964) viennent près de quarante ans après l'équation de Schrödinger.}

% ---- Question 7 (niveau 1, correct : C)
\qcmq[1]{Chap.~\ref{chap:dirac}\,$\cdot$\,\S\,\ref{sec:dirac}}
  {En notation de Dirac, comment appelle-t-on $\ket{\Psi}$ et comment le représente-t-on pour un qubit ?}
  {Un bra, représenté par un vecteur ligne à deux composantes complexes}
  {Un ket, représenté par un nombre complexe}
  {Un ket, représenté par un vecteur colonne à deux composantes complexes}
  {Un bra, représenté par une matrice $2\times2$}
  {C}
\qcmx
  {Le bra $\bra{\Psi}$ est le conjugué transposé du ket : un vecteur ligne, mais ce n'est pas la notation $\ket{\Psi}$.}
  {Un scalaire ne peut pas décrire une superposition : il faut deux amplitudes $\alpha$ et $\beta$.}
  {Le ket est le vecteur colonne $\binom{\alpha}{\beta}$ ; son bra $\bra{\Psi}$ est le vecteur ligne conjugué.}
  {Une matrice est un opérateur ; le bra est une matrice ligne $1\times2$.}

% ---- Question 8 (niveau 1, correct : B)
\qcmq[2]{Chap.~\ref{chap:dirac}\,$\cdot$\,\S\,\ref{sec:qubit}}
  {Quelle condition les amplitudes $\alpha$ et $\beta$ de $\alpha\ket{0}+\beta\ket{1}$ doivent-elles vérifier pour que l'état soit normalisé ?}
  {$\alpha+\beta=1$}
  {$|\alpha|^2+|\beta|^2=1$}
  {$\alpha^2+\beta^2=1$}
  {$|\alpha|+|\beta|=1$}
  {B}
\qcmx
  {On somme les amplitudes (complexes, parfois négatives) ; ce sont leurs modules au carré qui sont des probabilités.}
  {Les probabilités de mesure $P(0)=|\alpha|^2$ et $P(1)=|\beta|^2$ doivent sommer à $1$.}
  {Oubli du module : pour $\alpha=\frac{1}{\sqrt2}$ et $\beta=\frac{\ii}{\sqrt2}$, $\alpha^2+\beta^2=0$ alors que l'état est normalisé.}
  {Somme des modules au lieu des carrés : $\alpha=\beta=\frac{1}{\sqrt2}$ est normalisé, mais $|\alpha|+|\beta|=\sqrt2$.}

% ---- Question 9 (niveau 1, correct : C)
\qcmq[1]{Chap.~\ref{chap:dirac}\,$\cdot$\,\S\,\ref{sec:projecteurs}}
  {On mesure dans la base $Z$ un qubit préparé dans $\ket{+}=\frac{1}{\sqrt2}(\ket{0}+\ket{1})$. Que se passe-t-il ?}
  {On lit $+$ et le qubit reste dans $\ket{+}$}
  {On lit toujours $0$}
  {On lit $0$ ou $1$, chacun avec la probabilité $\frac12$, et l'état devient $\ket{0}$ ou $\ket{1}$}
  {On lit $0$ et $1$ à la fois}
  {C}
\qcmx
  {Une mesure en base $Z$ ne donne que $0$ ou $1$ ; elle modifie l'état.}
  {Les deux amplitudes ont le même module : $P(0)=P(1)=\frac12$.}
  {Postulat de la mesure : la superposition s'effondre sur l'état propre correspondant au résultat.}
  {Chaque mesure donne un seul résultat ; la superposition ne s'observe qu'à travers les statistiques.}

% ---- Question 10 (niveau 1, correct : B)
\qcmq[2]{Chap.~\ref{chap:dirac}\,$\cdot$\,\S\,\ref{sec:qubit}}
  {Que vaut le produit scalaire $\braket{0}{1}$ ?}
  {$1$}
  {$0$}
  {$\frac{1}{\sqrt2}$}
  {$\ket{01}$}
  {B}
\qcmx
  {C'est la valeur de $\braket{0}{0}$ et de $\braket{1}{1}$ (vecteurs normés), pas de $\braket{0}{1}$.}
  {$\ket{0}$ et $\ket{1}$ sont orthogonaux : $\begin{pmatrix}1&0\end{pmatrix}\begin{pmatrix}0\\1\end{pmatrix}=0$.}
  {C'est la valeur de $\braket{0}{+}$, pas de $\braket{0}{1}$.}
  {Confusion avec le produit tensoriel : un produit scalaire est un nombre, pas un ket.}

% ---- Question 11 (niveau 1, correct : D)
\qcmq[2]{Chap.~\ref{chap:dirac}\,$\cdot$\,\S\,\ref{sec:qubit}}
  {Un qubit est dans l'état $\ket{-}=\frac{1}{\sqrt2}(\ket{0}-\ket{1})$. Quelle est la probabilité de mesurer $1$ en base $Z$ ?}
  {$0$}
  {$-\frac12$}
  {$\frac{1}{\sqrt2}$}
  {$\frac12$}
  {D}
\qcmx
  {Le signe $-$ ne « retire » pas $\ket{1}$ de la superposition.}
  {Une probabilité est toujours positive : on prend le module au carré de l'amplitude.}
  {C'est l'amplitude en module, pas la probabilité (il faut élever au carré).}
  {$P(1)=\left|-\frac{1}{\sqrt2}\right|^2=\frac12$ : le signe est une phase relative, qui disparaît dans le module.}

% ---- Question 12 (niveau 1, correct : B)
\qcmq[2]{Chap.~\ref{chap:dirac}\,$\cdot$\,\S\,\ref{sec:projecteurs}}
  {Quel est le projecteur sur l'état $\ket{0}$ ?}
  {$\ket{0}\bra{1}=\begin{pmatrix}0&1\\0&0\end{pmatrix}$}
  {$\hat P_0=\ket{0}\bra{0}=\begin{pmatrix}1&0\\0&0\end{pmatrix}$}
  {$\braket{0}{0}=1$}
  {$\begin{pmatrix}0&0\\0&1\end{pmatrix}$}
  {B}
\qcmx
  {Cet opérateur envoie $\ket{1}$ sur $\ket{0}$ ; son carré est nul : ce n'est pas un projecteur.}
  {$\hat P_0\ket{\Psi}=\alpha\ket{0}$ : il garde la composante sur $\ket{0}$ ; $\hat P_0^2=\hat P_0$.}
  {C'est un nombre, pas un opérateur.}
  {C'est $\ket{1}\bra{1}$, le projecteur sur $\ket{1}$.}

% ---- Question 13 (niveau 1, correct : A)
\qcmq[1]{Chap.~\ref{chap:operateurs}\,$\cdot$\,\S\,\ref{sec:bloch}}
  {Où se trouvent les états $\ket{0}$ et $\ket{1}$ sur la sphère de Bloch ?}
  {Aux deux pôles, Nord et Sud}
  {Aux deux extrémités de l'axe $x$}
  {Aux deux extrémités de l'axe $y$}
  {En deux points voisins de l'équateur}
  {A}
\qcmx
  {$\ket{0}$ au pôle Nord ($\theta=0$), $\ket{1}$ au pôle Sud ($\theta=\pi$) : deux états orthogonaux sont diamétralement opposés.}
  {Ce sont les places de $\ket{+}$ et $\ket{-}$.}
  {Ce sont les places de $\ket{{+}\ii}$ et $\ket{{-}\ii}$.}
  {Des états orthogonaux sont antipodaux, jamais voisins.}

% ---- Question 14 (niveau 1, correct : A)
\qcmq[2]{Chap.~\ref{chap:operateurs}\,$\cdot$\,\S\,\ref{sec:tenseur}}
  {Quelle est la dimension de l'espace de Hilbert de $n$ qubits ?}
  {$2^n$}
  {$2n$}
  {$n^2$}
  {$2^n-1$}
  {A}
\qcmx
  {Le produit tensoriel multiplie les dimensions : $2\times2\times\cdots\times2$ ; deux qubits : $4$, trois qubits : $8$.}
  {On additionne les dimensions au lieu de les multiplier ; pour $n=3$ on trouverait $6$ au lieu de $8$.}
  {Juste pour $n=2$ par coïncidence seulement : pour $n=3$ on trouverait $9$ au lieu de $8$.}
  {Il y a bien $2^n$ vecteurs dans la base de calcul, $\ket{0\cdots0}$ compris.}

% ---- Question 15 (niveau 1, correct : C)
\qcmq[1]{Chap.~\ref{chap:operateurs}\,$\cdot$\,\S\,\ref{sec:rho}}
  {Que vaut toujours la trace d'une matrice de densité $\rho$ ?}
  {$0$}
  {$1$ pour un état pur, et moins de $1$ pour un état mixte}
  {$1$}
  {$2$, la dimension de l'espace}
  {C}
\qcmx
  {C'est la trace des matrices de Pauli, pas d'une matrice de densité.}
  {Confusion avec $\operatorname{Tr}(\rho^2)$ : c'est la pureté qui distingue les deux cas, la trace vaut toujours $1$.}
  {C'est la somme des probabilités de mesure (les populations), que l'état soit pur ou mixte.}
  {C'est la trace de l'identité $\operatorname{Tr}\mathrm{Id}=2$, pas celle de $\rho$.}

% ---- Question 16 (niveau 1, correct : B)
\qcmq[1]{Chap.~\ref{chap:operateurs}\,$\cdot$\,\S\,\ref{sec:pur-mixte}}
  {Quelle est la valeur de $\operatorname{Tr}(\rho^2)$ pour un état pur ?}
  {$0$}
  {$1$}
  {$\frac12$}
  {Une valeur strictement inférieure à $1$}
  {B}
\qcmx
  {$\rho^2=0$ imposerait $\rho=0$, ce qui contredit $\operatorname{Tr}\rho=1$.}
  {Un état pur s'écrit $\rho=\ket{\Psi}\bra{\Psi}$, donc $\rho^2=\rho$ et $\operatorname{Tr}\rho^2=\operatorname{Tr}\rho=1$.}
  {C'est la valeur minimale pour un qubit, atteinte par l'état maximalement mixte $\mathrm{Id}/2$.}
  {C'est la caractérisation d'un état mixte.}

% ---- Question 17 (niveau 1, correct : B)
\qcmq[1]{Chap.~\ref{chap:operateurs}\,$\cdot$\,\S\,\ref{sec:rho}}
  {Que représentent les éléments diagonaux $\rho_{00}$ et $\rho_{11}$ de la matrice de densité d'un qubit ?}
  {Les cohérences entre $\ket{0}$ et $\ket{1}$}
  {Les probabilités de mesurer $\ket{0}$ et $\ket{1}$ (les populations)}
  {Les valeurs propres de $\rho$}
  {Les amplitudes $\alpha$ et $\beta$}
  {B}
\qcmx
  {Les cohérences sont les éléments hors diagonale $\rho_{01}$ et $\rho_{10}$.}
  {$\rho_{00}=\bra{0}\rho\ket{0}=P(0)$ et $\rho_{11}=P(1)$, d'où $\operatorname{Tr}\rho=1$.}
  {Seulement si $\rho$ est diagonale dans la base de calcul : $\ket{+}\bra{+}$ a pour diagonale $(\frac12,\frac12)$ mais pour valeurs propres $(1,0)$.}
  {Pour un état pur, $\rho_{00}=|\alpha|^2$ : le module au carré, pas l'amplitude.}

% ---- Question 18 (niveau 1, correct : D)
\qcmq[1]{Chap.~\ref{chap:operateurs}\,$\cdot$\,\S\,\ref{sec:op-dirac}}
  {Pourquoi associe-t-on aux grandeurs mesurables (observables) des opérateurs hermitiens ?}
  {Parce qu'ils conservent la norme des vecteurs}
  {Parce qu'ils ont tous une trace nulle}
  {Parce que leurs valeurs propres valent toutes $\pm1$}
  {Parce que leurs valeurs propres sont réelles, comme les résultats d'une mesure}
  {D}
\qcmx
  {C'est la propriété des opérateurs unitaires (les portes), pas des opérateurs hermitiens.}
  {Vrai pour les matrices de Pauli, faux en général (l'identité est hermitienne, de trace $2$), et sans rapport avec le caractère mesurable.}
  {Vrai pour $X$, $Y$, $Z$ seulement ; l'hamiltonien du puits a les valeurs propres $E_n=n^2E_1$.}
  {Un opérateur hermitien ($A^\dagger=A$) a des valeurs propres réelles et des vecteurs propres orthogonaux ; les résultats possibles sont ses valeurs propres.}

% ---- Question 19 (niveau 1, correct : C)
\qcmq[2]{Chap.~\ref{chap:portes}\,$\cdot$\,\S\,\ref{sec:pauli}}
  {Quelle porte de Pauli réalise l'inversion de bit $\ket{0}\leftrightarrow\ket{1}$ ?}
  {$Z$}
  {$Y$}
  {$X$}
  {$H$}
  {C}
\qcmx
  {$Z$ laisse $\ket{0}$ et change $\ket{1}$ en $-\ket{1}$ : c'est une inversion de phase.}
  {$Y\ket{0}=\ii\ket{1}$ : elle inverse le bit mais ajoute une phase ; la porte d'inversion « pure » est $X$.}
  {$X\ket{0}=\ket{1}$ et $X\ket{1}=\ket{0}$ : c'est le NOT quantique.}
  {$H$ crée des superpositions : $H\ket{0}=\ket{+}$.}

% ---- Question 20 (niveau 1, correct : B)
\qcmq[2]{Chap.~\ref{chap:portes}\,$\cdot$\,\S\,\ref{sec:hadamard}}
  {Que donne la porte de Hadamard appliquée à $\ket{0}$ ?}
  {$\ket{1}$}
  {$\ket{+}=\frac{1}{\sqrt2}(\ket{0}+\ket{1})$}
  {$\ket{-}=\frac{1}{\sqrt2}(\ket{0}-\ket{1})$}
  {$\ket{0}$}
  {B}
\qcmx
  {C'est l'action de $X$.}
  {Superposition à poids égaux : $P(0)=P(1)=\frac12$ ; c'est la première étape de la plupart des algorithmes.}
  {C'est $H\ket{1}$ ; le signe $-$ n'apparaît qu'à partir de $\ket{1}$.}
  {$H$ n'est pas l'identité ; elle ne laisse inchangés ni $\ket{0}$ ni $\ket{1}$ (elle envoie $\ket{+}$ sur $\ket{0}$).}

% ---- Question 21 (niveau 1, correct : D)
\qcmq[2]{Chap.~\ref{chap:portes}\,$\cdot$\,\S\,\ref{sec:portes}}
  {Quelle propriété doit vérifier la matrice $U$ d'une porte quantique ?}
  {Elle est hermitienne : $U^\dagger=U$}
  {Sa trace vaut $1$}
  {Ses coefficients sont réels}
  {Elle est unitaire : $U^\dagger U=\mathrm{Id}$}
  {D}
\qcmx
  {Vrai pour $X$, $Y$, $Z$, $H$ mais faux pour $S$ et $T$ : ce n'est pas la condition générale.}
  {C'est une propriété des matrices de densité ; $\operatorname{Tr}X=0$ et $\operatorname{Tr}\mathrm{Id}=2$.}
  {$Y$ et $S$ sont des portes à coefficients complexes.}
  {Elle conserve la norme ($\braket{\Psi}{\Psi}=1$ reste vrai) et elle est réversible.}

% ---- Question 22 (niveau 1, correct : A)
\qcmq[2]{Chap.~\ref{chap:portes}\,$\cdot$\,\S\,\ref{sec:portes}}
  {Quel est l'inverse d'une porte quantique $U$ ?}
  {$U^{-1}=U^\dagger$ (transposée conjuguée)}
  {$U^{-1}=U^{\mathsf T}$ (transposée seule)}
  {$U^{-1}=U^{*}$ (conjuguée seule)}
  {$U^{-1}=-U$}
  {A}
\qcmx
  {Conséquence de $U^\dagger U=\mathrm{Id}$ ; par exemple $S^{-1}=S^\dagger=\mathrm{diag}(1,-\ii)$.}
  {Oubli de la conjugaison : $S^{\mathsf T}=S=\mathrm{diag}(1,\ii)$ alors que $S^{-1}=\mathrm{diag}(1,-\ii)$.}
  {Faux en général : $Y^*=-Y$ alors que $Y^{-1}=Y$.}
  {Faux : $X^{-1}=X\neq-X$.}

% ---- Question 23 (niveau 1, correct : C)
\qcmq[1]{Chap.~\ref{chap:portes}\,$\cdot$\,\S\,\ref{sec:cnot}}
  {Dans une porte CNOT, que devient le qubit cible lorsque le qubit de contrôle est dans l'état $\ket{0}$ ?}
  {Il est inversé par la porte $X$}
  {Il est remis à $\ket{0}$}
  {Il reste inchangé}
  {Il subit un changement de signe, comme avec $Z$}
  {C}
\qcmx
  {C'est ce qui se passe quand le contrôle vaut $\ket{1}$.}
  {Le CNOT est réversible : remettre la cible à zéro détruirait de l'information.}
  {Avec le contrôle à $\ket{0}$, la porte agit comme l'identité sur la cible.}
  {Ce serait la porte $Z$ contrôlée (CZ), pas le CNOT.}

% ---- Question 24 (niveau 1, correct : A)
\qcmq[2]{Chap.~\ref{chap:portes}\,$\cdot$\,\S\,\ref{sec:cnot}}
  {Le CNOT du cours a pour contrôle le second chiffre du ket et pour cible le premier : $\mathrm{CNOT}\ket{a,b}=\ket{a\oplus b,\,b}$. Que donne-t-il sur $\ket{0,1}$ ?}
  {$\ket{1,1}$}
  {$\ket{0,1}$}
  {$\ket{0,0}$}
  {$\ket{1,0}$}
  {A}
\qcmx
  {Le contrôle $b=1$ inverse la cible : $0\oplus1=1$.}
  {Le contrôle vaut $1$ : la cible est bien inversée.}
  {C'est le contrôle qui serait inversé : le contrôle n'est jamais modifié.}
  {Les deux qubits seraient inversés : seule la cible change.}

% ---- Question 25 (niveau 1, correct : D)
\qcmq[1]{Chap.~\ref{chap:bell}\,$\cdot$\,\S\,\ref{sec:bell}}
  {Quelle propriété caractérise les quatre états de Bell ?}
  {Ils sont séparables : chaque qubit a son propre état}
  {Ils décrivent un seul qubit}
  {Ils ne diffèrent que par une phase globale}
  {Ils sont maximalement intriqués et forment une base orthonormée de l'espace de deux qubits}
  {D}
\qcmx
  {$D=c_{00}c_{11}-c_{01}c_{10}=\pm\frac12\neq0$ : ils ne se factorisent pas.}
  {Ce sont des états de deux qubits (quatre amplitudes sur $\ket{00},\ket{01},\ket{10},\ket{11}$).}
  {Ils sont orthogonaux deux à deux : ce sont quatre états distincts, alors qu'une phase globale ne change pas l'état.}
  {Aucun n'est un produit d'états à un qubit ($|D|=\frac12$, valeur maximale) et ils sont orthogonaux deux à deux.}

% ---- Question 26 (niveau 1, correct : D)
\qcmq[1]{Chap.~\ref{chap:bell}\,$\cdot$\,\S\,\ref{sec:teleportation}}
  {Quel est l'objectif du protocole de téléportation quantique ?}
  {Copier l'état d'un qubit inconnu sur le qubit de Bob}
  {Transmettre une information plus vite que la lumière}
  {Déplacer la particule physique d'Alice vers Bob}
  {Transférer l'état d'un qubit inconnu d'Alice à Bob, avec une paire intriquée et deux bits classiques}
  {D}
\qcmx
  {La copie est interdite (non-clonage) : après le protocole, l'état n'existe plus que chez Bob.}
  {Sans les deux bits classiques, le qubit de Bob est dans l'état $\mathrm{Id}/2$ : il n'a reçu aucune information.}
  {C'est l'état quantique qui est transféré, pas la particule qui le porte.}
  {Le qubit lui-même n'est pas envoyé : son état est reconstruit chez Bob, et détruit chez Alice par la mesure.}

% ---- Question 27 (niveau 1, correct : B)
\qcmq[1]{Chap.~\ref{chap:bell}\,$\cdot$\,\S\,\ref{sec:teleportation}}
  {Dans la téléportation, que doit envoyer Alice à Bob sur un canal classique pour qu'il retrouve l'état ?}
  {Un qubit intriqué}
  {Deux bits classiques, le résultat de sa mesure}
  {Les amplitudes exactes $\alpha$ et $\beta$ de l'état}
  {Un seul bit}
  {B}
\qcmx
  {La paire intriquée est partagée avant le protocole ; on ne l'envoie pas après.}
  {Les quatre résultats $(a,q)$ correspondent aux quatre corrections possibles ($I$, $Z$, $X$, $X$ puis $Z$) : il faut deux bits.}
  {Alice ne connaît pas l'état à téléporter, et ces nombres complexes n'ont pas de description finie en bits.}
  {Un bit ne distingue que deux cas, alors que Bob doit choisir parmi quatre corrections.}

% ---- Question 28 (niveau 1, correct : A)
\qcmq[1]{Chap.~\ref{chap:bell}\,$\cdot$\,\S\,\ref{sec:bell} ; ex.\,\ref{ex:superdense}}
  {Quel est le but du codage superdense ?}
  {Transmettre deux bits classiques en envoyant un seul qubit, grâce à une paire intriquée partagée}
  {Transmettre l'état d'un qubit avec deux bits classiques}
  {Compresser plusieurs qubits en un seul}
  {Chiffrer un message de façon inviolable}
  {A}
\qcmx
  {Alice applique l'une des quatre portes $I$, $Z$, $X$, $XZ$ à son qubit de la paire, l'envoie, et Bob identifie l'état de Bell obtenu.}
  {C'est la téléportation, le protocole inverse.}
  {Le protocole ne compresse pas de qubits : il fait passer deux bits classiques par un canal à un qubit.}
  {Ce n'est pas un protocole de cryptographie : la paire intriquée sert à doubler la capacité du canal.}

% ---- Question 29 (niveau 1, correct : A)
\qcmq[1]{Chap.~\ref{chap:bell}\,$\cdot$\,\S\,\ref{sec:teleportation}}
  {Que dit le théorème de non-clonage ?}
  {Aucune opération ne peut copier parfaitement un état quantique inconnu quelconque}
  {Il est impossible de mesurer un qubit sans le perturber}
  {Il est impossible de créer de l'intrication}
  {Il est impossible de copier un bit classique}
  {A}
\qcmx
  {Une porte $U$ qui copierait $\ket{0}$ et $\ket{1}$ enverrait $\ket{+}\ket{0}$ sur un état intriqué, par linéarité, et non sur $\ket{+}\ket{+}$.}
  {C'est l'effondrement de l'état, un autre principe.}
  {L'intrication se crée très bien (Hadamard puis CNOT).}
  {Le CNOT copie un bit classique : $\ket{0,b}\to\ket{b,b}$.}

% ---- Question 30 (niveau 1, correct : B)
\qcmq[1]{Chap.~\ref{chap:bell}\,$\cdot$\,\S\,\ref{sec:circuit-bell}}
  {Quel enchaînement de portes transforme $\ket{00}$ en l'état de Bell $\ket{\Phi^+}=\frac{1}{\sqrt2}(\ket{00}+\ket{11})$ ?}
  {CNOT, puis Hadamard sur le contrôle}
  {Hadamard sur l'un des qubits, puis CNOT dont le contrôle est ce même qubit}
  {Hadamard sur chacun des deux qubits}
  {SWAP, puis $X$ sur un qubit}
  {B}
\qcmx
  {Le CNOT ne fait rien sur $\ket{00}$ (contrôle à $0$) ; $H$ donne ensuite un état séparable.}
  {$H$ crée $\ket{0}\otimes\ket{+}$, puis le CNOT recopie la superposition : $\frac{1}{\sqrt2}(\ket{00}+\ket{11})$.}
  {On obtient $\ket{+}\otimes\ket{+}$, séparable : des portes locales ne créent pas d'intrication.}
  {$\ket{00}$ devient $\ket{01}$ ou $\ket{10}$, un état de base séparable.}

\clearpage
%% ======================================================================
\subsection{Niveau 2 : intermédiaire (questions 31 à 60)}\label{sec:qcm-n2}
%% ======================================================================
\noindent{\color{insagris}\small Niveau de difficulté : Intermédiaire\quad $\bullet\bullet\circ\circ$}\par\smallskip
% ---- Question 31 (niveau 2, correct : A)
\qcmq[1]{Chap.~\ref{chap:histoire}\,$\cdot$\,\S\,\ref{sec:born}}
  {Quelle condition la fonction d'onde d'une particule doit-elle vérifier sur tout l'espace ?}
  {$\displaystyle\int_{-\infty}^{+\infty}|\Psi(x,t)|^2\,\dd x=1$}
  {$\displaystyle\int_{-\infty}^{+\infty}\Psi(x,t)\,\dd x=1$}
  {$|\Psi(x,t)|^2=1$ pour tout $x$}
  {$\partial_x\Psi(x,t)=0$ en tout point}
  {A}
\qcmx
  {La particule se trouve forcément quelque part : la probabilité totale vaut $1$ (condition de normalisation).}
  {Sans le module au carré : $\Psi$ est complexe et son intégrale n'est pas une probabilité (elle peut être nulle ou complexe).}
  {$|\Psi|^2$ est une densité : si elle valait $1$ partout, son intégrale divergerait.}
  {Une fonction constante n'est pas normalisable ($\int|\Psi|^2\dd x$ diverge) ; la condition porte sur l'intégrale de $|\Psi|^2$.}

% ---- Question 32 (niveau 2, correct : D)
\qcmq[2]{Chap.~\ref{chap:histoire}\,$\cdot$\,\S\,\ref{sec:operateurs}}
  {Quelle est l'expression de l'opérateur quantité de mouvement $\hat p$ en une dimension ?}
  {$\hat p=m\,v$}
  {$\hat p=-\dfrac{\hbar^2}{2m}\dfrac{\partial^2}{\partial x^2}$}
  {$\hat p=+\ii\hbar\,\dfrac{\partial}{\partial x}$}
  {$\hat p=-\ii\hbar\,\dfrac{\partial}{\partial x}$}
  {D}
\qcmx
  {C'est la grandeur classique ; en mécanique quantique $p$ est un opérateur différentiel.}
  {C'est l'énergie cinétique $\hat p^2/2m$.}
  {Erreur de signe : on trouverait $p=-\hbar k$ pour l'onde $\ee^{\ii kx}$.}
  {Convention du cours : $\hat p\,\ee^{\ii kx}=\hbar k\,\ee^{\ii kx}$, donc une onde vers les $x$ croissants a $p=\hbar k>0$.}

% ---- Question 33 (niveau 2, correct : B)
\qcmq[2]{Chap.~\ref{chap:histoire}\,$\cdot$\,\S\,\ref{sec:solutions}}
  {Pour une particule dans un puits de potentiel infini de largeur $a$, comment varie l'énergie $E_n$ avec le numéro $n$ du niveau ?}
  {$E_n=E_1/n$}
  {$E_n=n^2E_1$, avec $E_1=\dfrac{\pi^2\hbar^2}{2ma^2}$}
  {$E_n=nE_1$}
  {$E_n=E_1$ pour tout $n$}
  {B}
\qcmx
  {Les niveaux se resserreraient quand $n$ augmente ; c'est l'inverse pour le puits infini.}
  {$k_n=\frac{n\pi}{a}$ et $E_n=\frac{\hbar^2k_n^2}{2m}$ : les niveaux s'écartent quand $n$ croît ($E_{n+1}-E_n=(2n+1)E_1$).}
  {Niveaux équidistants : ce n'est pas le puits infini, où $E_2=4E_1$ et $E_3=9E_1$.}
  {Il n'y aurait pas de quantification : tous les états auraient la même énergie.}

% ---- Question 34 (niveau 2, correct : C)
\qcmq[2]{Chap.~\ref{chap:histoire}\,$\cdot$\,\S\,\ref{sec:solutions}}
  {Une particule est dans un état stationnaire $\Psi_n$ de l'hamiltonien ($\Hop\Psi_n=E_n\Psi_n$). Que vaut l'écart-type $\sigma_E$ de son énergie ?}
  {$E_n$}
  {$E_n^2$}
  {$0$}
  {$\sqrt{E_n}$}
  {C}
\qcmx
  {C'est la valeur moyenne de l'énergie, pas son écart-type.}
  {C'est la valeur de $\langle\Hop^2\rangle$, avant soustraction de $\langle\Hop\rangle^2$.}
  {$\langle\Hop\rangle=E_n$ et $\langle\Hop^2\rangle=E_n^2$, donc $\sigma_E=\sqrt{E_n^2-E_n^2}=0$ : l'énergie est parfaitement déterminée.}
  {Mélange avec la formule $\sigma=\sqrt{\langle\Hop^2\rangle-\langle\Hop\rangle^2}$ : les deux termes sont égaux à $E_n^2$.}

% ---- Question 35 (niveau 2, correct : C)
\qcmq[2]{Chap.~\ref{chap:histoire}\,$\cdot$\,\S\,\ref{sec:solutions}}
  {Pour un état stationnaire $\Psi_n$ du puits infini $]0,a[$, que vaut la position moyenne $\langle\hat x\rangle$ ?}
  {$a/n$}
  {$a/4$}
  {$a/2$ pour tout $n$}
  {$0$}
  {C}
\qcmx
  {Elle dépendrait de $n$, ce que la symétrie de $|\Psi_n|^2$ interdit.}
  {C'est la position d'un maximum de $|\Psi_2|^2$ ; la moyenne des deux maxima $\frac a4$ et $\frac{3a}{4}$ donne $\frac a2$.}
  {$|\Psi_n|^2=\frac2a\sin^2\frac{n\pi x}{a}$ est symétrique par rapport au centre du puits.}
  {$\Psi_n$ s'annule sur la paroi, mais $\langle\hat x\rangle$ est une moyenne pondérée sur tout le puits.}

% ---- Question 36 (niveau 2, correct : D)
\qcmq[2]{Chap.~\ref{chap:dirac}\,$\cdot$\,\S\,\ref{sec:qubit}}
  {Un qubit est dans l'état $\ket{\Psi}=\frac{1}{\sqrt5}\ket{0}+\frac{2\ii}{\sqrt5}\ket{1}$. Quelle est la probabilité de mesurer $1$ en base $Z$ ?}
  {$\frac25$}
  {$-\frac45$}
  {$\frac15$}
  {$\frac45$}
  {D}
\qcmx
  {On divise le coefficient $2$ par $5$ : on oublie à la fois la racine et le carré.}
  {On élève $\frac{2\ii}{\sqrt5}$ au carré sans conjuguer : $\left(\frac{2\ii}{\sqrt5}\right)^2=-\frac45$ ; une probabilité est $|\cdot|^2$.}
  {C'est $P(0)=\left|\frac{1}{\sqrt5}\right|^2$.}
  {$P(1)=\left|\frac{2\ii}{\sqrt5}\right|^2=\frac45$ : le facteur $\ii$ est une phase relative, invisible en base $Z$.}

% ---- Question 37 (niveau 2, correct : C)
\qcmq[2]{Chap.~\ref{chap:dirac}\,$\cdot$\,\S\,\ref{sec:projecteurs}}
  {Un qubit préparé dans $\ket{0}$ est mesuré dans la base $X$ $\{\ket{+},\ket{-}\}$. Quelle est la probabilité d'obtenir $-$ ?}
  {$0$}
  {$1$}
  {$\frac12$}
  {$\frac{1}{\sqrt2}$}
  {C}
\qcmx
  {$\ket{0}$ est certain en base $Z$, mais pas en base $X$ : une autre base donne une autre loi.}
  {Même confusion dans l'autre sens : $\ket{0}$ n'est pas l'état $\ket{-}$.}
  {$P(-)=|\braket{-}{0}|^2=\left|\frac{1}{\sqrt2}\right|^2=\frac12$.}
  {C'est l'amplitude $\braket{-}{0}$ ; il faut élever son module au carré.}

% ---- Question 38 (niveau 2, correct : C)
\qcmq[2]{Chap.~\ref{chap:dirac}\,$\cdot$\,\S\,\ref{sec:qubit}}
  {Quel facteur $N>0$ normalise l'état $N\big(\ket{0}+(1-\ii)\ket{1}\big)$ ?}
  {$N=\frac{1}{\sqrt2}$}
  {$N=\frac13$}
  {$N=\frac{1}{\sqrt3}$}
  {$N=\frac{1}{\sqrt{1-2\ii}}$}
  {C}
\qcmx
  {On oublie la composante $\ket{0}$ : $|1-\ii|^2=2$, mais la norme au carré est $1+2=3$.}
  {On oublie la racine carrée : $3N^2=1$ donne $N=\frac{1}{\sqrt3}$, pas $\frac13$.}
  {$\braket{\Psi}{\Psi}=N^2\big(1+|1-\ii|^2\big)=N^2(1+2)=3N^2$.}
  {On élève $1-\ii$ au carré sans prendre le module : $(1-\ii)^2=-2\ii$ n'est pas réel ; il faut $|1-\ii|^2=2$.}

% ---- Question 39 (niveau 2, correct : D)
\qcmq[2]{Chap.~\ref{chap:dirac}\,$\cdot$\,\S\,\ref{sec:hilbert}}
  {Soient $\ket{\psi}=\frac{1}{\sqrt2}(\ket{0}+\ii\ket{1})$ et $\ket{\varphi}=\frac{1}{\sqrt2}(\ket{0}-\ket{1})$. Que vaut $\braket{\psi}{\varphi}$ ?}
  {$\frac{1-\ii}{2}$}
  {$\frac12$}
  {$0$}
  {$\frac{1+\ii}{2}$}
  {D}
\qcmx
  {On oublie de conjuguer l'amplitude $\ii$ du bra : on calcule $\frac12\big(1+\ii\cdot(-1)\big)$.}
  {Seul le terme en $\ket{0}$ est compté ; le terme en $\ket{1}$ vaut $\frac12(-\ii)(-1)=\frac{\ii}{2}$, non nul.}
  {Confusion avec $\braket{+}{-}=0$ ; ici $|\braket{\psi}{\varphi}|^2=\frac12$.}
  {Le bra conjugue les amplitudes : $\bra{\psi}=\frac{1}{\sqrt2}(\bra{0}-\ii\bra{1})$, d'où $\frac12\big(1+(-\ii)(-1)\big)=\frac{1+\ii}{2}$.}

% ---- Question 40 (niveau 2, correct : A)
\qcmq[2]{Chap.~\ref{chap:dirac}\,$\cdot$\,\S\,\ref{sec:qubit}}
  {Lequel de ces états est orthogonal à $\ket{{+}\ii}=\frac{1}{\sqrt2}(\ket{0}+\ii\ket{1})$ ?}
  {$\ket{{-}\ii}=\frac{1}{\sqrt2}(\ket{0}-\ii\ket{1})$}
  {$\ket{-}=\frac{1}{\sqrt2}(\ket{0}-\ket{1})$}
  {$\ket{+}=\frac{1}{\sqrt2}(\ket{0}+\ket{1})$}
  {$\ket{1}$}
  {A}
\qcmx
  {$\braket{{+}\ii}{{-}\ii}=\frac12\big(1+(-\ii)(-\ii)\big)=\frac12(1-1)=0$, en conjuguant le bra.}
  {$\braket{{+}\ii}{-}=\frac{1+\ii}{2}\neq0$ : $\ket{-}$ est orthogonal à $\ket{+}$, pas à $\ket{{+}\ii}$.}
  {$\braket{{+}\ii}{+}=\frac{1-\ii}{2}\neq0$.}
  {$\braket{{+}\ii}{1}=\frac{-\ii}{\sqrt2}\neq0$ : $\ket{1}$ est orthogonal à $\ket{0}$ seulement.}

% ---- Question 41 (niveau 2, correct : D)
\qcmq[1]{Chap.~\ref{chap:operateurs}\,$\cdot$\,\S\,\ref{sec:pur-mixte}}
  {Une matrice de densité de qubit a pour pureté $\operatorname{Tr}(\rho^2)=\frac58$. Que peut-on en conclure ?}
  {L'état est pur}
  {La matrice n'est pas une matrice de densité valide}
  {L'état est maximalement mixte, $\rho=\mathrm{Id}/2$}
  {L'état est mixte (mélange statistique)}
  {D}
\qcmx
  {Il faudrait $\operatorname{Tr}\rho^2=1$.}
  {$\frac58$ est compris entre $\frac12$ et $1$ : valeur admissible (exemple du cours : $\mathrm{diag}(\frac34,\frac14)$).}
  {Ce cas donnerait $\operatorname{Tr}\rho^2=\frac12$.}
  {$\operatorname{Tr}\rho^2<1$ ; pour un qubit, $\frac12\leq\operatorname{Tr}\rho^2\leq1$ et $\frac58$ est une valeur admissible.}

% ---- Question 42 (niveau 2, correct : B)
\qcmq[2]{Chap.~\ref{chap:operateurs}\,$\cdot$\,\S\,\ref{sec:rho}}
  {Que vaut $\langle Z\rangle=\operatorname{Tr}(\rho Z)$ pour un qubit dans l'état $\ket{+}$ ?}
  {$1$}
  {$0$}
  {$-1$}
  {$\frac{1}{\sqrt2}$}
  {B}
\qcmx
  {C'est la valeur pour $\ket{0}$.}
  {$\rho=\frac12\begin{pmatrix}1&1\\1&1\end{pmatrix}$ et $\operatorname{Tr}(\rho Z)=\frac12-\frac12=0$ : les résultats $+1$ et $-1$ sont équiprobables.}
  {C'est la valeur pour $\ket{1}$.}
  {C'est l'amplitude de $\ket{0}$ dans $\ket{+}$, pas une valeur moyenne de $Z$.}

% ---- Question 43 (niveau 2, correct : A)
\qcmq[1]{Chap.~\ref{chap:operateurs}\,$\cdot$\,\S\,\ref{sec:bloch}}
  {Sur la boule de Bloch, que représente un point situé à l'intérieur de la sphère ($|\vec r|<1$) ?}
  {Un état mixte, de pureté $\operatorname{Tr}\rho^2=\frac{1+|\vec r|^2}{2}<1$}
  {Un état pur dont la norme n'est pas $1$}
  {Un état pur en superposition}
  {Un point sans signification physique}
  {A}
\qcmx
  {Les états purs sont sur la sphère ; le centre $\vec r=\vec0$ est l'état maximalement mixte $\mathrm{Id}/2$.}
  {Tout état pur est normalisé et situé sur la sphère.}
  {Les états purs, superpositions comprises, sont tous \emph{sur} la sphère.}
  {Tout point de la boule est une matrice de densité valide.}

% ---- Question 44 (niveau 2, correct : B)
\qcmq[1]{Chap.~\ref{chap:operateurs}\,$\cdot$\,\S\,\ref{sec:pur-mixte}}
  {Par quoi la matrice de densité distingue-t-elle l'état pur $\ket{+}$ du mélange équiprobable de $\ket{0}$ et $\ket{1}$ ?}
  {Par leur trace, égale à $1$ pour l'état pur seulement}
  {Par les éléments hors diagonale (cohérences), non nuls pour $\ket{+}$ et nuls pour le mélange}
  {Par leurs éléments diagonaux, qui sont différents}
  {Ils ne se distinguent pas : c'est la même matrice}
  {B}
\qcmx
  {La trace vaut $1$ pour toute matrice de densité.}
  {$\ket{+}\bra{+}=\frac12\begin{pmatrix}1&1\\1&1\end{pmatrix}$ et le mélange vaut $\frac12\mathrm{Id}$ ; une mesure en base $X$ les distingue.}
  {Les deux ont la diagonale $(\frac12,\frac12)$ : $P(0)=P(1)=\frac12$ dans les deux cas.}
  {Les matrices sont différentes : en base $X$, $\ket{+}$ donne $+$ avec certitude, le mélange donne $\pm$ à pile ou face.}

% ---- Question 45 (niveau 2, correct : C)
\qcmq[2]{Chap.~\ref{chap:operateurs}\,$\cdot$\,\S\,\ref{sec:rho}}
  {Un qubit est préparé dans $\ket{0}$ avec la probabilité $\frac34$ et dans $\ket{1}$ avec la probabilité $\frac14$. Que vaut $\langle Z\rangle=\operatorname{Tr}(\rho Z)$ ?}
  {$\frac34$}
  {$1$}
  {$\frac12$}
  {$0$}
  {C}
\qcmx
  {C'est la probabilité $\rho_{00}$ de lire $0$ ; or $Z$ vaut $+1$ pour $0$ et $-1$ pour $1$.}
  {Valeur pour $\rho=\ket{0}\bra{0}$ (préparation certaine de $\ket{0}$).}
  {$\rho=\mathrm{diag}(\frac34,\frac14)$, donc $\operatorname{Tr}(\rho Z)=\frac34\cdot1+\frac14\cdot(-1)=\frac12$ : c'est la coordonnée $r_z$ de la sphère de Bloch.}
  {Valeur pour le mélange équiprobable $\mathrm{Id}/2$.}

% ---- Question 46 (niveau 2, correct : B)
\qcmq[2]{Chap.~\ref{chap:operateurs}\,$\cdot$\,\S\,\ref{sec:tenseur}}
  {Dans la base ordonnée $\{\ket{00},\ket{01},\ket{10},\ket{11}\}$, quel vecteur colonne représente $\ket{1}\otimes\ket{0}=\ket{10}$ ?}
  {$(0,1,0,0)^{\mathsf T}$}
  {$(0,0,1,0)^{\mathsf T}$}
  {$(1,0,0,0)^{\mathsf T}$}
  {$(0,0,0,1)^{\mathsf T}$}
  {B}
\qcmx
  {C'est $\ket{01}$ : l'ordre des facteurs est inversé.}
  {$\binom{0}{1}\otimes\binom{1}{0}=\big(0\binom{1}{0},\,1\binom{1}{0}\big)=(0,0,1,0)^{\mathsf T}$ : c'est le troisième vecteur de la base.}
  {C'est $\ket{00}$.}
  {C'est $\ket{11}$.}

% ---- Question 47 (niveau 2, correct : D)
\qcmq[2]{Chap.~\ref{chap:portes}\,$\cdot$\,\S\,\ref{sec:phase}}
  {Quelle relation lie la porte de phase $S=\mathrm{diag}(1,\ii)$ à la porte $Z$ ?}
  {$S=Z$}
  {$S^2=X$}
  {$S^4=Z$}
  {$S^2=Z$}
  {D}
\qcmx
  {$Z=\mathrm{diag}(1,-1)$ : $S$ est un quart de tour, $Z$ un demi-tour.}
  {$S^2$ est diagonale alors que $X$ ne l'est pas ; c'est $HSH$ qui est une racine carrée de $X$.}
  {$S^4=\mathrm{diag}(1,\ii^4)=\mathrm{Id}$.}
  {$S^2=\mathrm{diag}(1,\ii^2)=\mathrm{diag}(1,-1)=Z$ : deux quarts de tour autour de $z$ font un demi-tour.}

% ---- Question 48 (niveau 2, correct : A)
\qcmq[2]{Chap.~\ref{chap:portes}\,$\cdot$\,\S\,\ref{sec:controlees}}
  {Dans la porte de Toffoli (CCNOT) à trois qubits $x,y,z$, à quelle condition la cible $z$ est-elle inversée ?}
  {Seulement si $x=1$ et $y=1$}
  {Si $x=1$ ou $y=1$}
  {Si $x\neq y$}
  {Si $x=0$ et $y=0$}
  {A}
\qcmx
  {$\ket{x,y,z}\to\ket{x,y,z\oplus xy}$ : la cible est inversée si et seulement si le produit $xy$ vaut $1$.}
  {Un seul contrôle à $1$ ne suffit pas : pour $(x,y)=(1,0)$, $xy=0$ et la cible ne change pas.}
  {Ce serait deux CNOT ($z\to z\oplus x\oplus y$) ; pour $(1,0)$ le Toffoli n'inverse rien.}
  {Pour $(0,0)$ on a $xy=0$ : la cible ne change pas.}

% ---- Question 49 (niveau 2, correct : A)
\qcmq[1]{Chap.~\ref{chap:portes}\,$\cdot$\,\S\,\ref{sec:portes}}
  {Pour un hamiltonien $\Hop$ indépendant du temps, quel opérateur d'évolution $U(t,t_0)$ vérifie $\ket{\Psi(t)}=U(t,t_0)\ket{\Psi(t_0)}$ ?}
  {$U(t,t_0)=\ee^{-\ii\Hop(t-t_0)/\hbar}$}
  {$U(t,t_0)=\ee^{+\ii\Hop(t-t_0)/\hbar}$}
  {$U(t,t_0)=-\ii\hbar\,\partial_t$}
  {$U(t,t_0)=\Hop\,(t-t_0)/\hbar$}
  {A}
\qcmx
  {Solution de $\ii\hbar\,\partial_t\ket{\Psi}=\Hop\ket{\Psi}$ ; comme $\Hop$ est hermitien, $U^\dagger U=\mathrm{Id}$ : l'évolution est unitaire.}
  {Erreur de signe : c'est $U^\dagger$, qui remonte le temps.}
  {C'est un opérateur différentiel, pas un opérateur d'évolution ; l'équation de Schrödinger s'écrit $\ii\hbar\,\partial_t\ket{\Psi}=\Hop\ket{\Psi}$.}
  {On oublie l'exponentielle et le facteur $\ii$ : cet opérateur est hermitien, pas unitaire.}

% ---- Question 50 (niveau 2, correct : D)
\qcmq[2]{Chap.~\ref{chap:portes}\,$\cdot$\,\S\,\ref{sec:controlees}}
  {La porte de Fredkin (SWAP contrôlé) a pour contrôle le premier qubit. Que donne-t-elle sur $\ket{1,0,1}$ ?}
  {$\ket{1,0,1}$}
  {$\ket{0,1,1}$}
  {$\ket{1,1,1}$}
  {$\ket{1,1,0}$}
  {D}
\qcmx
  {C'est le résultat si le contrôle valait $0$ ; ici il vaut $1$ et les deux autres qubits sont différents.}
  {Le SWAP porte sur les qubits $2$ et $3$, pas sur le contrôle.}
  {Un SWAP ne change pas le nombre de $1$, il permute seulement ; $\ket{1,1,1}$ viendrait d'un NOT.}
  {Le contrôle vaut $1$ : les deux autres qubits sont échangés, $\ket{0,1}\to\ket{1,0}$.}

% ---- Question 51 (niveau 2, correct : A)
\qcmq[1]{Chap.~\ref{chap:portes}\,$\cdot$\,\S\,\ref{sec:op2}}
  {Que donne $(H\otimes I)\ket{00}$, où $H$ agit sur le premier facteur du produit tensoriel $\ket{a}\otimes\ket{b}$ ?}
  {$\ket{+}\otimes\ket{0}=\frac{1}{\sqrt2}(\ket{00}+\ket{10})$}
  {$\frac12(\ket{00}+\ket{01}+\ket{10}+\ket{11})$}
  {$\ket{0}\otimes\ket{+}=\frac{1}{\sqrt2}(\ket{00}+\ket{01})$}
  {$\frac{1}{\sqrt2}(\ket{00}+\ket{11})$}
  {A}
\qcmx
  {$(A\otimes B)(\ket{u}\otimes\ket{v})=A\ket{u}\otimes B\ket{v}$ : $H\ket{0}=\ket{+}$ sur le premier facteur, $I\ket{0}=\ket{0}$ sur le second.}
  {C'est $(H\otimes H)\ket{00}$ : ici le second qubit n'est pas touché.}
  {$H$ agit sur le premier qubit, pas sur le second : les facteurs sont inversés.}
  {C'est $\ket{\Phi^+}$ : il faut un CNOT pour intriquer ; des portes locales laissent l'état séparable.}

% ---- Question 52 (niveau 2, correct : C)
\qcmq[1]{Chap.~\ref{chap:bell}\,$\cdot$\,\S\,\ref{sec:separable}}
  {Pourquoi l'état de Bell $\ket{\Phi^+}=\frac{1}{\sqrt2}(\ket{00}+\ket{11})$ est-il dit intriqué ?}
  {Parce que $P(01)=0$}
  {Parce que chaque qubit, pris seul, est dans un état pur}
  {Il ne peut pas s'écrire comme un produit $\ket{\varphi}_A\otimes\ket{\chi}_B$ : $D=c_{00}c_{11}-c_{01}c_{10}=\frac12\neq0$}
  {Parce qu'il est orthogonal à $\ket{00}$}
  {C}
\qcmx
  {C'est vrai, mais $\ket{00}$ a aussi $P(01)=0$ et il est séparable : cette condition ne prouve rien.}
  {C'est l'inverse : pris seul, chaque qubit est complètement mixte ($\rho_A=\mathrm{Id}/2$).}
  {Un produit imposerait $D=0$ ; aucun des deux qubits n'a d'état propre.}
  {$\braket{00}{\Phi^+}=\frac{1}{\sqrt2}\neq0$ ; l'orthogonalité n'a pas de lien avec l'intrication.}

% ---- Question 53 (niveau 2, correct : B)
\qcmq[2]{Chap.~\ref{chap:bell}\,$\cdot$\,\S\,\ref{sec:teleportation}}
  {Dans la téléportation du cours, Alice obtient les résultats $a=0$ et $q=1$. Quelle correction Bob doit-il appliquer à son qubit ?}
  {La porte $X$}
  {La porte $Z$}
  {$X$ puis $Z$}
  {Aucune, l'état est déjà $\ket{\varphi}$}
  {B}
\qcmx
  {$X$ corrige le cas $a=1$ (inversion de bit) ; ici $a=0$.}
  {Bob détient $\alpha\ket{0}-\beta\ket{1}=Z\ket{\varphi}$ ; $Z$ rétablit le signe : $Z(\alpha\ket{0}-\beta\ket{1})=\alpha\ket{0}+\beta\ket{1}$.}
  {Les deux corrections ne sont nécessaires que pour $(a,q)=(1,1)$.}
  {Ce n'est vrai que pour $(a,q)=(0,0)$.}

% ---- Question 54 (niveau 2, correct : B)
\qcmq[2]{Chap.~\ref{chap:bell}\,$\cdot$\,\S\,\ref{sec:bell} ; ex.\,\ref{ex:superdense}}
  {Dans le codage superdense, Alice et Bob partagent $\ket{\Phi^+}$ et Alice agit sur son qubit. Quelle porte transforme $\ket{\Phi^+}$ en $\ket{\Psi^+}=\frac{1}{\sqrt2}(\ket{01}+\ket{10})$ ?}
  {La porte $Z$}
  {La porte $X$}
  {La porte identité}
  {La porte $H$}
  {B}
\qcmx
  {$Z$ donne $\ket{\Phi^-}=\frac{1}{\sqrt2}(\ket{00}-\ket{11})$ : elle change un signe, pas les chiffres.}
  {$X$ échange $\ket{0}$ et $\ket{1}$ sur un qubit : $\ket{00}+\ket{11}\to\ket{01}+\ket{10}$ ; c'est le message $10$.}
  {Elle laisse $\ket{\Phi^+}$ : c'est le message $00$.}
  {$(H\otimes I)\ket{\Phi^+}$ est une superposition de deux états de Bell, pas un état de Bell.}

% ---- Question 55 (niveau 2, correct : D)
\qcmq[2]{Chap.~\ref{chap:bell}\,$\cdot$\,\S\,\ref{sec:info}}
  {Pour deux qubits dans l'état de Bell $\ket{\Phi^+}$, que vaut l'entropie conditionnelle $H(A|B)=H(A,B)-H(B)$ ?}
  {$0$ bit}
  {$+1$ bit}
  {$+2$ bits}
  {$-1$ bit}
  {D}
\qcmx
  {C'est la valeur classique quand $A$ est entièrement déterminé par $B$ (deux bits identiques) ; ici $H(A,B)=0$, pas $H(B)$.}
  {C'est la valeur quand $A$ et $B$ sont indépendants : $H(A|B)=H(A)=1$.}
  {C'est l'information mutuelle $I(A;B)=1+1-0=2$ bits, pas l'entropie conditionnelle.}
  {$H(A,B)=S(\rho_{AB})=0$ (état pur) et $H(B)=1$ bit ($\rho_B=\mathrm{Id}/2$) : $0-1=-1$, valeur impossible en classique.}

% ---- Question 56 (niveau 2, correct : B)
\qcmq[1]{Chap.~\ref{chap:bell}\,$\cdot$\,\S\,\ref{sec:trace}}
  {Quelle est la matrice de densité réduite $\rho_A=\operatorname{Tr}_B\rho_{AB}$ d'un qubit d'une paire dans l'état $\ket{\Phi^+}$ ?}
  {$\rho_A=\ket{0}\bra{0}$}
  {$\rho_A=\frac12\mathrm{Id}$, de pureté $\operatorname{Tr}\rho_A^2=\frac12$}
  {$\rho_A=\ket{+}\bra{+}=\frac12\begin{pmatrix}1&1\\1&1\end{pmatrix}$}
  {$\rho_A=\ket{\Phi^+}\bra{\Phi^+}$}
  {B}
\qcmx
  {Le qubit aurait un état défini, ce qui est faux pour un état intriqué : $P(0)=P(1)=\frac12$.}
  {$\rho_A=\frac12\big(\ket{0}\bra{0}+\ket{1}\bra{1}\big)$ : les termes croisés $\ket{00}\bra{11}$ s'annulent ($\braket{1}{0}=0$). Le couple est pur, mais chaque qubit est complètement mixte.}
  {Les termes croisés sont multipliés par $\braket{1}{0}=0$ dans la trace partielle : ils disparaissent.}
  {C'est la matrice $4\times4$ du couple ; la trace partielle donne une matrice $2\times2$.}

% ---- Question 57 (niveau 2, correct : C)
\qcmq[1]{Chap.~\ref{chap:bell}\,$\cdot$\,\S\,\ref{sec:separable}}
  {Pour $\ket{\Psi}=c_{00}\ket{00}+c_{01}\ket{01}+c_{10}\ket{10}+c_{11}\ket{11}$, quel est le critère de séparabilité ?}
  {$c_{00}c_{11}+c_{01}c_{10}=0$}
  {$c_{00}c_{01}-c_{10}c_{11}=0$}
  {$D=c_{00}c_{11}-c_{01}c_{10}=0$}
  {$c_{00}=c_{11}$ et $c_{01}=c_{10}$}
  {C}
\qcmx
  {Mauvais signe : $\ket{+}\otimes\ket{-}$ est séparable mais donnerait $-\frac12\neq0$.}
  {Mauvais appariement : $\ket{\Psi^+}$ (intriqué) la vérifie, puisque $c_{00}=c_{11}=0$.}
  {Un produit $(a_0\ket{0}+a_1\ket{1})\otimes(b_0\ket{0}+b_1\ket{1})$ a $c_{ij}=a_ib_j$, donc $c_{00}c_{11}=a_0a_1b_0b_1=c_{01}c_{10}$ ; la réciproque est vraie aussi.}
  {Ni nécessaire ($\ket{00}$ est séparable avec $c_{00}\neq c_{11}$) ni suffisant ($\ket{\Phi^+}$ la vérifie et est intriqué).}

% ---- Question 58 (niveau 2, correct : B)
\qcmq[2]{Chap.~\ref{chap:bell}\,$\cdot$\,\S\,\ref{sec:bell}}
  {Quelle est la valeur propre de $X\otimes X$ pour l'état $\ket{\Phi^+}$ ?}
  {$-1$}
  {$+1$}
  {$0$}
  {$\ket{\Phi^+}$ n'est pas un état propre de $X\otimes X$}
  {B}
\qcmx
  {C'est la valeur pour $\ket{\Phi^-}$ et $\ket{\Psi^-}$ : le signe $-$ de l'état se retrouve dans la valeur propre.}
  {$(X\otimes X)\ket{00}=\ket{11}$ et $(X\otimes X)\ket{11}=\ket{00}$, donc $(X\otimes X)\ket{\Phi^+}=\ket{\Phi^+}$.}
  {$X\otimes X$ est unitaire : ses valeurs propres sont de module $1$, jamais nulles.}
  {Les quatre états de Bell sont des états propres communs de $Z\otimes Z$ et de $X\otimes X$.}

% ---- Question 59 (niveau 2, correct : A)
\qcmq[2]{Chap.~\ref{chap:bell}\,$\cdot$\,\S\,\ref{sec:info}}
  {Que vaut l'entropie de von Neumann $S(\rho)=-\sum_i\lambda_i\log_2\lambda_i$ d'un état pur ?}
  {$0$ bit}
  {$1$ bit}
  {$-1$ bit}
  {$+\infty$}
  {A}
\qcmx
  {Les valeurs propres sont $1$ et $0$ : $1\cdot\log_21=0$ et $0\cdot\log_20=0$ (par continuité, $x\log_2x\to0$).}
  {C'est la valeur pour l'état maximalement mixte $\mathrm{Id}/2$ d'un qubit.}
  {L'entropie de $\rho$ est positive ou nulle ; seule l'entropie conditionnelle quantique peut être négative.}
  {$\log_20$ diverge, mais le terme $0\cdot\log_20$ vaut $0$ par continuité.}

% ---- Question 60 (niveau 2, correct : A)
\qcmq[2]{Chap.~\ref{chap:bell}\,$\cdot$\,\S\,\ref{sec:bell}}
  {Lequel des états de Bell a pour valeur propre $-1$ à la fois pour $Z\otimes Z$ et pour $X\otimes X$ ?}
  {$\ket{\Psi^-}=\frac{1}{\sqrt2}(\ket{01}-\ket{10})$}
  {$\ket{\Phi^+}$}
  {$\ket{\Phi^-}$}
  {$\ket{\Psi^+}$}
  {A}
\qcmx
  {Résultats opposés dans les deux bases : $Z\otimes Z=-1$ (chiffres différents) et $X\otimes X=-1$ (signe $-$ dans l'état).}
  {Valeurs propres $(Z\otimes Z,\,X\otimes X)=(+1,+1)$.}
  {Valeurs propres $(Z\otimes Z,\,X\otimes X)=(+1,-1)$.}
  {Valeurs propres $(Z\otimes Z,\,X\otimes X)=(-1,+1)$.}

\clearpage
%% ======================================================================
\subsection{Niveau 3 : avancé (questions 61 à 90)}\label{sec:qcm-n3}
%% ======================================================================
\noindent{\color{insagris}\small Niveau de difficulté : Avancé\quad $\bullet\bullet\bullet\circ$}\par\smallskip
% ---- Question 61 (niveau 3, correct : A)
\qcmq[2]{Chap.~\ref{chap:histoire}\,$\cdot$\,\S\,\ref{sec:solutions}}
  {Un électron d'un puits infini passe du niveau $n=3$ au niveau $n=2$ en émettant un photon. Quelle est l'énergie du photon, en fonction de $E_1$ ?}
  {$5E_1$}
  {$E_1$}
  {$\frac94E_1$}
  {$13E_1$}
  {A}
\qcmx
  {$E_3-E_2=9E_1-4E_1=5E_1$ ; le photon a la fréquence $f=(E_3-E_2)/h$.}
  {On a soustrait les numéros de niveau ($3-2=1$) au lieu de leurs carrés.}
  {C'est le rapport $E_3/E_2$, pas la différence $E_3-E_2$.}
  {On a additionné $9E_1+4E_1$ : l'énergie du photon est une différence de niveaux.}

% ---- Question 62 (niveau 3, correct : A)
\qcmq[2]{Chap.~\ref{chap:histoire}\,$\cdot$\,\S\,\ref{sec:superposition}}
  {À $t=0$, une particule du puits infini est dans l'état $\Psi=\frac{\sqrt3}{2}\,\psi_1+\frac12\,\psi_3$ ($\psi_n$ : états stationnaires normalisés, d'énergies $E_n=n^2E_1$). Quelle est son énergie moyenne $\langle E\rangle$ ?}
  {$3E_1$}
  {$5E_1$}
  {$\left(\frac{\sqrt3}{2}+\frac92\right)E_1$}
  {$7E_1$}
  {A}
\qcmx
  {$P(E_1)=\frac34$ et $P(E_3)=\frac14$, donc $\langle E\rangle=\frac34E_1+\frac14\cdot9E_1=3E_1$.}
  {On a pondéré également les deux niveaux, $\frac{E_1+9E_1}{2}$ : les poids sont les $|c_i|^2$, pas $\frac12$.}
  {On a pondéré par les amplitudes $c_i$ au lieu de leurs carrés $|c_i|^2$.}
  {Poids permutés : $\frac14E_1+\frac34\cdot9E_1=7E_1$ ; la probabilité de $E_1$ est $\left|\frac{\sqrt3}{2}\right|^2=\frac34$.}

% ---- Question 63 (niveau 3, correct : C)
\qcmq[1]{Chap.~\ref{chap:histoire}\,$\cdot$\,\S\,\ref{sec:solutions}}
  {Pour l'état stationnaire $n=2$ du puits infini $]0,a[$, que valent la densité de probabilité de présence en $x=a/2$ et la position moyenne $\langle\hat x\rangle$ ?}
  {Densité maximale en $a/2$, et $\langle\hat x\rangle=a/2$}
  {Densité nulle en $a/2$, et $\langle\hat x\rangle=0$}
  {Densité nulle en $a/2$, et $\langle\hat x\rangle=a/2$}
  {Densité maximale en $a/2$, et $\langle\hat x\rangle=a/4$}
  {C}
\qcmx
  {C'est le cas de $n=1$ ; pour $n=2$ il y a un nœud au centre.}
  {Une densité nulle au centre ne donne pas une moyenne nulle : $\langle\hat x\rangle$ pondère tout le puits et vaut $\frac a2$ par symétrie.}
  {$|\Psi_2|^2=\frac2a\sin^2\frac{2\pi x}{a}$ s'annule en $x=\frac a2$ (nœud), alors que la symétrie impose $\langle\hat x\rangle=\frac a2$ : la moyenne n'est pas la position la plus probable.}
  {Ni l'un ni l'autre : maxima en $\frac a4$ et $\frac{3a}4$, nœud en $\frac a2$, moyenne $\frac a2$.}

% ---- Question 64 (niveau 3, correct : C)
\qcmq[2]{Chap.~\ref{chap:dirac}\,$\cdot$\,\S\,\ref{sec:projecteurs}}
  {Quelle est la matrice du projecteur $\hat P_+=\ket{+}\bra{+}$ dans la base $\{\ket{0},\ket{1}\}$ ?}
  {$\frac12\begin{pmatrix}1&-1\\-1&1\end{pmatrix}$}
  {$\begin{pmatrix}1&0\\0&0\end{pmatrix}$}
  {$\frac12\begin{pmatrix}1&1\\1&1\end{pmatrix}$}
  {$\frac{1}{\sqrt2}\begin{pmatrix}1&1\\1&-1\end{pmatrix}$}
  {C}
\qcmx
  {C'est $\ket{-}\bra{-}$, le projecteur sur $\ket{-}$.}
  {C'est $\ket{0}\bra{0}$ : projecteur de la base $Z$, pas de la base $X$.}
  {$\ket{+}\bra{+}=\frac12(\ket0+\ket1)(\bra0+\bra1)$ ; on vérifie $\hat P_+^2=\hat P_+$ et $\operatorname{Tr}\hat P_+=1$.}
  {C'est la matrice de Hadamard, unitaire ($H^2=\mathrm{Id}$), et non un projecteur (un projecteur vérifie $P^2=P$).}

% ---- Question 65 (niveau 3, correct : B)
\qcmq[2]{Chap.~\ref{chap:dirac}\,$\cdot$\,\S\,\ref{sec:projecteurs}}
  {Un qubit préparé dans $\ket{0}$ est mesuré en base $X$, puis aussitôt en base $Z$. Quelle est la probabilité d'obtenir $1$ à la seconde mesure ?}
  {$0$}
  {$\frac12$}
  {$1$}
  {$\frac14$}
  {B}
\qcmx
  {On garde $\ket{0}$ en mémoire : la mesure en base $X$ a effondré l'état sur $\ket{\pm}$ et effacé $\ket{0}$.}
  {La première mesure donne $\ket{+}$ ou $\ket{-}$ (probabilité $\frac12$ chacun) ; dans les deux cas $P(1)=\frac12$ en base $Z$ : $\frac12\cdot\frac12+\frac12\cdot\frac12=\frac12$.}
  {Sans rapport : $\ket{0}$ n'est plus l'état après la mesure en base $X$.}
  {On n'a compté qu'une branche de la première mesure ; il faut sommer les deux ($\frac14+\frac14$).}

% ---- Question 66 (niveau 3, correct : D)
\qcmq[2]{Chap.~\ref{chap:dirac}\,$\cdot$\,\S\,\ref{sec:projecteurs}}
  {Un qubit est dans l'état $\ket{\psi}=\frac{3}{5}\ket{0}+\frac{4}{5}\ket{1}$. Quelle est la probabilité d'obtenir $-$ en mesurant dans la base $X$ ?}
  {$\frac{16}{25}$}
  {$\frac{1}{25}$}
  {$\frac{49}{50}$}
  {$\frac{1}{50}$}
  {D}
\qcmx
  {C'est $P(1)$ en base $Z$ ($|\frac45|^2$) : on a mesuré dans la mauvaise base.}
  {On a oublié le facteur $\frac{1}{\sqrt2}$ de $\ket{-}$ : $\left|\frac35-\frac45\right|^2=\frac1{25}$, il manque le $\frac12$.}
  {C'est $P(+)$ ; la probabilité demandée est celle du résultat $-$.}
  {$\braket{-}{\psi}=\frac{1}{\sqrt2}\left(\frac35-\frac45\right)=-\frac{1}{5\sqrt2}$, donc $P(-)=\frac1{50}$ ; et $P(+)=\frac{49}{50}$.}

% ---- Question 67 (niveau 3, correct : B)
\qcmq[1]{Chap.~\ref{chap:dirac}\,$\cdot$\,\S\,\ref{sec:projecteurs}}
  {Lequel de ces kets décrit le même état physique que $\ket{+}=\frac{1}{\sqrt2}(\ket{0}+\ket{1})$ ?}
  {$\frac{1}{\sqrt2}\big(\ket{0}+\ii\ket{1}\big)$}
  {$\ee^{\ii\pi/3}\ket{+}$}
  {$\frac{1}{\sqrt2}\big(\ket{0}+\ee^{\ii\pi}\ket{1}\big)$}
  {$\frac{1}{\sqrt2}\big(\ket{0}+\ee^{\ii\pi/3}\ket{1}\big)$}
  {B}
\qcmx
  {La phase $\ii$ ne porte que sur $\ket{1}$ : c'est une phase \emph{relative}, et l'état est $\ket{{+}\ii}$ (base $Y$).}
  {Un facteur de phase \emph{globale} multiplie tout le ket : toutes les probabilités de mesure (et $\rho=\ket{+}\bra{+}$) sont inchangées.}
  {$\ee^{\ii\pi}=-1$ : c'est $\ket{-}$, orthogonal à $\ket{+}$.}
  {Phase relative $\frac\pi3$ : c'est un autre état, sur l'équateur de la sphère de Bloch, à l'azimut $\frac\pi3$.}

% ---- Question 68 (niveau 3, correct : A)
\qcmq[2]{Chap.~\ref{chap:dirac}\,$\cdot$\,\S\,\ref{sec:qubit}}
  {Laquelle de ces paires est une base orthonormée de $\mathbb{C}^2$ ?}
  {$\{\ket{{+}\ii},\ket{{-}\ii}\}$}
  {$\{\ket{0},\ket{+}\}$}
  {$\{\ket{+},\ket{{-}\ii}\}$}
  {$\{\ket{0},\ket{0}+\ket{1}\}$}
  {A}
\qcmx
  {Ces deux états sont normés et orthogonaux : $\braket{{+}\ii}{{-}\ii}=\frac12(1-1)=0$.}
  {Normés, mais $\braket{0}{+}=\frac{1}{\sqrt2}\neq0$ : pas orthogonaux.}
  {$\braket{+}{{-}\ii}=\frac12(1-\ii)\neq0$ : pas orthogonaux.}
  {Le second ket n'est pas normé (norme au carré $2$) ni orthogonal à $\ket{0}$.}

% ---- Question 69 (niveau 3, correct : C)
\qcmq[2]{Chap.~\ref{chap:operateurs}\,$\cdot$\,\S\,\ref{sec:rho}}
  {Laquelle de ces matrices peut être la matrice de densité d'un qubit ?}
  {$\begin{pmatrix}\frac12&\frac34\\\frac34&\frac12\end{pmatrix}$}
  {$\begin{pmatrix}\frac12&\frac{\ii}{2}\\\frac{\ii}{2}&\frac12\end{pmatrix}$}
  {$\begin{pmatrix}\frac23&\frac{\ii}{3}\\-\frac{\ii}{3}&\frac13\end{pmatrix}$}
  {$\begin{pmatrix}\frac34&0\\0&\frac12\end{pmatrix}$}
  {C}
\qcmx
  {Hermitienne et de trace $1$, mais pas positive : $\det=\frac14-\frac9{16}<0$, valeurs propres $\frac54$ et $-\frac14$ (une « probabilité » négative).}
  {De trace $1$ mais pas hermitienne : il faudrait $\rho_{10}=\overline{\rho_{01}}=-\frac{\ii}{2}$.}
  {Hermitienne ($\rho_{10}=\overline{\rho_{01}}$), de trace $1$ et positive : $\det\rho=\frac29-\frac19=\frac19>0$ avec $\operatorname{Tr}\rho>0$ donne deux valeurs propres positives, $\frac{3\pm\sqrt5}{6}$.}
  {Hermitienne et positive, mais de trace $\frac54\neq1$ : les probabilités diagonales ne somment pas à $1$.}

% ---- Question 70 (niveau 3, correct : B)
\qcmq[2]{Chap.~\ref{chap:operateurs}\,$\cdot$\,\S\,\ref{sec:op-dirac}}
  {Un qubit est dans l'état $\ket{\psi}=\frac{1}{\sqrt5}\big(\ket{0}+2\ii\ket{1}\big)$. Que vaut la valeur moyenne $\langle Y\rangle=\bra{\psi}Y\ket{\psi}$, avec $Y=\begin{pmatrix}0&-\ii\\\ii&0\end{pmatrix}$ ?}
  {$0$}
  {$\frac45$}
  {$-\frac45$}
  {$-\frac35$}
  {B}
\qcmx
  {On a oublié de conjuguer les amplitudes du bra : $\psi^{\mathsf T}Y\psi=\frac15(2+2\ii\cdot\ii)=0$ (c'est d'ailleurs la valeur de $\langle X\rangle$ pour cet état).}
  {$Y\ket\psi=\frac1{\sqrt5}(2\ket0+\ii\ket1)$, puis $\bra\psi Y\ket\psi=\frac15\big(1\cdot2+(-2\ii)(\ii)\big)=\frac45$ ; en général $\langle Y\rangle=2\,\mathrm{Im}(\bar\alpha\beta)$.}
  {Erreur de signe : on a utilisé la transposée $\begin{pmatrix}0&\ii\\-\ii&0\end{pmatrix}=-Y$.}
  {C'est $\langle Z\rangle=\frac15-\frac45$ : on a calculé la valeur moyenne d'une autre observable.}

% ---- Question 71 (niveau 3, correct : D)
\qcmq[1]{Chap.~\ref{chap:operateurs}\,$\cdot$\,\S\,\ref{sec:spectral}}
  {On mesure l'observable $A=\begin{pmatrix}2&1\\1&2\end{pmatrix}$ sur un qubit préparé dans $\ket{0}$. Quelle affirmation est correcte ?}
  {On obtient toujours $2$}
  {On obtient $+1$ ou $-1$, avec la probabilité $\frac12$ chacun}
  {On obtient $2$ ou $1$, avec les probabilités $\frac45$ et $\frac15$}
  {On obtient $3$ ou $1$, avec la probabilité $\frac12$ chacun (moyenne $2$)}
  {D}
\qcmx
  {On a lu le coefficient diagonal : $A\ket0=2\ket0+\ket1$ n'est pas proportionnel à $\ket0$, donc $\ket0$ n'est pas vecteur propre de $A$.}
  {Ce sont les valeurs propres de $X$, pas de $A$ : le décalage $2\,\mathrm{Id}$ translate le spectre ($2\pm1$).}
  {On a lu $A\ket0=2\ket0+\ket1$ comme un état : $A$ n'est pas unitaire, ce n'est pas une évolution ; les résultats possibles d'une mesure sont les valeurs propres de $A$.}
  {$A=2\,\mathrm{Id}+X$ a les vecteurs propres de $X$ : valeur propre $3$ pour $\ket+$, $1$ pour $\ket-$. Comme $\ket0=\frac{1}{\sqrt2}(\ket++\ket-)$, les deux résultats sont équiprobables, et $\langle A\rangle=\bra0A\ket0=2$.}

% ---- Question 72 (niveau 3, correct : D)
\qcmq[2]{Chap.~\ref{chap:operateurs}\,$\cdot$\,\S\,\ref{sec:pur-mixte}}
  {Un qubit est préparé dans $\ket{0}$ avec la probabilité $\frac12$ et dans $\ket{+}$ avec la probabilité $\frac12$. Quelle est la pureté $\operatorname{Tr}(\rho^2)$ de son état ?}
  {$1$}
  {$\frac12$}
  {$\frac58$}
  {$\frac34$}
  {D}
\qcmx
  {Un mélange de deux états purs n'est pur que s'ils sont identiques ; ici $\ket0\neq\ket+$, donc $\operatorname{Tr}\rho^2<1$.}
  {Valeur de l'état maximalement mixte $\frac{\mathrm{Id}}2$, atteinte pour deux états \emph{orthogonaux} ; ici ils ne le sont pas (on a aussi pu oublier les termes croisés).}
  {On n'a additionné que les carrés des éléments diagonaux ($\frac9{16}+\frac1{16}$) : il manque $2\rho_{01}\rho_{10}=\frac2{16}$.}
  {$\rho=\frac12\ket0\bra0+\frac12\ket+\bra+=\begin{pmatrix}\frac34&\frac14\\\frac14&\frac14\end{pmatrix}$, donc $\operatorname{Tr}\rho^2=\frac9{16}+\frac1{16}+2\cdot\frac1{16}=\frac34$ : état mixte, mais moins que $\frac{\mathrm{Id}}2$ car $\ket0$ et $\ket+$ ne sont pas orthogonaux.}

% ---- Question 73 (niveau 3, correct : B)
\qcmq[2]{Chap.~\ref{chap:operateurs}\,$\cdot$\,\S\,\ref{sec:bloch}}
  {Quel est le vecteur de Bloch $(r_x,r_y,r_z)$ de l'état $\ket{\psi}=\cos\frac\pi6\,\ket0+\sin\frac\pi6\,\ket1$ ?}
  {$\left(\frac12,\,0,\,\frac{\sqrt3}{2}\right)$}
  {$\left(\frac{\sqrt3}{2},\,0,\,\frac12\right)$}
  {$\left(\frac{\sqrt3}{4},\,0,\,\frac12\right)$}
  {$\left(\frac{\sqrt3}{2},\,0,\,\frac34\right)$}
  {B}
\qcmx
  {On a pris $\theta=\frac\pi6$ : l'angle de Bloch est le \emph{double} de l'angle qui figure dans les amplitudes.}
  {Ici $\frac\theta2=\frac\pi6$, donc $\theta=\frac\pi3$ et $\varphi=0$ : $\vec r=(\sin\theta\cos\varphi,\sin\theta\sin\varphi,\cos\theta)$. Contrôle : $r_z=|\alpha|^2-|\beta|^2=\frac34-\frac14$ et $r_x=2\alpha\beta=\frac{\sqrt3}2$.}
  {On a écrit $r_x=\alpha\beta$ en oubliant le facteur $2$ ; on le voit à $|\vec r|^2=\frac7{16}<1$, impossible pour un état pur.}
  {On a pris $r_z=|\alpha|^2$ au lieu de $|\alpha|^2-|\beta|^2$ ; or $|\vec r|^2=\frac{21}{16}>1$ est impossible.}

% ---- Question 74 (niveau 3, correct : A)
\qcmq[2]{Chap.~\ref{chap:operateurs}\,$\cdot$\,\S\,\ref{sec:bloch}}
  {Un qubit a pour matrice de densité $\rho=\frac12\begin{pmatrix}1&-\ii\\\ii&1\end{pmatrix}$. Quel est son vecteur de Bloch ?}
  {$(0,\,1,\,0)$ : c'est l'état pur $\ket{{+}\ii}$}
  {$(0,\,-1,\,0)$}
  {$\left(0,\,\frac12,\,0\right)$}
  {$(0,\,0,\,0)$}
  {A}
\qcmx
  {Avec $\rho=\frac12\begin{pmatrix}1+r_z&r_x-\ii r_y\\r_x+\ii r_y&1-r_z\end{pmatrix}$ : $r_z=0$ et $r_x-\ii r_y=-\ii$, donc $r_y=1$, $r_x=0$. Comme $|\vec r|=1$, l'état est pur : $\ket{{+}\ii}\bra{{+}\ii}$.}
  {Erreur de signe : $\rho_{01}=\frac{r_x-\ii r_y}{2}$, donc $\rho_{01}=-\frac\ii2$ correspond à $r_y=+1$ ($r_y=-1$ pour $\ket{{-}\ii}$).}
  {On a lu $\rho_{01}=-\frac\ii2$ comme $-\ii r_y$ : il faut multiplier par $2$ ($\rho=\frac12(\mathrm{Id}+\vec r\cdot\vec\sigma)$).}
  {Les diagonales $\frac12,\frac12$ donnent $r_z=0$, mais les éléments non diagonaux (cohérences) ne sont pas nuls : l'état est pur, et non $\frac{\mathrm{Id}}2$.}

% ---- Question 75 (niveau 3, correct : C)
\qcmq[2]{Chap.~\ref{chap:operateurs}\,$\cdot$\,\S\,\ref{sec:tenseur}}
  {Que donne $(X\otimes Z)\,\big(\ket{+}\otimes\ket{-}\big)$ ?}
  {$-\ket{-}\otimes\ket{-}$}
  {$\ket{-}\otimes\ket{+}$}
  {$\ket{+}\otimes\ket{+}$}
  {$\ket{+}\otimes\ket{-}$}
  {C}
\qcmx
  {On a échangé les facteurs ($Z\otimes X$) : $Z\ket+=\ket-$ et $X\ket-=-\ket-$. Le premier opérateur agit sur le premier qubit.}
  {On a cru que $X$ transforme $\ket+$ en $\ket-$ ; or $X$ échange $\ket0$ et $\ket1$ et laisse $\ket+$ inchangé.}
  {$(A\otimes B)(\ket u\otimes\ket v)=A\ket u\otimes B\ket v$ : $X\ket+=\ket+$ ($\ket+$ est propre pour $X$, valeur propre $+1$) et $Z\ket-=\ket+$ ($Z$ échange $\ket+$ et $\ket-$).}
  {On a cru que $Z$ laisse $\ket-$ inchangé ; $Z$ ne laisse invariants que $\ket0$ et $\ket1$, et échange $\ket+$ et $\ket-$.}

% ---- Question 76 (niveau 3, correct : C)
\qcmq[2]{Chap.~\ref{chap:portes}\,$\cdot$\,\S\,\ref{sec:phase}}
  {Le qubit $\ket{+}$ traverse la porte $T=\mathrm{diag}\big(1,\ee^{\ii\pi/4}\big)$, puis est mesuré dans la base $X$ $\{\ket+,\ket-\}$. Quelle est la probabilité d'obtenir $+$ ?}
  {$\frac12$}
  {$\frac{2-\sqrt2}{4}\approx0{,}15$}
  {$\frac{2+\sqrt2}{4}\approx0{,}85$}
  {$\frac{\sqrt2}{2}\approx0{,}71$}
  {C}
\qcmx
  {On a cru que la phase ne change pas les probabilités : c'est vrai en base $Z$, mais la base $X$ est sensible à la phase \emph{relative} (le point de Bloch est à l'azimut $\frac\pi4$ de l'axe $+x$).}
  {C'est $P(-)=\sin^2\frac\pi8$ : on a permuté les résultats $+$ et $-$.}
  {$T\ket+=\frac1{\sqrt2}(\ket0+\ee^{\ii\pi/4}\ket1)$, donc $\braket{+}{\psi}=\frac12\big(1+\ee^{\ii\pi/4}\big)$ et $P(+)=\frac14\big|1+\ee^{\ii\pi/4}\big|^2=\frac{2+2\cos\frac\pi4}{4}=\frac{2+\sqrt2}{4}$ (soit $\cos^2\frac\pi8$).}
  {On a pris $\cos\frac\pi4$ comme probabilité ; il faut le \emph{demi}-angle et le carré : $P=\cos^2\frac{\pi/4}{2}=\cos^2\frac\pi8$.}

% ---- Question 77 (niveau 3, correct : B)
\qcmq[2]{Chap.~\ref{chap:portes}\,$\cdot$\,\S\,\ref{sec:hadamard}}
  {Que vaut le produit de matrices $HYH$, où $Y=\begin{pmatrix}0&-\ii\\\ii&0\end{pmatrix}$ ?}
  {$Y$}
  {$-Y$}
  {$Z$}
  {$X$}
  {B}
\qcmx
  {Par analogie avec $HXH=Z$, on a cru que $H$ ne modifie pas $Y$ ; en fait l'axe $y$ est retourné (signe $-$).}
  {$H$ échange les axes $x$ et $z$ et retourne l'axe $y$ : $HXH=Z$, $HZH=X$, $HYH=-Y$. Direct : $HY=\frac{1}{\sqrt2}\begin{pmatrix}\ii&-\ii\\-\ii&-\ii\end{pmatrix}$, puis $HYH=\begin{pmatrix}0&\ii\\-\ii&0\end{pmatrix}=-Y$.}
  {C'est $HXH$ : $H$ n'échange que les axes $x$ et $z$, l'axe $y$ reste l'axe $y$.}
  {C'est $HZH$ : $Y$ ne peut pas devenir $X$, car $H$ laisse l'axe $y$ à sa place (au signe près).}

% ---- Question 78 (niveau 3, correct : D)
\qcmq[1]{Chap.~\ref{chap:portes}\,$\cdot$\,\S\,\ref{sec:op2}}
  {Que donne $(H\otimes H)\ket{01}$, où $\ket{01}=\ket0\otimes\ket1$ ?}
  {$\frac12\big(\ket{00}+\ket{01}-\ket{10}-\ket{11}\big)$}
  {$\frac12\big(\ket{00}-\ket{01}-\ket{10}+\ket{11}\big)$}
  {$\frac12\big(\ket{00}+\ket{01}+\ket{10}+\ket{11}\big)$}
  {$\frac12\big(\ket{00}-\ket{01}+\ket{10}-\ket{11}\big)$}
  {D}
\qcmx
  {C'est $\ket-\otimes\ket+$ : on a attribué le $\ket-$ au premier qubit, alors que c'est le second qui vaut $\ket1$.}
  {C'est $\ket-\otimes\ket-$ : $H\ket0=\ket+$, et non $\ket-$.}
  {C'est $\ket+\otimes\ket+$, le résultat pour $\ket{00}$ : on a oublié que $H\ket1=\ket-$ introduit des signes.}
  {$H\ket0\otimes H\ket1=\ket+\otimes\ket-=\frac12(\ket0+\ket1)\otimes(\ket0-\ket1)$. Le signe $-$ apparaît quand le \emph{second} chiffre vaut $1$ (formule générale $\frac12\sum_y(-1)^{x\cdot y}\ket y$ avec $x=01$).}

% ---- Question 79 (niveau 3, correct : D)
\qcmq[2]{Chap.~\ref{chap:portes}\,$\cdot$\,\S\,\ref{sec:cnot}}
  {Avec le CNOT du cours (contrôle : second qubit $B$ ; cible : premier qubit $A$ ; $\mathrm{CNOT}\ket{a,b}=\ket{a\oplus b,\,b}$), quel état obtient-on à partir de $\ket{1}_A\otimes\ket{+}_B$ ?}
  {$\ket{1}\otimes\ket{+}$ (inchangé)}
  {$\frac{1}{\sqrt2}\big(\ket{00}+\ket{11}\big)=\ket{\Phi^+}$}
  {$\frac{1}{\sqrt2}\big(\ket{01}-\ket{10}\big)=\ket{\Psi^-}$}
  {$\frac{1}{\sqrt2}\big(\ket{01}+\ket{10}\big)=\ket{\Psi^+}$}
  {D}
\qcmx
  {On a pris le premier qubit comme contrôle : le CNOT agirait alors sur $B$, et $X\ket+=\ket+$ laisserait l'état inchangé. Mais ici le contrôle est $B$.}
  {C'est le résultat pour $\ket{0}_A\ket{+}_B$ ; ici $A$ vaut $\ket1$, ce qui échange les deux termes.}
  {Aucun signe $-$ n'apparaît : le CNOT permute des états de base sans leur ajouter de phase.}
  {$\ket1\ket+=\frac1{\sqrt2}(\ket{10}+\ket{11})$ ; $\ket{10}\mapsto\ket{1\oplus0,0}=\ket{10}$ et $\ket{11}\mapsto\ket{1\oplus1,1}=\ket{01}$. C'est un état de Bell, donc intriqué.}

% ---- Question 80 (niveau 3, correct : C)
\qcmq[2]{Chap.~\ref{chap:portes}\,$\cdot$\,\S\,\ref{sec:swap}}
  {Que vaut $\mathrm{SWAP}\,(X\otimes\mathrm{Id})\,\mathrm{SWAP}$ ?}
  {$X\otimes\mathrm{Id}$}
  {$\mathrm{Id}\otimes\mathrm{Id}$}
  {$\mathrm{Id}\otimes X$}
  {$X\otimes X$}
  {C}
\qcmx
  {On a supposé que $\mathrm{SWAP}$ commute avec $X\otimes\mathrm{Id}$ ; ce n'est pas le cas, puisque cet opérateur n'agit pas de la même façon sur les deux qubits.}
  {On a simplifié $\mathrm{SWAP}^2=\mathrm{Id}$ en oubliant que $X\otimes\mathrm{Id}$ se trouve \emph{entre} les deux SWAP : on ne peut pas les rapprocher.}
  {$\mathrm{SWAP}\,(X\otimes\mathrm{Id})\,\mathrm{SWAP}\ket{a,b}=\mathrm{SWAP}\,(X\otimes\mathrm{Id})\ket{b,a}=\mathrm{SWAP}\ket{b\oplus1,a}=\ket{a,b\oplus1}$ : le SWAP transporte l'opérateur d'un qubit à l'autre (même mécanisme que $\mathrm{CNOT}_{21}=\mathrm{SWAP}\,\mathrm{CNOT}_{12}\,\mathrm{SWAP}$).}
  {Cet opérateur inverserait les deux bits ; ici un seul bit est inversé, celui du \emph{second} qubit.}

% ---- Question 81 (niveau 3, correct : C)
\qcmq[1]{Chap.~\ref{chap:portes}\,$\cdot$\,\S\,\ref{sec:controlees}}
  {La porte $Z$ contrôlée ($\mathrm{CZ}=\mathrm{diag}(1,1,1,-1)$, contrôle : premier qubit) est appliquée à $\ket{+}\otimes\ket{+}$. Quel état obtient-on ?}
  {$\ket{+}\otimes\ket{+}$}
  {$\ket{-}\otimes\ket{-}$}
  {$\frac12\big(\ket{00}+\ket{01}+\ket{10}-\ket{11}\big)$, un état intriqué}
  {$\ket{+}\otimes\ket{-}$}
  {C}
\qcmx
  {On a cru qu'un simple signe est une phase globale ; mais le signe $-$ ne touche que $\ket{11}$ : c'est une phase \emph{relative}.}
  {C'est $Z\otimes Z$ ($Z$ sur les deux qubits), et non une porte contrôlée.}
  {CZ ne change que le signe de $\ket{11}$. Critère de séparabilité : $c_{00}c_{11}-c_{01}c_{10}=-\frac14-\frac14=-\frac12\neq0$ ; on peut l'écrire $\frac{1}{\sqrt2}\big(\ket0\ket++\ket1\ket-\big)$.}
  {On a appliqué $Z$ au second qubit sans condition ; le contrôle exige que le premier qubit soit dans $\ket1$.}

% ---- Question 82 (niveau 3, correct : D)
\qcmq[2]{Chap.~\ref{chap:portes}\,$\cdot$\,\S\,\ref{sec:phase}, \ref{sec:portes-bloch}}
  {Un qubit préparé dans $\ket{0}$ traverse successivement $H$, puis $S$, puis $H$ (dans cet ordre). Quel est l'état final, à une phase globale près ?}
  {$\ket{{+}\ii}$}
  {$\ket{0}$}
  {$\ket{1}$}
  {$\ket{{-}\ii}=\frac{1}{\sqrt2}\big(\ket0-\ii\ket1\big)$}
  {D}
\qcmx
  {Mauvais sens de rotation : on a perdu un signe en calculant $H\ket{{+}\ii}$ (la phase relative vaut $-\ii$, pas $+\ii$).}
  {On a « simplifié » $HSH$ en $S$ (puis $S\ket0=\ket0$) : mais $H$ ne commute pas avec $S$, on ne peut pas supprimer les deux $H$.}
  {C'est le résultat pour $HZH=X$ ; or $S$ n'est qu'un \emph{quart} de tour autour de $z$ (et non un demi-tour comme $Z$), donc $HSH$ est un quart de tour autour de $x$.}
  {$HSH=\sqrt X$ est un quart de tour autour de $x$ (§\,\ref{sec:phase}). Directement : $H\ket0=\ket+$, $S\ket+=\ket{{+}\ii}$, $H\ket{{+}\ii}=\frac{1+\ii}{2}\ket0+\frac{1-\ii}{2}\ket1=\frac{1+\ii}{2}\big(\ket0-\ii\ket1\big)$.}

% ---- Question 83 (niveau 3, correct : A)
\qcmq[1]{Chap.~\ref{chap:bell}\,$\cdot$\,\S\,\ref{sec:mesure2}}
  {Deux qubits sont dans l'état $\ket{\Psi}=\frac{1}{\sqrt3}\big(\ket{00}+\ket{01}+\ket{11}\big)$. On mesure le \emph{second} qubit $B$ en base $Z$ et on obtient $1$. Que valait la probabilité de ce résultat, et dans quel état se trouve alors $A$ ?}
  {$P(B{=}1)=\frac23$ ; $A$ est dans l'état $\ket{+}$}
  {$P(B{=}1)=\frac13$ ; $A$ est dans l'état $\ket{0}$}
  {$P(B{=}1)=\frac23$ ; $A$ est dans l'état $\ket{1}$}
  {$P(B{=}1)=\frac13$ ; $A$ est dans l'état $\ket{+}$}
  {A}
\qcmx
  {$P(B{=}1)=|c_{01}|^2+|c_{11}|^2=\frac13+\frac13$. On garde les termes dont le second chiffre vaut $1$ : $\frac{1}{\sqrt3}\big(\ket{01}+\ket{11}\big)=\frac{1}{\sqrt3}(\ket0+\ket1)\otimes\ket1$, qui renormalisé donne $\ket{+}_A\otimes\ket{1}_B$.}
  {C'est le résultat $B=0$ : seul le terme $\ket{00}$ a un second chiffre nul, d'où la probabilité $\frac13$ et $A$ dans $\ket0$.}
  {On a projeté sur le seul terme $\ket{11}$ ; or $\ket{01}$ a aussi son second chiffre égal à $1$ : $A$ vaut $0$ ou $1$ (état $\ket+$).}
  {L'état de $A$ est juste, mais on n'a compté qu'un seul terme pour la probabilité ; il faut sommer $|c_{01}|^2+|c_{11}|^2$.}

% ---- Question 84 (niveau 3, correct : A)
\qcmq[2]{Chap.~\ref{chap:bell}\,$\cdot$\,\S\,\ref{sec:separable}}
  {Pour quelle valeur de $x$ l'état (non normalisé) $\ket{00}+2\ket{01}+3\ket{10}+x\ket{11}$ est-il séparable ?}
  {$x=6$}
  {$x=-6$}
  {$x=\frac16$}
  {$x=5$}
  {A}
\qcmx
  {Critère $D=c_{00}c_{11}-c_{01}c_{10}=x-2\cdot3=0$. On obtient bien un produit : $(\ket0+3\ket1)\otimes(\ket0+2\ket1)=\ket{00}+2\ket{01}+3\ket{10}+6\ket{11}$.}
  {Erreur de signe : $D=x-6$ s'annule pour $x=+6$ ; pour $x=-6$, $D=-12\neq0$ et l'état est intriqué.}
  {On a inversé le rapport : $c_{00}c_{11}=c_{01}c_{10}$ donne $x=\frac{c_{01}c_{10}}{c_{00}}=6$, et non $\frac{c_{00}}{c_{01}c_{10}}$.}
  {On a additionné ($2+3$) au lieu de multiplier : le critère de séparabilité porte sur des \emph{produits} de coefficients.}

% ---- Question 85 (niveau 3, correct : B)
\qcmq[2]{Chap.~\ref{chap:bell}\,$\cdot$\,\S\,\ref{sec:bell}}
  {Comment s'écrit $\ket{00}$ dans la base de Bell ?}
  {$\frac{1}{\sqrt2}\big(\ket{\Phi^+}-\ket{\Phi^-}\big)$}
  {$\frac{1}{\sqrt2}\big(\ket{\Phi^+}+\ket{\Phi^-}\big)$}
  {$\frac{1}{\sqrt2}\big(\ket{\Phi^+}+\ket{\Psi^+}\big)$}
  {$\ket{\Phi^+}$}
  {B}
\qcmx
  {Erreur de signe : cette combinaison vaut $\ket{11}$, puisque $\ket{\Phi^+}-\ket{\Phi^-}=\sqrt2\,\ket{11}$.}
  {$\ket{\Phi^+}+\ket{\Phi^-}=\frac{1}{\sqrt2}\cdot2\ket{00}=\sqrt2\,\ket{00}$. Une mesure de Bell sur $\ket{00}$ donne donc $\Phi^+$ ou $\Phi^-$ avec la probabilité $\frac12$ chacun (jamais $\Psi^\pm$).}
  {Cette combinaison vaut $\ket{+}\otimes\ket{+}=\frac12(\ket{00}+\ket{01}+\ket{10}+\ket{11})$ : les deux états de Bell « pairs » et « impairs » s'y ajoutent.}
  {On a confondu $\ket{00}$, état séparable, avec $\ket{\Phi^+}=\frac{1}{\sqrt2}(\ket{00}+\ket{11})$, qui contient aussi $\ket{11}$ et est intriqué.}

% ---- Question 86 (niveau 3, correct : A)
\qcmq[1]{Chap.~\ref{chap:bell}\,$\cdot$\,\S\,\ref{sec:circuit-bell}}
  {Le générateur de Bell du cours est $G=\mathrm{CNOT}\cdot(\mathrm{Id}\otimes H)$ : $H$ sur le second qubit $B$, puis CNOT de contrôle $B$ et de cible $A$. Quel circuit permet de \emph{mesurer} une paire dans la base de Bell (on le place avant une mesure en base $Z$) ?}
  {D'abord le CNOT (contrôle $B$, cible $A$), puis $H$ sur $B$}
  {D'abord $H$ sur $B$, puis le CNOT}
  {D'abord le CNOT, puis $H$ sur $A$ (la cible)}
  {Aucun circuit : on mesure directement les deux qubits en base $Z$}
  {A}
\qcmx
  {C'est l'inverse $G^\dagger=(\mathrm{Id}\otimes H)\,\mathrm{CNOT}$ : les portes sont leurs propres inverses et l'ordre s'inverse. Il envoie $\Phi^+,\Phi^-,\Psi^+,\Psi^-$ sur $\ket{00},\ket{01},\ket{10},\ket{11}$ (à une phase près).}
  {C'est $G$ lui-même, qui \emph{fabrique} des états de Bell : on garde le même ordre au lieu de l'inverser, et la mesure ne distingue plus les quatre états.}
  {Mauvais qubit : pour $\Phi^+$ on trouve $\ket{+}\otimes\ket{+}$, dont les deux bits sont aléatoires. Le $H$ doit agir sur $B$, qui a reçu le $H$ dans $G$.}
  {$\Phi^+$ et $\Phi^-$ (comme $\Psi^+$ et $\Psi^-$) donnent les mêmes résultats en base $Z$ : cette base ne voit pas le signe relatif, il faut d'abord défaire l'intrication.}

% ---- Question 87 (niveau 3, correct : B)
\qcmq[2]{Chap.~\ref{chap:bell}\,$\cdot$\,\S\,\ref{sec:trace}}
  {Quelle est la matrice de densité réduite $\rho_A=\operatorname{Tr}_B\ket\Psi\bra\Psi$ du premier qubit pour $\ket{\Psi}=\frac12\ket{00}+\frac{\sqrt3}{2}\ket{11}$ ?}
  {$\begin{pmatrix}\frac14&\frac{\sqrt3}{4}\\\frac{\sqrt3}{4}&\frac34\end{pmatrix}$}
  {$\begin{pmatrix}\frac14&0\\0&\frac34\end{pmatrix}$}
  {$\begin{pmatrix}\frac12&0\\0&\frac12\end{pmatrix}$}
  {$\begin{pmatrix}\frac12&0\\0&\frac{\sqrt3}{2}\end{pmatrix}$}
  {B}
\qcmx
  {C'est la matrice $\ket\chi\bra\chi$ de l'état pur $\ket\chi=\frac12\ket0+\frac{\sqrt3}2\ket1$ : on a gardé les cohérences, comme si $A$ était seul ; l'intrication avec $B$ les détruit.}
  {Avec $C=\begin{pmatrix}c_{00}&c_{01}\\c_{10}&c_{11}\end{pmatrix}=\begin{pmatrix}\frac12&0\\0&\frac{\sqrt3}2\end{pmatrix}$, $\rho_A=CC^\dagger=\mathrm{diag}\big(\frac14,\frac34\big)$ : état mixte, de pureté $\frac{10}{16}=1-2|D|^2$ avec $D=\frac{\sqrt3}{4}$.}
  {C'est la réponse pour un état de Bell (intrication maximale) ; ici les poids $\frac14$ et $\frac34$ diffèrent : l'intrication est partielle.}
  {On a mis les amplitudes au lieu de leurs carrés : la trace vaudrait $\frac{1+\sqrt3}{2}\neq1$.}

% ---- Question 88 (niveau 3, correct : A)
\qcmq[2]{Chap.~\ref{chap:bell}\,$\cdot$\,\S\,\ref{sec:info}}
  {Deux qubits sont dans le mélange statistique $\rho_{AB}=\frac12\ket{00}\bra{00}+\frac12\ket{11}\bra{11}$ (corrélations purement classiques). Que vaut $H(A|B)=S(\rho_{AB})-S(\rho_B)$, en bits ?}
  {$0$}
  {$-1$}
  {$1$}
  {$2$}
  {A}
\qcmx
  {$\rho_{AB}$ a pour valeurs propres $\frac12,\frac12,0,0$, donc $S(\rho_{AB})=1$ ; et $\rho_B=\frac12\mathrm{Id}$ donne $S(\rho_B)=1$. Ainsi $H(A|B)=0$ : connaître $B$ fixe $A$. Une valeur négative est propre à l'intrication.}
  {C'est la valeur pour l'état de Bell \emph{pur} $\ket{\Phi^+}$ : mêmes résultats en base $Z$, mais $S(\rho_{AB})=0$ car l'état global est pur. Ici l'état est mixte.}
  {On a pris $H(A|B)=H(A)$, comme si $B$ ne disait rien sur $A$ ; or les deux bits sont toujours égaux.}
  {On a additionné les entropies des deux qubits ($1+1$), valeur de deux bits indépendants : ce n'est pas une entropie conditionnelle.}

% ---- Question 89 (niveau 3, correct : B)
\qcmq[1]{Chap.~\ref{chap:bell}\,$\cdot$\,\S\,\ref{sec:teleportation}}
  {Dans la téléportation, avant de recevoir les deux bits $(a,q)$ d'Alice, dans quel état se trouve le qubit de Bob ?}
  {Dans l'état pur $\alpha\ket0+\beta\ket1$, déjà téléporté}
  {Dans l'état complètement mélangé $\rho_B=\frac12\mathrm{Id}$, indépendant de $\alpha$ et $\beta$}
  {Dans le mélange $|\alpha|^2\ket0\bra0+|\beta|^2\ket1\bra1$}
  {Dans l'état pur $\ket{0}$, son état initial}
  {B}
\qcmx
  {Ce serait une communication instantanée. L'état $\alpha\ket0+\beta\ket1$ n'apparaît qu'\emph{après} la correction, qui dépend de $(a,q)$.}
  {Les quatre kets $\ket{\phi_{aq}}$ sont équiprobables ($\frac14$) et $\frac14\sum_{a,q}\ket{\phi_{aq}}\bra{\phi_{aq}}=\frac12\mathrm{Id}$. Bob n'apprend rien avant le message classique : rien ne voyage plus vite que la lumière.}
  {Bob connaîtrait alors les statistiques de $\ket\varphi$ sans aucun message : or les probabilités d'Alice ($\frac14$ pour chaque résultat) ne dépendent pas de $\alpha,\beta$.}
  {Le qubit de Bob est la moitié de la paire $\ket{\Phi^+}$ : pris seul, il n'a pas d'état pur ($\rho_B=\frac12\mathrm{Id}$ avant comme après la mesure d'Alice, tant que Bob ignore le résultat).}

% ---- Question 90 (niveau 3, correct : D)
\qcmq[1]{Chap.~\ref{chap:bell}\,$\cdot$\,\S\,\ref{sec:chsh}}
  {Au jeu CHSH (gain si $a\oplus b=x\cdot y$), quelles sont les meilleures probabilités de gain avec une stratégie classique, même avec du hasard partagé, et avec une paire $\ket{\Phi^+}$ mesurée dans des bases bien choisies ?}
  {$\frac12$ en classique ; $1$ avec l'intrication}
  {$\frac34$ dans les deux cas}
  {$1$ en classique si Alice et Bob s'accordent à l'avance ; aucun avantage quantique}
  {$\frac34$ en classique ; $\cos^2\frac\pi8\approx0{,}85$ avec l'intrication}
  {D}
\qcmx
  {Classique : répondre toujours $0$ donne déjà $\frac34$. Quantique : l'intrication ne permet pas de gagner à tous les coups (borne de Tsirelson $\cos^2\frac\pi8$).}
  {Si l'intrication n'apportait rien, on n'aurait pas de violation de l'inégalité de Bell ; l'avantage quantique ($0{,}85>0{,}75$) est précisément ce que l'expérience de Bell met en évidence.}
  {S'accorder à l'avance (ou partager du hasard) ne suffit pas : aucune stratégie déterministe ne gagne les quatre parties, d'après l'argument de parité.}
  {Une stratégie déterministe perd au moins une des quatre parties (la somme modulo $2$ des conditions donne $0=1$) ; avec $\Phi^+$, chaque partie est gagnée avec la probabilité $\cos^2\frac\pi8=\frac{2+\sqrt2}{4}$.}

\clearpage
%% ======================================================================
\subsection{Niveau 4 : maîtrise (questions 91 à 100)}\label{sec:qcm-n4}
%% ======================================================================
\noindent{\color{insagris}\small Niveau de difficulté : Maîtrise\quad $\bullet\bullet\bullet\bullet$}\par\smallskip
% ---- Question 91 (niveau 4, correct : C)
\qcmq[2]{Chap.~\ref{chap:histoire}\,$\cdot$\,\S\,\ref{sec:superposition}}
  {Une particule du puits infini est préparée à $t=0$ dans l'état $\Psi=\frac{1}{\sqrt2}\big(\psi_1+\psi_2\big)$, avec $E_n=n^2E_1$. Au bout de quelle durée minimale la densité de probabilité $|\Psi(x,t)|^2$ retrouve-t-elle sa forme initiale ?}
  {$T=\dfrac{2\pi\hbar}{E_1}$}
  {$T=\dfrac{2\pi\hbar}{5E_1}$}
  {$T=\dfrac{2\pi\hbar}{3E_1}$}
  {$T=\dfrac{\pi\hbar}{3E_1}$}
  {C}
\qcmx
  {C'est la période de la phase $\ee^{-\ii E_1t/\hbar}$ de $\psi_1$ seul ; or les phases globales ne se voient pas dans $|\Psi|^2$ : seul compte l'écart $E_2-E_1=3E_1$.}
  {On a pris la somme $E_1+E_2=5E_1$ ; c'est la \emph{différence} des énergies (phase relative) qui intervient.}
  {$|\Psi|^2=\frac12(\psi_1^2+\psi_2^2)+\psi_1\psi_2\cos\frac{(E_2-E_1)t}{\hbar}$ : seul le terme croisé dépend du temps, à la pulsation $\omega=\frac{E_2-E_1}{\hbar}=\frac{3E_1}{\hbar}$, donc $T=\frac{2\pi}{\omega}$.}
  {C'est la demi-période : à $t=\frac T2$ le cosinus vaut $-1$ et la densité est le symétrique de l'état initial ($x\to a-x$), pas encore sa forme de départ.}

% ---- Question 92 (niveau 4, correct : D)
\qcmq[1]{Chap.~\ref{chap:dirac}\,$\cdot$\,\S\,\ref{sec:projecteurs}, \ref{sec:qubit}}
  {Un expérimentateur reçoit un qubit préparé soit dans $\ket{0}$, soit dans $\ket{+}$. Laquelle de ces affirmations est correcte ?}
  {Une mesure en base $Z$ les distingue : $\ket0$ donne $0$ et $\ket+$ donne $1$}
  {Une mesure en base $X$ les distingue : $\ket+$ donne $+$ et $\ket0$ donne $-$}
  {On les distingue en répétant la mesure sur le même qubit}
  {Aucune mesure ne permet de les distinguer à coup sûr, car $\braket{0}{+}=\frac{1}{\sqrt2}\neq0$}
  {D}
\qcmx
  {$\ket+$ donne $0$ ou $1$ avec la probabilité $\frac12$ : le résultat $0$ est possible pour les deux états.}
  {$\ket0=\frac{1}{\sqrt2}(\ket++\ket-)$ donne $+$ ou $-$ avec la probabilité $\frac12$ : le résultat $+$ est possible pour les deux états.}
  {Après une première mesure l'état est effondré : répéter ne donne aucune information nouvelle, et on ne peut pas copier le qubit pour recommencer (non-clonage).}
  {Seuls des états \emph{orthogonaux} sont parfaitement discernables. En base $Z$, $\ket0$ donne toujours $0$ et $\ket+$ donne $0$ une fois sur deux : obtenir $0$ ne tranche pas (obtenir $1$ prouve seulement qu'on avait $\ket+$).}

% ---- Question 93 (niveau 4, correct : C)
\qcmq[1]{Chap.~\ref{chap:operateurs}\,$\cdot$\,\S\,\ref{sec:pur-mixte}}
  {Une source envoie au hasard $\ket{0}$ ou $\ket{1}$ (probabilité $\frac12$ chacun) ; une autre envoie au hasard $\ket{+}$ ou $\ket{-}$ (probabilité $\frac12$ chacun). Peut-on distinguer ces deux sources en mesurant les qubits émis ?}
  {Oui : en base $Z$, la première source donne des résultats prévisibles}
  {Oui : en base $X$, la seconde source donne des résultats prévisibles}
  {Non : les deux ont la même matrice de densité $\frac{\mathrm{Id}}{2}$, donc les mêmes statistiques pour toute mesure}
  {Oui : les états purs envoyés ne sont pas les mêmes}
  {C}
\qcmx
  {Chaque qubit est bien $\ket0$ ou $\ket1$, mais on ignore lequel : le résultat vaut $0$ ou $1$ avec la probabilité $\frac12$, exactement comme avec la seconde source ($|\braket{0}{\pm}|^2=\frac12$).}
  {Même erreur : chaque qubit est $\ket+$ ou $\ket-$, mais on ignore lequel ; le résultat en base $X$ est aléatoire, et la première source donne aussi $+$ ou $-$ avec la probabilité $\frac12$.}
  {$\frac12\ket0\bra0+\frac12\ket1\bra1=\frac12\ket+\bra++\frac12\ket-\bra-=\frac{\mathrm{Id}}2$ (relation de fermeture dans chaque base). Les probabilités de toute mesure ne dépendent que de $\rho$.}
  {Ce qui compte n'est pas la liste des états purs envoyés mais la matrice de densité moyenne : des préparations différentes peuvent avoir la même $\rho$.}

% ---- Question 94 (niveau 4, correct : B)
\qcmq[1]{Chap.~\ref{chap:operateurs}\,$\cdot$\,\S\,\ref{sec:pur-mixte} ; Chap.~\ref{chap:portes}\,$\cdot$\,\S\,\ref{sec:portes-bloch}}
  {Un qubit est dans l'état complètement mélangé $\rho=\frac{\mathrm{Id}}{2}$. Existe-t-il une porte quantique $U$ qui le transforme en un état pur, par exemple $\ket{0}\bra{0}$ ?}
  {Oui : la porte de Hadamard, qui crée des superpositions}
  {Non : $\rho\mapsto U\rho U^\dagger$ conserve les valeurs propres, donc la pureté ; ici $U\frac{\mathrm{Id}}{2}U^\dagger=\frac{\mathrm{Id}}{2}$}
  {Oui : une rotation de Bloch qui amène le centre de la boule au pôle nord}
  {Oui : une mesure en base $Z$, qui donne $\ket0$}
  {B}
\qcmx
  {$H$ envoie un état pur sur un état pur ($\ket0\mapsto\ket+$) et laisse $\frac{\mathrm{Id}}{2}$ invariant : $H\frac{\mathrm{Id}}2H=\frac{\mathrm{Id}}2$ car $H^2=\mathrm{Id}$.}
  {$\operatorname{Tr}\big[(U\rho U^\dagger)^2\big]=\operatorname{Tr}\rho^2$ ($=\frac12$ ici). Sur la boule de Bloch, une porte est une rotation : elle laisse le centre fixe.}
  {Une rotation tourne la boule autour de son centre, qui reste fixe : elle ne peut pas l'amener sur la sphère.}
  {Une mesure n'est pas une porte (elle est non unitaire) ; de plus elle donne $\ket0$ ou $\ket1$ au hasard, pas $\ket0$ à coup sûr.}

% ---- Question 95 (niveau 4, correct : D)
\qcmq[1]{Chap.~\ref{chap:portes}\,$\cdot$\,\S\,\ref{sec:cnot}}
  {Que vaut $(H\otimes H)\,\mathrm{CNOT}_{12}\,(H\otimes H)$, où $\mathrm{CNOT}_{12}$ a pour contrôle le \emph{premier} qubit et pour cible le second ?}
  {$\mathrm{CNOT}_{12}$, inchangé}
  {$\mathrm{CZ}$ (porte $Z$ contrôlée)}
  {$\mathrm{SWAP}$}
  {$\mathrm{CNOT}_{21}$ (contrôle : second qubit ; cible : premier qubit)}
  {D}
\qcmx
  {On a cru que les deux $H$ s'annulent ($H^2=\mathrm{Id}$) ; mais ils encadrent le CNOT au lieu d'être côte à côte : on ne peut pas les simplifier.}
  {C'est le résultat quand $H$ n'encadre que la \emph{cible} : $(\mathrm{Id}\otimes H)\,\mathrm{CNOT}_{12}\,(\mathrm{Id}\otimes H)=\mathrm{CZ}$, car $HXH=Z$. Ici $H$ agit aussi sur le contrôle.}
  {Le SWAP s'obtient avec \emph{trois} CNOT alternés, pas avec un seul CNOT conjugué par $H\otimes H$ ; il n'y a pas de contrôle dans un SWAP.}
  {$H\ket0\bra0H=\ket+\bra+$, $H\ket1\bra1H=\ket-\bra-$ et $HXH=Z$ donnent $\ket+\bra+\otimes I+\ket-\bra-\otimes Z=I\otimes\ket0\bra0+X\otimes\ket1\bra1=\mathrm{CNOT}_{21}$ : dans la base $\{\ket+,\ket-\}$, contrôle et cible échangent leurs rôles (rebond de phase).}

% ---- Question 96 (niveau 4, correct : A)
\qcmq[1]{Chap.~\ref{chap:portes}\,$\cdot$\,\S\,\ref{sec:controlees}}
  {On applique la porte de Toffoli (contrôles : qubits 1 et 2 ; cible : qubit 3) à $\ket{+}\otimes\ket{+}\otimes\ket{0}$, puis on mesure le troisième qubit en base $Z$. Que peut-on dire si l'on obtient $1$ ?}
  {Ce résultat a la probabilité $\frac14$, et les deux premiers qubits sont alors dans l'état $\ket{11}$}
  {Ce résultat a la probabilité $\frac12$, et les deux premiers qubits restent dans $\ket{+}\ket{+}$}
  {Ce résultat a la probabilité $\frac14$, mais les deux premiers qubits restent dans $\ket{+}\ket{+}$}
  {Ce résultat a la probabilité $\frac34$, et les deux premiers qubits sont dans l'état $\frac{1}{\sqrt3}\big(\ket{00}+\ket{01}+\ket{10}\big)$}
  {A}
\qcmx
  {$\mathrm{CCNOT}\ket{x,y,0}=\ket{x,y,xy}$ donne $\frac12\big(\ket{000}+\ket{010}+\ket{100}+\ket{111}\big)$ : seul $\ket{111}$ a un troisième chiffre égal à $1$ (probabilité $\frac14$), et il impose $x=y=1$.}
  {La cible n'est inversée que si $x=y=1$ (probabilité $\frac14$) et non une fois sur deux ; de plus l'état est intriqué, donc mesurer le qubit 3 modifie les qubits 1 et 2.}
  {La probabilité est juste, mais la mesure de la cible n'est pas sans effet sur les contrôles : l'état après la Toffoli est intriqué, donc le résultat $1$ les projette sur $\ket{11}$.}
  {C'est le cas du résultat $0$ (probabilité $\frac34$ : $x,y\in\{00,01,10\}$). Le résultat $1$ est le cas rare, celui du ET vrai ($x=y=1$).}

% ---- Question 97 (niveau 4, correct : A)
\qcmq[1]{Chap.~\ref{chap:bell}\,$\cdot$\,\S\,\ref{sec:bell}}
  {Que donne $(H\otimes H)\ket{\Psi^-}$, avec $\ket{\Psi^-}=\frac{1}{\sqrt2}\big(\ket{01}-\ket{10}\big)$ ?}
  {$-\ket{\Psi^-}$, c'est-à-dire le même état physique}
  {$\ket{\Phi^-}$}
  {$\ket{\Psi^+}$}
  {$\ket{\Phi^+}$}
  {A}
\qcmx
  {Le singulet est invariant par $U\otimes U$ à un facteur près : $(U\otimes U)\ket{\Psi^-}=\det(U)\ket{\Psi^-}$ avec $\det H=-1$. Directement : $\frac{1}{\sqrt2}\big(\ket+\ket--\ket-\ket+\big)=-\ket{\Psi^-}$.}
  {C'est l'image de $\ket{\Psi^+}$, pas de $\ket{\Psi^-}$ : $(H\otimes H)\ket{\Psi^+}=\ket{\Phi^-}$. Le signe $-$ du singulet change le calcul (les termes $\ket{00}$ et $\ket{11}$ se compensent).}
  {$\ket{\Psi^+}$ est symétrique par échange des deux qubits, $\ket{\Psi^-}$ antisymétrique ; $H\otimes H$ commute avec le SWAP et conserve donc cette symétrie.}
  {C'est l'état que $H\otimes H$ laisse invariant ($\ket{\Phi^+}=\frac{1}{\sqrt2}(\ket{++}+\ket{--})$) ; le singulet n'a pas ce comportement.}

% ---- Question 98 (niveau 4, correct : B)
\qcmq[1]{Chap.~\ref{chap:bell}\,$\cdot$\,\S\,\ref{sec:chsh}}
  {Une expérience de type CHSH (avec $S=E_{00}+E_{01}+E_{10}-E_{11}$) mesure $S=2{,}6$. Que peut-on en conclure ?}
  {C'est impossible : on a toujours $|S|\leq2$}
  {Aucun modèle local à résultats prédéterminés ne l'explique, mais c'est compatible avec la mécanique quantique}
  {C'est impossible en mécanique quantique, car $2\sqrt2\approx2{,}4<2{,}6$}
  {Les particules s'échangent de l'information plus vite que la lumière}
  {B}
\qcmx
  {Cette borne ne vaut que pour les modèles classiques (locaux) ; l'expérience et la mécanique quantique la dépassent, jusqu'à $2\sqrt2$.}
  {Classique : $|S|\leq2$ ; quantique : $|S|\leq2\sqrt2\approx2{,}83$ (borne de Tsirelson). Comme $2<2{,}6<2\sqrt2$, l'inégalité de Bell est violée sans dépasser le maximum quantique ; le gain correspondant est $\frac12+\frac{2{,}6}{8}\approx0{,}825$.}
  {Erreur numérique : $2\sqrt2\approx2{,}83$, donc $2{,}6$ est permis.}
  {Les corrélations violent l'inégalité de Bell mais ne permettent aucun signal : la réponse de Bob est uniforme quelle que soit la base choisie par Alice.}

% ---- Question 99 (niveau 4, correct : A)
\qcmq[2]{Chap.~\ref{chap:bell}\,$\cdot$\,\S\,\ref{sec:info}}
  {Pour l'état $\ket{\Psi}=\frac{1}{\sqrt3}\big(\ket{00}+\sqrt2\,\ket{11}\big)$, quelle est l'entropie de von Neumann $S(\rho_A)$ du premier qubit, en bits ?}
  {$\log_2 3-\frac23\approx0{,}92$}
  {$1$}
  {$0$}
  {$0{,}64$}
  {A}
\qcmx
  {$\rho_A=\mathrm{diag}\big(\frac13,\frac23\big)$ (trace partielle), donc $S=-\frac13\log_2\frac13-\frac23\log_2\frac23=\log_23-\frac23\approx1{,}585-0{,}667$. C'est moins que le maximum d'un bit (état de Bell) : l'intrication est partielle.}
  {C'est la valeur pour un état de Bell ($\rho_A=\frac{\mathrm{Id}}{2}$) ; ici les poids $\frac13$ et $\frac23$ diffèrent, donc $S<1$.}
  {$S(\rho_A)=0$ vaut pour un état séparable ; ici $D=\frac{\sqrt2}{3}\neq0$ : l'état est intriqué, $\rho_A$ est mixte et $S>0$ (le couple $AB$ est pur, pas $A$ seul).}
  {On a utilisé le logarithme \emph{népérien} (résultat en nats) ; en bits il faut $\log_2$, ce qui donne $0{,}92$.}

% ---- Question 100 (niveau 4, correct : B)
\qcmq[1]{Chap.~\ref{chap:bell}\,$\cdot$\,\S\,\ref{sec:bell}}
  {Pourquoi peut-on mesurer simultanément les deux « parités » $Z\otimes Z$ et $X\otimes X$ de deux qubits, alors qu'on ne peut pas mesurer simultanément $Z$ et $X$ sur un seul qubit ?}
  {Parce que l'incertitude de Heisenberg ne s'applique pas à deux qubits}
  {Parce que $Z\otimes Z$ et $X\otimes X$ commutent : les deux signes $-1$ de $ZX=-XZ$ se compensent}
  {Parce que les deux qubits sont indépendants}
  {Parce que $Z\otimes Z$ et $X\otimes X$ agissent sur des qubits différents}
  {B}
\qcmx
  {Elle s'applique toujours : $Z\otimes\mathrm{Id}$ et $X\otimes\mathrm{Id}$ ne commutent pas. C'est la commutation de ces deux opérateurs précis qui permet la mesure simultanée.}
  {$(Z\otimes Z)(X\otimes X)=ZX\otimes ZX=(-XZ)\otimes(-XZ)=XZ\otimes XZ=(X\otimes X)(Z\otimes Z)$. Ils ont donc une base propre commune : la base de Bell.}
  {Dans un état de Bell les qubits sont intriqués, donc loin d'être indépendants ; la commutation est une propriété des opérateurs, pas de l'état.}
  {Chacun des deux opérateurs agit sur \emph{les deux} qubits à la fois. (Ce sont $Z\otimes\mathrm{Id}$ et $\mathrm{Id}\otimes X$, sur des qubits différents, qui commutent pour cette raison.)}

\label{qcm:fin}

\clearpage
\ifqcmcorrige

% ======================================================================
\subsection{Clé des réponses et fiche de score}\label{sec:qcm-cle}
% ======================================================================
\noindent\textbf{Cette partie donne les réponses : ne la lisez qu'après avoir répondu à toutes les
questions.}

\paragraph{Clé des réponses.} La lettre de la bonne réponse de la question $n$ se lit à la ligne
qui contient $n$, dans la colonne du rang de $n$ dans la ligne (question~24 : ligne 21--30,
colonne 4).
\begin{center}
\renewcommand{\arraystretch}{1.25}
\setlength{\tabcolsep}{5.5pt}
\begin{tabular}{@{}l*{10}{c}@{}}
\toprule
 & 1 & 2 & 3 & 4 & 5 & 6 & 7 & 8 & 9 & 10 \\
\midrule
1--10 & \qcmkey{1} & \qcmkey{2} & \qcmkey{3} & \qcmkey{4} & \qcmkey{5} & \qcmkey{6} & \qcmkey{7} & \qcmkey{8} & \qcmkey{9} & \qcmkey{10} \\
11--20 & \qcmkey{11} & \qcmkey{12} & \qcmkey{13} & \qcmkey{14} & \qcmkey{15} & \qcmkey{16} & \qcmkey{17} & \qcmkey{18} & \qcmkey{19} & \qcmkey{20} \\
21--30 & \qcmkey{21} & \qcmkey{22} & \qcmkey{23} & \qcmkey{24} & \qcmkey{25} & \qcmkey{26} & \qcmkey{27} & \qcmkey{28} & \qcmkey{29} & \qcmkey{30} \\
31--40 & \qcmkey{31} & \qcmkey{32} & \qcmkey{33} & \qcmkey{34} & \qcmkey{35} & \qcmkey{36} & \qcmkey{37} & \qcmkey{38} & \qcmkey{39} & \qcmkey{40} \\
41--50 & \qcmkey{41} & \qcmkey{42} & \qcmkey{43} & \qcmkey{44} & \qcmkey{45} & \qcmkey{46} & \qcmkey{47} & \qcmkey{48} & \qcmkey{49} & \qcmkey{50} \\
51--60 & \qcmkey{51} & \qcmkey{52} & \qcmkey{53} & \qcmkey{54} & \qcmkey{55} & \qcmkey{56} & \qcmkey{57} & \qcmkey{58} & \qcmkey{59} & \qcmkey{60} \\
61--70 & \qcmkey{61} & \qcmkey{62} & \qcmkey{63} & \qcmkey{64} & \qcmkey{65} & \qcmkey{66} & \qcmkey{67} & \qcmkey{68} & \qcmkey{69} & \qcmkey{70} \\
71--80 & \qcmkey{71} & \qcmkey{72} & \qcmkey{73} & \qcmkey{74} & \qcmkey{75} & \qcmkey{76} & \qcmkey{77} & \qcmkey{78} & \qcmkey{79} & \qcmkey{80} \\
81--90 & \qcmkey{81} & \qcmkey{82} & \qcmkey{83} & \qcmkey{84} & \qcmkey{85} & \qcmkey{86} & \qcmkey{87} & \qcmkey{88} & \qcmkey{89} & \qcmkey{90} \\
91--100 & \qcmkey{91} & \qcmkey{92} & \qcmkey{93} & \qcmkey{94} & \qcmkey{95} & \qcmkey{96} & \qcmkey{97} & \qcmkey{98} & \qcmkey{99} & \qcmkey{100} \\
\bottomrule
\end{tabular}
\end{center}

\paragraph{Fiche de score.} Reportez le nombre de bonnes réponses de chaque niveau.
\begin{center}
\renewcommand{\arraystretch}{1.6}
\begin{tabular}{@{}l l c c c@{}}
\toprule
\textbf{Niveau} & \textbf{Questions} & \textbf{Bonnes réponses} & \textbf{Sur} & \textbf{Pourcentage}\\
\midrule
1. Débutant & 1--30 & \makebox[2.6cm]{\dotfill} & 30 & \makebox[2cm]{\dotfill}\,\%\\
2. Intermédiaire & 31--60 & \makebox[2.6cm]{\dotfill} & 30 & \makebox[2cm]{\dotfill}\,\%\\
3. Avancé & 61--90 & \makebox[2.6cm]{\dotfill} & 30 & \makebox[2cm]{\dotfill}\,\%\\
4. Maîtrise & 91--100 & \makebox[2.6cm]{\dotfill} & 10 & \makebox[2cm]{\dotfill}\,\%\\
\midrule
\textbf{Total} & 1--100 & \makebox[2.6cm]{\dotfill} & 100 & \makebox[2cm]{\dotfill}\,\%\\
\bottomrule
\end{tabular}
\end{center}
Repères : au moins $80\,\%$ à un niveau, il est acquis et on peut passer au suivant ; entre $50$ et
$80\,\%$, revoyez les sections des questions manquées ; en dessous, relisez le chapitre concerné et
son récapitulatif avant de recommencer.

\paragraph{Où revoir ?} Questions classées par chapitre du cours.
\begin{center}
\renewcommand{\arraystretch}{1.3}
\begin{tabular}{@{}l >{\raggedright\arraybackslash}p{5.9cm} >{\raggedright\arraybackslash}p{6.9cm}@{}}
\toprule
\textbf{Chapitre} & \textbf{Titre} & \textbf{Questions}\\
\midrule
Chap.~\ref{chap:histoire} & Histoire et mécanique quantique & 1, 2, 3, 4, 5, 6, 31, 32, 33, 34, 35, 61, 62, 63, 91 \\
Chap.~\ref{chap:dirac} & Notation de Dirac et espace de Hilbert & 7, 8, 9, 10, 11, 12, 36, 37, 38, 39, 40, 64, 65, 66, 67, 68, 92 \\
Chap.~\ref{chap:operateurs} & Opérateurs, matrice de densité, Bloch, produit tensoriel & 13, 14, 15, 16, 17, 18, 41, 42, 43, 44, 45, 46, 69, 70, 71, 72, 73, 74, 75, 93, 94 \\
Chap.~\ref{chap:portes} & Portes quantiques et systèmes multi-qubits & 19, 20, 21, 22, 23, 24, 47, 48, 49, 50, 51, 76, 77, 78, 79, 80, 81, 82, 95, 96 \\
Chap.~\ref{chap:bell} & États quantiques, intrication, Bell & 25, 26, 27, 28, 29, 30, 52, 53, 54, 55, 56, 57, 58, 59, 60, 83, 84, 85, 86, 87, 88, 89, 90, 97, 98, 99, 100 \\
\bottomrule
\end{tabular}
\end{center}

\clearpage
\subsection{Corrigé du niveau 1 : débutant (questions 1 à 30)}\label{sec:qcm-c1}
\qcmcorr{1}
\qcmcorr{2}
\qcmcorr{3}
\qcmcorr{4}
\qcmcorr{5}
\qcmcorr{6}
\qcmcorr{7}
\qcmcorr{8}
\qcmcorr{9}
\qcmcorr{10}
\qcmcorr{11}
\qcmcorr{12}
\qcmcorr{13}
\qcmcorr{14}
\qcmcorr{15}
\qcmcorr{16}
\qcmcorr{17}
\qcmcorr{18}
\qcmcorr{19}
\qcmcorr{20}
\qcmcorr{21}
\qcmcorr{22}
\qcmcorr{23}
\qcmcorr{24}
\qcmcorr{25}
\qcmcorr{26}
\qcmcorr{27}
\qcmcorr{28}
\qcmcorr{29}
\qcmcorr{30}

\clearpage
\subsection{Corrigé du niveau 2 : intermédiaire (questions 31 à 60)}\label{sec:qcm-c2}
\qcmcorr{31}
\qcmcorr{32}
\qcmcorr{33}
\qcmcorr{34}
\qcmcorr{35}
\qcmcorr{36}
\qcmcorr{37}
\qcmcorr{38}
\qcmcorr{39}
\qcmcorr{40}
\qcmcorr{41}
\qcmcorr{42}
\qcmcorr{43}
\qcmcorr{44}
\qcmcorr{45}
\qcmcorr{46}
\qcmcorr{47}
\qcmcorr{48}
\qcmcorr{49}
\qcmcorr{50}
\qcmcorr{51}
\qcmcorr{52}
\qcmcorr{53}
\qcmcorr{54}
\qcmcorr{55}
\qcmcorr{56}
\qcmcorr{57}
\qcmcorr{58}
\qcmcorr{59}
\qcmcorr{60}

\clearpage
\subsection{Corrigé du niveau 3 : avancé (questions 61 à 90)}\label{sec:qcm-c3}
\qcmcorr{61}
\qcmcorr{62}
\qcmcorr{63}
\qcmcorr{64}
\qcmcorr{65}
\qcmcorr{66}
\qcmcorr{67}
\qcmcorr{68}
\qcmcorr{69}
\qcmcorr{70}
\qcmcorr{71}
\qcmcorr{72}
\qcmcorr{73}
\qcmcorr{74}
\qcmcorr{75}
\qcmcorr{76}
\qcmcorr{77}
\qcmcorr{78}
\qcmcorr{79}
\qcmcorr{80}
\qcmcorr{81}
\qcmcorr{82}
\qcmcorr{83}
\qcmcorr{84}
\qcmcorr{85}
\qcmcorr{86}
\qcmcorr{87}
\qcmcorr{88}
\qcmcorr{89}
\qcmcorr{90}

\clearpage
\subsection{Corrigé du niveau 4 : maîtrise (questions 91 à 100)}\label{sec:qcm-c4}
\qcmcorr{91}
\qcmcorr{92}
\qcmcorr{93}
\qcmcorr{94}
\qcmcorr{95}
\qcmcorr{96}
\qcmcorr{97}
\qcmcorr{98}
\qcmcorr{99}
\qcmcorr{100}
\fi
. Un chapitre = une \section.
%
%  Organisation :
%    - les énoncés (niveaux 1 à 4) occupent des pages à part ;
%    - le corrigé (clé des réponses, explications, pièges, fiche de score)
%      commence sur une nouvelle page : on peut imprimer les énoncés seuls ;
%    - chaque question est écrite avec \qcmq puis \qcmx (voir plus bas) : le même
%      code produit l'énoncé (cases à cocher) et, plus loin, le corrigé.
%
%  Bascules (définies dans main.tex ; à placer avant % =====================================================================
%  Chapitre 7 : QCM sur le cours (100 questions, 4 niveaux).
%  Fichier importé par main.tex avec % =====================================================================
%  Chapitre 7 : QCM sur le cours (100 questions, 4 niveaux).
%  Fichier importé par main.tex avec \include{qcm}. Un chapitre = une \section.
%
%  Organisation :
%    - les énoncés (niveaux 1 à 4) occupent des pages à part ;
%    - le corrigé (clé des réponses, explications, pièges, fiche de score)
%      commence sur une nouvelle page : on peut imprimer les énoncés seuls ;
%    - chaque question est écrite avec \qcmq puis \qcmx (voir plus bas) : le même
%      code produit l'énoncé (cases à cocher) et, plus loin, le corrigé.
%
%  Bascules (définies dans main.tex ; à placer avant \include{qcm}) :
%    \qcmrefsfalse     masque la référence « Chap. · § » sous chaque énoncé ;
%    \qcmcorrigefalse  supprime le corrigé du document (QCM à distribuer seul).
%
%  Ce fichier est produit par outils_qcm/gen_qcm.py à partir de la banque de
%  questions (fichiers texte), mais il reste un fichier LaTeX ordinaire que l'on
%  peut modifier à la main. Pour ajouter une question à la main : insérer un
%  couple \qcmq ... \qcmx dans le bon niveau (la numérotation est automatique),
%  puis mettre à jour (1) les bornes « questions a à b » des titres de niveaux,
%  (2) les lignes \qcmcorr{n} du corrigé, (3) le tableau « Où revoir ? ».
% =====================================================================

% Macros de Dirac (déjà définies dans main.tex ; ce bloc évite l'erreur
% « Undefined control sequence \ket » si main.tex est une ancienne version).
\providecommand{\ket}[1]{|#1\rangle}
\providecommand{\bra}[1]{\langle #1|}
\providecommand{\braket}[2]{\langle #1|#2\rangle}

\makeatletter
% bascules : normalement déjà définies dans main.tex (valeurs par défaut : vrai)
\@ifundefined{ifqcmrefs}{\newif\ifqcmrefs\qcmrefstrue}{}
\@ifundefined{ifqcmcorrige}{\newif\ifqcmcorrige\qcmcorrigetrue}{}
\newcounter{qcmq}

% case à cocher (3 mm) pour le crayon
\newcommand{\qcmbox}{{\setlength{\fboxsep}{0pt}\setlength{\fboxrule}{0.5pt}%
  \raisebox{-0.4mm}{\fbox{\rule{0pt}{3mm}\rule{3mm}{0pt}}}}}

% une proposition : case, lettre, texte (retrait suspendu)
\newcommand{\qcm@opt}[2]{\noindent\hangindent=3.2em \hangafter=1
  \hspace*{0.3em}\qcmbox\hspace{0.45em}\textbf{#1)}\hspace{0.4em}#2\par\vspace{2.5pt}}

% \qcmq[colonnes]{ref}{énoncé}{A}{B}{C}{D}{bonne lettre}   (TeX limite les macros à 9 paramètres)
% \qcmx{com. A}{com. B}{com. C}{com. D}  : commentaires, placés juste après \qcmq
%   - com. de la bonne lettre = explication ; com. des autres = piège ;
%   - colonnes = 1 (une proposition par ligne) ou 2 (grille 2x2).
\newcommand{\qcmq}[8][1]{%
  \stepcounter{qcmq}%
  \expandafter\gdef\csname qcm@ref@\theqcmq\endcsname{#2}%
  \expandafter\gdef\csname qcm@rep@\theqcmq\endcsname{#8}%
  \expandafter\gdef\csname qcm@o@\theqcmq @A\endcsname{#4}%
  \expandafter\gdef\csname qcm@o@\theqcmq @B\endcsname{#5}%
  \expandafter\gdef\csname qcm@o@\theqcmq @C\endcsname{#6}%
  \expandafter\gdef\csname qcm@o@\theqcmq @D\endcsname{#7}%
  \par\vspace{7pt plus 3pt minus 2pt}\noindent
  \begin{minipage}{\linewidth}%
    \noindent{\color{insarouge}\bfseries Question~\theqcmq}%
    \ifqcmrefs\hfill{\footnotesize\color{insagris}#2}\fi\par\nopagebreak
    \vspace{2pt}\noindent #3\par\nopagebreak\vspace{3pt}%
    \ifnum#1=2
      \begin{minipage}[t]{0.495\linewidth}\qcm@opt{A}{#4}\end{minipage}\hfill
      \begin{minipage}[t]{0.495\linewidth}\qcm@opt{B}{#5}\end{minipage}\par\vspace{2pt}
      \begin{minipage}[t]{0.495\linewidth}\qcm@opt{C}{#6}\end{minipage}\hfill
      \begin{minipage}[t]{0.495\linewidth}\qcm@opt{D}{#7}\end{minipage}\par
    \else
      \qcm@opt{A}{#4}\qcm@opt{B}{#5}\qcm@opt{C}{#6}\qcm@opt{D}{#7}%
    \fi
  \end{minipage}\par}

\newcommand{\qcmx}[4]{%
  \expandafter\gdef\csname qcm@c@\theqcmq @A\endcsname{#1}%
  \expandafter\gdef\csname qcm@c@\theqcmq @B\endcsname{#2}%
  \expandafter\gdef\csname qcm@c@\theqcmq @C\endcsname{#3}%
  \expandafter\gdef\csname qcm@c@\theqcmq @D\endcsname{#4}}

% lettre de la bonne réponse de la question n (« ? » si la question n'existe pas)
\newcommand{\qcmkey}[1]{\@ifundefined{qcm@rep@#1}{?}{\csname qcm@rep@#1\endcsname}}

% une mauvaise réponse du corrigé (n° de question, lettre) : ignorée si c'est la bonne.
% (pas de \@for : « := » ne fonctionne pas avec le « : » rendu actif par babel-french)
\newcommand{\qcm@wrongone}[2]{%
  \edef\qcm@a{#2}\edef\qcm@b{\csname qcm@rep@#1\endcsname}%
  \ifx\qcm@a\qcm@b\else
    \noindent\hangindent=1.4em\hangafter=1 {\color{insagris}\textbf{#2)}}\enspace
    \csname qcm@o@#1@#2\endcsname\enspace--\enspace\emph{\csname qcm@c@#1@#2\endcsname}\par
  \fi}

% bloc de corrigé de la question n
\newcommand{\qcmcorr}[1]{%
  \par\vspace{6pt plus 2pt minus 2pt}\noindent
  \begin{minipage}{\linewidth}\small
    \noindent{\color{insarouge}\bfseries Question~#1}\enspace
    \textbf{Réponse : \csname qcm@rep@#1\endcsname}\hfill
    {\footnotesize\color{insagris}\csname qcm@ref@#1\endcsname}\par\nopagebreak
    \vspace{1pt}\noindent\hangindent=1.4em\hangafter=1
    \textbf{\csname qcm@rep@#1\endcsname)}\enspace
    \csname qcm@o@#1@\csname qcm@rep@#1\endcsname\endcsname\enspace---\enspace
    \csname qcm@c@#1@\csname qcm@rep@#1\endcsname\endcsname\par
    \qcm@wrongone{#1}{A}\qcm@wrongone{#1}{B}\qcm@wrongone{#1}{C}\qcm@wrongone{#1}{D}%
  \end{minipage}\par}
\makeatother

% =====================================================================
\section{QCM sur le cours}\label{chap:qcm}
% =====================================================================

\subsection{Mode d'emploi}\label{sec:qcm-mode}

Ce QCM compte \textbf{100 questions} sur les chapitres~\ref{chap:histoire} à~\ref{chap:bell},
classées en quatre niveaux. Chaque question propose quatre réponses dont \textbf{une seule} est
exacte : cochez au crayon la case de la proposition choisie.
\ifqcmrefs À droite du numéro de chaque question, une référence (chapitre, section) indique la
partie du cours concernée.\else\ifqcmcorrige{} La référence de chaque question (chapitre, section)
est donnée dans le corrigé.\fi\fi

\begin{center}
\renewcommand{\arraystretch}{1.3}
\begin{tabular}{@{}llc>{\raggedright\arraybackslash}p{8.2cm}@{}}
\toprule
\textbf{Niveau} & \textbf{Questions} & \textbf{Nombre} & \textbf{Ce qui est demandé}\\
\midrule
1. Débutant & 1 à 30 & 30 & définitions, notations, résultats du cours en une étape\\
2. Intermédiaire & 31 à 60 & 30 & petits calculs, propriétés, application directe d'un résultat\\
3. Avancé & 61 à 90 & 30 & calculs en plusieurs étapes, portes, circuits et protocoles\\
4. Maîtrise & 91 à 100 & 10 & raisonnements fins et pièges conceptuels, qui combinent plusieurs notions\\
\bottomrule
\end{tabular}
\end{center}

\begin{itemize}[nosep]
  \item \textbf{Sur papier} : imprimez seulement les pages~\pageref{chap:qcm} à~\pageref{qcm:fin}
        (mode d'emploi et énoncés, sections~\ref{sec:qcm-n1} à~\ref{sec:qcm-n4}).%
        \ifqcmcorrige{} Le corrigé commence sur une page distincte, à la page~\pageref{sec:qcm-cle}.\fi
  \item \textbf{Barème} : 1 point par bonne réponse, 0 sinon (pas de point négatif). Calculatrice
        autorisée, mais rarement utile (quelques logarithmes et racines).
  \item \textbf{Correction} :%
        \ifqcmcorrige{} comparez avec la clé des réponses, puis lisez l'explication de chaque
        question, même réussie. Les « pièges » décrivent l'erreur qui mène à chaque mauvaise réponse.%
        \else{} le corrigé n'est pas inclus dans cette version du document.\fi
\end{itemize}

\bigskip
\noindent Nom : \makebox[6cm]{\dotfill}\hfill Date : \makebox[3.5cm]{\dotfill}\hfill Score : \makebox[1.6cm]{\dotfill}\,/\,100
\bigskip

\begin{remarque}[Conventions des énoncés]
\begin{itemize}[nosep]
  \item $\ket{ab}=\ket{a}_A\otimes\ket{b}_B$ : le premier chiffre est le qubit $A$, le second le
        qubit $B$ ; dans un circuit, le fil du haut est le dernier chiffre.
  \item Le CNOT de la section~\ref{sec:cnot} a pour contrôle le \emph{second} chiffre ; les portes
        contrôlées de la section~\ref{sec:controlees} ont pour contrôle le \emph{premier}. Chaque
        énoncé précise le contrôle.
  \item Sans précision, une mesure est faite dans la base de calcul $Z$. Les logarithmes des
        entropies sont en base $2$ (résultats en bits).
  \item $\ket{\pm}=\frac{1}{\sqrt2}(\ket0\pm\ket1)$ et $\ket{{\pm}\ii}=\frac{1}{\sqrt2}(\ket0\pm\ii\ket1)$.
\end{itemize}
\end{remarque}

\clearpage
%% ======================================================================
\subsection{Niveau 1 : débutant (questions 1 à 30)}\label{sec:qcm-n1}
%% ======================================================================
\noindent{\color{insagris}\small Niveau de difficulté : Débutant\quad $\bullet\circ\circ\circ$}\par\smallskip
% ---- Question 1 (niveau 1, correct : D)
\qcmq[2]{Chap.~\ref{chap:histoire}\,$\cdot$\,\S\,\ref{sec:dualite}}
  {En 1905, qui explique l'effet photo-électrique en postulant que la lumière est faite de quanta d'énergie $E=hf$, les photons ?}
  {Louis de Broglie}
  {Erwin Schrödinger}
  {Max Planck}
  {Albert Einstein}
  {D}
\qcmx
  {Il propose en 1924 une onde de matière ($\lambda=h/p$) pour les particules comme l'électron ; l'idée du photon date de 1905.}
  {Il écrit en 1926 l'équation d'onde de la matière, plus de vingt ans après l'article sur les photons.}
  {Il introduit le quantum d'énergie en 1900 pour le rayonnement du corps noir ; c'est Einstein qui en fait une propriété de la lumière elle-même.}
  {Il postule que la lumière échange son énergie par paquets $E=hf=\hbar\omega$ : c'est la première face de la dualité onde--corpuscule.}

% ---- Question 2 (niveau 1, correct : A)
\qcmq[1]{Chap.~\ref{chap:histoire}\,$\cdot$\,\S\,\ref{sec:born}}
  {Selon l'interprétation de Max Born, que représente $|\Psi(x,t)|^2$ ?}
  {La densité de probabilité de présence de la particule en $x$ à l'instant $t$}
  {L'énergie cinétique de la particule}
  {L'amplitude de l'onde de matière}
  {La vitesse de la particule}
  {A}
\qcmx
  {$|\Psi(x,t)|^2\,\dd x$ est la probabilité de trouver la particule entre $x$ et $x+\dd x$ ; $\Psi$ lui-même n'est pas mesurable.}
  {L'énergie est associée à l'opérateur hamiltonien $\Hop$ ; $|\Psi|^2$ n'a pas la dimension d'une énergie.}
  {$\Psi$ est l'amplitude (complexe) ; $|\Psi|^2$ en est le carré du module.}
  {La vitesse s'obtient avec l'opérateur $\hat p=-\ii\hbar\,\partial_x$, pas avec $|\Psi|^2$.}

% ---- Question 3 (niveau 1, correct : D)
\qcmq[2]{Chap.~\ref{chap:histoire}\,$\cdot$\,\S\,\ref{sec:dualite}}
  {Une particule de quantité de mouvement $p$ est associée, d'après de Broglie, à une onde de longueur d'onde $\lambda$. Quelle relation est correcte ?}
  {$\lambda=hp$}
  {$\lambda=p/h$}
  {$\lambda=\hbar/p$}
  {$\lambda=h/p$}
  {D}
\qcmx
  {Produit au lieu du quotient : $\lambda$ grandirait avec $p$, le contraire de ce qu'on observe.}
  {Quotient inversé : $p/h=1/\lambda$ s'exprime en m$^{-1}$, pas en mètres.}
  {Confusion entre $h$ et $\hbar=h/2\pi$ : $\hbar/p=\lambda/2\pi$ est l'inverse du nombre d'onde $k$.}
  {Relation de de Broglie : plus la particule est rapide ou lourde, plus $\lambda$ est petite (balle de $0{,}1$~kg à $10$~m/s : $\lambda\approx10^{-34}$~m, inobservable).}

% ---- Question 4 (niveau 1, correct : C)
\qcmq[2]{Chap.~\ref{chap:histoire}\,$\cdot$\,\S\,\ref{sec:construction}}
  {Quelle est l'équation de Schrödinger dépendante du temps pour un état $\Psi$ et un hamiltonien $\Hop$ ?}
  {$-\ii\hbar\,\partial_t\Psi=\Hop\Psi$}
  {$\hbar\,\partial_t\Psi=\Hop\Psi$}
  {$\ii\hbar\,\partial_t\Psi=\Hop\Psi$}
  {$\ii\hbar\,\partial_x\Psi=\Hop\Psi$}
  {C}
\qcmx
  {Erreur de signe : la solution serait $\ee^{+\ii\Hop t/\hbar}\Psi(0)$, qui décrit l'évolution à rebours du temps.}
  {Sans le facteur $\ii$, les solutions sont des exponentielles réelles (croissantes ou décroissantes) et la norme n'est pas conservée.}
  {Elle relie l'évolution de l'état à l'énergie ; pour un $\Hop$ indépendant du temps, sa solution est $\Psi(t)=\ee^{-\ii\Hop t/\hbar}\Psi(0)$ (chapitre~\ref{chap:portes}).}
  {C'est la dérivée en $t$ qui décrit l'évolution ; la dérivée en $x$ sert à construire $\hat p=-\ii\hbar\,\partial_x$.}

% ---- Question 5 (niveau 1, correct : A)
\qcmq[1]{Chap.~\ref{chap:histoire}\,$\cdot$\,\S\,\ref{sec:construction}}
  {Quel opérateur représente l'énergie totale d'un système quantique ?}
  {L'hamiltonien $\Hop$}
  {L'opérateur quantité de mouvement $\hat p$}
  {L'opérateur position $\hat x$}
  {L'opérateur identité}
  {A}
\qcmx
  {Somme de l'énergie cinétique $-\frac{\hbar^2}{2m}\partial_x^2$ et de l'énergie potentielle $V$ ; ses valeurs propres sont les énergies mesurables.}
  {Il représente l'impulsion ; l'énergie cinétique est $\hat p^2/2m$, et il faut y ajouter $V$.}
  {C'est la multiplication par $x$ : il donne la position, pas l'énergie.}
  {Il vaut $1$ pour tout état : il ne contient aucune information sur l'énergie.}

% ---- Question 6 (niveau 1, correct : C)
\qcmq[1]{Chap.~\ref{chap:histoire}\,$\cdot$\,\S\,\ref{sec:histoire}}
  {Dans quel ordre chronologique ces avancées ont-elles eu lieu ?}
  {de Broglie $\to$ Planck $\to$ Schrödinger $\to$ Bell}
  {Planck $\to$ Schrödinger $\to$ de Broglie $\to$ Bell}
  {Planck (quantum d'énergie) $\to$ de Broglie (onde de matière) $\to$ Schrödinger (équation d'onde) $\to$ Bell (inégalités)}
  {Planck $\to$ de Broglie $\to$ Bell $\to$ Schrödinger}
  {C}
\qcmx
  {Le quantum d'énergie de Planck (1900) précède l'onde de matière de de Broglie (1924).}
  {L'équation de Schrödinger (1926) généralise l'idée d'onde de matière de de Broglie (1924) : elle vient après.}
  {Planck 1900, de Broglie 1924, Schrödinger 1926, Bell 1964.}
  {Les inégalités de Bell (1964) viennent près de quarante ans après l'équation de Schrödinger.}

% ---- Question 7 (niveau 1, correct : C)
\qcmq[1]{Chap.~\ref{chap:dirac}\,$\cdot$\,\S\,\ref{sec:dirac}}
  {En notation de Dirac, comment appelle-t-on $\ket{\Psi}$ et comment le représente-t-on pour un qubit ?}
  {Un bra, représenté par un vecteur ligne à deux composantes complexes}
  {Un ket, représenté par un nombre complexe}
  {Un ket, représenté par un vecteur colonne à deux composantes complexes}
  {Un bra, représenté par une matrice $2\times2$}
  {C}
\qcmx
  {Le bra $\bra{\Psi}$ est le conjugué transposé du ket : un vecteur ligne, mais ce n'est pas la notation $\ket{\Psi}$.}
  {Un scalaire ne peut pas décrire une superposition : il faut deux amplitudes $\alpha$ et $\beta$.}
  {Le ket est le vecteur colonne $\binom{\alpha}{\beta}$ ; son bra $\bra{\Psi}$ est le vecteur ligne conjugué.}
  {Une matrice est un opérateur ; le bra est une matrice ligne $1\times2$.}

% ---- Question 8 (niveau 1, correct : B)
\qcmq[2]{Chap.~\ref{chap:dirac}\,$\cdot$\,\S\,\ref{sec:qubit}}
  {Quelle condition les amplitudes $\alpha$ et $\beta$ de $\alpha\ket{0}+\beta\ket{1}$ doivent-elles vérifier pour que l'état soit normalisé ?}
  {$\alpha+\beta=1$}
  {$|\alpha|^2+|\beta|^2=1$}
  {$\alpha^2+\beta^2=1$}
  {$|\alpha|+|\beta|=1$}
  {B}
\qcmx
  {On somme les amplitudes (complexes, parfois négatives) ; ce sont leurs modules au carré qui sont des probabilités.}
  {Les probabilités de mesure $P(0)=|\alpha|^2$ et $P(1)=|\beta|^2$ doivent sommer à $1$.}
  {Oubli du module : pour $\alpha=\frac{1}{\sqrt2}$ et $\beta=\frac{\ii}{\sqrt2}$, $\alpha^2+\beta^2=0$ alors que l'état est normalisé.}
  {Somme des modules au lieu des carrés : $\alpha=\beta=\frac{1}{\sqrt2}$ est normalisé, mais $|\alpha|+|\beta|=\sqrt2$.}

% ---- Question 9 (niveau 1, correct : C)
\qcmq[1]{Chap.~\ref{chap:dirac}\,$\cdot$\,\S\,\ref{sec:projecteurs}}
  {On mesure dans la base $Z$ un qubit préparé dans $\ket{+}=\frac{1}{\sqrt2}(\ket{0}+\ket{1})$. Que se passe-t-il ?}
  {On lit $+$ et le qubit reste dans $\ket{+}$}
  {On lit toujours $0$}
  {On lit $0$ ou $1$, chacun avec la probabilité $\frac12$, et l'état devient $\ket{0}$ ou $\ket{1}$}
  {On lit $0$ et $1$ à la fois}
  {C}
\qcmx
  {Une mesure en base $Z$ ne donne que $0$ ou $1$ ; elle modifie l'état.}
  {Les deux amplitudes ont le même module : $P(0)=P(1)=\frac12$.}
  {Postulat de la mesure : la superposition s'effondre sur l'état propre correspondant au résultat.}
  {Chaque mesure donne un seul résultat ; la superposition ne s'observe qu'à travers les statistiques.}

% ---- Question 10 (niveau 1, correct : B)
\qcmq[2]{Chap.~\ref{chap:dirac}\,$\cdot$\,\S\,\ref{sec:qubit}}
  {Que vaut le produit scalaire $\braket{0}{1}$ ?}
  {$1$}
  {$0$}
  {$\frac{1}{\sqrt2}$}
  {$\ket{01}$}
  {B}
\qcmx
  {C'est la valeur de $\braket{0}{0}$ et de $\braket{1}{1}$ (vecteurs normés), pas de $\braket{0}{1}$.}
  {$\ket{0}$ et $\ket{1}$ sont orthogonaux : $\begin{pmatrix}1&0\end{pmatrix}\begin{pmatrix}0\\1\end{pmatrix}=0$.}
  {C'est la valeur de $\braket{0}{+}$, pas de $\braket{0}{1}$.}
  {Confusion avec le produit tensoriel : un produit scalaire est un nombre, pas un ket.}

% ---- Question 11 (niveau 1, correct : D)
\qcmq[2]{Chap.~\ref{chap:dirac}\,$\cdot$\,\S\,\ref{sec:qubit}}
  {Un qubit est dans l'état $\ket{-}=\frac{1}{\sqrt2}(\ket{0}-\ket{1})$. Quelle est la probabilité de mesurer $1$ en base $Z$ ?}
  {$0$}
  {$-\frac12$}
  {$\frac{1}{\sqrt2}$}
  {$\frac12$}
  {D}
\qcmx
  {Le signe $-$ ne « retire » pas $\ket{1}$ de la superposition.}
  {Une probabilité est toujours positive : on prend le module au carré de l'amplitude.}
  {C'est l'amplitude en module, pas la probabilité (il faut élever au carré).}
  {$P(1)=\left|-\frac{1}{\sqrt2}\right|^2=\frac12$ : le signe est une phase relative, qui disparaît dans le module.}

% ---- Question 12 (niveau 1, correct : B)
\qcmq[2]{Chap.~\ref{chap:dirac}\,$\cdot$\,\S\,\ref{sec:projecteurs}}
  {Quel est le projecteur sur l'état $\ket{0}$ ?}
  {$\ket{0}\bra{1}=\begin{pmatrix}0&1\\0&0\end{pmatrix}$}
  {$\hat P_0=\ket{0}\bra{0}=\begin{pmatrix}1&0\\0&0\end{pmatrix}$}
  {$\braket{0}{0}=1$}
  {$\begin{pmatrix}0&0\\0&1\end{pmatrix}$}
  {B}
\qcmx
  {Cet opérateur envoie $\ket{1}$ sur $\ket{0}$ ; son carré est nul : ce n'est pas un projecteur.}
  {$\hat P_0\ket{\Psi}=\alpha\ket{0}$ : il garde la composante sur $\ket{0}$ ; $\hat P_0^2=\hat P_0$.}
  {C'est un nombre, pas un opérateur.}
  {C'est $\ket{1}\bra{1}$, le projecteur sur $\ket{1}$.}

% ---- Question 13 (niveau 1, correct : A)
\qcmq[1]{Chap.~\ref{chap:operateurs}\,$\cdot$\,\S\,\ref{sec:bloch}}
  {Où se trouvent les états $\ket{0}$ et $\ket{1}$ sur la sphère de Bloch ?}
  {Aux deux pôles, Nord et Sud}
  {Aux deux extrémités de l'axe $x$}
  {Aux deux extrémités de l'axe $y$}
  {En deux points voisins de l'équateur}
  {A}
\qcmx
  {$\ket{0}$ au pôle Nord ($\theta=0$), $\ket{1}$ au pôle Sud ($\theta=\pi$) : deux états orthogonaux sont diamétralement opposés.}
  {Ce sont les places de $\ket{+}$ et $\ket{-}$.}
  {Ce sont les places de $\ket{{+}\ii}$ et $\ket{{-}\ii}$.}
  {Des états orthogonaux sont antipodaux, jamais voisins.}

% ---- Question 14 (niveau 1, correct : A)
\qcmq[2]{Chap.~\ref{chap:operateurs}\,$\cdot$\,\S\,\ref{sec:tenseur}}
  {Quelle est la dimension de l'espace de Hilbert de $n$ qubits ?}
  {$2^n$}
  {$2n$}
  {$n^2$}
  {$2^n-1$}
  {A}
\qcmx
  {Le produit tensoriel multiplie les dimensions : $2\times2\times\cdots\times2$ ; deux qubits : $4$, trois qubits : $8$.}
  {On additionne les dimensions au lieu de les multiplier ; pour $n=3$ on trouverait $6$ au lieu de $8$.}
  {Juste pour $n=2$ par coïncidence seulement : pour $n=3$ on trouverait $9$ au lieu de $8$.}
  {Il y a bien $2^n$ vecteurs dans la base de calcul, $\ket{0\cdots0}$ compris.}

% ---- Question 15 (niveau 1, correct : C)
\qcmq[1]{Chap.~\ref{chap:operateurs}\,$\cdot$\,\S\,\ref{sec:rho}}
  {Que vaut toujours la trace d'une matrice de densité $\rho$ ?}
  {$0$}
  {$1$ pour un état pur, et moins de $1$ pour un état mixte}
  {$1$}
  {$2$, la dimension de l'espace}
  {C}
\qcmx
  {C'est la trace des matrices de Pauli, pas d'une matrice de densité.}
  {Confusion avec $\operatorname{Tr}(\rho^2)$ : c'est la pureté qui distingue les deux cas, la trace vaut toujours $1$.}
  {C'est la somme des probabilités de mesure (les populations), que l'état soit pur ou mixte.}
  {C'est la trace de l'identité $\operatorname{Tr}\mathrm{Id}=2$, pas celle de $\rho$.}

% ---- Question 16 (niveau 1, correct : B)
\qcmq[1]{Chap.~\ref{chap:operateurs}\,$\cdot$\,\S\,\ref{sec:pur-mixte}}
  {Quelle est la valeur de $\operatorname{Tr}(\rho^2)$ pour un état pur ?}
  {$0$}
  {$1$}
  {$\frac12$}
  {Une valeur strictement inférieure à $1$}
  {B}
\qcmx
  {$\rho^2=0$ imposerait $\rho=0$, ce qui contredit $\operatorname{Tr}\rho=1$.}
  {Un état pur s'écrit $\rho=\ket{\Psi}\bra{\Psi}$, donc $\rho^2=\rho$ et $\operatorname{Tr}\rho^2=\operatorname{Tr}\rho=1$.}
  {C'est la valeur minimale pour un qubit, atteinte par l'état maximalement mixte $\mathrm{Id}/2$.}
  {C'est la caractérisation d'un état mixte.}

% ---- Question 17 (niveau 1, correct : B)
\qcmq[1]{Chap.~\ref{chap:operateurs}\,$\cdot$\,\S\,\ref{sec:rho}}
  {Que représentent les éléments diagonaux $\rho_{00}$ et $\rho_{11}$ de la matrice de densité d'un qubit ?}
  {Les cohérences entre $\ket{0}$ et $\ket{1}$}
  {Les probabilités de mesurer $\ket{0}$ et $\ket{1}$ (les populations)}
  {Les valeurs propres de $\rho$}
  {Les amplitudes $\alpha$ et $\beta$}
  {B}
\qcmx
  {Les cohérences sont les éléments hors diagonale $\rho_{01}$ et $\rho_{10}$.}
  {$\rho_{00}=\bra{0}\rho\ket{0}=P(0)$ et $\rho_{11}=P(1)$, d'où $\operatorname{Tr}\rho=1$.}
  {Seulement si $\rho$ est diagonale dans la base de calcul : $\ket{+}\bra{+}$ a pour diagonale $(\frac12,\frac12)$ mais pour valeurs propres $(1,0)$.}
  {Pour un état pur, $\rho_{00}=|\alpha|^2$ : le module au carré, pas l'amplitude.}

% ---- Question 18 (niveau 1, correct : D)
\qcmq[1]{Chap.~\ref{chap:operateurs}\,$\cdot$\,\S\,\ref{sec:op-dirac}}
  {Pourquoi associe-t-on aux grandeurs mesurables (observables) des opérateurs hermitiens ?}
  {Parce qu'ils conservent la norme des vecteurs}
  {Parce qu'ils ont tous une trace nulle}
  {Parce que leurs valeurs propres valent toutes $\pm1$}
  {Parce que leurs valeurs propres sont réelles, comme les résultats d'une mesure}
  {D}
\qcmx
  {C'est la propriété des opérateurs unitaires (les portes), pas des opérateurs hermitiens.}
  {Vrai pour les matrices de Pauli, faux en général (l'identité est hermitienne, de trace $2$), et sans rapport avec le caractère mesurable.}
  {Vrai pour $X$, $Y$, $Z$ seulement ; l'hamiltonien du puits a les valeurs propres $E_n=n^2E_1$.}
  {Un opérateur hermitien ($A^\dagger=A$) a des valeurs propres réelles et des vecteurs propres orthogonaux ; les résultats possibles sont ses valeurs propres.}

% ---- Question 19 (niveau 1, correct : C)
\qcmq[2]{Chap.~\ref{chap:portes}\,$\cdot$\,\S\,\ref{sec:pauli}}
  {Quelle porte de Pauli réalise l'inversion de bit $\ket{0}\leftrightarrow\ket{1}$ ?}
  {$Z$}
  {$Y$}
  {$X$}
  {$H$}
  {C}
\qcmx
  {$Z$ laisse $\ket{0}$ et change $\ket{1}$ en $-\ket{1}$ : c'est une inversion de phase.}
  {$Y\ket{0}=\ii\ket{1}$ : elle inverse le bit mais ajoute une phase ; la porte d'inversion « pure » est $X$.}
  {$X\ket{0}=\ket{1}$ et $X\ket{1}=\ket{0}$ : c'est le NOT quantique.}
  {$H$ crée des superpositions : $H\ket{0}=\ket{+}$.}

% ---- Question 20 (niveau 1, correct : B)
\qcmq[2]{Chap.~\ref{chap:portes}\,$\cdot$\,\S\,\ref{sec:hadamard}}
  {Que donne la porte de Hadamard appliquée à $\ket{0}$ ?}
  {$\ket{1}$}
  {$\ket{+}=\frac{1}{\sqrt2}(\ket{0}+\ket{1})$}
  {$\ket{-}=\frac{1}{\sqrt2}(\ket{0}-\ket{1})$}
  {$\ket{0}$}
  {B}
\qcmx
  {C'est l'action de $X$.}
  {Superposition à poids égaux : $P(0)=P(1)=\frac12$ ; c'est la première étape de la plupart des algorithmes.}
  {C'est $H\ket{1}$ ; le signe $-$ n'apparaît qu'à partir de $\ket{1}$.}
  {$H$ n'est pas l'identité ; elle ne laisse inchangés ni $\ket{0}$ ni $\ket{1}$ (elle envoie $\ket{+}$ sur $\ket{0}$).}

% ---- Question 21 (niveau 1, correct : D)
\qcmq[2]{Chap.~\ref{chap:portes}\,$\cdot$\,\S\,\ref{sec:portes}}
  {Quelle propriété doit vérifier la matrice $U$ d'une porte quantique ?}
  {Elle est hermitienne : $U^\dagger=U$}
  {Sa trace vaut $1$}
  {Ses coefficients sont réels}
  {Elle est unitaire : $U^\dagger U=\mathrm{Id}$}
  {D}
\qcmx
  {Vrai pour $X$, $Y$, $Z$, $H$ mais faux pour $S$ et $T$ : ce n'est pas la condition générale.}
  {C'est une propriété des matrices de densité ; $\operatorname{Tr}X=0$ et $\operatorname{Tr}\mathrm{Id}=2$.}
  {$Y$ et $S$ sont des portes à coefficients complexes.}
  {Elle conserve la norme ($\braket{\Psi}{\Psi}=1$ reste vrai) et elle est réversible.}

% ---- Question 22 (niveau 1, correct : A)
\qcmq[2]{Chap.~\ref{chap:portes}\,$\cdot$\,\S\,\ref{sec:portes}}
  {Quel est l'inverse d'une porte quantique $U$ ?}
  {$U^{-1}=U^\dagger$ (transposée conjuguée)}
  {$U^{-1}=U^{\mathsf T}$ (transposée seule)}
  {$U^{-1}=U^{*}$ (conjuguée seule)}
  {$U^{-1}=-U$}
  {A}
\qcmx
  {Conséquence de $U^\dagger U=\mathrm{Id}$ ; par exemple $S^{-1}=S^\dagger=\mathrm{diag}(1,-\ii)$.}
  {Oubli de la conjugaison : $S^{\mathsf T}=S=\mathrm{diag}(1,\ii)$ alors que $S^{-1}=\mathrm{diag}(1,-\ii)$.}
  {Faux en général : $Y^*=-Y$ alors que $Y^{-1}=Y$.}
  {Faux : $X^{-1}=X\neq-X$.}

% ---- Question 23 (niveau 1, correct : C)
\qcmq[1]{Chap.~\ref{chap:portes}\,$\cdot$\,\S\,\ref{sec:cnot}}
  {Dans une porte CNOT, que devient le qubit cible lorsque le qubit de contrôle est dans l'état $\ket{0}$ ?}
  {Il est inversé par la porte $X$}
  {Il est remis à $\ket{0}$}
  {Il reste inchangé}
  {Il subit un changement de signe, comme avec $Z$}
  {C}
\qcmx
  {C'est ce qui se passe quand le contrôle vaut $\ket{1}$.}
  {Le CNOT est réversible : remettre la cible à zéro détruirait de l'information.}
  {Avec le contrôle à $\ket{0}$, la porte agit comme l'identité sur la cible.}
  {Ce serait la porte $Z$ contrôlée (CZ), pas le CNOT.}

% ---- Question 24 (niveau 1, correct : A)
\qcmq[2]{Chap.~\ref{chap:portes}\,$\cdot$\,\S\,\ref{sec:cnot}}
  {Le CNOT du cours a pour contrôle le second chiffre du ket et pour cible le premier : $\mathrm{CNOT}\ket{a,b}=\ket{a\oplus b,\,b}$. Que donne-t-il sur $\ket{0,1}$ ?}
  {$\ket{1,1}$}
  {$\ket{0,1}$}
  {$\ket{0,0}$}
  {$\ket{1,0}$}
  {A}
\qcmx
  {Le contrôle $b=1$ inverse la cible : $0\oplus1=1$.}
  {Le contrôle vaut $1$ : la cible est bien inversée.}
  {C'est le contrôle qui serait inversé : le contrôle n'est jamais modifié.}
  {Les deux qubits seraient inversés : seule la cible change.}

% ---- Question 25 (niveau 1, correct : D)
\qcmq[1]{Chap.~\ref{chap:bell}\,$\cdot$\,\S\,\ref{sec:bell}}
  {Quelle propriété caractérise les quatre états de Bell ?}
  {Ils sont séparables : chaque qubit a son propre état}
  {Ils décrivent un seul qubit}
  {Ils ne diffèrent que par une phase globale}
  {Ils sont maximalement intriqués et forment une base orthonormée de l'espace de deux qubits}
  {D}
\qcmx
  {$D=c_{00}c_{11}-c_{01}c_{10}=\pm\frac12\neq0$ : ils ne se factorisent pas.}
  {Ce sont des états de deux qubits (quatre amplitudes sur $\ket{00},\ket{01},\ket{10},\ket{11}$).}
  {Ils sont orthogonaux deux à deux : ce sont quatre états distincts, alors qu'une phase globale ne change pas l'état.}
  {Aucun n'est un produit d'états à un qubit ($|D|=\frac12$, valeur maximale) et ils sont orthogonaux deux à deux.}

% ---- Question 26 (niveau 1, correct : D)
\qcmq[1]{Chap.~\ref{chap:bell}\,$\cdot$\,\S\,\ref{sec:teleportation}}
  {Quel est l'objectif du protocole de téléportation quantique ?}
  {Copier l'état d'un qubit inconnu sur le qubit de Bob}
  {Transmettre une information plus vite que la lumière}
  {Déplacer la particule physique d'Alice vers Bob}
  {Transférer l'état d'un qubit inconnu d'Alice à Bob, avec une paire intriquée et deux bits classiques}
  {D}
\qcmx
  {La copie est interdite (non-clonage) : après le protocole, l'état n'existe plus que chez Bob.}
  {Sans les deux bits classiques, le qubit de Bob est dans l'état $\mathrm{Id}/2$ : il n'a reçu aucune information.}
  {C'est l'état quantique qui est transféré, pas la particule qui le porte.}
  {Le qubit lui-même n'est pas envoyé : son état est reconstruit chez Bob, et détruit chez Alice par la mesure.}

% ---- Question 27 (niveau 1, correct : B)
\qcmq[1]{Chap.~\ref{chap:bell}\,$\cdot$\,\S\,\ref{sec:teleportation}}
  {Dans la téléportation, que doit envoyer Alice à Bob sur un canal classique pour qu'il retrouve l'état ?}
  {Un qubit intriqué}
  {Deux bits classiques, le résultat de sa mesure}
  {Les amplitudes exactes $\alpha$ et $\beta$ de l'état}
  {Un seul bit}
  {B}
\qcmx
  {La paire intriquée est partagée avant le protocole ; on ne l'envoie pas après.}
  {Les quatre résultats $(a,q)$ correspondent aux quatre corrections possibles ($I$, $Z$, $X$, $X$ puis $Z$) : il faut deux bits.}
  {Alice ne connaît pas l'état à téléporter, et ces nombres complexes n'ont pas de description finie en bits.}
  {Un bit ne distingue que deux cas, alors que Bob doit choisir parmi quatre corrections.}

% ---- Question 28 (niveau 1, correct : A)
\qcmq[1]{Chap.~\ref{chap:bell}\,$\cdot$\,\S\,\ref{sec:bell} ; ex.\,\ref{ex:superdense}}
  {Quel est le but du codage superdense ?}
  {Transmettre deux bits classiques en envoyant un seul qubit, grâce à une paire intriquée partagée}
  {Transmettre l'état d'un qubit avec deux bits classiques}
  {Compresser plusieurs qubits en un seul}
  {Chiffrer un message de façon inviolable}
  {A}
\qcmx
  {Alice applique l'une des quatre portes $I$, $Z$, $X$, $XZ$ à son qubit de la paire, l'envoie, et Bob identifie l'état de Bell obtenu.}
  {C'est la téléportation, le protocole inverse.}
  {Le protocole ne compresse pas de qubits : il fait passer deux bits classiques par un canal à un qubit.}
  {Ce n'est pas un protocole de cryptographie : la paire intriquée sert à doubler la capacité du canal.}

% ---- Question 29 (niveau 1, correct : A)
\qcmq[1]{Chap.~\ref{chap:bell}\,$\cdot$\,\S\,\ref{sec:teleportation}}
  {Que dit le théorème de non-clonage ?}
  {Aucune opération ne peut copier parfaitement un état quantique inconnu quelconque}
  {Il est impossible de mesurer un qubit sans le perturber}
  {Il est impossible de créer de l'intrication}
  {Il est impossible de copier un bit classique}
  {A}
\qcmx
  {Une porte $U$ qui copierait $\ket{0}$ et $\ket{1}$ enverrait $\ket{+}\ket{0}$ sur un état intriqué, par linéarité, et non sur $\ket{+}\ket{+}$.}
  {C'est l'effondrement de l'état, un autre principe.}
  {L'intrication se crée très bien (Hadamard puis CNOT).}
  {Le CNOT copie un bit classique : $\ket{0,b}\to\ket{b,b}$.}

% ---- Question 30 (niveau 1, correct : B)
\qcmq[1]{Chap.~\ref{chap:bell}\,$\cdot$\,\S\,\ref{sec:circuit-bell}}
  {Quel enchaînement de portes transforme $\ket{00}$ en l'état de Bell $\ket{\Phi^+}=\frac{1}{\sqrt2}(\ket{00}+\ket{11})$ ?}
  {CNOT, puis Hadamard sur le contrôle}
  {Hadamard sur l'un des qubits, puis CNOT dont le contrôle est ce même qubit}
  {Hadamard sur chacun des deux qubits}
  {SWAP, puis $X$ sur un qubit}
  {B}
\qcmx
  {Le CNOT ne fait rien sur $\ket{00}$ (contrôle à $0$) ; $H$ donne ensuite un état séparable.}
  {$H$ crée $\ket{0}\otimes\ket{+}$, puis le CNOT recopie la superposition : $\frac{1}{\sqrt2}(\ket{00}+\ket{11})$.}
  {On obtient $\ket{+}\otimes\ket{+}$, séparable : des portes locales ne créent pas d'intrication.}
  {$\ket{00}$ devient $\ket{01}$ ou $\ket{10}$, un état de base séparable.}

\clearpage
%% ======================================================================
\subsection{Niveau 2 : intermédiaire (questions 31 à 60)}\label{sec:qcm-n2}
%% ======================================================================
\noindent{\color{insagris}\small Niveau de difficulté : Intermédiaire\quad $\bullet\bullet\circ\circ$}\par\smallskip
% ---- Question 31 (niveau 2, correct : A)
\qcmq[1]{Chap.~\ref{chap:histoire}\,$\cdot$\,\S\,\ref{sec:born}}
  {Quelle condition la fonction d'onde d'une particule doit-elle vérifier sur tout l'espace ?}
  {$\displaystyle\int_{-\infty}^{+\infty}|\Psi(x,t)|^2\,\dd x=1$}
  {$\displaystyle\int_{-\infty}^{+\infty}\Psi(x,t)\,\dd x=1$}
  {$|\Psi(x,t)|^2=1$ pour tout $x$}
  {$\partial_x\Psi(x,t)=0$ en tout point}
  {A}
\qcmx
  {La particule se trouve forcément quelque part : la probabilité totale vaut $1$ (condition de normalisation).}
  {Sans le module au carré : $\Psi$ est complexe et son intégrale n'est pas une probabilité (elle peut être nulle ou complexe).}
  {$|\Psi|^2$ est une densité : si elle valait $1$ partout, son intégrale divergerait.}
  {Une fonction constante n'est pas normalisable ($\int|\Psi|^2\dd x$ diverge) ; la condition porte sur l'intégrale de $|\Psi|^2$.}

% ---- Question 32 (niveau 2, correct : D)
\qcmq[2]{Chap.~\ref{chap:histoire}\,$\cdot$\,\S\,\ref{sec:operateurs}}
  {Quelle est l'expression de l'opérateur quantité de mouvement $\hat p$ en une dimension ?}
  {$\hat p=m\,v$}
  {$\hat p=-\dfrac{\hbar^2}{2m}\dfrac{\partial^2}{\partial x^2}$}
  {$\hat p=+\ii\hbar\,\dfrac{\partial}{\partial x}$}
  {$\hat p=-\ii\hbar\,\dfrac{\partial}{\partial x}$}
  {D}
\qcmx
  {C'est la grandeur classique ; en mécanique quantique $p$ est un opérateur différentiel.}
  {C'est l'énergie cinétique $\hat p^2/2m$.}
  {Erreur de signe : on trouverait $p=-\hbar k$ pour l'onde $\ee^{\ii kx}$.}
  {Convention du cours : $\hat p\,\ee^{\ii kx}=\hbar k\,\ee^{\ii kx}$, donc une onde vers les $x$ croissants a $p=\hbar k>0$.}

% ---- Question 33 (niveau 2, correct : B)
\qcmq[2]{Chap.~\ref{chap:histoire}\,$\cdot$\,\S\,\ref{sec:solutions}}
  {Pour une particule dans un puits de potentiel infini de largeur $a$, comment varie l'énergie $E_n$ avec le numéro $n$ du niveau ?}
  {$E_n=E_1/n$}
  {$E_n=n^2E_1$, avec $E_1=\dfrac{\pi^2\hbar^2}{2ma^2}$}
  {$E_n=nE_1$}
  {$E_n=E_1$ pour tout $n$}
  {B}
\qcmx
  {Les niveaux se resserreraient quand $n$ augmente ; c'est l'inverse pour le puits infini.}
  {$k_n=\frac{n\pi}{a}$ et $E_n=\frac{\hbar^2k_n^2}{2m}$ : les niveaux s'écartent quand $n$ croît ($E_{n+1}-E_n=(2n+1)E_1$).}
  {Niveaux équidistants : ce n'est pas le puits infini, où $E_2=4E_1$ et $E_3=9E_1$.}
  {Il n'y aurait pas de quantification : tous les états auraient la même énergie.}

% ---- Question 34 (niveau 2, correct : C)
\qcmq[2]{Chap.~\ref{chap:histoire}\,$\cdot$\,\S\,\ref{sec:solutions}}
  {Une particule est dans un état stationnaire $\Psi_n$ de l'hamiltonien ($\Hop\Psi_n=E_n\Psi_n$). Que vaut l'écart-type $\sigma_E$ de son énergie ?}
  {$E_n$}
  {$E_n^2$}
  {$0$}
  {$\sqrt{E_n}$}
  {C}
\qcmx
  {C'est la valeur moyenne de l'énergie, pas son écart-type.}
  {C'est la valeur de $\langle\Hop^2\rangle$, avant soustraction de $\langle\Hop\rangle^2$.}
  {$\langle\Hop\rangle=E_n$ et $\langle\Hop^2\rangle=E_n^2$, donc $\sigma_E=\sqrt{E_n^2-E_n^2}=0$ : l'énergie est parfaitement déterminée.}
  {Mélange avec la formule $\sigma=\sqrt{\langle\Hop^2\rangle-\langle\Hop\rangle^2}$ : les deux termes sont égaux à $E_n^2$.}

% ---- Question 35 (niveau 2, correct : C)
\qcmq[2]{Chap.~\ref{chap:histoire}\,$\cdot$\,\S\,\ref{sec:solutions}}
  {Pour un état stationnaire $\Psi_n$ du puits infini $]0,a[$, que vaut la position moyenne $\langle\hat x\rangle$ ?}
  {$a/n$}
  {$a/4$}
  {$a/2$ pour tout $n$}
  {$0$}
  {C}
\qcmx
  {Elle dépendrait de $n$, ce que la symétrie de $|\Psi_n|^2$ interdit.}
  {C'est la position d'un maximum de $|\Psi_2|^2$ ; la moyenne des deux maxima $\frac a4$ et $\frac{3a}{4}$ donne $\frac a2$.}
  {$|\Psi_n|^2=\frac2a\sin^2\frac{n\pi x}{a}$ est symétrique par rapport au centre du puits.}
  {$\Psi_n$ s'annule sur la paroi, mais $\langle\hat x\rangle$ est une moyenne pondérée sur tout le puits.}

% ---- Question 36 (niveau 2, correct : D)
\qcmq[2]{Chap.~\ref{chap:dirac}\,$\cdot$\,\S\,\ref{sec:qubit}}
  {Un qubit est dans l'état $\ket{\Psi}=\frac{1}{\sqrt5}\ket{0}+\frac{2\ii}{\sqrt5}\ket{1}$. Quelle est la probabilité de mesurer $1$ en base $Z$ ?}
  {$\frac25$}
  {$-\frac45$}
  {$\frac15$}
  {$\frac45$}
  {D}
\qcmx
  {On divise le coefficient $2$ par $5$ : on oublie à la fois la racine et le carré.}
  {On élève $\frac{2\ii}{\sqrt5}$ au carré sans conjuguer : $\left(\frac{2\ii}{\sqrt5}\right)^2=-\frac45$ ; une probabilité est $|\cdot|^2$.}
  {C'est $P(0)=\left|\frac{1}{\sqrt5}\right|^2$.}
  {$P(1)=\left|\frac{2\ii}{\sqrt5}\right|^2=\frac45$ : le facteur $\ii$ est une phase relative, invisible en base $Z$.}

% ---- Question 37 (niveau 2, correct : C)
\qcmq[2]{Chap.~\ref{chap:dirac}\,$\cdot$\,\S\,\ref{sec:projecteurs}}
  {Un qubit préparé dans $\ket{0}$ est mesuré dans la base $X$ $\{\ket{+},\ket{-}\}$. Quelle est la probabilité d'obtenir $-$ ?}
  {$0$}
  {$1$}
  {$\frac12$}
  {$\frac{1}{\sqrt2}$}
  {C}
\qcmx
  {$\ket{0}$ est certain en base $Z$, mais pas en base $X$ : une autre base donne une autre loi.}
  {Même confusion dans l'autre sens : $\ket{0}$ n'est pas l'état $\ket{-}$.}
  {$P(-)=|\braket{-}{0}|^2=\left|\frac{1}{\sqrt2}\right|^2=\frac12$.}
  {C'est l'amplitude $\braket{-}{0}$ ; il faut élever son module au carré.}

% ---- Question 38 (niveau 2, correct : C)
\qcmq[2]{Chap.~\ref{chap:dirac}\,$\cdot$\,\S\,\ref{sec:qubit}}
  {Quel facteur $N>0$ normalise l'état $N\big(\ket{0}+(1-\ii)\ket{1}\big)$ ?}
  {$N=\frac{1}{\sqrt2}$}
  {$N=\frac13$}
  {$N=\frac{1}{\sqrt3}$}
  {$N=\frac{1}{\sqrt{1-2\ii}}$}
  {C}
\qcmx
  {On oublie la composante $\ket{0}$ : $|1-\ii|^2=2$, mais la norme au carré est $1+2=3$.}
  {On oublie la racine carrée : $3N^2=1$ donne $N=\frac{1}{\sqrt3}$, pas $\frac13$.}
  {$\braket{\Psi}{\Psi}=N^2\big(1+|1-\ii|^2\big)=N^2(1+2)=3N^2$.}
  {On élève $1-\ii$ au carré sans prendre le module : $(1-\ii)^2=-2\ii$ n'est pas réel ; il faut $|1-\ii|^2=2$.}

% ---- Question 39 (niveau 2, correct : D)
\qcmq[2]{Chap.~\ref{chap:dirac}\,$\cdot$\,\S\,\ref{sec:hilbert}}
  {Soient $\ket{\psi}=\frac{1}{\sqrt2}(\ket{0}+\ii\ket{1})$ et $\ket{\varphi}=\frac{1}{\sqrt2}(\ket{0}-\ket{1})$. Que vaut $\braket{\psi}{\varphi}$ ?}
  {$\frac{1-\ii}{2}$}
  {$\frac12$}
  {$0$}
  {$\frac{1+\ii}{2}$}
  {D}
\qcmx
  {On oublie de conjuguer l'amplitude $\ii$ du bra : on calcule $\frac12\big(1+\ii\cdot(-1)\big)$.}
  {Seul le terme en $\ket{0}$ est compté ; le terme en $\ket{1}$ vaut $\frac12(-\ii)(-1)=\frac{\ii}{2}$, non nul.}
  {Confusion avec $\braket{+}{-}=0$ ; ici $|\braket{\psi}{\varphi}|^2=\frac12$.}
  {Le bra conjugue les amplitudes : $\bra{\psi}=\frac{1}{\sqrt2}(\bra{0}-\ii\bra{1})$, d'où $\frac12\big(1+(-\ii)(-1)\big)=\frac{1+\ii}{2}$.}

% ---- Question 40 (niveau 2, correct : A)
\qcmq[2]{Chap.~\ref{chap:dirac}\,$\cdot$\,\S\,\ref{sec:qubit}}
  {Lequel de ces états est orthogonal à $\ket{{+}\ii}=\frac{1}{\sqrt2}(\ket{0}+\ii\ket{1})$ ?}
  {$\ket{{-}\ii}=\frac{1}{\sqrt2}(\ket{0}-\ii\ket{1})$}
  {$\ket{-}=\frac{1}{\sqrt2}(\ket{0}-\ket{1})$}
  {$\ket{+}=\frac{1}{\sqrt2}(\ket{0}+\ket{1})$}
  {$\ket{1}$}
  {A}
\qcmx
  {$\braket{{+}\ii}{{-}\ii}=\frac12\big(1+(-\ii)(-\ii)\big)=\frac12(1-1)=0$, en conjuguant le bra.}
  {$\braket{{+}\ii}{-}=\frac{1+\ii}{2}\neq0$ : $\ket{-}$ est orthogonal à $\ket{+}$, pas à $\ket{{+}\ii}$.}
  {$\braket{{+}\ii}{+}=\frac{1-\ii}{2}\neq0$.}
  {$\braket{{+}\ii}{1}=\frac{-\ii}{\sqrt2}\neq0$ : $\ket{1}$ est orthogonal à $\ket{0}$ seulement.}

% ---- Question 41 (niveau 2, correct : D)
\qcmq[1]{Chap.~\ref{chap:operateurs}\,$\cdot$\,\S\,\ref{sec:pur-mixte}}
  {Une matrice de densité de qubit a pour pureté $\operatorname{Tr}(\rho^2)=\frac58$. Que peut-on en conclure ?}
  {L'état est pur}
  {La matrice n'est pas une matrice de densité valide}
  {L'état est maximalement mixte, $\rho=\mathrm{Id}/2$}
  {L'état est mixte (mélange statistique)}
  {D}
\qcmx
  {Il faudrait $\operatorname{Tr}\rho^2=1$.}
  {$\frac58$ est compris entre $\frac12$ et $1$ : valeur admissible (exemple du cours : $\mathrm{diag}(\frac34,\frac14)$).}
  {Ce cas donnerait $\operatorname{Tr}\rho^2=\frac12$.}
  {$\operatorname{Tr}\rho^2<1$ ; pour un qubit, $\frac12\leq\operatorname{Tr}\rho^2\leq1$ et $\frac58$ est une valeur admissible.}

% ---- Question 42 (niveau 2, correct : B)
\qcmq[2]{Chap.~\ref{chap:operateurs}\,$\cdot$\,\S\,\ref{sec:rho}}
  {Que vaut $\langle Z\rangle=\operatorname{Tr}(\rho Z)$ pour un qubit dans l'état $\ket{+}$ ?}
  {$1$}
  {$0$}
  {$-1$}
  {$\frac{1}{\sqrt2}$}
  {B}
\qcmx
  {C'est la valeur pour $\ket{0}$.}
  {$\rho=\frac12\begin{pmatrix}1&1\\1&1\end{pmatrix}$ et $\operatorname{Tr}(\rho Z)=\frac12-\frac12=0$ : les résultats $+1$ et $-1$ sont équiprobables.}
  {C'est la valeur pour $\ket{1}$.}
  {C'est l'amplitude de $\ket{0}$ dans $\ket{+}$, pas une valeur moyenne de $Z$.}

% ---- Question 43 (niveau 2, correct : A)
\qcmq[1]{Chap.~\ref{chap:operateurs}\,$\cdot$\,\S\,\ref{sec:bloch}}
  {Sur la boule de Bloch, que représente un point situé à l'intérieur de la sphère ($|\vec r|<1$) ?}
  {Un état mixte, de pureté $\operatorname{Tr}\rho^2=\frac{1+|\vec r|^2}{2}<1$}
  {Un état pur dont la norme n'est pas $1$}
  {Un état pur en superposition}
  {Un point sans signification physique}
  {A}
\qcmx
  {Les états purs sont sur la sphère ; le centre $\vec r=\vec0$ est l'état maximalement mixte $\mathrm{Id}/2$.}
  {Tout état pur est normalisé et situé sur la sphère.}
  {Les états purs, superpositions comprises, sont tous \emph{sur} la sphère.}
  {Tout point de la boule est une matrice de densité valide.}

% ---- Question 44 (niveau 2, correct : B)
\qcmq[1]{Chap.~\ref{chap:operateurs}\,$\cdot$\,\S\,\ref{sec:pur-mixte}}
  {Par quoi la matrice de densité distingue-t-elle l'état pur $\ket{+}$ du mélange équiprobable de $\ket{0}$ et $\ket{1}$ ?}
  {Par leur trace, égale à $1$ pour l'état pur seulement}
  {Par les éléments hors diagonale (cohérences), non nuls pour $\ket{+}$ et nuls pour le mélange}
  {Par leurs éléments diagonaux, qui sont différents}
  {Ils ne se distinguent pas : c'est la même matrice}
  {B}
\qcmx
  {La trace vaut $1$ pour toute matrice de densité.}
  {$\ket{+}\bra{+}=\frac12\begin{pmatrix}1&1\\1&1\end{pmatrix}$ et le mélange vaut $\frac12\mathrm{Id}$ ; une mesure en base $X$ les distingue.}
  {Les deux ont la diagonale $(\frac12,\frac12)$ : $P(0)=P(1)=\frac12$ dans les deux cas.}
  {Les matrices sont différentes : en base $X$, $\ket{+}$ donne $+$ avec certitude, le mélange donne $\pm$ à pile ou face.}

% ---- Question 45 (niveau 2, correct : C)
\qcmq[2]{Chap.~\ref{chap:operateurs}\,$\cdot$\,\S\,\ref{sec:rho}}
  {Un qubit est préparé dans $\ket{0}$ avec la probabilité $\frac34$ et dans $\ket{1}$ avec la probabilité $\frac14$. Que vaut $\langle Z\rangle=\operatorname{Tr}(\rho Z)$ ?}
  {$\frac34$}
  {$1$}
  {$\frac12$}
  {$0$}
  {C}
\qcmx
  {C'est la probabilité $\rho_{00}$ de lire $0$ ; or $Z$ vaut $+1$ pour $0$ et $-1$ pour $1$.}
  {Valeur pour $\rho=\ket{0}\bra{0}$ (préparation certaine de $\ket{0}$).}
  {$\rho=\mathrm{diag}(\frac34,\frac14)$, donc $\operatorname{Tr}(\rho Z)=\frac34\cdot1+\frac14\cdot(-1)=\frac12$ : c'est la coordonnée $r_z$ de la sphère de Bloch.}
  {Valeur pour le mélange équiprobable $\mathrm{Id}/2$.}

% ---- Question 46 (niveau 2, correct : B)
\qcmq[2]{Chap.~\ref{chap:operateurs}\,$\cdot$\,\S\,\ref{sec:tenseur}}
  {Dans la base ordonnée $\{\ket{00},\ket{01},\ket{10},\ket{11}\}$, quel vecteur colonne représente $\ket{1}\otimes\ket{0}=\ket{10}$ ?}
  {$(0,1,0,0)^{\mathsf T}$}
  {$(0,0,1,0)^{\mathsf T}$}
  {$(1,0,0,0)^{\mathsf T}$}
  {$(0,0,0,1)^{\mathsf T}$}
  {B}
\qcmx
  {C'est $\ket{01}$ : l'ordre des facteurs est inversé.}
  {$\binom{0}{1}\otimes\binom{1}{0}=\big(0\binom{1}{0},\,1\binom{1}{0}\big)=(0,0,1,0)^{\mathsf T}$ : c'est le troisième vecteur de la base.}
  {C'est $\ket{00}$.}
  {C'est $\ket{11}$.}

% ---- Question 47 (niveau 2, correct : D)
\qcmq[2]{Chap.~\ref{chap:portes}\,$\cdot$\,\S\,\ref{sec:phase}}
  {Quelle relation lie la porte de phase $S=\mathrm{diag}(1,\ii)$ à la porte $Z$ ?}
  {$S=Z$}
  {$S^2=X$}
  {$S^4=Z$}
  {$S^2=Z$}
  {D}
\qcmx
  {$Z=\mathrm{diag}(1,-1)$ : $S$ est un quart de tour, $Z$ un demi-tour.}
  {$S^2$ est diagonale alors que $X$ ne l'est pas ; c'est $HSH$ qui est une racine carrée de $X$.}
  {$S^4=\mathrm{diag}(1,\ii^4)=\mathrm{Id}$.}
  {$S^2=\mathrm{diag}(1,\ii^2)=\mathrm{diag}(1,-1)=Z$ : deux quarts de tour autour de $z$ font un demi-tour.}

% ---- Question 48 (niveau 2, correct : A)
\qcmq[2]{Chap.~\ref{chap:portes}\,$\cdot$\,\S\,\ref{sec:controlees}}
  {Dans la porte de Toffoli (CCNOT) à trois qubits $x,y,z$, à quelle condition la cible $z$ est-elle inversée ?}
  {Seulement si $x=1$ et $y=1$}
  {Si $x=1$ ou $y=1$}
  {Si $x\neq y$}
  {Si $x=0$ et $y=0$}
  {A}
\qcmx
  {$\ket{x,y,z}\to\ket{x,y,z\oplus xy}$ : la cible est inversée si et seulement si le produit $xy$ vaut $1$.}
  {Un seul contrôle à $1$ ne suffit pas : pour $(x,y)=(1,0)$, $xy=0$ et la cible ne change pas.}
  {Ce serait deux CNOT ($z\to z\oplus x\oplus y$) ; pour $(1,0)$ le Toffoli n'inverse rien.}
  {Pour $(0,0)$ on a $xy=0$ : la cible ne change pas.}

% ---- Question 49 (niveau 2, correct : A)
\qcmq[1]{Chap.~\ref{chap:portes}\,$\cdot$\,\S\,\ref{sec:portes}}
  {Pour un hamiltonien $\Hop$ indépendant du temps, quel opérateur d'évolution $U(t,t_0)$ vérifie $\ket{\Psi(t)}=U(t,t_0)\ket{\Psi(t_0)}$ ?}
  {$U(t,t_0)=\ee^{-\ii\Hop(t-t_0)/\hbar}$}
  {$U(t,t_0)=\ee^{+\ii\Hop(t-t_0)/\hbar}$}
  {$U(t,t_0)=-\ii\hbar\,\partial_t$}
  {$U(t,t_0)=\Hop\,(t-t_0)/\hbar$}
  {A}
\qcmx
  {Solution de $\ii\hbar\,\partial_t\ket{\Psi}=\Hop\ket{\Psi}$ ; comme $\Hop$ est hermitien, $U^\dagger U=\mathrm{Id}$ : l'évolution est unitaire.}
  {Erreur de signe : c'est $U^\dagger$, qui remonte le temps.}
  {C'est un opérateur différentiel, pas un opérateur d'évolution ; l'équation de Schrödinger s'écrit $\ii\hbar\,\partial_t\ket{\Psi}=\Hop\ket{\Psi}$.}
  {On oublie l'exponentielle et le facteur $\ii$ : cet opérateur est hermitien, pas unitaire.}

% ---- Question 50 (niveau 2, correct : D)
\qcmq[2]{Chap.~\ref{chap:portes}\,$\cdot$\,\S\,\ref{sec:controlees}}
  {La porte de Fredkin (SWAP contrôlé) a pour contrôle le premier qubit. Que donne-t-elle sur $\ket{1,0,1}$ ?}
  {$\ket{1,0,1}$}
  {$\ket{0,1,1}$}
  {$\ket{1,1,1}$}
  {$\ket{1,1,0}$}
  {D}
\qcmx
  {C'est le résultat si le contrôle valait $0$ ; ici il vaut $1$ et les deux autres qubits sont différents.}
  {Le SWAP porte sur les qubits $2$ et $3$, pas sur le contrôle.}
  {Un SWAP ne change pas le nombre de $1$, il permute seulement ; $\ket{1,1,1}$ viendrait d'un NOT.}
  {Le contrôle vaut $1$ : les deux autres qubits sont échangés, $\ket{0,1}\to\ket{1,0}$.}

% ---- Question 51 (niveau 2, correct : A)
\qcmq[1]{Chap.~\ref{chap:portes}\,$\cdot$\,\S\,\ref{sec:op2}}
  {Que donne $(H\otimes I)\ket{00}$, où $H$ agit sur le premier facteur du produit tensoriel $\ket{a}\otimes\ket{b}$ ?}
  {$\ket{+}\otimes\ket{0}=\frac{1}{\sqrt2}(\ket{00}+\ket{10})$}
  {$\frac12(\ket{00}+\ket{01}+\ket{10}+\ket{11})$}
  {$\ket{0}\otimes\ket{+}=\frac{1}{\sqrt2}(\ket{00}+\ket{01})$}
  {$\frac{1}{\sqrt2}(\ket{00}+\ket{11})$}
  {A}
\qcmx
  {$(A\otimes B)(\ket{u}\otimes\ket{v})=A\ket{u}\otimes B\ket{v}$ : $H\ket{0}=\ket{+}$ sur le premier facteur, $I\ket{0}=\ket{0}$ sur le second.}
  {C'est $(H\otimes H)\ket{00}$ : ici le second qubit n'est pas touché.}
  {$H$ agit sur le premier qubit, pas sur le second : les facteurs sont inversés.}
  {C'est $\ket{\Phi^+}$ : il faut un CNOT pour intriquer ; des portes locales laissent l'état séparable.}

% ---- Question 52 (niveau 2, correct : C)
\qcmq[1]{Chap.~\ref{chap:bell}\,$\cdot$\,\S\,\ref{sec:separable}}
  {Pourquoi l'état de Bell $\ket{\Phi^+}=\frac{1}{\sqrt2}(\ket{00}+\ket{11})$ est-il dit intriqué ?}
  {Parce que $P(01)=0$}
  {Parce que chaque qubit, pris seul, est dans un état pur}
  {Il ne peut pas s'écrire comme un produit $\ket{\varphi}_A\otimes\ket{\chi}_B$ : $D=c_{00}c_{11}-c_{01}c_{10}=\frac12\neq0$}
  {Parce qu'il est orthogonal à $\ket{00}$}
  {C}
\qcmx
  {C'est vrai, mais $\ket{00}$ a aussi $P(01)=0$ et il est séparable : cette condition ne prouve rien.}
  {C'est l'inverse : pris seul, chaque qubit est complètement mixte ($\rho_A=\mathrm{Id}/2$).}
  {Un produit imposerait $D=0$ ; aucun des deux qubits n'a d'état propre.}
  {$\braket{00}{\Phi^+}=\frac{1}{\sqrt2}\neq0$ ; l'orthogonalité n'a pas de lien avec l'intrication.}

% ---- Question 53 (niveau 2, correct : B)
\qcmq[2]{Chap.~\ref{chap:bell}\,$\cdot$\,\S\,\ref{sec:teleportation}}
  {Dans la téléportation du cours, Alice obtient les résultats $a=0$ et $q=1$. Quelle correction Bob doit-il appliquer à son qubit ?}
  {La porte $X$}
  {La porte $Z$}
  {$X$ puis $Z$}
  {Aucune, l'état est déjà $\ket{\varphi}$}
  {B}
\qcmx
  {$X$ corrige le cas $a=1$ (inversion de bit) ; ici $a=0$.}
  {Bob détient $\alpha\ket{0}-\beta\ket{1}=Z\ket{\varphi}$ ; $Z$ rétablit le signe : $Z(\alpha\ket{0}-\beta\ket{1})=\alpha\ket{0}+\beta\ket{1}$.}
  {Les deux corrections ne sont nécessaires que pour $(a,q)=(1,1)$.}
  {Ce n'est vrai que pour $(a,q)=(0,0)$.}

% ---- Question 54 (niveau 2, correct : B)
\qcmq[2]{Chap.~\ref{chap:bell}\,$\cdot$\,\S\,\ref{sec:bell} ; ex.\,\ref{ex:superdense}}
  {Dans le codage superdense, Alice et Bob partagent $\ket{\Phi^+}$ et Alice agit sur son qubit. Quelle porte transforme $\ket{\Phi^+}$ en $\ket{\Psi^+}=\frac{1}{\sqrt2}(\ket{01}+\ket{10})$ ?}
  {La porte $Z$}
  {La porte $X$}
  {La porte identité}
  {La porte $H$}
  {B}
\qcmx
  {$Z$ donne $\ket{\Phi^-}=\frac{1}{\sqrt2}(\ket{00}-\ket{11})$ : elle change un signe, pas les chiffres.}
  {$X$ échange $\ket{0}$ et $\ket{1}$ sur un qubit : $\ket{00}+\ket{11}\to\ket{01}+\ket{10}$ ; c'est le message $10$.}
  {Elle laisse $\ket{\Phi^+}$ : c'est le message $00$.}
  {$(H\otimes I)\ket{\Phi^+}$ est une superposition de deux états de Bell, pas un état de Bell.}

% ---- Question 55 (niveau 2, correct : D)
\qcmq[2]{Chap.~\ref{chap:bell}\,$\cdot$\,\S\,\ref{sec:info}}
  {Pour deux qubits dans l'état de Bell $\ket{\Phi^+}$, que vaut l'entropie conditionnelle $H(A|B)=H(A,B)-H(B)$ ?}
  {$0$ bit}
  {$+1$ bit}
  {$+2$ bits}
  {$-1$ bit}
  {D}
\qcmx
  {C'est la valeur classique quand $A$ est entièrement déterminé par $B$ (deux bits identiques) ; ici $H(A,B)=0$, pas $H(B)$.}
  {C'est la valeur quand $A$ et $B$ sont indépendants : $H(A|B)=H(A)=1$.}
  {C'est l'information mutuelle $I(A;B)=1+1-0=2$ bits, pas l'entropie conditionnelle.}
  {$H(A,B)=S(\rho_{AB})=0$ (état pur) et $H(B)=1$ bit ($\rho_B=\mathrm{Id}/2$) : $0-1=-1$, valeur impossible en classique.}

% ---- Question 56 (niveau 2, correct : B)
\qcmq[1]{Chap.~\ref{chap:bell}\,$\cdot$\,\S\,\ref{sec:trace}}
  {Quelle est la matrice de densité réduite $\rho_A=\operatorname{Tr}_B\rho_{AB}$ d'un qubit d'une paire dans l'état $\ket{\Phi^+}$ ?}
  {$\rho_A=\ket{0}\bra{0}$}
  {$\rho_A=\frac12\mathrm{Id}$, de pureté $\operatorname{Tr}\rho_A^2=\frac12$}
  {$\rho_A=\ket{+}\bra{+}=\frac12\begin{pmatrix}1&1\\1&1\end{pmatrix}$}
  {$\rho_A=\ket{\Phi^+}\bra{\Phi^+}$}
  {B}
\qcmx
  {Le qubit aurait un état défini, ce qui est faux pour un état intriqué : $P(0)=P(1)=\frac12$.}
  {$\rho_A=\frac12\big(\ket{0}\bra{0}+\ket{1}\bra{1}\big)$ : les termes croisés $\ket{00}\bra{11}$ s'annulent ($\braket{1}{0}=0$). Le couple est pur, mais chaque qubit est complètement mixte.}
  {Les termes croisés sont multipliés par $\braket{1}{0}=0$ dans la trace partielle : ils disparaissent.}
  {C'est la matrice $4\times4$ du couple ; la trace partielle donne une matrice $2\times2$.}

% ---- Question 57 (niveau 2, correct : C)
\qcmq[1]{Chap.~\ref{chap:bell}\,$\cdot$\,\S\,\ref{sec:separable}}
  {Pour $\ket{\Psi}=c_{00}\ket{00}+c_{01}\ket{01}+c_{10}\ket{10}+c_{11}\ket{11}$, quel est le critère de séparabilité ?}
  {$c_{00}c_{11}+c_{01}c_{10}=0$}
  {$c_{00}c_{01}-c_{10}c_{11}=0$}
  {$D=c_{00}c_{11}-c_{01}c_{10}=0$}
  {$c_{00}=c_{11}$ et $c_{01}=c_{10}$}
  {C}
\qcmx
  {Mauvais signe : $\ket{+}\otimes\ket{-}$ est séparable mais donnerait $-\frac12\neq0$.}
  {Mauvais appariement : $\ket{\Psi^+}$ (intriqué) la vérifie, puisque $c_{00}=c_{11}=0$.}
  {Un produit $(a_0\ket{0}+a_1\ket{1})\otimes(b_0\ket{0}+b_1\ket{1})$ a $c_{ij}=a_ib_j$, donc $c_{00}c_{11}=a_0a_1b_0b_1=c_{01}c_{10}$ ; la réciproque est vraie aussi.}
  {Ni nécessaire ($\ket{00}$ est séparable avec $c_{00}\neq c_{11}$) ni suffisant ($\ket{\Phi^+}$ la vérifie et est intriqué).}

% ---- Question 58 (niveau 2, correct : B)
\qcmq[2]{Chap.~\ref{chap:bell}\,$\cdot$\,\S\,\ref{sec:bell}}
  {Quelle est la valeur propre de $X\otimes X$ pour l'état $\ket{\Phi^+}$ ?}
  {$-1$}
  {$+1$}
  {$0$}
  {$\ket{\Phi^+}$ n'est pas un état propre de $X\otimes X$}
  {B}
\qcmx
  {C'est la valeur pour $\ket{\Phi^-}$ et $\ket{\Psi^-}$ : le signe $-$ de l'état se retrouve dans la valeur propre.}
  {$(X\otimes X)\ket{00}=\ket{11}$ et $(X\otimes X)\ket{11}=\ket{00}$, donc $(X\otimes X)\ket{\Phi^+}=\ket{\Phi^+}$.}
  {$X\otimes X$ est unitaire : ses valeurs propres sont de module $1$, jamais nulles.}
  {Les quatre états de Bell sont des états propres communs de $Z\otimes Z$ et de $X\otimes X$.}

% ---- Question 59 (niveau 2, correct : A)
\qcmq[2]{Chap.~\ref{chap:bell}\,$\cdot$\,\S\,\ref{sec:info}}
  {Que vaut l'entropie de von Neumann $S(\rho)=-\sum_i\lambda_i\log_2\lambda_i$ d'un état pur ?}
  {$0$ bit}
  {$1$ bit}
  {$-1$ bit}
  {$+\infty$}
  {A}
\qcmx
  {Les valeurs propres sont $1$ et $0$ : $1\cdot\log_21=0$ et $0\cdot\log_20=0$ (par continuité, $x\log_2x\to0$).}
  {C'est la valeur pour l'état maximalement mixte $\mathrm{Id}/2$ d'un qubit.}
  {L'entropie de $\rho$ est positive ou nulle ; seule l'entropie conditionnelle quantique peut être négative.}
  {$\log_20$ diverge, mais le terme $0\cdot\log_20$ vaut $0$ par continuité.}

% ---- Question 60 (niveau 2, correct : A)
\qcmq[2]{Chap.~\ref{chap:bell}\,$\cdot$\,\S\,\ref{sec:bell}}
  {Lequel des états de Bell a pour valeur propre $-1$ à la fois pour $Z\otimes Z$ et pour $X\otimes X$ ?}
  {$\ket{\Psi^-}=\frac{1}{\sqrt2}(\ket{01}-\ket{10})$}
  {$\ket{\Phi^+}$}
  {$\ket{\Phi^-}$}
  {$\ket{\Psi^+}$}
  {A}
\qcmx
  {Résultats opposés dans les deux bases : $Z\otimes Z=-1$ (chiffres différents) et $X\otimes X=-1$ (signe $-$ dans l'état).}
  {Valeurs propres $(Z\otimes Z,\,X\otimes X)=(+1,+1)$.}
  {Valeurs propres $(Z\otimes Z,\,X\otimes X)=(+1,-1)$.}
  {Valeurs propres $(Z\otimes Z,\,X\otimes X)=(-1,+1)$.}

\clearpage
%% ======================================================================
\subsection{Niveau 3 : avancé (questions 61 à 90)}\label{sec:qcm-n3}
%% ======================================================================
\noindent{\color{insagris}\small Niveau de difficulté : Avancé\quad $\bullet\bullet\bullet\circ$}\par\smallskip
% ---- Question 61 (niveau 3, correct : A)
\qcmq[2]{Chap.~\ref{chap:histoire}\,$\cdot$\,\S\,\ref{sec:solutions}}
  {Un électron d'un puits infini passe du niveau $n=3$ au niveau $n=2$ en émettant un photon. Quelle est l'énergie du photon, en fonction de $E_1$ ?}
  {$5E_1$}
  {$E_1$}
  {$\frac94E_1$}
  {$13E_1$}
  {A}
\qcmx
  {$E_3-E_2=9E_1-4E_1=5E_1$ ; le photon a la fréquence $f=(E_3-E_2)/h$.}
  {On a soustrait les numéros de niveau ($3-2=1$) au lieu de leurs carrés.}
  {C'est le rapport $E_3/E_2$, pas la différence $E_3-E_2$.}
  {On a additionné $9E_1+4E_1$ : l'énergie du photon est une différence de niveaux.}

% ---- Question 62 (niveau 3, correct : A)
\qcmq[2]{Chap.~\ref{chap:histoire}\,$\cdot$\,\S\,\ref{sec:superposition}}
  {À $t=0$, une particule du puits infini est dans l'état $\Psi=\frac{\sqrt3}{2}\,\psi_1+\frac12\,\psi_3$ ($\psi_n$ : états stationnaires normalisés, d'énergies $E_n=n^2E_1$). Quelle est son énergie moyenne $\langle E\rangle$ ?}
  {$3E_1$}
  {$5E_1$}
  {$\left(\frac{\sqrt3}{2}+\frac92\right)E_1$}
  {$7E_1$}
  {A}
\qcmx
  {$P(E_1)=\frac34$ et $P(E_3)=\frac14$, donc $\langle E\rangle=\frac34E_1+\frac14\cdot9E_1=3E_1$.}
  {On a pondéré également les deux niveaux, $\frac{E_1+9E_1}{2}$ : les poids sont les $|c_i|^2$, pas $\frac12$.}
  {On a pondéré par les amplitudes $c_i$ au lieu de leurs carrés $|c_i|^2$.}
  {Poids permutés : $\frac14E_1+\frac34\cdot9E_1=7E_1$ ; la probabilité de $E_1$ est $\left|\frac{\sqrt3}{2}\right|^2=\frac34$.}

% ---- Question 63 (niveau 3, correct : C)
\qcmq[1]{Chap.~\ref{chap:histoire}\,$\cdot$\,\S\,\ref{sec:solutions}}
  {Pour l'état stationnaire $n=2$ du puits infini $]0,a[$, que valent la densité de probabilité de présence en $x=a/2$ et la position moyenne $\langle\hat x\rangle$ ?}
  {Densité maximale en $a/2$, et $\langle\hat x\rangle=a/2$}
  {Densité nulle en $a/2$, et $\langle\hat x\rangle=0$}
  {Densité nulle en $a/2$, et $\langle\hat x\rangle=a/2$}
  {Densité maximale en $a/2$, et $\langle\hat x\rangle=a/4$}
  {C}
\qcmx
  {C'est le cas de $n=1$ ; pour $n=2$ il y a un nœud au centre.}
  {Une densité nulle au centre ne donne pas une moyenne nulle : $\langle\hat x\rangle$ pondère tout le puits et vaut $\frac a2$ par symétrie.}
  {$|\Psi_2|^2=\frac2a\sin^2\frac{2\pi x}{a}$ s'annule en $x=\frac a2$ (nœud), alors que la symétrie impose $\langle\hat x\rangle=\frac a2$ : la moyenne n'est pas la position la plus probable.}
  {Ni l'un ni l'autre : maxima en $\frac a4$ et $\frac{3a}4$, nœud en $\frac a2$, moyenne $\frac a2$.}

% ---- Question 64 (niveau 3, correct : C)
\qcmq[2]{Chap.~\ref{chap:dirac}\,$\cdot$\,\S\,\ref{sec:projecteurs}}
  {Quelle est la matrice du projecteur $\hat P_+=\ket{+}\bra{+}$ dans la base $\{\ket{0},\ket{1}\}$ ?}
  {$\frac12\begin{pmatrix}1&-1\\-1&1\end{pmatrix}$}
  {$\begin{pmatrix}1&0\\0&0\end{pmatrix}$}
  {$\frac12\begin{pmatrix}1&1\\1&1\end{pmatrix}$}
  {$\frac{1}{\sqrt2}\begin{pmatrix}1&1\\1&-1\end{pmatrix}$}
  {C}
\qcmx
  {C'est $\ket{-}\bra{-}$, le projecteur sur $\ket{-}$.}
  {C'est $\ket{0}\bra{0}$ : projecteur de la base $Z$, pas de la base $X$.}
  {$\ket{+}\bra{+}=\frac12(\ket0+\ket1)(\bra0+\bra1)$ ; on vérifie $\hat P_+^2=\hat P_+$ et $\operatorname{Tr}\hat P_+=1$.}
  {C'est la matrice de Hadamard, unitaire ($H^2=\mathrm{Id}$), et non un projecteur (un projecteur vérifie $P^2=P$).}

% ---- Question 65 (niveau 3, correct : B)
\qcmq[2]{Chap.~\ref{chap:dirac}\,$\cdot$\,\S\,\ref{sec:projecteurs}}
  {Un qubit préparé dans $\ket{0}$ est mesuré en base $X$, puis aussitôt en base $Z$. Quelle est la probabilité d'obtenir $1$ à la seconde mesure ?}
  {$0$}
  {$\frac12$}
  {$1$}
  {$\frac14$}
  {B}
\qcmx
  {On garde $\ket{0}$ en mémoire : la mesure en base $X$ a effondré l'état sur $\ket{\pm}$ et effacé $\ket{0}$.}
  {La première mesure donne $\ket{+}$ ou $\ket{-}$ (probabilité $\frac12$ chacun) ; dans les deux cas $P(1)=\frac12$ en base $Z$ : $\frac12\cdot\frac12+\frac12\cdot\frac12=\frac12$.}
  {Sans rapport : $\ket{0}$ n'est plus l'état après la mesure en base $X$.}
  {On n'a compté qu'une branche de la première mesure ; il faut sommer les deux ($\frac14+\frac14$).}

% ---- Question 66 (niveau 3, correct : D)
\qcmq[2]{Chap.~\ref{chap:dirac}\,$\cdot$\,\S\,\ref{sec:projecteurs}}
  {Un qubit est dans l'état $\ket{\psi}=\frac{3}{5}\ket{0}+\frac{4}{5}\ket{1}$. Quelle est la probabilité d'obtenir $-$ en mesurant dans la base $X$ ?}
  {$\frac{16}{25}$}
  {$\frac{1}{25}$}
  {$\frac{49}{50}$}
  {$\frac{1}{50}$}
  {D}
\qcmx
  {C'est $P(1)$ en base $Z$ ($|\frac45|^2$) : on a mesuré dans la mauvaise base.}
  {On a oublié le facteur $\frac{1}{\sqrt2}$ de $\ket{-}$ : $\left|\frac35-\frac45\right|^2=\frac1{25}$, il manque le $\frac12$.}
  {C'est $P(+)$ ; la probabilité demandée est celle du résultat $-$.}
  {$\braket{-}{\psi}=\frac{1}{\sqrt2}\left(\frac35-\frac45\right)=-\frac{1}{5\sqrt2}$, donc $P(-)=\frac1{50}$ ; et $P(+)=\frac{49}{50}$.}

% ---- Question 67 (niveau 3, correct : B)
\qcmq[1]{Chap.~\ref{chap:dirac}\,$\cdot$\,\S\,\ref{sec:projecteurs}}
  {Lequel de ces kets décrit le même état physique que $\ket{+}=\frac{1}{\sqrt2}(\ket{0}+\ket{1})$ ?}
  {$\frac{1}{\sqrt2}\big(\ket{0}+\ii\ket{1}\big)$}
  {$\ee^{\ii\pi/3}\ket{+}$}
  {$\frac{1}{\sqrt2}\big(\ket{0}+\ee^{\ii\pi}\ket{1}\big)$}
  {$\frac{1}{\sqrt2}\big(\ket{0}+\ee^{\ii\pi/3}\ket{1}\big)$}
  {B}
\qcmx
  {La phase $\ii$ ne porte que sur $\ket{1}$ : c'est une phase \emph{relative}, et l'état est $\ket{{+}\ii}$ (base $Y$).}
  {Un facteur de phase \emph{globale} multiplie tout le ket : toutes les probabilités de mesure (et $\rho=\ket{+}\bra{+}$) sont inchangées.}
  {$\ee^{\ii\pi}=-1$ : c'est $\ket{-}$, orthogonal à $\ket{+}$.}
  {Phase relative $\frac\pi3$ : c'est un autre état, sur l'équateur de la sphère de Bloch, à l'azimut $\frac\pi3$.}

% ---- Question 68 (niveau 3, correct : A)
\qcmq[2]{Chap.~\ref{chap:dirac}\,$\cdot$\,\S\,\ref{sec:qubit}}
  {Laquelle de ces paires est une base orthonormée de $\mathbb{C}^2$ ?}
  {$\{\ket{{+}\ii},\ket{{-}\ii}\}$}
  {$\{\ket{0},\ket{+}\}$}
  {$\{\ket{+},\ket{{-}\ii}\}$}
  {$\{\ket{0},\ket{0}+\ket{1}\}$}
  {A}
\qcmx
  {Ces deux états sont normés et orthogonaux : $\braket{{+}\ii}{{-}\ii}=\frac12(1-1)=0$.}
  {Normés, mais $\braket{0}{+}=\frac{1}{\sqrt2}\neq0$ : pas orthogonaux.}
  {$\braket{+}{{-}\ii}=\frac12(1-\ii)\neq0$ : pas orthogonaux.}
  {Le second ket n'est pas normé (norme au carré $2$) ni orthogonal à $\ket{0}$.}

% ---- Question 69 (niveau 3, correct : C)
\qcmq[2]{Chap.~\ref{chap:operateurs}\,$\cdot$\,\S\,\ref{sec:rho}}
  {Laquelle de ces matrices peut être la matrice de densité d'un qubit ?}
  {$\begin{pmatrix}\frac12&\frac34\\\frac34&\frac12\end{pmatrix}$}
  {$\begin{pmatrix}\frac12&\frac{\ii}{2}\\\frac{\ii}{2}&\frac12\end{pmatrix}$}
  {$\begin{pmatrix}\frac23&\frac{\ii}{3}\\-\frac{\ii}{3}&\frac13\end{pmatrix}$}
  {$\begin{pmatrix}\frac34&0\\0&\frac12\end{pmatrix}$}
  {C}
\qcmx
  {Hermitienne et de trace $1$, mais pas positive : $\det=\frac14-\frac9{16}<0$, valeurs propres $\frac54$ et $-\frac14$ (une « probabilité » négative).}
  {De trace $1$ mais pas hermitienne : il faudrait $\rho_{10}=\overline{\rho_{01}}=-\frac{\ii}{2}$.}
  {Hermitienne ($\rho_{10}=\overline{\rho_{01}}$), de trace $1$ et positive : $\det\rho=\frac29-\frac19=\frac19>0$ avec $\operatorname{Tr}\rho>0$ donne deux valeurs propres positives, $\frac{3\pm\sqrt5}{6}$.}
  {Hermitienne et positive, mais de trace $\frac54\neq1$ : les probabilités diagonales ne somment pas à $1$.}

% ---- Question 70 (niveau 3, correct : B)
\qcmq[2]{Chap.~\ref{chap:operateurs}\,$\cdot$\,\S\,\ref{sec:op-dirac}}
  {Un qubit est dans l'état $\ket{\psi}=\frac{1}{\sqrt5}\big(\ket{0}+2\ii\ket{1}\big)$. Que vaut la valeur moyenne $\langle Y\rangle=\bra{\psi}Y\ket{\psi}$, avec $Y=\begin{pmatrix}0&-\ii\\\ii&0\end{pmatrix}$ ?}
  {$0$}
  {$\frac45$}
  {$-\frac45$}
  {$-\frac35$}
  {B}
\qcmx
  {On a oublié de conjuguer les amplitudes du bra : $\psi^{\mathsf T}Y\psi=\frac15(2+2\ii\cdot\ii)=0$ (c'est d'ailleurs la valeur de $\langle X\rangle$ pour cet état).}
  {$Y\ket\psi=\frac1{\sqrt5}(2\ket0+\ii\ket1)$, puis $\bra\psi Y\ket\psi=\frac15\big(1\cdot2+(-2\ii)(\ii)\big)=\frac45$ ; en général $\langle Y\rangle=2\,\mathrm{Im}(\bar\alpha\beta)$.}
  {Erreur de signe : on a utilisé la transposée $\begin{pmatrix}0&\ii\\-\ii&0\end{pmatrix}=-Y$.}
  {C'est $\langle Z\rangle=\frac15-\frac45$ : on a calculé la valeur moyenne d'une autre observable.}

% ---- Question 71 (niveau 3, correct : D)
\qcmq[1]{Chap.~\ref{chap:operateurs}\,$\cdot$\,\S\,\ref{sec:spectral}}
  {On mesure l'observable $A=\begin{pmatrix}2&1\\1&2\end{pmatrix}$ sur un qubit préparé dans $\ket{0}$. Quelle affirmation est correcte ?}
  {On obtient toujours $2$}
  {On obtient $+1$ ou $-1$, avec la probabilité $\frac12$ chacun}
  {On obtient $2$ ou $1$, avec les probabilités $\frac45$ et $\frac15$}
  {On obtient $3$ ou $1$, avec la probabilité $\frac12$ chacun (moyenne $2$)}
  {D}
\qcmx
  {On a lu le coefficient diagonal : $A\ket0=2\ket0+\ket1$ n'est pas proportionnel à $\ket0$, donc $\ket0$ n'est pas vecteur propre de $A$.}
  {Ce sont les valeurs propres de $X$, pas de $A$ : le décalage $2\,\mathrm{Id}$ translate le spectre ($2\pm1$).}
  {On a lu $A\ket0=2\ket0+\ket1$ comme un état : $A$ n'est pas unitaire, ce n'est pas une évolution ; les résultats possibles d'une mesure sont les valeurs propres de $A$.}
  {$A=2\,\mathrm{Id}+X$ a les vecteurs propres de $X$ : valeur propre $3$ pour $\ket+$, $1$ pour $\ket-$. Comme $\ket0=\frac{1}{\sqrt2}(\ket++\ket-)$, les deux résultats sont équiprobables, et $\langle A\rangle=\bra0A\ket0=2$.}

% ---- Question 72 (niveau 3, correct : D)
\qcmq[2]{Chap.~\ref{chap:operateurs}\,$\cdot$\,\S\,\ref{sec:pur-mixte}}
  {Un qubit est préparé dans $\ket{0}$ avec la probabilité $\frac12$ et dans $\ket{+}$ avec la probabilité $\frac12$. Quelle est la pureté $\operatorname{Tr}(\rho^2)$ de son état ?}
  {$1$}
  {$\frac12$}
  {$\frac58$}
  {$\frac34$}
  {D}
\qcmx
  {Un mélange de deux états purs n'est pur que s'ils sont identiques ; ici $\ket0\neq\ket+$, donc $\operatorname{Tr}\rho^2<1$.}
  {Valeur de l'état maximalement mixte $\frac{\mathrm{Id}}2$, atteinte pour deux états \emph{orthogonaux} ; ici ils ne le sont pas (on a aussi pu oublier les termes croisés).}
  {On n'a additionné que les carrés des éléments diagonaux ($\frac9{16}+\frac1{16}$) : il manque $2\rho_{01}\rho_{10}=\frac2{16}$.}
  {$\rho=\frac12\ket0\bra0+\frac12\ket+\bra+=\begin{pmatrix}\frac34&\frac14\\\frac14&\frac14\end{pmatrix}$, donc $\operatorname{Tr}\rho^2=\frac9{16}+\frac1{16}+2\cdot\frac1{16}=\frac34$ : état mixte, mais moins que $\frac{\mathrm{Id}}2$ car $\ket0$ et $\ket+$ ne sont pas orthogonaux.}

% ---- Question 73 (niveau 3, correct : B)
\qcmq[2]{Chap.~\ref{chap:operateurs}\,$\cdot$\,\S\,\ref{sec:bloch}}
  {Quel est le vecteur de Bloch $(r_x,r_y,r_z)$ de l'état $\ket{\psi}=\cos\frac\pi6\,\ket0+\sin\frac\pi6\,\ket1$ ?}
  {$\left(\frac12,\,0,\,\frac{\sqrt3}{2}\right)$}
  {$\left(\frac{\sqrt3}{2},\,0,\,\frac12\right)$}
  {$\left(\frac{\sqrt3}{4},\,0,\,\frac12\right)$}
  {$\left(\frac{\sqrt3}{2},\,0,\,\frac34\right)$}
  {B}
\qcmx
  {On a pris $\theta=\frac\pi6$ : l'angle de Bloch est le \emph{double} de l'angle qui figure dans les amplitudes.}
  {Ici $\frac\theta2=\frac\pi6$, donc $\theta=\frac\pi3$ et $\varphi=0$ : $\vec r=(\sin\theta\cos\varphi,\sin\theta\sin\varphi,\cos\theta)$. Contrôle : $r_z=|\alpha|^2-|\beta|^2=\frac34-\frac14$ et $r_x=2\alpha\beta=\frac{\sqrt3}2$.}
  {On a écrit $r_x=\alpha\beta$ en oubliant le facteur $2$ ; on le voit à $|\vec r|^2=\frac7{16}<1$, impossible pour un état pur.}
  {On a pris $r_z=|\alpha|^2$ au lieu de $|\alpha|^2-|\beta|^2$ ; or $|\vec r|^2=\frac{21}{16}>1$ est impossible.}

% ---- Question 74 (niveau 3, correct : A)
\qcmq[2]{Chap.~\ref{chap:operateurs}\,$\cdot$\,\S\,\ref{sec:bloch}}
  {Un qubit a pour matrice de densité $\rho=\frac12\begin{pmatrix}1&-\ii\\\ii&1\end{pmatrix}$. Quel est son vecteur de Bloch ?}
  {$(0,\,1,\,0)$ : c'est l'état pur $\ket{{+}\ii}$}
  {$(0,\,-1,\,0)$}
  {$\left(0,\,\frac12,\,0\right)$}
  {$(0,\,0,\,0)$}
  {A}
\qcmx
  {Avec $\rho=\frac12\begin{pmatrix}1+r_z&r_x-\ii r_y\\r_x+\ii r_y&1-r_z\end{pmatrix}$ : $r_z=0$ et $r_x-\ii r_y=-\ii$, donc $r_y=1$, $r_x=0$. Comme $|\vec r|=1$, l'état est pur : $\ket{{+}\ii}\bra{{+}\ii}$.}
  {Erreur de signe : $\rho_{01}=\frac{r_x-\ii r_y}{2}$, donc $\rho_{01}=-\frac\ii2$ correspond à $r_y=+1$ ($r_y=-1$ pour $\ket{{-}\ii}$).}
  {On a lu $\rho_{01}=-\frac\ii2$ comme $-\ii r_y$ : il faut multiplier par $2$ ($\rho=\frac12(\mathrm{Id}+\vec r\cdot\vec\sigma)$).}
  {Les diagonales $\frac12,\frac12$ donnent $r_z=0$, mais les éléments non diagonaux (cohérences) ne sont pas nuls : l'état est pur, et non $\frac{\mathrm{Id}}2$.}

% ---- Question 75 (niveau 3, correct : C)
\qcmq[2]{Chap.~\ref{chap:operateurs}\,$\cdot$\,\S\,\ref{sec:tenseur}}
  {Que donne $(X\otimes Z)\,\big(\ket{+}\otimes\ket{-}\big)$ ?}
  {$-\ket{-}\otimes\ket{-}$}
  {$\ket{-}\otimes\ket{+}$}
  {$\ket{+}\otimes\ket{+}$}
  {$\ket{+}\otimes\ket{-}$}
  {C}
\qcmx
  {On a échangé les facteurs ($Z\otimes X$) : $Z\ket+=\ket-$ et $X\ket-=-\ket-$. Le premier opérateur agit sur le premier qubit.}
  {On a cru que $X$ transforme $\ket+$ en $\ket-$ ; or $X$ échange $\ket0$ et $\ket1$ et laisse $\ket+$ inchangé.}
  {$(A\otimes B)(\ket u\otimes\ket v)=A\ket u\otimes B\ket v$ : $X\ket+=\ket+$ ($\ket+$ est propre pour $X$, valeur propre $+1$) et $Z\ket-=\ket+$ ($Z$ échange $\ket+$ et $\ket-$).}
  {On a cru que $Z$ laisse $\ket-$ inchangé ; $Z$ ne laisse invariants que $\ket0$ et $\ket1$, et échange $\ket+$ et $\ket-$.}

% ---- Question 76 (niveau 3, correct : C)
\qcmq[2]{Chap.~\ref{chap:portes}\,$\cdot$\,\S\,\ref{sec:phase}}
  {Le qubit $\ket{+}$ traverse la porte $T=\mathrm{diag}\big(1,\ee^{\ii\pi/4}\big)$, puis est mesuré dans la base $X$ $\{\ket+,\ket-\}$. Quelle est la probabilité d'obtenir $+$ ?}
  {$\frac12$}
  {$\frac{2-\sqrt2}{4}\approx0{,}15$}
  {$\frac{2+\sqrt2}{4}\approx0{,}85$}
  {$\frac{\sqrt2}{2}\approx0{,}71$}
  {C}
\qcmx
  {On a cru que la phase ne change pas les probabilités : c'est vrai en base $Z$, mais la base $X$ est sensible à la phase \emph{relative} (le point de Bloch est à l'azimut $\frac\pi4$ de l'axe $+x$).}
  {C'est $P(-)=\sin^2\frac\pi8$ : on a permuté les résultats $+$ et $-$.}
  {$T\ket+=\frac1{\sqrt2}(\ket0+\ee^{\ii\pi/4}\ket1)$, donc $\braket{+}{\psi}=\frac12\big(1+\ee^{\ii\pi/4}\big)$ et $P(+)=\frac14\big|1+\ee^{\ii\pi/4}\big|^2=\frac{2+2\cos\frac\pi4}{4}=\frac{2+\sqrt2}{4}$ (soit $\cos^2\frac\pi8$).}
  {On a pris $\cos\frac\pi4$ comme probabilité ; il faut le \emph{demi}-angle et le carré : $P=\cos^2\frac{\pi/4}{2}=\cos^2\frac\pi8$.}

% ---- Question 77 (niveau 3, correct : B)
\qcmq[2]{Chap.~\ref{chap:portes}\,$\cdot$\,\S\,\ref{sec:hadamard}}
  {Que vaut le produit de matrices $HYH$, où $Y=\begin{pmatrix}0&-\ii\\\ii&0\end{pmatrix}$ ?}
  {$Y$}
  {$-Y$}
  {$Z$}
  {$X$}
  {B}
\qcmx
  {Par analogie avec $HXH=Z$, on a cru que $H$ ne modifie pas $Y$ ; en fait l'axe $y$ est retourné (signe $-$).}
  {$H$ échange les axes $x$ et $z$ et retourne l'axe $y$ : $HXH=Z$, $HZH=X$, $HYH=-Y$. Direct : $HY=\frac{1}{\sqrt2}\begin{pmatrix}\ii&-\ii\\-\ii&-\ii\end{pmatrix}$, puis $HYH=\begin{pmatrix}0&\ii\\-\ii&0\end{pmatrix}=-Y$.}
  {C'est $HXH$ : $H$ n'échange que les axes $x$ et $z$, l'axe $y$ reste l'axe $y$.}
  {C'est $HZH$ : $Y$ ne peut pas devenir $X$, car $H$ laisse l'axe $y$ à sa place (au signe près).}

% ---- Question 78 (niveau 3, correct : D)
\qcmq[1]{Chap.~\ref{chap:portes}\,$\cdot$\,\S\,\ref{sec:op2}}
  {Que donne $(H\otimes H)\ket{01}$, où $\ket{01}=\ket0\otimes\ket1$ ?}
  {$\frac12\big(\ket{00}+\ket{01}-\ket{10}-\ket{11}\big)$}
  {$\frac12\big(\ket{00}-\ket{01}-\ket{10}+\ket{11}\big)$}
  {$\frac12\big(\ket{00}+\ket{01}+\ket{10}+\ket{11}\big)$}
  {$\frac12\big(\ket{00}-\ket{01}+\ket{10}-\ket{11}\big)$}
  {D}
\qcmx
  {C'est $\ket-\otimes\ket+$ : on a attribué le $\ket-$ au premier qubit, alors que c'est le second qui vaut $\ket1$.}
  {C'est $\ket-\otimes\ket-$ : $H\ket0=\ket+$, et non $\ket-$.}
  {C'est $\ket+\otimes\ket+$, le résultat pour $\ket{00}$ : on a oublié que $H\ket1=\ket-$ introduit des signes.}
  {$H\ket0\otimes H\ket1=\ket+\otimes\ket-=\frac12(\ket0+\ket1)\otimes(\ket0-\ket1)$. Le signe $-$ apparaît quand le \emph{second} chiffre vaut $1$ (formule générale $\frac12\sum_y(-1)^{x\cdot y}\ket y$ avec $x=01$).}

% ---- Question 79 (niveau 3, correct : D)
\qcmq[2]{Chap.~\ref{chap:portes}\,$\cdot$\,\S\,\ref{sec:cnot}}
  {Avec le CNOT du cours (contrôle : second qubit $B$ ; cible : premier qubit $A$ ; $\mathrm{CNOT}\ket{a,b}=\ket{a\oplus b,\,b}$), quel état obtient-on à partir de $\ket{1}_A\otimes\ket{+}_B$ ?}
  {$\ket{1}\otimes\ket{+}$ (inchangé)}
  {$\frac{1}{\sqrt2}\big(\ket{00}+\ket{11}\big)=\ket{\Phi^+}$}
  {$\frac{1}{\sqrt2}\big(\ket{01}-\ket{10}\big)=\ket{\Psi^-}$}
  {$\frac{1}{\sqrt2}\big(\ket{01}+\ket{10}\big)=\ket{\Psi^+}$}
  {D}
\qcmx
  {On a pris le premier qubit comme contrôle : le CNOT agirait alors sur $B$, et $X\ket+=\ket+$ laisserait l'état inchangé. Mais ici le contrôle est $B$.}
  {C'est le résultat pour $\ket{0}_A\ket{+}_B$ ; ici $A$ vaut $\ket1$, ce qui échange les deux termes.}
  {Aucun signe $-$ n'apparaît : le CNOT permute des états de base sans leur ajouter de phase.}
  {$\ket1\ket+=\frac1{\sqrt2}(\ket{10}+\ket{11})$ ; $\ket{10}\mapsto\ket{1\oplus0,0}=\ket{10}$ et $\ket{11}\mapsto\ket{1\oplus1,1}=\ket{01}$. C'est un état de Bell, donc intriqué.}

% ---- Question 80 (niveau 3, correct : C)
\qcmq[2]{Chap.~\ref{chap:portes}\,$\cdot$\,\S\,\ref{sec:swap}}
  {Que vaut $\mathrm{SWAP}\,(X\otimes\mathrm{Id})\,\mathrm{SWAP}$ ?}
  {$X\otimes\mathrm{Id}$}
  {$\mathrm{Id}\otimes\mathrm{Id}$}
  {$\mathrm{Id}\otimes X$}
  {$X\otimes X$}
  {C}
\qcmx
  {On a supposé que $\mathrm{SWAP}$ commute avec $X\otimes\mathrm{Id}$ ; ce n'est pas le cas, puisque cet opérateur n'agit pas de la même façon sur les deux qubits.}
  {On a simplifié $\mathrm{SWAP}^2=\mathrm{Id}$ en oubliant que $X\otimes\mathrm{Id}$ se trouve \emph{entre} les deux SWAP : on ne peut pas les rapprocher.}
  {$\mathrm{SWAP}\,(X\otimes\mathrm{Id})\,\mathrm{SWAP}\ket{a,b}=\mathrm{SWAP}\,(X\otimes\mathrm{Id})\ket{b,a}=\mathrm{SWAP}\ket{b\oplus1,a}=\ket{a,b\oplus1}$ : le SWAP transporte l'opérateur d'un qubit à l'autre (même mécanisme que $\mathrm{CNOT}_{21}=\mathrm{SWAP}\,\mathrm{CNOT}_{12}\,\mathrm{SWAP}$).}
  {Cet opérateur inverserait les deux bits ; ici un seul bit est inversé, celui du \emph{second} qubit.}

% ---- Question 81 (niveau 3, correct : C)
\qcmq[1]{Chap.~\ref{chap:portes}\,$\cdot$\,\S\,\ref{sec:controlees}}
  {La porte $Z$ contrôlée ($\mathrm{CZ}=\mathrm{diag}(1,1,1,-1)$, contrôle : premier qubit) est appliquée à $\ket{+}\otimes\ket{+}$. Quel état obtient-on ?}
  {$\ket{+}\otimes\ket{+}$}
  {$\ket{-}\otimes\ket{-}$}
  {$\frac12\big(\ket{00}+\ket{01}+\ket{10}-\ket{11}\big)$, un état intriqué}
  {$\ket{+}\otimes\ket{-}$}
  {C}
\qcmx
  {On a cru qu'un simple signe est une phase globale ; mais le signe $-$ ne touche que $\ket{11}$ : c'est une phase \emph{relative}.}
  {C'est $Z\otimes Z$ ($Z$ sur les deux qubits), et non une porte contrôlée.}
  {CZ ne change que le signe de $\ket{11}$. Critère de séparabilité : $c_{00}c_{11}-c_{01}c_{10}=-\frac14-\frac14=-\frac12\neq0$ ; on peut l'écrire $\frac{1}{\sqrt2}\big(\ket0\ket++\ket1\ket-\big)$.}
  {On a appliqué $Z$ au second qubit sans condition ; le contrôle exige que le premier qubit soit dans $\ket1$.}

% ---- Question 82 (niveau 3, correct : D)
\qcmq[2]{Chap.~\ref{chap:portes}\,$\cdot$\,\S\,\ref{sec:phase}, \ref{sec:portes-bloch}}
  {Un qubit préparé dans $\ket{0}$ traverse successivement $H$, puis $S$, puis $H$ (dans cet ordre). Quel est l'état final, à une phase globale près ?}
  {$\ket{{+}\ii}$}
  {$\ket{0}$}
  {$\ket{1}$}
  {$\ket{{-}\ii}=\frac{1}{\sqrt2}\big(\ket0-\ii\ket1\big)$}
  {D}
\qcmx
  {Mauvais sens de rotation : on a perdu un signe en calculant $H\ket{{+}\ii}$ (la phase relative vaut $-\ii$, pas $+\ii$).}
  {On a « simplifié » $HSH$ en $S$ (puis $S\ket0=\ket0$) : mais $H$ ne commute pas avec $S$, on ne peut pas supprimer les deux $H$.}
  {C'est le résultat pour $HZH=X$ ; or $S$ n'est qu'un \emph{quart} de tour autour de $z$ (et non un demi-tour comme $Z$), donc $HSH$ est un quart de tour autour de $x$.}
  {$HSH=\sqrt X$ est un quart de tour autour de $x$ (§\,\ref{sec:phase}). Directement : $H\ket0=\ket+$, $S\ket+=\ket{{+}\ii}$, $H\ket{{+}\ii}=\frac{1+\ii}{2}\ket0+\frac{1-\ii}{2}\ket1=\frac{1+\ii}{2}\big(\ket0-\ii\ket1\big)$.}

% ---- Question 83 (niveau 3, correct : A)
\qcmq[1]{Chap.~\ref{chap:bell}\,$\cdot$\,\S\,\ref{sec:mesure2}}
  {Deux qubits sont dans l'état $\ket{\Psi}=\frac{1}{\sqrt3}\big(\ket{00}+\ket{01}+\ket{11}\big)$. On mesure le \emph{second} qubit $B$ en base $Z$ et on obtient $1$. Que valait la probabilité de ce résultat, et dans quel état se trouve alors $A$ ?}
  {$P(B{=}1)=\frac23$ ; $A$ est dans l'état $\ket{+}$}
  {$P(B{=}1)=\frac13$ ; $A$ est dans l'état $\ket{0}$}
  {$P(B{=}1)=\frac23$ ; $A$ est dans l'état $\ket{1}$}
  {$P(B{=}1)=\frac13$ ; $A$ est dans l'état $\ket{+}$}
  {A}
\qcmx
  {$P(B{=}1)=|c_{01}|^2+|c_{11}|^2=\frac13+\frac13$. On garde les termes dont le second chiffre vaut $1$ : $\frac{1}{\sqrt3}\big(\ket{01}+\ket{11}\big)=\frac{1}{\sqrt3}(\ket0+\ket1)\otimes\ket1$, qui renormalisé donne $\ket{+}_A\otimes\ket{1}_B$.}
  {C'est le résultat $B=0$ : seul le terme $\ket{00}$ a un second chiffre nul, d'où la probabilité $\frac13$ et $A$ dans $\ket0$.}
  {On a projeté sur le seul terme $\ket{11}$ ; or $\ket{01}$ a aussi son second chiffre égal à $1$ : $A$ vaut $0$ ou $1$ (état $\ket+$).}
  {L'état de $A$ est juste, mais on n'a compté qu'un seul terme pour la probabilité ; il faut sommer $|c_{01}|^2+|c_{11}|^2$.}

% ---- Question 84 (niveau 3, correct : A)
\qcmq[2]{Chap.~\ref{chap:bell}\,$\cdot$\,\S\,\ref{sec:separable}}
  {Pour quelle valeur de $x$ l'état (non normalisé) $\ket{00}+2\ket{01}+3\ket{10}+x\ket{11}$ est-il séparable ?}
  {$x=6$}
  {$x=-6$}
  {$x=\frac16$}
  {$x=5$}
  {A}
\qcmx
  {Critère $D=c_{00}c_{11}-c_{01}c_{10}=x-2\cdot3=0$. On obtient bien un produit : $(\ket0+3\ket1)\otimes(\ket0+2\ket1)=\ket{00}+2\ket{01}+3\ket{10}+6\ket{11}$.}
  {Erreur de signe : $D=x-6$ s'annule pour $x=+6$ ; pour $x=-6$, $D=-12\neq0$ et l'état est intriqué.}
  {On a inversé le rapport : $c_{00}c_{11}=c_{01}c_{10}$ donne $x=\frac{c_{01}c_{10}}{c_{00}}=6$, et non $\frac{c_{00}}{c_{01}c_{10}}$.}
  {On a additionné ($2+3$) au lieu de multiplier : le critère de séparabilité porte sur des \emph{produits} de coefficients.}

% ---- Question 85 (niveau 3, correct : B)
\qcmq[2]{Chap.~\ref{chap:bell}\,$\cdot$\,\S\,\ref{sec:bell}}
  {Comment s'écrit $\ket{00}$ dans la base de Bell ?}
  {$\frac{1}{\sqrt2}\big(\ket{\Phi^+}-\ket{\Phi^-}\big)$}
  {$\frac{1}{\sqrt2}\big(\ket{\Phi^+}+\ket{\Phi^-}\big)$}
  {$\frac{1}{\sqrt2}\big(\ket{\Phi^+}+\ket{\Psi^+}\big)$}
  {$\ket{\Phi^+}$}
  {B}
\qcmx
  {Erreur de signe : cette combinaison vaut $\ket{11}$, puisque $\ket{\Phi^+}-\ket{\Phi^-}=\sqrt2\,\ket{11}$.}
  {$\ket{\Phi^+}+\ket{\Phi^-}=\frac{1}{\sqrt2}\cdot2\ket{00}=\sqrt2\,\ket{00}$. Une mesure de Bell sur $\ket{00}$ donne donc $\Phi^+$ ou $\Phi^-$ avec la probabilité $\frac12$ chacun (jamais $\Psi^\pm$).}
  {Cette combinaison vaut $\ket{+}\otimes\ket{+}=\frac12(\ket{00}+\ket{01}+\ket{10}+\ket{11})$ : les deux états de Bell « pairs » et « impairs » s'y ajoutent.}
  {On a confondu $\ket{00}$, état séparable, avec $\ket{\Phi^+}=\frac{1}{\sqrt2}(\ket{00}+\ket{11})$, qui contient aussi $\ket{11}$ et est intriqué.}

% ---- Question 86 (niveau 3, correct : A)
\qcmq[1]{Chap.~\ref{chap:bell}\,$\cdot$\,\S\,\ref{sec:circuit-bell}}
  {Le générateur de Bell du cours est $G=\mathrm{CNOT}\cdot(\mathrm{Id}\otimes H)$ : $H$ sur le second qubit $B$, puis CNOT de contrôle $B$ et de cible $A$. Quel circuit permet de \emph{mesurer} une paire dans la base de Bell (on le place avant une mesure en base $Z$) ?}
  {D'abord le CNOT (contrôle $B$, cible $A$), puis $H$ sur $B$}
  {D'abord $H$ sur $B$, puis le CNOT}
  {D'abord le CNOT, puis $H$ sur $A$ (la cible)}
  {Aucun circuit : on mesure directement les deux qubits en base $Z$}
  {A}
\qcmx
  {C'est l'inverse $G^\dagger=(\mathrm{Id}\otimes H)\,\mathrm{CNOT}$ : les portes sont leurs propres inverses et l'ordre s'inverse. Il envoie $\Phi^+,\Phi^-,\Psi^+,\Psi^-$ sur $\ket{00},\ket{01},\ket{10},\ket{11}$ (à une phase près).}
  {C'est $G$ lui-même, qui \emph{fabrique} des états de Bell : on garde le même ordre au lieu de l'inverser, et la mesure ne distingue plus les quatre états.}
  {Mauvais qubit : pour $\Phi^+$ on trouve $\ket{+}\otimes\ket{+}$, dont les deux bits sont aléatoires. Le $H$ doit agir sur $B$, qui a reçu le $H$ dans $G$.}
  {$\Phi^+$ et $\Phi^-$ (comme $\Psi^+$ et $\Psi^-$) donnent les mêmes résultats en base $Z$ : cette base ne voit pas le signe relatif, il faut d'abord défaire l'intrication.}

% ---- Question 87 (niveau 3, correct : B)
\qcmq[2]{Chap.~\ref{chap:bell}\,$\cdot$\,\S\,\ref{sec:trace}}
  {Quelle est la matrice de densité réduite $\rho_A=\operatorname{Tr}_B\ket\Psi\bra\Psi$ du premier qubit pour $\ket{\Psi}=\frac12\ket{00}+\frac{\sqrt3}{2}\ket{11}$ ?}
  {$\begin{pmatrix}\frac14&\frac{\sqrt3}{4}\\\frac{\sqrt3}{4}&\frac34\end{pmatrix}$}
  {$\begin{pmatrix}\frac14&0\\0&\frac34\end{pmatrix}$}
  {$\begin{pmatrix}\frac12&0\\0&\frac12\end{pmatrix}$}
  {$\begin{pmatrix}\frac12&0\\0&\frac{\sqrt3}{2}\end{pmatrix}$}
  {B}
\qcmx
  {C'est la matrice $\ket\chi\bra\chi$ de l'état pur $\ket\chi=\frac12\ket0+\frac{\sqrt3}2\ket1$ : on a gardé les cohérences, comme si $A$ était seul ; l'intrication avec $B$ les détruit.}
  {Avec $C=\begin{pmatrix}c_{00}&c_{01}\\c_{10}&c_{11}\end{pmatrix}=\begin{pmatrix}\frac12&0\\0&\frac{\sqrt3}2\end{pmatrix}$, $\rho_A=CC^\dagger=\mathrm{diag}\big(\frac14,\frac34\big)$ : état mixte, de pureté $\frac{10}{16}=1-2|D|^2$ avec $D=\frac{\sqrt3}{4}$.}
  {C'est la réponse pour un état de Bell (intrication maximale) ; ici les poids $\frac14$ et $\frac34$ diffèrent : l'intrication est partielle.}
  {On a mis les amplitudes au lieu de leurs carrés : la trace vaudrait $\frac{1+\sqrt3}{2}\neq1$.}

% ---- Question 88 (niveau 3, correct : A)
\qcmq[2]{Chap.~\ref{chap:bell}\,$\cdot$\,\S\,\ref{sec:info}}
  {Deux qubits sont dans le mélange statistique $\rho_{AB}=\frac12\ket{00}\bra{00}+\frac12\ket{11}\bra{11}$ (corrélations purement classiques). Que vaut $H(A|B)=S(\rho_{AB})-S(\rho_B)$, en bits ?}
  {$0$}
  {$-1$}
  {$1$}
  {$2$}
  {A}
\qcmx
  {$\rho_{AB}$ a pour valeurs propres $\frac12,\frac12,0,0$, donc $S(\rho_{AB})=1$ ; et $\rho_B=\frac12\mathrm{Id}$ donne $S(\rho_B)=1$. Ainsi $H(A|B)=0$ : connaître $B$ fixe $A$. Une valeur négative est propre à l'intrication.}
  {C'est la valeur pour l'état de Bell \emph{pur} $\ket{\Phi^+}$ : mêmes résultats en base $Z$, mais $S(\rho_{AB})=0$ car l'état global est pur. Ici l'état est mixte.}
  {On a pris $H(A|B)=H(A)$, comme si $B$ ne disait rien sur $A$ ; or les deux bits sont toujours égaux.}
  {On a additionné les entropies des deux qubits ($1+1$), valeur de deux bits indépendants : ce n'est pas une entropie conditionnelle.}

% ---- Question 89 (niveau 3, correct : B)
\qcmq[1]{Chap.~\ref{chap:bell}\,$\cdot$\,\S\,\ref{sec:teleportation}}
  {Dans la téléportation, avant de recevoir les deux bits $(a,q)$ d'Alice, dans quel état se trouve le qubit de Bob ?}
  {Dans l'état pur $\alpha\ket0+\beta\ket1$, déjà téléporté}
  {Dans l'état complètement mélangé $\rho_B=\frac12\mathrm{Id}$, indépendant de $\alpha$ et $\beta$}
  {Dans le mélange $|\alpha|^2\ket0\bra0+|\beta|^2\ket1\bra1$}
  {Dans l'état pur $\ket{0}$, son état initial}
  {B}
\qcmx
  {Ce serait une communication instantanée. L'état $\alpha\ket0+\beta\ket1$ n'apparaît qu'\emph{après} la correction, qui dépend de $(a,q)$.}
  {Les quatre kets $\ket{\phi_{aq}}$ sont équiprobables ($\frac14$) et $\frac14\sum_{a,q}\ket{\phi_{aq}}\bra{\phi_{aq}}=\frac12\mathrm{Id}$. Bob n'apprend rien avant le message classique : rien ne voyage plus vite que la lumière.}
  {Bob connaîtrait alors les statistiques de $\ket\varphi$ sans aucun message : or les probabilités d'Alice ($\frac14$ pour chaque résultat) ne dépendent pas de $\alpha,\beta$.}
  {Le qubit de Bob est la moitié de la paire $\ket{\Phi^+}$ : pris seul, il n'a pas d'état pur ($\rho_B=\frac12\mathrm{Id}$ avant comme après la mesure d'Alice, tant que Bob ignore le résultat).}

% ---- Question 90 (niveau 3, correct : D)
\qcmq[1]{Chap.~\ref{chap:bell}\,$\cdot$\,\S\,\ref{sec:chsh}}
  {Au jeu CHSH (gain si $a\oplus b=x\cdot y$), quelles sont les meilleures probabilités de gain avec une stratégie classique, même avec du hasard partagé, et avec une paire $\ket{\Phi^+}$ mesurée dans des bases bien choisies ?}
  {$\frac12$ en classique ; $1$ avec l'intrication}
  {$\frac34$ dans les deux cas}
  {$1$ en classique si Alice et Bob s'accordent à l'avance ; aucun avantage quantique}
  {$\frac34$ en classique ; $\cos^2\frac\pi8\approx0{,}85$ avec l'intrication}
  {D}
\qcmx
  {Classique : répondre toujours $0$ donne déjà $\frac34$. Quantique : l'intrication ne permet pas de gagner à tous les coups (borne de Tsirelson $\cos^2\frac\pi8$).}
  {Si l'intrication n'apportait rien, on n'aurait pas de violation de l'inégalité de Bell ; l'avantage quantique ($0{,}85>0{,}75$) est précisément ce que l'expérience de Bell met en évidence.}
  {S'accorder à l'avance (ou partager du hasard) ne suffit pas : aucune stratégie déterministe ne gagne les quatre parties, d'après l'argument de parité.}
  {Une stratégie déterministe perd au moins une des quatre parties (la somme modulo $2$ des conditions donne $0=1$) ; avec $\Phi^+$, chaque partie est gagnée avec la probabilité $\cos^2\frac\pi8=\frac{2+\sqrt2}{4}$.}

\clearpage
%% ======================================================================
\subsection{Niveau 4 : maîtrise (questions 91 à 100)}\label{sec:qcm-n4}
%% ======================================================================
\noindent{\color{insagris}\small Niveau de difficulté : Maîtrise\quad $\bullet\bullet\bullet\bullet$}\par\smallskip
% ---- Question 91 (niveau 4, correct : C)
\qcmq[2]{Chap.~\ref{chap:histoire}\,$\cdot$\,\S\,\ref{sec:superposition}}
  {Une particule du puits infini est préparée à $t=0$ dans l'état $\Psi=\frac{1}{\sqrt2}\big(\psi_1+\psi_2\big)$, avec $E_n=n^2E_1$. Au bout de quelle durée minimale la densité de probabilité $|\Psi(x,t)|^2$ retrouve-t-elle sa forme initiale ?}
  {$T=\dfrac{2\pi\hbar}{E_1}$}
  {$T=\dfrac{2\pi\hbar}{5E_1}$}
  {$T=\dfrac{2\pi\hbar}{3E_1}$}
  {$T=\dfrac{\pi\hbar}{3E_1}$}
  {C}
\qcmx
  {C'est la période de la phase $\ee^{-\ii E_1t/\hbar}$ de $\psi_1$ seul ; or les phases globales ne se voient pas dans $|\Psi|^2$ : seul compte l'écart $E_2-E_1=3E_1$.}
  {On a pris la somme $E_1+E_2=5E_1$ ; c'est la \emph{différence} des énergies (phase relative) qui intervient.}
  {$|\Psi|^2=\frac12(\psi_1^2+\psi_2^2)+\psi_1\psi_2\cos\frac{(E_2-E_1)t}{\hbar}$ : seul le terme croisé dépend du temps, à la pulsation $\omega=\frac{E_2-E_1}{\hbar}=\frac{3E_1}{\hbar}$, donc $T=\frac{2\pi}{\omega}$.}
  {C'est la demi-période : à $t=\frac T2$ le cosinus vaut $-1$ et la densité est le symétrique de l'état initial ($x\to a-x$), pas encore sa forme de départ.}

% ---- Question 92 (niveau 4, correct : D)
\qcmq[1]{Chap.~\ref{chap:dirac}\,$\cdot$\,\S\,\ref{sec:projecteurs}, \ref{sec:qubit}}
  {Un expérimentateur reçoit un qubit préparé soit dans $\ket{0}$, soit dans $\ket{+}$. Laquelle de ces affirmations est correcte ?}
  {Une mesure en base $Z$ les distingue : $\ket0$ donne $0$ et $\ket+$ donne $1$}
  {Une mesure en base $X$ les distingue : $\ket+$ donne $+$ et $\ket0$ donne $-$}
  {On les distingue en répétant la mesure sur le même qubit}
  {Aucune mesure ne permet de les distinguer à coup sûr, car $\braket{0}{+}=\frac{1}{\sqrt2}\neq0$}
  {D}
\qcmx
  {$\ket+$ donne $0$ ou $1$ avec la probabilité $\frac12$ : le résultat $0$ est possible pour les deux états.}
  {$\ket0=\frac{1}{\sqrt2}(\ket++\ket-)$ donne $+$ ou $-$ avec la probabilité $\frac12$ : le résultat $+$ est possible pour les deux états.}
  {Après une première mesure l'état est effondré : répéter ne donne aucune information nouvelle, et on ne peut pas copier le qubit pour recommencer (non-clonage).}
  {Seuls des états \emph{orthogonaux} sont parfaitement discernables. En base $Z$, $\ket0$ donne toujours $0$ et $\ket+$ donne $0$ une fois sur deux : obtenir $0$ ne tranche pas (obtenir $1$ prouve seulement qu'on avait $\ket+$).}

% ---- Question 93 (niveau 4, correct : C)
\qcmq[1]{Chap.~\ref{chap:operateurs}\,$\cdot$\,\S\,\ref{sec:pur-mixte}}
  {Une source envoie au hasard $\ket{0}$ ou $\ket{1}$ (probabilité $\frac12$ chacun) ; une autre envoie au hasard $\ket{+}$ ou $\ket{-}$ (probabilité $\frac12$ chacun). Peut-on distinguer ces deux sources en mesurant les qubits émis ?}
  {Oui : en base $Z$, la première source donne des résultats prévisibles}
  {Oui : en base $X$, la seconde source donne des résultats prévisibles}
  {Non : les deux ont la même matrice de densité $\frac{\mathrm{Id}}{2}$, donc les mêmes statistiques pour toute mesure}
  {Oui : les états purs envoyés ne sont pas les mêmes}
  {C}
\qcmx
  {Chaque qubit est bien $\ket0$ ou $\ket1$, mais on ignore lequel : le résultat vaut $0$ ou $1$ avec la probabilité $\frac12$, exactement comme avec la seconde source ($|\braket{0}{\pm}|^2=\frac12$).}
  {Même erreur : chaque qubit est $\ket+$ ou $\ket-$, mais on ignore lequel ; le résultat en base $X$ est aléatoire, et la première source donne aussi $+$ ou $-$ avec la probabilité $\frac12$.}
  {$\frac12\ket0\bra0+\frac12\ket1\bra1=\frac12\ket+\bra++\frac12\ket-\bra-=\frac{\mathrm{Id}}2$ (relation de fermeture dans chaque base). Les probabilités de toute mesure ne dépendent que de $\rho$.}
  {Ce qui compte n'est pas la liste des états purs envoyés mais la matrice de densité moyenne : des préparations différentes peuvent avoir la même $\rho$.}

% ---- Question 94 (niveau 4, correct : B)
\qcmq[1]{Chap.~\ref{chap:operateurs}\,$\cdot$\,\S\,\ref{sec:pur-mixte} ; Chap.~\ref{chap:portes}\,$\cdot$\,\S\,\ref{sec:portes-bloch}}
  {Un qubit est dans l'état complètement mélangé $\rho=\frac{\mathrm{Id}}{2}$. Existe-t-il une porte quantique $U$ qui le transforme en un état pur, par exemple $\ket{0}\bra{0}$ ?}
  {Oui : la porte de Hadamard, qui crée des superpositions}
  {Non : $\rho\mapsto U\rho U^\dagger$ conserve les valeurs propres, donc la pureté ; ici $U\frac{\mathrm{Id}}{2}U^\dagger=\frac{\mathrm{Id}}{2}$}
  {Oui : une rotation de Bloch qui amène le centre de la boule au pôle nord}
  {Oui : une mesure en base $Z$, qui donne $\ket0$}
  {B}
\qcmx
  {$H$ envoie un état pur sur un état pur ($\ket0\mapsto\ket+$) et laisse $\frac{\mathrm{Id}}{2}$ invariant : $H\frac{\mathrm{Id}}2H=\frac{\mathrm{Id}}2$ car $H^2=\mathrm{Id}$.}
  {$\operatorname{Tr}\big[(U\rho U^\dagger)^2\big]=\operatorname{Tr}\rho^2$ ($=\frac12$ ici). Sur la boule de Bloch, une porte est une rotation : elle laisse le centre fixe.}
  {Une rotation tourne la boule autour de son centre, qui reste fixe : elle ne peut pas l'amener sur la sphère.}
  {Une mesure n'est pas une porte (elle est non unitaire) ; de plus elle donne $\ket0$ ou $\ket1$ au hasard, pas $\ket0$ à coup sûr.}

% ---- Question 95 (niveau 4, correct : D)
\qcmq[1]{Chap.~\ref{chap:portes}\,$\cdot$\,\S\,\ref{sec:cnot}}
  {Que vaut $(H\otimes H)\,\mathrm{CNOT}_{12}\,(H\otimes H)$, où $\mathrm{CNOT}_{12}$ a pour contrôle le \emph{premier} qubit et pour cible le second ?}
  {$\mathrm{CNOT}_{12}$, inchangé}
  {$\mathrm{CZ}$ (porte $Z$ contrôlée)}
  {$\mathrm{SWAP}$}
  {$\mathrm{CNOT}_{21}$ (contrôle : second qubit ; cible : premier qubit)}
  {D}
\qcmx
  {On a cru que les deux $H$ s'annulent ($H^2=\mathrm{Id}$) ; mais ils encadrent le CNOT au lieu d'être côte à côte : on ne peut pas les simplifier.}
  {C'est le résultat quand $H$ n'encadre que la \emph{cible} : $(\mathrm{Id}\otimes H)\,\mathrm{CNOT}_{12}\,(\mathrm{Id}\otimes H)=\mathrm{CZ}$, car $HXH=Z$. Ici $H$ agit aussi sur le contrôle.}
  {Le SWAP s'obtient avec \emph{trois} CNOT alternés, pas avec un seul CNOT conjugué par $H\otimes H$ ; il n'y a pas de contrôle dans un SWAP.}
  {$H\ket0\bra0H=\ket+\bra+$, $H\ket1\bra1H=\ket-\bra-$ et $HXH=Z$ donnent $\ket+\bra+\otimes I+\ket-\bra-\otimes Z=I\otimes\ket0\bra0+X\otimes\ket1\bra1=\mathrm{CNOT}_{21}$ : dans la base $\{\ket+,\ket-\}$, contrôle et cible échangent leurs rôles (rebond de phase).}

% ---- Question 96 (niveau 4, correct : A)
\qcmq[1]{Chap.~\ref{chap:portes}\,$\cdot$\,\S\,\ref{sec:controlees}}
  {On applique la porte de Toffoli (contrôles : qubits 1 et 2 ; cible : qubit 3) à $\ket{+}\otimes\ket{+}\otimes\ket{0}$, puis on mesure le troisième qubit en base $Z$. Que peut-on dire si l'on obtient $1$ ?}
  {Ce résultat a la probabilité $\frac14$, et les deux premiers qubits sont alors dans l'état $\ket{11}$}
  {Ce résultat a la probabilité $\frac12$, et les deux premiers qubits restent dans $\ket{+}\ket{+}$}
  {Ce résultat a la probabilité $\frac14$, mais les deux premiers qubits restent dans $\ket{+}\ket{+}$}
  {Ce résultat a la probabilité $\frac34$, et les deux premiers qubits sont dans l'état $\frac{1}{\sqrt3}\big(\ket{00}+\ket{01}+\ket{10}\big)$}
  {A}
\qcmx
  {$\mathrm{CCNOT}\ket{x,y,0}=\ket{x,y,xy}$ donne $\frac12\big(\ket{000}+\ket{010}+\ket{100}+\ket{111}\big)$ : seul $\ket{111}$ a un troisième chiffre égal à $1$ (probabilité $\frac14$), et il impose $x=y=1$.}
  {La cible n'est inversée que si $x=y=1$ (probabilité $\frac14$) et non une fois sur deux ; de plus l'état est intriqué, donc mesurer le qubit 3 modifie les qubits 1 et 2.}
  {La probabilité est juste, mais la mesure de la cible n'est pas sans effet sur les contrôles : l'état après la Toffoli est intriqué, donc le résultat $1$ les projette sur $\ket{11}$.}
  {C'est le cas du résultat $0$ (probabilité $\frac34$ : $x,y\in\{00,01,10\}$). Le résultat $1$ est le cas rare, celui du ET vrai ($x=y=1$).}

% ---- Question 97 (niveau 4, correct : A)
\qcmq[1]{Chap.~\ref{chap:bell}\,$\cdot$\,\S\,\ref{sec:bell}}
  {Que donne $(H\otimes H)\ket{\Psi^-}$, avec $\ket{\Psi^-}=\frac{1}{\sqrt2}\big(\ket{01}-\ket{10}\big)$ ?}
  {$-\ket{\Psi^-}$, c'est-à-dire le même état physique}
  {$\ket{\Phi^-}$}
  {$\ket{\Psi^+}$}
  {$\ket{\Phi^+}$}
  {A}
\qcmx
  {Le singulet est invariant par $U\otimes U$ à un facteur près : $(U\otimes U)\ket{\Psi^-}=\det(U)\ket{\Psi^-}$ avec $\det H=-1$. Directement : $\frac{1}{\sqrt2}\big(\ket+\ket--\ket-\ket+\big)=-\ket{\Psi^-}$.}
  {C'est l'image de $\ket{\Psi^+}$, pas de $\ket{\Psi^-}$ : $(H\otimes H)\ket{\Psi^+}=\ket{\Phi^-}$. Le signe $-$ du singulet change le calcul (les termes $\ket{00}$ et $\ket{11}$ se compensent).}
  {$\ket{\Psi^+}$ est symétrique par échange des deux qubits, $\ket{\Psi^-}$ antisymétrique ; $H\otimes H$ commute avec le SWAP et conserve donc cette symétrie.}
  {C'est l'état que $H\otimes H$ laisse invariant ($\ket{\Phi^+}=\frac{1}{\sqrt2}(\ket{++}+\ket{--})$) ; le singulet n'a pas ce comportement.}

% ---- Question 98 (niveau 4, correct : B)
\qcmq[1]{Chap.~\ref{chap:bell}\,$\cdot$\,\S\,\ref{sec:chsh}}
  {Une expérience de type CHSH (avec $S=E_{00}+E_{01}+E_{10}-E_{11}$) mesure $S=2{,}6$. Que peut-on en conclure ?}
  {C'est impossible : on a toujours $|S|\leq2$}
  {Aucun modèle local à résultats prédéterminés ne l'explique, mais c'est compatible avec la mécanique quantique}
  {C'est impossible en mécanique quantique, car $2\sqrt2\approx2{,}4<2{,}6$}
  {Les particules s'échangent de l'information plus vite que la lumière}
  {B}
\qcmx
  {Cette borne ne vaut que pour les modèles classiques (locaux) ; l'expérience et la mécanique quantique la dépassent, jusqu'à $2\sqrt2$.}
  {Classique : $|S|\leq2$ ; quantique : $|S|\leq2\sqrt2\approx2{,}83$ (borne de Tsirelson). Comme $2<2{,}6<2\sqrt2$, l'inégalité de Bell est violée sans dépasser le maximum quantique ; le gain correspondant est $\frac12+\frac{2{,}6}{8}\approx0{,}825$.}
  {Erreur numérique : $2\sqrt2\approx2{,}83$, donc $2{,}6$ est permis.}
  {Les corrélations violent l'inégalité de Bell mais ne permettent aucun signal : la réponse de Bob est uniforme quelle que soit la base choisie par Alice.}

% ---- Question 99 (niveau 4, correct : A)
\qcmq[2]{Chap.~\ref{chap:bell}\,$\cdot$\,\S\,\ref{sec:info}}
  {Pour l'état $\ket{\Psi}=\frac{1}{\sqrt3}\big(\ket{00}+\sqrt2\,\ket{11}\big)$, quelle est l'entropie de von Neumann $S(\rho_A)$ du premier qubit, en bits ?}
  {$\log_2 3-\frac23\approx0{,}92$}
  {$1$}
  {$0$}
  {$0{,}64$}
  {A}
\qcmx
  {$\rho_A=\mathrm{diag}\big(\frac13,\frac23\big)$ (trace partielle), donc $S=-\frac13\log_2\frac13-\frac23\log_2\frac23=\log_23-\frac23\approx1{,}585-0{,}667$. C'est moins que le maximum d'un bit (état de Bell) : l'intrication est partielle.}
  {C'est la valeur pour un état de Bell ($\rho_A=\frac{\mathrm{Id}}{2}$) ; ici les poids $\frac13$ et $\frac23$ diffèrent, donc $S<1$.}
  {$S(\rho_A)=0$ vaut pour un état séparable ; ici $D=\frac{\sqrt2}{3}\neq0$ : l'état est intriqué, $\rho_A$ est mixte et $S>0$ (le couple $AB$ est pur, pas $A$ seul).}
  {On a utilisé le logarithme \emph{népérien} (résultat en nats) ; en bits il faut $\log_2$, ce qui donne $0{,}92$.}

% ---- Question 100 (niveau 4, correct : B)
\qcmq[1]{Chap.~\ref{chap:bell}\,$\cdot$\,\S\,\ref{sec:bell}}
  {Pourquoi peut-on mesurer simultanément les deux « parités » $Z\otimes Z$ et $X\otimes X$ de deux qubits, alors qu'on ne peut pas mesurer simultanément $Z$ et $X$ sur un seul qubit ?}
  {Parce que l'incertitude de Heisenberg ne s'applique pas à deux qubits}
  {Parce que $Z\otimes Z$ et $X\otimes X$ commutent : les deux signes $-1$ de $ZX=-XZ$ se compensent}
  {Parce que les deux qubits sont indépendants}
  {Parce que $Z\otimes Z$ et $X\otimes X$ agissent sur des qubits différents}
  {B}
\qcmx
  {Elle s'applique toujours : $Z\otimes\mathrm{Id}$ et $X\otimes\mathrm{Id}$ ne commutent pas. C'est la commutation de ces deux opérateurs précis qui permet la mesure simultanée.}
  {$(Z\otimes Z)(X\otimes X)=ZX\otimes ZX=(-XZ)\otimes(-XZ)=XZ\otimes XZ=(X\otimes X)(Z\otimes Z)$. Ils ont donc une base propre commune : la base de Bell.}
  {Dans un état de Bell les qubits sont intriqués, donc loin d'être indépendants ; la commutation est une propriété des opérateurs, pas de l'état.}
  {Chacun des deux opérateurs agit sur \emph{les deux} qubits à la fois. (Ce sont $Z\otimes\mathrm{Id}$ et $\mathrm{Id}\otimes X$, sur des qubits différents, qui commutent pour cette raison.)}

\label{qcm:fin}

\clearpage
\ifqcmcorrige

% ======================================================================
\subsection{Clé des réponses et fiche de score}\label{sec:qcm-cle}
% ======================================================================
\noindent\textbf{Cette partie donne les réponses : ne la lisez qu'après avoir répondu à toutes les
questions.}

\paragraph{Clé des réponses.} La lettre de la bonne réponse de la question $n$ se lit à la ligne
qui contient $n$, dans la colonne du rang de $n$ dans la ligne (question~24 : ligne 21--30,
colonne 4).
\begin{center}
\renewcommand{\arraystretch}{1.25}
\setlength{\tabcolsep}{5.5pt}
\begin{tabular}{@{}l*{10}{c}@{}}
\toprule
 & 1 & 2 & 3 & 4 & 5 & 6 & 7 & 8 & 9 & 10 \\
\midrule
1--10 & \qcmkey{1} & \qcmkey{2} & \qcmkey{3} & \qcmkey{4} & \qcmkey{5} & \qcmkey{6} & \qcmkey{7} & \qcmkey{8} & \qcmkey{9} & \qcmkey{10} \\
11--20 & \qcmkey{11} & \qcmkey{12} & \qcmkey{13} & \qcmkey{14} & \qcmkey{15} & \qcmkey{16} & \qcmkey{17} & \qcmkey{18} & \qcmkey{19} & \qcmkey{20} \\
21--30 & \qcmkey{21} & \qcmkey{22} & \qcmkey{23} & \qcmkey{24} & \qcmkey{25} & \qcmkey{26} & \qcmkey{27} & \qcmkey{28} & \qcmkey{29} & \qcmkey{30} \\
31--40 & \qcmkey{31} & \qcmkey{32} & \qcmkey{33} & \qcmkey{34} & \qcmkey{35} & \qcmkey{36} & \qcmkey{37} & \qcmkey{38} & \qcmkey{39} & \qcmkey{40} \\
41--50 & \qcmkey{41} & \qcmkey{42} & \qcmkey{43} & \qcmkey{44} & \qcmkey{45} & \qcmkey{46} & \qcmkey{47} & \qcmkey{48} & \qcmkey{49} & \qcmkey{50} \\
51--60 & \qcmkey{51} & \qcmkey{52} & \qcmkey{53} & \qcmkey{54} & \qcmkey{55} & \qcmkey{56} & \qcmkey{57} & \qcmkey{58} & \qcmkey{59} & \qcmkey{60} \\
61--70 & \qcmkey{61} & \qcmkey{62} & \qcmkey{63} & \qcmkey{64} & \qcmkey{65} & \qcmkey{66} & \qcmkey{67} & \qcmkey{68} & \qcmkey{69} & \qcmkey{70} \\
71--80 & \qcmkey{71} & \qcmkey{72} & \qcmkey{73} & \qcmkey{74} & \qcmkey{75} & \qcmkey{76} & \qcmkey{77} & \qcmkey{78} & \qcmkey{79} & \qcmkey{80} \\
81--90 & \qcmkey{81} & \qcmkey{82} & \qcmkey{83} & \qcmkey{84} & \qcmkey{85} & \qcmkey{86} & \qcmkey{87} & \qcmkey{88} & \qcmkey{89} & \qcmkey{90} \\
91--100 & \qcmkey{91} & \qcmkey{92} & \qcmkey{93} & \qcmkey{94} & \qcmkey{95} & \qcmkey{96} & \qcmkey{97} & \qcmkey{98} & \qcmkey{99} & \qcmkey{100} \\
\bottomrule
\end{tabular}
\end{center}

\paragraph{Fiche de score.} Reportez le nombre de bonnes réponses de chaque niveau.
\begin{center}
\renewcommand{\arraystretch}{1.6}
\begin{tabular}{@{}l l c c c@{}}
\toprule
\textbf{Niveau} & \textbf{Questions} & \textbf{Bonnes réponses} & \textbf{Sur} & \textbf{Pourcentage}\\
\midrule
1. Débutant & 1--30 & \makebox[2.6cm]{\dotfill} & 30 & \makebox[2cm]{\dotfill}\,\%\\
2. Intermédiaire & 31--60 & \makebox[2.6cm]{\dotfill} & 30 & \makebox[2cm]{\dotfill}\,\%\\
3. Avancé & 61--90 & \makebox[2.6cm]{\dotfill} & 30 & \makebox[2cm]{\dotfill}\,\%\\
4. Maîtrise & 91--100 & \makebox[2.6cm]{\dotfill} & 10 & \makebox[2cm]{\dotfill}\,\%\\
\midrule
\textbf{Total} & 1--100 & \makebox[2.6cm]{\dotfill} & 100 & \makebox[2cm]{\dotfill}\,\%\\
\bottomrule
\end{tabular}
\end{center}
Repères : au moins $80\,\%$ à un niveau, il est acquis et on peut passer au suivant ; entre $50$ et
$80\,\%$, revoyez les sections des questions manquées ; en dessous, relisez le chapitre concerné et
son récapitulatif avant de recommencer.

\paragraph{Où revoir ?} Questions classées par chapitre du cours.
\begin{center}
\renewcommand{\arraystretch}{1.3}
\begin{tabular}{@{}l >{\raggedright\arraybackslash}p{5.9cm} >{\raggedright\arraybackslash}p{6.9cm}@{}}
\toprule
\textbf{Chapitre} & \textbf{Titre} & \textbf{Questions}\\
\midrule
Chap.~\ref{chap:histoire} & Histoire et mécanique quantique & 1, 2, 3, 4, 5, 6, 31, 32, 33, 34, 35, 61, 62, 63, 91 \\
Chap.~\ref{chap:dirac} & Notation de Dirac et espace de Hilbert & 7, 8, 9, 10, 11, 12, 36, 37, 38, 39, 40, 64, 65, 66, 67, 68, 92 \\
Chap.~\ref{chap:operateurs} & Opérateurs, matrice de densité, Bloch, produit tensoriel & 13, 14, 15, 16, 17, 18, 41, 42, 43, 44, 45, 46, 69, 70, 71, 72, 73, 74, 75, 93, 94 \\
Chap.~\ref{chap:portes} & Portes quantiques et systèmes multi-qubits & 19, 20, 21, 22, 23, 24, 47, 48, 49, 50, 51, 76, 77, 78, 79, 80, 81, 82, 95, 96 \\
Chap.~\ref{chap:bell} & États quantiques, intrication, Bell & 25, 26, 27, 28, 29, 30, 52, 53, 54, 55, 56, 57, 58, 59, 60, 83, 84, 85, 86, 87, 88, 89, 90, 97, 98, 99, 100 \\
\bottomrule
\end{tabular}
\end{center}

\clearpage
\subsection{Corrigé du niveau 1 : débutant (questions 1 à 30)}\label{sec:qcm-c1}
\qcmcorr{1}
\qcmcorr{2}
\qcmcorr{3}
\qcmcorr{4}
\qcmcorr{5}
\qcmcorr{6}
\qcmcorr{7}
\qcmcorr{8}
\qcmcorr{9}
\qcmcorr{10}
\qcmcorr{11}
\qcmcorr{12}
\qcmcorr{13}
\qcmcorr{14}
\qcmcorr{15}
\qcmcorr{16}
\qcmcorr{17}
\qcmcorr{18}
\qcmcorr{19}
\qcmcorr{20}
\qcmcorr{21}
\qcmcorr{22}
\qcmcorr{23}
\qcmcorr{24}
\qcmcorr{25}
\qcmcorr{26}
\qcmcorr{27}
\qcmcorr{28}
\qcmcorr{29}
\qcmcorr{30}

\clearpage
\subsection{Corrigé du niveau 2 : intermédiaire (questions 31 à 60)}\label{sec:qcm-c2}
\qcmcorr{31}
\qcmcorr{32}
\qcmcorr{33}
\qcmcorr{34}
\qcmcorr{35}
\qcmcorr{36}
\qcmcorr{37}
\qcmcorr{38}
\qcmcorr{39}
\qcmcorr{40}
\qcmcorr{41}
\qcmcorr{42}
\qcmcorr{43}
\qcmcorr{44}
\qcmcorr{45}
\qcmcorr{46}
\qcmcorr{47}
\qcmcorr{48}
\qcmcorr{49}
\qcmcorr{50}
\qcmcorr{51}
\qcmcorr{52}
\qcmcorr{53}
\qcmcorr{54}
\qcmcorr{55}
\qcmcorr{56}
\qcmcorr{57}
\qcmcorr{58}
\qcmcorr{59}
\qcmcorr{60}

\clearpage
\subsection{Corrigé du niveau 3 : avancé (questions 61 à 90)}\label{sec:qcm-c3}
\qcmcorr{61}
\qcmcorr{62}
\qcmcorr{63}
\qcmcorr{64}
\qcmcorr{65}
\qcmcorr{66}
\qcmcorr{67}
\qcmcorr{68}
\qcmcorr{69}
\qcmcorr{70}
\qcmcorr{71}
\qcmcorr{72}
\qcmcorr{73}
\qcmcorr{74}
\qcmcorr{75}
\qcmcorr{76}
\qcmcorr{77}
\qcmcorr{78}
\qcmcorr{79}
\qcmcorr{80}
\qcmcorr{81}
\qcmcorr{82}
\qcmcorr{83}
\qcmcorr{84}
\qcmcorr{85}
\qcmcorr{86}
\qcmcorr{87}
\qcmcorr{88}
\qcmcorr{89}
\qcmcorr{90}

\clearpage
\subsection{Corrigé du niveau 4 : maîtrise (questions 91 à 100)}\label{sec:qcm-c4}
\qcmcorr{91}
\qcmcorr{92}
\qcmcorr{93}
\qcmcorr{94}
\qcmcorr{95}
\qcmcorr{96}
\qcmcorr{97}
\qcmcorr{98}
\qcmcorr{99}
\qcmcorr{100}
\fi
. Un chapitre = une \section.
%
%  Organisation :
%    - les énoncés (niveaux 1 à 4) occupent des pages à part ;
%    - le corrigé (clé des réponses, explications, pièges, fiche de score)
%      commence sur une nouvelle page : on peut imprimer les énoncés seuls ;
%    - chaque question est écrite avec \qcmq puis \qcmx (voir plus bas) : le même
%      code produit l'énoncé (cases à cocher) et, plus loin, le corrigé.
%
%  Bascules (définies dans main.tex ; à placer avant % =====================================================================
%  Chapitre 7 : QCM sur le cours (100 questions, 4 niveaux).
%  Fichier importé par main.tex avec \include{qcm}. Un chapitre = une \section.
%
%  Organisation :
%    - les énoncés (niveaux 1 à 4) occupent des pages à part ;
%    - le corrigé (clé des réponses, explications, pièges, fiche de score)
%      commence sur une nouvelle page : on peut imprimer les énoncés seuls ;
%    - chaque question est écrite avec \qcmq puis \qcmx (voir plus bas) : le même
%      code produit l'énoncé (cases à cocher) et, plus loin, le corrigé.
%
%  Bascules (définies dans main.tex ; à placer avant \include{qcm}) :
%    \qcmrefsfalse     masque la référence « Chap. · § » sous chaque énoncé ;
%    \qcmcorrigefalse  supprime le corrigé du document (QCM à distribuer seul).
%
%  Ce fichier est produit par outils_qcm/gen_qcm.py à partir de la banque de
%  questions (fichiers texte), mais il reste un fichier LaTeX ordinaire que l'on
%  peut modifier à la main. Pour ajouter une question à la main : insérer un
%  couple \qcmq ... \qcmx dans le bon niveau (la numérotation est automatique),
%  puis mettre à jour (1) les bornes « questions a à b » des titres de niveaux,
%  (2) les lignes \qcmcorr{n} du corrigé, (3) le tableau « Où revoir ? ».
% =====================================================================

% Macros de Dirac (déjà définies dans main.tex ; ce bloc évite l'erreur
% « Undefined control sequence \ket » si main.tex est une ancienne version).
\providecommand{\ket}[1]{|#1\rangle}
\providecommand{\bra}[1]{\langle #1|}
\providecommand{\braket}[2]{\langle #1|#2\rangle}

\makeatletter
% bascules : normalement déjà définies dans main.tex (valeurs par défaut : vrai)
\@ifundefined{ifqcmrefs}{\newif\ifqcmrefs\qcmrefstrue}{}
\@ifundefined{ifqcmcorrige}{\newif\ifqcmcorrige\qcmcorrigetrue}{}
\newcounter{qcmq}

% case à cocher (3 mm) pour le crayon
\newcommand{\qcmbox}{{\setlength{\fboxsep}{0pt}\setlength{\fboxrule}{0.5pt}%
  \raisebox{-0.4mm}{\fbox{\rule{0pt}{3mm}\rule{3mm}{0pt}}}}}

% une proposition : case, lettre, texte (retrait suspendu)
\newcommand{\qcm@opt}[2]{\noindent\hangindent=3.2em \hangafter=1
  \hspace*{0.3em}\qcmbox\hspace{0.45em}\textbf{#1)}\hspace{0.4em}#2\par\vspace{2.5pt}}

% \qcmq[colonnes]{ref}{énoncé}{A}{B}{C}{D}{bonne lettre}   (TeX limite les macros à 9 paramètres)
% \qcmx{com. A}{com. B}{com. C}{com. D}  : commentaires, placés juste après \qcmq
%   - com. de la bonne lettre = explication ; com. des autres = piège ;
%   - colonnes = 1 (une proposition par ligne) ou 2 (grille 2x2).
\newcommand{\qcmq}[8][1]{%
  \stepcounter{qcmq}%
  \expandafter\gdef\csname qcm@ref@\theqcmq\endcsname{#2}%
  \expandafter\gdef\csname qcm@rep@\theqcmq\endcsname{#8}%
  \expandafter\gdef\csname qcm@o@\theqcmq @A\endcsname{#4}%
  \expandafter\gdef\csname qcm@o@\theqcmq @B\endcsname{#5}%
  \expandafter\gdef\csname qcm@o@\theqcmq @C\endcsname{#6}%
  \expandafter\gdef\csname qcm@o@\theqcmq @D\endcsname{#7}%
  \par\vspace{7pt plus 3pt minus 2pt}\noindent
  \begin{minipage}{\linewidth}%
    \noindent{\color{insarouge}\bfseries Question~\theqcmq}%
    \ifqcmrefs\hfill{\footnotesize\color{insagris}#2}\fi\par\nopagebreak
    \vspace{2pt}\noindent #3\par\nopagebreak\vspace{3pt}%
    \ifnum#1=2
      \begin{minipage}[t]{0.495\linewidth}\qcm@opt{A}{#4}\end{minipage}\hfill
      \begin{minipage}[t]{0.495\linewidth}\qcm@opt{B}{#5}\end{minipage}\par\vspace{2pt}
      \begin{minipage}[t]{0.495\linewidth}\qcm@opt{C}{#6}\end{minipage}\hfill
      \begin{minipage}[t]{0.495\linewidth}\qcm@opt{D}{#7}\end{minipage}\par
    \else
      \qcm@opt{A}{#4}\qcm@opt{B}{#5}\qcm@opt{C}{#6}\qcm@opt{D}{#7}%
    \fi
  \end{minipage}\par}

\newcommand{\qcmx}[4]{%
  \expandafter\gdef\csname qcm@c@\theqcmq @A\endcsname{#1}%
  \expandafter\gdef\csname qcm@c@\theqcmq @B\endcsname{#2}%
  \expandafter\gdef\csname qcm@c@\theqcmq @C\endcsname{#3}%
  \expandafter\gdef\csname qcm@c@\theqcmq @D\endcsname{#4}}

% lettre de la bonne réponse de la question n (« ? » si la question n'existe pas)
\newcommand{\qcmkey}[1]{\@ifundefined{qcm@rep@#1}{?}{\csname qcm@rep@#1\endcsname}}

% une mauvaise réponse du corrigé (n° de question, lettre) : ignorée si c'est la bonne.
% (pas de \@for : « := » ne fonctionne pas avec le « : » rendu actif par babel-french)
\newcommand{\qcm@wrongone}[2]{%
  \edef\qcm@a{#2}\edef\qcm@b{\csname qcm@rep@#1\endcsname}%
  \ifx\qcm@a\qcm@b\else
    \noindent\hangindent=1.4em\hangafter=1 {\color{insagris}\textbf{#2)}}\enspace
    \csname qcm@o@#1@#2\endcsname\enspace--\enspace\emph{\csname qcm@c@#1@#2\endcsname}\par
  \fi}

% bloc de corrigé de la question n
\newcommand{\qcmcorr}[1]{%
  \par\vspace{6pt plus 2pt minus 2pt}\noindent
  \begin{minipage}{\linewidth}\small
    \noindent{\color{insarouge}\bfseries Question~#1}\enspace
    \textbf{Réponse : \csname qcm@rep@#1\endcsname}\hfill
    {\footnotesize\color{insagris}\csname qcm@ref@#1\endcsname}\par\nopagebreak
    \vspace{1pt}\noindent\hangindent=1.4em\hangafter=1
    \textbf{\csname qcm@rep@#1\endcsname)}\enspace
    \csname qcm@o@#1@\csname qcm@rep@#1\endcsname\endcsname\enspace---\enspace
    \csname qcm@c@#1@\csname qcm@rep@#1\endcsname\endcsname\par
    \qcm@wrongone{#1}{A}\qcm@wrongone{#1}{B}\qcm@wrongone{#1}{C}\qcm@wrongone{#1}{D}%
  \end{minipage}\par}
\makeatother

% =====================================================================
\section{QCM sur le cours}\label{chap:qcm}
% =====================================================================

\subsection{Mode d'emploi}\label{sec:qcm-mode}

Ce QCM compte \textbf{100 questions} sur les chapitres~\ref{chap:histoire} à~\ref{chap:bell},
classées en quatre niveaux. Chaque question propose quatre réponses dont \textbf{une seule} est
exacte : cochez au crayon la case de la proposition choisie.
\ifqcmrefs À droite du numéro de chaque question, une référence (chapitre, section) indique la
partie du cours concernée.\else\ifqcmcorrige{} La référence de chaque question (chapitre, section)
est donnée dans le corrigé.\fi\fi

\begin{center}
\renewcommand{\arraystretch}{1.3}
\begin{tabular}{@{}llc>{\raggedright\arraybackslash}p{8.2cm}@{}}
\toprule
\textbf{Niveau} & \textbf{Questions} & \textbf{Nombre} & \textbf{Ce qui est demandé}\\
\midrule
1. Débutant & 1 à 30 & 30 & définitions, notations, résultats du cours en une étape\\
2. Intermédiaire & 31 à 60 & 30 & petits calculs, propriétés, application directe d'un résultat\\
3. Avancé & 61 à 90 & 30 & calculs en plusieurs étapes, portes, circuits et protocoles\\
4. Maîtrise & 91 à 100 & 10 & raisonnements fins et pièges conceptuels, qui combinent plusieurs notions\\
\bottomrule
\end{tabular}
\end{center}

\begin{itemize}[nosep]
  \item \textbf{Sur papier} : imprimez seulement les pages~\pageref{chap:qcm} à~\pageref{qcm:fin}
        (mode d'emploi et énoncés, sections~\ref{sec:qcm-n1} à~\ref{sec:qcm-n4}).%
        \ifqcmcorrige{} Le corrigé commence sur une page distincte, à la page~\pageref{sec:qcm-cle}.\fi
  \item \textbf{Barème} : 1 point par bonne réponse, 0 sinon (pas de point négatif). Calculatrice
        autorisée, mais rarement utile (quelques logarithmes et racines).
  \item \textbf{Correction} :%
        \ifqcmcorrige{} comparez avec la clé des réponses, puis lisez l'explication de chaque
        question, même réussie. Les « pièges » décrivent l'erreur qui mène à chaque mauvaise réponse.%
        \else{} le corrigé n'est pas inclus dans cette version du document.\fi
\end{itemize}

\bigskip
\noindent Nom : \makebox[6cm]{\dotfill}\hfill Date : \makebox[3.5cm]{\dotfill}\hfill Score : \makebox[1.6cm]{\dotfill}\,/\,100
\bigskip

\begin{remarque}[Conventions des énoncés]
\begin{itemize}[nosep]
  \item $\ket{ab}=\ket{a}_A\otimes\ket{b}_B$ : le premier chiffre est le qubit $A$, le second le
        qubit $B$ ; dans un circuit, le fil du haut est le dernier chiffre.
  \item Le CNOT de la section~\ref{sec:cnot} a pour contrôle le \emph{second} chiffre ; les portes
        contrôlées de la section~\ref{sec:controlees} ont pour contrôle le \emph{premier}. Chaque
        énoncé précise le contrôle.
  \item Sans précision, une mesure est faite dans la base de calcul $Z$. Les logarithmes des
        entropies sont en base $2$ (résultats en bits).
  \item $\ket{\pm}=\frac{1}{\sqrt2}(\ket0\pm\ket1)$ et $\ket{{\pm}\ii}=\frac{1}{\sqrt2}(\ket0\pm\ii\ket1)$.
\end{itemize}
\end{remarque}

\clearpage
%% ======================================================================
\subsection{Niveau 1 : débutant (questions 1 à 30)}\label{sec:qcm-n1}
%% ======================================================================
\noindent{\color{insagris}\small Niveau de difficulté : Débutant\quad $\bullet\circ\circ\circ$}\par\smallskip
% ---- Question 1 (niveau 1, correct : D)
\qcmq[2]{Chap.~\ref{chap:histoire}\,$\cdot$\,\S\,\ref{sec:dualite}}
  {En 1905, qui explique l'effet photo-électrique en postulant que la lumière est faite de quanta d'énergie $E=hf$, les photons ?}
  {Louis de Broglie}
  {Erwin Schrödinger}
  {Max Planck}
  {Albert Einstein}
  {D}
\qcmx
  {Il propose en 1924 une onde de matière ($\lambda=h/p$) pour les particules comme l'électron ; l'idée du photon date de 1905.}
  {Il écrit en 1926 l'équation d'onde de la matière, plus de vingt ans après l'article sur les photons.}
  {Il introduit le quantum d'énergie en 1900 pour le rayonnement du corps noir ; c'est Einstein qui en fait une propriété de la lumière elle-même.}
  {Il postule que la lumière échange son énergie par paquets $E=hf=\hbar\omega$ : c'est la première face de la dualité onde--corpuscule.}

% ---- Question 2 (niveau 1, correct : A)
\qcmq[1]{Chap.~\ref{chap:histoire}\,$\cdot$\,\S\,\ref{sec:born}}
  {Selon l'interprétation de Max Born, que représente $|\Psi(x,t)|^2$ ?}
  {La densité de probabilité de présence de la particule en $x$ à l'instant $t$}
  {L'énergie cinétique de la particule}
  {L'amplitude de l'onde de matière}
  {La vitesse de la particule}
  {A}
\qcmx
  {$|\Psi(x,t)|^2\,\dd x$ est la probabilité de trouver la particule entre $x$ et $x+\dd x$ ; $\Psi$ lui-même n'est pas mesurable.}
  {L'énergie est associée à l'opérateur hamiltonien $\Hop$ ; $|\Psi|^2$ n'a pas la dimension d'une énergie.}
  {$\Psi$ est l'amplitude (complexe) ; $|\Psi|^2$ en est le carré du module.}
  {La vitesse s'obtient avec l'opérateur $\hat p=-\ii\hbar\,\partial_x$, pas avec $|\Psi|^2$.}

% ---- Question 3 (niveau 1, correct : D)
\qcmq[2]{Chap.~\ref{chap:histoire}\,$\cdot$\,\S\,\ref{sec:dualite}}
  {Une particule de quantité de mouvement $p$ est associée, d'après de Broglie, à une onde de longueur d'onde $\lambda$. Quelle relation est correcte ?}
  {$\lambda=hp$}
  {$\lambda=p/h$}
  {$\lambda=\hbar/p$}
  {$\lambda=h/p$}
  {D}
\qcmx
  {Produit au lieu du quotient : $\lambda$ grandirait avec $p$, le contraire de ce qu'on observe.}
  {Quotient inversé : $p/h=1/\lambda$ s'exprime en m$^{-1}$, pas en mètres.}
  {Confusion entre $h$ et $\hbar=h/2\pi$ : $\hbar/p=\lambda/2\pi$ est l'inverse du nombre d'onde $k$.}
  {Relation de de Broglie : plus la particule est rapide ou lourde, plus $\lambda$ est petite (balle de $0{,}1$~kg à $10$~m/s : $\lambda\approx10^{-34}$~m, inobservable).}

% ---- Question 4 (niveau 1, correct : C)
\qcmq[2]{Chap.~\ref{chap:histoire}\,$\cdot$\,\S\,\ref{sec:construction}}
  {Quelle est l'équation de Schrödinger dépendante du temps pour un état $\Psi$ et un hamiltonien $\Hop$ ?}
  {$-\ii\hbar\,\partial_t\Psi=\Hop\Psi$}
  {$\hbar\,\partial_t\Psi=\Hop\Psi$}
  {$\ii\hbar\,\partial_t\Psi=\Hop\Psi$}
  {$\ii\hbar\,\partial_x\Psi=\Hop\Psi$}
  {C}
\qcmx
  {Erreur de signe : la solution serait $\ee^{+\ii\Hop t/\hbar}\Psi(0)$, qui décrit l'évolution à rebours du temps.}
  {Sans le facteur $\ii$, les solutions sont des exponentielles réelles (croissantes ou décroissantes) et la norme n'est pas conservée.}
  {Elle relie l'évolution de l'état à l'énergie ; pour un $\Hop$ indépendant du temps, sa solution est $\Psi(t)=\ee^{-\ii\Hop t/\hbar}\Psi(0)$ (chapitre~\ref{chap:portes}).}
  {C'est la dérivée en $t$ qui décrit l'évolution ; la dérivée en $x$ sert à construire $\hat p=-\ii\hbar\,\partial_x$.}

% ---- Question 5 (niveau 1, correct : A)
\qcmq[1]{Chap.~\ref{chap:histoire}\,$\cdot$\,\S\,\ref{sec:construction}}
  {Quel opérateur représente l'énergie totale d'un système quantique ?}
  {L'hamiltonien $\Hop$}
  {L'opérateur quantité de mouvement $\hat p$}
  {L'opérateur position $\hat x$}
  {L'opérateur identité}
  {A}
\qcmx
  {Somme de l'énergie cinétique $-\frac{\hbar^2}{2m}\partial_x^2$ et de l'énergie potentielle $V$ ; ses valeurs propres sont les énergies mesurables.}
  {Il représente l'impulsion ; l'énergie cinétique est $\hat p^2/2m$, et il faut y ajouter $V$.}
  {C'est la multiplication par $x$ : il donne la position, pas l'énergie.}
  {Il vaut $1$ pour tout état : il ne contient aucune information sur l'énergie.}

% ---- Question 6 (niveau 1, correct : C)
\qcmq[1]{Chap.~\ref{chap:histoire}\,$\cdot$\,\S\,\ref{sec:histoire}}
  {Dans quel ordre chronologique ces avancées ont-elles eu lieu ?}
  {de Broglie $\to$ Planck $\to$ Schrödinger $\to$ Bell}
  {Planck $\to$ Schrödinger $\to$ de Broglie $\to$ Bell}
  {Planck (quantum d'énergie) $\to$ de Broglie (onde de matière) $\to$ Schrödinger (équation d'onde) $\to$ Bell (inégalités)}
  {Planck $\to$ de Broglie $\to$ Bell $\to$ Schrödinger}
  {C}
\qcmx
  {Le quantum d'énergie de Planck (1900) précède l'onde de matière de de Broglie (1924).}
  {L'équation de Schrödinger (1926) généralise l'idée d'onde de matière de de Broglie (1924) : elle vient après.}
  {Planck 1900, de Broglie 1924, Schrödinger 1926, Bell 1964.}
  {Les inégalités de Bell (1964) viennent près de quarante ans après l'équation de Schrödinger.}

% ---- Question 7 (niveau 1, correct : C)
\qcmq[1]{Chap.~\ref{chap:dirac}\,$\cdot$\,\S\,\ref{sec:dirac}}
  {En notation de Dirac, comment appelle-t-on $\ket{\Psi}$ et comment le représente-t-on pour un qubit ?}
  {Un bra, représenté par un vecteur ligne à deux composantes complexes}
  {Un ket, représenté par un nombre complexe}
  {Un ket, représenté par un vecteur colonne à deux composantes complexes}
  {Un bra, représenté par une matrice $2\times2$}
  {C}
\qcmx
  {Le bra $\bra{\Psi}$ est le conjugué transposé du ket : un vecteur ligne, mais ce n'est pas la notation $\ket{\Psi}$.}
  {Un scalaire ne peut pas décrire une superposition : il faut deux amplitudes $\alpha$ et $\beta$.}
  {Le ket est le vecteur colonne $\binom{\alpha}{\beta}$ ; son bra $\bra{\Psi}$ est le vecteur ligne conjugué.}
  {Une matrice est un opérateur ; le bra est une matrice ligne $1\times2$.}

% ---- Question 8 (niveau 1, correct : B)
\qcmq[2]{Chap.~\ref{chap:dirac}\,$\cdot$\,\S\,\ref{sec:qubit}}
  {Quelle condition les amplitudes $\alpha$ et $\beta$ de $\alpha\ket{0}+\beta\ket{1}$ doivent-elles vérifier pour que l'état soit normalisé ?}
  {$\alpha+\beta=1$}
  {$|\alpha|^2+|\beta|^2=1$}
  {$\alpha^2+\beta^2=1$}
  {$|\alpha|+|\beta|=1$}
  {B}
\qcmx
  {On somme les amplitudes (complexes, parfois négatives) ; ce sont leurs modules au carré qui sont des probabilités.}
  {Les probabilités de mesure $P(0)=|\alpha|^2$ et $P(1)=|\beta|^2$ doivent sommer à $1$.}
  {Oubli du module : pour $\alpha=\frac{1}{\sqrt2}$ et $\beta=\frac{\ii}{\sqrt2}$, $\alpha^2+\beta^2=0$ alors que l'état est normalisé.}
  {Somme des modules au lieu des carrés : $\alpha=\beta=\frac{1}{\sqrt2}$ est normalisé, mais $|\alpha|+|\beta|=\sqrt2$.}

% ---- Question 9 (niveau 1, correct : C)
\qcmq[1]{Chap.~\ref{chap:dirac}\,$\cdot$\,\S\,\ref{sec:projecteurs}}
  {On mesure dans la base $Z$ un qubit préparé dans $\ket{+}=\frac{1}{\sqrt2}(\ket{0}+\ket{1})$. Que se passe-t-il ?}
  {On lit $+$ et le qubit reste dans $\ket{+}$}
  {On lit toujours $0$}
  {On lit $0$ ou $1$, chacun avec la probabilité $\frac12$, et l'état devient $\ket{0}$ ou $\ket{1}$}
  {On lit $0$ et $1$ à la fois}
  {C}
\qcmx
  {Une mesure en base $Z$ ne donne que $0$ ou $1$ ; elle modifie l'état.}
  {Les deux amplitudes ont le même module : $P(0)=P(1)=\frac12$.}
  {Postulat de la mesure : la superposition s'effondre sur l'état propre correspondant au résultat.}
  {Chaque mesure donne un seul résultat ; la superposition ne s'observe qu'à travers les statistiques.}

% ---- Question 10 (niveau 1, correct : B)
\qcmq[2]{Chap.~\ref{chap:dirac}\,$\cdot$\,\S\,\ref{sec:qubit}}
  {Que vaut le produit scalaire $\braket{0}{1}$ ?}
  {$1$}
  {$0$}
  {$\frac{1}{\sqrt2}$}
  {$\ket{01}$}
  {B}
\qcmx
  {C'est la valeur de $\braket{0}{0}$ et de $\braket{1}{1}$ (vecteurs normés), pas de $\braket{0}{1}$.}
  {$\ket{0}$ et $\ket{1}$ sont orthogonaux : $\begin{pmatrix}1&0\end{pmatrix}\begin{pmatrix}0\\1\end{pmatrix}=0$.}
  {C'est la valeur de $\braket{0}{+}$, pas de $\braket{0}{1}$.}
  {Confusion avec le produit tensoriel : un produit scalaire est un nombre, pas un ket.}

% ---- Question 11 (niveau 1, correct : D)
\qcmq[2]{Chap.~\ref{chap:dirac}\,$\cdot$\,\S\,\ref{sec:qubit}}
  {Un qubit est dans l'état $\ket{-}=\frac{1}{\sqrt2}(\ket{0}-\ket{1})$. Quelle est la probabilité de mesurer $1$ en base $Z$ ?}
  {$0$}
  {$-\frac12$}
  {$\frac{1}{\sqrt2}$}
  {$\frac12$}
  {D}
\qcmx
  {Le signe $-$ ne « retire » pas $\ket{1}$ de la superposition.}
  {Une probabilité est toujours positive : on prend le module au carré de l'amplitude.}
  {C'est l'amplitude en module, pas la probabilité (il faut élever au carré).}
  {$P(1)=\left|-\frac{1}{\sqrt2}\right|^2=\frac12$ : le signe est une phase relative, qui disparaît dans le module.}

% ---- Question 12 (niveau 1, correct : B)
\qcmq[2]{Chap.~\ref{chap:dirac}\,$\cdot$\,\S\,\ref{sec:projecteurs}}
  {Quel est le projecteur sur l'état $\ket{0}$ ?}
  {$\ket{0}\bra{1}=\begin{pmatrix}0&1\\0&0\end{pmatrix}$}
  {$\hat P_0=\ket{0}\bra{0}=\begin{pmatrix}1&0\\0&0\end{pmatrix}$}
  {$\braket{0}{0}=1$}
  {$\begin{pmatrix}0&0\\0&1\end{pmatrix}$}
  {B}
\qcmx
  {Cet opérateur envoie $\ket{1}$ sur $\ket{0}$ ; son carré est nul : ce n'est pas un projecteur.}
  {$\hat P_0\ket{\Psi}=\alpha\ket{0}$ : il garde la composante sur $\ket{0}$ ; $\hat P_0^2=\hat P_0$.}
  {C'est un nombre, pas un opérateur.}
  {C'est $\ket{1}\bra{1}$, le projecteur sur $\ket{1}$.}

% ---- Question 13 (niveau 1, correct : A)
\qcmq[1]{Chap.~\ref{chap:operateurs}\,$\cdot$\,\S\,\ref{sec:bloch}}
  {Où se trouvent les états $\ket{0}$ et $\ket{1}$ sur la sphère de Bloch ?}
  {Aux deux pôles, Nord et Sud}
  {Aux deux extrémités de l'axe $x$}
  {Aux deux extrémités de l'axe $y$}
  {En deux points voisins de l'équateur}
  {A}
\qcmx
  {$\ket{0}$ au pôle Nord ($\theta=0$), $\ket{1}$ au pôle Sud ($\theta=\pi$) : deux états orthogonaux sont diamétralement opposés.}
  {Ce sont les places de $\ket{+}$ et $\ket{-}$.}
  {Ce sont les places de $\ket{{+}\ii}$ et $\ket{{-}\ii}$.}
  {Des états orthogonaux sont antipodaux, jamais voisins.}

% ---- Question 14 (niveau 1, correct : A)
\qcmq[2]{Chap.~\ref{chap:operateurs}\,$\cdot$\,\S\,\ref{sec:tenseur}}
  {Quelle est la dimension de l'espace de Hilbert de $n$ qubits ?}
  {$2^n$}
  {$2n$}
  {$n^2$}
  {$2^n-1$}
  {A}
\qcmx
  {Le produit tensoriel multiplie les dimensions : $2\times2\times\cdots\times2$ ; deux qubits : $4$, trois qubits : $8$.}
  {On additionne les dimensions au lieu de les multiplier ; pour $n=3$ on trouverait $6$ au lieu de $8$.}
  {Juste pour $n=2$ par coïncidence seulement : pour $n=3$ on trouverait $9$ au lieu de $8$.}
  {Il y a bien $2^n$ vecteurs dans la base de calcul, $\ket{0\cdots0}$ compris.}

% ---- Question 15 (niveau 1, correct : C)
\qcmq[1]{Chap.~\ref{chap:operateurs}\,$\cdot$\,\S\,\ref{sec:rho}}
  {Que vaut toujours la trace d'une matrice de densité $\rho$ ?}
  {$0$}
  {$1$ pour un état pur, et moins de $1$ pour un état mixte}
  {$1$}
  {$2$, la dimension de l'espace}
  {C}
\qcmx
  {C'est la trace des matrices de Pauli, pas d'une matrice de densité.}
  {Confusion avec $\operatorname{Tr}(\rho^2)$ : c'est la pureté qui distingue les deux cas, la trace vaut toujours $1$.}
  {C'est la somme des probabilités de mesure (les populations), que l'état soit pur ou mixte.}
  {C'est la trace de l'identité $\operatorname{Tr}\mathrm{Id}=2$, pas celle de $\rho$.}

% ---- Question 16 (niveau 1, correct : B)
\qcmq[1]{Chap.~\ref{chap:operateurs}\,$\cdot$\,\S\,\ref{sec:pur-mixte}}
  {Quelle est la valeur de $\operatorname{Tr}(\rho^2)$ pour un état pur ?}
  {$0$}
  {$1$}
  {$\frac12$}
  {Une valeur strictement inférieure à $1$}
  {B}
\qcmx
  {$\rho^2=0$ imposerait $\rho=0$, ce qui contredit $\operatorname{Tr}\rho=1$.}
  {Un état pur s'écrit $\rho=\ket{\Psi}\bra{\Psi}$, donc $\rho^2=\rho$ et $\operatorname{Tr}\rho^2=\operatorname{Tr}\rho=1$.}
  {C'est la valeur minimale pour un qubit, atteinte par l'état maximalement mixte $\mathrm{Id}/2$.}
  {C'est la caractérisation d'un état mixte.}

% ---- Question 17 (niveau 1, correct : B)
\qcmq[1]{Chap.~\ref{chap:operateurs}\,$\cdot$\,\S\,\ref{sec:rho}}
  {Que représentent les éléments diagonaux $\rho_{00}$ et $\rho_{11}$ de la matrice de densité d'un qubit ?}
  {Les cohérences entre $\ket{0}$ et $\ket{1}$}
  {Les probabilités de mesurer $\ket{0}$ et $\ket{1}$ (les populations)}
  {Les valeurs propres de $\rho$}
  {Les amplitudes $\alpha$ et $\beta$}
  {B}
\qcmx
  {Les cohérences sont les éléments hors diagonale $\rho_{01}$ et $\rho_{10}$.}
  {$\rho_{00}=\bra{0}\rho\ket{0}=P(0)$ et $\rho_{11}=P(1)$, d'où $\operatorname{Tr}\rho=1$.}
  {Seulement si $\rho$ est diagonale dans la base de calcul : $\ket{+}\bra{+}$ a pour diagonale $(\frac12,\frac12)$ mais pour valeurs propres $(1,0)$.}
  {Pour un état pur, $\rho_{00}=|\alpha|^2$ : le module au carré, pas l'amplitude.}

% ---- Question 18 (niveau 1, correct : D)
\qcmq[1]{Chap.~\ref{chap:operateurs}\,$\cdot$\,\S\,\ref{sec:op-dirac}}
  {Pourquoi associe-t-on aux grandeurs mesurables (observables) des opérateurs hermitiens ?}
  {Parce qu'ils conservent la norme des vecteurs}
  {Parce qu'ils ont tous une trace nulle}
  {Parce que leurs valeurs propres valent toutes $\pm1$}
  {Parce que leurs valeurs propres sont réelles, comme les résultats d'une mesure}
  {D}
\qcmx
  {C'est la propriété des opérateurs unitaires (les portes), pas des opérateurs hermitiens.}
  {Vrai pour les matrices de Pauli, faux en général (l'identité est hermitienne, de trace $2$), et sans rapport avec le caractère mesurable.}
  {Vrai pour $X$, $Y$, $Z$ seulement ; l'hamiltonien du puits a les valeurs propres $E_n=n^2E_1$.}
  {Un opérateur hermitien ($A^\dagger=A$) a des valeurs propres réelles et des vecteurs propres orthogonaux ; les résultats possibles sont ses valeurs propres.}

% ---- Question 19 (niveau 1, correct : C)
\qcmq[2]{Chap.~\ref{chap:portes}\,$\cdot$\,\S\,\ref{sec:pauli}}
  {Quelle porte de Pauli réalise l'inversion de bit $\ket{0}\leftrightarrow\ket{1}$ ?}
  {$Z$}
  {$Y$}
  {$X$}
  {$H$}
  {C}
\qcmx
  {$Z$ laisse $\ket{0}$ et change $\ket{1}$ en $-\ket{1}$ : c'est une inversion de phase.}
  {$Y\ket{0}=\ii\ket{1}$ : elle inverse le bit mais ajoute une phase ; la porte d'inversion « pure » est $X$.}
  {$X\ket{0}=\ket{1}$ et $X\ket{1}=\ket{0}$ : c'est le NOT quantique.}
  {$H$ crée des superpositions : $H\ket{0}=\ket{+}$.}

% ---- Question 20 (niveau 1, correct : B)
\qcmq[2]{Chap.~\ref{chap:portes}\,$\cdot$\,\S\,\ref{sec:hadamard}}
  {Que donne la porte de Hadamard appliquée à $\ket{0}$ ?}
  {$\ket{1}$}
  {$\ket{+}=\frac{1}{\sqrt2}(\ket{0}+\ket{1})$}
  {$\ket{-}=\frac{1}{\sqrt2}(\ket{0}-\ket{1})$}
  {$\ket{0}$}
  {B}
\qcmx
  {C'est l'action de $X$.}
  {Superposition à poids égaux : $P(0)=P(1)=\frac12$ ; c'est la première étape de la plupart des algorithmes.}
  {C'est $H\ket{1}$ ; le signe $-$ n'apparaît qu'à partir de $\ket{1}$.}
  {$H$ n'est pas l'identité ; elle ne laisse inchangés ni $\ket{0}$ ni $\ket{1}$ (elle envoie $\ket{+}$ sur $\ket{0}$).}

% ---- Question 21 (niveau 1, correct : D)
\qcmq[2]{Chap.~\ref{chap:portes}\,$\cdot$\,\S\,\ref{sec:portes}}
  {Quelle propriété doit vérifier la matrice $U$ d'une porte quantique ?}
  {Elle est hermitienne : $U^\dagger=U$}
  {Sa trace vaut $1$}
  {Ses coefficients sont réels}
  {Elle est unitaire : $U^\dagger U=\mathrm{Id}$}
  {D}
\qcmx
  {Vrai pour $X$, $Y$, $Z$, $H$ mais faux pour $S$ et $T$ : ce n'est pas la condition générale.}
  {C'est une propriété des matrices de densité ; $\operatorname{Tr}X=0$ et $\operatorname{Tr}\mathrm{Id}=2$.}
  {$Y$ et $S$ sont des portes à coefficients complexes.}
  {Elle conserve la norme ($\braket{\Psi}{\Psi}=1$ reste vrai) et elle est réversible.}

% ---- Question 22 (niveau 1, correct : A)
\qcmq[2]{Chap.~\ref{chap:portes}\,$\cdot$\,\S\,\ref{sec:portes}}
  {Quel est l'inverse d'une porte quantique $U$ ?}
  {$U^{-1}=U^\dagger$ (transposée conjuguée)}
  {$U^{-1}=U^{\mathsf T}$ (transposée seule)}
  {$U^{-1}=U^{*}$ (conjuguée seule)}
  {$U^{-1}=-U$}
  {A}
\qcmx
  {Conséquence de $U^\dagger U=\mathrm{Id}$ ; par exemple $S^{-1}=S^\dagger=\mathrm{diag}(1,-\ii)$.}
  {Oubli de la conjugaison : $S^{\mathsf T}=S=\mathrm{diag}(1,\ii)$ alors que $S^{-1}=\mathrm{diag}(1,-\ii)$.}
  {Faux en général : $Y^*=-Y$ alors que $Y^{-1}=Y$.}
  {Faux : $X^{-1}=X\neq-X$.}

% ---- Question 23 (niveau 1, correct : C)
\qcmq[1]{Chap.~\ref{chap:portes}\,$\cdot$\,\S\,\ref{sec:cnot}}
  {Dans une porte CNOT, que devient le qubit cible lorsque le qubit de contrôle est dans l'état $\ket{0}$ ?}
  {Il est inversé par la porte $X$}
  {Il est remis à $\ket{0}$}
  {Il reste inchangé}
  {Il subit un changement de signe, comme avec $Z$}
  {C}
\qcmx
  {C'est ce qui se passe quand le contrôle vaut $\ket{1}$.}
  {Le CNOT est réversible : remettre la cible à zéro détruirait de l'information.}
  {Avec le contrôle à $\ket{0}$, la porte agit comme l'identité sur la cible.}
  {Ce serait la porte $Z$ contrôlée (CZ), pas le CNOT.}

% ---- Question 24 (niveau 1, correct : A)
\qcmq[2]{Chap.~\ref{chap:portes}\,$\cdot$\,\S\,\ref{sec:cnot}}
  {Le CNOT du cours a pour contrôle le second chiffre du ket et pour cible le premier : $\mathrm{CNOT}\ket{a,b}=\ket{a\oplus b,\,b}$. Que donne-t-il sur $\ket{0,1}$ ?}
  {$\ket{1,1}$}
  {$\ket{0,1}$}
  {$\ket{0,0}$}
  {$\ket{1,0}$}
  {A}
\qcmx
  {Le contrôle $b=1$ inverse la cible : $0\oplus1=1$.}
  {Le contrôle vaut $1$ : la cible est bien inversée.}
  {C'est le contrôle qui serait inversé : le contrôle n'est jamais modifié.}
  {Les deux qubits seraient inversés : seule la cible change.}

% ---- Question 25 (niveau 1, correct : D)
\qcmq[1]{Chap.~\ref{chap:bell}\,$\cdot$\,\S\,\ref{sec:bell}}
  {Quelle propriété caractérise les quatre états de Bell ?}
  {Ils sont séparables : chaque qubit a son propre état}
  {Ils décrivent un seul qubit}
  {Ils ne diffèrent que par une phase globale}
  {Ils sont maximalement intriqués et forment une base orthonormée de l'espace de deux qubits}
  {D}
\qcmx
  {$D=c_{00}c_{11}-c_{01}c_{10}=\pm\frac12\neq0$ : ils ne se factorisent pas.}
  {Ce sont des états de deux qubits (quatre amplitudes sur $\ket{00},\ket{01},\ket{10},\ket{11}$).}
  {Ils sont orthogonaux deux à deux : ce sont quatre états distincts, alors qu'une phase globale ne change pas l'état.}
  {Aucun n'est un produit d'états à un qubit ($|D|=\frac12$, valeur maximale) et ils sont orthogonaux deux à deux.}

% ---- Question 26 (niveau 1, correct : D)
\qcmq[1]{Chap.~\ref{chap:bell}\,$\cdot$\,\S\,\ref{sec:teleportation}}
  {Quel est l'objectif du protocole de téléportation quantique ?}
  {Copier l'état d'un qubit inconnu sur le qubit de Bob}
  {Transmettre une information plus vite que la lumière}
  {Déplacer la particule physique d'Alice vers Bob}
  {Transférer l'état d'un qubit inconnu d'Alice à Bob, avec une paire intriquée et deux bits classiques}
  {D}
\qcmx
  {La copie est interdite (non-clonage) : après le protocole, l'état n'existe plus que chez Bob.}
  {Sans les deux bits classiques, le qubit de Bob est dans l'état $\mathrm{Id}/2$ : il n'a reçu aucune information.}
  {C'est l'état quantique qui est transféré, pas la particule qui le porte.}
  {Le qubit lui-même n'est pas envoyé : son état est reconstruit chez Bob, et détruit chez Alice par la mesure.}

% ---- Question 27 (niveau 1, correct : B)
\qcmq[1]{Chap.~\ref{chap:bell}\,$\cdot$\,\S\,\ref{sec:teleportation}}
  {Dans la téléportation, que doit envoyer Alice à Bob sur un canal classique pour qu'il retrouve l'état ?}
  {Un qubit intriqué}
  {Deux bits classiques, le résultat de sa mesure}
  {Les amplitudes exactes $\alpha$ et $\beta$ de l'état}
  {Un seul bit}
  {B}
\qcmx
  {La paire intriquée est partagée avant le protocole ; on ne l'envoie pas après.}
  {Les quatre résultats $(a,q)$ correspondent aux quatre corrections possibles ($I$, $Z$, $X$, $X$ puis $Z$) : il faut deux bits.}
  {Alice ne connaît pas l'état à téléporter, et ces nombres complexes n'ont pas de description finie en bits.}
  {Un bit ne distingue que deux cas, alors que Bob doit choisir parmi quatre corrections.}

% ---- Question 28 (niveau 1, correct : A)
\qcmq[1]{Chap.~\ref{chap:bell}\,$\cdot$\,\S\,\ref{sec:bell} ; ex.\,\ref{ex:superdense}}
  {Quel est le but du codage superdense ?}
  {Transmettre deux bits classiques en envoyant un seul qubit, grâce à une paire intriquée partagée}
  {Transmettre l'état d'un qubit avec deux bits classiques}
  {Compresser plusieurs qubits en un seul}
  {Chiffrer un message de façon inviolable}
  {A}
\qcmx
  {Alice applique l'une des quatre portes $I$, $Z$, $X$, $XZ$ à son qubit de la paire, l'envoie, et Bob identifie l'état de Bell obtenu.}
  {C'est la téléportation, le protocole inverse.}
  {Le protocole ne compresse pas de qubits : il fait passer deux bits classiques par un canal à un qubit.}
  {Ce n'est pas un protocole de cryptographie : la paire intriquée sert à doubler la capacité du canal.}

% ---- Question 29 (niveau 1, correct : A)
\qcmq[1]{Chap.~\ref{chap:bell}\,$\cdot$\,\S\,\ref{sec:teleportation}}
  {Que dit le théorème de non-clonage ?}
  {Aucune opération ne peut copier parfaitement un état quantique inconnu quelconque}
  {Il est impossible de mesurer un qubit sans le perturber}
  {Il est impossible de créer de l'intrication}
  {Il est impossible de copier un bit classique}
  {A}
\qcmx
  {Une porte $U$ qui copierait $\ket{0}$ et $\ket{1}$ enverrait $\ket{+}\ket{0}$ sur un état intriqué, par linéarité, et non sur $\ket{+}\ket{+}$.}
  {C'est l'effondrement de l'état, un autre principe.}
  {L'intrication se crée très bien (Hadamard puis CNOT).}
  {Le CNOT copie un bit classique : $\ket{0,b}\to\ket{b,b}$.}

% ---- Question 30 (niveau 1, correct : B)
\qcmq[1]{Chap.~\ref{chap:bell}\,$\cdot$\,\S\,\ref{sec:circuit-bell}}
  {Quel enchaînement de portes transforme $\ket{00}$ en l'état de Bell $\ket{\Phi^+}=\frac{1}{\sqrt2}(\ket{00}+\ket{11})$ ?}
  {CNOT, puis Hadamard sur le contrôle}
  {Hadamard sur l'un des qubits, puis CNOT dont le contrôle est ce même qubit}
  {Hadamard sur chacun des deux qubits}
  {SWAP, puis $X$ sur un qubit}
  {B}
\qcmx
  {Le CNOT ne fait rien sur $\ket{00}$ (contrôle à $0$) ; $H$ donne ensuite un état séparable.}
  {$H$ crée $\ket{0}\otimes\ket{+}$, puis le CNOT recopie la superposition : $\frac{1}{\sqrt2}(\ket{00}+\ket{11})$.}
  {On obtient $\ket{+}\otimes\ket{+}$, séparable : des portes locales ne créent pas d'intrication.}
  {$\ket{00}$ devient $\ket{01}$ ou $\ket{10}$, un état de base séparable.}

\clearpage
%% ======================================================================
\subsection{Niveau 2 : intermédiaire (questions 31 à 60)}\label{sec:qcm-n2}
%% ======================================================================
\noindent{\color{insagris}\small Niveau de difficulté : Intermédiaire\quad $\bullet\bullet\circ\circ$}\par\smallskip
% ---- Question 31 (niveau 2, correct : A)
\qcmq[1]{Chap.~\ref{chap:histoire}\,$\cdot$\,\S\,\ref{sec:born}}
  {Quelle condition la fonction d'onde d'une particule doit-elle vérifier sur tout l'espace ?}
  {$\displaystyle\int_{-\infty}^{+\infty}|\Psi(x,t)|^2\,\dd x=1$}
  {$\displaystyle\int_{-\infty}^{+\infty}\Psi(x,t)\,\dd x=1$}
  {$|\Psi(x,t)|^2=1$ pour tout $x$}
  {$\partial_x\Psi(x,t)=0$ en tout point}
  {A}
\qcmx
  {La particule se trouve forcément quelque part : la probabilité totale vaut $1$ (condition de normalisation).}
  {Sans le module au carré : $\Psi$ est complexe et son intégrale n'est pas une probabilité (elle peut être nulle ou complexe).}
  {$|\Psi|^2$ est une densité : si elle valait $1$ partout, son intégrale divergerait.}
  {Une fonction constante n'est pas normalisable ($\int|\Psi|^2\dd x$ diverge) ; la condition porte sur l'intégrale de $|\Psi|^2$.}

% ---- Question 32 (niveau 2, correct : D)
\qcmq[2]{Chap.~\ref{chap:histoire}\,$\cdot$\,\S\,\ref{sec:operateurs}}
  {Quelle est l'expression de l'opérateur quantité de mouvement $\hat p$ en une dimension ?}
  {$\hat p=m\,v$}
  {$\hat p=-\dfrac{\hbar^2}{2m}\dfrac{\partial^2}{\partial x^2}$}
  {$\hat p=+\ii\hbar\,\dfrac{\partial}{\partial x}$}
  {$\hat p=-\ii\hbar\,\dfrac{\partial}{\partial x}$}
  {D}
\qcmx
  {C'est la grandeur classique ; en mécanique quantique $p$ est un opérateur différentiel.}
  {C'est l'énergie cinétique $\hat p^2/2m$.}
  {Erreur de signe : on trouverait $p=-\hbar k$ pour l'onde $\ee^{\ii kx}$.}
  {Convention du cours : $\hat p\,\ee^{\ii kx}=\hbar k\,\ee^{\ii kx}$, donc une onde vers les $x$ croissants a $p=\hbar k>0$.}

% ---- Question 33 (niveau 2, correct : B)
\qcmq[2]{Chap.~\ref{chap:histoire}\,$\cdot$\,\S\,\ref{sec:solutions}}
  {Pour une particule dans un puits de potentiel infini de largeur $a$, comment varie l'énergie $E_n$ avec le numéro $n$ du niveau ?}
  {$E_n=E_1/n$}
  {$E_n=n^2E_1$, avec $E_1=\dfrac{\pi^2\hbar^2}{2ma^2}$}
  {$E_n=nE_1$}
  {$E_n=E_1$ pour tout $n$}
  {B}
\qcmx
  {Les niveaux se resserreraient quand $n$ augmente ; c'est l'inverse pour le puits infini.}
  {$k_n=\frac{n\pi}{a}$ et $E_n=\frac{\hbar^2k_n^2}{2m}$ : les niveaux s'écartent quand $n$ croît ($E_{n+1}-E_n=(2n+1)E_1$).}
  {Niveaux équidistants : ce n'est pas le puits infini, où $E_2=4E_1$ et $E_3=9E_1$.}
  {Il n'y aurait pas de quantification : tous les états auraient la même énergie.}

% ---- Question 34 (niveau 2, correct : C)
\qcmq[2]{Chap.~\ref{chap:histoire}\,$\cdot$\,\S\,\ref{sec:solutions}}
  {Une particule est dans un état stationnaire $\Psi_n$ de l'hamiltonien ($\Hop\Psi_n=E_n\Psi_n$). Que vaut l'écart-type $\sigma_E$ de son énergie ?}
  {$E_n$}
  {$E_n^2$}
  {$0$}
  {$\sqrt{E_n}$}
  {C}
\qcmx
  {C'est la valeur moyenne de l'énergie, pas son écart-type.}
  {C'est la valeur de $\langle\Hop^2\rangle$, avant soustraction de $\langle\Hop\rangle^2$.}
  {$\langle\Hop\rangle=E_n$ et $\langle\Hop^2\rangle=E_n^2$, donc $\sigma_E=\sqrt{E_n^2-E_n^2}=0$ : l'énergie est parfaitement déterminée.}
  {Mélange avec la formule $\sigma=\sqrt{\langle\Hop^2\rangle-\langle\Hop\rangle^2}$ : les deux termes sont égaux à $E_n^2$.}

% ---- Question 35 (niveau 2, correct : C)
\qcmq[2]{Chap.~\ref{chap:histoire}\,$\cdot$\,\S\,\ref{sec:solutions}}
  {Pour un état stationnaire $\Psi_n$ du puits infini $]0,a[$, que vaut la position moyenne $\langle\hat x\rangle$ ?}
  {$a/n$}
  {$a/4$}
  {$a/2$ pour tout $n$}
  {$0$}
  {C}
\qcmx
  {Elle dépendrait de $n$, ce que la symétrie de $|\Psi_n|^2$ interdit.}
  {C'est la position d'un maximum de $|\Psi_2|^2$ ; la moyenne des deux maxima $\frac a4$ et $\frac{3a}{4}$ donne $\frac a2$.}
  {$|\Psi_n|^2=\frac2a\sin^2\frac{n\pi x}{a}$ est symétrique par rapport au centre du puits.}
  {$\Psi_n$ s'annule sur la paroi, mais $\langle\hat x\rangle$ est une moyenne pondérée sur tout le puits.}

% ---- Question 36 (niveau 2, correct : D)
\qcmq[2]{Chap.~\ref{chap:dirac}\,$\cdot$\,\S\,\ref{sec:qubit}}
  {Un qubit est dans l'état $\ket{\Psi}=\frac{1}{\sqrt5}\ket{0}+\frac{2\ii}{\sqrt5}\ket{1}$. Quelle est la probabilité de mesurer $1$ en base $Z$ ?}
  {$\frac25$}
  {$-\frac45$}
  {$\frac15$}
  {$\frac45$}
  {D}
\qcmx
  {On divise le coefficient $2$ par $5$ : on oublie à la fois la racine et le carré.}
  {On élève $\frac{2\ii}{\sqrt5}$ au carré sans conjuguer : $\left(\frac{2\ii}{\sqrt5}\right)^2=-\frac45$ ; une probabilité est $|\cdot|^2$.}
  {C'est $P(0)=\left|\frac{1}{\sqrt5}\right|^2$.}
  {$P(1)=\left|\frac{2\ii}{\sqrt5}\right|^2=\frac45$ : le facteur $\ii$ est une phase relative, invisible en base $Z$.}

% ---- Question 37 (niveau 2, correct : C)
\qcmq[2]{Chap.~\ref{chap:dirac}\,$\cdot$\,\S\,\ref{sec:projecteurs}}
  {Un qubit préparé dans $\ket{0}$ est mesuré dans la base $X$ $\{\ket{+},\ket{-}\}$. Quelle est la probabilité d'obtenir $-$ ?}
  {$0$}
  {$1$}
  {$\frac12$}
  {$\frac{1}{\sqrt2}$}
  {C}
\qcmx
  {$\ket{0}$ est certain en base $Z$, mais pas en base $X$ : une autre base donne une autre loi.}
  {Même confusion dans l'autre sens : $\ket{0}$ n'est pas l'état $\ket{-}$.}
  {$P(-)=|\braket{-}{0}|^2=\left|\frac{1}{\sqrt2}\right|^2=\frac12$.}
  {C'est l'amplitude $\braket{-}{0}$ ; il faut élever son module au carré.}

% ---- Question 38 (niveau 2, correct : C)
\qcmq[2]{Chap.~\ref{chap:dirac}\,$\cdot$\,\S\,\ref{sec:qubit}}
  {Quel facteur $N>0$ normalise l'état $N\big(\ket{0}+(1-\ii)\ket{1}\big)$ ?}
  {$N=\frac{1}{\sqrt2}$}
  {$N=\frac13$}
  {$N=\frac{1}{\sqrt3}$}
  {$N=\frac{1}{\sqrt{1-2\ii}}$}
  {C}
\qcmx
  {On oublie la composante $\ket{0}$ : $|1-\ii|^2=2$, mais la norme au carré est $1+2=3$.}
  {On oublie la racine carrée : $3N^2=1$ donne $N=\frac{1}{\sqrt3}$, pas $\frac13$.}
  {$\braket{\Psi}{\Psi}=N^2\big(1+|1-\ii|^2\big)=N^2(1+2)=3N^2$.}
  {On élève $1-\ii$ au carré sans prendre le module : $(1-\ii)^2=-2\ii$ n'est pas réel ; il faut $|1-\ii|^2=2$.}

% ---- Question 39 (niveau 2, correct : D)
\qcmq[2]{Chap.~\ref{chap:dirac}\,$\cdot$\,\S\,\ref{sec:hilbert}}
  {Soient $\ket{\psi}=\frac{1}{\sqrt2}(\ket{0}+\ii\ket{1})$ et $\ket{\varphi}=\frac{1}{\sqrt2}(\ket{0}-\ket{1})$. Que vaut $\braket{\psi}{\varphi}$ ?}
  {$\frac{1-\ii}{2}$}
  {$\frac12$}
  {$0$}
  {$\frac{1+\ii}{2}$}
  {D}
\qcmx
  {On oublie de conjuguer l'amplitude $\ii$ du bra : on calcule $\frac12\big(1+\ii\cdot(-1)\big)$.}
  {Seul le terme en $\ket{0}$ est compté ; le terme en $\ket{1}$ vaut $\frac12(-\ii)(-1)=\frac{\ii}{2}$, non nul.}
  {Confusion avec $\braket{+}{-}=0$ ; ici $|\braket{\psi}{\varphi}|^2=\frac12$.}
  {Le bra conjugue les amplitudes : $\bra{\psi}=\frac{1}{\sqrt2}(\bra{0}-\ii\bra{1})$, d'où $\frac12\big(1+(-\ii)(-1)\big)=\frac{1+\ii}{2}$.}

% ---- Question 40 (niveau 2, correct : A)
\qcmq[2]{Chap.~\ref{chap:dirac}\,$\cdot$\,\S\,\ref{sec:qubit}}
  {Lequel de ces états est orthogonal à $\ket{{+}\ii}=\frac{1}{\sqrt2}(\ket{0}+\ii\ket{1})$ ?}
  {$\ket{{-}\ii}=\frac{1}{\sqrt2}(\ket{0}-\ii\ket{1})$}
  {$\ket{-}=\frac{1}{\sqrt2}(\ket{0}-\ket{1})$}
  {$\ket{+}=\frac{1}{\sqrt2}(\ket{0}+\ket{1})$}
  {$\ket{1}$}
  {A}
\qcmx
  {$\braket{{+}\ii}{{-}\ii}=\frac12\big(1+(-\ii)(-\ii)\big)=\frac12(1-1)=0$, en conjuguant le bra.}
  {$\braket{{+}\ii}{-}=\frac{1+\ii}{2}\neq0$ : $\ket{-}$ est orthogonal à $\ket{+}$, pas à $\ket{{+}\ii}$.}
  {$\braket{{+}\ii}{+}=\frac{1-\ii}{2}\neq0$.}
  {$\braket{{+}\ii}{1}=\frac{-\ii}{\sqrt2}\neq0$ : $\ket{1}$ est orthogonal à $\ket{0}$ seulement.}

% ---- Question 41 (niveau 2, correct : D)
\qcmq[1]{Chap.~\ref{chap:operateurs}\,$\cdot$\,\S\,\ref{sec:pur-mixte}}
  {Une matrice de densité de qubit a pour pureté $\operatorname{Tr}(\rho^2)=\frac58$. Que peut-on en conclure ?}
  {L'état est pur}
  {La matrice n'est pas une matrice de densité valide}
  {L'état est maximalement mixte, $\rho=\mathrm{Id}/2$}
  {L'état est mixte (mélange statistique)}
  {D}
\qcmx
  {Il faudrait $\operatorname{Tr}\rho^2=1$.}
  {$\frac58$ est compris entre $\frac12$ et $1$ : valeur admissible (exemple du cours : $\mathrm{diag}(\frac34,\frac14)$).}
  {Ce cas donnerait $\operatorname{Tr}\rho^2=\frac12$.}
  {$\operatorname{Tr}\rho^2<1$ ; pour un qubit, $\frac12\leq\operatorname{Tr}\rho^2\leq1$ et $\frac58$ est une valeur admissible.}

% ---- Question 42 (niveau 2, correct : B)
\qcmq[2]{Chap.~\ref{chap:operateurs}\,$\cdot$\,\S\,\ref{sec:rho}}
  {Que vaut $\langle Z\rangle=\operatorname{Tr}(\rho Z)$ pour un qubit dans l'état $\ket{+}$ ?}
  {$1$}
  {$0$}
  {$-1$}
  {$\frac{1}{\sqrt2}$}
  {B}
\qcmx
  {C'est la valeur pour $\ket{0}$.}
  {$\rho=\frac12\begin{pmatrix}1&1\\1&1\end{pmatrix}$ et $\operatorname{Tr}(\rho Z)=\frac12-\frac12=0$ : les résultats $+1$ et $-1$ sont équiprobables.}
  {C'est la valeur pour $\ket{1}$.}
  {C'est l'amplitude de $\ket{0}$ dans $\ket{+}$, pas une valeur moyenne de $Z$.}

% ---- Question 43 (niveau 2, correct : A)
\qcmq[1]{Chap.~\ref{chap:operateurs}\,$\cdot$\,\S\,\ref{sec:bloch}}
  {Sur la boule de Bloch, que représente un point situé à l'intérieur de la sphère ($|\vec r|<1$) ?}
  {Un état mixte, de pureté $\operatorname{Tr}\rho^2=\frac{1+|\vec r|^2}{2}<1$}
  {Un état pur dont la norme n'est pas $1$}
  {Un état pur en superposition}
  {Un point sans signification physique}
  {A}
\qcmx
  {Les états purs sont sur la sphère ; le centre $\vec r=\vec0$ est l'état maximalement mixte $\mathrm{Id}/2$.}
  {Tout état pur est normalisé et situé sur la sphère.}
  {Les états purs, superpositions comprises, sont tous \emph{sur} la sphère.}
  {Tout point de la boule est une matrice de densité valide.}

% ---- Question 44 (niveau 2, correct : B)
\qcmq[1]{Chap.~\ref{chap:operateurs}\,$\cdot$\,\S\,\ref{sec:pur-mixte}}
  {Par quoi la matrice de densité distingue-t-elle l'état pur $\ket{+}$ du mélange équiprobable de $\ket{0}$ et $\ket{1}$ ?}
  {Par leur trace, égale à $1$ pour l'état pur seulement}
  {Par les éléments hors diagonale (cohérences), non nuls pour $\ket{+}$ et nuls pour le mélange}
  {Par leurs éléments diagonaux, qui sont différents}
  {Ils ne se distinguent pas : c'est la même matrice}
  {B}
\qcmx
  {La trace vaut $1$ pour toute matrice de densité.}
  {$\ket{+}\bra{+}=\frac12\begin{pmatrix}1&1\\1&1\end{pmatrix}$ et le mélange vaut $\frac12\mathrm{Id}$ ; une mesure en base $X$ les distingue.}
  {Les deux ont la diagonale $(\frac12,\frac12)$ : $P(0)=P(1)=\frac12$ dans les deux cas.}
  {Les matrices sont différentes : en base $X$, $\ket{+}$ donne $+$ avec certitude, le mélange donne $\pm$ à pile ou face.}

% ---- Question 45 (niveau 2, correct : C)
\qcmq[2]{Chap.~\ref{chap:operateurs}\,$\cdot$\,\S\,\ref{sec:rho}}
  {Un qubit est préparé dans $\ket{0}$ avec la probabilité $\frac34$ et dans $\ket{1}$ avec la probabilité $\frac14$. Que vaut $\langle Z\rangle=\operatorname{Tr}(\rho Z)$ ?}
  {$\frac34$}
  {$1$}
  {$\frac12$}
  {$0$}
  {C}
\qcmx
  {C'est la probabilité $\rho_{00}$ de lire $0$ ; or $Z$ vaut $+1$ pour $0$ et $-1$ pour $1$.}
  {Valeur pour $\rho=\ket{0}\bra{0}$ (préparation certaine de $\ket{0}$).}
  {$\rho=\mathrm{diag}(\frac34,\frac14)$, donc $\operatorname{Tr}(\rho Z)=\frac34\cdot1+\frac14\cdot(-1)=\frac12$ : c'est la coordonnée $r_z$ de la sphère de Bloch.}
  {Valeur pour le mélange équiprobable $\mathrm{Id}/2$.}

% ---- Question 46 (niveau 2, correct : B)
\qcmq[2]{Chap.~\ref{chap:operateurs}\,$\cdot$\,\S\,\ref{sec:tenseur}}
  {Dans la base ordonnée $\{\ket{00},\ket{01},\ket{10},\ket{11}\}$, quel vecteur colonne représente $\ket{1}\otimes\ket{0}=\ket{10}$ ?}
  {$(0,1,0,0)^{\mathsf T}$}
  {$(0,0,1,0)^{\mathsf T}$}
  {$(1,0,0,0)^{\mathsf T}$}
  {$(0,0,0,1)^{\mathsf T}$}
  {B}
\qcmx
  {C'est $\ket{01}$ : l'ordre des facteurs est inversé.}
  {$\binom{0}{1}\otimes\binom{1}{0}=\big(0\binom{1}{0},\,1\binom{1}{0}\big)=(0,0,1,0)^{\mathsf T}$ : c'est le troisième vecteur de la base.}
  {C'est $\ket{00}$.}
  {C'est $\ket{11}$.}

% ---- Question 47 (niveau 2, correct : D)
\qcmq[2]{Chap.~\ref{chap:portes}\,$\cdot$\,\S\,\ref{sec:phase}}
  {Quelle relation lie la porte de phase $S=\mathrm{diag}(1,\ii)$ à la porte $Z$ ?}
  {$S=Z$}
  {$S^2=X$}
  {$S^4=Z$}
  {$S^2=Z$}
  {D}
\qcmx
  {$Z=\mathrm{diag}(1,-1)$ : $S$ est un quart de tour, $Z$ un demi-tour.}
  {$S^2$ est diagonale alors que $X$ ne l'est pas ; c'est $HSH$ qui est une racine carrée de $X$.}
  {$S^4=\mathrm{diag}(1,\ii^4)=\mathrm{Id}$.}
  {$S^2=\mathrm{diag}(1,\ii^2)=\mathrm{diag}(1,-1)=Z$ : deux quarts de tour autour de $z$ font un demi-tour.}

% ---- Question 48 (niveau 2, correct : A)
\qcmq[2]{Chap.~\ref{chap:portes}\,$\cdot$\,\S\,\ref{sec:controlees}}
  {Dans la porte de Toffoli (CCNOT) à trois qubits $x,y,z$, à quelle condition la cible $z$ est-elle inversée ?}
  {Seulement si $x=1$ et $y=1$}
  {Si $x=1$ ou $y=1$}
  {Si $x\neq y$}
  {Si $x=0$ et $y=0$}
  {A}
\qcmx
  {$\ket{x,y,z}\to\ket{x,y,z\oplus xy}$ : la cible est inversée si et seulement si le produit $xy$ vaut $1$.}
  {Un seul contrôle à $1$ ne suffit pas : pour $(x,y)=(1,0)$, $xy=0$ et la cible ne change pas.}
  {Ce serait deux CNOT ($z\to z\oplus x\oplus y$) ; pour $(1,0)$ le Toffoli n'inverse rien.}
  {Pour $(0,0)$ on a $xy=0$ : la cible ne change pas.}

% ---- Question 49 (niveau 2, correct : A)
\qcmq[1]{Chap.~\ref{chap:portes}\,$\cdot$\,\S\,\ref{sec:portes}}
  {Pour un hamiltonien $\Hop$ indépendant du temps, quel opérateur d'évolution $U(t,t_0)$ vérifie $\ket{\Psi(t)}=U(t,t_0)\ket{\Psi(t_0)}$ ?}
  {$U(t,t_0)=\ee^{-\ii\Hop(t-t_0)/\hbar}$}
  {$U(t,t_0)=\ee^{+\ii\Hop(t-t_0)/\hbar}$}
  {$U(t,t_0)=-\ii\hbar\,\partial_t$}
  {$U(t,t_0)=\Hop\,(t-t_0)/\hbar$}
  {A}
\qcmx
  {Solution de $\ii\hbar\,\partial_t\ket{\Psi}=\Hop\ket{\Psi}$ ; comme $\Hop$ est hermitien, $U^\dagger U=\mathrm{Id}$ : l'évolution est unitaire.}
  {Erreur de signe : c'est $U^\dagger$, qui remonte le temps.}
  {C'est un opérateur différentiel, pas un opérateur d'évolution ; l'équation de Schrödinger s'écrit $\ii\hbar\,\partial_t\ket{\Psi}=\Hop\ket{\Psi}$.}
  {On oublie l'exponentielle et le facteur $\ii$ : cet opérateur est hermitien, pas unitaire.}

% ---- Question 50 (niveau 2, correct : D)
\qcmq[2]{Chap.~\ref{chap:portes}\,$\cdot$\,\S\,\ref{sec:controlees}}
  {La porte de Fredkin (SWAP contrôlé) a pour contrôle le premier qubit. Que donne-t-elle sur $\ket{1,0,1}$ ?}
  {$\ket{1,0,1}$}
  {$\ket{0,1,1}$}
  {$\ket{1,1,1}$}
  {$\ket{1,1,0}$}
  {D}
\qcmx
  {C'est le résultat si le contrôle valait $0$ ; ici il vaut $1$ et les deux autres qubits sont différents.}
  {Le SWAP porte sur les qubits $2$ et $3$, pas sur le contrôle.}
  {Un SWAP ne change pas le nombre de $1$, il permute seulement ; $\ket{1,1,1}$ viendrait d'un NOT.}
  {Le contrôle vaut $1$ : les deux autres qubits sont échangés, $\ket{0,1}\to\ket{1,0}$.}

% ---- Question 51 (niveau 2, correct : A)
\qcmq[1]{Chap.~\ref{chap:portes}\,$\cdot$\,\S\,\ref{sec:op2}}
  {Que donne $(H\otimes I)\ket{00}$, où $H$ agit sur le premier facteur du produit tensoriel $\ket{a}\otimes\ket{b}$ ?}
  {$\ket{+}\otimes\ket{0}=\frac{1}{\sqrt2}(\ket{00}+\ket{10})$}
  {$\frac12(\ket{00}+\ket{01}+\ket{10}+\ket{11})$}
  {$\ket{0}\otimes\ket{+}=\frac{1}{\sqrt2}(\ket{00}+\ket{01})$}
  {$\frac{1}{\sqrt2}(\ket{00}+\ket{11})$}
  {A}
\qcmx
  {$(A\otimes B)(\ket{u}\otimes\ket{v})=A\ket{u}\otimes B\ket{v}$ : $H\ket{0}=\ket{+}$ sur le premier facteur, $I\ket{0}=\ket{0}$ sur le second.}
  {C'est $(H\otimes H)\ket{00}$ : ici le second qubit n'est pas touché.}
  {$H$ agit sur le premier qubit, pas sur le second : les facteurs sont inversés.}
  {C'est $\ket{\Phi^+}$ : il faut un CNOT pour intriquer ; des portes locales laissent l'état séparable.}

% ---- Question 52 (niveau 2, correct : C)
\qcmq[1]{Chap.~\ref{chap:bell}\,$\cdot$\,\S\,\ref{sec:separable}}
  {Pourquoi l'état de Bell $\ket{\Phi^+}=\frac{1}{\sqrt2}(\ket{00}+\ket{11})$ est-il dit intriqué ?}
  {Parce que $P(01)=0$}
  {Parce que chaque qubit, pris seul, est dans un état pur}
  {Il ne peut pas s'écrire comme un produit $\ket{\varphi}_A\otimes\ket{\chi}_B$ : $D=c_{00}c_{11}-c_{01}c_{10}=\frac12\neq0$}
  {Parce qu'il est orthogonal à $\ket{00}$}
  {C}
\qcmx
  {C'est vrai, mais $\ket{00}$ a aussi $P(01)=0$ et il est séparable : cette condition ne prouve rien.}
  {C'est l'inverse : pris seul, chaque qubit est complètement mixte ($\rho_A=\mathrm{Id}/2$).}
  {Un produit imposerait $D=0$ ; aucun des deux qubits n'a d'état propre.}
  {$\braket{00}{\Phi^+}=\frac{1}{\sqrt2}\neq0$ ; l'orthogonalité n'a pas de lien avec l'intrication.}

% ---- Question 53 (niveau 2, correct : B)
\qcmq[2]{Chap.~\ref{chap:bell}\,$\cdot$\,\S\,\ref{sec:teleportation}}
  {Dans la téléportation du cours, Alice obtient les résultats $a=0$ et $q=1$. Quelle correction Bob doit-il appliquer à son qubit ?}
  {La porte $X$}
  {La porte $Z$}
  {$X$ puis $Z$}
  {Aucune, l'état est déjà $\ket{\varphi}$}
  {B}
\qcmx
  {$X$ corrige le cas $a=1$ (inversion de bit) ; ici $a=0$.}
  {Bob détient $\alpha\ket{0}-\beta\ket{1}=Z\ket{\varphi}$ ; $Z$ rétablit le signe : $Z(\alpha\ket{0}-\beta\ket{1})=\alpha\ket{0}+\beta\ket{1}$.}
  {Les deux corrections ne sont nécessaires que pour $(a,q)=(1,1)$.}
  {Ce n'est vrai que pour $(a,q)=(0,0)$.}

% ---- Question 54 (niveau 2, correct : B)
\qcmq[2]{Chap.~\ref{chap:bell}\,$\cdot$\,\S\,\ref{sec:bell} ; ex.\,\ref{ex:superdense}}
  {Dans le codage superdense, Alice et Bob partagent $\ket{\Phi^+}$ et Alice agit sur son qubit. Quelle porte transforme $\ket{\Phi^+}$ en $\ket{\Psi^+}=\frac{1}{\sqrt2}(\ket{01}+\ket{10})$ ?}
  {La porte $Z$}
  {La porte $X$}
  {La porte identité}
  {La porte $H$}
  {B}
\qcmx
  {$Z$ donne $\ket{\Phi^-}=\frac{1}{\sqrt2}(\ket{00}-\ket{11})$ : elle change un signe, pas les chiffres.}
  {$X$ échange $\ket{0}$ et $\ket{1}$ sur un qubit : $\ket{00}+\ket{11}\to\ket{01}+\ket{10}$ ; c'est le message $10$.}
  {Elle laisse $\ket{\Phi^+}$ : c'est le message $00$.}
  {$(H\otimes I)\ket{\Phi^+}$ est une superposition de deux états de Bell, pas un état de Bell.}

% ---- Question 55 (niveau 2, correct : D)
\qcmq[2]{Chap.~\ref{chap:bell}\,$\cdot$\,\S\,\ref{sec:info}}
  {Pour deux qubits dans l'état de Bell $\ket{\Phi^+}$, que vaut l'entropie conditionnelle $H(A|B)=H(A,B)-H(B)$ ?}
  {$0$ bit}
  {$+1$ bit}
  {$+2$ bits}
  {$-1$ bit}
  {D}
\qcmx
  {C'est la valeur classique quand $A$ est entièrement déterminé par $B$ (deux bits identiques) ; ici $H(A,B)=0$, pas $H(B)$.}
  {C'est la valeur quand $A$ et $B$ sont indépendants : $H(A|B)=H(A)=1$.}
  {C'est l'information mutuelle $I(A;B)=1+1-0=2$ bits, pas l'entropie conditionnelle.}
  {$H(A,B)=S(\rho_{AB})=0$ (état pur) et $H(B)=1$ bit ($\rho_B=\mathrm{Id}/2$) : $0-1=-1$, valeur impossible en classique.}

% ---- Question 56 (niveau 2, correct : B)
\qcmq[1]{Chap.~\ref{chap:bell}\,$\cdot$\,\S\,\ref{sec:trace}}
  {Quelle est la matrice de densité réduite $\rho_A=\operatorname{Tr}_B\rho_{AB}$ d'un qubit d'une paire dans l'état $\ket{\Phi^+}$ ?}
  {$\rho_A=\ket{0}\bra{0}$}
  {$\rho_A=\frac12\mathrm{Id}$, de pureté $\operatorname{Tr}\rho_A^2=\frac12$}
  {$\rho_A=\ket{+}\bra{+}=\frac12\begin{pmatrix}1&1\\1&1\end{pmatrix}$}
  {$\rho_A=\ket{\Phi^+}\bra{\Phi^+}$}
  {B}
\qcmx
  {Le qubit aurait un état défini, ce qui est faux pour un état intriqué : $P(0)=P(1)=\frac12$.}
  {$\rho_A=\frac12\big(\ket{0}\bra{0}+\ket{1}\bra{1}\big)$ : les termes croisés $\ket{00}\bra{11}$ s'annulent ($\braket{1}{0}=0$). Le couple est pur, mais chaque qubit est complètement mixte.}
  {Les termes croisés sont multipliés par $\braket{1}{0}=0$ dans la trace partielle : ils disparaissent.}
  {C'est la matrice $4\times4$ du couple ; la trace partielle donne une matrice $2\times2$.}

% ---- Question 57 (niveau 2, correct : C)
\qcmq[1]{Chap.~\ref{chap:bell}\,$\cdot$\,\S\,\ref{sec:separable}}
  {Pour $\ket{\Psi}=c_{00}\ket{00}+c_{01}\ket{01}+c_{10}\ket{10}+c_{11}\ket{11}$, quel est le critère de séparabilité ?}
  {$c_{00}c_{11}+c_{01}c_{10}=0$}
  {$c_{00}c_{01}-c_{10}c_{11}=0$}
  {$D=c_{00}c_{11}-c_{01}c_{10}=0$}
  {$c_{00}=c_{11}$ et $c_{01}=c_{10}$}
  {C}
\qcmx
  {Mauvais signe : $\ket{+}\otimes\ket{-}$ est séparable mais donnerait $-\frac12\neq0$.}
  {Mauvais appariement : $\ket{\Psi^+}$ (intriqué) la vérifie, puisque $c_{00}=c_{11}=0$.}
  {Un produit $(a_0\ket{0}+a_1\ket{1})\otimes(b_0\ket{0}+b_1\ket{1})$ a $c_{ij}=a_ib_j$, donc $c_{00}c_{11}=a_0a_1b_0b_1=c_{01}c_{10}$ ; la réciproque est vraie aussi.}
  {Ni nécessaire ($\ket{00}$ est séparable avec $c_{00}\neq c_{11}$) ni suffisant ($\ket{\Phi^+}$ la vérifie et est intriqué).}

% ---- Question 58 (niveau 2, correct : B)
\qcmq[2]{Chap.~\ref{chap:bell}\,$\cdot$\,\S\,\ref{sec:bell}}
  {Quelle est la valeur propre de $X\otimes X$ pour l'état $\ket{\Phi^+}$ ?}
  {$-1$}
  {$+1$}
  {$0$}
  {$\ket{\Phi^+}$ n'est pas un état propre de $X\otimes X$}
  {B}
\qcmx
  {C'est la valeur pour $\ket{\Phi^-}$ et $\ket{\Psi^-}$ : le signe $-$ de l'état se retrouve dans la valeur propre.}
  {$(X\otimes X)\ket{00}=\ket{11}$ et $(X\otimes X)\ket{11}=\ket{00}$, donc $(X\otimes X)\ket{\Phi^+}=\ket{\Phi^+}$.}
  {$X\otimes X$ est unitaire : ses valeurs propres sont de module $1$, jamais nulles.}
  {Les quatre états de Bell sont des états propres communs de $Z\otimes Z$ et de $X\otimes X$.}

% ---- Question 59 (niveau 2, correct : A)
\qcmq[2]{Chap.~\ref{chap:bell}\,$\cdot$\,\S\,\ref{sec:info}}
  {Que vaut l'entropie de von Neumann $S(\rho)=-\sum_i\lambda_i\log_2\lambda_i$ d'un état pur ?}
  {$0$ bit}
  {$1$ bit}
  {$-1$ bit}
  {$+\infty$}
  {A}
\qcmx
  {Les valeurs propres sont $1$ et $0$ : $1\cdot\log_21=0$ et $0\cdot\log_20=0$ (par continuité, $x\log_2x\to0$).}
  {C'est la valeur pour l'état maximalement mixte $\mathrm{Id}/2$ d'un qubit.}
  {L'entropie de $\rho$ est positive ou nulle ; seule l'entropie conditionnelle quantique peut être négative.}
  {$\log_20$ diverge, mais le terme $0\cdot\log_20$ vaut $0$ par continuité.}

% ---- Question 60 (niveau 2, correct : A)
\qcmq[2]{Chap.~\ref{chap:bell}\,$\cdot$\,\S\,\ref{sec:bell}}
  {Lequel des états de Bell a pour valeur propre $-1$ à la fois pour $Z\otimes Z$ et pour $X\otimes X$ ?}
  {$\ket{\Psi^-}=\frac{1}{\sqrt2}(\ket{01}-\ket{10})$}
  {$\ket{\Phi^+}$}
  {$\ket{\Phi^-}$}
  {$\ket{\Psi^+}$}
  {A}
\qcmx
  {Résultats opposés dans les deux bases : $Z\otimes Z=-1$ (chiffres différents) et $X\otimes X=-1$ (signe $-$ dans l'état).}
  {Valeurs propres $(Z\otimes Z,\,X\otimes X)=(+1,+1)$.}
  {Valeurs propres $(Z\otimes Z,\,X\otimes X)=(+1,-1)$.}
  {Valeurs propres $(Z\otimes Z,\,X\otimes X)=(-1,+1)$.}

\clearpage
%% ======================================================================
\subsection{Niveau 3 : avancé (questions 61 à 90)}\label{sec:qcm-n3}
%% ======================================================================
\noindent{\color{insagris}\small Niveau de difficulté : Avancé\quad $\bullet\bullet\bullet\circ$}\par\smallskip
% ---- Question 61 (niveau 3, correct : A)
\qcmq[2]{Chap.~\ref{chap:histoire}\,$\cdot$\,\S\,\ref{sec:solutions}}
  {Un électron d'un puits infini passe du niveau $n=3$ au niveau $n=2$ en émettant un photon. Quelle est l'énergie du photon, en fonction de $E_1$ ?}
  {$5E_1$}
  {$E_1$}
  {$\frac94E_1$}
  {$13E_1$}
  {A}
\qcmx
  {$E_3-E_2=9E_1-4E_1=5E_1$ ; le photon a la fréquence $f=(E_3-E_2)/h$.}
  {On a soustrait les numéros de niveau ($3-2=1$) au lieu de leurs carrés.}
  {C'est le rapport $E_3/E_2$, pas la différence $E_3-E_2$.}
  {On a additionné $9E_1+4E_1$ : l'énergie du photon est une différence de niveaux.}

% ---- Question 62 (niveau 3, correct : A)
\qcmq[2]{Chap.~\ref{chap:histoire}\,$\cdot$\,\S\,\ref{sec:superposition}}
  {À $t=0$, une particule du puits infini est dans l'état $\Psi=\frac{\sqrt3}{2}\,\psi_1+\frac12\,\psi_3$ ($\psi_n$ : états stationnaires normalisés, d'énergies $E_n=n^2E_1$). Quelle est son énergie moyenne $\langle E\rangle$ ?}
  {$3E_1$}
  {$5E_1$}
  {$\left(\frac{\sqrt3}{2}+\frac92\right)E_1$}
  {$7E_1$}
  {A}
\qcmx
  {$P(E_1)=\frac34$ et $P(E_3)=\frac14$, donc $\langle E\rangle=\frac34E_1+\frac14\cdot9E_1=3E_1$.}
  {On a pondéré également les deux niveaux, $\frac{E_1+9E_1}{2}$ : les poids sont les $|c_i|^2$, pas $\frac12$.}
  {On a pondéré par les amplitudes $c_i$ au lieu de leurs carrés $|c_i|^2$.}
  {Poids permutés : $\frac14E_1+\frac34\cdot9E_1=7E_1$ ; la probabilité de $E_1$ est $\left|\frac{\sqrt3}{2}\right|^2=\frac34$.}

% ---- Question 63 (niveau 3, correct : C)
\qcmq[1]{Chap.~\ref{chap:histoire}\,$\cdot$\,\S\,\ref{sec:solutions}}
  {Pour l'état stationnaire $n=2$ du puits infini $]0,a[$, que valent la densité de probabilité de présence en $x=a/2$ et la position moyenne $\langle\hat x\rangle$ ?}
  {Densité maximale en $a/2$, et $\langle\hat x\rangle=a/2$}
  {Densité nulle en $a/2$, et $\langle\hat x\rangle=0$}
  {Densité nulle en $a/2$, et $\langle\hat x\rangle=a/2$}
  {Densité maximale en $a/2$, et $\langle\hat x\rangle=a/4$}
  {C}
\qcmx
  {C'est le cas de $n=1$ ; pour $n=2$ il y a un nœud au centre.}
  {Une densité nulle au centre ne donne pas une moyenne nulle : $\langle\hat x\rangle$ pondère tout le puits et vaut $\frac a2$ par symétrie.}
  {$|\Psi_2|^2=\frac2a\sin^2\frac{2\pi x}{a}$ s'annule en $x=\frac a2$ (nœud), alors que la symétrie impose $\langle\hat x\rangle=\frac a2$ : la moyenne n'est pas la position la plus probable.}
  {Ni l'un ni l'autre : maxima en $\frac a4$ et $\frac{3a}4$, nœud en $\frac a2$, moyenne $\frac a2$.}

% ---- Question 64 (niveau 3, correct : C)
\qcmq[2]{Chap.~\ref{chap:dirac}\,$\cdot$\,\S\,\ref{sec:projecteurs}}
  {Quelle est la matrice du projecteur $\hat P_+=\ket{+}\bra{+}$ dans la base $\{\ket{0},\ket{1}\}$ ?}
  {$\frac12\begin{pmatrix}1&-1\\-1&1\end{pmatrix}$}
  {$\begin{pmatrix}1&0\\0&0\end{pmatrix}$}
  {$\frac12\begin{pmatrix}1&1\\1&1\end{pmatrix}$}
  {$\frac{1}{\sqrt2}\begin{pmatrix}1&1\\1&-1\end{pmatrix}$}
  {C}
\qcmx
  {C'est $\ket{-}\bra{-}$, le projecteur sur $\ket{-}$.}
  {C'est $\ket{0}\bra{0}$ : projecteur de la base $Z$, pas de la base $X$.}
  {$\ket{+}\bra{+}=\frac12(\ket0+\ket1)(\bra0+\bra1)$ ; on vérifie $\hat P_+^2=\hat P_+$ et $\operatorname{Tr}\hat P_+=1$.}
  {C'est la matrice de Hadamard, unitaire ($H^2=\mathrm{Id}$), et non un projecteur (un projecteur vérifie $P^2=P$).}

% ---- Question 65 (niveau 3, correct : B)
\qcmq[2]{Chap.~\ref{chap:dirac}\,$\cdot$\,\S\,\ref{sec:projecteurs}}
  {Un qubit préparé dans $\ket{0}$ est mesuré en base $X$, puis aussitôt en base $Z$. Quelle est la probabilité d'obtenir $1$ à la seconde mesure ?}
  {$0$}
  {$\frac12$}
  {$1$}
  {$\frac14$}
  {B}
\qcmx
  {On garde $\ket{0}$ en mémoire : la mesure en base $X$ a effondré l'état sur $\ket{\pm}$ et effacé $\ket{0}$.}
  {La première mesure donne $\ket{+}$ ou $\ket{-}$ (probabilité $\frac12$ chacun) ; dans les deux cas $P(1)=\frac12$ en base $Z$ : $\frac12\cdot\frac12+\frac12\cdot\frac12=\frac12$.}
  {Sans rapport : $\ket{0}$ n'est plus l'état après la mesure en base $X$.}
  {On n'a compté qu'une branche de la première mesure ; il faut sommer les deux ($\frac14+\frac14$).}

% ---- Question 66 (niveau 3, correct : D)
\qcmq[2]{Chap.~\ref{chap:dirac}\,$\cdot$\,\S\,\ref{sec:projecteurs}}
  {Un qubit est dans l'état $\ket{\psi}=\frac{3}{5}\ket{0}+\frac{4}{5}\ket{1}$. Quelle est la probabilité d'obtenir $-$ en mesurant dans la base $X$ ?}
  {$\frac{16}{25}$}
  {$\frac{1}{25}$}
  {$\frac{49}{50}$}
  {$\frac{1}{50}$}
  {D}
\qcmx
  {C'est $P(1)$ en base $Z$ ($|\frac45|^2$) : on a mesuré dans la mauvaise base.}
  {On a oublié le facteur $\frac{1}{\sqrt2}$ de $\ket{-}$ : $\left|\frac35-\frac45\right|^2=\frac1{25}$, il manque le $\frac12$.}
  {C'est $P(+)$ ; la probabilité demandée est celle du résultat $-$.}
  {$\braket{-}{\psi}=\frac{1}{\sqrt2}\left(\frac35-\frac45\right)=-\frac{1}{5\sqrt2}$, donc $P(-)=\frac1{50}$ ; et $P(+)=\frac{49}{50}$.}

% ---- Question 67 (niveau 3, correct : B)
\qcmq[1]{Chap.~\ref{chap:dirac}\,$\cdot$\,\S\,\ref{sec:projecteurs}}
  {Lequel de ces kets décrit le même état physique que $\ket{+}=\frac{1}{\sqrt2}(\ket{0}+\ket{1})$ ?}
  {$\frac{1}{\sqrt2}\big(\ket{0}+\ii\ket{1}\big)$}
  {$\ee^{\ii\pi/3}\ket{+}$}
  {$\frac{1}{\sqrt2}\big(\ket{0}+\ee^{\ii\pi}\ket{1}\big)$}
  {$\frac{1}{\sqrt2}\big(\ket{0}+\ee^{\ii\pi/3}\ket{1}\big)$}
  {B}
\qcmx
  {La phase $\ii$ ne porte que sur $\ket{1}$ : c'est une phase \emph{relative}, et l'état est $\ket{{+}\ii}$ (base $Y$).}
  {Un facteur de phase \emph{globale} multiplie tout le ket : toutes les probabilités de mesure (et $\rho=\ket{+}\bra{+}$) sont inchangées.}
  {$\ee^{\ii\pi}=-1$ : c'est $\ket{-}$, orthogonal à $\ket{+}$.}
  {Phase relative $\frac\pi3$ : c'est un autre état, sur l'équateur de la sphère de Bloch, à l'azimut $\frac\pi3$.}

% ---- Question 68 (niveau 3, correct : A)
\qcmq[2]{Chap.~\ref{chap:dirac}\,$\cdot$\,\S\,\ref{sec:qubit}}
  {Laquelle de ces paires est une base orthonormée de $\mathbb{C}^2$ ?}
  {$\{\ket{{+}\ii},\ket{{-}\ii}\}$}
  {$\{\ket{0},\ket{+}\}$}
  {$\{\ket{+},\ket{{-}\ii}\}$}
  {$\{\ket{0},\ket{0}+\ket{1}\}$}
  {A}
\qcmx
  {Ces deux états sont normés et orthogonaux : $\braket{{+}\ii}{{-}\ii}=\frac12(1-1)=0$.}
  {Normés, mais $\braket{0}{+}=\frac{1}{\sqrt2}\neq0$ : pas orthogonaux.}
  {$\braket{+}{{-}\ii}=\frac12(1-\ii)\neq0$ : pas orthogonaux.}
  {Le second ket n'est pas normé (norme au carré $2$) ni orthogonal à $\ket{0}$.}

% ---- Question 69 (niveau 3, correct : C)
\qcmq[2]{Chap.~\ref{chap:operateurs}\,$\cdot$\,\S\,\ref{sec:rho}}
  {Laquelle de ces matrices peut être la matrice de densité d'un qubit ?}
  {$\begin{pmatrix}\frac12&\frac34\\\frac34&\frac12\end{pmatrix}$}
  {$\begin{pmatrix}\frac12&\frac{\ii}{2}\\\frac{\ii}{2}&\frac12\end{pmatrix}$}
  {$\begin{pmatrix}\frac23&\frac{\ii}{3}\\-\frac{\ii}{3}&\frac13\end{pmatrix}$}
  {$\begin{pmatrix}\frac34&0\\0&\frac12\end{pmatrix}$}
  {C}
\qcmx
  {Hermitienne et de trace $1$, mais pas positive : $\det=\frac14-\frac9{16}<0$, valeurs propres $\frac54$ et $-\frac14$ (une « probabilité » négative).}
  {De trace $1$ mais pas hermitienne : il faudrait $\rho_{10}=\overline{\rho_{01}}=-\frac{\ii}{2}$.}
  {Hermitienne ($\rho_{10}=\overline{\rho_{01}}$), de trace $1$ et positive : $\det\rho=\frac29-\frac19=\frac19>0$ avec $\operatorname{Tr}\rho>0$ donne deux valeurs propres positives, $\frac{3\pm\sqrt5}{6}$.}
  {Hermitienne et positive, mais de trace $\frac54\neq1$ : les probabilités diagonales ne somment pas à $1$.}

% ---- Question 70 (niveau 3, correct : B)
\qcmq[2]{Chap.~\ref{chap:operateurs}\,$\cdot$\,\S\,\ref{sec:op-dirac}}
  {Un qubit est dans l'état $\ket{\psi}=\frac{1}{\sqrt5}\big(\ket{0}+2\ii\ket{1}\big)$. Que vaut la valeur moyenne $\langle Y\rangle=\bra{\psi}Y\ket{\psi}$, avec $Y=\begin{pmatrix}0&-\ii\\\ii&0\end{pmatrix}$ ?}
  {$0$}
  {$\frac45$}
  {$-\frac45$}
  {$-\frac35$}
  {B}
\qcmx
  {On a oublié de conjuguer les amplitudes du bra : $\psi^{\mathsf T}Y\psi=\frac15(2+2\ii\cdot\ii)=0$ (c'est d'ailleurs la valeur de $\langle X\rangle$ pour cet état).}
  {$Y\ket\psi=\frac1{\sqrt5}(2\ket0+\ii\ket1)$, puis $\bra\psi Y\ket\psi=\frac15\big(1\cdot2+(-2\ii)(\ii)\big)=\frac45$ ; en général $\langle Y\rangle=2\,\mathrm{Im}(\bar\alpha\beta)$.}
  {Erreur de signe : on a utilisé la transposée $\begin{pmatrix}0&\ii\\-\ii&0\end{pmatrix}=-Y$.}
  {C'est $\langle Z\rangle=\frac15-\frac45$ : on a calculé la valeur moyenne d'une autre observable.}

% ---- Question 71 (niveau 3, correct : D)
\qcmq[1]{Chap.~\ref{chap:operateurs}\,$\cdot$\,\S\,\ref{sec:spectral}}
  {On mesure l'observable $A=\begin{pmatrix}2&1\\1&2\end{pmatrix}$ sur un qubit préparé dans $\ket{0}$. Quelle affirmation est correcte ?}
  {On obtient toujours $2$}
  {On obtient $+1$ ou $-1$, avec la probabilité $\frac12$ chacun}
  {On obtient $2$ ou $1$, avec les probabilités $\frac45$ et $\frac15$}
  {On obtient $3$ ou $1$, avec la probabilité $\frac12$ chacun (moyenne $2$)}
  {D}
\qcmx
  {On a lu le coefficient diagonal : $A\ket0=2\ket0+\ket1$ n'est pas proportionnel à $\ket0$, donc $\ket0$ n'est pas vecteur propre de $A$.}
  {Ce sont les valeurs propres de $X$, pas de $A$ : le décalage $2\,\mathrm{Id}$ translate le spectre ($2\pm1$).}
  {On a lu $A\ket0=2\ket0+\ket1$ comme un état : $A$ n'est pas unitaire, ce n'est pas une évolution ; les résultats possibles d'une mesure sont les valeurs propres de $A$.}
  {$A=2\,\mathrm{Id}+X$ a les vecteurs propres de $X$ : valeur propre $3$ pour $\ket+$, $1$ pour $\ket-$. Comme $\ket0=\frac{1}{\sqrt2}(\ket++\ket-)$, les deux résultats sont équiprobables, et $\langle A\rangle=\bra0A\ket0=2$.}

% ---- Question 72 (niveau 3, correct : D)
\qcmq[2]{Chap.~\ref{chap:operateurs}\,$\cdot$\,\S\,\ref{sec:pur-mixte}}
  {Un qubit est préparé dans $\ket{0}$ avec la probabilité $\frac12$ et dans $\ket{+}$ avec la probabilité $\frac12$. Quelle est la pureté $\operatorname{Tr}(\rho^2)$ de son état ?}
  {$1$}
  {$\frac12$}
  {$\frac58$}
  {$\frac34$}
  {D}
\qcmx
  {Un mélange de deux états purs n'est pur que s'ils sont identiques ; ici $\ket0\neq\ket+$, donc $\operatorname{Tr}\rho^2<1$.}
  {Valeur de l'état maximalement mixte $\frac{\mathrm{Id}}2$, atteinte pour deux états \emph{orthogonaux} ; ici ils ne le sont pas (on a aussi pu oublier les termes croisés).}
  {On n'a additionné que les carrés des éléments diagonaux ($\frac9{16}+\frac1{16}$) : il manque $2\rho_{01}\rho_{10}=\frac2{16}$.}
  {$\rho=\frac12\ket0\bra0+\frac12\ket+\bra+=\begin{pmatrix}\frac34&\frac14\\\frac14&\frac14\end{pmatrix}$, donc $\operatorname{Tr}\rho^2=\frac9{16}+\frac1{16}+2\cdot\frac1{16}=\frac34$ : état mixte, mais moins que $\frac{\mathrm{Id}}2$ car $\ket0$ et $\ket+$ ne sont pas orthogonaux.}

% ---- Question 73 (niveau 3, correct : B)
\qcmq[2]{Chap.~\ref{chap:operateurs}\,$\cdot$\,\S\,\ref{sec:bloch}}
  {Quel est le vecteur de Bloch $(r_x,r_y,r_z)$ de l'état $\ket{\psi}=\cos\frac\pi6\,\ket0+\sin\frac\pi6\,\ket1$ ?}
  {$\left(\frac12,\,0,\,\frac{\sqrt3}{2}\right)$}
  {$\left(\frac{\sqrt3}{2},\,0,\,\frac12\right)$}
  {$\left(\frac{\sqrt3}{4},\,0,\,\frac12\right)$}
  {$\left(\frac{\sqrt3}{2},\,0,\,\frac34\right)$}
  {B}
\qcmx
  {On a pris $\theta=\frac\pi6$ : l'angle de Bloch est le \emph{double} de l'angle qui figure dans les amplitudes.}
  {Ici $\frac\theta2=\frac\pi6$, donc $\theta=\frac\pi3$ et $\varphi=0$ : $\vec r=(\sin\theta\cos\varphi,\sin\theta\sin\varphi,\cos\theta)$. Contrôle : $r_z=|\alpha|^2-|\beta|^2=\frac34-\frac14$ et $r_x=2\alpha\beta=\frac{\sqrt3}2$.}
  {On a écrit $r_x=\alpha\beta$ en oubliant le facteur $2$ ; on le voit à $|\vec r|^2=\frac7{16}<1$, impossible pour un état pur.}
  {On a pris $r_z=|\alpha|^2$ au lieu de $|\alpha|^2-|\beta|^2$ ; or $|\vec r|^2=\frac{21}{16}>1$ est impossible.}

% ---- Question 74 (niveau 3, correct : A)
\qcmq[2]{Chap.~\ref{chap:operateurs}\,$\cdot$\,\S\,\ref{sec:bloch}}
  {Un qubit a pour matrice de densité $\rho=\frac12\begin{pmatrix}1&-\ii\\\ii&1\end{pmatrix}$. Quel est son vecteur de Bloch ?}
  {$(0,\,1,\,0)$ : c'est l'état pur $\ket{{+}\ii}$}
  {$(0,\,-1,\,0)$}
  {$\left(0,\,\frac12,\,0\right)$}
  {$(0,\,0,\,0)$}
  {A}
\qcmx
  {Avec $\rho=\frac12\begin{pmatrix}1+r_z&r_x-\ii r_y\\r_x+\ii r_y&1-r_z\end{pmatrix}$ : $r_z=0$ et $r_x-\ii r_y=-\ii$, donc $r_y=1$, $r_x=0$. Comme $|\vec r|=1$, l'état est pur : $\ket{{+}\ii}\bra{{+}\ii}$.}
  {Erreur de signe : $\rho_{01}=\frac{r_x-\ii r_y}{2}$, donc $\rho_{01}=-\frac\ii2$ correspond à $r_y=+1$ ($r_y=-1$ pour $\ket{{-}\ii}$).}
  {On a lu $\rho_{01}=-\frac\ii2$ comme $-\ii r_y$ : il faut multiplier par $2$ ($\rho=\frac12(\mathrm{Id}+\vec r\cdot\vec\sigma)$).}
  {Les diagonales $\frac12,\frac12$ donnent $r_z=0$, mais les éléments non diagonaux (cohérences) ne sont pas nuls : l'état est pur, et non $\frac{\mathrm{Id}}2$.}

% ---- Question 75 (niveau 3, correct : C)
\qcmq[2]{Chap.~\ref{chap:operateurs}\,$\cdot$\,\S\,\ref{sec:tenseur}}
  {Que donne $(X\otimes Z)\,\big(\ket{+}\otimes\ket{-}\big)$ ?}
  {$-\ket{-}\otimes\ket{-}$}
  {$\ket{-}\otimes\ket{+}$}
  {$\ket{+}\otimes\ket{+}$}
  {$\ket{+}\otimes\ket{-}$}
  {C}
\qcmx
  {On a échangé les facteurs ($Z\otimes X$) : $Z\ket+=\ket-$ et $X\ket-=-\ket-$. Le premier opérateur agit sur le premier qubit.}
  {On a cru que $X$ transforme $\ket+$ en $\ket-$ ; or $X$ échange $\ket0$ et $\ket1$ et laisse $\ket+$ inchangé.}
  {$(A\otimes B)(\ket u\otimes\ket v)=A\ket u\otimes B\ket v$ : $X\ket+=\ket+$ ($\ket+$ est propre pour $X$, valeur propre $+1$) et $Z\ket-=\ket+$ ($Z$ échange $\ket+$ et $\ket-$).}
  {On a cru que $Z$ laisse $\ket-$ inchangé ; $Z$ ne laisse invariants que $\ket0$ et $\ket1$, et échange $\ket+$ et $\ket-$.}

% ---- Question 76 (niveau 3, correct : C)
\qcmq[2]{Chap.~\ref{chap:portes}\,$\cdot$\,\S\,\ref{sec:phase}}
  {Le qubit $\ket{+}$ traverse la porte $T=\mathrm{diag}\big(1,\ee^{\ii\pi/4}\big)$, puis est mesuré dans la base $X$ $\{\ket+,\ket-\}$. Quelle est la probabilité d'obtenir $+$ ?}
  {$\frac12$}
  {$\frac{2-\sqrt2}{4}\approx0{,}15$}
  {$\frac{2+\sqrt2}{4}\approx0{,}85$}
  {$\frac{\sqrt2}{2}\approx0{,}71$}
  {C}
\qcmx
  {On a cru que la phase ne change pas les probabilités : c'est vrai en base $Z$, mais la base $X$ est sensible à la phase \emph{relative} (le point de Bloch est à l'azimut $\frac\pi4$ de l'axe $+x$).}
  {C'est $P(-)=\sin^2\frac\pi8$ : on a permuté les résultats $+$ et $-$.}
  {$T\ket+=\frac1{\sqrt2}(\ket0+\ee^{\ii\pi/4}\ket1)$, donc $\braket{+}{\psi}=\frac12\big(1+\ee^{\ii\pi/4}\big)$ et $P(+)=\frac14\big|1+\ee^{\ii\pi/4}\big|^2=\frac{2+2\cos\frac\pi4}{4}=\frac{2+\sqrt2}{4}$ (soit $\cos^2\frac\pi8$).}
  {On a pris $\cos\frac\pi4$ comme probabilité ; il faut le \emph{demi}-angle et le carré : $P=\cos^2\frac{\pi/4}{2}=\cos^2\frac\pi8$.}

% ---- Question 77 (niveau 3, correct : B)
\qcmq[2]{Chap.~\ref{chap:portes}\,$\cdot$\,\S\,\ref{sec:hadamard}}
  {Que vaut le produit de matrices $HYH$, où $Y=\begin{pmatrix}0&-\ii\\\ii&0\end{pmatrix}$ ?}
  {$Y$}
  {$-Y$}
  {$Z$}
  {$X$}
  {B}
\qcmx
  {Par analogie avec $HXH=Z$, on a cru que $H$ ne modifie pas $Y$ ; en fait l'axe $y$ est retourné (signe $-$).}
  {$H$ échange les axes $x$ et $z$ et retourne l'axe $y$ : $HXH=Z$, $HZH=X$, $HYH=-Y$. Direct : $HY=\frac{1}{\sqrt2}\begin{pmatrix}\ii&-\ii\\-\ii&-\ii\end{pmatrix}$, puis $HYH=\begin{pmatrix}0&\ii\\-\ii&0\end{pmatrix}=-Y$.}
  {C'est $HXH$ : $H$ n'échange que les axes $x$ et $z$, l'axe $y$ reste l'axe $y$.}
  {C'est $HZH$ : $Y$ ne peut pas devenir $X$, car $H$ laisse l'axe $y$ à sa place (au signe près).}

% ---- Question 78 (niveau 3, correct : D)
\qcmq[1]{Chap.~\ref{chap:portes}\,$\cdot$\,\S\,\ref{sec:op2}}
  {Que donne $(H\otimes H)\ket{01}$, où $\ket{01}=\ket0\otimes\ket1$ ?}
  {$\frac12\big(\ket{00}+\ket{01}-\ket{10}-\ket{11}\big)$}
  {$\frac12\big(\ket{00}-\ket{01}-\ket{10}+\ket{11}\big)$}
  {$\frac12\big(\ket{00}+\ket{01}+\ket{10}+\ket{11}\big)$}
  {$\frac12\big(\ket{00}-\ket{01}+\ket{10}-\ket{11}\big)$}
  {D}
\qcmx
  {C'est $\ket-\otimes\ket+$ : on a attribué le $\ket-$ au premier qubit, alors que c'est le second qui vaut $\ket1$.}
  {C'est $\ket-\otimes\ket-$ : $H\ket0=\ket+$, et non $\ket-$.}
  {C'est $\ket+\otimes\ket+$, le résultat pour $\ket{00}$ : on a oublié que $H\ket1=\ket-$ introduit des signes.}
  {$H\ket0\otimes H\ket1=\ket+\otimes\ket-=\frac12(\ket0+\ket1)\otimes(\ket0-\ket1)$. Le signe $-$ apparaît quand le \emph{second} chiffre vaut $1$ (formule générale $\frac12\sum_y(-1)^{x\cdot y}\ket y$ avec $x=01$).}

% ---- Question 79 (niveau 3, correct : D)
\qcmq[2]{Chap.~\ref{chap:portes}\,$\cdot$\,\S\,\ref{sec:cnot}}
  {Avec le CNOT du cours (contrôle : second qubit $B$ ; cible : premier qubit $A$ ; $\mathrm{CNOT}\ket{a,b}=\ket{a\oplus b,\,b}$), quel état obtient-on à partir de $\ket{1}_A\otimes\ket{+}_B$ ?}
  {$\ket{1}\otimes\ket{+}$ (inchangé)}
  {$\frac{1}{\sqrt2}\big(\ket{00}+\ket{11}\big)=\ket{\Phi^+}$}
  {$\frac{1}{\sqrt2}\big(\ket{01}-\ket{10}\big)=\ket{\Psi^-}$}
  {$\frac{1}{\sqrt2}\big(\ket{01}+\ket{10}\big)=\ket{\Psi^+}$}
  {D}
\qcmx
  {On a pris le premier qubit comme contrôle : le CNOT agirait alors sur $B$, et $X\ket+=\ket+$ laisserait l'état inchangé. Mais ici le contrôle est $B$.}
  {C'est le résultat pour $\ket{0}_A\ket{+}_B$ ; ici $A$ vaut $\ket1$, ce qui échange les deux termes.}
  {Aucun signe $-$ n'apparaît : le CNOT permute des états de base sans leur ajouter de phase.}
  {$\ket1\ket+=\frac1{\sqrt2}(\ket{10}+\ket{11})$ ; $\ket{10}\mapsto\ket{1\oplus0,0}=\ket{10}$ et $\ket{11}\mapsto\ket{1\oplus1,1}=\ket{01}$. C'est un état de Bell, donc intriqué.}

% ---- Question 80 (niveau 3, correct : C)
\qcmq[2]{Chap.~\ref{chap:portes}\,$\cdot$\,\S\,\ref{sec:swap}}
  {Que vaut $\mathrm{SWAP}\,(X\otimes\mathrm{Id})\,\mathrm{SWAP}$ ?}
  {$X\otimes\mathrm{Id}$}
  {$\mathrm{Id}\otimes\mathrm{Id}$}
  {$\mathrm{Id}\otimes X$}
  {$X\otimes X$}
  {C}
\qcmx
  {On a supposé que $\mathrm{SWAP}$ commute avec $X\otimes\mathrm{Id}$ ; ce n'est pas le cas, puisque cet opérateur n'agit pas de la même façon sur les deux qubits.}
  {On a simplifié $\mathrm{SWAP}^2=\mathrm{Id}$ en oubliant que $X\otimes\mathrm{Id}$ se trouve \emph{entre} les deux SWAP : on ne peut pas les rapprocher.}
  {$\mathrm{SWAP}\,(X\otimes\mathrm{Id})\,\mathrm{SWAP}\ket{a,b}=\mathrm{SWAP}\,(X\otimes\mathrm{Id})\ket{b,a}=\mathrm{SWAP}\ket{b\oplus1,a}=\ket{a,b\oplus1}$ : le SWAP transporte l'opérateur d'un qubit à l'autre (même mécanisme que $\mathrm{CNOT}_{21}=\mathrm{SWAP}\,\mathrm{CNOT}_{12}\,\mathrm{SWAP}$).}
  {Cet opérateur inverserait les deux bits ; ici un seul bit est inversé, celui du \emph{second} qubit.}

% ---- Question 81 (niveau 3, correct : C)
\qcmq[1]{Chap.~\ref{chap:portes}\,$\cdot$\,\S\,\ref{sec:controlees}}
  {La porte $Z$ contrôlée ($\mathrm{CZ}=\mathrm{diag}(1,1,1,-1)$, contrôle : premier qubit) est appliquée à $\ket{+}\otimes\ket{+}$. Quel état obtient-on ?}
  {$\ket{+}\otimes\ket{+}$}
  {$\ket{-}\otimes\ket{-}$}
  {$\frac12\big(\ket{00}+\ket{01}+\ket{10}-\ket{11}\big)$, un état intriqué}
  {$\ket{+}\otimes\ket{-}$}
  {C}
\qcmx
  {On a cru qu'un simple signe est une phase globale ; mais le signe $-$ ne touche que $\ket{11}$ : c'est une phase \emph{relative}.}
  {C'est $Z\otimes Z$ ($Z$ sur les deux qubits), et non une porte contrôlée.}
  {CZ ne change que le signe de $\ket{11}$. Critère de séparabilité : $c_{00}c_{11}-c_{01}c_{10}=-\frac14-\frac14=-\frac12\neq0$ ; on peut l'écrire $\frac{1}{\sqrt2}\big(\ket0\ket++\ket1\ket-\big)$.}
  {On a appliqué $Z$ au second qubit sans condition ; le contrôle exige que le premier qubit soit dans $\ket1$.}

% ---- Question 82 (niveau 3, correct : D)
\qcmq[2]{Chap.~\ref{chap:portes}\,$\cdot$\,\S\,\ref{sec:phase}, \ref{sec:portes-bloch}}
  {Un qubit préparé dans $\ket{0}$ traverse successivement $H$, puis $S$, puis $H$ (dans cet ordre). Quel est l'état final, à une phase globale près ?}
  {$\ket{{+}\ii}$}
  {$\ket{0}$}
  {$\ket{1}$}
  {$\ket{{-}\ii}=\frac{1}{\sqrt2}\big(\ket0-\ii\ket1\big)$}
  {D}
\qcmx
  {Mauvais sens de rotation : on a perdu un signe en calculant $H\ket{{+}\ii}$ (la phase relative vaut $-\ii$, pas $+\ii$).}
  {On a « simplifié » $HSH$ en $S$ (puis $S\ket0=\ket0$) : mais $H$ ne commute pas avec $S$, on ne peut pas supprimer les deux $H$.}
  {C'est le résultat pour $HZH=X$ ; or $S$ n'est qu'un \emph{quart} de tour autour de $z$ (et non un demi-tour comme $Z$), donc $HSH$ est un quart de tour autour de $x$.}
  {$HSH=\sqrt X$ est un quart de tour autour de $x$ (§\,\ref{sec:phase}). Directement : $H\ket0=\ket+$, $S\ket+=\ket{{+}\ii}$, $H\ket{{+}\ii}=\frac{1+\ii}{2}\ket0+\frac{1-\ii}{2}\ket1=\frac{1+\ii}{2}\big(\ket0-\ii\ket1\big)$.}

% ---- Question 83 (niveau 3, correct : A)
\qcmq[1]{Chap.~\ref{chap:bell}\,$\cdot$\,\S\,\ref{sec:mesure2}}
  {Deux qubits sont dans l'état $\ket{\Psi}=\frac{1}{\sqrt3}\big(\ket{00}+\ket{01}+\ket{11}\big)$. On mesure le \emph{second} qubit $B$ en base $Z$ et on obtient $1$. Que valait la probabilité de ce résultat, et dans quel état se trouve alors $A$ ?}
  {$P(B{=}1)=\frac23$ ; $A$ est dans l'état $\ket{+}$}
  {$P(B{=}1)=\frac13$ ; $A$ est dans l'état $\ket{0}$}
  {$P(B{=}1)=\frac23$ ; $A$ est dans l'état $\ket{1}$}
  {$P(B{=}1)=\frac13$ ; $A$ est dans l'état $\ket{+}$}
  {A}
\qcmx
  {$P(B{=}1)=|c_{01}|^2+|c_{11}|^2=\frac13+\frac13$. On garde les termes dont le second chiffre vaut $1$ : $\frac{1}{\sqrt3}\big(\ket{01}+\ket{11}\big)=\frac{1}{\sqrt3}(\ket0+\ket1)\otimes\ket1$, qui renormalisé donne $\ket{+}_A\otimes\ket{1}_B$.}
  {C'est le résultat $B=0$ : seul le terme $\ket{00}$ a un second chiffre nul, d'où la probabilité $\frac13$ et $A$ dans $\ket0$.}
  {On a projeté sur le seul terme $\ket{11}$ ; or $\ket{01}$ a aussi son second chiffre égal à $1$ : $A$ vaut $0$ ou $1$ (état $\ket+$).}
  {L'état de $A$ est juste, mais on n'a compté qu'un seul terme pour la probabilité ; il faut sommer $|c_{01}|^2+|c_{11}|^2$.}

% ---- Question 84 (niveau 3, correct : A)
\qcmq[2]{Chap.~\ref{chap:bell}\,$\cdot$\,\S\,\ref{sec:separable}}
  {Pour quelle valeur de $x$ l'état (non normalisé) $\ket{00}+2\ket{01}+3\ket{10}+x\ket{11}$ est-il séparable ?}
  {$x=6$}
  {$x=-6$}
  {$x=\frac16$}
  {$x=5$}
  {A}
\qcmx
  {Critère $D=c_{00}c_{11}-c_{01}c_{10}=x-2\cdot3=0$. On obtient bien un produit : $(\ket0+3\ket1)\otimes(\ket0+2\ket1)=\ket{00}+2\ket{01}+3\ket{10}+6\ket{11}$.}
  {Erreur de signe : $D=x-6$ s'annule pour $x=+6$ ; pour $x=-6$, $D=-12\neq0$ et l'état est intriqué.}
  {On a inversé le rapport : $c_{00}c_{11}=c_{01}c_{10}$ donne $x=\frac{c_{01}c_{10}}{c_{00}}=6$, et non $\frac{c_{00}}{c_{01}c_{10}}$.}
  {On a additionné ($2+3$) au lieu de multiplier : le critère de séparabilité porte sur des \emph{produits} de coefficients.}

% ---- Question 85 (niveau 3, correct : B)
\qcmq[2]{Chap.~\ref{chap:bell}\,$\cdot$\,\S\,\ref{sec:bell}}
  {Comment s'écrit $\ket{00}$ dans la base de Bell ?}
  {$\frac{1}{\sqrt2}\big(\ket{\Phi^+}-\ket{\Phi^-}\big)$}
  {$\frac{1}{\sqrt2}\big(\ket{\Phi^+}+\ket{\Phi^-}\big)$}
  {$\frac{1}{\sqrt2}\big(\ket{\Phi^+}+\ket{\Psi^+}\big)$}
  {$\ket{\Phi^+}$}
  {B}
\qcmx
  {Erreur de signe : cette combinaison vaut $\ket{11}$, puisque $\ket{\Phi^+}-\ket{\Phi^-}=\sqrt2\,\ket{11}$.}
  {$\ket{\Phi^+}+\ket{\Phi^-}=\frac{1}{\sqrt2}\cdot2\ket{00}=\sqrt2\,\ket{00}$. Une mesure de Bell sur $\ket{00}$ donne donc $\Phi^+$ ou $\Phi^-$ avec la probabilité $\frac12$ chacun (jamais $\Psi^\pm$).}
  {Cette combinaison vaut $\ket{+}\otimes\ket{+}=\frac12(\ket{00}+\ket{01}+\ket{10}+\ket{11})$ : les deux états de Bell « pairs » et « impairs » s'y ajoutent.}
  {On a confondu $\ket{00}$, état séparable, avec $\ket{\Phi^+}=\frac{1}{\sqrt2}(\ket{00}+\ket{11})$, qui contient aussi $\ket{11}$ et est intriqué.}

% ---- Question 86 (niveau 3, correct : A)
\qcmq[1]{Chap.~\ref{chap:bell}\,$\cdot$\,\S\,\ref{sec:circuit-bell}}
  {Le générateur de Bell du cours est $G=\mathrm{CNOT}\cdot(\mathrm{Id}\otimes H)$ : $H$ sur le second qubit $B$, puis CNOT de contrôle $B$ et de cible $A$. Quel circuit permet de \emph{mesurer} une paire dans la base de Bell (on le place avant une mesure en base $Z$) ?}
  {D'abord le CNOT (contrôle $B$, cible $A$), puis $H$ sur $B$}
  {D'abord $H$ sur $B$, puis le CNOT}
  {D'abord le CNOT, puis $H$ sur $A$ (la cible)}
  {Aucun circuit : on mesure directement les deux qubits en base $Z$}
  {A}
\qcmx
  {C'est l'inverse $G^\dagger=(\mathrm{Id}\otimes H)\,\mathrm{CNOT}$ : les portes sont leurs propres inverses et l'ordre s'inverse. Il envoie $\Phi^+,\Phi^-,\Psi^+,\Psi^-$ sur $\ket{00},\ket{01},\ket{10},\ket{11}$ (à une phase près).}
  {C'est $G$ lui-même, qui \emph{fabrique} des états de Bell : on garde le même ordre au lieu de l'inverser, et la mesure ne distingue plus les quatre états.}
  {Mauvais qubit : pour $\Phi^+$ on trouve $\ket{+}\otimes\ket{+}$, dont les deux bits sont aléatoires. Le $H$ doit agir sur $B$, qui a reçu le $H$ dans $G$.}
  {$\Phi^+$ et $\Phi^-$ (comme $\Psi^+$ et $\Psi^-$) donnent les mêmes résultats en base $Z$ : cette base ne voit pas le signe relatif, il faut d'abord défaire l'intrication.}

% ---- Question 87 (niveau 3, correct : B)
\qcmq[2]{Chap.~\ref{chap:bell}\,$\cdot$\,\S\,\ref{sec:trace}}
  {Quelle est la matrice de densité réduite $\rho_A=\operatorname{Tr}_B\ket\Psi\bra\Psi$ du premier qubit pour $\ket{\Psi}=\frac12\ket{00}+\frac{\sqrt3}{2}\ket{11}$ ?}
  {$\begin{pmatrix}\frac14&\frac{\sqrt3}{4}\\\frac{\sqrt3}{4}&\frac34\end{pmatrix}$}
  {$\begin{pmatrix}\frac14&0\\0&\frac34\end{pmatrix}$}
  {$\begin{pmatrix}\frac12&0\\0&\frac12\end{pmatrix}$}
  {$\begin{pmatrix}\frac12&0\\0&\frac{\sqrt3}{2}\end{pmatrix}$}
  {B}
\qcmx
  {C'est la matrice $\ket\chi\bra\chi$ de l'état pur $\ket\chi=\frac12\ket0+\frac{\sqrt3}2\ket1$ : on a gardé les cohérences, comme si $A$ était seul ; l'intrication avec $B$ les détruit.}
  {Avec $C=\begin{pmatrix}c_{00}&c_{01}\\c_{10}&c_{11}\end{pmatrix}=\begin{pmatrix}\frac12&0\\0&\frac{\sqrt3}2\end{pmatrix}$, $\rho_A=CC^\dagger=\mathrm{diag}\big(\frac14,\frac34\big)$ : état mixte, de pureté $\frac{10}{16}=1-2|D|^2$ avec $D=\frac{\sqrt3}{4}$.}
  {C'est la réponse pour un état de Bell (intrication maximale) ; ici les poids $\frac14$ et $\frac34$ diffèrent : l'intrication est partielle.}
  {On a mis les amplitudes au lieu de leurs carrés : la trace vaudrait $\frac{1+\sqrt3}{2}\neq1$.}

% ---- Question 88 (niveau 3, correct : A)
\qcmq[2]{Chap.~\ref{chap:bell}\,$\cdot$\,\S\,\ref{sec:info}}
  {Deux qubits sont dans le mélange statistique $\rho_{AB}=\frac12\ket{00}\bra{00}+\frac12\ket{11}\bra{11}$ (corrélations purement classiques). Que vaut $H(A|B)=S(\rho_{AB})-S(\rho_B)$, en bits ?}
  {$0$}
  {$-1$}
  {$1$}
  {$2$}
  {A}
\qcmx
  {$\rho_{AB}$ a pour valeurs propres $\frac12,\frac12,0,0$, donc $S(\rho_{AB})=1$ ; et $\rho_B=\frac12\mathrm{Id}$ donne $S(\rho_B)=1$. Ainsi $H(A|B)=0$ : connaître $B$ fixe $A$. Une valeur négative est propre à l'intrication.}
  {C'est la valeur pour l'état de Bell \emph{pur} $\ket{\Phi^+}$ : mêmes résultats en base $Z$, mais $S(\rho_{AB})=0$ car l'état global est pur. Ici l'état est mixte.}
  {On a pris $H(A|B)=H(A)$, comme si $B$ ne disait rien sur $A$ ; or les deux bits sont toujours égaux.}
  {On a additionné les entropies des deux qubits ($1+1$), valeur de deux bits indépendants : ce n'est pas une entropie conditionnelle.}

% ---- Question 89 (niveau 3, correct : B)
\qcmq[1]{Chap.~\ref{chap:bell}\,$\cdot$\,\S\,\ref{sec:teleportation}}
  {Dans la téléportation, avant de recevoir les deux bits $(a,q)$ d'Alice, dans quel état se trouve le qubit de Bob ?}
  {Dans l'état pur $\alpha\ket0+\beta\ket1$, déjà téléporté}
  {Dans l'état complètement mélangé $\rho_B=\frac12\mathrm{Id}$, indépendant de $\alpha$ et $\beta$}
  {Dans le mélange $|\alpha|^2\ket0\bra0+|\beta|^2\ket1\bra1$}
  {Dans l'état pur $\ket{0}$, son état initial}
  {B}
\qcmx
  {Ce serait une communication instantanée. L'état $\alpha\ket0+\beta\ket1$ n'apparaît qu'\emph{après} la correction, qui dépend de $(a,q)$.}
  {Les quatre kets $\ket{\phi_{aq}}$ sont équiprobables ($\frac14$) et $\frac14\sum_{a,q}\ket{\phi_{aq}}\bra{\phi_{aq}}=\frac12\mathrm{Id}$. Bob n'apprend rien avant le message classique : rien ne voyage plus vite que la lumière.}
  {Bob connaîtrait alors les statistiques de $\ket\varphi$ sans aucun message : or les probabilités d'Alice ($\frac14$ pour chaque résultat) ne dépendent pas de $\alpha,\beta$.}
  {Le qubit de Bob est la moitié de la paire $\ket{\Phi^+}$ : pris seul, il n'a pas d'état pur ($\rho_B=\frac12\mathrm{Id}$ avant comme après la mesure d'Alice, tant que Bob ignore le résultat).}

% ---- Question 90 (niveau 3, correct : D)
\qcmq[1]{Chap.~\ref{chap:bell}\,$\cdot$\,\S\,\ref{sec:chsh}}
  {Au jeu CHSH (gain si $a\oplus b=x\cdot y$), quelles sont les meilleures probabilités de gain avec une stratégie classique, même avec du hasard partagé, et avec une paire $\ket{\Phi^+}$ mesurée dans des bases bien choisies ?}
  {$\frac12$ en classique ; $1$ avec l'intrication}
  {$\frac34$ dans les deux cas}
  {$1$ en classique si Alice et Bob s'accordent à l'avance ; aucun avantage quantique}
  {$\frac34$ en classique ; $\cos^2\frac\pi8\approx0{,}85$ avec l'intrication}
  {D}
\qcmx
  {Classique : répondre toujours $0$ donne déjà $\frac34$. Quantique : l'intrication ne permet pas de gagner à tous les coups (borne de Tsirelson $\cos^2\frac\pi8$).}
  {Si l'intrication n'apportait rien, on n'aurait pas de violation de l'inégalité de Bell ; l'avantage quantique ($0{,}85>0{,}75$) est précisément ce que l'expérience de Bell met en évidence.}
  {S'accorder à l'avance (ou partager du hasard) ne suffit pas : aucune stratégie déterministe ne gagne les quatre parties, d'après l'argument de parité.}
  {Une stratégie déterministe perd au moins une des quatre parties (la somme modulo $2$ des conditions donne $0=1$) ; avec $\Phi^+$, chaque partie est gagnée avec la probabilité $\cos^2\frac\pi8=\frac{2+\sqrt2}{4}$.}

\clearpage
%% ======================================================================
\subsection{Niveau 4 : maîtrise (questions 91 à 100)}\label{sec:qcm-n4}
%% ======================================================================
\noindent{\color{insagris}\small Niveau de difficulté : Maîtrise\quad $\bullet\bullet\bullet\bullet$}\par\smallskip
% ---- Question 91 (niveau 4, correct : C)
\qcmq[2]{Chap.~\ref{chap:histoire}\,$\cdot$\,\S\,\ref{sec:superposition}}
  {Une particule du puits infini est préparée à $t=0$ dans l'état $\Psi=\frac{1}{\sqrt2}\big(\psi_1+\psi_2\big)$, avec $E_n=n^2E_1$. Au bout de quelle durée minimale la densité de probabilité $|\Psi(x,t)|^2$ retrouve-t-elle sa forme initiale ?}
  {$T=\dfrac{2\pi\hbar}{E_1}$}
  {$T=\dfrac{2\pi\hbar}{5E_1}$}
  {$T=\dfrac{2\pi\hbar}{3E_1}$}
  {$T=\dfrac{\pi\hbar}{3E_1}$}
  {C}
\qcmx
  {C'est la période de la phase $\ee^{-\ii E_1t/\hbar}$ de $\psi_1$ seul ; or les phases globales ne se voient pas dans $|\Psi|^2$ : seul compte l'écart $E_2-E_1=3E_1$.}
  {On a pris la somme $E_1+E_2=5E_1$ ; c'est la \emph{différence} des énergies (phase relative) qui intervient.}
  {$|\Psi|^2=\frac12(\psi_1^2+\psi_2^2)+\psi_1\psi_2\cos\frac{(E_2-E_1)t}{\hbar}$ : seul le terme croisé dépend du temps, à la pulsation $\omega=\frac{E_2-E_1}{\hbar}=\frac{3E_1}{\hbar}$, donc $T=\frac{2\pi}{\omega}$.}
  {C'est la demi-période : à $t=\frac T2$ le cosinus vaut $-1$ et la densité est le symétrique de l'état initial ($x\to a-x$), pas encore sa forme de départ.}

% ---- Question 92 (niveau 4, correct : D)
\qcmq[1]{Chap.~\ref{chap:dirac}\,$\cdot$\,\S\,\ref{sec:projecteurs}, \ref{sec:qubit}}
  {Un expérimentateur reçoit un qubit préparé soit dans $\ket{0}$, soit dans $\ket{+}$. Laquelle de ces affirmations est correcte ?}
  {Une mesure en base $Z$ les distingue : $\ket0$ donne $0$ et $\ket+$ donne $1$}
  {Une mesure en base $X$ les distingue : $\ket+$ donne $+$ et $\ket0$ donne $-$}
  {On les distingue en répétant la mesure sur le même qubit}
  {Aucune mesure ne permet de les distinguer à coup sûr, car $\braket{0}{+}=\frac{1}{\sqrt2}\neq0$}
  {D}
\qcmx
  {$\ket+$ donne $0$ ou $1$ avec la probabilité $\frac12$ : le résultat $0$ est possible pour les deux états.}
  {$\ket0=\frac{1}{\sqrt2}(\ket++\ket-)$ donne $+$ ou $-$ avec la probabilité $\frac12$ : le résultat $+$ est possible pour les deux états.}
  {Après une première mesure l'état est effondré : répéter ne donne aucune information nouvelle, et on ne peut pas copier le qubit pour recommencer (non-clonage).}
  {Seuls des états \emph{orthogonaux} sont parfaitement discernables. En base $Z$, $\ket0$ donne toujours $0$ et $\ket+$ donne $0$ une fois sur deux : obtenir $0$ ne tranche pas (obtenir $1$ prouve seulement qu'on avait $\ket+$).}

% ---- Question 93 (niveau 4, correct : C)
\qcmq[1]{Chap.~\ref{chap:operateurs}\,$\cdot$\,\S\,\ref{sec:pur-mixte}}
  {Une source envoie au hasard $\ket{0}$ ou $\ket{1}$ (probabilité $\frac12$ chacun) ; une autre envoie au hasard $\ket{+}$ ou $\ket{-}$ (probabilité $\frac12$ chacun). Peut-on distinguer ces deux sources en mesurant les qubits émis ?}
  {Oui : en base $Z$, la première source donne des résultats prévisibles}
  {Oui : en base $X$, la seconde source donne des résultats prévisibles}
  {Non : les deux ont la même matrice de densité $\frac{\mathrm{Id}}{2}$, donc les mêmes statistiques pour toute mesure}
  {Oui : les états purs envoyés ne sont pas les mêmes}
  {C}
\qcmx
  {Chaque qubit est bien $\ket0$ ou $\ket1$, mais on ignore lequel : le résultat vaut $0$ ou $1$ avec la probabilité $\frac12$, exactement comme avec la seconde source ($|\braket{0}{\pm}|^2=\frac12$).}
  {Même erreur : chaque qubit est $\ket+$ ou $\ket-$, mais on ignore lequel ; le résultat en base $X$ est aléatoire, et la première source donne aussi $+$ ou $-$ avec la probabilité $\frac12$.}
  {$\frac12\ket0\bra0+\frac12\ket1\bra1=\frac12\ket+\bra++\frac12\ket-\bra-=\frac{\mathrm{Id}}2$ (relation de fermeture dans chaque base). Les probabilités de toute mesure ne dépendent que de $\rho$.}
  {Ce qui compte n'est pas la liste des états purs envoyés mais la matrice de densité moyenne : des préparations différentes peuvent avoir la même $\rho$.}

% ---- Question 94 (niveau 4, correct : B)
\qcmq[1]{Chap.~\ref{chap:operateurs}\,$\cdot$\,\S\,\ref{sec:pur-mixte} ; Chap.~\ref{chap:portes}\,$\cdot$\,\S\,\ref{sec:portes-bloch}}
  {Un qubit est dans l'état complètement mélangé $\rho=\frac{\mathrm{Id}}{2}$. Existe-t-il une porte quantique $U$ qui le transforme en un état pur, par exemple $\ket{0}\bra{0}$ ?}
  {Oui : la porte de Hadamard, qui crée des superpositions}
  {Non : $\rho\mapsto U\rho U^\dagger$ conserve les valeurs propres, donc la pureté ; ici $U\frac{\mathrm{Id}}{2}U^\dagger=\frac{\mathrm{Id}}{2}$}
  {Oui : une rotation de Bloch qui amène le centre de la boule au pôle nord}
  {Oui : une mesure en base $Z$, qui donne $\ket0$}
  {B}
\qcmx
  {$H$ envoie un état pur sur un état pur ($\ket0\mapsto\ket+$) et laisse $\frac{\mathrm{Id}}{2}$ invariant : $H\frac{\mathrm{Id}}2H=\frac{\mathrm{Id}}2$ car $H^2=\mathrm{Id}$.}
  {$\operatorname{Tr}\big[(U\rho U^\dagger)^2\big]=\operatorname{Tr}\rho^2$ ($=\frac12$ ici). Sur la boule de Bloch, une porte est une rotation : elle laisse le centre fixe.}
  {Une rotation tourne la boule autour de son centre, qui reste fixe : elle ne peut pas l'amener sur la sphère.}
  {Une mesure n'est pas une porte (elle est non unitaire) ; de plus elle donne $\ket0$ ou $\ket1$ au hasard, pas $\ket0$ à coup sûr.}

% ---- Question 95 (niveau 4, correct : D)
\qcmq[1]{Chap.~\ref{chap:portes}\,$\cdot$\,\S\,\ref{sec:cnot}}
  {Que vaut $(H\otimes H)\,\mathrm{CNOT}_{12}\,(H\otimes H)$, où $\mathrm{CNOT}_{12}$ a pour contrôle le \emph{premier} qubit et pour cible le second ?}
  {$\mathrm{CNOT}_{12}$, inchangé}
  {$\mathrm{CZ}$ (porte $Z$ contrôlée)}
  {$\mathrm{SWAP}$}
  {$\mathrm{CNOT}_{21}$ (contrôle : second qubit ; cible : premier qubit)}
  {D}
\qcmx
  {On a cru que les deux $H$ s'annulent ($H^2=\mathrm{Id}$) ; mais ils encadrent le CNOT au lieu d'être côte à côte : on ne peut pas les simplifier.}
  {C'est le résultat quand $H$ n'encadre que la \emph{cible} : $(\mathrm{Id}\otimes H)\,\mathrm{CNOT}_{12}\,(\mathrm{Id}\otimes H)=\mathrm{CZ}$, car $HXH=Z$. Ici $H$ agit aussi sur le contrôle.}
  {Le SWAP s'obtient avec \emph{trois} CNOT alternés, pas avec un seul CNOT conjugué par $H\otimes H$ ; il n'y a pas de contrôle dans un SWAP.}
  {$H\ket0\bra0H=\ket+\bra+$, $H\ket1\bra1H=\ket-\bra-$ et $HXH=Z$ donnent $\ket+\bra+\otimes I+\ket-\bra-\otimes Z=I\otimes\ket0\bra0+X\otimes\ket1\bra1=\mathrm{CNOT}_{21}$ : dans la base $\{\ket+,\ket-\}$, contrôle et cible échangent leurs rôles (rebond de phase).}

% ---- Question 96 (niveau 4, correct : A)
\qcmq[1]{Chap.~\ref{chap:portes}\,$\cdot$\,\S\,\ref{sec:controlees}}
  {On applique la porte de Toffoli (contrôles : qubits 1 et 2 ; cible : qubit 3) à $\ket{+}\otimes\ket{+}\otimes\ket{0}$, puis on mesure le troisième qubit en base $Z$. Que peut-on dire si l'on obtient $1$ ?}
  {Ce résultat a la probabilité $\frac14$, et les deux premiers qubits sont alors dans l'état $\ket{11}$}
  {Ce résultat a la probabilité $\frac12$, et les deux premiers qubits restent dans $\ket{+}\ket{+}$}
  {Ce résultat a la probabilité $\frac14$, mais les deux premiers qubits restent dans $\ket{+}\ket{+}$}
  {Ce résultat a la probabilité $\frac34$, et les deux premiers qubits sont dans l'état $\frac{1}{\sqrt3}\big(\ket{00}+\ket{01}+\ket{10}\big)$}
  {A}
\qcmx
  {$\mathrm{CCNOT}\ket{x,y,0}=\ket{x,y,xy}$ donne $\frac12\big(\ket{000}+\ket{010}+\ket{100}+\ket{111}\big)$ : seul $\ket{111}$ a un troisième chiffre égal à $1$ (probabilité $\frac14$), et il impose $x=y=1$.}
  {La cible n'est inversée que si $x=y=1$ (probabilité $\frac14$) et non une fois sur deux ; de plus l'état est intriqué, donc mesurer le qubit 3 modifie les qubits 1 et 2.}
  {La probabilité est juste, mais la mesure de la cible n'est pas sans effet sur les contrôles : l'état après la Toffoli est intriqué, donc le résultat $1$ les projette sur $\ket{11}$.}
  {C'est le cas du résultat $0$ (probabilité $\frac34$ : $x,y\in\{00,01,10\}$). Le résultat $1$ est le cas rare, celui du ET vrai ($x=y=1$).}

% ---- Question 97 (niveau 4, correct : A)
\qcmq[1]{Chap.~\ref{chap:bell}\,$\cdot$\,\S\,\ref{sec:bell}}
  {Que donne $(H\otimes H)\ket{\Psi^-}$, avec $\ket{\Psi^-}=\frac{1}{\sqrt2}\big(\ket{01}-\ket{10}\big)$ ?}
  {$-\ket{\Psi^-}$, c'est-à-dire le même état physique}
  {$\ket{\Phi^-}$}
  {$\ket{\Psi^+}$}
  {$\ket{\Phi^+}$}
  {A}
\qcmx
  {Le singulet est invariant par $U\otimes U$ à un facteur près : $(U\otimes U)\ket{\Psi^-}=\det(U)\ket{\Psi^-}$ avec $\det H=-1$. Directement : $\frac{1}{\sqrt2}\big(\ket+\ket--\ket-\ket+\big)=-\ket{\Psi^-}$.}
  {C'est l'image de $\ket{\Psi^+}$, pas de $\ket{\Psi^-}$ : $(H\otimes H)\ket{\Psi^+}=\ket{\Phi^-}$. Le signe $-$ du singulet change le calcul (les termes $\ket{00}$ et $\ket{11}$ se compensent).}
  {$\ket{\Psi^+}$ est symétrique par échange des deux qubits, $\ket{\Psi^-}$ antisymétrique ; $H\otimes H$ commute avec le SWAP et conserve donc cette symétrie.}
  {C'est l'état que $H\otimes H$ laisse invariant ($\ket{\Phi^+}=\frac{1}{\sqrt2}(\ket{++}+\ket{--})$) ; le singulet n'a pas ce comportement.}

% ---- Question 98 (niveau 4, correct : B)
\qcmq[1]{Chap.~\ref{chap:bell}\,$\cdot$\,\S\,\ref{sec:chsh}}
  {Une expérience de type CHSH (avec $S=E_{00}+E_{01}+E_{10}-E_{11}$) mesure $S=2{,}6$. Que peut-on en conclure ?}
  {C'est impossible : on a toujours $|S|\leq2$}
  {Aucun modèle local à résultats prédéterminés ne l'explique, mais c'est compatible avec la mécanique quantique}
  {C'est impossible en mécanique quantique, car $2\sqrt2\approx2{,}4<2{,}6$}
  {Les particules s'échangent de l'information plus vite que la lumière}
  {B}
\qcmx
  {Cette borne ne vaut que pour les modèles classiques (locaux) ; l'expérience et la mécanique quantique la dépassent, jusqu'à $2\sqrt2$.}
  {Classique : $|S|\leq2$ ; quantique : $|S|\leq2\sqrt2\approx2{,}83$ (borne de Tsirelson). Comme $2<2{,}6<2\sqrt2$, l'inégalité de Bell est violée sans dépasser le maximum quantique ; le gain correspondant est $\frac12+\frac{2{,}6}{8}\approx0{,}825$.}
  {Erreur numérique : $2\sqrt2\approx2{,}83$, donc $2{,}6$ est permis.}
  {Les corrélations violent l'inégalité de Bell mais ne permettent aucun signal : la réponse de Bob est uniforme quelle que soit la base choisie par Alice.}

% ---- Question 99 (niveau 4, correct : A)
\qcmq[2]{Chap.~\ref{chap:bell}\,$\cdot$\,\S\,\ref{sec:info}}
  {Pour l'état $\ket{\Psi}=\frac{1}{\sqrt3}\big(\ket{00}+\sqrt2\,\ket{11}\big)$, quelle est l'entropie de von Neumann $S(\rho_A)$ du premier qubit, en bits ?}
  {$\log_2 3-\frac23\approx0{,}92$}
  {$1$}
  {$0$}
  {$0{,}64$}
  {A}
\qcmx
  {$\rho_A=\mathrm{diag}\big(\frac13,\frac23\big)$ (trace partielle), donc $S=-\frac13\log_2\frac13-\frac23\log_2\frac23=\log_23-\frac23\approx1{,}585-0{,}667$. C'est moins que le maximum d'un bit (état de Bell) : l'intrication est partielle.}
  {C'est la valeur pour un état de Bell ($\rho_A=\frac{\mathrm{Id}}{2}$) ; ici les poids $\frac13$ et $\frac23$ diffèrent, donc $S<1$.}
  {$S(\rho_A)=0$ vaut pour un état séparable ; ici $D=\frac{\sqrt2}{3}\neq0$ : l'état est intriqué, $\rho_A$ est mixte et $S>0$ (le couple $AB$ est pur, pas $A$ seul).}
  {On a utilisé le logarithme \emph{népérien} (résultat en nats) ; en bits il faut $\log_2$, ce qui donne $0{,}92$.}

% ---- Question 100 (niveau 4, correct : B)
\qcmq[1]{Chap.~\ref{chap:bell}\,$\cdot$\,\S\,\ref{sec:bell}}
  {Pourquoi peut-on mesurer simultanément les deux « parités » $Z\otimes Z$ et $X\otimes X$ de deux qubits, alors qu'on ne peut pas mesurer simultanément $Z$ et $X$ sur un seul qubit ?}
  {Parce que l'incertitude de Heisenberg ne s'applique pas à deux qubits}
  {Parce que $Z\otimes Z$ et $X\otimes X$ commutent : les deux signes $-1$ de $ZX=-XZ$ se compensent}
  {Parce que les deux qubits sont indépendants}
  {Parce que $Z\otimes Z$ et $X\otimes X$ agissent sur des qubits différents}
  {B}
\qcmx
  {Elle s'applique toujours : $Z\otimes\mathrm{Id}$ et $X\otimes\mathrm{Id}$ ne commutent pas. C'est la commutation de ces deux opérateurs précis qui permet la mesure simultanée.}
  {$(Z\otimes Z)(X\otimes X)=ZX\otimes ZX=(-XZ)\otimes(-XZ)=XZ\otimes XZ=(X\otimes X)(Z\otimes Z)$. Ils ont donc une base propre commune : la base de Bell.}
  {Dans un état de Bell les qubits sont intriqués, donc loin d'être indépendants ; la commutation est une propriété des opérateurs, pas de l'état.}
  {Chacun des deux opérateurs agit sur \emph{les deux} qubits à la fois. (Ce sont $Z\otimes\mathrm{Id}$ et $\mathrm{Id}\otimes X$, sur des qubits différents, qui commutent pour cette raison.)}

\label{qcm:fin}

\clearpage
\ifqcmcorrige

% ======================================================================
\subsection{Clé des réponses et fiche de score}\label{sec:qcm-cle}
% ======================================================================
\noindent\textbf{Cette partie donne les réponses : ne la lisez qu'après avoir répondu à toutes les
questions.}

\paragraph{Clé des réponses.} La lettre de la bonne réponse de la question $n$ se lit à la ligne
qui contient $n$, dans la colonne du rang de $n$ dans la ligne (question~24 : ligne 21--30,
colonne 4).
\begin{center}
\renewcommand{\arraystretch}{1.25}
\setlength{\tabcolsep}{5.5pt}
\begin{tabular}{@{}l*{10}{c}@{}}
\toprule
 & 1 & 2 & 3 & 4 & 5 & 6 & 7 & 8 & 9 & 10 \\
\midrule
1--10 & \qcmkey{1} & \qcmkey{2} & \qcmkey{3} & \qcmkey{4} & \qcmkey{5} & \qcmkey{6} & \qcmkey{7} & \qcmkey{8} & \qcmkey{9} & \qcmkey{10} \\
11--20 & \qcmkey{11} & \qcmkey{12} & \qcmkey{13} & \qcmkey{14} & \qcmkey{15} & \qcmkey{16} & \qcmkey{17} & \qcmkey{18} & \qcmkey{19} & \qcmkey{20} \\
21--30 & \qcmkey{21} & \qcmkey{22} & \qcmkey{23} & \qcmkey{24} & \qcmkey{25} & \qcmkey{26} & \qcmkey{27} & \qcmkey{28} & \qcmkey{29} & \qcmkey{30} \\
31--40 & \qcmkey{31} & \qcmkey{32} & \qcmkey{33} & \qcmkey{34} & \qcmkey{35} & \qcmkey{36} & \qcmkey{37} & \qcmkey{38} & \qcmkey{39} & \qcmkey{40} \\
41--50 & \qcmkey{41} & \qcmkey{42} & \qcmkey{43} & \qcmkey{44} & \qcmkey{45} & \qcmkey{46} & \qcmkey{47} & \qcmkey{48} & \qcmkey{49} & \qcmkey{50} \\
51--60 & \qcmkey{51} & \qcmkey{52} & \qcmkey{53} & \qcmkey{54} & \qcmkey{55} & \qcmkey{56} & \qcmkey{57} & \qcmkey{58} & \qcmkey{59} & \qcmkey{60} \\
61--70 & \qcmkey{61} & \qcmkey{62} & \qcmkey{63} & \qcmkey{64} & \qcmkey{65} & \qcmkey{66} & \qcmkey{67} & \qcmkey{68} & \qcmkey{69} & \qcmkey{70} \\
71--80 & \qcmkey{71} & \qcmkey{72} & \qcmkey{73} & \qcmkey{74} & \qcmkey{75} & \qcmkey{76} & \qcmkey{77} & \qcmkey{78} & \qcmkey{79} & \qcmkey{80} \\
81--90 & \qcmkey{81} & \qcmkey{82} & \qcmkey{83} & \qcmkey{84} & \qcmkey{85} & \qcmkey{86} & \qcmkey{87} & \qcmkey{88} & \qcmkey{89} & \qcmkey{90} \\
91--100 & \qcmkey{91} & \qcmkey{92} & \qcmkey{93} & \qcmkey{94} & \qcmkey{95} & \qcmkey{96} & \qcmkey{97} & \qcmkey{98} & \qcmkey{99} & \qcmkey{100} \\
\bottomrule
\end{tabular}
\end{center}

\paragraph{Fiche de score.} Reportez le nombre de bonnes réponses de chaque niveau.
\begin{center}
\renewcommand{\arraystretch}{1.6}
\begin{tabular}{@{}l l c c c@{}}
\toprule
\textbf{Niveau} & \textbf{Questions} & \textbf{Bonnes réponses} & \textbf{Sur} & \textbf{Pourcentage}\\
\midrule
1. Débutant & 1--30 & \makebox[2.6cm]{\dotfill} & 30 & \makebox[2cm]{\dotfill}\,\%\\
2. Intermédiaire & 31--60 & \makebox[2.6cm]{\dotfill} & 30 & \makebox[2cm]{\dotfill}\,\%\\
3. Avancé & 61--90 & \makebox[2.6cm]{\dotfill} & 30 & \makebox[2cm]{\dotfill}\,\%\\
4. Maîtrise & 91--100 & \makebox[2.6cm]{\dotfill} & 10 & \makebox[2cm]{\dotfill}\,\%\\
\midrule
\textbf{Total} & 1--100 & \makebox[2.6cm]{\dotfill} & 100 & \makebox[2cm]{\dotfill}\,\%\\
\bottomrule
\end{tabular}
\end{center}
Repères : au moins $80\,\%$ à un niveau, il est acquis et on peut passer au suivant ; entre $50$ et
$80\,\%$, revoyez les sections des questions manquées ; en dessous, relisez le chapitre concerné et
son récapitulatif avant de recommencer.

\paragraph{Où revoir ?} Questions classées par chapitre du cours.
\begin{center}
\renewcommand{\arraystretch}{1.3}
\begin{tabular}{@{}l >{\raggedright\arraybackslash}p{5.9cm} >{\raggedright\arraybackslash}p{6.9cm}@{}}
\toprule
\textbf{Chapitre} & \textbf{Titre} & \textbf{Questions}\\
\midrule
Chap.~\ref{chap:histoire} & Histoire et mécanique quantique & 1, 2, 3, 4, 5, 6, 31, 32, 33, 34, 35, 61, 62, 63, 91 \\
Chap.~\ref{chap:dirac} & Notation de Dirac et espace de Hilbert & 7, 8, 9, 10, 11, 12, 36, 37, 38, 39, 40, 64, 65, 66, 67, 68, 92 \\
Chap.~\ref{chap:operateurs} & Opérateurs, matrice de densité, Bloch, produit tensoriel & 13, 14, 15, 16, 17, 18, 41, 42, 43, 44, 45, 46, 69, 70, 71, 72, 73, 74, 75, 93, 94 \\
Chap.~\ref{chap:portes} & Portes quantiques et systèmes multi-qubits & 19, 20, 21, 22, 23, 24, 47, 48, 49, 50, 51, 76, 77, 78, 79, 80, 81, 82, 95, 96 \\
Chap.~\ref{chap:bell} & États quantiques, intrication, Bell & 25, 26, 27, 28, 29, 30, 52, 53, 54, 55, 56, 57, 58, 59, 60, 83, 84, 85, 86, 87, 88, 89, 90, 97, 98, 99, 100 \\
\bottomrule
\end{tabular}
\end{center}

\clearpage
\subsection{Corrigé du niveau 1 : débutant (questions 1 à 30)}\label{sec:qcm-c1}
\qcmcorr{1}
\qcmcorr{2}
\qcmcorr{3}
\qcmcorr{4}
\qcmcorr{5}
\qcmcorr{6}
\qcmcorr{7}
\qcmcorr{8}
\qcmcorr{9}
\qcmcorr{10}
\qcmcorr{11}
\qcmcorr{12}
\qcmcorr{13}
\qcmcorr{14}
\qcmcorr{15}
\qcmcorr{16}
\qcmcorr{17}
\qcmcorr{18}
\qcmcorr{19}
\qcmcorr{20}
\qcmcorr{21}
\qcmcorr{22}
\qcmcorr{23}
\qcmcorr{24}
\qcmcorr{25}
\qcmcorr{26}
\qcmcorr{27}
\qcmcorr{28}
\qcmcorr{29}
\qcmcorr{30}

\clearpage
\subsection{Corrigé du niveau 2 : intermédiaire (questions 31 à 60)}\label{sec:qcm-c2}
\qcmcorr{31}
\qcmcorr{32}
\qcmcorr{33}
\qcmcorr{34}
\qcmcorr{35}
\qcmcorr{36}
\qcmcorr{37}
\qcmcorr{38}
\qcmcorr{39}
\qcmcorr{40}
\qcmcorr{41}
\qcmcorr{42}
\qcmcorr{43}
\qcmcorr{44}
\qcmcorr{45}
\qcmcorr{46}
\qcmcorr{47}
\qcmcorr{48}
\qcmcorr{49}
\qcmcorr{50}
\qcmcorr{51}
\qcmcorr{52}
\qcmcorr{53}
\qcmcorr{54}
\qcmcorr{55}
\qcmcorr{56}
\qcmcorr{57}
\qcmcorr{58}
\qcmcorr{59}
\qcmcorr{60}

\clearpage
\subsection{Corrigé du niveau 3 : avancé (questions 61 à 90)}\label{sec:qcm-c3}
\qcmcorr{61}
\qcmcorr{62}
\qcmcorr{63}
\qcmcorr{64}
\qcmcorr{65}
\qcmcorr{66}
\qcmcorr{67}
\qcmcorr{68}
\qcmcorr{69}
\qcmcorr{70}
\qcmcorr{71}
\qcmcorr{72}
\qcmcorr{73}
\qcmcorr{74}
\qcmcorr{75}
\qcmcorr{76}
\qcmcorr{77}
\qcmcorr{78}
\qcmcorr{79}
\qcmcorr{80}
\qcmcorr{81}
\qcmcorr{82}
\qcmcorr{83}
\qcmcorr{84}
\qcmcorr{85}
\qcmcorr{86}
\qcmcorr{87}
\qcmcorr{88}
\qcmcorr{89}
\qcmcorr{90}

\clearpage
\subsection{Corrigé du niveau 4 : maîtrise (questions 91 à 100)}\label{sec:qcm-c4}
\qcmcorr{91}
\qcmcorr{92}
\qcmcorr{93}
\qcmcorr{94}
\qcmcorr{95}
\qcmcorr{96}
\qcmcorr{97}
\qcmcorr{98}
\qcmcorr{99}
\qcmcorr{100}
\fi
) :
%    \qcmrefsfalse     masque la référence « Chap. · § » sous chaque énoncé ;
%    \qcmcorrigefalse  supprime le corrigé du document (QCM à distribuer seul).
%
%  Ce fichier est produit par outils_qcm/gen_qcm.py à partir de la banque de
%  questions (fichiers texte), mais il reste un fichier LaTeX ordinaire que l'on
%  peut modifier à la main. Pour ajouter une question à la main : insérer un
%  couple \qcmq ... \qcmx dans le bon niveau (la numérotation est automatique),
%  puis mettre à jour (1) les bornes « questions a à b » des titres de niveaux,
%  (2) les lignes \qcmcorr{n} du corrigé, (3) le tableau « Où revoir ? ».
% =====================================================================

% Macros de Dirac (déjà définies dans main.tex ; ce bloc évite l'erreur
% « Undefined control sequence \ket » si main.tex est une ancienne version).
\providecommand{\ket}[1]{|#1\rangle}
\providecommand{\bra}[1]{\langle #1|}
\providecommand{\braket}[2]{\langle #1|#2\rangle}

\makeatletter
% bascules : normalement déjà définies dans main.tex (valeurs par défaut : vrai)
\@ifundefined{ifqcmrefs}{\newif\ifqcmrefs\qcmrefstrue}{}
\@ifundefined{ifqcmcorrige}{\newif\ifqcmcorrige\qcmcorrigetrue}{}
\newcounter{qcmq}

% case à cocher (3 mm) pour le crayon
\newcommand{\qcmbox}{{\setlength{\fboxsep}{0pt}\setlength{\fboxrule}{0.5pt}%
  \raisebox{-0.4mm}{\fbox{\rule{0pt}{3mm}\rule{3mm}{0pt}}}}}

% une proposition : case, lettre, texte (retrait suspendu)
\newcommand{\qcm@opt}[2]{\noindent\hangindent=3.2em \hangafter=1
  \hspace*{0.3em}\qcmbox\hspace{0.45em}\textbf{#1)}\hspace{0.4em}#2\par\vspace{2.5pt}}

% \qcmq[colonnes]{ref}{énoncé}{A}{B}{C}{D}{bonne lettre}   (TeX limite les macros à 9 paramètres)
% \qcmx{com. A}{com. B}{com. C}{com. D}  : commentaires, placés juste après \qcmq
%   - com. de la bonne lettre = explication ; com. des autres = piège ;
%   - colonnes = 1 (une proposition par ligne) ou 2 (grille 2x2).
\newcommand{\qcmq}[8][1]{%
  \stepcounter{qcmq}%
  \expandafter\gdef\csname qcm@ref@\theqcmq\endcsname{#2}%
  \expandafter\gdef\csname qcm@rep@\theqcmq\endcsname{#8}%
  \expandafter\gdef\csname qcm@o@\theqcmq @A\endcsname{#4}%
  \expandafter\gdef\csname qcm@o@\theqcmq @B\endcsname{#5}%
  \expandafter\gdef\csname qcm@o@\theqcmq @C\endcsname{#6}%
  \expandafter\gdef\csname qcm@o@\theqcmq @D\endcsname{#7}%
  \par\vspace{7pt plus 3pt minus 2pt}\noindent
  \begin{minipage}{\linewidth}%
    \noindent{\color{insarouge}\bfseries Question~\theqcmq}%
    \ifqcmrefs\hfill{\footnotesize\color{insagris}#2}\fi\par\nopagebreak
    \vspace{2pt}\noindent #3\par\nopagebreak\vspace{3pt}%
    \ifnum#1=2
      \begin{minipage}[t]{0.495\linewidth}\qcm@opt{A}{#4}\end{minipage}\hfill
      \begin{minipage}[t]{0.495\linewidth}\qcm@opt{B}{#5}\end{minipage}\par\vspace{2pt}
      \begin{minipage}[t]{0.495\linewidth}\qcm@opt{C}{#6}\end{minipage}\hfill
      \begin{minipage}[t]{0.495\linewidth}\qcm@opt{D}{#7}\end{minipage}\par
    \else
      \qcm@opt{A}{#4}\qcm@opt{B}{#5}\qcm@opt{C}{#6}\qcm@opt{D}{#7}%
    \fi
  \end{minipage}\par}

\newcommand{\qcmx}[4]{%
  \expandafter\gdef\csname qcm@c@\theqcmq @A\endcsname{#1}%
  \expandafter\gdef\csname qcm@c@\theqcmq @B\endcsname{#2}%
  \expandafter\gdef\csname qcm@c@\theqcmq @C\endcsname{#3}%
  \expandafter\gdef\csname qcm@c@\theqcmq @D\endcsname{#4}}

% lettre de la bonne réponse de la question n (« ? » si la question n'existe pas)
\newcommand{\qcmkey}[1]{\@ifundefined{qcm@rep@#1}{?}{\csname qcm@rep@#1\endcsname}}

% une mauvaise réponse du corrigé (n° de question, lettre) : ignorée si c'est la bonne.
% (pas de \@for : « := » ne fonctionne pas avec le « : » rendu actif par babel-french)
\newcommand{\qcm@wrongone}[2]{%
  \edef\qcm@a{#2}\edef\qcm@b{\csname qcm@rep@#1\endcsname}%
  \ifx\qcm@a\qcm@b\else
    \noindent\hangindent=1.4em\hangafter=1 {\color{insagris}\textbf{#2)}}\enspace
    \csname qcm@o@#1@#2\endcsname\enspace--\enspace\emph{\csname qcm@c@#1@#2\endcsname}\par
  \fi}

% bloc de corrigé de la question n
\newcommand{\qcmcorr}[1]{%
  \par\vspace{6pt plus 2pt minus 2pt}\noindent
  \begin{minipage}{\linewidth}\small
    \noindent{\color{insarouge}\bfseries Question~#1}\enspace
    \textbf{Réponse : \csname qcm@rep@#1\endcsname}\hfill
    {\footnotesize\color{insagris}\csname qcm@ref@#1\endcsname}\par\nopagebreak
    \vspace{1pt}\noindent\hangindent=1.4em\hangafter=1
    \textbf{\csname qcm@rep@#1\endcsname)}\enspace
    \csname qcm@o@#1@\csname qcm@rep@#1\endcsname\endcsname\enspace---\enspace
    \csname qcm@c@#1@\csname qcm@rep@#1\endcsname\endcsname\par
    \qcm@wrongone{#1}{A}\qcm@wrongone{#1}{B}\qcm@wrongone{#1}{C}\qcm@wrongone{#1}{D}%
  \end{minipage}\par}
\makeatother

% =====================================================================
\section{QCM sur le cours}\label{chap:qcm}
% =====================================================================

\subsection{Mode d'emploi}\label{sec:qcm-mode}

Ce QCM compte \textbf{100 questions} sur les chapitres~\ref{chap:histoire} à~\ref{chap:bell},
classées en quatre niveaux. Chaque question propose quatre réponses dont \textbf{une seule} est
exacte : cochez au crayon la case de la proposition choisie.
\ifqcmrefs À droite du numéro de chaque question, une référence (chapitre, section) indique la
partie du cours concernée.\else\ifqcmcorrige{} La référence de chaque question (chapitre, section)
est donnée dans le corrigé.\fi\fi

\begin{center}
\renewcommand{\arraystretch}{1.3}
\begin{tabular}{@{}llc>{\raggedright\arraybackslash}p{8.2cm}@{}}
\toprule
\textbf{Niveau} & \textbf{Questions} & \textbf{Nombre} & \textbf{Ce qui est demandé}\\
\midrule
1. Débutant & 1 à 30 & 30 & définitions, notations, résultats du cours en une étape\\
2. Intermédiaire & 31 à 60 & 30 & petits calculs, propriétés, application directe d'un résultat\\
3. Avancé & 61 à 90 & 30 & calculs en plusieurs étapes, portes, circuits et protocoles\\
4. Maîtrise & 91 à 100 & 10 & raisonnements fins et pièges conceptuels, qui combinent plusieurs notions\\
\bottomrule
\end{tabular}
\end{center}

\begin{itemize}[nosep]
  \item \textbf{Sur papier} : imprimez seulement les pages~\pageref{chap:qcm} à~\pageref{qcm:fin}
        (mode d'emploi et énoncés, sections~\ref{sec:qcm-n1} à~\ref{sec:qcm-n4}).%
        \ifqcmcorrige{} Le corrigé commence sur une page distincte, à la page~\pageref{sec:qcm-cle}.\fi
  \item \textbf{Barème} : 1 point par bonne réponse, 0 sinon (pas de point négatif). Calculatrice
        autorisée, mais rarement utile (quelques logarithmes et racines).
  \item \textbf{Correction} :%
        \ifqcmcorrige{} comparez avec la clé des réponses, puis lisez l'explication de chaque
        question, même réussie. Les « pièges » décrivent l'erreur qui mène à chaque mauvaise réponse.%
        \else{} le corrigé n'est pas inclus dans cette version du document.\fi
\end{itemize}

\bigskip
\noindent Nom : \makebox[6cm]{\dotfill}\hfill Date : \makebox[3.5cm]{\dotfill}\hfill Score : \makebox[1.6cm]{\dotfill}\,/\,100
\bigskip

\begin{remarque}[Conventions des énoncés]
\begin{itemize}[nosep]
  \item $\ket{ab}=\ket{a}_A\otimes\ket{b}_B$ : le premier chiffre est le qubit $A$, le second le
        qubit $B$ ; dans un circuit, le fil du haut est le dernier chiffre.
  \item Le CNOT de la section~\ref{sec:cnot} a pour contrôle le \emph{second} chiffre ; les portes
        contrôlées de la section~\ref{sec:controlees} ont pour contrôle le \emph{premier}. Chaque
        énoncé précise le contrôle.
  \item Sans précision, une mesure est faite dans la base de calcul $Z$. Les logarithmes des
        entropies sont en base $2$ (résultats en bits).
  \item $\ket{\pm}=\frac{1}{\sqrt2}(\ket0\pm\ket1)$ et $\ket{{\pm}\ii}=\frac{1}{\sqrt2}(\ket0\pm\ii\ket1)$.
\end{itemize}
\end{remarque}

\clearpage
%% ======================================================================
\subsection{Niveau 1 : débutant (questions 1 à 30)}\label{sec:qcm-n1}
%% ======================================================================
\noindent{\color{insagris}\small Niveau de difficulté : Débutant\quad $\bullet\circ\circ\circ$}\par\smallskip
% ---- Question 1 (niveau 1, correct : D)
\qcmq[2]{Chap.~\ref{chap:histoire}\,$\cdot$\,\S\,\ref{sec:dualite}}
  {En 1905, qui explique l'effet photo-électrique en postulant que la lumière est faite de quanta d'énergie $E=hf$, les photons ?}
  {Louis de Broglie}
  {Erwin Schrödinger}
  {Max Planck}
  {Albert Einstein}
  {D}
\qcmx
  {Il propose en 1924 une onde de matière ($\lambda=h/p$) pour les particules comme l'électron ; l'idée du photon date de 1905.}
  {Il écrit en 1926 l'équation d'onde de la matière, plus de vingt ans après l'article sur les photons.}
  {Il introduit le quantum d'énergie en 1900 pour le rayonnement du corps noir ; c'est Einstein qui en fait une propriété de la lumière elle-même.}
  {Il postule que la lumière échange son énergie par paquets $E=hf=\hbar\omega$ : c'est la première face de la dualité onde--corpuscule.}

% ---- Question 2 (niveau 1, correct : A)
\qcmq[1]{Chap.~\ref{chap:histoire}\,$\cdot$\,\S\,\ref{sec:born}}
  {Selon l'interprétation de Max Born, que représente $|\Psi(x,t)|^2$ ?}
  {La densité de probabilité de présence de la particule en $x$ à l'instant $t$}
  {L'énergie cinétique de la particule}
  {L'amplitude de l'onde de matière}
  {La vitesse de la particule}
  {A}
\qcmx
  {$|\Psi(x,t)|^2\,\dd x$ est la probabilité de trouver la particule entre $x$ et $x+\dd x$ ; $\Psi$ lui-même n'est pas mesurable.}
  {L'énergie est associée à l'opérateur hamiltonien $\Hop$ ; $|\Psi|^2$ n'a pas la dimension d'une énergie.}
  {$\Psi$ est l'amplitude (complexe) ; $|\Psi|^2$ en est le carré du module.}
  {La vitesse s'obtient avec l'opérateur $\hat p=-\ii\hbar\,\partial_x$, pas avec $|\Psi|^2$.}

% ---- Question 3 (niveau 1, correct : D)
\qcmq[2]{Chap.~\ref{chap:histoire}\,$\cdot$\,\S\,\ref{sec:dualite}}
  {Une particule de quantité de mouvement $p$ est associée, d'après de Broglie, à une onde de longueur d'onde $\lambda$. Quelle relation est correcte ?}
  {$\lambda=hp$}
  {$\lambda=p/h$}
  {$\lambda=\hbar/p$}
  {$\lambda=h/p$}
  {D}
\qcmx
  {Produit au lieu du quotient : $\lambda$ grandirait avec $p$, le contraire de ce qu'on observe.}
  {Quotient inversé : $p/h=1/\lambda$ s'exprime en m$^{-1}$, pas en mètres.}
  {Confusion entre $h$ et $\hbar=h/2\pi$ : $\hbar/p=\lambda/2\pi$ est l'inverse du nombre d'onde $k$.}
  {Relation de de Broglie : plus la particule est rapide ou lourde, plus $\lambda$ est petite (balle de $0{,}1$~kg à $10$~m/s : $\lambda\approx10^{-34}$~m, inobservable).}

% ---- Question 4 (niveau 1, correct : C)
\qcmq[2]{Chap.~\ref{chap:histoire}\,$\cdot$\,\S\,\ref{sec:construction}}
  {Quelle est l'équation de Schrödinger dépendante du temps pour un état $\Psi$ et un hamiltonien $\Hop$ ?}
  {$-\ii\hbar\,\partial_t\Psi=\Hop\Psi$}
  {$\hbar\,\partial_t\Psi=\Hop\Psi$}
  {$\ii\hbar\,\partial_t\Psi=\Hop\Psi$}
  {$\ii\hbar\,\partial_x\Psi=\Hop\Psi$}
  {C}
\qcmx
  {Erreur de signe : la solution serait $\ee^{+\ii\Hop t/\hbar}\Psi(0)$, qui décrit l'évolution à rebours du temps.}
  {Sans le facteur $\ii$, les solutions sont des exponentielles réelles (croissantes ou décroissantes) et la norme n'est pas conservée.}
  {Elle relie l'évolution de l'état à l'énergie ; pour un $\Hop$ indépendant du temps, sa solution est $\Psi(t)=\ee^{-\ii\Hop t/\hbar}\Psi(0)$ (chapitre~\ref{chap:portes}).}
  {C'est la dérivée en $t$ qui décrit l'évolution ; la dérivée en $x$ sert à construire $\hat p=-\ii\hbar\,\partial_x$.}

% ---- Question 5 (niveau 1, correct : A)
\qcmq[1]{Chap.~\ref{chap:histoire}\,$\cdot$\,\S\,\ref{sec:construction}}
  {Quel opérateur représente l'énergie totale d'un système quantique ?}
  {L'hamiltonien $\Hop$}
  {L'opérateur quantité de mouvement $\hat p$}
  {L'opérateur position $\hat x$}
  {L'opérateur identité}
  {A}
\qcmx
  {Somme de l'énergie cinétique $-\frac{\hbar^2}{2m}\partial_x^2$ et de l'énergie potentielle $V$ ; ses valeurs propres sont les énergies mesurables.}
  {Il représente l'impulsion ; l'énergie cinétique est $\hat p^2/2m$, et il faut y ajouter $V$.}
  {C'est la multiplication par $x$ : il donne la position, pas l'énergie.}
  {Il vaut $1$ pour tout état : il ne contient aucune information sur l'énergie.}

% ---- Question 6 (niveau 1, correct : C)
\qcmq[1]{Chap.~\ref{chap:histoire}\,$\cdot$\,\S\,\ref{sec:histoire}}
  {Dans quel ordre chronologique ces avancées ont-elles eu lieu ?}
  {de Broglie $\to$ Planck $\to$ Schrödinger $\to$ Bell}
  {Planck $\to$ Schrödinger $\to$ de Broglie $\to$ Bell}
  {Planck (quantum d'énergie) $\to$ de Broglie (onde de matière) $\to$ Schrödinger (équation d'onde) $\to$ Bell (inégalités)}
  {Planck $\to$ de Broglie $\to$ Bell $\to$ Schrödinger}
  {C}
\qcmx
  {Le quantum d'énergie de Planck (1900) précède l'onde de matière de de Broglie (1924).}
  {L'équation de Schrödinger (1926) généralise l'idée d'onde de matière de de Broglie (1924) : elle vient après.}
  {Planck 1900, de Broglie 1924, Schrödinger 1926, Bell 1964.}
  {Les inégalités de Bell (1964) viennent près de quarante ans après l'équation de Schrödinger.}

% ---- Question 7 (niveau 1, correct : C)
\qcmq[1]{Chap.~\ref{chap:dirac}\,$\cdot$\,\S\,\ref{sec:dirac}}
  {En notation de Dirac, comment appelle-t-on $\ket{\Psi}$ et comment le représente-t-on pour un qubit ?}
  {Un bra, représenté par un vecteur ligne à deux composantes complexes}
  {Un ket, représenté par un nombre complexe}
  {Un ket, représenté par un vecteur colonne à deux composantes complexes}
  {Un bra, représenté par une matrice $2\times2$}
  {C}
\qcmx
  {Le bra $\bra{\Psi}$ est le conjugué transposé du ket : un vecteur ligne, mais ce n'est pas la notation $\ket{\Psi}$.}
  {Un scalaire ne peut pas décrire une superposition : il faut deux amplitudes $\alpha$ et $\beta$.}
  {Le ket est le vecteur colonne $\binom{\alpha}{\beta}$ ; son bra $\bra{\Psi}$ est le vecteur ligne conjugué.}
  {Une matrice est un opérateur ; le bra est une matrice ligne $1\times2$.}

% ---- Question 8 (niveau 1, correct : B)
\qcmq[2]{Chap.~\ref{chap:dirac}\,$\cdot$\,\S\,\ref{sec:qubit}}
  {Quelle condition les amplitudes $\alpha$ et $\beta$ de $\alpha\ket{0}+\beta\ket{1}$ doivent-elles vérifier pour que l'état soit normalisé ?}
  {$\alpha+\beta=1$}
  {$|\alpha|^2+|\beta|^2=1$}
  {$\alpha^2+\beta^2=1$}
  {$|\alpha|+|\beta|=1$}
  {B}
\qcmx
  {On somme les amplitudes (complexes, parfois négatives) ; ce sont leurs modules au carré qui sont des probabilités.}
  {Les probabilités de mesure $P(0)=|\alpha|^2$ et $P(1)=|\beta|^2$ doivent sommer à $1$.}
  {Oubli du module : pour $\alpha=\frac{1}{\sqrt2}$ et $\beta=\frac{\ii}{\sqrt2}$, $\alpha^2+\beta^2=0$ alors que l'état est normalisé.}
  {Somme des modules au lieu des carrés : $\alpha=\beta=\frac{1}{\sqrt2}$ est normalisé, mais $|\alpha|+|\beta|=\sqrt2$.}

% ---- Question 9 (niveau 1, correct : C)
\qcmq[1]{Chap.~\ref{chap:dirac}\,$\cdot$\,\S\,\ref{sec:projecteurs}}
  {On mesure dans la base $Z$ un qubit préparé dans $\ket{+}=\frac{1}{\sqrt2}(\ket{0}+\ket{1})$. Que se passe-t-il ?}
  {On lit $+$ et le qubit reste dans $\ket{+}$}
  {On lit toujours $0$}
  {On lit $0$ ou $1$, chacun avec la probabilité $\frac12$, et l'état devient $\ket{0}$ ou $\ket{1}$}
  {On lit $0$ et $1$ à la fois}
  {C}
\qcmx
  {Une mesure en base $Z$ ne donne que $0$ ou $1$ ; elle modifie l'état.}
  {Les deux amplitudes ont le même module : $P(0)=P(1)=\frac12$.}
  {Postulat de la mesure : la superposition s'effondre sur l'état propre correspondant au résultat.}
  {Chaque mesure donne un seul résultat ; la superposition ne s'observe qu'à travers les statistiques.}

% ---- Question 10 (niveau 1, correct : B)
\qcmq[2]{Chap.~\ref{chap:dirac}\,$\cdot$\,\S\,\ref{sec:qubit}}
  {Que vaut le produit scalaire $\braket{0}{1}$ ?}
  {$1$}
  {$0$}
  {$\frac{1}{\sqrt2}$}
  {$\ket{01}$}
  {B}
\qcmx
  {C'est la valeur de $\braket{0}{0}$ et de $\braket{1}{1}$ (vecteurs normés), pas de $\braket{0}{1}$.}
  {$\ket{0}$ et $\ket{1}$ sont orthogonaux : $\begin{pmatrix}1&0\end{pmatrix}\begin{pmatrix}0\\1\end{pmatrix}=0$.}
  {C'est la valeur de $\braket{0}{+}$, pas de $\braket{0}{1}$.}
  {Confusion avec le produit tensoriel : un produit scalaire est un nombre, pas un ket.}

% ---- Question 11 (niveau 1, correct : D)
\qcmq[2]{Chap.~\ref{chap:dirac}\,$\cdot$\,\S\,\ref{sec:qubit}}
  {Un qubit est dans l'état $\ket{-}=\frac{1}{\sqrt2}(\ket{0}-\ket{1})$. Quelle est la probabilité de mesurer $1$ en base $Z$ ?}
  {$0$}
  {$-\frac12$}
  {$\frac{1}{\sqrt2}$}
  {$\frac12$}
  {D}
\qcmx
  {Le signe $-$ ne « retire » pas $\ket{1}$ de la superposition.}
  {Une probabilité est toujours positive : on prend le module au carré de l'amplitude.}
  {C'est l'amplitude en module, pas la probabilité (il faut élever au carré).}
  {$P(1)=\left|-\frac{1}{\sqrt2}\right|^2=\frac12$ : le signe est une phase relative, qui disparaît dans le module.}

% ---- Question 12 (niveau 1, correct : B)
\qcmq[2]{Chap.~\ref{chap:dirac}\,$\cdot$\,\S\,\ref{sec:projecteurs}}
  {Quel est le projecteur sur l'état $\ket{0}$ ?}
  {$\ket{0}\bra{1}=\begin{pmatrix}0&1\\0&0\end{pmatrix}$}
  {$\hat P_0=\ket{0}\bra{0}=\begin{pmatrix}1&0\\0&0\end{pmatrix}$}
  {$\braket{0}{0}=1$}
  {$\begin{pmatrix}0&0\\0&1\end{pmatrix}$}
  {B}
\qcmx
  {Cet opérateur envoie $\ket{1}$ sur $\ket{0}$ ; son carré est nul : ce n'est pas un projecteur.}
  {$\hat P_0\ket{\Psi}=\alpha\ket{0}$ : il garde la composante sur $\ket{0}$ ; $\hat P_0^2=\hat P_0$.}
  {C'est un nombre, pas un opérateur.}
  {C'est $\ket{1}\bra{1}$, le projecteur sur $\ket{1}$.}

% ---- Question 13 (niveau 1, correct : A)
\qcmq[1]{Chap.~\ref{chap:operateurs}\,$\cdot$\,\S\,\ref{sec:bloch}}
  {Où se trouvent les états $\ket{0}$ et $\ket{1}$ sur la sphère de Bloch ?}
  {Aux deux pôles, Nord et Sud}
  {Aux deux extrémités de l'axe $x$}
  {Aux deux extrémités de l'axe $y$}
  {En deux points voisins de l'équateur}
  {A}
\qcmx
  {$\ket{0}$ au pôle Nord ($\theta=0$), $\ket{1}$ au pôle Sud ($\theta=\pi$) : deux états orthogonaux sont diamétralement opposés.}
  {Ce sont les places de $\ket{+}$ et $\ket{-}$.}
  {Ce sont les places de $\ket{{+}\ii}$ et $\ket{{-}\ii}$.}
  {Des états orthogonaux sont antipodaux, jamais voisins.}

% ---- Question 14 (niveau 1, correct : A)
\qcmq[2]{Chap.~\ref{chap:operateurs}\,$\cdot$\,\S\,\ref{sec:tenseur}}
  {Quelle est la dimension de l'espace de Hilbert de $n$ qubits ?}
  {$2^n$}
  {$2n$}
  {$n^2$}
  {$2^n-1$}
  {A}
\qcmx
  {Le produit tensoriel multiplie les dimensions : $2\times2\times\cdots\times2$ ; deux qubits : $4$, trois qubits : $8$.}
  {On additionne les dimensions au lieu de les multiplier ; pour $n=3$ on trouverait $6$ au lieu de $8$.}
  {Juste pour $n=2$ par coïncidence seulement : pour $n=3$ on trouverait $9$ au lieu de $8$.}
  {Il y a bien $2^n$ vecteurs dans la base de calcul, $\ket{0\cdots0}$ compris.}

% ---- Question 15 (niveau 1, correct : C)
\qcmq[1]{Chap.~\ref{chap:operateurs}\,$\cdot$\,\S\,\ref{sec:rho}}
  {Que vaut toujours la trace d'une matrice de densité $\rho$ ?}
  {$0$}
  {$1$ pour un état pur, et moins de $1$ pour un état mixte}
  {$1$}
  {$2$, la dimension de l'espace}
  {C}
\qcmx
  {C'est la trace des matrices de Pauli, pas d'une matrice de densité.}
  {Confusion avec $\operatorname{Tr}(\rho^2)$ : c'est la pureté qui distingue les deux cas, la trace vaut toujours $1$.}
  {C'est la somme des probabilités de mesure (les populations), que l'état soit pur ou mixte.}
  {C'est la trace de l'identité $\operatorname{Tr}\mathrm{Id}=2$, pas celle de $\rho$.}

% ---- Question 16 (niveau 1, correct : B)
\qcmq[1]{Chap.~\ref{chap:operateurs}\,$\cdot$\,\S\,\ref{sec:pur-mixte}}
  {Quelle est la valeur de $\operatorname{Tr}(\rho^2)$ pour un état pur ?}
  {$0$}
  {$1$}
  {$\frac12$}
  {Une valeur strictement inférieure à $1$}
  {B}
\qcmx
  {$\rho^2=0$ imposerait $\rho=0$, ce qui contredit $\operatorname{Tr}\rho=1$.}
  {Un état pur s'écrit $\rho=\ket{\Psi}\bra{\Psi}$, donc $\rho^2=\rho$ et $\operatorname{Tr}\rho^2=\operatorname{Tr}\rho=1$.}
  {C'est la valeur minimale pour un qubit, atteinte par l'état maximalement mixte $\mathrm{Id}/2$.}
  {C'est la caractérisation d'un état mixte.}

% ---- Question 17 (niveau 1, correct : B)
\qcmq[1]{Chap.~\ref{chap:operateurs}\,$\cdot$\,\S\,\ref{sec:rho}}
  {Que représentent les éléments diagonaux $\rho_{00}$ et $\rho_{11}$ de la matrice de densité d'un qubit ?}
  {Les cohérences entre $\ket{0}$ et $\ket{1}$}
  {Les probabilités de mesurer $\ket{0}$ et $\ket{1}$ (les populations)}
  {Les valeurs propres de $\rho$}
  {Les amplitudes $\alpha$ et $\beta$}
  {B}
\qcmx
  {Les cohérences sont les éléments hors diagonale $\rho_{01}$ et $\rho_{10}$.}
  {$\rho_{00}=\bra{0}\rho\ket{0}=P(0)$ et $\rho_{11}=P(1)$, d'où $\operatorname{Tr}\rho=1$.}
  {Seulement si $\rho$ est diagonale dans la base de calcul : $\ket{+}\bra{+}$ a pour diagonale $(\frac12,\frac12)$ mais pour valeurs propres $(1,0)$.}
  {Pour un état pur, $\rho_{00}=|\alpha|^2$ : le module au carré, pas l'amplitude.}

% ---- Question 18 (niveau 1, correct : D)
\qcmq[1]{Chap.~\ref{chap:operateurs}\,$\cdot$\,\S\,\ref{sec:op-dirac}}
  {Pourquoi associe-t-on aux grandeurs mesurables (observables) des opérateurs hermitiens ?}
  {Parce qu'ils conservent la norme des vecteurs}
  {Parce qu'ils ont tous une trace nulle}
  {Parce que leurs valeurs propres valent toutes $\pm1$}
  {Parce que leurs valeurs propres sont réelles, comme les résultats d'une mesure}
  {D}
\qcmx
  {C'est la propriété des opérateurs unitaires (les portes), pas des opérateurs hermitiens.}
  {Vrai pour les matrices de Pauli, faux en général (l'identité est hermitienne, de trace $2$), et sans rapport avec le caractère mesurable.}
  {Vrai pour $X$, $Y$, $Z$ seulement ; l'hamiltonien du puits a les valeurs propres $E_n=n^2E_1$.}
  {Un opérateur hermitien ($A^\dagger=A$) a des valeurs propres réelles et des vecteurs propres orthogonaux ; les résultats possibles sont ses valeurs propres.}

% ---- Question 19 (niveau 1, correct : C)
\qcmq[2]{Chap.~\ref{chap:portes}\,$\cdot$\,\S\,\ref{sec:pauli}}
  {Quelle porte de Pauli réalise l'inversion de bit $\ket{0}\leftrightarrow\ket{1}$ ?}
  {$Z$}
  {$Y$}
  {$X$}
  {$H$}
  {C}
\qcmx
  {$Z$ laisse $\ket{0}$ et change $\ket{1}$ en $-\ket{1}$ : c'est une inversion de phase.}
  {$Y\ket{0}=\ii\ket{1}$ : elle inverse le bit mais ajoute une phase ; la porte d'inversion « pure » est $X$.}
  {$X\ket{0}=\ket{1}$ et $X\ket{1}=\ket{0}$ : c'est le NOT quantique.}
  {$H$ crée des superpositions : $H\ket{0}=\ket{+}$.}

% ---- Question 20 (niveau 1, correct : B)
\qcmq[2]{Chap.~\ref{chap:portes}\,$\cdot$\,\S\,\ref{sec:hadamard}}
  {Que donne la porte de Hadamard appliquée à $\ket{0}$ ?}
  {$\ket{1}$}
  {$\ket{+}=\frac{1}{\sqrt2}(\ket{0}+\ket{1})$}
  {$\ket{-}=\frac{1}{\sqrt2}(\ket{0}-\ket{1})$}
  {$\ket{0}$}
  {B}
\qcmx
  {C'est l'action de $X$.}
  {Superposition à poids égaux : $P(0)=P(1)=\frac12$ ; c'est la première étape de la plupart des algorithmes.}
  {C'est $H\ket{1}$ ; le signe $-$ n'apparaît qu'à partir de $\ket{1}$.}
  {$H$ n'est pas l'identité ; elle ne laisse inchangés ni $\ket{0}$ ni $\ket{1}$ (elle envoie $\ket{+}$ sur $\ket{0}$).}

% ---- Question 21 (niveau 1, correct : D)
\qcmq[2]{Chap.~\ref{chap:portes}\,$\cdot$\,\S\,\ref{sec:portes}}
  {Quelle propriété doit vérifier la matrice $U$ d'une porte quantique ?}
  {Elle est hermitienne : $U^\dagger=U$}
  {Sa trace vaut $1$}
  {Ses coefficients sont réels}
  {Elle est unitaire : $U^\dagger U=\mathrm{Id}$}
  {D}
\qcmx
  {Vrai pour $X$, $Y$, $Z$, $H$ mais faux pour $S$ et $T$ : ce n'est pas la condition générale.}
  {C'est une propriété des matrices de densité ; $\operatorname{Tr}X=0$ et $\operatorname{Tr}\mathrm{Id}=2$.}
  {$Y$ et $S$ sont des portes à coefficients complexes.}
  {Elle conserve la norme ($\braket{\Psi}{\Psi}=1$ reste vrai) et elle est réversible.}

% ---- Question 22 (niveau 1, correct : A)
\qcmq[2]{Chap.~\ref{chap:portes}\,$\cdot$\,\S\,\ref{sec:portes}}
  {Quel est l'inverse d'une porte quantique $U$ ?}
  {$U^{-1}=U^\dagger$ (transposée conjuguée)}
  {$U^{-1}=U^{\mathsf T}$ (transposée seule)}
  {$U^{-1}=U^{*}$ (conjuguée seule)}
  {$U^{-1}=-U$}
  {A}
\qcmx
  {Conséquence de $U^\dagger U=\mathrm{Id}$ ; par exemple $S^{-1}=S^\dagger=\mathrm{diag}(1,-\ii)$.}
  {Oubli de la conjugaison : $S^{\mathsf T}=S=\mathrm{diag}(1,\ii)$ alors que $S^{-1}=\mathrm{diag}(1,-\ii)$.}
  {Faux en général : $Y^*=-Y$ alors que $Y^{-1}=Y$.}
  {Faux : $X^{-1}=X\neq-X$.}

% ---- Question 23 (niveau 1, correct : C)
\qcmq[1]{Chap.~\ref{chap:portes}\,$\cdot$\,\S\,\ref{sec:cnot}}
  {Dans une porte CNOT, que devient le qubit cible lorsque le qubit de contrôle est dans l'état $\ket{0}$ ?}
  {Il est inversé par la porte $X$}
  {Il est remis à $\ket{0}$}
  {Il reste inchangé}
  {Il subit un changement de signe, comme avec $Z$}
  {C}
\qcmx
  {C'est ce qui se passe quand le contrôle vaut $\ket{1}$.}
  {Le CNOT est réversible : remettre la cible à zéro détruirait de l'information.}
  {Avec le contrôle à $\ket{0}$, la porte agit comme l'identité sur la cible.}
  {Ce serait la porte $Z$ contrôlée (CZ), pas le CNOT.}

% ---- Question 24 (niveau 1, correct : A)
\qcmq[2]{Chap.~\ref{chap:portes}\,$\cdot$\,\S\,\ref{sec:cnot}}
  {Le CNOT du cours a pour contrôle le second chiffre du ket et pour cible le premier : $\mathrm{CNOT}\ket{a,b}=\ket{a\oplus b,\,b}$. Que donne-t-il sur $\ket{0,1}$ ?}
  {$\ket{1,1}$}
  {$\ket{0,1}$}
  {$\ket{0,0}$}
  {$\ket{1,0}$}
  {A}
\qcmx
  {Le contrôle $b=1$ inverse la cible : $0\oplus1=1$.}
  {Le contrôle vaut $1$ : la cible est bien inversée.}
  {C'est le contrôle qui serait inversé : le contrôle n'est jamais modifié.}
  {Les deux qubits seraient inversés : seule la cible change.}

% ---- Question 25 (niveau 1, correct : D)
\qcmq[1]{Chap.~\ref{chap:bell}\,$\cdot$\,\S\,\ref{sec:bell}}
  {Quelle propriété caractérise les quatre états de Bell ?}
  {Ils sont séparables : chaque qubit a son propre état}
  {Ils décrivent un seul qubit}
  {Ils ne diffèrent que par une phase globale}
  {Ils sont maximalement intriqués et forment une base orthonormée de l'espace de deux qubits}
  {D}
\qcmx
  {$D=c_{00}c_{11}-c_{01}c_{10}=\pm\frac12\neq0$ : ils ne se factorisent pas.}
  {Ce sont des états de deux qubits (quatre amplitudes sur $\ket{00},\ket{01},\ket{10},\ket{11}$).}
  {Ils sont orthogonaux deux à deux : ce sont quatre états distincts, alors qu'une phase globale ne change pas l'état.}
  {Aucun n'est un produit d'états à un qubit ($|D|=\frac12$, valeur maximale) et ils sont orthogonaux deux à deux.}

% ---- Question 26 (niveau 1, correct : D)
\qcmq[1]{Chap.~\ref{chap:bell}\,$\cdot$\,\S\,\ref{sec:teleportation}}
  {Quel est l'objectif du protocole de téléportation quantique ?}
  {Copier l'état d'un qubit inconnu sur le qubit de Bob}
  {Transmettre une information plus vite que la lumière}
  {Déplacer la particule physique d'Alice vers Bob}
  {Transférer l'état d'un qubit inconnu d'Alice à Bob, avec une paire intriquée et deux bits classiques}
  {D}
\qcmx
  {La copie est interdite (non-clonage) : après le protocole, l'état n'existe plus que chez Bob.}
  {Sans les deux bits classiques, le qubit de Bob est dans l'état $\mathrm{Id}/2$ : il n'a reçu aucune information.}
  {C'est l'état quantique qui est transféré, pas la particule qui le porte.}
  {Le qubit lui-même n'est pas envoyé : son état est reconstruit chez Bob, et détruit chez Alice par la mesure.}

% ---- Question 27 (niveau 1, correct : B)
\qcmq[1]{Chap.~\ref{chap:bell}\,$\cdot$\,\S\,\ref{sec:teleportation}}
  {Dans la téléportation, que doit envoyer Alice à Bob sur un canal classique pour qu'il retrouve l'état ?}
  {Un qubit intriqué}
  {Deux bits classiques, le résultat de sa mesure}
  {Les amplitudes exactes $\alpha$ et $\beta$ de l'état}
  {Un seul bit}
  {B}
\qcmx
  {La paire intriquée est partagée avant le protocole ; on ne l'envoie pas après.}
  {Les quatre résultats $(a,q)$ correspondent aux quatre corrections possibles ($I$, $Z$, $X$, $X$ puis $Z$) : il faut deux bits.}
  {Alice ne connaît pas l'état à téléporter, et ces nombres complexes n'ont pas de description finie en bits.}
  {Un bit ne distingue que deux cas, alors que Bob doit choisir parmi quatre corrections.}

% ---- Question 28 (niveau 1, correct : A)
\qcmq[1]{Chap.~\ref{chap:bell}\,$\cdot$\,\S\,\ref{sec:bell} ; ex.\,\ref{ex:superdense}}
  {Quel est le but du codage superdense ?}
  {Transmettre deux bits classiques en envoyant un seul qubit, grâce à une paire intriquée partagée}
  {Transmettre l'état d'un qubit avec deux bits classiques}
  {Compresser plusieurs qubits en un seul}
  {Chiffrer un message de façon inviolable}
  {A}
\qcmx
  {Alice applique l'une des quatre portes $I$, $Z$, $X$, $XZ$ à son qubit de la paire, l'envoie, et Bob identifie l'état de Bell obtenu.}
  {C'est la téléportation, le protocole inverse.}
  {Le protocole ne compresse pas de qubits : il fait passer deux bits classiques par un canal à un qubit.}
  {Ce n'est pas un protocole de cryptographie : la paire intriquée sert à doubler la capacité du canal.}

% ---- Question 29 (niveau 1, correct : A)
\qcmq[1]{Chap.~\ref{chap:bell}\,$\cdot$\,\S\,\ref{sec:teleportation}}
  {Que dit le théorème de non-clonage ?}
  {Aucune opération ne peut copier parfaitement un état quantique inconnu quelconque}
  {Il est impossible de mesurer un qubit sans le perturber}
  {Il est impossible de créer de l'intrication}
  {Il est impossible de copier un bit classique}
  {A}
\qcmx
  {Une porte $U$ qui copierait $\ket{0}$ et $\ket{1}$ enverrait $\ket{+}\ket{0}$ sur un état intriqué, par linéarité, et non sur $\ket{+}\ket{+}$.}
  {C'est l'effondrement de l'état, un autre principe.}
  {L'intrication se crée très bien (Hadamard puis CNOT).}
  {Le CNOT copie un bit classique : $\ket{0,b}\to\ket{b,b}$.}

% ---- Question 30 (niveau 1, correct : B)
\qcmq[1]{Chap.~\ref{chap:bell}\,$\cdot$\,\S\,\ref{sec:circuit-bell}}
  {Quel enchaînement de portes transforme $\ket{00}$ en l'état de Bell $\ket{\Phi^+}=\frac{1}{\sqrt2}(\ket{00}+\ket{11})$ ?}
  {CNOT, puis Hadamard sur le contrôle}
  {Hadamard sur l'un des qubits, puis CNOT dont le contrôle est ce même qubit}
  {Hadamard sur chacun des deux qubits}
  {SWAP, puis $X$ sur un qubit}
  {B}
\qcmx
  {Le CNOT ne fait rien sur $\ket{00}$ (contrôle à $0$) ; $H$ donne ensuite un état séparable.}
  {$H$ crée $\ket{0}\otimes\ket{+}$, puis le CNOT recopie la superposition : $\frac{1}{\sqrt2}(\ket{00}+\ket{11})$.}
  {On obtient $\ket{+}\otimes\ket{+}$, séparable : des portes locales ne créent pas d'intrication.}
  {$\ket{00}$ devient $\ket{01}$ ou $\ket{10}$, un état de base séparable.}

\clearpage
%% ======================================================================
\subsection{Niveau 2 : intermédiaire (questions 31 à 60)}\label{sec:qcm-n2}
%% ======================================================================
\noindent{\color{insagris}\small Niveau de difficulté : Intermédiaire\quad $\bullet\bullet\circ\circ$}\par\smallskip
% ---- Question 31 (niveau 2, correct : A)
\qcmq[1]{Chap.~\ref{chap:histoire}\,$\cdot$\,\S\,\ref{sec:born}}
  {Quelle condition la fonction d'onde d'une particule doit-elle vérifier sur tout l'espace ?}
  {$\displaystyle\int_{-\infty}^{+\infty}|\Psi(x,t)|^2\,\dd x=1$}
  {$\displaystyle\int_{-\infty}^{+\infty}\Psi(x,t)\,\dd x=1$}
  {$|\Psi(x,t)|^2=1$ pour tout $x$}
  {$\partial_x\Psi(x,t)=0$ en tout point}
  {A}
\qcmx
  {La particule se trouve forcément quelque part : la probabilité totale vaut $1$ (condition de normalisation).}
  {Sans le module au carré : $\Psi$ est complexe et son intégrale n'est pas une probabilité (elle peut être nulle ou complexe).}
  {$|\Psi|^2$ est une densité : si elle valait $1$ partout, son intégrale divergerait.}
  {Une fonction constante n'est pas normalisable ($\int|\Psi|^2\dd x$ diverge) ; la condition porte sur l'intégrale de $|\Psi|^2$.}

% ---- Question 32 (niveau 2, correct : D)
\qcmq[2]{Chap.~\ref{chap:histoire}\,$\cdot$\,\S\,\ref{sec:operateurs}}
  {Quelle est l'expression de l'opérateur quantité de mouvement $\hat p$ en une dimension ?}
  {$\hat p=m\,v$}
  {$\hat p=-\dfrac{\hbar^2}{2m}\dfrac{\partial^2}{\partial x^2}$}
  {$\hat p=+\ii\hbar\,\dfrac{\partial}{\partial x}$}
  {$\hat p=-\ii\hbar\,\dfrac{\partial}{\partial x}$}
  {D}
\qcmx
  {C'est la grandeur classique ; en mécanique quantique $p$ est un opérateur différentiel.}
  {C'est l'énergie cinétique $\hat p^2/2m$.}
  {Erreur de signe : on trouverait $p=-\hbar k$ pour l'onde $\ee^{\ii kx}$.}
  {Convention du cours : $\hat p\,\ee^{\ii kx}=\hbar k\,\ee^{\ii kx}$, donc une onde vers les $x$ croissants a $p=\hbar k>0$.}

% ---- Question 33 (niveau 2, correct : B)
\qcmq[2]{Chap.~\ref{chap:histoire}\,$\cdot$\,\S\,\ref{sec:solutions}}
  {Pour une particule dans un puits de potentiel infini de largeur $a$, comment varie l'énergie $E_n$ avec le numéro $n$ du niveau ?}
  {$E_n=E_1/n$}
  {$E_n=n^2E_1$, avec $E_1=\dfrac{\pi^2\hbar^2}{2ma^2}$}
  {$E_n=nE_1$}
  {$E_n=E_1$ pour tout $n$}
  {B}
\qcmx
  {Les niveaux se resserreraient quand $n$ augmente ; c'est l'inverse pour le puits infini.}
  {$k_n=\frac{n\pi}{a}$ et $E_n=\frac{\hbar^2k_n^2}{2m}$ : les niveaux s'écartent quand $n$ croît ($E_{n+1}-E_n=(2n+1)E_1$).}
  {Niveaux équidistants : ce n'est pas le puits infini, où $E_2=4E_1$ et $E_3=9E_1$.}
  {Il n'y aurait pas de quantification : tous les états auraient la même énergie.}

% ---- Question 34 (niveau 2, correct : C)
\qcmq[2]{Chap.~\ref{chap:histoire}\,$\cdot$\,\S\,\ref{sec:solutions}}
  {Une particule est dans un état stationnaire $\Psi_n$ de l'hamiltonien ($\Hop\Psi_n=E_n\Psi_n$). Que vaut l'écart-type $\sigma_E$ de son énergie ?}
  {$E_n$}
  {$E_n^2$}
  {$0$}
  {$\sqrt{E_n}$}
  {C}
\qcmx
  {C'est la valeur moyenne de l'énergie, pas son écart-type.}
  {C'est la valeur de $\langle\Hop^2\rangle$, avant soustraction de $\langle\Hop\rangle^2$.}
  {$\langle\Hop\rangle=E_n$ et $\langle\Hop^2\rangle=E_n^2$, donc $\sigma_E=\sqrt{E_n^2-E_n^2}=0$ : l'énergie est parfaitement déterminée.}
  {Mélange avec la formule $\sigma=\sqrt{\langle\Hop^2\rangle-\langle\Hop\rangle^2}$ : les deux termes sont égaux à $E_n^2$.}

% ---- Question 35 (niveau 2, correct : C)
\qcmq[2]{Chap.~\ref{chap:histoire}\,$\cdot$\,\S\,\ref{sec:solutions}}
  {Pour un état stationnaire $\Psi_n$ du puits infini $]0,a[$, que vaut la position moyenne $\langle\hat x\rangle$ ?}
  {$a/n$}
  {$a/4$}
  {$a/2$ pour tout $n$}
  {$0$}
  {C}
\qcmx
  {Elle dépendrait de $n$, ce que la symétrie de $|\Psi_n|^2$ interdit.}
  {C'est la position d'un maximum de $|\Psi_2|^2$ ; la moyenne des deux maxima $\frac a4$ et $\frac{3a}{4}$ donne $\frac a2$.}
  {$|\Psi_n|^2=\frac2a\sin^2\frac{n\pi x}{a}$ est symétrique par rapport au centre du puits.}
  {$\Psi_n$ s'annule sur la paroi, mais $\langle\hat x\rangle$ est une moyenne pondérée sur tout le puits.}

% ---- Question 36 (niveau 2, correct : D)
\qcmq[2]{Chap.~\ref{chap:dirac}\,$\cdot$\,\S\,\ref{sec:qubit}}
  {Un qubit est dans l'état $\ket{\Psi}=\frac{1}{\sqrt5}\ket{0}+\frac{2\ii}{\sqrt5}\ket{1}$. Quelle est la probabilité de mesurer $1$ en base $Z$ ?}
  {$\frac25$}
  {$-\frac45$}
  {$\frac15$}
  {$\frac45$}
  {D}
\qcmx
  {On divise le coefficient $2$ par $5$ : on oublie à la fois la racine et le carré.}
  {On élève $\frac{2\ii}{\sqrt5}$ au carré sans conjuguer : $\left(\frac{2\ii}{\sqrt5}\right)^2=-\frac45$ ; une probabilité est $|\cdot|^2$.}
  {C'est $P(0)=\left|\frac{1}{\sqrt5}\right|^2$.}
  {$P(1)=\left|\frac{2\ii}{\sqrt5}\right|^2=\frac45$ : le facteur $\ii$ est une phase relative, invisible en base $Z$.}

% ---- Question 37 (niveau 2, correct : C)
\qcmq[2]{Chap.~\ref{chap:dirac}\,$\cdot$\,\S\,\ref{sec:projecteurs}}
  {Un qubit préparé dans $\ket{0}$ est mesuré dans la base $X$ $\{\ket{+},\ket{-}\}$. Quelle est la probabilité d'obtenir $-$ ?}
  {$0$}
  {$1$}
  {$\frac12$}
  {$\frac{1}{\sqrt2}$}
  {C}
\qcmx
  {$\ket{0}$ est certain en base $Z$, mais pas en base $X$ : une autre base donne une autre loi.}
  {Même confusion dans l'autre sens : $\ket{0}$ n'est pas l'état $\ket{-}$.}
  {$P(-)=|\braket{-}{0}|^2=\left|\frac{1}{\sqrt2}\right|^2=\frac12$.}
  {C'est l'amplitude $\braket{-}{0}$ ; il faut élever son module au carré.}

% ---- Question 38 (niveau 2, correct : C)
\qcmq[2]{Chap.~\ref{chap:dirac}\,$\cdot$\,\S\,\ref{sec:qubit}}
  {Quel facteur $N>0$ normalise l'état $N\big(\ket{0}+(1-\ii)\ket{1}\big)$ ?}
  {$N=\frac{1}{\sqrt2}$}
  {$N=\frac13$}
  {$N=\frac{1}{\sqrt3}$}
  {$N=\frac{1}{\sqrt{1-2\ii}}$}
  {C}
\qcmx
  {On oublie la composante $\ket{0}$ : $|1-\ii|^2=2$, mais la norme au carré est $1+2=3$.}
  {On oublie la racine carrée : $3N^2=1$ donne $N=\frac{1}{\sqrt3}$, pas $\frac13$.}
  {$\braket{\Psi}{\Psi}=N^2\big(1+|1-\ii|^2\big)=N^2(1+2)=3N^2$.}
  {On élève $1-\ii$ au carré sans prendre le module : $(1-\ii)^2=-2\ii$ n'est pas réel ; il faut $|1-\ii|^2=2$.}

% ---- Question 39 (niveau 2, correct : D)
\qcmq[2]{Chap.~\ref{chap:dirac}\,$\cdot$\,\S\,\ref{sec:hilbert}}
  {Soient $\ket{\psi}=\frac{1}{\sqrt2}(\ket{0}+\ii\ket{1})$ et $\ket{\varphi}=\frac{1}{\sqrt2}(\ket{0}-\ket{1})$. Que vaut $\braket{\psi}{\varphi}$ ?}
  {$\frac{1-\ii}{2}$}
  {$\frac12$}
  {$0$}
  {$\frac{1+\ii}{2}$}
  {D}
\qcmx
  {On oublie de conjuguer l'amplitude $\ii$ du bra : on calcule $\frac12\big(1+\ii\cdot(-1)\big)$.}
  {Seul le terme en $\ket{0}$ est compté ; le terme en $\ket{1}$ vaut $\frac12(-\ii)(-1)=\frac{\ii}{2}$, non nul.}
  {Confusion avec $\braket{+}{-}=0$ ; ici $|\braket{\psi}{\varphi}|^2=\frac12$.}
  {Le bra conjugue les amplitudes : $\bra{\psi}=\frac{1}{\sqrt2}(\bra{0}-\ii\bra{1})$, d'où $\frac12\big(1+(-\ii)(-1)\big)=\frac{1+\ii}{2}$.}

% ---- Question 40 (niveau 2, correct : A)
\qcmq[2]{Chap.~\ref{chap:dirac}\,$\cdot$\,\S\,\ref{sec:qubit}}
  {Lequel de ces états est orthogonal à $\ket{{+}\ii}=\frac{1}{\sqrt2}(\ket{0}+\ii\ket{1})$ ?}
  {$\ket{{-}\ii}=\frac{1}{\sqrt2}(\ket{0}-\ii\ket{1})$}
  {$\ket{-}=\frac{1}{\sqrt2}(\ket{0}-\ket{1})$}
  {$\ket{+}=\frac{1}{\sqrt2}(\ket{0}+\ket{1})$}
  {$\ket{1}$}
  {A}
\qcmx
  {$\braket{{+}\ii}{{-}\ii}=\frac12\big(1+(-\ii)(-\ii)\big)=\frac12(1-1)=0$, en conjuguant le bra.}
  {$\braket{{+}\ii}{-}=\frac{1+\ii}{2}\neq0$ : $\ket{-}$ est orthogonal à $\ket{+}$, pas à $\ket{{+}\ii}$.}
  {$\braket{{+}\ii}{+}=\frac{1-\ii}{2}\neq0$.}
  {$\braket{{+}\ii}{1}=\frac{-\ii}{\sqrt2}\neq0$ : $\ket{1}$ est orthogonal à $\ket{0}$ seulement.}

% ---- Question 41 (niveau 2, correct : D)
\qcmq[1]{Chap.~\ref{chap:operateurs}\,$\cdot$\,\S\,\ref{sec:pur-mixte}}
  {Une matrice de densité de qubit a pour pureté $\operatorname{Tr}(\rho^2)=\frac58$. Que peut-on en conclure ?}
  {L'état est pur}
  {La matrice n'est pas une matrice de densité valide}
  {L'état est maximalement mixte, $\rho=\mathrm{Id}/2$}
  {L'état est mixte (mélange statistique)}
  {D}
\qcmx
  {Il faudrait $\operatorname{Tr}\rho^2=1$.}
  {$\frac58$ est compris entre $\frac12$ et $1$ : valeur admissible (exemple du cours : $\mathrm{diag}(\frac34,\frac14)$).}
  {Ce cas donnerait $\operatorname{Tr}\rho^2=\frac12$.}
  {$\operatorname{Tr}\rho^2<1$ ; pour un qubit, $\frac12\leq\operatorname{Tr}\rho^2\leq1$ et $\frac58$ est une valeur admissible.}

% ---- Question 42 (niveau 2, correct : B)
\qcmq[2]{Chap.~\ref{chap:operateurs}\,$\cdot$\,\S\,\ref{sec:rho}}
  {Que vaut $\langle Z\rangle=\operatorname{Tr}(\rho Z)$ pour un qubit dans l'état $\ket{+}$ ?}
  {$1$}
  {$0$}
  {$-1$}
  {$\frac{1}{\sqrt2}$}
  {B}
\qcmx
  {C'est la valeur pour $\ket{0}$.}
  {$\rho=\frac12\begin{pmatrix}1&1\\1&1\end{pmatrix}$ et $\operatorname{Tr}(\rho Z)=\frac12-\frac12=0$ : les résultats $+1$ et $-1$ sont équiprobables.}
  {C'est la valeur pour $\ket{1}$.}
  {C'est l'amplitude de $\ket{0}$ dans $\ket{+}$, pas une valeur moyenne de $Z$.}

% ---- Question 43 (niveau 2, correct : A)
\qcmq[1]{Chap.~\ref{chap:operateurs}\,$\cdot$\,\S\,\ref{sec:bloch}}
  {Sur la boule de Bloch, que représente un point situé à l'intérieur de la sphère ($|\vec r|<1$) ?}
  {Un état mixte, de pureté $\operatorname{Tr}\rho^2=\frac{1+|\vec r|^2}{2}<1$}
  {Un état pur dont la norme n'est pas $1$}
  {Un état pur en superposition}
  {Un point sans signification physique}
  {A}
\qcmx
  {Les états purs sont sur la sphère ; le centre $\vec r=\vec0$ est l'état maximalement mixte $\mathrm{Id}/2$.}
  {Tout état pur est normalisé et situé sur la sphère.}
  {Les états purs, superpositions comprises, sont tous \emph{sur} la sphère.}
  {Tout point de la boule est une matrice de densité valide.}

% ---- Question 44 (niveau 2, correct : B)
\qcmq[1]{Chap.~\ref{chap:operateurs}\,$\cdot$\,\S\,\ref{sec:pur-mixte}}
  {Par quoi la matrice de densité distingue-t-elle l'état pur $\ket{+}$ du mélange équiprobable de $\ket{0}$ et $\ket{1}$ ?}
  {Par leur trace, égale à $1$ pour l'état pur seulement}
  {Par les éléments hors diagonale (cohérences), non nuls pour $\ket{+}$ et nuls pour le mélange}
  {Par leurs éléments diagonaux, qui sont différents}
  {Ils ne se distinguent pas : c'est la même matrice}
  {B}
\qcmx
  {La trace vaut $1$ pour toute matrice de densité.}
  {$\ket{+}\bra{+}=\frac12\begin{pmatrix}1&1\\1&1\end{pmatrix}$ et le mélange vaut $\frac12\mathrm{Id}$ ; une mesure en base $X$ les distingue.}
  {Les deux ont la diagonale $(\frac12,\frac12)$ : $P(0)=P(1)=\frac12$ dans les deux cas.}
  {Les matrices sont différentes : en base $X$, $\ket{+}$ donne $+$ avec certitude, le mélange donne $\pm$ à pile ou face.}

% ---- Question 45 (niveau 2, correct : C)
\qcmq[2]{Chap.~\ref{chap:operateurs}\,$\cdot$\,\S\,\ref{sec:rho}}
  {Un qubit est préparé dans $\ket{0}$ avec la probabilité $\frac34$ et dans $\ket{1}$ avec la probabilité $\frac14$. Que vaut $\langle Z\rangle=\operatorname{Tr}(\rho Z)$ ?}
  {$\frac34$}
  {$1$}
  {$\frac12$}
  {$0$}
  {C}
\qcmx
  {C'est la probabilité $\rho_{00}$ de lire $0$ ; or $Z$ vaut $+1$ pour $0$ et $-1$ pour $1$.}
  {Valeur pour $\rho=\ket{0}\bra{0}$ (préparation certaine de $\ket{0}$).}
  {$\rho=\mathrm{diag}(\frac34,\frac14)$, donc $\operatorname{Tr}(\rho Z)=\frac34\cdot1+\frac14\cdot(-1)=\frac12$ : c'est la coordonnée $r_z$ de la sphère de Bloch.}
  {Valeur pour le mélange équiprobable $\mathrm{Id}/2$.}

% ---- Question 46 (niveau 2, correct : B)
\qcmq[2]{Chap.~\ref{chap:operateurs}\,$\cdot$\,\S\,\ref{sec:tenseur}}
  {Dans la base ordonnée $\{\ket{00},\ket{01},\ket{10},\ket{11}\}$, quel vecteur colonne représente $\ket{1}\otimes\ket{0}=\ket{10}$ ?}
  {$(0,1,0,0)^{\mathsf T}$}
  {$(0,0,1,0)^{\mathsf T}$}
  {$(1,0,0,0)^{\mathsf T}$}
  {$(0,0,0,1)^{\mathsf T}$}
  {B}
\qcmx
  {C'est $\ket{01}$ : l'ordre des facteurs est inversé.}
  {$\binom{0}{1}\otimes\binom{1}{0}=\big(0\binom{1}{0},\,1\binom{1}{0}\big)=(0,0,1,0)^{\mathsf T}$ : c'est le troisième vecteur de la base.}
  {C'est $\ket{00}$.}
  {C'est $\ket{11}$.}

% ---- Question 47 (niveau 2, correct : D)
\qcmq[2]{Chap.~\ref{chap:portes}\,$\cdot$\,\S\,\ref{sec:phase}}
  {Quelle relation lie la porte de phase $S=\mathrm{diag}(1,\ii)$ à la porte $Z$ ?}
  {$S=Z$}
  {$S^2=X$}
  {$S^4=Z$}
  {$S^2=Z$}
  {D}
\qcmx
  {$Z=\mathrm{diag}(1,-1)$ : $S$ est un quart de tour, $Z$ un demi-tour.}
  {$S^2$ est diagonale alors que $X$ ne l'est pas ; c'est $HSH$ qui est une racine carrée de $X$.}
  {$S^4=\mathrm{diag}(1,\ii^4)=\mathrm{Id}$.}
  {$S^2=\mathrm{diag}(1,\ii^2)=\mathrm{diag}(1,-1)=Z$ : deux quarts de tour autour de $z$ font un demi-tour.}

% ---- Question 48 (niveau 2, correct : A)
\qcmq[2]{Chap.~\ref{chap:portes}\,$\cdot$\,\S\,\ref{sec:controlees}}
  {Dans la porte de Toffoli (CCNOT) à trois qubits $x,y,z$, à quelle condition la cible $z$ est-elle inversée ?}
  {Seulement si $x=1$ et $y=1$}
  {Si $x=1$ ou $y=1$}
  {Si $x\neq y$}
  {Si $x=0$ et $y=0$}
  {A}
\qcmx
  {$\ket{x,y,z}\to\ket{x,y,z\oplus xy}$ : la cible est inversée si et seulement si le produit $xy$ vaut $1$.}
  {Un seul contrôle à $1$ ne suffit pas : pour $(x,y)=(1,0)$, $xy=0$ et la cible ne change pas.}
  {Ce serait deux CNOT ($z\to z\oplus x\oplus y$) ; pour $(1,0)$ le Toffoli n'inverse rien.}
  {Pour $(0,0)$ on a $xy=0$ : la cible ne change pas.}

% ---- Question 49 (niveau 2, correct : A)
\qcmq[1]{Chap.~\ref{chap:portes}\,$\cdot$\,\S\,\ref{sec:portes}}
  {Pour un hamiltonien $\Hop$ indépendant du temps, quel opérateur d'évolution $U(t,t_0)$ vérifie $\ket{\Psi(t)}=U(t,t_0)\ket{\Psi(t_0)}$ ?}
  {$U(t,t_0)=\ee^{-\ii\Hop(t-t_0)/\hbar}$}
  {$U(t,t_0)=\ee^{+\ii\Hop(t-t_0)/\hbar}$}
  {$U(t,t_0)=-\ii\hbar\,\partial_t$}
  {$U(t,t_0)=\Hop\,(t-t_0)/\hbar$}
  {A}
\qcmx
  {Solution de $\ii\hbar\,\partial_t\ket{\Psi}=\Hop\ket{\Psi}$ ; comme $\Hop$ est hermitien, $U^\dagger U=\mathrm{Id}$ : l'évolution est unitaire.}
  {Erreur de signe : c'est $U^\dagger$, qui remonte le temps.}
  {C'est un opérateur différentiel, pas un opérateur d'évolution ; l'équation de Schrödinger s'écrit $\ii\hbar\,\partial_t\ket{\Psi}=\Hop\ket{\Psi}$.}
  {On oublie l'exponentielle et le facteur $\ii$ : cet opérateur est hermitien, pas unitaire.}

% ---- Question 50 (niveau 2, correct : D)
\qcmq[2]{Chap.~\ref{chap:portes}\,$\cdot$\,\S\,\ref{sec:controlees}}
  {La porte de Fredkin (SWAP contrôlé) a pour contrôle le premier qubit. Que donne-t-elle sur $\ket{1,0,1}$ ?}
  {$\ket{1,0,1}$}
  {$\ket{0,1,1}$}
  {$\ket{1,1,1}$}
  {$\ket{1,1,0}$}
  {D}
\qcmx
  {C'est le résultat si le contrôle valait $0$ ; ici il vaut $1$ et les deux autres qubits sont différents.}
  {Le SWAP porte sur les qubits $2$ et $3$, pas sur le contrôle.}
  {Un SWAP ne change pas le nombre de $1$, il permute seulement ; $\ket{1,1,1}$ viendrait d'un NOT.}
  {Le contrôle vaut $1$ : les deux autres qubits sont échangés, $\ket{0,1}\to\ket{1,0}$.}

% ---- Question 51 (niveau 2, correct : A)
\qcmq[1]{Chap.~\ref{chap:portes}\,$\cdot$\,\S\,\ref{sec:op2}}
  {Que donne $(H\otimes I)\ket{00}$, où $H$ agit sur le premier facteur du produit tensoriel $\ket{a}\otimes\ket{b}$ ?}
  {$\ket{+}\otimes\ket{0}=\frac{1}{\sqrt2}(\ket{00}+\ket{10})$}
  {$\frac12(\ket{00}+\ket{01}+\ket{10}+\ket{11})$}
  {$\ket{0}\otimes\ket{+}=\frac{1}{\sqrt2}(\ket{00}+\ket{01})$}
  {$\frac{1}{\sqrt2}(\ket{00}+\ket{11})$}
  {A}
\qcmx
  {$(A\otimes B)(\ket{u}\otimes\ket{v})=A\ket{u}\otimes B\ket{v}$ : $H\ket{0}=\ket{+}$ sur le premier facteur, $I\ket{0}=\ket{0}$ sur le second.}
  {C'est $(H\otimes H)\ket{00}$ : ici le second qubit n'est pas touché.}
  {$H$ agit sur le premier qubit, pas sur le second : les facteurs sont inversés.}
  {C'est $\ket{\Phi^+}$ : il faut un CNOT pour intriquer ; des portes locales laissent l'état séparable.}

% ---- Question 52 (niveau 2, correct : C)
\qcmq[1]{Chap.~\ref{chap:bell}\,$\cdot$\,\S\,\ref{sec:separable}}
  {Pourquoi l'état de Bell $\ket{\Phi^+}=\frac{1}{\sqrt2}(\ket{00}+\ket{11})$ est-il dit intriqué ?}
  {Parce que $P(01)=0$}
  {Parce que chaque qubit, pris seul, est dans un état pur}
  {Il ne peut pas s'écrire comme un produit $\ket{\varphi}_A\otimes\ket{\chi}_B$ : $D=c_{00}c_{11}-c_{01}c_{10}=\frac12\neq0$}
  {Parce qu'il est orthogonal à $\ket{00}$}
  {C}
\qcmx
  {C'est vrai, mais $\ket{00}$ a aussi $P(01)=0$ et il est séparable : cette condition ne prouve rien.}
  {C'est l'inverse : pris seul, chaque qubit est complètement mixte ($\rho_A=\mathrm{Id}/2$).}
  {Un produit imposerait $D=0$ ; aucun des deux qubits n'a d'état propre.}
  {$\braket{00}{\Phi^+}=\frac{1}{\sqrt2}\neq0$ ; l'orthogonalité n'a pas de lien avec l'intrication.}

% ---- Question 53 (niveau 2, correct : B)
\qcmq[2]{Chap.~\ref{chap:bell}\,$\cdot$\,\S\,\ref{sec:teleportation}}
  {Dans la téléportation du cours, Alice obtient les résultats $a=0$ et $q=1$. Quelle correction Bob doit-il appliquer à son qubit ?}
  {La porte $X$}
  {La porte $Z$}
  {$X$ puis $Z$}
  {Aucune, l'état est déjà $\ket{\varphi}$}
  {B}
\qcmx
  {$X$ corrige le cas $a=1$ (inversion de bit) ; ici $a=0$.}
  {Bob détient $\alpha\ket{0}-\beta\ket{1}=Z\ket{\varphi}$ ; $Z$ rétablit le signe : $Z(\alpha\ket{0}-\beta\ket{1})=\alpha\ket{0}+\beta\ket{1}$.}
  {Les deux corrections ne sont nécessaires que pour $(a,q)=(1,1)$.}
  {Ce n'est vrai que pour $(a,q)=(0,0)$.}

% ---- Question 54 (niveau 2, correct : B)
\qcmq[2]{Chap.~\ref{chap:bell}\,$\cdot$\,\S\,\ref{sec:bell} ; ex.\,\ref{ex:superdense}}
  {Dans le codage superdense, Alice et Bob partagent $\ket{\Phi^+}$ et Alice agit sur son qubit. Quelle porte transforme $\ket{\Phi^+}$ en $\ket{\Psi^+}=\frac{1}{\sqrt2}(\ket{01}+\ket{10})$ ?}
  {La porte $Z$}
  {La porte $X$}
  {La porte identité}
  {La porte $H$}
  {B}
\qcmx
  {$Z$ donne $\ket{\Phi^-}=\frac{1}{\sqrt2}(\ket{00}-\ket{11})$ : elle change un signe, pas les chiffres.}
  {$X$ échange $\ket{0}$ et $\ket{1}$ sur un qubit : $\ket{00}+\ket{11}\to\ket{01}+\ket{10}$ ; c'est le message $10$.}
  {Elle laisse $\ket{\Phi^+}$ : c'est le message $00$.}
  {$(H\otimes I)\ket{\Phi^+}$ est une superposition de deux états de Bell, pas un état de Bell.}

% ---- Question 55 (niveau 2, correct : D)
\qcmq[2]{Chap.~\ref{chap:bell}\,$\cdot$\,\S\,\ref{sec:info}}
  {Pour deux qubits dans l'état de Bell $\ket{\Phi^+}$, que vaut l'entropie conditionnelle $H(A|B)=H(A,B)-H(B)$ ?}
  {$0$ bit}
  {$+1$ bit}
  {$+2$ bits}
  {$-1$ bit}
  {D}
\qcmx
  {C'est la valeur classique quand $A$ est entièrement déterminé par $B$ (deux bits identiques) ; ici $H(A,B)=0$, pas $H(B)$.}
  {C'est la valeur quand $A$ et $B$ sont indépendants : $H(A|B)=H(A)=1$.}
  {C'est l'information mutuelle $I(A;B)=1+1-0=2$ bits, pas l'entropie conditionnelle.}
  {$H(A,B)=S(\rho_{AB})=0$ (état pur) et $H(B)=1$ bit ($\rho_B=\mathrm{Id}/2$) : $0-1=-1$, valeur impossible en classique.}

% ---- Question 56 (niveau 2, correct : B)
\qcmq[1]{Chap.~\ref{chap:bell}\,$\cdot$\,\S\,\ref{sec:trace}}
  {Quelle est la matrice de densité réduite $\rho_A=\operatorname{Tr}_B\rho_{AB}$ d'un qubit d'une paire dans l'état $\ket{\Phi^+}$ ?}
  {$\rho_A=\ket{0}\bra{0}$}
  {$\rho_A=\frac12\mathrm{Id}$, de pureté $\operatorname{Tr}\rho_A^2=\frac12$}
  {$\rho_A=\ket{+}\bra{+}=\frac12\begin{pmatrix}1&1\\1&1\end{pmatrix}$}
  {$\rho_A=\ket{\Phi^+}\bra{\Phi^+}$}
  {B}
\qcmx
  {Le qubit aurait un état défini, ce qui est faux pour un état intriqué : $P(0)=P(1)=\frac12$.}
  {$\rho_A=\frac12\big(\ket{0}\bra{0}+\ket{1}\bra{1}\big)$ : les termes croisés $\ket{00}\bra{11}$ s'annulent ($\braket{1}{0}=0$). Le couple est pur, mais chaque qubit est complètement mixte.}
  {Les termes croisés sont multipliés par $\braket{1}{0}=0$ dans la trace partielle : ils disparaissent.}
  {C'est la matrice $4\times4$ du couple ; la trace partielle donne une matrice $2\times2$.}

% ---- Question 57 (niveau 2, correct : C)
\qcmq[1]{Chap.~\ref{chap:bell}\,$\cdot$\,\S\,\ref{sec:separable}}
  {Pour $\ket{\Psi}=c_{00}\ket{00}+c_{01}\ket{01}+c_{10}\ket{10}+c_{11}\ket{11}$, quel est le critère de séparabilité ?}
  {$c_{00}c_{11}+c_{01}c_{10}=0$}
  {$c_{00}c_{01}-c_{10}c_{11}=0$}
  {$D=c_{00}c_{11}-c_{01}c_{10}=0$}
  {$c_{00}=c_{11}$ et $c_{01}=c_{10}$}
  {C}
\qcmx
  {Mauvais signe : $\ket{+}\otimes\ket{-}$ est séparable mais donnerait $-\frac12\neq0$.}
  {Mauvais appariement : $\ket{\Psi^+}$ (intriqué) la vérifie, puisque $c_{00}=c_{11}=0$.}
  {Un produit $(a_0\ket{0}+a_1\ket{1})\otimes(b_0\ket{0}+b_1\ket{1})$ a $c_{ij}=a_ib_j$, donc $c_{00}c_{11}=a_0a_1b_0b_1=c_{01}c_{10}$ ; la réciproque est vraie aussi.}
  {Ni nécessaire ($\ket{00}$ est séparable avec $c_{00}\neq c_{11}$) ni suffisant ($\ket{\Phi^+}$ la vérifie et est intriqué).}

% ---- Question 58 (niveau 2, correct : B)
\qcmq[2]{Chap.~\ref{chap:bell}\,$\cdot$\,\S\,\ref{sec:bell}}
  {Quelle est la valeur propre de $X\otimes X$ pour l'état $\ket{\Phi^+}$ ?}
  {$-1$}
  {$+1$}
  {$0$}
  {$\ket{\Phi^+}$ n'est pas un état propre de $X\otimes X$}
  {B}
\qcmx
  {C'est la valeur pour $\ket{\Phi^-}$ et $\ket{\Psi^-}$ : le signe $-$ de l'état se retrouve dans la valeur propre.}
  {$(X\otimes X)\ket{00}=\ket{11}$ et $(X\otimes X)\ket{11}=\ket{00}$, donc $(X\otimes X)\ket{\Phi^+}=\ket{\Phi^+}$.}
  {$X\otimes X$ est unitaire : ses valeurs propres sont de module $1$, jamais nulles.}
  {Les quatre états de Bell sont des états propres communs de $Z\otimes Z$ et de $X\otimes X$.}

% ---- Question 59 (niveau 2, correct : A)
\qcmq[2]{Chap.~\ref{chap:bell}\,$\cdot$\,\S\,\ref{sec:info}}
  {Que vaut l'entropie de von Neumann $S(\rho)=-\sum_i\lambda_i\log_2\lambda_i$ d'un état pur ?}
  {$0$ bit}
  {$1$ bit}
  {$-1$ bit}
  {$+\infty$}
  {A}
\qcmx
  {Les valeurs propres sont $1$ et $0$ : $1\cdot\log_21=0$ et $0\cdot\log_20=0$ (par continuité, $x\log_2x\to0$).}
  {C'est la valeur pour l'état maximalement mixte $\mathrm{Id}/2$ d'un qubit.}
  {L'entropie de $\rho$ est positive ou nulle ; seule l'entropie conditionnelle quantique peut être négative.}
  {$\log_20$ diverge, mais le terme $0\cdot\log_20$ vaut $0$ par continuité.}

% ---- Question 60 (niveau 2, correct : A)
\qcmq[2]{Chap.~\ref{chap:bell}\,$\cdot$\,\S\,\ref{sec:bell}}
  {Lequel des états de Bell a pour valeur propre $-1$ à la fois pour $Z\otimes Z$ et pour $X\otimes X$ ?}
  {$\ket{\Psi^-}=\frac{1}{\sqrt2}(\ket{01}-\ket{10})$}
  {$\ket{\Phi^+}$}
  {$\ket{\Phi^-}$}
  {$\ket{\Psi^+}$}
  {A}
\qcmx
  {Résultats opposés dans les deux bases : $Z\otimes Z=-1$ (chiffres différents) et $X\otimes X=-1$ (signe $-$ dans l'état).}
  {Valeurs propres $(Z\otimes Z,\,X\otimes X)=(+1,+1)$.}
  {Valeurs propres $(Z\otimes Z,\,X\otimes X)=(+1,-1)$.}
  {Valeurs propres $(Z\otimes Z,\,X\otimes X)=(-1,+1)$.}

\clearpage
%% ======================================================================
\subsection{Niveau 3 : avancé (questions 61 à 90)}\label{sec:qcm-n3}
%% ======================================================================
\noindent{\color{insagris}\small Niveau de difficulté : Avancé\quad $\bullet\bullet\bullet\circ$}\par\smallskip
% ---- Question 61 (niveau 3, correct : A)
\qcmq[2]{Chap.~\ref{chap:histoire}\,$\cdot$\,\S\,\ref{sec:solutions}}
  {Un électron d'un puits infini passe du niveau $n=3$ au niveau $n=2$ en émettant un photon. Quelle est l'énergie du photon, en fonction de $E_1$ ?}
  {$5E_1$}
  {$E_1$}
  {$\frac94E_1$}
  {$13E_1$}
  {A}
\qcmx
  {$E_3-E_2=9E_1-4E_1=5E_1$ ; le photon a la fréquence $f=(E_3-E_2)/h$.}
  {On a soustrait les numéros de niveau ($3-2=1$) au lieu de leurs carrés.}
  {C'est le rapport $E_3/E_2$, pas la différence $E_3-E_2$.}
  {On a additionné $9E_1+4E_1$ : l'énergie du photon est une différence de niveaux.}

% ---- Question 62 (niveau 3, correct : A)
\qcmq[2]{Chap.~\ref{chap:histoire}\,$\cdot$\,\S\,\ref{sec:superposition}}
  {À $t=0$, une particule du puits infini est dans l'état $\Psi=\frac{\sqrt3}{2}\,\psi_1+\frac12\,\psi_3$ ($\psi_n$ : états stationnaires normalisés, d'énergies $E_n=n^2E_1$). Quelle est son énergie moyenne $\langle E\rangle$ ?}
  {$3E_1$}
  {$5E_1$}
  {$\left(\frac{\sqrt3}{2}+\frac92\right)E_1$}
  {$7E_1$}
  {A}
\qcmx
  {$P(E_1)=\frac34$ et $P(E_3)=\frac14$, donc $\langle E\rangle=\frac34E_1+\frac14\cdot9E_1=3E_1$.}
  {On a pondéré également les deux niveaux, $\frac{E_1+9E_1}{2}$ : les poids sont les $|c_i|^2$, pas $\frac12$.}
  {On a pondéré par les amplitudes $c_i$ au lieu de leurs carrés $|c_i|^2$.}
  {Poids permutés : $\frac14E_1+\frac34\cdot9E_1=7E_1$ ; la probabilité de $E_1$ est $\left|\frac{\sqrt3}{2}\right|^2=\frac34$.}

% ---- Question 63 (niveau 3, correct : C)
\qcmq[1]{Chap.~\ref{chap:histoire}\,$\cdot$\,\S\,\ref{sec:solutions}}
  {Pour l'état stationnaire $n=2$ du puits infini $]0,a[$, que valent la densité de probabilité de présence en $x=a/2$ et la position moyenne $\langle\hat x\rangle$ ?}
  {Densité maximale en $a/2$, et $\langle\hat x\rangle=a/2$}
  {Densité nulle en $a/2$, et $\langle\hat x\rangle=0$}
  {Densité nulle en $a/2$, et $\langle\hat x\rangle=a/2$}
  {Densité maximale en $a/2$, et $\langle\hat x\rangle=a/4$}
  {C}
\qcmx
  {C'est le cas de $n=1$ ; pour $n=2$ il y a un nœud au centre.}
  {Une densité nulle au centre ne donne pas une moyenne nulle : $\langle\hat x\rangle$ pondère tout le puits et vaut $\frac a2$ par symétrie.}
  {$|\Psi_2|^2=\frac2a\sin^2\frac{2\pi x}{a}$ s'annule en $x=\frac a2$ (nœud), alors que la symétrie impose $\langle\hat x\rangle=\frac a2$ : la moyenne n'est pas la position la plus probable.}
  {Ni l'un ni l'autre : maxima en $\frac a4$ et $\frac{3a}4$, nœud en $\frac a2$, moyenne $\frac a2$.}

% ---- Question 64 (niveau 3, correct : C)
\qcmq[2]{Chap.~\ref{chap:dirac}\,$\cdot$\,\S\,\ref{sec:projecteurs}}
  {Quelle est la matrice du projecteur $\hat P_+=\ket{+}\bra{+}$ dans la base $\{\ket{0},\ket{1}\}$ ?}
  {$\frac12\begin{pmatrix}1&-1\\-1&1\end{pmatrix}$}
  {$\begin{pmatrix}1&0\\0&0\end{pmatrix}$}
  {$\frac12\begin{pmatrix}1&1\\1&1\end{pmatrix}$}
  {$\frac{1}{\sqrt2}\begin{pmatrix}1&1\\1&-1\end{pmatrix}$}
  {C}
\qcmx
  {C'est $\ket{-}\bra{-}$, le projecteur sur $\ket{-}$.}
  {C'est $\ket{0}\bra{0}$ : projecteur de la base $Z$, pas de la base $X$.}
  {$\ket{+}\bra{+}=\frac12(\ket0+\ket1)(\bra0+\bra1)$ ; on vérifie $\hat P_+^2=\hat P_+$ et $\operatorname{Tr}\hat P_+=1$.}
  {C'est la matrice de Hadamard, unitaire ($H^2=\mathrm{Id}$), et non un projecteur (un projecteur vérifie $P^2=P$).}

% ---- Question 65 (niveau 3, correct : B)
\qcmq[2]{Chap.~\ref{chap:dirac}\,$\cdot$\,\S\,\ref{sec:projecteurs}}
  {Un qubit préparé dans $\ket{0}$ est mesuré en base $X$, puis aussitôt en base $Z$. Quelle est la probabilité d'obtenir $1$ à la seconde mesure ?}
  {$0$}
  {$\frac12$}
  {$1$}
  {$\frac14$}
  {B}
\qcmx
  {On garde $\ket{0}$ en mémoire : la mesure en base $X$ a effondré l'état sur $\ket{\pm}$ et effacé $\ket{0}$.}
  {La première mesure donne $\ket{+}$ ou $\ket{-}$ (probabilité $\frac12$ chacun) ; dans les deux cas $P(1)=\frac12$ en base $Z$ : $\frac12\cdot\frac12+\frac12\cdot\frac12=\frac12$.}
  {Sans rapport : $\ket{0}$ n'est plus l'état après la mesure en base $X$.}
  {On n'a compté qu'une branche de la première mesure ; il faut sommer les deux ($\frac14+\frac14$).}

% ---- Question 66 (niveau 3, correct : D)
\qcmq[2]{Chap.~\ref{chap:dirac}\,$\cdot$\,\S\,\ref{sec:projecteurs}}
  {Un qubit est dans l'état $\ket{\psi}=\frac{3}{5}\ket{0}+\frac{4}{5}\ket{1}$. Quelle est la probabilité d'obtenir $-$ en mesurant dans la base $X$ ?}
  {$\frac{16}{25}$}
  {$\frac{1}{25}$}
  {$\frac{49}{50}$}
  {$\frac{1}{50}$}
  {D}
\qcmx
  {C'est $P(1)$ en base $Z$ ($|\frac45|^2$) : on a mesuré dans la mauvaise base.}
  {On a oublié le facteur $\frac{1}{\sqrt2}$ de $\ket{-}$ : $\left|\frac35-\frac45\right|^2=\frac1{25}$, il manque le $\frac12$.}
  {C'est $P(+)$ ; la probabilité demandée est celle du résultat $-$.}
  {$\braket{-}{\psi}=\frac{1}{\sqrt2}\left(\frac35-\frac45\right)=-\frac{1}{5\sqrt2}$, donc $P(-)=\frac1{50}$ ; et $P(+)=\frac{49}{50}$.}

% ---- Question 67 (niveau 3, correct : B)
\qcmq[1]{Chap.~\ref{chap:dirac}\,$\cdot$\,\S\,\ref{sec:projecteurs}}
  {Lequel de ces kets décrit le même état physique que $\ket{+}=\frac{1}{\sqrt2}(\ket{0}+\ket{1})$ ?}
  {$\frac{1}{\sqrt2}\big(\ket{0}+\ii\ket{1}\big)$}
  {$\ee^{\ii\pi/3}\ket{+}$}
  {$\frac{1}{\sqrt2}\big(\ket{0}+\ee^{\ii\pi}\ket{1}\big)$}
  {$\frac{1}{\sqrt2}\big(\ket{0}+\ee^{\ii\pi/3}\ket{1}\big)$}
  {B}
\qcmx
  {La phase $\ii$ ne porte que sur $\ket{1}$ : c'est une phase \emph{relative}, et l'état est $\ket{{+}\ii}$ (base $Y$).}
  {Un facteur de phase \emph{globale} multiplie tout le ket : toutes les probabilités de mesure (et $\rho=\ket{+}\bra{+}$) sont inchangées.}
  {$\ee^{\ii\pi}=-1$ : c'est $\ket{-}$, orthogonal à $\ket{+}$.}
  {Phase relative $\frac\pi3$ : c'est un autre état, sur l'équateur de la sphère de Bloch, à l'azimut $\frac\pi3$.}

% ---- Question 68 (niveau 3, correct : A)
\qcmq[2]{Chap.~\ref{chap:dirac}\,$\cdot$\,\S\,\ref{sec:qubit}}
  {Laquelle de ces paires est une base orthonormée de $\mathbb{C}^2$ ?}
  {$\{\ket{{+}\ii},\ket{{-}\ii}\}$}
  {$\{\ket{0},\ket{+}\}$}
  {$\{\ket{+},\ket{{-}\ii}\}$}
  {$\{\ket{0},\ket{0}+\ket{1}\}$}
  {A}
\qcmx
  {Ces deux états sont normés et orthogonaux : $\braket{{+}\ii}{{-}\ii}=\frac12(1-1)=0$.}
  {Normés, mais $\braket{0}{+}=\frac{1}{\sqrt2}\neq0$ : pas orthogonaux.}
  {$\braket{+}{{-}\ii}=\frac12(1-\ii)\neq0$ : pas orthogonaux.}
  {Le second ket n'est pas normé (norme au carré $2$) ni orthogonal à $\ket{0}$.}

% ---- Question 69 (niveau 3, correct : C)
\qcmq[2]{Chap.~\ref{chap:operateurs}\,$\cdot$\,\S\,\ref{sec:rho}}
  {Laquelle de ces matrices peut être la matrice de densité d'un qubit ?}
  {$\begin{pmatrix}\frac12&\frac34\\\frac34&\frac12\end{pmatrix}$}
  {$\begin{pmatrix}\frac12&\frac{\ii}{2}\\\frac{\ii}{2}&\frac12\end{pmatrix}$}
  {$\begin{pmatrix}\frac23&\frac{\ii}{3}\\-\frac{\ii}{3}&\frac13\end{pmatrix}$}
  {$\begin{pmatrix}\frac34&0\\0&\frac12\end{pmatrix}$}
  {C}
\qcmx
  {Hermitienne et de trace $1$, mais pas positive : $\det=\frac14-\frac9{16}<0$, valeurs propres $\frac54$ et $-\frac14$ (une « probabilité » négative).}
  {De trace $1$ mais pas hermitienne : il faudrait $\rho_{10}=\overline{\rho_{01}}=-\frac{\ii}{2}$.}
  {Hermitienne ($\rho_{10}=\overline{\rho_{01}}$), de trace $1$ et positive : $\det\rho=\frac29-\frac19=\frac19>0$ avec $\operatorname{Tr}\rho>0$ donne deux valeurs propres positives, $\frac{3\pm\sqrt5}{6}$.}
  {Hermitienne et positive, mais de trace $\frac54\neq1$ : les probabilités diagonales ne somment pas à $1$.}

% ---- Question 70 (niveau 3, correct : B)
\qcmq[2]{Chap.~\ref{chap:operateurs}\,$\cdot$\,\S\,\ref{sec:op-dirac}}
  {Un qubit est dans l'état $\ket{\psi}=\frac{1}{\sqrt5}\big(\ket{0}+2\ii\ket{1}\big)$. Que vaut la valeur moyenne $\langle Y\rangle=\bra{\psi}Y\ket{\psi}$, avec $Y=\begin{pmatrix}0&-\ii\\\ii&0\end{pmatrix}$ ?}
  {$0$}
  {$\frac45$}
  {$-\frac45$}
  {$-\frac35$}
  {B}
\qcmx
  {On a oublié de conjuguer les amplitudes du bra : $\psi^{\mathsf T}Y\psi=\frac15(2+2\ii\cdot\ii)=0$ (c'est d'ailleurs la valeur de $\langle X\rangle$ pour cet état).}
  {$Y\ket\psi=\frac1{\sqrt5}(2\ket0+\ii\ket1)$, puis $\bra\psi Y\ket\psi=\frac15\big(1\cdot2+(-2\ii)(\ii)\big)=\frac45$ ; en général $\langle Y\rangle=2\,\mathrm{Im}(\bar\alpha\beta)$.}
  {Erreur de signe : on a utilisé la transposée $\begin{pmatrix}0&\ii\\-\ii&0\end{pmatrix}=-Y$.}
  {C'est $\langle Z\rangle=\frac15-\frac45$ : on a calculé la valeur moyenne d'une autre observable.}

% ---- Question 71 (niveau 3, correct : D)
\qcmq[1]{Chap.~\ref{chap:operateurs}\,$\cdot$\,\S\,\ref{sec:spectral}}
  {On mesure l'observable $A=\begin{pmatrix}2&1\\1&2\end{pmatrix}$ sur un qubit préparé dans $\ket{0}$. Quelle affirmation est correcte ?}
  {On obtient toujours $2$}
  {On obtient $+1$ ou $-1$, avec la probabilité $\frac12$ chacun}
  {On obtient $2$ ou $1$, avec les probabilités $\frac45$ et $\frac15$}
  {On obtient $3$ ou $1$, avec la probabilité $\frac12$ chacun (moyenne $2$)}
  {D}
\qcmx
  {On a lu le coefficient diagonal : $A\ket0=2\ket0+\ket1$ n'est pas proportionnel à $\ket0$, donc $\ket0$ n'est pas vecteur propre de $A$.}
  {Ce sont les valeurs propres de $X$, pas de $A$ : le décalage $2\,\mathrm{Id}$ translate le spectre ($2\pm1$).}
  {On a lu $A\ket0=2\ket0+\ket1$ comme un état : $A$ n'est pas unitaire, ce n'est pas une évolution ; les résultats possibles d'une mesure sont les valeurs propres de $A$.}
  {$A=2\,\mathrm{Id}+X$ a les vecteurs propres de $X$ : valeur propre $3$ pour $\ket+$, $1$ pour $\ket-$. Comme $\ket0=\frac{1}{\sqrt2}(\ket++\ket-)$, les deux résultats sont équiprobables, et $\langle A\rangle=\bra0A\ket0=2$.}

% ---- Question 72 (niveau 3, correct : D)
\qcmq[2]{Chap.~\ref{chap:operateurs}\,$\cdot$\,\S\,\ref{sec:pur-mixte}}
  {Un qubit est préparé dans $\ket{0}$ avec la probabilité $\frac12$ et dans $\ket{+}$ avec la probabilité $\frac12$. Quelle est la pureté $\operatorname{Tr}(\rho^2)$ de son état ?}
  {$1$}
  {$\frac12$}
  {$\frac58$}
  {$\frac34$}
  {D}
\qcmx
  {Un mélange de deux états purs n'est pur que s'ils sont identiques ; ici $\ket0\neq\ket+$, donc $\operatorname{Tr}\rho^2<1$.}
  {Valeur de l'état maximalement mixte $\frac{\mathrm{Id}}2$, atteinte pour deux états \emph{orthogonaux} ; ici ils ne le sont pas (on a aussi pu oublier les termes croisés).}
  {On n'a additionné que les carrés des éléments diagonaux ($\frac9{16}+\frac1{16}$) : il manque $2\rho_{01}\rho_{10}=\frac2{16}$.}
  {$\rho=\frac12\ket0\bra0+\frac12\ket+\bra+=\begin{pmatrix}\frac34&\frac14\\\frac14&\frac14\end{pmatrix}$, donc $\operatorname{Tr}\rho^2=\frac9{16}+\frac1{16}+2\cdot\frac1{16}=\frac34$ : état mixte, mais moins que $\frac{\mathrm{Id}}2$ car $\ket0$ et $\ket+$ ne sont pas orthogonaux.}

% ---- Question 73 (niveau 3, correct : B)
\qcmq[2]{Chap.~\ref{chap:operateurs}\,$\cdot$\,\S\,\ref{sec:bloch}}
  {Quel est le vecteur de Bloch $(r_x,r_y,r_z)$ de l'état $\ket{\psi}=\cos\frac\pi6\,\ket0+\sin\frac\pi6\,\ket1$ ?}
  {$\left(\frac12,\,0,\,\frac{\sqrt3}{2}\right)$}
  {$\left(\frac{\sqrt3}{2},\,0,\,\frac12\right)$}
  {$\left(\frac{\sqrt3}{4},\,0,\,\frac12\right)$}
  {$\left(\frac{\sqrt3}{2},\,0,\,\frac34\right)$}
  {B}
\qcmx
  {On a pris $\theta=\frac\pi6$ : l'angle de Bloch est le \emph{double} de l'angle qui figure dans les amplitudes.}
  {Ici $\frac\theta2=\frac\pi6$, donc $\theta=\frac\pi3$ et $\varphi=0$ : $\vec r=(\sin\theta\cos\varphi,\sin\theta\sin\varphi,\cos\theta)$. Contrôle : $r_z=|\alpha|^2-|\beta|^2=\frac34-\frac14$ et $r_x=2\alpha\beta=\frac{\sqrt3}2$.}
  {On a écrit $r_x=\alpha\beta$ en oubliant le facteur $2$ ; on le voit à $|\vec r|^2=\frac7{16}<1$, impossible pour un état pur.}
  {On a pris $r_z=|\alpha|^2$ au lieu de $|\alpha|^2-|\beta|^2$ ; or $|\vec r|^2=\frac{21}{16}>1$ est impossible.}

% ---- Question 74 (niveau 3, correct : A)
\qcmq[2]{Chap.~\ref{chap:operateurs}\,$\cdot$\,\S\,\ref{sec:bloch}}
  {Un qubit a pour matrice de densité $\rho=\frac12\begin{pmatrix}1&-\ii\\\ii&1\end{pmatrix}$. Quel est son vecteur de Bloch ?}
  {$(0,\,1,\,0)$ : c'est l'état pur $\ket{{+}\ii}$}
  {$(0,\,-1,\,0)$}
  {$\left(0,\,\frac12,\,0\right)$}
  {$(0,\,0,\,0)$}
  {A}
\qcmx
  {Avec $\rho=\frac12\begin{pmatrix}1+r_z&r_x-\ii r_y\\r_x+\ii r_y&1-r_z\end{pmatrix}$ : $r_z=0$ et $r_x-\ii r_y=-\ii$, donc $r_y=1$, $r_x=0$. Comme $|\vec r|=1$, l'état est pur : $\ket{{+}\ii}\bra{{+}\ii}$.}
  {Erreur de signe : $\rho_{01}=\frac{r_x-\ii r_y}{2}$, donc $\rho_{01}=-\frac\ii2$ correspond à $r_y=+1$ ($r_y=-1$ pour $\ket{{-}\ii}$).}
  {On a lu $\rho_{01}=-\frac\ii2$ comme $-\ii r_y$ : il faut multiplier par $2$ ($\rho=\frac12(\mathrm{Id}+\vec r\cdot\vec\sigma)$).}
  {Les diagonales $\frac12,\frac12$ donnent $r_z=0$, mais les éléments non diagonaux (cohérences) ne sont pas nuls : l'état est pur, et non $\frac{\mathrm{Id}}2$.}

% ---- Question 75 (niveau 3, correct : C)
\qcmq[2]{Chap.~\ref{chap:operateurs}\,$\cdot$\,\S\,\ref{sec:tenseur}}
  {Que donne $(X\otimes Z)\,\big(\ket{+}\otimes\ket{-}\big)$ ?}
  {$-\ket{-}\otimes\ket{-}$}
  {$\ket{-}\otimes\ket{+}$}
  {$\ket{+}\otimes\ket{+}$}
  {$\ket{+}\otimes\ket{-}$}
  {C}
\qcmx
  {On a échangé les facteurs ($Z\otimes X$) : $Z\ket+=\ket-$ et $X\ket-=-\ket-$. Le premier opérateur agit sur le premier qubit.}
  {On a cru que $X$ transforme $\ket+$ en $\ket-$ ; or $X$ échange $\ket0$ et $\ket1$ et laisse $\ket+$ inchangé.}
  {$(A\otimes B)(\ket u\otimes\ket v)=A\ket u\otimes B\ket v$ : $X\ket+=\ket+$ ($\ket+$ est propre pour $X$, valeur propre $+1$) et $Z\ket-=\ket+$ ($Z$ échange $\ket+$ et $\ket-$).}
  {On a cru que $Z$ laisse $\ket-$ inchangé ; $Z$ ne laisse invariants que $\ket0$ et $\ket1$, et échange $\ket+$ et $\ket-$.}

% ---- Question 76 (niveau 3, correct : C)
\qcmq[2]{Chap.~\ref{chap:portes}\,$\cdot$\,\S\,\ref{sec:phase}}
  {Le qubit $\ket{+}$ traverse la porte $T=\mathrm{diag}\big(1,\ee^{\ii\pi/4}\big)$, puis est mesuré dans la base $X$ $\{\ket+,\ket-\}$. Quelle est la probabilité d'obtenir $+$ ?}
  {$\frac12$}
  {$\frac{2-\sqrt2}{4}\approx0{,}15$}
  {$\frac{2+\sqrt2}{4}\approx0{,}85$}
  {$\frac{\sqrt2}{2}\approx0{,}71$}
  {C}
\qcmx
  {On a cru que la phase ne change pas les probabilités : c'est vrai en base $Z$, mais la base $X$ est sensible à la phase \emph{relative} (le point de Bloch est à l'azimut $\frac\pi4$ de l'axe $+x$).}
  {C'est $P(-)=\sin^2\frac\pi8$ : on a permuté les résultats $+$ et $-$.}
  {$T\ket+=\frac1{\sqrt2}(\ket0+\ee^{\ii\pi/4}\ket1)$, donc $\braket{+}{\psi}=\frac12\big(1+\ee^{\ii\pi/4}\big)$ et $P(+)=\frac14\big|1+\ee^{\ii\pi/4}\big|^2=\frac{2+2\cos\frac\pi4}{4}=\frac{2+\sqrt2}{4}$ (soit $\cos^2\frac\pi8$).}
  {On a pris $\cos\frac\pi4$ comme probabilité ; il faut le \emph{demi}-angle et le carré : $P=\cos^2\frac{\pi/4}{2}=\cos^2\frac\pi8$.}

% ---- Question 77 (niveau 3, correct : B)
\qcmq[2]{Chap.~\ref{chap:portes}\,$\cdot$\,\S\,\ref{sec:hadamard}}
  {Que vaut le produit de matrices $HYH$, où $Y=\begin{pmatrix}0&-\ii\\\ii&0\end{pmatrix}$ ?}
  {$Y$}
  {$-Y$}
  {$Z$}
  {$X$}
  {B}
\qcmx
  {Par analogie avec $HXH=Z$, on a cru que $H$ ne modifie pas $Y$ ; en fait l'axe $y$ est retourné (signe $-$).}
  {$H$ échange les axes $x$ et $z$ et retourne l'axe $y$ : $HXH=Z$, $HZH=X$, $HYH=-Y$. Direct : $HY=\frac{1}{\sqrt2}\begin{pmatrix}\ii&-\ii\\-\ii&-\ii\end{pmatrix}$, puis $HYH=\begin{pmatrix}0&\ii\\-\ii&0\end{pmatrix}=-Y$.}
  {C'est $HXH$ : $H$ n'échange que les axes $x$ et $z$, l'axe $y$ reste l'axe $y$.}
  {C'est $HZH$ : $Y$ ne peut pas devenir $X$, car $H$ laisse l'axe $y$ à sa place (au signe près).}

% ---- Question 78 (niveau 3, correct : D)
\qcmq[1]{Chap.~\ref{chap:portes}\,$\cdot$\,\S\,\ref{sec:op2}}
  {Que donne $(H\otimes H)\ket{01}$, où $\ket{01}=\ket0\otimes\ket1$ ?}
  {$\frac12\big(\ket{00}+\ket{01}-\ket{10}-\ket{11}\big)$}
  {$\frac12\big(\ket{00}-\ket{01}-\ket{10}+\ket{11}\big)$}
  {$\frac12\big(\ket{00}+\ket{01}+\ket{10}+\ket{11}\big)$}
  {$\frac12\big(\ket{00}-\ket{01}+\ket{10}-\ket{11}\big)$}
  {D}
\qcmx
  {C'est $\ket-\otimes\ket+$ : on a attribué le $\ket-$ au premier qubit, alors que c'est le second qui vaut $\ket1$.}
  {C'est $\ket-\otimes\ket-$ : $H\ket0=\ket+$, et non $\ket-$.}
  {C'est $\ket+\otimes\ket+$, le résultat pour $\ket{00}$ : on a oublié que $H\ket1=\ket-$ introduit des signes.}
  {$H\ket0\otimes H\ket1=\ket+\otimes\ket-=\frac12(\ket0+\ket1)\otimes(\ket0-\ket1)$. Le signe $-$ apparaît quand le \emph{second} chiffre vaut $1$ (formule générale $\frac12\sum_y(-1)^{x\cdot y}\ket y$ avec $x=01$).}

% ---- Question 79 (niveau 3, correct : D)
\qcmq[2]{Chap.~\ref{chap:portes}\,$\cdot$\,\S\,\ref{sec:cnot}}
  {Avec le CNOT du cours (contrôle : second qubit $B$ ; cible : premier qubit $A$ ; $\mathrm{CNOT}\ket{a,b}=\ket{a\oplus b,\,b}$), quel état obtient-on à partir de $\ket{1}_A\otimes\ket{+}_B$ ?}
  {$\ket{1}\otimes\ket{+}$ (inchangé)}
  {$\frac{1}{\sqrt2}\big(\ket{00}+\ket{11}\big)=\ket{\Phi^+}$}
  {$\frac{1}{\sqrt2}\big(\ket{01}-\ket{10}\big)=\ket{\Psi^-}$}
  {$\frac{1}{\sqrt2}\big(\ket{01}+\ket{10}\big)=\ket{\Psi^+}$}
  {D}
\qcmx
  {On a pris le premier qubit comme contrôle : le CNOT agirait alors sur $B$, et $X\ket+=\ket+$ laisserait l'état inchangé. Mais ici le contrôle est $B$.}
  {C'est le résultat pour $\ket{0}_A\ket{+}_B$ ; ici $A$ vaut $\ket1$, ce qui échange les deux termes.}
  {Aucun signe $-$ n'apparaît : le CNOT permute des états de base sans leur ajouter de phase.}
  {$\ket1\ket+=\frac1{\sqrt2}(\ket{10}+\ket{11})$ ; $\ket{10}\mapsto\ket{1\oplus0,0}=\ket{10}$ et $\ket{11}\mapsto\ket{1\oplus1,1}=\ket{01}$. C'est un état de Bell, donc intriqué.}

% ---- Question 80 (niveau 3, correct : C)
\qcmq[2]{Chap.~\ref{chap:portes}\,$\cdot$\,\S\,\ref{sec:swap}}
  {Que vaut $\mathrm{SWAP}\,(X\otimes\mathrm{Id})\,\mathrm{SWAP}$ ?}
  {$X\otimes\mathrm{Id}$}
  {$\mathrm{Id}\otimes\mathrm{Id}$}
  {$\mathrm{Id}\otimes X$}
  {$X\otimes X$}
  {C}
\qcmx
  {On a supposé que $\mathrm{SWAP}$ commute avec $X\otimes\mathrm{Id}$ ; ce n'est pas le cas, puisque cet opérateur n'agit pas de la même façon sur les deux qubits.}
  {On a simplifié $\mathrm{SWAP}^2=\mathrm{Id}$ en oubliant que $X\otimes\mathrm{Id}$ se trouve \emph{entre} les deux SWAP : on ne peut pas les rapprocher.}
  {$\mathrm{SWAP}\,(X\otimes\mathrm{Id})\,\mathrm{SWAP}\ket{a,b}=\mathrm{SWAP}\,(X\otimes\mathrm{Id})\ket{b,a}=\mathrm{SWAP}\ket{b\oplus1,a}=\ket{a,b\oplus1}$ : le SWAP transporte l'opérateur d'un qubit à l'autre (même mécanisme que $\mathrm{CNOT}_{21}=\mathrm{SWAP}\,\mathrm{CNOT}_{12}\,\mathrm{SWAP}$).}
  {Cet opérateur inverserait les deux bits ; ici un seul bit est inversé, celui du \emph{second} qubit.}

% ---- Question 81 (niveau 3, correct : C)
\qcmq[1]{Chap.~\ref{chap:portes}\,$\cdot$\,\S\,\ref{sec:controlees}}
  {La porte $Z$ contrôlée ($\mathrm{CZ}=\mathrm{diag}(1,1,1,-1)$, contrôle : premier qubit) est appliquée à $\ket{+}\otimes\ket{+}$. Quel état obtient-on ?}
  {$\ket{+}\otimes\ket{+}$}
  {$\ket{-}\otimes\ket{-}$}
  {$\frac12\big(\ket{00}+\ket{01}+\ket{10}-\ket{11}\big)$, un état intriqué}
  {$\ket{+}\otimes\ket{-}$}
  {C}
\qcmx
  {On a cru qu'un simple signe est une phase globale ; mais le signe $-$ ne touche que $\ket{11}$ : c'est une phase \emph{relative}.}
  {C'est $Z\otimes Z$ ($Z$ sur les deux qubits), et non une porte contrôlée.}
  {CZ ne change que le signe de $\ket{11}$. Critère de séparabilité : $c_{00}c_{11}-c_{01}c_{10}=-\frac14-\frac14=-\frac12\neq0$ ; on peut l'écrire $\frac{1}{\sqrt2}\big(\ket0\ket++\ket1\ket-\big)$.}
  {On a appliqué $Z$ au second qubit sans condition ; le contrôle exige que le premier qubit soit dans $\ket1$.}

% ---- Question 82 (niveau 3, correct : D)
\qcmq[2]{Chap.~\ref{chap:portes}\,$\cdot$\,\S\,\ref{sec:phase}, \ref{sec:portes-bloch}}
  {Un qubit préparé dans $\ket{0}$ traverse successivement $H$, puis $S$, puis $H$ (dans cet ordre). Quel est l'état final, à une phase globale près ?}
  {$\ket{{+}\ii}$}
  {$\ket{0}$}
  {$\ket{1}$}
  {$\ket{{-}\ii}=\frac{1}{\sqrt2}\big(\ket0-\ii\ket1\big)$}
  {D}
\qcmx
  {Mauvais sens de rotation : on a perdu un signe en calculant $H\ket{{+}\ii}$ (la phase relative vaut $-\ii$, pas $+\ii$).}
  {On a « simplifié » $HSH$ en $S$ (puis $S\ket0=\ket0$) : mais $H$ ne commute pas avec $S$, on ne peut pas supprimer les deux $H$.}
  {C'est le résultat pour $HZH=X$ ; or $S$ n'est qu'un \emph{quart} de tour autour de $z$ (et non un demi-tour comme $Z$), donc $HSH$ est un quart de tour autour de $x$.}
  {$HSH=\sqrt X$ est un quart de tour autour de $x$ (§\,\ref{sec:phase}). Directement : $H\ket0=\ket+$, $S\ket+=\ket{{+}\ii}$, $H\ket{{+}\ii}=\frac{1+\ii}{2}\ket0+\frac{1-\ii}{2}\ket1=\frac{1+\ii}{2}\big(\ket0-\ii\ket1\big)$.}

% ---- Question 83 (niveau 3, correct : A)
\qcmq[1]{Chap.~\ref{chap:bell}\,$\cdot$\,\S\,\ref{sec:mesure2}}
  {Deux qubits sont dans l'état $\ket{\Psi}=\frac{1}{\sqrt3}\big(\ket{00}+\ket{01}+\ket{11}\big)$. On mesure le \emph{second} qubit $B$ en base $Z$ et on obtient $1$. Que valait la probabilité de ce résultat, et dans quel état se trouve alors $A$ ?}
  {$P(B{=}1)=\frac23$ ; $A$ est dans l'état $\ket{+}$}
  {$P(B{=}1)=\frac13$ ; $A$ est dans l'état $\ket{0}$}
  {$P(B{=}1)=\frac23$ ; $A$ est dans l'état $\ket{1}$}
  {$P(B{=}1)=\frac13$ ; $A$ est dans l'état $\ket{+}$}
  {A}
\qcmx
  {$P(B{=}1)=|c_{01}|^2+|c_{11}|^2=\frac13+\frac13$. On garde les termes dont le second chiffre vaut $1$ : $\frac{1}{\sqrt3}\big(\ket{01}+\ket{11}\big)=\frac{1}{\sqrt3}(\ket0+\ket1)\otimes\ket1$, qui renormalisé donne $\ket{+}_A\otimes\ket{1}_B$.}
  {C'est le résultat $B=0$ : seul le terme $\ket{00}$ a un second chiffre nul, d'où la probabilité $\frac13$ et $A$ dans $\ket0$.}
  {On a projeté sur le seul terme $\ket{11}$ ; or $\ket{01}$ a aussi son second chiffre égal à $1$ : $A$ vaut $0$ ou $1$ (état $\ket+$).}
  {L'état de $A$ est juste, mais on n'a compté qu'un seul terme pour la probabilité ; il faut sommer $|c_{01}|^2+|c_{11}|^2$.}

% ---- Question 84 (niveau 3, correct : A)
\qcmq[2]{Chap.~\ref{chap:bell}\,$\cdot$\,\S\,\ref{sec:separable}}
  {Pour quelle valeur de $x$ l'état (non normalisé) $\ket{00}+2\ket{01}+3\ket{10}+x\ket{11}$ est-il séparable ?}
  {$x=6$}
  {$x=-6$}
  {$x=\frac16$}
  {$x=5$}
  {A}
\qcmx
  {Critère $D=c_{00}c_{11}-c_{01}c_{10}=x-2\cdot3=0$. On obtient bien un produit : $(\ket0+3\ket1)\otimes(\ket0+2\ket1)=\ket{00}+2\ket{01}+3\ket{10}+6\ket{11}$.}
  {Erreur de signe : $D=x-6$ s'annule pour $x=+6$ ; pour $x=-6$, $D=-12\neq0$ et l'état est intriqué.}
  {On a inversé le rapport : $c_{00}c_{11}=c_{01}c_{10}$ donne $x=\frac{c_{01}c_{10}}{c_{00}}=6$, et non $\frac{c_{00}}{c_{01}c_{10}}$.}
  {On a additionné ($2+3$) au lieu de multiplier : le critère de séparabilité porte sur des \emph{produits} de coefficients.}

% ---- Question 85 (niveau 3, correct : B)
\qcmq[2]{Chap.~\ref{chap:bell}\,$\cdot$\,\S\,\ref{sec:bell}}
  {Comment s'écrit $\ket{00}$ dans la base de Bell ?}
  {$\frac{1}{\sqrt2}\big(\ket{\Phi^+}-\ket{\Phi^-}\big)$}
  {$\frac{1}{\sqrt2}\big(\ket{\Phi^+}+\ket{\Phi^-}\big)$}
  {$\frac{1}{\sqrt2}\big(\ket{\Phi^+}+\ket{\Psi^+}\big)$}
  {$\ket{\Phi^+}$}
  {B}
\qcmx
  {Erreur de signe : cette combinaison vaut $\ket{11}$, puisque $\ket{\Phi^+}-\ket{\Phi^-}=\sqrt2\,\ket{11}$.}
  {$\ket{\Phi^+}+\ket{\Phi^-}=\frac{1}{\sqrt2}\cdot2\ket{00}=\sqrt2\,\ket{00}$. Une mesure de Bell sur $\ket{00}$ donne donc $\Phi^+$ ou $\Phi^-$ avec la probabilité $\frac12$ chacun (jamais $\Psi^\pm$).}
  {Cette combinaison vaut $\ket{+}\otimes\ket{+}=\frac12(\ket{00}+\ket{01}+\ket{10}+\ket{11})$ : les deux états de Bell « pairs » et « impairs » s'y ajoutent.}
  {On a confondu $\ket{00}$, état séparable, avec $\ket{\Phi^+}=\frac{1}{\sqrt2}(\ket{00}+\ket{11})$, qui contient aussi $\ket{11}$ et est intriqué.}

% ---- Question 86 (niveau 3, correct : A)
\qcmq[1]{Chap.~\ref{chap:bell}\,$\cdot$\,\S\,\ref{sec:circuit-bell}}
  {Le générateur de Bell du cours est $G=\mathrm{CNOT}\cdot(\mathrm{Id}\otimes H)$ : $H$ sur le second qubit $B$, puis CNOT de contrôle $B$ et de cible $A$. Quel circuit permet de \emph{mesurer} une paire dans la base de Bell (on le place avant une mesure en base $Z$) ?}
  {D'abord le CNOT (contrôle $B$, cible $A$), puis $H$ sur $B$}
  {D'abord $H$ sur $B$, puis le CNOT}
  {D'abord le CNOT, puis $H$ sur $A$ (la cible)}
  {Aucun circuit : on mesure directement les deux qubits en base $Z$}
  {A}
\qcmx
  {C'est l'inverse $G^\dagger=(\mathrm{Id}\otimes H)\,\mathrm{CNOT}$ : les portes sont leurs propres inverses et l'ordre s'inverse. Il envoie $\Phi^+,\Phi^-,\Psi^+,\Psi^-$ sur $\ket{00},\ket{01},\ket{10},\ket{11}$ (à une phase près).}
  {C'est $G$ lui-même, qui \emph{fabrique} des états de Bell : on garde le même ordre au lieu de l'inverser, et la mesure ne distingue plus les quatre états.}
  {Mauvais qubit : pour $\Phi^+$ on trouve $\ket{+}\otimes\ket{+}$, dont les deux bits sont aléatoires. Le $H$ doit agir sur $B$, qui a reçu le $H$ dans $G$.}
  {$\Phi^+$ et $\Phi^-$ (comme $\Psi^+$ et $\Psi^-$) donnent les mêmes résultats en base $Z$ : cette base ne voit pas le signe relatif, il faut d'abord défaire l'intrication.}

% ---- Question 87 (niveau 3, correct : B)
\qcmq[2]{Chap.~\ref{chap:bell}\,$\cdot$\,\S\,\ref{sec:trace}}
  {Quelle est la matrice de densité réduite $\rho_A=\operatorname{Tr}_B\ket\Psi\bra\Psi$ du premier qubit pour $\ket{\Psi}=\frac12\ket{00}+\frac{\sqrt3}{2}\ket{11}$ ?}
  {$\begin{pmatrix}\frac14&\frac{\sqrt3}{4}\\\frac{\sqrt3}{4}&\frac34\end{pmatrix}$}
  {$\begin{pmatrix}\frac14&0\\0&\frac34\end{pmatrix}$}
  {$\begin{pmatrix}\frac12&0\\0&\frac12\end{pmatrix}$}
  {$\begin{pmatrix}\frac12&0\\0&\frac{\sqrt3}{2}\end{pmatrix}$}
  {B}
\qcmx
  {C'est la matrice $\ket\chi\bra\chi$ de l'état pur $\ket\chi=\frac12\ket0+\frac{\sqrt3}2\ket1$ : on a gardé les cohérences, comme si $A$ était seul ; l'intrication avec $B$ les détruit.}
  {Avec $C=\begin{pmatrix}c_{00}&c_{01}\\c_{10}&c_{11}\end{pmatrix}=\begin{pmatrix}\frac12&0\\0&\frac{\sqrt3}2\end{pmatrix}$, $\rho_A=CC^\dagger=\mathrm{diag}\big(\frac14,\frac34\big)$ : état mixte, de pureté $\frac{10}{16}=1-2|D|^2$ avec $D=\frac{\sqrt3}{4}$.}
  {C'est la réponse pour un état de Bell (intrication maximale) ; ici les poids $\frac14$ et $\frac34$ diffèrent : l'intrication est partielle.}
  {On a mis les amplitudes au lieu de leurs carrés : la trace vaudrait $\frac{1+\sqrt3}{2}\neq1$.}

% ---- Question 88 (niveau 3, correct : A)
\qcmq[2]{Chap.~\ref{chap:bell}\,$\cdot$\,\S\,\ref{sec:info}}
  {Deux qubits sont dans le mélange statistique $\rho_{AB}=\frac12\ket{00}\bra{00}+\frac12\ket{11}\bra{11}$ (corrélations purement classiques). Que vaut $H(A|B)=S(\rho_{AB})-S(\rho_B)$, en bits ?}
  {$0$}
  {$-1$}
  {$1$}
  {$2$}
  {A}
\qcmx
  {$\rho_{AB}$ a pour valeurs propres $\frac12,\frac12,0,0$, donc $S(\rho_{AB})=1$ ; et $\rho_B=\frac12\mathrm{Id}$ donne $S(\rho_B)=1$. Ainsi $H(A|B)=0$ : connaître $B$ fixe $A$. Une valeur négative est propre à l'intrication.}
  {C'est la valeur pour l'état de Bell \emph{pur} $\ket{\Phi^+}$ : mêmes résultats en base $Z$, mais $S(\rho_{AB})=0$ car l'état global est pur. Ici l'état est mixte.}
  {On a pris $H(A|B)=H(A)$, comme si $B$ ne disait rien sur $A$ ; or les deux bits sont toujours égaux.}
  {On a additionné les entropies des deux qubits ($1+1$), valeur de deux bits indépendants : ce n'est pas une entropie conditionnelle.}

% ---- Question 89 (niveau 3, correct : B)
\qcmq[1]{Chap.~\ref{chap:bell}\,$\cdot$\,\S\,\ref{sec:teleportation}}
  {Dans la téléportation, avant de recevoir les deux bits $(a,q)$ d'Alice, dans quel état se trouve le qubit de Bob ?}
  {Dans l'état pur $\alpha\ket0+\beta\ket1$, déjà téléporté}
  {Dans l'état complètement mélangé $\rho_B=\frac12\mathrm{Id}$, indépendant de $\alpha$ et $\beta$}
  {Dans le mélange $|\alpha|^2\ket0\bra0+|\beta|^2\ket1\bra1$}
  {Dans l'état pur $\ket{0}$, son état initial}
  {B}
\qcmx
  {Ce serait une communication instantanée. L'état $\alpha\ket0+\beta\ket1$ n'apparaît qu'\emph{après} la correction, qui dépend de $(a,q)$.}
  {Les quatre kets $\ket{\phi_{aq}}$ sont équiprobables ($\frac14$) et $\frac14\sum_{a,q}\ket{\phi_{aq}}\bra{\phi_{aq}}=\frac12\mathrm{Id}$. Bob n'apprend rien avant le message classique : rien ne voyage plus vite que la lumière.}
  {Bob connaîtrait alors les statistiques de $\ket\varphi$ sans aucun message : or les probabilités d'Alice ($\frac14$ pour chaque résultat) ne dépendent pas de $\alpha,\beta$.}
  {Le qubit de Bob est la moitié de la paire $\ket{\Phi^+}$ : pris seul, il n'a pas d'état pur ($\rho_B=\frac12\mathrm{Id}$ avant comme après la mesure d'Alice, tant que Bob ignore le résultat).}

% ---- Question 90 (niveau 3, correct : D)
\qcmq[1]{Chap.~\ref{chap:bell}\,$\cdot$\,\S\,\ref{sec:chsh}}
  {Au jeu CHSH (gain si $a\oplus b=x\cdot y$), quelles sont les meilleures probabilités de gain avec une stratégie classique, même avec du hasard partagé, et avec une paire $\ket{\Phi^+}$ mesurée dans des bases bien choisies ?}
  {$\frac12$ en classique ; $1$ avec l'intrication}
  {$\frac34$ dans les deux cas}
  {$1$ en classique si Alice et Bob s'accordent à l'avance ; aucun avantage quantique}
  {$\frac34$ en classique ; $\cos^2\frac\pi8\approx0{,}85$ avec l'intrication}
  {D}
\qcmx
  {Classique : répondre toujours $0$ donne déjà $\frac34$. Quantique : l'intrication ne permet pas de gagner à tous les coups (borne de Tsirelson $\cos^2\frac\pi8$).}
  {Si l'intrication n'apportait rien, on n'aurait pas de violation de l'inégalité de Bell ; l'avantage quantique ($0{,}85>0{,}75$) est précisément ce que l'expérience de Bell met en évidence.}
  {S'accorder à l'avance (ou partager du hasard) ne suffit pas : aucune stratégie déterministe ne gagne les quatre parties, d'après l'argument de parité.}
  {Une stratégie déterministe perd au moins une des quatre parties (la somme modulo $2$ des conditions donne $0=1$) ; avec $\Phi^+$, chaque partie est gagnée avec la probabilité $\cos^2\frac\pi8=\frac{2+\sqrt2}{4}$.}

\clearpage
%% ======================================================================
\subsection{Niveau 4 : maîtrise (questions 91 à 100)}\label{sec:qcm-n4}
%% ======================================================================
\noindent{\color{insagris}\small Niveau de difficulté : Maîtrise\quad $\bullet\bullet\bullet\bullet$}\par\smallskip
% ---- Question 91 (niveau 4, correct : C)
\qcmq[2]{Chap.~\ref{chap:histoire}\,$\cdot$\,\S\,\ref{sec:superposition}}
  {Une particule du puits infini est préparée à $t=0$ dans l'état $\Psi=\frac{1}{\sqrt2}\big(\psi_1+\psi_2\big)$, avec $E_n=n^2E_1$. Au bout de quelle durée minimale la densité de probabilité $|\Psi(x,t)|^2$ retrouve-t-elle sa forme initiale ?}
  {$T=\dfrac{2\pi\hbar}{E_1}$}
  {$T=\dfrac{2\pi\hbar}{5E_1}$}
  {$T=\dfrac{2\pi\hbar}{3E_1}$}
  {$T=\dfrac{\pi\hbar}{3E_1}$}
  {C}
\qcmx
  {C'est la période de la phase $\ee^{-\ii E_1t/\hbar}$ de $\psi_1$ seul ; or les phases globales ne se voient pas dans $|\Psi|^2$ : seul compte l'écart $E_2-E_1=3E_1$.}
  {On a pris la somme $E_1+E_2=5E_1$ ; c'est la \emph{différence} des énergies (phase relative) qui intervient.}
  {$|\Psi|^2=\frac12(\psi_1^2+\psi_2^2)+\psi_1\psi_2\cos\frac{(E_2-E_1)t}{\hbar}$ : seul le terme croisé dépend du temps, à la pulsation $\omega=\frac{E_2-E_1}{\hbar}=\frac{3E_1}{\hbar}$, donc $T=\frac{2\pi}{\omega}$.}
  {C'est la demi-période : à $t=\frac T2$ le cosinus vaut $-1$ et la densité est le symétrique de l'état initial ($x\to a-x$), pas encore sa forme de départ.}

% ---- Question 92 (niveau 4, correct : D)
\qcmq[1]{Chap.~\ref{chap:dirac}\,$\cdot$\,\S\,\ref{sec:projecteurs}, \ref{sec:qubit}}
  {Un expérimentateur reçoit un qubit préparé soit dans $\ket{0}$, soit dans $\ket{+}$. Laquelle de ces affirmations est correcte ?}
  {Une mesure en base $Z$ les distingue : $\ket0$ donne $0$ et $\ket+$ donne $1$}
  {Une mesure en base $X$ les distingue : $\ket+$ donne $+$ et $\ket0$ donne $-$}
  {On les distingue en répétant la mesure sur le même qubit}
  {Aucune mesure ne permet de les distinguer à coup sûr, car $\braket{0}{+}=\frac{1}{\sqrt2}\neq0$}
  {D}
\qcmx
  {$\ket+$ donne $0$ ou $1$ avec la probabilité $\frac12$ : le résultat $0$ est possible pour les deux états.}
  {$\ket0=\frac{1}{\sqrt2}(\ket++\ket-)$ donne $+$ ou $-$ avec la probabilité $\frac12$ : le résultat $+$ est possible pour les deux états.}
  {Après une première mesure l'état est effondré : répéter ne donne aucune information nouvelle, et on ne peut pas copier le qubit pour recommencer (non-clonage).}
  {Seuls des états \emph{orthogonaux} sont parfaitement discernables. En base $Z$, $\ket0$ donne toujours $0$ et $\ket+$ donne $0$ une fois sur deux : obtenir $0$ ne tranche pas (obtenir $1$ prouve seulement qu'on avait $\ket+$).}

% ---- Question 93 (niveau 4, correct : C)
\qcmq[1]{Chap.~\ref{chap:operateurs}\,$\cdot$\,\S\,\ref{sec:pur-mixte}}
  {Une source envoie au hasard $\ket{0}$ ou $\ket{1}$ (probabilité $\frac12$ chacun) ; une autre envoie au hasard $\ket{+}$ ou $\ket{-}$ (probabilité $\frac12$ chacun). Peut-on distinguer ces deux sources en mesurant les qubits émis ?}
  {Oui : en base $Z$, la première source donne des résultats prévisibles}
  {Oui : en base $X$, la seconde source donne des résultats prévisibles}
  {Non : les deux ont la même matrice de densité $\frac{\mathrm{Id}}{2}$, donc les mêmes statistiques pour toute mesure}
  {Oui : les états purs envoyés ne sont pas les mêmes}
  {C}
\qcmx
  {Chaque qubit est bien $\ket0$ ou $\ket1$, mais on ignore lequel : le résultat vaut $0$ ou $1$ avec la probabilité $\frac12$, exactement comme avec la seconde source ($|\braket{0}{\pm}|^2=\frac12$).}
  {Même erreur : chaque qubit est $\ket+$ ou $\ket-$, mais on ignore lequel ; le résultat en base $X$ est aléatoire, et la première source donne aussi $+$ ou $-$ avec la probabilité $\frac12$.}
  {$\frac12\ket0\bra0+\frac12\ket1\bra1=\frac12\ket+\bra++\frac12\ket-\bra-=\frac{\mathrm{Id}}2$ (relation de fermeture dans chaque base). Les probabilités de toute mesure ne dépendent que de $\rho$.}
  {Ce qui compte n'est pas la liste des états purs envoyés mais la matrice de densité moyenne : des préparations différentes peuvent avoir la même $\rho$.}

% ---- Question 94 (niveau 4, correct : B)
\qcmq[1]{Chap.~\ref{chap:operateurs}\,$\cdot$\,\S\,\ref{sec:pur-mixte} ; Chap.~\ref{chap:portes}\,$\cdot$\,\S\,\ref{sec:portes-bloch}}
  {Un qubit est dans l'état complètement mélangé $\rho=\frac{\mathrm{Id}}{2}$. Existe-t-il une porte quantique $U$ qui le transforme en un état pur, par exemple $\ket{0}\bra{0}$ ?}
  {Oui : la porte de Hadamard, qui crée des superpositions}
  {Non : $\rho\mapsto U\rho U^\dagger$ conserve les valeurs propres, donc la pureté ; ici $U\frac{\mathrm{Id}}{2}U^\dagger=\frac{\mathrm{Id}}{2}$}
  {Oui : une rotation de Bloch qui amène le centre de la boule au pôle nord}
  {Oui : une mesure en base $Z$, qui donne $\ket0$}
  {B}
\qcmx
  {$H$ envoie un état pur sur un état pur ($\ket0\mapsto\ket+$) et laisse $\frac{\mathrm{Id}}{2}$ invariant : $H\frac{\mathrm{Id}}2H=\frac{\mathrm{Id}}2$ car $H^2=\mathrm{Id}$.}
  {$\operatorname{Tr}\big[(U\rho U^\dagger)^2\big]=\operatorname{Tr}\rho^2$ ($=\frac12$ ici). Sur la boule de Bloch, une porte est une rotation : elle laisse le centre fixe.}
  {Une rotation tourne la boule autour de son centre, qui reste fixe : elle ne peut pas l'amener sur la sphère.}
  {Une mesure n'est pas une porte (elle est non unitaire) ; de plus elle donne $\ket0$ ou $\ket1$ au hasard, pas $\ket0$ à coup sûr.}

% ---- Question 95 (niveau 4, correct : D)
\qcmq[1]{Chap.~\ref{chap:portes}\,$\cdot$\,\S\,\ref{sec:cnot}}
  {Que vaut $(H\otimes H)\,\mathrm{CNOT}_{12}\,(H\otimes H)$, où $\mathrm{CNOT}_{12}$ a pour contrôle le \emph{premier} qubit et pour cible le second ?}
  {$\mathrm{CNOT}_{12}$, inchangé}
  {$\mathrm{CZ}$ (porte $Z$ contrôlée)}
  {$\mathrm{SWAP}$}
  {$\mathrm{CNOT}_{21}$ (contrôle : second qubit ; cible : premier qubit)}
  {D}
\qcmx
  {On a cru que les deux $H$ s'annulent ($H^2=\mathrm{Id}$) ; mais ils encadrent le CNOT au lieu d'être côte à côte : on ne peut pas les simplifier.}
  {C'est le résultat quand $H$ n'encadre que la \emph{cible} : $(\mathrm{Id}\otimes H)\,\mathrm{CNOT}_{12}\,(\mathrm{Id}\otimes H)=\mathrm{CZ}$, car $HXH=Z$. Ici $H$ agit aussi sur le contrôle.}
  {Le SWAP s'obtient avec \emph{trois} CNOT alternés, pas avec un seul CNOT conjugué par $H\otimes H$ ; il n'y a pas de contrôle dans un SWAP.}
  {$H\ket0\bra0H=\ket+\bra+$, $H\ket1\bra1H=\ket-\bra-$ et $HXH=Z$ donnent $\ket+\bra+\otimes I+\ket-\bra-\otimes Z=I\otimes\ket0\bra0+X\otimes\ket1\bra1=\mathrm{CNOT}_{21}$ : dans la base $\{\ket+,\ket-\}$, contrôle et cible échangent leurs rôles (rebond de phase).}

% ---- Question 96 (niveau 4, correct : A)
\qcmq[1]{Chap.~\ref{chap:portes}\,$\cdot$\,\S\,\ref{sec:controlees}}
  {On applique la porte de Toffoli (contrôles : qubits 1 et 2 ; cible : qubit 3) à $\ket{+}\otimes\ket{+}\otimes\ket{0}$, puis on mesure le troisième qubit en base $Z$. Que peut-on dire si l'on obtient $1$ ?}
  {Ce résultat a la probabilité $\frac14$, et les deux premiers qubits sont alors dans l'état $\ket{11}$}
  {Ce résultat a la probabilité $\frac12$, et les deux premiers qubits restent dans $\ket{+}\ket{+}$}
  {Ce résultat a la probabilité $\frac14$, mais les deux premiers qubits restent dans $\ket{+}\ket{+}$}
  {Ce résultat a la probabilité $\frac34$, et les deux premiers qubits sont dans l'état $\frac{1}{\sqrt3}\big(\ket{00}+\ket{01}+\ket{10}\big)$}
  {A}
\qcmx
  {$\mathrm{CCNOT}\ket{x,y,0}=\ket{x,y,xy}$ donne $\frac12\big(\ket{000}+\ket{010}+\ket{100}+\ket{111}\big)$ : seul $\ket{111}$ a un troisième chiffre égal à $1$ (probabilité $\frac14$), et il impose $x=y=1$.}
  {La cible n'est inversée que si $x=y=1$ (probabilité $\frac14$) et non une fois sur deux ; de plus l'état est intriqué, donc mesurer le qubit 3 modifie les qubits 1 et 2.}
  {La probabilité est juste, mais la mesure de la cible n'est pas sans effet sur les contrôles : l'état après la Toffoli est intriqué, donc le résultat $1$ les projette sur $\ket{11}$.}
  {C'est le cas du résultat $0$ (probabilité $\frac34$ : $x,y\in\{00,01,10\}$). Le résultat $1$ est le cas rare, celui du ET vrai ($x=y=1$).}

% ---- Question 97 (niveau 4, correct : A)
\qcmq[1]{Chap.~\ref{chap:bell}\,$\cdot$\,\S\,\ref{sec:bell}}
  {Que donne $(H\otimes H)\ket{\Psi^-}$, avec $\ket{\Psi^-}=\frac{1}{\sqrt2}\big(\ket{01}-\ket{10}\big)$ ?}
  {$-\ket{\Psi^-}$, c'est-à-dire le même état physique}
  {$\ket{\Phi^-}$}
  {$\ket{\Psi^+}$}
  {$\ket{\Phi^+}$}
  {A}
\qcmx
  {Le singulet est invariant par $U\otimes U$ à un facteur près : $(U\otimes U)\ket{\Psi^-}=\det(U)\ket{\Psi^-}$ avec $\det H=-1$. Directement : $\frac{1}{\sqrt2}\big(\ket+\ket--\ket-\ket+\big)=-\ket{\Psi^-}$.}
  {C'est l'image de $\ket{\Psi^+}$, pas de $\ket{\Psi^-}$ : $(H\otimes H)\ket{\Psi^+}=\ket{\Phi^-}$. Le signe $-$ du singulet change le calcul (les termes $\ket{00}$ et $\ket{11}$ se compensent).}
  {$\ket{\Psi^+}$ est symétrique par échange des deux qubits, $\ket{\Psi^-}$ antisymétrique ; $H\otimes H$ commute avec le SWAP et conserve donc cette symétrie.}
  {C'est l'état que $H\otimes H$ laisse invariant ($\ket{\Phi^+}=\frac{1}{\sqrt2}(\ket{++}+\ket{--})$) ; le singulet n'a pas ce comportement.}

% ---- Question 98 (niveau 4, correct : B)
\qcmq[1]{Chap.~\ref{chap:bell}\,$\cdot$\,\S\,\ref{sec:chsh}}
  {Une expérience de type CHSH (avec $S=E_{00}+E_{01}+E_{10}-E_{11}$) mesure $S=2{,}6$. Que peut-on en conclure ?}
  {C'est impossible : on a toujours $|S|\leq2$}
  {Aucun modèle local à résultats prédéterminés ne l'explique, mais c'est compatible avec la mécanique quantique}
  {C'est impossible en mécanique quantique, car $2\sqrt2\approx2{,}4<2{,}6$}
  {Les particules s'échangent de l'information plus vite que la lumière}
  {B}
\qcmx
  {Cette borne ne vaut que pour les modèles classiques (locaux) ; l'expérience et la mécanique quantique la dépassent, jusqu'à $2\sqrt2$.}
  {Classique : $|S|\leq2$ ; quantique : $|S|\leq2\sqrt2\approx2{,}83$ (borne de Tsirelson). Comme $2<2{,}6<2\sqrt2$, l'inégalité de Bell est violée sans dépasser le maximum quantique ; le gain correspondant est $\frac12+\frac{2{,}6}{8}\approx0{,}825$.}
  {Erreur numérique : $2\sqrt2\approx2{,}83$, donc $2{,}6$ est permis.}
  {Les corrélations violent l'inégalité de Bell mais ne permettent aucun signal : la réponse de Bob est uniforme quelle que soit la base choisie par Alice.}

% ---- Question 99 (niveau 4, correct : A)
\qcmq[2]{Chap.~\ref{chap:bell}\,$\cdot$\,\S\,\ref{sec:info}}
  {Pour l'état $\ket{\Psi}=\frac{1}{\sqrt3}\big(\ket{00}+\sqrt2\,\ket{11}\big)$, quelle est l'entropie de von Neumann $S(\rho_A)$ du premier qubit, en bits ?}
  {$\log_2 3-\frac23\approx0{,}92$}
  {$1$}
  {$0$}
  {$0{,}64$}
  {A}
\qcmx
  {$\rho_A=\mathrm{diag}\big(\frac13,\frac23\big)$ (trace partielle), donc $S=-\frac13\log_2\frac13-\frac23\log_2\frac23=\log_23-\frac23\approx1{,}585-0{,}667$. C'est moins que le maximum d'un bit (état de Bell) : l'intrication est partielle.}
  {C'est la valeur pour un état de Bell ($\rho_A=\frac{\mathrm{Id}}{2}$) ; ici les poids $\frac13$ et $\frac23$ diffèrent, donc $S<1$.}
  {$S(\rho_A)=0$ vaut pour un état séparable ; ici $D=\frac{\sqrt2}{3}\neq0$ : l'état est intriqué, $\rho_A$ est mixte et $S>0$ (le couple $AB$ est pur, pas $A$ seul).}
  {On a utilisé le logarithme \emph{népérien} (résultat en nats) ; en bits il faut $\log_2$, ce qui donne $0{,}92$.}

% ---- Question 100 (niveau 4, correct : B)
\qcmq[1]{Chap.~\ref{chap:bell}\,$\cdot$\,\S\,\ref{sec:bell}}
  {Pourquoi peut-on mesurer simultanément les deux « parités » $Z\otimes Z$ et $X\otimes X$ de deux qubits, alors qu'on ne peut pas mesurer simultanément $Z$ et $X$ sur un seul qubit ?}
  {Parce que l'incertitude de Heisenberg ne s'applique pas à deux qubits}
  {Parce que $Z\otimes Z$ et $X\otimes X$ commutent : les deux signes $-1$ de $ZX=-XZ$ se compensent}
  {Parce que les deux qubits sont indépendants}
  {Parce que $Z\otimes Z$ et $X\otimes X$ agissent sur des qubits différents}
  {B}
\qcmx
  {Elle s'applique toujours : $Z\otimes\mathrm{Id}$ et $X\otimes\mathrm{Id}$ ne commutent pas. C'est la commutation de ces deux opérateurs précis qui permet la mesure simultanée.}
  {$(Z\otimes Z)(X\otimes X)=ZX\otimes ZX=(-XZ)\otimes(-XZ)=XZ\otimes XZ=(X\otimes X)(Z\otimes Z)$. Ils ont donc une base propre commune : la base de Bell.}
  {Dans un état de Bell les qubits sont intriqués, donc loin d'être indépendants ; la commutation est une propriété des opérateurs, pas de l'état.}
  {Chacun des deux opérateurs agit sur \emph{les deux} qubits à la fois. (Ce sont $Z\otimes\mathrm{Id}$ et $\mathrm{Id}\otimes X$, sur des qubits différents, qui commutent pour cette raison.)}

\label{qcm:fin}

\clearpage
\ifqcmcorrige

% ======================================================================
\subsection{Clé des réponses et fiche de score}\label{sec:qcm-cle}
% ======================================================================
\noindent\textbf{Cette partie donne les réponses : ne la lisez qu'après avoir répondu à toutes les
questions.}

\paragraph{Clé des réponses.} La lettre de la bonne réponse de la question $n$ se lit à la ligne
qui contient $n$, dans la colonne du rang de $n$ dans la ligne (question~24 : ligne 21--30,
colonne 4).
\begin{center}
\renewcommand{\arraystretch}{1.25}
\setlength{\tabcolsep}{5.5pt}
\begin{tabular}{@{}l*{10}{c}@{}}
\toprule
 & 1 & 2 & 3 & 4 & 5 & 6 & 7 & 8 & 9 & 10 \\
\midrule
1--10 & \qcmkey{1} & \qcmkey{2} & \qcmkey{3} & \qcmkey{4} & \qcmkey{5} & \qcmkey{6} & \qcmkey{7} & \qcmkey{8} & \qcmkey{9} & \qcmkey{10} \\
11--20 & \qcmkey{11} & \qcmkey{12} & \qcmkey{13} & \qcmkey{14} & \qcmkey{15} & \qcmkey{16} & \qcmkey{17} & \qcmkey{18} & \qcmkey{19} & \qcmkey{20} \\
21--30 & \qcmkey{21} & \qcmkey{22} & \qcmkey{23} & \qcmkey{24} & \qcmkey{25} & \qcmkey{26} & \qcmkey{27} & \qcmkey{28} & \qcmkey{29} & \qcmkey{30} \\
31--40 & \qcmkey{31} & \qcmkey{32} & \qcmkey{33} & \qcmkey{34} & \qcmkey{35} & \qcmkey{36} & \qcmkey{37} & \qcmkey{38} & \qcmkey{39} & \qcmkey{40} \\
41--50 & \qcmkey{41} & \qcmkey{42} & \qcmkey{43} & \qcmkey{44} & \qcmkey{45} & \qcmkey{46} & \qcmkey{47} & \qcmkey{48} & \qcmkey{49} & \qcmkey{50} \\
51--60 & \qcmkey{51} & \qcmkey{52} & \qcmkey{53} & \qcmkey{54} & \qcmkey{55} & \qcmkey{56} & \qcmkey{57} & \qcmkey{58} & \qcmkey{59} & \qcmkey{60} \\
61--70 & \qcmkey{61} & \qcmkey{62} & \qcmkey{63} & \qcmkey{64} & \qcmkey{65} & \qcmkey{66} & \qcmkey{67} & \qcmkey{68} & \qcmkey{69} & \qcmkey{70} \\
71--80 & \qcmkey{71} & \qcmkey{72} & \qcmkey{73} & \qcmkey{74} & \qcmkey{75} & \qcmkey{76} & \qcmkey{77} & \qcmkey{78} & \qcmkey{79} & \qcmkey{80} \\
81--90 & \qcmkey{81} & \qcmkey{82} & \qcmkey{83} & \qcmkey{84} & \qcmkey{85} & \qcmkey{86} & \qcmkey{87} & \qcmkey{88} & \qcmkey{89} & \qcmkey{90} \\
91--100 & \qcmkey{91} & \qcmkey{92} & \qcmkey{93} & \qcmkey{94} & \qcmkey{95} & \qcmkey{96} & \qcmkey{97} & \qcmkey{98} & \qcmkey{99} & \qcmkey{100} \\
\bottomrule
\end{tabular}
\end{center}

\paragraph{Fiche de score.} Reportez le nombre de bonnes réponses de chaque niveau.
\begin{center}
\renewcommand{\arraystretch}{1.6}
\begin{tabular}{@{}l l c c c@{}}
\toprule
\textbf{Niveau} & \textbf{Questions} & \textbf{Bonnes réponses} & \textbf{Sur} & \textbf{Pourcentage}\\
\midrule
1. Débutant & 1--30 & \makebox[2.6cm]{\dotfill} & 30 & \makebox[2cm]{\dotfill}\,\%\\
2. Intermédiaire & 31--60 & \makebox[2.6cm]{\dotfill} & 30 & \makebox[2cm]{\dotfill}\,\%\\
3. Avancé & 61--90 & \makebox[2.6cm]{\dotfill} & 30 & \makebox[2cm]{\dotfill}\,\%\\
4. Maîtrise & 91--100 & \makebox[2.6cm]{\dotfill} & 10 & \makebox[2cm]{\dotfill}\,\%\\
\midrule
\textbf{Total} & 1--100 & \makebox[2.6cm]{\dotfill} & 100 & \makebox[2cm]{\dotfill}\,\%\\
\bottomrule
\end{tabular}
\end{center}
Repères : au moins $80\,\%$ à un niveau, il est acquis et on peut passer au suivant ; entre $50$ et
$80\,\%$, revoyez les sections des questions manquées ; en dessous, relisez le chapitre concerné et
son récapitulatif avant de recommencer.

\paragraph{Où revoir ?} Questions classées par chapitre du cours.
\begin{center}
\renewcommand{\arraystretch}{1.3}
\begin{tabular}{@{}l >{\raggedright\arraybackslash}p{5.9cm} >{\raggedright\arraybackslash}p{6.9cm}@{}}
\toprule
\textbf{Chapitre} & \textbf{Titre} & \textbf{Questions}\\
\midrule
Chap.~\ref{chap:histoire} & Histoire et mécanique quantique & 1, 2, 3, 4, 5, 6, 31, 32, 33, 34, 35, 61, 62, 63, 91 \\
Chap.~\ref{chap:dirac} & Notation de Dirac et espace de Hilbert & 7, 8, 9, 10, 11, 12, 36, 37, 38, 39, 40, 64, 65, 66, 67, 68, 92 \\
Chap.~\ref{chap:operateurs} & Opérateurs, matrice de densité, Bloch, produit tensoriel & 13, 14, 15, 16, 17, 18, 41, 42, 43, 44, 45, 46, 69, 70, 71, 72, 73, 74, 75, 93, 94 \\
Chap.~\ref{chap:portes} & Portes quantiques et systèmes multi-qubits & 19, 20, 21, 22, 23, 24, 47, 48, 49, 50, 51, 76, 77, 78, 79, 80, 81, 82, 95, 96 \\
Chap.~\ref{chap:bell} & États quantiques, intrication, Bell & 25, 26, 27, 28, 29, 30, 52, 53, 54, 55, 56, 57, 58, 59, 60, 83, 84, 85, 86, 87, 88, 89, 90, 97, 98, 99, 100 \\
\bottomrule
\end{tabular}
\end{center}

\clearpage
\subsection{Corrigé du niveau 1 : débutant (questions 1 à 30)}\label{sec:qcm-c1}
\qcmcorr{1}
\qcmcorr{2}
\qcmcorr{3}
\qcmcorr{4}
\qcmcorr{5}
\qcmcorr{6}
\qcmcorr{7}
\qcmcorr{8}
\qcmcorr{9}
\qcmcorr{10}
\qcmcorr{11}
\qcmcorr{12}
\qcmcorr{13}
\qcmcorr{14}
\qcmcorr{15}
\qcmcorr{16}
\qcmcorr{17}
\qcmcorr{18}
\qcmcorr{19}
\qcmcorr{20}
\qcmcorr{21}
\qcmcorr{22}
\qcmcorr{23}
\qcmcorr{24}
\qcmcorr{25}
\qcmcorr{26}
\qcmcorr{27}
\qcmcorr{28}
\qcmcorr{29}
\qcmcorr{30}

\clearpage
\subsection{Corrigé du niveau 2 : intermédiaire (questions 31 à 60)}\label{sec:qcm-c2}
\qcmcorr{31}
\qcmcorr{32}
\qcmcorr{33}
\qcmcorr{34}
\qcmcorr{35}
\qcmcorr{36}
\qcmcorr{37}
\qcmcorr{38}
\qcmcorr{39}
\qcmcorr{40}
\qcmcorr{41}
\qcmcorr{42}
\qcmcorr{43}
\qcmcorr{44}
\qcmcorr{45}
\qcmcorr{46}
\qcmcorr{47}
\qcmcorr{48}
\qcmcorr{49}
\qcmcorr{50}
\qcmcorr{51}
\qcmcorr{52}
\qcmcorr{53}
\qcmcorr{54}
\qcmcorr{55}
\qcmcorr{56}
\qcmcorr{57}
\qcmcorr{58}
\qcmcorr{59}
\qcmcorr{60}

\clearpage
\subsection{Corrigé du niveau 3 : avancé (questions 61 à 90)}\label{sec:qcm-c3}
\qcmcorr{61}
\qcmcorr{62}
\qcmcorr{63}
\qcmcorr{64}
\qcmcorr{65}
\qcmcorr{66}
\qcmcorr{67}
\qcmcorr{68}
\qcmcorr{69}
\qcmcorr{70}
\qcmcorr{71}
\qcmcorr{72}
\qcmcorr{73}
\qcmcorr{74}
\qcmcorr{75}
\qcmcorr{76}
\qcmcorr{77}
\qcmcorr{78}
\qcmcorr{79}
\qcmcorr{80}
\qcmcorr{81}
\qcmcorr{82}
\qcmcorr{83}
\qcmcorr{84}
\qcmcorr{85}
\qcmcorr{86}
\qcmcorr{87}
\qcmcorr{88}
\qcmcorr{89}
\qcmcorr{90}

\clearpage
\subsection{Corrigé du niveau 4 : maîtrise (questions 91 à 100)}\label{sec:qcm-c4}
\qcmcorr{91}
\qcmcorr{92}
\qcmcorr{93}
\qcmcorr{94}
\qcmcorr{95}
\qcmcorr{96}
\qcmcorr{97}
\qcmcorr{98}
\qcmcorr{99}
\qcmcorr{100}
\fi
) :
%    \qcmrefsfalse     masque la référence « Chap. · § » sous chaque énoncé ;
%    \qcmcorrigefalse  supprime le corrigé du document (QCM à distribuer seul).
%
%  Ce fichier est produit par outils_qcm/gen_qcm.py à partir de la banque de
%  questions (fichiers texte), mais il reste un fichier LaTeX ordinaire que l'on
%  peut modifier à la main. Pour ajouter une question à la main : insérer un
%  couple \qcmq ... \qcmx dans le bon niveau (la numérotation est automatique),
%  puis mettre à jour (1) les bornes « questions a à b » des titres de niveaux,
%  (2) les lignes \qcmcorr{n} du corrigé, (3) le tableau « Où revoir ? ».
% =====================================================================

% Macros de Dirac (déjà définies dans main.tex ; ce bloc évite l'erreur
% « Undefined control sequence \ket » si main.tex est une ancienne version).
\providecommand{\ket}[1]{|#1\rangle}
\providecommand{\bra}[1]{\langle #1|}
\providecommand{\braket}[2]{\langle #1|#2\rangle}

\makeatletter
% bascules : normalement déjà définies dans main.tex (valeurs par défaut : vrai)
\@ifundefined{ifqcmrefs}{\newif\ifqcmrefs\qcmrefstrue}{}
\@ifundefined{ifqcmcorrige}{\newif\ifqcmcorrige\qcmcorrigetrue}{}
\newcounter{qcmq}

% case à cocher (3 mm) pour le crayon
\newcommand{\qcmbox}{{\setlength{\fboxsep}{0pt}\setlength{\fboxrule}{0.5pt}%
  \raisebox{-0.4mm}{\fbox{\rule{0pt}{3mm}\rule{3mm}{0pt}}}}}

% une proposition : case, lettre, texte (retrait suspendu)
\newcommand{\qcm@opt}[2]{\noindent\hangindent=3.2em \hangafter=1
  \hspace*{0.3em}\qcmbox\hspace{0.45em}\textbf{#1)}\hspace{0.4em}#2\par\vspace{2.5pt}}

% \qcmq[colonnes]{ref}{énoncé}{A}{B}{C}{D}{bonne lettre}   (TeX limite les macros à 9 paramètres)
% \qcmx{com. A}{com. B}{com. C}{com. D}  : commentaires, placés juste après \qcmq
%   - com. de la bonne lettre = explication ; com. des autres = piège ;
%   - colonnes = 1 (une proposition par ligne) ou 2 (grille 2x2).
\newcommand{\qcmq}[8][1]{%
  \stepcounter{qcmq}%
  \expandafter\gdef\csname qcm@ref@\theqcmq\endcsname{#2}%
  \expandafter\gdef\csname qcm@rep@\theqcmq\endcsname{#8}%
  \expandafter\gdef\csname qcm@o@\theqcmq @A\endcsname{#4}%
  \expandafter\gdef\csname qcm@o@\theqcmq @B\endcsname{#5}%
  \expandafter\gdef\csname qcm@o@\theqcmq @C\endcsname{#6}%
  \expandafter\gdef\csname qcm@o@\theqcmq @D\endcsname{#7}%
  \par\vspace{7pt plus 3pt minus 2pt}\noindent
  \begin{minipage}{\linewidth}%
    \noindent{\color{insarouge}\bfseries Question~\theqcmq}%
    \ifqcmrefs\hfill{\footnotesize\color{insagris}#2}\fi\par\nopagebreak
    \vspace{2pt}\noindent #3\par\nopagebreak\vspace{3pt}%
    \ifnum#1=2
      \begin{minipage}[t]{0.495\linewidth}\qcm@opt{A}{#4}\end{minipage}\hfill
      \begin{minipage}[t]{0.495\linewidth}\qcm@opt{B}{#5}\end{minipage}\par\vspace{2pt}
      \begin{minipage}[t]{0.495\linewidth}\qcm@opt{C}{#6}\end{minipage}\hfill
      \begin{minipage}[t]{0.495\linewidth}\qcm@opt{D}{#7}\end{minipage}\par
    \else
      \qcm@opt{A}{#4}\qcm@opt{B}{#5}\qcm@opt{C}{#6}\qcm@opt{D}{#7}%
    \fi
  \end{minipage}\par}

\newcommand{\qcmx}[4]{%
  \expandafter\gdef\csname qcm@c@\theqcmq @A\endcsname{#1}%
  \expandafter\gdef\csname qcm@c@\theqcmq @B\endcsname{#2}%
  \expandafter\gdef\csname qcm@c@\theqcmq @C\endcsname{#3}%
  \expandafter\gdef\csname qcm@c@\theqcmq @D\endcsname{#4}}

% lettre de la bonne réponse de la question n (« ? » si la question n'existe pas)
\newcommand{\qcmkey}[1]{\@ifundefined{qcm@rep@#1}{?}{\csname qcm@rep@#1\endcsname}}

% une mauvaise réponse du corrigé (n° de question, lettre) : ignorée si c'est la bonne.
% (pas de \@for : « := » ne fonctionne pas avec le « : » rendu actif par babel-french)
\newcommand{\qcm@wrongone}[2]{%
  \edef\qcm@a{#2}\edef\qcm@b{\csname qcm@rep@#1\endcsname}%
  \ifx\qcm@a\qcm@b\else
    \noindent\hangindent=1.4em\hangafter=1 {\color{insagris}\textbf{#2)}}\enspace
    \csname qcm@o@#1@#2\endcsname\enspace--\enspace\emph{\csname qcm@c@#1@#2\endcsname}\par
  \fi}

% bloc de corrigé de la question n
\newcommand{\qcmcorr}[1]{%
  \par\vspace{6pt plus 2pt minus 2pt}\noindent
  \begin{minipage}{\linewidth}\small
    \noindent{\color{insarouge}\bfseries Question~#1}\enspace
    \textbf{Réponse : \csname qcm@rep@#1\endcsname}\hfill
    {\footnotesize\color{insagris}\csname qcm@ref@#1\endcsname}\par\nopagebreak
    \vspace{1pt}\noindent\hangindent=1.4em\hangafter=1
    \textbf{\csname qcm@rep@#1\endcsname)}\enspace
    \csname qcm@o@#1@\csname qcm@rep@#1\endcsname\endcsname\enspace---\enspace
    \csname qcm@c@#1@\csname qcm@rep@#1\endcsname\endcsname\par
    \qcm@wrongone{#1}{A}\qcm@wrongone{#1}{B}\qcm@wrongone{#1}{C}\qcm@wrongone{#1}{D}%
  \end{minipage}\par}
\makeatother

% =====================================================================
\section{QCM sur le cours}\label{chap:qcm}
% =====================================================================

\subsection{Mode d'emploi}\label{sec:qcm-mode}

Ce QCM compte \textbf{100 questions} sur les chapitres~\ref{chap:histoire} à~\ref{chap:bell},
classées en quatre niveaux. Chaque question propose quatre réponses dont \textbf{une seule} est
exacte : cochez au crayon la case de la proposition choisie.
\ifqcmrefs À droite du numéro de chaque question, une référence (chapitre, section) indique la
partie du cours concernée.\else\ifqcmcorrige{} La référence de chaque question (chapitre, section)
est donnée dans le corrigé.\fi\fi

\begin{center}
\renewcommand{\arraystretch}{1.3}
\begin{tabular}{@{}llc>{\raggedright\arraybackslash}p{8.2cm}@{}}
\toprule
\textbf{Niveau} & \textbf{Questions} & \textbf{Nombre} & \textbf{Ce qui est demandé}\\
\midrule
1. Débutant & 1 à 30 & 30 & définitions, notations, résultats du cours en une étape\\
2. Intermédiaire & 31 à 60 & 30 & petits calculs, propriétés, application directe d'un résultat\\
3. Avancé & 61 à 90 & 30 & calculs en plusieurs étapes, portes, circuits et protocoles\\
4. Maîtrise & 91 à 100 & 10 & raisonnements fins et pièges conceptuels, qui combinent plusieurs notions\\
\bottomrule
\end{tabular}
\end{center}

\begin{itemize}[nosep]
  \item \textbf{Sur papier} : imprimez seulement les pages~\pageref{chap:qcm} à~\pageref{qcm:fin}
        (mode d'emploi et énoncés, sections~\ref{sec:qcm-n1} à~\ref{sec:qcm-n4}).%
        \ifqcmcorrige{} Le corrigé commence sur une page distincte, à la page~\pageref{sec:qcm-cle}.\fi
  \item \textbf{Barème} : 1 point par bonne réponse, 0 sinon (pas de point négatif). Calculatrice
        autorisée, mais rarement utile (quelques logarithmes et racines).
  \item \textbf{Correction} :%
        \ifqcmcorrige{} comparez avec la clé des réponses, puis lisez l'explication de chaque
        question, même réussie. Les « pièges » décrivent l'erreur qui mène à chaque mauvaise réponse.%
        \else{} le corrigé n'est pas inclus dans cette version du document.\fi
\end{itemize}

\bigskip
\noindent Nom : \makebox[6cm]{\dotfill}\hfill Date : \makebox[3.5cm]{\dotfill}\hfill Score : \makebox[1.6cm]{\dotfill}\,/\,100
\bigskip

\begin{remarque}[Conventions des énoncés]
\begin{itemize}[nosep]
  \item $\ket{ab}=\ket{a}_A\otimes\ket{b}_B$ : le premier chiffre est le qubit $A$, le second le
        qubit $B$ ; dans un circuit, le fil du haut est le dernier chiffre.
  \item Le CNOT de la section~\ref{sec:cnot} a pour contrôle le \emph{second} chiffre ; les portes
        contrôlées de la section~\ref{sec:controlees} ont pour contrôle le \emph{premier}. Chaque
        énoncé précise le contrôle.
  \item Sans précision, une mesure est faite dans la base de calcul $Z$. Les logarithmes des
        entropies sont en base $2$ (résultats en bits).
  \item $\ket{\pm}=\frac{1}{\sqrt2}(\ket0\pm\ket1)$ et $\ket{{\pm}\ii}=\frac{1}{\sqrt2}(\ket0\pm\ii\ket1)$.
\end{itemize}
\end{remarque}

\clearpage
%% ======================================================================
\subsection{Niveau 1 : débutant (questions 1 à 30)}\label{sec:qcm-n1}
%% ======================================================================
\noindent{\color{insagris}\small Niveau de difficulté : Débutant\quad $\bullet\circ\circ\circ$}\par\smallskip
% ---- Question 1 (niveau 1, correct : D)
\qcmq[2]{Chap.~\ref{chap:histoire}\,$\cdot$\,\S\,\ref{sec:dualite}}
  {En 1905, qui explique l'effet photo-électrique en postulant que la lumière est faite de quanta d'énergie $E=hf$, les photons ?}
  {Louis de Broglie}
  {Erwin Schrödinger}
  {Max Planck}
  {Albert Einstein}
  {D}
\qcmx
  {Il propose en 1924 une onde de matière ($\lambda=h/p$) pour les particules comme l'électron ; l'idée du photon date de 1905.}
  {Il écrit en 1926 l'équation d'onde de la matière, plus de vingt ans après l'article sur les photons.}
  {Il introduit le quantum d'énergie en 1900 pour le rayonnement du corps noir ; c'est Einstein qui en fait une propriété de la lumière elle-même.}
  {Il postule que la lumière échange son énergie par paquets $E=hf=\hbar\omega$ : c'est la première face de la dualité onde--corpuscule.}

% ---- Question 2 (niveau 1, correct : A)
\qcmq[1]{Chap.~\ref{chap:histoire}\,$\cdot$\,\S\,\ref{sec:born}}
  {Selon l'interprétation de Max Born, que représente $|\Psi(x,t)|^2$ ?}
  {La densité de probabilité de présence de la particule en $x$ à l'instant $t$}
  {L'énergie cinétique de la particule}
  {L'amplitude de l'onde de matière}
  {La vitesse de la particule}
  {A}
\qcmx
  {$|\Psi(x,t)|^2\,\dd x$ est la probabilité de trouver la particule entre $x$ et $x+\dd x$ ; $\Psi$ lui-même n'est pas mesurable.}
  {L'énergie est associée à l'opérateur hamiltonien $\Hop$ ; $|\Psi|^2$ n'a pas la dimension d'une énergie.}
  {$\Psi$ est l'amplitude (complexe) ; $|\Psi|^2$ en est le carré du module.}
  {La vitesse s'obtient avec l'opérateur $\hat p=-\ii\hbar\,\partial_x$, pas avec $|\Psi|^2$.}

% ---- Question 3 (niveau 1, correct : D)
\qcmq[2]{Chap.~\ref{chap:histoire}\,$\cdot$\,\S\,\ref{sec:dualite}}
  {Une particule de quantité de mouvement $p$ est associée, d'après de Broglie, à une onde de longueur d'onde $\lambda$. Quelle relation est correcte ?}
  {$\lambda=hp$}
  {$\lambda=p/h$}
  {$\lambda=\hbar/p$}
  {$\lambda=h/p$}
  {D}
\qcmx
  {Produit au lieu du quotient : $\lambda$ grandirait avec $p$, le contraire de ce qu'on observe.}
  {Quotient inversé : $p/h=1/\lambda$ s'exprime en m$^{-1}$, pas en mètres.}
  {Confusion entre $h$ et $\hbar=h/2\pi$ : $\hbar/p=\lambda/2\pi$ est l'inverse du nombre d'onde $k$.}
  {Relation de de Broglie : plus la particule est rapide ou lourde, plus $\lambda$ est petite (balle de $0{,}1$~kg à $10$~m/s : $\lambda\approx10^{-34}$~m, inobservable).}

% ---- Question 4 (niveau 1, correct : C)
\qcmq[2]{Chap.~\ref{chap:histoire}\,$\cdot$\,\S\,\ref{sec:construction}}
  {Quelle est l'équation de Schrödinger dépendante du temps pour un état $\Psi$ et un hamiltonien $\Hop$ ?}
  {$-\ii\hbar\,\partial_t\Psi=\Hop\Psi$}
  {$\hbar\,\partial_t\Psi=\Hop\Psi$}
  {$\ii\hbar\,\partial_t\Psi=\Hop\Psi$}
  {$\ii\hbar\,\partial_x\Psi=\Hop\Psi$}
  {C}
\qcmx
  {Erreur de signe : la solution serait $\ee^{+\ii\Hop t/\hbar}\Psi(0)$, qui décrit l'évolution à rebours du temps.}
  {Sans le facteur $\ii$, les solutions sont des exponentielles réelles (croissantes ou décroissantes) et la norme n'est pas conservée.}
  {Elle relie l'évolution de l'état à l'énergie ; pour un $\Hop$ indépendant du temps, sa solution est $\Psi(t)=\ee^{-\ii\Hop t/\hbar}\Psi(0)$ (chapitre~\ref{chap:portes}).}
  {C'est la dérivée en $t$ qui décrit l'évolution ; la dérivée en $x$ sert à construire $\hat p=-\ii\hbar\,\partial_x$.}

% ---- Question 5 (niveau 1, correct : A)
\qcmq[1]{Chap.~\ref{chap:histoire}\,$\cdot$\,\S\,\ref{sec:construction}}
  {Quel opérateur représente l'énergie totale d'un système quantique ?}
  {L'hamiltonien $\Hop$}
  {L'opérateur quantité de mouvement $\hat p$}
  {L'opérateur position $\hat x$}
  {L'opérateur identité}
  {A}
\qcmx
  {Somme de l'énergie cinétique $-\frac{\hbar^2}{2m}\partial_x^2$ et de l'énergie potentielle $V$ ; ses valeurs propres sont les énergies mesurables.}
  {Il représente l'impulsion ; l'énergie cinétique est $\hat p^2/2m$, et il faut y ajouter $V$.}
  {C'est la multiplication par $x$ : il donne la position, pas l'énergie.}
  {Il vaut $1$ pour tout état : il ne contient aucune information sur l'énergie.}

% ---- Question 6 (niveau 1, correct : C)
\qcmq[1]{Chap.~\ref{chap:histoire}\,$\cdot$\,\S\,\ref{sec:histoire}}
  {Dans quel ordre chronologique ces avancées ont-elles eu lieu ?}
  {de Broglie $\to$ Planck $\to$ Schrödinger $\to$ Bell}
  {Planck $\to$ Schrödinger $\to$ de Broglie $\to$ Bell}
  {Planck (quantum d'énergie) $\to$ de Broglie (onde de matière) $\to$ Schrödinger (équation d'onde) $\to$ Bell (inégalités)}
  {Planck $\to$ de Broglie $\to$ Bell $\to$ Schrödinger}
  {C}
\qcmx
  {Le quantum d'énergie de Planck (1900) précède l'onde de matière de de Broglie (1924).}
  {L'équation de Schrödinger (1926) généralise l'idée d'onde de matière de de Broglie (1924) : elle vient après.}
  {Planck 1900, de Broglie 1924, Schrödinger 1926, Bell 1964.}
  {Les inégalités de Bell (1964) viennent près de quarante ans après l'équation de Schrödinger.}

% ---- Question 7 (niveau 1, correct : C)
\qcmq[1]{Chap.~\ref{chap:dirac}\,$\cdot$\,\S\,\ref{sec:dirac}}
  {En notation de Dirac, comment appelle-t-on $\ket{\Psi}$ et comment le représente-t-on pour un qubit ?}
  {Un bra, représenté par un vecteur ligne à deux composantes complexes}
  {Un ket, représenté par un nombre complexe}
  {Un ket, représenté par un vecteur colonne à deux composantes complexes}
  {Un bra, représenté par une matrice $2\times2$}
  {C}
\qcmx
  {Le bra $\bra{\Psi}$ est le conjugué transposé du ket : un vecteur ligne, mais ce n'est pas la notation $\ket{\Psi}$.}
  {Un scalaire ne peut pas décrire une superposition : il faut deux amplitudes $\alpha$ et $\beta$.}
  {Le ket est le vecteur colonne $\binom{\alpha}{\beta}$ ; son bra $\bra{\Psi}$ est le vecteur ligne conjugué.}
  {Une matrice est un opérateur ; le bra est une matrice ligne $1\times2$.}

% ---- Question 8 (niveau 1, correct : B)
\qcmq[2]{Chap.~\ref{chap:dirac}\,$\cdot$\,\S\,\ref{sec:qubit}}
  {Quelle condition les amplitudes $\alpha$ et $\beta$ de $\alpha\ket{0}+\beta\ket{1}$ doivent-elles vérifier pour que l'état soit normalisé ?}
  {$\alpha+\beta=1$}
  {$|\alpha|^2+|\beta|^2=1$}
  {$\alpha^2+\beta^2=1$}
  {$|\alpha|+|\beta|=1$}
  {B}
\qcmx
  {On somme les amplitudes (complexes, parfois négatives) ; ce sont leurs modules au carré qui sont des probabilités.}
  {Les probabilités de mesure $P(0)=|\alpha|^2$ et $P(1)=|\beta|^2$ doivent sommer à $1$.}
  {Oubli du module : pour $\alpha=\frac{1}{\sqrt2}$ et $\beta=\frac{\ii}{\sqrt2}$, $\alpha^2+\beta^2=0$ alors que l'état est normalisé.}
  {Somme des modules au lieu des carrés : $\alpha=\beta=\frac{1}{\sqrt2}$ est normalisé, mais $|\alpha|+|\beta|=\sqrt2$.}

% ---- Question 9 (niveau 1, correct : C)
\qcmq[1]{Chap.~\ref{chap:dirac}\,$\cdot$\,\S\,\ref{sec:projecteurs}}
  {On mesure dans la base $Z$ un qubit préparé dans $\ket{+}=\frac{1}{\sqrt2}(\ket{0}+\ket{1})$. Que se passe-t-il ?}
  {On lit $+$ et le qubit reste dans $\ket{+}$}
  {On lit toujours $0$}
  {On lit $0$ ou $1$, chacun avec la probabilité $\frac12$, et l'état devient $\ket{0}$ ou $\ket{1}$}
  {On lit $0$ et $1$ à la fois}
  {C}
\qcmx
  {Une mesure en base $Z$ ne donne que $0$ ou $1$ ; elle modifie l'état.}
  {Les deux amplitudes ont le même module : $P(0)=P(1)=\frac12$.}
  {Postulat de la mesure : la superposition s'effondre sur l'état propre correspondant au résultat.}
  {Chaque mesure donne un seul résultat ; la superposition ne s'observe qu'à travers les statistiques.}

% ---- Question 10 (niveau 1, correct : B)
\qcmq[2]{Chap.~\ref{chap:dirac}\,$\cdot$\,\S\,\ref{sec:qubit}}
  {Que vaut le produit scalaire $\braket{0}{1}$ ?}
  {$1$}
  {$0$}
  {$\frac{1}{\sqrt2}$}
  {$\ket{01}$}
  {B}
\qcmx
  {C'est la valeur de $\braket{0}{0}$ et de $\braket{1}{1}$ (vecteurs normés), pas de $\braket{0}{1}$.}
  {$\ket{0}$ et $\ket{1}$ sont orthogonaux : $\begin{pmatrix}1&0\end{pmatrix}\begin{pmatrix}0\\1\end{pmatrix}=0$.}
  {C'est la valeur de $\braket{0}{+}$, pas de $\braket{0}{1}$.}
  {Confusion avec le produit tensoriel : un produit scalaire est un nombre, pas un ket.}

% ---- Question 11 (niveau 1, correct : D)
\qcmq[2]{Chap.~\ref{chap:dirac}\,$\cdot$\,\S\,\ref{sec:qubit}}
  {Un qubit est dans l'état $\ket{-}=\frac{1}{\sqrt2}(\ket{0}-\ket{1})$. Quelle est la probabilité de mesurer $1$ en base $Z$ ?}
  {$0$}
  {$-\frac12$}
  {$\frac{1}{\sqrt2}$}
  {$\frac12$}
  {D}
\qcmx
  {Le signe $-$ ne « retire » pas $\ket{1}$ de la superposition.}
  {Une probabilité est toujours positive : on prend le module au carré de l'amplitude.}
  {C'est l'amplitude en module, pas la probabilité (il faut élever au carré).}
  {$P(1)=\left|-\frac{1}{\sqrt2}\right|^2=\frac12$ : le signe est une phase relative, qui disparaît dans le module.}

% ---- Question 12 (niveau 1, correct : B)
\qcmq[2]{Chap.~\ref{chap:dirac}\,$\cdot$\,\S\,\ref{sec:projecteurs}}
  {Quel est le projecteur sur l'état $\ket{0}$ ?}
  {$\ket{0}\bra{1}=\begin{pmatrix}0&1\\0&0\end{pmatrix}$}
  {$\hat P_0=\ket{0}\bra{0}=\begin{pmatrix}1&0\\0&0\end{pmatrix}$}
  {$\braket{0}{0}=1$}
  {$\begin{pmatrix}0&0\\0&1\end{pmatrix}$}
  {B}
\qcmx
  {Cet opérateur envoie $\ket{1}$ sur $\ket{0}$ ; son carré est nul : ce n'est pas un projecteur.}
  {$\hat P_0\ket{\Psi}=\alpha\ket{0}$ : il garde la composante sur $\ket{0}$ ; $\hat P_0^2=\hat P_0$.}
  {C'est un nombre, pas un opérateur.}
  {C'est $\ket{1}\bra{1}$, le projecteur sur $\ket{1}$.}

% ---- Question 13 (niveau 1, correct : A)
\qcmq[1]{Chap.~\ref{chap:operateurs}\,$\cdot$\,\S\,\ref{sec:bloch}}
  {Où se trouvent les états $\ket{0}$ et $\ket{1}$ sur la sphère de Bloch ?}
  {Aux deux pôles, Nord et Sud}
  {Aux deux extrémités de l'axe $x$}
  {Aux deux extrémités de l'axe $y$}
  {En deux points voisins de l'équateur}
  {A}
\qcmx
  {$\ket{0}$ au pôle Nord ($\theta=0$), $\ket{1}$ au pôle Sud ($\theta=\pi$) : deux états orthogonaux sont diamétralement opposés.}
  {Ce sont les places de $\ket{+}$ et $\ket{-}$.}
  {Ce sont les places de $\ket{{+}\ii}$ et $\ket{{-}\ii}$.}
  {Des états orthogonaux sont antipodaux, jamais voisins.}

% ---- Question 14 (niveau 1, correct : A)
\qcmq[2]{Chap.~\ref{chap:operateurs}\,$\cdot$\,\S\,\ref{sec:tenseur}}
  {Quelle est la dimension de l'espace de Hilbert de $n$ qubits ?}
  {$2^n$}
  {$2n$}
  {$n^2$}
  {$2^n-1$}
  {A}
\qcmx
  {Le produit tensoriel multiplie les dimensions : $2\times2\times\cdots\times2$ ; deux qubits : $4$, trois qubits : $8$.}
  {On additionne les dimensions au lieu de les multiplier ; pour $n=3$ on trouverait $6$ au lieu de $8$.}
  {Juste pour $n=2$ par coïncidence seulement : pour $n=3$ on trouverait $9$ au lieu de $8$.}
  {Il y a bien $2^n$ vecteurs dans la base de calcul, $\ket{0\cdots0}$ compris.}

% ---- Question 15 (niveau 1, correct : C)
\qcmq[1]{Chap.~\ref{chap:operateurs}\,$\cdot$\,\S\,\ref{sec:rho}}
  {Que vaut toujours la trace d'une matrice de densité $\rho$ ?}
  {$0$}
  {$1$ pour un état pur, et moins de $1$ pour un état mixte}
  {$1$}
  {$2$, la dimension de l'espace}
  {C}
\qcmx
  {C'est la trace des matrices de Pauli, pas d'une matrice de densité.}
  {Confusion avec $\operatorname{Tr}(\rho^2)$ : c'est la pureté qui distingue les deux cas, la trace vaut toujours $1$.}
  {C'est la somme des probabilités de mesure (les populations), que l'état soit pur ou mixte.}
  {C'est la trace de l'identité $\operatorname{Tr}\mathrm{Id}=2$, pas celle de $\rho$.}

% ---- Question 16 (niveau 1, correct : B)
\qcmq[1]{Chap.~\ref{chap:operateurs}\,$\cdot$\,\S\,\ref{sec:pur-mixte}}
  {Quelle est la valeur de $\operatorname{Tr}(\rho^2)$ pour un état pur ?}
  {$0$}
  {$1$}
  {$\frac12$}
  {Une valeur strictement inférieure à $1$}
  {B}
\qcmx
  {$\rho^2=0$ imposerait $\rho=0$, ce qui contredit $\operatorname{Tr}\rho=1$.}
  {Un état pur s'écrit $\rho=\ket{\Psi}\bra{\Psi}$, donc $\rho^2=\rho$ et $\operatorname{Tr}\rho^2=\operatorname{Tr}\rho=1$.}
  {C'est la valeur minimale pour un qubit, atteinte par l'état maximalement mixte $\mathrm{Id}/2$.}
  {C'est la caractérisation d'un état mixte.}

% ---- Question 17 (niveau 1, correct : B)
\qcmq[1]{Chap.~\ref{chap:operateurs}\,$\cdot$\,\S\,\ref{sec:rho}}
  {Que représentent les éléments diagonaux $\rho_{00}$ et $\rho_{11}$ de la matrice de densité d'un qubit ?}
  {Les cohérences entre $\ket{0}$ et $\ket{1}$}
  {Les probabilités de mesurer $\ket{0}$ et $\ket{1}$ (les populations)}
  {Les valeurs propres de $\rho$}
  {Les amplitudes $\alpha$ et $\beta$}
  {B}
\qcmx
  {Les cohérences sont les éléments hors diagonale $\rho_{01}$ et $\rho_{10}$.}
  {$\rho_{00}=\bra{0}\rho\ket{0}=P(0)$ et $\rho_{11}=P(1)$, d'où $\operatorname{Tr}\rho=1$.}
  {Seulement si $\rho$ est diagonale dans la base de calcul : $\ket{+}\bra{+}$ a pour diagonale $(\frac12,\frac12)$ mais pour valeurs propres $(1,0)$.}
  {Pour un état pur, $\rho_{00}=|\alpha|^2$ : le module au carré, pas l'amplitude.}

% ---- Question 18 (niveau 1, correct : D)
\qcmq[1]{Chap.~\ref{chap:operateurs}\,$\cdot$\,\S\,\ref{sec:op-dirac}}
  {Pourquoi associe-t-on aux grandeurs mesurables (observables) des opérateurs hermitiens ?}
  {Parce qu'ils conservent la norme des vecteurs}
  {Parce qu'ils ont tous une trace nulle}
  {Parce que leurs valeurs propres valent toutes $\pm1$}
  {Parce que leurs valeurs propres sont réelles, comme les résultats d'une mesure}
  {D}
\qcmx
  {C'est la propriété des opérateurs unitaires (les portes), pas des opérateurs hermitiens.}
  {Vrai pour les matrices de Pauli, faux en général (l'identité est hermitienne, de trace $2$), et sans rapport avec le caractère mesurable.}
  {Vrai pour $X$, $Y$, $Z$ seulement ; l'hamiltonien du puits a les valeurs propres $E_n=n^2E_1$.}
  {Un opérateur hermitien ($A^\dagger=A$) a des valeurs propres réelles et des vecteurs propres orthogonaux ; les résultats possibles sont ses valeurs propres.}

% ---- Question 19 (niveau 1, correct : C)
\qcmq[2]{Chap.~\ref{chap:portes}\,$\cdot$\,\S\,\ref{sec:pauli}}
  {Quelle porte de Pauli réalise l'inversion de bit $\ket{0}\leftrightarrow\ket{1}$ ?}
  {$Z$}
  {$Y$}
  {$X$}
  {$H$}
  {C}
\qcmx
  {$Z$ laisse $\ket{0}$ et change $\ket{1}$ en $-\ket{1}$ : c'est une inversion de phase.}
  {$Y\ket{0}=\ii\ket{1}$ : elle inverse le bit mais ajoute une phase ; la porte d'inversion « pure » est $X$.}
  {$X\ket{0}=\ket{1}$ et $X\ket{1}=\ket{0}$ : c'est le NOT quantique.}
  {$H$ crée des superpositions : $H\ket{0}=\ket{+}$.}

% ---- Question 20 (niveau 1, correct : B)
\qcmq[2]{Chap.~\ref{chap:portes}\,$\cdot$\,\S\,\ref{sec:hadamard}}
  {Que donne la porte de Hadamard appliquée à $\ket{0}$ ?}
  {$\ket{1}$}
  {$\ket{+}=\frac{1}{\sqrt2}(\ket{0}+\ket{1})$}
  {$\ket{-}=\frac{1}{\sqrt2}(\ket{0}-\ket{1})$}
  {$\ket{0}$}
  {B}
\qcmx
  {C'est l'action de $X$.}
  {Superposition à poids égaux : $P(0)=P(1)=\frac12$ ; c'est la première étape de la plupart des algorithmes.}
  {C'est $H\ket{1}$ ; le signe $-$ n'apparaît qu'à partir de $\ket{1}$.}
  {$H$ n'est pas l'identité ; elle ne laisse inchangés ni $\ket{0}$ ni $\ket{1}$ (elle envoie $\ket{+}$ sur $\ket{0}$).}

% ---- Question 21 (niveau 1, correct : D)
\qcmq[2]{Chap.~\ref{chap:portes}\,$\cdot$\,\S\,\ref{sec:portes}}
  {Quelle propriété doit vérifier la matrice $U$ d'une porte quantique ?}
  {Elle est hermitienne : $U^\dagger=U$}
  {Sa trace vaut $1$}
  {Ses coefficients sont réels}
  {Elle est unitaire : $U^\dagger U=\mathrm{Id}$}
  {D}
\qcmx
  {Vrai pour $X$, $Y$, $Z$, $H$ mais faux pour $S$ et $T$ : ce n'est pas la condition générale.}
  {C'est une propriété des matrices de densité ; $\operatorname{Tr}X=0$ et $\operatorname{Tr}\mathrm{Id}=2$.}
  {$Y$ et $S$ sont des portes à coefficients complexes.}
  {Elle conserve la norme ($\braket{\Psi}{\Psi}=1$ reste vrai) et elle est réversible.}

% ---- Question 22 (niveau 1, correct : A)
\qcmq[2]{Chap.~\ref{chap:portes}\,$\cdot$\,\S\,\ref{sec:portes}}
  {Quel est l'inverse d'une porte quantique $U$ ?}
  {$U^{-1}=U^\dagger$ (transposée conjuguée)}
  {$U^{-1}=U^{\mathsf T}$ (transposée seule)}
  {$U^{-1}=U^{*}$ (conjuguée seule)}
  {$U^{-1}=-U$}
  {A}
\qcmx
  {Conséquence de $U^\dagger U=\mathrm{Id}$ ; par exemple $S^{-1}=S^\dagger=\mathrm{diag}(1,-\ii)$.}
  {Oubli de la conjugaison : $S^{\mathsf T}=S=\mathrm{diag}(1,\ii)$ alors que $S^{-1}=\mathrm{diag}(1,-\ii)$.}
  {Faux en général : $Y^*=-Y$ alors que $Y^{-1}=Y$.}
  {Faux : $X^{-1}=X\neq-X$.}

% ---- Question 23 (niveau 1, correct : C)
\qcmq[1]{Chap.~\ref{chap:portes}\,$\cdot$\,\S\,\ref{sec:cnot}}
  {Dans une porte CNOT, que devient le qubit cible lorsque le qubit de contrôle est dans l'état $\ket{0}$ ?}
  {Il est inversé par la porte $X$}
  {Il est remis à $\ket{0}$}
  {Il reste inchangé}
  {Il subit un changement de signe, comme avec $Z$}
  {C}
\qcmx
  {C'est ce qui se passe quand le contrôle vaut $\ket{1}$.}
  {Le CNOT est réversible : remettre la cible à zéro détruirait de l'information.}
  {Avec le contrôle à $\ket{0}$, la porte agit comme l'identité sur la cible.}
  {Ce serait la porte $Z$ contrôlée (CZ), pas le CNOT.}

% ---- Question 24 (niveau 1, correct : A)
\qcmq[2]{Chap.~\ref{chap:portes}\,$\cdot$\,\S\,\ref{sec:cnot}}
  {Le CNOT du cours a pour contrôle le second chiffre du ket et pour cible le premier : $\mathrm{CNOT}\ket{a,b}=\ket{a\oplus b,\,b}$. Que donne-t-il sur $\ket{0,1}$ ?}
  {$\ket{1,1}$}
  {$\ket{0,1}$}
  {$\ket{0,0}$}
  {$\ket{1,0}$}
  {A}
\qcmx
  {Le contrôle $b=1$ inverse la cible : $0\oplus1=1$.}
  {Le contrôle vaut $1$ : la cible est bien inversée.}
  {C'est le contrôle qui serait inversé : le contrôle n'est jamais modifié.}
  {Les deux qubits seraient inversés : seule la cible change.}

% ---- Question 25 (niveau 1, correct : D)
\qcmq[1]{Chap.~\ref{chap:bell}\,$\cdot$\,\S\,\ref{sec:bell}}
  {Quelle propriété caractérise les quatre états de Bell ?}
  {Ils sont séparables : chaque qubit a son propre état}
  {Ils décrivent un seul qubit}
  {Ils ne diffèrent que par une phase globale}
  {Ils sont maximalement intriqués et forment une base orthonormée de l'espace de deux qubits}
  {D}
\qcmx
  {$D=c_{00}c_{11}-c_{01}c_{10}=\pm\frac12\neq0$ : ils ne se factorisent pas.}
  {Ce sont des états de deux qubits (quatre amplitudes sur $\ket{00},\ket{01},\ket{10},\ket{11}$).}
  {Ils sont orthogonaux deux à deux : ce sont quatre états distincts, alors qu'une phase globale ne change pas l'état.}
  {Aucun n'est un produit d'états à un qubit ($|D|=\frac12$, valeur maximale) et ils sont orthogonaux deux à deux.}

% ---- Question 26 (niveau 1, correct : D)
\qcmq[1]{Chap.~\ref{chap:bell}\,$\cdot$\,\S\,\ref{sec:teleportation}}
  {Quel est l'objectif du protocole de téléportation quantique ?}
  {Copier l'état d'un qubit inconnu sur le qubit de Bob}
  {Transmettre une information plus vite que la lumière}
  {Déplacer la particule physique d'Alice vers Bob}
  {Transférer l'état d'un qubit inconnu d'Alice à Bob, avec une paire intriquée et deux bits classiques}
  {D}
\qcmx
  {La copie est interdite (non-clonage) : après le protocole, l'état n'existe plus que chez Bob.}
  {Sans les deux bits classiques, le qubit de Bob est dans l'état $\mathrm{Id}/2$ : il n'a reçu aucune information.}
  {C'est l'état quantique qui est transféré, pas la particule qui le porte.}
  {Le qubit lui-même n'est pas envoyé : son état est reconstruit chez Bob, et détruit chez Alice par la mesure.}

% ---- Question 27 (niveau 1, correct : B)
\qcmq[1]{Chap.~\ref{chap:bell}\,$\cdot$\,\S\,\ref{sec:teleportation}}
  {Dans la téléportation, que doit envoyer Alice à Bob sur un canal classique pour qu'il retrouve l'état ?}
  {Un qubit intriqué}
  {Deux bits classiques, le résultat de sa mesure}
  {Les amplitudes exactes $\alpha$ et $\beta$ de l'état}
  {Un seul bit}
  {B}
\qcmx
  {La paire intriquée est partagée avant le protocole ; on ne l'envoie pas après.}
  {Les quatre résultats $(a,q)$ correspondent aux quatre corrections possibles ($I$, $Z$, $X$, $X$ puis $Z$) : il faut deux bits.}
  {Alice ne connaît pas l'état à téléporter, et ces nombres complexes n'ont pas de description finie en bits.}
  {Un bit ne distingue que deux cas, alors que Bob doit choisir parmi quatre corrections.}

% ---- Question 28 (niveau 1, correct : A)
\qcmq[1]{Chap.~\ref{chap:bell}\,$\cdot$\,\S\,\ref{sec:bell} ; ex.\,\ref{ex:superdense}}
  {Quel est le but du codage superdense ?}
  {Transmettre deux bits classiques en envoyant un seul qubit, grâce à une paire intriquée partagée}
  {Transmettre l'état d'un qubit avec deux bits classiques}
  {Compresser plusieurs qubits en un seul}
  {Chiffrer un message de façon inviolable}
  {A}
\qcmx
  {Alice applique l'une des quatre portes $I$, $Z$, $X$, $XZ$ à son qubit de la paire, l'envoie, et Bob identifie l'état de Bell obtenu.}
  {C'est la téléportation, le protocole inverse.}
  {Le protocole ne compresse pas de qubits : il fait passer deux bits classiques par un canal à un qubit.}
  {Ce n'est pas un protocole de cryptographie : la paire intriquée sert à doubler la capacité du canal.}

% ---- Question 29 (niveau 1, correct : A)
\qcmq[1]{Chap.~\ref{chap:bell}\,$\cdot$\,\S\,\ref{sec:teleportation}}
  {Que dit le théorème de non-clonage ?}
  {Aucune opération ne peut copier parfaitement un état quantique inconnu quelconque}
  {Il est impossible de mesurer un qubit sans le perturber}
  {Il est impossible de créer de l'intrication}
  {Il est impossible de copier un bit classique}
  {A}
\qcmx
  {Une porte $U$ qui copierait $\ket{0}$ et $\ket{1}$ enverrait $\ket{+}\ket{0}$ sur un état intriqué, par linéarité, et non sur $\ket{+}\ket{+}$.}
  {C'est l'effondrement de l'état, un autre principe.}
  {L'intrication se crée très bien (Hadamard puis CNOT).}
  {Le CNOT copie un bit classique : $\ket{0,b}\to\ket{b,b}$.}

% ---- Question 30 (niveau 1, correct : B)
\qcmq[1]{Chap.~\ref{chap:bell}\,$\cdot$\,\S\,\ref{sec:circuit-bell}}
  {Quel enchaînement de portes transforme $\ket{00}$ en l'état de Bell $\ket{\Phi^+}=\frac{1}{\sqrt2}(\ket{00}+\ket{11})$ ?}
  {CNOT, puis Hadamard sur le contrôle}
  {Hadamard sur l'un des qubits, puis CNOT dont le contrôle est ce même qubit}
  {Hadamard sur chacun des deux qubits}
  {SWAP, puis $X$ sur un qubit}
  {B}
\qcmx
  {Le CNOT ne fait rien sur $\ket{00}$ (contrôle à $0$) ; $H$ donne ensuite un état séparable.}
  {$H$ crée $\ket{0}\otimes\ket{+}$, puis le CNOT recopie la superposition : $\frac{1}{\sqrt2}(\ket{00}+\ket{11})$.}
  {On obtient $\ket{+}\otimes\ket{+}$, séparable : des portes locales ne créent pas d'intrication.}
  {$\ket{00}$ devient $\ket{01}$ ou $\ket{10}$, un état de base séparable.}

\clearpage
%% ======================================================================
\subsection{Niveau 2 : intermédiaire (questions 31 à 60)}\label{sec:qcm-n2}
%% ======================================================================
\noindent{\color{insagris}\small Niveau de difficulté : Intermédiaire\quad $\bullet\bullet\circ\circ$}\par\smallskip
% ---- Question 31 (niveau 2, correct : A)
\qcmq[1]{Chap.~\ref{chap:histoire}\,$\cdot$\,\S\,\ref{sec:born}}
  {Quelle condition la fonction d'onde d'une particule doit-elle vérifier sur tout l'espace ?}
  {$\displaystyle\int_{-\infty}^{+\infty}|\Psi(x,t)|^2\,\dd x=1$}
  {$\displaystyle\int_{-\infty}^{+\infty}\Psi(x,t)\,\dd x=1$}
  {$|\Psi(x,t)|^2=1$ pour tout $x$}
  {$\partial_x\Psi(x,t)=0$ en tout point}
  {A}
\qcmx
  {La particule se trouve forcément quelque part : la probabilité totale vaut $1$ (condition de normalisation).}
  {Sans le module au carré : $\Psi$ est complexe et son intégrale n'est pas une probabilité (elle peut être nulle ou complexe).}
  {$|\Psi|^2$ est une densité : si elle valait $1$ partout, son intégrale divergerait.}
  {Une fonction constante n'est pas normalisable ($\int|\Psi|^2\dd x$ diverge) ; la condition porte sur l'intégrale de $|\Psi|^2$.}

% ---- Question 32 (niveau 2, correct : D)
\qcmq[2]{Chap.~\ref{chap:histoire}\,$\cdot$\,\S\,\ref{sec:operateurs}}
  {Quelle est l'expression de l'opérateur quantité de mouvement $\hat p$ en une dimension ?}
  {$\hat p=m\,v$}
  {$\hat p=-\dfrac{\hbar^2}{2m}\dfrac{\partial^2}{\partial x^2}$}
  {$\hat p=+\ii\hbar\,\dfrac{\partial}{\partial x}$}
  {$\hat p=-\ii\hbar\,\dfrac{\partial}{\partial x}$}
  {D}
\qcmx
  {C'est la grandeur classique ; en mécanique quantique $p$ est un opérateur différentiel.}
  {C'est l'énergie cinétique $\hat p^2/2m$.}
  {Erreur de signe : on trouverait $p=-\hbar k$ pour l'onde $\ee^{\ii kx}$.}
  {Convention du cours : $\hat p\,\ee^{\ii kx}=\hbar k\,\ee^{\ii kx}$, donc une onde vers les $x$ croissants a $p=\hbar k>0$.}

% ---- Question 33 (niveau 2, correct : B)
\qcmq[2]{Chap.~\ref{chap:histoire}\,$\cdot$\,\S\,\ref{sec:solutions}}
  {Pour une particule dans un puits de potentiel infini de largeur $a$, comment varie l'énergie $E_n$ avec le numéro $n$ du niveau ?}
  {$E_n=E_1/n$}
  {$E_n=n^2E_1$, avec $E_1=\dfrac{\pi^2\hbar^2}{2ma^2}$}
  {$E_n=nE_1$}
  {$E_n=E_1$ pour tout $n$}
  {B}
\qcmx
  {Les niveaux se resserreraient quand $n$ augmente ; c'est l'inverse pour le puits infini.}
  {$k_n=\frac{n\pi}{a}$ et $E_n=\frac{\hbar^2k_n^2}{2m}$ : les niveaux s'écartent quand $n$ croît ($E_{n+1}-E_n=(2n+1)E_1$).}
  {Niveaux équidistants : ce n'est pas le puits infini, où $E_2=4E_1$ et $E_3=9E_1$.}
  {Il n'y aurait pas de quantification : tous les états auraient la même énergie.}

% ---- Question 34 (niveau 2, correct : C)
\qcmq[2]{Chap.~\ref{chap:histoire}\,$\cdot$\,\S\,\ref{sec:solutions}}
  {Une particule est dans un état stationnaire $\Psi_n$ de l'hamiltonien ($\Hop\Psi_n=E_n\Psi_n$). Que vaut l'écart-type $\sigma_E$ de son énergie ?}
  {$E_n$}
  {$E_n^2$}
  {$0$}
  {$\sqrt{E_n}$}
  {C}
\qcmx
  {C'est la valeur moyenne de l'énergie, pas son écart-type.}
  {C'est la valeur de $\langle\Hop^2\rangle$, avant soustraction de $\langle\Hop\rangle^2$.}
  {$\langle\Hop\rangle=E_n$ et $\langle\Hop^2\rangle=E_n^2$, donc $\sigma_E=\sqrt{E_n^2-E_n^2}=0$ : l'énergie est parfaitement déterminée.}
  {Mélange avec la formule $\sigma=\sqrt{\langle\Hop^2\rangle-\langle\Hop\rangle^2}$ : les deux termes sont égaux à $E_n^2$.}

% ---- Question 35 (niveau 2, correct : C)
\qcmq[2]{Chap.~\ref{chap:histoire}\,$\cdot$\,\S\,\ref{sec:solutions}}
  {Pour un état stationnaire $\Psi_n$ du puits infini $]0,a[$, que vaut la position moyenne $\langle\hat x\rangle$ ?}
  {$a/n$}
  {$a/4$}
  {$a/2$ pour tout $n$}
  {$0$}
  {C}
\qcmx
  {Elle dépendrait de $n$, ce que la symétrie de $|\Psi_n|^2$ interdit.}
  {C'est la position d'un maximum de $|\Psi_2|^2$ ; la moyenne des deux maxima $\frac a4$ et $\frac{3a}{4}$ donne $\frac a2$.}
  {$|\Psi_n|^2=\frac2a\sin^2\frac{n\pi x}{a}$ est symétrique par rapport au centre du puits.}
  {$\Psi_n$ s'annule sur la paroi, mais $\langle\hat x\rangle$ est une moyenne pondérée sur tout le puits.}

% ---- Question 36 (niveau 2, correct : D)
\qcmq[2]{Chap.~\ref{chap:dirac}\,$\cdot$\,\S\,\ref{sec:qubit}}
  {Un qubit est dans l'état $\ket{\Psi}=\frac{1}{\sqrt5}\ket{0}+\frac{2\ii}{\sqrt5}\ket{1}$. Quelle est la probabilité de mesurer $1$ en base $Z$ ?}
  {$\frac25$}
  {$-\frac45$}
  {$\frac15$}
  {$\frac45$}
  {D}
\qcmx
  {On divise le coefficient $2$ par $5$ : on oublie à la fois la racine et le carré.}
  {On élève $\frac{2\ii}{\sqrt5}$ au carré sans conjuguer : $\left(\frac{2\ii}{\sqrt5}\right)^2=-\frac45$ ; une probabilité est $|\cdot|^2$.}
  {C'est $P(0)=\left|\frac{1}{\sqrt5}\right|^2$.}
  {$P(1)=\left|\frac{2\ii}{\sqrt5}\right|^2=\frac45$ : le facteur $\ii$ est une phase relative, invisible en base $Z$.}

% ---- Question 37 (niveau 2, correct : C)
\qcmq[2]{Chap.~\ref{chap:dirac}\,$\cdot$\,\S\,\ref{sec:projecteurs}}
  {Un qubit préparé dans $\ket{0}$ est mesuré dans la base $X$ $\{\ket{+},\ket{-}\}$. Quelle est la probabilité d'obtenir $-$ ?}
  {$0$}
  {$1$}
  {$\frac12$}
  {$\frac{1}{\sqrt2}$}
  {C}
\qcmx
  {$\ket{0}$ est certain en base $Z$, mais pas en base $X$ : une autre base donne une autre loi.}
  {Même confusion dans l'autre sens : $\ket{0}$ n'est pas l'état $\ket{-}$.}
  {$P(-)=|\braket{-}{0}|^2=\left|\frac{1}{\sqrt2}\right|^2=\frac12$.}
  {C'est l'amplitude $\braket{-}{0}$ ; il faut élever son module au carré.}

% ---- Question 38 (niveau 2, correct : C)
\qcmq[2]{Chap.~\ref{chap:dirac}\,$\cdot$\,\S\,\ref{sec:qubit}}
  {Quel facteur $N>0$ normalise l'état $N\big(\ket{0}+(1-\ii)\ket{1}\big)$ ?}
  {$N=\frac{1}{\sqrt2}$}
  {$N=\frac13$}
  {$N=\frac{1}{\sqrt3}$}
  {$N=\frac{1}{\sqrt{1-2\ii}}$}
  {C}
\qcmx
  {On oublie la composante $\ket{0}$ : $|1-\ii|^2=2$, mais la norme au carré est $1+2=3$.}
  {On oublie la racine carrée : $3N^2=1$ donne $N=\frac{1}{\sqrt3}$, pas $\frac13$.}
  {$\braket{\Psi}{\Psi}=N^2\big(1+|1-\ii|^2\big)=N^2(1+2)=3N^2$.}
  {On élève $1-\ii$ au carré sans prendre le module : $(1-\ii)^2=-2\ii$ n'est pas réel ; il faut $|1-\ii|^2=2$.}

% ---- Question 39 (niveau 2, correct : D)
\qcmq[2]{Chap.~\ref{chap:dirac}\,$\cdot$\,\S\,\ref{sec:hilbert}}
  {Soient $\ket{\psi}=\frac{1}{\sqrt2}(\ket{0}+\ii\ket{1})$ et $\ket{\varphi}=\frac{1}{\sqrt2}(\ket{0}-\ket{1})$. Que vaut $\braket{\psi}{\varphi}$ ?}
  {$\frac{1-\ii}{2}$}
  {$\frac12$}
  {$0$}
  {$\frac{1+\ii}{2}$}
  {D}
\qcmx
  {On oublie de conjuguer l'amplitude $\ii$ du bra : on calcule $\frac12\big(1+\ii\cdot(-1)\big)$.}
  {Seul le terme en $\ket{0}$ est compté ; le terme en $\ket{1}$ vaut $\frac12(-\ii)(-1)=\frac{\ii}{2}$, non nul.}
  {Confusion avec $\braket{+}{-}=0$ ; ici $|\braket{\psi}{\varphi}|^2=\frac12$.}
  {Le bra conjugue les amplitudes : $\bra{\psi}=\frac{1}{\sqrt2}(\bra{0}-\ii\bra{1})$, d'où $\frac12\big(1+(-\ii)(-1)\big)=\frac{1+\ii}{2}$.}

% ---- Question 40 (niveau 2, correct : A)
\qcmq[2]{Chap.~\ref{chap:dirac}\,$\cdot$\,\S\,\ref{sec:qubit}}
  {Lequel de ces états est orthogonal à $\ket{{+}\ii}=\frac{1}{\sqrt2}(\ket{0}+\ii\ket{1})$ ?}
  {$\ket{{-}\ii}=\frac{1}{\sqrt2}(\ket{0}-\ii\ket{1})$}
  {$\ket{-}=\frac{1}{\sqrt2}(\ket{0}-\ket{1})$}
  {$\ket{+}=\frac{1}{\sqrt2}(\ket{0}+\ket{1})$}
  {$\ket{1}$}
  {A}
\qcmx
  {$\braket{{+}\ii}{{-}\ii}=\frac12\big(1+(-\ii)(-\ii)\big)=\frac12(1-1)=0$, en conjuguant le bra.}
  {$\braket{{+}\ii}{-}=\frac{1+\ii}{2}\neq0$ : $\ket{-}$ est orthogonal à $\ket{+}$, pas à $\ket{{+}\ii}$.}
  {$\braket{{+}\ii}{+}=\frac{1-\ii}{2}\neq0$.}
  {$\braket{{+}\ii}{1}=\frac{-\ii}{\sqrt2}\neq0$ : $\ket{1}$ est orthogonal à $\ket{0}$ seulement.}

% ---- Question 41 (niveau 2, correct : D)
\qcmq[1]{Chap.~\ref{chap:operateurs}\,$\cdot$\,\S\,\ref{sec:pur-mixte}}
  {Une matrice de densité de qubit a pour pureté $\operatorname{Tr}(\rho^2)=\frac58$. Que peut-on en conclure ?}
  {L'état est pur}
  {La matrice n'est pas une matrice de densité valide}
  {L'état est maximalement mixte, $\rho=\mathrm{Id}/2$}
  {L'état est mixte (mélange statistique)}
  {D}
\qcmx
  {Il faudrait $\operatorname{Tr}\rho^2=1$.}
  {$\frac58$ est compris entre $\frac12$ et $1$ : valeur admissible (exemple du cours : $\mathrm{diag}(\frac34,\frac14)$).}
  {Ce cas donnerait $\operatorname{Tr}\rho^2=\frac12$.}
  {$\operatorname{Tr}\rho^2<1$ ; pour un qubit, $\frac12\leq\operatorname{Tr}\rho^2\leq1$ et $\frac58$ est une valeur admissible.}

% ---- Question 42 (niveau 2, correct : B)
\qcmq[2]{Chap.~\ref{chap:operateurs}\,$\cdot$\,\S\,\ref{sec:rho}}
  {Que vaut $\langle Z\rangle=\operatorname{Tr}(\rho Z)$ pour un qubit dans l'état $\ket{+}$ ?}
  {$1$}
  {$0$}
  {$-1$}
  {$\frac{1}{\sqrt2}$}
  {B}
\qcmx
  {C'est la valeur pour $\ket{0}$.}
  {$\rho=\frac12\begin{pmatrix}1&1\\1&1\end{pmatrix}$ et $\operatorname{Tr}(\rho Z)=\frac12-\frac12=0$ : les résultats $+1$ et $-1$ sont équiprobables.}
  {C'est la valeur pour $\ket{1}$.}
  {C'est l'amplitude de $\ket{0}$ dans $\ket{+}$, pas une valeur moyenne de $Z$.}

% ---- Question 43 (niveau 2, correct : A)
\qcmq[1]{Chap.~\ref{chap:operateurs}\,$\cdot$\,\S\,\ref{sec:bloch}}
  {Sur la boule de Bloch, que représente un point situé à l'intérieur de la sphère ($|\vec r|<1$) ?}
  {Un état mixte, de pureté $\operatorname{Tr}\rho^2=\frac{1+|\vec r|^2}{2}<1$}
  {Un état pur dont la norme n'est pas $1$}
  {Un état pur en superposition}
  {Un point sans signification physique}
  {A}
\qcmx
  {Les états purs sont sur la sphère ; le centre $\vec r=\vec0$ est l'état maximalement mixte $\mathrm{Id}/2$.}
  {Tout état pur est normalisé et situé sur la sphère.}
  {Les états purs, superpositions comprises, sont tous \emph{sur} la sphère.}
  {Tout point de la boule est une matrice de densité valide.}

% ---- Question 44 (niveau 2, correct : B)
\qcmq[1]{Chap.~\ref{chap:operateurs}\,$\cdot$\,\S\,\ref{sec:pur-mixte}}
  {Par quoi la matrice de densité distingue-t-elle l'état pur $\ket{+}$ du mélange équiprobable de $\ket{0}$ et $\ket{1}$ ?}
  {Par leur trace, égale à $1$ pour l'état pur seulement}
  {Par les éléments hors diagonale (cohérences), non nuls pour $\ket{+}$ et nuls pour le mélange}
  {Par leurs éléments diagonaux, qui sont différents}
  {Ils ne se distinguent pas : c'est la même matrice}
  {B}
\qcmx
  {La trace vaut $1$ pour toute matrice de densité.}
  {$\ket{+}\bra{+}=\frac12\begin{pmatrix}1&1\\1&1\end{pmatrix}$ et le mélange vaut $\frac12\mathrm{Id}$ ; une mesure en base $X$ les distingue.}
  {Les deux ont la diagonale $(\frac12,\frac12)$ : $P(0)=P(1)=\frac12$ dans les deux cas.}
  {Les matrices sont différentes : en base $X$, $\ket{+}$ donne $+$ avec certitude, le mélange donne $\pm$ à pile ou face.}

% ---- Question 45 (niveau 2, correct : C)
\qcmq[2]{Chap.~\ref{chap:operateurs}\,$\cdot$\,\S\,\ref{sec:rho}}
  {Un qubit est préparé dans $\ket{0}$ avec la probabilité $\frac34$ et dans $\ket{1}$ avec la probabilité $\frac14$. Que vaut $\langle Z\rangle=\operatorname{Tr}(\rho Z)$ ?}
  {$\frac34$}
  {$1$}
  {$\frac12$}
  {$0$}
  {C}
\qcmx
  {C'est la probabilité $\rho_{00}$ de lire $0$ ; or $Z$ vaut $+1$ pour $0$ et $-1$ pour $1$.}
  {Valeur pour $\rho=\ket{0}\bra{0}$ (préparation certaine de $\ket{0}$).}
  {$\rho=\mathrm{diag}(\frac34,\frac14)$, donc $\operatorname{Tr}(\rho Z)=\frac34\cdot1+\frac14\cdot(-1)=\frac12$ : c'est la coordonnée $r_z$ de la sphère de Bloch.}
  {Valeur pour le mélange équiprobable $\mathrm{Id}/2$.}

% ---- Question 46 (niveau 2, correct : B)
\qcmq[2]{Chap.~\ref{chap:operateurs}\,$\cdot$\,\S\,\ref{sec:tenseur}}
  {Dans la base ordonnée $\{\ket{00},\ket{01},\ket{10},\ket{11}\}$, quel vecteur colonne représente $\ket{1}\otimes\ket{0}=\ket{10}$ ?}
  {$(0,1,0,0)^{\mathsf T}$}
  {$(0,0,1,0)^{\mathsf T}$}
  {$(1,0,0,0)^{\mathsf T}$}
  {$(0,0,0,1)^{\mathsf T}$}
  {B}
\qcmx
  {C'est $\ket{01}$ : l'ordre des facteurs est inversé.}
  {$\binom{0}{1}\otimes\binom{1}{0}=\big(0\binom{1}{0},\,1\binom{1}{0}\big)=(0,0,1,0)^{\mathsf T}$ : c'est le troisième vecteur de la base.}
  {C'est $\ket{00}$.}
  {C'est $\ket{11}$.}

% ---- Question 47 (niveau 2, correct : D)
\qcmq[2]{Chap.~\ref{chap:portes}\,$\cdot$\,\S\,\ref{sec:phase}}
  {Quelle relation lie la porte de phase $S=\mathrm{diag}(1,\ii)$ à la porte $Z$ ?}
  {$S=Z$}
  {$S^2=X$}
  {$S^4=Z$}
  {$S^2=Z$}
  {D}
\qcmx
  {$Z=\mathrm{diag}(1,-1)$ : $S$ est un quart de tour, $Z$ un demi-tour.}
  {$S^2$ est diagonale alors que $X$ ne l'est pas ; c'est $HSH$ qui est une racine carrée de $X$.}
  {$S^4=\mathrm{diag}(1,\ii^4)=\mathrm{Id}$.}
  {$S^2=\mathrm{diag}(1,\ii^2)=\mathrm{diag}(1,-1)=Z$ : deux quarts de tour autour de $z$ font un demi-tour.}

% ---- Question 48 (niveau 2, correct : A)
\qcmq[2]{Chap.~\ref{chap:portes}\,$\cdot$\,\S\,\ref{sec:controlees}}
  {Dans la porte de Toffoli (CCNOT) à trois qubits $x,y,z$, à quelle condition la cible $z$ est-elle inversée ?}
  {Seulement si $x=1$ et $y=1$}
  {Si $x=1$ ou $y=1$}
  {Si $x\neq y$}
  {Si $x=0$ et $y=0$}
  {A}
\qcmx
  {$\ket{x,y,z}\to\ket{x,y,z\oplus xy}$ : la cible est inversée si et seulement si le produit $xy$ vaut $1$.}
  {Un seul contrôle à $1$ ne suffit pas : pour $(x,y)=(1,0)$, $xy=0$ et la cible ne change pas.}
  {Ce serait deux CNOT ($z\to z\oplus x\oplus y$) ; pour $(1,0)$ le Toffoli n'inverse rien.}
  {Pour $(0,0)$ on a $xy=0$ : la cible ne change pas.}

% ---- Question 49 (niveau 2, correct : A)
\qcmq[1]{Chap.~\ref{chap:portes}\,$\cdot$\,\S\,\ref{sec:portes}}
  {Pour un hamiltonien $\Hop$ indépendant du temps, quel opérateur d'évolution $U(t,t_0)$ vérifie $\ket{\Psi(t)}=U(t,t_0)\ket{\Psi(t_0)}$ ?}
  {$U(t,t_0)=\ee^{-\ii\Hop(t-t_0)/\hbar}$}
  {$U(t,t_0)=\ee^{+\ii\Hop(t-t_0)/\hbar}$}
  {$U(t,t_0)=-\ii\hbar\,\partial_t$}
  {$U(t,t_0)=\Hop\,(t-t_0)/\hbar$}
  {A}
\qcmx
  {Solution de $\ii\hbar\,\partial_t\ket{\Psi}=\Hop\ket{\Psi}$ ; comme $\Hop$ est hermitien, $U^\dagger U=\mathrm{Id}$ : l'évolution est unitaire.}
  {Erreur de signe : c'est $U^\dagger$, qui remonte le temps.}
  {C'est un opérateur différentiel, pas un opérateur d'évolution ; l'équation de Schrödinger s'écrit $\ii\hbar\,\partial_t\ket{\Psi}=\Hop\ket{\Psi}$.}
  {On oublie l'exponentielle et le facteur $\ii$ : cet opérateur est hermitien, pas unitaire.}

% ---- Question 50 (niveau 2, correct : D)
\qcmq[2]{Chap.~\ref{chap:portes}\,$\cdot$\,\S\,\ref{sec:controlees}}
  {La porte de Fredkin (SWAP contrôlé) a pour contrôle le premier qubit. Que donne-t-elle sur $\ket{1,0,1}$ ?}
  {$\ket{1,0,1}$}
  {$\ket{0,1,1}$}
  {$\ket{1,1,1}$}
  {$\ket{1,1,0}$}
  {D}
\qcmx
  {C'est le résultat si le contrôle valait $0$ ; ici il vaut $1$ et les deux autres qubits sont différents.}
  {Le SWAP porte sur les qubits $2$ et $3$, pas sur le contrôle.}
  {Un SWAP ne change pas le nombre de $1$, il permute seulement ; $\ket{1,1,1}$ viendrait d'un NOT.}
  {Le contrôle vaut $1$ : les deux autres qubits sont échangés, $\ket{0,1}\to\ket{1,0}$.}

% ---- Question 51 (niveau 2, correct : A)
\qcmq[1]{Chap.~\ref{chap:portes}\,$\cdot$\,\S\,\ref{sec:op2}}
  {Que donne $(H\otimes I)\ket{00}$, où $H$ agit sur le premier facteur du produit tensoriel $\ket{a}\otimes\ket{b}$ ?}
  {$\ket{+}\otimes\ket{0}=\frac{1}{\sqrt2}(\ket{00}+\ket{10})$}
  {$\frac12(\ket{00}+\ket{01}+\ket{10}+\ket{11})$}
  {$\ket{0}\otimes\ket{+}=\frac{1}{\sqrt2}(\ket{00}+\ket{01})$}
  {$\frac{1}{\sqrt2}(\ket{00}+\ket{11})$}
  {A}
\qcmx
  {$(A\otimes B)(\ket{u}\otimes\ket{v})=A\ket{u}\otimes B\ket{v}$ : $H\ket{0}=\ket{+}$ sur le premier facteur, $I\ket{0}=\ket{0}$ sur le second.}
  {C'est $(H\otimes H)\ket{00}$ : ici le second qubit n'est pas touché.}
  {$H$ agit sur le premier qubit, pas sur le second : les facteurs sont inversés.}
  {C'est $\ket{\Phi^+}$ : il faut un CNOT pour intriquer ; des portes locales laissent l'état séparable.}

% ---- Question 52 (niveau 2, correct : C)
\qcmq[1]{Chap.~\ref{chap:bell}\,$\cdot$\,\S\,\ref{sec:separable}}
  {Pourquoi l'état de Bell $\ket{\Phi^+}=\frac{1}{\sqrt2}(\ket{00}+\ket{11})$ est-il dit intriqué ?}
  {Parce que $P(01)=0$}
  {Parce que chaque qubit, pris seul, est dans un état pur}
  {Il ne peut pas s'écrire comme un produit $\ket{\varphi}_A\otimes\ket{\chi}_B$ : $D=c_{00}c_{11}-c_{01}c_{10}=\frac12\neq0$}
  {Parce qu'il est orthogonal à $\ket{00}$}
  {C}
\qcmx
  {C'est vrai, mais $\ket{00}$ a aussi $P(01)=0$ et il est séparable : cette condition ne prouve rien.}
  {C'est l'inverse : pris seul, chaque qubit est complètement mixte ($\rho_A=\mathrm{Id}/2$).}
  {Un produit imposerait $D=0$ ; aucun des deux qubits n'a d'état propre.}
  {$\braket{00}{\Phi^+}=\frac{1}{\sqrt2}\neq0$ ; l'orthogonalité n'a pas de lien avec l'intrication.}

% ---- Question 53 (niveau 2, correct : B)
\qcmq[2]{Chap.~\ref{chap:bell}\,$\cdot$\,\S\,\ref{sec:teleportation}}
  {Dans la téléportation du cours, Alice obtient les résultats $a=0$ et $q=1$. Quelle correction Bob doit-il appliquer à son qubit ?}
  {La porte $X$}
  {La porte $Z$}
  {$X$ puis $Z$}
  {Aucune, l'état est déjà $\ket{\varphi}$}
  {B}
\qcmx
  {$X$ corrige le cas $a=1$ (inversion de bit) ; ici $a=0$.}
  {Bob détient $\alpha\ket{0}-\beta\ket{1}=Z\ket{\varphi}$ ; $Z$ rétablit le signe : $Z(\alpha\ket{0}-\beta\ket{1})=\alpha\ket{0}+\beta\ket{1}$.}
  {Les deux corrections ne sont nécessaires que pour $(a,q)=(1,1)$.}
  {Ce n'est vrai que pour $(a,q)=(0,0)$.}

% ---- Question 54 (niveau 2, correct : B)
\qcmq[2]{Chap.~\ref{chap:bell}\,$\cdot$\,\S\,\ref{sec:bell} ; ex.\,\ref{ex:superdense}}
  {Dans le codage superdense, Alice et Bob partagent $\ket{\Phi^+}$ et Alice agit sur son qubit. Quelle porte transforme $\ket{\Phi^+}$ en $\ket{\Psi^+}=\frac{1}{\sqrt2}(\ket{01}+\ket{10})$ ?}
  {La porte $Z$}
  {La porte $X$}
  {La porte identité}
  {La porte $H$}
  {B}
\qcmx
  {$Z$ donne $\ket{\Phi^-}=\frac{1}{\sqrt2}(\ket{00}-\ket{11})$ : elle change un signe, pas les chiffres.}
  {$X$ échange $\ket{0}$ et $\ket{1}$ sur un qubit : $\ket{00}+\ket{11}\to\ket{01}+\ket{10}$ ; c'est le message $10$.}
  {Elle laisse $\ket{\Phi^+}$ : c'est le message $00$.}
  {$(H\otimes I)\ket{\Phi^+}$ est une superposition de deux états de Bell, pas un état de Bell.}

% ---- Question 55 (niveau 2, correct : D)
\qcmq[2]{Chap.~\ref{chap:bell}\,$\cdot$\,\S\,\ref{sec:info}}
  {Pour deux qubits dans l'état de Bell $\ket{\Phi^+}$, que vaut l'entropie conditionnelle $H(A|B)=H(A,B)-H(B)$ ?}
  {$0$ bit}
  {$+1$ bit}
  {$+2$ bits}
  {$-1$ bit}
  {D}
\qcmx
  {C'est la valeur classique quand $A$ est entièrement déterminé par $B$ (deux bits identiques) ; ici $H(A,B)=0$, pas $H(B)$.}
  {C'est la valeur quand $A$ et $B$ sont indépendants : $H(A|B)=H(A)=1$.}
  {C'est l'information mutuelle $I(A;B)=1+1-0=2$ bits, pas l'entropie conditionnelle.}
  {$H(A,B)=S(\rho_{AB})=0$ (état pur) et $H(B)=1$ bit ($\rho_B=\mathrm{Id}/2$) : $0-1=-1$, valeur impossible en classique.}

% ---- Question 56 (niveau 2, correct : B)
\qcmq[1]{Chap.~\ref{chap:bell}\,$\cdot$\,\S\,\ref{sec:trace}}
  {Quelle est la matrice de densité réduite $\rho_A=\operatorname{Tr}_B\rho_{AB}$ d'un qubit d'une paire dans l'état $\ket{\Phi^+}$ ?}
  {$\rho_A=\ket{0}\bra{0}$}
  {$\rho_A=\frac12\mathrm{Id}$, de pureté $\operatorname{Tr}\rho_A^2=\frac12$}
  {$\rho_A=\ket{+}\bra{+}=\frac12\begin{pmatrix}1&1\\1&1\end{pmatrix}$}
  {$\rho_A=\ket{\Phi^+}\bra{\Phi^+}$}
  {B}
\qcmx
  {Le qubit aurait un état défini, ce qui est faux pour un état intriqué : $P(0)=P(1)=\frac12$.}
  {$\rho_A=\frac12\big(\ket{0}\bra{0}+\ket{1}\bra{1}\big)$ : les termes croisés $\ket{00}\bra{11}$ s'annulent ($\braket{1}{0}=0$). Le couple est pur, mais chaque qubit est complètement mixte.}
  {Les termes croisés sont multipliés par $\braket{1}{0}=0$ dans la trace partielle : ils disparaissent.}
  {C'est la matrice $4\times4$ du couple ; la trace partielle donne une matrice $2\times2$.}

% ---- Question 57 (niveau 2, correct : C)
\qcmq[1]{Chap.~\ref{chap:bell}\,$\cdot$\,\S\,\ref{sec:separable}}
  {Pour $\ket{\Psi}=c_{00}\ket{00}+c_{01}\ket{01}+c_{10}\ket{10}+c_{11}\ket{11}$, quel est le critère de séparabilité ?}
  {$c_{00}c_{11}+c_{01}c_{10}=0$}
  {$c_{00}c_{01}-c_{10}c_{11}=0$}
  {$D=c_{00}c_{11}-c_{01}c_{10}=0$}
  {$c_{00}=c_{11}$ et $c_{01}=c_{10}$}
  {C}
\qcmx
  {Mauvais signe : $\ket{+}\otimes\ket{-}$ est séparable mais donnerait $-\frac12\neq0$.}
  {Mauvais appariement : $\ket{\Psi^+}$ (intriqué) la vérifie, puisque $c_{00}=c_{11}=0$.}
  {Un produit $(a_0\ket{0}+a_1\ket{1})\otimes(b_0\ket{0}+b_1\ket{1})$ a $c_{ij}=a_ib_j$, donc $c_{00}c_{11}=a_0a_1b_0b_1=c_{01}c_{10}$ ; la réciproque est vraie aussi.}
  {Ni nécessaire ($\ket{00}$ est séparable avec $c_{00}\neq c_{11}$) ni suffisant ($\ket{\Phi^+}$ la vérifie et est intriqué).}

% ---- Question 58 (niveau 2, correct : B)
\qcmq[2]{Chap.~\ref{chap:bell}\,$\cdot$\,\S\,\ref{sec:bell}}
  {Quelle est la valeur propre de $X\otimes X$ pour l'état $\ket{\Phi^+}$ ?}
  {$-1$}
  {$+1$}
  {$0$}
  {$\ket{\Phi^+}$ n'est pas un état propre de $X\otimes X$}
  {B}
\qcmx
  {C'est la valeur pour $\ket{\Phi^-}$ et $\ket{\Psi^-}$ : le signe $-$ de l'état se retrouve dans la valeur propre.}
  {$(X\otimes X)\ket{00}=\ket{11}$ et $(X\otimes X)\ket{11}=\ket{00}$, donc $(X\otimes X)\ket{\Phi^+}=\ket{\Phi^+}$.}
  {$X\otimes X$ est unitaire : ses valeurs propres sont de module $1$, jamais nulles.}
  {Les quatre états de Bell sont des états propres communs de $Z\otimes Z$ et de $X\otimes X$.}

% ---- Question 59 (niveau 2, correct : A)
\qcmq[2]{Chap.~\ref{chap:bell}\,$\cdot$\,\S\,\ref{sec:info}}
  {Que vaut l'entropie de von Neumann $S(\rho)=-\sum_i\lambda_i\log_2\lambda_i$ d'un état pur ?}
  {$0$ bit}
  {$1$ bit}
  {$-1$ bit}
  {$+\infty$}
  {A}
\qcmx
  {Les valeurs propres sont $1$ et $0$ : $1\cdot\log_21=0$ et $0\cdot\log_20=0$ (par continuité, $x\log_2x\to0$).}
  {C'est la valeur pour l'état maximalement mixte $\mathrm{Id}/2$ d'un qubit.}
  {L'entropie de $\rho$ est positive ou nulle ; seule l'entropie conditionnelle quantique peut être négative.}
  {$\log_20$ diverge, mais le terme $0\cdot\log_20$ vaut $0$ par continuité.}

% ---- Question 60 (niveau 2, correct : A)
\qcmq[2]{Chap.~\ref{chap:bell}\,$\cdot$\,\S\,\ref{sec:bell}}
  {Lequel des états de Bell a pour valeur propre $-1$ à la fois pour $Z\otimes Z$ et pour $X\otimes X$ ?}
  {$\ket{\Psi^-}=\frac{1}{\sqrt2}(\ket{01}-\ket{10})$}
  {$\ket{\Phi^+}$}
  {$\ket{\Phi^-}$}
  {$\ket{\Psi^+}$}
  {A}
\qcmx
  {Résultats opposés dans les deux bases : $Z\otimes Z=-1$ (chiffres différents) et $X\otimes X=-1$ (signe $-$ dans l'état).}
  {Valeurs propres $(Z\otimes Z,\,X\otimes X)=(+1,+1)$.}
  {Valeurs propres $(Z\otimes Z,\,X\otimes X)=(+1,-1)$.}
  {Valeurs propres $(Z\otimes Z,\,X\otimes X)=(-1,+1)$.}

\clearpage
%% ======================================================================
\subsection{Niveau 3 : avancé (questions 61 à 90)}\label{sec:qcm-n3}
%% ======================================================================
\noindent{\color{insagris}\small Niveau de difficulté : Avancé\quad $\bullet\bullet\bullet\circ$}\par\smallskip
% ---- Question 61 (niveau 3, correct : A)
\qcmq[2]{Chap.~\ref{chap:histoire}\,$\cdot$\,\S\,\ref{sec:solutions}}
  {Un électron d'un puits infini passe du niveau $n=3$ au niveau $n=2$ en émettant un photon. Quelle est l'énergie du photon, en fonction de $E_1$ ?}
  {$5E_1$}
  {$E_1$}
  {$\frac94E_1$}
  {$13E_1$}
  {A}
\qcmx
  {$E_3-E_2=9E_1-4E_1=5E_1$ ; le photon a la fréquence $f=(E_3-E_2)/h$.}
  {On a soustrait les numéros de niveau ($3-2=1$) au lieu de leurs carrés.}
  {C'est le rapport $E_3/E_2$, pas la différence $E_3-E_2$.}
  {On a additionné $9E_1+4E_1$ : l'énergie du photon est une différence de niveaux.}

% ---- Question 62 (niveau 3, correct : A)
\qcmq[2]{Chap.~\ref{chap:histoire}\,$\cdot$\,\S\,\ref{sec:superposition}}
  {À $t=0$, une particule du puits infini est dans l'état $\Psi=\frac{\sqrt3}{2}\,\psi_1+\frac12\,\psi_3$ ($\psi_n$ : états stationnaires normalisés, d'énergies $E_n=n^2E_1$). Quelle est son énergie moyenne $\langle E\rangle$ ?}
  {$3E_1$}
  {$5E_1$}
  {$\left(\frac{\sqrt3}{2}+\frac92\right)E_1$}
  {$7E_1$}
  {A}
\qcmx
  {$P(E_1)=\frac34$ et $P(E_3)=\frac14$, donc $\langle E\rangle=\frac34E_1+\frac14\cdot9E_1=3E_1$.}
  {On a pondéré également les deux niveaux, $\frac{E_1+9E_1}{2}$ : les poids sont les $|c_i|^2$, pas $\frac12$.}
  {On a pondéré par les amplitudes $c_i$ au lieu de leurs carrés $|c_i|^2$.}
  {Poids permutés : $\frac14E_1+\frac34\cdot9E_1=7E_1$ ; la probabilité de $E_1$ est $\left|\frac{\sqrt3}{2}\right|^2=\frac34$.}

% ---- Question 63 (niveau 3, correct : C)
\qcmq[1]{Chap.~\ref{chap:histoire}\,$\cdot$\,\S\,\ref{sec:solutions}}
  {Pour l'état stationnaire $n=2$ du puits infini $]0,a[$, que valent la densité de probabilité de présence en $x=a/2$ et la position moyenne $\langle\hat x\rangle$ ?}
  {Densité maximale en $a/2$, et $\langle\hat x\rangle=a/2$}
  {Densité nulle en $a/2$, et $\langle\hat x\rangle=0$}
  {Densité nulle en $a/2$, et $\langle\hat x\rangle=a/2$}
  {Densité maximale en $a/2$, et $\langle\hat x\rangle=a/4$}
  {C}
\qcmx
  {C'est le cas de $n=1$ ; pour $n=2$ il y a un nœud au centre.}
  {Une densité nulle au centre ne donne pas une moyenne nulle : $\langle\hat x\rangle$ pondère tout le puits et vaut $\frac a2$ par symétrie.}
  {$|\Psi_2|^2=\frac2a\sin^2\frac{2\pi x}{a}$ s'annule en $x=\frac a2$ (nœud), alors que la symétrie impose $\langle\hat x\rangle=\frac a2$ : la moyenne n'est pas la position la plus probable.}
  {Ni l'un ni l'autre : maxima en $\frac a4$ et $\frac{3a}4$, nœud en $\frac a2$, moyenne $\frac a2$.}

% ---- Question 64 (niveau 3, correct : C)
\qcmq[2]{Chap.~\ref{chap:dirac}\,$\cdot$\,\S\,\ref{sec:projecteurs}}
  {Quelle est la matrice du projecteur $\hat P_+=\ket{+}\bra{+}$ dans la base $\{\ket{0},\ket{1}\}$ ?}
  {$\frac12\begin{pmatrix}1&-1\\-1&1\end{pmatrix}$}
  {$\begin{pmatrix}1&0\\0&0\end{pmatrix}$}
  {$\frac12\begin{pmatrix}1&1\\1&1\end{pmatrix}$}
  {$\frac{1}{\sqrt2}\begin{pmatrix}1&1\\1&-1\end{pmatrix}$}
  {C}
\qcmx
  {C'est $\ket{-}\bra{-}$, le projecteur sur $\ket{-}$.}
  {C'est $\ket{0}\bra{0}$ : projecteur de la base $Z$, pas de la base $X$.}
  {$\ket{+}\bra{+}=\frac12(\ket0+\ket1)(\bra0+\bra1)$ ; on vérifie $\hat P_+^2=\hat P_+$ et $\operatorname{Tr}\hat P_+=1$.}
  {C'est la matrice de Hadamard, unitaire ($H^2=\mathrm{Id}$), et non un projecteur (un projecteur vérifie $P^2=P$).}

% ---- Question 65 (niveau 3, correct : B)
\qcmq[2]{Chap.~\ref{chap:dirac}\,$\cdot$\,\S\,\ref{sec:projecteurs}}
  {Un qubit préparé dans $\ket{0}$ est mesuré en base $X$, puis aussitôt en base $Z$. Quelle est la probabilité d'obtenir $1$ à la seconde mesure ?}
  {$0$}
  {$\frac12$}
  {$1$}
  {$\frac14$}
  {B}
\qcmx
  {On garde $\ket{0}$ en mémoire : la mesure en base $X$ a effondré l'état sur $\ket{\pm}$ et effacé $\ket{0}$.}
  {La première mesure donne $\ket{+}$ ou $\ket{-}$ (probabilité $\frac12$ chacun) ; dans les deux cas $P(1)=\frac12$ en base $Z$ : $\frac12\cdot\frac12+\frac12\cdot\frac12=\frac12$.}
  {Sans rapport : $\ket{0}$ n'est plus l'état après la mesure en base $X$.}
  {On n'a compté qu'une branche de la première mesure ; il faut sommer les deux ($\frac14+\frac14$).}

% ---- Question 66 (niveau 3, correct : D)
\qcmq[2]{Chap.~\ref{chap:dirac}\,$\cdot$\,\S\,\ref{sec:projecteurs}}
  {Un qubit est dans l'état $\ket{\psi}=\frac{3}{5}\ket{0}+\frac{4}{5}\ket{1}$. Quelle est la probabilité d'obtenir $-$ en mesurant dans la base $X$ ?}
  {$\frac{16}{25}$}
  {$\frac{1}{25}$}
  {$\frac{49}{50}$}
  {$\frac{1}{50}$}
  {D}
\qcmx
  {C'est $P(1)$ en base $Z$ ($|\frac45|^2$) : on a mesuré dans la mauvaise base.}
  {On a oublié le facteur $\frac{1}{\sqrt2}$ de $\ket{-}$ : $\left|\frac35-\frac45\right|^2=\frac1{25}$, il manque le $\frac12$.}
  {C'est $P(+)$ ; la probabilité demandée est celle du résultat $-$.}
  {$\braket{-}{\psi}=\frac{1}{\sqrt2}\left(\frac35-\frac45\right)=-\frac{1}{5\sqrt2}$, donc $P(-)=\frac1{50}$ ; et $P(+)=\frac{49}{50}$.}

% ---- Question 67 (niveau 3, correct : B)
\qcmq[1]{Chap.~\ref{chap:dirac}\,$\cdot$\,\S\,\ref{sec:projecteurs}}
  {Lequel de ces kets décrit le même état physique que $\ket{+}=\frac{1}{\sqrt2}(\ket{0}+\ket{1})$ ?}
  {$\frac{1}{\sqrt2}\big(\ket{0}+\ii\ket{1}\big)$}
  {$\ee^{\ii\pi/3}\ket{+}$}
  {$\frac{1}{\sqrt2}\big(\ket{0}+\ee^{\ii\pi}\ket{1}\big)$}
  {$\frac{1}{\sqrt2}\big(\ket{0}+\ee^{\ii\pi/3}\ket{1}\big)$}
  {B}
\qcmx
  {La phase $\ii$ ne porte que sur $\ket{1}$ : c'est une phase \emph{relative}, et l'état est $\ket{{+}\ii}$ (base $Y$).}
  {Un facteur de phase \emph{globale} multiplie tout le ket : toutes les probabilités de mesure (et $\rho=\ket{+}\bra{+}$) sont inchangées.}
  {$\ee^{\ii\pi}=-1$ : c'est $\ket{-}$, orthogonal à $\ket{+}$.}
  {Phase relative $\frac\pi3$ : c'est un autre état, sur l'équateur de la sphère de Bloch, à l'azimut $\frac\pi3$.}

% ---- Question 68 (niveau 3, correct : A)
\qcmq[2]{Chap.~\ref{chap:dirac}\,$\cdot$\,\S\,\ref{sec:qubit}}
  {Laquelle de ces paires est une base orthonormée de $\mathbb{C}^2$ ?}
  {$\{\ket{{+}\ii},\ket{{-}\ii}\}$}
  {$\{\ket{0},\ket{+}\}$}
  {$\{\ket{+},\ket{{-}\ii}\}$}
  {$\{\ket{0},\ket{0}+\ket{1}\}$}
  {A}
\qcmx
  {Ces deux états sont normés et orthogonaux : $\braket{{+}\ii}{{-}\ii}=\frac12(1-1)=0$.}
  {Normés, mais $\braket{0}{+}=\frac{1}{\sqrt2}\neq0$ : pas orthogonaux.}
  {$\braket{+}{{-}\ii}=\frac12(1-\ii)\neq0$ : pas orthogonaux.}
  {Le second ket n'est pas normé (norme au carré $2$) ni orthogonal à $\ket{0}$.}

% ---- Question 69 (niveau 3, correct : C)
\qcmq[2]{Chap.~\ref{chap:operateurs}\,$\cdot$\,\S\,\ref{sec:rho}}
  {Laquelle de ces matrices peut être la matrice de densité d'un qubit ?}
  {$\begin{pmatrix}\frac12&\frac34\\\frac34&\frac12\end{pmatrix}$}
  {$\begin{pmatrix}\frac12&\frac{\ii}{2}\\\frac{\ii}{2}&\frac12\end{pmatrix}$}
  {$\begin{pmatrix}\frac23&\frac{\ii}{3}\\-\frac{\ii}{3}&\frac13\end{pmatrix}$}
  {$\begin{pmatrix}\frac34&0\\0&\frac12\end{pmatrix}$}
  {C}
\qcmx
  {Hermitienne et de trace $1$, mais pas positive : $\det=\frac14-\frac9{16}<0$, valeurs propres $\frac54$ et $-\frac14$ (une « probabilité » négative).}
  {De trace $1$ mais pas hermitienne : il faudrait $\rho_{10}=\overline{\rho_{01}}=-\frac{\ii}{2}$.}
  {Hermitienne ($\rho_{10}=\overline{\rho_{01}}$), de trace $1$ et positive : $\det\rho=\frac29-\frac19=\frac19>0$ avec $\operatorname{Tr}\rho>0$ donne deux valeurs propres positives, $\frac{3\pm\sqrt5}{6}$.}
  {Hermitienne et positive, mais de trace $\frac54\neq1$ : les probabilités diagonales ne somment pas à $1$.}

% ---- Question 70 (niveau 3, correct : B)
\qcmq[2]{Chap.~\ref{chap:operateurs}\,$\cdot$\,\S\,\ref{sec:op-dirac}}
  {Un qubit est dans l'état $\ket{\psi}=\frac{1}{\sqrt5}\big(\ket{0}+2\ii\ket{1}\big)$. Que vaut la valeur moyenne $\langle Y\rangle=\bra{\psi}Y\ket{\psi}$, avec $Y=\begin{pmatrix}0&-\ii\\\ii&0\end{pmatrix}$ ?}
  {$0$}
  {$\frac45$}
  {$-\frac45$}
  {$-\frac35$}
  {B}
\qcmx
  {On a oublié de conjuguer les amplitudes du bra : $\psi^{\mathsf T}Y\psi=\frac15(2+2\ii\cdot\ii)=0$ (c'est d'ailleurs la valeur de $\langle X\rangle$ pour cet état).}
  {$Y\ket\psi=\frac1{\sqrt5}(2\ket0+\ii\ket1)$, puis $\bra\psi Y\ket\psi=\frac15\big(1\cdot2+(-2\ii)(\ii)\big)=\frac45$ ; en général $\langle Y\rangle=2\,\mathrm{Im}(\bar\alpha\beta)$.}
  {Erreur de signe : on a utilisé la transposée $\begin{pmatrix}0&\ii\\-\ii&0\end{pmatrix}=-Y$.}
  {C'est $\langle Z\rangle=\frac15-\frac45$ : on a calculé la valeur moyenne d'une autre observable.}

% ---- Question 71 (niveau 3, correct : D)
\qcmq[1]{Chap.~\ref{chap:operateurs}\,$\cdot$\,\S\,\ref{sec:spectral}}
  {On mesure l'observable $A=\begin{pmatrix}2&1\\1&2\end{pmatrix}$ sur un qubit préparé dans $\ket{0}$. Quelle affirmation est correcte ?}
  {On obtient toujours $2$}
  {On obtient $+1$ ou $-1$, avec la probabilité $\frac12$ chacun}
  {On obtient $2$ ou $1$, avec les probabilités $\frac45$ et $\frac15$}
  {On obtient $3$ ou $1$, avec la probabilité $\frac12$ chacun (moyenne $2$)}
  {D}
\qcmx
  {On a lu le coefficient diagonal : $A\ket0=2\ket0+\ket1$ n'est pas proportionnel à $\ket0$, donc $\ket0$ n'est pas vecteur propre de $A$.}
  {Ce sont les valeurs propres de $X$, pas de $A$ : le décalage $2\,\mathrm{Id}$ translate le spectre ($2\pm1$).}
  {On a lu $A\ket0=2\ket0+\ket1$ comme un état : $A$ n'est pas unitaire, ce n'est pas une évolution ; les résultats possibles d'une mesure sont les valeurs propres de $A$.}
  {$A=2\,\mathrm{Id}+X$ a les vecteurs propres de $X$ : valeur propre $3$ pour $\ket+$, $1$ pour $\ket-$. Comme $\ket0=\frac{1}{\sqrt2}(\ket++\ket-)$, les deux résultats sont équiprobables, et $\langle A\rangle=\bra0A\ket0=2$.}

% ---- Question 72 (niveau 3, correct : D)
\qcmq[2]{Chap.~\ref{chap:operateurs}\,$\cdot$\,\S\,\ref{sec:pur-mixte}}
  {Un qubit est préparé dans $\ket{0}$ avec la probabilité $\frac12$ et dans $\ket{+}$ avec la probabilité $\frac12$. Quelle est la pureté $\operatorname{Tr}(\rho^2)$ de son état ?}
  {$1$}
  {$\frac12$}
  {$\frac58$}
  {$\frac34$}
  {D}
\qcmx
  {Un mélange de deux états purs n'est pur que s'ils sont identiques ; ici $\ket0\neq\ket+$, donc $\operatorname{Tr}\rho^2<1$.}
  {Valeur de l'état maximalement mixte $\frac{\mathrm{Id}}2$, atteinte pour deux états \emph{orthogonaux} ; ici ils ne le sont pas (on a aussi pu oublier les termes croisés).}
  {On n'a additionné que les carrés des éléments diagonaux ($\frac9{16}+\frac1{16}$) : il manque $2\rho_{01}\rho_{10}=\frac2{16}$.}
  {$\rho=\frac12\ket0\bra0+\frac12\ket+\bra+=\begin{pmatrix}\frac34&\frac14\\\frac14&\frac14\end{pmatrix}$, donc $\operatorname{Tr}\rho^2=\frac9{16}+\frac1{16}+2\cdot\frac1{16}=\frac34$ : état mixte, mais moins que $\frac{\mathrm{Id}}2$ car $\ket0$ et $\ket+$ ne sont pas orthogonaux.}

% ---- Question 73 (niveau 3, correct : B)
\qcmq[2]{Chap.~\ref{chap:operateurs}\,$\cdot$\,\S\,\ref{sec:bloch}}
  {Quel est le vecteur de Bloch $(r_x,r_y,r_z)$ de l'état $\ket{\psi}=\cos\frac\pi6\,\ket0+\sin\frac\pi6\,\ket1$ ?}
  {$\left(\frac12,\,0,\,\frac{\sqrt3}{2}\right)$}
  {$\left(\frac{\sqrt3}{2},\,0,\,\frac12\right)$}
  {$\left(\frac{\sqrt3}{4},\,0,\,\frac12\right)$}
  {$\left(\frac{\sqrt3}{2},\,0,\,\frac34\right)$}
  {B}
\qcmx
  {On a pris $\theta=\frac\pi6$ : l'angle de Bloch est le \emph{double} de l'angle qui figure dans les amplitudes.}
  {Ici $\frac\theta2=\frac\pi6$, donc $\theta=\frac\pi3$ et $\varphi=0$ : $\vec r=(\sin\theta\cos\varphi,\sin\theta\sin\varphi,\cos\theta)$. Contrôle : $r_z=|\alpha|^2-|\beta|^2=\frac34-\frac14$ et $r_x=2\alpha\beta=\frac{\sqrt3}2$.}
  {On a écrit $r_x=\alpha\beta$ en oubliant le facteur $2$ ; on le voit à $|\vec r|^2=\frac7{16}<1$, impossible pour un état pur.}
  {On a pris $r_z=|\alpha|^2$ au lieu de $|\alpha|^2-|\beta|^2$ ; or $|\vec r|^2=\frac{21}{16}>1$ est impossible.}

% ---- Question 74 (niveau 3, correct : A)
\qcmq[2]{Chap.~\ref{chap:operateurs}\,$\cdot$\,\S\,\ref{sec:bloch}}
  {Un qubit a pour matrice de densité $\rho=\frac12\begin{pmatrix}1&-\ii\\\ii&1\end{pmatrix}$. Quel est son vecteur de Bloch ?}
  {$(0,\,1,\,0)$ : c'est l'état pur $\ket{{+}\ii}$}
  {$(0,\,-1,\,0)$}
  {$\left(0,\,\frac12,\,0\right)$}
  {$(0,\,0,\,0)$}
  {A}
\qcmx
  {Avec $\rho=\frac12\begin{pmatrix}1+r_z&r_x-\ii r_y\\r_x+\ii r_y&1-r_z\end{pmatrix}$ : $r_z=0$ et $r_x-\ii r_y=-\ii$, donc $r_y=1$, $r_x=0$. Comme $|\vec r|=1$, l'état est pur : $\ket{{+}\ii}\bra{{+}\ii}$.}
  {Erreur de signe : $\rho_{01}=\frac{r_x-\ii r_y}{2}$, donc $\rho_{01}=-\frac\ii2$ correspond à $r_y=+1$ ($r_y=-1$ pour $\ket{{-}\ii}$).}
  {On a lu $\rho_{01}=-\frac\ii2$ comme $-\ii r_y$ : il faut multiplier par $2$ ($\rho=\frac12(\mathrm{Id}+\vec r\cdot\vec\sigma)$).}
  {Les diagonales $\frac12,\frac12$ donnent $r_z=0$, mais les éléments non diagonaux (cohérences) ne sont pas nuls : l'état est pur, et non $\frac{\mathrm{Id}}2$.}

% ---- Question 75 (niveau 3, correct : C)
\qcmq[2]{Chap.~\ref{chap:operateurs}\,$\cdot$\,\S\,\ref{sec:tenseur}}
  {Que donne $(X\otimes Z)\,\big(\ket{+}\otimes\ket{-}\big)$ ?}
  {$-\ket{-}\otimes\ket{-}$}
  {$\ket{-}\otimes\ket{+}$}
  {$\ket{+}\otimes\ket{+}$}
  {$\ket{+}\otimes\ket{-}$}
  {C}
\qcmx
  {On a échangé les facteurs ($Z\otimes X$) : $Z\ket+=\ket-$ et $X\ket-=-\ket-$. Le premier opérateur agit sur le premier qubit.}
  {On a cru que $X$ transforme $\ket+$ en $\ket-$ ; or $X$ échange $\ket0$ et $\ket1$ et laisse $\ket+$ inchangé.}
  {$(A\otimes B)(\ket u\otimes\ket v)=A\ket u\otimes B\ket v$ : $X\ket+=\ket+$ ($\ket+$ est propre pour $X$, valeur propre $+1$) et $Z\ket-=\ket+$ ($Z$ échange $\ket+$ et $\ket-$).}
  {On a cru que $Z$ laisse $\ket-$ inchangé ; $Z$ ne laisse invariants que $\ket0$ et $\ket1$, et échange $\ket+$ et $\ket-$.}

% ---- Question 76 (niveau 3, correct : C)
\qcmq[2]{Chap.~\ref{chap:portes}\,$\cdot$\,\S\,\ref{sec:phase}}
  {Le qubit $\ket{+}$ traverse la porte $T=\mathrm{diag}\big(1,\ee^{\ii\pi/4}\big)$, puis est mesuré dans la base $X$ $\{\ket+,\ket-\}$. Quelle est la probabilité d'obtenir $+$ ?}
  {$\frac12$}
  {$\frac{2-\sqrt2}{4}\approx0{,}15$}
  {$\frac{2+\sqrt2}{4}\approx0{,}85$}
  {$\frac{\sqrt2}{2}\approx0{,}71$}
  {C}
\qcmx
  {On a cru que la phase ne change pas les probabilités : c'est vrai en base $Z$, mais la base $X$ est sensible à la phase \emph{relative} (le point de Bloch est à l'azimut $\frac\pi4$ de l'axe $+x$).}
  {C'est $P(-)=\sin^2\frac\pi8$ : on a permuté les résultats $+$ et $-$.}
  {$T\ket+=\frac1{\sqrt2}(\ket0+\ee^{\ii\pi/4}\ket1)$, donc $\braket{+}{\psi}=\frac12\big(1+\ee^{\ii\pi/4}\big)$ et $P(+)=\frac14\big|1+\ee^{\ii\pi/4}\big|^2=\frac{2+2\cos\frac\pi4}{4}=\frac{2+\sqrt2}{4}$ (soit $\cos^2\frac\pi8$).}
  {On a pris $\cos\frac\pi4$ comme probabilité ; il faut le \emph{demi}-angle et le carré : $P=\cos^2\frac{\pi/4}{2}=\cos^2\frac\pi8$.}

% ---- Question 77 (niveau 3, correct : B)
\qcmq[2]{Chap.~\ref{chap:portes}\,$\cdot$\,\S\,\ref{sec:hadamard}}
  {Que vaut le produit de matrices $HYH$, où $Y=\begin{pmatrix}0&-\ii\\\ii&0\end{pmatrix}$ ?}
  {$Y$}
  {$-Y$}
  {$Z$}
  {$X$}
  {B}
\qcmx
  {Par analogie avec $HXH=Z$, on a cru que $H$ ne modifie pas $Y$ ; en fait l'axe $y$ est retourné (signe $-$).}
  {$H$ échange les axes $x$ et $z$ et retourne l'axe $y$ : $HXH=Z$, $HZH=X$, $HYH=-Y$. Direct : $HY=\frac{1}{\sqrt2}\begin{pmatrix}\ii&-\ii\\-\ii&-\ii\end{pmatrix}$, puis $HYH=\begin{pmatrix}0&\ii\\-\ii&0\end{pmatrix}=-Y$.}
  {C'est $HXH$ : $H$ n'échange que les axes $x$ et $z$, l'axe $y$ reste l'axe $y$.}
  {C'est $HZH$ : $Y$ ne peut pas devenir $X$, car $H$ laisse l'axe $y$ à sa place (au signe près).}

% ---- Question 78 (niveau 3, correct : D)
\qcmq[1]{Chap.~\ref{chap:portes}\,$\cdot$\,\S\,\ref{sec:op2}}
  {Que donne $(H\otimes H)\ket{01}$, où $\ket{01}=\ket0\otimes\ket1$ ?}
  {$\frac12\big(\ket{00}+\ket{01}-\ket{10}-\ket{11}\big)$}
  {$\frac12\big(\ket{00}-\ket{01}-\ket{10}+\ket{11}\big)$}
  {$\frac12\big(\ket{00}+\ket{01}+\ket{10}+\ket{11}\big)$}
  {$\frac12\big(\ket{00}-\ket{01}+\ket{10}-\ket{11}\big)$}
  {D}
\qcmx
  {C'est $\ket-\otimes\ket+$ : on a attribué le $\ket-$ au premier qubit, alors que c'est le second qui vaut $\ket1$.}
  {C'est $\ket-\otimes\ket-$ : $H\ket0=\ket+$, et non $\ket-$.}
  {C'est $\ket+\otimes\ket+$, le résultat pour $\ket{00}$ : on a oublié que $H\ket1=\ket-$ introduit des signes.}
  {$H\ket0\otimes H\ket1=\ket+\otimes\ket-=\frac12(\ket0+\ket1)\otimes(\ket0-\ket1)$. Le signe $-$ apparaît quand le \emph{second} chiffre vaut $1$ (formule générale $\frac12\sum_y(-1)^{x\cdot y}\ket y$ avec $x=01$).}

% ---- Question 79 (niveau 3, correct : D)
\qcmq[2]{Chap.~\ref{chap:portes}\,$\cdot$\,\S\,\ref{sec:cnot}}
  {Avec le CNOT du cours (contrôle : second qubit $B$ ; cible : premier qubit $A$ ; $\mathrm{CNOT}\ket{a,b}=\ket{a\oplus b,\,b}$), quel état obtient-on à partir de $\ket{1}_A\otimes\ket{+}_B$ ?}
  {$\ket{1}\otimes\ket{+}$ (inchangé)}
  {$\frac{1}{\sqrt2}\big(\ket{00}+\ket{11}\big)=\ket{\Phi^+}$}
  {$\frac{1}{\sqrt2}\big(\ket{01}-\ket{10}\big)=\ket{\Psi^-}$}
  {$\frac{1}{\sqrt2}\big(\ket{01}+\ket{10}\big)=\ket{\Psi^+}$}
  {D}
\qcmx
  {On a pris le premier qubit comme contrôle : le CNOT agirait alors sur $B$, et $X\ket+=\ket+$ laisserait l'état inchangé. Mais ici le contrôle est $B$.}
  {C'est le résultat pour $\ket{0}_A\ket{+}_B$ ; ici $A$ vaut $\ket1$, ce qui échange les deux termes.}
  {Aucun signe $-$ n'apparaît : le CNOT permute des états de base sans leur ajouter de phase.}
  {$\ket1\ket+=\frac1{\sqrt2}(\ket{10}+\ket{11})$ ; $\ket{10}\mapsto\ket{1\oplus0,0}=\ket{10}$ et $\ket{11}\mapsto\ket{1\oplus1,1}=\ket{01}$. C'est un état de Bell, donc intriqué.}

% ---- Question 80 (niveau 3, correct : C)
\qcmq[2]{Chap.~\ref{chap:portes}\,$\cdot$\,\S\,\ref{sec:swap}}
  {Que vaut $\mathrm{SWAP}\,(X\otimes\mathrm{Id})\,\mathrm{SWAP}$ ?}
  {$X\otimes\mathrm{Id}$}
  {$\mathrm{Id}\otimes\mathrm{Id}$}
  {$\mathrm{Id}\otimes X$}
  {$X\otimes X$}
  {C}
\qcmx
  {On a supposé que $\mathrm{SWAP}$ commute avec $X\otimes\mathrm{Id}$ ; ce n'est pas le cas, puisque cet opérateur n'agit pas de la même façon sur les deux qubits.}
  {On a simplifié $\mathrm{SWAP}^2=\mathrm{Id}$ en oubliant que $X\otimes\mathrm{Id}$ se trouve \emph{entre} les deux SWAP : on ne peut pas les rapprocher.}
  {$\mathrm{SWAP}\,(X\otimes\mathrm{Id})\,\mathrm{SWAP}\ket{a,b}=\mathrm{SWAP}\,(X\otimes\mathrm{Id})\ket{b,a}=\mathrm{SWAP}\ket{b\oplus1,a}=\ket{a,b\oplus1}$ : le SWAP transporte l'opérateur d'un qubit à l'autre (même mécanisme que $\mathrm{CNOT}_{21}=\mathrm{SWAP}\,\mathrm{CNOT}_{12}\,\mathrm{SWAP}$).}
  {Cet opérateur inverserait les deux bits ; ici un seul bit est inversé, celui du \emph{second} qubit.}

% ---- Question 81 (niveau 3, correct : C)
\qcmq[1]{Chap.~\ref{chap:portes}\,$\cdot$\,\S\,\ref{sec:controlees}}
  {La porte $Z$ contrôlée ($\mathrm{CZ}=\mathrm{diag}(1,1,1,-1)$, contrôle : premier qubit) est appliquée à $\ket{+}\otimes\ket{+}$. Quel état obtient-on ?}
  {$\ket{+}\otimes\ket{+}$}
  {$\ket{-}\otimes\ket{-}$}
  {$\frac12\big(\ket{00}+\ket{01}+\ket{10}-\ket{11}\big)$, un état intriqué}
  {$\ket{+}\otimes\ket{-}$}
  {C}
\qcmx
  {On a cru qu'un simple signe est une phase globale ; mais le signe $-$ ne touche que $\ket{11}$ : c'est une phase \emph{relative}.}
  {C'est $Z\otimes Z$ ($Z$ sur les deux qubits), et non une porte contrôlée.}
  {CZ ne change que le signe de $\ket{11}$. Critère de séparabilité : $c_{00}c_{11}-c_{01}c_{10}=-\frac14-\frac14=-\frac12\neq0$ ; on peut l'écrire $\frac{1}{\sqrt2}\big(\ket0\ket++\ket1\ket-\big)$.}
  {On a appliqué $Z$ au second qubit sans condition ; le contrôle exige que le premier qubit soit dans $\ket1$.}

% ---- Question 82 (niveau 3, correct : D)
\qcmq[2]{Chap.~\ref{chap:portes}\,$\cdot$\,\S\,\ref{sec:phase}, \ref{sec:portes-bloch}}
  {Un qubit préparé dans $\ket{0}$ traverse successivement $H$, puis $S$, puis $H$ (dans cet ordre). Quel est l'état final, à une phase globale près ?}
  {$\ket{{+}\ii}$}
  {$\ket{0}$}
  {$\ket{1}$}
  {$\ket{{-}\ii}=\frac{1}{\sqrt2}\big(\ket0-\ii\ket1\big)$}
  {D}
\qcmx
  {Mauvais sens de rotation : on a perdu un signe en calculant $H\ket{{+}\ii}$ (la phase relative vaut $-\ii$, pas $+\ii$).}
  {On a « simplifié » $HSH$ en $S$ (puis $S\ket0=\ket0$) : mais $H$ ne commute pas avec $S$, on ne peut pas supprimer les deux $H$.}
  {C'est le résultat pour $HZH=X$ ; or $S$ n'est qu'un \emph{quart} de tour autour de $z$ (et non un demi-tour comme $Z$), donc $HSH$ est un quart de tour autour de $x$.}
  {$HSH=\sqrt X$ est un quart de tour autour de $x$ (§\,\ref{sec:phase}). Directement : $H\ket0=\ket+$, $S\ket+=\ket{{+}\ii}$, $H\ket{{+}\ii}=\frac{1+\ii}{2}\ket0+\frac{1-\ii}{2}\ket1=\frac{1+\ii}{2}\big(\ket0-\ii\ket1\big)$.}

% ---- Question 83 (niveau 3, correct : A)
\qcmq[1]{Chap.~\ref{chap:bell}\,$\cdot$\,\S\,\ref{sec:mesure2}}
  {Deux qubits sont dans l'état $\ket{\Psi}=\frac{1}{\sqrt3}\big(\ket{00}+\ket{01}+\ket{11}\big)$. On mesure le \emph{second} qubit $B$ en base $Z$ et on obtient $1$. Que valait la probabilité de ce résultat, et dans quel état se trouve alors $A$ ?}
  {$P(B{=}1)=\frac23$ ; $A$ est dans l'état $\ket{+}$}
  {$P(B{=}1)=\frac13$ ; $A$ est dans l'état $\ket{0}$}
  {$P(B{=}1)=\frac23$ ; $A$ est dans l'état $\ket{1}$}
  {$P(B{=}1)=\frac13$ ; $A$ est dans l'état $\ket{+}$}
  {A}
\qcmx
  {$P(B{=}1)=|c_{01}|^2+|c_{11}|^2=\frac13+\frac13$. On garde les termes dont le second chiffre vaut $1$ : $\frac{1}{\sqrt3}\big(\ket{01}+\ket{11}\big)=\frac{1}{\sqrt3}(\ket0+\ket1)\otimes\ket1$, qui renormalisé donne $\ket{+}_A\otimes\ket{1}_B$.}
  {C'est le résultat $B=0$ : seul le terme $\ket{00}$ a un second chiffre nul, d'où la probabilité $\frac13$ et $A$ dans $\ket0$.}
  {On a projeté sur le seul terme $\ket{11}$ ; or $\ket{01}$ a aussi son second chiffre égal à $1$ : $A$ vaut $0$ ou $1$ (état $\ket+$).}
  {L'état de $A$ est juste, mais on n'a compté qu'un seul terme pour la probabilité ; il faut sommer $|c_{01}|^2+|c_{11}|^2$.}

% ---- Question 84 (niveau 3, correct : A)
\qcmq[2]{Chap.~\ref{chap:bell}\,$\cdot$\,\S\,\ref{sec:separable}}
  {Pour quelle valeur de $x$ l'état (non normalisé) $\ket{00}+2\ket{01}+3\ket{10}+x\ket{11}$ est-il séparable ?}
  {$x=6$}
  {$x=-6$}
  {$x=\frac16$}
  {$x=5$}
  {A}
\qcmx
  {Critère $D=c_{00}c_{11}-c_{01}c_{10}=x-2\cdot3=0$. On obtient bien un produit : $(\ket0+3\ket1)\otimes(\ket0+2\ket1)=\ket{00}+2\ket{01}+3\ket{10}+6\ket{11}$.}
  {Erreur de signe : $D=x-6$ s'annule pour $x=+6$ ; pour $x=-6$, $D=-12\neq0$ et l'état est intriqué.}
  {On a inversé le rapport : $c_{00}c_{11}=c_{01}c_{10}$ donne $x=\frac{c_{01}c_{10}}{c_{00}}=6$, et non $\frac{c_{00}}{c_{01}c_{10}}$.}
  {On a additionné ($2+3$) au lieu de multiplier : le critère de séparabilité porte sur des \emph{produits} de coefficients.}

% ---- Question 85 (niveau 3, correct : B)
\qcmq[2]{Chap.~\ref{chap:bell}\,$\cdot$\,\S\,\ref{sec:bell}}
  {Comment s'écrit $\ket{00}$ dans la base de Bell ?}
  {$\frac{1}{\sqrt2}\big(\ket{\Phi^+}-\ket{\Phi^-}\big)$}
  {$\frac{1}{\sqrt2}\big(\ket{\Phi^+}+\ket{\Phi^-}\big)$}
  {$\frac{1}{\sqrt2}\big(\ket{\Phi^+}+\ket{\Psi^+}\big)$}
  {$\ket{\Phi^+}$}
  {B}
\qcmx
  {Erreur de signe : cette combinaison vaut $\ket{11}$, puisque $\ket{\Phi^+}-\ket{\Phi^-}=\sqrt2\,\ket{11}$.}
  {$\ket{\Phi^+}+\ket{\Phi^-}=\frac{1}{\sqrt2}\cdot2\ket{00}=\sqrt2\,\ket{00}$. Une mesure de Bell sur $\ket{00}$ donne donc $\Phi^+$ ou $\Phi^-$ avec la probabilité $\frac12$ chacun (jamais $\Psi^\pm$).}
  {Cette combinaison vaut $\ket{+}\otimes\ket{+}=\frac12(\ket{00}+\ket{01}+\ket{10}+\ket{11})$ : les deux états de Bell « pairs » et « impairs » s'y ajoutent.}
  {On a confondu $\ket{00}$, état séparable, avec $\ket{\Phi^+}=\frac{1}{\sqrt2}(\ket{00}+\ket{11})$, qui contient aussi $\ket{11}$ et est intriqué.}

% ---- Question 86 (niveau 3, correct : A)
\qcmq[1]{Chap.~\ref{chap:bell}\,$\cdot$\,\S\,\ref{sec:circuit-bell}}
  {Le générateur de Bell du cours est $G=\mathrm{CNOT}\cdot(\mathrm{Id}\otimes H)$ : $H$ sur le second qubit $B$, puis CNOT de contrôle $B$ et de cible $A$. Quel circuit permet de \emph{mesurer} une paire dans la base de Bell (on le place avant une mesure en base $Z$) ?}
  {D'abord le CNOT (contrôle $B$, cible $A$), puis $H$ sur $B$}
  {D'abord $H$ sur $B$, puis le CNOT}
  {D'abord le CNOT, puis $H$ sur $A$ (la cible)}
  {Aucun circuit : on mesure directement les deux qubits en base $Z$}
  {A}
\qcmx
  {C'est l'inverse $G^\dagger=(\mathrm{Id}\otimes H)\,\mathrm{CNOT}$ : les portes sont leurs propres inverses et l'ordre s'inverse. Il envoie $\Phi^+,\Phi^-,\Psi^+,\Psi^-$ sur $\ket{00},\ket{01},\ket{10},\ket{11}$ (à une phase près).}
  {C'est $G$ lui-même, qui \emph{fabrique} des états de Bell : on garde le même ordre au lieu de l'inverser, et la mesure ne distingue plus les quatre états.}
  {Mauvais qubit : pour $\Phi^+$ on trouve $\ket{+}\otimes\ket{+}$, dont les deux bits sont aléatoires. Le $H$ doit agir sur $B$, qui a reçu le $H$ dans $G$.}
  {$\Phi^+$ et $\Phi^-$ (comme $\Psi^+$ et $\Psi^-$) donnent les mêmes résultats en base $Z$ : cette base ne voit pas le signe relatif, il faut d'abord défaire l'intrication.}

% ---- Question 87 (niveau 3, correct : B)
\qcmq[2]{Chap.~\ref{chap:bell}\,$\cdot$\,\S\,\ref{sec:trace}}
  {Quelle est la matrice de densité réduite $\rho_A=\operatorname{Tr}_B\ket\Psi\bra\Psi$ du premier qubit pour $\ket{\Psi}=\frac12\ket{00}+\frac{\sqrt3}{2}\ket{11}$ ?}
  {$\begin{pmatrix}\frac14&\frac{\sqrt3}{4}\\\frac{\sqrt3}{4}&\frac34\end{pmatrix}$}
  {$\begin{pmatrix}\frac14&0\\0&\frac34\end{pmatrix}$}
  {$\begin{pmatrix}\frac12&0\\0&\frac12\end{pmatrix}$}
  {$\begin{pmatrix}\frac12&0\\0&\frac{\sqrt3}{2}\end{pmatrix}$}
  {B}
\qcmx
  {C'est la matrice $\ket\chi\bra\chi$ de l'état pur $\ket\chi=\frac12\ket0+\frac{\sqrt3}2\ket1$ : on a gardé les cohérences, comme si $A$ était seul ; l'intrication avec $B$ les détruit.}
  {Avec $C=\begin{pmatrix}c_{00}&c_{01}\\c_{10}&c_{11}\end{pmatrix}=\begin{pmatrix}\frac12&0\\0&\frac{\sqrt3}2\end{pmatrix}$, $\rho_A=CC^\dagger=\mathrm{diag}\big(\frac14,\frac34\big)$ : état mixte, de pureté $\frac{10}{16}=1-2|D|^2$ avec $D=\frac{\sqrt3}{4}$.}
  {C'est la réponse pour un état de Bell (intrication maximale) ; ici les poids $\frac14$ et $\frac34$ diffèrent : l'intrication est partielle.}
  {On a mis les amplitudes au lieu de leurs carrés : la trace vaudrait $\frac{1+\sqrt3}{2}\neq1$.}

% ---- Question 88 (niveau 3, correct : A)
\qcmq[2]{Chap.~\ref{chap:bell}\,$\cdot$\,\S\,\ref{sec:info}}
  {Deux qubits sont dans le mélange statistique $\rho_{AB}=\frac12\ket{00}\bra{00}+\frac12\ket{11}\bra{11}$ (corrélations purement classiques). Que vaut $H(A|B)=S(\rho_{AB})-S(\rho_B)$, en bits ?}
  {$0$}
  {$-1$}
  {$1$}
  {$2$}
  {A}
\qcmx
  {$\rho_{AB}$ a pour valeurs propres $\frac12,\frac12,0,0$, donc $S(\rho_{AB})=1$ ; et $\rho_B=\frac12\mathrm{Id}$ donne $S(\rho_B)=1$. Ainsi $H(A|B)=0$ : connaître $B$ fixe $A$. Une valeur négative est propre à l'intrication.}
  {C'est la valeur pour l'état de Bell \emph{pur} $\ket{\Phi^+}$ : mêmes résultats en base $Z$, mais $S(\rho_{AB})=0$ car l'état global est pur. Ici l'état est mixte.}
  {On a pris $H(A|B)=H(A)$, comme si $B$ ne disait rien sur $A$ ; or les deux bits sont toujours égaux.}
  {On a additionné les entropies des deux qubits ($1+1$), valeur de deux bits indépendants : ce n'est pas une entropie conditionnelle.}

% ---- Question 89 (niveau 3, correct : B)
\qcmq[1]{Chap.~\ref{chap:bell}\,$\cdot$\,\S\,\ref{sec:teleportation}}
  {Dans la téléportation, avant de recevoir les deux bits $(a,q)$ d'Alice, dans quel état se trouve le qubit de Bob ?}
  {Dans l'état pur $\alpha\ket0+\beta\ket1$, déjà téléporté}
  {Dans l'état complètement mélangé $\rho_B=\frac12\mathrm{Id}$, indépendant de $\alpha$ et $\beta$}
  {Dans le mélange $|\alpha|^2\ket0\bra0+|\beta|^2\ket1\bra1$}
  {Dans l'état pur $\ket{0}$, son état initial}
  {B}
\qcmx
  {Ce serait une communication instantanée. L'état $\alpha\ket0+\beta\ket1$ n'apparaît qu'\emph{après} la correction, qui dépend de $(a,q)$.}
  {Les quatre kets $\ket{\phi_{aq}}$ sont équiprobables ($\frac14$) et $\frac14\sum_{a,q}\ket{\phi_{aq}}\bra{\phi_{aq}}=\frac12\mathrm{Id}$. Bob n'apprend rien avant le message classique : rien ne voyage plus vite que la lumière.}
  {Bob connaîtrait alors les statistiques de $\ket\varphi$ sans aucun message : or les probabilités d'Alice ($\frac14$ pour chaque résultat) ne dépendent pas de $\alpha,\beta$.}
  {Le qubit de Bob est la moitié de la paire $\ket{\Phi^+}$ : pris seul, il n'a pas d'état pur ($\rho_B=\frac12\mathrm{Id}$ avant comme après la mesure d'Alice, tant que Bob ignore le résultat).}

% ---- Question 90 (niveau 3, correct : D)
\qcmq[1]{Chap.~\ref{chap:bell}\,$\cdot$\,\S\,\ref{sec:chsh}}
  {Au jeu CHSH (gain si $a\oplus b=x\cdot y$), quelles sont les meilleures probabilités de gain avec une stratégie classique, même avec du hasard partagé, et avec une paire $\ket{\Phi^+}$ mesurée dans des bases bien choisies ?}
  {$\frac12$ en classique ; $1$ avec l'intrication}
  {$\frac34$ dans les deux cas}
  {$1$ en classique si Alice et Bob s'accordent à l'avance ; aucun avantage quantique}
  {$\frac34$ en classique ; $\cos^2\frac\pi8\approx0{,}85$ avec l'intrication}
  {D}
\qcmx
  {Classique : répondre toujours $0$ donne déjà $\frac34$. Quantique : l'intrication ne permet pas de gagner à tous les coups (borne de Tsirelson $\cos^2\frac\pi8$).}
  {Si l'intrication n'apportait rien, on n'aurait pas de violation de l'inégalité de Bell ; l'avantage quantique ($0{,}85>0{,}75$) est précisément ce que l'expérience de Bell met en évidence.}
  {S'accorder à l'avance (ou partager du hasard) ne suffit pas : aucune stratégie déterministe ne gagne les quatre parties, d'après l'argument de parité.}
  {Une stratégie déterministe perd au moins une des quatre parties (la somme modulo $2$ des conditions donne $0=1$) ; avec $\Phi^+$, chaque partie est gagnée avec la probabilité $\cos^2\frac\pi8=\frac{2+\sqrt2}{4}$.}

\clearpage
%% ======================================================================
\subsection{Niveau 4 : maîtrise (questions 91 à 100)}\label{sec:qcm-n4}
%% ======================================================================
\noindent{\color{insagris}\small Niveau de difficulté : Maîtrise\quad $\bullet\bullet\bullet\bullet$}\par\smallskip
% ---- Question 91 (niveau 4, correct : C)
\qcmq[2]{Chap.~\ref{chap:histoire}\,$\cdot$\,\S\,\ref{sec:superposition}}
  {Une particule du puits infini est préparée à $t=0$ dans l'état $\Psi=\frac{1}{\sqrt2}\big(\psi_1+\psi_2\big)$, avec $E_n=n^2E_1$. Au bout de quelle durée minimale la densité de probabilité $|\Psi(x,t)|^2$ retrouve-t-elle sa forme initiale ?}
  {$T=\dfrac{2\pi\hbar}{E_1}$}
  {$T=\dfrac{2\pi\hbar}{5E_1}$}
  {$T=\dfrac{2\pi\hbar}{3E_1}$}
  {$T=\dfrac{\pi\hbar}{3E_1}$}
  {C}
\qcmx
  {C'est la période de la phase $\ee^{-\ii E_1t/\hbar}$ de $\psi_1$ seul ; or les phases globales ne se voient pas dans $|\Psi|^2$ : seul compte l'écart $E_2-E_1=3E_1$.}
  {On a pris la somme $E_1+E_2=5E_1$ ; c'est la \emph{différence} des énergies (phase relative) qui intervient.}
  {$|\Psi|^2=\frac12(\psi_1^2+\psi_2^2)+\psi_1\psi_2\cos\frac{(E_2-E_1)t}{\hbar}$ : seul le terme croisé dépend du temps, à la pulsation $\omega=\frac{E_2-E_1}{\hbar}=\frac{3E_1}{\hbar}$, donc $T=\frac{2\pi}{\omega}$.}
  {C'est la demi-période : à $t=\frac T2$ le cosinus vaut $-1$ et la densité est le symétrique de l'état initial ($x\to a-x$), pas encore sa forme de départ.}

% ---- Question 92 (niveau 4, correct : D)
\qcmq[1]{Chap.~\ref{chap:dirac}\,$\cdot$\,\S\,\ref{sec:projecteurs}, \ref{sec:qubit}}
  {Un expérimentateur reçoit un qubit préparé soit dans $\ket{0}$, soit dans $\ket{+}$. Laquelle de ces affirmations est correcte ?}
  {Une mesure en base $Z$ les distingue : $\ket0$ donne $0$ et $\ket+$ donne $1$}
  {Une mesure en base $X$ les distingue : $\ket+$ donne $+$ et $\ket0$ donne $-$}
  {On les distingue en répétant la mesure sur le même qubit}
  {Aucune mesure ne permet de les distinguer à coup sûr, car $\braket{0}{+}=\frac{1}{\sqrt2}\neq0$}
  {D}
\qcmx
  {$\ket+$ donne $0$ ou $1$ avec la probabilité $\frac12$ : le résultat $0$ est possible pour les deux états.}
  {$\ket0=\frac{1}{\sqrt2}(\ket++\ket-)$ donne $+$ ou $-$ avec la probabilité $\frac12$ : le résultat $+$ est possible pour les deux états.}
  {Après une première mesure l'état est effondré : répéter ne donne aucune information nouvelle, et on ne peut pas copier le qubit pour recommencer (non-clonage).}
  {Seuls des états \emph{orthogonaux} sont parfaitement discernables. En base $Z$, $\ket0$ donne toujours $0$ et $\ket+$ donne $0$ une fois sur deux : obtenir $0$ ne tranche pas (obtenir $1$ prouve seulement qu'on avait $\ket+$).}

% ---- Question 93 (niveau 4, correct : C)
\qcmq[1]{Chap.~\ref{chap:operateurs}\,$\cdot$\,\S\,\ref{sec:pur-mixte}}
  {Une source envoie au hasard $\ket{0}$ ou $\ket{1}$ (probabilité $\frac12$ chacun) ; une autre envoie au hasard $\ket{+}$ ou $\ket{-}$ (probabilité $\frac12$ chacun). Peut-on distinguer ces deux sources en mesurant les qubits émis ?}
  {Oui : en base $Z$, la première source donne des résultats prévisibles}
  {Oui : en base $X$, la seconde source donne des résultats prévisibles}
  {Non : les deux ont la même matrice de densité $\frac{\mathrm{Id}}{2}$, donc les mêmes statistiques pour toute mesure}
  {Oui : les états purs envoyés ne sont pas les mêmes}
  {C}
\qcmx
  {Chaque qubit est bien $\ket0$ ou $\ket1$, mais on ignore lequel : le résultat vaut $0$ ou $1$ avec la probabilité $\frac12$, exactement comme avec la seconde source ($|\braket{0}{\pm}|^2=\frac12$).}
  {Même erreur : chaque qubit est $\ket+$ ou $\ket-$, mais on ignore lequel ; le résultat en base $X$ est aléatoire, et la première source donne aussi $+$ ou $-$ avec la probabilité $\frac12$.}
  {$\frac12\ket0\bra0+\frac12\ket1\bra1=\frac12\ket+\bra++\frac12\ket-\bra-=\frac{\mathrm{Id}}2$ (relation de fermeture dans chaque base). Les probabilités de toute mesure ne dépendent que de $\rho$.}
  {Ce qui compte n'est pas la liste des états purs envoyés mais la matrice de densité moyenne : des préparations différentes peuvent avoir la même $\rho$.}

% ---- Question 94 (niveau 4, correct : B)
\qcmq[1]{Chap.~\ref{chap:operateurs}\,$\cdot$\,\S\,\ref{sec:pur-mixte} ; Chap.~\ref{chap:portes}\,$\cdot$\,\S\,\ref{sec:portes-bloch}}
  {Un qubit est dans l'état complètement mélangé $\rho=\frac{\mathrm{Id}}{2}$. Existe-t-il une porte quantique $U$ qui le transforme en un état pur, par exemple $\ket{0}\bra{0}$ ?}
  {Oui : la porte de Hadamard, qui crée des superpositions}
  {Non : $\rho\mapsto U\rho U^\dagger$ conserve les valeurs propres, donc la pureté ; ici $U\frac{\mathrm{Id}}{2}U^\dagger=\frac{\mathrm{Id}}{2}$}
  {Oui : une rotation de Bloch qui amène le centre de la boule au pôle nord}
  {Oui : une mesure en base $Z$, qui donne $\ket0$}
  {B}
\qcmx
  {$H$ envoie un état pur sur un état pur ($\ket0\mapsto\ket+$) et laisse $\frac{\mathrm{Id}}{2}$ invariant : $H\frac{\mathrm{Id}}2H=\frac{\mathrm{Id}}2$ car $H^2=\mathrm{Id}$.}
  {$\operatorname{Tr}\big[(U\rho U^\dagger)^2\big]=\operatorname{Tr}\rho^2$ ($=\frac12$ ici). Sur la boule de Bloch, une porte est une rotation : elle laisse le centre fixe.}
  {Une rotation tourne la boule autour de son centre, qui reste fixe : elle ne peut pas l'amener sur la sphère.}
  {Une mesure n'est pas une porte (elle est non unitaire) ; de plus elle donne $\ket0$ ou $\ket1$ au hasard, pas $\ket0$ à coup sûr.}

% ---- Question 95 (niveau 4, correct : D)
\qcmq[1]{Chap.~\ref{chap:portes}\,$\cdot$\,\S\,\ref{sec:cnot}}
  {Que vaut $(H\otimes H)\,\mathrm{CNOT}_{12}\,(H\otimes H)$, où $\mathrm{CNOT}_{12}$ a pour contrôle le \emph{premier} qubit et pour cible le second ?}
  {$\mathrm{CNOT}_{12}$, inchangé}
  {$\mathrm{CZ}$ (porte $Z$ contrôlée)}
  {$\mathrm{SWAP}$}
  {$\mathrm{CNOT}_{21}$ (contrôle : second qubit ; cible : premier qubit)}
  {D}
\qcmx
  {On a cru que les deux $H$ s'annulent ($H^2=\mathrm{Id}$) ; mais ils encadrent le CNOT au lieu d'être côte à côte : on ne peut pas les simplifier.}
  {C'est le résultat quand $H$ n'encadre que la \emph{cible} : $(\mathrm{Id}\otimes H)\,\mathrm{CNOT}_{12}\,(\mathrm{Id}\otimes H)=\mathrm{CZ}$, car $HXH=Z$. Ici $H$ agit aussi sur le contrôle.}
  {Le SWAP s'obtient avec \emph{trois} CNOT alternés, pas avec un seul CNOT conjugué par $H\otimes H$ ; il n'y a pas de contrôle dans un SWAP.}
  {$H\ket0\bra0H=\ket+\bra+$, $H\ket1\bra1H=\ket-\bra-$ et $HXH=Z$ donnent $\ket+\bra+\otimes I+\ket-\bra-\otimes Z=I\otimes\ket0\bra0+X\otimes\ket1\bra1=\mathrm{CNOT}_{21}$ : dans la base $\{\ket+,\ket-\}$, contrôle et cible échangent leurs rôles (rebond de phase).}

% ---- Question 96 (niveau 4, correct : A)
\qcmq[1]{Chap.~\ref{chap:portes}\,$\cdot$\,\S\,\ref{sec:controlees}}
  {On applique la porte de Toffoli (contrôles : qubits 1 et 2 ; cible : qubit 3) à $\ket{+}\otimes\ket{+}\otimes\ket{0}$, puis on mesure le troisième qubit en base $Z$. Que peut-on dire si l'on obtient $1$ ?}
  {Ce résultat a la probabilité $\frac14$, et les deux premiers qubits sont alors dans l'état $\ket{11}$}
  {Ce résultat a la probabilité $\frac12$, et les deux premiers qubits restent dans $\ket{+}\ket{+}$}
  {Ce résultat a la probabilité $\frac14$, mais les deux premiers qubits restent dans $\ket{+}\ket{+}$}
  {Ce résultat a la probabilité $\frac34$, et les deux premiers qubits sont dans l'état $\frac{1}{\sqrt3}\big(\ket{00}+\ket{01}+\ket{10}\big)$}
  {A}
\qcmx
  {$\mathrm{CCNOT}\ket{x,y,0}=\ket{x,y,xy}$ donne $\frac12\big(\ket{000}+\ket{010}+\ket{100}+\ket{111}\big)$ : seul $\ket{111}$ a un troisième chiffre égal à $1$ (probabilité $\frac14$), et il impose $x=y=1$.}
  {La cible n'est inversée que si $x=y=1$ (probabilité $\frac14$) et non une fois sur deux ; de plus l'état est intriqué, donc mesurer le qubit 3 modifie les qubits 1 et 2.}
  {La probabilité est juste, mais la mesure de la cible n'est pas sans effet sur les contrôles : l'état après la Toffoli est intriqué, donc le résultat $1$ les projette sur $\ket{11}$.}
  {C'est le cas du résultat $0$ (probabilité $\frac34$ : $x,y\in\{00,01,10\}$). Le résultat $1$ est le cas rare, celui du ET vrai ($x=y=1$).}

% ---- Question 97 (niveau 4, correct : A)
\qcmq[1]{Chap.~\ref{chap:bell}\,$\cdot$\,\S\,\ref{sec:bell}}
  {Que donne $(H\otimes H)\ket{\Psi^-}$, avec $\ket{\Psi^-}=\frac{1}{\sqrt2}\big(\ket{01}-\ket{10}\big)$ ?}
  {$-\ket{\Psi^-}$, c'est-à-dire le même état physique}
  {$\ket{\Phi^-}$}
  {$\ket{\Psi^+}$}
  {$\ket{\Phi^+}$}
  {A}
\qcmx
  {Le singulet est invariant par $U\otimes U$ à un facteur près : $(U\otimes U)\ket{\Psi^-}=\det(U)\ket{\Psi^-}$ avec $\det H=-1$. Directement : $\frac{1}{\sqrt2}\big(\ket+\ket--\ket-\ket+\big)=-\ket{\Psi^-}$.}
  {C'est l'image de $\ket{\Psi^+}$, pas de $\ket{\Psi^-}$ : $(H\otimes H)\ket{\Psi^+}=\ket{\Phi^-}$. Le signe $-$ du singulet change le calcul (les termes $\ket{00}$ et $\ket{11}$ se compensent).}
  {$\ket{\Psi^+}$ est symétrique par échange des deux qubits, $\ket{\Psi^-}$ antisymétrique ; $H\otimes H$ commute avec le SWAP et conserve donc cette symétrie.}
  {C'est l'état que $H\otimes H$ laisse invariant ($\ket{\Phi^+}=\frac{1}{\sqrt2}(\ket{++}+\ket{--})$) ; le singulet n'a pas ce comportement.}

% ---- Question 98 (niveau 4, correct : B)
\qcmq[1]{Chap.~\ref{chap:bell}\,$\cdot$\,\S\,\ref{sec:chsh}}
  {Une expérience de type CHSH (avec $S=E_{00}+E_{01}+E_{10}-E_{11}$) mesure $S=2{,}6$. Que peut-on en conclure ?}
  {C'est impossible : on a toujours $|S|\leq2$}
  {Aucun modèle local à résultats prédéterminés ne l'explique, mais c'est compatible avec la mécanique quantique}
  {C'est impossible en mécanique quantique, car $2\sqrt2\approx2{,}4<2{,}6$}
  {Les particules s'échangent de l'information plus vite que la lumière}
  {B}
\qcmx
  {Cette borne ne vaut que pour les modèles classiques (locaux) ; l'expérience et la mécanique quantique la dépassent, jusqu'à $2\sqrt2$.}
  {Classique : $|S|\leq2$ ; quantique : $|S|\leq2\sqrt2\approx2{,}83$ (borne de Tsirelson). Comme $2<2{,}6<2\sqrt2$, l'inégalité de Bell est violée sans dépasser le maximum quantique ; le gain correspondant est $\frac12+\frac{2{,}6}{8}\approx0{,}825$.}
  {Erreur numérique : $2\sqrt2\approx2{,}83$, donc $2{,}6$ est permis.}
  {Les corrélations violent l'inégalité de Bell mais ne permettent aucun signal : la réponse de Bob est uniforme quelle que soit la base choisie par Alice.}

% ---- Question 99 (niveau 4, correct : A)
\qcmq[2]{Chap.~\ref{chap:bell}\,$\cdot$\,\S\,\ref{sec:info}}
  {Pour l'état $\ket{\Psi}=\frac{1}{\sqrt3}\big(\ket{00}+\sqrt2\,\ket{11}\big)$, quelle est l'entropie de von Neumann $S(\rho_A)$ du premier qubit, en bits ?}
  {$\log_2 3-\frac23\approx0{,}92$}
  {$1$}
  {$0$}
  {$0{,}64$}
  {A}
\qcmx
  {$\rho_A=\mathrm{diag}\big(\frac13,\frac23\big)$ (trace partielle), donc $S=-\frac13\log_2\frac13-\frac23\log_2\frac23=\log_23-\frac23\approx1{,}585-0{,}667$. C'est moins que le maximum d'un bit (état de Bell) : l'intrication est partielle.}
  {C'est la valeur pour un état de Bell ($\rho_A=\frac{\mathrm{Id}}{2}$) ; ici les poids $\frac13$ et $\frac23$ diffèrent, donc $S<1$.}
  {$S(\rho_A)=0$ vaut pour un état séparable ; ici $D=\frac{\sqrt2}{3}\neq0$ : l'état est intriqué, $\rho_A$ est mixte et $S>0$ (le couple $AB$ est pur, pas $A$ seul).}
  {On a utilisé le logarithme \emph{népérien} (résultat en nats) ; en bits il faut $\log_2$, ce qui donne $0{,}92$.}

% ---- Question 100 (niveau 4, correct : B)
\qcmq[1]{Chap.~\ref{chap:bell}\,$\cdot$\,\S\,\ref{sec:bell}}
  {Pourquoi peut-on mesurer simultanément les deux « parités » $Z\otimes Z$ et $X\otimes X$ de deux qubits, alors qu'on ne peut pas mesurer simultanément $Z$ et $X$ sur un seul qubit ?}
  {Parce que l'incertitude de Heisenberg ne s'applique pas à deux qubits}
  {Parce que $Z\otimes Z$ et $X\otimes X$ commutent : les deux signes $-1$ de $ZX=-XZ$ se compensent}
  {Parce que les deux qubits sont indépendants}
  {Parce que $Z\otimes Z$ et $X\otimes X$ agissent sur des qubits différents}
  {B}
\qcmx
  {Elle s'applique toujours : $Z\otimes\mathrm{Id}$ et $X\otimes\mathrm{Id}$ ne commutent pas. C'est la commutation de ces deux opérateurs précis qui permet la mesure simultanée.}
  {$(Z\otimes Z)(X\otimes X)=ZX\otimes ZX=(-XZ)\otimes(-XZ)=XZ\otimes XZ=(X\otimes X)(Z\otimes Z)$. Ils ont donc une base propre commune : la base de Bell.}
  {Dans un état de Bell les qubits sont intriqués, donc loin d'être indépendants ; la commutation est une propriété des opérateurs, pas de l'état.}
  {Chacun des deux opérateurs agit sur \emph{les deux} qubits à la fois. (Ce sont $Z\otimes\mathrm{Id}$ et $\mathrm{Id}\otimes X$, sur des qubits différents, qui commutent pour cette raison.)}

\label{qcm:fin}

\clearpage
\ifqcmcorrige

% ======================================================================
\subsection{Clé des réponses et fiche de score}\label{sec:qcm-cle}
% ======================================================================
\noindent\textbf{Cette partie donne les réponses : ne la lisez qu'après avoir répondu à toutes les
questions.}

\paragraph{Clé des réponses.} La lettre de la bonne réponse de la question $n$ se lit à la ligne
qui contient $n$, dans la colonne du rang de $n$ dans la ligne (question~24 : ligne 21--30,
colonne 4).
\begin{center}
\renewcommand{\arraystretch}{1.25}
\setlength{\tabcolsep}{5.5pt}
\begin{tabular}{@{}l*{10}{c}@{}}
\toprule
 & 1 & 2 & 3 & 4 & 5 & 6 & 7 & 8 & 9 & 10 \\
\midrule
1--10 & \qcmkey{1} & \qcmkey{2} & \qcmkey{3} & \qcmkey{4} & \qcmkey{5} & \qcmkey{6} & \qcmkey{7} & \qcmkey{8} & \qcmkey{9} & \qcmkey{10} \\
11--20 & \qcmkey{11} & \qcmkey{12} & \qcmkey{13} & \qcmkey{14} & \qcmkey{15} & \qcmkey{16} & \qcmkey{17} & \qcmkey{18} & \qcmkey{19} & \qcmkey{20} \\
21--30 & \qcmkey{21} & \qcmkey{22} & \qcmkey{23} & \qcmkey{24} & \qcmkey{25} & \qcmkey{26} & \qcmkey{27} & \qcmkey{28} & \qcmkey{29} & \qcmkey{30} \\
31--40 & \qcmkey{31} & \qcmkey{32} & \qcmkey{33} & \qcmkey{34} & \qcmkey{35} & \qcmkey{36} & \qcmkey{37} & \qcmkey{38} & \qcmkey{39} & \qcmkey{40} \\
41--50 & \qcmkey{41} & \qcmkey{42} & \qcmkey{43} & \qcmkey{44} & \qcmkey{45} & \qcmkey{46} & \qcmkey{47} & \qcmkey{48} & \qcmkey{49} & \qcmkey{50} \\
51--60 & \qcmkey{51} & \qcmkey{52} & \qcmkey{53} & \qcmkey{54} & \qcmkey{55} & \qcmkey{56} & \qcmkey{57} & \qcmkey{58} & \qcmkey{59} & \qcmkey{60} \\
61--70 & \qcmkey{61} & \qcmkey{62} & \qcmkey{63} & \qcmkey{64} & \qcmkey{65} & \qcmkey{66} & \qcmkey{67} & \qcmkey{68} & \qcmkey{69} & \qcmkey{70} \\
71--80 & \qcmkey{71} & \qcmkey{72} & \qcmkey{73} & \qcmkey{74} & \qcmkey{75} & \qcmkey{76} & \qcmkey{77} & \qcmkey{78} & \qcmkey{79} & \qcmkey{80} \\
81--90 & \qcmkey{81} & \qcmkey{82} & \qcmkey{83} & \qcmkey{84} & \qcmkey{85} & \qcmkey{86} & \qcmkey{87} & \qcmkey{88} & \qcmkey{89} & \qcmkey{90} \\
91--100 & \qcmkey{91} & \qcmkey{92} & \qcmkey{93} & \qcmkey{94} & \qcmkey{95} & \qcmkey{96} & \qcmkey{97} & \qcmkey{98} & \qcmkey{99} & \qcmkey{100} \\
\bottomrule
\end{tabular}
\end{center}

\paragraph{Fiche de score.} Reportez le nombre de bonnes réponses de chaque niveau.
\begin{center}
\renewcommand{\arraystretch}{1.6}
\begin{tabular}{@{}l l c c c@{}}
\toprule
\textbf{Niveau} & \textbf{Questions} & \textbf{Bonnes réponses} & \textbf{Sur} & \textbf{Pourcentage}\\
\midrule
1. Débutant & 1--30 & \makebox[2.6cm]{\dotfill} & 30 & \makebox[2cm]{\dotfill}\,\%\\
2. Intermédiaire & 31--60 & \makebox[2.6cm]{\dotfill} & 30 & \makebox[2cm]{\dotfill}\,\%\\
3. Avancé & 61--90 & \makebox[2.6cm]{\dotfill} & 30 & \makebox[2cm]{\dotfill}\,\%\\
4. Maîtrise & 91--100 & \makebox[2.6cm]{\dotfill} & 10 & \makebox[2cm]{\dotfill}\,\%\\
\midrule
\textbf{Total} & 1--100 & \makebox[2.6cm]{\dotfill} & 100 & \makebox[2cm]{\dotfill}\,\%\\
\bottomrule
\end{tabular}
\end{center}
Repères : au moins $80\,\%$ à un niveau, il est acquis et on peut passer au suivant ; entre $50$ et
$80\,\%$, revoyez les sections des questions manquées ; en dessous, relisez le chapitre concerné et
son récapitulatif avant de recommencer.

\paragraph{Où revoir ?} Questions classées par chapitre du cours.
\begin{center}
\renewcommand{\arraystretch}{1.3}
\begin{tabular}{@{}l >{\raggedright\arraybackslash}p{5.9cm} >{\raggedright\arraybackslash}p{6.9cm}@{}}
\toprule
\textbf{Chapitre} & \textbf{Titre} & \textbf{Questions}\\
\midrule
Chap.~\ref{chap:histoire} & Histoire et mécanique quantique & 1, 2, 3, 4, 5, 6, 31, 32, 33, 34, 35, 61, 62, 63, 91 \\
Chap.~\ref{chap:dirac} & Notation de Dirac et espace de Hilbert & 7, 8, 9, 10, 11, 12, 36, 37, 38, 39, 40, 64, 65, 66, 67, 68, 92 \\
Chap.~\ref{chap:operateurs} & Opérateurs, matrice de densité, Bloch, produit tensoriel & 13, 14, 15, 16, 17, 18, 41, 42, 43, 44, 45, 46, 69, 70, 71, 72, 73, 74, 75, 93, 94 \\
Chap.~\ref{chap:portes} & Portes quantiques et systèmes multi-qubits & 19, 20, 21, 22, 23, 24, 47, 48, 49, 50, 51, 76, 77, 78, 79, 80, 81, 82, 95, 96 \\
Chap.~\ref{chap:bell} & États quantiques, intrication, Bell & 25, 26, 27, 28, 29, 30, 52, 53, 54, 55, 56, 57, 58, 59, 60, 83, 84, 85, 86, 87, 88, 89, 90, 97, 98, 99, 100 \\
\bottomrule
\end{tabular}
\end{center}

\clearpage
\subsection{Corrigé du niveau 1 : débutant (questions 1 à 30)}\label{sec:qcm-c1}
\qcmcorr{1}
\qcmcorr{2}
\qcmcorr{3}
\qcmcorr{4}
\qcmcorr{5}
\qcmcorr{6}
\qcmcorr{7}
\qcmcorr{8}
\qcmcorr{9}
\qcmcorr{10}
\qcmcorr{11}
\qcmcorr{12}
\qcmcorr{13}
\qcmcorr{14}
\qcmcorr{15}
\qcmcorr{16}
\qcmcorr{17}
\qcmcorr{18}
\qcmcorr{19}
\qcmcorr{20}
\qcmcorr{21}
\qcmcorr{22}
\qcmcorr{23}
\qcmcorr{24}
\qcmcorr{25}
\qcmcorr{26}
\qcmcorr{27}
\qcmcorr{28}
\qcmcorr{29}
\qcmcorr{30}

\clearpage
\subsection{Corrigé du niveau 2 : intermédiaire (questions 31 à 60)}\label{sec:qcm-c2}
\qcmcorr{31}
\qcmcorr{32}
\qcmcorr{33}
\qcmcorr{34}
\qcmcorr{35}
\qcmcorr{36}
\qcmcorr{37}
\qcmcorr{38}
\qcmcorr{39}
\qcmcorr{40}
\qcmcorr{41}
\qcmcorr{42}
\qcmcorr{43}
\qcmcorr{44}
\qcmcorr{45}
\qcmcorr{46}
\qcmcorr{47}
\qcmcorr{48}
\qcmcorr{49}
\qcmcorr{50}
\qcmcorr{51}
\qcmcorr{52}
\qcmcorr{53}
\qcmcorr{54}
\qcmcorr{55}
\qcmcorr{56}
\qcmcorr{57}
\qcmcorr{58}
\qcmcorr{59}
\qcmcorr{60}

\clearpage
\subsection{Corrigé du niveau 3 : avancé (questions 61 à 90)}\label{sec:qcm-c3}
\qcmcorr{61}
\qcmcorr{62}
\qcmcorr{63}
\qcmcorr{64}
\qcmcorr{65}
\qcmcorr{66}
\qcmcorr{67}
\qcmcorr{68}
\qcmcorr{69}
\qcmcorr{70}
\qcmcorr{71}
\qcmcorr{72}
\qcmcorr{73}
\qcmcorr{74}
\qcmcorr{75}
\qcmcorr{76}
\qcmcorr{77}
\qcmcorr{78}
\qcmcorr{79}
\qcmcorr{80}
\qcmcorr{81}
\qcmcorr{82}
\qcmcorr{83}
\qcmcorr{84}
\qcmcorr{85}
\qcmcorr{86}
\qcmcorr{87}
\qcmcorr{88}
\qcmcorr{89}
\qcmcorr{90}

\clearpage
\subsection{Corrigé du niveau 4 : maîtrise (questions 91 à 100)}\label{sec:qcm-c4}
\qcmcorr{91}
\qcmcorr{92}
\qcmcorr{93}
\qcmcorr{94}
\qcmcorr{95}
\qcmcorr{96}
\qcmcorr{97}
\qcmcorr{98}
\qcmcorr{99}
\qcmcorr{100}
\fi
       % 7. QCM sur le cours (100 questions, énoncés puis corrigé)
%% =====================================================================
%  Chapitre 7 : portes multi-qubits, circuits quantiques et téléportation
%  Fichier importé par main.tex avec % =====================================================================
%  Chapitre 7 : portes multi-qubits, circuits quantiques et téléportation
%  Fichier importé par main.tex avec % =====================================================================
%  Chapitre 7 : portes multi-qubits, circuits quantiques et téléportation
%  Fichier importé par main.tex avec \include{chapitre7}.
%  Comme chapitre4.tex, chapitre5.tex et chapitre6.tex, il ouvre une NOUVELLE
%  SECTION (numéro 5). Pour le rattacher à la section précédente, remplacer
%  \section par \subsection et chaque \subsection par \subsubsection.
%  Pas de préambule : environnements (definition, resultat, remarque),
%  macros (\ii, \ee, \dd, \ket, \bra, \braket) et couleurs (insarouge,
%  insagris) sont définis dans main.tex. Les symboles de circuits (contrôle,
%  cible, croix, mesure) sont définis juste après avec \tikzset.
% =====================================================================

% Macros de Dirac (déjà définies dans main.tex récent ; ce bloc évite l'erreur
% « Undefined control sequence \ket » si main.tex est une ancienne version).
\providecommand{\ket}[1]{|#1\rangle}
\providecommand{\bra}[1]{\langle #1|}
\providecommand{\braket}[2]{\langle #1|#2\rangle}

% Symboles de circuits pour les figures de ce chapitre.
\tikzset{
  ctl/.pic   = {\fill[insarouge] (0,0) circle (2.6pt);},
  cible/.pic = {\draw[insarouge, thick, fill=white] (0,0) circle (0.28);
                \draw[insarouge, thick] (-0.28,0) -- (0.28,0);
                \draw[insarouge, thick] (0,-0.28) -- (0,0.28);},
  croix/.pic = {\draw[insarouge, thick] (-0.17,-0.17) -- (0.17,0.17);
                \draw[insarouge, thick] (-0.17,0.17) -- (0.17,-0.17);},
  mesure/.pic = {\draw[thick, fill=white] (-0.3,-0.27) rectangle (0.3,0.27);
                 \draw[thick] (-0.19,-0.13) arc[start angle=155, end angle=25, radius=0.21];
                 \draw[thick, -{Stealth[length=1.3mm]}] (0,-0.17) -- (0.14,0.12);},
  porte/.style = {draw=insarouge, fill=insarouge!8, minimum size=7mm, thick},
}

% =====================================================================
\section{Portes multi-qubits, circuits et téléportation}\label{sec:circuits}
% =====================================================================

La section précédente a décrit les états de deux qubits. On étudie ici les portes qui
agissent sur plusieurs qubits : le CNOT et le SWAP, puis les portes contrôlées à deux et à
trois qubits (contrôlé-$U$, Fredkin, Toffoli). On assemble ensuite des portes dans un
circuit qui fabrique un état de Bell, on mesure l'information contenue dans un tel état,
on distingue les états purs des états mélangés, et on démontre le protocole de
téléportation. Tous les calculs sont détaillés.

\begin{remarque}[Convention pour les circuits]
Dans un circuit, le temps va de gauche à droite et chaque qubit a un fil. Les fils sont
dessinés dans l'ordre \emph{inverse} de l'écriture des kets : le fil du \emph{haut} porte le
\emph{dernier} chiffre du ket, le fil du \emph{bas} le premier. Par exemple, si le qubit $A$
est dessiné au-dessus du qubit $B$, le ket s'écrit $\ket{b,a}=\ket{b}_B\otimes\ket{a}_A$.
La figure~\ref{fig:bell} de la section précédente suit la même convention.
\end{remarque}

% ---------------------------------------------------------------------
\subsection{Porte CNOT}\label{sec:cnot}
% ---------------------------------------------------------------------

\begin{definition}{Porte CNOT (NOT contrôlé)}
La porte CNOT agit sur deux qubits : le qubit $A$ est le \emph{contrôle}, le qubit $B$ est la
\emph{cible}. La cible est inversée (porte $X$) si le contrôle vaut $1$ et reste inchangée
s'il vaut $0$. Avec le système écrit $(B,A)$, soit $\ket{b,a}=\ket{b}_B\otimes\ket{a}_A$ :
\[
  \mathrm{CNOT}\,\ket{b,a}=\ket{b\oplus a,\,a},
  \qquad
  \mathrm{CNOT}=I\otimes\ket{0}\bra{0}+X\otimes\ket{1}\bra{1},
\]
où $\oplus$ est l'addition modulo $2$ ($0\oplus0=1\oplus1=0$ et $0\oplus1=1\oplus0=1$).
\end{definition}

Sur le circuit de la figure~\ref{fig:cnot}, le point noir posé sur le fil de $A$ marque le
contrôle et le symbole $\oplus$ posé sur le fil de $B$ marque la cible.

\begin{figure}[!ht]
\centering
\begin{tikzpicture}[>=Stealth, font=\small]
  \draw (0,1.2) -- (6,1.2);
  \draw (0,0) -- (6,0);
  \node[anchor=east] at (0,1.2) {$\ket{a}_A$};
  \node[anchor=east] at (0,0) {$\ket{b}_B$};
  \node[anchor=west] at (6,1.2) {$\ket{a}_A$};
  \node[anchor=west] at (6,0) {$\ket{b\oplus a}_B$};
  \draw[insarouge, thick] (3,1.2) -- (3,0);
  \pic at (3,1.2) {ctl};
  \pic at (3,0) {cible};
  \node[anchor=south, text=insagris] at (3,1.5) {qubit de contrôle};
  \node[anchor=north, text=insagris] at (3,-0.5) {qubit cible};
\end{tikzpicture}
\caption{Porte CNOT : le fil du haut ($A$) est le contrôle, le fil du bas ($B$) est la
cible. Le ket s'écrit $\ket{b,a}$ (le fil du haut donne le dernier chiffre).}
\label{fig:cnot}
\end{figure}

\subsubsection{Table de vérité}

\begin{center}
\renewcommand{\arraystretch}{1.3}
\begin{tabular}{@{}cccc@{}}
\toprule
Entrée $\ket{b,a}$ & Contrôle $a$ & Sortie $\ket{b\oplus a,a}$ & Effet sur la cible $B$ \\
\midrule
$\ket{00}$ & $0$ & $\ket{00}$ & aucun \\
$\ket{01}$ & $1$ & $\ket{11}$ & $0\to1$ \\
$\ket{10}$ & $0$ & $\ket{10}$ & aucun \\
$\ket{11}$ & $1$ & $\ket{01}$ & $1\to0$ \\
\bottomrule
\end{tabular}
\end{center}

\subsubsection{Matrice du CNOT}

La $j$-ième colonne d'une matrice $4\times4$ est l'image du $j$-ième vecteur de base
($\ket{00},\ket{01},\ket{10},\ket{11}$, dans cet ordre). La table donne :
\begin{align*}
  \mathrm{CNOT}\ket{00}&=\ket{00}=\begin{pmatrix}1\\0\\0\\0\end{pmatrix},&
  \mathrm{CNOT}\ket{01}&=\ket{11}=\begin{pmatrix}0\\0\\0\\1\end{pmatrix},\\[1.5ex]
  \mathrm{CNOT}\ket{10}&=\ket{10}=\begin{pmatrix}0\\0\\1\\0\end{pmatrix},&
  \mathrm{CNOT}\ket{11}&=\ket{01}=\begin{pmatrix}0\\1\\0\\0\end{pmatrix},
\end{align*}
d'où, en juxtaposant ces quatre colonnes,
\[
  \mathrm{CNOT}=\begin{pmatrix}
    1 & 0 & 0 & 0\\ 0 & 0 & 0 & 1\\ 0 & 0 & 1 & 0\\ 0 & 1 & 0 & 0
  \end{pmatrix}.
\]

On retrouve cette matrice avec la formule $I\otimes\ket{0}\bra{0}+X\otimes\ket{1}\bra{1}$ et
la règle $A\otimes B=(A_{ij}B)$ (section~\ref{sec:op2}), où
$\ket{0}\bra{0}=\begin{pmatrix}1&0\\0&0\end{pmatrix}$ et
$\ket{1}\bra{1}=\begin{pmatrix}0&0\\0&1\end{pmatrix}$ :
\[
  I\otimes\ket{0}\bra{0}=\begin{pmatrix}1&0&0&0\\0&0&0&0\\0&0&1&0\\0&0&0&0\end{pmatrix},
  \qquad
  X\otimes\ket{1}\bra{1}=\begin{pmatrix}0&0&0&0\\0&0&0&1\\0&0&0&0\\0&1&0&0\end{pmatrix},
\]
et la somme de ces deux matrices est la matrice du CNOT.

\subsubsection{Action sur un état quelconque}

Pour $\ket{\Psi}=c_{00}\ket{00}+c_{01}\ket{01}+c_{10}\ket{10}+c_{11}\ket{11}$ :
\[
  \mathrm{CNOT}\begin{pmatrix}c_{00}\\c_{01}\\c_{10}\\c_{11}\end{pmatrix}
  =\begin{pmatrix}
    1\cdot c_{00}+0\cdot c_{01}+0\cdot c_{10}+0\cdot c_{11}\\
    0\cdot c_{00}+0\cdot c_{01}+0\cdot c_{10}+1\cdot c_{11}\\
    0\cdot c_{00}+0\cdot c_{01}+1\cdot c_{10}+0\cdot c_{11}\\
    0\cdot c_{00}+1\cdot c_{01}+0\cdot c_{10}+0\cdot c_{11}
  \end{pmatrix}
  =\begin{pmatrix}c_{00}\\c_{11}\\c_{10}\\c_{01}\end{pmatrix}.
\]
Le CNOT échange les amplitudes de $\ket{01}$ et de $\ket{11}$, c'est-à-dire des deux états
dont le contrôle vaut $1$, et laisse les autres.

\paragraph{Exemple : contrôle en superposition.} Soit $\ket{0}_B\otimes\ket{+}_A
=\frac{1}{\sqrt2}\big(\ket{00}+\ket{01}\big)=\frac{1}{\sqrt2}(1,1,0,0)^{\mathsf T}$,
un état séparable. Alors
\[
  \mathrm{CNOT}\,\frac{1}{\sqrt2}\begin{pmatrix}1\\1\\0\\0\end{pmatrix}
  =\frac{1}{\sqrt2}\begin{pmatrix}1\\0\\0\\1\end{pmatrix}
  =\frac{1}{\sqrt2}\big(\ket{00}+\ket{11}\big)=\ket{\Phi^+}.
\]
Le résultat est intriqué ($D=c_{00}c_{11}-c_{01}c_{10}=\frac12\neq0$ ; section~\ref{sec:separable}) :
un CNOT dont le contrôle est en superposition crée de l'intrication.

\paragraph{Réversibilité.} Deux CNOT successifs redonnent l'état initial :
$\ket{b,a}\to\ket{b\oplus a,a}\to\ket{b\oplus a\oplus a,a}=\ket{b,a}$, donc
$\mathrm{CNOT}^2=\mathrm{Id}$. La matrice est réelle et symétrique, donc
$\mathrm{CNOT}^\dagger\mathrm{CNOT}=\mathrm{CNOT}^2=\mathrm{Id}$ : elle est bien unitaire.

\begin{remarque}[Contrôle en dernier ou en premier]
La matrice dépend de l'ordre d'écriture du ket. Notons $\mathrm{CNOT}_{21}$ le CNOT dont le
contrôle est le \emph{second} chiffre du ket (c'est la matrice ci-dessus) et
$\mathrm{CNOT}_{12}$ celui dont le contrôle est le \emph{premier} chiffre, utilisé à la
section précédente :
\[
  \mathrm{CNOT}_{12}=\ket{0}\bra{0}\otimes I+\ket{1}\bra{1}\otimes X
  =\begin{pmatrix}1&0&0&0\\0&1&0&0\\0&0&0&1\\0&0&1&0\end{pmatrix}.
\]
Un même circuit (contrôle $A$, cible $B$) a donc deux matrices, selon que l'on écrit
$A$ en dernier ($\mathrm{CNOT}_{21}$) ou en premier ($\mathrm{CNOT}_{12}$). Les deux sont
reliées par le SWAP (section suivante).
\end{remarque}

% ---------------------------------------------------------------------
\subsection{Porte SWAP}\label{sec:swap}
% ---------------------------------------------------------------------

\begin{definition}{Porte SWAP (échange)}
La porte SWAP échange les états des deux qubits :
\[
  \mathrm{SWAP}\big(\ket{\psi}\otimes\ket{\phi}\big)=\ket{\phi}\otimes\ket{\psi}.
\]
Sur les états de base, $\mathrm{SWAP}\ket{a,b}=\ket{b,a}$ ; on prolonge par linéarité.
\end{definition}

Son symbole est formé de deux croix reliées par un trait vertical
(figure~\ref{fig:swap}).

\begin{figure}[!ht]
\centering
\begin{tikzpicture}[>=Stealth, font=\small]
  \draw (0,1.2) -- (6,1.2);
  \draw (0,0) -- (6,0);
  \node[anchor=east] at (0,1.2) {$\ket{a}$};
  \node[anchor=east] at (0,0) {$\ket{b}$};
  \node[anchor=west] at (6,1.2) {$\ket{b}$};
  \node[anchor=west] at (6,0) {$\ket{a}$};
  \draw[insarouge, thick] (3,1.2) -- (3,0);
  \pic at (3,1.2) {croix};
  \pic at (3,0) {croix};
\end{tikzpicture}
\caption{Porte SWAP : les deux qubits échangent leurs états.}
\label{fig:swap}
\end{figure}

\subsubsection{Table et matrice}

\begin{center}
\renewcommand{\arraystretch}{1.3}
\begin{tabular}{@{}cc@{}}
\toprule
Entrée & Sortie \\
\midrule
$\ket{00}$ & $\ket{00}$ \\
$\ket{01}$ & $\ket{10}$ \\
$\ket{10}$ & $\ket{01}$ \\
$\ket{11}$ & $\ket{11}$ \\
\bottomrule
\end{tabular}
\end{center}

Les images des quatre vecteurs de base sont les colonnes de la matrice :
\[
  \mathrm{SWAP}\ket{00}=\begin{pmatrix}1\\0\\0\\0\end{pmatrix},\quad
  \mathrm{SWAP}\ket{01}=\begin{pmatrix}0\\0\\1\\0\end{pmatrix},\quad
  \mathrm{SWAP}\ket{10}=\begin{pmatrix}0\\1\\0\\0\end{pmatrix},\quad
  \mathrm{SWAP}\ket{11}=\begin{pmatrix}0\\0\\0\\1\end{pmatrix},
\]
\[
  \mathrm{SWAP}=\begin{pmatrix}
    1 & 0 & 0 & 0\\ 0 & 0 & 1 & 0\\ 0 & 1 & 0 & 0\\ 0 & 0 & 0 & 1
  \end{pmatrix}.
\]
Appliquée à $\ket{01}=(0,1,0,0)^{\mathsf T}$, cette matrice sélectionne sa deuxième colonne :
\[
  \begin{pmatrix}
    1 & 0 & 0 & 0\\ 0 & 0 & 1 & 0\\ 0 & 1 & 0 & 0\\ 0 & 0 & 0 & 1
  \end{pmatrix}
  \begin{pmatrix}0\\1\\0\\0\end{pmatrix}
  =\begin{pmatrix}0\\0\\1\\0\end{pmatrix}=\ket{10}.
\]
Sur un état quelconque, $\mathrm{SWAP}\,(c_{00},c_{01},c_{10},c_{11})^{\mathsf T}
=(c_{00},c_{10},c_{01},c_{11})^{\mathsf T}$ : les amplitudes de $\ket{01}$ et $\ket{10}$
sont échangées.

\subsubsection{Accord avec la définition}

Pour deux états $\ket{\psi}=(a_0,a_1)^{\mathsf T}$ et $\ket{\phi}=(b_0,b_1)^{\mathsf T}$ :
\[
  \mathrm{SWAP}\,(\ket{\psi}\otimes\ket{\phi})
  =\mathrm{SWAP}\begin{pmatrix}a_0b_0\\a_0b_1\\a_1b_0\\a_1b_1\end{pmatrix}
  =\begin{pmatrix}a_0b_0\\a_1b_0\\a_0b_1\\a_1b_1\end{pmatrix}
  =\begin{pmatrix}b_0\\b_1\end{pmatrix}\otimes\begin{pmatrix}a_0\\a_1\end{pmatrix}
  =\ket{\phi}\otimes\ket{\psi},
\]
ce qui est bien la définition. La porte s'écrit aussi avec des projecteurs :
\[
  \mathrm{SWAP}=\sum_{a,b\in\{0,1\}}\ket{a}\bra{b}\otimes\ket{b}\bra{a}.
\]
En effet, pour $\ket{c}\otimes\ket{d}$ on trouve
$\sum_{a,b}\ket{a}\braket{b}{c}\otimes\ket{b}\braket{a}{d}$, et seul le terme $b=c$, $a=d$
est non nul : le résultat est $\ket{d}\otimes\ket{c}$.

\paragraph{Exemples.}
\begin{itemize}[nosep]
  \item Deux états différents : $\mathrm{SWAP}\big(\ket{0}\otimes\ket{+}\big)
        =\mathrm{SWAP}\,\frac{1}{\sqrt2}(1,1,0,0)^{\mathsf T}=\frac{1}{\sqrt2}(1,0,1,0)^{\mathsf T}
        =\ket{+}\otimes\ket{0}$.
  \item L'état $\ket{\psi}=\frac{1}{\sqrt2}\ket{00}+\frac{1}{\sqrt2}\ket{11}$ ne change pas :
        $\mathrm{SWAP}\ket{\psi}=\frac{1}{\sqrt2}\ket{00}+\frac{1}{\sqrt2}\ket{11}=\ket{\psi}$,
        car $\ket{00}$ et $\ket{11}$ sont invariants. De même
        $\mathrm{SWAP}\ket{\Phi^-}=\ket{\Phi^-}$ et $\mathrm{SWAP}\ket{\Psi^+}=\ket{\Psi^+}$, alors que
        $\mathrm{SWAP}\ket{\Psi^-}=\frac{1}{\sqrt2}(\ket{10}-\ket{01})=-\ket{\Psi^-}$.
\end{itemize}

\paragraph{Propriétés.} Échanger deux fois redonne l'état de départ :
$\mathrm{SWAP}^2=\mathrm{Id}$ ; la matrice est réelle et symétrique, donc
$\mathrm{SWAP}^\dagger\mathrm{SWAP}=\mathrm{SWAP}^2=\mathrm{Id}$ : elle est unitaire.

\subsubsection{Lien avec le CNOT}

\begin{resultat}{SWAP échange les rôles de contrôle et de cible}
\[
  \mathrm{CNOT}_{21}=\mathrm{SWAP}\;\mathrm{CNOT}_{12}\;\mathrm{SWAP}.
\]
\end{resultat}

\emph{Sur les états de base.}
$\mathrm{SWAP}\,\mathrm{CNOT}_{12}\,\mathrm{SWAP}\ket{b,a}
=\mathrm{SWAP}\,\mathrm{CNOT}_{12}\ket{a,b}
=\mathrm{SWAP}\ket{a,\,a\oplus b}=\ket{a\oplus b,\,a}=\mathrm{CNOT}_{21}\ket{b,a}$.

\emph{Par les matrices.} Multiplier à gauche par $\mathrm{SWAP}$ échange les lignes $2$ et
$3$, multiplier à droite l'échange les colonnes $2$ et $3$ :
\[
  \mathrm{SWAP}\,\mathrm{CNOT}_{12}
  =\begin{pmatrix}1&0&0&0\\0&0&0&1\\0&1&0&0\\0&0&1&0\end{pmatrix},
  \qquad
  \mathrm{SWAP}\,\mathrm{CNOT}_{12}\,\mathrm{SWAP}
  =\begin{pmatrix}1&0&0&0\\0&0&0&1\\0&0&1&0\\0&1&0&0\end{pmatrix}=\mathrm{CNOT}_{21}.
\]

\begin{remarque}[SWAP avec trois CNOT]
$\mathrm{SWAP}=\mathrm{CNOT}_{12}\,\mathrm{CNOT}_{21}\,\mathrm{CNOT}_{12}$. En effet, sur
$\ket{a,b}$ (on applique les portes de droite à gauche) :
\[
  \ket{a,b}\ \xrightarrow{\ \mathrm{CNOT}_{12}\ }\ \ket{a,\,a\oplus b}
  \ \xrightarrow{\ \mathrm{CNOT}_{21}\ }\ \ket{a\oplus(a\oplus b),\,a\oplus b}=\ket{b,\,a\oplus b}
  \ \xrightarrow{\ \mathrm{CNOT}_{12}\ }\ \ket{b,\,b\oplus(a\oplus b)}=\ket{b,a}.
\]
Les deux membres coïncident sur les états de base, donc sur tous les états par linéarité.
\end{remarque}

% ---------------------------------------------------------------------
\subsection{Un circuit quantique : le générateur d'états de Bell}\label{sec:circuit-bell}
% ---------------------------------------------------------------------

On assemble une porte de Hadamard et un CNOT. Le système est $(A,B)$ avec le ket
$\ket{a,b}=\ket{a}_A\otimes\ket{b}_B$ ; d'après la convention de la section, $B$ est dessiné
en haut et $A$ en bas (figure~\ref{fig:circuit-bell}). Le circuit a deux étapes :
\begin{enumerate}[nosep]
  \item la porte $I\otimes H$ : identité sur $A$, Hadamard sur $B$ ;
  \item un CNOT de contrôle $B$ et de cible $A$ (la matrice $\mathrm{CNOT}_{21}$ de la
        section~\ref{sec:cnot}, puisque le contrôle est le second chiffre).
\end{enumerate}
Les traits pointillés délimitent les états $\ket{\Pi_0}$, $\ket{\Pi_1}$ et $\ket{\Pi_2}$ :
avant les portes, après $I\otimes H$, après le CNOT.

\begin{figure}[!ht]
\centering
\begin{tikzpicture}[>=Stealth, font=\small]
  % fils : B en haut, A en bas
  \draw (0,1.4) -- (7.6,1.4);
  \draw (0,0) -- (7.6,0);
  \node[anchor=east] at (0,1.4) {$\ket{0}_B$};
  \node[anchor=east] at (0,0) {$\ket{0}_A$};
  % première étape : H sur B, identité (pointillée) sur A
  \node[porte] at (2.2,1.4) {$H$};
  \node[draw=insagris, dashed, fill=white, minimum size=7mm] at (2.2,0) {$I$};
  % CNOT : contrôle B (haut), cible A (bas)
  \draw[insarouge, thick] (5,1.4) -- (5,0);
  \pic at (5,1.4) {ctl};
  \pic at (5,0) {cible};
  % accolade de sortie
  \draw[insagris, thick] (7.75,1.6) -- (7.9,1.6) -- (7.9,-0.2) -- (7.75,-0.2);
  \node[anchor=west] at (7.9,0.7) {$\ket{\beta^{00}}_{AB}$};
  % états intermédiaires
  \foreach \x in {0.7,3.6,6.4}
    \draw[dashed, insagris!70] (\x,1.9) -- (\x,-0.7);
  \node[anchor=north, font=\footnotesize] at (0.7,-0.7) {$\ket{\Pi_0}=\ket{00}$};
  \node[anchor=north, font=\footnotesize] at (3.6,-0.7) {$\ket{\Pi_1}=\ket{0}\otimes\ket{+}$};
  \node[anchor=north, font=\footnotesize] at (6.4,-0.7) {$\ket{\Pi_2}=\ket{\beta^{00}}$};
\end{tikzpicture}
\caption{Générateur d'états de Bell : $I\otimes H$, puis un CNOT de contrôle $B$ (fil du haut)
et de cible $A$ (fil du bas). L'entrée représentée est $\ket{00}$.}
\label{fig:circuit-bell}
\end{figure}

\subsubsection{Entrée \texorpdfstring{$\ket{00}$}{|00>} : calcul pas à pas}

\begin{align*}
  \ket{\Pi_0}&=\ket{0}_A\otimes\ket{0}_B=\ket{00},\\[1ex]
  \ket{\Pi_1}&=(I\otimes H)\ket{\Pi_0}=I\ket{0}_A\otimes H\ket{0}_B=\ket{0}\otimes\ket{+}
    =\frac{\ket{00}+\ket{01}}{\sqrt2},\\[1ex]
  \ket{\Pi_2}&=\mathrm{CNOT}\ket{\Pi_1}=\frac{\mathrm{CNOT}\ket{00}+\mathrm{CNOT}\ket{01}}{\sqrt2}
    =\frac{\ket{00}+\ket{11}}{\sqrt2}=\ket{\beta^{00}}=\ket{\Phi^+}.
\end{align*}
Dans $\ket{\Pi_1}$, le premier chiffre est celui de $A$ et le second celui de $B$. Le CNOT a
pour contrôle $B$ : $\ket{00}$ ($B=0$) ne change pas, $\ket{01}$ ($B=1$) devient $\ket{11}$
puisque $A$ est inversé. Le résultat est un \emph{état de Bell}.

Avec les matrices, la règle $A\otimes B=(A_{ij}B)$ donne
\[
  I\otimes H=\begin{pmatrix}1\cdot H&0\\0&1\cdot H\end{pmatrix}
  =\frac{1}{\sqrt2}\begin{pmatrix}1&1&0&0\\1&-1&0&0\\0&0&1&1\\0&0&1&-1\end{pmatrix}.
\]
Sa première colonne est $\ket{\Pi_1}=\frac{1}{\sqrt2}(1,1,0,0)^{\mathsf T}$ (image de $\ket{00}$).
Le CNOT échange les amplitudes de $\ket{01}$ et de $\ket{11}$ (section~\ref{sec:cnot}), donc
\[
  \ket{\Pi_2}=\mathrm{CNOT}\,\frac{1}{\sqrt2}\begin{pmatrix}1\\1\\0\\0\end{pmatrix}
  =\frac{1}{\sqrt2}\begin{pmatrix}1\\0\\0\\1\end{pmatrix}.
\]

\subsubsection{Les quatre entrées}

Prenons $B$ dans l'état $\ket{x}$ et $A$ dans l'état $\ket{y}$, avec $x,y\in\{0,1\}$. Alors
$\ket{\Pi_0}=\ket{y}_A\otimes\ket{x}_B=\ket{y,x}$. Comme
$H\ket{x}=\frac{1}{\sqrt2}\big(\ket{0}+(-1)^x\ket{1}\big)$ (c'est $\ket{+}$ pour $x=0$ et
$\ket{-}$ pour $x=1$),
\[
  \ket{\Pi_1}=\ket{y}\otimes H\ket{x}=\frac{1}{\sqrt2}\Big(\ket{y,0}+(-1)^x\ket{y,1}\Big).
\]
Le CNOT (contrôle $B$) ne change pas $\ket{y,0}$ et transforme $\ket{y,1}$ en
$\ket{\bar y,1}$, avec $\bar y=1-y$ :
\[
  \ket{\Pi_2}=\ket{\beta^{xy}}=\frac{1}{\sqrt2}\Big(\ket{y,0}+(-1)^x\ket{\bar y,1}\Big).
\]
Le premier indice de $\ket{\beta^{xy}}$ est l'entrée de $B$ (le qubit qui reçoit $H$), le
second celle de $A$ (la cible). Les quatre cas :

\begin{center}
\renewcommand{\arraystretch}{1.7}
\begin{tabular}{@{}ccccc@{}}
\toprule
$(B,A)=(x,y)$ & $\ket{\Pi_0}=\ket{y,x}$ & $\ket{\Pi_1}$ & $\ket{\Pi_2}=\ket{\beta^{xy}}$ & Nom \\
\midrule
$(0,0)$ & $\ket{00}$ & $\frac{1}{\sqrt2}(\ket{00}+\ket{01})$ & $\frac{1}{\sqrt2}(\ket{00}+\ket{11})$ & $\ket{\Phi^+}$ \\
$(1,0)$ & $\ket{01}$ & $\frac{1}{\sqrt2}(\ket{00}-\ket{01})$ & $\frac{1}{\sqrt2}(\ket{00}-\ket{11})$ & $\ket{\Phi^-}$ \\
$(0,1)$ & $\ket{10}$ & $\frac{1}{\sqrt2}(\ket{10}+\ket{11})$ & $\frac{1}{\sqrt2}(\ket{10}+\ket{01})$ & $\ket{\Psi^+}$ \\
$(1,1)$ & $\ket{11}$ & $\frac{1}{\sqrt2}(\ket{10}-\ket{11})$ & $\frac{1}{\sqrt2}(\ket{10}-\ket{01})$ & $-\ket{\Psi^-}$ \\
\bottomrule
\end{tabular}
\end{center}
Détail de la dernière ligne : $\ket{\Pi_1}=\ket{1}\otimes\ket{-}=\frac{1}{\sqrt2}(\ket{10}-\ket{11})$ ;
le CNOT laisse $\ket{10}$ et envoie $\ket{11}$ sur $\ket{01}$, d'où
$\frac{1}{\sqrt2}(\ket{10}-\ket{01})=-\frac{1}{\sqrt2}(\ket{01}-\ket{10})=-\ket{\Psi^-}$.

\paragraph{Forme matricielle.} $\mathrm{CNOT}_{21}$ échange les lignes $2$ et $4$ de
$I\otimes H$ (il échange les amplitudes de $\ket{01}$ et $\ket{11}$), donc
\[
  \mathrm{CNOT}_{21}\,(I\otimes H)=\frac{1}{\sqrt2}\begin{pmatrix}
    1 & 1 & 0 & 0\\ 0 & 0 & 1 & -1\\ 0 & 0 & 1 & 1\\ 1 & -1 & 0 & 0
  \end{pmatrix}.
\]
Ses colonnes sont les images de $\ket{00},\ket{01},\ket{10},\ket{11}$, c'est-à-dire de
$(B,A)=(0,0),(1,0),(0,1),(1,1)$ : ce sont $\ket{\beta^{00}}$, $\ket{\beta^{10}}$,
$\ket{\beta^{01}}$, $\ket{\beta^{11}}$, soit $\ket{\Phi^+}$, $\ket{\Phi^-}$, $\ket{\Psi^+}$
et $-\ket{\Psi^-}$.

\begin{remarque}[Comparaison avec la section précédente]
Le circuit de la figure~\ref{fig:bell} applique $H$ au \emph{premier} chiffre et prend celui-ci
pour contrôle ; ici $H$ et le contrôle portent sur le \emph{dernier} chiffre. Les deux
circuits se déduisent l'un de l'autre en échangeant les deux fils (conjugaison par le SWAP,
section~\ref{sec:swap}). Les quatre états de Bell obtenus sont les mêmes, avec les mêmes
indices, sauf le dernier : $-\ket{\Psi^-}$ ici, $+\ket{\Psi^-}$ là-bas. Ce signe est une phase
globale, sans effet sur les mesures (elle disparaît de $\rho$, section~\ref{sec:rho}).
\end{remarque}

\subsubsection{Matrice de densité de l'état de sortie}

Pour $\ket{\Psi}=\ket{\beta^{00}}=\frac{1}{\sqrt2}\ket{00}+\frac{1}{\sqrt2}\ket{11}$ :
\[
  \rho_{AB}=\ket{\Psi}\bra{\Psi}
  =\begin{pmatrix}\frac{1}{\sqrt2}\\0\\0\\\frac{1}{\sqrt2}\end{pmatrix}
   \begin{pmatrix}\frac{1}{\sqrt2}&0&0&\frac{1}{\sqrt2}\end{pmatrix}
  =\frac12\begin{pmatrix}1&0&0&1\\0&0&0&0\\0&0&0&0\\1&0&0&1\end{pmatrix}.
\]
Sous forme de ket-bra :
$\rho_{AB}=\frac12\big(\ket{00}\bra{00}+\ket{00}\bra{11}+\ket{11}\bra{00}+\ket{11}\bra{11}\big)$.
On vérifie $\operatorname{Tr}\rho_{AB}=\frac12(1+1)=1$ et
\[
  \rho_{AB}^2=\frac14\begin{pmatrix}2&0&0&2\\0&0&0&0\\0&0&0&0\\2&0&0&2\end{pmatrix}=\rho_{AB},
\]
donc $\operatorname{Tr}\rho_{AB}^2=1$ : le couple $(A,B)$ est dans un état \emph{pur}.

% ---------------------------------------------------------------------
\subsection{Information : variables aléatoires et états quantiques}\label{sec:info}
% ---------------------------------------------------------------------

\begin{definition}{Entropie, information mutuelle, entropie conditionnelle}
Soit $A$ une variable aléatoire de loi $p(a)$. Son \emph{entropie de Shannon}, en bits, est
\[
  H(A)=-\sum_a p(a)\log_2p(a)\qquad(\text{avec }0\log_20=0).
\]
Elle mesure l'incertitude moyenne sur $A$ : une pièce équilibrée a $H=1$ bit. Pour un couple
$(A,B)$ de loi $p(a,b)$, $H(A,B)=-\sum_{a,b}p(a,b)\log_2p(a,b)$, et
\begin{align*}
  I(A;B)&=H(A)+H(B)-H(A,B) &&\text{(information mutuelle)},\\
  H(A|B)&=H(A,B)-H(B)=H(A)-I(A;B) &&\text{(entropie conditionnelle)}.
\end{align*}
\end{definition}

Les deux écritures de $H(A|B)$ sont égales : la définition de $I(A;B)$ donne
$H(A,B)=H(A)+H(B)-I(A;B)$, donc $H(A,B)-H(B)=H(A)-I(A;B)$. L'information mutuelle $I(A;B)$
mesure ce que la connaissance de $B$ apprend sur $A$ ; $H(A|B)$ est l'incertitude qui reste sur
$A$ quand on connaît $B$.

\subsubsection{Deux variables aléatoires indépendantes}

Alice et Bob lancent chacun une pièce équilibrée, sans lien entre les deux lancers :
$p(a,b)=\frac14$ pour les quatre couples. Chaque pièce donne $0$ ou $1$ avec la probabilité $\frac12$, donc
\begin{align*}
  H(A)&=-\Big(\tfrac12\log_2\tfrac12+\tfrac12\log_2\tfrac12\Big)=1\text{ bit},\\
  H(B)&=1\text{ bit (même calcul)},\\
  H(A,B)&=-4\cdot\tfrac14\log_2\tfrac14=2\text{ bits},\\
  I(A;B)&=H(A)+H(B)-H(A,B)=1+1-2=0,\\
  H(A|B)&=H(A)-I(A;B)=1-0=H(A)=1\text{ bit}.
\end{align*}
Connaître $B$ n'apprend rien sur $A$.

\subsubsection{Deux variables solidaires}

Les deux pièces tombent toujours du même côté : $p(0,0)=p(1,1)=\frac12$ et
$p(0,1)=p(1,0)=0$. Les lois de chaque pièce sont inchangées ($p(a{=}0)=p(0,0)+p(0,1)=\frac12$,
etc.), donc $H(A)=H(B)=1$ bit, mais
\begin{align*}
  H(A,B)&=-2\cdot\tfrac12\log_2\tfrac12-2\cdot0\log_20=1\text{ bit},\\
  I(A;B)&=1+1-1=1\text{ bit},\qquad H(A|B)=1-1=0.
\end{align*}
Connaître $B$ détermine $A$ : plus aucune incertitude ne reste.

\subsubsection{Deux qubits dans un état de Bell}

Pour un système quantique, on remplace la loi de probabilité par la matrice de densité.

\begin{definition}{Trace partielle et entropie de von Neumann}
L'état du seul qubit $A$ d'un système $(A,B)$ de matrice de densité $\rho_{AB}$ est obtenu en
\emph{traçant sur $B$} : $\rho_A=\operatorname{Tr}_B(\rho_{AB})$. La trace partielle est
linéaire et remplace la partie $B$ de chaque terme par sa trace :
\[
  \operatorname{Tr}_B\big(\ket{a}\bra{a'}\otimes\ket{b}\bra{b'}\big)=\ket{a}\bra{a'}\;\braket{b'}{b}.
\]
L'\emph{entropie de von Neumann} d'un état $\rho$, de valeurs propres $\lambda_i$, est
$S(\rho)=-\operatorname{Tr}(\rho\log_2\rho)=-\sum_i\lambda_i\log_2\lambda_i$. En quantique,
$H(A)=S(\rho_A)$, $H(B)=S(\rho_B)$ et $H(A,B)=S(\rho_{AB})$ ; les définitions de
$I(A;B)$ et de $H(A|B)$ sont les mêmes.
\end{definition}

On part de l'état de Bell
$\rho_{AB}=\frac12\big(\ket{00}\bra{00}+\ket{00}\bra{11}+\ket{11}\bra{00}+\ket{11}\bra{11}\big)$
produit par le circuit de la section~\ref{sec:circuit-bell}. Chaque terme est un produit
d'opérateurs à un qubit :
$\ket{00}\bra{11}=\ket{0}\bra{1}\otimes\ket{0}\bra{1}$, etc. Alors
\begin{align*}
  \rho_A=\operatorname{Tr}_B(\rho_{AB})
  &=\tfrac12\Big(\ket{0}\bra{0}\,\braket{0}{0}+\ket{0}\bra{1}\,\braket{1}{0}
      +\ket{1}\bra{0}\,\braket{0}{1}+\ket{1}\bra{1}\,\braket{1}{1}\Big)\\
  &=\tfrac12\Big(\ket{0}\bra{0}\cdot1+\ket{0}\bra{1}\cdot0+\ket{1}\bra{0}\cdot0+\ket{1}\bra{1}\cdot1\Big)\\
  &=\tfrac12\big(\ket{0}\bra{0}+\ket{1}\bra{1}\big)=\tfrac12\,\mathrm{Id}
  =\begin{pmatrix}\frac12&0\\0&\frac12\end{pmatrix}.
\end{align*}
La formule des composantes de la section~\ref{sec:separable},
$(\rho_A)_{ij}=\rho_{(i0),(j0)}+\rho_{(i1),(j1)}$, donne le même résultat. En traçant sur $A$
au lieu de $B$, la règle est $\operatorname{Tr}_A(\ket{a}\bra{a'}\otimes\ket{b}\bra{b'})
=\braket{a'}{a}\,\ket{b}\bra{b'}$, et on trouve de même $\rho_B=\frac12\mathrm{Id}$.

Calcul des entropies :
\begin{itemize}[nosep]
  \item $\rho_{AB}$ est pur : ses valeurs propres sont $1,0,0,0$, donc
        $H(A,B)=S(\rho_{AB})=-1\cdot\log_21=0$ ;
  \item $\rho_A$ et $\rho_B$ ont pour valeurs propres $\frac12,\frac12$, donc
        $H(A)=H(B)=-\big(\frac12\log_2\frac12+\frac12\log_2\frac12\big)=1$ bit.
\end{itemize}
D'où
\begin{align*}
  I(A;B)&=H(A)+H(B)-H(A,B)=1+1-0=2\text{ bits},\\
  H(A|B)&=H(A,B)-H(B)=0-1=-1\text{ bit}\qquad(=H(A)-I(A;B)=1-2).
\end{align*}

\begin{remarque}[Une entropie conditionnelle négative]
Pour des variables classiques, $p(a,b)\leq p(b)$, donc $-\log_2p(a,b)\geq-\log_2p(b)$ et
$H(A,B)\geq H(B)$ : on a toujours $H(A|B)\geq0$, et $I(A;B)\leq H(A)$. Ici $H(A|B)=-1$ et
$I(A;B)=2>H(A)$ : ces valeurs n'ont pas d'équivalent classique, c'est une signature de
l'intrication. Le résultat est surprenant mais mathématiquement correct.
\end{remarque}

\begin{center}
\renewcommand{\arraystretch}{1.4}
\begin{tabular}{@{}lccc@{}}
\toprule
 & indépendantes & solidaires & état de Bell \\
\midrule
$H(A)$ & $1$ & $1$ & $1$ \\
$H(B)$ & $1$ & $1$ & $1$ \\
$H(A,B)$ & $2$ & $1$ & $0$ \\
$I(A;B)$ & $0$ & $1$ & $2$ \\
$H(A|B)$ & $1$ & $0$ & $-1$ \\
\bottomrule
\end{tabular}
\end{center}
(en bits ; les deux premières colonnes sont classiques, la troisième est quantique).

% ---------------------------------------------------------------------
\subsection{Système pur et système mélangé}\label{sec:pur-melange}
% ---------------------------------------------------------------------

Rappel (section~\ref{sec:pur-mixte}) : si $\lambda_i$ sont les valeurs propres de la matrice de
densité $\rho$, alors $\operatorname{Tr}\rho^2=\sum_i\lambda_i^2\leq1$. Le système est
\emph{pur} si $\operatorname{Tr}\rho^2=1$ (alors $\rho=\ket{\psi}\bra{\psi}$ et $\rho^2=\rho$),
et \emph{mélangé} si $\operatorname{Tr}\rho^2<1$. Pour un qubit,
$\operatorname{Tr}\rho^2=\frac{1+|\vec r|^2}{2}$ (section~\ref{sec:bloch}) est compris entre
$\frac12$ et $1$.

\subsubsection{Un système pur : exercice}

\emph{Énoncé.} Soit $\ket{\psi}=\frac{1}{\sqrt2}\ket{0}+\frac{1}{\sqrt2}\ket{1}$.
\begin{enumerate}[label=\alph*), nosep]
  \item Calculer la matrice de densité $\rho$.
  \item Calculer la pureté du système quantique.
  \item Calculer $\langle\sigma_x\rangle$ et $\langle\sigma_z\rangle$.
  \item Écrire $\ket{\psi}$ dans la base $X$.
\end{enumerate}

\paragraph{a) Matrice de densité.} $\rho=\ket{\psi}\bra{\psi}$ est le produit d'un vecteur
colonne par un vecteur ligne (les amplitudes sont réelles) :
\[
  \rho=\begin{pmatrix}\frac{1}{\sqrt2}\\[0.4ex]\frac{1}{\sqrt2}\end{pmatrix}
       \begin{pmatrix}\frac{1}{\sqrt2}&\frac{1}{\sqrt2}\end{pmatrix}
  =\begin{pmatrix}\frac12&\frac12\\[0.4ex]\frac12&\frac12\end{pmatrix}.
\]

\paragraph{b) Pureté.}
\[
  \rho^2=\begin{pmatrix}\frac12&\frac12\\[0.4ex]\frac12&\frac12\end{pmatrix}
         \begin{pmatrix}\frac12&\frac12\\[0.4ex]\frac12&\frac12\end{pmatrix}
  =\begin{pmatrix}\frac14+\frac14&\frac14+\frac14\\[0.4ex]\frac14+\frac14&\frac14+\frac14\end{pmatrix}
  =\begin{pmatrix}\frac12&\frac12\\[0.4ex]\frac12&\frac12\end{pmatrix}=\rho,
\]
donc $\operatorname{Tr}\rho^2=\operatorname{Tr}\rho=\frac12+\frac12=1$ : le système est \emph{pur}.

\paragraph{c) Valeurs moyennes.} Avec $\sigma_x=X=\begin{pmatrix}0&1\\1&0\end{pmatrix}$ et
$\sigma_z=Z=\begin{pmatrix}1&0\\0&-1\end{pmatrix}$ (section~\ref{sec:op-dirac}), la valeur
moyenne est $\langle\sigma\rangle=\operatorname{Tr}(\rho\,\sigma)$ (section~\ref{sec:rho}).
\begin{align*}
  \rho\,\sigma_x&=\begin{pmatrix}\frac12&\frac12\\[0.4ex]\frac12&\frac12\end{pmatrix}
     \begin{pmatrix}0&1\\1&0\end{pmatrix}
     =\begin{pmatrix}\frac12&\frac12\\[0.4ex]\frac12&\frac12\end{pmatrix},
     &\langle\sigma_x\rangle&=\operatorname{Tr}(\rho\,\sigma_x)=\tfrac12+\tfrac12=1,\\[1.5ex]
  \rho\,\sigma_z&=\begin{pmatrix}\frac12&\frac12\\[0.4ex]\frac12&\frac12\end{pmatrix}
     \begin{pmatrix}1&0\\0&-1\end{pmatrix}
     =\begin{pmatrix}\frac12&-\frac12\\[0.4ex]\frac12&-\frac12\end{pmatrix},
     &\langle\sigma_z\rangle&=\operatorname{Tr}(\rho\,\sigma_z)=\tfrac12-\tfrac12=0.
\end{align*}
On retrouve $\langle\sigma_z\rangle$ avec les probabilités de mesure : dans la base $Z$,
$\sigma_z$ vaut $+1$ pour $\ket{0}$ et $-1$ pour $\ket{1}$, et
$|\braket{0}{\psi}|^2=|\braket{1}{\psi}|^2=\frac12$, donc
$\langle\sigma_z\rangle=\frac12\times(+1)+\frac12\times(-1)=0$.

\paragraph{d) Écriture dans la base $X$.} La base $X$ est $\{\ket{+},\ket{-}\}$, orthonormée, donc
$\ket{+}\bra{+}+\ket{-}\bra{-}=\mathrm{Id}$ et
$\ket{\psi}=\ket{+}\braket{+}{\psi}+\ket{-}\braket{-}{\psi}$, avec
\[
  \braket{+}{\psi}=\tfrac{1}{\sqrt2}\Big(\tfrac{1}{\sqrt2}+\tfrac{1}{\sqrt2}\Big)=1,
  \qquad
  \braket{-}{\psi}=\tfrac{1}{\sqrt2}\Big(\tfrac{1}{\sqrt2}-\tfrac{1}{\sqrt2}\Big)=0 .
\]
Donc $\ket{\psi}=1\cdot\ket{+}+0\cdot\ket{-}=\ket{+}$ : dans la base $X$, l'état n'a qu'une
seule composante. Comme $\sigma_x\ket{+}=+\ket{+}$ et $\sigma_x\ket{-}=-\ket{-}$, une
mesure de $\sigma_x$ donne $+1$ avec certitude, et
$\langle\sigma_x\rangle=1\cdot(+1)+0\cdot(-1)=1$ : c'est le résultat de la question c).
Le vecteur de Bloch de $\ket{+}$ est $(\langle\sigma_x\rangle,\langle\sigma_y\rangle,\langle\sigma_z\rangle)=(1,0,0)$,
de norme $1$ : l'état est sur la sphère, en accord avec $\operatorname{Tr}\rho^2=\frac{1+1}{2}=1$.

\subsubsection{Un système mélangé : un qubit de la paire de Bell}

À la section~\ref{sec:info}, la trace sur $B$ de l'état de Bell a donné
$\rho_A=\frac12\mathrm{Id}=\begin{pmatrix}\frac12&0\\0&\frac12\end{pmatrix}$. Alors
\[
  \rho_A^2=\begin{pmatrix}\frac14&0\\0&\frac14\end{pmatrix},
  \qquad
  \operatorname{Tr}\big(\rho_A^2\big)=\frac14+\frac14=\frac12<1 :
\]
le qubit $A$ est dans un état \emph{mélangé}, et c'est le plus mélangé possible pour un qubit
(centre de la boule de Bloch, $\vec r=\vec0$). Pourtant le couple $(A,B)$ est pur
($\operatorname{Tr}\rho_{AB}^2=1$) : on connaît parfaitement l'état du couple, mais pas celui d'un
qubit pris seul. L'information est entièrement dans les corrélations entre $A$ et $B$, et
$S(\rho_A)=1$ bit.

\subsubsection{Un système mélangé : le mélange statistique de Bob}

Bob prépare un qubit dans l'état $\ket{0}$ avec la probabilité $p_1=\frac34$ (75\,\%) et dans
l'état $\ket{1}$ avec la probabilité $p_2=\frac14$ (25\,\%), sans que l'on sache lequel des deux a
été préparé. Ce n'est pas un état pur $\ket{\psi}$ : on décrit le système par la matrice de
densité $\rho=\sum_ip_i\ket{\psi_i}\bra{\psi_i}$ (section~\ref{sec:pur-mixte}) :
\[
  \rho=\underbrace{\tfrac34}_{p_1}\ket{0}\bra{0}+\underbrace{\tfrac14}_{p_2}\ket{1}\bra{1}
  =\begin{pmatrix}\frac34&0\\0&\frac14\end{pmatrix}.
\]
Alors
\[
  \rho^2=\begin{pmatrix}\frac{9}{16}&0\\0&\frac{1}{16}\end{pmatrix},
  \qquad
  \operatorname{Tr}\big(\rho^2\big)=\frac{9}{16}+\frac{1}{16}=\frac{10}{16}=\frac58<1 :
\]
le système est \emph{mélangé}.

Les valeurs moyennes sont les moyennes pondérées sur la préparation :
$\langle\sigma_z\rangle=\operatorname{Tr}(\rho\,\sigma_z)=\frac34-\frac14=\frac12
=p_1\cdot(+1)+p_2\cdot(-1)$, et $\langle\sigma_x\rangle=\operatorname{Tr}(\rho\,\sigma_x)=0$
(les termes hors diagonale de $\rho$ sont nuls). Le vecteur de Bloch est $(0,0,\frac12)$,
à l'intérieur de la boule, et $\operatorname{Tr}\rho^2=\frac{1+1/4}{2}=\frac58$. Son entropie
est $S(\rho)=-\frac34\log_2\frac34-\frac14\log_2\frac14=2-\frac34\log_23\approx0{,}811$ bit.

\begin{remarque}[Mélange et superposition]
Considérons l'état pur $\ket{\psi'}=\frac{\sqrt3}{2}\ket{0}+\frac12\ket{1}$, qui a les mêmes
probabilités en base $Z$ ($P(0)=\frac34$, $P(1)=\frac14$). Sa matrice de densité est
\[
  \rho'=\begin{pmatrix}\frac34&\frac{\sqrt3}{4}\\[0.4ex]\frac{\sqrt3}{4}&\frac14\end{pmatrix},
  \qquad
  \operatorname{Tr}\rho'^2=\frac{9}{16}+\frac1{16}+2\cdot\frac{3}{16}=1,
  \qquad
  \langle\sigma_x\rangle=2\cdot\frac{\sqrt3}{4}=\frac{\sqrt3}{2}.
\]
Le mélange de Bob a la même diagonale mais aucun terme hors diagonale ($\operatorname{Tr}\rho^2=\frac58$,
$\langle\sigma_x\rangle=0$) : une superposition n'est pas un tirage au hasard entre $\ket{0}$ et
$\ket{1}$. De plus, pour un état pur il existe une mesure au résultat certain (ici $\sigma_x$ pour
$\ket{+}$) ; pour le mélange de Bob, la meilleure probabilité d'obtenir un même résultat est
$75\,\%$ (mesure en base $Z$, résultat $\ket{0}$).
\end{remarque}

\subsubsection{Bilan}

\begin{center}
\renewcommand{\arraystretch}{1.5}
\begin{tabular}{@{}lcccl@{}}
\toprule
\textbf{Système} & $\rho$ & $\operatorname{Tr}\rho^2$ & $S(\rho)$ & \textbf{Nature} \\
\midrule
$\ket{\psi}=\ket{+}$ & ${\renewcommand{\arraystretch}{1}\begin{pmatrix}\frac12&\frac12\\\frac12&\frac12\end{pmatrix}}$ & $1$ & $0$ & pur \\[2ex]
Qubit $A$ d'une paire de Bell & ${\renewcommand{\arraystretch}{1}\begin{pmatrix}\frac12&0\\0&\frac12\end{pmatrix}}$ & $\frac12$ & $1$ bit & mélangé \\[2ex]
Source de Bob (75\,\%/25\,\%) & ${\renewcommand{\arraystretch}{1}\begin{pmatrix}\frac34&0\\0&\frac14\end{pmatrix}}$ & $\frac58$ & $\approx0{,}811$ bit & mélangé \\
\bottomrule
\end{tabular}
\end{center}

% ---------------------------------------------------------------------
\subsection{Algorithme quantique : la téléportation}\label{sec:teleportation}
% ---------------------------------------------------------------------

\emph{Problème.} Alice détient un qubit $Q$ dans un état qu'elle ne connaît pas,
$\ket{\varphi}=\alpha\ket{0}+\beta\ket{1}$ avec $|\alpha|^2+|\beta|^2=1$, et elle veut que Bob
obtienne cet état sur l'un de ses qubits. Elle ne peut pas mesurer $Q$ pour connaître
$\alpha$ et $\beta$ (une mesure ne donne qu'un bit et détruit l'état), elle ne peut pas le
copier (théorème de non-clonage) et elle n'envoie pas le qubit lui-même.

\begin{definition}{Protocole de téléportation quantique}
\textbf{Ressources.}
\begin{itemize}[nosep]
  \item Un \emph{e-bit} : Alice et Bob partagent une paire de qubits $(A,B)$ dans l'état de Bell
        $\ket{\Phi^+}_{BA}=\frac{1}{\sqrt2}\big(\ket{00}+\ket{11}\big)$ ; Alice détient $A$,
        Bob détient $B$.
  \item Le qubit $Q$ à téléporter, chez Alice.
  \item Un canal classique (deux bits).
\end{itemize}
Le système est écrit $(B,A,Q)$ : $\ket{b,a,q}=\ket{b}_B\otimes\ket{a}_A\otimes\ket{q}_Q$.

\textbf{Étapes.}
\begin{enumerate}[nosep]
  \item Alice applique un CNOT de contrôle $Q$ et de cible $A$.
  \item Alice applique une porte $H$ sur $Q$.
  \item Alice mesure $A$ et $Q$ dans la base de calcul, obtient deux bits $(a,q)$ et les envoie
        à Bob.
  \item Bob applique $X$ si $a=1$, puis $Z$ si $q=1$, à son qubit $B$.
\end{enumerate}
\end{definition}

\begin{figure}[!ht]
\centering
\begin{adjustbox}{max width=\linewidth}
\begin{tikzpicture}[>=Stealth, font=\small]
  % fils quantiques : Q en haut, A au milieu, B en bas
  \draw (0,2.4) -- (5.3,2.4);
  \draw (0,1.2) -- (5.3,1.2);
  \draw (0,0) -- (10.6,0);
  \node[anchor=east] at (0,2.4) {$\ket{\varphi}_Q$};
  \node[anchor=south west, text=insagris, font=\footnotesize] at (0.05,1.2) {$A$};
  \node[anchor=south west, text=insagris, font=\footnotesize] at (0.05,0) {$B$};
  % paire de Bell partagée
  \draw[insagris, thick] (-0.1,1.5) -- (-0.25,1.5) -- (-0.25,-0.3) -- (-0.1,-0.3);
  \node[anchor=east] at (-0.35,0.6) {$\ket{\Phi^+}_{BA}$};
  % CNOT : contrôle Q, cible A
  \draw[insarouge, thick] (2.6,2.4) -- (2.6,1.2);
  \pic at (2.6,2.4) {ctl};
  \pic at (2.6,1.2) {cible};
  % Hadamard sur Q
  \node[porte] at (4.2,2.4) {$H$};
  % mesures d'Alice
  \pic at (5.6,2.4) {mesure};
  \pic at (5.6,1.2) {mesure};
  % fils classiques (doubles)
  \draw[double, double distance=1.2pt, insagris] (5.9,1.2) -- (7.6,1.2) -- (7.6,0.35);
  \draw[double, double distance=1.2pt, insagris] (5.9,2.4) -- (9.6,2.4) -- (9.6,0.35);
  \node[above, font=\footnotesize] at (6.5,1.2) {$a$};
  \node[above, font=\footnotesize] at (6.5,2.4) {$q$};
  % corrections de Bob
  \node[porte] at (7.6,0) {$X$};
  \node[porte] at (9.6,0) {$Z$};
  \node[anchor=west] at (10.6,0) {$\ket{\varphi}_B$};
  % états intermédiaires
  \foreach \x in {1.4,3.4,4.9}
    \draw[dashed, insagris!70] (\x,2.8) -- (\x,-0.5);
  \node[anchor=north, font=\footnotesize] at (1.4,-0.5) {$\ket{\Pi_0}$};
  \node[anchor=north, font=\footnotesize] at (3.4,-0.5) {$\ket{\Pi_1}$};
  \node[anchor=north, font=\footnotesize] at (4.9,-0.5) {$\ket{\Pi_2}$};
\end{tikzpicture}
\end{adjustbox}
\caption{Circuit de téléportation. Alice détient $Q$ (dans $\ket{\varphi}=\alpha\ket{0}+\beta\ket{1}$)
et $A$, Bob détient $B$ ; $A$ et $B$ forment la paire $\ket{\Phi^+}_{BA}$. Les traits doubles
sont des fils \emph{classiques} qui portent les résultats $a$ (mesure de $A$) et $q$ (mesure de
$Q$) ; $X$ est appliqué si $a=1$, $Z$ si $q=1$. Le ket s'écrit $\ket{b,a,q}$ : le fil du haut
($Q$) est le dernier chiffre.}
\label{fig:teleportation}
\end{figure}

\subsubsection{Calcul pas à pas}

\paragraph{État initial $\ket{\Pi_0}$.} Alice et Bob partagent la paire, et $Q$ est indépendant
d'elle :
\begin{align*}
  \ket{\Pi_0}&=\ket{\Phi^+}_{BA}\otimes\ket{\varphi}_Q
    =\frac{1}{\sqrt2}\big(\ket{00}+\ket{11}\big)\otimes\big(\alpha\ket{0}+\beta\ket{1}\big)\\
  &=\frac{1}{\sqrt2}\Big(\alpha\ket{00}\ket{0}+\beta\ket{00}\ket{1}+\alpha\ket{11}\ket{0}+\beta\ket{11}\ket{1}\Big)\\
  &=\frac{\alpha\ket{000}+\beta\ket{001}+\alpha\ket{110}+\beta\ket{111}}{\sqrt2}.
\end{align*}

\paragraph{Après le CNOT $Q\to A$ : $\ket{\Pi_1}$.} Le CNOT (contrôle $Q$, dernier chiffre ;
cible $A$, chiffre du milieu) agit ainsi : $\ket{b,a,q}\to\ket{b,\,a\oplus q,\,q}$. Terme par terme :
\begin{center}
\renewcommand{\arraystretch}{1.3}
\begin{tabular}{@{}cccc@{}}
\toprule
Terme & $q$ & Effet sur $A$ & Résultat \\
\midrule
$\alpha\ket{000}$ & $0$ & aucun & $\alpha\ket{000}$ \\
$\beta\ket{001}$ & $1$ & $0\to1$ & $\beta\ket{011}$ \\
$\alpha\ket{110}$ & $0$ & aucun & $\alpha\ket{110}$ \\
$\beta\ket{111}$ & $1$ & $1\to0$ & $\beta\ket{101}$ \\
\bottomrule
\end{tabular}
\end{center}
donc
\[
  \ket{\Pi_1}=\frac{\alpha\ket{000}+\beta\ket{011}+\alpha\ket{110}+\beta\ket{101}}{\sqrt2}.
\]
On isole ensuite le dernier chiffre, celui de $Q$, qui va changer à l'étape suivante
(le premier ket est celui du couple $\ket{b,a}$) :
\[
  \ket{\Pi_1}=\frac{\alpha\ket{00}\otimes\ket{0}+\beta\ket{01}\otimes\ket{1}
    +\alpha\ket{11}\otimes\ket{0}+\beta\ket{10}\otimes\ket{1}}{\sqrt2}.
\]

\paragraph{Après la porte $H$ sur $Q$ : $\ket{\Pi_2}$.} On remplace $\ket{0}$ par
$H\ket{0}=\ket{+}$ et $\ket{1}$ par $H\ket{1}=\ket{-}$ dans le facteur $Q$ :
\[
  \ket{\Pi_2}=\frac{1}{\sqrt2}\Big(\alpha\ket{00}\otimes\ket{+}+\beta\ket{01}\otimes\ket{-}
    +\alpha\ket{11}\otimes\ket{+}+\beta\ket{10}\otimes\ket{-}\Big).
\]
Avec $\ket{\pm}=\frac{1}{\sqrt2}(\ket{0}\pm\ket{1})$, le préfacteur devient
$\frac{1}{\sqrt2}\cdot\frac{1}{\sqrt2}=\frac12$ :
\begin{align*}
  \ket{\Pi_2}&=\frac12\Big(\alpha\ket{00}\big(\ket{0}+\ket{1}\big)+\beta\ket{01}\big(\ket{0}-\ket{1}\big)
     +\alpha\ket{11}\big(\ket{0}+\ket{1}\big)+\beta\ket{10}\big(\ket{0}-\ket{1}\big)\Big)\\
  &=\frac12\Big(\alpha\ket{000}+\alpha\ket{001}+\beta\ket{010}-\beta\ket{011}
     +\alpha\ket{110}+\alpha\ket{111}+\beta\ket{100}-\beta\ket{101}\Big).
\end{align*}

\paragraph{Regroupement selon le résultat d'Alice.} Alice mesurera $A$ et $Q$, donc les deux
derniers chiffres. On regroupe les huit termes qui ont les mêmes chiffres $(a,q)$, en mettant
le ket de Bob en premier :
\begin{center}
\renewcommand{\arraystretch}{1.3}
\begin{tabular}{@{}cll@{}}
\toprule
$(a,q)$ & Termes de $\ket{\Pi_2}$ (facteur $\frac12$ omis) & Ket de Bob \\
\midrule
$(0,0)$ & $\alpha\ket{000}+\beta\ket{100}$ & $(\alpha\ket{0}+\beta\ket{1})\ket{00}$ \\
$(0,1)$ & $\alpha\ket{001}-\beta\ket{101}$ & $(\alpha\ket{0}-\beta\ket{1})\ket{01}$ \\
$(1,0)$ & $\beta\ket{010}+\alpha\ket{110}$ & $(\beta\ket{0}+\alpha\ket{1})\ket{10}$ \\
$(1,1)$ & $-\beta\ket{011}+\alpha\ket{111}$ & $(-\beta\ket{0}+\alpha\ket{1})\ket{11}$ \\
\bottomrule
\end{tabular}
\end{center}
d'où
\[
  \ket{\Pi_2}=\frac12\Big[\big(\alpha\ket{0}+\beta\ket{1}\big)\ket{00}
    +\big(\alpha\ket{0}-\beta\ket{1}\big)\ket{01}
    +\big(\beta\ket{0}+\alpha\ket{1}\big)\ket{10}
    +\big(\alpha\ket{1}-\beta\ket{0}\big)\ket{11}\Big].
\]
Dans chaque terme, le ket de gauche est l'état de $B$ et le ket de droite $\ket{a,q}$ celui du
couple $(A,Q)$ d'Alice.

\subsubsection{Mesure d'Alice et corrections de Bob}

Alice mesure $A$ et $Q$. Comme à la section~\ref{sec:mesure2}, la probabilité d'un résultat
$(a,q)$ est le carré de la norme du terme correspondant, et l'état de Bob devient le ket
associé (déjà normé). Ici
\[
  P(a,q)=\Big\|\tfrac12\ket{\phi_{aq}}\Big\|^2=\tfrac14\big(|\alpha|^2+|\beta|^2\big)=\tfrac14
  \quad\text{pour les quatre résultats,}
\]
qui sont donc équiprobables et \emph{indépendants} de $\alpha$ et $\beta$ : les deux bits
envoyés ne révèlent rien sur $\ket{\varphi}$. Les kets de Bob sont
$\ket{\phi_{00}}=\alpha\ket{0}+\beta\ket{1}$, $\ket{\phi_{01}}=\alpha\ket{0}-\beta\ket{1}$,
$\ket{\phi_{10}}=\beta\ket{0}+\alpha\ket{1}$ et $\ket{\phi_{11}}=-\beta\ket{0}+\alpha\ket{1}$.

\begin{center}
\renewcommand{\arraystretch}{1.5}
\begin{tabular}{@{}cccll@{}}
\toprule
Résultat $(a,q)$ & Probabilité & État de Bob & Lien avec $\ket{\varphi}$ & Correction de Bob \\
\midrule
$(0,0)$ & $\frac14$ & $\alpha\ket{0}+\beta\ket{1}$ & $\ket{\varphi}$ & $I$ (rien) \\
$(0,1)$ & $\frac14$ & $\alpha\ket{0}-\beta\ket{1}$ & $Z\ket{\varphi}$ & $Z$ \\
$(1,0)$ & $\frac14$ & $\beta\ket{0}+\alpha\ket{1}$ & $X\ket{\varphi}$ & $X$ \\
$(1,1)$ & $\frac14$ & $-\beta\ket{0}+\alpha\ket{1}$ & $XZ\ket{\varphi}$ & $X$ puis $Z$ (opérateur $ZX$) \\
\bottomrule
\end{tabular}
\end{center}

\paragraph{Vérification des quatre corrections.} Avec $X\ket{0}=\ket{1}$, $X\ket{1}=\ket{0}$,
$Z\ket{0}=\ket{0}$ et $Z\ket{1}=-\ket{1}$ :
\begin{align*}
  (0,0):\quad & I\big(\alpha\ket{0}+\beta\ket{1}\big)=\alpha\ket{0}+\beta\ket{1},\\
  (0,1):\quad & Z\big(\alpha\ket{0}-\beta\ket{1}\big)=\alpha\ket{0}-\beta Z\ket{1}=\alpha\ket{0}+\beta\ket{1},\\
  (1,0):\quad & X\big(\beta\ket{0}+\alpha\ket{1}\big)=\beta\ket{1}+\alpha\ket{0}=\alpha\ket{0}+\beta\ket{1},\\
  (1,1):\quad & X\big(-\beta\ket{0}+\alpha\ket{1}\big)=-\beta\ket{1}+\alpha\ket{0}=\alpha\ket{0}-\beta\ket{1},
     \ \text{puis}\ Z\big(\alpha\ket{0}-\beta\ket{1}\big)=\alpha\ket{0}+\beta\ket{1}.
\end{align*}
Dans le dernier cas, avec les matrices,
$ZX=\begin{pmatrix}1&0\\0&-1\end{pmatrix}\begin{pmatrix}0&1\\1&0\end{pmatrix}
=\begin{pmatrix}0&1\\-1&0\end{pmatrix}$ et
$ZX\begin{pmatrix}-\beta\\ \alpha\end{pmatrix}=\begin{pmatrix}\alpha\\ \beta\end{pmatrix}$.
L'ordre compte : l'état de Bob est $XZ\ket{\varphi}$ ($Z\ket{\varphi}=\alpha\ket{0}-\beta\ket{1}$,
puis $X$ donne $\alpha\ket{1}-\beta\ket{0}$), et on l'annule avec
$(XZ)^{-1}=Z^{-1}X^{-1}=ZX$ : on applique d'abord $X$, puis $Z$, comme dans le circuit.
Dans les quatre cas, Bob retrouve exactement $\ket{\varphi}=\alpha\ket{0}+\beta\ket{1}$.

\subsubsection{Ce que Bob sait avant de recevoir les bits}

Tant que Bob ignore $(a,q)$, son qubit est dans le mélange des quatre kets $\ket{\phi_{aq}}$, chacun
avec la probabilité $\frac14$ :
\begin{align*}
  \rho_B&=\frac14\sum_{a,q}\ket{\phi_{aq}}\bra{\phi_{aq}}\\
  &=\frac14\Bigg[\begin{pmatrix}|\alpha|^2&\alpha\beta^*\\ \alpha^*\beta&|\beta|^2\end{pmatrix}
    +\begin{pmatrix}|\alpha|^2&-\alpha\beta^*\\ -\alpha^*\beta&|\beta|^2\end{pmatrix}
    +\begin{pmatrix}|\beta|^2&\alpha^*\beta\\ \alpha\beta^*&|\alpha|^2\end{pmatrix}
    +\begin{pmatrix}|\beta|^2&-\alpha^*\beta\\ -\alpha\beta^*&|\alpha|^2\end{pmatrix}\Bigg]\\
  &=\frac14\begin{pmatrix}2(|\alpha|^2+|\beta|^2)&0\\0&2(|\alpha|^2+|\beta|^2)\end{pmatrix}
   =\frac12\begin{pmatrix}1&0\\0&1\end{pmatrix}=\frac12\,\mathrm{Id}.
\end{align*}
Sans les deux bits, Bob a donc un qubit complètement mélangé, qui ne dépend pas de $\alpha$ ni de
$\beta$ : il n'a reçu aucune information sur $\ket{\varphi}$. C'est l'envoi classique de $(a,q)$
qui permet de la retrouver, donc rien ne voyage plus vite que la lumière.

\begin{remarque}[Bilan de la téléportation]
\begin{itemize}[nosep]
  \item Ressources : $1$ e-bit (la paire de Bell) et $2$ bits classiques pour transmettre
        \emph{un qubit} d'état quelconque, sans qu'Alice connaisse $\alpha$ ni $\beta$.
  \item Pas de clonage : après la mesure, $A$ et $Q$ sont dans l'état de base $\ket{a,q}$.
        Alice n'a plus $\ket{\varphi}$ : l'état a été \emph{déplacé}, pas copié. L'intrication
        entre $A$ et $B$ a été consommée.
  \item Les quatre corrections $I,Z,X,ZX$ sont des portes contrôlées par des bits classiques
        (section~\ref{sec:controlees}).
\end{itemize}
\end{remarque}

% ---------------------------------------------------------------------
\subsection{Portes contrôlées : contrôlé-\texorpdfstring{$U$}{U}, Fredkin et Toffoli}\label{sec:controlees}
% ---------------------------------------------------------------------

Le CNOT est le cas le plus simple d'une idée générale : appliquer une opération à un ou
plusieurs qubits \emph{cibles} seulement si un ou plusieurs qubits de \emph{contrôle} valent $1$.
Dans cette section, comme à la section~\ref{sec:op2}, le contrôle est le \emph{premier} chiffre
du ket. Les qubits sont notés $x$, $y$, $z$ (le qubit $x$ est le contrôle) et
$\ket{x,y,z}=\ket{x}\otimes\ket{y}\otimes\ket{z}$ ; d'après la convention de la section, le
contrôle est dessiné en bas (figure~\ref{fig:controlees}). Pour éviter la confusion avec les
matrices de Pauli, on écrit ici $\sigma_x=X$ pour la porte NOT.

\begin{figure}[!ht]
\centering
\begin{adjustbox}{max width=\linewidth}
\begin{tikzpicture}[>=Stealth, font=\small]
  % (a) contrôlé-U : contrôle x en bas, cible y en haut
  \draw (0,1.2) -- (3.4,1.2);
  \draw (0,0) -- (3.4,0);
  \node[anchor=east] at (0,1.2) {$y$};
  \node[anchor=east] at (0,0) {$x$};
  \draw[insarouge, thick] (1.7,0) -- (1.7,0.85);
  \pic at (1.7,0) {ctl};
  \node[porte] at (1.7,1.2) {$U$};
  \node[anchor=north, font=\footnotesize] at (1.7,-0.5) {(a) contrôlé-$U$};
  % (b) contrôlé-SWAP (Fredkin) : x contrôle, y et z échangés
  \draw (5,2.4) -- (8.4,2.4);
  \draw (5,1.2) -- (8.4,1.2);
  \draw (5,0) -- (8.4,0);
  \node[anchor=east] at (5,2.4) {$z$};
  \node[anchor=east] at (5,1.2) {$y$};
  \node[anchor=east] at (5,0) {$x$};
  \draw[insarouge, thick] (6.7,0) -- (6.7,2.4);
  \pic at (6.7,0) {ctl};
  \pic at (6.7,1.2) {croix};
  \pic at (6.7,2.4) {croix};
  \node[anchor=north, font=\footnotesize] at (6.7,-0.5) {(b) Fredkin (C-SWAP)};
  % (c) Toffoli (CCNOT) : x et y contrôlent, z cible
  \draw (10,2.4) -- (13.4,2.4);
  \draw (10,1.2) -- (13.4,1.2);
  \draw (10,0) -- (13.4,0);
  \node[anchor=east] at (10,2.4) {$z$};
  \node[anchor=east] at (10,1.2) {$y$};
  \node[anchor=east] at (10,0) {$x$};
  \draw[insarouge, thick] (11.7,0) -- (11.7,2.4);
  \pic at (11.7,0) {ctl};
  \pic at (11.7,1.2) {ctl};
  \pic at (11.7,2.4) {cible};
  \node[anchor=north, font=\footnotesize] at (11.7,-0.5) {(c) Toffoli (CCNOT)};
\end{tikzpicture}
\end{adjustbox}
\caption{Portes contrôlées. Le contrôle $x$ est le premier chiffre du ket, donc le fil du bas.
(a) Le qubit $y$ subit $U$ si $x=1$ ; pour $U=X$ on retrouve le CNOT. (b) Si $x=1$, les qubits
$y$ et $z$ sont échangés. (c) Si $x=y=1$, le qubit $z$ est inversé.}
\label{fig:controlees}
\end{figure}

\subsubsection{Contrôlé-\texorpdfstring{$U$}{U}}

\begin{definition}{Porte contrôlée-$U$}
Soit $U$ une matrice unitaire $2\times2$. Sur le système $(x,y)$, de contrôle $x$ et de cible $y$,
\[
  \mathrm{C}U=\ket{0}\bra{0}_x\otimes I_y+\ket{1}\bra{1}_x\otimes U_y .
\]
Le premier terme agit quand le contrôle est $\ket{0}$ : la cible ne change pas ($I$). Le second
agit quand le contrôle est $\ket{1}$ : la cible subit $U$.
\end{definition}

Comme $\ket{0}\bra{0}\ket{x}=\delta_{x0}\ket{0}$ et $\ket{1}\bra{1}\ket{x}=\delta_{x1}\ket{1}$,
\[
  \mathrm{C}U\big(\ket{x}\otimes\ket{y}\big)=\ket{x}\otimes U^x\ket{y},\qquad U^0=I,\ U^1=U .
\]
Par la règle $A\otimes B=(A_{ij}B)$, $\ket{0}\bra{0}\otimes I=\begin{pmatrix}I&0\\0&0\end{pmatrix}$
et $\ket{1}\bra{1}\otimes U=\begin{pmatrix}0&0\\0&U\end{pmatrix}$ (blocs $2\times2$), donc
\[
  \mathrm{C}U=\begin{pmatrix}I&0\\0&U\end{pmatrix}
  =\begin{pmatrix}1&0&0&0\\0&1&0&0\\0&0&u_{00}&u_{01}\\0&0&u_{10}&u_{11}\end{pmatrix}
  \qquad\text{pour }U=\begin{pmatrix}u_{00}&u_{01}\\u_{10}&u_{11}\end{pmatrix}.
\]
Cette matrice est unitaire car $I^\dagger I=I$ et $U^\dagger U=I$.

\paragraph{Le CNOT est un cas particulier.} Pour $U=\sigma_x$,
$\mathrm{C}\sigma_x=\ket{0}\bra{0}_x\otimes I_y+\ket{1}\bra{1}_x\otimes\sigma_x$ est le CNOT de
contrôle $x$ (premier chiffre), c'est-à-dire $\mathrm{CNOT}_{12}$ de la section~\ref{sec:cnot}.
On calcule l'action sur les quatre états de base, avec $\braket{0}{0}=\braket{1}{1}=1$,
$\braket{0}{1}=\braket{1}{0}=0$, $\sigma_x\ket{0}=\ket{1}$ et $\sigma_x\ket{1}=\ket{0}$ :
\begin{align*}
  \mathrm{C}\sigma_x\ket{00}&=\big[\ket{0}\bra{0}_x\otimes I_y\big]\ket{0}_x\ket{0}_y
      +\big[\ket{1}\bra{1}_x\otimes\sigma_x\big]\ket{0}_x\ket{0}_y\\
    &=\ket{0}\braket{0}{0}\otimes I\ket{0}+\ket{1}\braket{1}{0}\otimes\sigma_x\ket{0}
     =\ket{0}\otimes\ket{0}+0=\ket{00},\\[1.2ex]
  \mathrm{C}\sigma_x\ket{01}&=\ket{0}\braket{0}{0}\otimes I\ket{1}+\ket{1}\braket{1}{0}\otimes\sigma_x\ket{1}
     =\ket{0}\otimes\ket{1}+0=\ket{01},\\[1.2ex]
  \mathrm{C}\sigma_x\ket{10}&=\ket{0}\braket{0}{1}\otimes I\ket{0}+\ket{1}\braket{1}{1}\otimes\sigma_x\ket{0}
     =0+\ket{1}\otimes\ket{1}=\ket{11},\\[1.2ex]
  \mathrm{C}\sigma_x\ket{11}&=\ket{0}\braket{0}{1}\otimes I\ket{1}+\ket{1}\braket{1}{1}\otimes\sigma_x\ket{1}
     =0+\ket{1}\otimes\ket{0}=\ket{10}.
\end{align*}
On retrouve la table du CNOT : la cible n'est inversée que si le contrôle vaut $1$.

\paragraph{Un autre exemple : $U=Z$.} Ici
\[
  \mathrm{C}Z=\begin{pmatrix}I&0\\0&Z\end{pmatrix}=\mathrm{diag}(1,1,1,-1) :
\]
la porte multiplie $\ket{11}$ par $-1$ et ne touche pas les autres états de base. Elle est
symétrique entre les deux qubits. Elle se déduit du CNOT par
$\mathrm{C}Z=(I\otimes H)\,\mathrm{C}X\,(I\otimes H)$, où $\mathrm{C}X$ est le contrôlé-$X$.
En effet, $I\otimes H=\mathrm{diag}(H,H)$ (deux blocs $2\times2$ égaux à $H$), et le produit de
matrices diagonales par blocs se fait bloc par bloc :
\[
  (I\otimes H)\begin{pmatrix}I&0\\0&X\end{pmatrix}(I\otimes H)
  =\begin{pmatrix}H\,I\,H&0\\0&H\,X\,H\end{pmatrix}
  =\begin{pmatrix}I&0\\0&HXH\end{pmatrix},
\]
car $H^2=I$, et
\[
  HXH=\frac12\begin{pmatrix}1&1\\1&-1\end{pmatrix}\begin{pmatrix}0&1\\1&0\end{pmatrix}\begin{pmatrix}1&1\\1&-1\end{pmatrix}
  =\frac12\begin{pmatrix}1&1\\1&-1\end{pmatrix}\begin{pmatrix}1&-1\\1&1\end{pmatrix}
  =\frac12\begin{pmatrix}2&0\\0&-2\end{pmatrix}=Z
\]
(section~\ref{sec:hadamard} : $H$ échange les axes $x$ et $z$).

\subsubsection{Contrôlé-SWAP : la porte de Fredkin}

\begin{definition}{Porte de Fredkin (C-SWAP)}
Sur trois qubits $(x,y,z)$ (espace de dimension $2^3=8$), avec $x$ pour contrôle,
\[
  \mathrm{C\text{-}SWAP}=\ket{0}\bra{0}_x\otimes I_y\otimes I_z+\ket{1}\bra{1}_x\otimes\mathrm{SWAP}_{yz}.
\]
Si $x=0$, rien ne change ; si $x=1$, les qubits $y$ et $z$ sont échangés.
\end{definition}

Calcul sur quatre états de base (avec $I\ket{ab}=\ket{ab}$ et $\mathrm{SWAP}\ket{ab}=\ket{ba}$) :
\begin{align*}
  \mathrm{C\text{-}SWAP}\ket{000}&=\ket{0}\braket{0}{0}\otimes I\ket{00}+\ket{1}\braket{1}{0}\otimes\mathrm{SWAP}\ket{00}
     =\ket{0}\otimes\ket{00}+0=\ket{000},\\[1.2ex]
  \mathrm{C\text{-}SWAP}\ket{101}&=\ket{0}\braket{0}{1}\otimes I\ket{01}+\ket{1}\braket{1}{1}\otimes\mathrm{SWAP}\ket{01}
     =0+\ket{1}\otimes\ket{10}=\ket{110},\\[1.2ex]
  \mathrm{C\text{-}SWAP}\ket{110}&=\ket{0}\braket{0}{1}\otimes I\ket{10}+\ket{1}\braket{1}{1}\otimes\mathrm{SWAP}\ket{10}
     =0+\ket{1}\otimes\ket{01}=\ket{101},\\[1.2ex]
  \mathrm{C\text{-}SWAP}\ket{010}&=\ket{0}\braket{0}{0}\otimes I\ket{10}+\ket{1}\braket{1}{0}\otimes\mathrm{SWAP}\ket{10}
     =\ket{0}\otimes\ket{10}+0=\ket{010}.
\end{align*}
Les huit états de base donnent :
\begin{center}
\renewcommand{\arraystretch}{1.3}
\begin{tabular}{@{}ccccccccc@{}}
\toprule
Entrée & $\ket{000}$ & $\ket{001}$ & $\ket{010}$ & $\ket{011}$ & $\ket{100}$ & $\ket{101}$ & $\ket{110}$ & $\ket{111}$ \\
\midrule
Sortie & $\ket{000}$ & $\ket{001}$ & $\ket{010}$ & $\ket{011}$ & $\ket{100}$ & $\ket{110}$ & $\ket{101}$ & $\ket{111}$ \\
\bottomrule
\end{tabular}
\end{center}
Seuls $\ket{101}$ et $\ket{110}$ changent (contrôle à $1$ et $y\neq z$) ; pour $\ket{100}$ et
$\ket{111}$ l'échange de deux qubits égaux ne change rien. Dans la base
$\ket{000},\ket{001},\dots,\ket{111}$ (ordre binaire), la matrice est
$\mathrm{diag}(I_4,\mathrm{SWAP})$ (deux blocs $4\times4$) :
\[
  \mathrm{C\text{-}SWAP}=\begingroup\setlength{\arraycolsep}{3pt}\begin{pmatrix}
    1&0&0&0&0&0&0&0\\
    0&1&0&0&0&0&0&0\\
    0&0&1&0&0&0&0&0\\
    0&0&0&1&0&0&0&0\\
    0&0&0&0&1&0&0&0\\
    0&0&0&0&0&0&1&0\\
    0&0&0&0&0&1&0&0\\
    0&0&0&0&0&0&0&1
  \end{pmatrix}\endgroup .
\]
C'est l'identité dont les lignes $6$ et $7$ sont échangées.

\paragraph{Contrôle en superposition.} Avec $\ket{+}_x\otimes\ket{0}_y\otimes\ket{1}_z
=\frac{1}{\sqrt2}\big(\ket{001}+\ket{101}\big)$ :
\[
  \mathrm{C\text{-}SWAP}\,\frac{1}{\sqrt2}\big(\ket{001}+\ket{101}\big)
  =\frac{1}{\sqrt2}\big(\ket{001}+\ket{110}\big).
\]
Le résultat n'est pas un produit : les coefficients de $\ket{x}\otimes\ket{yz}$ forment la
matrice (lignes $x=0,1$, colonnes $yz=00,01,10,11$)
\[
  \frac{1}{\sqrt2}\begin{pmatrix}0&1&0&0\\0&0&1&0\end{pmatrix},
\]
dont les deux lignes ne sont pas proportionnelles. Le qubit $x$ est intriqué avec la paire
$(y,z)$.

\subsubsection{Contrôlé-contrôlé-NOT : la porte de Toffoli}

\begin{definition}{Porte de Toffoli (CCNOT)}
Sur trois qubits $(x,y,z)$, avec $x$ et $y$ pour contrôles et $z$ pour cible,
\[
  \mathrm{CCNOT}=\ket{0}\bra{0}_x\otimes I_y\otimes I_z
   +\ket{1}\bra{1}_x\otimes\Big[\ket{0}\bra{0}_y\otimes I_z+\ket{1}\bra{1}_y\otimes\sigma_x\Big].
\]
Le crochet est un CNOT (contrôle $y$, cible $z$), appliqué seulement si $x=1$ : si $x=0$, rien ;
si $x=1$ et $y=0$, rien ; si $x=1$ et $y=1$, le qubit $z$ est inversé.
\end{definition}

Calcul de $\mathrm{CCNOT}\ket{110}$ :
\begin{align*}
  \mathrm{CCNOT}\ket{110}&=\ket{0}\braket{0}{1}\otimes I\ket{1}\otimes I\ket{0}
     +\ket{1}\braket{1}{1}\otimes\Big[\ket{0}\braket{0}{1}\otimes I\ket{0}
       +\ket{1}\braket{1}{1}\otimes\sigma_x\ket{0}\Big]\\
  &=0+\ket{1}\otimes\Big[0+\ket{1}\otimes\ket{1}\Big]=\ket{111}.
\end{align*}
Pour les huit états de base, $\mathrm{CCNOT}\ket{x,y,z}=\ket{x,\,y,\,z\oplus xy}$ :
\begin{center}
\renewcommand{\arraystretch}{1.3}
\begin{tabular}{@{}ccccccccc@{}}
\toprule
Entrée & $\ket{000}$ & $\ket{001}$ & $\ket{010}$ & $\ket{011}$ & $\ket{100}$ & $\ket{101}$ & $\ket{110}$ & $\ket{111}$ \\
\midrule
Sortie & $\ket{000}$ & $\ket{001}$ & $\ket{010}$ & $\ket{011}$ & $\ket{100}$ & $\ket{101}$ & $\ket{111}$ & $\ket{110}$ \\
\bottomrule
\end{tabular}
\end{center}
Seuls $\ket{110}$ et $\ket{111}$ sont échangés : la matrice est $\mathrm{diag}(I_6,X)$
(un bloc $6\times6$ égal à l'identité, puis $X$),
\[
  \mathrm{CCNOT}=\begingroup\setlength{\arraycolsep}{3pt}\begin{pmatrix}
    1&0&0&0&0&0&0&0\\
    0&1&0&0&0&0&0&0\\
    0&0&1&0&0&0&0&0\\
    0&0&0&1&0&0&0&0\\
    0&0&0&0&1&0&0&0\\
    0&0&0&0&0&1&0&0\\
    0&0&0&0&0&0&0&1\\
    0&0&0&0&0&0&1&0
  \end{pmatrix}\endgroup .
\]

\begin{remarque}[Portes réversibles]
Les portes de Fredkin et de Toffoli sont des permutations des états de base qui échangent
seulement deux d'entre eux : elles sont leurs propres inverses ($\mathrm{C\text{-}SWAP}^2=\mathrm{CCNOT}^2=\mathrm{Id}$)
et unitaires. Avec $z=0$, la Toffoli calcule le ET logique $x\wedge y$ sans rien effacer :
$\ket{x,y,0}\to\ket{x,y,xy}$ ; avec $z=1$, elle calcule le NON-ET. Elle sert à écrire tout calcul
classique de façon réversible.
\end{remarque}

% ---------------------------------------------------------------------
\subsection{Récapitulatif du chapitre}
% ---------------------------------------------------------------------

\begin{center}
\renewcommand{\arraystretch}{1.5}
\begin{tabular}{@{}lp{12cm}@{}}
\toprule
\textbf{Notion} & \textbf{Formule ou résultat} \\
\midrule
CNOT & $\ket{b,a}\to\ket{b\oplus a,a}$ (contrôle $a$ écrit en dernier), $\mathrm{CNOT}^2=\mathrm{Id}$ \\
SWAP & $\mathrm{SWAP}\ket{a,b}=\ket{b,a}$, \quad $\mathrm{CNOT}_{21}=\mathrm{SWAP}\,\mathrm{CNOT}_{12}\,\mathrm{SWAP}$ \\
Circuit de Bell & $(I\otimes H)$ puis CNOT : $\ket{\beta^{xy}}=\frac{1}{\sqrt2}\big(\ket{y,0}+(-1)^x\ket{\bar y,1}\big)$, $\ket{\beta^{00}}=\ket{\Phi^+}$ \\
Information & $I(A;B)=H(A)+H(B)-H(A,B)$, \quad $H(A|B)=H(A,B)-H(B)$ \\
Trace partielle & $\operatorname{Tr}_B\big(\ket{a}\bra{a'}\otimes\ket{b}\bra{b'}\big)=\ket{a}\bra{a'}\braket{b'}{b}$ ; état de Bell : $\rho_A=\frac12\mathrm{Id}$ \\
État de Bell (bits) & $H(A,B)=0$, \quad $H(A)=H(B)=1$, \quad $I(A;B)=2$, \quad $H(A|B)=-1$ \\
Pur / mélangé & $\operatorname{Tr}\rho^2=1$ / $<1$ ; $\ket{+}$ : $1$ ; qubit d'une paire de Bell : $\frac12$ ; source de Bob : $\frac58$ \\
Téléportation & $1$ e-bit $+\ 2$ bits classiques ; correction : $X$ si $a=1$, puis $Z$ si $q=1$ ; $P(a,q)=\frac14$ \\
Contrôlé-$U$ & $\ket{0}\bra{0}\otimes I+\ket{1}\bra{1}\otimes U=\mathrm{diag}(I,U)$ ; $\ket{x,y}\to\ket{x}\otimes U^x\ket{y}$ \\
Fredkin & $\ket{0}\bra{0}\otimes I\otimes I+\ket{1}\bra{1}\otimes\mathrm{SWAP}$ ; $\ket{101}\leftrightarrow\ket{110}$ \\
Toffoli & $\ket{x,y,z}\to\ket{x,y,z\oplus xy}$ ; $\ket{110}\leftrightarrow\ket{111}$ \\
\bottomrule
\end{tabular}
\end{center}
.
%  Comme chapitre4.tex, chapitre5.tex et chapitre6.tex, il ouvre une NOUVELLE
%  SECTION (numéro 5). Pour le rattacher à la section précédente, remplacer
%  \section par \subsection et chaque \subsection par \subsubsection.
%  Pas de préambule : environnements (definition, resultat, remarque),
%  macros (\ii, \ee, \dd, \ket, \bra, \braket) et couleurs (insarouge,
%  insagris) sont définis dans main.tex. Les symboles de circuits (contrôle,
%  cible, croix, mesure) sont définis juste après avec \tikzset.
% =====================================================================

% Macros de Dirac (déjà définies dans main.tex récent ; ce bloc évite l'erreur
% « Undefined control sequence \ket » si main.tex est une ancienne version).
\providecommand{\ket}[1]{|#1\rangle}
\providecommand{\bra}[1]{\langle #1|}
\providecommand{\braket}[2]{\langle #1|#2\rangle}

% Symboles de circuits pour les figures de ce chapitre.
\tikzset{
  ctl/.pic   = {\fill[insarouge] (0,0) circle (2.6pt);},
  cible/.pic = {\draw[insarouge, thick, fill=white] (0,0) circle (0.28);
                \draw[insarouge, thick] (-0.28,0) -- (0.28,0);
                \draw[insarouge, thick] (0,-0.28) -- (0,0.28);},
  croix/.pic = {\draw[insarouge, thick] (-0.17,-0.17) -- (0.17,0.17);
                \draw[insarouge, thick] (-0.17,0.17) -- (0.17,-0.17);},
  mesure/.pic = {\draw[thick, fill=white] (-0.3,-0.27) rectangle (0.3,0.27);
                 \draw[thick] (-0.19,-0.13) arc[start angle=155, end angle=25, radius=0.21];
                 \draw[thick, -{Stealth[length=1.3mm]}] (0,-0.17) -- (0.14,0.12);},
  porte/.style = {draw=insarouge, fill=insarouge!8, minimum size=7mm, thick},
}

% =====================================================================
\section{Portes multi-qubits, circuits et téléportation}\label{sec:circuits}
% =====================================================================

La section précédente a décrit les états de deux qubits. On étudie ici les portes qui
agissent sur plusieurs qubits : le CNOT et le SWAP, puis les portes contrôlées à deux et à
trois qubits (contrôlé-$U$, Fredkin, Toffoli). On assemble ensuite des portes dans un
circuit qui fabrique un état de Bell, on mesure l'information contenue dans un tel état,
on distingue les états purs des états mélangés, et on démontre le protocole de
téléportation. Tous les calculs sont détaillés.

\begin{remarque}[Convention pour les circuits]
Dans un circuit, le temps va de gauche à droite et chaque qubit a un fil. Les fils sont
dessinés dans l'ordre \emph{inverse} de l'écriture des kets : le fil du \emph{haut} porte le
\emph{dernier} chiffre du ket, le fil du \emph{bas} le premier. Par exemple, si le qubit $A$
est dessiné au-dessus du qubit $B$, le ket s'écrit $\ket{b,a}=\ket{b}_B\otimes\ket{a}_A$.
La figure~\ref{fig:bell} de la section précédente suit la même convention.
\end{remarque}

% ---------------------------------------------------------------------
\subsection{Porte CNOT}\label{sec:cnot}
% ---------------------------------------------------------------------

\begin{definition}{Porte CNOT (NOT contrôlé)}
La porte CNOT agit sur deux qubits : le qubit $A$ est le \emph{contrôle}, le qubit $B$ est la
\emph{cible}. La cible est inversée (porte $X$) si le contrôle vaut $1$ et reste inchangée
s'il vaut $0$. Avec le système écrit $(B,A)$, soit $\ket{b,a}=\ket{b}_B\otimes\ket{a}_A$ :
\[
  \mathrm{CNOT}\,\ket{b,a}=\ket{b\oplus a,\,a},
  \qquad
  \mathrm{CNOT}=I\otimes\ket{0}\bra{0}+X\otimes\ket{1}\bra{1},
\]
où $\oplus$ est l'addition modulo $2$ ($0\oplus0=1\oplus1=0$ et $0\oplus1=1\oplus0=1$).
\end{definition}

Sur le circuit de la figure~\ref{fig:cnot}, le point noir posé sur le fil de $A$ marque le
contrôle et le symbole $\oplus$ posé sur le fil de $B$ marque la cible.

\begin{figure}[!ht]
\centering
\begin{tikzpicture}[>=Stealth, font=\small]
  \draw (0,1.2) -- (6,1.2);
  \draw (0,0) -- (6,0);
  \node[anchor=east] at (0,1.2) {$\ket{a}_A$};
  \node[anchor=east] at (0,0) {$\ket{b}_B$};
  \node[anchor=west] at (6,1.2) {$\ket{a}_A$};
  \node[anchor=west] at (6,0) {$\ket{b\oplus a}_B$};
  \draw[insarouge, thick] (3,1.2) -- (3,0);
  \pic at (3,1.2) {ctl};
  \pic at (3,0) {cible};
  \node[anchor=south, text=insagris] at (3,1.5) {qubit de contrôle};
  \node[anchor=north, text=insagris] at (3,-0.5) {qubit cible};
\end{tikzpicture}
\caption{Porte CNOT : le fil du haut ($A$) est le contrôle, le fil du bas ($B$) est la
cible. Le ket s'écrit $\ket{b,a}$ (le fil du haut donne le dernier chiffre).}
\label{fig:cnot}
\end{figure}

\subsubsection{Table de vérité}

\begin{center}
\renewcommand{\arraystretch}{1.3}
\begin{tabular}{@{}cccc@{}}
\toprule
Entrée $\ket{b,a}$ & Contrôle $a$ & Sortie $\ket{b\oplus a,a}$ & Effet sur la cible $B$ \\
\midrule
$\ket{00}$ & $0$ & $\ket{00}$ & aucun \\
$\ket{01}$ & $1$ & $\ket{11}$ & $0\to1$ \\
$\ket{10}$ & $0$ & $\ket{10}$ & aucun \\
$\ket{11}$ & $1$ & $\ket{01}$ & $1\to0$ \\
\bottomrule
\end{tabular}
\end{center}

\subsubsection{Matrice du CNOT}

La $j$-ième colonne d'une matrice $4\times4$ est l'image du $j$-ième vecteur de base
($\ket{00},\ket{01},\ket{10},\ket{11}$, dans cet ordre). La table donne :
\begin{align*}
  \mathrm{CNOT}\ket{00}&=\ket{00}=\begin{pmatrix}1\\0\\0\\0\end{pmatrix},&
  \mathrm{CNOT}\ket{01}&=\ket{11}=\begin{pmatrix}0\\0\\0\\1\end{pmatrix},\\[1.5ex]
  \mathrm{CNOT}\ket{10}&=\ket{10}=\begin{pmatrix}0\\0\\1\\0\end{pmatrix},&
  \mathrm{CNOT}\ket{11}&=\ket{01}=\begin{pmatrix}0\\1\\0\\0\end{pmatrix},
\end{align*}
d'où, en juxtaposant ces quatre colonnes,
\[
  \mathrm{CNOT}=\begin{pmatrix}
    1 & 0 & 0 & 0\\ 0 & 0 & 0 & 1\\ 0 & 0 & 1 & 0\\ 0 & 1 & 0 & 0
  \end{pmatrix}.
\]

On retrouve cette matrice avec la formule $I\otimes\ket{0}\bra{0}+X\otimes\ket{1}\bra{1}$ et
la règle $A\otimes B=(A_{ij}B)$ (section~\ref{sec:op2}), où
$\ket{0}\bra{0}=\begin{pmatrix}1&0\\0&0\end{pmatrix}$ et
$\ket{1}\bra{1}=\begin{pmatrix}0&0\\0&1\end{pmatrix}$ :
\[
  I\otimes\ket{0}\bra{0}=\begin{pmatrix}1&0&0&0\\0&0&0&0\\0&0&1&0\\0&0&0&0\end{pmatrix},
  \qquad
  X\otimes\ket{1}\bra{1}=\begin{pmatrix}0&0&0&0\\0&0&0&1\\0&0&0&0\\0&1&0&0\end{pmatrix},
\]
et la somme de ces deux matrices est la matrice du CNOT.

\subsubsection{Action sur un état quelconque}

Pour $\ket{\Psi}=c_{00}\ket{00}+c_{01}\ket{01}+c_{10}\ket{10}+c_{11}\ket{11}$ :
\[
  \mathrm{CNOT}\begin{pmatrix}c_{00}\\c_{01}\\c_{10}\\c_{11}\end{pmatrix}
  =\begin{pmatrix}
    1\cdot c_{00}+0\cdot c_{01}+0\cdot c_{10}+0\cdot c_{11}\\
    0\cdot c_{00}+0\cdot c_{01}+0\cdot c_{10}+1\cdot c_{11}\\
    0\cdot c_{00}+0\cdot c_{01}+1\cdot c_{10}+0\cdot c_{11}\\
    0\cdot c_{00}+1\cdot c_{01}+0\cdot c_{10}+0\cdot c_{11}
  \end{pmatrix}
  =\begin{pmatrix}c_{00}\\c_{11}\\c_{10}\\c_{01}\end{pmatrix}.
\]
Le CNOT échange les amplitudes de $\ket{01}$ et de $\ket{11}$, c'est-à-dire des deux états
dont le contrôle vaut $1$, et laisse les autres.

\paragraph{Exemple : contrôle en superposition.} Soit $\ket{0}_B\otimes\ket{+}_A
=\frac{1}{\sqrt2}\big(\ket{00}+\ket{01}\big)=\frac{1}{\sqrt2}(1,1,0,0)^{\mathsf T}$,
un état séparable. Alors
\[
  \mathrm{CNOT}\,\frac{1}{\sqrt2}\begin{pmatrix}1\\1\\0\\0\end{pmatrix}
  =\frac{1}{\sqrt2}\begin{pmatrix}1\\0\\0\\1\end{pmatrix}
  =\frac{1}{\sqrt2}\big(\ket{00}+\ket{11}\big)=\ket{\Phi^+}.
\]
Le résultat est intriqué ($D=c_{00}c_{11}-c_{01}c_{10}=\frac12\neq0$ ; section~\ref{sec:separable}) :
un CNOT dont le contrôle est en superposition crée de l'intrication.

\paragraph{Réversibilité.} Deux CNOT successifs redonnent l'état initial :
$\ket{b,a}\to\ket{b\oplus a,a}\to\ket{b\oplus a\oplus a,a}=\ket{b,a}$, donc
$\mathrm{CNOT}^2=\mathrm{Id}$. La matrice est réelle et symétrique, donc
$\mathrm{CNOT}^\dagger\mathrm{CNOT}=\mathrm{CNOT}^2=\mathrm{Id}$ : elle est bien unitaire.

\begin{remarque}[Contrôle en dernier ou en premier]
La matrice dépend de l'ordre d'écriture du ket. Notons $\mathrm{CNOT}_{21}$ le CNOT dont le
contrôle est le \emph{second} chiffre du ket (c'est la matrice ci-dessus) et
$\mathrm{CNOT}_{12}$ celui dont le contrôle est le \emph{premier} chiffre, utilisé à la
section précédente :
\[
  \mathrm{CNOT}_{12}=\ket{0}\bra{0}\otimes I+\ket{1}\bra{1}\otimes X
  =\begin{pmatrix}1&0&0&0\\0&1&0&0\\0&0&0&1\\0&0&1&0\end{pmatrix}.
\]
Un même circuit (contrôle $A$, cible $B$) a donc deux matrices, selon que l'on écrit
$A$ en dernier ($\mathrm{CNOT}_{21}$) ou en premier ($\mathrm{CNOT}_{12}$). Les deux sont
reliées par le SWAP (section suivante).
\end{remarque}

% ---------------------------------------------------------------------
\subsection{Porte SWAP}\label{sec:swap}
% ---------------------------------------------------------------------

\begin{definition}{Porte SWAP (échange)}
La porte SWAP échange les états des deux qubits :
\[
  \mathrm{SWAP}\big(\ket{\psi}\otimes\ket{\phi}\big)=\ket{\phi}\otimes\ket{\psi}.
\]
Sur les états de base, $\mathrm{SWAP}\ket{a,b}=\ket{b,a}$ ; on prolonge par linéarité.
\end{definition}

Son symbole est formé de deux croix reliées par un trait vertical
(figure~\ref{fig:swap}).

\begin{figure}[!ht]
\centering
\begin{tikzpicture}[>=Stealth, font=\small]
  \draw (0,1.2) -- (6,1.2);
  \draw (0,0) -- (6,0);
  \node[anchor=east] at (0,1.2) {$\ket{a}$};
  \node[anchor=east] at (0,0) {$\ket{b}$};
  \node[anchor=west] at (6,1.2) {$\ket{b}$};
  \node[anchor=west] at (6,0) {$\ket{a}$};
  \draw[insarouge, thick] (3,1.2) -- (3,0);
  \pic at (3,1.2) {croix};
  \pic at (3,0) {croix};
\end{tikzpicture}
\caption{Porte SWAP : les deux qubits échangent leurs états.}
\label{fig:swap}
\end{figure}

\subsubsection{Table et matrice}

\begin{center}
\renewcommand{\arraystretch}{1.3}
\begin{tabular}{@{}cc@{}}
\toprule
Entrée & Sortie \\
\midrule
$\ket{00}$ & $\ket{00}$ \\
$\ket{01}$ & $\ket{10}$ \\
$\ket{10}$ & $\ket{01}$ \\
$\ket{11}$ & $\ket{11}$ \\
\bottomrule
\end{tabular}
\end{center}

Les images des quatre vecteurs de base sont les colonnes de la matrice :
\[
  \mathrm{SWAP}\ket{00}=\begin{pmatrix}1\\0\\0\\0\end{pmatrix},\quad
  \mathrm{SWAP}\ket{01}=\begin{pmatrix}0\\0\\1\\0\end{pmatrix},\quad
  \mathrm{SWAP}\ket{10}=\begin{pmatrix}0\\1\\0\\0\end{pmatrix},\quad
  \mathrm{SWAP}\ket{11}=\begin{pmatrix}0\\0\\0\\1\end{pmatrix},
\]
\[
  \mathrm{SWAP}=\begin{pmatrix}
    1 & 0 & 0 & 0\\ 0 & 0 & 1 & 0\\ 0 & 1 & 0 & 0\\ 0 & 0 & 0 & 1
  \end{pmatrix}.
\]
Appliquée à $\ket{01}=(0,1,0,0)^{\mathsf T}$, cette matrice sélectionne sa deuxième colonne :
\[
  \begin{pmatrix}
    1 & 0 & 0 & 0\\ 0 & 0 & 1 & 0\\ 0 & 1 & 0 & 0\\ 0 & 0 & 0 & 1
  \end{pmatrix}
  \begin{pmatrix}0\\1\\0\\0\end{pmatrix}
  =\begin{pmatrix}0\\0\\1\\0\end{pmatrix}=\ket{10}.
\]
Sur un état quelconque, $\mathrm{SWAP}\,(c_{00},c_{01},c_{10},c_{11})^{\mathsf T}
=(c_{00},c_{10},c_{01},c_{11})^{\mathsf T}$ : les amplitudes de $\ket{01}$ et $\ket{10}$
sont échangées.

\subsubsection{Accord avec la définition}

Pour deux états $\ket{\psi}=(a_0,a_1)^{\mathsf T}$ et $\ket{\phi}=(b_0,b_1)^{\mathsf T}$ :
\[
  \mathrm{SWAP}\,(\ket{\psi}\otimes\ket{\phi})
  =\mathrm{SWAP}\begin{pmatrix}a_0b_0\\a_0b_1\\a_1b_0\\a_1b_1\end{pmatrix}
  =\begin{pmatrix}a_0b_0\\a_1b_0\\a_0b_1\\a_1b_1\end{pmatrix}
  =\begin{pmatrix}b_0\\b_1\end{pmatrix}\otimes\begin{pmatrix}a_0\\a_1\end{pmatrix}
  =\ket{\phi}\otimes\ket{\psi},
\]
ce qui est bien la définition. La porte s'écrit aussi avec des projecteurs :
\[
  \mathrm{SWAP}=\sum_{a,b\in\{0,1\}}\ket{a}\bra{b}\otimes\ket{b}\bra{a}.
\]
En effet, pour $\ket{c}\otimes\ket{d}$ on trouve
$\sum_{a,b}\ket{a}\braket{b}{c}\otimes\ket{b}\braket{a}{d}$, et seul le terme $b=c$, $a=d$
est non nul : le résultat est $\ket{d}\otimes\ket{c}$.

\paragraph{Exemples.}
\begin{itemize}[nosep]
  \item Deux états différents : $\mathrm{SWAP}\big(\ket{0}\otimes\ket{+}\big)
        =\mathrm{SWAP}\,\frac{1}{\sqrt2}(1,1,0,0)^{\mathsf T}=\frac{1}{\sqrt2}(1,0,1,0)^{\mathsf T}
        =\ket{+}\otimes\ket{0}$.
  \item L'état $\ket{\psi}=\frac{1}{\sqrt2}\ket{00}+\frac{1}{\sqrt2}\ket{11}$ ne change pas :
        $\mathrm{SWAP}\ket{\psi}=\frac{1}{\sqrt2}\ket{00}+\frac{1}{\sqrt2}\ket{11}=\ket{\psi}$,
        car $\ket{00}$ et $\ket{11}$ sont invariants. De même
        $\mathrm{SWAP}\ket{\Phi^-}=\ket{\Phi^-}$ et $\mathrm{SWAP}\ket{\Psi^+}=\ket{\Psi^+}$, alors que
        $\mathrm{SWAP}\ket{\Psi^-}=\frac{1}{\sqrt2}(\ket{10}-\ket{01})=-\ket{\Psi^-}$.
\end{itemize}

\paragraph{Propriétés.} Échanger deux fois redonne l'état de départ :
$\mathrm{SWAP}^2=\mathrm{Id}$ ; la matrice est réelle et symétrique, donc
$\mathrm{SWAP}^\dagger\mathrm{SWAP}=\mathrm{SWAP}^2=\mathrm{Id}$ : elle est unitaire.

\subsubsection{Lien avec le CNOT}

\begin{resultat}{SWAP échange les rôles de contrôle et de cible}
\[
  \mathrm{CNOT}_{21}=\mathrm{SWAP}\;\mathrm{CNOT}_{12}\;\mathrm{SWAP}.
\]
\end{resultat}

\emph{Sur les états de base.}
$\mathrm{SWAP}\,\mathrm{CNOT}_{12}\,\mathrm{SWAP}\ket{b,a}
=\mathrm{SWAP}\,\mathrm{CNOT}_{12}\ket{a,b}
=\mathrm{SWAP}\ket{a,\,a\oplus b}=\ket{a\oplus b,\,a}=\mathrm{CNOT}_{21}\ket{b,a}$.

\emph{Par les matrices.} Multiplier à gauche par $\mathrm{SWAP}$ échange les lignes $2$ et
$3$, multiplier à droite l'échange les colonnes $2$ et $3$ :
\[
  \mathrm{SWAP}\,\mathrm{CNOT}_{12}
  =\begin{pmatrix}1&0&0&0\\0&0&0&1\\0&1&0&0\\0&0&1&0\end{pmatrix},
  \qquad
  \mathrm{SWAP}\,\mathrm{CNOT}_{12}\,\mathrm{SWAP}
  =\begin{pmatrix}1&0&0&0\\0&0&0&1\\0&0&1&0\\0&1&0&0\end{pmatrix}=\mathrm{CNOT}_{21}.
\]

\begin{remarque}[SWAP avec trois CNOT]
$\mathrm{SWAP}=\mathrm{CNOT}_{12}\,\mathrm{CNOT}_{21}\,\mathrm{CNOT}_{12}$. En effet, sur
$\ket{a,b}$ (on applique les portes de droite à gauche) :
\[
  \ket{a,b}\ \xrightarrow{\ \mathrm{CNOT}_{12}\ }\ \ket{a,\,a\oplus b}
  \ \xrightarrow{\ \mathrm{CNOT}_{21}\ }\ \ket{a\oplus(a\oplus b),\,a\oplus b}=\ket{b,\,a\oplus b}
  \ \xrightarrow{\ \mathrm{CNOT}_{12}\ }\ \ket{b,\,b\oplus(a\oplus b)}=\ket{b,a}.
\]
Les deux membres coïncident sur les états de base, donc sur tous les états par linéarité.
\end{remarque}

% ---------------------------------------------------------------------
\subsection{Un circuit quantique : le générateur d'états de Bell}\label{sec:circuit-bell}
% ---------------------------------------------------------------------

On assemble une porte de Hadamard et un CNOT. Le système est $(A,B)$ avec le ket
$\ket{a,b}=\ket{a}_A\otimes\ket{b}_B$ ; d'après la convention de la section, $B$ est dessiné
en haut et $A$ en bas (figure~\ref{fig:circuit-bell}). Le circuit a deux étapes :
\begin{enumerate}[nosep]
  \item la porte $I\otimes H$ : identité sur $A$, Hadamard sur $B$ ;
  \item un CNOT de contrôle $B$ et de cible $A$ (la matrice $\mathrm{CNOT}_{21}$ de la
        section~\ref{sec:cnot}, puisque le contrôle est le second chiffre).
\end{enumerate}
Les traits pointillés délimitent les états $\ket{\Pi_0}$, $\ket{\Pi_1}$ et $\ket{\Pi_2}$ :
avant les portes, après $I\otimes H$, après le CNOT.

\begin{figure}[!ht]
\centering
\begin{tikzpicture}[>=Stealth, font=\small]
  % fils : B en haut, A en bas
  \draw (0,1.4) -- (7.6,1.4);
  \draw (0,0) -- (7.6,0);
  \node[anchor=east] at (0,1.4) {$\ket{0}_B$};
  \node[anchor=east] at (0,0) {$\ket{0}_A$};
  % première étape : H sur B, identité (pointillée) sur A
  \node[porte] at (2.2,1.4) {$H$};
  \node[draw=insagris, dashed, fill=white, minimum size=7mm] at (2.2,0) {$I$};
  % CNOT : contrôle B (haut), cible A (bas)
  \draw[insarouge, thick] (5,1.4) -- (5,0);
  \pic at (5,1.4) {ctl};
  \pic at (5,0) {cible};
  % accolade de sortie
  \draw[insagris, thick] (7.75,1.6) -- (7.9,1.6) -- (7.9,-0.2) -- (7.75,-0.2);
  \node[anchor=west] at (7.9,0.7) {$\ket{\beta^{00}}_{AB}$};
  % états intermédiaires
  \foreach \x in {0.7,3.6,6.4}
    \draw[dashed, insagris!70] (\x,1.9) -- (\x,-0.7);
  \node[anchor=north, font=\footnotesize] at (0.7,-0.7) {$\ket{\Pi_0}=\ket{00}$};
  \node[anchor=north, font=\footnotesize] at (3.6,-0.7) {$\ket{\Pi_1}=\ket{0}\otimes\ket{+}$};
  \node[anchor=north, font=\footnotesize] at (6.4,-0.7) {$\ket{\Pi_2}=\ket{\beta^{00}}$};
\end{tikzpicture}
\caption{Générateur d'états de Bell : $I\otimes H$, puis un CNOT de contrôle $B$ (fil du haut)
et de cible $A$ (fil du bas). L'entrée représentée est $\ket{00}$.}
\label{fig:circuit-bell}
\end{figure}

\subsubsection{Entrée \texorpdfstring{$\ket{00}$}{|00>} : calcul pas à pas}

\begin{align*}
  \ket{\Pi_0}&=\ket{0}_A\otimes\ket{0}_B=\ket{00},\\[1ex]
  \ket{\Pi_1}&=(I\otimes H)\ket{\Pi_0}=I\ket{0}_A\otimes H\ket{0}_B=\ket{0}\otimes\ket{+}
    =\frac{\ket{00}+\ket{01}}{\sqrt2},\\[1ex]
  \ket{\Pi_2}&=\mathrm{CNOT}\ket{\Pi_1}=\frac{\mathrm{CNOT}\ket{00}+\mathrm{CNOT}\ket{01}}{\sqrt2}
    =\frac{\ket{00}+\ket{11}}{\sqrt2}=\ket{\beta^{00}}=\ket{\Phi^+}.
\end{align*}
Dans $\ket{\Pi_1}$, le premier chiffre est celui de $A$ et le second celui de $B$. Le CNOT a
pour contrôle $B$ : $\ket{00}$ ($B=0$) ne change pas, $\ket{01}$ ($B=1$) devient $\ket{11}$
puisque $A$ est inversé. Le résultat est un \emph{état de Bell}.

Avec les matrices, la règle $A\otimes B=(A_{ij}B)$ donne
\[
  I\otimes H=\begin{pmatrix}1\cdot H&0\\0&1\cdot H\end{pmatrix}
  =\frac{1}{\sqrt2}\begin{pmatrix}1&1&0&0\\1&-1&0&0\\0&0&1&1\\0&0&1&-1\end{pmatrix}.
\]
Sa première colonne est $\ket{\Pi_1}=\frac{1}{\sqrt2}(1,1,0,0)^{\mathsf T}$ (image de $\ket{00}$).
Le CNOT échange les amplitudes de $\ket{01}$ et de $\ket{11}$ (section~\ref{sec:cnot}), donc
\[
  \ket{\Pi_2}=\mathrm{CNOT}\,\frac{1}{\sqrt2}\begin{pmatrix}1\\1\\0\\0\end{pmatrix}
  =\frac{1}{\sqrt2}\begin{pmatrix}1\\0\\0\\1\end{pmatrix}.
\]

\subsubsection{Les quatre entrées}

Prenons $B$ dans l'état $\ket{x}$ et $A$ dans l'état $\ket{y}$, avec $x,y\in\{0,1\}$. Alors
$\ket{\Pi_0}=\ket{y}_A\otimes\ket{x}_B=\ket{y,x}$. Comme
$H\ket{x}=\frac{1}{\sqrt2}\big(\ket{0}+(-1)^x\ket{1}\big)$ (c'est $\ket{+}$ pour $x=0$ et
$\ket{-}$ pour $x=1$),
\[
  \ket{\Pi_1}=\ket{y}\otimes H\ket{x}=\frac{1}{\sqrt2}\Big(\ket{y,0}+(-1)^x\ket{y,1}\Big).
\]
Le CNOT (contrôle $B$) ne change pas $\ket{y,0}$ et transforme $\ket{y,1}$ en
$\ket{\bar y,1}$, avec $\bar y=1-y$ :
\[
  \ket{\Pi_2}=\ket{\beta^{xy}}=\frac{1}{\sqrt2}\Big(\ket{y,0}+(-1)^x\ket{\bar y,1}\Big).
\]
Le premier indice de $\ket{\beta^{xy}}$ est l'entrée de $B$ (le qubit qui reçoit $H$), le
second celle de $A$ (la cible). Les quatre cas :

\begin{center}
\renewcommand{\arraystretch}{1.7}
\begin{tabular}{@{}ccccc@{}}
\toprule
$(B,A)=(x,y)$ & $\ket{\Pi_0}=\ket{y,x}$ & $\ket{\Pi_1}$ & $\ket{\Pi_2}=\ket{\beta^{xy}}$ & Nom \\
\midrule
$(0,0)$ & $\ket{00}$ & $\frac{1}{\sqrt2}(\ket{00}+\ket{01})$ & $\frac{1}{\sqrt2}(\ket{00}+\ket{11})$ & $\ket{\Phi^+}$ \\
$(1,0)$ & $\ket{01}$ & $\frac{1}{\sqrt2}(\ket{00}-\ket{01})$ & $\frac{1}{\sqrt2}(\ket{00}-\ket{11})$ & $\ket{\Phi^-}$ \\
$(0,1)$ & $\ket{10}$ & $\frac{1}{\sqrt2}(\ket{10}+\ket{11})$ & $\frac{1}{\sqrt2}(\ket{10}+\ket{01})$ & $\ket{\Psi^+}$ \\
$(1,1)$ & $\ket{11}$ & $\frac{1}{\sqrt2}(\ket{10}-\ket{11})$ & $\frac{1}{\sqrt2}(\ket{10}-\ket{01})$ & $-\ket{\Psi^-}$ \\
\bottomrule
\end{tabular}
\end{center}
Détail de la dernière ligne : $\ket{\Pi_1}=\ket{1}\otimes\ket{-}=\frac{1}{\sqrt2}(\ket{10}-\ket{11})$ ;
le CNOT laisse $\ket{10}$ et envoie $\ket{11}$ sur $\ket{01}$, d'où
$\frac{1}{\sqrt2}(\ket{10}-\ket{01})=-\frac{1}{\sqrt2}(\ket{01}-\ket{10})=-\ket{\Psi^-}$.

\paragraph{Forme matricielle.} $\mathrm{CNOT}_{21}$ échange les lignes $2$ et $4$ de
$I\otimes H$ (il échange les amplitudes de $\ket{01}$ et $\ket{11}$), donc
\[
  \mathrm{CNOT}_{21}\,(I\otimes H)=\frac{1}{\sqrt2}\begin{pmatrix}
    1 & 1 & 0 & 0\\ 0 & 0 & 1 & -1\\ 0 & 0 & 1 & 1\\ 1 & -1 & 0 & 0
  \end{pmatrix}.
\]
Ses colonnes sont les images de $\ket{00},\ket{01},\ket{10},\ket{11}$, c'est-à-dire de
$(B,A)=(0,0),(1,0),(0,1),(1,1)$ : ce sont $\ket{\beta^{00}}$, $\ket{\beta^{10}}$,
$\ket{\beta^{01}}$, $\ket{\beta^{11}}$, soit $\ket{\Phi^+}$, $\ket{\Phi^-}$, $\ket{\Psi^+}$
et $-\ket{\Psi^-}$.

\begin{remarque}[Comparaison avec la section précédente]
Le circuit de la figure~\ref{fig:bell} applique $H$ au \emph{premier} chiffre et prend celui-ci
pour contrôle ; ici $H$ et le contrôle portent sur le \emph{dernier} chiffre. Les deux
circuits se déduisent l'un de l'autre en échangeant les deux fils (conjugaison par le SWAP,
section~\ref{sec:swap}). Les quatre états de Bell obtenus sont les mêmes, avec les mêmes
indices, sauf le dernier : $-\ket{\Psi^-}$ ici, $+\ket{\Psi^-}$ là-bas. Ce signe est une phase
globale, sans effet sur les mesures (elle disparaît de $\rho$, section~\ref{sec:rho}).
\end{remarque}

\subsubsection{Matrice de densité de l'état de sortie}

Pour $\ket{\Psi}=\ket{\beta^{00}}=\frac{1}{\sqrt2}\ket{00}+\frac{1}{\sqrt2}\ket{11}$ :
\[
  \rho_{AB}=\ket{\Psi}\bra{\Psi}
  =\begin{pmatrix}\frac{1}{\sqrt2}\\0\\0\\\frac{1}{\sqrt2}\end{pmatrix}
   \begin{pmatrix}\frac{1}{\sqrt2}&0&0&\frac{1}{\sqrt2}\end{pmatrix}
  =\frac12\begin{pmatrix}1&0&0&1\\0&0&0&0\\0&0&0&0\\1&0&0&1\end{pmatrix}.
\]
Sous forme de ket-bra :
$\rho_{AB}=\frac12\big(\ket{00}\bra{00}+\ket{00}\bra{11}+\ket{11}\bra{00}+\ket{11}\bra{11}\big)$.
On vérifie $\operatorname{Tr}\rho_{AB}=\frac12(1+1)=1$ et
\[
  \rho_{AB}^2=\frac14\begin{pmatrix}2&0&0&2\\0&0&0&0\\0&0&0&0\\2&0&0&2\end{pmatrix}=\rho_{AB},
\]
donc $\operatorname{Tr}\rho_{AB}^2=1$ : le couple $(A,B)$ est dans un état \emph{pur}.

% ---------------------------------------------------------------------
\subsection{Information : variables aléatoires et états quantiques}\label{sec:info}
% ---------------------------------------------------------------------

\begin{definition}{Entropie, information mutuelle, entropie conditionnelle}
Soit $A$ une variable aléatoire de loi $p(a)$. Son \emph{entropie de Shannon}, en bits, est
\[
  H(A)=-\sum_a p(a)\log_2p(a)\qquad(\text{avec }0\log_20=0).
\]
Elle mesure l'incertitude moyenne sur $A$ : une pièce équilibrée a $H=1$ bit. Pour un couple
$(A,B)$ de loi $p(a,b)$, $H(A,B)=-\sum_{a,b}p(a,b)\log_2p(a,b)$, et
\begin{align*}
  I(A;B)&=H(A)+H(B)-H(A,B) &&\text{(information mutuelle)},\\
  H(A|B)&=H(A,B)-H(B)=H(A)-I(A;B) &&\text{(entropie conditionnelle)}.
\end{align*}
\end{definition}

Les deux écritures de $H(A|B)$ sont égales : la définition de $I(A;B)$ donne
$H(A,B)=H(A)+H(B)-I(A;B)$, donc $H(A,B)-H(B)=H(A)-I(A;B)$. L'information mutuelle $I(A;B)$
mesure ce que la connaissance de $B$ apprend sur $A$ ; $H(A|B)$ est l'incertitude qui reste sur
$A$ quand on connaît $B$.

\subsubsection{Deux variables aléatoires indépendantes}

Alice et Bob lancent chacun une pièce équilibrée, sans lien entre les deux lancers :
$p(a,b)=\frac14$ pour les quatre couples. Chaque pièce donne $0$ ou $1$ avec la probabilité $\frac12$, donc
\begin{align*}
  H(A)&=-\Big(\tfrac12\log_2\tfrac12+\tfrac12\log_2\tfrac12\Big)=1\text{ bit},\\
  H(B)&=1\text{ bit (même calcul)},\\
  H(A,B)&=-4\cdot\tfrac14\log_2\tfrac14=2\text{ bits},\\
  I(A;B)&=H(A)+H(B)-H(A,B)=1+1-2=0,\\
  H(A|B)&=H(A)-I(A;B)=1-0=H(A)=1\text{ bit}.
\end{align*}
Connaître $B$ n'apprend rien sur $A$.

\subsubsection{Deux variables solidaires}

Les deux pièces tombent toujours du même côté : $p(0,0)=p(1,1)=\frac12$ et
$p(0,1)=p(1,0)=0$. Les lois de chaque pièce sont inchangées ($p(a{=}0)=p(0,0)+p(0,1)=\frac12$,
etc.), donc $H(A)=H(B)=1$ bit, mais
\begin{align*}
  H(A,B)&=-2\cdot\tfrac12\log_2\tfrac12-2\cdot0\log_20=1\text{ bit},\\
  I(A;B)&=1+1-1=1\text{ bit},\qquad H(A|B)=1-1=0.
\end{align*}
Connaître $B$ détermine $A$ : plus aucune incertitude ne reste.

\subsubsection{Deux qubits dans un état de Bell}

Pour un système quantique, on remplace la loi de probabilité par la matrice de densité.

\begin{definition}{Trace partielle et entropie de von Neumann}
L'état du seul qubit $A$ d'un système $(A,B)$ de matrice de densité $\rho_{AB}$ est obtenu en
\emph{traçant sur $B$} : $\rho_A=\operatorname{Tr}_B(\rho_{AB})$. La trace partielle est
linéaire et remplace la partie $B$ de chaque terme par sa trace :
\[
  \operatorname{Tr}_B\big(\ket{a}\bra{a'}\otimes\ket{b}\bra{b'}\big)=\ket{a}\bra{a'}\;\braket{b'}{b}.
\]
L'\emph{entropie de von Neumann} d'un état $\rho$, de valeurs propres $\lambda_i$, est
$S(\rho)=-\operatorname{Tr}(\rho\log_2\rho)=-\sum_i\lambda_i\log_2\lambda_i$. En quantique,
$H(A)=S(\rho_A)$, $H(B)=S(\rho_B)$ et $H(A,B)=S(\rho_{AB})$ ; les définitions de
$I(A;B)$ et de $H(A|B)$ sont les mêmes.
\end{definition}

On part de l'état de Bell
$\rho_{AB}=\frac12\big(\ket{00}\bra{00}+\ket{00}\bra{11}+\ket{11}\bra{00}+\ket{11}\bra{11}\big)$
produit par le circuit de la section~\ref{sec:circuit-bell}. Chaque terme est un produit
d'opérateurs à un qubit :
$\ket{00}\bra{11}=\ket{0}\bra{1}\otimes\ket{0}\bra{1}$, etc. Alors
\begin{align*}
  \rho_A=\operatorname{Tr}_B(\rho_{AB})
  &=\tfrac12\Big(\ket{0}\bra{0}\,\braket{0}{0}+\ket{0}\bra{1}\,\braket{1}{0}
      +\ket{1}\bra{0}\,\braket{0}{1}+\ket{1}\bra{1}\,\braket{1}{1}\Big)\\
  &=\tfrac12\Big(\ket{0}\bra{0}\cdot1+\ket{0}\bra{1}\cdot0+\ket{1}\bra{0}\cdot0+\ket{1}\bra{1}\cdot1\Big)\\
  &=\tfrac12\big(\ket{0}\bra{0}+\ket{1}\bra{1}\big)=\tfrac12\,\mathrm{Id}
  =\begin{pmatrix}\frac12&0\\0&\frac12\end{pmatrix}.
\end{align*}
La formule des composantes de la section~\ref{sec:separable},
$(\rho_A)_{ij}=\rho_{(i0),(j0)}+\rho_{(i1),(j1)}$, donne le même résultat. En traçant sur $A$
au lieu de $B$, la règle est $\operatorname{Tr}_A(\ket{a}\bra{a'}\otimes\ket{b}\bra{b'})
=\braket{a'}{a}\,\ket{b}\bra{b'}$, et on trouve de même $\rho_B=\frac12\mathrm{Id}$.

Calcul des entropies :
\begin{itemize}[nosep]
  \item $\rho_{AB}$ est pur : ses valeurs propres sont $1,0,0,0$, donc
        $H(A,B)=S(\rho_{AB})=-1\cdot\log_21=0$ ;
  \item $\rho_A$ et $\rho_B$ ont pour valeurs propres $\frac12,\frac12$, donc
        $H(A)=H(B)=-\big(\frac12\log_2\frac12+\frac12\log_2\frac12\big)=1$ bit.
\end{itemize}
D'où
\begin{align*}
  I(A;B)&=H(A)+H(B)-H(A,B)=1+1-0=2\text{ bits},\\
  H(A|B)&=H(A,B)-H(B)=0-1=-1\text{ bit}\qquad(=H(A)-I(A;B)=1-2).
\end{align*}

\begin{remarque}[Une entropie conditionnelle négative]
Pour des variables classiques, $p(a,b)\leq p(b)$, donc $-\log_2p(a,b)\geq-\log_2p(b)$ et
$H(A,B)\geq H(B)$ : on a toujours $H(A|B)\geq0$, et $I(A;B)\leq H(A)$. Ici $H(A|B)=-1$ et
$I(A;B)=2>H(A)$ : ces valeurs n'ont pas d'équivalent classique, c'est une signature de
l'intrication. Le résultat est surprenant mais mathématiquement correct.
\end{remarque}

\begin{center}
\renewcommand{\arraystretch}{1.4}
\begin{tabular}{@{}lccc@{}}
\toprule
 & indépendantes & solidaires & état de Bell \\
\midrule
$H(A)$ & $1$ & $1$ & $1$ \\
$H(B)$ & $1$ & $1$ & $1$ \\
$H(A,B)$ & $2$ & $1$ & $0$ \\
$I(A;B)$ & $0$ & $1$ & $2$ \\
$H(A|B)$ & $1$ & $0$ & $-1$ \\
\bottomrule
\end{tabular}
\end{center}
(en bits ; les deux premières colonnes sont classiques, la troisième est quantique).

% ---------------------------------------------------------------------
\subsection{Système pur et système mélangé}\label{sec:pur-melange}
% ---------------------------------------------------------------------

Rappel (section~\ref{sec:pur-mixte}) : si $\lambda_i$ sont les valeurs propres de la matrice de
densité $\rho$, alors $\operatorname{Tr}\rho^2=\sum_i\lambda_i^2\leq1$. Le système est
\emph{pur} si $\operatorname{Tr}\rho^2=1$ (alors $\rho=\ket{\psi}\bra{\psi}$ et $\rho^2=\rho$),
et \emph{mélangé} si $\operatorname{Tr}\rho^2<1$. Pour un qubit,
$\operatorname{Tr}\rho^2=\frac{1+|\vec r|^2}{2}$ (section~\ref{sec:bloch}) est compris entre
$\frac12$ et $1$.

\subsubsection{Un système pur : exercice}

\emph{Énoncé.} Soit $\ket{\psi}=\frac{1}{\sqrt2}\ket{0}+\frac{1}{\sqrt2}\ket{1}$.
\begin{enumerate}[label=\alph*), nosep]
  \item Calculer la matrice de densité $\rho$.
  \item Calculer la pureté du système quantique.
  \item Calculer $\langle\sigma_x\rangle$ et $\langle\sigma_z\rangle$.
  \item Écrire $\ket{\psi}$ dans la base $X$.
\end{enumerate}

\paragraph{a) Matrice de densité.} $\rho=\ket{\psi}\bra{\psi}$ est le produit d'un vecteur
colonne par un vecteur ligne (les amplitudes sont réelles) :
\[
  \rho=\begin{pmatrix}\frac{1}{\sqrt2}\\[0.4ex]\frac{1}{\sqrt2}\end{pmatrix}
       \begin{pmatrix}\frac{1}{\sqrt2}&\frac{1}{\sqrt2}\end{pmatrix}
  =\begin{pmatrix}\frac12&\frac12\\[0.4ex]\frac12&\frac12\end{pmatrix}.
\]

\paragraph{b) Pureté.}
\[
  \rho^2=\begin{pmatrix}\frac12&\frac12\\[0.4ex]\frac12&\frac12\end{pmatrix}
         \begin{pmatrix}\frac12&\frac12\\[0.4ex]\frac12&\frac12\end{pmatrix}
  =\begin{pmatrix}\frac14+\frac14&\frac14+\frac14\\[0.4ex]\frac14+\frac14&\frac14+\frac14\end{pmatrix}
  =\begin{pmatrix}\frac12&\frac12\\[0.4ex]\frac12&\frac12\end{pmatrix}=\rho,
\]
donc $\operatorname{Tr}\rho^2=\operatorname{Tr}\rho=\frac12+\frac12=1$ : le système est \emph{pur}.

\paragraph{c) Valeurs moyennes.} Avec $\sigma_x=X=\begin{pmatrix}0&1\\1&0\end{pmatrix}$ et
$\sigma_z=Z=\begin{pmatrix}1&0\\0&-1\end{pmatrix}$ (section~\ref{sec:op-dirac}), la valeur
moyenne est $\langle\sigma\rangle=\operatorname{Tr}(\rho\,\sigma)$ (section~\ref{sec:rho}).
\begin{align*}
  \rho\,\sigma_x&=\begin{pmatrix}\frac12&\frac12\\[0.4ex]\frac12&\frac12\end{pmatrix}
     \begin{pmatrix}0&1\\1&0\end{pmatrix}
     =\begin{pmatrix}\frac12&\frac12\\[0.4ex]\frac12&\frac12\end{pmatrix},
     &\langle\sigma_x\rangle&=\operatorname{Tr}(\rho\,\sigma_x)=\tfrac12+\tfrac12=1,\\[1.5ex]
  \rho\,\sigma_z&=\begin{pmatrix}\frac12&\frac12\\[0.4ex]\frac12&\frac12\end{pmatrix}
     \begin{pmatrix}1&0\\0&-1\end{pmatrix}
     =\begin{pmatrix}\frac12&-\frac12\\[0.4ex]\frac12&-\frac12\end{pmatrix},
     &\langle\sigma_z\rangle&=\operatorname{Tr}(\rho\,\sigma_z)=\tfrac12-\tfrac12=0.
\end{align*}
On retrouve $\langle\sigma_z\rangle$ avec les probabilités de mesure : dans la base $Z$,
$\sigma_z$ vaut $+1$ pour $\ket{0}$ et $-1$ pour $\ket{1}$, et
$|\braket{0}{\psi}|^2=|\braket{1}{\psi}|^2=\frac12$, donc
$\langle\sigma_z\rangle=\frac12\times(+1)+\frac12\times(-1)=0$.

\paragraph{d) Écriture dans la base $X$.} La base $X$ est $\{\ket{+},\ket{-}\}$, orthonormée, donc
$\ket{+}\bra{+}+\ket{-}\bra{-}=\mathrm{Id}$ et
$\ket{\psi}=\ket{+}\braket{+}{\psi}+\ket{-}\braket{-}{\psi}$, avec
\[
  \braket{+}{\psi}=\tfrac{1}{\sqrt2}\Big(\tfrac{1}{\sqrt2}+\tfrac{1}{\sqrt2}\Big)=1,
  \qquad
  \braket{-}{\psi}=\tfrac{1}{\sqrt2}\Big(\tfrac{1}{\sqrt2}-\tfrac{1}{\sqrt2}\Big)=0 .
\]
Donc $\ket{\psi}=1\cdot\ket{+}+0\cdot\ket{-}=\ket{+}$ : dans la base $X$, l'état n'a qu'une
seule composante. Comme $\sigma_x\ket{+}=+\ket{+}$ et $\sigma_x\ket{-}=-\ket{-}$, une
mesure de $\sigma_x$ donne $+1$ avec certitude, et
$\langle\sigma_x\rangle=1\cdot(+1)+0\cdot(-1)=1$ : c'est le résultat de la question c).
Le vecteur de Bloch de $\ket{+}$ est $(\langle\sigma_x\rangle,\langle\sigma_y\rangle,\langle\sigma_z\rangle)=(1,0,0)$,
de norme $1$ : l'état est sur la sphère, en accord avec $\operatorname{Tr}\rho^2=\frac{1+1}{2}=1$.

\subsubsection{Un système mélangé : un qubit de la paire de Bell}

À la section~\ref{sec:info}, la trace sur $B$ de l'état de Bell a donné
$\rho_A=\frac12\mathrm{Id}=\begin{pmatrix}\frac12&0\\0&\frac12\end{pmatrix}$. Alors
\[
  \rho_A^2=\begin{pmatrix}\frac14&0\\0&\frac14\end{pmatrix},
  \qquad
  \operatorname{Tr}\big(\rho_A^2\big)=\frac14+\frac14=\frac12<1 :
\]
le qubit $A$ est dans un état \emph{mélangé}, et c'est le plus mélangé possible pour un qubit
(centre de la boule de Bloch, $\vec r=\vec0$). Pourtant le couple $(A,B)$ est pur
($\operatorname{Tr}\rho_{AB}^2=1$) : on connaît parfaitement l'état du couple, mais pas celui d'un
qubit pris seul. L'information est entièrement dans les corrélations entre $A$ et $B$, et
$S(\rho_A)=1$ bit.

\subsubsection{Un système mélangé : le mélange statistique de Bob}

Bob prépare un qubit dans l'état $\ket{0}$ avec la probabilité $p_1=\frac34$ (75\,\%) et dans
l'état $\ket{1}$ avec la probabilité $p_2=\frac14$ (25\,\%), sans que l'on sache lequel des deux a
été préparé. Ce n'est pas un état pur $\ket{\psi}$ : on décrit le système par la matrice de
densité $\rho=\sum_ip_i\ket{\psi_i}\bra{\psi_i}$ (section~\ref{sec:pur-mixte}) :
\[
  \rho=\underbrace{\tfrac34}_{p_1}\ket{0}\bra{0}+\underbrace{\tfrac14}_{p_2}\ket{1}\bra{1}
  =\begin{pmatrix}\frac34&0\\0&\frac14\end{pmatrix}.
\]
Alors
\[
  \rho^2=\begin{pmatrix}\frac{9}{16}&0\\0&\frac{1}{16}\end{pmatrix},
  \qquad
  \operatorname{Tr}\big(\rho^2\big)=\frac{9}{16}+\frac{1}{16}=\frac{10}{16}=\frac58<1 :
\]
le système est \emph{mélangé}.

Les valeurs moyennes sont les moyennes pondérées sur la préparation :
$\langle\sigma_z\rangle=\operatorname{Tr}(\rho\,\sigma_z)=\frac34-\frac14=\frac12
=p_1\cdot(+1)+p_2\cdot(-1)$, et $\langle\sigma_x\rangle=\operatorname{Tr}(\rho\,\sigma_x)=0$
(les termes hors diagonale de $\rho$ sont nuls). Le vecteur de Bloch est $(0,0,\frac12)$,
à l'intérieur de la boule, et $\operatorname{Tr}\rho^2=\frac{1+1/4}{2}=\frac58$. Son entropie
est $S(\rho)=-\frac34\log_2\frac34-\frac14\log_2\frac14=2-\frac34\log_23\approx0{,}811$ bit.

\begin{remarque}[Mélange et superposition]
Considérons l'état pur $\ket{\psi'}=\frac{\sqrt3}{2}\ket{0}+\frac12\ket{1}$, qui a les mêmes
probabilités en base $Z$ ($P(0)=\frac34$, $P(1)=\frac14$). Sa matrice de densité est
\[
  \rho'=\begin{pmatrix}\frac34&\frac{\sqrt3}{4}\\[0.4ex]\frac{\sqrt3}{4}&\frac14\end{pmatrix},
  \qquad
  \operatorname{Tr}\rho'^2=\frac{9}{16}+\frac1{16}+2\cdot\frac{3}{16}=1,
  \qquad
  \langle\sigma_x\rangle=2\cdot\frac{\sqrt3}{4}=\frac{\sqrt3}{2}.
\]
Le mélange de Bob a la même diagonale mais aucun terme hors diagonale ($\operatorname{Tr}\rho^2=\frac58$,
$\langle\sigma_x\rangle=0$) : une superposition n'est pas un tirage au hasard entre $\ket{0}$ et
$\ket{1}$. De plus, pour un état pur il existe une mesure au résultat certain (ici $\sigma_x$ pour
$\ket{+}$) ; pour le mélange de Bob, la meilleure probabilité d'obtenir un même résultat est
$75\,\%$ (mesure en base $Z$, résultat $\ket{0}$).
\end{remarque}

\subsubsection{Bilan}

\begin{center}
\renewcommand{\arraystretch}{1.5}
\begin{tabular}{@{}lcccl@{}}
\toprule
\textbf{Système} & $\rho$ & $\operatorname{Tr}\rho^2$ & $S(\rho)$ & \textbf{Nature} \\
\midrule
$\ket{\psi}=\ket{+}$ & ${\renewcommand{\arraystretch}{1}\begin{pmatrix}\frac12&\frac12\\\frac12&\frac12\end{pmatrix}}$ & $1$ & $0$ & pur \\[2ex]
Qubit $A$ d'une paire de Bell & ${\renewcommand{\arraystretch}{1}\begin{pmatrix}\frac12&0\\0&\frac12\end{pmatrix}}$ & $\frac12$ & $1$ bit & mélangé \\[2ex]
Source de Bob (75\,\%/25\,\%) & ${\renewcommand{\arraystretch}{1}\begin{pmatrix}\frac34&0\\0&\frac14\end{pmatrix}}$ & $\frac58$ & $\approx0{,}811$ bit & mélangé \\
\bottomrule
\end{tabular}
\end{center}

% ---------------------------------------------------------------------
\subsection{Algorithme quantique : la téléportation}\label{sec:teleportation}
% ---------------------------------------------------------------------

\emph{Problème.} Alice détient un qubit $Q$ dans un état qu'elle ne connaît pas,
$\ket{\varphi}=\alpha\ket{0}+\beta\ket{1}$ avec $|\alpha|^2+|\beta|^2=1$, et elle veut que Bob
obtienne cet état sur l'un de ses qubits. Elle ne peut pas mesurer $Q$ pour connaître
$\alpha$ et $\beta$ (une mesure ne donne qu'un bit et détruit l'état), elle ne peut pas le
copier (théorème de non-clonage) et elle n'envoie pas le qubit lui-même.

\begin{definition}{Protocole de téléportation quantique}
\textbf{Ressources.}
\begin{itemize}[nosep]
  \item Un \emph{e-bit} : Alice et Bob partagent une paire de qubits $(A,B)$ dans l'état de Bell
        $\ket{\Phi^+}_{BA}=\frac{1}{\sqrt2}\big(\ket{00}+\ket{11}\big)$ ; Alice détient $A$,
        Bob détient $B$.
  \item Le qubit $Q$ à téléporter, chez Alice.
  \item Un canal classique (deux bits).
\end{itemize}
Le système est écrit $(B,A,Q)$ : $\ket{b,a,q}=\ket{b}_B\otimes\ket{a}_A\otimes\ket{q}_Q$.

\textbf{Étapes.}
\begin{enumerate}[nosep]
  \item Alice applique un CNOT de contrôle $Q$ et de cible $A$.
  \item Alice applique une porte $H$ sur $Q$.
  \item Alice mesure $A$ et $Q$ dans la base de calcul, obtient deux bits $(a,q)$ et les envoie
        à Bob.
  \item Bob applique $X$ si $a=1$, puis $Z$ si $q=1$, à son qubit $B$.
\end{enumerate}
\end{definition}

\begin{figure}[!ht]
\centering
\begin{adjustbox}{max width=\linewidth}
\begin{tikzpicture}[>=Stealth, font=\small]
  % fils quantiques : Q en haut, A au milieu, B en bas
  \draw (0,2.4) -- (5.3,2.4);
  \draw (0,1.2) -- (5.3,1.2);
  \draw (0,0) -- (10.6,0);
  \node[anchor=east] at (0,2.4) {$\ket{\varphi}_Q$};
  \node[anchor=south west, text=insagris, font=\footnotesize] at (0.05,1.2) {$A$};
  \node[anchor=south west, text=insagris, font=\footnotesize] at (0.05,0) {$B$};
  % paire de Bell partagée
  \draw[insagris, thick] (-0.1,1.5) -- (-0.25,1.5) -- (-0.25,-0.3) -- (-0.1,-0.3);
  \node[anchor=east] at (-0.35,0.6) {$\ket{\Phi^+}_{BA}$};
  % CNOT : contrôle Q, cible A
  \draw[insarouge, thick] (2.6,2.4) -- (2.6,1.2);
  \pic at (2.6,2.4) {ctl};
  \pic at (2.6,1.2) {cible};
  % Hadamard sur Q
  \node[porte] at (4.2,2.4) {$H$};
  % mesures d'Alice
  \pic at (5.6,2.4) {mesure};
  \pic at (5.6,1.2) {mesure};
  % fils classiques (doubles)
  \draw[double, double distance=1.2pt, insagris] (5.9,1.2) -- (7.6,1.2) -- (7.6,0.35);
  \draw[double, double distance=1.2pt, insagris] (5.9,2.4) -- (9.6,2.4) -- (9.6,0.35);
  \node[above, font=\footnotesize] at (6.5,1.2) {$a$};
  \node[above, font=\footnotesize] at (6.5,2.4) {$q$};
  % corrections de Bob
  \node[porte] at (7.6,0) {$X$};
  \node[porte] at (9.6,0) {$Z$};
  \node[anchor=west] at (10.6,0) {$\ket{\varphi}_B$};
  % états intermédiaires
  \foreach \x in {1.4,3.4,4.9}
    \draw[dashed, insagris!70] (\x,2.8) -- (\x,-0.5);
  \node[anchor=north, font=\footnotesize] at (1.4,-0.5) {$\ket{\Pi_0}$};
  \node[anchor=north, font=\footnotesize] at (3.4,-0.5) {$\ket{\Pi_1}$};
  \node[anchor=north, font=\footnotesize] at (4.9,-0.5) {$\ket{\Pi_2}$};
\end{tikzpicture}
\end{adjustbox}
\caption{Circuit de téléportation. Alice détient $Q$ (dans $\ket{\varphi}=\alpha\ket{0}+\beta\ket{1}$)
et $A$, Bob détient $B$ ; $A$ et $B$ forment la paire $\ket{\Phi^+}_{BA}$. Les traits doubles
sont des fils \emph{classiques} qui portent les résultats $a$ (mesure de $A$) et $q$ (mesure de
$Q$) ; $X$ est appliqué si $a=1$, $Z$ si $q=1$. Le ket s'écrit $\ket{b,a,q}$ : le fil du haut
($Q$) est le dernier chiffre.}
\label{fig:teleportation}
\end{figure}

\subsubsection{Calcul pas à pas}

\paragraph{État initial $\ket{\Pi_0}$.} Alice et Bob partagent la paire, et $Q$ est indépendant
d'elle :
\begin{align*}
  \ket{\Pi_0}&=\ket{\Phi^+}_{BA}\otimes\ket{\varphi}_Q
    =\frac{1}{\sqrt2}\big(\ket{00}+\ket{11}\big)\otimes\big(\alpha\ket{0}+\beta\ket{1}\big)\\
  &=\frac{1}{\sqrt2}\Big(\alpha\ket{00}\ket{0}+\beta\ket{00}\ket{1}+\alpha\ket{11}\ket{0}+\beta\ket{11}\ket{1}\Big)\\
  &=\frac{\alpha\ket{000}+\beta\ket{001}+\alpha\ket{110}+\beta\ket{111}}{\sqrt2}.
\end{align*}

\paragraph{Après le CNOT $Q\to A$ : $\ket{\Pi_1}$.} Le CNOT (contrôle $Q$, dernier chiffre ;
cible $A$, chiffre du milieu) agit ainsi : $\ket{b,a,q}\to\ket{b,\,a\oplus q,\,q}$. Terme par terme :
\begin{center}
\renewcommand{\arraystretch}{1.3}
\begin{tabular}{@{}cccc@{}}
\toprule
Terme & $q$ & Effet sur $A$ & Résultat \\
\midrule
$\alpha\ket{000}$ & $0$ & aucun & $\alpha\ket{000}$ \\
$\beta\ket{001}$ & $1$ & $0\to1$ & $\beta\ket{011}$ \\
$\alpha\ket{110}$ & $0$ & aucun & $\alpha\ket{110}$ \\
$\beta\ket{111}$ & $1$ & $1\to0$ & $\beta\ket{101}$ \\
\bottomrule
\end{tabular}
\end{center}
donc
\[
  \ket{\Pi_1}=\frac{\alpha\ket{000}+\beta\ket{011}+\alpha\ket{110}+\beta\ket{101}}{\sqrt2}.
\]
On isole ensuite le dernier chiffre, celui de $Q$, qui va changer à l'étape suivante
(le premier ket est celui du couple $\ket{b,a}$) :
\[
  \ket{\Pi_1}=\frac{\alpha\ket{00}\otimes\ket{0}+\beta\ket{01}\otimes\ket{1}
    +\alpha\ket{11}\otimes\ket{0}+\beta\ket{10}\otimes\ket{1}}{\sqrt2}.
\]

\paragraph{Après la porte $H$ sur $Q$ : $\ket{\Pi_2}$.} On remplace $\ket{0}$ par
$H\ket{0}=\ket{+}$ et $\ket{1}$ par $H\ket{1}=\ket{-}$ dans le facteur $Q$ :
\[
  \ket{\Pi_2}=\frac{1}{\sqrt2}\Big(\alpha\ket{00}\otimes\ket{+}+\beta\ket{01}\otimes\ket{-}
    +\alpha\ket{11}\otimes\ket{+}+\beta\ket{10}\otimes\ket{-}\Big).
\]
Avec $\ket{\pm}=\frac{1}{\sqrt2}(\ket{0}\pm\ket{1})$, le préfacteur devient
$\frac{1}{\sqrt2}\cdot\frac{1}{\sqrt2}=\frac12$ :
\begin{align*}
  \ket{\Pi_2}&=\frac12\Big(\alpha\ket{00}\big(\ket{0}+\ket{1}\big)+\beta\ket{01}\big(\ket{0}-\ket{1}\big)
     +\alpha\ket{11}\big(\ket{0}+\ket{1}\big)+\beta\ket{10}\big(\ket{0}-\ket{1}\big)\Big)\\
  &=\frac12\Big(\alpha\ket{000}+\alpha\ket{001}+\beta\ket{010}-\beta\ket{011}
     +\alpha\ket{110}+\alpha\ket{111}+\beta\ket{100}-\beta\ket{101}\Big).
\end{align*}

\paragraph{Regroupement selon le résultat d'Alice.} Alice mesurera $A$ et $Q$, donc les deux
derniers chiffres. On regroupe les huit termes qui ont les mêmes chiffres $(a,q)$, en mettant
le ket de Bob en premier :
\begin{center}
\renewcommand{\arraystretch}{1.3}
\begin{tabular}{@{}cll@{}}
\toprule
$(a,q)$ & Termes de $\ket{\Pi_2}$ (facteur $\frac12$ omis) & Ket de Bob \\
\midrule
$(0,0)$ & $\alpha\ket{000}+\beta\ket{100}$ & $(\alpha\ket{0}+\beta\ket{1})\ket{00}$ \\
$(0,1)$ & $\alpha\ket{001}-\beta\ket{101}$ & $(\alpha\ket{0}-\beta\ket{1})\ket{01}$ \\
$(1,0)$ & $\beta\ket{010}+\alpha\ket{110}$ & $(\beta\ket{0}+\alpha\ket{1})\ket{10}$ \\
$(1,1)$ & $-\beta\ket{011}+\alpha\ket{111}$ & $(-\beta\ket{0}+\alpha\ket{1})\ket{11}$ \\
\bottomrule
\end{tabular}
\end{center}
d'où
\[
  \ket{\Pi_2}=\frac12\Big[\big(\alpha\ket{0}+\beta\ket{1}\big)\ket{00}
    +\big(\alpha\ket{0}-\beta\ket{1}\big)\ket{01}
    +\big(\beta\ket{0}+\alpha\ket{1}\big)\ket{10}
    +\big(\alpha\ket{1}-\beta\ket{0}\big)\ket{11}\Big].
\]
Dans chaque terme, le ket de gauche est l'état de $B$ et le ket de droite $\ket{a,q}$ celui du
couple $(A,Q)$ d'Alice.

\subsubsection{Mesure d'Alice et corrections de Bob}

Alice mesure $A$ et $Q$. Comme à la section~\ref{sec:mesure2}, la probabilité d'un résultat
$(a,q)$ est le carré de la norme du terme correspondant, et l'état de Bob devient le ket
associé (déjà normé). Ici
\[
  P(a,q)=\Big\|\tfrac12\ket{\phi_{aq}}\Big\|^2=\tfrac14\big(|\alpha|^2+|\beta|^2\big)=\tfrac14
  \quad\text{pour les quatre résultats,}
\]
qui sont donc équiprobables et \emph{indépendants} de $\alpha$ et $\beta$ : les deux bits
envoyés ne révèlent rien sur $\ket{\varphi}$. Les kets de Bob sont
$\ket{\phi_{00}}=\alpha\ket{0}+\beta\ket{1}$, $\ket{\phi_{01}}=\alpha\ket{0}-\beta\ket{1}$,
$\ket{\phi_{10}}=\beta\ket{0}+\alpha\ket{1}$ et $\ket{\phi_{11}}=-\beta\ket{0}+\alpha\ket{1}$.

\begin{center}
\renewcommand{\arraystretch}{1.5}
\begin{tabular}{@{}cccll@{}}
\toprule
Résultat $(a,q)$ & Probabilité & État de Bob & Lien avec $\ket{\varphi}$ & Correction de Bob \\
\midrule
$(0,0)$ & $\frac14$ & $\alpha\ket{0}+\beta\ket{1}$ & $\ket{\varphi}$ & $I$ (rien) \\
$(0,1)$ & $\frac14$ & $\alpha\ket{0}-\beta\ket{1}$ & $Z\ket{\varphi}$ & $Z$ \\
$(1,0)$ & $\frac14$ & $\beta\ket{0}+\alpha\ket{1}$ & $X\ket{\varphi}$ & $X$ \\
$(1,1)$ & $\frac14$ & $-\beta\ket{0}+\alpha\ket{1}$ & $XZ\ket{\varphi}$ & $X$ puis $Z$ (opérateur $ZX$) \\
\bottomrule
\end{tabular}
\end{center}

\paragraph{Vérification des quatre corrections.} Avec $X\ket{0}=\ket{1}$, $X\ket{1}=\ket{0}$,
$Z\ket{0}=\ket{0}$ et $Z\ket{1}=-\ket{1}$ :
\begin{align*}
  (0,0):\quad & I\big(\alpha\ket{0}+\beta\ket{1}\big)=\alpha\ket{0}+\beta\ket{1},\\
  (0,1):\quad & Z\big(\alpha\ket{0}-\beta\ket{1}\big)=\alpha\ket{0}-\beta Z\ket{1}=\alpha\ket{0}+\beta\ket{1},\\
  (1,0):\quad & X\big(\beta\ket{0}+\alpha\ket{1}\big)=\beta\ket{1}+\alpha\ket{0}=\alpha\ket{0}+\beta\ket{1},\\
  (1,1):\quad & X\big(-\beta\ket{0}+\alpha\ket{1}\big)=-\beta\ket{1}+\alpha\ket{0}=\alpha\ket{0}-\beta\ket{1},
     \ \text{puis}\ Z\big(\alpha\ket{0}-\beta\ket{1}\big)=\alpha\ket{0}+\beta\ket{1}.
\end{align*}
Dans le dernier cas, avec les matrices,
$ZX=\begin{pmatrix}1&0\\0&-1\end{pmatrix}\begin{pmatrix}0&1\\1&0\end{pmatrix}
=\begin{pmatrix}0&1\\-1&0\end{pmatrix}$ et
$ZX\begin{pmatrix}-\beta\\ \alpha\end{pmatrix}=\begin{pmatrix}\alpha\\ \beta\end{pmatrix}$.
L'ordre compte : l'état de Bob est $XZ\ket{\varphi}$ ($Z\ket{\varphi}=\alpha\ket{0}-\beta\ket{1}$,
puis $X$ donne $\alpha\ket{1}-\beta\ket{0}$), et on l'annule avec
$(XZ)^{-1}=Z^{-1}X^{-1}=ZX$ : on applique d'abord $X$, puis $Z$, comme dans le circuit.
Dans les quatre cas, Bob retrouve exactement $\ket{\varphi}=\alpha\ket{0}+\beta\ket{1}$.

\subsubsection{Ce que Bob sait avant de recevoir les bits}

Tant que Bob ignore $(a,q)$, son qubit est dans le mélange des quatre kets $\ket{\phi_{aq}}$, chacun
avec la probabilité $\frac14$ :
\begin{align*}
  \rho_B&=\frac14\sum_{a,q}\ket{\phi_{aq}}\bra{\phi_{aq}}\\
  &=\frac14\Bigg[\begin{pmatrix}|\alpha|^2&\alpha\beta^*\\ \alpha^*\beta&|\beta|^2\end{pmatrix}
    +\begin{pmatrix}|\alpha|^2&-\alpha\beta^*\\ -\alpha^*\beta&|\beta|^2\end{pmatrix}
    +\begin{pmatrix}|\beta|^2&\alpha^*\beta\\ \alpha\beta^*&|\alpha|^2\end{pmatrix}
    +\begin{pmatrix}|\beta|^2&-\alpha^*\beta\\ -\alpha\beta^*&|\alpha|^2\end{pmatrix}\Bigg]\\
  &=\frac14\begin{pmatrix}2(|\alpha|^2+|\beta|^2)&0\\0&2(|\alpha|^2+|\beta|^2)\end{pmatrix}
   =\frac12\begin{pmatrix}1&0\\0&1\end{pmatrix}=\frac12\,\mathrm{Id}.
\end{align*}
Sans les deux bits, Bob a donc un qubit complètement mélangé, qui ne dépend pas de $\alpha$ ni de
$\beta$ : il n'a reçu aucune information sur $\ket{\varphi}$. C'est l'envoi classique de $(a,q)$
qui permet de la retrouver, donc rien ne voyage plus vite que la lumière.

\begin{remarque}[Bilan de la téléportation]
\begin{itemize}[nosep]
  \item Ressources : $1$ e-bit (la paire de Bell) et $2$ bits classiques pour transmettre
        \emph{un qubit} d'état quelconque, sans qu'Alice connaisse $\alpha$ ni $\beta$.
  \item Pas de clonage : après la mesure, $A$ et $Q$ sont dans l'état de base $\ket{a,q}$.
        Alice n'a plus $\ket{\varphi}$ : l'état a été \emph{déplacé}, pas copié. L'intrication
        entre $A$ et $B$ a été consommée.
  \item Les quatre corrections $I,Z,X,ZX$ sont des portes contrôlées par des bits classiques
        (section~\ref{sec:controlees}).
\end{itemize}
\end{remarque}

% ---------------------------------------------------------------------
\subsection{Portes contrôlées : contrôlé-\texorpdfstring{$U$}{U}, Fredkin et Toffoli}\label{sec:controlees}
% ---------------------------------------------------------------------

Le CNOT est le cas le plus simple d'une idée générale : appliquer une opération à un ou
plusieurs qubits \emph{cibles} seulement si un ou plusieurs qubits de \emph{contrôle} valent $1$.
Dans cette section, comme à la section~\ref{sec:op2}, le contrôle est le \emph{premier} chiffre
du ket. Les qubits sont notés $x$, $y$, $z$ (le qubit $x$ est le contrôle) et
$\ket{x,y,z}=\ket{x}\otimes\ket{y}\otimes\ket{z}$ ; d'après la convention de la section, le
contrôle est dessiné en bas (figure~\ref{fig:controlees}). Pour éviter la confusion avec les
matrices de Pauli, on écrit ici $\sigma_x=X$ pour la porte NOT.

\begin{figure}[!ht]
\centering
\begin{adjustbox}{max width=\linewidth}
\begin{tikzpicture}[>=Stealth, font=\small]
  % (a) contrôlé-U : contrôle x en bas, cible y en haut
  \draw (0,1.2) -- (3.4,1.2);
  \draw (0,0) -- (3.4,0);
  \node[anchor=east] at (0,1.2) {$y$};
  \node[anchor=east] at (0,0) {$x$};
  \draw[insarouge, thick] (1.7,0) -- (1.7,0.85);
  \pic at (1.7,0) {ctl};
  \node[porte] at (1.7,1.2) {$U$};
  \node[anchor=north, font=\footnotesize] at (1.7,-0.5) {(a) contrôlé-$U$};
  % (b) contrôlé-SWAP (Fredkin) : x contrôle, y et z échangés
  \draw (5,2.4) -- (8.4,2.4);
  \draw (5,1.2) -- (8.4,1.2);
  \draw (5,0) -- (8.4,0);
  \node[anchor=east] at (5,2.4) {$z$};
  \node[anchor=east] at (5,1.2) {$y$};
  \node[anchor=east] at (5,0) {$x$};
  \draw[insarouge, thick] (6.7,0) -- (6.7,2.4);
  \pic at (6.7,0) {ctl};
  \pic at (6.7,1.2) {croix};
  \pic at (6.7,2.4) {croix};
  \node[anchor=north, font=\footnotesize] at (6.7,-0.5) {(b) Fredkin (C-SWAP)};
  % (c) Toffoli (CCNOT) : x et y contrôlent, z cible
  \draw (10,2.4) -- (13.4,2.4);
  \draw (10,1.2) -- (13.4,1.2);
  \draw (10,0) -- (13.4,0);
  \node[anchor=east] at (10,2.4) {$z$};
  \node[anchor=east] at (10,1.2) {$y$};
  \node[anchor=east] at (10,0) {$x$};
  \draw[insarouge, thick] (11.7,0) -- (11.7,2.4);
  \pic at (11.7,0) {ctl};
  \pic at (11.7,1.2) {ctl};
  \pic at (11.7,2.4) {cible};
  \node[anchor=north, font=\footnotesize] at (11.7,-0.5) {(c) Toffoli (CCNOT)};
\end{tikzpicture}
\end{adjustbox}
\caption{Portes contrôlées. Le contrôle $x$ est le premier chiffre du ket, donc le fil du bas.
(a) Le qubit $y$ subit $U$ si $x=1$ ; pour $U=X$ on retrouve le CNOT. (b) Si $x=1$, les qubits
$y$ et $z$ sont échangés. (c) Si $x=y=1$, le qubit $z$ est inversé.}
\label{fig:controlees}
\end{figure}

\subsubsection{Contrôlé-\texorpdfstring{$U$}{U}}

\begin{definition}{Porte contrôlée-$U$}
Soit $U$ une matrice unitaire $2\times2$. Sur le système $(x,y)$, de contrôle $x$ et de cible $y$,
\[
  \mathrm{C}U=\ket{0}\bra{0}_x\otimes I_y+\ket{1}\bra{1}_x\otimes U_y .
\]
Le premier terme agit quand le contrôle est $\ket{0}$ : la cible ne change pas ($I$). Le second
agit quand le contrôle est $\ket{1}$ : la cible subit $U$.
\end{definition}

Comme $\ket{0}\bra{0}\ket{x}=\delta_{x0}\ket{0}$ et $\ket{1}\bra{1}\ket{x}=\delta_{x1}\ket{1}$,
\[
  \mathrm{C}U\big(\ket{x}\otimes\ket{y}\big)=\ket{x}\otimes U^x\ket{y},\qquad U^0=I,\ U^1=U .
\]
Par la règle $A\otimes B=(A_{ij}B)$, $\ket{0}\bra{0}\otimes I=\begin{pmatrix}I&0\\0&0\end{pmatrix}$
et $\ket{1}\bra{1}\otimes U=\begin{pmatrix}0&0\\0&U\end{pmatrix}$ (blocs $2\times2$), donc
\[
  \mathrm{C}U=\begin{pmatrix}I&0\\0&U\end{pmatrix}
  =\begin{pmatrix}1&0&0&0\\0&1&0&0\\0&0&u_{00}&u_{01}\\0&0&u_{10}&u_{11}\end{pmatrix}
  \qquad\text{pour }U=\begin{pmatrix}u_{00}&u_{01}\\u_{10}&u_{11}\end{pmatrix}.
\]
Cette matrice est unitaire car $I^\dagger I=I$ et $U^\dagger U=I$.

\paragraph{Le CNOT est un cas particulier.} Pour $U=\sigma_x$,
$\mathrm{C}\sigma_x=\ket{0}\bra{0}_x\otimes I_y+\ket{1}\bra{1}_x\otimes\sigma_x$ est le CNOT de
contrôle $x$ (premier chiffre), c'est-à-dire $\mathrm{CNOT}_{12}$ de la section~\ref{sec:cnot}.
On calcule l'action sur les quatre états de base, avec $\braket{0}{0}=\braket{1}{1}=1$,
$\braket{0}{1}=\braket{1}{0}=0$, $\sigma_x\ket{0}=\ket{1}$ et $\sigma_x\ket{1}=\ket{0}$ :
\begin{align*}
  \mathrm{C}\sigma_x\ket{00}&=\big[\ket{0}\bra{0}_x\otimes I_y\big]\ket{0}_x\ket{0}_y
      +\big[\ket{1}\bra{1}_x\otimes\sigma_x\big]\ket{0}_x\ket{0}_y\\
    &=\ket{0}\braket{0}{0}\otimes I\ket{0}+\ket{1}\braket{1}{0}\otimes\sigma_x\ket{0}
     =\ket{0}\otimes\ket{0}+0=\ket{00},\\[1.2ex]
  \mathrm{C}\sigma_x\ket{01}&=\ket{0}\braket{0}{0}\otimes I\ket{1}+\ket{1}\braket{1}{0}\otimes\sigma_x\ket{1}
     =\ket{0}\otimes\ket{1}+0=\ket{01},\\[1.2ex]
  \mathrm{C}\sigma_x\ket{10}&=\ket{0}\braket{0}{1}\otimes I\ket{0}+\ket{1}\braket{1}{1}\otimes\sigma_x\ket{0}
     =0+\ket{1}\otimes\ket{1}=\ket{11},\\[1.2ex]
  \mathrm{C}\sigma_x\ket{11}&=\ket{0}\braket{0}{1}\otimes I\ket{1}+\ket{1}\braket{1}{1}\otimes\sigma_x\ket{1}
     =0+\ket{1}\otimes\ket{0}=\ket{10}.
\end{align*}
On retrouve la table du CNOT : la cible n'est inversée que si le contrôle vaut $1$.

\paragraph{Un autre exemple : $U=Z$.} Ici
\[
  \mathrm{C}Z=\begin{pmatrix}I&0\\0&Z\end{pmatrix}=\mathrm{diag}(1,1,1,-1) :
\]
la porte multiplie $\ket{11}$ par $-1$ et ne touche pas les autres états de base. Elle est
symétrique entre les deux qubits. Elle se déduit du CNOT par
$\mathrm{C}Z=(I\otimes H)\,\mathrm{C}X\,(I\otimes H)$, où $\mathrm{C}X$ est le contrôlé-$X$.
En effet, $I\otimes H=\mathrm{diag}(H,H)$ (deux blocs $2\times2$ égaux à $H$), et le produit de
matrices diagonales par blocs se fait bloc par bloc :
\[
  (I\otimes H)\begin{pmatrix}I&0\\0&X\end{pmatrix}(I\otimes H)
  =\begin{pmatrix}H\,I\,H&0\\0&H\,X\,H\end{pmatrix}
  =\begin{pmatrix}I&0\\0&HXH\end{pmatrix},
\]
car $H^2=I$, et
\[
  HXH=\frac12\begin{pmatrix}1&1\\1&-1\end{pmatrix}\begin{pmatrix}0&1\\1&0\end{pmatrix}\begin{pmatrix}1&1\\1&-1\end{pmatrix}
  =\frac12\begin{pmatrix}1&1\\1&-1\end{pmatrix}\begin{pmatrix}1&-1\\1&1\end{pmatrix}
  =\frac12\begin{pmatrix}2&0\\0&-2\end{pmatrix}=Z
\]
(section~\ref{sec:hadamard} : $H$ échange les axes $x$ et $z$).

\subsubsection{Contrôlé-SWAP : la porte de Fredkin}

\begin{definition}{Porte de Fredkin (C-SWAP)}
Sur trois qubits $(x,y,z)$ (espace de dimension $2^3=8$), avec $x$ pour contrôle,
\[
  \mathrm{C\text{-}SWAP}=\ket{0}\bra{0}_x\otimes I_y\otimes I_z+\ket{1}\bra{1}_x\otimes\mathrm{SWAP}_{yz}.
\]
Si $x=0$, rien ne change ; si $x=1$, les qubits $y$ et $z$ sont échangés.
\end{definition}

Calcul sur quatre états de base (avec $I\ket{ab}=\ket{ab}$ et $\mathrm{SWAP}\ket{ab}=\ket{ba}$) :
\begin{align*}
  \mathrm{C\text{-}SWAP}\ket{000}&=\ket{0}\braket{0}{0}\otimes I\ket{00}+\ket{1}\braket{1}{0}\otimes\mathrm{SWAP}\ket{00}
     =\ket{0}\otimes\ket{00}+0=\ket{000},\\[1.2ex]
  \mathrm{C\text{-}SWAP}\ket{101}&=\ket{0}\braket{0}{1}\otimes I\ket{01}+\ket{1}\braket{1}{1}\otimes\mathrm{SWAP}\ket{01}
     =0+\ket{1}\otimes\ket{10}=\ket{110},\\[1.2ex]
  \mathrm{C\text{-}SWAP}\ket{110}&=\ket{0}\braket{0}{1}\otimes I\ket{10}+\ket{1}\braket{1}{1}\otimes\mathrm{SWAP}\ket{10}
     =0+\ket{1}\otimes\ket{01}=\ket{101},\\[1.2ex]
  \mathrm{C\text{-}SWAP}\ket{010}&=\ket{0}\braket{0}{0}\otimes I\ket{10}+\ket{1}\braket{1}{0}\otimes\mathrm{SWAP}\ket{10}
     =\ket{0}\otimes\ket{10}+0=\ket{010}.
\end{align*}
Les huit états de base donnent :
\begin{center}
\renewcommand{\arraystretch}{1.3}
\begin{tabular}{@{}ccccccccc@{}}
\toprule
Entrée & $\ket{000}$ & $\ket{001}$ & $\ket{010}$ & $\ket{011}$ & $\ket{100}$ & $\ket{101}$ & $\ket{110}$ & $\ket{111}$ \\
\midrule
Sortie & $\ket{000}$ & $\ket{001}$ & $\ket{010}$ & $\ket{011}$ & $\ket{100}$ & $\ket{110}$ & $\ket{101}$ & $\ket{111}$ \\
\bottomrule
\end{tabular}
\end{center}
Seuls $\ket{101}$ et $\ket{110}$ changent (contrôle à $1$ et $y\neq z$) ; pour $\ket{100}$ et
$\ket{111}$ l'échange de deux qubits égaux ne change rien. Dans la base
$\ket{000},\ket{001},\dots,\ket{111}$ (ordre binaire), la matrice est
$\mathrm{diag}(I_4,\mathrm{SWAP})$ (deux blocs $4\times4$) :
\[
  \mathrm{C\text{-}SWAP}=\begingroup\setlength{\arraycolsep}{3pt}\begin{pmatrix}
    1&0&0&0&0&0&0&0\\
    0&1&0&0&0&0&0&0\\
    0&0&1&0&0&0&0&0\\
    0&0&0&1&0&0&0&0\\
    0&0&0&0&1&0&0&0\\
    0&0&0&0&0&0&1&0\\
    0&0&0&0&0&1&0&0\\
    0&0&0&0&0&0&0&1
  \end{pmatrix}\endgroup .
\]
C'est l'identité dont les lignes $6$ et $7$ sont échangées.

\paragraph{Contrôle en superposition.} Avec $\ket{+}_x\otimes\ket{0}_y\otimes\ket{1}_z
=\frac{1}{\sqrt2}\big(\ket{001}+\ket{101}\big)$ :
\[
  \mathrm{C\text{-}SWAP}\,\frac{1}{\sqrt2}\big(\ket{001}+\ket{101}\big)
  =\frac{1}{\sqrt2}\big(\ket{001}+\ket{110}\big).
\]
Le résultat n'est pas un produit : les coefficients de $\ket{x}\otimes\ket{yz}$ forment la
matrice (lignes $x=0,1$, colonnes $yz=00,01,10,11$)
\[
  \frac{1}{\sqrt2}\begin{pmatrix}0&1&0&0\\0&0&1&0\end{pmatrix},
\]
dont les deux lignes ne sont pas proportionnelles. Le qubit $x$ est intriqué avec la paire
$(y,z)$.

\subsubsection{Contrôlé-contrôlé-NOT : la porte de Toffoli}

\begin{definition}{Porte de Toffoli (CCNOT)}
Sur trois qubits $(x,y,z)$, avec $x$ et $y$ pour contrôles et $z$ pour cible,
\[
  \mathrm{CCNOT}=\ket{0}\bra{0}_x\otimes I_y\otimes I_z
   +\ket{1}\bra{1}_x\otimes\Big[\ket{0}\bra{0}_y\otimes I_z+\ket{1}\bra{1}_y\otimes\sigma_x\Big].
\]
Le crochet est un CNOT (contrôle $y$, cible $z$), appliqué seulement si $x=1$ : si $x=0$, rien ;
si $x=1$ et $y=0$, rien ; si $x=1$ et $y=1$, le qubit $z$ est inversé.
\end{definition}

Calcul de $\mathrm{CCNOT}\ket{110}$ :
\begin{align*}
  \mathrm{CCNOT}\ket{110}&=\ket{0}\braket{0}{1}\otimes I\ket{1}\otimes I\ket{0}
     +\ket{1}\braket{1}{1}\otimes\Big[\ket{0}\braket{0}{1}\otimes I\ket{0}
       +\ket{1}\braket{1}{1}\otimes\sigma_x\ket{0}\Big]\\
  &=0+\ket{1}\otimes\Big[0+\ket{1}\otimes\ket{1}\Big]=\ket{111}.
\end{align*}
Pour les huit états de base, $\mathrm{CCNOT}\ket{x,y,z}=\ket{x,\,y,\,z\oplus xy}$ :
\begin{center}
\renewcommand{\arraystretch}{1.3}
\begin{tabular}{@{}ccccccccc@{}}
\toprule
Entrée & $\ket{000}$ & $\ket{001}$ & $\ket{010}$ & $\ket{011}$ & $\ket{100}$ & $\ket{101}$ & $\ket{110}$ & $\ket{111}$ \\
\midrule
Sortie & $\ket{000}$ & $\ket{001}$ & $\ket{010}$ & $\ket{011}$ & $\ket{100}$ & $\ket{101}$ & $\ket{111}$ & $\ket{110}$ \\
\bottomrule
\end{tabular}
\end{center}
Seuls $\ket{110}$ et $\ket{111}$ sont échangés : la matrice est $\mathrm{diag}(I_6,X)$
(un bloc $6\times6$ égal à l'identité, puis $X$),
\[
  \mathrm{CCNOT}=\begingroup\setlength{\arraycolsep}{3pt}\begin{pmatrix}
    1&0&0&0&0&0&0&0\\
    0&1&0&0&0&0&0&0\\
    0&0&1&0&0&0&0&0\\
    0&0&0&1&0&0&0&0\\
    0&0&0&0&1&0&0&0\\
    0&0&0&0&0&1&0&0\\
    0&0&0&0&0&0&0&1\\
    0&0&0&0&0&0&1&0
  \end{pmatrix}\endgroup .
\]

\begin{remarque}[Portes réversibles]
Les portes de Fredkin et de Toffoli sont des permutations des états de base qui échangent
seulement deux d'entre eux : elles sont leurs propres inverses ($\mathrm{C\text{-}SWAP}^2=\mathrm{CCNOT}^2=\mathrm{Id}$)
et unitaires. Avec $z=0$, la Toffoli calcule le ET logique $x\wedge y$ sans rien effacer :
$\ket{x,y,0}\to\ket{x,y,xy}$ ; avec $z=1$, elle calcule le NON-ET. Elle sert à écrire tout calcul
classique de façon réversible.
\end{remarque}

% ---------------------------------------------------------------------
\subsection{Récapitulatif du chapitre}
% ---------------------------------------------------------------------

\begin{center}
\renewcommand{\arraystretch}{1.5}
\begin{tabular}{@{}lp{12cm}@{}}
\toprule
\textbf{Notion} & \textbf{Formule ou résultat} \\
\midrule
CNOT & $\ket{b,a}\to\ket{b\oplus a,a}$ (contrôle $a$ écrit en dernier), $\mathrm{CNOT}^2=\mathrm{Id}$ \\
SWAP & $\mathrm{SWAP}\ket{a,b}=\ket{b,a}$, \quad $\mathrm{CNOT}_{21}=\mathrm{SWAP}\,\mathrm{CNOT}_{12}\,\mathrm{SWAP}$ \\
Circuit de Bell & $(I\otimes H)$ puis CNOT : $\ket{\beta^{xy}}=\frac{1}{\sqrt2}\big(\ket{y,0}+(-1)^x\ket{\bar y,1}\big)$, $\ket{\beta^{00}}=\ket{\Phi^+}$ \\
Information & $I(A;B)=H(A)+H(B)-H(A,B)$, \quad $H(A|B)=H(A,B)-H(B)$ \\
Trace partielle & $\operatorname{Tr}_B\big(\ket{a}\bra{a'}\otimes\ket{b}\bra{b'}\big)=\ket{a}\bra{a'}\braket{b'}{b}$ ; état de Bell : $\rho_A=\frac12\mathrm{Id}$ \\
État de Bell (bits) & $H(A,B)=0$, \quad $H(A)=H(B)=1$, \quad $I(A;B)=2$, \quad $H(A|B)=-1$ \\
Pur / mélangé & $\operatorname{Tr}\rho^2=1$ / $<1$ ; $\ket{+}$ : $1$ ; qubit d'une paire de Bell : $\frac12$ ; source de Bob : $\frac58$ \\
Téléportation & $1$ e-bit $+\ 2$ bits classiques ; correction : $X$ si $a=1$, puis $Z$ si $q=1$ ; $P(a,q)=\frac14$ \\
Contrôlé-$U$ & $\ket{0}\bra{0}\otimes I+\ket{1}\bra{1}\otimes U=\mathrm{diag}(I,U)$ ; $\ket{x,y}\to\ket{x}\otimes U^x\ket{y}$ \\
Fredkin & $\ket{0}\bra{0}\otimes I\otimes I+\ket{1}\bra{1}\otimes\mathrm{SWAP}$ ; $\ket{101}\leftrightarrow\ket{110}$ \\
Toffoli & $\ket{x,y,z}\to\ket{x,y,z\oplus xy}$ ; $\ket{110}\leftrightarrow\ket{111}$ \\
\bottomrule
\end{tabular}
\end{center}
.
%  Comme chapitre4.tex, chapitre5.tex et chapitre6.tex, il ouvre une NOUVELLE
%  SECTION (numéro 5). Pour le rattacher à la section précédente, remplacer
%  \section par \subsection et chaque \subsection par \subsubsection.
%  Pas de préambule : environnements (definition, resultat, remarque),
%  macros (\ii, \ee, \dd, \ket, \bra, \braket) et couleurs (insarouge,
%  insagris) sont définis dans main.tex. Les symboles de circuits (contrôle,
%  cible, croix, mesure) sont définis juste après avec \tikzset.
% =====================================================================

% Macros de Dirac (déjà définies dans main.tex récent ; ce bloc évite l'erreur
% « Undefined control sequence \ket » si main.tex est une ancienne version).
\providecommand{\ket}[1]{|#1\rangle}
\providecommand{\bra}[1]{\langle #1|}
\providecommand{\braket}[2]{\langle #1|#2\rangle}

% Symboles de circuits pour les figures de ce chapitre.
\tikzset{
  ctl/.pic   = {\fill[insarouge] (0,0) circle (2.6pt);},
  cible/.pic = {\draw[insarouge, thick, fill=white] (0,0) circle (0.28);
                \draw[insarouge, thick] (-0.28,0) -- (0.28,0);
                \draw[insarouge, thick] (0,-0.28) -- (0,0.28);},
  croix/.pic = {\draw[insarouge, thick] (-0.17,-0.17) -- (0.17,0.17);
                \draw[insarouge, thick] (-0.17,0.17) -- (0.17,-0.17);},
  mesure/.pic = {\draw[thick, fill=white] (-0.3,-0.27) rectangle (0.3,0.27);
                 \draw[thick] (-0.19,-0.13) arc[start angle=155, end angle=25, radius=0.21];
                 \draw[thick, -{Stealth[length=1.3mm]}] (0,-0.17) -- (0.14,0.12);},
  porte/.style = {draw=insarouge, fill=insarouge!8, minimum size=7mm, thick},
}

% =====================================================================
\section{Portes multi-qubits, circuits et téléportation}\label{sec:circuits}
% =====================================================================

La section précédente a décrit les états de deux qubits. On étudie ici les portes qui
agissent sur plusieurs qubits : le CNOT et le SWAP, puis les portes contrôlées à deux et à
trois qubits (contrôlé-$U$, Fredkin, Toffoli). On assemble ensuite des portes dans un
circuit qui fabrique un état de Bell, on mesure l'information contenue dans un tel état,
on distingue les états purs des états mélangés, et on démontre le protocole de
téléportation. Tous les calculs sont détaillés.

\begin{remarque}[Convention pour les circuits]
Dans un circuit, le temps va de gauche à droite et chaque qubit a un fil. Les fils sont
dessinés dans l'ordre \emph{inverse} de l'écriture des kets : le fil du \emph{haut} porte le
\emph{dernier} chiffre du ket, le fil du \emph{bas} le premier. Par exemple, si le qubit $A$
est dessiné au-dessus du qubit $B$, le ket s'écrit $\ket{b,a}=\ket{b}_B\otimes\ket{a}_A$.
La figure~\ref{fig:bell} de la section précédente suit la même convention.
\end{remarque}

% ---------------------------------------------------------------------
\subsection{Porte CNOT}\label{sec:cnot}
% ---------------------------------------------------------------------

\begin{definition}{Porte CNOT (NOT contrôlé)}
La porte CNOT agit sur deux qubits : le qubit $A$ est le \emph{contrôle}, le qubit $B$ est la
\emph{cible}. La cible est inversée (porte $X$) si le contrôle vaut $1$ et reste inchangée
s'il vaut $0$. Avec le système écrit $(B,A)$, soit $\ket{b,a}=\ket{b}_B\otimes\ket{a}_A$ :
\[
  \mathrm{CNOT}\,\ket{b,a}=\ket{b\oplus a,\,a},
  \qquad
  \mathrm{CNOT}=I\otimes\ket{0}\bra{0}+X\otimes\ket{1}\bra{1},
\]
où $\oplus$ est l'addition modulo $2$ ($0\oplus0=1\oplus1=0$ et $0\oplus1=1\oplus0=1$).
\end{definition}

Sur le circuit de la figure~\ref{fig:cnot}, le point noir posé sur le fil de $A$ marque le
contrôle et le symbole $\oplus$ posé sur le fil de $B$ marque la cible.

\begin{figure}[!ht]
\centering
\begin{tikzpicture}[>=Stealth, font=\small]
  \draw (0,1.2) -- (6,1.2);
  \draw (0,0) -- (6,0);
  \node[anchor=east] at (0,1.2) {$\ket{a}_A$};
  \node[anchor=east] at (0,0) {$\ket{b}_B$};
  \node[anchor=west] at (6,1.2) {$\ket{a}_A$};
  \node[anchor=west] at (6,0) {$\ket{b\oplus a}_B$};
  \draw[insarouge, thick] (3,1.2) -- (3,0);
  \pic at (3,1.2) {ctl};
  \pic at (3,0) {cible};
  \node[anchor=south, text=insagris] at (3,1.5) {qubit de contrôle};
  \node[anchor=north, text=insagris] at (3,-0.5) {qubit cible};
\end{tikzpicture}
\caption{Porte CNOT : le fil du haut ($A$) est le contrôle, le fil du bas ($B$) est la
cible. Le ket s'écrit $\ket{b,a}$ (le fil du haut donne le dernier chiffre).}
\label{fig:cnot}
\end{figure}

\subsubsection{Table de vérité}

\begin{center}
\renewcommand{\arraystretch}{1.3}
\begin{tabular}{@{}cccc@{}}
\toprule
Entrée $\ket{b,a}$ & Contrôle $a$ & Sortie $\ket{b\oplus a,a}$ & Effet sur la cible $B$ \\
\midrule
$\ket{00}$ & $0$ & $\ket{00}$ & aucun \\
$\ket{01}$ & $1$ & $\ket{11}$ & $0\to1$ \\
$\ket{10}$ & $0$ & $\ket{10}$ & aucun \\
$\ket{11}$ & $1$ & $\ket{01}$ & $1\to0$ \\
\bottomrule
\end{tabular}
\end{center}

\subsubsection{Matrice du CNOT}

La $j$-ième colonne d'une matrice $4\times4$ est l'image du $j$-ième vecteur de base
($\ket{00},\ket{01},\ket{10},\ket{11}$, dans cet ordre). La table donne :
\begin{align*}
  \mathrm{CNOT}\ket{00}&=\ket{00}=\begin{pmatrix}1\\0\\0\\0\end{pmatrix},&
  \mathrm{CNOT}\ket{01}&=\ket{11}=\begin{pmatrix}0\\0\\0\\1\end{pmatrix},\\[1.5ex]
  \mathrm{CNOT}\ket{10}&=\ket{10}=\begin{pmatrix}0\\0\\1\\0\end{pmatrix},&
  \mathrm{CNOT}\ket{11}&=\ket{01}=\begin{pmatrix}0\\1\\0\\0\end{pmatrix},
\end{align*}
d'où, en juxtaposant ces quatre colonnes,
\[
  \mathrm{CNOT}=\begin{pmatrix}
    1 & 0 & 0 & 0\\ 0 & 0 & 0 & 1\\ 0 & 0 & 1 & 0\\ 0 & 1 & 0 & 0
  \end{pmatrix}.
\]

On retrouve cette matrice avec la formule $I\otimes\ket{0}\bra{0}+X\otimes\ket{1}\bra{1}$ et
la règle $A\otimes B=(A_{ij}B)$ (section~\ref{sec:op2}), où
$\ket{0}\bra{0}=\begin{pmatrix}1&0\\0&0\end{pmatrix}$ et
$\ket{1}\bra{1}=\begin{pmatrix}0&0\\0&1\end{pmatrix}$ :
\[
  I\otimes\ket{0}\bra{0}=\begin{pmatrix}1&0&0&0\\0&0&0&0\\0&0&1&0\\0&0&0&0\end{pmatrix},
  \qquad
  X\otimes\ket{1}\bra{1}=\begin{pmatrix}0&0&0&0\\0&0&0&1\\0&0&0&0\\0&1&0&0\end{pmatrix},
\]
et la somme de ces deux matrices est la matrice du CNOT.

\subsubsection{Action sur un état quelconque}

Pour $\ket{\Psi}=c_{00}\ket{00}+c_{01}\ket{01}+c_{10}\ket{10}+c_{11}\ket{11}$ :
\[
  \mathrm{CNOT}\begin{pmatrix}c_{00}\\c_{01}\\c_{10}\\c_{11}\end{pmatrix}
  =\begin{pmatrix}
    1\cdot c_{00}+0\cdot c_{01}+0\cdot c_{10}+0\cdot c_{11}\\
    0\cdot c_{00}+0\cdot c_{01}+0\cdot c_{10}+1\cdot c_{11}\\
    0\cdot c_{00}+0\cdot c_{01}+1\cdot c_{10}+0\cdot c_{11}\\
    0\cdot c_{00}+1\cdot c_{01}+0\cdot c_{10}+0\cdot c_{11}
  \end{pmatrix}
  =\begin{pmatrix}c_{00}\\c_{11}\\c_{10}\\c_{01}\end{pmatrix}.
\]
Le CNOT échange les amplitudes de $\ket{01}$ et de $\ket{11}$, c'est-à-dire des deux états
dont le contrôle vaut $1$, et laisse les autres.

\paragraph{Exemple : contrôle en superposition.} Soit $\ket{0}_B\otimes\ket{+}_A
=\frac{1}{\sqrt2}\big(\ket{00}+\ket{01}\big)=\frac{1}{\sqrt2}(1,1,0,0)^{\mathsf T}$,
un état séparable. Alors
\[
  \mathrm{CNOT}\,\frac{1}{\sqrt2}\begin{pmatrix}1\\1\\0\\0\end{pmatrix}
  =\frac{1}{\sqrt2}\begin{pmatrix}1\\0\\0\\1\end{pmatrix}
  =\frac{1}{\sqrt2}\big(\ket{00}+\ket{11}\big)=\ket{\Phi^+}.
\]
Le résultat est intriqué ($D=c_{00}c_{11}-c_{01}c_{10}=\frac12\neq0$ ; section~\ref{sec:separable}) :
un CNOT dont le contrôle est en superposition crée de l'intrication.

\paragraph{Réversibilité.} Deux CNOT successifs redonnent l'état initial :
$\ket{b,a}\to\ket{b\oplus a,a}\to\ket{b\oplus a\oplus a,a}=\ket{b,a}$, donc
$\mathrm{CNOT}^2=\mathrm{Id}$. La matrice est réelle et symétrique, donc
$\mathrm{CNOT}^\dagger\mathrm{CNOT}=\mathrm{CNOT}^2=\mathrm{Id}$ : elle est bien unitaire.

\begin{remarque}[Contrôle en dernier ou en premier]
La matrice dépend de l'ordre d'écriture du ket. Notons $\mathrm{CNOT}_{21}$ le CNOT dont le
contrôle est le \emph{second} chiffre du ket (c'est la matrice ci-dessus) et
$\mathrm{CNOT}_{12}$ celui dont le contrôle est le \emph{premier} chiffre, utilisé à la
section précédente :
\[
  \mathrm{CNOT}_{12}=\ket{0}\bra{0}\otimes I+\ket{1}\bra{1}\otimes X
  =\begin{pmatrix}1&0&0&0\\0&1&0&0\\0&0&0&1\\0&0&1&0\end{pmatrix}.
\]
Un même circuit (contrôle $A$, cible $B$) a donc deux matrices, selon que l'on écrit
$A$ en dernier ($\mathrm{CNOT}_{21}$) ou en premier ($\mathrm{CNOT}_{12}$). Les deux sont
reliées par le SWAP (section suivante).
\end{remarque}

% ---------------------------------------------------------------------
\subsection{Porte SWAP}\label{sec:swap}
% ---------------------------------------------------------------------

\begin{definition}{Porte SWAP (échange)}
La porte SWAP échange les états des deux qubits :
\[
  \mathrm{SWAP}\big(\ket{\psi}\otimes\ket{\phi}\big)=\ket{\phi}\otimes\ket{\psi}.
\]
Sur les états de base, $\mathrm{SWAP}\ket{a,b}=\ket{b,a}$ ; on prolonge par linéarité.
\end{definition}

Son symbole est formé de deux croix reliées par un trait vertical
(figure~\ref{fig:swap}).

\begin{figure}[!ht]
\centering
\begin{tikzpicture}[>=Stealth, font=\small]
  \draw (0,1.2) -- (6,1.2);
  \draw (0,0) -- (6,0);
  \node[anchor=east] at (0,1.2) {$\ket{a}$};
  \node[anchor=east] at (0,0) {$\ket{b}$};
  \node[anchor=west] at (6,1.2) {$\ket{b}$};
  \node[anchor=west] at (6,0) {$\ket{a}$};
  \draw[insarouge, thick] (3,1.2) -- (3,0);
  \pic at (3,1.2) {croix};
  \pic at (3,0) {croix};
\end{tikzpicture}
\caption{Porte SWAP : les deux qubits échangent leurs états.}
\label{fig:swap}
\end{figure}

\subsubsection{Table et matrice}

\begin{center}
\renewcommand{\arraystretch}{1.3}
\begin{tabular}{@{}cc@{}}
\toprule
Entrée & Sortie \\
\midrule
$\ket{00}$ & $\ket{00}$ \\
$\ket{01}$ & $\ket{10}$ \\
$\ket{10}$ & $\ket{01}$ \\
$\ket{11}$ & $\ket{11}$ \\
\bottomrule
\end{tabular}
\end{center}

Les images des quatre vecteurs de base sont les colonnes de la matrice :
\[
  \mathrm{SWAP}\ket{00}=\begin{pmatrix}1\\0\\0\\0\end{pmatrix},\quad
  \mathrm{SWAP}\ket{01}=\begin{pmatrix}0\\0\\1\\0\end{pmatrix},\quad
  \mathrm{SWAP}\ket{10}=\begin{pmatrix}0\\1\\0\\0\end{pmatrix},\quad
  \mathrm{SWAP}\ket{11}=\begin{pmatrix}0\\0\\0\\1\end{pmatrix},
\]
\[
  \mathrm{SWAP}=\begin{pmatrix}
    1 & 0 & 0 & 0\\ 0 & 0 & 1 & 0\\ 0 & 1 & 0 & 0\\ 0 & 0 & 0 & 1
  \end{pmatrix}.
\]
Appliquée à $\ket{01}=(0,1,0,0)^{\mathsf T}$, cette matrice sélectionne sa deuxième colonne :
\[
  \begin{pmatrix}
    1 & 0 & 0 & 0\\ 0 & 0 & 1 & 0\\ 0 & 1 & 0 & 0\\ 0 & 0 & 0 & 1
  \end{pmatrix}
  \begin{pmatrix}0\\1\\0\\0\end{pmatrix}
  =\begin{pmatrix}0\\0\\1\\0\end{pmatrix}=\ket{10}.
\]
Sur un état quelconque, $\mathrm{SWAP}\,(c_{00},c_{01},c_{10},c_{11})^{\mathsf T}
=(c_{00},c_{10},c_{01},c_{11})^{\mathsf T}$ : les amplitudes de $\ket{01}$ et $\ket{10}$
sont échangées.

\subsubsection{Accord avec la définition}

Pour deux états $\ket{\psi}=(a_0,a_1)^{\mathsf T}$ et $\ket{\phi}=(b_0,b_1)^{\mathsf T}$ :
\[
  \mathrm{SWAP}\,(\ket{\psi}\otimes\ket{\phi})
  =\mathrm{SWAP}\begin{pmatrix}a_0b_0\\a_0b_1\\a_1b_0\\a_1b_1\end{pmatrix}
  =\begin{pmatrix}a_0b_0\\a_1b_0\\a_0b_1\\a_1b_1\end{pmatrix}
  =\begin{pmatrix}b_0\\b_1\end{pmatrix}\otimes\begin{pmatrix}a_0\\a_1\end{pmatrix}
  =\ket{\phi}\otimes\ket{\psi},
\]
ce qui est bien la définition. La porte s'écrit aussi avec des projecteurs :
\[
  \mathrm{SWAP}=\sum_{a,b\in\{0,1\}}\ket{a}\bra{b}\otimes\ket{b}\bra{a}.
\]
En effet, pour $\ket{c}\otimes\ket{d}$ on trouve
$\sum_{a,b}\ket{a}\braket{b}{c}\otimes\ket{b}\braket{a}{d}$, et seul le terme $b=c$, $a=d$
est non nul : le résultat est $\ket{d}\otimes\ket{c}$.

\paragraph{Exemples.}
\begin{itemize}[nosep]
  \item Deux états différents : $\mathrm{SWAP}\big(\ket{0}\otimes\ket{+}\big)
        =\mathrm{SWAP}\,\frac{1}{\sqrt2}(1,1,0,0)^{\mathsf T}=\frac{1}{\sqrt2}(1,0,1,0)^{\mathsf T}
        =\ket{+}\otimes\ket{0}$.
  \item L'état $\ket{\psi}=\frac{1}{\sqrt2}\ket{00}+\frac{1}{\sqrt2}\ket{11}$ ne change pas :
        $\mathrm{SWAP}\ket{\psi}=\frac{1}{\sqrt2}\ket{00}+\frac{1}{\sqrt2}\ket{11}=\ket{\psi}$,
        car $\ket{00}$ et $\ket{11}$ sont invariants. De même
        $\mathrm{SWAP}\ket{\Phi^-}=\ket{\Phi^-}$ et $\mathrm{SWAP}\ket{\Psi^+}=\ket{\Psi^+}$, alors que
        $\mathrm{SWAP}\ket{\Psi^-}=\frac{1}{\sqrt2}(\ket{10}-\ket{01})=-\ket{\Psi^-}$.
\end{itemize}

\paragraph{Propriétés.} Échanger deux fois redonne l'état de départ :
$\mathrm{SWAP}^2=\mathrm{Id}$ ; la matrice est réelle et symétrique, donc
$\mathrm{SWAP}^\dagger\mathrm{SWAP}=\mathrm{SWAP}^2=\mathrm{Id}$ : elle est unitaire.

\subsubsection{Lien avec le CNOT}

\begin{resultat}{SWAP échange les rôles de contrôle et de cible}
\[
  \mathrm{CNOT}_{21}=\mathrm{SWAP}\;\mathrm{CNOT}_{12}\;\mathrm{SWAP}.
\]
\end{resultat}

\emph{Sur les états de base.}
$\mathrm{SWAP}\,\mathrm{CNOT}_{12}\,\mathrm{SWAP}\ket{b,a}
=\mathrm{SWAP}\,\mathrm{CNOT}_{12}\ket{a,b}
=\mathrm{SWAP}\ket{a,\,a\oplus b}=\ket{a\oplus b,\,a}=\mathrm{CNOT}_{21}\ket{b,a}$.

\emph{Par les matrices.} Multiplier à gauche par $\mathrm{SWAP}$ échange les lignes $2$ et
$3$, multiplier à droite l'échange les colonnes $2$ et $3$ :
\[
  \mathrm{SWAP}\,\mathrm{CNOT}_{12}
  =\begin{pmatrix}1&0&0&0\\0&0&0&1\\0&1&0&0\\0&0&1&0\end{pmatrix},
  \qquad
  \mathrm{SWAP}\,\mathrm{CNOT}_{12}\,\mathrm{SWAP}
  =\begin{pmatrix}1&0&0&0\\0&0&0&1\\0&0&1&0\\0&1&0&0\end{pmatrix}=\mathrm{CNOT}_{21}.
\]

\begin{remarque}[SWAP avec trois CNOT]
$\mathrm{SWAP}=\mathrm{CNOT}_{12}\,\mathrm{CNOT}_{21}\,\mathrm{CNOT}_{12}$. En effet, sur
$\ket{a,b}$ (on applique les portes de droite à gauche) :
\[
  \ket{a,b}\ \xrightarrow{\ \mathrm{CNOT}_{12}\ }\ \ket{a,\,a\oplus b}
  \ \xrightarrow{\ \mathrm{CNOT}_{21}\ }\ \ket{a\oplus(a\oplus b),\,a\oplus b}=\ket{b,\,a\oplus b}
  \ \xrightarrow{\ \mathrm{CNOT}_{12}\ }\ \ket{b,\,b\oplus(a\oplus b)}=\ket{b,a}.
\]
Les deux membres coïncident sur les états de base, donc sur tous les états par linéarité.
\end{remarque}

% ---------------------------------------------------------------------
\subsection{Un circuit quantique : le générateur d'états de Bell}\label{sec:circuit-bell}
% ---------------------------------------------------------------------

On assemble une porte de Hadamard et un CNOT. Le système est $(A,B)$ avec le ket
$\ket{a,b}=\ket{a}_A\otimes\ket{b}_B$ ; d'après la convention de la section, $B$ est dessiné
en haut et $A$ en bas (figure~\ref{fig:circuit-bell}). Le circuit a deux étapes :
\begin{enumerate}[nosep]
  \item la porte $I\otimes H$ : identité sur $A$, Hadamard sur $B$ ;
  \item un CNOT de contrôle $B$ et de cible $A$ (la matrice $\mathrm{CNOT}_{21}$ de la
        section~\ref{sec:cnot}, puisque le contrôle est le second chiffre).
\end{enumerate}
Les traits pointillés délimitent les états $\ket{\Pi_0}$, $\ket{\Pi_1}$ et $\ket{\Pi_2}$ :
avant les portes, après $I\otimes H$, après le CNOT.

\begin{figure}[!ht]
\centering
\begin{tikzpicture}[>=Stealth, font=\small]
  % fils : B en haut, A en bas
  \draw (0,1.4) -- (7.6,1.4);
  \draw (0,0) -- (7.6,0);
  \node[anchor=east] at (0,1.4) {$\ket{0}_B$};
  \node[anchor=east] at (0,0) {$\ket{0}_A$};
  % première étape : H sur B, identité (pointillée) sur A
  \node[porte] at (2.2,1.4) {$H$};
  \node[draw=insagris, dashed, fill=white, minimum size=7mm] at (2.2,0) {$I$};
  % CNOT : contrôle B (haut), cible A (bas)
  \draw[insarouge, thick] (5,1.4) -- (5,0);
  \pic at (5,1.4) {ctl};
  \pic at (5,0) {cible};
  % accolade de sortie
  \draw[insagris, thick] (7.75,1.6) -- (7.9,1.6) -- (7.9,-0.2) -- (7.75,-0.2);
  \node[anchor=west] at (7.9,0.7) {$\ket{\beta^{00}}_{AB}$};
  % états intermédiaires
  \foreach \x in {0.7,3.6,6.4}
    \draw[dashed, insagris!70] (\x,1.9) -- (\x,-0.7);
  \node[anchor=north, font=\footnotesize] at (0.7,-0.7) {$\ket{\Pi_0}=\ket{00}$};
  \node[anchor=north, font=\footnotesize] at (3.6,-0.7) {$\ket{\Pi_1}=\ket{0}\otimes\ket{+}$};
  \node[anchor=north, font=\footnotesize] at (6.4,-0.7) {$\ket{\Pi_2}=\ket{\beta^{00}}$};
\end{tikzpicture}
\caption{Générateur d'états de Bell : $I\otimes H$, puis un CNOT de contrôle $B$ (fil du haut)
et de cible $A$ (fil du bas). L'entrée représentée est $\ket{00}$.}
\label{fig:circuit-bell}
\end{figure}

\subsubsection{Entrée \texorpdfstring{$\ket{00}$}{|00>} : calcul pas à pas}

\begin{align*}
  \ket{\Pi_0}&=\ket{0}_A\otimes\ket{0}_B=\ket{00},\\[1ex]
  \ket{\Pi_1}&=(I\otimes H)\ket{\Pi_0}=I\ket{0}_A\otimes H\ket{0}_B=\ket{0}\otimes\ket{+}
    =\frac{\ket{00}+\ket{01}}{\sqrt2},\\[1ex]
  \ket{\Pi_2}&=\mathrm{CNOT}\ket{\Pi_1}=\frac{\mathrm{CNOT}\ket{00}+\mathrm{CNOT}\ket{01}}{\sqrt2}
    =\frac{\ket{00}+\ket{11}}{\sqrt2}=\ket{\beta^{00}}=\ket{\Phi^+}.
\end{align*}
Dans $\ket{\Pi_1}$, le premier chiffre est celui de $A$ et le second celui de $B$. Le CNOT a
pour contrôle $B$ : $\ket{00}$ ($B=0$) ne change pas, $\ket{01}$ ($B=1$) devient $\ket{11}$
puisque $A$ est inversé. Le résultat est un \emph{état de Bell}.

Avec les matrices, la règle $A\otimes B=(A_{ij}B)$ donne
\[
  I\otimes H=\begin{pmatrix}1\cdot H&0\\0&1\cdot H\end{pmatrix}
  =\frac{1}{\sqrt2}\begin{pmatrix}1&1&0&0\\1&-1&0&0\\0&0&1&1\\0&0&1&-1\end{pmatrix}.
\]
Sa première colonne est $\ket{\Pi_1}=\frac{1}{\sqrt2}(1,1,0,0)^{\mathsf T}$ (image de $\ket{00}$).
Le CNOT échange les amplitudes de $\ket{01}$ et de $\ket{11}$ (section~\ref{sec:cnot}), donc
\[
  \ket{\Pi_2}=\mathrm{CNOT}\,\frac{1}{\sqrt2}\begin{pmatrix}1\\1\\0\\0\end{pmatrix}
  =\frac{1}{\sqrt2}\begin{pmatrix}1\\0\\0\\1\end{pmatrix}.
\]

\subsubsection{Les quatre entrées}

Prenons $B$ dans l'état $\ket{x}$ et $A$ dans l'état $\ket{y}$, avec $x,y\in\{0,1\}$. Alors
$\ket{\Pi_0}=\ket{y}_A\otimes\ket{x}_B=\ket{y,x}$. Comme
$H\ket{x}=\frac{1}{\sqrt2}\big(\ket{0}+(-1)^x\ket{1}\big)$ (c'est $\ket{+}$ pour $x=0$ et
$\ket{-}$ pour $x=1$),
\[
  \ket{\Pi_1}=\ket{y}\otimes H\ket{x}=\frac{1}{\sqrt2}\Big(\ket{y,0}+(-1)^x\ket{y,1}\Big).
\]
Le CNOT (contrôle $B$) ne change pas $\ket{y,0}$ et transforme $\ket{y,1}$ en
$\ket{\bar y,1}$, avec $\bar y=1-y$ :
\[
  \ket{\Pi_2}=\ket{\beta^{xy}}=\frac{1}{\sqrt2}\Big(\ket{y,0}+(-1)^x\ket{\bar y,1}\Big).
\]
Le premier indice de $\ket{\beta^{xy}}$ est l'entrée de $B$ (le qubit qui reçoit $H$), le
second celle de $A$ (la cible). Les quatre cas :

\begin{center}
\renewcommand{\arraystretch}{1.7}
\begin{tabular}{@{}ccccc@{}}
\toprule
$(B,A)=(x,y)$ & $\ket{\Pi_0}=\ket{y,x}$ & $\ket{\Pi_1}$ & $\ket{\Pi_2}=\ket{\beta^{xy}}$ & Nom \\
\midrule
$(0,0)$ & $\ket{00}$ & $\frac{1}{\sqrt2}(\ket{00}+\ket{01})$ & $\frac{1}{\sqrt2}(\ket{00}+\ket{11})$ & $\ket{\Phi^+}$ \\
$(1,0)$ & $\ket{01}$ & $\frac{1}{\sqrt2}(\ket{00}-\ket{01})$ & $\frac{1}{\sqrt2}(\ket{00}-\ket{11})$ & $\ket{\Phi^-}$ \\
$(0,1)$ & $\ket{10}$ & $\frac{1}{\sqrt2}(\ket{10}+\ket{11})$ & $\frac{1}{\sqrt2}(\ket{10}+\ket{01})$ & $\ket{\Psi^+}$ \\
$(1,1)$ & $\ket{11}$ & $\frac{1}{\sqrt2}(\ket{10}-\ket{11})$ & $\frac{1}{\sqrt2}(\ket{10}-\ket{01})$ & $-\ket{\Psi^-}$ \\
\bottomrule
\end{tabular}
\end{center}
Détail de la dernière ligne : $\ket{\Pi_1}=\ket{1}\otimes\ket{-}=\frac{1}{\sqrt2}(\ket{10}-\ket{11})$ ;
le CNOT laisse $\ket{10}$ et envoie $\ket{11}$ sur $\ket{01}$, d'où
$\frac{1}{\sqrt2}(\ket{10}-\ket{01})=-\frac{1}{\sqrt2}(\ket{01}-\ket{10})=-\ket{\Psi^-}$.

\paragraph{Forme matricielle.} $\mathrm{CNOT}_{21}$ échange les lignes $2$ et $4$ de
$I\otimes H$ (il échange les amplitudes de $\ket{01}$ et $\ket{11}$), donc
\[
  \mathrm{CNOT}_{21}\,(I\otimes H)=\frac{1}{\sqrt2}\begin{pmatrix}
    1 & 1 & 0 & 0\\ 0 & 0 & 1 & -1\\ 0 & 0 & 1 & 1\\ 1 & -1 & 0 & 0
  \end{pmatrix}.
\]
Ses colonnes sont les images de $\ket{00},\ket{01},\ket{10},\ket{11}$, c'est-à-dire de
$(B,A)=(0,0),(1,0),(0,1),(1,1)$ : ce sont $\ket{\beta^{00}}$, $\ket{\beta^{10}}$,
$\ket{\beta^{01}}$, $\ket{\beta^{11}}$, soit $\ket{\Phi^+}$, $\ket{\Phi^-}$, $\ket{\Psi^+}$
et $-\ket{\Psi^-}$.

\begin{remarque}[Comparaison avec la section précédente]
Le circuit de la figure~\ref{fig:bell} applique $H$ au \emph{premier} chiffre et prend celui-ci
pour contrôle ; ici $H$ et le contrôle portent sur le \emph{dernier} chiffre. Les deux
circuits se déduisent l'un de l'autre en échangeant les deux fils (conjugaison par le SWAP,
section~\ref{sec:swap}). Les quatre états de Bell obtenus sont les mêmes, avec les mêmes
indices, sauf le dernier : $-\ket{\Psi^-}$ ici, $+\ket{\Psi^-}$ là-bas. Ce signe est une phase
globale, sans effet sur les mesures (elle disparaît de $\rho$, section~\ref{sec:rho}).
\end{remarque}

\subsubsection{Matrice de densité de l'état de sortie}

Pour $\ket{\Psi}=\ket{\beta^{00}}=\frac{1}{\sqrt2}\ket{00}+\frac{1}{\sqrt2}\ket{11}$ :
\[
  \rho_{AB}=\ket{\Psi}\bra{\Psi}
  =\begin{pmatrix}\frac{1}{\sqrt2}\\0\\0\\\frac{1}{\sqrt2}\end{pmatrix}
   \begin{pmatrix}\frac{1}{\sqrt2}&0&0&\frac{1}{\sqrt2}\end{pmatrix}
  =\frac12\begin{pmatrix}1&0&0&1\\0&0&0&0\\0&0&0&0\\1&0&0&1\end{pmatrix}.
\]
Sous forme de ket-bra :
$\rho_{AB}=\frac12\big(\ket{00}\bra{00}+\ket{00}\bra{11}+\ket{11}\bra{00}+\ket{11}\bra{11}\big)$.
On vérifie $\operatorname{Tr}\rho_{AB}=\frac12(1+1)=1$ et
\[
  \rho_{AB}^2=\frac14\begin{pmatrix}2&0&0&2\\0&0&0&0\\0&0&0&0\\2&0&0&2\end{pmatrix}=\rho_{AB},
\]
donc $\operatorname{Tr}\rho_{AB}^2=1$ : le couple $(A,B)$ est dans un état \emph{pur}.

% ---------------------------------------------------------------------
\subsection{Information : variables aléatoires et états quantiques}\label{sec:info}
% ---------------------------------------------------------------------

\begin{definition}{Entropie, information mutuelle, entropie conditionnelle}
Soit $A$ une variable aléatoire de loi $p(a)$. Son \emph{entropie de Shannon}, en bits, est
\[
  H(A)=-\sum_a p(a)\log_2p(a)\qquad(\text{avec }0\log_20=0).
\]
Elle mesure l'incertitude moyenne sur $A$ : une pièce équilibrée a $H=1$ bit. Pour un couple
$(A,B)$ de loi $p(a,b)$, $H(A,B)=-\sum_{a,b}p(a,b)\log_2p(a,b)$, et
\begin{align*}
  I(A;B)&=H(A)+H(B)-H(A,B) &&\text{(information mutuelle)},\\
  H(A|B)&=H(A,B)-H(B)=H(A)-I(A;B) &&\text{(entropie conditionnelle)}.
\end{align*}
\end{definition}

Les deux écritures de $H(A|B)$ sont égales : la définition de $I(A;B)$ donne
$H(A,B)=H(A)+H(B)-I(A;B)$, donc $H(A,B)-H(B)=H(A)-I(A;B)$. L'information mutuelle $I(A;B)$
mesure ce que la connaissance de $B$ apprend sur $A$ ; $H(A|B)$ est l'incertitude qui reste sur
$A$ quand on connaît $B$.

\subsubsection{Deux variables aléatoires indépendantes}

Alice et Bob lancent chacun une pièce équilibrée, sans lien entre les deux lancers :
$p(a,b)=\frac14$ pour les quatre couples. Chaque pièce donne $0$ ou $1$ avec la probabilité $\frac12$, donc
\begin{align*}
  H(A)&=-\Big(\tfrac12\log_2\tfrac12+\tfrac12\log_2\tfrac12\Big)=1\text{ bit},\\
  H(B)&=1\text{ bit (même calcul)},\\
  H(A,B)&=-4\cdot\tfrac14\log_2\tfrac14=2\text{ bits},\\
  I(A;B)&=H(A)+H(B)-H(A,B)=1+1-2=0,\\
  H(A|B)&=H(A)-I(A;B)=1-0=H(A)=1\text{ bit}.
\end{align*}
Connaître $B$ n'apprend rien sur $A$.

\subsubsection{Deux variables solidaires}

Les deux pièces tombent toujours du même côté : $p(0,0)=p(1,1)=\frac12$ et
$p(0,1)=p(1,0)=0$. Les lois de chaque pièce sont inchangées ($p(a{=}0)=p(0,0)+p(0,1)=\frac12$,
etc.), donc $H(A)=H(B)=1$ bit, mais
\begin{align*}
  H(A,B)&=-2\cdot\tfrac12\log_2\tfrac12-2\cdot0\log_20=1\text{ bit},\\
  I(A;B)&=1+1-1=1\text{ bit},\qquad H(A|B)=1-1=0.
\end{align*}
Connaître $B$ détermine $A$ : plus aucune incertitude ne reste.

\subsubsection{Deux qubits dans un état de Bell}

Pour un système quantique, on remplace la loi de probabilité par la matrice de densité.

\begin{definition}{Trace partielle et entropie de von Neumann}
L'état du seul qubit $A$ d'un système $(A,B)$ de matrice de densité $\rho_{AB}$ est obtenu en
\emph{traçant sur $B$} : $\rho_A=\operatorname{Tr}_B(\rho_{AB})$. La trace partielle est
linéaire et remplace la partie $B$ de chaque terme par sa trace :
\[
  \operatorname{Tr}_B\big(\ket{a}\bra{a'}\otimes\ket{b}\bra{b'}\big)=\ket{a}\bra{a'}\;\braket{b'}{b}.
\]
L'\emph{entropie de von Neumann} d'un état $\rho$, de valeurs propres $\lambda_i$, est
$S(\rho)=-\operatorname{Tr}(\rho\log_2\rho)=-\sum_i\lambda_i\log_2\lambda_i$. En quantique,
$H(A)=S(\rho_A)$, $H(B)=S(\rho_B)$ et $H(A,B)=S(\rho_{AB})$ ; les définitions de
$I(A;B)$ et de $H(A|B)$ sont les mêmes.
\end{definition}

On part de l'état de Bell
$\rho_{AB}=\frac12\big(\ket{00}\bra{00}+\ket{00}\bra{11}+\ket{11}\bra{00}+\ket{11}\bra{11}\big)$
produit par le circuit de la section~\ref{sec:circuit-bell}. Chaque terme est un produit
d'opérateurs à un qubit :
$\ket{00}\bra{11}=\ket{0}\bra{1}\otimes\ket{0}\bra{1}$, etc. Alors
\begin{align*}
  \rho_A=\operatorname{Tr}_B(\rho_{AB})
  &=\tfrac12\Big(\ket{0}\bra{0}\,\braket{0}{0}+\ket{0}\bra{1}\,\braket{1}{0}
      +\ket{1}\bra{0}\,\braket{0}{1}+\ket{1}\bra{1}\,\braket{1}{1}\Big)\\
  &=\tfrac12\Big(\ket{0}\bra{0}\cdot1+\ket{0}\bra{1}\cdot0+\ket{1}\bra{0}\cdot0+\ket{1}\bra{1}\cdot1\Big)\\
  &=\tfrac12\big(\ket{0}\bra{0}+\ket{1}\bra{1}\big)=\tfrac12\,\mathrm{Id}
  =\begin{pmatrix}\frac12&0\\0&\frac12\end{pmatrix}.
\end{align*}
La formule des composantes de la section~\ref{sec:separable},
$(\rho_A)_{ij}=\rho_{(i0),(j0)}+\rho_{(i1),(j1)}$, donne le même résultat. En traçant sur $A$
au lieu de $B$, la règle est $\operatorname{Tr}_A(\ket{a}\bra{a'}\otimes\ket{b}\bra{b'})
=\braket{a'}{a}\,\ket{b}\bra{b'}$, et on trouve de même $\rho_B=\frac12\mathrm{Id}$.

Calcul des entropies :
\begin{itemize}[nosep]
  \item $\rho_{AB}$ est pur : ses valeurs propres sont $1,0,0,0$, donc
        $H(A,B)=S(\rho_{AB})=-1\cdot\log_21=0$ ;
  \item $\rho_A$ et $\rho_B$ ont pour valeurs propres $\frac12,\frac12$, donc
        $H(A)=H(B)=-\big(\frac12\log_2\frac12+\frac12\log_2\frac12\big)=1$ bit.
\end{itemize}
D'où
\begin{align*}
  I(A;B)&=H(A)+H(B)-H(A,B)=1+1-0=2\text{ bits},\\
  H(A|B)&=H(A,B)-H(B)=0-1=-1\text{ bit}\qquad(=H(A)-I(A;B)=1-2).
\end{align*}

\begin{remarque}[Une entropie conditionnelle négative]
Pour des variables classiques, $p(a,b)\leq p(b)$, donc $-\log_2p(a,b)\geq-\log_2p(b)$ et
$H(A,B)\geq H(B)$ : on a toujours $H(A|B)\geq0$, et $I(A;B)\leq H(A)$. Ici $H(A|B)=-1$ et
$I(A;B)=2>H(A)$ : ces valeurs n'ont pas d'équivalent classique, c'est une signature de
l'intrication. Le résultat est surprenant mais mathématiquement correct.
\end{remarque}

\begin{center}
\renewcommand{\arraystretch}{1.4}
\begin{tabular}{@{}lccc@{}}
\toprule
 & indépendantes & solidaires & état de Bell \\
\midrule
$H(A)$ & $1$ & $1$ & $1$ \\
$H(B)$ & $1$ & $1$ & $1$ \\
$H(A,B)$ & $2$ & $1$ & $0$ \\
$I(A;B)$ & $0$ & $1$ & $2$ \\
$H(A|B)$ & $1$ & $0$ & $-1$ \\
\bottomrule
\end{tabular}
\end{center}
(en bits ; les deux premières colonnes sont classiques, la troisième est quantique).

% ---------------------------------------------------------------------
\subsection{Système pur et système mélangé}\label{sec:pur-melange}
% ---------------------------------------------------------------------

Rappel (section~\ref{sec:pur-mixte}) : si $\lambda_i$ sont les valeurs propres de la matrice de
densité $\rho$, alors $\operatorname{Tr}\rho^2=\sum_i\lambda_i^2\leq1$. Le système est
\emph{pur} si $\operatorname{Tr}\rho^2=1$ (alors $\rho=\ket{\psi}\bra{\psi}$ et $\rho^2=\rho$),
et \emph{mélangé} si $\operatorname{Tr}\rho^2<1$. Pour un qubit,
$\operatorname{Tr}\rho^2=\frac{1+|\vec r|^2}{2}$ (section~\ref{sec:bloch}) est compris entre
$\frac12$ et $1$.

\subsubsection{Un système pur : exercice}

\emph{Énoncé.} Soit $\ket{\psi}=\frac{1}{\sqrt2}\ket{0}+\frac{1}{\sqrt2}\ket{1}$.
\begin{enumerate}[label=\alph*), nosep]
  \item Calculer la matrice de densité $\rho$.
  \item Calculer la pureté du système quantique.
  \item Calculer $\langle\sigma_x\rangle$ et $\langle\sigma_z\rangle$.
  \item Écrire $\ket{\psi}$ dans la base $X$.
\end{enumerate}

\paragraph{a) Matrice de densité.} $\rho=\ket{\psi}\bra{\psi}$ est le produit d'un vecteur
colonne par un vecteur ligne (les amplitudes sont réelles) :
\[
  \rho=\begin{pmatrix}\frac{1}{\sqrt2}\\[0.4ex]\frac{1}{\sqrt2}\end{pmatrix}
       \begin{pmatrix}\frac{1}{\sqrt2}&\frac{1}{\sqrt2}\end{pmatrix}
  =\begin{pmatrix}\frac12&\frac12\\[0.4ex]\frac12&\frac12\end{pmatrix}.
\]

\paragraph{b) Pureté.}
\[
  \rho^2=\begin{pmatrix}\frac12&\frac12\\[0.4ex]\frac12&\frac12\end{pmatrix}
         \begin{pmatrix}\frac12&\frac12\\[0.4ex]\frac12&\frac12\end{pmatrix}
  =\begin{pmatrix}\frac14+\frac14&\frac14+\frac14\\[0.4ex]\frac14+\frac14&\frac14+\frac14\end{pmatrix}
  =\begin{pmatrix}\frac12&\frac12\\[0.4ex]\frac12&\frac12\end{pmatrix}=\rho,
\]
donc $\operatorname{Tr}\rho^2=\operatorname{Tr}\rho=\frac12+\frac12=1$ : le système est \emph{pur}.

\paragraph{c) Valeurs moyennes.} Avec $\sigma_x=X=\begin{pmatrix}0&1\\1&0\end{pmatrix}$ et
$\sigma_z=Z=\begin{pmatrix}1&0\\0&-1\end{pmatrix}$ (section~\ref{sec:op-dirac}), la valeur
moyenne est $\langle\sigma\rangle=\operatorname{Tr}(\rho\,\sigma)$ (section~\ref{sec:rho}).
\begin{align*}
  \rho\,\sigma_x&=\begin{pmatrix}\frac12&\frac12\\[0.4ex]\frac12&\frac12\end{pmatrix}
     \begin{pmatrix}0&1\\1&0\end{pmatrix}
     =\begin{pmatrix}\frac12&\frac12\\[0.4ex]\frac12&\frac12\end{pmatrix},
     &\langle\sigma_x\rangle&=\operatorname{Tr}(\rho\,\sigma_x)=\tfrac12+\tfrac12=1,\\[1.5ex]
  \rho\,\sigma_z&=\begin{pmatrix}\frac12&\frac12\\[0.4ex]\frac12&\frac12\end{pmatrix}
     \begin{pmatrix}1&0\\0&-1\end{pmatrix}
     =\begin{pmatrix}\frac12&-\frac12\\[0.4ex]\frac12&-\frac12\end{pmatrix},
     &\langle\sigma_z\rangle&=\operatorname{Tr}(\rho\,\sigma_z)=\tfrac12-\tfrac12=0.
\end{align*}
On retrouve $\langle\sigma_z\rangle$ avec les probabilités de mesure : dans la base $Z$,
$\sigma_z$ vaut $+1$ pour $\ket{0}$ et $-1$ pour $\ket{1}$, et
$|\braket{0}{\psi}|^2=|\braket{1}{\psi}|^2=\frac12$, donc
$\langle\sigma_z\rangle=\frac12\times(+1)+\frac12\times(-1)=0$.

\paragraph{d) Écriture dans la base $X$.} La base $X$ est $\{\ket{+},\ket{-}\}$, orthonormée, donc
$\ket{+}\bra{+}+\ket{-}\bra{-}=\mathrm{Id}$ et
$\ket{\psi}=\ket{+}\braket{+}{\psi}+\ket{-}\braket{-}{\psi}$, avec
\[
  \braket{+}{\psi}=\tfrac{1}{\sqrt2}\Big(\tfrac{1}{\sqrt2}+\tfrac{1}{\sqrt2}\Big)=1,
  \qquad
  \braket{-}{\psi}=\tfrac{1}{\sqrt2}\Big(\tfrac{1}{\sqrt2}-\tfrac{1}{\sqrt2}\Big)=0 .
\]
Donc $\ket{\psi}=1\cdot\ket{+}+0\cdot\ket{-}=\ket{+}$ : dans la base $X$, l'état n'a qu'une
seule composante. Comme $\sigma_x\ket{+}=+\ket{+}$ et $\sigma_x\ket{-}=-\ket{-}$, une
mesure de $\sigma_x$ donne $+1$ avec certitude, et
$\langle\sigma_x\rangle=1\cdot(+1)+0\cdot(-1)=1$ : c'est le résultat de la question c).
Le vecteur de Bloch de $\ket{+}$ est $(\langle\sigma_x\rangle,\langle\sigma_y\rangle,\langle\sigma_z\rangle)=(1,0,0)$,
de norme $1$ : l'état est sur la sphère, en accord avec $\operatorname{Tr}\rho^2=\frac{1+1}{2}=1$.

\subsubsection{Un système mélangé : un qubit de la paire de Bell}

À la section~\ref{sec:info}, la trace sur $B$ de l'état de Bell a donné
$\rho_A=\frac12\mathrm{Id}=\begin{pmatrix}\frac12&0\\0&\frac12\end{pmatrix}$. Alors
\[
  \rho_A^2=\begin{pmatrix}\frac14&0\\0&\frac14\end{pmatrix},
  \qquad
  \operatorname{Tr}\big(\rho_A^2\big)=\frac14+\frac14=\frac12<1 :
\]
le qubit $A$ est dans un état \emph{mélangé}, et c'est le plus mélangé possible pour un qubit
(centre de la boule de Bloch, $\vec r=\vec0$). Pourtant le couple $(A,B)$ est pur
($\operatorname{Tr}\rho_{AB}^2=1$) : on connaît parfaitement l'état du couple, mais pas celui d'un
qubit pris seul. L'information est entièrement dans les corrélations entre $A$ et $B$, et
$S(\rho_A)=1$ bit.

\subsubsection{Un système mélangé : le mélange statistique de Bob}

Bob prépare un qubit dans l'état $\ket{0}$ avec la probabilité $p_1=\frac34$ (75\,\%) et dans
l'état $\ket{1}$ avec la probabilité $p_2=\frac14$ (25\,\%), sans que l'on sache lequel des deux a
été préparé. Ce n'est pas un état pur $\ket{\psi}$ : on décrit le système par la matrice de
densité $\rho=\sum_ip_i\ket{\psi_i}\bra{\psi_i}$ (section~\ref{sec:pur-mixte}) :
\[
  \rho=\underbrace{\tfrac34}_{p_1}\ket{0}\bra{0}+\underbrace{\tfrac14}_{p_2}\ket{1}\bra{1}
  =\begin{pmatrix}\frac34&0\\0&\frac14\end{pmatrix}.
\]
Alors
\[
  \rho^2=\begin{pmatrix}\frac{9}{16}&0\\0&\frac{1}{16}\end{pmatrix},
  \qquad
  \operatorname{Tr}\big(\rho^2\big)=\frac{9}{16}+\frac{1}{16}=\frac{10}{16}=\frac58<1 :
\]
le système est \emph{mélangé}.

Les valeurs moyennes sont les moyennes pondérées sur la préparation :
$\langle\sigma_z\rangle=\operatorname{Tr}(\rho\,\sigma_z)=\frac34-\frac14=\frac12
=p_1\cdot(+1)+p_2\cdot(-1)$, et $\langle\sigma_x\rangle=\operatorname{Tr}(\rho\,\sigma_x)=0$
(les termes hors diagonale de $\rho$ sont nuls). Le vecteur de Bloch est $(0,0,\frac12)$,
à l'intérieur de la boule, et $\operatorname{Tr}\rho^2=\frac{1+1/4}{2}=\frac58$. Son entropie
est $S(\rho)=-\frac34\log_2\frac34-\frac14\log_2\frac14=2-\frac34\log_23\approx0{,}811$ bit.

\begin{remarque}[Mélange et superposition]
Considérons l'état pur $\ket{\psi'}=\frac{\sqrt3}{2}\ket{0}+\frac12\ket{1}$, qui a les mêmes
probabilités en base $Z$ ($P(0)=\frac34$, $P(1)=\frac14$). Sa matrice de densité est
\[
  \rho'=\begin{pmatrix}\frac34&\frac{\sqrt3}{4}\\[0.4ex]\frac{\sqrt3}{4}&\frac14\end{pmatrix},
  \qquad
  \operatorname{Tr}\rho'^2=\frac{9}{16}+\frac1{16}+2\cdot\frac{3}{16}=1,
  \qquad
  \langle\sigma_x\rangle=2\cdot\frac{\sqrt3}{4}=\frac{\sqrt3}{2}.
\]
Le mélange de Bob a la même diagonale mais aucun terme hors diagonale ($\operatorname{Tr}\rho^2=\frac58$,
$\langle\sigma_x\rangle=0$) : une superposition n'est pas un tirage au hasard entre $\ket{0}$ et
$\ket{1}$. De plus, pour un état pur il existe une mesure au résultat certain (ici $\sigma_x$ pour
$\ket{+}$) ; pour le mélange de Bob, la meilleure probabilité d'obtenir un même résultat est
$75\,\%$ (mesure en base $Z$, résultat $\ket{0}$).
\end{remarque}

\subsubsection{Bilan}

\begin{center}
\renewcommand{\arraystretch}{1.5}
\begin{tabular}{@{}lcccl@{}}
\toprule
\textbf{Système} & $\rho$ & $\operatorname{Tr}\rho^2$ & $S(\rho)$ & \textbf{Nature} \\
\midrule
$\ket{\psi}=\ket{+}$ & ${\renewcommand{\arraystretch}{1}\begin{pmatrix}\frac12&\frac12\\\frac12&\frac12\end{pmatrix}}$ & $1$ & $0$ & pur \\[2ex]
Qubit $A$ d'une paire de Bell & ${\renewcommand{\arraystretch}{1}\begin{pmatrix}\frac12&0\\0&\frac12\end{pmatrix}}$ & $\frac12$ & $1$ bit & mélangé \\[2ex]
Source de Bob (75\,\%/25\,\%) & ${\renewcommand{\arraystretch}{1}\begin{pmatrix}\frac34&0\\0&\frac14\end{pmatrix}}$ & $\frac58$ & $\approx0{,}811$ bit & mélangé \\
\bottomrule
\end{tabular}
\end{center}

% ---------------------------------------------------------------------
\subsection{Algorithme quantique : la téléportation}\label{sec:teleportation}
% ---------------------------------------------------------------------

\emph{Problème.} Alice détient un qubit $Q$ dans un état qu'elle ne connaît pas,
$\ket{\varphi}=\alpha\ket{0}+\beta\ket{1}$ avec $|\alpha|^2+|\beta|^2=1$, et elle veut que Bob
obtienne cet état sur l'un de ses qubits. Elle ne peut pas mesurer $Q$ pour connaître
$\alpha$ et $\beta$ (une mesure ne donne qu'un bit et détruit l'état), elle ne peut pas le
copier (théorème de non-clonage) et elle n'envoie pas le qubit lui-même.

\begin{definition}{Protocole de téléportation quantique}
\textbf{Ressources.}
\begin{itemize}[nosep]
  \item Un \emph{e-bit} : Alice et Bob partagent une paire de qubits $(A,B)$ dans l'état de Bell
        $\ket{\Phi^+}_{BA}=\frac{1}{\sqrt2}\big(\ket{00}+\ket{11}\big)$ ; Alice détient $A$,
        Bob détient $B$.
  \item Le qubit $Q$ à téléporter, chez Alice.
  \item Un canal classique (deux bits).
\end{itemize}
Le système est écrit $(B,A,Q)$ : $\ket{b,a,q}=\ket{b}_B\otimes\ket{a}_A\otimes\ket{q}_Q$.

\textbf{Étapes.}
\begin{enumerate}[nosep]
  \item Alice applique un CNOT de contrôle $Q$ et de cible $A$.
  \item Alice applique une porte $H$ sur $Q$.
  \item Alice mesure $A$ et $Q$ dans la base de calcul, obtient deux bits $(a,q)$ et les envoie
        à Bob.
  \item Bob applique $X$ si $a=1$, puis $Z$ si $q=1$, à son qubit $B$.
\end{enumerate}
\end{definition}

\begin{figure}[!ht]
\centering
\begin{adjustbox}{max width=\linewidth}
\begin{tikzpicture}[>=Stealth, font=\small]
  % fils quantiques : Q en haut, A au milieu, B en bas
  \draw (0,2.4) -- (5.3,2.4);
  \draw (0,1.2) -- (5.3,1.2);
  \draw (0,0) -- (10.6,0);
  \node[anchor=east] at (0,2.4) {$\ket{\varphi}_Q$};
  \node[anchor=south west, text=insagris, font=\footnotesize] at (0.05,1.2) {$A$};
  \node[anchor=south west, text=insagris, font=\footnotesize] at (0.05,0) {$B$};
  % paire de Bell partagée
  \draw[insagris, thick] (-0.1,1.5) -- (-0.25,1.5) -- (-0.25,-0.3) -- (-0.1,-0.3);
  \node[anchor=east] at (-0.35,0.6) {$\ket{\Phi^+}_{BA}$};
  % CNOT : contrôle Q, cible A
  \draw[insarouge, thick] (2.6,2.4) -- (2.6,1.2);
  \pic at (2.6,2.4) {ctl};
  \pic at (2.6,1.2) {cible};
  % Hadamard sur Q
  \node[porte] at (4.2,2.4) {$H$};
  % mesures d'Alice
  \pic at (5.6,2.4) {mesure};
  \pic at (5.6,1.2) {mesure};
  % fils classiques (doubles)
  \draw[double, double distance=1.2pt, insagris] (5.9,1.2) -- (7.6,1.2) -- (7.6,0.35);
  \draw[double, double distance=1.2pt, insagris] (5.9,2.4) -- (9.6,2.4) -- (9.6,0.35);
  \node[above, font=\footnotesize] at (6.5,1.2) {$a$};
  \node[above, font=\footnotesize] at (6.5,2.4) {$q$};
  % corrections de Bob
  \node[porte] at (7.6,0) {$X$};
  \node[porte] at (9.6,0) {$Z$};
  \node[anchor=west] at (10.6,0) {$\ket{\varphi}_B$};
  % états intermédiaires
  \foreach \x in {1.4,3.4,4.9}
    \draw[dashed, insagris!70] (\x,2.8) -- (\x,-0.5);
  \node[anchor=north, font=\footnotesize] at (1.4,-0.5) {$\ket{\Pi_0}$};
  \node[anchor=north, font=\footnotesize] at (3.4,-0.5) {$\ket{\Pi_1}$};
  \node[anchor=north, font=\footnotesize] at (4.9,-0.5) {$\ket{\Pi_2}$};
\end{tikzpicture}
\end{adjustbox}
\caption{Circuit de téléportation. Alice détient $Q$ (dans $\ket{\varphi}=\alpha\ket{0}+\beta\ket{1}$)
et $A$, Bob détient $B$ ; $A$ et $B$ forment la paire $\ket{\Phi^+}_{BA}$. Les traits doubles
sont des fils \emph{classiques} qui portent les résultats $a$ (mesure de $A$) et $q$ (mesure de
$Q$) ; $X$ est appliqué si $a=1$, $Z$ si $q=1$. Le ket s'écrit $\ket{b,a,q}$ : le fil du haut
($Q$) est le dernier chiffre.}
\label{fig:teleportation}
\end{figure}

\subsubsection{Calcul pas à pas}

\paragraph{État initial $\ket{\Pi_0}$.} Alice et Bob partagent la paire, et $Q$ est indépendant
d'elle :
\begin{align*}
  \ket{\Pi_0}&=\ket{\Phi^+}_{BA}\otimes\ket{\varphi}_Q
    =\frac{1}{\sqrt2}\big(\ket{00}+\ket{11}\big)\otimes\big(\alpha\ket{0}+\beta\ket{1}\big)\\
  &=\frac{1}{\sqrt2}\Big(\alpha\ket{00}\ket{0}+\beta\ket{00}\ket{1}+\alpha\ket{11}\ket{0}+\beta\ket{11}\ket{1}\Big)\\
  &=\frac{\alpha\ket{000}+\beta\ket{001}+\alpha\ket{110}+\beta\ket{111}}{\sqrt2}.
\end{align*}

\paragraph{Après le CNOT $Q\to A$ : $\ket{\Pi_1}$.} Le CNOT (contrôle $Q$, dernier chiffre ;
cible $A$, chiffre du milieu) agit ainsi : $\ket{b,a,q}\to\ket{b,\,a\oplus q,\,q}$. Terme par terme :
\begin{center}
\renewcommand{\arraystretch}{1.3}
\begin{tabular}{@{}cccc@{}}
\toprule
Terme & $q$ & Effet sur $A$ & Résultat \\
\midrule
$\alpha\ket{000}$ & $0$ & aucun & $\alpha\ket{000}$ \\
$\beta\ket{001}$ & $1$ & $0\to1$ & $\beta\ket{011}$ \\
$\alpha\ket{110}$ & $0$ & aucun & $\alpha\ket{110}$ \\
$\beta\ket{111}$ & $1$ & $1\to0$ & $\beta\ket{101}$ \\
\bottomrule
\end{tabular}
\end{center}
donc
\[
  \ket{\Pi_1}=\frac{\alpha\ket{000}+\beta\ket{011}+\alpha\ket{110}+\beta\ket{101}}{\sqrt2}.
\]
On isole ensuite le dernier chiffre, celui de $Q$, qui va changer à l'étape suivante
(le premier ket est celui du couple $\ket{b,a}$) :
\[
  \ket{\Pi_1}=\frac{\alpha\ket{00}\otimes\ket{0}+\beta\ket{01}\otimes\ket{1}
    +\alpha\ket{11}\otimes\ket{0}+\beta\ket{10}\otimes\ket{1}}{\sqrt2}.
\]

\paragraph{Après la porte $H$ sur $Q$ : $\ket{\Pi_2}$.} On remplace $\ket{0}$ par
$H\ket{0}=\ket{+}$ et $\ket{1}$ par $H\ket{1}=\ket{-}$ dans le facteur $Q$ :
\[
  \ket{\Pi_2}=\frac{1}{\sqrt2}\Big(\alpha\ket{00}\otimes\ket{+}+\beta\ket{01}\otimes\ket{-}
    +\alpha\ket{11}\otimes\ket{+}+\beta\ket{10}\otimes\ket{-}\Big).
\]
Avec $\ket{\pm}=\frac{1}{\sqrt2}(\ket{0}\pm\ket{1})$, le préfacteur devient
$\frac{1}{\sqrt2}\cdot\frac{1}{\sqrt2}=\frac12$ :
\begin{align*}
  \ket{\Pi_2}&=\frac12\Big(\alpha\ket{00}\big(\ket{0}+\ket{1}\big)+\beta\ket{01}\big(\ket{0}-\ket{1}\big)
     +\alpha\ket{11}\big(\ket{0}+\ket{1}\big)+\beta\ket{10}\big(\ket{0}-\ket{1}\big)\Big)\\
  &=\frac12\Big(\alpha\ket{000}+\alpha\ket{001}+\beta\ket{010}-\beta\ket{011}
     +\alpha\ket{110}+\alpha\ket{111}+\beta\ket{100}-\beta\ket{101}\Big).
\end{align*}

\paragraph{Regroupement selon le résultat d'Alice.} Alice mesurera $A$ et $Q$, donc les deux
derniers chiffres. On regroupe les huit termes qui ont les mêmes chiffres $(a,q)$, en mettant
le ket de Bob en premier :
\begin{center}
\renewcommand{\arraystretch}{1.3}
\begin{tabular}{@{}cll@{}}
\toprule
$(a,q)$ & Termes de $\ket{\Pi_2}$ (facteur $\frac12$ omis) & Ket de Bob \\
\midrule
$(0,0)$ & $\alpha\ket{000}+\beta\ket{100}$ & $(\alpha\ket{0}+\beta\ket{1})\ket{00}$ \\
$(0,1)$ & $\alpha\ket{001}-\beta\ket{101}$ & $(\alpha\ket{0}-\beta\ket{1})\ket{01}$ \\
$(1,0)$ & $\beta\ket{010}+\alpha\ket{110}$ & $(\beta\ket{0}+\alpha\ket{1})\ket{10}$ \\
$(1,1)$ & $-\beta\ket{011}+\alpha\ket{111}$ & $(-\beta\ket{0}+\alpha\ket{1})\ket{11}$ \\
\bottomrule
\end{tabular}
\end{center}
d'où
\[
  \ket{\Pi_2}=\frac12\Big[\big(\alpha\ket{0}+\beta\ket{1}\big)\ket{00}
    +\big(\alpha\ket{0}-\beta\ket{1}\big)\ket{01}
    +\big(\beta\ket{0}+\alpha\ket{1}\big)\ket{10}
    +\big(\alpha\ket{1}-\beta\ket{0}\big)\ket{11}\Big].
\]
Dans chaque terme, le ket de gauche est l'état de $B$ et le ket de droite $\ket{a,q}$ celui du
couple $(A,Q)$ d'Alice.

\subsubsection{Mesure d'Alice et corrections de Bob}

Alice mesure $A$ et $Q$. Comme à la section~\ref{sec:mesure2}, la probabilité d'un résultat
$(a,q)$ est le carré de la norme du terme correspondant, et l'état de Bob devient le ket
associé (déjà normé). Ici
\[
  P(a,q)=\Big\|\tfrac12\ket{\phi_{aq}}\Big\|^2=\tfrac14\big(|\alpha|^2+|\beta|^2\big)=\tfrac14
  \quad\text{pour les quatre résultats,}
\]
qui sont donc équiprobables et \emph{indépendants} de $\alpha$ et $\beta$ : les deux bits
envoyés ne révèlent rien sur $\ket{\varphi}$. Les kets de Bob sont
$\ket{\phi_{00}}=\alpha\ket{0}+\beta\ket{1}$, $\ket{\phi_{01}}=\alpha\ket{0}-\beta\ket{1}$,
$\ket{\phi_{10}}=\beta\ket{0}+\alpha\ket{1}$ et $\ket{\phi_{11}}=-\beta\ket{0}+\alpha\ket{1}$.

\begin{center}
\renewcommand{\arraystretch}{1.5}
\begin{tabular}{@{}cccll@{}}
\toprule
Résultat $(a,q)$ & Probabilité & État de Bob & Lien avec $\ket{\varphi}$ & Correction de Bob \\
\midrule
$(0,0)$ & $\frac14$ & $\alpha\ket{0}+\beta\ket{1}$ & $\ket{\varphi}$ & $I$ (rien) \\
$(0,1)$ & $\frac14$ & $\alpha\ket{0}-\beta\ket{1}$ & $Z\ket{\varphi}$ & $Z$ \\
$(1,0)$ & $\frac14$ & $\beta\ket{0}+\alpha\ket{1}$ & $X\ket{\varphi}$ & $X$ \\
$(1,1)$ & $\frac14$ & $-\beta\ket{0}+\alpha\ket{1}$ & $XZ\ket{\varphi}$ & $X$ puis $Z$ (opérateur $ZX$) \\
\bottomrule
\end{tabular}
\end{center}

\paragraph{Vérification des quatre corrections.} Avec $X\ket{0}=\ket{1}$, $X\ket{1}=\ket{0}$,
$Z\ket{0}=\ket{0}$ et $Z\ket{1}=-\ket{1}$ :
\begin{align*}
  (0,0):\quad & I\big(\alpha\ket{0}+\beta\ket{1}\big)=\alpha\ket{0}+\beta\ket{1},\\
  (0,1):\quad & Z\big(\alpha\ket{0}-\beta\ket{1}\big)=\alpha\ket{0}-\beta Z\ket{1}=\alpha\ket{0}+\beta\ket{1},\\
  (1,0):\quad & X\big(\beta\ket{0}+\alpha\ket{1}\big)=\beta\ket{1}+\alpha\ket{0}=\alpha\ket{0}+\beta\ket{1},\\
  (1,1):\quad & X\big(-\beta\ket{0}+\alpha\ket{1}\big)=-\beta\ket{1}+\alpha\ket{0}=\alpha\ket{0}-\beta\ket{1},
     \ \text{puis}\ Z\big(\alpha\ket{0}-\beta\ket{1}\big)=\alpha\ket{0}+\beta\ket{1}.
\end{align*}
Dans le dernier cas, avec les matrices,
$ZX=\begin{pmatrix}1&0\\0&-1\end{pmatrix}\begin{pmatrix}0&1\\1&0\end{pmatrix}
=\begin{pmatrix}0&1\\-1&0\end{pmatrix}$ et
$ZX\begin{pmatrix}-\beta\\ \alpha\end{pmatrix}=\begin{pmatrix}\alpha\\ \beta\end{pmatrix}$.
L'ordre compte : l'état de Bob est $XZ\ket{\varphi}$ ($Z\ket{\varphi}=\alpha\ket{0}-\beta\ket{1}$,
puis $X$ donne $\alpha\ket{1}-\beta\ket{0}$), et on l'annule avec
$(XZ)^{-1}=Z^{-1}X^{-1}=ZX$ : on applique d'abord $X$, puis $Z$, comme dans le circuit.
Dans les quatre cas, Bob retrouve exactement $\ket{\varphi}=\alpha\ket{0}+\beta\ket{1}$.

\subsubsection{Ce que Bob sait avant de recevoir les bits}

Tant que Bob ignore $(a,q)$, son qubit est dans le mélange des quatre kets $\ket{\phi_{aq}}$, chacun
avec la probabilité $\frac14$ :
\begin{align*}
  \rho_B&=\frac14\sum_{a,q}\ket{\phi_{aq}}\bra{\phi_{aq}}\\
  &=\frac14\Bigg[\begin{pmatrix}|\alpha|^2&\alpha\beta^*\\ \alpha^*\beta&|\beta|^2\end{pmatrix}
    +\begin{pmatrix}|\alpha|^2&-\alpha\beta^*\\ -\alpha^*\beta&|\beta|^2\end{pmatrix}
    +\begin{pmatrix}|\beta|^2&\alpha^*\beta\\ \alpha\beta^*&|\alpha|^2\end{pmatrix}
    +\begin{pmatrix}|\beta|^2&-\alpha^*\beta\\ -\alpha\beta^*&|\alpha|^2\end{pmatrix}\Bigg]\\
  &=\frac14\begin{pmatrix}2(|\alpha|^2+|\beta|^2)&0\\0&2(|\alpha|^2+|\beta|^2)\end{pmatrix}
   =\frac12\begin{pmatrix}1&0\\0&1\end{pmatrix}=\frac12\,\mathrm{Id}.
\end{align*}
Sans les deux bits, Bob a donc un qubit complètement mélangé, qui ne dépend pas de $\alpha$ ni de
$\beta$ : il n'a reçu aucune information sur $\ket{\varphi}$. C'est l'envoi classique de $(a,q)$
qui permet de la retrouver, donc rien ne voyage plus vite que la lumière.

\begin{remarque}[Bilan de la téléportation]
\begin{itemize}[nosep]
  \item Ressources : $1$ e-bit (la paire de Bell) et $2$ bits classiques pour transmettre
        \emph{un qubit} d'état quelconque, sans qu'Alice connaisse $\alpha$ ni $\beta$.
  \item Pas de clonage : après la mesure, $A$ et $Q$ sont dans l'état de base $\ket{a,q}$.
        Alice n'a plus $\ket{\varphi}$ : l'état a été \emph{déplacé}, pas copié. L'intrication
        entre $A$ et $B$ a été consommée.
  \item Les quatre corrections $I,Z,X,ZX$ sont des portes contrôlées par des bits classiques
        (section~\ref{sec:controlees}).
\end{itemize}
\end{remarque}

% ---------------------------------------------------------------------
\subsection{Portes contrôlées : contrôlé-\texorpdfstring{$U$}{U}, Fredkin et Toffoli}\label{sec:controlees}
% ---------------------------------------------------------------------

Le CNOT est le cas le plus simple d'une idée générale : appliquer une opération à un ou
plusieurs qubits \emph{cibles} seulement si un ou plusieurs qubits de \emph{contrôle} valent $1$.
Dans cette section, comme à la section~\ref{sec:op2}, le contrôle est le \emph{premier} chiffre
du ket. Les qubits sont notés $x$, $y$, $z$ (le qubit $x$ est le contrôle) et
$\ket{x,y,z}=\ket{x}\otimes\ket{y}\otimes\ket{z}$ ; d'après la convention de la section, le
contrôle est dessiné en bas (figure~\ref{fig:controlees}). Pour éviter la confusion avec les
matrices de Pauli, on écrit ici $\sigma_x=X$ pour la porte NOT.

\begin{figure}[!ht]
\centering
\begin{adjustbox}{max width=\linewidth}
\begin{tikzpicture}[>=Stealth, font=\small]
  % (a) contrôlé-U : contrôle x en bas, cible y en haut
  \draw (0,1.2) -- (3.4,1.2);
  \draw (0,0) -- (3.4,0);
  \node[anchor=east] at (0,1.2) {$y$};
  \node[anchor=east] at (0,0) {$x$};
  \draw[insarouge, thick] (1.7,0) -- (1.7,0.85);
  \pic at (1.7,0) {ctl};
  \node[porte] at (1.7,1.2) {$U$};
  \node[anchor=north, font=\footnotesize] at (1.7,-0.5) {(a) contrôlé-$U$};
  % (b) contrôlé-SWAP (Fredkin) : x contrôle, y et z échangés
  \draw (5,2.4) -- (8.4,2.4);
  \draw (5,1.2) -- (8.4,1.2);
  \draw (5,0) -- (8.4,0);
  \node[anchor=east] at (5,2.4) {$z$};
  \node[anchor=east] at (5,1.2) {$y$};
  \node[anchor=east] at (5,0) {$x$};
  \draw[insarouge, thick] (6.7,0) -- (6.7,2.4);
  \pic at (6.7,0) {ctl};
  \pic at (6.7,1.2) {croix};
  \pic at (6.7,2.4) {croix};
  \node[anchor=north, font=\footnotesize] at (6.7,-0.5) {(b) Fredkin (C-SWAP)};
  % (c) Toffoli (CCNOT) : x et y contrôlent, z cible
  \draw (10,2.4) -- (13.4,2.4);
  \draw (10,1.2) -- (13.4,1.2);
  \draw (10,0) -- (13.4,0);
  \node[anchor=east] at (10,2.4) {$z$};
  \node[anchor=east] at (10,1.2) {$y$};
  \node[anchor=east] at (10,0) {$x$};
  \draw[insarouge, thick] (11.7,0) -- (11.7,2.4);
  \pic at (11.7,0) {ctl};
  \pic at (11.7,1.2) {ctl};
  \pic at (11.7,2.4) {cible};
  \node[anchor=north, font=\footnotesize] at (11.7,-0.5) {(c) Toffoli (CCNOT)};
\end{tikzpicture}
\end{adjustbox}
\caption{Portes contrôlées. Le contrôle $x$ est le premier chiffre du ket, donc le fil du bas.
(a) Le qubit $y$ subit $U$ si $x=1$ ; pour $U=X$ on retrouve le CNOT. (b) Si $x=1$, les qubits
$y$ et $z$ sont échangés. (c) Si $x=y=1$, le qubit $z$ est inversé.}
\label{fig:controlees}
\end{figure}

\subsubsection{Contrôlé-\texorpdfstring{$U$}{U}}

\begin{definition}{Porte contrôlée-$U$}
Soit $U$ une matrice unitaire $2\times2$. Sur le système $(x,y)$, de contrôle $x$ et de cible $y$,
\[
  \mathrm{C}U=\ket{0}\bra{0}_x\otimes I_y+\ket{1}\bra{1}_x\otimes U_y .
\]
Le premier terme agit quand le contrôle est $\ket{0}$ : la cible ne change pas ($I$). Le second
agit quand le contrôle est $\ket{1}$ : la cible subit $U$.
\end{definition}

Comme $\ket{0}\bra{0}\ket{x}=\delta_{x0}\ket{0}$ et $\ket{1}\bra{1}\ket{x}=\delta_{x1}\ket{1}$,
\[
  \mathrm{C}U\big(\ket{x}\otimes\ket{y}\big)=\ket{x}\otimes U^x\ket{y},\qquad U^0=I,\ U^1=U .
\]
Par la règle $A\otimes B=(A_{ij}B)$, $\ket{0}\bra{0}\otimes I=\begin{pmatrix}I&0\\0&0\end{pmatrix}$
et $\ket{1}\bra{1}\otimes U=\begin{pmatrix}0&0\\0&U\end{pmatrix}$ (blocs $2\times2$), donc
\[
  \mathrm{C}U=\begin{pmatrix}I&0\\0&U\end{pmatrix}
  =\begin{pmatrix}1&0&0&0\\0&1&0&0\\0&0&u_{00}&u_{01}\\0&0&u_{10}&u_{11}\end{pmatrix}
  \qquad\text{pour }U=\begin{pmatrix}u_{00}&u_{01}\\u_{10}&u_{11}\end{pmatrix}.
\]
Cette matrice est unitaire car $I^\dagger I=I$ et $U^\dagger U=I$.

\paragraph{Le CNOT est un cas particulier.} Pour $U=\sigma_x$,
$\mathrm{C}\sigma_x=\ket{0}\bra{0}_x\otimes I_y+\ket{1}\bra{1}_x\otimes\sigma_x$ est le CNOT de
contrôle $x$ (premier chiffre), c'est-à-dire $\mathrm{CNOT}_{12}$ de la section~\ref{sec:cnot}.
On calcule l'action sur les quatre états de base, avec $\braket{0}{0}=\braket{1}{1}=1$,
$\braket{0}{1}=\braket{1}{0}=0$, $\sigma_x\ket{0}=\ket{1}$ et $\sigma_x\ket{1}=\ket{0}$ :
\begin{align*}
  \mathrm{C}\sigma_x\ket{00}&=\big[\ket{0}\bra{0}_x\otimes I_y\big]\ket{0}_x\ket{0}_y
      +\big[\ket{1}\bra{1}_x\otimes\sigma_x\big]\ket{0}_x\ket{0}_y\\
    &=\ket{0}\braket{0}{0}\otimes I\ket{0}+\ket{1}\braket{1}{0}\otimes\sigma_x\ket{0}
     =\ket{0}\otimes\ket{0}+0=\ket{00},\\[1.2ex]
  \mathrm{C}\sigma_x\ket{01}&=\ket{0}\braket{0}{0}\otimes I\ket{1}+\ket{1}\braket{1}{0}\otimes\sigma_x\ket{1}
     =\ket{0}\otimes\ket{1}+0=\ket{01},\\[1.2ex]
  \mathrm{C}\sigma_x\ket{10}&=\ket{0}\braket{0}{1}\otimes I\ket{0}+\ket{1}\braket{1}{1}\otimes\sigma_x\ket{0}
     =0+\ket{1}\otimes\ket{1}=\ket{11},\\[1.2ex]
  \mathrm{C}\sigma_x\ket{11}&=\ket{0}\braket{0}{1}\otimes I\ket{1}+\ket{1}\braket{1}{1}\otimes\sigma_x\ket{1}
     =0+\ket{1}\otimes\ket{0}=\ket{10}.
\end{align*}
On retrouve la table du CNOT : la cible n'est inversée que si le contrôle vaut $1$.

\paragraph{Un autre exemple : $U=Z$.} Ici
\[
  \mathrm{C}Z=\begin{pmatrix}I&0\\0&Z\end{pmatrix}=\mathrm{diag}(1,1,1,-1) :
\]
la porte multiplie $\ket{11}$ par $-1$ et ne touche pas les autres états de base. Elle est
symétrique entre les deux qubits. Elle se déduit du CNOT par
$\mathrm{C}Z=(I\otimes H)\,\mathrm{C}X\,(I\otimes H)$, où $\mathrm{C}X$ est le contrôlé-$X$.
En effet, $I\otimes H=\mathrm{diag}(H,H)$ (deux blocs $2\times2$ égaux à $H$), et le produit de
matrices diagonales par blocs se fait bloc par bloc :
\[
  (I\otimes H)\begin{pmatrix}I&0\\0&X\end{pmatrix}(I\otimes H)
  =\begin{pmatrix}H\,I\,H&0\\0&H\,X\,H\end{pmatrix}
  =\begin{pmatrix}I&0\\0&HXH\end{pmatrix},
\]
car $H^2=I$, et
\[
  HXH=\frac12\begin{pmatrix}1&1\\1&-1\end{pmatrix}\begin{pmatrix}0&1\\1&0\end{pmatrix}\begin{pmatrix}1&1\\1&-1\end{pmatrix}
  =\frac12\begin{pmatrix}1&1\\1&-1\end{pmatrix}\begin{pmatrix}1&-1\\1&1\end{pmatrix}
  =\frac12\begin{pmatrix}2&0\\0&-2\end{pmatrix}=Z
\]
(section~\ref{sec:hadamard} : $H$ échange les axes $x$ et $z$).

\subsubsection{Contrôlé-SWAP : la porte de Fredkin}

\begin{definition}{Porte de Fredkin (C-SWAP)}
Sur trois qubits $(x,y,z)$ (espace de dimension $2^3=8$), avec $x$ pour contrôle,
\[
  \mathrm{C\text{-}SWAP}=\ket{0}\bra{0}_x\otimes I_y\otimes I_z+\ket{1}\bra{1}_x\otimes\mathrm{SWAP}_{yz}.
\]
Si $x=0$, rien ne change ; si $x=1$, les qubits $y$ et $z$ sont échangés.
\end{definition}

Calcul sur quatre états de base (avec $I\ket{ab}=\ket{ab}$ et $\mathrm{SWAP}\ket{ab}=\ket{ba}$) :
\begin{align*}
  \mathrm{C\text{-}SWAP}\ket{000}&=\ket{0}\braket{0}{0}\otimes I\ket{00}+\ket{1}\braket{1}{0}\otimes\mathrm{SWAP}\ket{00}
     =\ket{0}\otimes\ket{00}+0=\ket{000},\\[1.2ex]
  \mathrm{C\text{-}SWAP}\ket{101}&=\ket{0}\braket{0}{1}\otimes I\ket{01}+\ket{1}\braket{1}{1}\otimes\mathrm{SWAP}\ket{01}
     =0+\ket{1}\otimes\ket{10}=\ket{110},\\[1.2ex]
  \mathrm{C\text{-}SWAP}\ket{110}&=\ket{0}\braket{0}{1}\otimes I\ket{10}+\ket{1}\braket{1}{1}\otimes\mathrm{SWAP}\ket{10}
     =0+\ket{1}\otimes\ket{01}=\ket{101},\\[1.2ex]
  \mathrm{C\text{-}SWAP}\ket{010}&=\ket{0}\braket{0}{0}\otimes I\ket{10}+\ket{1}\braket{1}{0}\otimes\mathrm{SWAP}\ket{10}
     =\ket{0}\otimes\ket{10}+0=\ket{010}.
\end{align*}
Les huit états de base donnent :
\begin{center}
\renewcommand{\arraystretch}{1.3}
\begin{tabular}{@{}ccccccccc@{}}
\toprule
Entrée & $\ket{000}$ & $\ket{001}$ & $\ket{010}$ & $\ket{011}$ & $\ket{100}$ & $\ket{101}$ & $\ket{110}$ & $\ket{111}$ \\
\midrule
Sortie & $\ket{000}$ & $\ket{001}$ & $\ket{010}$ & $\ket{011}$ & $\ket{100}$ & $\ket{110}$ & $\ket{101}$ & $\ket{111}$ \\
\bottomrule
\end{tabular}
\end{center}
Seuls $\ket{101}$ et $\ket{110}$ changent (contrôle à $1$ et $y\neq z$) ; pour $\ket{100}$ et
$\ket{111}$ l'échange de deux qubits égaux ne change rien. Dans la base
$\ket{000},\ket{001},\dots,\ket{111}$ (ordre binaire), la matrice est
$\mathrm{diag}(I_4,\mathrm{SWAP})$ (deux blocs $4\times4$) :
\[
  \mathrm{C\text{-}SWAP}=\begingroup\setlength{\arraycolsep}{3pt}\begin{pmatrix}
    1&0&0&0&0&0&0&0\\
    0&1&0&0&0&0&0&0\\
    0&0&1&0&0&0&0&0\\
    0&0&0&1&0&0&0&0\\
    0&0&0&0&1&0&0&0\\
    0&0&0&0&0&0&1&0\\
    0&0&0&0&0&1&0&0\\
    0&0&0&0&0&0&0&1
  \end{pmatrix}\endgroup .
\]
C'est l'identité dont les lignes $6$ et $7$ sont échangées.

\paragraph{Contrôle en superposition.} Avec $\ket{+}_x\otimes\ket{0}_y\otimes\ket{1}_z
=\frac{1}{\sqrt2}\big(\ket{001}+\ket{101}\big)$ :
\[
  \mathrm{C\text{-}SWAP}\,\frac{1}{\sqrt2}\big(\ket{001}+\ket{101}\big)
  =\frac{1}{\sqrt2}\big(\ket{001}+\ket{110}\big).
\]
Le résultat n'est pas un produit : les coefficients de $\ket{x}\otimes\ket{yz}$ forment la
matrice (lignes $x=0,1$, colonnes $yz=00,01,10,11$)
\[
  \frac{1}{\sqrt2}\begin{pmatrix}0&1&0&0\\0&0&1&0\end{pmatrix},
\]
dont les deux lignes ne sont pas proportionnelles. Le qubit $x$ est intriqué avec la paire
$(y,z)$.

\subsubsection{Contrôlé-contrôlé-NOT : la porte de Toffoli}

\begin{definition}{Porte de Toffoli (CCNOT)}
Sur trois qubits $(x,y,z)$, avec $x$ et $y$ pour contrôles et $z$ pour cible,
\[
  \mathrm{CCNOT}=\ket{0}\bra{0}_x\otimes I_y\otimes I_z
   +\ket{1}\bra{1}_x\otimes\Big[\ket{0}\bra{0}_y\otimes I_z+\ket{1}\bra{1}_y\otimes\sigma_x\Big].
\]
Le crochet est un CNOT (contrôle $y$, cible $z$), appliqué seulement si $x=1$ : si $x=0$, rien ;
si $x=1$ et $y=0$, rien ; si $x=1$ et $y=1$, le qubit $z$ est inversé.
\end{definition}

Calcul de $\mathrm{CCNOT}\ket{110}$ :
\begin{align*}
  \mathrm{CCNOT}\ket{110}&=\ket{0}\braket{0}{1}\otimes I\ket{1}\otimes I\ket{0}
     +\ket{1}\braket{1}{1}\otimes\Big[\ket{0}\braket{0}{1}\otimes I\ket{0}
       +\ket{1}\braket{1}{1}\otimes\sigma_x\ket{0}\Big]\\
  &=0+\ket{1}\otimes\Big[0+\ket{1}\otimes\ket{1}\Big]=\ket{111}.
\end{align*}
Pour les huit états de base, $\mathrm{CCNOT}\ket{x,y,z}=\ket{x,\,y,\,z\oplus xy}$ :
\begin{center}
\renewcommand{\arraystretch}{1.3}
\begin{tabular}{@{}ccccccccc@{}}
\toprule
Entrée & $\ket{000}$ & $\ket{001}$ & $\ket{010}$ & $\ket{011}$ & $\ket{100}$ & $\ket{101}$ & $\ket{110}$ & $\ket{111}$ \\
\midrule
Sortie & $\ket{000}$ & $\ket{001}$ & $\ket{010}$ & $\ket{011}$ & $\ket{100}$ & $\ket{101}$ & $\ket{111}$ & $\ket{110}$ \\
\bottomrule
\end{tabular}
\end{center}
Seuls $\ket{110}$ et $\ket{111}$ sont échangés : la matrice est $\mathrm{diag}(I_6,X)$
(un bloc $6\times6$ égal à l'identité, puis $X$),
\[
  \mathrm{CCNOT}=\begingroup\setlength{\arraycolsep}{3pt}\begin{pmatrix}
    1&0&0&0&0&0&0&0\\
    0&1&0&0&0&0&0&0\\
    0&0&1&0&0&0&0&0\\
    0&0&0&1&0&0&0&0\\
    0&0&0&0&1&0&0&0\\
    0&0&0&0&0&1&0&0\\
    0&0&0&0&0&0&0&1\\
    0&0&0&0&0&0&1&0
  \end{pmatrix}\endgroup .
\]

\begin{remarque}[Portes réversibles]
Les portes de Fredkin et de Toffoli sont des permutations des états de base qui échangent
seulement deux d'entre eux : elles sont leurs propres inverses ($\mathrm{C\text{-}SWAP}^2=\mathrm{CCNOT}^2=\mathrm{Id}$)
et unitaires. Avec $z=0$, la Toffoli calcule le ET logique $x\wedge y$ sans rien effacer :
$\ket{x,y,0}\to\ket{x,y,xy}$ ; avec $z=1$, elle calcule le NON-ET. Elle sert à écrire tout calcul
classique de façon réversible.
\end{remarque}

% ---------------------------------------------------------------------
\subsection{Récapitulatif du chapitre}
% ---------------------------------------------------------------------

\begin{center}
\renewcommand{\arraystretch}{1.5}
\begin{tabular}{@{}lp{12cm}@{}}
\toprule
\textbf{Notion} & \textbf{Formule ou résultat} \\
\midrule
CNOT & $\ket{b,a}\to\ket{b\oplus a,a}$ (contrôle $a$ écrit en dernier), $\mathrm{CNOT}^2=\mathrm{Id}$ \\
SWAP & $\mathrm{SWAP}\ket{a,b}=\ket{b,a}$, \quad $\mathrm{CNOT}_{21}=\mathrm{SWAP}\,\mathrm{CNOT}_{12}\,\mathrm{SWAP}$ \\
Circuit de Bell & $(I\otimes H)$ puis CNOT : $\ket{\beta^{xy}}=\frac{1}{\sqrt2}\big(\ket{y,0}+(-1)^x\ket{\bar y,1}\big)$, $\ket{\beta^{00}}=\ket{\Phi^+}$ \\
Information & $I(A;B)=H(A)+H(B)-H(A,B)$, \quad $H(A|B)=H(A,B)-H(B)$ \\
Trace partielle & $\operatorname{Tr}_B\big(\ket{a}\bra{a'}\otimes\ket{b}\bra{b'}\big)=\ket{a}\bra{a'}\braket{b'}{b}$ ; état de Bell : $\rho_A=\frac12\mathrm{Id}$ \\
État de Bell (bits) & $H(A,B)=0$, \quad $H(A)=H(B)=1$, \quad $I(A;B)=2$, \quad $H(A|B)=-1$ \\
Pur / mélangé & $\operatorname{Tr}\rho^2=1$ / $<1$ ; $\ket{+}$ : $1$ ; qubit d'une paire de Bell : $\frac12$ ; source de Bob : $\frac58$ \\
Téléportation & $1$ e-bit $+\ 2$ bits classiques ; correction : $X$ si $a=1$, puis $Z$ si $q=1$ ; $P(a,q)=\frac14$ \\
Contrôlé-$U$ & $\ket{0}\bra{0}\otimes I+\ket{1}\bra{1}\otimes U=\mathrm{diag}(I,U)$ ; $\ket{x,y}\to\ket{x}\otimes U^x\ket{y}$ \\
Fredkin & $\ket{0}\bra{0}\otimes I\otimes I+\ket{1}\bra{1}\otimes\mathrm{SWAP}$ ; $\ket{101}\leftrightarrow\ket{110}$ \\
Toffoli & $\ket{x,y,z}\to\ket{x,y,z\oplus xy}$ ; $\ket{110}\leftrightarrow\ket{111}$ \\
\bottomrule
\end{tabular}
\end{center}
 % (optionnel) fichier supplémentaire, qui serait numéroté 8

%%%%% Conclusion (section non numérotée, ajoutée au sommaire) %%%%%
% =====================================================================
%  Conclusion : importée par main.tex avec % =====================================================================
%  Conclusion : importée par main.tex avec % =====================================================================
%  Conclusion : importée par main.tex avec \include{conclusion}.
%  Section non numérotée, ajoutée au sommaire.
% =====================================================================

% Macros de Dirac (déjà définies dans main.tex ; ce bloc évite l'erreur
% « Undefined control sequence \ket » si main.tex est une ancienne version).
\providecommand{\ket}[1]{|#1\rangle}
\providecommand{\bra}[1]{\langle #1|}
\providecommand{\braket}[2]{\langle #1|#2\rangle}

\phantomsection
\section*{Conclusion}
\addcontentsline{toc}{section}{Conclusion}

On est parti de l'onde de matière de de Broglie et de l'équation de Schrödinger pour
aboutir à un modèle de calcul complet. Le tableau ci-dessous résume la traduction de la
physique ondulatoire vers l'informatique quantique.

\begin{center}
\renewcommand{\arraystretch}{1.45}
\begin{tabular}{@{}p{3.6cm}p{5.4cm}p{5.6cm}@{}}
\toprule
\textbf{Notion} & \textbf{Mécanique ondulatoire} & \textbf{Informatique quantique} \\
\midrule
État & fonction d'onde $\Psi(x,t)$ & vecteur $\ket{\Psi}$ d'un espace de Hilbert, ou $\rho$ \\
Grandeur physique & opérateur $\hat p$, $\hat H$ & matrice hermitienne \\
Évolution & équation de Schrödinger & porte unitaire $U$ \\
Mesure & règle de Born $|\Psi|^2$ & $P(i)=|\braket{i}{\Psi}|^2$ et projection \\
Plusieurs systèmes & fonction de plusieurs variables & produit tensoriel $\otimes$ \\
Corrélations & --- & intrication, états de Bell \\
\bottomrule
\end{tabular}
\end{center}

L'intrication est la ressource qui distingue les systèmes quantiques : le protocole de
téléportation du chapitre~\ref{chap:bell} la consomme pour transmettre un état inconnu, sans le copier.

\paragraph{Pour aller plus loin.} Ce document ne traite ni les algorithmes (Deutsch, Grover,
Shor \parencite{shor1997}), ni les inégalités de Bell et CHSH \parencite{bell1964}, ni la
correction d'erreurs quantiques, qui prolongent directement les chapitres~\ref{chap:portes}
et~\ref{chap:bell}. Les ouvrages de \textcite{nielsen2010} et de \textcite{mermin2007} en
donnent une présentation complète.
.
%  Section non numérotée, ajoutée au sommaire.
% =====================================================================

% Macros de Dirac (déjà définies dans main.tex ; ce bloc évite l'erreur
% « Undefined control sequence \ket » si main.tex est une ancienne version).
\providecommand{\ket}[1]{|#1\rangle}
\providecommand{\bra}[1]{\langle #1|}
\providecommand{\braket}[2]{\langle #1|#2\rangle}

\phantomsection
\section*{Conclusion}
\addcontentsline{toc}{section}{Conclusion}

On est parti de l'onde de matière de de Broglie et de l'équation de Schrödinger pour
aboutir à un modèle de calcul complet. Le tableau ci-dessous résume la traduction de la
physique ondulatoire vers l'informatique quantique.

\begin{center}
\renewcommand{\arraystretch}{1.45}
\begin{tabular}{@{}p{3.6cm}p{5.4cm}p{5.6cm}@{}}
\toprule
\textbf{Notion} & \textbf{Mécanique ondulatoire} & \textbf{Informatique quantique} \\
\midrule
État & fonction d'onde $\Psi(x,t)$ & vecteur $\ket{\Psi}$ d'un espace de Hilbert, ou $\rho$ \\
Grandeur physique & opérateur $\hat p$, $\hat H$ & matrice hermitienne \\
Évolution & équation de Schrödinger & porte unitaire $U$ \\
Mesure & règle de Born $|\Psi|^2$ & $P(i)=|\braket{i}{\Psi}|^2$ et projection \\
Plusieurs systèmes & fonction de plusieurs variables & produit tensoriel $\otimes$ \\
Corrélations & --- & intrication, états de Bell \\
\bottomrule
\end{tabular}
\end{center}

L'intrication est la ressource qui distingue les systèmes quantiques : le protocole de
téléportation du chapitre~\ref{chap:bell} la consomme pour transmettre un état inconnu, sans le copier.

\paragraph{Pour aller plus loin.} Ce document ne traite ni les algorithmes (Deutsch, Grover,
Shor \parencite{shor1997}), ni les inégalités de Bell et CHSH \parencite{bell1964}, ni la
correction d'erreurs quantiques, qui prolongent directement les chapitres~\ref{chap:portes}
et~\ref{chap:bell}. Les ouvrages de \textcite{nielsen2010} et de \textcite{mermin2007} en
donnent une présentation complète.
.
%  Section non numérotée, ajoutée au sommaire.
% =====================================================================

% Macros de Dirac (déjà définies dans main.tex ; ce bloc évite l'erreur
% « Undefined control sequence \ket » si main.tex est une ancienne version).
\providecommand{\ket}[1]{|#1\rangle}
\providecommand{\bra}[1]{\langle #1|}
\providecommand{\braket}[2]{\langle #1|#2\rangle}

\phantomsection
\section*{Conclusion}
\addcontentsline{toc}{section}{Conclusion}

On est parti de l'onde de matière de de Broglie et de l'équation de Schrödinger pour
aboutir à un modèle de calcul complet. Le tableau ci-dessous résume la traduction de la
physique ondulatoire vers l'informatique quantique.

\begin{center}
\renewcommand{\arraystretch}{1.45}
\begin{tabular}{@{}p{3.6cm}p{5.4cm}p{5.6cm}@{}}
\toprule
\textbf{Notion} & \textbf{Mécanique ondulatoire} & \textbf{Informatique quantique} \\
\midrule
État & fonction d'onde $\Psi(x,t)$ & vecteur $\ket{\Psi}$ d'un espace de Hilbert, ou $\rho$ \\
Grandeur physique & opérateur $\hat p$, $\hat H$ & matrice hermitienne \\
Évolution & équation de Schrödinger & porte unitaire $U$ \\
Mesure & règle de Born $|\Psi|^2$ & $P(i)=|\braket{i}{\Psi}|^2$ et projection \\
Plusieurs systèmes & fonction de plusieurs variables & produit tensoriel $\otimes$ \\
Corrélations & --- & intrication, états de Bell \\
\bottomrule
\end{tabular}
\end{center}

L'intrication est la ressource qui distingue les systèmes quantiques : le protocole de
téléportation du chapitre~\ref{chap:bell} la consomme pour transmettre un état inconnu, sans le copier.

\paragraph{Pour aller plus loin.} Ce document ne traite ni les algorithmes (Deutsch, Grover,
Shor \parencite{shor1997}), ni les inégalités de Bell et CHSH \parencite{bell1964}, ni la
correction d'erreurs quantiques, qui prolongent directement les chapitres~\ref{chap:portes}
et~\ref{chap:bell}. Les ouvrages de \textcite{nielsen2010} et de \textcite{mermin2007} en
donnent une présentation complète.


%%%%% Bibliographie %%%%%
\clearpage
\printbibliography[title=Bibliographie, heading=bibintoc]

%%%%% Table des illustrations %%%%%
%\newpage
%\renewcommand{\listfigurename}{ }
%\section*{Table des illustrations}
%\addcontentsline{toc}{section}{Table des illustrations}
%\listoffigures

%%%%% Glossaire %%%%%
\clearpage
\phantomsection
\section*{Glossaire}
\addcontentsline{toc}{section}{Glossaire}
\glsaddall
\begingroup
\renewcommand*{\glossarysection}[2][]{}% titre déjà donné ci-dessus
\printnoidxglossary
\endgroup

%%%%% Annexes %%%%%
%\newpage
%\section{Annexes}
%\input{annexes}

\end{document}
